\documentclass[11pt]{article}
\usepackage[margin=1in]{geometry}
\usepackage[T1]{fontenc}
\usepackage{lmodern,microtype}
\usepackage{amsmath,amssymb,amsthm,mathtools,bm,mathrsfs}
\usepackage{booktabs,longtable,array,enumitem}
\usepackage{xurl}
\usepackage[colorlinks=true,linkcolor=blue,citecolor=blue,urlcolor=blue]{hyperref}
\usepackage[nameinlink,noabbrev]{cleveref}
\hypersetup{pdftitle={Perturbative Einstein--Standard Model EFT Continuum Programme: V71},pdfauthor={Research notes},pdfsubject={Two quartic vertices, generated marginal coupling, and the retained six-field return},pdfkeywords={Einstein Standard Model, effective field theory, measured renormalization, gravitational shell, BV}}
\allowdisplaybreaks
\setlength{\emergencystretch}{3em}
\setlength{\parskip}{4pt}
\setlength{\parindent}{0pt}
\setlist{itemsep=2pt,topsep=4pt}
\newtheorem{theorem}{Theorem}[section]
\newtheorem{proposition}[theorem]{Proposition}
\newtheorem{lemma}[theorem]{Lemma}
\newtheorem{corollary}[theorem]{Corollary}
\theoremstyle{definition}
\newtheorem{definition}[theorem]{Definition}
\newtheorem{hypothesis}[theorem]{Hypothesis}
\crefname{hypothesis}{Hypothesis}{Hypotheses}
\theoremstyle{remark}
\newtheorem{remark}[theorem]{Remark}
\newcommand{\R}{\mathbb R}
\newcommand{\C}{\mathbb C}
\newcommand{\E}{\mathbb E}
\newcommand{\Tr}{\operatorname{Tr}}
\newcommand{\Cov}{\operatorname{Cov}}
\newcommand{\Cum}{\operatorname{Cum}}
\newcommand{\dd}{\mathrm d}
\newcommand{\Id}{\mathrm{id}}
\newcommand{\loc}{\mathrm{loc}}
\newcommand{\ext}{\mathrm{ext}}
\newcommand{\op}{\mathrm{op}}
\newcommand{\wt}{\operatorname{wt}}
\newcommand{\norm}[1]{\left\lVert#1\right\rVert}
\newcommand{\abs}[1]{\left|#1\right|}
\newcommand{\ip}[2]{\left\langle#1,#2\right\rangle}
\newcommand{\jet}{\mathrm J}
\newcommand{\srcid}[1]{\nolinkurl{#1}}
\newcommand{\PKD}{\mathscr P_{K,d}}
\usepackage{needspace}
\usepackage{etoolbox}
\makeatletter
\patchcmd{\l@section}{1.5em}{2.1em}{}{}
\patchcmd{\l@subsection}{2.3em}{3.1em}{}{}
\makeatother
\title{Perturbative Einstein--Standard Model\\[3pt]
Effective Field Theory Continuum Programme\\[4pt]
\large Cumulative notes, V71\\
Two quartic vertices and the retained six-field return}
\author{Research notes}
\date{20 September 2026}
\begin{document}
\maketitle
\noindent\textbf{Version history.} V1--V70 (Sections 1--543), including
all mathematical text, labels and continuation markers, are retained
verbatim below. V71 begins at Section~\ref{sec:v71-start}. It evaluates
the connected two-quartic coefficient, supplies its one-loop four-field
normalization and tests the induced marginal update under exact
Gaussian composition. The remaining source and EFT conditions are
recorded in Section~\ref{sec:v71-ledgers}. The original abstract
summarizes V1.

\begin{abstract}
This first installment of a separate perturbative continuum lineage addresses
which information an actual measured shell must retain before it can be
iterated. A quadratic local source in the inherited positive Gaussian gauge
shell is shown to generate a nonzero source-quadratic, field-independent
coefficient. Such a term is not a removable vacuum normalization. A finite
source cumulant tower, rather than a linear source map alone, is therefore
necessary. The inherited two-loop field-degree packet is extended to every
prescribed finite loop order by a weighted-jet theorem for nonlinear measured
pushforwards. Its proof identifies all required action and constraint jets,
proves composition at the coefficient level, and gives connected spatial
bounds without an extensive volume loss. The same weighted truncation determines
the finite BV master-equation defect under stated algebraic hypotheses.
Finally, physical kinetic extraction distinguishes the leading measured
one-loop shift from its two-loop correction. These results close finite-order
algebraic interfaces in P1 and part of P7, with an algebraic P3 compatibility
statement. They do not establish an invariant analytic class for the full
compact integral, the sign of an internally generated YM flow, an all-order
uniform remainder, or a continuum limit. The missing complete-shell term is
retained explicitly, and the last section records the next necessary estimate.
\end{abstract}
\clearpage
\tableofcontents
\clearpage

\section{Context, source audit, and conventions}
\label{sec:v1-context}

The supplied programme \cite{Programme} asks for the chain
\[
\begin{gathered}
\text{measured shell}\ \longrightarrow\ \text{invariant RG class}
\ \longrightarrow\ \text{generated couplings}\\
\longrightarrow\ \text{continuum source functionals}.
\end{gathered}
\]
The present filename starts at V1 as requested; it does not renumber the
attached Yang--Mills or Standard-Model lineages. Subsequent installments are
to be appended before the continuation marker at the end, preserving the
mathematical body and its labels.

Two endpoints remain distinct. For a coefficientwise perturbative continuum, each prescribed finite-order
coefficient must have a renormalized limit as the cutoff is removed. An
analytic trajectory additionally requires estimates for the actual
interacting law over the claimed scale interval. Formal expansion alone
does not identify these endpoints. No nonperturbative mass-gap claim,
positive continuum measure, or ordinary Standard-Model ultraviolet
completion is made here. Gravity is not integrated in this installment.

\subsection{What is inherited, and what is not repeated}

The dependency audit uses the attached YM closure monograph V7 \cite{YM7}
and SM--Einstein continuum monograph \cite{SM}, with their exact internal
labels. The accompanying Grand Tensor and mass-gap files are background
sources, not substitutes for a perturbative full-shell estimate. In particular:

\begin{center}
\small
\begin{tabular}{>{\raggedright\arraybackslash}p{0.32\textwidth}>{\raggedright\arraybackslash}p{0.61\textwidth}}
\toprule
Inherited input & Source location and limitation\\
\midrule
Exact blocking and measured redundancy & YM V7,
\srcid{thm:measured-redundancy}; the complete divergence/Jacobian tangent is
retained. No new proof of this result is claimed.\\
Positive eliminated Gaussian and root chart & YM V7,
\srcid{sec:nonlinear-root-block} and \srcid{eq:moving-minimum}; positivity
concerns the eliminated fibre, not the retained massless gauge field.\\
Local measured two-loop expansion & YM V7,
\srcid{thm:measured-two-loop}; its remainder is controlled through four marks
on the local box, not on the full compact complement.\\
Finite two-loop composition & YM V7,
\srcid{thm:measured-loop-composition} and
\srcid{thm:closed-loop-jet-packet}; the packet
$(\jet^8 S,\jet^6 Q,\jet^4 B;\jet^7 F)$ is already established.\\
Conditional stable and source bounds & YM V7,
\srcid{eq:reciprocal}, \srcid{cor:two-loop-summability}; SM,
\srcid{thm:stable-graph} and \srcid{thm:localized-shell}.
These retain their supplied-trajectory and moment hypotheses.\\
BV and determinant interfaces & SM, sections ``BRST, BV, and exact
elimination'' and ``Spatial, scale, and reflection compatibility'';
\srcid{thm:coherent-line}. Algebraic BV pushforward and conditional line
scalarization are not analytic full-shell closure.\\
\bottomrule
\end{tabular}
\end{center}

The programme refers to an earlier perturbative V1. A separate manuscript
with that title is not among the eight files in this attachment set. The
inherited statements used here are therefore anchored to the displayed
monograph locations rather than to an uninspected Section~7.
The two new substantive extensions are the source-complete coefficient carrier
and the arbitrary-finite-order weighted-jet theorem. The two-loop special case
is cited, not presented as a new finding.

Perturbative renormalization in BV and background-field settings is already
studied in other constructions \cite{Costello,Anselmi}. Those results provide
context, not an imported estimate for this particular compact blocking map.
The proofs below are finite-regulator algebra and explicitly stated spatial
estimates. They are not claims to the first general perturbative
renormalization theorem.

\subsection{Loop order, coupling order, field degree, and canonical dimension}

Write the dimensionless negative logarithm of a density as
\begin{equation}
 A_h(x,J)=h^{-1}S_0(x,J)+S_1(x,J)+hS_2(x,J)+\cdots,
 \qquad \mathcal S_h=hA_h=\sum_{\ell\ge0}h^\ell S_\ell.
 \label{eq:v1-action-convention}
\end{equation}
Here $h$ is a formal loop-counting variable for coefficient statements, and a
small positive parameter only when an actual integral is explicitly used.
For the local YM convention it is $\beta^{-1}$ before physical recalibration.
Identifying it with a running physical coupling after every step is a further
normalization operation. It is not built into the formal calculus.

Thus $S_1$ is the one-loop term in the negative logarithm and $hS_2$ is its
two-loop term. Sources $J$ are not integrated. An ordinary source factor
$e^{\sum_a J_a O_a}$ contributes $-h\sum_aJ_aO_a$ to
$\mathcal S_h$. Field and source Taylor degrees will be counted together;
canonical dimensions and external-momentum orders will not. No canonical
power-counting conclusion follows solely from the degree count below.

At a finite regulator, subtracting an additive constant independent of all
retained fields \emph{and} sources is harmless for normalized source
functionals. A term that is field independent but source dependent is not
such a constant. This distinction is needed before a recursive class can be
specified.

\section{The first necessary enlargement: a source cumulant tower}
\label{sec:v1-source-tower}

\subsection{Exact conditional source transport}

Let $B:X\to Y$ be the actual finite blocking map and let $\mu(\dd x\mid y)$
be its normalized conditional law at zero external source. For a fixed finite
family of observables $O_1,\ldots,O_m$, set
\begin{equation}
 H_B(y,J)=\log\E\!\left[
      e^{\sum_aJ_aO_a}\mid B=y\right],\qquad H_B(y,0)=0.
 \label{eq:v1-conditional-source-log}
\end{equation}
For actual identities assume a nonzero moment-generating function on the
source neighbourhood in question. For formal identities, take Taylor
coefficients at zero when the required moments exist. The new action is
\begin{equation}
 A_+(y,J)=A_+(y,0)-H_B(y,J).
 \label{eq:v1-source-action-update}
\end{equation}
Its source coefficients are the conditional cumulants
\[
 H_B(y,J)=\sum_{r\ge1}\frac1{r!}
       \sum_{a_1,\ldots,a_r}J_{a_1}\cdots J_{a_r}
       C^{(r)}_{a_1\ldots a_r}(y).
\]
The first coefficient is the usual conditional observable; the second is a
conditional covariance. Higher coefficients are part of the effective action,
not an optional correction to a linear transport.

\begin{theorem}[Exact finite-source composition]
\label[theorem]{thm:v1-total-cumulance}
Consider $X\xrightarrow{B_1}Y\xrightarrow{B_2}Z$ with both conditional laws
derived from the same original zero-source law. Let $C_I(Y)$ denote the
conditional cumulant of $(O_i)_{i\in I}$ given $Y$. Then, conditional on $Z$,
\begin{equation}
 \Cum(O_i:i\in I\mid Z)
 =\sum_{\pi\in\mathfrak P(I)}
   \Cum\bigl(C_B(Y):B\in\pi\mid Z\bigr).
 \label{eq:v1-total-cumulance}
\end{equation}
The source tower through order $r$ closes under a finite sequence of exact
blockings: no source coefficient of order greater than $r$ contributes to a
coefficient of order at most $r$. This statement alone gives no norm bound
uniform in the number of blockings.
\end{theorem}
\begin{proof}
The tower property gives
\[
 \E[e^{J\cdot O}\mid Z]
 =\E[e^{H_{B_1}(Y,J)}\mid Z].
\]
Expand the logarithm on the right in outer conditional cumulants of
$H_{B_1}$. A derivative with one distinct mark for each $i\in I$ assigns
these marks to nonempty groups, each group differentiating one copy of
$H_{B_1}$. The groups form a set partition $\pi$; the inner derivative is
$C_B$, and the outer derivative is the cumulant displayed in
\eqref{eq:v1-total-cumulance}. The exponential and cumulant factorials cancel
the orderings of those groups. All groups have size at most $|I|$, proving
the finite-order assertion. Repeated marks follow by polarization.
\end{proof}

For example, at second order the two contributions are
\[
 \E[C^{(2)}_{ab}(Y)\mid Z]
 +\Cov(C^{(1)}_a(Y),C^{(1)}_b(Y)\mid Z).
\]
Keeping only a composed linear observable map drops the first term. This is
not repaired by a later scalar partition-function normalization.

\subsection{A nonzero contact coefficient in the inherited gauge shell}

The following calculation uses the same kind of positive eliminated Gaussian
as the local root construction. It is not an alternative fixed-reference
model for the interacting law.

\begin{theorem}[Quadratic gauge sources generate source-only terms]
\label[theorem]{thm:v1-gaussian-contact}
Let $\eta$ be a real Gaussian vector of covariance $hC$ on the eliminated
fibre, with $C>0$. For symmetric matrices $A_a$ define
\[
 O_a(\eta)=\tfrac12\eta^{\mathsf T}A_a\eta,
 \qquad A(J)=\sum_a J_aA_a.
\]
For real $J$ such that $I-hC^{1/2}A(J)C^{1/2}>0$,
\begin{align}
 H(J)&=\log\E e^{\sum_aJ_aO_a}
      =-\tfrac12\Tr\log\bigl(I-hC^{1/2}A(J)C^{1/2}\bigr),
 \label{eq:v1-quadratic-log}\\
 D_aD_bH(0)&=\frac{h^2}{2}\Tr(CA_aCA_b).
 \label{eq:v1-quadratic-contact}
\end{align}
In particular $D_a^2H(0)>0$ whenever $A_a\ne0$.
For $A_p=E^*d_p^*d_pE$, where $E$ is the eliminated-field inclusion and
$d_p$ is a linearized plaquette curl, at least one such coefficient is
strictly positive whenever
$\sum_p A_p=E^*d^*dE>0$ on a nonzero eliminated fibre.
\end{theorem}
\begin{proof}
Completing the Gaussian quadratic form gives the determinant in
\eqref{eq:v1-quadratic-log}. On a smaller neighbourhood its convergent series is
\[
 H(J)=\frac12\sum_{n\ge1}\frac{h^n}{n}\Tr\bigl((CA(J))^n\bigr).
\]
The $n=2$ term yields \eqref{eq:v1-quadratic-contact} by cyclicity of the
trace. For $a=b$, the trace is the squared Hilbert--Schmidt norm of the
nonzero real symmetric matrix $C^{1/2}A_aC^{1/2}$. Positivity of
$\sum_p A_p$ on a nonzero space implies that at least one $A_p$ is nonzero.
These $A_p$ are precisely restrictions of the linearized local curvature
quadratics to the Gaussian fibre.
\end{proof}

The theorem evaluates a coefficient of the inherited shell, not an assumed
cofinal interacting covariance. The negative effective action contains
$-H(J)$, so the corresponding sign in that action is negative. Subtraction
of $H(0)=0$ does not remove the $J^2$ term. Normal ordering the individual
quadratics changes their first cumulants and leaves the second cumulant
unchanged. In \eqref{eq:v1-action-convention}, this particular $J^2$ term
has order $h^2$ in the negative logarithm and order $h^3$ in
$\mathcal S_h$. A packet for these unscaled sources must therefore use
$K\ge3$ to include it. Ordinary insertion sources and the inherited
background-field marks are not assigned interchangeable loop orders.

For separated plaquette labels the coefficient can be quasilocal rather
than supported exactly on the diagonal. Its local part is a source contact
term; its remaining kernel must also be retained. Thus a source-complete
state contains source-polynomial interaction kernels, including kernels
with no retained-field legs. A vector of wave-function factors and a linear
map on $J$ alone is not closed even at this Gaussian stage.

\section{An arbitrary-order weighted-jet compiler}
\label{sec:v1-compiler}

\subsection{The coefficient carrier}

Fix integers $K\ge1$ and $d\ge2$, and put $D=d+2K$. All retained fields,
eliminated fields, and nonintegrated source marks are assigned Taylor weight
one; $h$ has weight two. For a monomial $h^\ell x^\alpha J^\gamma$ set
\[
 \wt(h^\ell x^\alpha J^\gamma)=2\ell+|\alpha|+|\gamma|.
\]
Source variables are never contracted. Define the packet
\begin{equation}
 \PKD(\mathcal S)
   =\bigl(\jet^{D}S_0,\jet^{D-2}S_1,\ldots,
                   \jet^{D-2K}S_K\bigr).
 \label{eq:v1-packet}
\end{equation}
It is the quotient by terms with loop index greater than $K$ or weight
greater than $D$. This is a finite Taylor-degree carrier at each regulator;
it is not a replacement for the complete analytic action.

\begin{hypothesis}[Finite measured stationary branch]
\label[hypothesis]{hyp:v1-branch}
The fields are finite dimensional. $S_0$ and the blocking map $F$ are analytic
near zero, $S_0(0)=0$, $DS_0(0)=0$, $F(0)=0$, and $DF(0)$ is surjective.
The quadratic form $D^2S_0(0)$ is strictly positive on $\ker DF(0)$.
The density coefficients $S_\ell$ are analytic germs. Ordinary sources occur
in $S_\ell$, $\ell\ge1$, and are fixed during blocking. A branch of the
nonzero reference density is specified. For an iterated statement these
conditions hold at each intermediate stationary branch.
\end{hypothesis}

Choose a right inverse $R$ of $L=DF(0)$ and an inclusion $E$ of $\ker L$.
There is an analytic fibre chart
\begin{equation}
 X(y,c)=Ey+Rv(y,c),\qquad F(X(y,c))=c.
 \label{eq:v1-chart}
\end{equation}
The normalized measured action in this chart is
\begin{equation}
 \widetilde{\mathcal S}_h(y,c,J)
   =\mathcal S_h(X(y,c),J)-h\log|\det D_{(y,c)}X|.
 \label{eq:v1-measured-chart}
\end{equation}
Any additional prescribed incoming reference density is included in $S_1$;
a nonconstant outgoing reference is added once after elimination. The
affine quadratic shift in $y$ removes the mixed quadratic form. It uses
only $D^2S_0(0)$ and $DF(0)$.

\begin{lemma}[Graph weight and finiteness]
\label[lemma]{lem:v1-graph-weight}
After this quadratic shift, separate the Gaussian and the field-independent
constants. An interaction vertex from $h^{\ell_v}S_{\ell_v}$ has total
field/source valence $n_v$ and contributes $h^{\ell_v-1}$ to the exponent.
For a connected Wick graph with $v$ vertices, $e$ internal edges, and $r$
uncontracted field/source legs, its contribution to $-h\log Z$ has loop index
\begin{equation}
 L_\Gamma=\sum_{v}\ell_v+e-v+1\ge0,
 \qquad
 r+2L_\Gamma=2+\sum_v(n_v+2\ell_v-2).
 \label{eq:v1-graph-weight}
\end{equation}
Every nonconstant interaction vertex has $n_v+2\ell_v-2\ge1$. Consequently,
only finitely many graphs contribute to $L_\Gamma\le K$ and
$r+2L_\Gamma\le D$. A primitive coefficient of weight greater than $D$ or
loop index greater than $K$ cannot contribute to this output packet.
\end{lemma}
\begin{proof}
Each covariance contributes $h$, each exponent vertex contributes
$h^{\ell_v-1}$, and the final multiplication by $-h$ contributes one power.
This gives the first expression. Connectedness implies $e\ge v-1$.
Counting legs gives $r=\sum_v n_v-2e$, which proves the second expression.
The non-Gaussian classical vertices have valence at least three. For
$\ell_v=1$ a nonconstant vertex has valence at least one; for
$\ell_v\ge2$ its weight excess is positive. Constants have no edges and
contribute directly as one-vertex terms.

For a nonconstant graph in the packet, the sum of positive vertex excesses
is at most $D-2$, so the vertex number is at most $D-2$. Every vertex weight
is at most $D$, which bounds all valences and the edge number; the finite
field species then give only finitely many contraction patterns. If a
vertex has weight $w>D$, \eqref{eq:v1-graph-weight} gives an output weight
at least $w$ because all other vertex excesses are nonnegative. A vertex
with $\ell_v>K$ forces $L_\Gamma>K$ by connectedness. This proves the final
assertion, including disconnected-free constants.
\end{proof}

\begin{theorem}[Source-complete measured weighted-jet closure]
\label[theorem]{thm:v1-all-order-packet}
Under \cref{hyp:v1-branch}, write the formal measured pushforward, after
separating its source-independent Gaussian scalar normalization, as
\[
 -\log F_*\bigl(e^{-\mathcal S_h/h}\dd x\bigr)
      =h^{-1}G_0+G_1+hG_2+\cdots.
\]
The input
\begin{equation}
 \boxed{\bigl(\jet^{d+2K}S_0,\jet^{d+2K-2}S_1,\ldots,
                   \jet^d S_K;\ \jet^{d+2K-1}F\bigr)}
 \label{eq:v1-main-packet}
\end{equation}
determines exactly
\[
 (\jet^{d+2K}G_0,\jet^{d+2K-2}G_1,\ldots,\jet^dG_K).
\]
The statement includes all mixed source coefficients in those Taylor
polynomials and all source-only coefficients. On compatible stationary
branches the coefficient maps compose:
\begin{equation}
 \mathfrak R_{F_2}^{K,d}\mathfrak R_{F_1}^{K,d}
     =\mathfrak R_{F_2\circ F_1}^{K,d}.
 \label{eq:v1-compiler-composition}
\end{equation}
No assertion of an $h^K$ remainder bound uniform in regulator, volume,
or number of steps is part of this theorem.
\end{theorem}
\begin{proof}
Set $D=d+2K$. In \eqref{eq:v1-chart}, the homogeneous degree-$q$ equation
for $v$ has the identity as its linear coefficient. Induction therefore
determines $\jet^{D-1}X$ from $\jet^{D-1}F$. Changing the unknown higher
part of $X$ by $O(|(y,c)|^D)$ changes $S_0\circ X$ by
$O(|(y,c)|^{D+1})$, since $DS_0(0)=0$. Thus $\jet^D S_0$ and this chart
jet determine $\jet^D(S_0\circ X)$.

For $\ell\ge1$, an error in $S_\ell$ of Taylor degree greater than
$D-2\ell$ stays of greater degree after substitution into a zero-preserving
chart. The unknown part of $X$ cannot affect these smaller jets. The
Jacobian uses one derivative of $X$: its jet through degree $D-2$ is fixed
by $\jet^{D-1}X$. Since it is multiplied by $h$, this is exactly its required
budget in $S_1$. This proves that \eqref{eq:v1-main-packet} determines the
complete measured interaction packet before Gaussian elimination.

After the affine harmonic shift, expand the formal exponential and its
logarithm. To see explicitly why only connected contractions remain,
partition every Wick graph into connected components. The exponential
formula sums unordered collections of components with their factorials;
taking the logarithm leaves one component. This is a finite combinatorial
identity at each weight by \cref{lem:v1-graph-weight}. That lemma then
shows that no omitted primitive coefficient can influence an output
coefficient in the packet. The Gaussian normalization contributes only
the declared scalar and its measured quadratic-reference determinant.

For composition, integrate the two eliminated Gaussian coordinate sets in
stages. Their covariance blocks and Schur complement are those of the
single joint Gaussian; the determinant factors multiply, and the complete
chart Jacobians obey the chain rule. Wick contraction of a polynomial is
unchanged by the order of these finite Gaussian integrations. Sorting the
same connected diagrams by their first elimination therefore gives the
staged and joint formal logarithms identically. The exclusion statement
of \cref{lem:v1-graph-weight} allows the weighted packet to be taken between
steps. Equivalently, this is formal Fubini on the compatible stationary
branch, with no interchange of infinite cutoff limits. The classical
minimum and its positive second Schur complement ensure that the two
expansions describe that same branch.
\end{proof}

For $K=2$ and $d=4$, \eqref{eq:v1-main-packet} is exactly the inherited
$(\jet^8S,\jet^6Q,\jet^4B;\jet^7F)$ packet. At $K=3,d=4$ it becomes
\[
 (\jet^{10}S_0,\jet^8S_1,\jet^6S_2,\jet^4S_3;\jet^9F).
\]
The extension concerns arbitrary finite loop order, not merely another
calculation of the inherited two-loop coefficient. For a requested mixed
field/source derivative of total order at most $d$, it includes that
coefficient at every loop index through $K$.

\subsection{Why the higher input degrees cannot be discarded}

\begin{proposition}[Primitive-degree and constraint-degree tests]
\label[proposition]{prop:v1-sharpness}
Fix an even $d\ge2$. For $0\le\ell\le K$, an input monomial of field
degree $d+2(K-\ell)$ in $S_\ell$ can affect the degree-$d$ coefficient of
$G_K$. A degree-$(d+2K-1)$ change of the blocking constraint can do so as
well. These tests rule out a common degree-$d$ truncation of every input.
\end{proposition}
\begin{proof}
Let $S_0(c,y)=\tfrac12(c^2+y^2)$ and add
$\lambda h^\ell c^d y^{2(K-\ell)}$ to $\mathcal S_h$.
The derivative at $\lambda=0$ of its normalized effective action is
\[
 h^\ell c^d\E[y^{2(K-\ell)}]
   =(2(K-\ell)-1)!!\,h^Kc^d,
\]
where $y$ has variance $h$ and $(-1)!!=1$. It is nonzero. This is a formal
coefficient identity; for the displayed even powers the first derivative
also follows from a stable one-sided $\lambda\ge0$ integral.

For the constraint test use the same quadratic action in $(x,y)$ and
\[
 F_\varepsilon(x,y)=x+\varepsilon x^{d-1}y^{2K}.
\]
Solving $F_\varepsilon=c$ gives
$x=c-\varepsilon c^{d-1}y^{2K}+O(\varepsilon^2)$.
The order-$\varepsilon$ measured action variation in this chart is
\[
 -c^d y^{2K}+h(d-1)c^{d-2}y^{2K}.
\]
Its Gaussian mean is
\[
 -(2K-1)!!h^Kc^d
  +(d-1)(2K-1)!!h^{K+1}c^{d-2}.
\]
The first term is in the target $G_K$ and is nonzero. The second belongs
to loop index $K+1$, so it cannot cancel the first. The Jacobian has been
included; this is a change of the actual blocking constraint, not a pure
reparametrization of an unchanged fibre.
\end{proof}

\begin{remark}[Fermions and the finite coefficient statement]
The graph argument also applies to a finite nondegenerate graded Gaussian
whose contractions each carry one power of $h$, with a regular measured
superchart and its Berezinian. Wick signs change coefficients but not
\eqref{eq:v1-graph-weight}. This extension presupposes a fixed nonzero
Gaussian determinant/Pfaffian branch. It neither crosses a zero mode nor
constructs a chiral line over gauge or cutoff families. Those are separate
parts of P8.
\end{remark}

\section{Rooted spatial norms and connected, not extensive, products}
\label{sec:v1-locality}

Let $\Lambda$ be a finite connected cell graph with a fixed finite set of
field/source components per cell. Distances are measured in current-shell
cell units. For $n\ge1$, let $\tau(x_1,\ldots,x_n)$ be the minimum length
of a connected graph joining the marked sites, with repeated marks allowed.
For a symmetric tensor kernel $K_n$ set
\begin{equation}
 \norm{K_n}_{a}
 =\sup_{x_1}\sum_{x_2,\ldots,x_n}
       e^{a\tau(x_1,\ldots,x_n)}\abs{K_n(x_1,\ldots,x_n)}.
 \label{eq:v1-kernel-norm}
\end{equation}
Finite component sums are included. For a constant use
$|\Lambda|^{-1}|K_0|$. Source-only kernels of positive source degree use
\eqref{eq:v1-kernel-norm}, not the constant convention. For nonsymmetric
tensors take the maximum over the possible anchored legs.

Write
$S_\ell=\sum_n(n!)^{-1}\langle K_{\ell,n},\varphi^{\otimes n}\rangle$,
where $\varphi$ denotes the combined field/source list. The finite direct
sum of these norms, with any fixed positive degree weights, is a Banach
norm for the packet. For an analytic completion one may use
\begin{equation}
 \norm{S_\ell}_{a,R}
 =\frac{|K_{\ell,0}|}{|\Lambda|}
  +\sum_{n\ge1}\frac{R^n}{n!}\norm{K_{\ell,n}}_a.
 \label{eq:v1-analytic-norm}
\end{equation}
A separate source radius is obtained by replacing $R^n$ by
$R^{n_{\rm field}}\rho_J^{n_{\rm source}}$. Only the finite-degree norm is
used in the next theorem; no convergence of the full analytic sum is inferred.

\begin{lemma}[An extensive action norm is not an algebra norm]
\label[lemma]{lem:v1-not-algebra}
Even at $a=0$, the rooted derivative norm of an extensive action need not
control the norm of its ordinary square uniformly in $|\Lambda|$.
\end{lemma}
\begin{proof}
For scalar variables let $F(\varphi)=\sum_{x\in\Lambda}\varphi_x$.
Its only nonzero kernel is $K_1=1$, of rooted norm one. The Hessian of
$F^2$ is $K_2(x,y)=2$, so its rooted norm is $2|\Lambda|$.
\end{proof}

Thus the recursive action space must be closed under \emph{connected}
contractions. One can form exponentials in a separate formal symmetric
algebra of disconnected components; one must not claim that the exponential
has the same small extensive-action norm. The formal logarithm returns to
connected kernels.

\begin{theorem}[Connected contraction without a volume factor]
\label[theorem]{thm:v1-connected-kernel}
Let a connected contraction graph have vertex tensors with finite norms
\eqref{eq:v1-kernel-norm}. Suppose its covariance obeys
\begin{equation}
 \begin{gathered}
 G_a=\max\left\{\sup_x\sum_y e^{a d(x,y)}\abs{C(x,y)},
                    \sup_y\sum_x e^{a d(x,y)}\abs{C(x,y)}\right\}<\infty,\\
 M_a=\sup_{x,y} e^{a d(x,y)}\abs{C(x,y)}\le G_a.
 \end{gathered}
 \label{eq:v1-covariance-norm}
\end{equation}
If $T$ is a spanning tree of its vertex graph and $e$ is the number of
all internal edges, its external tensor satisfies
\begin{equation}
 \norm{K_\Gamma}_a
 \le G_a^{|T|}M_a^{e-|T|}\prod_v\norm{K_v}_a,
 \label{eq:v1-connected-bound}
\end{equation}
for a fixed contraction pattern. For a vacuum graph the same estimate
holds after dividing its absolute value by $|\Lambda|$. Finite symmetry
and component factors are to be included when diagrams are summed.
\end{theorem}
\begin{proof}
The union of a tree joining the legs within each vertex and the internal
covariance edges is connected and meets all external legs. Therefore
\[
 \tau(\text{external legs})
 \le\sum_v\tau(\text{legs of }v)
       +\sum_{e'}d(\text{endpoints of }e').
\]
Allocate this exponential weight to the corresponding vertex and edge
factors. Bound every covariance not in $T$, including self-contractions,
by $M_a$. Anchor the root vertex at one fixed external leg and orient
$T$ away from it. A leaf vertex is anchored at its incoming covariance
leg; summing its other leg coordinates costs its rooted tensor norm.
Summing the incoming anchor against the parent covariance costs $G_a$.
Remove the leaf and repeat. At the root the remaining sum costs the
rooted norm of its tensor. This gives \eqref{eq:v1-connected-bound}.
For a vacuum graph first choose an arbitrary leg at a root vertex, run
the same argument with its coordinate fixed, and then sum that coordinate.
The final sum is at most $|\Lambda|$. A zero-valence one-vertex constant
is handled directly by the stated intensive convention.
\end{proof}

\begin{corollary}[Finite-order one-shell locality]
\label[corollary]{cor:v1-finite-locality}
At fixed $K,d$, assume uniform bounds in volume for the measured chart
and harmonic-lift jets required in \eqref{eq:v1-main-packet}, their
transformed vertex kernels, and \eqref{eq:v1-covariance-norm}. Then every
nonconstant coefficient in the output packet, and every interaction-generated
constant after intensive normalization, has a norm bounded by a finite
polynomial of these stocks. The polynomial and its degree depend on
$K,d$ and the fixed field species, not on $|\Lambda|$. An intensive
bound on the constant Gaussian log determinant additionally follows from
uniform upper and lower eigenvalue bounds and a bounded number of fibre
components per cell.
\end{corollary}
\begin{proof}
By \cref{lem:v1-graph-weight}, the packet has finitely many diagrams.
Apply \cref{thm:v1-connected-kernel} to each and sum their finite symmetry
factors. Nonconstant diagrams have at most $D-2$ vertices and at most
$D(D-2)/2$ internal edges, a sufficient crude bound for the degree of
the resulting polynomial. Constants and the Gaussian reference are
added separately. The spectral bound gives
$|\log\det H_0|\le\dim(H_0)\max\{|\log\lambda_{\min}|,
|\log\lambda_{\max}|\}$.
\end{proof}

This supplies a concrete coefficient-level spatial estimate rather than
an unspecified ``small remainder.'' It applies to a shell whose stated
primitive stocks are bounded. It does not prove that those stocks survive
arbitrarily many actual shell steps. Nor does it establish that a local
counterterm of each coefficient belongs to a renormalizable finite sector;
that question requires canonical power counting and Ward-compatible
subtraction, not just rooted summability.

\section{A compatible finite BV defect, without a false counterterm claim}
\label{sec:v1-bv}

Work in a finite graded-commutative formal algebra with Darboux field and
antifield coordinates, each of weight one, and central sources of weight one.
Let the antibracket have its constant Darboux coefficients, and let $\Delta$
be the corresponding constant-coefficient second-order BV Laplacian. Each
antibracket contraction lowers field weight by two, while $h\Delta$
preserves total weight. Use the Euclidean sign convention
\begin{equation}
 \mathcal Q(\mathcal S)=\tfrac12(\mathcal S,\mathcal S)-h\Delta\mathcal S.
 \label{eq:v1-qme-defect}
\end{equation}
The overall sign convention does not change the degree statement.

\begin{theorem}[The master-equation defect descends to the packet]
\label[theorem]{thm:v1-bv-filtration}
Suppose $\mathcal S$ is even, has nonnegative powers of $h$, and has no monomial of total
weight less than two. Let $\mathcal I_{K,D}$ consist of terms of loop index
greater than $K$ or total weight greater than $D$. For any other even
action $\widetilde{\mathcal S}$ in the same algebra,
\[
 \mathcal S\equiv\widetilde{\mathcal S}\pmod{\mathcal I_{K,D}}
 \quad\Longrightarrow\quad
 \mathcal Q(\mathcal S)\equiv\mathcal Q(\widetilde{\mathcal S})
                       \pmod{\mathcal I_{K,D}}.
\]
Consequently, on an admissible finite BV fibration for which the inherited
exact/formal BV pushforward preserves the QME, its coefficient compiler
preserves the QME modulo this packet ideal. More generally, the packet
determines its own QME defect, including any nonzero retained coefficient;
a nonzero defect is not asserted to be unchanged by pushforward.
\end{theorem}
\begin{proof}
Write $\widetilde{\mathcal S}=\mathcal S+E$ with $E\in\mathcal I_{K,D}$.
The difference of defects is
\[
 (\mathcal S,E)+\tfrac12(E,E)-h\Delta E.
\]
An antibracket of a term of weight greater than $D$ with one of weight
at least two still has weight greater than $D$. The operator $h\Delta$
subtracts two field weights and adds two $h$ weights, so also preserves
the ideal. A term of loop index greater than $K$ remains of such loop
index under a bracket with nonnegative powers, and $h\Delta$ increases
that index. These observations handle all summands, including the
quadratic term in $E$. Apply this identity to an effective action satisfying
the inherited full formal QME to obtain the last assertion.
\end{proof}

The theorem is not P3 in its entirety. Weighted Taylor truncation is not
local renormalization. It neither constructs a bounded local counterterm
projector commuting with BRST nor proves absence of a chiral anomaly.
Moreover, the compact root block has not been identified here with a
complete BV fibration, including its dual antifield projection. The claim
is exactly that the finite packet does not itself introduce an artificial
master-equation defect when that algebraic structure has been supplied.
A shifted nonstationary background must first be recentered: a classical
linear term would violate the weight hypothesis.

\section{The measured marginal coordinate and its degree provenance}
\label{sec:v1-marginal}

The supplied programme emphasizes $\mathcal B_1$, but its own convention
\[
 \widehat W_\beta=\beta G+Q_0+\beta^{-1}\mathcal B_1+O(\beta^{-2})
\]
puts the first inverse-coupling shift in $Q_0$. This is a degree issue,
not a proposed numerical beta coefficient. The density, mixed-shell, and
field-normalization terms belonging to $Q_0$ must be assembled first.

Let $I$ be the chosen physical kinetic action and let $\mathfrak b$ be a
fixed bounded linear kinetic extractor in the already calibrated background
convention, with
\[
 \mathfrak b(I)=1,\qquad \mathfrak b(\text{field/source-independent constants})=0.
\]
It is evaluated at zero external source. Its existence and its interpretation
on the physical gauge/redundant quotient are normalization data, not
consequences of an arbitrary Hessian. The monograph's measured
field-redefinition theorem is retained whenever such a quotient is used.

\begin{theorem}[Kinetic extraction through two loops]
\label[theorem]{thm:v1-marginal-extraction}
Suppose the complete real action in that same coordinate convention satisfies
\begin{equation}
 A_+(u)=u^{-1}G_0+G_1+uG_2+R_u,
 \quad \mathfrak b(G_0)=1,
 \quad\norm{R_u}\le C u^2,
 \label{eq:v1-calibrated-expansion}
\end{equation}
and $q_0=\mathfrak b(G_1)$, $q_1=\mathfrak b(G_2)$ are uniformly bounded
on the stated parameter domain. Then its extracted inverse kinetic
coefficient and physical coupling are
\begin{align}
 k_+(u)&=u^{-1}+q_0+uq_1+r_u,
       &|r_u|&\le\norm{\mathfrak b}Cu^2,\\
 u_+(u)&=k_+(u)^{-1}
       =u-q_0u^2+(q_0^2-q_1)u^3+O(u^4).
 \label{eq:v1-marginal-series}
\end{align}
The error constant is uniform in any variables in which all the displayed
input bounds are uniform. The formula is oriented in the direction of
this measured pushforward; it is not automatically the ultraviolet
inverse recursion.
\end{theorem}
\begin{proof}
Apply the bounded linear extractor to \eqref{eq:v1-calibrated-expansion}.
Factor $u^{-1}$ from $k_+$ and expand
$(1+uq_0+u^2q_1+ur_u)^{-1}$ on a domain where the modulus of its
nonconstant part is at most $1/2$. Since $ur_u=O(u^3)$, multiplication
by the leading $u$ gives \eqref{eq:v1-marginal-series} and the stated
uniform error. Positivity of the real kinetic coefficient holds for
sufficiently small positive $u$ under these same bounds.
\end{proof}

The novelty of the present application is the provenance of each degree,
not the reciprocal sensitivity estimate already in YM V7. In its notation,
$G_1$ contains the complete measured one-loop term
\begin{equation}
 Q_0(c)=\tfrac12\log\det H_c-\ell(z_*(c),c)+\log j_c(c)
 \label{eq:v1-q0-owner}
\end{equation}
plus every incoming one-loop term transported through that shell. The
coefficient $G_2$ contains the transported incoming two-loop term and the
complete measured contractions in
\srcid{eq:measured-loop-step} of \cite{YM7}. This includes the $q_y$ and
$q_{yy}$ returns of the incoming density. Taking a separate absolute
value before those signed terms are assembled changes the quantity being
extracted.

\begin{proposition}[One-loop information cannot be reconstructed from
 the two-loop coefficient alone]
\label[proposition]{prop:v1-two-loop-insufficient}
In general, knowing $G_0$ and $G_2$ does not determine $q_0$ in
\eqref{eq:v1-marginal-series}.
\end{proposition}
\begin{proof}
For a Gaussian eliminated variable $y$ let
$\mathcal S_h(c,y)=I(c)+\tfrac12 y^2+h\alpha I(c)$ and block to $c$.
The integral is exact, with $G_0=I$, $G_1=\alpha I$, and $G_2=0$ for
every $\alpha$. Thus $q_0=\alpha$ varies while the other two pieces are
unchanged. This simply tests the independent incoming measured density;
it does not assign an arbitrary counterterm to the physical YM theory.
\end{proof}

To use \cref{thm:v1-marginal-extraction} for the actual compact theory,
$A_+$ must mean the full pushforward, not its root-chart approximation.
The local theorem establishes its local analogue but gives neither the
sign nor the value of the complete $q_0$. In addition, a downward coarse
map changes the cutoff in the opposite direction from the upward UV
index of the programme. Matching and inverting the \emph{joint} coupling
and stable-coordinate map is necessary before declaring an asymptotically
free branch. No known continuum beta coefficient is inserted here to
bypass that step.

\section{A typed exact carrier and the first unclosed estimate}
\label{sec:v1-carrier}

\subsection{All discarded Taylor information has an owner}

For fixed $K,d$, a usable state at a finite scale records the following
objects, in addition to the reference quadratic form and the actual block:
\begin{equation}
 \mathfrak S_j=
 \bigl(u_j,z_j,R^{\rm ord}_j,\mathcal J_j,
        \mathcal L_j,\Phi_j,\mathcal T_j,
        C^{\ext}_j,\mathcal D^{\rm BV}_j\bigr).
 \label{eq:v1-exact-state}
\end{equation}
Here $\mathcal J_j$ is the source-log tower of
\cref{sec:v1-source-tower}, not just its linear part.
$\mathcal D^{\rm BV}_j$ is the actual defect
\eqref{eq:v1-qme-defect}. $\mathcal L_j$ is the determinant/Pfaffian line,
$\Phi_j$ its specified local frame/phase and any adjacent line map,
and $\mathcal T_j$ contains the actual field normalization and source-tower
transport. A collection of line maps without its cocycle is not a coherent
line datum; a cocycle is retained as a condition to be checked.

For the action component choose a declared, bounded local extraction
operator $\Pi_{\loc,j}$ onto a finite list of relevant/marginal and
normalization coordinates, including the required local source-contact
coordinates at the chosen source order. At this stage its Ward compatibility
is a separate recorded condition, not a theorem. Set
\begin{equation}
 \lambda_j=\Pi_{\loc,j} A_j,\qquad
 A_j=\iota_j \lambda_j+(I-\iota_j\Pi_{\loc,j})A_j,
 \qquad \Pi_{\loc,j}\iota_j=I.
 \label{eq:v1-retained-projection}
\end{equation}
The complementary term is retained in $z_j$ or in its explicitly labelled
remainder. Formula \eqref{eq:v1-retained-projection} is a coordinate split,
not a subtraction of the second summand from the theory. The numerical
coupling convention is a separate, generally nonlinear normalization
$u_j=\mathcal N_j(\lambda_j)$. In particular, the linear kinetic extractor
returns $k_j$, whereas the physical YM coordinate is $u_j^{\rm YM}=k_j^{-1}$.
A linear projector is not being identified with this reciprocal map.

In particular, writing $D_\ell=d+2(K-\ell)$, every full coefficient has
an exact Taylor split
\begin{equation}
 S_{\ell,j}=\jet^{D_\ell}S_{\ell,j}+Z_{\ell,j},
 \qquad \jet^{D_\ell}Z_{\ell,j}=0.
 \label{eq:v1-taylor-tail}
\end{equation}
The $Z_{\ell,j}$ are genuine analytic interaction tails, not zero.
Their nonparticipation in a selected finite packet is proved in
\cref{thm:v1-all-order-packet}; their smallness for a finite-amplitude
integral is \emph{not} proved by that statement.

For a local expansion through $G_K$, define the actual order remainder by
\begin{equation}
 R^{\rm ord}_{j,h}
  =A^{\loc}_{j,h}-\sum_{\ell=0}^K h^{\ell-1}S_{\ell,j}
       -c_{j,h}^{\rm vac}.
 \label{eq:v1-order-remainder}
\end{equation}
The source-independent Gaussian scalar $c_{j,h}^{\rm vac}$ is displayed
separately. An $h^K$ bound on \eqref{eq:v1-order-remainder} requires an
actual moment/differentiation theorem; the finite graph compiler does not
supply it. The inherited $K=2$, four-mark theorem is an example with
its own local hypotheses.

\subsection{The full compact ratio is a separate interaction}

Use a compatible nonnegative local cutoff or local chart contribution in
the actual compact fibre, so that $Z_j^{\loc}$ and $Z_j^{\rm full}$ are
integrals of the same measured action and source convention. On a simply
connected local domain where neither partition function vanishes, choose
the logarithmic branch fixed at the reference point and define
\begin{equation}
 C_j^{\ext}(c,J)
 =-\log\frac{Z_j^{\rm full}(c,J)}{Z_j^{\loc}(c,J)}.
 \label{eq:v1-exterior}
\end{equation}
The complete action is exactly
\begin{equation}
 A_j^{\rm full}=A_j^{\loc}+C_j^{\ext}.
 \label{eq:v1-full-local}
\end{equation}
Only $C_j^{\ext}(0,0)$ can be removed as a field- and source-independent
scalar. Its other local jets, including its kinetic and source-contact
jets, enter \eqref{eq:v1-retained-projection}; its complement stays in
$z_j$. The full compact estimates in the attached bare-Wilson notes do
not identify this ratio for a generated interacting action, as expressly
recorded in YM V7, \srcid{sec:compact-local-interface}.

For example, at a fixed background in the positive YM law, with the same
observable in both conditional integrals, normalized differentiation gives
\begin{align}
 D_{J_a}C_j^{\ext}(c,0)
 &=\E_{\loc,c}O_a-\E_{{\rm full},c}O_a,\\
 D_{J_a}D_{J_b}C_j^{\ext}(c,0)
 &=\Cov_{\loc,c}(O_a,O_b)-\Cov_{{\rm full},c}(O_a,O_b).
 \label{eq:v1-exterior-source-jets}
\end{align}
These are exact identities, not bounds inferred from the probability of
an exceptional set. If the local cutoff depends on sources, its complete
source derivative must additionally be included. The displayed formulas
use a source-independent cutoff.

A quantitative candidate class can now specify, without an unspecified
``small remainder,'' budgets such as
\begin{align}
 &\norm{S_{\ell,j}}_{a,R}\le M_\ell,\qquad
   \norm{Z_{\ell,j}}_{a,R'}\le Z_\ell\quad (R'<R),\\
 &\sup_{0<h\le\rho}h^{-K}\norm{R^{\rm ord}_{j,h}}_{a,r}\le M_{K,r},
   \qquad
   \norm{C^{\ext}_{j,h}-C^{\ext}_{j,h}(0,0)}_{a,r}
                                  \le E_r(h),
 \label{eq:v1-class-budgets}
\end{align}
with source-tower, BV-defect, frame, and adjacent-transport norms on the
same declared domains. Here $\norm{\cdot}_{a,r}$ is the rooted marked
norm of YM V7, or its mixed field/source version; order zero has the
intensive normalization. The functions $E_r$ are budgets to be established,
not numbers already furnished by the local-box theorem. Restricting to a
fixed finite mark order is permitted; an assertion for all mark orders
requires the corresponding family of estimates.

\begin{proposition}[Exact bookkeeping and finite-packet self-map]
\label[proposition]{prop:v1-bookkeeping}
At any finite step for which the indicated integrals, branch, and coordinate
extractions are defined, equations
\eqref{eq:v1-source-action-update}, \eqref{eq:v1-retained-projection},
\eqref{eq:v1-taylor-tail}, and \eqref{eq:v1-full-local} give an exact update
with no deleted complementary interaction. The local formal component
maps to the same weighted packet by \cref{thm:v1-all-order-packet} and
obeys the spatial estimate of \cref{cor:v1-finite-locality} when its
primitive stocks satisfy that corollary. These facts do not imply that
the fixed analytic budgets \eqref{eq:v1-class-budgets} are preserved.
\end{proposition}
\begin{proof}
The source identity is the conditional definition of the pushed density.
The projection and Taylor splits are exact additive identities, and
\eqref{eq:v1-full-local} follows by taking the logarithm of the same two
partition functions. Their sum reconstructs the entire action. The formal
and spatial assertions are the cited theorems. None of these algebraic
identities bounds the new exterior ratio or the finite-amplitude tails,
so none proves preservation of the displayed budgets.
\end{proof}

The notation $\mathfrak C_{\rm pert}$ will therefore be reserved for an
\emph{invariant} analytic class after those budgets have actually been
proved. At present \eqref{eq:v1-exact-state} is a typed carrier, and
\eqref{eq:v1-main-packet} is its closed finite coefficient quotient. This
separation prevents finite-order algebraic closure from being renamed a
continuum construction.

\subsection{The next load-bearing theorem}

The smallest missing analytic interface exposed by this installment is
control of the \emph{complete normalized shell}, starting with the actual
kinetic and source jets of \eqref{eq:v1-exterior}, on the interaction class
produced by the previous shell. The estimate must be uniform over that
class for $0<h\le\rho$, before a trajectory is assumed. Its measured local
part must be assigned to the running couplings and source contacts before
its complementary part is bounded. More explicitly, the next attack must
supply a common domain and quantitative estimates for
\begin{equation}
 \mathfrak b_j C_j^{\ext},\qquad
 D_J^sD_c^t C_j^{\ext}\big|_{J=0},\qquad
 (I-\iota_j\Pi_{\loc,j})C_j^{\ext},
 \label{eq:v1-next-estimates}
\end{equation}
along with the corresponding finite-amplitude local-order remainder.
The scalar $\mathfrak b_jC_j^{\ext}$ is not set to zero merely because
it is absent from the local-box expansion.

An exponentially small exterior estimate would be sufficient on a branch
where it is true, but it is not imposed as a fact: additional compact wells,
sector multiplicities, or degeneracies can require retaining more reference
sectors instead. A proof must either control their normalized contribution
or enlarge the carrier to keep them. A global union bound with a volume
factor is not a substitute for the rooted source estimates in
\eqref{eq:v1-next-estimates}.

Only after this complete-shell interface and physical calibration are paid
can $q_{0,j}$ in \eqref{eq:v1-marginal-series} be called the internally
generated physical marginal shift. The current V1 neither supplies its
YM sign nor assumes an external $u_j\sim1/(bj)$ to prove it.

\section{Gate assessment, anti-circularity audit, and uniformity}
\label{sec:v1-assessment}

\begin{center}
\small
\begin{tabular}{>{\raggedright\arraybackslash}p{0.14\textwidth}>{\raggedright\arraybackslash}p{0.78\textwidth}}
\toprule
Gate & V1 assessment\\
\midrule
P1 & A necessary source-complete carrier, explicit rooted coefficient norms,
and an arbitrary-finite-order formal self-map are established. A fixed
invariant \emph{analytic} ball for the complete action is not established.\\
P2 & The complete compact correction is retained in
\eqref{eq:v1-exterior}. Its derivative estimates for generated interactions
remain the first analytic gap.\\
P3 & Weighted-jet QME compatibility is proved under the finite BV
hypotheses. An actual Ward-compatible renormalization projector and the
compact-block/BV identification remain open.\\
P4--P6 & The measured degree and normalization interface is explicit.
No signed physical YM marginal flow, generated stable graph, or uniformly
bounded composed physical source transport is claimed.\\
P7 & The finite coefficient compiler exists for every $K,d$, with spatial
bounds under explicit primitive stocks. Arbitrary-order uniform remainders
and renormalizable counterterm closure are not proved.\\
P8--P14 & Chiral line descent, derived adjacent-cutoff budgets, continuum
source limits, and scheme independence are not closed by V1.\\
\bottomrule
\end{tabular}
\end{center}

No running trajectory is used in the proof of the packet theorem. No bound
on composed transports is used to derive the source tower. No continuum Ward
identity is used in the finite BV weight calculation. No infinite-volume
locality is assumed to prove \eqref{eq:v1-connected-bound}. The absolute
bounds are applied only after the formal logarithm has selected connected
terms; the marginal extraction acts on the complete signed coefficient.
A source-only term is never hidden in a source-independent vacuum constant.
The formal graph theorem is never used as an analytic remainder theorem.

\begin{center}
\small
\begin{tabular}{>{\raggedright\arraybackslash}p{0.26\textwidth}>{\raggedright\arraybackslash}p{0.66\textwidth}}
\toprule
Variable & Exact scope of uniformity\\
\midrule
Volume & The source and compiler identities hold at every finite volume.
The connected norm bound has no volume factor under common primitive
stocks; vacuum terms use the explicit intensive normalization.\\
Cutoff & No cutoff-uniform bound is inferred unless the covariance,
chart, lift, and vertex stocks are themselves cutoff uniform. These are
not verified for the complete generated action here.\\
Perturbative order & Each fixed $K$ is covered. Combinatorial constants
may grow with $K$; no convergent or Borel-summable perturbation series
is asserted.\\
Source order & Every fixed finite source order fits a sufficiently large
$d$. The analytic inherited two-loop remainder remains restricted to
its proven mark order, at most four.\\
RG steps & Formal composition holds for every finite compatible sequence.
No norm estimate uniform in the number of steps is established.\\
Coupling domain & The quadratic-source logarithm is analytic on its specified
nonvanishing branch, with the stated positive real domain.
The marginal-series error is uniform only on domains with the displayed
bounded coefficients and remainder.\\
\bottomrule
\end{tabular}
\end{center}

\section{Final ledgers for V1}
\label{sec:v1-ledgers}

\subsection*{Proved}
\addcontentsline{toc}{subsection}{Proved}
The conditional source tower obeys the exact finite-order composition law
\eqref{eq:v1-total-cumulance}. The inherited positive Gaussian curvature
shell generates a strictly nonzero source-quadratic coefficient whenever
its corresponding quadratic observable is nonzero
(\cref{thm:v1-gaussian-contact}). The measured weighted packet
\eqref{eq:v1-main-packet} closes and composes at arbitrary prescribed finite
loop order (\cref{thm:v1-all-order-packet}); the primitive-degree and
constraint-degree tests exhibit why its higher jets are needed. Connected
spatial contractions satisfy \eqref{eq:v1-connected-bound}, yielding a
volume-independent finite-order coefficient estimate under its explicit
stocks. The finite BV defect descends to the packet under the stated
Darboux/weight hypotheses. Kinetic extraction gives
\eqref{eq:v1-marginal-series} from a complete calibrated expansion and
places the leading shift in its one-loop term.

\subsection*{Inherited}
\addcontentsline{toc}{subsection}{Inherited}
The exact compact pushforward and measured redundancy identities, positive
eliminated Gaussian, nonlinear root chart and inverse-Haar density, local
measured two-loop formula and four-mark local remainder, and the existing
two-loop composition packet are inherited from the cited locations in
\cite{YM7}. The finite BV pushforward principle and the gauge--matter,
stable-graph, and determinant-line interfaces retain the hypotheses and
scope stated in \cite{SM}. The supplied mass-gap notes are not used to
assert a perturbative continuum conclusion.

\subsection*{Conditional}
\addcontentsline{toc}{subsection}{Conditional}
A cutoff-uniform coefficient bound follows if the primitive stocks in
\cref{cor:v1-finite-locality} remain cutoff uniform on the generated class.
A physical marginal update through the displayed order follows from an
actual complete-shell calibrated expansion, not merely the local-box one.
A BV-compatible coefficient flow requires an admissible full field/antifield
fibration. An analytic action class requires preserved budgets for all
finite-amplitude tails, the normalized compact ratio, source coefficients,
and determinant/phase data. These assumptions have not been proved for the
full iterated SM--YM theory in this installment.

\subsection*{Open}
\addcontentsline{toc}{subsection}{Open}
The first load-bearing open theorem is the complete normalized one-shell
estimate \eqref{eq:v1-next-estimates} for the \emph{generated} interaction
class, including the physical kinetic jet of the compact correction and
all source contacts at the claimed order. It must be established before
claiming a signed internally generated YM trajectory or an invariant
analytic class. Canonical counterterm closure, arbitrary-order analytic
remainders, source-transport products, the chiral line cocycle on the
actual cutoff family, and continuum/scheme independence remain downstream.
No one of these is inferred from the formal packet theorem.

\begin{thebibliography}{9}
\bibitem{Programme}
Supplied programme, \emph{Perturbative SM--YM Continuum Closure Programme},
attachment \srcid{Pasted markdown(20260919-030347).md}.
The P1--P14 gates and four-ledger convention are taken from this document.

\bibitem{YM7}
\emph{Renewal Yang--Mills Closure Monograph}, V7,
attachment \srcid{Renewal_Yang_Mills_Closure_Monograph_V7(1).tex}.
Internal theorem/equation labels used in this note are given in
\cref{sec:v1-context} and at the point of use.

\bibitem{SM}
\emph{SM--Einstein Continuum Monograph}, updated supplied version,
attachment \srcid{SM_Einstein_Continuum_Monograph_Updated(3)(8).tex}.

\bibitem{Companions}
Supplied companion context:
\emph{Renewal Grand Tensor Monograph} V073(6);
\emph{Renewal Yang--Mills Mass Gap Research Notes} V59(1), V68(1),
V100\_f14(7), and V100\_f15(7).
These lineages are not reclassified as a proof of the complete perturbative
shell estimates required here.

\bibitem{Costello}
K.~J.~Costello, \emph{Renormalisation and the Batalin--Vilkovisky formalism},
arXiv:0706.1533 (2007). External methodological context only.

\bibitem{Anselmi}
D.~Anselmi, \emph{Background field method, Batalin--Vilkovisky formalism and
parametric completeness of renormalization},
Physical Review D \textbf{89}, 045004 (2014), arXiv:1311.2704.
External methodological context only.
\end{thebibliography}

\clearpage
\part*{V2: Global completion and exact defect-sector coordinates}
\addcontentsline{toc}{part}{V2: Global completion and exact defect-sector coordinates}

\section{The analytic interface, rather than another coefficient calculation}
\label{sec:v2-interface}

This installment preserves Sections~1--9 of V1. It attacks P1 and the first
complete-shell interface of P2. The finite coefficient theorem
\cref{thm:v1-all-order-packet}, its source tower, and the inherited two-loop
formula are not reproved. The issue is now which \emph{global} information
must accompany the local analytic germ, and how that information can be
converted into the actual compact correction.

There are three separate results. First, an explicit compact-group example
proves that the full-shell kinetic coefficient is not controlled by a small
root-chart analytic norm, even when an arbitrarily long finite jet and the
positive local Gaussian remain unchanged. Second, a two-norm interaction
space is constructed, together with exact connected-density multiplication
and reblocking. Third, the normalized compact ratio is given a canonical
hard-core representation by \emph{defect-sector cumulants}. Its activities
are derived from the actual integrals rather than identified with bad-set
probabilities. Their dependence on remote retained fields and sources is
retained by a support decomposition. A rooted norm of these decorated
activities controls the compact correction, its kinetic jet, and its
source contacts.

The last implication uses the connected hard-core estimate already present
in \cite{YM7}, \srcid{thm:four-mark-clustering}; it is not advertised as a
new cluster-expansion principle. What is new in this lineage is the exact
choice of activities for the ratio in \eqref{eq:v1-exterior}, the explicit
support information they require, and the global/local carrier on which
these operations make sense. Normed-algebra and localization constructions
in other rigorous RG methods provide methodological context
\cite{V2BrydgesSladeI,V2BrydgesSladeII}, not estimates for this gauge shell.

\subsection{A global-completion obstruction inside a compact group fibre}

The following is a test of the proposed \emph{class}, not a claim that its
adversarial perturbation has already been produced by Wilson blocking.
It shows why local information alone cannot imply the desired theorem on
a class allowing arbitrary globally completed analytic interactions.

Let $U\in\mathrm{SU}(2)$, with normalized Haar measure, and write
\[
 v(U)=1-\tfrac12\Tr U\in[0,2],\qquad q(U)=v(U)/2.
\]
Fix $\gamma>2$, $\eta>0$, and an integer $n\ge2$. For a retained scalar
$c$ put
\begin{equation}
 S_n(c,U)=\tfrac12c^2+v(U)-\gamma q(U)^n
                          +\gamma\eta c^2q(U)^n.
 \label{eq:v2-adversarial-action}
\end{equation}
Block by retaining $c$ and integrating $U$. Let $\chi$ be a fixed,
source- and $c$-independent nonnegative cutoff supported sufficiently
near $I$, positive on a neighbourhood of $I$. Define the full and local
integrals with weights $1$ and $\chi$, respectively, using the same source
factor $e^{Jq(U)}$ and the same action $S_n/h$.

\begin{theorem}[Local analytic smallness does not control the full kinetic jet]
\label[theorem]{thm:v2-global-obstruction}
For each prescribed finite Taylor degree $D$, there are arbitrarily large
$n$ for which the perturbation in \eqref{eq:v2-adversarial-action} has zero
Taylor jet through degree $D$ at $(c,U)=(0,I)$ and tends to zero in a fixed
small root-chart analytic coefficient norm as $n\to\infty$. The positive
eliminated Hessian at that point is unchanged. Nevertheless, with
$C_{n,h}^{\ext}=-\log(Z_{n,h}^{\rm full}/Z_{n,h}^{\loc})$,
\begin{equation}
 \lim_{h\downarrow0}h\,\partial_c^2 C_{n,h}^{\ext}(0,0)
       =2\gamma\eta,
 \qquad
 \lim_{h\downarrow0}\partial_J C_{n,h}^{\ext}(0,0)=-1.
 \label{eq:v2-obstruction-limits}
\end{equation}
These limits hold for each sufficiently large fixed $n$. Consequently no
bound on these full-shell jets that vanishes with the local perturbation
norm can hold uniformly over this class and all $0<h\le h_0$.
\end{theorem}
\begin{proof}
Use $U=\exp(i\sum_{a=1}^3 a_a\sigma_a)$, so that
$v=1-\cos\sqrt{a_1^2+a_2^2+a_3^2}$. In the coefficient $\ell^1$ norm on
$|a_a|\le R$, its power series obeys
\[
 \|q\|_R\le q_R:=\tfrac12\{\cosh(\sqrt3 R)-1\}.
\]
Choose $R$ with $q_R<1$. On $|c|\le R_c$ the perturbation has norm at most
$\gamma(1+\eta R_c^2)q_R^n$. Its first monomial in $a$ has degree $2n$.
Taking $2n>D$ proves the jet assertion; $n\ge2$ leaves the Hessian unchanged.
This coefficient bound also controls the fixed finite-component rooted
local norm used in V1.

At $c=0$ the action is $f_n(v)=v-\gamma(v/2)^n$. This function is strictly
concave for $v>0$, with endpoint values $0$ and $2-\gamma<0$. Its unique
global minimizing group element is $U=-I$. For a fixed local support
$v\le v_0<2$, take $n$ large enough that
$(\gamma/2)(v_0/2)^{n-1}\le1/2$. Then $f_n(v)\ge v/2$ there, and the
unique local minimum is $I$.

For completeness, concentration at a unique minimizer of a continuous
function on this compact space follows directly from its energy gap off
any neighbourhood of that minimizer. A smaller positive-measure
neighbourhood has energy within half that gap; comparison of its integral
with the outside integral gives an exponentially vanishing outside
probability. The same argument applies to the fixed local cutoff.
Thus, for fixed $n$, the full expectations of $q$ and $q^n$ tend to one,
while their local expectations tend to zero.

At $c=0$, $\partial_c S_n=0$ and
$\partial_c^2S_n=1+2\gamma\eta q^n$. Differentiating the finite positive
integrals therefore gives, for either integral,
\[
 h\,\partial_c^2(-\log Z)(0,0)=1+2\gamma\eta\E[q^n].
\]
Subtract the two equations. Similarly
$\partial_J C^{\ext}=\E_{\loc}q-\E_{\rm full}q$.
This proves \eqref{eq:v2-obstruction-limits}. Finally, first choose $n$
so that the root norm is as small as desired and then $h$ sufficiently
small. A claimed uniform vanishing bound is contradicted.
\end{proof}

The example respects compactness, conjugation invariance, a positive local
Gaussian, and the retained blocking coordinate. It does not contradict
V1: V1's formal compiler follows the declared local stationary branch.
Here the global integral follows a different branch. In particular the
missing compact term can change the \emph{classical} physical kinetic
coefficient, not only an exponentially small remainder. A norm of a root
germ cannot certify the global branch.

\begin{remark}[The analogous noncompact stability issue]
For a real Higgs coordinate, $\lambda H^4-\varepsilon H^6$ with
$\lambda,\varepsilon>0$ is arbitrarily close to $\lambda H^4$ on any fixed
small complex disc as $\varepsilon\downarrow0$, but
$\int_{\R}\exp[-\lambda H^4+\varepsilon H^6]\dd H=\infty$.
For sufficiently large $|H|$, the exponent is at least
$(\varepsilon/2)H^6$, proving divergence. Thus the SM carrier also needs
a global stability domain, not an unrestricted small ball of local
analytic action germs. This does not replace the inherited stable
one-sided Higgs result; it specifies which perturbations that result
cannot admit without additional global control.
\end{remark}

\section{A two-norm admissible interaction space}
\label{sec:v2-two-norm}

\subsection{Local algebras and support labels}

Fix a finite connected block graph $\Lambda$, with graph distance measured
in current-shell block units. Each block has a fixed finite number of
field and source species. For every nonempty finite support $X$, let
$\tau(X)$ be the length of a shortest connected graph containing $X$;
$X$ itself need not be connected. Supports are part of the representation,
so equality of two total actions does not automatically identify their
support decompositions. Redundant decompositions may be quotiented by the
infimum norm after a specific extraction convention is fixed.

For compact gauge variables on $X$, define the global source algebra
$\mathcal G_X(\sigma,\omega)$ by expansions
\[
 F_X(U,J,\theta)=\sum_{\beta,\zeta}
       F_{X,\beta,\zeta}(U)J^\beta\theta^\zeta,
 \qquad
 \|F_X\|_{\mathcal G_X}
 =\sum_{\beta,\zeta}\sigma^{|\beta|}\omega^{|\zeta|}
                         \sup_U|F_{X,\beta,\zeta}(U)|.
\]
The coefficients are continuous on the compact configuration space.
The $\theta$'s generate a finite exterior algebra with a fixed ordered
basis; the action and density activities are even. The source algebra
can instead be quotiented by source degree greater than a fixed $r$.
That quotient retains precisely source derivatives through order $r$;
it does not assert a full analytic source functional.

Let $\mathcal A_X(R,\sigma,\omega)$ be the local coefficient algebra in a
specified root chart $U=U(a)$, with norm
\[
 \|f_X\|_{\mathcal A_X}
 =\sum_{\alpha,\beta,\zeta}|f_{X,\alpha,\beta,\zeta}|
                   R^{|\alpha|}\sigma^{|\beta|}\omega^{|\zeta|}.
\]
Both algebras are unital Banach algebras. Multiplication is ordinary
coefficient convolution with the exterior signs retained before taking
norms. The triangle inequality and the factorization of the degree
weights prove submultiplicativity; completeness is that of weighted
$\ell^1$ with complete coefficient spaces.

A \emph{globally completed local interaction} is a compatible pair
$\Psi_X=(F_X,f_X)$ with $f_X(a,J,\theta)=F_X(U(a),J,\theta)$ on the real
root chart. Give this closed subspace the norm
\begin{equation}
 \|\Psi_X\|_{\mathcal B_X}
   =\max\{\|F_X\|_{\mathcal G_X},\|f_X\|_{\mathcal A_X}\}.
 \label{eq:v2-local-global-norm}
\end{equation}
Compatibility is closed because convergence in either norm implies
uniform convergence on every smaller real chart and source polydisc.
The product is compatible and submultiplicative. The global part of
\eqref{eq:v2-local-global-norm} excludes the perturbations of
\cref{thm:v2-global-obstruction} from a small global ball, even though
their local parts converge to zero.

Define the rooted interaction norm
\begin{equation}
 \|\Psi\|_{a,\mu}
  =\sup_{x\in\Lambda}\sum_{X\ni x}
        e^{a|X|+\mu\tau(X)}\|\Psi_X\|_{\mathcal B_X},
 \qquad a>0,\quad\mu\ge0.
 \label{eq:v2-interaction-norm}
\end{equation}
A scalar independent of all fields and sources may be stored separately
with intensive norm. A source-only coefficient still has its source
support and is governed by \eqref{eq:v2-interaction-norm}.

\begin{proposition}[Completeness and local derivative control]
\label[proposition]{prop:v2-complete-space}
The compatible support families with finite
\eqref{eq:v2-interaction-norm} form a Banach space. Its structural norm
inequalities do not depend on $|\Lambda|$. Taylor projection onto any
fixed set of multi-degrees is bounded by one on the local coefficient
part. For $R'<R$, $\sigma'<\sigma$, an order-$m$ field derivative and
order-$s$ source derivative are bounded on the local coefficient part by
\begin{equation}
 m!s!(R-R')^{-m}(\sigma-\sigma')^{-s}
 \label{eq:v2-derivative-loss}
\end{equation}
times the original norm, with the component/direction convention fixed.
No derivative bound at unchanged radii is asserted.
\end{proposition}
\begin{proof}
A Cauchy family converges in each $\mathcal B_X$. Fatou's inequality for
the positive support sums shows that the limit has finite rooted norm;
applying the same argument to the differences proves norm convergence.
The compatibility subspaces were shown to be closed. Taylor deletion
removes nonnegative terms from the local coefficient norm. To prove the
derivative bound, for a single variable use the coefficient inequality
\[
 \binom{k}{m}(R')^{k-m}(R-R')^m\le R^k,
\]
which is one summand of the binomial theorem. Multiply by $m!$, sum
coefficients, and apply the analogous argument to source variables.
Directional derivatives follow by multilinearity with the declared
$\ell^1$ component normalization. None of these steps sums over an
unanchored copy of $\Lambda$.
\end{proof}

The support norm is an interaction norm, not a norm for the unrestricted
ordinary product of two extensive actions. V1's counterexample
\cref{lem:v1-not-algebra} still applies. Density operations are performed
in the supported algebras below, and their connected output is then
measured in \eqref{eq:v2-interaction-norm}.

\subsection{Actual actions, formal coefficients, and the two remainders}

Let $h\in(0,h_0]$, without requiring holomorphy at $h=0$. At a fixed scale
record an actual reference density, actual blocking map, and globally
completed residual interaction $\Psi_h$. The full action
$A_h^{\rm full}$ is primary. On its root chart impose the exact matching
identity
\begin{equation}
 \operatorname{res}A_h^{\rm full}
 =c_h^{\rm vac}+\sum_{\ell=0}^{K}h^{\ell-1}S_\ell
                    +R_h^{\rm ord}+C_h^{\ext}.
 \label{eq:v2-exact-matching}
\end{equation}
Here $R_h^{\rm ord}$ is the local-order remainder, and $C_h^{\ext}$ is the
actual ratio correction for the declared local integral. They are not
identified. In particular a bound $R_h^{\rm ord}=O(h^K)$ does not assign
that order to $C_h^{\ext}$. Both are determined by
\eqref{eq:v2-exact-matching} and their defining integrals; they are not
extra factors to be multiplied into an already complete density.

The coefficient $S_\ell$ retains its V1 packet and its entire Taylor tail
$Z_\ell$. Its local norm is an analytic norm, not the global norm of the
full $h$-dependent density. The finite coefficients may contain relevant,
marginal, redundant, source-contact, and irrelevant monomials. A fixed
bounded extraction $\Pi_{\loc}$ with a specified right inverse $\iota$
gives the exact split
\[
 \lambda=\Pi_{\loc}A^{\rm full},\qquad
 A^{\rm full}=\iota\lambda+(I-\iota\Pi_{\loc})A^{\rm full}.
\]
Its target contains the physical action normalizations and the required
source contacts. The nonlinear normalization $u=\mathcal N(\lambda)$
is separate. Existence of a \emph{Ward-intertwining} local extraction is
not assumed merely because a bounded extraction has been fixed.

For the remainder functions a typical component norm is
\[
 \sup_{0<h\le h_0}h^{-K}\|R_h^{\rm ord}\|_{a,\mu;R',\sigma'},
\]
with its mark orders declared. For $C_h^{\ext}$ we instead use the
sector-activity norm of \cref{sec:v2-exterior-bound}. It carries an
explicit scale function $\epsilon(h)$ rather than a presumed power of $h$.

\subsection{Gauge--matter and line components}

For noncompact Higgs variables, replace the global supremum by a
\emph{graded} family of regulator norms
\[
 \|F_X\|_{{\rm gl},t}
 =\sum_{\beta,\zeta}\sigma^{|\beta|}\omega^{|\zeta|}
   \sup_{U,H}e^{-t\mathcal V_X(H)}|F_{X,\beta,\zeta}(U,H)|,
 \qquad \mathcal V_X=\sum_{x\in X}(1+|H_x|^2+|H_x|^4).
\]
For $t,s\ge0$, multiplication sends grades $(t,s)$ to $t+s$, because
$t\mathcal V_X+s\mathcal V_Y\le(t+s)\mathcal V_{X\cup Y}$.
This is a Banach \emph{scale}, not a single fixed-grade algebra for
arbitrary exponentiation. Actual bosonic integrability is a separate
condition, for example
\[
 \operatorname{Re}A_h^{\rm body}\ge
       \kappa_h\sum_x |H_x|^4-B_h|\Lambda|,
 \qquad \kappa_h>0,
\]
together with domination of the finitely many insertion coefficients at
the declared source order. An exponential is admitted only when this
density bound or another explicit integrable bound is proved. No estimate
for a general exponential follows from the local chart norm.

Keep the finite fermionic/ghost coefficients in the exterior algebra,
and keep each nonzero hard determinant/Pfaffian pivot with its inverse
margin. On a chosen perturbative component, line frames and adjacent maps
are recorded with their logarithmic field/source jets and inverse bounds.
A change of frame changes the scalar density representative and the line
map together. The cocycle, orientation, and Ward compatibility are exact
constraints or separately budgeted defects. They are not supplied by an
absolute determinant bound. Positivity-based comparisons later in this
installment apply to the real compact YM law, not to a general chiral
Berezin-integrated weight.

\subsection{An admissible carrier with explicit analytic and sector norms}

For fixed $K$, source order or radius, field radius, decay weights, and
structural reference, let $\mathfrak E_j^{K,r}(\mathbf M)$ consist of
states satisfying the following finite or graded family of requirements.
All bounds refer to their actual support decompositions.
\begin{center}
\small
\begin{tabular}{>{\raggedright\arraybackslash}p{.25\textwidth}>{\raggedright\arraybackslash}p{.66\textwidth}}
\toprule
Component & Type and required norm or constraint\\
\midrule
$u_j,\lambda_j$ & Physically normalized coordinates and their extracted
kinetic/relevant coefficients, in a declared finite-dimensional domain.\\
$z_j,S_{\ell,j},Z_{\ell,j}$ & Complete complementary support kernels;
local coefficient norms, V1 weighted degrees, and retained Taylor tails.\\
$\Psi_{j,h}$ & Actual full globally completed interaction, with
\eqref{eq:v2-interaction-norm}, or the stated matter regulator grades
and stability cone.\\
$R^{\rm ord}_{j,h}$ & Exact local-order remainder with the displayed
$h^{-K}$ norm and specified mark order.\\
$\mathcal J_j,\mathcal T_j$ & Full source coefficient tower, including
source-only contacts; one-step field/source maps with explicit radius
inclusion and their kernel norms. No product bound is stipulated.\\
$\mathcal L_j,\Phi_j$ & Fixed-component line data, nonzero pivot/inverse
margins, frame jets, and actual cocycle/phase defects.\\
$\mathcal D_j^{\rm BV}$ & Calculated master-equation defect, in the local
coefficient/marked norm; not an independently adjustable field.\\
$\kappa_{j,h}(D;E)$ & Exact decorated sector cumulants from the next
measured fibre, as defined below; norm
\eqref{eq:v2-sector-epsilon} bounded by a declared $\epsilon_j(h)$.\\
\bottomrule
\end{tabular}
\end{center}
The sector coordinates are constrained functions of the full density,
not free auxiliary activities. Local integral nonvanishing is included
in their domain of definition. The budgets $\mathbf M$ and
$\epsilon_j(h)$ must be proved for a proposed input; writing them in this
definition does not prove membership of the generated gauge theory.

The complete transformation is defined on the finite-regulator
nonvanishing/integrable domain and all of its outputs have specified
owners. The notation $\mathfrak C_{\rm pert}$ remains reserved, as in V1,
for a class whose \emph{fixed budgets} have been proved invariant. What
is constructed here is the admissible two-norm carrier
$\mathfrak E_j^{K,r}$ and its algebraic operations.

\section{Connected density coordinates and reblocking}
\label{sec:v2-density}

The reference action, including its kinetic and determinant normalization,
is kept outside the following small residual interaction. Smallness is
measured for the actual negative-log-density perturbation $\Psi_h$, not
silently for $h\Psi_h$. This distinction matters when a classical
interaction is divided by $h$.

\subsection{An exact conversion before integrating any field}

For a finite support family of even interactions define
\[
 f_X=e^{-\Psi_X}-1.
\]
Let a finite family of supports be overlap-connected when its graph of
nonempty intersections is connected. Put
\begin{equation}
 K(Y)=\sum_{\substack{\mathcal F\ne\varnothing\ \text{of distinct supports}\\
                  \mathcal F\ \text{overlap-connected},\ \cup\mathcal F=Y}}
                  \prod_{X\in\mathcal F} f_X.
 \label{eq:v2-mayer-activity}
\end{equation}
The word ``distinct'' refers to the support labels of this particular
interaction decomposition. Multiple interaction species on one support
can first be summed or kept as separate labels. For a finite volume this
is an exact finite expression, although the number of terms can be large.

\begin{theorem}[Exact density coordinates with a rooted bound]
\label[theorem]{thm:v2-density-conversion}
At fixed fields and sources,
\begin{equation}
 e^{-\sum_X\Psi_X}
   =\sum_{\Gamma\ \text{pairwise disjoint}}\prod_{Y\in\Gamma}K(Y).
 \label{eq:v2-density-gas}
\end{equation}
Suppose
\begin{equation}
 \epsilon:=\sup_x\sum_{X\ni x}
 e^{(a+1)|X|+\mu\tau(X)}\|f_X\|_{\mathcal B_X}<1.
 \label{eq:v2-mayer-small}
\end{equation}
Then the connected density family satisfies
\begin{equation}
 \sup_x\sum_{Y\ni x}e^{a|Y|+\mu\tau(Y)}
                             \|K(Y)\|_{\mathcal B_Y}\le\epsilon.
 \label{eq:v2-density-bound}
\end{equation}
In particular, if $\delta=\|\Psi\|_{a+1,\mu}$ and
$e^\delta\delta<1$, the right side may be taken to be $e^\delta\delta$.
All constants are independent of total volume. This theorem performs no
Gaussian or compact fibre integration.
\end{theorem}
\begin{proof}
Expand $\prod_X(1+f_X)$. Every chosen support family has a unique partition
into overlap-connected components. Distinct components have disjoint
unions, yielding \eqref{eq:v2-density-gas} with no factorization hypothesis
on any measure.

Write $w_X=e^{a|X|+\mu\tau(X)}\|f_X\|$. If a family is
intersection-connected, the union of trees inside its supports is
connected, whence
$\tau(\cup\mathcal F)\le\sum_{X\in\mathcal F}\tau(X)$ and
$|\cup\mathcal F|\le\sum |X|$. Its weighted norm is therefore at most
$\prod w_X$. Root an overlap spanning tree at a support containing a
fixed $x$. Counting all such trees overcounts each connected family.
A nonnegative rooted-tree majorant is given by increasing finite-depth
iterates of
\[
 T_X=w_X\exp\left(\sum_{Y:Y\cap X\ne\varnothing}T_Y\right).
\]
Start at $T_X^{(0)}=w_X$. If $T_Y\le w_Ye^{|Y|}$, then the sum in the
exponent is at most $|X|\epsilon\le|X|$ by
\eqref{eq:v2-mayer-small}. Induction bounds every iterate by
$w_Xe^{|X|}$. The exponential counts unordered children; allowing repeated
labels only increases the majorant. For a family of distinct labels, the
labeling factorials cancel, so each rooted spanning tree is counted with
at least its required unit weight. Summing roots through $x$ gives
\eqref{eq:v2-density-bound}.

Finally $\|f_X\|\le e^{\|\Psi_X\|}-1\le e^\delta\|\Psi_X\|$
by the local algebra bound and $\|\Psi_X\|\le\delta$.
\end{proof}

This is a volume-independent conversion of an actual globally bounded
small residual density perturbation. It does not show that the full
interacting reference, or all of its outgoing residuals, is small.
For the noncompact matter scale it applies at grade zero; other grades
need the explicit integrable regulator estimate described above.

\subsection{Products and geometrical reblocking have explicit costs}

Two densities in the form \eqref{eq:v2-density-gas} can be multiplied
without discarding interactions: color the polymers by their input
family and collect connected components of the combined overlap graph.
The same rooted-tree proof bounds the new activity at weight $a$ by
$\|K_1\|_{a+1,\mu}+\|K_2\|_{a+1,\mu}$ when this sum is less than one.
Within each color only disjoint configurations are allowed; dropping
that restriction in the nonnegative tree bound is an upper bound.
There is a size-weight reserve, not a same-radius contraction.

Here is the precise reblocking statement at $\mu=0$, which avoids any
unstated comparison between fine and coarse distances. Partition fine
blocks into coarse blocks containing at most $q$ fine blocks. Let $b(Y)$
be the coarse support of $Y$. Group any disjoint fine polymer
configuration by connected components of the overlap graph of $b(Y)$.
Let $K_b(X)$ be the sum of products in one such component of coarse union
$X$. Then the density is unchanged, now expressed as disjoint coarse
polymers, and
\begin{equation}
 q\|K\|_{a+1,0}<1
 \quad\Longrightarrow\quad
 \|K_b\|_{a,0}\le q\|K\|_{a+1,0}.
 \label{eq:v2-reblock}
\end{equation}
Indeed $|b(Y)|\le|Y|$, and anchoring a coarse block costs at most $q$
fine anchors. The proof of \cref{thm:v2-density-conversion} then applies
to the coarse-overlap graph. An estimate with spatial weight $\mu>0$
requires a stated comparison of the two metrics. This operation has not
yet integrated the fine fields inside a coarse polymer or rescaled them.
The factor $q$ and the weight reserve cannot be hidden in a claim of
uniformity over arbitrarily many steps.

\subsection{Why integrated scalar activities do not simply multiply}

For two disjoint field sets, even a Gaussian reference is generally not a
product. For centered jointly Gaussian $\xi_x,\xi_y$ of covariance $hC$,
\begin{equation}
 \Cov(e^{t\xi_x},e^{s\xi_y})
 =\exp\!\left[\tfrac h2(t^2C_{xx}+s^2C_{yy})\right]
                       \{e^{htsC_{xy}}-1\}.
 \label{eq:v2-gaussian-cross-support}
\end{equation}
The Gaussian generating function proves this equality. It is nonzero in
general despite the disjoint supports. Thus
$\E\prod K(Y)$ cannot be replaced by $\prod\E K(Y)$ in
\eqref{eq:v2-density-gas}. Integrating a polymer density creates connected
interactions between different incoming polymers. V1 determines their
finite formal coefficients; it does not supply the full analytic sum.

An exact two-support bound identifies the first required connecting
operation. Let $F$ and $G$ be smooth bounded functions on disjoint
Gaussian coordinate sets $X$ and $Y$, with bounded first derivatives.
Multiply the off-diagonal blocks $C_{XY},C_{YX}$ by $t\in[0,1]$ and
keep the diagonal blocks fixed. These interpolating covariances are
positive semidefinite because they are convex combinations of the full
and block-diagonal covariances. Gaussian differentiation yields
\begin{equation}
 \Cov(F,G)=h\int_0^1\sum_{x\in X,y\in Y}C_{xy}
                 \E_t[(\partial_xF)(\partial_yG)]\,\dd t.
 \label{eq:v2-cross-covariance}
\end{equation}
For nondegenerate covariances this follows by integration by parts from
$\partial_t\E_t H=\frac h2\sum_{ij}\dot C_{ij}\E_t\partial_i\partial_jH$;
degenerate cases follow by adding a positive diagonal regularizer and
taking its limit. Consequently its absolute value is at most
$h\sum_{x,y}|C_{xy}|\|\partial_xF\|_\infty\|\partial_yG\|_\infty$.
Higher connected integrations and the non-Gaussian reference need their
own bounds. Formula \eqref{eq:v2-cross-covariance} is not a uniform
all-order cumulant estimate.

\section{Canonical activities of the complete compact ratio}
\label{sec:v2-sector}

\subsection{A partition of the actual measured fibre}

At a fixed retained background and source, use the actual completed fibre
coordinates and its full density, including Haar and determinant factors.
Choose measurable local cutoffs $0\le\chi_x\le1$, independent of the
sources, such that the declared local contribution is
$\prod_x\chi_x$. They are fixed in the fibre coordinates when retained
parameters are varied. No global holomorphic hard cutoff is asserted.
If the chosen local integral does not have this factorization, it must
first be replaced by an explicitly compared local integral; that
comparison is another retained ratio, not a free change of scheme.

For $D\subseteq\Lambda$ define
\begin{equation}
 Z_D(c,J)=\int e^{-A(c,y,J)}
        \prod_{x\in D}(1-\chi_x(y))
        \prod_{x\notin D}\chi_x(y)\,\dd\nu_c(y).
 \label{eq:v2-sector-integrals}
\end{equation}
All measured background dependence of $\nu_c$ belongs to this formula.
For positive real YM data the integrals are nonnegative. In the
complex/fermionic setting the algebraic formulas require their defined
scalar representatives and a nonzero $Z_\varnothing$; they do not use
positivity. The exact identities are
\[
 Z^{\loc}=Z_\varnothing,\qquad Z^{\rm full}=\sum_{D\subseteq\Lambda}Z_D.
\]
Assume a common local field/source domain on which $Z_\varnothing\ne0$,
and put $r(D)=Z_D/Z_\varnothing$, $r(\varnothing)=1$.
For the analytic estimates below, holomorphy of these sector integrals
on that common domain is an additional requirement. Global continuity
plus analyticity near the all-good chart does not by itself prove
analytic dependence of sectors containing remote configurations.
For smooth cutoffs these are labelled partition contributions rather
than probabilities of disjoint events. For indicator cutoffs they are
literal sector integrals. The construction is valid in either case.

\begin{definition}[Defect-sector cumulants]
For $D\ne\varnothing$ define
\begin{equation}
 \kappa(D;c,J)=\sum_{\pi\in\mathfrak P(D)}
    (-1)^{|\pi|-1}(|\pi|-1)!\prod_{B\in\pi}r(B;c,J).
 \label{eq:v2-sector-cumulants}
\end{equation}
They are cumulants of the normalized \emph{sector array} $r$, not
ordinary cumulants of bad-set indicators under the full measure.
\end{definition}

\begin{theorem}[Exact normalized sector-to-polymer identity]
\label[theorem]{thm:v2-sector-identity}
For every finite regulator and every point of the stated nonzero local
integral domain,
\begin{equation}
 \boxed{\frac{Z^{\rm full}}{Z^{\loc}}
       =\sum_{\substack{\Gamma\ \text{pairwise disjoint}\\
                         D\ne\varnothing\ \text{for all }D\in\Gamma}}
                           \prod_{D\in\Gamma}\kappa(D).}
 \label{eq:v2-sector-gas}
\end{equation}
All nonempty subsets $D$ are allowed. The construction does not assume
that the fibre measure factorizes or that the activities vanish on
geometrically disconnected sets. The source dependence of every
normalizing denominator is retained.
\end{theorem}
\begin{proof}
In the commutative finite algebra with $t_x^2=0$, set
$R(t)=\sum_Dr(D)\prod_{x\in D}t_x$. Its constant coefficient is one.
Expanding its finite nilpotent logarithm shows that the coefficient of
$t_D$ is \eqref{eq:v2-sector-cumulants}: a partition into $k$ nonempty
blocks appears $k!$ times in $(R-1)^k$, and division by $k$ leaves
$(-1)^{k-1}(k-1)!$. Conversely exponentiating proves the coefficient
identity
\begin{equation}
 r(D)=\sum_{\pi\in\mathfrak P(D)}\prod_{B\in\pi}\kappa(B).
 \label{eq:v2-sector-inversion}
\end{equation}
Sum this identity over all $D$. A partition of an arbitrary subset is
exactly a family of pairwise disjoint nonempty subsets, proving
\eqref{eq:v2-sector-gas}. Only finite algebra was used.
\end{proof}

It is important \emph{not} to replace
$\log(Z^{\rm full}/Z^{\loc})$ by $\sum_D\kappa(D)$. Evaluation at
$t_x=1$ is not a homomorphism from the nilpotent algebra. Already with
one block, $\kappa=r(\{x\})$ but the ordinary logarithm is
$\log(1+\kappa)$, not $\kappa$. The ordinary hard-core logarithm, with
its repeated incompatible polymers, is needed in the next section.

\subsection{Remote source dependence is another retained coordinate}

The integrals in \eqref{eq:v2-sector-integrals} share all the remaining
interacting fields. There is no exact reason for $\kappa(D;c,J)$ to
depend only on $(c,J)$ in $D$. Even the derivative
\[
 \partial_{J_x}r(D)=r(D)
 \{\E_D O_x-\E_\varnothing O_x\}
\]
when the source is a linear insertion, need not vanish for $x\notin D$.
This formula is used only where $Z_D\ne0$; differentiating the quotient
directly covers zero sector integrals as well. Thus the missing
information is not repaired by assigning the label $D$ as the activity's
entire source support.

On a smaller common analytic domain, expand $\kappa(D)$ in the retained
field, source, and odd variables around the chosen reference point.
Group its Taylor monomials by their exact occupied block set $E$:
\begin{equation}
 \kappa(D;c,J,\theta)=\sum_{E\subseteq\Lambda}
                     \kappa(D;E;c_E,J_E,\theta_E).
 \label{eq:v2-decorated-activity}
\end{equation}
The $E=\varnothing$ term is the constant coefficient. Absolute local
coefficient convergence makes this an exact grouping; for fixed source
order the same construction is made in the source quotient algebra.
Each activity is labelled by $(D,E)$ and is carried on
$X(D,E)=D\cup E$. Its incompatibility rule is \emph{intersection of the
defect sets $D$}, not necessarily of $X(D,E)$. Different activities may
share source variables while remaining compatible. Expanding the sum
over $E$ in \eqref{eq:v2-sector-gas} gives this decorated hard-core gas
exactly. It cannot select two decorations of the same $D$, since that
would violate the defect incompatibility rule.

The decoration is not cosmetic. It records a generated boundary/source
interaction connecting the defects to the remaining retained theory.
Finite-volume analyticity gives this decomposition at a sufficiently
small radius; it does not give a cutoff-independent radius or spatial
norm. Those quantitative requirements are explicit below.

\section{A rooted norm for the actual exterior correction}
\label{sec:v2-exterior-bound}

Write $\|\kappa(D;E)\|_{R,\sigma,\omega}$ for its absolute analytic
coefficient norm on the stated local domain. Define
\begin{equation}
 \epsilon_{a,\mu}(h)
 :=\sup_{x\in\Lambda}\sum_{\substack{D\ne\varnothing,\ E\subseteq\Lambda\\
                                  x\in D\cup E}}
 e^{(a+1)|D\cup E|+\mu\tau(D\cup E)}
                  \|\kappa_h(D;E)\|_{R,\sigma,\omega}.
 \label{eq:v2-sector-epsilon}
\end{equation}
The additional one in the size weight pays the rooted tree summation.
The weighted span records both defect separation and remote retained
support. This is an input norm of explicitly constructed quantities,
not a bound inferred from a rare-event probability.

\begin{theorem}[Decorated sector norm controls the complete compact jet]
\label[theorem]{thm:v2-exterior-control}
Suppose the exact activities above have
$\epsilon_{a,\mu}(h)<1$ on the declared local field/source domain.
Then $Z^{\rm full}/Z^{\loc}$ has a nonvanishing branch connected to one
at zero activity. There is a support decomposition
\begin{equation}
 C_h^{\ext}=-\log(Z^{\rm full}/Z^{\loc})=\sum_{X\ne\varnothing}P_{X,h}
 \label{eq:v2-exterior-support}
\end{equation}
whose local coefficient norm satisfies
\begin{equation}
 \boxed{\sup_x\sum_{X\ni x}e^{a|X|+\mu\tau(X)}
                           \|P_{X,h}\|_{R,\sigma,\omega}
                 \le\epsilon_{a,\mu}(h).}
 \label{eq:v2-exterior-bound}
\end{equation}
The source- and field-independent constant is at most
$|\Lambda|\epsilon_{a,\mu}(h)$ in absolute value. Removing that constant
support by support does not increase the local coefficient norm.
Every fixed mixed kinetic/source jet at the reference point is bounded
by its coefficient-extraction norm times $\epsilon_{a,\mu}(h)$.
In particular, for the fixed kinetic extractor of V1,
\begin{equation}
 |\mathfrak b C_h^{\ext}|\le
               \|\mathfrak b\|\epsilon_{a,\mu}(h),\qquad
 \|(I-\iota\Pi_{\loc})C_h^{\ext}\|
 \le\|I-\iota\Pi_{\loc}\|\epsilon_{a,\mu}(h).
 \label{eq:v2-exterior-projections}
\end{equation}
The latter projections use the same local support norm and a declared
bounded section. The theorem holds uniformly in volume and in $h$
where its activity bounds and structural operator norms are uniform.
It is not asserted that the generated YM/SM states satisfy this bound.
\end{theorem}
\begin{proof}
Let $\alpha=(D,E)$, $X_\alpha=D\cup E$, and define
$q_\alpha=e^{a|X_\alpha|+\mu\tau(X_\alpha)}\|\kappa_\alpha\|$.
The hard-core logarithm is its connected Ursell series
\[
 \log\mathcal Z
 =\sum_{n\ge1}\frac1{n!}\sum_{\alpha_1,\ldots,\alpha_n}
       \varphi^T(\alpha_1,\ldots,\alpha_n)
                         \prod_{i=1}^n\kappa_{\alpha_i}.
\]
Repeated activities occur here. The hard-core connected-graph tree
bound used in \cite{YM7}, \srcid{thm:four-mark-clustering}, applies to
this incompatibility relation: the absolute Ursell coefficient is at
most the sum of trees whose edges are incompatible pairs. It depends
on the incompatibility graph, not on a claim that each $D$ is a connected
lattice set.

If $D_\alpha\cap D_\beta\ne\varnothing$, then
$X_\alpha\cap X_\beta\ne\varnothing$. Therefore
\[
 \sum_{\beta:\,D_\beta\cap D_\alpha\ne\varnothing}
             q_\beta e^{|X_\beta|}
 \le |D_\alpha|\epsilon_{a,\mu}\le |X_\alpha|.
\]
The same rooted-tree recursion as in
\cref{thm:v2-density-conversion} bounds the sum rooted at $\alpha$ by
$q_\alpha e^{|X_\alpha|}$. A connected incompatibility cluster has
intersecting full supports along its tree; hence the weight of their
union is at most the product of their weights. Assign that union as
the support $X$ of its coefficient in \eqref{eq:v2-exterior-support}.
To bound all clusters whose union contains a fixed $x$, select a member
whose full support contains $x$. The resulting overcount replaces
$1/n!$ by the rooted $1/(n-1)!$ and is bounded by
\[
 \sum_{\alpha:x\in X_\alpha}q_\alpha e^{|X_\alpha|}
                      \le\epsilon_{a,\mu}.
\]
This proves normal convergence in the support coefficient algebra and
\eqref{eq:v2-exterior-bound}.

Because the inequality is strict, multiply all activities by a complex
parameter $t$ of modulus slightly larger than one while keeping the
tree bound. Near $t=0$, the exponential of the convergent Ursell series
is the exact finite hard-core partition polynomial. Analytic
continuation establishes the equality throughout this disc; the
exponential cannot vanish. This fixes the logarithmic branch at $t=1$.
For the real positive original ratio it is the real branch, by
continuity from $t=0$ along the nonzero real polynomial.

Summing support coefficients without an anchor costs at most
$|\Lambda|$. Constant deletion contracts the absolute coefficient norm.
A mixed derivative at zero is bounded by its factorial and inverse
radius weights; at a smaller domain use
\eqref{eq:v2-derivative-loss}. The projection estimates are bounded
linear-operator inequalities applied to this same complete, signed
correction.
\end{proof}

The proof reuses the inherited hard-core bound. It does not reuse the
unproved hypothesis that suitable source-local activities exist with
small norm: \cref{thm:v2-sector-identity} constructs the activities,
\eqref{eq:v2-decorated-activity} constructs their support dependence,
and \eqref{eq:v2-sector-epsilon} is the remaining quantitative test.
In particular the geometry constants for a bare filling-contour graph
cannot simply be substituted into this test. The actual sets and their
retained support decorations must be counted.

For a source derivative with one fixed anchor, the bound contains the
local support norm and source radii, not the norm of an extensive action.
A source-only monomial with $E\ne\varnothing$ survives constant deletion.
The normalized functional $W[J]=\log Z[J]-\log Z[0]$ therefore retains
all of the contact coefficients required by V1.

\section{What global perturbation control does, and does not, buy}
\label{sec:v2-repair}

The following localized comparison transfers a \emph{given} real
reference repair to a new globally controlled interaction. It isolates
one useful consequence of the new global norm. It does not construct
the reference repair or the connected source estimate.

\begin{proposition}[Local-support stability of a sector comparison]
\label[proposition]{prop:v2-repair-tilt}
Use indicator sectors and a positive reference fibre measure.
For each $D$, suppose there is a measurable map $T_D$ from its sector
to the all-good sector such that
\[
 (T_D)_*\nu_D^0\le q^{|D|}\nu_\varnothing^0.
\]
This is domination of the \emph{unnormalized} measures, including the
actual repair Jacobian and reference action. Suppose $T_D$ changes
variables only in a fixed set $R_D$ with $|R_D|\le b|D|$.
For a real perturbation $\Psi=\sum_X\Psi_X$, put
\[
 \Omega(\Psi)=\sup_x\sum_{X\ni x}\operatorname{osc}\Psi_X,
\]
where each oscillation is over its complete global variable domain at
the fixed retained parameters. Then the normalized perturbed sector
weights satisfy
\begin{equation}
 \frac{Z_D^\Psi}{Z_\varnothing^\Psi}
                  \le\{q e^{b\Omega(\Psi)}\}^{|D|}.
 \label{eq:v2-repair-bound}
\end{equation}
There is no total-volume factor. For globally bounded interactions
$\Omega(\Psi)\le2\|\Psi\|_{a,\mu}$.
\end{proposition}
\begin{proof}
Only supports meeting $R_D$ can change under repair, so
\[
 |\Psi(y)-\Psi(T_Dy)|
 \le\sum_{X:X\cap R_D\ne\varnothing}\operatorname{osc}\Psi_X
 \le |R_D|\Omega(\Psi).
\]
Multiply $e^{-\Psi(T_Dy)}$ by the resulting exponential comparison,
integrate against $\nu_D^0$, and apply the pushforward domination. The
right side is $e^{b|D|\Omega}q^{|D|}Z_\varnothing^\Psi$.
Division proves the result. The last inequality follows from
$\operatorname{osc}\Psi_X\le2\sup|\Psi_X|$ and the weights being at
least one.
\end{proof}

This is the support-local counterpart of bounded-tilt comparison, not a
new entropy inequality. The companion notes already contain normalized
bounded-tilt estimates, for example
\srcid{f15v24:tilt-lemma} in \cite{Companions}. The point here is that a
localized repair pays only supports it changes, even for a nonlocal
interaction family with finite rooted global norm. The existence of a
suitable same-law reference repair for the next generated shell remains
an input to this proposition.

\begin{proposition}[Sector rarity alone does not bound connected activities]
\label[proposition]{prop:v2-rarity-not-clustering}
There are strictly positive sector arrays with $r(D)\le q^{|D|}$ for
every nonempty $D$ whose two-defect cumulants have no spatial decay and
whose rooted cumulant norm diverges with volume.
\end{proposition}
\begin{proof}
Fix $0<p<1$ and $q>0$, and set
$r(\varnothing)=1$ and $r(D)=p q^{|D|}$ for $D\ne\varnothing$.
These are positive finite sector masses after division by their sum,
with $Z_\varnothing$ restored as the normalizing reference mass.
For distinct $x,y$, \eqref{eq:v2-sector-cumulants} gives
\[
 \kappa(\{x,y\})=p(1-p)q^2.
\]
At zero spatial weight, the sum over pairs containing a fixed $x$ is
$(|\Lambda|-1)p(1-p)q^2$. It is unbounded as volume grows, despite the
stated smallness of every individual sector. At positive span weight
the failure is stronger. The corresponding exact ratio is
$1-p+p(1+q)^{|\Lambda|}$, showing explicitly the shared sector dependence.
\end{proof}

The last test concerns a nonlocal sector array; it is not a counterexample
to a proved spatial mixing theorem for the actual gauge action. It shows
why \eqref{eq:v2-repair-bound}, or a rooted bad-set probability estimate,
cannot by itself be called the activity estimate
\eqref{eq:v2-sector-epsilon}. The supplied F15 discussion likewise
separates bad-region probabilities from constructed normalized polymer
activities. V2 now supplies an exact representation, but the spatial
norm of its signed cumulants remains a substantive requirement.

\section{The complete update and the next measured estimate}
\label{sec:v2-update}

At a finite step in the declared domain, form the exact pushed density,
not a product of separately averaged incoming polymers. Its local
coefficient part is obtained from the inherited measured compiler;
the full compact correction is obtained from
\eqref{eq:v2-sector-integrals}--\eqref{eq:v2-sector-gas} with the same
sources and reference. Apply the retained extraction to the sum
\[
 A_+^{\rm full}=A_+^{\loc}+C_+^{\ext}.
\]
The source tower is extracted from this full expression. The determinant
frame and adjacent map are updated with their actual measured factors;
the BV defect is recomputed, not set to zero. Complementary interactions,
source contacts, and local-order remainders retain the owners specified
in \eqref{eq:v2-exact-matching}.

\subsection{An exterior budget feeds the physical normalization without a supplied flow}

Suppose the calibrated local two-loop expansion has
\[
 \mathfrak b A_+^{\loc}=u^{-1}+q_0+u q_1+r_{\loc},
 \qquad |r_{\loc}|\le C u^2,
\]
in the same coordinate convention as V1, and suppose
\cref{thm:v2-exterior-control} applies with $h=u$ at this particular
input. The actual extracted update is then
\begin{equation}
 k_+^{\rm full}=u^{-1}+q_0+u q_1+r_{\loc}+e_{\ext},
 \qquad |e_{\ext}|\le\|\mathfrak b\|\epsilon_{a,\mu}(u).
 \label{eq:v2-full-kinetic}
\end{equation}
The signed $q_0$ still comes from the complete measured one-loop local
coefficient; it is not replaced by the two-loop coefficient. If
$|k_+^{\loc}|\ge(2u)^{-1}$ and
$|e_{\ext}|\le(4u)^{-1}$, the exact reciprocal identity gives
\begin{equation}
 |u_+^{\rm full}-u_+^{\loc}|
 =\frac{|e_{\ext}|}{|k_+^{\loc}(k_+^{\loc}+e_{\ext})|}
 \le8u^2\|\mathfrak b\|\epsilon_{a,\mu}(u).
 \label{eq:v2-exterior-kinetic-shield}
\end{equation}
Thus an $O(u^2)$ sector budget is sufficient to keep this correction
inside the $O(u^4)$ physical-coupling error of the displayed two-loop
update. More generally its physical error is $O(u^2\epsilon(u))$.
This is the inherited reciprocal sensitivity applied to the newly
specified \emph{actual} compact contribution. Neither its smallness
nor a trajectory is being assumed to prove the local coefficient.

\subsection{The remaining estimate is now a concrete norm of exact integrals}

The next load-bearing analytic task is to derive, on the generated
interaction class and before choosing a marginal trajectory, a common
field/source domain and a bound
\begin{equation}
 \sup_x\sum_{D\ne\varnothing,E:\,x\in D\cup E}
 e^{(a+1)|D\cup E|+\mu\tau(D\cup E)}
 \left\|\left[\sum_{\pi\in\mathfrak P(D)}
 (-1)^{|\pi|-1}(|\pi|-1)!
 \prod_{B\in\pi}\frac{Z_B}{Z_\varnothing}\right]_E\right\|
 \le\epsilon(h)<1.
 \label{eq:v2-next-target}
\end{equation}
The bracket subscript means the exact retained-support Taylor grouping.
The $Z_B$ are the complete sector integrals, with the existing global
interaction, density, and sources left intact. This estimate must control
connected defect interactions and their remote retained/source tails;
a bound for each $Z_B/Z_\varnothing$ separately is not enough.

A proof of \eqref{eq:v2-next-target} with, for example,
$\epsilon(h)=O(h^2)$ would close the specific exterior kinetic/source
interface through the inherited two-loop order. It would still have to
be combined with preservation of the global interaction/stability
budgets, local-order remainders, Ward-compatible extraction, and
one-step source/line estimates to establish a complete self-map.
The finite coefficient compiler and the sector identity do not supply
these analytic inequalities.

The upstream options are now precise. One may control the connected
mixed-support cumulants in the normalized good reference, or construct
an exact local repair/decoupling expansion whose signed connected
activities satisfy \eqref{eq:v2-next-target}. A repair paying only a
rare-sector probability is not sufficient. No continuum beta coefficient,
externally prescribed $u_j$, or assumed product bound on composed source
transports has been inserted.

\subsection{Gate status and uniformity}

P1 has been strengthened from local coefficient bookkeeping to a
compatible global/local Banach-scale carrier, exact supported density
operations, and exact sector coordinates. P1's \emph{invariant-budget}
claim remains coupled to P2; it is not declared closed. P2 now has an
explicit exact activity norm that controls its compact ratio, but that
norm has not been bounded for the complete generated theory. The
nonfactorization and global-completion tests prevent two incorrect
routes to such a bound. P3--P14 retain the unclosed interfaces recorded
in V1; in particular no signed YM trajectory or SM ultraviolet
completion is asserted.

\begin{center}
\small
\begin{tabular}{>{\raggedright\arraybackslash}p{.24\textwidth}>{\raggedright\arraybackslash}p{.67\textwidth}}
\toprule
Variable & Scope of V2's estimates\\
\midrule
Volume & Rooted density, sector-log, and localized-repair bounds have no
volume factor under their displayed stocks. The removable constant
has an explicit $|\Lambda|$ factor. Identities hold at each finite volume.\\
Cutoff & Algebraic identities hold at every finite regulator. Analytic
bounds are cutoff uniform only when graph/species, radii, reference,
activity, and extraction stocks are cutoff uniform; this is unproved
for the generated theory.\\
Perturbative order & The sector and density identities are exact, not
truncated in $h$. Matching uses each fixed inherited packet $K$.
No all-order remainder or uniform bound in $K$ is obtained.\\
Source order & The full analytic coefficient theorem applies on the
stated common radius. Finite source quotients give exactly their fixed
order. Derivative constants contain factorials and radius losses.\\
RG steps & Product/reblocking have explicit reserve losses and a factor
$q$ in \eqref{eq:v2-reblock}. No fixed-budget estimate uniform in the
number of shell steps follows.\\
Coupling domain & Actual functions live on $(0,h_0]$. Bounds are uniform
only where $\epsilon(h)<1$ and all declared real stability/nonzero
reference conditions hold. Holomorphy at $h=0$ is not assumed.\\
\bottomrule
\end{tabular}
\end{center}

\begingroup
\renewcommand{\refname}{External methodological references for V2}
\begin{thebibliography}{9}
\bibitem[7]{V2BrydgesSladeI}
D.~C.~Brydges and G.~Slade, \emph{A renormalisation group method.
I. Gaussian integration and normed algebras}, arXiv:1403.7244 (2014),
Journal of Statistical Physics \textbf{159} (2015), 421--460.
Context for the separation of local algebras, Gaussian integration,
and interaction estimates; no gauge-shell result is imported.
\bibitem[8]{V2BrydgesSladeII}
D.~C.~Brydges and G.~Slade, \emph{A renormalisation group method.
II. Approximation by local polynomials}, arXiv:1403.7253 (2014).
Context for localization as a separate operation, not a substitute
for a Ward-compatible extraction in the present construction.
\end{thebibliography}
\endgroup

\section{Final ledgers for V2}
\label{sec:v2-ledgers}

\subsection*{Proved}
\addcontentsline{toc}{subsection}{V2: Proved}
The compact-group obstruction
\cref{thm:v2-global-obstruction} proves that arbitrary local analytic
smallness and unchanged finite jets do not control the complete kinetic
or source response. The compatible global/local support space is Banach
and has explicit derivative-radius bounds
(\cref{prop:v2-complete-space}). Small globally bounded residual
interactions have exact connected density coordinates with the
volume-independent bound \eqref{eq:v2-density-bound}; products and
geometrical reblocking have the stated reserve costs. Normalized sector
integrals admit the exact cumulant and decorated-support representation
\cref{thm:v2-sector-identity}. Applying the inherited connected hard-core
bound to those exact activities gives
\cref{thm:v2-exterior-control}, including kinetic, complementary, and
source-contact estimates. A localized real repair comparison survives
a global rooted perturbation with the cost
\eqref{eq:v2-repair-bound}. The explicit sector-array test proves that
rarity is not connected locality. Each analytic theorem retains its
stated hypotheses; none verifies the generated-theory activity stock.

\subsection*{Inherited}
\addcontentsline{toc}{subsection}{V2: Inherited}
V1's weighted coefficient compiler, exact source tower, connected formal
kernel estimates, finite BV filtration, and kinetic normalization are
unchanged. The actual compact block, root chart, local measured two-loop
expansion, four-mark hard-core clustering argument, and reciprocal
sensitivity are taken from their cited monograph locations. The
companion notes' distinctions between probability bounds, normalized
activities, and retained conditional correlations remain in force.
The stable Higgs, Berezin, and determinant-line interfaces retain their
original hypotheses.

\subsection*{Conditional}
\addcontentsline{toc}{subsection}{V2: Conditional}
The complete exterior correction is small in the required local topology
if its \emph{exact decorated sector cumulants} satisfy
\eqref{eq:v2-next-target} on a common domain. The real perturbation
comparison requires a same-law, locally supported reference repair;
it is not a chiral complex-measure theorem. Its sector-rarity bound alone
does not imply the activity estimate. A complete kinetic update follows
from the inherited local theorem plus that activity bound and physical
calibration. An invariant analytic RG class still requires preservation
of the global/stability budgets, the local-order remainder, and the
Ward/source/line components. Those requirements have not been proved
for the full generated SM--YM theory in V2.

\subsection*{Open}
\addcontentsline{toc}{subsection}{V2: Open}
The first remaining estimate is \eqref{eq:v2-next-target} for the actual
normalized shell and generated global interaction, with controlled
remote retained/source decorations. It must be proved without assuming
a running trajectory, independent scalar defect activities, or a
small root-chart norm as a global substitute. The next version should
attack the first connected two-defect/source term and its decoupling
on the same normalized fibre, then the higher connected expansion
needed to sum the decorated norm. Preservation of the global carrier
and physical one-loop normalization remain adjoining tasks. No
perturbative continuum limit, infinite-step analytic trajectory,
nonperturbative mass gap, or ordinary SM ultraviolet completion has
been claimed.

\clearpage
\part*{V3: Global fibre coercivity and normalized exterior insertions}
\addcontentsline{toc}{part}{V3: Global fibre coercivity and normalized exterior insertions}

\section{Upstream from the sector norm}
\label{sec:v3-start}

V2 gives an exact definition of the exterior activities but does not bound
those activities on a generated action. Assuming their smallness again
would not advance P2. This installment instead returns to the \emph{actual}
completed root block of \cite{YM7}, proves a global deterministic estimate
for that block, and uses it to construct and estimate normalized compact
insertions. The principal specialization is the positive real
$\mathrm{SU}(2)$ Wilson fibre. No chiral positivity assumption is made.

Three distinctions are essential. First, a coercive \emph{action} is not a
uniform diffusion gap. Second, a compact defect well outside a narrow core
has a different conditional price from a fluctuation immediately outside
that same core. Third, estimates at each fixed number of defects and source
marks do not sum an unrestricted analytic activity norm. All three
restrictions enter the statements, not only their final interpretation.

The new results are: a global plaquette-to-fibre inequality for the
completed block; a normalized, volume-independent price for arbitrary
sets of \emph{guarded} exterior defects, stable under explicitly bounded
global perturbations; a common-reference representation of separated
insertions; and actual connected two-defect bounds. For exponentially
quasilocal $C^2$ residual interactions the last bound is proved at zero
mark order, together with the remote first jet of a single insertion.
For finite-range smooth residuals, every fixed finite mark order is
covered. The transition collar between the narrow core and the outer
analytic box is retained explicitly; its resummation is the next gap.

\subsection{Source locations and nonduplication}

The completed product carrier, its positive local inverse-Haar factors,
and their fixed stencils are inherited from \cite{YM7},
\srcid{thm:global-pivot} and \srcid{cor:compact-product}. The local Hessian
floor and the common small-field germ come from
\srcid{sec:actual-one-loop}. The monograph expressly distinguishes the
auxiliary convex \emph{pivot-solving} potential from the Wilson action;
we do not use the former as a lower Hessian bound for the latter.

The companion result \srcid{f15v44:full-colour} controls a rooted
plaquette-coherence observable on a different boundary-pinned Wilson
problem; its roots can have length comparable to that box. The linear
quotient result \srcid{f15v11:coercivity} also has a different carrier.
Neither is relabelled as the local, nonlinear link-deviation inequality
proved below. The Witten/Helffer--Sj\"ostrand covariance method is standard
\cite{V3Helffer,V3Lo}; a self-contained finite-box proof of the particular
weighted estimate used here is included. It is applied to an interacting
\emph{normalized good law}, not merely to its Gaussian Hessian.

Throughout V3, $h=\beta^{-1}\in(0,1]$. This is the parameter of the input
measured action. No running trajectory, ultraviolet orientation, or
physical beta-function coefficient is assumed. All distances between
cells are in the current coarse-cell lattice units.

\section{A global inequality for the completed nonlinear root fibre}
\label{sec:v3-global-geometry}

\subsection{Fixed paths and a finite plaquette filling}

Use the parity cells and the tree clearing the highest active parity bit
from \cite{YM7}. The path from a cell root to
$\varepsilon\in\{0,1\}^4$ adds its active directions in increasing order.
Denote the corresponding direction word by $r(\varepsilon)$. On the tree
slice every transport along this tree is the identity. Each cell has
seventeen internal non-tree groups and twenty-eight nonpivot boundary
groups: these are its forty-five independent groups $Y_x$.

For a coarse edge $e=(x,i)$ let $B_e$ be its pivot and $C_e$ its fixed
coarse output. A boundary group in this family is indexed by
$\varepsilon_i=1$; the pivot has $\varepsilon=e_i$. The completed equation
is exactly
\begin{equation}
 C_e=B_e\exp X_e,\qquad
 X_e=\frac1{16}\sum_{\varepsilon\in\{0,1\}^4}
        g_\rho(B_e^{-1}P_{e,\varepsilon}),
 \label{eq:v3-pivot-equation}
\end{equation}
where the two duplicate root paths have $P_{e,0}=P_{e,e_i}=B_e$
and hence zero summands. The other fourteen products are the inherited
chronological two-link products, dressed by the root trees.

Let $d_G$ be the bi-invariant Hilbert--Schmidt geodesic distance on
$\mathrm{SU}(2)$. A fixed positive plaquette has holonomy $U_p$ and Wilson
energy $w_p=1-\frac12\Tr U_p\ge0$. Set $S_{\rm W}=\sum_p w_p$.
For this normalization,
\begin{equation}
 d_G(I,U_p)^2\le\pi^2 w_p.
 \label{eq:v3-metric-wilson}
\end{equation}
Indeed, for principal angle $0\le\theta\le\pi$ the two sides reduce to
$2\theta^2$ and $\pi^2(1-\cos\theta)$, respectively; use
$\sin(\theta/2)\ge\theta/\pi$.

\begin{lemma}[Finite path filling, with no commutativity assumption]
\label[lemma]{lem:v3-fillings}
There are nonnegative, fixed-stencil plaquette coefficients $a_{lp}$,
one row for each independent fine group $l$, such that
\begin{equation}
 \delta_l\le\sum_p a_{lp}d_G(I,U_p),\qquad
 \delta_l=\begin{cases}
 d_G(Y_l,I),&l\text{ internal},\\
 d_G(Y_l,C_e),&l\text{ a boundary group of }e.
 \end{cases}
 \label{eq:v3-link-rows}
\end{equation}
They can be chosen with
\begin{equation}
 \sup_l\sum_p a_{lp}\le9,\qquad
 \sup_p\sum_l a_{lp}\le180.
 \label{eq:v3-incidence-stocks}
\end{equation}
The sums count multiplicities and include plaquettes crossing cell
boundaries. The constants are independent of volume and of the group
configuration and coarse output.
\end{lemma}
\begin{proof}
Swapping two adjacent distinct positive coordinate steps replaces the
path transport by a plaquette holonomy conjugated by a prefix transport.
Successive swaps therefore express the discrepancy of the two path
transports as an ordered product of conjugates of plaquettes or their
inverses. Bi-invariance and the triangle inequality bound its distance
from the identity by the sum of their distances. No Abelian Stokes
identity or independent-plaquette representation is used.

For an internal link in direction $i$, $\varepsilon_i=0$. Move the last
$i$ in $r(\varepsilon)i$ into increasing order. There are
$k_>(\varepsilon,i)=\#\{j>i:\varepsilon_j=1\}\le3$ swaps.
The comparison path is the root tree to $\varepsilon+e_i$, so the
resulting discrepancy is the internal link itself.

For a boundary link, compare $r(\varepsilon)i$ with
$ii\,r(\varepsilon-e_i)$. If
$k=\sum_{j\ne i}\varepsilon_j$ and
$k_<=\sum_{j<i}\varepsilon_j$, at most $k+k_<\le6$ swaps are
needed. The target transport is $B_e$, because its other links are on
the two root trees. Thus $d_G(Y_l,B_e)$ has such a plaquette bound.
For a chronological path compare $r(\varepsilon)ii$ with
$ii\,r(\varepsilon)$. Moving the two added steps left costs at most
$2k\le6$ swaps and bounds $d_G(P_{e,\varepsilon},B_e)$.
All swaps stay in the fixed rectangular union of the two adjacent cells.

The actual cutoff satisfies $\|g_\rho(V)\|\le d_G(I,V)$ on the
whole group, including where it is zero. Hence
\[
 d_G(B_e,C_e)\le\|X_e\|
       \le\frac1{16}\sum_\varepsilon
                            d_G(P_{e,\varepsilon},B_e).
\]
For fixed $i$, $\sum_\varepsilon 2k=48$, so this contributes at most
three to each boundary row. Adding the direct boundary comparison
bounds that row sum by nine; the internal row sum is at most three.
This proves \eqref{eq:v3-link-rows} and its first stock.

For clarity, a volume-independent column bound can be obtained without
optimizing any individual column. Per coarse cell the internal direct
row masses sum to
$4(3+2+1+0)=24$. For boundary direct rows, in direction $i$ the sum is
$12+4(i-1)$, so the four directions total $72$. The twenty-eight
nonpivot boundary rows each receive the additional mass at most three,
for a further $84$. Their total is at most $180$. Translate this fixed
list by the coarse lattice. A given literal plaquette occurrence in
that list can land at a specified fine plaquette for at most one coarse
translation modulo the period. Summing its weights bounds every column
by the total list mass $180$. Coincident plaquettes on a periodic
quotient are counted with multiplicity and obey the same bound.
\end{proof}

\begin{theorem}[All-configuration nonlinear fibre coercivity]
\label[theorem]{thm:v3-global-coercivity}
For every compact input $Y$ and every coarse output $C$ of the actual
completed block,
\begin{equation}
 \boxed{\sum_{l\ \mathrm{independent}}\delta_l(Y,C)^2
          \le K_{\rm geom} S_{\rm W}(Y,B(Y,C)),\qquad
          K_{\rm geom}=1620\pi^2.}
 \label{eq:v3-global-coercivity}
\end{equation}
In particular, at $C=I$ the full compact Wilson action on the tree-fixed
fibre has a unique zero, $Y=I$, and
$S_{\rm W}\ge K_{\rm geom}^{-1}\sum_l d_G(Y_l,I)^2$.
This is a deterministic action estimate, not a diffusion spectral gap
or a normalized probability estimate.
\end{theorem}
\begin{proof}
Weighted Cauchy--Schwarz and \cref{lem:v3-fillings} give
\[
 \sum_l\delta_l^2
 \le9\sum_p\Bigl(\sum_l a_{lp}\Bigr)d_G(I,U_p)^2
 \le1620\sum_p d_G(I,U_p)^2.
\]
Apply \eqref{eq:v3-metric-wilson}. If $C=I$ and the action vanishes,
all independent groups are the identity. The unique solved pivots are
then also the identity by \eqref{eq:v3-pivot-equation} and the inherited
pivot uniqueness. Conversely this configuration has zero action.
\end{proof}

The proof is global in the group variables; no logarithm of an arbitrary
plaquette was inserted into a local expansion. For nontrivial $C$ the
boundary deviations are from their own $C_e$, not from a fictitious
simultaneously flat coarse connection. The estimate does \emph{not}
subtract the nonzero coarse minimum. It is consequently compatible with
nonzero coarse curvature and with positive-energy metastable compact
wells. No assertion about a retained massless field follows.

\subsection{An independently reproducible finite incidence check}

The companion script \srcid{check_root_fibre_fillings_V3.py} implements
only the direction-word swaps in \cref{lem:v3-fillings}, using integers.
It checks seventeen internal and twenty-eight boundary rows, the
$3,6,6$ upper counts for the three types of fillings, row mass at most
nine, and the analytical column bound $180$. This is a combinatorial
consistency check, not numerical evidence for an interacting estimate.
The proof above uses the explicit analytical bound, not any improved
constant printed by the script.

\section{An actual normalized price for guarded exterior sectors}
\label{sec:v3-sector-price}

\subsection{Two radii and every affected interaction}

Choose a small fixed outer radius $R$ inside the common root chart.
For one cell, let $G_s$ be the image of the product Lie-coordinate cube
of radius $s$ in its forty-five independent groups. Choose
$0<\varepsilon<R$ and define three disjoint measurable regions
\begin{equation}
 G=G_\varepsilon,\qquad T=G_R\setminus G_\varepsilon,
       \qquad B=(G_R)^c.
 \label{eq:v3-three-regions}
\end{equation}
Boundaries of these regions can be assigned arbitrarily, being Haar null.
Each cell in $B$ has at least one independent group at distance at least
$R$ from the identity. A cell in $G$ has all such distances at most
$c_G\varepsilon$, for a fixed coordinate constant $c_G$.
All coarse parameters below lie in a fixed coordinate cube with
$d_G(C_e,I)\le c_G\varepsilon$.

Write the complete positive real density as
\begin{equation}
 e^{-S_{\rm W}(Y,B(Y,C))/h-\Psi_h(Y,C,J)}j(Y,C)\dd H(Y),
       \qquad \Psi_h=\sum_X\Psi_{X,h}.
 \label{eq:v3-actual-density}
\end{equation}
This uses precisely the inherited global product Haar carrier and its
inverse-pivot density. The residual $\Psi_h$ is an actual negative-log
interaction, not a formal Taylor tail. Define its complete compact
oscillation stock, uniformly over the real parameter domain, by
\begin{equation}
 \Omega_h=\sup_x\sum_{X\ni x}
             \operatorname{osc}_{Y_X}\Psi_{X,h}.
 \label{eq:v3-osc-stock}
\end{equation}
No total-volume supremum of $\sum_X\Psi_{X,h}$ is used.

For a cell set $D$, let $\mathcal P(D)$ be \emph{all} plaquette terms of
the composed action whose actual independent-variable stencil meets $D$.
This includes every affected exterior plaquette. Put
$S_D=\sum_{p\in\mathcal P(D)}w_p$; the remaining action does not depend
on $Y_D$. There is a structural bound
$|\mathcal P(D)|\le b_P|D|$. There is likewise a constant $J_0$ such
that changing $Y_D$ changes $\log j$ by at most $J_0|D|$. This follows
by retaining all local inverse-Haar factors meeting $D$, using their
fixed incidence and individual positive upper/lower bounds.

\begin{lemma}[Guarded compact energy separation]
\label[lemma]{lem:v3-energy-price}
There is a structural constant $C_0<\infty$ such that, for all finite
$D$, all outside values $Y_{D^c}\in G^{D^c}$, and the stated small
coarse fields,
\begin{equation}
 \inf_{Y_D\in B^D}S_D(Y)
       -\sup_{Y_D\in G^D}S_D(Y)
 \ge |D|\left\{\frac{R^2}{4K_{\rm geom}}-C_0\varepsilon^2\right\}.
 \label{eq:v3-energy-separation}
\end{equation}
Here $\varepsilon$ is first restricted so that
$c_G\varepsilon\le R/2$. In particular it can be fixed, independently
of volume and $h$, so that the brace is a positive number
$\Delta\ge R^2/(8K_{\rm geom})$.
\end{lemma}
\begin{proof}
Apply the rows of \eqref{eq:v3-link-rows} only to the independent groups
in $D$. For a deep cell at least one of its deviations $\delta_l$ is
at least $R-c_G\varepsilon\ge R/2$. Thus the sum of their squares is
at least $R^2|D|/4$.

Some plaquettes in these rows need not belong to $\mathcal P(D)$.
Every such term is independent of $Y_D$; its stencil sees only small
outside groups and small coarse output. Smoothness of the solved
pivots on their fixed stencils, and their value $I$ at zero input,
give $d_G(I,U_p)^2\le C\varepsilon^2$. The sum of the row weights
in question is at most $9\cdot45|D|$. The same weighted
Cauchy--Schwarz argument as before therefore gives
\[
 R^2|D|/4\le K_{\rm geom} S_D(Y)+C_1\varepsilon^2|D|.
\]
This step has retained the unaffected plaquettes as an explicit debit;
it has not asserted that $D$ is action closed.

When $Y_D\in G^D$, every input in an affected stencil is small.
The same smooth local bound and $|\mathcal P(D)|\le b_P|D|$ show
$S_D(Y)\le C_2\varepsilon^2|D|$. Subtract this upper bound from
the preceding lower bound and absorb the fixed constants into $C_0$.
Finally take $\varepsilon/R$ small enough to obtain the claimed margin.
\end{proof}

\begin{theorem}[Normalized deep-sector price on the complete compact law]
\label[theorem]{thm:v3-deep-price}
Fix these radii and $\Delta>0$. Let
$b_\varepsilon=H^{\otimes45}(G_\varepsilon)>0$ be the product Haar mass of
one core cell. At every fixed outside core configuration the ratio of
the conditional integrals over $B^D$ and $G^D$, with the same action,
sources, and measured density, is at most $q(h)^{|D|}$, where
\begin{equation}
 \boxed{q(h)=b_\varepsilon^{-1}
              \exp\{J_0+\Omega_h-\Delta/h\}.}
 \label{eq:v3-q}
\end{equation}
Consequently, if $Z_{\varnothing,D}$ denotes the integral with exactly
$D$ in $B$ and every other cell in $G$, then
\begin{equation}
 0\le r^{\rm out}(D):=\frac{Z_{\varnothing,D}}{Z_{\varnothing,\varnothing}}
             \le q(h)^{|D|}.
 \label{eq:v3-actual-sector-price}
\end{equation}
This holds for arbitrary $D$, without a geometric connectedness
restriction. If $h\Omega_h\le\omega<\Delta$, the price is
$q(h)\le q_0e^{-(\Delta-\omega)/h}$, with
$q_0=b_\varepsilon^{-1}e^{J_0}$. The constants have no volume factor.
\end{theorem}
\begin{proof}
Hold $Y_{D^c}$ fixed in its core domain. All factors not meeting $D$
cancel in the ratio. The oscillation of the sum of the residual terms
meeting $D$ is at most $|D|\Omega_h$, by summing over their possible
anchors in $D$. The corresponding logarithmic density oscillation is
at most $J_0|D|$.

The numerator has product Haar mass at most one. Lower-bound the
denominator by its full core product mass $b_\varepsilon^{|D|}$
times the infimum of its density. Comparing the two extreme densities
using \cref{lem:v3-energy-price} proves the conditional ratio bound.
Multiplication by the identical outside density and integration over
$G^{D^c}$ proves \eqref{eq:v3-actual-sector-price}. Every denominator
has been kept; there is no change to an unnormalized probability law.
\end{proof}

This proves a real reference-sector estimate rather than assuming the
repair domination in \cref{prop:v2-repair-tilt}. It also admits genuinely
nonlocal residual interactions through \eqref{eq:v3-osc-stock}. The
residual may have size $O(h^{-1})$, provided its \emph{normalized-action}
oscillation $h\Omega_h$ stays below the fixed geometric margin. This
condition is not inferred for future generated actions.

\subsection{Why the two radii cannot be silently identified}

\begin{proposition}[A positive conditional Hessian does not give same-radius odds]
\label[proposition]{prop:v3-guard-necessary}
A uniformly positive finite-dimensional Hessian does not imply a
uniformly small ratio of outside-core to inside-core conditional masses
when the boundary variables range over that same core.
\end{proposition}
\begin{proof}
Take $V_h(x,y)=\{(x-2y)^2+y^2\}/(2h)$ on a fixed square containing
$[-3r,3r]^2$. Its Hessian is
$h^{-1}\left(\begin{smallmatrix}1&-2\\-2&5\end{smallmatrix}\right)$,
whose smaller eigenvalue is $(3-2\sqrt2)/h>0$. Fix the interior core
value $y=3r/4$. The conditional minimum is at $x=3r/2$, outside
$[-r,r]$. For small $h$, an interval of length $\sqrt h$ around that
minimum is inside the exterior region and has conditional integral
at least $c\sqrt h$ after the common $y$ factor is cancelled. The
core integral is at most $2r\exp[-r^2/(8h)]$. Their ratio diverges
exponentially. The example tests exactly the proposed conditional
inference; it is not substituted for the actual gauge measure.
\end{proof}

The guard in \cref{thm:v3-deep-price} is a mathematical condition, not a
notational choice. The transition region $T$ in
\eqref{eq:v3-three-regions} has not been assigned zero mass. Its exact
reintroduction is given in \cref{sec:v3-collar}.

\section{Connected integration in the actual normalized core law}
\label{sec:v3-good-law}

\subsection{Primitive smooth stocks, rather than assumed mixing}

Let $\Gamma_h=h\Psi_h$ be the normalized-action residual. For this
section its global compact coefficients have two continuous derivatives
in the independent groups and in the declared small \emph{real}
retained/source parameters $\theta$. Uniform bounds are taken over that
whole compact/small-parameter domain. A convenient stock is
\begin{equation}
 \mathcal M_{\nu,2}(\Gamma_h)=\sup_x\sum_{X\ni x}e^{\nu\tau(X)}
 \left[\operatorname{osc}_{Y_X}\Gamma_{X,h}
       +\sum_{k=1}^2\sum_{z_1,\ldots,z_k\in X}
                    \sup\|D_{z_1}\cdots D_{z_k}\Gamma_{X,h}\|\right].
 \label{eq:v3-smooth-stock}
\end{equation}
A derivative index comprises a cell and one of its fixed finite species;
compact derivatives are taken in fixed orthonormal invariant frames,
and parameter derivatives in the stated coordinate chart. The norm on
mixed tensors includes those finite species sums. A field/source-independent
constant is irrelevant to this stock and is stored separately. Source-only
coefficients are not deleted from the action. Coordinate conversions to
the small Lie chart have fixed bounds and are included in the constants.

Assume $\mathcal M_{\nu,2}(\Gamma_h)\le\zeta$ for a sufficiently small
fixed $\zeta$. This permits exponentially quasilocal interactions; no
fixed range is imposed here. It is an explicit primitive input
condition, satisfied by $\Psi_h=0$, not an assumed bound for any output
activity. For fixed higher marked orders we will additionally assume
bounded higher derivative stocks, as specified below.

The normalized core measure $\mu_G$ is the probability obtained from
\eqref{eq:v3-actual-density} by restricting every cell to $G$. In its
product Lie cube its full potential relative to Lebesgue measure is
\begin{equation}
 V_h(z,\theta)=h^{-1}f(z,c)-\ell(z,c)+\Psi_h(z,\theta).
 \label{eq:v3-good-potential}
\end{equation}
Here, specifically in the free-group coordinates $Y_l=\exp z_l$,
\[
 f(z,c)=S_{\rm W}(\exp z,B(\exp z,\exp c)),\qquad
 \ell(z,c)=\log j(\exp z,\exp c)+\sum_{l\ {\rm independent}}\log j_G(z_l).
\]
Thus every free-group Haar factor and the inverse-pivot Haar factor
is included once. No additional coarse-Haar factor is inserted: the
completed law already has density relative to coarse Haar. A change
from the inherited local fibre coordinates includes its ordinary
coordinate Jacobian; its zero-field tangent has fixed norm-equivalence
constants. The potential displayed here is the one in the actual
free-group product cube used to define $\mu_G$.

\begin{lemma}[A primitive core Hessian and commutator margin]
\label[lemma]{lem:v3-good-hessian}
Choose the radii and then $\zeta,h_0>0$ sufficiently small. There are
$m>0$ and $\mu\in(0,\nu)$, independent of volume and of
$0<h\le h_0$, such that on the entire core cube
\begin{equation}
 D_z^2V_h\succeq (m/h)I,\qquad
 s_\mu:=\sup_x\sum_y(e^{\mu d(x,y)}-1)
                 \sup_z\|h(V_h)_{xy}(z)\|<m/2.
 \label{eq:v3-hessian-stocks}
\end{equation}
These bounds are consequences of the inherited zero-field fibre floor,
finite stencils, and \eqref{eq:v3-smooth-stock}. They do not assume a
covariance estimate for $\mu_G$.
\end{lemma}
\begin{proof}
In the inherited local normalization the zero-field fibre Hessian has
a fixed positive lower bound (the stated value is $24/2327$ in its
orthonormal convention). If a fixed coordinate conversion is used,
reduce this lower bound by its fixed equivalence factor and call it
$m_0>0$. The finite-stencil third derivative row bounds imply
$\|f_{zz}(z,c)-f_{zz}(0,0)\|_{\op}\le C(\varepsilon+\|c\|_\infty)$.
Take the core and coarse radii so this is at most $m_0/4$.
The local density has a fixed second-derivative row bound.
The residual chart Hessian has row and column bounds at most
$C\zeta/h$. Taking $C\zeta\le m_0/8$ and $h$ sufficiently small
therefore gives the first inequality with, for example, $m=m_0/2$
after any further harmless reduction.

For the second inequality the finite-range part of $hV_{zz}$ is
uniformly bounded and its contribution tends to zero as $\mu\downarrow0$.
The remaining Hessian has an exponential row bound at some fixed
positive exponent no larger than $\nu$. Its commutator stock also tends
to zero uniformly: multiply the exponential row summands by
$(e^{\mu t}-1)e^{-\nu' t}$ with fixed $\nu'>0$ and
$0<\mu<\nu'$, whose supremum tends to zero. The symmetric block Hessian
gives the identical column control. Choose a single $\mu$ with the
stated strict margin. All estimates are in current-cell units.
\end{proof}

\begin{theorem}[Weighted covariance for the interacting core]
\label[theorem]{thm:v3-core-covariance}
Under \eqref{eq:v3-hessian-stocks}, for $C^1$ functions on the product
cube and cell anchors $a,b$, define
\[
 L_a(F)=\left[\sum_x e^{2\mu d(a,x)}
                         \sup\|\nabla_xF\|^2\right]^{1/2}.
\]
Then
\begin{equation}
 \boxed{|\Cov_{\mu_G}(F,G)|
       \le\frac{h}{m-s_\mu}e^{-\mu d(a,b)}L_a(F)L_b(G).}
 \label{eq:v3-core-covariance}
\end{equation}
For functions supported on sets $A,B$, the corresponding estimate is
\[
 |\Cov_{\mu_G}(F,G)|\le \frac{h}{m-s_\mu}e^{-\mu d(A,B)}
       \|\nabla F\|_{L^2(\mu_G)}\|\nabla G\|_{L^2(\mu_G)}.
\]
No Gaussian replacement of the core law occurs.
\end{theorem}
\begin{proof}
Use the nonnegative Neumann operator
$L=-\Delta+\nabla V_h\cdot\nabla$ in $L^2(\mu_G)$.
For centered $G$, solve $L\phi=G$ with Neumann boundary condition.
At each fixed finite cube, comparison of its strictly positive density
with normalized Lebesgue measure transfers the ordinary cube Poincare
inequality with a finite constant. Lax--Milgram on mean-zero functions
therefore gives the solution; no uniform constant from that existence
argument is used. Differentiating the equation gives the one-form
operator
\[
 \mathcal H=L\otimes I+D^2V_h,\qquad
 \nabla\phi=\mathcal H^{-1}\nabla G.
\]
Here a vector component normal to a cube face has zero trace, while its
tangential components have the corresponding Neumann condition. The
quadratic form is
\[
 \int\left\{\sum_{i,k}|\partial_k u_i|^2
                       +\langle u,D^2V_h u\rangle\right\}\dd\mu_G.
\]
It is coercive by $m/h$. Integration by parts on the flat faces gives
the differentiated identity and
$\Cov(F,G)=\langle\nabla F,\mathcal H^{-1}\nabla G\rangle$.
These are weak-form identities: smooth Neumann functions and the
associated vector form give them by density; edges of the cube have
zero face measure. Thus no unaccounted boundary curvature term is
introduced.

Let $W$ multiply the vector components in cell $x$ by
$e^{\mu d(x,b)}$. It is constant in the configuration variables,
preserves the vector boundary conditions, and commutes with $L$.
The difference $W\mathcal HW^{-1}-\mathcal H$ consists only of the
Hessian blocks multiplied by
$e^{\mu(d(x,b)-d(y,b))}-1$. The row and column Schur bounds in
\eqref{eq:v3-hessian-stocks} give its operator norm at most $s_\mu/h$.
A bounded-perturbation resolvent identity therefore yields
\[
 \|W\mathcal H^{-1}W^{-1}\|_{\op}\le h/(m-s_\mu).
\]
Finally
$\|W^{-1}\nabla F\|_2\le e^{-\mu d(a,b)}L_a(F)$ and
$\|W\nabla G\|_2\le L_b(G)$ by the metric triangle inequality.
This proves \eqref{eq:v3-core-covariance}. For supports $A,B$, use
$W_x=e^{\mu d(x,B)}$ instead. Approximation proves the stated $C^1$
version. The only inverse margin used in the estimate is the explicit
one-form margin, not a volume-dependent scalar existence constant.
\end{proof}

\subsection{A finite-order consequence, with its decay loss visible}

\begin{lemma}[Finite cumulants from the product-covariance estimate]
\label[lemma]{lem:v3-finite-cumulants}
Let $F_1,\ldots,F_n$, $n\ge2$, be bounded $C^1$ functions supported in
fixed-radius cell balls about $x_1,\ldots,x_n$. Write
$\|F_i\|_\infty\le A_i$ and
$\|\nabla F_i\|_{L^\infty(\ell^2)}\le A_i L_i$, with $A_i>0$.
Then
\begin{equation}
 |\Cum_{\mu_G}(F_1,\ldots,F_n)|
 \le C_n h\left(\prod_i A_i\right)\left(\sum_i L_i\right)^2
       e^{-\mu\tau(x_1,\ldots,x_n)/(n-1)}.
 \label{eq:v3-finite-cumulants}
\end{equation}
The constant depends on $n$, the support radius, and the margins in
\eqref{eq:v3-hessian-stocks}, but not on volume. No bound uniform in
$n$, and no decay exponent independent of $n$, is asserted.
\end{lemma}
\begin{proof}
Choose a minimum spanning tree for the anchors and remove one of its
longest edges. This partitions the anchors into nonempty groups $A,B$.
Every cross-group distance is at least that edge length, by the
exchange property of a minimum spanning tree. The length is at least
the spanning-tree length divided by $n-1$, which is at least
$\tau/(n-1)$. Enlarging anchors to their support balls changes the
separation by at most twice their radius.

Write the cumulant as its finite partition polynomial in moments.
Replace, for every index set $I$, its moment by
\[
 \E\prod_{i\in I\cap A}F_i\;\E\prod_{i\in I\cap B}F_i.
\]
The resulting cumulant is zero: these are the moments of the product
of the two marginal laws of the two groups. The difference of an old
and a new mixed moment is a covariance of two products. Their suprema
are bounded by the corresponding products of $A_i$, and their gradient
norms by those products times the sums of their $L_i$. Apply
\cref{thm:v3-core-covariance} across the chosen cut. Telescope the finite
moment products in each partition term and sum the absolute values of
the finite partition coefficients. This gives
\eqref{eq:v3-finite-cumulants}, including its stated loss of exponent.
\end{proof}

This does not contradict the inherited warning that a two-entry estimate
for a restricted observable list need not control all higher cumulants.
Here the covariance estimate holds for \emph{all products} of the
bounded local functions in question, with explicit derivative costs.
Even so, its finite-order consequence is too weak to sum the V2 activity
norm over an unbounded number of defects at one fixed decay exponent.

\section{Exact exterior insertions and their connected bounds}
\label{sec:v3-insertions}

\subsection{An ordinary cumulant representation when cores are separated}

For a single cell $x$, remove all action terms not depending on $Y_x$
from its conditional integral, and denote the remaining measured
conditional factors integrated over $G$ and $B$ by $K_x^G(v)$ and
$K_x^B(v)$. Here $v=Y_{x^c}$; all unchanged density factors cancel.
Define the actual conditional odds
\begin{equation}
 a_x(v;\theta)=K_x^B(v;\theta)/K_x^G(v;\theta).
 \label{eq:v3-conditional-odds}
\end{equation}
The numerator is positive and the denominator is nonzero for real data.
By \cref{thm:v3-deep-price}, $0\le a_x\le q(h)$ on every outside core
configuration. These odds are derived from the complete compact
conditional, not chosen as independent scalar activities.

\begin{proposition}[Same-reference identities for separated deep defects]
\label[proposition]{prop:v3-common-reference}
Suppose the full measured action has fixed finite interaction range
$R_s$. If no interaction stencil meets two distinct cells in a finite
set $D$, then
\begin{equation}
 r^{\rm out}(D)=\E_{\mu_G}\prod_{x\in D}a_x.
 \label{eq:v3-moment-identity}
\end{equation}
Thus, for such a pair $x,y$, its sector cumulant is exactly
\begin{equation}
 \kappa^{\rm out}(\{x,y\})
   :=r^{\rm out}(\{x,y\})-r^{\rm out}(\{x\})r^{\rm out}(\{y\})
   =\Cov_{\mu_G}(a_x,a_y).
 \label{eq:v3-exact-pair}
\end{equation}
For a separated set, its higher sector cumulant is likewise the
ordinary joint cumulant of the odds under this same core law.
\end{proposition}
\begin{proof}
Condition on $Y_{D^c}$. Under the stated stencil condition the
conditional action splits as a sum of terms, one for each $Y_x$,
plus an outside term. The product Haar measure and the complete local
density factors have the same split. Replace $K_x^G$ by $K_x^B$ by
multiplying by its exact ratio, and then integrate the unchanged
outside law. This proves \eqref{eq:v3-moment-identity}; partition
inversion proves the cumulant statements. In particular, the good
variables are not independent after the outside integration. Their
connected contribution is precisely the covariance in
\eqref{eq:v3-exact-pair}.
\end{proof}

\begin{lemma}[Fixed derivatives of the actual odds]
\label[lemma]{lem:v3-odds-jets}
In the finite-range case, assume uniform global/local smooth stocks
through order $s+1$ for the normalized action and the measured factors.
For any fixed collection of $k\le s+1$ outside-field or real parameter
derivatives,
\begin{equation}
 |D^k a_x|\le C_k q(h)h^{-k}.
 \label{eq:v3-odds-jets}
\end{equation}
The derivative vanishes outside a fixed neighbourhood of $x$. For two
nearby cells, the joint conditional odds
$a_{xy}=K_{xy}^{BB}/K_{xy}^{GG}$ instead satisfy
$|D^k a_{xy}|\le C_k q(h)^2h^{-k}$, with support in a fixed
neighbourhood of the pair. These constants are independent of volume.
\end{lemma}
\begin{proof}
Only finitely many incident stencils appear in each conditional action.
Every fixed derivative of that action has supremum at most $C/h$.
Differentiating a logarithmic conditional normalizer gives a finite
sum of cumulants of these score derivatives. At order $k$ their
absolute values are at most $C_k h^{-k}$ by the finite partition
formula and boundedness of the scores. Apply this to
$\log a_x=\log K_x^B-\log K_x^G$, and then use the finite
exponential differentiation formula and $a_x\le q(h)$.
The domains are fixed and compact, so no boundary derivative is missing.
The same proof uses the conditional $|D|=2$ price for $a_{xy}$.
Finite range proves the stated support assertions.
\end{proof}

\begin{theorem}[Guarded marked one- and two-defect estimates: finite range]
\label[theorem]{thm:v3-marked-pairs}
Fix a finite mark order $r\ge0$. On the primitive domain above, suppose
the residual has fixed finite range and fixed smooth stocks through
order $r+2$. Let $I$ be $r$ labelled real local parameter derivatives
at cells $z_1,\ldots,z_r$, including permitted coarse and insertion
source derivatives. There exist $C_r<\infty$, $\mu_r>0$ independent
of volume and $h\in(0,h_0]$ such that, for $x\ne y$,
\begin{align}
 |D_I r^{\rm out}(\{x\})|
   &\le C_r q(h)h^{-P_r}
                     e^{-\mu_r\tau(x,z_1,\ldots,z_r)},\\
 |D_I\kappa^{\rm out}(\{x,y\})|
   &\le C_r q(h)^2h^{-P_r}
                     e^{-\mu_r\tau(x,y,z_1,\ldots,z_r)},
 \label{eq:v3-marked-pair-bound}
\end{align}
where the nonoptimal choice $P_r=2r+4$ suffices. One may take a
positive $\mu_r$ no larger than $\mu/[2(r+2)]$. The corresponding
rooted sums over the other cell and the mark cells are bounded by
$C'_r q(h)^2h^{-P_r}$, or $C'_r q(h)h^{-P_r}$ for one defect,
at every smaller fixed spatial exponent.
\end{theorem}
\begin{proof}
For a parameter-dependent potential $V$ and observables $F_1,\ldots,F_n$
on a fixed domain, labelled differentiation of a cumulant gives
\begin{equation}
 \begin{split}
 D_I\Cum(F_1,\ldots,F_n)
  =\sum_{I_1\sqcup\cdots\sqcup I_n\sqcup R'=I}
     \sum_{\pi\in\mathfrak P(R')}(-1)^{|\pi|}
       \Cum\bigl(D_{I_1}F_1,\ldots,D_{I_n}F_n,
                                      D_B V:B\in\pi\bigr).
 \end{split}
 \label{eq:v3-cumulant-jets}
\end{equation}
For $n=1$ a one-entry cumulant means expectation. The empty partition
has one term. To prove the formula, apply the derivatives to
$\log\int\exp[-V(\theta)+\sum_i t_iF_i(\theta)]$ and extract the
coefficient with one mark of each $t_i$. Each original observable
receives its assigned parameter set; all remaining parameter marks
are partitioned into differentiated potential insertions. This also
proves the formula by induction from the normalized first-derivative
identity. The signs are fixed by $e^{-V}$.

For a separated pair apply \eqref{eq:v3-cumulant-jets} to
\eqref{eq:v3-exact-pair}. Every term has between two and $r+2$ entries
and contains two odds derivatives. The latter have amplitudes
$q h^{-k}$ and one further gradient cost at most $C/h$ by
\cref{lem:v3-odds-jets}. Each potential-score entry has amplitude
$C/h$ and bounded-range support near all of the marks differentiating
it. Apply \cref{lem:v3-finite-cumulants}. The product of amplitudes
costs at most $C_r q^2h^{-r}$, and the covariance-gradient factor
costs at most $C_r/h$. All the omitted individual marks are a bounded
distance from their parent entry, so reinstating them in the support
tree changes only $C_r$. The claimed weaker power and exponent follow.

For a single defect use $r^{\rm out}(\{x\})=\E a_x$ and the
$n=1$ formula. A term with no inserted score is supported near $x$
and has the direct odds bound; every other term is handled by the
same finite-cumulant estimate.

For a nonseparated pair use
$r^{\rm out}(\{x,y\})=\E_{\mu_G} a_{xy}$ and the just-proved
one-insertion argument for that bounded-size cluster. Differentiate
the subtracted product of the two single-defect expectations by
Leibniz. The two trees in a product can be joined by the bounded edge
from $x$ to $y$. The largest required crude power is at most
$h^{-(r+2)}$, hence is covered by $h^{-P_r}$. This proves the second
bound also at short separation.

Finally, on a four-dimensional periodic cell graph, balls have a
volume bound $C(1+t)^4$ independent of the period. For a tree spanning
an anchor and $k$ other marks, each of their distances from the anchor
is at most its tree length; consequently its length is at least the
sum of those distances divided by $k$. The spare positive exponent
therefore sums all the $k$ sites by a product of convergent exponential
row sums. This proves the rooted statements. Source species contribute
only fixed finite factors.
\end{proof}

The theorem is an actual compact-sector estimate for the bare completed
Wilson shell and for the specified finite-range deformations, uniformly
on its real small-parameter domain. It is not an estimate for a formal
Gaussian defect proxy. The constants may grow with $r$ and use the
fixed completion cutoff $\rho$; no limit $\rho\downarrow0$ is taken.

\subsection{A nonlocal residual does not require discarding the pair interaction}

The finite-range assumption is inappropriate for a general generated
interaction. At zero mark order it can be removed with a direct
comparison that keeps the same core measure throughout.

\begin{theorem}[Guarded two-defect locality with quasilocal global residuals]
\label[theorem]{thm:v3-nonlocal-pair}
Assume \eqref{eq:v3-smooth-stock} with
$\zeta\le\Delta/4$, reduced further if needed for
\cref{lem:v3-good-hessian}. Define
\begin{equation}
 q_*(h)=q_0e^{-\Delta/(2h)},\qquad
                    q_0=b_\varepsilon^{-1}e^{J_0}.
 \label{eq:v3-qstar}
\end{equation}
Then, for the exact sector array in
\eqref{eq:v3-actual-sector-price}, there are $C<\infty$ and
$\mu_*>0$, independent of volume and $h\in(0,h_0]$, such that
\begin{equation}
 \boxed{|\kappa^{\rm out}(\{x,y\})|
       \le C h^{-1}q_*(h)^2 e^{-\mu_*d(x,y)}\quad(x\ne y).}
 \label{eq:v3-nonlocal-pair}
\end{equation}
In particular its rooted pair sum is exponentially small in $h^{-1}$.
The primitive domain permits all support sets with the stated
exponentially summable global $C^2$ stock. It is not restricted to a
finite list of finite-range couplings.
\end{theorem}
\begin{proof}
At separation larger than the fixed bare stencil diameter, no Wilson
or measured-density term contains both cells. Split the residual into
terms meeting both cells and the rest. With all other core variables
$v$ held fixed, call the sum of the former terms in the negative-log
action $W_{xy}$, and let $U^0$ be the conditional action after removing
only those terms. Its $x,y$ variables factor. Define its two one-cell
odds $a_x^0(v),a_y^0(v)$. Removing a subset of the residual support
family does not increase the stock, so both odds are bounded by
\[
 q(h)\le q_0e^{-(\Delta-\zeta)/h}.
\]
Their weighted gradient norms about $x$ and $y$ are at most
$Cq(h)/h$. Indeed logarithmic differentiation of either conditional
ratio is a difference of two conditional scores. The bare incident
scores have a fixed stencil and size $C/h$; the weighted sum of the
residual scores is bounded by \eqref{eq:v3-smooth-stock}. This bounds
the weighted $\ell^1$ gradient, and therefore also the $\ell^2$
gradient used in \cref{thm:v3-core-covariance}.

Use the \emph{actual} full core law $\mu_G$, not the decoupled law,
in that covariance theorem. It follows that
\[
 |\Cov_{\mu_G}(a_x^0,a_y^0)|
                         \le Cq(h)^2h^{-1}e^{-\mu d(x,y)}.
\]
It remains to compare this covariance with the actual sector cumulant.
For fixed $v$, set
\[
 w(v)=\operatorname{osc}_{Y_x,Y_y}W_{xy}
            \le (\zeta/h)e^{-\nu d(x,y)}.
\]
Subtract a $v$-dependent scalar from $W_{xy}$ so that it lies in
$[0,w]$; this scalar cancels in every conditional ratio. Each of the
four conditional integrals $K^{GG},K^{BG},K^{GB},K^{BB}$ is between
$e^{-w}$ and $1$ times its decoupled value, apart from that common
scalar. Hence their three ratios with denominator $K^{GG}$ differ
multiplicatively from $a_x^0,a_y^0,a_x^0a_y^0$ by factors in
$[e^{-w},e^w]$.

Their expectations under the actual core marginal are exactly
$r^{\rm out}(\{x\}),r^{\rm out}(\{y\})$ and
$r^{\rm out}(\{x,y\})$. Taking the signed difference first yields
\begin{equation}
 |\kappa^{\rm out}(\{x,y\})-\Cov_{\mu_G}(a_x^0,a_y^0)|
       \le q(h)^2\{e^{w_*}+e^{2w_*}-2\}
       \le3q(h)^2 w_*e^{2w_*},
 \label{eq:v3-direct-pair-debit}
\end{equation}
where $w_*=(\zeta/h)e^{-\nu d(x,y)}$. Since $\zeta\le\Delta/4$,
$q(h)^2e^{2\zeta/h}\le q_*(h)^2$. Combine the two bounds and reduce
the spatial exponent if necessary.

For the bounded set of shorter separations, positivity and
\cref{thm:v3-deep-price} give
$|\kappa^{\rm out}(\{x,y\})|\le2q(h)^2\le2q_*(h)^2$.
Increase $C$ to cover these distances and use $h\le1$.
The rooted sum follows from the fixed-dimensional exponential row
sum, not from a total-volume supremum.
\end{proof}

\begin{corollary}[A remote first field/source jet on the same nonlocal domain]
\label[corollary]{cor:v3-first-jet}
When the stock \eqref{eq:v3-smooth-stock} includes the mixed compact
and real parameter derivatives, each declared local parameter component
$\theta_z$ satisfies
\begin{equation}
 |\partial_{\theta_z}r^{\rm out}(\{x\})|
                \le C h^{-1}q_*(h)e^{-\mu_*d(x,z)}.
 \label{eq:v3-first-jet}
\end{equation}
This includes coarse-field and ordinary real insertion-source
components with the stated bounded local primitive scores.
\end{corollary}
\begin{proof}
The single-cell identity $r^{\rm out}(\{x\})=\E_{\mu_G}a_x$
holds without a finite-range restriction. Exact differentiation gives
\[
 \partial_{\theta_z}r^{\rm out}(\{x\})
   =\E_{\mu_G}\partial_{\theta_z}a_x
                  -\Cov_{\mu_G}(a_x,\partial_{\theta_z}V_h).
\]
Terms independent of $Y_x$ cancel inside the conditional log-odds
score. The direct derivative is therefore bounded by
$Cq(h)h^{-1}e^{-\nu' d(x,z)}$, using the terms incident to both $x$
and $z$, with the fixed bare stencil included. The weighted gradient
of $a_x$ is at most $Cq(h)/h$. The weighted gradient of
$\partial_{\theta_z}V_h$ about $z$ is at most $C/h$ by the mixed
second-derivative stock. Apply \cref{thm:v3-core-covariance}, then
use $q(h)\le q_*(h)$. Every denominator and density score was
included before the estimate.
\end{proof}

The nonlocal theorem proves a genuine two-defect input and a first
remote jet on an explicit globally controlled class. It does not yet
give the nonlocal \emph{second} coarse-field jet needed for kinetic
extraction, all mixed higher jets, or the full analytic decoration
sum in \eqref{eq:v2-sector-epsilon}. The finite-range higher-mark result
is not silently applied to the nonlocal case.

\section{The transition collar and the exact remaining ratio}
\label{sec:v3-collar}

\subsection{No missing middle region}

Let $T,D\subseteq\Lambda$ be disjoint. In this subsection the letter
$T$ in a subscript denotes the set of cells assigned to the transition
region of \eqref{eq:v3-three-regions}. Define the actual ternary sector
integral
\begin{equation}
 Z_{T,D}(\theta)=\int e^{-S_{\rm W}/h-\Psi_h}j\,
       \prod_{x\in T}\mathbf1_{G_R\setminus G_\varepsilon}(Y_x)
       \prod_{x\in D}\mathbf1_{G_R^c}(Y_x)
       \prod_{x\notin T\cup D}\mathbf1_{G_\varepsilon}(Y_x)
                                                    \dd H(Y).
 \label{eq:v3-ternary-sector}
\end{equation}
All factors, including their parameter dependence, are those of
\eqref{eq:v3-actual-density}. The integration regions are fixed in the
independent product coordinates. Write
\[
 Z_G=Z_{\varnothing,\varnothing},\qquad
 Z_R=\sum_T Z_{T,\varnothing},\qquad
 Z_{\rm full}=\sum_{T\cap D=\varnothing}Z_{T,D}.
\]
Thus $Z_R$ is the same-law integral over the outer product box, not a
new Gaussian reference. The results in the preceding sections concern
$r^{\rm out}(D)=Z_{\varnothing,D}/Z_G$, with \emph{no} transition
cells elsewhere.

\begin{proposition}[Exact collar dressing of the binary sector array]
\label[proposition]{prop:v3-collar-dressing}
On the positive real domain, define
\begin{equation}
 F_D=\sum_{T\subseteq D^c}\frac{Z_{T,D}}{Z_G},\qquad
 \widehat r(D)=\frac{F_D}{F_\varnothing}.
 \label{eq:v3-dressed-array}
\end{equation}
Then $F_\varnothing=Z_R/Z_G$, $\widehat r(\varnothing)=1$, and
\begin{align}
 \frac{Z_{\rm full}}{Z_G}
       &=\frac{Z_R}{Z_G}\frac{Z_{\rm full}}{Z_R},
       &\frac{Z_{\rm full}}{Z_R}&=\sum_D\widehat r(D),
                                                        \label{eq:v3-two-ratios}\\
 C_G^{\ext}&=C_{G\to R}+C_R^{\ext},
       &C_{G\to R}&=-\log(Z_R/Z_G),\qquad
        C_R^{\ext}=-\log(Z_{\rm full}/Z_R).
                                                        \label{eq:v3-collar-split}
\end{align}
The cumulants of $\widehat r$ are exactly V2's sector cumulants for
indicator core $G_R$. In particular,
\begin{equation}
 \boxed{\widehat\kappa(\{x,y\})
     =\frac{F_{\{x,y\}}}{F_\varnothing}
       -\frac{F_{\{x\}}F_{\{y\}}}{F_\varnothing^2}.}
 \label{eq:v3-dressed-pair}
\end{equation}
Every mixed derivative of these identities differentiates all displayed
normalizers. Source-dependent constants are not removed.
\end{proposition}
\begin{proof}
The three regions partition each compact cell. Expanding their product
therefore gives the displayed finite sums for $Z_R$ and $Z_{\rm full}$.
For fixed $D$, summing over $T\subseteq D^c$ integrates all remaining
cells over $G_R$. Dividing by $Z_R$ is exactly the binary normalization
with core $G_R$, proving \eqref{eq:v3-dressed-array} and
\eqref{eq:v3-two-ratios}. Positivity permits the real logarithms and
proves \eqref{eq:v3-collar-split}. Applying the finite partition
inversion of \cref{thm:v2-sector-identity} gives
\eqref{eq:v3-dressed-pair}. Fixed finite compact domains justify each
of the finitely many differentiated identities under its stated smooth
stocks.
\end{proof}

This is not a claim that collar dressing is norm preserving. In
particular, $r^{\rm out}(D)$ is only the $T=\varnothing$ summand in
$F_D$. Deleting all other summands changes both the numerator and the
reference law. Positivity gives a lower bound from that deletion, not
the upper bound needed here. Even a bound on $F_\varnothing$ alone
does not control the mixed $T,D$ terms in $F_D$ or their signed
connected combination \eqref{eq:v3-dressed-pair}.

The first ratio in \eqref{eq:v3-collar-split} lies entirely inside the
analytic root neighbourhood. The inherited normalized wall comparison
is the appropriate available mechanism for that ratio, with its own
source-order, domain, and convex-extension hypotheses. V3 does not
silently upgrade the inherited four-mark wall theorem to an arbitrary
source-order theorem on all quasilocal perturbations. The ratio $Z_R/Z_G$ is not asserted to be close to one without the
stationary-section interior margin required by that wall comparison.
The second ratio still samples the compact exterior and requires the
dressed activities.
The same-radius obstruction in \cref{prop:v3-guard-necessary} prevents
reusing the guarded conditional price with the outside cube enlarged
from $G_\varepsilon$ to $G_R$ without a new argument.

\subsection{Which part of P2 has actually advanced}

The previous load-bearing assumption was an unspecified estimate on
exact sector cumulants. The following more primitive implication is
now proved for the actual completed positive Wilson shell:
\begin{equation}
 \begin{gathered}
 \text{fixed root-path geometry and small global smooth residual stock}
 \\ \Longrightarrow\quad
 \text{guarded normalized defect price and connected pair locality}.
 \end{gathered}
 \label{eq:v3-proved-implication}
\end{equation}
It includes all affected plaquette and density terms, allows nonlocal
residual interactions in the pair theorem, and does not assume a
running coupling, a rare-event polymer representation, or a mixing
estimate for the core. The core mixing estimate was obtained from the
primitive Hessian and weighted row bounds. This is a genuine upstream
input for a collar/contour expansion, rather than the assertion that
such an expansion is already small.

There are two subsequent analytic tasks, with neither implied by
\eqref{eq:v3-proved-implication}. First, the $T$ sums in
\eqref{eq:v3-dressed-array} must be reorganized and estimated so that
the complete signed cumulants, beginning with
\eqref{eq:v3-dressed-pair}, retain a useful decay weight and a compact
energy price. Second, those estimates must extend to an unrestricted
number of defects with degree-faithful constants and retained/source
support control sufficient for \eqref{eq:v2-sector-epsilon}.
The finite-cumulant bound \eqref{eq:v3-finite-cumulants} deliberately
loses its spatial exponent as the number of entries increases and
cannot supply this second step by summation.

For each fixed power $N$ and fixed $P$, the proved guarded factors
satisfy
\[
 \sup_{0<h\le h_0}h^{-N-P}e^{-c/h}<\infty\quad(c>0).
\]
Indeed, after $t=c/h$, the supremum is bounded by a constant times
$\sup_{t>0}t^{N+P}e^{-t}<\infty$ (the case of a negative exponent is
handled on $t\ge c/h_0$). Thus the guarded contributions just bounded
are smaller than every fixed perturbative power. This does \emph{not}
prove the same statement for $C_G^{\ext}$, $C_R^{\ext}$, or the
physical kinetic jet of either full ratio. The undischarged collar
and higher-decoration estimates cannot be replaced by the displayed
scalar exponential calculation.

The next local target may be stated without introducing a trajectory:
construct a same-law expansion of \eqref{eq:v3-dressed-pair} and its
second coarse-field and required source derivatives, using the
complete $Z_{T,D}$ and a fixed declared global residual domain. It
must retain the effect of arbitrarily many transition cells, rather
than put the surrounding field deterministically in the smaller core.
Such a theorem would permit the physical kinetic extractor of V1 to
act on a genuinely complete, not guarded, correction. Higher connected
orders and the preservation of the primitive global stocks remain
necessary for the full P2 self-map.

\section{Scope of the new domain and uniformity}
\label{sec:v3-scope}

The input domain of V3 is nonempty: it contains the bare completed
$\mathrm{SU}(2)$ Wilson shell and a sufficiently small neighbourhood
in the global normalized-action smooth stock
\eqref{eq:v3-smooth-stock}. Its quasilocal support family is not
restricted to the original Wilson plaquettes. However, we have not
proved that a previous measured shell lands inside the same small
stock domain, nor that the budgets of V2's larger global/local carrier
are equivalent to these smooth stocks. This is a specified class of
actual inputs, not an invariant class renamed after the proof.

The deterministic global estimate needs only the inherited compact
completion and the Wilson action. The positive conditional and
covariance theorems additionally use real weights. A determinant or
Pfaffian with an uncontrolled phase does not satisfy those hypotheses.
No extension to the chiral Standard-Model measure or to its line
transport is claimed. The separate Higgs stability and line-cocycle
requirements of V2 are unchanged.

The completion uses a smooth cutoff. A fixed number of real derivatives
of the compact sectors exists and is sufficient for the marked results
proved here. It is not a proof that every exterior sector extends
holomorphically to a common complex field/source polydisc. Consequently
none of the finite smooth-jet bounds is identified with the full
analytic decoration norm of V2. A real Taylor jet, an analytic radius,
and a uniform bound over an unbounded number of source marks are three
different assertions.

\begin{center}
\small
\begin{tabular}{>{\raggedright\arraybackslash}p{.23\textwidth}>{\raggedright\arraybackslash}p{.68\textwidth}}
\toprule
Variable & Precise scope in V3\\
\midrule
Volume & Global geometry, guarded price, core covariance, and rooted pair
bounds have no total-volume factor under their stated primitive stocks.
Every normalization is the actual finite-volume one.\\
Cutoff & The root geometry is fixed in current-cell units. The analytic
bounds are cutoff uniform only over families with common radii, species,
completion, Hessian, and global smooth stocks. Their preservation under
generated cutoff changes is not proved.\\
Perturbative order & Geometry and sector identities are exact in $h$.
Every fixed power is dominated by the guarded exponential price, but
this is not a full-shell remainder or all-orders continuum theorem.\\
Source order & At finite interaction range, any fixed $r$ is covered
under the required higher smooth stocks; constants and decay exponents
depend on $r$. For quasilocal $C^2$ residuals, the established results
are the unmarked pair and the single-defect first remote jet.\\
Number of defects & The conditional rarity bound holds for arbitrary
$D$. Connected estimates cover the stated fixed orders. A uniform
connected activity sum over all orders has not been proved.\\
RG steps & These are one-shell estimates on a specified input domain.
No repeated-step norm bound, invariant budget, or source-transport
product bound follows.\\
Coupling domain & $0<h\le h_0$, with fixed positive geometric guard
$\Delta$ and global normalized-action stock below its stated margin.
No holomorphy at $h=0$ and no supplied marginal recursion is used.\\
\bottomrule
\end{tabular}
\end{center}

P1 receives a concrete globally controlled smooth subdomain and a
new deterministic stability estimate. P2 receives actual normalized
insertion and connected-pair bounds on that domain. The complete
collar-dressed, source-decorated self-map remains unproved. The one-loop
normalization of P4 is unchanged: even a small full exterior correction
would not by itself compute the signed local marginal coefficient or
identify the upward ultraviolet recursion. P3 and P5--P14 retain their
previous unclosed Ward, stability, source, line, continuum, and scheme
interfaces. No nonperturbative mass-gap inference is made.

\begingroup
\renewcommand{\refname}{External methodological references for V3}
\begin{thebibliography}{10}
\bibitem[9]{V3Helffer}
B.~Helffer, \emph{Remarks on decay of correlations and Witten Laplacians:
Brascamp--Lieb inequalities and semiclassical limit}, Journal of
Functional Analysis \textbf{155} (1998), 571--586.
DOI: \srcid{10.1006/jfan.1997.3239}.
Methodological context for the one-form covariance representation;
no estimate for the completed compact root block is imported.
\bibitem[10]{V3Lo}
A.~Lo, \emph{Witten Laplacian Methods for the Decay of Correlations},
arXiv:\srcid{math-ph/0611002} (2006).
Methodological context for weighted covariance estimates; the finite
product-cube boundary conditions and the actual core law used in V3
are specified in \cref{thm:v3-core-covariance}.
\end{thebibliography}
\endgroup

\section{Final ledgers for V3}
\label{sec:v3-ledgers}

\subsection*{Proved}
\addcontentsline{toc}{subsection}{V3: Proved}
The actual completed root block satisfies the all-configuration
plaquette-to-fibre estimate \cref{thm:v3-global-coercivity}, with the
explicit structural constant $1620\pi^2$. The guarded affected-action
estimate and normalized sector price
\cref{lem:v3-energy-price,thm:v3-deep-price} hold for arbitrary defect
sets with all boundary and inverse-Haar factors included. A positive
Hessian alone does not justify a same-radius conditional price
(\cref{prop:v3-guard-necessary}). The primitive global smooth domain
has the stated core Hessian and weighted commutator margins; they
imply the interacting-core covariance theorem and its finite-cumulant
consequence. Actual compact conditional odds represent separated
sector moments under the same normalized core law. They yield
fixed-order marked pair locality for finite-range smooth residuals,
and the unmarked pair and single first remote jet for exponentially
quasilocal global $C^2$ residuals
(\cref{thm:v3-marked-pairs,thm:v3-nonlocal-pair,cor:v3-first-jet}).
The ternary sector sum and collar dressing are exact, including every
normalization in \eqref{eq:v3-dressed-pair}. Each assertion has the
scope stated in \cref{sec:v3-scope}; the guarded sectors are not the
complete compact ratio.

\subsection*{Inherited}
\addcontentsline{toc}{subsection}{V3: Inherited}
V1's source-complete coefficient compiler, source tower, and measured
one-loop kinetic convention, and V2's global/local carrier, exact
sector inversion, decorated activity norm, and conditional exterior
estimate are unchanged. The product Haar carrier, actual solved
pivots, inverse-Haar factors, local root germ, zero-field eliminated
Hessian floor, and local measured expansion are inherited from the
identified sections of \cite{YM7}. The local wall estimates keep their
original hypotheses and marked orders. The covariance method has the
external methodological provenance stated above; the application to
the present interacting core and its norm margins is proved here.
No theorem on a different boundary-pinned or massive-curl carrier is
substituted for this one.

\subsection*{Conditional}
\addcontentsline{toc}{subsection}{V3: Conditional}
Application to a generated effective action requires that action to
satisfy the stated global smooth, oscillation, and Hessian margins;
this membership is not established for arbitrary prior shell outputs.
The full binary compact correction is controlled by V2 only after
collar-dressed cumulants and their retained/source decorations satisfy
its activity norm. The new guarded pair estimate is an input toward,
not a proof of, that condition. Higher marked quasilocal estimates,
a common complex analytic domain, and the all-defect connected sum
remain separate requirements. A physical coupling update still needs
the complete normalized kinetic jet, physical calibration, and the
signed local one-loop term. Positive-real comparison gives no general
chiral phase or determinant-line estimate.

\subsection*{Open}
\addcontentsline{toc}{subsection}{V3: Open}
The smallest new unresolved interface is collar dressing of the actual
signed connected pair \eqref{eq:v3-dressed-pair}, with the second
coarse-field and required source derivatives, on a common globally
controlled quasilocal input class. The proof must sum the mixed
transition/deep sectors and their normalizer instead of replacing the
surrounding configuration by the narrow core. It must then extend to
higher connected activities with sufficient degree and spatial control
to sum \eqref{eq:v2-sector-epsilon}. Preservation of the primitive
global stocks under the measured shell, Ward-compatible local
subtraction, the generated marginal trajectory, composed source and
line transport, and continuum/scheme independence remain downstream.
A proof of the guarded estimate has not been reclassified as closure
of any of these interfaces.


\clearpage
\part*{V4: A time-dependent one-loop interface from the native Wilson kernel}
\addcontentsline{toc}{part}{V4: A time-dependent one-loop interface from the native Wilson kernel}

\section{What the V70 companion changes, and what it does not}
\label{sec:v4-start}

The newly supplied cumulative companion \cite{V4Companion} is used at its
actual scope. Its V68 identifies the horizontally reduced Gaussian reference
at stationary flat backgrounds; V69 computes a static plane-wave spectral
response; V70 identifies the frozen source Gram and explicitly leaves the
electric response and a controlled shell elimination open. These statements
do not estimate the mixed transition/deep sectors in
\eqref{eq:v3-dressed-array}. Sections~1--25 are preserved, and the collar
problem in \cref{sec:v3-collar} is not declared solved.

There is, however, a useful parallel attack on P4. The static trace cannot
be differentiated in time without specifying the two-endpoint kinetic
quadratic form and the moving hard subspace. This installment derives those
objects from the companion's \emph{native Wilson time-link action}, rather
than replacing it by an instantaneous spectrum. Its principal result is an
explicit two-jet formula for the time-dependent hard-fibre Gaussian
determinant, including the kinetic vertices. It proves an exact
finite-regulator coefficient statement for the specified local conditional
fibre. Identification with the full gauge-projected measured shell, and the
analytic complete-shell remainder, remain separate obligations.

\subsection{Source audit and limits on imported conclusions}

The source locations in \cite{V4Companion} are
\srcid{f16v68:reduction}, \srcid{f16v69:definition},
\srcid{f16v69:wellposed}, \srcid{f16v69:blocks},
\srcid{f16v69:erratum}, and \srcid{f16v70:decrement}.
The source's finite-momentum fit is not an exact beta coefficient, and its
three-dimensional transfer parameter $\gamma$ is not silently identified
with a four-dimensional root-block physical coupling. In this installment
$h=\gamma^{-1}$ only within the explicitly defined transfer-fibre integral.
No numeric slope, continuum Ward identity, or marginal trajectory is used.

Two restrictions matter. First, the horizontal Hessian identity at a
stationary flat background does not by itself identify the constrained
Hessian of a different nonlinear block at a nonstationary background. The
constrained Hessian and measured-density terms of \cite{YM7},
\srcid{eq:actual-constrained-Hessian} and
\srcid{eq:root-local-density}, remain required. We verify the needed
recentering order for the particular plane wave below; we do not assume
that a projection always removes it. Second, the reflection inequality in
\srcid{f16v69:erratum} does not support its subsequent asserted nonvanishing
threshold merely by substitution. At
$V/\gamma\asymp(\log V)^{-2}$ its displayed certified lower bound has order
$e^{-O(1)}/\log V$, which tends to zero, not order one. This is an audit of
that consequence, not a refutation of the scalar reflection inequality.
Neither the claimed threshold nor a large-volume operator-norm obstruction
is used here. The source's higher connected-cumulant assertion retains its
conditional status.

The new results below are the analytic hard-bundle construction, a verified
third-order force cancellation for this plane wave, the native temporal
quadratic form, and its exact frequency response. General finite Gaussian
integration, the one-mode Mehler identity, measured Schur ownership, and
V1's weighted compiler are inherited algebra, not claimed as new principles.

\section{An analytic hard bundle and a verified recentering cancellation}
\label{sec:v4-hard-bundle}

\subsection{Resolving the apparent rank change without numerical mode counts}

Work on a three-dimensional periodic lattice of side $L$, $V=L^3$, with
real link tangent space $\mathcal X=\R^{9V}$. Fix a transverse polarization
$\nu\in\{2,3\}$, color $n=e_3$, and
$k=2\pi m/L$ with $0<k<\pi$. The excluded endpoints ensure that the real
sine and cosine modes are independent. Put
\[
 b_{x\mu}=\mathbf1_{\mu=\nu}\cos(kx_1)n,
 \qquad \overline U_e(\alpha)=\exp(\alpha b_e).
\]
Quaternion coordinates are used, so that $[v,w]=2v\times w$ and
$1-\tfrac12\Tr e^v=1-\cos|v|$. The right-coordinate vertical derivative is
\[
 (G_\alpha\eta)_{x\mu}
   =\operatorname{Ad}(\overline U_{x\mu}(\alpha))^{\mathsf T}\eta_x
       -\eta_{x+\widehat\mu}.
\]
Let $G_\alpha^b$ be its restriction to vertex fields with $\eta_o=0$.
For the two constant vertex fields $e_1,e_2$ define
\[
 Z_a(\alpha)=\alpha^{-1}G_\alpha e_a\quad(\alpha\ne0),\qquad
 (Z_a(0))_e=-2\cos(kx_1)\mathbf1_{\mu=\nu}(n\times e_a).
\]
These divided columns have analytic continuations at zero. Let $E$ contain
the nine constant link fields, the cosine mode $b$, and its sine partner in
the same polarization and color. Form the rectangular frame
\begin{equation}
 F(\alpha)=[\,G_\alpha^b\mid Z_1(\alpha)\mid Z_2(\alpha)\mid E\,].
 \label{eq:v4-excluded-frame}
\end{equation}

\begin{theorem}[Analytic completion of the plane-wave hard subspace]
\label[theorem]{thm:v4-analytic-bundle}
For each fixed $L,k$ there is $a_L>0$ such that
\eqref{eq:v4-excluded-frame} has full column rank $3V+10$ on
$|\alpha|<a_L$. Hence
\begin{equation}
 P(\alpha)=I-F(\alpha)(F(\alpha)^*F(\alpha))^{-1}F(\alpha)^*
 \label{eq:v4-hard-projector}
\end{equation}
is a real-analytic orthogonal projector of rank $d=6V-10$.
For $\alpha\ne0$, its kernel is precisely the span of the full spatial
gauge directions and the eleven columns of $E$. At zero it additionally
retains in the excluded space the two limiting color-orbit modes specified
above. If $U_0:\R^d\to\mathcal X$ is any orthonormal frame of
$\operatorname{Ran}P(0)$, then
\begin{equation}
 U(\alpha)=P(\alpha)U_0
             (U_0^*P(\alpha)U_0)^{-1/2}
 \label{eq:v4-polar-frame}
\end{equation}
is analytic on a smaller interval and satisfies
\begin{equation}
 U^*U=I,\quad U_0^*U'(0)=0,\quad
 N:=U'(0)=P'(0)U_0,\quad U_0^*U''(0)=-N^*N.
 \label{eq:v4-frame-jets}
\end{equation}
No lower bound on $a_L$ independent of $L$ is asserted.
\end{theorem}
\begin{proof}
At zero, $G_0^b$ is the incidence derivative on based vertex fields.
A zero gradient is constant on the connected graph and its based value is
zero, so these $3V-3$ columns are independent. They are orthogonal to the
nine constant link fields. Each of the four remaining plane-wave columns
is divergence free, has mean zero, and is orthogonal to the other three:
two use distinct transverse colors and the last two use the sine and cosine
in color $n$. Their norms are nonzero for the stated $k$. Thus the frame
has rank $(3V-3)+9+4=3V+10$ at zero. Analyticity and the strictly positive
finite Gram determinant prove full rank on an interval.

Every vertex field is its based part plus a constant. The constant
$n$ field gives zero under $G_\alpha$, and the other two constants are
$\alpha Z_a(\alpha)$. This proves the stated span for nonzero $\alpha$.
Formula \eqref{eq:v4-hard-projector} is therefore analytic and has the
claimed rank, including at zero. The positive matrix
$U_0^*P(\alpha)U_0$ is close to the identity, so its inverse square root
is analytic. Direct multiplication proves $U^*U=I$ and
$U_0^*U(\alpha)=(U_0^*P(\alpha)U_0)^{1/2}$, which is symmetric.
Differentiating $P^2=P$ gives $P(0)P'(0)P(0)=0$, and hence the first
three identities in \eqref{eq:v4-frame-jets}. Differentiating $U^*U=I$
twice, using the symmetry of $U_0^*U''(0)$, gives the last identity.
\end{proof}

Let $H^{\rm amb}(\alpha)$ be the right-exponential Hessian of the spatial
Wilson action at $\overline U(\alpha)$ and define
\begin{equation}
 H(\alpha)=U(\alpha)^*H^{\rm amb}(\alpha)U(\alpha),\qquad
 H_0=H(0),\quad A=H'(0),\quad D_H=H''(0).
 \label{eq:v4-magnetic-jets}
\end{equation}
All these finite matrices are specified by plaquette derivatives and
\eqref{eq:v4-hard-projector}; a finite-difference definition of their jets
is unnecessary. At zero the curl Hessian on this hard space is positive,
with $H_0\succeq4\sin^2(\pi/L)I$. This follows by the transverse Fourier
decomposition of the discrete curl, after removal of constants and
gradients; removing additional plane waves can only increase its minimum
Rayleigh quotient. Thus $H(\alpha)$ stays positive on a smaller interval.
The small eigenvalue scale is displayed rather than called volume uniform.

\subsection{Why freezing at this plane wave gives the correct quadratic jet}

For a general background, a small but nonzero force in the eliminated
direction can change the quadratic one-loop response. The next statement
verifies its absence to the required order in the particular source used
in the companion.

\begin{theorem}[Third-order hard force for the abelian plane wave]
\label[theorem]{thm:v4-recentering}
Define the local spatial high-fibre action by
\[
 f(q,\alpha)=S_{\rm W}\bigl(\overline U_e(\alpha)
                                  \exp((U(\alpha)q)_e)\bigr).
\]
Then $f_q(0,\alpha)=O(\alpha^3)$ at each fixed regulator. Its nearby
stationary section satisfies $q_*(\alpha)=O(\alpha^3)$. Consequently the
two-jets at $\alpha=0$ of $f_{qq}(q_*(\alpha),\alpha)$ and $H(\alpha)$
agree. The same statement holds on a fixed periodic time cycle for the
native action
\begin{equation}
 \mathscr S(U)=\sum_t S_{\rm W}(U_t)
  +\sum_{t,e}\left(1-\tfrac12\Tr U_{t,e}U_{t+1,e}^{-1}\right),
 \label{eq:v4-native-cycle-action}
\end{equation}
with background $\overline U_t=\overline U(\alpha_t)$ and frame
$U(\alpha_t)$ on each slice: its hard force has no terms of total source
degree one or two in $(\alpha_t)_t$. The stationary shift has total degree
at least three. Thus the degree-two Gaussian determinant is computed from
the Hessian at zero hard fluctuation, with no unproved recentering omission.
\end{theorem}
\begin{proof}
At the abelian background all spatial plaquette angles are
$\alpha(db)_p$ in color $n$. A right variation perpendicular to $n$ has
zero first trace derivative. In color $n$, the exact gradient is
$d^*\sin(\alpha db)$, componentwise. Its linear part is
$\alpha d^*db=\alpha\lambda_k b$, and its remainder has degree at least
three because sine is odd. Since $U(\alpha)^*b=0$ exactly, the hard force
is $O(\alpha^3)$. The analytic implicit function theorem applies with
invertible hard Hessian $H_0$. Its degree-one and degree-two equations give
zero coefficients of $q_*$; equivalently $q_*=O(\alpha^3)$. Taylor expansion
of $f_{qq}$ proves the two-jet assertion.

For a time link set $\delta_{t,e}=(\alpha_t-\alpha_{t+1})b_e^0$, where
$b_e=b_e^0n$. The first derivative of its action is
$\sin\delta_{t,e}\,n\cdot(v_t-v_{t+1})$.
At a fixed time, its linear source part is a scalar multiple of the same
link vector $b$, namely $(2\alpha_t-\alpha_{t-1}-\alpha_{t+1})b$;
projection by $U(\alpha_t)^*$ annihilates it. The sine remainder starts
at degree three. The spatial force has already been treated. The reference
time-cycle Hessian is $H_0$ plus the nonnegative discrete time Laplacian,
so it is invertible. The finite-dimensional multivariable implicit
function theorem completes the proof.
\end{proof}

This proof is about a \emph{specified conditional high fibre}. It does not
identify it with the four-dimensional nonlinear root block, nor with an
integration over temporal gauge variables. For a measured local integral
with smooth leading density $\rho(\boldsymbol\alpha,q)$, the degree-two
one-loop coefficient is therefore
\begin{equation}
 \Gamma_1[\boldsymbol\alpha]
  =\tfrac12\log\det Q(\boldsymbol\alpha)
       -\log\rho(\boldsymbol\alpha,0)
       \quad\text{through total source degree two},
 \label{eq:v4-measured-coefficient}
\end{equation}
up to a source-independent Gaussian scalar. Indeed evaluation at $q_*$
changes neither displayed two-jet. The density term is \emph{not} set to
zero. The finite-regulator stationary expansion used here is V1's inherited
measured Gaussian coefficient rule, not a new uniform remainder theorem.

\section{The time-link quadratic form and the static normalization}
\label{sec:v4-native-quadratic}

\subsection{A two-endpoint calculation that an instantaneous spectrum misses}

Let $Jv=n\times v$, so $J^*=-J$ and $J^*J=I-nn^*$. For one native time
link at abelian relative angle $\delta$ write the endpoint fluctuations as
imaginary quaternions $v,w$.

\begin{lemma}[Exact quadratic jet of a native temporal plaquette]
\label[lemma]{lem:v4-time-link}
The Taylor polynomial through degree two in $(v,w)$ is
\begin{equation}
\begin{split}
 1-\operatorname{Re}\{e^{\delta n}e^ve^{-w}\}
  ={}&1-\cos\delta+\sin\delta\,n\cdot(v-w)\\
     &+\tfrac12\cos\delta\,(|v|^2+|w|^2)
         -v^*(\cos\delta I-\sin\delta J)w+O((|v|+|w|)^3).
\end{split}
 \label{eq:v4-time-link-jet}
\end{equation}
The remainder bound is uniform for $\delta$ in a fixed compact interval.
In particular the mixed kinetic vertex proportional to $\sin\delta J$ is
present before any Gaussian integration.
\end{lemma}
\begin{proof}
The scalar and vector parts of $e^ve^{-w}$ through degree two are
$1-|v-w|^2/2$ and $v-w-v\times w$. Multiplication by
$(\cos\delta,\sin\delta n)$ gives the displayed scalar part, because
$n\cdot(v\times w)=-v^*Jw$. Analyticity on the compact $\delta$ interval
bounds the third-order remainder. This is a finite quaternion identity,
not an expansion of a gauge-averaged kernel.
\end{proof}

Let $U_t=U(\alpha_t)$. On the periodic cycle of length $T\ge3$ define
ambient block-diagonal matrices
\[
 C_t=\operatorname{diag}_e\cos((\alpha_t-\alpha_{t+1})b_e^0)\otimes I_3,
 \quad
 S_t=\operatorname{diag}_e\sin((\alpha_t-\alpha_{t+1})b_e^0)\otimes J.
\]
The complete Hessian in \cref{thm:v4-recentering} is the real symmetric block
matrix assembled by
\begin{equation}
\begin{split}
 Q_{tt}&=H(\alpha_t)+U_t^*(C_{t-1}+C_t)U_t,\\
 Q_{t,t+1}&=U_t^*(-C_t+S_t)U_{t+1},\qquad
 Q_{t+1,t}=Q_{t,t+1}^*.
\end{split}
 \label{eq:v4-exact-temporal-Q}
\end{equation}
At zero $Q_0=H_0\otimes I_T+I_d\otimes(-\Delta_T)$, with temporal symbol
$H_0+2-2\cos\omega$. For a sufficiently small real background the
finite matrix remains positive. This yields an exact, defined Gaussian
integral and a logarithmic determinant, not merely a formal eigenvalue
prescription.

Changing frames by $U_t\mapsto U_tO_t$, $O_t$ orthogonal, conjugates $Q$ by
$\operatorname{diag}(O_t)$, provided $H_t$ is conjugated at the same time.
Thus its determinant is frame independent. In a nonorthonormal frame the
associated Jacobian belongs to $\rho$ in
\eqref{eq:v4-measured-coefficient}; it cannot be discarded before changing
coordinates. This is consistent with, but does not repeat or extend, the
stationary horizontal cancellation in the companion's V68.

\subsection{The temporal determinant has the companion's half-exponent}

Diagonalize $H_0$ with eigenvalues $\lambda_i>0$ and put
\begin{equation}
 \xi_i=\operatorname{arcosh}(1+\lambda_i/2),\qquad
 c_i(r)=\frac{e^{-\xi_i|r|}}{2\sinh\xi_i}.
 \label{eq:v4-xi-covariance}
\end{equation}

\begin{proposition}[Cyclic determinant and infinite-time covariance]
\label[proposition]{prop:v4-cyclic-mehler}
For each hard mode,
\begin{align}
 \prod_{j=0}^{T-1}(\lambda_i+2-2\cos(2\pi j/T))
    &=e^{T\xi_i}(1-e^{-T\xi_i})^2,\\
 \frac1{2T}\log\det Q_0
    &=\frac12\sum_i\xi_i+
           \frac1T\sum_i\log(1-e^{-T\xi_i}).
 \label{eq:v4-static-normalization}
\end{align}
The infinite-time covariance is diagonal in this basis with kernel
$c_i(r)$, and its Fourier symbol is
$C(\omega)=(H_0+2-2\cos\omega)^{-1}$.
In particular the one-mode vacuum energy is $\xi_i/2$, not $\xi_i$.
\end{proposition}
\begin{proof}
Factor
$2\cosh\xi-2\cos\omega
=e^\xi(1-e^{-\xi+i\omega})(1-e^{-\xi-i\omega})$.
The roots-of-unity product gives the first identity and then the second.
Summing the geometric series
$\sum_{r\in\mathbb Z}e^{-\xi|r|}e^{-ir\omega}
=\sinh\xi/(\cosh\xi-\cos\omega)$ proves the covariance formula.
Equivalently, the normalized Mehler kernel has positive Gaussian ground
state $e^{-\sinh\xi\,q^2/2}$ with eigenvalue $e^{-\xi/2}$ by a single
Gaussian integration. This recovers the companion normalization through
the exact temporal determinant.
\end{proof}

Only infinite \emph{time} at a fixed spatial regulator is taken here.
It is not a spatial continuum limit. The frozen source covariance in the
companion's \srcid{f16v70:decrement} is consistent with precisely
\eqref{eq:v4-xi-covariance}; no new source-capture conclusion is inferred.

\section{Exact temporal response of the native hard-fibre determinant}
\label{sec:v4-temporal-response}

This section gives scalar background response. Mixed temporal-source
derivatives for this fixed spatial profile follow from the resulting
quadratic functional. Other spatial or color profiles require their own
bundle and hard-force checks. The statement needs only the matrix data explicitly
derived above. All traces are finite hard-space traces.

Define ambient matrices and their hard compressions by
\begin{equation}
 \mathsf B=\operatorname{diag}_e(b_e^0)\otimes J,
 \qquad \mathsf D=\operatorname{diag}_e((b_e^0)^2)\otimes I_3,
 \qquad B=U_0^*\mathsf B U_0,\quad D=U_0^*\mathsf D U_0.
 \label{eq:v4-native-BD}
\end{equation}
Thus $B^*=-B$. Let
\[
 C_0=\operatorname{diag}_i\frac1{2\sinh\xi_i},\qquad
 C_1=\operatorname{diag}_i\frac{e^{-\xi_i}}{2\sinh\xi_i},
 \qquad \widehat\Omega^2=2-2\cos\Omega.
\]
The quadratic part of $W=\tfrac12\log\det Q$ is written
$\tfrac12\sum_\Omega|\widehat\alpha(\Omega)|^2\Pi(\Omega)$, with unitary
Fourier normalization. The infinite-time symbol below is the limit of the
corresponding finite-cycle response.

\begin{theorem}[Native one-loop temporal response]
\label[theorem]{thm:v4-native-response}
Let $H(\alpha)$ and $U(\alpha)$ have the jets
\eqref{eq:v4-magnetic-jets} and \eqref{eq:v4-frame-jets}, with $H_0>0$.
Set
\begin{align}
 \mathcal T={}&\Tr(C_1N^*N)+2\Tr(C_1N^*\mathsf B U_0)
                         -\Tr((C_0-C_1)D),
 \label{eq:v4-tadpole-T}\\
 V(\omega,\Omega)={}&A+
                   2\{\cos\omega-\cos(\omega+\Omega)\}B.
 \label{eq:v4-frequency-vertex}
\end{align}
Then the complete determinant response of
\eqref{eq:v4-exact-temporal-Q} is
\begin{equation}
\boxed{\begin{aligned}
 \Pi(\Omega)={}&\tfrac12\Tr(C_0D_H)+\widehat\Omega^2\mathcal T\\
 &-\tfrac12\int_{-\pi}^{\pi}\frac{\dd\omega}{2\pi}
   \Tr\!\left[C(\omega+\Omega)V(\omega,\Omega)
                    C(\omega)V(\omega,\Omega)^*\right].
\end{aligned}}
 \label{eq:v4-native-response}
\end{equation}
At finite $T$, the same formula holds for $\Omega\in(2\pi/T)\mathbb Z$
with the integral replaced by the discrete mean, $C_0$ by the discrete
mean of $C(\omega)$, and $C_1$ by the discrete mean of
$\cos\omega\,C(\omega)$. There is no small-frequency approximation in
this identity.
\end{theorem}
\begin{proof}
For a real direction of the background, normalized Gaussian differentiation
is the exact matrix identity
\[
 D^2W=\tfrac12\Tr(CQ'')-\tfrac12\Tr(CQ'CQ'),\qquad C=Q_0^{-1}.
\]
We compute both terms before taking absolute values. Write
$U_t=U_0+\alpha_tN+\alpha_t^2U''(0)/2+O(\alpha_t^3)$ and
$\delta\alpha_t=\alpha_t-\alpha_{t+1}$. From
\eqref{eq:v4-frame-jets},
\[
 U_t^*U_{t+1}=I-\tfrac12(\delta\alpha_t)^2N^*N
                         +O(|\boldsymbol\alpha|^3).
\]
Also $C_t=I-(\delta\alpha_t)^2\mathsf D/2+O(|\boldsymbol\alpha|^4)$ and
$S_t=\delta\alpha_t\mathsf B+O(|\boldsymbol\alpha|^3)$.
The degree-one off-diagonal block of $Q$ is consequently
$\delta\alpha_t B$; its transpose in the opposite block is
$-\delta\alpha_t B$. Acting on a mode $e^{i\omega t}$ with a background
$e^{i\Omega t}$ gives exactly \eqref{eq:v4-frequency-vertex}.
The diagonal vertex is $A$. This yields the bubble integral in
\eqref{eq:v4-native-response}, including the two time-link vertices.

For clarity, the degree-two terms of the quadratic action, per time edge,
are the following three contributions:
\begin{align*}
 &\tfrac12(\delta\alpha_t)^2q_t^*N^*Nq_{t+1},\\
 &-\tfrac14(\delta\alpha_t)^2
       (q_t^*Dq_t+q_{t+1}^*Dq_{t+1})
                  +\tfrac12(\delta\alpha_t)^2q_t^*Dq_{t+1},\\
 &\delta\alpha_t\,q_t^*
     (\alpha_tN^*\mathsf B U_0+\alpha_{t+1}U_0^*\mathsf B N)q_{t+1}.
\end{align*}
The first is frame motion, the second is the cosine curvature of the
native time link, and the third is their sine/frame cross term.
Since $U_0^*\mathsf B N=-(N^*\mathsf B U_0)^*$ and the one-step covariance
$C_1$ is symmetric, their Gaussian expectations sum to
$\tfrac12(\delta\alpha_t)^2\mathcal T$.
The magnetic quadratic term has expectation
$\alpha_t^2\Tr(C_0D_H)/4$. Fourier transformation of
$\sum_t(\delta\alpha_t)^2$ gives $\widehat\Omega^2$, proving the tadpole.
The same calculation uses the periodic covariance at finite $T$, and all
sums are finite. This proves both versions of the theorem.
\end{proof}

The source-independent Gaussian scalar cancels from $W[\alpha]-W[0]$;
the displayed degree-two source terms do not. Neither the $B$ vertex nor
the cosine curvature term can be read off from the instantaneous eigenvalues
of $H(\alpha)$. Replacing \eqref{eq:v4-exact-temporal-Q} by a product of
instantaneous Mehler eigenvalues would retain only $\Pi(0)$ and delete
these temporal contributions.

\subsection{A closed spectral reduction, including the signed cancellations}

In the eigenbasis of $H_0$, define
\begin{equation}
 I_{ij}(\Omega)=
 \frac{\sinh(\xi_i+\xi_j)}
 {4\sinh\xi_i\sinh\xi_j
           \{\cosh(\xi_i+\xi_j)-\cos\Omega\}},\qquad
 a_{ij}=\cosh(\xi_i+\xi_j)-1.
 \label{eq:v4-bubble-I}
\end{equation}
Introduce the symmetric matrix and scalar
\begin{equation}
 A_{\rm cov}=A-[H_0,B],\qquad
 \mathcal T_{\rm cov}=\mathcal T-\Tr(C_1B^*B).
 \label{eq:v4-covariant-data}
\end{equation}
The subscript ``cov'' denotes this proved commutator combination only;
it is not a claim that the complete gauge Ward identity has been closed.

\begin{theorem}[Exact dispersive reduction and the electric two-jet]
\label[theorem]{thm:v4-electric}
For the same determinant,
\begin{align}
 \Pi(0)&=\tfrac12\Tr(C_0D_H)
                     -\tfrac12\sum_{i,j}A_{ij}^2I_{ij}(0),
 \label{eq:v4-static-Pi}\\
 \Pi(\Omega)-\Pi(0)
   &=\widehat\Omega^2\left\{
       \mathcal T_{\rm cov}+
       \frac14\sum_{i,j}
       \frac{|(A_{\rm cov})_{ij}|^2 I_{ij}(0)}
                                  {a_{ij}+\widehat\Omega^2/2}\right\}.
 \label{eq:v4-exact-dispersion}
\end{align}
In particular the coefficient of $\widehat\Omega^2$ is
\begin{equation}
 \boxed{\mathcal E_{\rm det}
 =\mathcal T_{\rm cov}+
       \frac14\sum_{i,j}
                \frac{|(A_{\rm cov})_{ij}|^2I_{ij}(0)}{a_{ij}}.}
 \label{eq:v4-electric-coefficient}
\end{equation}
This is the electric \emph{candidate coefficient of this specified hard-fibre
determinant}, not an identification with the physical running gauge
coupling. The second term is nonnegative. The complete coefficient need
not have a fixed sign, since $\mathcal T_{\rm cov}$ is retained with its
native sign.
\end{theorem}
\begin{proof}
The Fourier transform of $c_i(r)c_j(r)$ from
\eqref{eq:v4-xi-covariance} is the geometric-series expression
\eqref{eq:v4-bubble-I}. It equals
$\int c_i(\omega+\Omega)c_j(\omega)\dd\omega/(2\pi)$, where now
$c_i(\omega)=(\lambda_i+2-2\cos\omega)^{-1}$.
Let $d(\omega,\Omega)=2(\cos\omega-\cos(\omega+\Omega))$ and
$\Delta_{ij}=\lambda_i-\lambda_j$. Direct subtraction of denominators gives
\begin{align}
 \int d\,c_i(\omega+\Omega)c_j(\omega)\frac{\dd\omega}{2\pi}
     &=\Delta_{ij}\{I_{ij}(0)-I_{ij}(\Omega)\},
 \label{eq:v4-bubble-first}\\
 \int d^2 c_i(\omega+\Omega)c_j(\omega)\frac{\dd\omega}{2\pi}
     &=\widehat\Omega^2\{(C_1)_{ii}+(C_1)_{jj}\}
           -\Delta_{ij}^2\{I_{ij}(0)-I_{ij}(\Omega)\}.
 \label{eq:v4-bubble-second}
\end{align}
Here is a direct verification of the second identity. Put
$q_i=\lambda_i+2-2\cos(\omega+\Omega)$ and
$q_j=\lambda_j+2-2\cos\omega$, so $d=q_i-q_j-\Delta_{ij}$.
Expand $d^2/(q_iq_j)$. The difference of the one-denominator means is
$(C_0)_{jj}-(C_0)_{ii}=\Delta_{ij}I_{ij}(0)$.
The remaining trigonometric integrals are
$\int d\,c_j=\widehat\Omega^2(C_1)_{jj}$ and
$\int d\,c_i(\omega+\Omega)=-\widehat\Omega^2(C_1)_{ii}$,
which give \eqref{eq:v4-bubble-second}. The first identity follows by the
same denominator subtraction without squaring.

The integrand trace in \eqref{eq:v4-native-response} is
$\sum_{ij}c_i(\omega+\Omega)c_j(\omega)(A_{ij}+dB_{ij})^2$.
Subtracting its value at zero frequency and using
\eqref{eq:v4-bubble-first}--\eqref{eq:v4-bubble-second} gives
\[
 -\sum_{ij}(A_{ij}-\Delta_{ij}B_{ij})^2
                  \{I_{ij}(0)-I_{ij}(\Omega)\}
        +2\widehat\Omega^2\Tr(C_1B^*B).
\]
This is the signed combination that must be formed before estimating.
Since $[H_0,B]_{ij}=\Delta_{ij}B_{ij}$ and
\[
 I_{ij}(0)-I_{ij}(\Omega)
    =\frac{\widehat\Omega^2 I_{ij}(0)}
                         {2(a_{ij}+\widehat\Omega^2/2)},
\]
substitution proves \eqref{eq:v4-exact-dispersion} and its limit
\eqref{eq:v4-electric-coefficient}. Degenerate eigenvalues are harmless;
all displayed functions are continuous and $a_{ij}>0$.
\end{proof}

The geometric scalar also has an equivalent form that exhibits all three
native contributions. Put $S=(I-P(0))\mathsf B U_0$ and
$D_\parallel=U_0^*(\operatorname{diag}_e((b_e^0)^2)\otimes nn^*)U_0$.
Then
\begin{equation}
 \mathcal T_{\rm cov}
 =\Tr\{C_1(N+S)^*(N+S)\}
                       -\Tr(C_0D)+\Tr(C_1D_\parallel).
 \label{eq:v4-geometric-completion}
\end{equation}
Indeed $\mathsf B^*\mathsf B=\mathsf D-
\operatorname{diag}((b_e^0)^2)\otimes nn^*$ and
$\mathsf B U_0=U_0B+S$. Expanding the square proves the identity.
The negative curvature trace in \eqref{eq:v4-geometric-completion} is not
removed by naming the other terms geometric. Any cancellation with gauge
integration or density must be proved in that same measured law.

\section{Quantitative frequency control and exact shell ownership}
\label{sec:v4-frequency-bounds}

\subsection{An explicit remainder and its gap dependence}

Assume for this subsection that $H_0\succeq m^2I$ and put
$\xi_*=\operatorname{arcosh}(1+m^2/2)>0$. Hilbert--Schmidt norms are denoted
by $\|\cdot\|_{\rm HS}$.

\begin{corollary}[Uniform temporal Taylor bound under an explicit hard gap]
\label[corollary]{cor:v4-electric-remainder}
The exact response satisfies
\begin{equation}
 \Pi(\Omega)=\Pi(0)+\widehat\Omega^2\mathcal E_{\rm det}
                                      +\mathcal R(\Omega),
 \qquad
 |\mathcal R(\Omega)|\le
 \frac{\coth\xi_*}{128\sinh^6\xi_*}
       \|A_{\rm cov}\|_{\rm HS}^2\widehat\Omega^4.
 \label{eq:v4-electric-remainder}
\end{equation}
The nonnegative dispersive part of $\mathcal E_{\rm det}$ is at most
\begin{equation}
 \frac{\coth\xi_*}{32\sinh^4\xi_*}\|A_{\rm cov}\|_{\rm HS}^2.
 \label{eq:v4-electric-upper}
\end{equation}
These estimates have no additional matrix-dimension or time-length factor.
Their uniformity requires a common $m$ and bounds on the displayed vertex
norms; it does not follow for the companion's entire hard space when
$L\to\infty$, since its lowest frequency tends to zero.
\end{corollary}
\begin{proof}
Subtract the zero-frequency denominator in
\eqref{eq:v4-exact-dispersion}. The remainder is exactly
\[
 -\frac{\widehat\Omega^4}{8}\sum_{ij}
  \frac{|(A_{\rm cov})_{ij}|^2I_{ij}(0)}
          {a_{ij}(a_{ij}+\widehat\Omega^2/2)}.
\]
Use
$a_{ij}\ge2\sinh^2\xi_*$ and
$I_{ij}(0)=\coth((\xi_i+\xi_j)/2)/(4\sinh\xi_i\sinh\xi_j)
\le\coth\xi_*/(4\sinh^2\xi_*)$.
Summing the squared entries gives both inequalities. The bilinear extension of this dispersive expression is bounded by the
product of the corresponding Hilbert--Schmidt norms by Cauchy--Schwarz;
this algebraic observation does not identify different background bundles.
All constants and their hard-gap dependence are explicit.
\end{proof}

At fixed positive gap the temporal response also has an exponentially
summable time kernel: the nonlocal part before frequency reduction is made
from products $c_i(r)c_j(r)$ and their finitely shifted copies. They are
bounded by constants times $e^{-2\xi_*|r|}$. In particular every fixed
temporal moment exists. For a fixed finite local source, its primitive
vertex ranks or Hilbert--Schmidt norms can be volume independent; for a
plane wave they generally scale with volume and must be normalized by its
physical quadratic form. No locality of a sharp spatial momentum projector
has been inferred from a temporal gap.

\begin{proposition}[The finite-cycle limit is controlled, not an extra cutoff limit]
\label[proposition]{prop:v4-time-limit}
Fix $H_0>0$ and the finite vertex matrices. For every
$0<\eta<\xi_*$, the discrete-mean response in
\cref{thm:v4-native-response} differs from its infinite-time expression at
any allowed real $\Omega$ by at most $C_\eta e^{-\eta T}$, after increasing
$C_\eta$ to cover the finitely many small $T$. The constant depends on the
hard gap and matrix stocks, not $T$ or the allowed frequency.
\end{proposition}
\begin{proof}
On $|\operatorname{Im}z|\le\eta$,
\[
 |\lambda_i+2-2\cos z|
 \ge\operatorname{Re}(\lambda_i+2-2\cos z)
 \ge2(\cosh\xi_* -\cosh\eta)>0.
\]
The response integrands and the means defining $C_0,C_1$ are therefore
periodic holomorphic functions in this strip with uniform bounds. The
transpose vertex is continued algebraically as $A-dB$, not by conjugating
the complex frequency. Real
shifts by $\Omega$ leave the bound unchanged. Their Fourier coefficients
have size at most $M_\eta e^{-\eta|n|}$ by shifting the contour. The
$T$-point average keeps exactly the coefficients with $n$ a multiple of
$T$, so its difference from the integral is at most
$2M_\eta e^{-\eta T}/(1-e^{-\eta T})$. The same reasoning applies to each
finite trace term and proves the assertion.
\end{proof}

\subsection{A spatial shell restriction must keep cross-shell bubbles}

The companion proposes a shell-restricted static calculation. It is exactly
additive when the complete background-dependent quadratic form reduces the
chosen shell subspaces, as happens for a preserved transverse-translation
label in its restricted background. It is not additive for arbitrary local
source insertions merely because $H_0$ is diagonal at zero.

\begin{proposition}[Two-jet shell composition with its return term]
\label[proposition]{prop:v4-shell-return}
At finite $T$, partition the hard coordinates into two spectral subspaces of
$H_0$, denoted $H$ and $L$. Write the full positive temporal matrix as
\[
 Q(a)=\begin{pmatrix}Q_{HH}(a)&Q_{HL}(a)\\Q_{LH}(a)&Q_{LL}(a)\end{pmatrix},
 \qquad Q_{HL}(0)=0.
\]
Exact integration of the $H$ variables produces the determinant
$\tfrac12\log\det Q_{HH}$ and the retained quadratic form
\begin{equation}
 Q^{\rm eff}_{LL}=Q_{LL}-Q_{LH}Q_{HH}^{-1}Q_{HL}.
 \label{eq:v4-temporal-Schur}
\end{equation}
For any real source direction $a$, its second derivative at zero is
\begin{equation}
 (Q^{\rm eff}_{LL})''(0)
   =Q_{LL}''(0)-2Q_{LH}'(0)Q_{HH}(0)^{-1}Q_{HL}'(0).
 \label{eq:v4-return-second}
\end{equation}
The later $L$ determinant therefore receives
$-\Tr(C_LQ_{LH}'C_HQ_{HL}')$, exactly the two cross-shell terms in the
full bubble. Discarding this return changes the one-loop electric and
magnetic response in general. Repeated elimination assigns such a pair to
one Schur return once; it does not justify dropping it.
\end{proposition}
\begin{proof}
Complete the finite Gaussian square, retaining its determinant. The matrix
factorization gives
$\det Q=\det Q_{HH}\det Q^{\rm eff}_{LL}$ exactly. At zero the two
cross blocks vanish, so the second derivative of their product is exactly
the term in \eqref{eq:v4-return-second}; differentiating its inverse factor
would require a nonzero zero-order cross block and contributes nothing at
this order. Half-tracing the second derivative against $C_L$ gives the
stated return. The full log-determinant bubble contains both orientations
$H\to L\to H$ and $L\to H\to L$ with coefficient $-1/2$ each, so they
agree. Iteration uses this exact factorization at every elimination.
\end{proof}

The Schur identity itself is inherited from the monographs. Its concrete
application here is to the native time-dependent vertex
\eqref{eq:v4-frequency-vertex} and the coefficient
\eqref{eq:v4-electric-coefficient}. A sharp shell may destroy spatial
quasilocality even when this finite identity is exact. A shell used in P1/P4
must therefore supply its own spatial kernel and source-transport bounds;
a list of eigenvalues is not a substitute.

\section{The precise transfer to measured marginal coordinates}
\label{sec:v4-transfer}

The preceding formulas give computable, rigorously specified finite
coefficients from $H_0,A,D_H,U_0,N$ and the diagonal profile $b_e^0$. They
replace finite-amplitude subtraction of eigenvalue sums by exact matrix
jets, remove the artificial hard-rank jump at zero, and include the native
temporal vertices. They do not require a known beta function. This is the
specific P4 interface advanced by the companion, not another source-bank
calibration.

Three distinct statements must not be conflated. The local conditional
high-fibre determinant in \eqref{eq:v4-exact-temporal-Q} is fully specified
and its two-jet is computed. The corresponding \emph{measured} local
one-loop coefficient is \eqref{eq:v4-measured-coefficient}, which also
requires the actual induced density. Finally the gauge-projected native
transfer, or the four-dimensional root-block pushforward of the original
programme, may have extra Gaussian variables and a different constraint.
Those require their actual constrained Hessian and density, not just the
projector \eqref{eq:v4-hard-projector}.

For example, if temporal gauge or other eliminated coordinates $z$ remain,
the actual quadratic matrix has blocks in $(q,z)$. Eliminating $z$ produces
its log determinant and the Schur complement of its block; both enter the
one-loop coefficient. Applying \cref{prop:v4-shell-return} keeps all mixed
kinetic returns. The existence of the stationary horizontal cancellation
in the companion's V68 does not prove that these terms vanish for
$\alpha_t\ne\alpha_{t+1}$. The ambient time-link calculation above is
explicitly prior to that extra gauge integration. It has not been called
a complete gauge-invariant electric renormalization.

Once the same measured convention has been identified, fix two physical
linear extractions, $\mathfrak b_B$ for the magnetic quadratic term and
$\mathfrak b_E$ for the temporal-gradient term, with their classical
normalizations specified. Their one-loop outputs are separate coordinates,
\[
 q_B=\mathfrak b_B\Gamma_1,\qquad q_E=\mathfrak b_E\Gamma_1.
\]
The determinant part of the second is given by
\eqref{eq:v4-electric-coefficient}; its density and extra-variable Schur
parts are added in the same signed expression. Equality of $q_B$ and $q_E$
is not imposed. An anisotropy coordinate must remain in the state until
its normalization or symmetry identity is proved. This is the P1 coupling
content required by the supplied programme, rather than an identification
of the static magnetic fit with a scalar physical coupling.

There is no new physical reciprocal-flow theorem hidden in this notation.
After complete extraction, V1's reciprocal rule applies to each positive
kinetic coefficient in its declared convention. Without the measured
identification, Ward-compatible local extraction, the full compact ratio,
and their preserved budgets, these numbers do not generate the claimed
ultraviolet trajectory. In particular the analytic collar task of V3 and
the chiral line task for the SM are unchanged.

\subsection{Reproducible checks and the distinction between proof and diagnostics}

The accompanying standalone script
\srcid{check_native_temporal_response_V4.py} performs an exact rational
quaternion check of \eqref{eq:v4-time-link-jet}. Its numerical tests compare
the finite-cycle determinant Hessian with
\eqref{eq:v4-native-response}, compare the signed spectral reduction with
the unreduced frequency sum, check the bubble integral against
\eqref{eq:v4-bubble-I}, and verify the finite-cycle determinant identity.
It also constructs the excluded frame on a small lattice and checks its
predicted rank at zero and nonzero amplitude. All tests have declared
finite tolerances and fixed seeds. They are independent diagnostics, not
proofs of a continuum limit or interval-certified lattice beta coefficients.
The theorems are proved above without relying on those floating-point tests.

\subsection{Uniformity and the next load-bearing calculation}

\begin{center}
\small
\begin{tabular}{>{\raggedright\arraybackslash}p{.23\textwidth}>{\raggedright\arraybackslash}p{.68\textwidth}}
\toprule
Variable & Exact scope in V4\\
\midrule
Spatial volume & Bundle analyticity and stationary recentering are proved at
every fixed finite $L$. A common positive gap and bounded local vertex
stocks give dimension-free frequency estimates. The entire companion
hard space does not have a volume-independent lower gap.\\
Time length & The determinant identity and response hold for every finite
$T\ge3$. The infinite-time response and its exponential comparison are
proved at fixed positive hard gap and fixed vertex stocks.\\
Cutoff and shell & No preservation of the gap, spatial localization,
induced-density stocks, or source normalization through changing spatial
cutoffs is proved. Spectral additivity needs reducing vertices; otherwise
the Schur returns are retained.\\
Loop/source order & This installment computes the one-loop coefficient
through total background degree two. The Gaussian frequency dependence at
that degree is exact. It does not supply an arbitrary-order nonlinear
remainder or an all-source analytic bound.\\
RG steps & Exact Gaussian Schur composition is finite-step algebra. A bound
uniform over arbitrarily many interacting RG steps is not inferred.\\
Coupling & The coefficient calculation is independent of a supplied running
$h$. Matching it to an actual small-positive-$h$ full-shell expansion still
requires the local remainder and compact-sector estimates.\\
\bottomrule
\end{tabular}
\end{center}

There are now two explicit adjacent tasks, not a claim that either is paid
by the other. On the P4 coefficient branch, compute the actual gauge/constraint
Schur and induced-density two-jets for a time-dependent retained source,
then insert them in the signed formula above. The source must be the one
transported by the declared block; replacing it by a static spectral
parameter is not allowed. On P2's analytic branch, prove the
collar-dressed pair estimate \eqref{eq:v3-dressed-pair} and its higher
activity extension. The new companion supplies no estimate for that ratio.
The next iteration should attack whichever of these explicit primitive
interfaces admits a same-law proof, not repeat a Gaussian trace calculation.

\begingroup
\renewcommand{\refname}{Additional source for V4}
\begin{thebibliography}{99}
\bibitem[11]{V4Companion}
\emph{Renewal Yang--Mills Mass Gap Research Notes}, cumulative V70,
attachment \srcid{Renewal_Yang_Mills_Mass_Gap_Research_Notes_V70(2).tex}.
The relevant increments are V66 (native temporal Wilson kernel),
V68 (stationary horizontal Gaussian reduction), V69 (static plane-wave
response and cubic-chaos erratum), and V70 (frozen source covariance and
explicit remaining one-shell/electric tasks). Numerical tables and
unprovided verifier files are not promoted to rigorous uniform estimates.
\end{thebibliography}
\endgroup

\section{Final ledgers for V4}
\label{sec:v4-ledgers}

\subsection*{Proved}
\addcontentsline{toc}{subsection}{V4: Proved}
The divided gauge columns yield an analytic fixed-rank hard bundle for the
companion's plane wave, with the precise excluded-mode count and explicit
frame jets (\cref{thm:v4-analytic-bundle}). Its spatial and native temporal
hard force begins at total source degree three, so stationary recentering
does not change the determinant or measured-density two-jet
(\cref{thm:v4-recentering}). The native temporal plaquette has the exact
quadratic form \eqref{eq:v4-time-link-jet}, leading to the fully specified
conditional high-fibre matrix \eqref{eq:v4-exact-temporal-Q}.
Its Gaussian determinant obeys the exact temporal response
\eqref{eq:v4-native-response}, the signed commutator reduction
\eqref{eq:v4-exact-dispersion}, and the electric candidate formula
\eqref{eq:v4-electric-coefficient}. Its temporal Taylor remainder and
finite-cycle limit have explicit hard-gap-dependent bounds. Its
cross-shell two-jet is assigned exactly by
\eqref{eq:v4-return-second}. These are coefficient statements about the
specified local fibre, with the scope in the uniformity table; no complete
nonlinear shell bound is included.

\subsection*{Inherited}
\addcontentsline{toc}{subsection}{V4: Inherited}
V1's measured finite-order compiler, source tower, and reciprocal
normalization; V2's exact sector representation and global/local carrier;
and V3's compact geometry, guarded sector estimates, and unclosed collar
identity remain unchanged. The actual Wilson temporal kernel, Mehler
normalization, and stationary horizontal Gaussian identities are taken from
the cited companion locations at their stated scope. The constrained-Hessian,
measured-density, and Schur principles remain the monographs' inputs. The
companion's frozen source-decrement calculation is not repeated as a new
physical capture estimate.

\subsection*{Conditional}
\addcontentsline{toc}{subsection}{V4: Conditional}
A physical electric one-loop coefficient follows only after the conditional
fibre is identified with the actual measured gauge/constraint elimination,
including the induced density and any extra Gaussian Schur returns.
Cutoff-uniform frequency bounds require uniform hard gaps and local vertex
stocks; a sharp spectral shell alone supplies neither the spatial locality
nor the source normalization. An actual marginal update also requires the
complete local-order and compact-correction estimates. None of these
requirements is replaced by the companion's static numerical slope or by
the exact Gaussian time limit.

\subsection*{Open}
\addcontentsline{toc}{subsection}{V4: Open}
The newly isolated coefficient task is the time-dependent measured
constraint/gauge and density two-jet for the actual block, in a single
physical source convention. V3's collar-dressed connected sectors remain
the first unclosed analytic compact interface, followed by their higher
connected activity norm and preservation of the global residual class.
Ward-intertwining local subtraction, a generated joint marginal/stable
trajectory, composed source transport, chiral line coherence, and the
perturbative continuum/scheme limits remain unproved here. These are not
reclassified as consequences of a static transfer spectrum. The separate
YM mass-gap problem and nonperturbative SM ultraviolet completion remain
outside this lineage.


\clearpage
\part*{V5: Relative gauge integration, the measured electric two-jet, and its volume bounds}
\addcontentsline{toc}{part}{V5: Relative gauge integration and the measured electric two-jet}

\section{A measured continuation of the native transfer calculation}
\label{sec:v5-start}

Sections~1--32 are preserved. The calculation in V4 explicitly stops before
integrating relative temporal gauge transformations. The present installment
performs their \emph{based} integration in the native Wilson transfer, constructs
the associated source-coordinate density, and evaluates both missing
contributions to the electric two-jet. The result is a measured local
coefficient, not a determinant evaluated with an unspecified density. The
residual spatially constant gauge transformations and all low coordinates
remain retained; they are not regularized away or counted as nonzero Gaussian
modes. This qualification is essential at the trivial background.

The inherited starting point is the exact based-gauge tube and half-density in
\cite{V4Companion}, \srcid{f16v44:tube}, together with its native Wilson kernel
\srcid{f16v66:structure}. The stationary horizontal cancellation in
\srcid{f16v68:reduction} treats equal backgrounds. Its nonequal-time extension
is not an inherited theorem. We compute it below. V4's hard projector,
quaternion convention, spatial plane wave, and hard covariance retain their
precise definitions. In particular, $h=\gamma^{-1}$ is the transfer parameter
in this local calculation; it has not been identified with the physical
coupling of the four-dimensional root block.

The new conclusion is
\[
 \Pi_{\rm based}(\Omega)=\Pi_{\rm V4}(\Omega)
      +\widehat\Omega^2\bigl(\mathcal E_{\rm orb}
                                    +\mathcal E_{\rm ret}\bigr).
\]
Both added coefficients are explicit, nonpositive, and separately bounded
below after normalization by the spatial profile norm. Their sum obeys a
logarithmic, rather than a power-law, infrared bound in three spatial
dimensions. This bound is uniform in the spatial volume at fixed nonzero
external momentum. No statement about the sign of the \emph{total} electric
coefficient or of a beta function follows from the sign of this addition.

The work addresses a necessary measured part of P4. It does not replace P2's
collar-dressed compact-sector estimate. Background-field renormalization based
on exact lattice BRS, background-gauge, and shift identities is established
in a different formulation \cite{V5LuscherWeisz}; no counterterm, Ward identity,
or running coupling from that formulation is imported into the present
transfer fibre.

\subsection{The actual quotient and the retained variables}

Let $\mathcal X=\mathrm{SU}(2)^{3V}$, $V=L^3$, $L\ge3$, and let
$\mathcal G_o=\{g:g_o=I\}$ be the based vertex group. Its action is free.
Write
\[
 \mathcal K_h(X,X')=N_h^{-3V}
  \exp\left[-h^{-1}\sum_e
       \{1-\operatorname{Re}(X_eX_e'^{-1})\}\right].
\]
The spatial magnetic factors on the two endpoints are
$\exp[-S_{\rm W}(X)/(2h)]$ and $\exp[-S_{\rm W}(X')/(2h)]$.
Their product with $\mathcal K_h$ is the native transfer kernel. Here
$\operatorname{Re}$ means half the fundamental trace. Its based-averaged
kinetic kernel is exactly
\begin{equation}
 \mathcal K_h^o(X,X')=\int_{\mathcal G_o}
                           \mathcal K_h(X,g\cdot X')\,\dd H_o(g).
 \label{eq:v5-based-kernel}
\end{equation}
This is the operator on based-invariant functions obtained by the based
orthogonal projector, which commutes with the native transfer. All group
integrals use probability Haar. There is no adjustable gauge-fixing Gaussian
in \eqref{eq:v5-based-kernel}.

We use V4's profile $b_e=b_e^0n$, $n=e_3$, with $b_e^0=\cos(kx_1)$ on
polarization $\nu\ne1$ and zero otherwise, where $0<k=2\pi m/L<\pi$.
Let $U(a)$ denote the \emph{hard frame} of \eqref{eq:v4-polar-frame}, not a
group-valued configuration; write $\overline X(a)_e=\exp(a b_e)$ for the
background configuration. The $d=6V-10$ columns of $U(a)$ are orthogonal to
all spatial gauge directions. The based horizontal space has dimension
$6V+3$, so its complementary low space has dimension thirteen. These
thirteen coordinates, including the amplitude $a$, are retained.

\begin{proposition}[A source chart with a specified leading density]
\label{prop:v5-source-chart}
At each fixed $L,k$, V4's hard fibre extends to a local chart of the based
quotient with coordinates $(a,z,q)$, $z\in\R^{12}$, $q\in\R^d$, of the form
\begin{equation}
 X(a,z,q)_e=\overline X(a)_e
           \exp\{(V_\perp(a)z+U(a)q)_e\}.
 \label{eq:v5-full-slice}
\end{equation}
The columns $V_\perp(a)$ can be chosen analytic, orthonormal, based-horizontal,
and orthogonal to both $U(a)$ and $b$. Relative to
$\|b\|\,\dd a\,\dd z\,\dd q$, the quotient Haar density $\mathfrak j$
satisfies, up to a positive source-independent constant,
\begin{equation}
 \mathfrak j(a,0,0)=\det\mathcal D(a)^{1/2},\qquad
 \mathcal D(a)=(G_a^b)^*G_a^b.
 \label{eq:v5-source-density}
\end{equation}
This is the density of the declared source chart, not the assertion that the
fixed transverse slice of the companion has unit horizontal Jacobian.
\end{proposition}
\begin{proof}
The based vertical Gram is positive because the based action is free.
The vector $b$ is exactly based-horizontal at every background: rotations
about $n$ fix its color, and its lattice divergence in direction $\nu$
is zero. It is orthogonal to $U(a)$ by construction. The orthogonal
complement of $\operatorname{Ran}G_a^b$, $\operatorname{Ran}U(a)$, and $b$
has constant dimension twelve. Projecting a basis at zero onto this
complement and multiplying by its inverse Gram square root constructs
$V_\perp(a)$ analytically on a smaller interval.

At $(z,q)=(0,0)$, the right-trivialized derivative columns of the slice are
$[b,V_\perp(a),U(a)]$. After replacing $b$ by $b/\|b\|$, they are an
orthonormal frame of the based horizontal space. The derivative of
$(g,a,z,q)\mapsto g\cdot X(a,z,q)$ is therefore an invertible square
matrix with columns $[G_a^b,b,V_\perp(a),U(a)]$. The inverse function theorem
and the inherited based tube give a local quotient chart. The squared
absolute volume determinant at the reference section is
$\|b\|^2\det((G_a^b)^*G_a^b)$, since vertical and horizontal columns are
orthogonal. Right translation is an isometry and preserves fine Haar.
The normalization constants of fine and based Haar are independent of
$a$. This proves \eqref{eq:v5-source-density}. The full density away from
this section remains the exact pulled-back Haar density.
\end{proof}

In a symmetric kernel on flat quotient-coordinate measure, each endpoint
contributes $\mathfrak j^{1/2}$. On a periodic cycle their product is exactly
one factor $\mathfrak j$ per slice. Equivalently, use the nonsymmetrized
kernel with measure $\mathfrak j\,\dd a\,\dd z\,\dd q$ from the outset.
These give the same cycle integral; inserting both densities would count the
same Jacobian twice. We hold $z=0$ while evaluating the amplitude two-jet.
Other retained coordinates are not integrated or replaced by their means.

The full vertex group decomposes uniquely as $g_x=h k_x$, $k_o=I$, with
product Haar $\dd h\,\dd H_o(k)$. Thus the remaining global projector can
still be applied as an exact compact operation. In this installment it is
retained, just as the residual simultaneous adjoint action is retained in
\srcid{f16v44:tube}. A trace on based-invariant functions is not asserted to
be a trace on the fully invariant Hilbert space without that operation.

\section{The relative-gauge Hessian from the native temporal plaquette}
\label{sec:v5-gauge-hessian}

Let $\mathcal V_o=\{\eta\in(\R^3)^V:\eta_o=0\}$. For an oriented spatial
edge $e=(x,x+\widehat\mu)$ define maps from $\mathcal V_o$ to link vectors by
\[
 (\mathsf T_a\eta)_e=
       \operatorname{Ad}(\overline X_e(a))^{\mathsf T}\eta_x,
 \qquad (\mathsf H\eta)_e=\eta_{x+\widehat\mu},
 \qquad G_a=\mathsf T_a-\mathsf H.
\]
Thus $G_a=G_a^b$. In this section $\mathsf H$ is the head-evaluation map;
$H(a)$ without a sans-serif font remains V4's magnetic hard Hessian.
Set
\[
 \mathsf T=\mathsf T_0,\quad G=\mathsf T-\mathsf H,\quad
 \mathsf F=\mathsf T+\mathsf H,\quad
 \mathcal D=G^*G,\quad \mathcal G=\mathcal D^{-1}.
\]
Use the ambient matrices $\mathsf B,\mathsf D$ of
\eqref{eq:v4-native-BD}; in particular $\mathsf B^*=-\mathsf B$.
Then $\mathsf T_a=e^{-2a\mathsf B}\mathsf T$.

\subsection{The four-factor identity}

For endpoint backgrounds $a,b$, write $v=U(a)q$, $w=U(b)q'$, and
$g_x=e^{\eta_x}$. By cyclicity of the trace, the relative-gauge temporal
word on edge $e$ is exactly
\begin{equation}
 \operatorname{Re}\left\{
  e^{(a-b)b_e^0n}e^{v_e}e^{\eta_{x+\widehat\mu}}
           e^{-w_e}e^{-(\mathsf T_b\eta)_e}\right\}.
 \label{eq:v5-four-factor-word}
\end{equation}
The use of the \emph{second} endpoint in $\mathsf T_b$ follows from moving
$\overline X(b)^{-1}$ through the trace, not from a convention for an
independent temporal field.

For four small imaginary quaternions $x_1,\ldots,x_4$, multiplication gives
\begin{equation}
\begin{split}
 1-\operatorname{Re}\{e^{\delta n}\prod_{i=1}^4e^{x_i}\}
 ={}&1-\cos\delta+\sin\delta\,n\cdot\sum_i x_i
        +\frac{\cos\delta}{2}\left|\sum_i x_i\right|^2\\
    &+\sin\delta\,n\cdot\sum_{i<j}x_i\times x_j+O(|x|^3).
\end{split}
 \label{eq:v5-four-factor-jet}
\end{equation}
Indeed the scalar and vector parts of the product through degree two are
$1-|\sum_i x_i|^2/2$ and $\sum_i x_i+\sum_{i<j}x_i\times x_j$.
The remainder is bounded on every fixed compact $\delta$ interval.
This proves the identity directly with the native group curvature retained.

Let
\[
 C(a,b)=\operatorname{diag}_e\cos((a-b)b_e^0)\otimes I_3,
 \qquad S(a,b)=\operatorname{diag}_e\sin((a-b)b_e^0)\otimes J,
 \quad Jv=n\times v.
\]
Thus $S^*=-S$. Substitution of
$(x_1,x_2,x_3,x_4)=(v,\mathsf H\eta,-w,-\mathsf T_b\eta)$
gives the following exact quadratic Hessian blocks.

\begin{theorem}[Native vertical and mixed blocks]
\label{thm:v5-native-blocks}
The $\eta$ block and the two hard--vertical blocks of the summed temporal
plaquettes are
\begin{align}
 M(a,b)&=G_b^*C(a,b)G_b
       +\mathsf H^*S(a,b)\mathsf T_b
                  -\mathsf T_b^*S(a,b)\mathsf H,
             \label{eq:v5-M}\\
 L_-(a,b)&=U(a)^*\{-C(a,b)+S(a,b)\}G_b,
             \label{eq:v5-Lminus}\\
 L_+(a,b)&=U(b)^*\{C(a,b)G_b
                         -S(a,b)(\mathsf T_b+\mathsf H)\}.
             \label{eq:v5-Lplus}
\end{align}
The hard--hard block is exactly V4's native block, before this new integration.
The quadratic action contains
$q^*L_-\eta+q'^*L_+\eta+\eta^*M\eta/2$.
At equal backgrounds,
\begin{equation}
 M(a,a)=\mathcal D(a),\qquad L_-(a,a)=L_+(a,a)=0.
 \label{eq:v5-static-cancellation-data}
\end{equation}
All formulas are finite-dimensional analytic identities near $(a,b)=(0,0)$.
\end{theorem}
\begin{proof}
The cosine term in \eqref{eq:v5-four-factor-jet} is
$\frac12(v-w-G_b\eta)^*C(v-w-G_b\eta)$.
The sine cross products, using $n\cdot(x\times y)=-x^*Jy$, give
\[
 v^*Sw+v^*SG_b\eta-w^*S(\mathsf T_b+\mathsf H)\eta
                         +(\mathsf H\eta)^*S\mathsf T_b\eta.
\]
The last scalar has symmetric Hessian
$\mathsf H^*S\mathsf T_b-\mathsf T_b^*S\mathsf H$.
Reading the remaining bilinear terms proves
\eqref{eq:v5-M}--\eqref{eq:v5-Lplus}, including their signs.
At $a=b$, $C=I$, $S=0$, and $U(a)^*G_a=0$; this proves
\eqref{eq:v5-static-cancellation-data}. Analyticity follows from the actual
quaternion functions and V4's analytic frame.
\end{proof}

\subsection{The jets that enter a measured electric coefficient}

Define symmetric vertex matrices on $\mathcal V_o$ by
\begin{align}
 \mathcal R&=\mathsf H^*\mathsf B\mathsf T
                                -\mathsf T^*\mathsf B\mathsf H,\\
 \mathcal R_1&=-2(\mathsf H^*\mathsf B^2\mathsf T
                                +\mathsf T^*\mathsf B^2\mathsf H),\\
 \mathcal W&=G^*\mathsf D G,\qquad
 K=U_0^*\mathsf B\mathsf F.
 \label{eq:v5-gauge-vertices}
\end{align}
The letter $K$ in this display is a finite mixed vertex, not a loop-order
index or a blocking kernel.

\begin{lemma}[Symmetric vertical jet and the common endpoint return]
\label{lem:v5-jets}
At total amplitude degree at most two,
\begin{align}
 \mathcal D(a)&=\mathcal D+2a\mathcal R+a^2\mathcal R_1+O(a^3),
                                                   \label{eq:v5-Djet}\\
 M(a,b)&=\mathcal D+(a+b)\mathcal R+ab\mathcal R_1
                   -\tfrac12(a-b)^2\mathcal W+O((|a|+|b|)^3),
                                                   \label{eq:v5-Mjet}\\
 L_-(a,b)&=-(a-b)K+O((|a|+|b|)^2),\qquad
 L_+(a,b)=-(a-b)K+O((|a|+|b|)^2).
                                                   \label{eq:v5-Ljet}
\end{align}
The same sum $q+q'$ therefore couples to the vertical mode at first source
order. Replacing it by a difference would give the wrong return coefficient.
\end{lemma}
\begin{proof}
Expand $\mathsf T_a=\mathsf T-2a\mathsf B\mathsf T
+2a^2\mathsf B^2\mathsf T+O(a^3)$.
In $G_a^*G_a$ the terms containing
$\mathsf T^*\mathsf B\mathsf T$ cancel at first order; at second order the
$\mathsf T^*\mathsf B^2\mathsf T$ terms also cancel. This gives
\eqref{eq:v5-Djet}. Next use
$C=I-(a-b)^2\mathsf D/2+O((a-b)^4)$ and
$S=(a-b)\mathsf B+O((a-b)^3)$ in \eqref{eq:v5-M}.
The sine part is $(a-b)\mathcal R+b(a-b)\mathcal R_1$ to this order.
Combining it with \eqref{eq:v5-Djet} at $b$ gives \eqref{eq:v5-Mjet}.

Differentiate $U(a)^*G_a=0$ at zero:
$N^*G=-U_0^*G'_0$, where $G'_0=-2\mathsf B\mathsf T$.
The linear part of $L_-$ is
$(a-b)U_0^*(G'_0+\mathsf B G)=-(a-b)U_0^*\mathsf B\mathsf F$.
For $L_+$ the term $U(b)^*G_b$ vanishes exactly; its sine term has this same
linear part. This proves \eqref{eq:v5-Ljet} without using a numerical
frame derivative.
\end{proof}

\section{The stationary section and the complete local based coefficient}
\label{sec:v5-measured-coefficient}

Consider a periodic time cycle of length $T\ge3$. Use the exact based kernel
\eqref{eq:v5-based-kernel}, the chart \eqref{eq:v5-full-slice} on every slice,
and relative based variables $g_t=e^{\eta_t}$ on every time edge. Retain
$(a_t,z_t)$ and set $z_t=0$ on the evaluated source family. There is an exact
local integral obtained by restricting $(q_t,\eta_t)$ to fixed small
neighbourhoods containing their reference zeros. Its exponent is the native
cycle action with the relative gauge transformations left in the temporal
plaquettes. Its density is
\[
 \prod_t\mathfrak j(a_t,0,q_t)\,j_o(\eta_t),
 \qquad j_o(0)=\text{positive constant},
\]
where $j_o$ is the based-Haar density in exponential coordinates. This local
integral is a specified contribution to the exact transfer, not its compact
complement. Cutoffs equal to one near the stationary section do not change
its formal stationary coefficients.

\begin{theorem}[No missing recentering term at source degree two]
\label{thm:v5-stationary}
For this actual local $(q,\eta)$ integral at fixed $L,T,k$, the force in
$\eta$ at $q=\eta=0$ vanishes exactly for every small $(a_t)$.
The force in $q$ starts at total source degree three. Its joint stationary
section consequently satisfies
$(q_*,\eta_*)=O(\|\boldsymbol a\|^3)$.
Let $Q(\boldsymbol a)$ be V4's hard temporal matrix,
$M=\operatorname{diag}_t M(a_t,a_{t+1})$, and let $L$ place $L_-$ and
$L_+$ in the two hard endpoint rows for each vertical edge variable.
The local measured one-loop coefficient through total source degree two is
\begin{equation}
 \boxed{\Gamma^o_1[\boldsymbol a]
  =\frac12\log\det\begin{pmatrix}Q&L\\L^*&M\end{pmatrix}
       -\frac12\sum_t\log\det\mathcal D(a_t)
       \quad\pmod{\text{source degree}\ge3},}
 \label{eq:v5-complete-measured}
\end{equation}
up to a source-independent scalar. Both the based gauge determinant and the
induced source density are thus specified.
\end{theorem}
\begin{proof}
At the abelian background, a temporal gauge variation has color-$n$ force
proportional to the spatial divergence of
$\sin((a_t-a_{t+1})b_e^0)n$. This field is supported in direction $\nu$
and is independent of $x_\nu$, so its divergence is identically zero.
The adjoint factor $\operatorname{Ad}(\overline X(b))$ fixes $n$, and the
based restriction cannot change a vanishing full gradient. The other
colors have zero first trace derivative. Thus the $\eta$ force is exactly
zero. At $\eta=0$, V4's hard-force calculation applies without change:
the spatial linear term and each temporal linear term lie in the retained
$b$ direction, while sine has no quadratic term. The $q$ force is therefore
$O(\|\boldsymbol a\|^3)$.

At zero source the joint Hessian is block diagonal, with positive blocks
$H_0-\Delta_T$ and $\operatorname{diag}_t\mathcal D$. Based freeness proves
the latter positivity. The analytic implicit function theorem gives the
joint stationary section, whose degree-one and degree-two equations have
zero inhomogeneous terms. Evaluating the Hessian or the smooth leading
density at that section therefore leaves their source two-jets unchanged.
V1's measured stationary coefficient rule gives half the log determinant
of this joint Hessian minus the log leading density. By
\eqref{eq:v5-source-density}, that density at the reference section has
source-dependent part $\prod_t\det\mathcal D(a_t)^{1/2}$.
The based-Haar density at $\eta=0$, the profile Jacobian $\|b\|$, and the
normalizations $N_h$ contribute only source-independent scalars.
This proves \eqref{eq:v5-complete-measured}.
\end{proof}

The formula applies to the bare native action and its stated local source
chart. An incoming effective density or a generated classical interaction
would contribute its own jets and could change the recentering argument.
Those terms are not covered by assigning their values to the bare matrices.
No full compact remainder bound is deduced from the stationary coefficient
rule.

Gaussian elimination now gives, within the same formula,
\begin{equation}
 \Gamma^o_1=\tfrac12\log\det(Q-LM^{-1}L^*)
        +\tfrac12\sum_t\log\det M(a_t,a_{t+1})
        -\tfrac12\sum_t\log\det\mathcal D(a_t).
 \label{eq:v5-measured-schur}
\end{equation}
The three terms must be formed before any separate norm bound. In
particular, a time-dependent $M$ is not the endpoint orbit Gram
$\mathcal D(a_t)$, despite their equality when $a_t=a_{t+1}$.

\section{The added electric coefficient and its signed structure}
\label{sec:v5-electric}

Use $C_{0,T}$ and $C_{1,T}$ for the equal-time and one-step covariances of
$Q_0=H_0-\Delta_T$. These are the discrete means specified in V4. In
infinite time at fixed $L$ they become $C_0,C_1$ from
\eqref{eq:v4-xi-covariance}. Define
\begin{align}
 \mathcal E_{\rm orb}
   &=\tfrac12\left\{
       \Tr(\mathcal G\mathcal R\mathcal G\mathcal R)
             -\Tr(\mathcal G(\mathcal R_1+\mathcal W))\right\},
                                           \label{eq:v5-orb}\\
 \mathcal E_{{\rm ret},T}
   &=-2\Tr\{(C_{0,T}+C_{1,T})K\mathcal G K^*\}.
                                           \label{eq:v5-return}
\end{align}
No inverse of an unremoved constant gauge mode appears in either definition:
$\mathcal G$ is the based inverse.

\begin{theorem}[Complete based-gauge correction to V4's electric two-jet]
\label{thm:v5-electric}
For the measured local source family of \cref{thm:v5-stationary}, write
its quadratic source functional as
$\frac12\sum_\Omega|\widehat a(\Omega)|^2\Pi^o_T(\Omega)$.
For every allowed $\Omega=2\pi n/T$, exactly at this source order,
\begin{equation}
 \boxed{\Pi^o_T(\Omega)-\Pi^{\rm V4}_T(\Omega)
       =\widehat\Omega^2
                  (\mathcal E_{\rm orb}+\mathcal E_{{\rm ret},T}).}
 \label{eq:v5-response}
\end{equation}
In infinite time at fixed spatial regulator,
\begin{equation}
 \mathcal E_{\rm measured}^o
       =\mathcal E_{\det}^{\rm V4}
                  +\mathcal E_{\rm orb}+\mathcal E_{\rm ret},
 \quad \mathcal E_{\rm ret}=-2\Tr\{(C_0+C_1)K\mathcal G K^*\}.
 \label{eq:v5-total-electric}
\end{equation}
There is no additional frequency remainder from this gauge/density
correction. The remaining dispersive remainder at source degree two is
precisely the one in V4. The static two-jet is unchanged.
\end{theorem}
\begin{proof}
First expand the last two terms of \eqref{eq:v5-measured-schur}. Put
$\delta_t=a_t-a_{t+1}$. From \eqref{eq:v5-Mjet},
\begin{align*}
 \log\det M(a,b)={}&\log\det\mathcal D
 +(a+b)\Tr(\mathcal G\mathcal R)
 +ab\Tr(\mathcal G\mathcal R_1)
 -\tfrac12(a-b)^2\Tr(\mathcal G\mathcal W)\\
 &-\tfrac12(a+b)^2\Tr(\mathcal G\mathcal R\mathcal G\mathcal R)
 +O((|a|+|b|)^3),\\
 \log\det\mathcal D(a)={}&\log\det\mathcal D
 +2a\Tr(\mathcal G\mathcal R)
 +a^2\Tr(\mathcal G\mathcal R_1)
 -2a^2\Tr(\mathcal G\mathcal R\mathcal G\mathcal R)+O(a^3).
\end{align*}
On the cycle the constant and linear terms cancel. Use
$\sum_t(a_ta_{t+1}-a_t^2)=-\frac12\sum_t\delta_t^2$.
Their quadratic difference, including the prefactors $1/2$, is
$\frac12\sum_t\delta_t^2\mathcal E_{\rm orb}$.

Next, $L(0)=0$, so only its first jet and $M(0)^{-1}$ enter the
source-quadratic Schur subtraction. By \eqref{eq:v5-Ljet}, this subtracts
from the hard quadratic action
\[
 \frac12\sum_t\delta_t^2
             (q_t+q_{t+1})^*K\mathcal G K^*(q_t+q_{t+1}).
\]
The first variation of half the log determinant is the Gaussian mean of
the quadratic-action variation. Since
$\E[(q_t+q_{t+1})(q_t+q_{t+1})^*]=2(C_{0,T}+C_{1,T})$,
its contribution to the free energy is
$-\sum_t\delta_t^2\Tr\{(C_{0,T}+C_{1,T})K\mathcal G K^*\}$.
Adding the orbit part and using
$\sum_t\delta_t^2=\sum_\Omega\widehat\Omega^2
|\widehat a(\Omega)|^2$ proves \eqref{eq:v5-response}.
There are no other degree-two terms: derivatives of $M^{-1}$ multiply
at least two first-order factors of $L$ and hence start at degree three.
V4's fixed-regulator time limit proves \eqref{eq:v5-total-electric}.
\end{proof}

\subsection{A negative-square identity before estimating the traces}

Define the based gradient projector and a color-parallel multiplier by
\[
 P_g=G\mathcal G G^*,\qquad
 \mathsf B_\parallel=
             \operatorname{diag}_e(b_e^0)\otimes nn^*.
\]
The plane wave has the exact divergence identity
\begin{equation}
 \mathsf T^*\mathsf B\mathsf T
                   =\mathsf H^*\mathsf B\mathsf H.
 \label{eq:v5-divergence-profile}
\end{equation}
Indeed the only nonzero edge direction is $\nu$, and its weight is unchanged
by a shift in that direction. Restricting the vertex variables to the based
subspace preserves this matrix identity.

\begin{theorem}[The measured gauge addition is nonpositive]
\label{thm:v5-sign}
For this transverse plane-wave profile,
\begin{equation}
 \boxed{\mathcal E_{\rm orb}
 =-\tfrac12\|(I-P_g)\mathsf B\mathsf F\mathcal G^{1/2}\|_{\rm HS}^2
  -\tfrac12\|\mathsf B_\parallel G\mathcal G^{1/2}\|_{\rm HS}^2
 \le0.}
 \label{eq:v5-negative-square}
\end{equation}
Moreover $\mathcal E_{{\rm ret},T}\le0$ for every $T\ge3$, and
$\mathcal E_{\rm ret}\le0$ in infinite time. Thus the added correction in
\eqref{eq:v5-total-electric} is nonpositive. This statement concerns the
added terms, not the sign of $\mathcal E_{\rm measured}^o$.
\end{theorem}
\begin{proof}
Expanding \eqref{eq:v5-divergence-profile} gives
\[
 \mathcal R=\mathsf F^*\mathsf B G=-G^*\mathsf B\mathsf F.
\]
Since $\mathsf B^*\mathsf B=\mathsf D-\mathsf B_\parallel^*
\mathsf B_\parallel$, direct expansion also gives
\[
 \mathcal R_1+\mathcal W
 =\mathsf F^*\mathsf B^*\mathsf B\mathsf F
                  +G^*\mathsf B_\parallel^*\mathsf B_\parallel G.
\]
Cyclicity of the trace and the symmetry of $\mathcal R$ imply
\[
 \Tr(\mathcal G\mathcal R\mathcal G\mathcal R)
 =\Tr(\mathcal G\mathsf F^*\mathsf B^*P_g\mathsf B\mathsf F).
\]
Subtract the preceding positive decomposition to obtain
\eqref{eq:v5-negative-square}. The matrix $P_g$ is an orthogonal projector.
The sum $C_{0,T}+C_{1,T}$ is positive: it is half the covariance of
$q_t+q_{t+1}$. Since $K\mathcal G K^*\succeq0$, its trace product is
nonnegative. This proves the other signs.
\end{proof}

The negative-square form is useful analytically: bounding the two original
log determinants before cancellation would have kept artificial static and
linear terms. The calculation neither imposes a continuum Ward identity nor
uses a continuum beta coefficient. It is a signed identity of the specified
native, based-quotient local integral.

\section{Volume bounds without a smallest-eigenvalue power loss}
\label{sec:v5-volume}

The formulas so far hold at fixed finite spatial regulator. A crude bound by
$\|\mathcal D^{-1}\|$ would introduce a large pinned-mode loss. For the
return this loss is unnecessary; V4's particular excluded modes remove it
exactly. The orbit term admits a different local Green-function bound.

\subsection{The mixed vertex annihilates the constant gauge modes}

Let $\mathsf T_f,\mathsf H_f$ be tail and head evaluation on all vertex
fields, and let
\[
 G_f=\mathsf T_f-\mathsf H_f,\quad
 \Delta=G_f^*G_f,\quad
 K_f=U_0^*\mathsf B(\mathsf T_f+\mathsf H_f).
\]
Write $\Delta^+$ for its inverse on mean-zero vertex fields and zero on
constants. This notation does not mean that the full constant gauge
integration has been performed.

\begin{lemma}[Exact removal of the pinned inverse from the return]
\label{lem:v5-pin}
The constant-color vertex fields satisfy $K_f\eta_{\rm const}=0$, and
\begin{equation}
 K\mathcal G K^*=K_f\Delta^+K_f^*.
 \label{eq:v5-pin-identity}
\end{equation}
The identity concerns this particular contracted inverse, not an operator
norm comparison of the based inverse and the pseudoinverse.
\end{lemma}
\begin{proof}
For a constant color vector $v$,
$\mathsf B(\mathsf T_f+\mathsf H_f)v=2b_e^0(n\times v)$.
It is zero for $v\parallel n$. The other two vectors are precisely the
cosine color-orbit modes excluded from V4's hard space at zero.
Thus $K_f$ annihilates constants.

Let $I_o$ embed based fields by their zero root value. Then
$K=K_fI_o$ and $\mathcal D=I_o^*\Delta I_o$.
For $f=K_f^*u$, $f$ has zero mean. The full solution
$v=\Delta^+f$ can be shifted by the constant $-v(o)$ to satisfy the based
condition; its nonroot entries solve
$\mathcal D I_o^*(v-v(o))=I_o^*f$.
Applying $K_f$ removes the constant shift. This proves
\eqref{eq:v5-pin-identity} for every $u$.
\end{proof}

This result explains why deleting the two limiting color-orbit columns from
V4's excluded frame would spoil a later estimate, even if one still counted
only nonzero scalar eigenvalues. They are needed for the \emph{source vertex},
not just the rank convention.

\subsection{Two elementary lattice sums}

For $p\in(2\pi/L)\mathbb Z^3/(2\pi\mathbb Z)^3$, put
$\lambda(p)=4\sum_{\mu=1}^3\sin^2(p_\mu/2)$ and let $|p|$ be the norm of
its nearest representative in $[-\pi,\pi]^3$. With $\kappa=|k e_1|=k$,
all constants below are absolute and independent of $L,k,T$ in their stated
ranges. Define $\log_+(x)=\max\{0,\log x\}$.

\begin{lemma}[Infrared convolution bounds on the periodic lattice]
\label{lem:v5-lattice-sums}
For $L\ge3$ and $0<k=2\pi m/L<\pi$,
\begin{align}
 g_L(0):=\frac1V\sum_{p\ne0}\lambda(p)^{-1}&\le C_g,
                                                    \label{eq:v5-g0}\\
 \frac1V\sum_{p\ne0,\ p+k e_1\ne0}
        \frac1{\lambda(p)\sqrt{\lambda(p+k e_1)}}
                       &\le C\{1+\log_+(1/\kappa)\},
                                                    \label{eq:v5-logsum}\\
 \frac1V\sum_{p\ne0,\ p+k e_1\ne0}
        \frac1{\lambda(p)\lambda(p+k e_1)}&\le C/\kappa.
                                                    \label{eq:v5-coulombsum}
\end{align}
Addition of momenta and the omitted zero conditions are periodic.
\end{lemma}
\begin{proof}
For the nearest representative,
$4|p|^2/\pi^2\le\lambda(p)\le|p|^2$.
A dyadic shell $r\le|p|<2r$ has at most $C(Lr)^3$ lattice points when
$r\ge2\pi/L$; enclosing it in a cube proves this bound. Consequently
\[
 V^{-1}\sum_{0<|p|\le r}|p|^{-2}\le Cr,\qquad
 V^{-1}\sum_{0<|p|\le r}|p|^{-1}\le Cr^2
\]
for $r$ above the mesh, and the sums are zero below the first mesh radius.
Summing $Cr$ over the dyadic shells up to the fixed torus diameter proves
\eqref{eq:v5-g0}.

For \eqref{eq:v5-logsum}, split the torus into the ball $|p|<\kappa/2$,
the ball $|p+k e_1|<\kappa/2$, and their complement. In the first ball the
other distance is at least $\kappa/2$, giving a bound
$C\kappa^{-1}\cdot\kappa=C$. In the second, use
$C\kappa^{-2}\cdot\kappa^2=C$. On the remaining region with
$|p|\le2\kappa$, both denominators have the lower distances
$\kappa/2$, while the number of terms is at most $CV\kappa^3$;
this contributes at most a constant. On $|p|>2\kappa$, the triangle
inequality on the flat torus gives $|p+k e_1|\ge|p|/2$.
Each remaining dyadic shell contributes at most a constant, and there are
at most $C(1+\log_+(1/\kappa))$ shells.
The same decomposition with powers $(2,2)$ gives $C/\kappa$ in each near
ball and in the middle region. Its far shells contribute $C/r$, whose
geometric sum is at most $C/\kappa$. This proves
\eqref{eq:v5-coulombsum}. The condition $\kappa\ge2\pi/L$ controls all
mesh-endpoint estimates uniformly.
\end{proof}

\subsection{Orbit and return estimates on the actual plane-wave family}

\begin{theorem}[Physically normalized volume bound for the added coefficient]
\label{thm:v5-volume-bound}
Let $H_0$ be the actual curl Hessian on V4's hard space at zero, and use the
bare native data above. Then, for an absolute constant $C$,
\begin{align}
 0\le-\mathcal E_{\rm orb}&\le(16C_g+\tfrac12)\|b\|^2,
                                                   \label{eq:v5-orb-bound}\\
 0\le-\mathcal E_{{\rm ret},T}
      &\le C\|b\|^2
       \left\{1+\log_+(1/\kappa)+\frac1{T\kappa}\right\},
                    \qquad T\ge3,                \label{eq:v5-return-bound}\\
 0\le-\mathcal E_{\rm ret}
      &\le C\|b\|^2\{1+\log_+(1/\kappa)\}.
                                                   \label{eq:v5-return-inf-bound}
\end{align}
Here $\|b\|^2=V/2$. Thus division by the classical temporal profile norm
removes the extensive factor. At fixed nonzero external momentum the
infinite-time added coefficient is bounded independently of spatial volume.
For a finite time cycle the additional condition $T\kappa\ge c>0$ gives
the corresponding common bound. No infrared-uniform bound as
$\kappa\downarrow0$ is asserted.
\end{theorem}
\begin{proof}
For a scalar vertex color, the based inverse kernel is
\[
 \mathcal G(x,y)=g_L(x-y)-g_L(x-o)-g_L(y-o)+g_L(0),
 \quad
 g_L(x)=V^{-1}\sum_{p\ne0}\frac{e^{ip\cdot x}}{\lambda(p)}.
\]
To verify the formula, apply $\Delta$ and note that it vanishes at the
root; the omitted root equation is exactly the based Dirichlet condition.
Since $|g_L(x)|\le g_L(0)$, its diagonal is at most $4g_L(0)$.
Cauchy--Schwarz for this positive covariance gives
$\E|\eta_x+\eta_y|^2\le16g_L(0)$ in one color. The first square in
\eqref{eq:v5-negative-square}, with $I-P_g$ bounded by $I$, therefore
costs at most $32g_L(0)\sum_e(b_e^0)^2$ before its prefactor $1/2$;
there are two transverse colors. The second square is
$\Tr(\mathsf B_\parallel P_g\mathsf B_\parallel^*)$,
which is at most $\sum_e(b_e^0)^2$, since $0\preceq P_g\preceq I$.
Equation \eqref{eq:v5-g0} proves \eqref{eq:v5-orb-bound}.

For the return use \eqref{eq:v5-pin-identity}. At zero, the hard projector
commutes with the curl Hessian: each removed plane wave is a transverse
curl eigenvector, and any mixing of its real sine/cosine representatives
is inside an equal-eigenvalue space. The extension of the hard covariance
$C_0+C_1$ to ambient link space is therefore bounded above by the Fourier
multiplier $\lambda(p)^{-1/2}I$ on nonzero momenta and zero on constants.
Indeed its scalar value is
\[
 c_+(\lambda)=\frac{1+e^{-\xi}}{2\sinh\xi}
             =\frac1{e^\xi-1}\le\lambda^{-1/2},
 \qquad \xi=\operatorname{arcosh}(1+\lambda/2).
\]
Removing additional hard directions can only decrease this positive
multiplier on each curl eigenspace.

The Fourier map $\mathsf B(\mathsf T_f+\mathsf H_f)$ sends a vertex mode
$p$ only to $p+k e_1$ and $p-k e_1$, in link direction $\nu$.
Each coefficient has norm at most one, because its cosine factor is
$1/2$ and the tail-plus-head symbol has norm at most two. Its color matrix
$J$ has squared Hilbert--Schmidt norm two. Hence positivity and the ambient
multiplier bound give
\[
 \Tr\{(C_0+C_1)K_f\Delta^+K_f^*\}
 \le2\sum_{s=\pm1}\sum_{p\ne0,\ p+s k e_1\ne0}
                  \frac1{\lambda(p)\sqrt{\lambda(p+s k e_1)}}.
\]
If the extra real-mode exclusions mix opposite Fourier labels, use the
positive ambient multiplier inequality \emph{before} expanding this trace;
no translation-invariance assertion for that projector is needed.
The $p=0$ input was removed exactly by \cref{lem:v5-pin}.
Apply \eqref{eq:v5-logsum}, the factor two in
\eqref{eq:v5-return}, and $V=2\|b\|^2$ to obtain
\eqref{eq:v5-return-inf-bound}.

For finite $T$, periodizing the positive scalar covariance gives
\[
 c_{+,T}(\lambda)=\frac1{e^\xi-1}
                +\frac{\coth(\xi/2)}{e^{T\xi}-1}
       \le\lambda^{-1/2}+\frac{16}{T\lambda},\qquad0<\lambda\le12.
\]
For the inequality, $e^{T\xi}-1\ge T\xi$,
$\coth(\xi/2)\le1+2/\xi$, and
$\xi=2\operatorname{arsinh}(\sqrt\lambda/2)\ge\sqrt\lambda/2$
on this range; thus $\lambda\coth(\xi/2)/\xi<16$.
The extra trace is bounded by \eqref{eq:v5-coulombsum}, proving
\eqref{eq:v5-return-bound}. The signs are from \cref{thm:v5-sign}.
\end{proof}

These estimates control the \emph{new} gauge/density correction, not every
term of V4's determinant response. They do not turn a plane-wave normalized
bound into a rooted spatial norm for arbitrary local sources. The logarithmic
allowance has not been replaced by a calculated logarithmic coefficient.
At the minimal momentum $k=2\pi/L$ the bound allows $O(\log L)$, plus
$O(L/T)$ for a finite time cycle, after profile normalization. This exposes
the two limit parameters rather than hiding them in a finite hard-gap bound.

\section{Measured identification and the next transfer}
\label{sec:v5-transfer}

The measured object calculated here is precise: the local conditional
high-fibre coefficient of the \emph{based-averaged native transfer}, in the
source chart \eqref{eq:v5-full-slice}, with all thirteen low coordinates and
the residual global projector retained. It includes the actual relative
based-group Gaussian, its native curvature vertices, the hard--vertical
Schur return, and the leading Haar density of that source chart. No
unspecified determinant or density remains in its source two-jet.
The passage from V4's determinant to \eqref{eq:v5-total-electric} is therefore
a proved measured identification on this explicitly limited branch.

Three further identifications would be invalid. First, based averaging is
not replacement of the full residual global integral by its value at the
identity. That operation remains in the low theory; at zero background its
constant modes do not have a nonsingular Gaussian Hessian. Second, the
thirteen-dimensional retained family used here is not the four-dimensional
nonlinear root block. The two constructions must be compared with their
actual retained source maps and induced densities. Third, a bare local
stationary coefficient is not an analytic estimate for the generated action
or its compact complement. In particular \eqref{eq:v3-dressed-pair} is not
estimated by the present transfer calculation.

If magnetic and electric coordinates are extracted on this branch, their
local one-loop magnetic two-jet agrees with V4 and their electric two-jet
is \eqref{eq:v5-total-electric}. Their equality is not imposed. Changing the
retained coordinates requires transforming the reference source density as
well as the action. The exact measured-redundancy theorem of \cite{YM7}
remains the appropriate rule; deleting the density or the Schur return would
change the defined source functional.

The useful next P4 task is now a source-resolved comparison on an actual
\emph{spatial shell}, not another static-spectrum fit. It must retain the
relative based and global data, the low normalization, and the cross-shell
returns when the V4 hard fibre is split. A uniformly local shell covariance
and its source map must be constructed before replacing the plane-wave
normalization by the rooted source norms of P1. The coefficient of a
logarithm, even if subsequently calculated, is not a substitute for that
measured comparison. On P2's parallel branch, the same-law collar-dressed
activity bound remains the first complete compact analytic task.

\subsection{Verification, limitations, and preservation}

The accompanying script \srcid{check_measured_gauge_response_V5.py}
verifies the four-factor quaternion identity in exact rational polynomial
algebra. It then assembles the based incidence maps, V4 hard space, and
actual curl Hessian at $L=3$. Its floating-point diagnostics check the
vertical and mixed jets, the negative-square identity, the contracted
pinned-inverse identity, and the finite-cycle response by direct positive
matrix determinants. The last check sets the higher magnetic source jets
to zero on \emph{both} compared matrices: the added gauge/density two-jet
is independent of those jets, so this isolates precisely the asserted
addition rather than reporting a lattice beta coefficient. It uses the
actual native temporal and gauge matrices and the actual zero-background
curl Hessian. Tolerances and source steps are recorded. These numerical
checks are not interval proofs or estimates uniform in a cutoff.

\begin{center}
\small
\begin{tabular}{>{\raggedright\arraybackslash}p{.23\textwidth}>{\raggedright\arraybackslash}p{.68\textwidth}}
\toprule
Variable & Exact scope in V5\\
\midrule
Spatial volume & Source-chart existence and local stationary sections are
proved at each fixed $L$. The added electric coefficient divided by $V/2$
has the volume-independent bounds in \cref{thm:v5-volume-bound}; no
volume-independent chart radius is asserted.\\
Time length & The measured two-jet formula holds at every $T\ge3$ on its
finite-regulator domain. Its bound records $1/(T k)$ explicitly. Infinite
time is taken first at fixed $L$ when using $C_0,C_1$.\\
Cutoff & No invariant source chart or rooted local kernel norm through
spatial cutoff changes is proved. At fixed nonzero lattice momentum the
stated normalized addition has no volume loss; the infrared allowance
remains logarithmic.\\
Loop and source order & One-loop order and total amplitude degree two on
one declared spatial profile. Frequency dependence at this degree is
exact. No arbitrary-order source, coupling, or nonlinear remainder theorem
is inferred.\\
RG steps & The relative based integration is an actual local operation.
There is no estimate uniform over arbitrarily many interacting steps and
no supplied running trajectory.\\
Coupling and sectors & Formal stationary coefficients of the positive bare
native law, with $h=\gamma^{-1}$. Its compact complement, generated
interactions, residual global integration, and SM chiral weights retain
separate obligations.\\
\bottomrule
\end{tabular}
\end{center}

\begingroup
\renewcommand{\refname}{External methodological context for V5}
\begin{thebibliography}{99}
\bibitem[12]{V5LuscherWeisz}
M.~L\"uscher and P.~Weisz,
\emph{Background field technique and renormalization in lattice gauge theory},
Nuclear Physics B \textbf{452} (1995), 213--233,
arXiv:hep-lat/9504006, doi:10.1016/0550-3213(95)00346-T.
The all-orders BRS/background-field framework is context only; no
renormalization or locality theorem for the present transfer map is imported.
\end{thebibliography}
\endgroup

\section{Final ledgers for V5}
\label{sec:v5-ledgers}

\subsection*{Proved}
\addcontentsline{toc}{subsection}{V5: Proved}
V4's hard fibre embeds in an actual local based-quotient source chart with
leading density \eqref{eq:v5-source-density}. The native relative-gauge
word gives the complete vertical and mixed blocks
\eqref{eq:v5-M}--\eqref{eq:v5-Lplus}, their symmetric endpoint jets, and a
joint stationary section whose displacement has source degree at least
three. The resulting measured local coefficient is
\eqref{eq:v5-complete-measured}. Its difference from V4 is exactly
\eqref{eq:v5-response}, with the explicit orbit/density and Schur-return
coefficients \eqref{eq:v5-orb}--\eqref{eq:v5-return}. Their signs follow from
the negative-square identity and positivity. The mixed vertex annihilates
constant gauge fields, allowing the exact contracted inverse replacement
\eqref{eq:v5-pin-identity}. Elementary lattice sums then prove the
physically normalized volume and finite-time bounds in
\cref{thm:v5-volume-bound}. These results have precisely the loop, source,
sector, and uniformity scope in the preceding table.

\subsection*{Inherited}
\addcontentsline{toc}{subsection}{V5: Inherited}
The native Wilson kernel, the based quotient with its exact half-density,
and the residual global action come from the stated companion locations.
V4 supplies the hard bundle, zero-background curl spectrum, hard recentering
calculation, temporal determinant response, and covariance. V1 supplies the
measured stationary coefficient rule and source tower. The monographs'
Schur and measured-redundancy identities are reused with the actual new
blocks, not asserted as new general principles. V2's global/local space
and exact sector coordinates and V3's guarded estimates remain unchanged.

\subsection*{Conditional}
\addcontentsline{toc}{subsection}{V5: Conditional}
A physical running electric coordinate requires a retained-source and
spatial-shell identification of this based-quotient coefficient, with the
residual global data kept and with Ward-compatible local extraction.
Extension to generated interactions requires their actual source chart,
force, Hessian, and density jets; bare formulas do not certify it.
An analytic one-step map further requires the local-order remainder and
complete compact-sector bound. The displayed logarithmic bound is not a
calculation of a beta coefficient or a proof of ultraviolet asymptotic
freedom. No result about chiral line transport is obtained by the positive
based-group integration.

\subsection*{Open}
\addcontentsline{toc}{subsection}{V5: Open}
The next coefficient interface is a uniformly local, source-complete
spatial-shell realization of the measured transfer and its retained
low/global data, with the mixed returns assigned before comparison to the
four-dimensional root source. The parallel first analytic compact task
remains the collar-dressed sector cumulant and its higher decorated norm
from V3. An invariant generated interaction class, internally generated
joint marginal/stable trajectory, composed renormalized source transport,
Ward-compatible subtraction, SM line coherence, and the perturbative
continuum and regulator-independence theorems are not yet established.
The new measured based coefficient closes the specific missing relative
based-gauge and density two-jet, not these remaining gates.


\clearpage
\part*{V6: Einstein--Standard Model EFT, weighted trajectories, and a native free gravitational shell}
\addcontentsline{toc}{part}{V6: Einstein--Standard Model EFT and a native gravitational reference}

\section{The enlarged endpoint and the nonduplication audit}
\label{sec:v6-start}

The updated programme \cite{V6Programme} replaces the SM--YM endpoint by a
source-complete \emph{Einstein--Standard Model effective field theory} (ESM EFT).
Sections~1--39 are preserved. They supply gauge--matter and measured-source
interfaces; they are not retrospectively promoted to a quantum gravitational
construction. Throughout this installment, $a$ denotes spatial mesh,
$\tau$ temporal mesh, and $\hbar$ loop-counting normalization. Neither is the
inverse gauge coefficient denoted by $h$ in the earlier installments.

The requested extension has an important previously available input. The
separate measured-continuum article \cite{V6CMP},
\srcid{hyp:rg-class}, \srcid{thm:coefficient-trajectory},
\srcid{thm:source-continuum}, and \srcid{prop:scheme}, already proves the
R1--R4 two-boundary coefficient theorem, ultraviolet boundary loss, source
landing, and matched-scheme comparison. Its definition of the local bank
\emph{already permits dependence on the retained degree and composite-source
panel}. Its \srcid{thm:slavnov-splitting} also already gives the exact
finite-dimensional criterion for a commuting local projection. We import these
results, including their hypotheses, rather than proving them again under an
ESM name. They describe an admissible shell class, not a verified arbitrary
compact Renewal regulator.

The actual new issue for that theorem is not merely the number of tensor
types. A dimension-$d>4$ Wilson coefficient fixed at the physical scale has
inverse dilation factor $b^{d-4}>1$. This violates the particular
$\|A_j^{-1}\|\le1$ normalization of \cite{V6CMP}. The first result below repairs
that interface by \emph{degreewise scale weights}, with an explicit gap between
local growth and complementary decay. We then construct and estimate a
nonzero-momentum free gravitational shell from the inherited finite action's
flat quadratic constraints. The linear gravitational ghosts and the
constraint/gauge Jacobian are evaluated, not merely named. The construction
stops short of a nonlinear coframe--connection quantum measure, and this
limitation is tested by a native high-momentum symmetry defect.

\subsection{Which geometric and matter data are actually inherited}

The new Lorentzian source \cite{V6Spacetime} supplies the calibrated coframe,
metric-compatible connection, Palatini--Holst branch, physical variation
spaces, and a particular phase-compatible canonical flat jet. The relevant
locations are \srcid{thm:main-gravitational-classification},
\srcid{thm:supp-palatini-torsion}, \srcid{lem:supp-initial-calculus}, and
\srcid{lem:supp-initial-balance}. Its constructive harmonic evolution and its
literal finite-action evolution are explicitly distinguished in that source.
They are not interchangeable quantum regulators.

The structural source \cite{V6Duality} supplies the same-carrier internal
packet and the explicit action \srcid{eq:finite-SM-action}: gauge curvature,
covariant Higgs differences, Higgs potential, and a finite spin--Dirac/Yukawa
operator with independent dual fermionic amplitude. Its gauge Ward, internal
BRST, and stress identities are \srcid{thm:finite-action-identities}. The
classical ESM handoff \srcid{prop:classical-compatibility} requires convergent
common first variations and stationarity. It is not a quantum statement.
The spectral source \cite{V6Spectral} supplies a compatible differential and
spin-discretization interface; its spectral reconstruction does not determine
a gravitational integration measure or its normalization.

The finite common action remains
\begin{equation}
 \mathscr S_X^{\rm common}(e,\omega,U,H,\Psi,\overline\Psi)
 =\mathscr S_{g,X}(e,\omega)+
   \mathscr S_{{\rm SM},X}^{L}(e,\omega,U,H,\Psi,\overline\Psi).
 \label{eq:v6-common-action}
\end{equation}
Its actual finite representatives, not continuum Feynman rules, determine the
vertices. At this stage the supplied classical action does not by itself fix a
unique quantum measure, contour, gravitational gauge-fixing, or nonlinear
finite diffeomorphism complex. These are explicit remaining data/estimates.
For the reference calculation we select the flat, zero cosmological-energy
branch, with vanishing matter fluctuations. Choosing a different classical
background requires its own stationary and Hessian analysis.

\subsection{Literature comparison and the meaning of the claim}

Gravity as a quantum EFT is established methodology \cite{V6Donoghue}.
The one-loop calculation of 't~Hooft--Veltman \cite{V6tHV}, the two-loop
calculations of Goroff--Sagnotti and van~de~Ven \cite{V6GS,V6vdV}, and the
locally covariant, renormalized BV construction of Brunetti--Fredenhagen--Rejzner
\cite{V6BFR} preclude presenting ``perturbative gravity exists'' as the new
result. Gravity coupled to gauge fields also has explicit early one-loop
calculations \cite{V6DN}. Lattice convergence and gauge-theory constructions
have separate regulator hypotheses \cite{V6Reisz,V6EGK}; none is inferred from
agreement of one quadratic symbol. These works are comparison and normalization
context here, not imported estimates for the finite action
\eqref{eq:v6-common-action}. No numerical curvature-counterterm coefficient is
borrowed to conceal an uncomputed native loop.

\section{A source-complete finite EFT bank and its ancestry}
\label{sec:v6-bank}

\subsection{Covariant tensor types, rather than finitely many graviton monomials}

Fix a spacetime dimension four, the finite structural SM representations,
a finite source panel, a perturbative bookkeeping cutoff $N$, and an EFT
weight cutoff $\Delta$. On a nondegenerate coframe chart take covariant jets
of curvature, torsion or contorsion, gauge curvature, Higgs and spinor fields.
Assign weights
\[
 [D]=1,\quad [R]=2,\quad [T]=[K]=1,\quad [F]=2,
 \quad[H]=1,\quad[\Psi]=[\overline\Psi]=3/2.
\]
Coframes, inverse coframes and invariant metrics are contraction data;
one volume density is included. They are not independent zero-weight scalar
generators of arbitrary functions. Gauge-covariant derivative indices and
representation indices must be completely contracted, or paired with the
specified covariant sources. A source at a fixed coordinate is not thereby a
relational, diffeomorphism-invariant observable: its transformation law belongs
to the enlarged source/BV packet.

Let $\mathcal U^{\rm ESM}_{N,\Delta}$ be the span of these allowed tensor
\emph{types}, with all mixed source contacts through the fixed source order.
Ghost and antifield monomials additionally have fixed positive bookkeeping
budgets and fixed ghost number. This additional budget is essential when a
physical source or an antifield has nonpositive canonical dimension. Quotient
only by the declared algebraic identities, integration by parts with the
chosen boundary completion, and represented measured redundancies. A
cohomological quotient and a bounded choice of representatives are different
operations. In a background-dependent symmetry-breaking bank, powers of the
coframe fluctuation additionally carry positive perturbative degree; a fixed
Taylor budget is retained. Arbitrary zero-cost functions of the reference and
dynamical metrics are not admitted as finitely many counterterms.

\begin{proposition}[Finite types at fixed order]
\label[proposition]{prop:v6-finite-bank}
With these conventions, $\mathcal U^{\rm ESM}_{N,\Delta}$ is finite dimensional.
There is a nested finite family when $N,\Delta$ and source budgets increase.
This does not require a uniform dimension bound as the order tends to infinity.
\end{proposition}
\begin{proof}
Every positive-weight covariant factor and every derivative consumes at least
one half-unit of the fixed weight. There are finitely many species and
representation components, hence finitely many words of the allowed lengths
and finitely many contractions of their indices. The fixed source/ghost and
antifield budgets leave finitely many further choices. Metric/coframe
contractions may be reduced using the inverse relations, leaving tensor
contractions of this finite list and the single density. Taking a quotient
of a finite span preserves finiteness. Enlarging the budgets includes the
previous types; dependent representatives may be removed after inclusion.
\end{proof}

This elementary finiteness statement does not prove that a regulator divergence
is local or symmetry allowed. It specifies a bank on which those questions
can be tested. A covariant type such as $\sqrt{|g|}R^2$ includes its full
geometric expansion; a finite Taylor packet of that expansion is a separate
quotient, as in V1.

The initial list contains volume, Palatini--Holst and Einstein terms, curvature
quadratics, and $R H^\dagger H$. At higher weights it admits $RF^2$,
$R(DH)^\dagger DH$, curvature-cubics and mixed matter operators. A schematic
$R\overline\psi\psi$ must respect the actual chiral representation: for a
charged SM Yukawa channel a scalar contraction instead includes the Higgs,
for example $R\overline L H e_R$ together with its adjoint. The optional
neutral Majorana channel has a different allowed type. These are type
requirements, not a computation of their Wilson coefficients.

In a first-order gravitational formulation, the torsion sector cannot be
removed using the spinless connection equation when fermions are present.
The inherited common-action equation is
\[
 -D(P(e\wedge e))-\Sigma=0,
\]
where $\Sigma$ is the spin current; see \cite{Companions}, Grand Tensor,
\srcid{thm:Palatini-torsion}. Eliminating contorsion then also generates
spin-current interactions and a quantum determinant if performed inside an
integral. They remain in the bank or in its uncomputed measure component.
Likewise the four-dimensional torsion-free Euler-density identity is used
only after the torsion-free and boundary assumptions are valid; it does not
justify deleting a source-dependent contact or a boundary term.

\subsection{A lower canonical degree can have a higher EFT ancestor}

Finite EFT filtration is not literal deletion by field degree. For example,
for $\eta\sim N(0,v)$,
\begin{equation}
 \E(\phi+\eta)^6
 =\phi^6+15v\phi^4+45v^2\phi^2+15v^3.
 \label{eq:v6-return-example}
\end{equation}
All four terms retain the incoming Wilson-coefficient tag, such as $M^{-2}$.
A hard-cutoff contraction can therefore return a higher-dimensional insertion
to a lower-dimensional operator without removing its EFT ancestry. The
upper-right block in the inherited R1 map has precisely this role. One must
not infer that projection by canonical dimension alone commutes with blocking.

We predeclare a finite \emph{ancestor envelope} of operators for each requested
order: it includes the primitive vertices, the retained source products, and
all local contraction descendants that can contribute through that order.
V1's weighted-jet theorem handles the finite field-degree part of this
construction; it is not reproved. Nonlocal descendants stay in the complement.

For comparison with an eventual power-counting proof, a connected graph whose
propagating fields have canonical dimensions $d_\phi$ and inverse propagator
orders $4-2d_\phi$ obeys
\begin{equation}
 \omega(\Gamma)+\sum_{\text{external legs}}d_\phi
       =4+\sum_{v}(d_v-4).
 \label{eq:v6-graph-dimension}
\end{equation}
Here $d_v$ includes derivatives and field dimensions, and $\omega$ is the
superficial momentum degree. Indeed each loop contributes four, each internal
line $2d_\phi-4$, and each vertex its derivative degree; use
$L=I-V+1$ and count the field dimensions at both ends of internal lines.
For a pure metric gravity graph with two-derivative propagators this gives the
familiar derivative budget
\begin{equation}
 d_{\partial}=2+2L+\sum_v(r_v-2),
 \label{eq:v6-gravity-degree}
\end{equation}
where $r_v$ is the derivative count of its vertex. These are graph identities,
not convergence theorems. Algebraic auxiliary connection variables must first
be eliminated with their actual determinant, or assigned their own propagator
orders, before \eqref{eq:v6-graph-dimension} is applied.

At fixed positive interaction/source bookkeeping degree only finitely many
insertions occur. At fixed gravity-loop and EFT order, covariant completion
organizes their local tensor types. A finite ancestor envelope and a finite
subtraction order can thus be specified, but proving that they control the
\emph{native} integrals still requires the shell estimates below and their
nonlinear extension. Increasing the EFT bank does not absorb an uncontrolled
nonlocal remainder.

\section{The filtered two-boundary extension: growth versus decay}
\label{sec:v6-filtered-flow}

\subsection{The exact point where the older hypotheses change}

We retain the scale orientation of \cite{V6CMP}: $j$ increases toward the
ultraviolet and elimination maps $j+1$ to $j$, with $\Lambda_j=\mu b^j$.
At each fixed positive total degree $n$ in the finite ancestor envelope, let
$\mathcal E_n=\mathcal U_n\oplus\mathcal X_n$. The first factor is the finite
local EFT/source bank; the second is its complete Banach kernel complement.
Both may depend on $n$, but not on the ultraviolet endpoint $J$ after the
prescribed identifications. The highest-degree equations have the inherited
triangular form
\begin{equation}
\begin{split}
 r_{j,n}&=B_{j,n}r_{j+1,n}+g_{j,n}(x_{j+1,<n}),\\
 u_{j,n}&=A_{j,n}u_{j+1,n}+C_{j,n}r_{j+1,n}
                                  +f_{j,n}(x_{j+1,<n}),\qquad x=(u,r).
\end{split}
\label{eq:v6-triangular}
\end{equation}
In particular there is no assumed same-degree local-to-remainder block.
Verifying this block structure in a symmetry-compatible finite bank is still
part of the shell hypothesis.

The new norms allow
\begin{equation}
 \|A_{j,n}^{-1}\|\le b^{s_n},\quad
 \|B_{j,n}\|\le b^{-t_n},\quad \|C_{j,n}\|\le C_n,
 \qquad s_n\ge0,
 \label{eq:v6-growth-decay}
\end{equation}
in place of inverse local contraction. Choose once, for each retained degree,
a real exponent $\sigma_n$ with
\begin{equation}
 \boxed{s_n\le\sigma_n<t_n.}
 \label{eq:v6-window}
\end{equation}
There is no claim of a common $\sigma$ for infinitely many EFT orders.

\begin{hypothesis}[Weighted filtered R1--R4 realization]
\label[hypothesis]{hyp:v6-filtered}
The measured maps and their complete source, density, ghost and phase factors
have the finite triangular coefficients \eqref{eq:v6-triangular}. Their local
banks and terminal physical evaluation are fixed before estimates are taken.
The bounds \eqref{eq:v6-growth-decay}--\eqref{eq:v6-window} hold. With
\begin{equation}
 y_{j,n}=b^{-\sigma_nj}x_{j,n},
 \label{eq:v6-weight}
\end{equation}
the normalized forcing maps
\begin{equation}
 b^{-\sigma_n j}(f_{j,n},g_{j,n})
       \bigl((b^{\sigma_m(j+1)}y_{j+1,m})_{m<n}\bigr)
 \label{eq:v6-normalized-forcing}
\end{equation}
and their multilinear first derivatives have polynomial growth in $1+j$ on
polynomial coefficient families, in the declared quasilocal source norms.
At the two boundaries,
\[
 u_{0,n}=u_n^{\rm ren},\qquad
 \|r_{J,n}\|\le K_n b^{\sigma_nJ}(1+J)^{p_n}.
\]
Terminal normalized Wick/Berezin evaluation at $\mu$ satisfies the older R4,
including a nonzero degree-zero body and a fixed infrared prescription.
All statements are degreewise and compatible with restriction to a smaller
ancestor envelope.
\end{hypothesis}

This is a test on the shell maps, not an assumed trajectory. A sufficient
check of \eqref{eq:v6-normalized-forcing} is useful: for each multilinear
monomial with coefficient norm at most
$K(1+j)^p b^{\gamma j}$ and incoming degrees $m_1,\ldots,m_v<n$, require
\begin{equation}
 \gamma+\sum_{i=1}^v\sigma_{m_i}\le\sigma_n.
 \label{eq:v6-grade-balance}
\end{equation}
Substitution in \eqref{eq:v6-normalized-forcing} then leaves a bounded
exponential factor and polynomial growth. Differentiation with respect to a
normalized lower coefficient leaves the same factor. These finitely many
inequalities are part of the EFT/subtraction design; finiteness of the bank
alone does not imply them.

\begin{theorem}[Filtered coefficient trajectory with physical EFT matching]
\label[theorem]{thm:v6-filtered-trajectory}
Under \cref{hyp:v6-filtered}, each finite cutoff and finite ancestor envelope
has a unique two-boundary coefficient trajectory. There is a unique infinite
trajectory with the same physical local data for which every
$y_{j,n}=b^{-\sigma_nj}x_{j,n}$ has polynomial growth. Put
\[
 \rho_*:=\max_{n\le N}b^{\sigma_n-t_n}<1.
\]
For every $\theta\in(\rho_*,1)$ there are $K_n,p_n,q_n$, independent of $J$,
such that
\begin{align}
 \|x_{j,n}^{(J)}\|&\le K_n b^{\sigma_nj}(1+j)^{p_n},\\
 \|x_{j,n}^{(J)}-x_{j,n}^{(\infty)}\|
   &\le K_n b^{\sigma_nj}(1+J)^{q_n}\theta^{J-j}.
 \label{eq:v6-filtered-rate}
\end{align}
The normalized connected local source coefficients at scale $\mu$ converge,
with error at most $K_{N,r}(1+J)^q\theta^J$ times the specified source test
seminorms. They are independent of the allowed ultraviolet complementary
boundary data, but depend on \emph{all} physical EFT Wilson coefficients and
source normalizations supplied at scale $\mu$.
\end{theorem}
\begin{proof}
Conjugate \eqref{eq:v6-triangular} by \eqref{eq:v6-weight}. Its linear blocks
become
\[
 \widetilde A_{j,n}=b^{\sigma_n}A_{j,n},\quad
 \widetilde B_{j,n}=b^{\sigma_n}B_{j,n},\quad
 \widetilde C_{j,n}=b^{\sigma_n}C_{j,n}.
\]
Their inverse-local, complementary and mixing bounds are respectively
$1$, $b^{\sigma_n-t_n}<1$, and $b^{\sigma_n}C_n$. The nonlinear maps are
exactly \eqref{eq:v6-normalized-forcing}; the endpoint remainder is now
polynomially bounded. At $j=0$ the conjugacy is the identity, so both physical
renormalization data and the terminal evaluation are unchanged.

The transformed system therefore satisfies the inherited R1--R4 theorem
\cite{V6CMP}, \srcid{thm:coefficient-trajectory} and
\srcid{thm:source-continuum}. That theorem permits finitely many
degree-dependent coefficient spaces. Apply it to the present finite envelope,
using the maximum $\rho_*$. Multiplying its coefficient estimates by
$b^{\sigma_nj}$ gives \eqref{eq:v6-filtered-rate}. At the fixed physical
endpoint no factor remains. Its finite polynomial terminal evaluation gives
the source estimate. Uniqueness follows by conjugating any other
weighted-polynomial trajectory and applying the same inherited uniqueness
statement. No all-degree uniform estimate or numerical evaluation of a formal
series is used.
\end{proof}

The new proof is the weighted reduction to the earlier theorem, not a second
proof of that theorem. Enlarging the finite bank at a subsequent EFT order
requires new constants and may require stronger subtraction. It does not
require one finite-dimensional local space at every order.

\begin{proposition}[The growth--decay inequality has real content]
\label[proposition]{prop:v6-window-test}
Neither a finite bank nor a strict unweighted complementary contraction alone
ensures the filtered conclusion. For
\[
 r_j=\rho r_{j+1}+b^{\sigma j},\qquad r_J=0,\qquad 0<\rho<1,
\]
$r_0^{(J)}$ diverges if $\rho b^\sigma\ge1$. Moreover the homogeneous infinite
equation has nonzero solutions of the allowed weighted growth when
$\rho b^\sigma\ge1$.
\end{proposition}
\begin{proof}
Backward substitution gives
$r_0^{(J)}=\sum_{k=0}^{J-1}(\rho b^\sigma)^k$, which grows linearly at equality
and exponentially above it. The homogeneous solution $r_j=c\rho^{-j}$ obeys
$b^{-\sigma j}|r_j|\le|c|$ whenever $\rho b^\sigma\ge1$. Thus neither
convergence nor uniqueness can be inferred at that threshold.
\end{proof}

\subsection{What follows for symmetry, schemes, and order compatibility}

If the local bank includes every descendant required by the finite envelope,
and lower-order maps, source conventions and boundary data agree under
truncation, the larger-order trajectory restricts to the smaller one by
uniqueness. This is the filtered compatibility statement. New independent
Wilson coefficients at higher order are not determined by the lower-order
trajectory.

If two weighted systems admit the inherited local/source conjugacies, with
polynomial bounds \emph{after} \eqref{eq:v6-weight}, and matching terminal
evaluation, the matched-scheme theorem \cite{V6CMP}, \srcid{prop:scheme},
applies to their conjugated maps. A vanishing comparison defect adds its
explicit physical-endpoint error. The gauge group and an agreement of one
low-order coefficient do not provide that conjugacy.

Likewise, if the complete combined BV identities hold through the retained
order, or are restored with a \emph{vanishing total} regulator defect in the
same source norm, they pass to the limiting coefficients. A merely summable
sequence of symmetry errors need not have zero sum. The independent anomaly,
projection and finite-regulator identity hypotheses are retained from
\cite{V6CMP}; the filtered trajectory theorem does not prove them.

\section{The combined symmetry problem is upstream of ghost notation}
\label{sec:v6-symmetry}

The intended continuum symmetry is diffeomorphisms acting on local Lorentz
and internal gauge transformations, with the actual matter representations.
It is semidirect: a vector field differentiates the local gauge parameters.
The direct-sum notation in the programme names the sectors; it is not a claim
that these transformations commute. The extended source bank includes their
ghosts, antighosts, auxiliary fields and antifields, and the geometry dependence
of all SM insertions.

The finite internal BRST theorem of \cite{V6Duality} is inherited. The
finite projection theorem of \cite{V6CMP}, \srcid{thm:slavnov-splitting},
may be used for a combined differential \emph{after} a genuine nilpotent
finite differential and an invariant local subcomplex have been constructed.
At each EFT order it asks for the vanishing of the corresponding cohomological
extension obstruction; it does not supply a cutoff-uniform projection norm.
Nothing in that theorem manufactures the missing diffeomorphism differential.

A source-specific obstruction prevents a naive construction. On the actual
odd-grid phase convention in \cite{V6Spacetime}, let
\begin{equation}
 \widehat{\delta_a f}(m)=i\kappa_a(m)\widehat f(m),\qquad
 \kappa_a(m)=2a^{-1}\sin(\pi a m),\qquad a=1/N.
 \label{eq:v6-phase}
\end{equation}
Products in the following test are ordinary pointwise products of lattice
records. This convention is stated because a different transported product
would require a separate calculation.

\begin{proposition}[The native derivative is not a cutoff-uniform Leibniz rule]
\label[proposition]{prop:v6-Leibniz}
For modes $e_m(x)=e^{2\pi i mx}$ and $e_n(x)$ with $m+n$ below the cutoff,
\begin{equation}
 \delta_a(e_me_n)-(\delta_a e_m)e_n-e_m\delta_a e_n
 =i\{\kappa_a(m+n)-\kappa_a(m)-\kappa_a(n)\}e_{m+n}.
 \label{eq:v6-Leibniz-defect}
\end{equation}
For $N=4m+1$ and $n=m$,
\begin{equation}
 \lim_{m\to\infty}a\,
 \|\delta_a(e_m^2)-2e_m\delta_a e_m\|_\infty
 =2(\sqrt2-1)>0.
 \label{eq:v6-Leibniz-ultraviolet}
\end{equation}
For each pair of fixed physical mode labels, by contrast, the defect is
$O(a^2)$.
\end{proposition}
\begin{proof}
The Fourier multiplier gives \eqref{eq:v6-Leibniz-defect}. In the displayed
odd sequence $2m=(N-1)/2$ is an admitted mode, so there is no aliasing in this
test. Multiplication by $a$ leaves
$|2\sin(2\pi m/N)-4\sin(\pi m/N)|$, whose limit is
$|2-2\sqrt2|$. At fixed labels, expand sine through cubic order; the linear
terms cancel and the cubic remainder is $O(a^2)$.
\end{proof}

Thus smooth-field consistency cannot be used as a uniform high-frequency
Slavnov estimate. Substitution of $\delta_a$ into continuum transformation
formulas that require the Leibniz identity leaves a defect already visible
in \eqref{eq:v6-Leibniz-defect}. This is not a no-go theorem for a corrected
cochain, homotopy-BV, transported-product or counterterm construction. It is a
proof that such a correction must be constructed, with its weight and source
bounds, instead of assumed.

The conclusion agrees with the source's independent exact-action audit:
\srcid{eq:supp-exact-homogeneity} states global multiplier homogeneity, not
local discrete diffeomorphism identities, and
\srcid{thm:supp-exact-constraint-rank} describes a regular branch with a
finite characteristic quotient different from a freely assumed continuum
first-class constraint system. That inherited theorem is not reproved.
The quadratic reference below uses only the linearized flat symmetry, for
which the necessary constraints do commute exactly.

\section{A native flat gravitational reference with evaluated linear ghosts}
\label{sec:v6-native-free}

\subsection{Connection reduction and the quantization convention}

We use the stationary-connection and Legendre reduction of the same finite
phase-compatible action, as provided in \cite{V6Spacetime},
\srcid{lem:supp-initial-calculus}. The coframe is represented on this reduced
flat branch by the symmetric spatial metric perturbation $u$ and its momentum
$\pi$. This is not substitution of a different continuum action: the source
identifies the following linear constraints and the transverse quadratic
Hamiltonian as jets of its actual finite action. It also specifies that the
remaining link terms start at higher order there.

Stationary classical reduction does \emph{not} authorize deleting a connection
determinant from a quantum integral. We therefore specify the present object
as the canonical reduced \emph{quadratic reference}. A quantum equivalence to
the original coframe--connection measure, the local Lorentz gauge factor and
any contorsion determinant remain in the measure/interaction interface.
For this reference choose canonical Liouville measure on $(u,\pi)$ and the
oriented linear Faddeev--Popov prescription below. These choices are explicit;
they are not claimed to follow uniquely from the spectral construction.

Let the spatial period be $\ell=Na$, with odd $N$, and let
$p\in(2\pi/\ell)\mathbb Z^3$ in the fundamental momentum cube. Write
\[
 \kappa_i(p)=2a^{-1}\sin(ap_i/2),\quad
 \lambda_a(p)=\sum_i\kappa_i(p)^2,
 \quad \Omega_a^2=-\sum_i\delta_i^2.
\]
The zero spatial mode is retained separately; no inverse below acts on it.
Use the orthonormal symmetric-tensor metric $\|u\|^2=\sum_{ij}|u_{ij}|^2$ and
the canonical Poisson form. The inherited constraints at first order are
\begin{equation}
 C_0=\delta_i\delta_j u_{ij}-\sum_i\delta_i^2\operatorname{tr}u,
 \qquad C_i=-2\delta_j\pi^{ij}.
 \label{eq:v6-linear-constraints}
\end{equation}
Choose linear gauge functions
\begin{equation}
 F_0=\operatorname{tr}\pi,\qquad
 F_i=\delta_j u_{ij}-\tfrac12\delta_i\operatorname{tr}u.
 \label{eq:v6-gauge}
\end{equation}
They are genuine conditions on the reduced reference, not claims about a
nonlinear finite gauge orbit.

\begin{theorem}[Exact constraint--ghost cancellation on each nonzero mode]
\label[theorem]{thm:v6-ghost-reduction}
The constraints \eqref{eq:v6-linear-constraints} commute in the linear
canonical Poisson algebra. On each real nonzero-mode block, the absolute
Faddeev--Popov determinant of \eqref{eq:v6-gauge} is
\begin{equation}
 |\det\{F,C\}|=2\lambda_a(p)^4.
 \label{eq:v6-ghost-det}
\end{equation}
Their joint zero set is exactly $u,\pi$ transverse and traceless with respect
to $\kappa(p)$. Let
\begin{equation}
 P_p=I-\frac{\kappa\kappa^*}{\lambda_a},\qquad
 \mathbb P_p A=P_pAP_p-\tfrac12 P_p\operatorname{tr}(P_pA).
 \label{eq:v6-TT-projector}
\end{equation}
Then $\mathbb P_p$ is an orthogonal projector of rank two on symmetric tensors.
With the induced orthonormal TT measures, the exact finite-dimensional identity
is
\begin{equation}
 \delta(C)\delta(F)|\det\{F,C\}|\,\dd u\,\dd\pi
                   =\dd u_{\rm TT}\,\dd\pi_{\rm TT}.
 \label{eq:v6-constraint-measure}
\end{equation}
A fixed ghost orientation realizes the determinant by four Grassmann
antighost--ghost pairs. The corresponding linear BRST differential is
nilpotent. No assertion about its nonlinear completion is included.
\end{theorem}
\begin{proof}
In momentum variables $C_0=\lambda\operatorname{tr}u-\kappa^*u\kappa$.
Its contraction with the symmetrized gradient generating $C_i$ vanishes;
this proves $\{C_0,C_i\}=0$, and the other constraint brackets vanish because
they involve only $u$ or only $\pi$. Using the real sine/cosine blocks, or
rotating the spatial gauge parameters to real symbol conventions, the
Faddeev--Popov matrix has diagonal entries $-2\lambda$ and three entries of
modulus $\lambda$. Its off-diagonal entries vanish. The signs of the spatial
entries depend on whether the factor $i$ is assigned to the gauge parameter;
its absolute determinant is \eqref{eq:v6-ghost-det}.

$F_i=0$ gives $\kappa_j u_{ij}=\kappa_i\operatorname{tr}u/2$.
Inserting this in $C_0=0$ gives $\operatorname{tr}u=0$ and hence
$\kappa_j u_{ij}=0$. The momentum constraints and $F_0=0$ give the analogous
conditions for $\pi$. Since $P_p$ has rank two, its symmetric range has
dimension three; removal of its one-dimensional trace leaves dimension two.
Direct substitution proves that \eqref{eq:v6-TT-projector} is the orthogonal
projection onto this space.

To compute the measure rather than assume a formal cancellation, rotate
$\kappa$ to $(0,0,k)$, $k=\sqrt\lambda$. The configuration normal coordinates
are
\[
 ((u_{11}+u_{22})/\sqrt2,u_{33},\sqrt2u_{13},\sqrt2u_{23}).
\]
The absolute Jacobian to $(C_0,F_1,F_2,F_3)$ is
$\lambda k^3/(2\sqrt2)$. The analogous momentum normal coordinates have
absolute Jacobian $4\sqrt2 k^3$ to $(F_0,C_1,C_2,C_3)$.
Their product is $2k^8=2\lambda^4$, precisely the factor in
\eqref{eq:v6-ghost-det}. Integration of the normal delta functions proves
\eqref{eq:v6-constraint-measure}. Real and imaginary Fourier components are
counted once in orthonormal real blocks, so no extra complex-mode factor is
introduced.

For the ghost statement, let $s_0 z=\{z,C_\alpha\}c^\alpha$,
$s_0c=0$, $s_0\overline c=B$, $s_0B=0$. The commuting linear Hamiltonian
vector fields give $s_0^2=0$. The gauge fermion
$\overline c_\alpha F_\alpha$ gives the auxiliary gauge constraint and the
ghost matrix $\{F,C\}$; its Berezin integral equals the oriented determinant.
The orientation is constant on each nonzero-mode block. This proves the
linear statement without invoking a continuum Ward identity.
\end{proof}

This theorem does not reproduce the source's rank calculation for the initial
constraint map. Its new content is the explicit gauge choice, determinant,
measure cancellation and resulting quantum reference on that actual flat jet.
The homogeneous constraint and global sectors are retained, not priced by a
small inverse eigenvalue or deleted as gauge.

\subsection{The propagating quadratic form and its finite time carrier}

On the TT sector, the source's canonical quadratic Hamiltonian is
\begin{equation}
 H^{[2]}=\|\pi_{\rm TT}\|^2+\frac{1}{4}
                   \langle u_{\rm TT},\Omega_a^2u_{\rm TT}\rangle.
 \label{eq:v6-native-H}
\end{equation}
The symplectic normalization $q=u/\sqrt2$, $p=\sqrt2\pi$ converts it to
$\tfrac12(\|p\|^2+\langle q,\Omega_a^2q\rangle)$ without a Jacobian.
We use its positive Euclidean oscillator reference. This is an explicitly
chosen analytic continuation for the reduced free modes, not a positive
Euclidean measure on all real coframes. The conformal obstruction is already
recorded in the SM--Einstein monograph, Chapter~6, and is not reproved here.

On a periodic time mesh of $T$ sites and step $\tau$, the chosen free reference
has action
\begin{equation}
 S^{[2]}_{a,\tau}(q)=\frac{\tau}{2}\sum_{t=0}^{T-1}
  \left\{\left\|\frac{q_{t+1}-q_t}{\tau}\right\|_{a}^2
                        +\langle q_t,\Omega_a^2q_t\rangle_a\right\}.
 \label{eq:v6-free-cycle}
\end{equation}
This finite time prescription is stated as part of the reference. Equality to
a particular nonlinear time discretization is not inferred from equality of
the canonical quadratic Hamiltonian. The covariance in the physical
$a^3\tau$ inner product is $\hbar\,\mathbb P_p/D_{a,\tau}$, where
\begin{equation}
 D_{a,\tau}(p,\omega)=\lambda_a(p)+
          4\tau^{-2}\sin^2(\tau\omega/2).
 \label{eq:v6-free-symbol}
\end{equation}
It is positive on every nonzero spatial TT mode, with no claim of a lower
bound uniform as spatial volume grows. The next shell construction avoids
that loss by keeping zero and soft modes outside the integrated covariance.

\section{A smooth native graviton shell and its quasilocal bound}
\label{sec:v6-smooth-shell}

\subsection{Positive covariance splitting instead of a sharp spectral deletion}

Fix $\chi\in C_c^\infty((1,4))$, $0\le\chi\le1$, a scale $\Lambda>0$, and
assume
\begin{equation}
 a\Lambda\le c_0,\qquad \tau\Lambda\le c_0,
 \label{eq:v6-mesh-domain}
\end{equation}
for a fixed sufficiently small numerical $c_0$, say $1/8$. Define
\begin{equation}
 \boxed{C_\Lambda(p,\omega)=
  \frac{\chi(\lambda_a(p)/\Lambda^2)}
       {D_{a,\tau}(p,\omega)}\,\mathbb P_p.}
 \label{eq:v6-shell-covariance}
\end{equation}
The zero spatial mode is assigned zero. No global polarization frame is
chosen; the matrix projector \eqref{eq:v6-TT-projector} is sufficient.

\begin{proposition}[Exact measured free split]
\label[proposition]{prop:v6-free-split}
Let $C$ be the full free inverse in \eqref{eq:v6-free-symbol}, on the retained
nonzero-mode space. Then $0\preceq C_\Lambda\preceq C$. Independent normalized
Gaussians $\eta\sim N(0,\hbar C_\Lambda)$ and
$q_<\sim N(0,\hbar(C-C_\Lambda))$ satisfy $q_<+\eta\sim N(0,\hbar C)$.
Consequently
\begin{equation}
 e^{-V_<(q_<,J)/\hbar}
  =\E_\eta e^{-V(q_<+\eta,J)/\hbar}
 \label{eq:v6-exact-free-push}
\end{equation}
is an exact conditional density map whenever the integral is defined, and a
well-defined formal map through every fixed positive interaction/source order.
It keeps every source contact. A covariance partition into several positive
shells composes by the same normalized Gaussian convolution.
\end{proposition}
\begin{proof}
The scalar cutoff lies between zero and one and commutes with the positive
TT inverse. Add the covariance matrices of the independent Gaussians. Their
characteristic functions multiply to that of covariance $\hbar C$, which
proves the identity of laws. Fubini proves \eqref{eq:v6-exact-free-push} for
actual integrable weights; finite Gaussian moments prove its coefficientwise
version. Source differentiation takes place in the same expression. Singular
covariances are interpreted on their supports, so no inverse of the retained
zero mode is used.
\end{proof}

This reuses the general measured Gaussian principle. The new specified input
is the native gravitational constraint reduction and covariance
\eqref{eq:v6-shell-covariance}. A smooth split is not an orthogonal hard-mode
projection: the two independent fields can have overlapping Fourier supports.
Their extra coordinates are auxiliary Gaussian variables, not extra physical
degrees of freedom. Their normalization must be preserved when this reference
is coupled to a nonconstant density.

\subsection{Uniform spatial and temporal moments of the shell kernel}

Let $C_\Lambda(x,t)$ denote the convolution kernel relative to counting weights
$a^3\tau$ on the periodic spacetime mesh, and let $d_{\rm per}$ be the shortest
product distance on that spacetime torus. Spatial and temporal periods are
arbitrary finite multiples of their meshes.

\begin{theorem}[Quasilocal shell estimate on the inherited flat symbol]
\label[theorem]{thm:v6-shell-locality}
For every fixed $s\ge0$, under \eqref{eq:v6-mesh-domain}, there is a constant
$K_s$ depending on $s,c_0,\chi$ and the fixed tensor components, but not on
$a,\tau,\Lambda$, spatial volume or time length, such that
\begin{equation}
 \boxed{a^3\tau\sum_{x,t}
 (1+\Lambda d_{\rm per}((x,t),0))^s
                         \|C_\Lambda(x,t)\|\le K_s\Lambda^{-2}.}
 \label{eq:v6-shell-row}
\end{equation}
On the infinite mesh before periodization, for every fixed integer $m$ the
kernel also satisfies
\begin{equation}
 \|C_\Lambda(x,t)\|
 \le K_m\Lambda^2(1+\Lambda|x|)^{-m}e^{-c\Lambda|t|}
 \label{eq:v6-shell-pointwise}
\end{equation}
with one $c>0$ independent of the same variables. The spatial decay is rapid
polynomial decay, not an asserted common exponential in space.
\end{theorem}
\begin{proof}
On the support of $\chi$, $\lambda_a\in(\Lambda^2,4\Lambda^2)$.
The mesh restriction places the support strictly inside the Brillouin cube,
away from the phase-coordinate seams, and $|p|$ is comparable to $\Lambda$.
The functions $\mathbb P_p$ and $\chi(\lambda_a/\Lambda^2)$ have derivative
bounds $K_\alpha\Lambda^{-|\alpha|}$ there. This follows by differentiating
$\kappa_i=2a^{-1}\sin(ap_i/2)$ and the displayed rational projector; its
denominators are bounded below by $\Lambda^2$.

For fixed $\lambda>0$, set
$\xi(\lambda)=\operatorname{arcosh}(1+\tau^2\lambda/2)$.
The inverse temporal Fourier transform in the physical time measure is
\begin{equation}
 c_\lambda(n\tau)=\frac{\tau}{2\sinh\xi(\lambda)}
                              e^{-\xi(\lambda)|n|}.
 \label{eq:v6-time-kernel}
\end{equation}
Indeed the inverse of
$\lambda+\tau^{-2}(2-2\cos v)$ in the discrete Fourier series is
$\tau^2e^{-\xi|n|}/(2\sinh\xi)$; division by the convolution weight $\tau$
gives \eqref{eq:v6-time-kernel}. On the annulus and mesh domain,
$\xi\asymp\tau\Lambda$ and $\tau/\sinh\xi\asymp\Lambda^{-1}$.
Repeated differentiation with respect to $p$ thus bounds the product of
\eqref{eq:v6-time-kernel}, the projector and the cutoff by
\[
 K_\alpha\Lambda^{-1-|\alpha|}
       (1+\Lambda|t|)^{|\alpha|}e^{-c_1\Lambda|t|}
 \le K'_\alpha\Lambda^{-1-|\alpha|}e^{-c_2\Lambda|t|}.
\]
The constants are uniform since the functions of $\tau\sqrt\lambda$ involved
here extend smoothly to its zero endpoint after these factors have been
removed. Spatial Fourier integration ranges over volume $O(\Lambda^3)$.
Integrating by parts in $p$ an arbitrary fixed number of times, without
boundary terms because the cutoff is compactly supported inside the cube,
yields \eqref{eq:v6-shell-pointwise}.

Choose $m>s+3$. The weighted spatial mesh sum of
$(1+\Lambda|x|)^{-m+s}$ is at most $K\Lambda^{-3}$ since $a\Lambda\le c_0$.
The temporal mesh sum of
$(1+\Lambda|t|)^s e^{-c\Lambda|t|}$ is at most $K_s\Lambda^{-1}$ since
$\tau\Lambda\le c_0$. Multiplying by the factor $\Lambda^2$ in
\eqref{eq:v6-shell-pointwise} proves the row estimate on the infinite mesh.

The periodic kernel is the sum over spatial and temporal period translates of
this absolutely summable infinite-mesh kernel; equality follows either by
Fourier sampling or by the finite Fourier series. For every translate, shortest
periodic distance is no larger than its unwrapped distance. Triangle inequality
therefore bounds the weighted periodic sum by the infinite-mesh sum. This
proves \eqref{eq:v6-shell-row} without a volume or time-length factor.
\end{proof}

The theorem is an actual free gravitational shell estimate, not an application
of the Yang--Mills numerical gap. It supplies the covariance part of the
filtered connected-kernel estimates. It does not prove bounds for arbitrary
interacting coframe vertices, a curved-background reference, local counterterm
projection, or the complete Lorentzian quantum contour. Sending the reference
scale into the infrared also loses the factor $\Lambda^{-2}$; no physical mass
gap is implied.

\section{The first mixed gravity--matter source contribution}
\label{sec:v6-mixed}

The geometric variation of the native matter action has to be used, rather than
a separately supplied stress tensor. Let $\mathsf T_{ij,X}$ be its spatial
metric stress in the convention
\begin{equation}
 D_\gamma\mathscr S_{{\rm SM},X}^{L}[v]
             =-\tfrac12\langle\mathsf T_X,v\rangle_X.
 \label{eq:v6-native-stress}
\end{equation}
This differentiates the Hodge weights, covariant spin operator, mass/Yukawa
and source-dependent terms of the same finite action. All other components,
including lapse, shift and spin-current variations, remain in the common
source packet. The TT source exchange below is not their replacement.

Use $\delta\gamma=\sqrt2\kappa_g q$ for the normalized propagating variable of
\eqref{eq:v6-native-H}; $\kappa_g$ denotes the chosen gravitational expansion
normalization. The linear TT coupling is
$-\kappa_g\langle q,\mathsf T\rangle/\sqrt2$. Let $J_g$ be an additional
scaled gravitational source, included in the action as $-\langle q,J_g\rangle$,
and put
\[
 \mathcal K(\varphi,J)=\kappa_g\mathsf T_X(\varphi,J)/\sqrt2+J_g.
\]
Here $\varphi$ denotes the remaining matter/soft variables; when fermions are
present, $\mathcal K$ is even in the finite Grassmann algebra. The pairings
in the following Gaussian source formulas are the complex-bilinear extensions
of the real field pairing, not sesquilinear absolute-value squares.

\begin{proposition}[Exact mixed current exchange and its source-only contact]
\label[proposition]{prop:v6-stress-exchange}
For the linear hard-field coupling in the reference
\eqref{eq:v6-exact-free-push}, the complete generated contribution is
\begin{equation}
 \boxed{\Delta\mathscr S_\Lambda
       =-\tfrac12\langle\mathcal K,C_\Lambda\mathcal K\rangle.}
 \label{eq:v6-exchange}
\end{equation}
It contains the matter term
$-\kappa_g^2\langle\mathsf T,C_\Lambda\mathsf T\rangle/4$, the mixed
$\mathsf T$--$J_g$ contact, and the source-only term
$-\langle J_g,C_\Lambda J_g\rangle/2$. Every finite source derivative of
$\mathsf T(\varphi,J)$ is included. If its local current vertices have rooted
norms $M_r$ through a declared order, the corresponding fixed-order connected
exchange kernels are bounded by $K_r\Lambda^{-2}$ times the appropriate
products of those stocks and $\kappa_g$ factors.
\end{proposition}
\begin{proof}
The normalized Gaussian of covariance $\hbar C_\Lambda$ has moment-generating
function
\[
 \E\exp(\hbar^{-1}\langle\eta,\mathcal K\rangle)
       =\exp\{(2\hbar)^{-1}\langle\mathcal K,C_\Lambda\mathcal K\rangle\}.
\]
Multiplying its logarithm by $-\hbar$ gives \eqref{eq:v6-exchange}.
Even Grassmann coefficients commute, so the same identity holds by finite
coefficient expansion. Expanding the square gives the three stated terms.
Differentiate this actual square, not a stress insertion with its geometric
source dependence frozen. The kernel bound follows from the single covariance
contraction and \cref{thm:v6-shell-locality}. Use polynomial rooted tree
weights $(1+\Lambda\,\tau(\text{marked sites}))^s$ for these exchange norms.
The connected contraction proof of V1 applies with this weight because
$1+\sum_i a_i\le\prod_i(1+a_i)$ for $a_i\ge0$; no spatial exponential
norm is being inferred. Fixed finite source order contributes only its
finite product-rule factors.
\end{proof}

This is the first explicit mixed contribution in the new direction. It is a
\emph{tree exchange of the specified gravitational reference}, exact for the
linear coupling, not a completed mixed one-loop renormalization. Quadratic and
higher metric dependence of the native matter action, connection/spin exchange,
nonpropagating constraints, and their quantum measure factors remain to be
integrated in the same source convention. A bound on $C_\Lambda$ alone does not
bound the native stress stocks $M_r$ over an arbitrary generated state.
No coefficient of $R^2$, $RF^2$ or a torsion contact is claimed computed here.

The tree-level classical consistency interface is nevertheless fixed. On an
actual stationary eliminated branch $z_*(y,J)$ of the common finite action,
\[
 D_{y,J}\mathscr S^{\rm tree}_{\rm eff}
     =D_{y,J}\mathscr S_X^{\rm common}(y,z_*(y,J),J),
\]
because its eliminated derivative vanishes. This is the inherited measured
stationary rule, not a new reconstruction theorem. With the classical
first-variation convergence hypotheses of \cite{V6Duality},
\srcid{prop:classical-compatibility}, its leading stress and connection rows
therefore have the same classical ESM target. Quantum counterterms and source
contacts are higher coefficients; classical convergence does not determine
them.

\section{What the first gravitational shell test closes}
\label{sec:v6-test}

The weighted theorem and the concrete reference are deliberately tested against
each other. The covariance \eqref{eq:v6-shell-covariance} is positive on its
supported TT modes and has uniform fixed polynomial spatial moments. The linear
gravitational gauge determinant is evaluated on the native flat constraints.
The first stress exchange is in the same measured source calculus. These are
nonempty explicit inputs, not an unrelated continuum graviton propagator.

However, the new shell does \emph{not} verify the full weighted R1--R4 system
for \eqref{eq:v6-common-action}. Three specific interfaces remain.
First, the coframe/connection quantum measure, contour and nonlinear gauge
completion must be related to the reduced reference, including all auxiliary
determinants and the retained homogeneous modes. Second, the complete native
nonlinear vertices must satisfy the ancestor-envelope locality, growth-balance
and subtraction estimates. Third, the combined symmetry breaking must have an
actual restoration or vanishing bound; \cref{prop:v6-Leibniz} shows why
continuum consistency alone is not that bound.

This gives a sharper next attack than a further static YM coefficient. At the
first nonlinear order, form the cubic native gravitational vertices and the
cubic BV defect on the same finite phase/coframe carrier. Determine whether the
defect is a local coboundary in the fixed EFT/source bank, and estimate a
restoring primitive with the same shell weights. If additional finite
geometric incidence or a transported product is required, retain it as a
change in the regulator data and compare it explicitly. Do not simply assign
the continuum Lie bracket to the finite fields.

The full compact gauge collar estimates of V3 and the measured gauge/source
matching of V5 are still required when consumed by this unified calculation.
They remain useful infrastructure, not the entire new endpoint. The first
complete interacting graviton/coframe-plus-ghost shell, its induced curvature
counterterms and mixed loops, and its realization of
\cref{hyp:v6-filtered} remain open; the reference result is not renamed that
theorem.

\subsection{Uniformity and reproducible checks}

\begin{center}
\small
\begin{tabular}{>{\raggedright\arraybackslash}p{.23\textwidth}>{\raggedright\arraybackslash}p{.68\textwidth}}
\toprule
Variable & Exact scope in V6\\
\midrule
EFT and perturbative order & Each finite predeclared ancestor envelope is covered.
Bank dimensions, scale weights and constants can grow with order. No convergent
all-orders series is asserted.\\
Spatial volume and time length & The supported free shell row bound has neither
factor. The free constraint identities hold modewise. Homogeneous modes are
retained, and no uniform interacting or curved-background gap is asserted.\\
Mesh and shell scale & \eqref{eq:v6-shell-row} is uniform for
$a\Lambda,\tau\Lambda\le c_0$, with the explicit $\Lambda^{-2}$ factor.
Weighted trajectory constants are cutoff uniform only under the stated
map-level hypotheses; their full native realization is open.\\
Source order & Fixed finite source order and tensor species. Polynomial moment
constants and current-vertex stocks appear explicitly. No common all-source
analytic radius is constructed for the gravitational sector.\\
Number of shells & The conditional filtered theorem has endpoint-independent
constants at each finite order. The explicit reference proves one-shell
covariance estimates and finite Gaussian composition, not interacting budget
invariance.\\
Coupling and reality & The EFT theorem is coefficientwise. The free oscillator
reference uses a declared positive reduced continuation. No positive full
Lorentzian/coframe law, numerical gravitational UV trajectory or unrestricted
SM coupling domain is inferred.\\
\bottomrule
\end{tabular}
\end{center}

The standalone file \srcid{check_ESM_filtered_shell_V6.py} verifies the linear
constraint brackets, TT projector and normal/ghost Jacobian cancellation in
exact symbolic arithmetic. It also checks the Gaussian current square and
Wick return exactly. Separately labelled numerical diagnostics test a weighted
triangular two-boundary system, its cutoff loss, the periodic time covariance
normalization, and the native high-mode product defect. Those tests do not
certify a nonlinear BV identity, an interacting shell estimate or a continuum
limit; the theorems above do not rely on numerical sampling.

\begingroup
\renewcommand{\refname}{Sources and literature for V6}
\begin{thebibliography}{99}
\bibitem[13]{V6Programme}
Supplied programme update, \emph{Enlarge the perturbative SM/YM continuum
programme to the full Einstein--Standard Model},
\srcid{Pasted markdown(20260919-055340).md}.
\bibitem[14]{V6CMP}
\emph{Measured Renormalization and the Perturbative Continuum}, supplied prior
Library article \srcid{Measured_Renormalization_Perturbative_Continuum_CMP_V3.tex}.
Relevant exact labels: \srcid{hyp:rg-class}, \srcid{thm:coefficient-trajectory},
\srcid{thm:source-continuum}, \srcid{prop:scheme},
\srcid{thm:slavnov-splitting}. These are inherited results, not new V6 theorems.
\bibitem[15]{V6Spacetime}
A.~P\'elissier, \emph{Renewal Geometry and the Emergence of Lorentzian
Spacetime}, supplied text \srcid{Pasted text (2)(20260919-055437).txt}.
\bibitem[16]{V6Duality}
A.~P\'elissier, \emph{A Duality Between Spacetime--Gauge--Incidence Actions
and Matter Multiplicity}, supplied text
\srcid{Pasted text (3)(20260919-055444).txt}.
\bibitem[17]{V6Spectral}
A.~P\'elissier, \emph{From Predictive Dynamics to Spectral Geometry}, supplied
text \srcid{Pasted text (4)(20260919-055451).txt}.
\bibitem[18]{V6Donoghue}
J.~F.~Donoghue, \emph{General relativity as an effective field theory:
The leading quantum corrections}, Physical Review D \textbf{50} (1994),
3874--3888, arXiv:gr-qc/9405057, doi:10.1103/PhysRevD.50.3874.
\bibitem[19]{V6tHV}
G.~'t~Hooft and M.~Veltman, \emph{One-loop divergencies in the theory of
gravitation}, Annales de l'Institut Henri Poincar\'e A \textbf{20} (1974),
69--94.
\bibitem[20]{V6GS}
M.~H.~Goroff and A.~Sagnotti, \emph{The ultraviolet behavior of Einstein
gravity}, Nuclear Physics B \textbf{266} (1986), 709--736,
doi:10.1016/0550-3213(86)90193-8.
\bibitem[21]{V6vdV}
A.~E.~M.~van~de~Ven, \emph{Two-loop quantum gravity}, Nuclear Physics B
\textbf{378} (1992), 309--366, doi:10.1016/0550-3213(92)90011-Y.
\bibitem[22]{V6BFR}
R.~Brunetti, K.~Fredenhagen and K.~Rejzner, \emph{Quantum gravity from the
point of view of locally covariant quantum field theory}, Communications in
Mathematical Physics \textbf{345} (2016), 741--779, arXiv:1306.1058,
doi:10.1007/s00220-016-2676-x.
\bibitem[23]{V6DN}
S.~Deser and P.~van~Nieuwenhuizen, \emph{One-loop divergences of quantized
Einstein--Maxwell fields}, Physical Review D \textbf{10} (1974), 401--410,
doi:10.1103/PhysRevD.10.401.
\bibitem[24]{V6Reisz}
T.~Reisz, \emph{A power counting theorem for Feynman integrals on the lattice},
Communications in Mathematical Physics \textbf{116} (1988), 81--126,
doi:10.1007/BF01239027.
\bibitem[25]{V6EGK}
A.~N.~Efremov, R.~Guida and C.~Kopper, \emph{Renormalization of SU(2)
Yang--Mills theory with Flow Equations}, Journal of Mathematical Physics
\textbf{58} (2017), 093503, arXiv:1704.06799, doi:10.1063/1.5000041.
\end{thebibliography}
\endgroup

\section{Final ledgers for V6}
\label{sec:v6-ledgers}

\subsection*{Proved}
\addcontentsline{toc}{subsection}{V6: Proved}
The fixed-order covariant EFT/source bank is finite under its explicit
tensor and positive-bookkeeping conventions. The weighted conjugacy
\cref{thm:v6-filtered-trajectory} extends the inherited coefficient theorem to
EFT local growth, with the explicit condition
$s_n\le\sigma_n<t_n$ and the cutoff error \eqref{eq:v6-filtered-rate}.
The scalar threshold test proves that the missing growth--decay inequality
cannot be replaced by bank finiteness. On the inherited native flat canonical
jet, the linear gravitational ghost determinant and normal constraint
Jacobian cancel exactly, leaving two TT polarizations per nonzero real mode
(\cref{thm:v6-ghost-reduction}). The specified smooth covariance shell is
positive and satisfies the mesh-, volume- and time-uniform fixed-moment estimate
\eqref{eq:v6-shell-row}. Its linear matter-stress/gravitational-source coupling
has the exact exchange and source-only contact \eqref{eq:v6-exchange}. The
native phase derivative has the explicit ultraviolet Leibniz defect
\eqref{eq:v6-Leibniz-ultraviolet}. These statements have the scopes recorded
above; a free reference is not a complete interacting shell.

\subsection*{Inherited}
\addcontentsline{toc}{subsection}{V6: Inherited}
V1--V5 remain unchanged. The prior measured-continuum article supplies the
R1--R4 coefficient trajectory, source convergence, matched schemes and finite
Slavnov splitting. Its already allowed degree-dependent local banks are not
claimed as new. The Lorentzian source supplies the native phase constraints,
TT canonical quadratic form, coframe/connection branch and classical
variation interfaces. The structural source supplies the full finite SM
representations and common matter action, its internal gauge identities and
geometric stress definition. The spectral source supplies its declared
compatible spin/differential interface, not a quantum gravity measure.

\subsection*{Literature input}
\addcontentsline{toc}{subsection}{V6: Literature input}
The EFT interpretation, historical gravitational loop calculations, curved-space
BV perturbation theory and regulator-specific gauge/lattice constructions are
external comparison inputs cited above. Their existence is not a new result of
this lineage. No external counterterm number or beta function is substituted
for a measured coefficient of the native ESM regulator.

\subsection*{Conditional}
\addcontentsline{toc}{subsection}{V6: Conditional}
The full source-complete ESM coefficientwise continuum follows from the
filtered theorem only for a family verifying
\cref{hyp:v6-filtered}, the combined symmetry-restoration hypotheses and the
matched physical/source conventions. The theorem proves their consequence;
it does not prove membership of the full native ESM system. The canonical
reduced reference requires a separate quantum measure/connection comparison
to represent the original first-order integral. Curved backgrounds, actual
mixed loop vertices, spin-current and auxiliary determinants, and the phase of
the chiral sector require their own same-family estimates. Classical ESM
stationarity follows at tree level only with the inherited common-action
convergence and variation hypotheses.

\subsection*{Open}
\addcontentsline{toc}{subsection}{V6: Open}
The next smallest gravitational interface is the first nonlinear native
BV/constraint defect and its local restoring primitive, with cutoff-uniform
weighted bounds on the actual coframe/connection carrier. It must address the
high-mode product defect rather than invoke a continuum Leibniz identity.
The first complete interacting graviton/coframe-plus-ghost shell remains to be
constructed, including its measure, curvature counterterms, mixed matter
loops and source contacts. Verification of the filtered growth-balance and
complement estimates, the gauge compact collar, the chiral line, and matched
ESM source landing remain subsequent or coupled tasks. No finite-parameter
all-orders Einstein--Hilbert renormalizability, asymptotic safety, numerical
quantum-gravity UV trajectory, nonperturbative ESM completion, or mass gap is
claimed.

\clearpage
\part*{V7: A pole-free mixed gravitational Ward repair on the finite Higgs carrier}
\addcontentsline{toc}{part}{V7: A pole-free mixed gravitational Ward repair}

\section{A first nonlinear ESM component that can be completed explicitly}
\label{sec:v7-start}

Sections~1--48 are preserved. The next question is not whether a free TT
covariance can be written again, but whether a nonlinear vertex of the common
matter--geometry action can satisfy a finite symmetry identity with controlled
locality. We address the first metric--Higgs component, including its matter
antifield partner. This is a necessary component of the full ESM deformation,
not a replacement for the pure gravitational cubic vertex.

The input from \cite{V6Duality}, \srcid{eq:finite-SM-action}, is the covariant
edge Higgs form $\frac12\langle D_UH,\star_1D_UH\rangle$ and its geometric
variation. That source does not identify this form with the particular
momentum-space vertex below. We therefore specify the local Hodge realization
used in this installment as additional regulator data. The comparison to its
unrepaired vertex is calculated explicitly; no equality to every other choice
of Hodge writer is asserted. The same finite field carrier, link differences,
phase symbol, and physical geometric sources are retained.

We restrict to the flat $U=I$, zero-Higgs background and to the coefficient
with two Higgs fields and one metric perturbation. There are $n_H=4$ real Higgs
components, with identical formulas for each. The positive scalar quartic,
Yukawa and gauge couplings remain in the ESM state but do not occur in this
coefficient. A mass $m^2\ge0$ can be included in the reference; $m=0$ is the
unbroken $v_H=0$ branch of the supplied Higgs action. A reference mass different
from the physical mass must have its compensating insertion retained.

The result is an exact cubic Ward/BV repair on a declared nonaliased source
window, a local EFT expansion of the repair, and its explicit first measured
Higgs-shell contribution. It does not supply a nonlinear diffeomorphism group
on the full lattice. In particular the original action's regular
characteristic quotient, described in \cite{V6Spacetime},
\srcid{thm:supp-exact-constraint-rank}, is not changed by a verbal declaration
of additional first-class constraints.

\subsection{The finite local Hodge writer and its measured metric vertex}

For clarity use equal spatial and temporal meshes $a$ and odd periods
$L_\mu/a$, $\mu=0,1,2,3$. This is a subfamily of the reference domain of V6.
The Euclidean scalar reference is positive. No unrestricted positive Euclidean
coframe measure or nonlinear Lorentzian contour is introduced. Set
\[
 D_\mu^+=(T_\mu-I)/a,\qquad D_\mu^-=(I-T_\mu^{-1})/a,
 \qquad \mathfrak v=a^4.
\]
All tensors in this installment have Euclidean contraction convention.
For a positive metric record $g_x$ define
\begin{equation}
 \begin{split}
 S_{H,a}[g,\phi]={}&\frac{\mathfrak v}{2^{5}}
   \sum_x\sum_{\epsilon\in\{+,-\}^{4}}
       \sqrt{\det g_x}\,g_x^{\mu\nu}
       (D_\mu^{\epsilon_\mu}\phi)_x
       (D_\nu^{\epsilon_\nu}\phi)_x\\
 &+\frac{\mathfrak v}{2}\sum_x\sqrt{\det g_x}\,m^2\phi_x^2.
 \end{split}
 \label{eq:v7-Hodge-writer}
\end{equation}
Repeated spacetime indices are summed; the real Higgs-component sum is
suppressed. Each of the sixteen terms is a positive metric quadratic form
and has fixed finite support. This is an orientation-averaged realization of
the positive edge Hodge form. At nonzero internal links its incoming and
outgoing differences can first be transported to the same site, preserving
internal gauge covariance of this \emph{seed}. The nonlinear spacetime repair
below is proved at $U=I$ only; its internal-gauge covariant extension is not
assumed.

At $g=I$ the diagonal sum gives exactly the usual nearest-neighbour scalar
kinetic form. With physical momenta in the fundamental cube put
\begin{equation}
 \kappa_\mu(p)=\frac2a\sin\frac{ap_\mu}{2},\qquad
 P_a(p)=m^2+\sum_\mu\kappa_\mu(p)^2,\qquad
 v_\mu(p)=\frac12\partial_{p_\mu}P_a(p)=\frac{\sin(ap_\mu)}a.
 \label{eq:v7-native-symbols}
\end{equation}
Thus $P_a$ is the scalar version of the V6 free denominator, and $i v_\mu$
is the genuinely local centred difference. The phase differential used in
the free metric gauge transformation still has symbol $i\kappa_\mu$.

Write $g=I+\varepsilon h$. The coefficient of $\varepsilon$ is
$S_{3,a}^{b}=\frac14\langle\phi,\mathcal V_h^b\phi\rangle$; the superscript
$b$ denotes the unrepaired Hodge seed, not a renormalized coefficient.
For scalar momenta $p,q$, total momentum $k=p+q$, and metric momentum $-k$,
its tensor is
\begin{align}
 B_{\mu\nu}^{b}(p,q)&=
 \begin{cases}
  \kappa_\mu(p)\kappa_\mu(q)\cos(ak_\mu/2),&\mu=\nu,\\
  v_\mu(p)v_\nu(q),&\mu\ne\nu,
 \end{cases}\notag\\
 V_{\mu\nu}^{b}(p,q)&=B_{\mu\nu}^{b}(p,q)+B_{\nu\mu}^{b}(p,q)
       -\delta_{\mu\nu}\left(\sum_\rho B_{\rho\rho}^{b}(p,q)-m^2\right).
 \label{eq:v7-seed-vertex}
\end{align}
Here and below a bilinear momentum tensor is paired with the symmetric metric
source; in position space all products retain the physical cell weights.

\begin{lemma}[The seed vertex is an actual finite geometric variation]
\label[lemma]{lem:v7-seed}
Equation~\eqref{eq:v7-seed-vertex} is the first metric variation of
\eqref{eq:v7-Hodge-writer}. It is symmetric in its metric indices and in its
two scalar arguments, is real under simultaneous momentum reversal, and
has the limit
\begin{equation}
 V_{\mu\nu}^{(0)}(p,q)=p_\mu q_\nu+p_\nu q_\mu
                     -\delta_{\mu\nu}(p\cdot q-m^2).
 \label{eq:v7-continuum-vertex}
\end{equation}
On $|p|,|q|,m\le C\Lambda$, $a\Lambda\le c$, its difference from
\eqref{eq:v7-continuum-vertex}, with any fixed number of scaled momentum
derivatives, is $O(a^2\Lambda^4)$.
\end{lemma}
\begin{proof}
Use $\sqrt{\det(I+\varepsilon h)}=1+\varepsilon\operatorname{tr}h/2+O(\varepsilon^2)$
and $g^{-1}=I-\varepsilon h+O(\varepsilon^2)$. The coefficient at one sign
choice is
\[
 -\tfrac12h_{\mu\nu}D_\mu^{\epsilon_\mu}\phi
                         D_\nu^{\epsilon_\nu}\phi
 +\tfrac14\operatorname{tr}h
       \left(\sum_\mu(D_\mu^{\epsilon_\mu}\phi)^2+m^2\phi^2\right).
\]
The symbols of $D_\mu^\pm$ are $i\kappa_\mu(p)e^{\pm iap_\mu/2}$.
For equal indices the same sign occurs twice, giving
$\cos(a(p_\mu+q_\mu)/2)$. For different indices the independent sign average
factorizes into $\cos(ap_\mu/2)\cos(aq_\nu/2)$, giving $v_\mu(p)v_\nu(q)$.
The two factors $i$ account for the signs in \eqref{eq:v7-seed-vertex}.
This proves the coefficient and its symmetries. Taylor expansion of sine and
cosine, on a fixed scaled complex neighbourhood if derivatives are desired,
proves the last assertion.
\end{proof}

\section{Why an action-only repair with the naive matter ghost fails}
\label{sec:v7-obstruction}

Let the free metric transformation be
$\delta_0 h_{\mu\nu}=\delta_\mu c_\nu+\delta_\nu c_\mu$, where the phase
symbols are those of \eqref{eq:v7-native-symbols}. Take a candidate first
matter transformation with momentum velocity $d_\nu(p)$, continuously differentiable and odd in $p$:
$\delta_1\phi=c^\nu d_\nu(\delta)\phi$, with differential symbol $i d_\nu$.
Before applying source windows, the $c\phi\phi$ Noether equation is exactly
\begin{equation}
 \boxed{\sum_\mu\kappa_\mu(k)V_{\mu\nu}(p,q)
       =d_\nu(p)P_a(q)+d_\nu(q)P_a(p),\qquad k=p+q.}
 \label{eq:v7-Ward}
\end{equation}
Indeed $\delta_0S_3$ has coefficient $-i\kappa_\mu(k)V_{\mu\nu}/2$,
whereas $\delta_1S_2$ has coefficient
$i[d_\nu(p)P_a(q)+d_\nu(q)P_a(p)]/2$.
This calculation uses finite differentiation of the action, not a continuum
Leibniz identity.

\begin{proposition}[A zero-transfer integrability obstruction]
\label[proposition]{prop:v7-integrability}
Fix the native $P_a$ and free phase metric transformation. If a symmetric
vertex $V_{\mu\nu}$, continuously differentiable in the total momentum near
$k=0$, solves \eqref{eq:v7-Ward} on an open momentum chart, then for every
relative momentum $t$ the matrix
\begin{equation}
 M_{\mu\nu}^{d}(t)
    =P_a(t)\partial_\mu d_\nu(t)
                      -d_\nu(t)\partial_\mu P_a(t)
 \label{eq:v7-Md}
\end{equation}
must be symmetric. For the naive choice $d=\kappa$ this condition fails at
generic $t$ with two unequal nonzero momentum components. In particular,
\begin{equation}
 M_{\mu\nu}^{\kappa}-M_{\nu\mu}^{\kappa}
  =2\kappa_\mu(t)\kappa_\nu(t)
       \left\{\cos(at_\nu/2)-\cos(at_\mu/2)\right\}.
 \label{eq:v7-skew-obstruction}
\end{equation}
Consequently changing only the regular symmetric metric--Higgs vertex cannot
solve this Ward equation while keeping that naive matter transformation.
\end{proposition}
\begin{proof}
Set $p=t+k/2$ and $q=-t+k/2$. At $k=0$ the right side vanishes because
$P_a$ is even and $d$ odd. Differentiate in $k_\mu$ at zero. Since
$D\kappa(0)=I$, the left derivative is $V_{\mu\nu}(t,-t)$; the derivative
of the right side is \eqref{eq:v7-Md}. Symmetry of the vertex is therefore
necessary. Now $\partial_\mu\kappa_\nu=\delta_{\mu\nu}\cos(at_\mu/2)$
and $\partial_\mu P_a=2\kappa_\mu\cos(at_\mu/2)$. Substitution gives
\eqref{eq:v7-skew-obstruction}, which is nonzero, for example, for distinct
small positive $at_\mu,at_\nu$. This remains an obstruction however many
regular symmetric action counterterms are added while $d$ is held fixed.
\end{proof}

This is a statement about an open symbol chart, and hence about a regular
local thermodynamic symbol family. Arbitrary interpolation on the finitely
many momenta of one fixed torus is not excluded; it need not have controlled
derivatives as the torus grows. The proposition also does not rule out other
matter transformations containing derivatives of the ghost. It rules out
the specified naive ansatz, and supplies a reason to correct the ghost
partner as well as the action. The general lattice Leibniz obstructions of
\cite{V7Kato} are related context, not the proof of this tensor condition.

A useful nonarbitrary correction is the native dispersion velocity
$d=v=\frac12\nabla P_a$. For this choice
\begin{equation}
 M_{\mu\nu}^{v}(t)
   =P_a(t)\delta_{\mu\nu}\cos(at_\mu)-2v_\mu(t)v_\nu(t),
 \label{eq:v7-Mv}
\end{equation}
which is symmetric. No inverse of $P_a(t)$ is used, so this assertion also
holds at $m=t=0$. Relative to the naive phase velocity,
\begin{equation}
 v_\mu(p)-\kappa_\mu(p)
                  =-\frac{a^2}{8}p_\mu^3+O(a^4|p_\mu|^5).
 \label{eq:v7-ghost-correction}
\end{equation}
This corrects a transformation coefficient at the same EFT order as the
vertex correction below. It is not a supplied continuum beta function or a
change of the free inverse propagator.

\section{An explicit symmetric vertex with no soft-momentum inverse}
\label{sec:v7-closed}

For $k=p+q$, $t=(p-q)/2$, introduce
\begin{equation}
 \begin{aligned}
 r_\mu&=\kappa_\mu(k),& w_\mu&=\frac{\sin(at_\mu)}a,\\
 c_\mu&=\cos(at_\mu),& C_\mu&=\cos(ak_\mu/2),\\
 W&=\sum_\rho w_\rho r_\rho,&
 \ell_\mu&=-\frac{a^2}{4(1+C_\mu)},\\
 A&=\frac{P_a(p)+P_a(q)}2
       =m^2+\frac2{a^2}\sum_\rho(1-c_\rho C_\rho).
 \end{aligned}
 \label{eq:v7-closed-data}
\end{equation}
Here $W$ is a scalar, not a generating functional. On $|ak_\mu|<\pi$ the
denominator in $\ell_\mu$ is bounded away from zero. The identities
\begin{equation}
 C_\mu=1+r_\mu^2\ell_\mu,\qquad
 P_a(p)=A+W,\quad P_a(q)=A-W
 \label{eq:v7-basic-identities}
\end{equation}
follow from elementary trigonometry. The expression for $\ell_\mu$ is its
analytic value also at $r_\mu=0$, not a division by $r_\mu^2$.

\begin{theorem}[Pole-free cubic metric--Higgs completion]
\label[theorem]{thm:v7-completion}
On the stated chart define
\begin{equation}
 \boxed{\begin{aligned}
 V^\sharp_{\mu\nu}(p,q)={}&
 \delta_{\mu\nu}\{c_\mu A-2w_\mu r_\mu\ell_\mu W\}
       -2w_\mu w_\nu\\
 &-\frac12\{ |r|^2\delta_{\mu\nu}-r_\mu r_\nu\}.
 \end{aligned}}
 \label{eq:v7-closed-vertex}
\end{equation}
There is no sum on $\mu$ inside the diagonal braces. This is a symmetric,
real, scalar-exchange-symmetric analytic tensor, and it obeys
\begin{equation}
 \boxed{r_\mu V^\sharp_{\mu\nu}(p,q)
        =v_\nu(p)P_a(q)+v_\nu(q)P_a(p).}
 \label{eq:v7-exact-repaired-Ward}
\end{equation}
Its limit at $a=0$ is exactly \eqref{eq:v7-continuum-vertex}.
The correction $\Delta V=V^\sharp-V^b$ is $O(a^2\Lambda^4)$ on every fixed
scaled domain $|p|,|q|,m\le C\Lambda$, $a\Lambda\le c_0$, with the same
bound after any fixed number of scaled momentum derivatives.
\end{theorem}
\begin{proof}
The two scalar velocities are
\[
 v_\nu(p)=w_\nu C_\nu+c_\nu r_\nu/2,\qquad
 v_\nu(q)=-w_\nu C_\nu+c_\nu r_\nu/2.
\]
Using \eqref{eq:v7-basic-identities}, the right side of
\eqref{eq:v7-exact-repaired-Ward} is
$r_\nu c_\nu A-2w_\nu C_\nu W$. Contract
\eqref{eq:v7-closed-vertex} with $r_\mu$. The last tensor is transverse,
and the other terms give
\[
 r_\nu c_\nu A-2w_\nu W-2w_\nu r_\nu^2\ell_\nu W
              =r_\nu c_\nu A-2w_\nu C_\nu W.
\]
This proves the exact identity without dividing by a momentum norm.
Metric symmetry is manifest. Scalar exchange sends $t$ to $-t$, reversing
both $w$ and $W$ but leaving the displayed tensor unchanged. Simultaneous
momentum reversal reverses $t,k,w,r$ and preserves it as well.

At $a=0$, $A=m^2+|t|^2+|k|^2/4$, $w=t$, $r=k$, and
$\ell_\mu=0$. Thus the limit is
\[
 \delta_{\mu\nu}(m^2+|t|^2-|k|^2/4)
             -2t_\mu t_\nu+k_\mu k_\nu/2,
\]
which is \eqref{eq:v7-continuum-vertex}. The transverse term in
\eqref{eq:v7-closed-vertex} is important: omitting it would add a finite
continuum curvature improvement, not an innocuous regulator error.
Both seed and repaired expressions have even analytic expansions in $a$ and
agree at $a=0$. Rescale momenta and $m$ by $\Lambda$; sine divided by its
argument and $1/(1+\cos)$ are analytic on a common neighbourhood when
$a\Lambda$ is small. Their difference and every fixed scaled derivative
are bounded by $K(a\Lambda)^2\Lambda^2$, proving the final assertion.
\end{proof}

The Ward equation does not determine all transverse physics. For every
symmetric tensor $T_{\mu\nu}$ with $r_\mu T_{\mu\nu}=0$, the vertex
$V^\sharp+T$ obeys the same identity. In particular an independent coefficient
multiplying $(|r|^2\delta_{\mu\nu}-r_\mu r_\nu)$ retains the $RH^\dagger H$
improvement channel. Its physical normalization remains an EFT datum.
The formula chooses the minimal continuum vertex when that independent
coefficient is zero; it does not determine it from symmetry.

\subsection{A finite source window makes the identity a literal regulator identity}

Fix a real even $\zeta\in C_c^\infty(\mathbb R^4)$, equal to one on a
neighbourhood of zero and supported in $|z|_\infty<2$, and set
\begin{equation}
 \Theta_\Lambda(k,p,q)=\zeta(k/\Lambda)\zeta(p/\Lambda)\zeta(q/\Lambda),
             \qquad a\Lambda\le1/32.
 \label{eq:v7-window}
\end{equation}
Use this same factor in the action and matter-ghost vertices. Nonzero
coefficients then have $k=p+q$ without Brillouin-zone wrapping. Multiplying
\eqref{eq:v7-exact-repaired-Ward} by $\Theta_\Lambda$ gives an exact identity
at every allowed finite-torus momentum; outside its support both sides vanish.
In particular the construction does not assume that a projected pointwise
product is associative. It has specified one trilinear coefficient with an
exact Ward identity. The window is part of the regulator, not a discarded
exceptional event.

\section{The local EFT primitive and its scale-uniform kernel bound}
\label{sec:v7-locality}

The exact repaired tensor is analytic near zero transfer and quasilocal after
\eqref{eq:v7-window}. It need not have finite lattice support. At each fixed
EFT order its expansion instead gives finitely many local differential
counterterm types, plus a controlled remainder.

Here is an explicit first coefficient, useful for independent verification.
Use $k,t$ as in \eqref{eq:v7-closed-data} and write
\[
 A_0=m^2+|t|^2+|k|^2/4,\qquad
 A_2=-\sum_\rho\left(t_\rho^4/12+t_\rho^2k_\rho^2/8+k_\rho^4/192\right).
\]
Then $V^\sharp=V^{(0)}+a^2V^{\sharp,(2)}+O(a^4\Lambda^6)$, where
\begin{equation}
 \begin{split}
 V^{\sharp,(2)}_{\mu\nu}={}&\delta_{\mu\nu}
 \left\{A_2-\frac{t_\mu^2A_0}{2}
       +\frac{t_\mu k_\mu(t\cdot k)}4
       +\frac1{24}\sum_\rho k_\rho^4\right\}\\
 &+\frac{t_\mu^3t_\nu+t_\mu t_\nu^3}{3}
              -\frac{k_\mu^3k_\nu+k_\mu k_\nu^3}{48}.
 \end{split}
 \label{eq:v7-four-derivative}
\end{equation}
The diagonal braces again have no additional sum on $\mu$.
For the seed define
\begin{equation}
 B_{\mu\nu}^{b,(2)}=
 \begin{cases}
 -\dfrac{p_\mu^3q_\mu+p_\mu q_\mu^3}{24}
             -\dfrac{p_\mu q_\mu k_\mu^2}{8},&\mu=\nu,\\[2mm]
 -\dfrac{p_\mu^3q_\nu+p_\mu q_\nu^3}{6},&\mu\ne\nu.
 \end{cases}
 \label{eq:v7-seed-four-derivative}
\end{equation}
Put $V^{b,(2)}=B^{b,(2)}+(B^{b,(2)})^{\mathsf T}
-\Id\sum_\rho B_{\rho\rho}^{b,(2)}$.
The actual first action counterterm is therefore
$a^2(V^{\sharp,(2)}-V^{b,(2)})$, not the entire tensor
\eqref{eq:v7-four-derivative}. Along with
\eqref{eq:v7-ghost-correction}, this specifies the first restoring
\emph{action and antifield} coefficient in elementary polynomials.

\begin{proposition}[Finite EFT expansion of the actual correction]
\label[proposition]{prop:v7-EFT}
For every fixed integer $R\ge1$ there are polynomials
$\Delta V^{(2)},\ldots,\Delta V^{(2R)}$, of joint momentum/mass weight
$2+2s$ at order $a^{2s}$, such that
\begin{equation}
 \left|\partial_p^\alpha\partial_q^\beta
  \left(\Delta V-\sum_{s=1}^{R}a^{2s}\Delta V^{(2s)}\right)\right|
 \le K_{R,\alpha,\beta}\,
     a^{2R+2}\Lambda^{2R+4-|\alpha|-|\beta|}.
 \label{eq:v7-EFT-error}
\end{equation}
For the velocity correction the coefficient weight is $1+2s$, and its
differentiated tail is bounded by
\[
 K a^{2R+2}\Lambda^{2R+3-|\alpha|}.
\] Retaining both expansions
through the same order leaves a Ward residual at most
$K a^{2R+2}\Lambda^{2R+5}$, with corresponding differentiated bounds.
\end{proposition}
\begin{proof}
All expressions in \eqref{eq:v7-closed-data}--\eqref{eq:v7-closed-vertex}
have convergent even series in $a$, with removable apparent negative powers.
Their coefficients at order $a^{2s}$ are polynomials of the stated weight;
this follows either by multiplying the sine/cosine series or by simultaneous
scaling of $p,q,m$. The same is true of \eqref{eq:v7-seed-vertex}.
On the scaled complex neighbourhood from the proof of
\cref{thm:v7-completion}, Taylor's integral estimate and Cauchy estimates
bound every fixed differentiated tail, giving \eqref{eq:v7-EFT-error}.
The velocity has one rather than two leading momentum powers. Multiply the
vertex tail by $r=O(\Lambda)$ and the velocity tail by $P_a=O(\Lambda^2)$
in the exact Ward identity; the stated residual follows. Formula
\eqref{eq:v7-four-derivative} is obtained by using
$\ell_\mu=-a^2/8+O(a^4k_\mu^2)$ and expanding the other factors through
$a^2$, retaining the transverse term before simplifying.
\end{proof}

These are differential coefficients in the local EFT expansion. The exact
regulator representatives use the smooth source window. A polynomial in the
continuous momenta is not here asserted to be a finite-range operator on all
Brillouin-zone modes; the windowed comparison is the one controlled below.
No invariant curvature classification is inferred for an individual
symmetry-restoring lattice polynomial. Covariant EFT operators and restoring
antifield/source terms are separate coordinates in the V6 bank.

\begin{theorem}[No inverse-momentum loss in the rooted cubic correction]
\label[theorem]{thm:v7-kernel}
Let $K_{\Lambda,a}(y,z)$ be the inverse Fourier kernel in the two relative
variables of $\Theta_\Lambda\Delta V$, with convolution weights
$\mathfrak v^2=a^8$. For every fixed $s\ge0$, if
$a\Lambda\le1/32$ and $0\le m\le M\Lambda$, then
\begin{equation}
 \boxed{a^8\sum_{y,z}
  \{1+\Lambda(d_{\rm per}(y,0)+d_{\rm per}(z,0))\}^{s}
             \|K_{\Lambda,a}(y,z)\|
       \le K_{s,M}a^2\Lambda^4.}
 \label{eq:v7-kernel-bound}
\end{equation}
The constant is independent of mesh, all four periods, and $\Lambda$ in
this domain. The analogous matter-antifield correction kernel is bounded by
$K_{s,M}a^2\Lambda^3$. For the tails in \eqref{eq:v7-EFT-error}, replace
the right sides by $K_{R,s,M}a^{2R+2}\Lambda^{2R+4}$ and
$K_{R,s,M}a^{2R+2}\Lambda^{2R+3}$ respectively.
\end{theorem}
\begin{proof}
Differentiating the cutoff and the explicit regular tensor gives, for every
multi-index $\gamma$ in the eight momentum variables,
\[
 \|\partial^\gamma(\Theta_\Lambda\Delta V)\|_\infty
                         \le K_\gamma a^2\Lambda^{4-|\gamma|}.
\]
The support has eight-dimensional volume $O(\Lambda^8)$ and lies strictly
inside the momentum cube. Integration by parts an arbitrary fixed number of
times therefore bounds the infinite-mesh kernel by
\[
 K_n a^2\Lambda^{12}
             (1+\Lambda(|y|+|z|))^{-n}.
\]
There are no seam boundary terms. Choose $n>s+8$ and sum with weights
$a^8$; the sum is at most $K\Lambda^{-8}$ because $a\Lambda\le1/32$.
This proves \eqref{eq:v7-kernel-bound} on the infinite mesh. Periodization
is absolutely convergent, and the shortest periodic distance is no larger
than the distance of any representative. Its weighted sum is consequently
bounded by the same infinite sum, independently of the periods. Apply the
same argument to the velocity difference and the Taylor tails. All tensor
and Higgs component counts are fixed.
\end{proof}

After the two-derivative action normalization, the relative correction stock
is $O((a\Lambda)^2)$; the antifield correction has the same relative stock
after one-derivative normalization. The bound includes the zero-total-momentum
region. A construction using $k_\mu/|k|^2$ to solve the Ward identity would
not give this conclusion. Higher fixed local source derivatives insert only
the finite coefficient and product-rule factors of these same kernels;
no all-source analytic radius or arbitrary-order uniformity is asserted.

\section{The first mixed BV equation and the measured Higgs-shell output}
\label{sec:v7-BV-shell}

\subsection{A finite cubic cocycle, not an assumed nonlinear gauge algebra}

Reinstate the unreduced metric components as fields/sources in the free
polynomial complex. The commuting phase differences define the usual
linear Fierz--Pauli Einstein operator
\begin{equation}
 \begin{split}
 (\mathcal G_a h)_{\mu\nu}=\tfrac12\{&
  \delta_\rho\delta_\mu h_{\rho\nu}
 +\delta_\rho\delta_\nu h_{\rho\mu}-\delta^2h_{\mu\nu}
 -\delta_\mu\delta_\nu\operatorname{tr}h\\
 &-\delta_{\mu\nu}(\delta_\rho\delta_\sigma h_{\rho\sigma}
                       -\delta^2\operatorname{tr}h)\}.
 \end{split}
 \label{eq:v7-linear-G}
\end{equation}
Commutativity of the finite differences proves $\mathcal G_a R_a=0$ for
$(R_ac)_{\mu\nu}=\delta_\mu c_\nu+\delta_\nu c_\mu$.
This is the free quadratic extension of the flat V6 constraint reference;
its TT action has the same normalization after $h_{\rm TT}=\sqrt2 q$.
No nonlinear identity follows from this extension, and the full quadratic
metric form is not being treated as a positive measure.

Let $\mathcal S_0$ be the sum of the free metric action
$\frac12\langle h,\mathcal G_a h\rangle$, the scalar action with inverse
$P_a$, and the canonical metric-antifield term $\langle h^*,R_ac\rangle$.
Use the canonical BV bracket with the convention that its first-order master
equation is the Noether sum $\delta_0S_3+\delta_1S_2=0$.
Equivalently, this is the standard right/left derivative bracket with the
antifield placed to the left of its odd transformation. Add nonminimal
contractible pairs only when a gauge-fixing calculation needs them.

For the cubic source window, define
\begin{equation}
 \mathcal S_1^\sharp=S_{3,a}^\sharp+
                    \sum_A\langle\phi_A^*,\mathcal R_v(c)\phi_A\rangle.
 \label{eq:v7-S1}
\end{equation}
The kernel of $\mathcal R_v(c)$ has symbol $i v_\nu(q)c_\nu(-p-q)$ times
$\Theta_\Lambda(p+q,p,q)$ when paired with an antifield of momentum $p$.
The seed $\mathcal S_1^b$ uses $V^b$ and $i\kappa_\nu(q)$ instead.
Antifields are independent coordinates, not suppressed source derivatives.

\begin{theorem}[Explicit first mixed classical and quantum master coefficient]
\label[theorem]{thm:v7-BV}
On every finite periodic regulator with \eqref{eq:v7-window},
\begin{equation}
 (\mathcal S_0,\mathcal S_1^\sharp)=0,\qquad
                 \Delta_{\rm BV}\mathcal S_1^\sharp=0.
 \label{eq:v7-master-first}
\end{equation}
Here $\Delta_{\rm BV}$ uses the specified flat scalar/antifield density.
Consequently $\mathcal S_0+\varepsilon\mathcal S_1^\sharp$ satisfies the
quantum master equation modulo $O(\varepsilon^2)$ in this free-metric/Higgs
component. More explicitly, its seed defect has the actual restoring primitive
\begin{equation}
 \boxed{\begin{aligned}
 (\mathcal S_0,\mathcal S_1^b)
       & =-(\mathcal S_0,\Delta\mathcal S_1),\\
 \Delta\mathcal S_1&=
  \tfrac14\langle\phi,\Delta\mathcal V_h\phi\rangle+
  \langle\phi^*,(\mathcal R_v-\mathcal R_\kappa)(c)\phi\rangle.
 \end{aligned}}
 \label{eq:v7-primitive}
\end{equation}
Both pieces of this primitive have the quasilocal bounds of
\cref{thm:v7-kernel} and the finite EFT expansions above.
\end{theorem}
\begin{proof}
The only nonzero contributions to the first mixed master coefficient are
the metric variation of $S_{3,a}^\sharp$ and the free scalar Euler covector
paired with $\mathcal R_v(c)\phi$. Their coefficients are the two opposite
sides of \eqref{eq:v7-exact-repaired-Ward}, with the common cutoff factor.
This proves the first identity. The free bracket squares to zero because
$\mathcal G_aR_a=0$ and the linear gauge transformations commute; the
scalar fields have zero free gauge transformation.

For the BV Laplacian, the action part has no antifields. The other part gives
$n_H\operatorname{Tr}\mathcal R_v(c)$. The diagonal Fourier contraction forces
the ghost momentum to zero. Its coefficient is proportional to
\[
  c_\nu(0)\sum_p\zeta(p/\Lambda)^2\,i v_\nu(p).
\]
The momenta occur in opposite pairs, the cutoff is even, and $v$ is odd.
Every pair cancels and the zero mode contributes zero. This is a finite
measure-divergence computation, not an appeal to anomaly cancellation in
the chiral SM. The same argument gives $\Delta_{\rm BV}\mathcal S_0=0$.
Expanding the quantum master equation in $\varepsilon$ proves the stated
order. Subtract the first identity from the seed expression to obtain
\eqref{eq:v7-primitive}. The kernel and EFT estimates have already been
proved for its two explicit differences.
\end{proof}

The theorem is restricted to the displayed mixed component, not the full
ESM master equation. In particular it contains no pure $hhh$ deformation,
no mixed order $\varepsilon g_{\rm SM}$, no Yukawa or spin-current loop,
and no local Lorentz/connection determinant. Flat Berezin-divergence zero
is not determinant-line coherence. The general first-order deformation
interpretation is established BV methodology \cite{V7BarnichHenneaux}; the
new content here is its finite native-dispersion vertex and the bounded
restoring primitive \eqref{eq:v7-primitive}.

\subsection{The corrected vertex is integrated, not dropped after repair}

Take an actual positive scalar shell $C_H$, for example the smooth Fourier
multiplier
\begin{equation}
 C_H(p)=\frac{\chi(\lambda_a(p)/\Lambda^2)}{P_a(p)},\qquad
       \lambda_a(p)=\sum_\mu\kappa_\mu(p)^2,\quad
       \chi\in C_c^\infty((1,4)),\quad 0\le\chi\le1.
 \label{eq:v7-Higgs-covariance}
\end{equation}
It is bounded above by the full scalar covariance on its support and excludes
the massless zero mode. It therefore defines a normalized Gaussian innovation
in exactly the same measured sense as \cref{prop:v6-free-split}.
Retain the metric source, matter antifields and ghosts, a scalar field $\phi$
and a source $J$ for the eliminated Higgs variable $\eta$.

\begin{proposition}[Source-complete first mixed shell coefficient]
\label[proposition]{prop:v7-shell-output}
For $\eta\sim N(0,\hbar C_H)$, the coefficient of $\varepsilon$ in
\begin{equation}
 -\hbar\log\E\exp\{\hbar^{-1}\langle J,\eta\rangle
       -\varepsilon\hbar^{-1}
            \mathcal S_1^\sharp(h,\phi+\eta,\phi^*,c)\}
 \label{eq:v7-measured-integral}
\end{equation}
is exactly
\begin{equation}
 \boxed{\mathcal S_1^\sharp(h,\phi+C_HJ,\phi^*,c)
           +\frac{\hbar}{4}\operatorname{Tr}(\mathcal V_h^\sharp C_H).}
 \label{eq:v7-shell-output}
\end{equation}
The degree-zero term is $-\langle J,C_HJ\rangle/2$.
In particular the output retains the metric--source-only contact
$\langle C_HJ,\mathcal V_h^\sharp C_HJ\rangle/4$, its mixed retained-field
partner, the antifield--ghost--$J$ contact, and the one-Higgs-loop metric
one-point coefficient. Replacing $\sharp$ by $b$ gives the actual seed output;
their difference is obtained by substituting the explicit correction, without
an unrecorded normalization change.
\end{proposition}
\begin{proof}
Differentiate the finite logarithm at $\varepsilon=0$. Completing the
Gaussian source square shifts $\eta$ by $C_HJ$ and contributes
$-\langle J,C_HJ\rangle/2$ to the negative logarithm. The quadratic Higgs
part has expectation its shifted value plus
$\hbar\operatorname{Tr}(\mathcal V_h^\sharp C_H)/4$. The antifield term is
linear in the Higgs field and only shifts. This proves
\eqref{eq:v7-shell-output}, also for the even Grassmann coefficient algebra.
The formula is an exact finite coefficient; a finite real neighbourhood of
$\varepsilon=0$ also exists whenever its Gaussian quadratic form remains
positive. No uniform higher-order interacting remainder is inferred.
\end{proof}

\subsection{The induced stress one-point coefficient and its cutoff price}

With normalized momentum sum
$\int_p^\ell=\operatorname{Vol}(\mathbb T_\ell)^{-1}\sum_p$, the one-point
coefficient per physical volume is
\begin{equation}
 t_{\mu\nu}^{\sharp}(a,\Lambda)
 =\frac{\hbar n_H}{4}\int_p^\ell
   \Theta_\Lambda(0,p,-p)V_{\mu\nu}^{\sharp}(p,-p)C_H(p).
 \label{eq:v7-tadpole}
\end{equation}
Zero total momentum gives an especially simple evaluated tensor:
\begin{equation}
 \boxed{V_{\mu\nu}^{\sharp}(p,-p)
       =P_a(p)\delta_{\mu\nu}\cos(ap_\mu)-2v_\mu(p)v_\nu(p).}
 \label{eq:v7-zero-transfer-vertex}
\end{equation}
For the specified Hodge seed, its off-diagonal entries at $(p,-p)$ are already
$-2v_\mu v_\nu$, while its diagonal entries are $P_a-2\kappa_\mu^2$.
Thus the actual difference is diagonal, with
\begin{equation}
 \Delta V_{\mu\mu}(p,-p)
      =P_a(p)\{\cos(ap_\mu)-1\}
                      +2\{\kappa_\mu(p)^2-v_\mu(p)^2\}.
 \label{eq:v7-tadpole-correction}
\end{equation}
The one-point term is a geometric source term, not a source-independent
vacuum constant. On an equal-period hypercubic family with a symmetric cutoff
its tensor is proportional to $\delta_{\mu\nu}$ by sign reversals and
permutations. That special symmetry does not license deleting its variation
with respect to the metric.

Define the explicit short-period factor
\[
 \mathfrak t_\Lambda=\prod_{\mu=0}^{3}\{1+(\Lambda L_\mu)^{-1}\}.
\]
Then
\begin{equation}
 |t^\sharp-t^b|\le K_M\hbar n_H a^2\Lambda^6\mathfrak t_\Lambda.
 \label{eq:v7-tadpole-bound}
\end{equation}
Indeed the shell contains at most
$K\operatorname{Vol}(\mathbb T_\ell)\Lambda^4\mathfrak t_\Lambda$ momenta,
$C_H\le\Lambda^{-2}$ there, and
$\|\Delta V\|\le K_Ma^2\Lambda^4$. The short-period factor is retained
because a diagonal trace estimate is stronger than an $\ell^1$ kernel row
estimate when one physical period is very small.

For a fixed infrared scale $\mu>0$ and periods with $\mu L_\nu\ge1$, put
$\Lambda_j=\mu b^j$, $b>1$, $0\le j\le J$. Summing the \emph{same explicit correction}
through $\Lambda_J$ gives
\begin{align}
 \sum_{j=0}^{J}\Lambda_j^{-2}\|\Delta K_j\|_{\rm rooted}
                      &\le K_b(a\Lambda_J)^2,\\
 \sum_{j=0}^{J}|t_j^\sharp-t_j^b|
                      &\le K_{b,M}\hbar n_Ha^2\Lambda_J^6.
 \label{eq:v7-summed-correction}
\end{align}
The first uses \eqref{eq:v7-kernel-bound}; the second uses
\eqref{eq:v7-tadpole-bound}. Both are geometric sums. For example the
separated cutoff choice
$\Lambda_J\le\mu(a\mu)^{-\nu}$ with $0<\nu<1/3$ makes even the
unsubtracted one-point correction tend to zero. This proves a cofinal
comparison for these coefficients only. It does not estimate all ESM diagrams
or replace local subtraction of their divergent parts. At a top scale of
order $a^{-1}$ the displayed bounds do not vanish.

\section{The next nonlinear equation and the scope of the advance}
\label{sec:v7-next}

The cubic construction has not secretly restored a full lattice Lie algebra.
Already in one direction, with $D=(T-T^{-1})/(2a)$ and
$L_\xi=M_\xi D$, define
\[
 A_+=\xi(T\eta)-\eta(T\xi),\qquad
 A_-=\xi(T^{-1}\eta)-\eta(T^{-1}\xi).
\]
Direct multiplication gives the exact identity
\begin{equation}
 [L_\xi,L_\eta]
  =\frac1{4a^2}\{M_{A_+}(T^2-I)+M_{A_-}(T^{-2}-I)\}.
 \label{eq:v7-next-algebra}
\end{equation}
For generic profiles on a periodic lattice of at least seven sites, the
$T^2$ and $T^{-2}$ entries are nonzero. Thus this commutator is not another
operator $M_\beta D$, whose off-diagonal support is at distance one. This
remains a check on the proposed nonlinear ansatz, not an obstruction to
field-dependent higher transformations. Smooth windows add their own
commutator contributions and must also be transported.

In BV form the next coefficient produced by the completed mixed cocycle is
\begin{equation}
 \mathcal O_2=\tfrac12(\mathcal S_1^\sharp,\mathcal S_1^\sharp),
 \qquad (\mathcal S_0,\mathcal O_2)=0.
 \label{eq:v7-next-cocycle}
\end{equation}
The second equality follows from the graded Jacobi identity and
\eqref{eq:v7-master-first}. Its cancellation requires an actual quartic
coefficient, higher transformations, and, for full ESM, the native pure
metric cubic and the internal-gauge and spin sectors. In particular one must
solve the combined equation
$(\mathcal S_0,\mathcal S_2)=-\mathcal O_2$ after all relevant cubic sectors
have been included, not set the right side to zero. The local cohomological
and uniform-norm issues emphasized in \cite{V7BarnichHenneaux} are distinct;
our explicit first primitive proves neither for this next equation.

This sets a concrete next interface. The scalar part of the first mixed
cocycle, its contact terms and its kernel bounds are now specified. It must
be coupled to the actual gravitational cubic vertex and its measure before
a complete graviton/coframe-plus-ghost shell is claimed. The $\varepsilon^2$
metric--metric--Higgs seagull and its antifield terms are the next mixed
coefficients; no curvature-squared loop coefficient can be inferred by
closing only the one-metric tadpole. The gauge collar and chiral-line tasks
from earlier versions remain where the full ESM map consumes them.

\subsection{Uniformity, nonduplication, and verification}

The general R1--R4, filtered trajectory, Slavnov-splitting, Gaussian
pushforward and source-contact principles are inherited. The new calculation
is the Hodge-seed vertex, its integrability obstruction, the explicit
pole-free repair, its fixed-order local primitive, and its measured scalar
shell coefficient. It does not relabel a standard continuum graviton vertex
as a native finite calculation. The actual $P_a$, phase gauge generator,
local sign-averaged Hodge writer, source window and flat measure have all been
specified.

\begin{center}
\small
\begin{tabular}{>{\raggedright\arraybackslash}p{.22\textwidth}>{\raggedright\arraybackslash}p{.69\textwidth}}
\toprule
Variable & Precise scope of V7\\
\midrule
Mesh and shell & Exact Ward repair on the stated chart; kernel estimates for
$a\Lambda\le1/32$. The full Brillouin-zone interacting theory is not repaired.\\
Volume and time & Polynomial kernel moments are uniform in all periods.
The one-point trace bound retains $\mathfrak t_\Lambda$ explicitly.\\
EFT order & Every fixed $R$ has finitely many restoring differential types.
Constants can grow with $R$; the exact analytic symbol is not a finite-support
all-orders lattice action.\\
Coupling and field degree & First gravitational coupling in the
$h\phi\phi$ and matter-antifield sector. Pure gravitational cubic, mixed gauge
orders and higher field degrees are not included in the solved equation.\\
Sources & All contacts in the displayed finite Gaussian coefficient are
retained, with fixed finite source derivatives. No complete diffeomorphism-
invariant relational observable algebra is constructed.\\
Measure & The BV divergence is evaluated for the stated flat scalar measure.
Nonlinear coframe, Lorentz, connection and chiral-line factors remain separate.\\
Shell number & The explicit correction sums have the bounds
\eqref{eq:v7-summed-correction}; they are not an invariant interacting ESM
RG class or a bound for its full many-shell source transport.\\
\bottomrule
\end{tabular}
\end{center}

The standalone script \srcid{check_ESM_cubic_Ward_V7.py} verifies the
closed Ward contraction as an exact polynomial identity using
$C_\mu=1+r_\mu^2\ell_\mu$. It independently checks the full four-dimensional
$a^2$ Ward polynomial, the seed correction, the transverse improvement and
paired-mode divergence in symbolic arithmetic. Separately labelled numerical
diagnostics check the unexpanded trigonometric formula, its $a^2$ comparison,
the finite Gaussian source derivative and \eqref{eq:v7-next-algebra}.
None of those floating tests is used as a proof of continuum convergence.

\begingroup
\renewcommand{\refname}{Additional methodological references for V7}
\begin{thebibliography}{99}
\bibitem[26]{V7BarnichHenneaux}
G.~Barnich and M.~Henneaux, \emph{Consistent couplings between fields with a
gauge freedom and deformations of the master equation}, Physics Letters B
\textbf{311} (1993), 123--129, arXiv:hep-th/9304057,
doi:10.1016/0370-2693(93)90544-R. The local-versus-unrestricted deformation
problem is methodological context, not a supplied primitive for this lattice.
\bibitem[27]{V7Kato}
M.~Kato, M.~Sakamoto and H.~So, \emph{Taming the Leibniz Rule on the Lattice},
arXiv:0803.3121 (2008). General locality/Leibniz obstruction as comparison;
the metric--Higgs tensor test and its repair are calculated above.
\end{thebibliography}
\endgroup

\section{Final ledgers for V7}
\label{sec:v7-ledgers}

\subsection*{Proved}
\addcontentsline{toc}{subsection}{V7: Proved}
The declared local Hodge writer has the explicit metric seed
\eqref{eq:v7-seed-vertex}. The naive phase-velocity Ward problem fails the
necessary symmetric zero-transfer test \cref{prop:v7-integrability}.
The dispersion velocity and \eqref{eq:v7-closed-vertex} give an exact
pole-free symmetric cubic completion, preserving the minimal continuum
vertex while retaining independent transverse improvements. Their actual
action/antifield differences form the explicit restoring primitive
\eqref{eq:v7-primitive}, with the finite EFT expansions and volume/mesh-uniform
rooted kernel bounds proved above. Its first finite scalar BV divergence is
zero. The measured Higgs shell yields \eqref{eq:v7-shell-output}, including
the geometric one-point coefficient and every displayed source contact.
The explicit regulator-correction sums have \eqref{eq:v7-summed-correction}.
The next nearest-neighbour transformation ansatz has the nonclosure
\eqref{eq:v7-next-algebra}. All assertions retain the scopes in the table.

\subsection*{Inherited}
\addcontentsline{toc}{subsection}{V7: Inherited}
V1--V6 and their labels are unchanged. The common finite matter carrier and
edge-Higgs form, internal gauge identities and geometric stress definition
come from \cite{V6Duality}. The native phase symbol, flat gravitational
quadratic/constraint data and distinction between exact-action and harmonic
writer come from \cite{V6Spacetime} and V6. Filtered trajectory and
source-complete measured elimination results are used only at their stated
hypotheses. The source does not prescribe the sign-averaged variable-metric
Hodge realization; that is an explicit regulator choice made in this version,
not an inherited equality between different source writers.

\subsection*{Literature input}
\addcontentsline{toc}{subsection}{V7: Literature input}
The BV deformation framework and the general lattice Leibniz obstruction
are established comparison results in \cite{V7BarnichHenneaux,V7Kato}.
They supply neither the displayed native Ward tensor nor its uniform kernel
estimate. Gravity-as-EFT and prior perturbative gravity constructions retain
the status recorded in V6; no novelty claim is made for their existence.

\subsection*{Conditional}
\addcontentsline{toc}{subsection}{V7: Conditional}
Using this primitive inside the full ESM shell requires its compatible
extension to the other geometric source writers and to nonzero internal
gauge fields. Fermionic spin currents, the nonlinear coframe and connection
measure, and the pure gravitational cubic sector must also be included. The first-order mixed BV identity does not
imply the quartic master equation. The separated-cutoff comparison controls
only the displayed coefficients, not all renormalized ESM amplitudes.
The full weighted R1--R4 realization and continuum conclusion of V6 therefore
remain conditional.

\subsection*{Open}
\addcontentsline{toc}{subsection}{V7: Open}
The next coefficient problem is to assemble the native pure-gravity cubic
and this explicit mixed cocycle on one coframe/connection measure, and then
solve the first coupled quartic master equation, including the
metric--metric--Higgs seagull and source/antifield partners. Its primitive
must have the same scale-local bounds and must not discard physical
transverse EFT operators. Complete gravitational loop subtraction, mixed
spin/gauge sectors, global/source-compatible regulator transport, the gauge
compact collar and chiral phase remain further requirements. No invariant
full ESM shell class, full perturbative continuum, finite-parameter quantum
gravity, asymptotic safety or nonperturbative mass gap is claimed.


\clearpage
\part*{V8: The complete Higgs-loop metric two-jet and its subtracted continuum limit}
\addcontentsline{toc}{part}{V8: Complete Higgs-loop metric two-jet and continuum limit}

\section{From the quartic question to a complete measured coefficient}
\label{sec:v8-start}

Sections~1--55 are preserved. The unresolved equation
\eqref{eq:v7-next-cocycle} cannot be solved by assigning a continuum
nonlinear diffeomorphism algebra to the finite action. In particular, the
native pure-gravity cubic and its ghost partner have not been computed in
this lineage. The present installment evaluates the missing
metric--metric--Higgs action vertex and then follows it, together with both
cubic insertions, through the actual normalized Higgs integral. It proves a
cutoff theorem for that complete two-metric, one-Higgs-loop coefficient.
The finite nonlinear master equation remains a separate, explicitly
calculated defect; it is not assumed in the cutoff proof.

There is an upstream simplification. The tensor of
\cref{thm:v7-completion} requires small \emph{total metric momentum}, but its
relative scalar momentum need not be small. Written with one unwrapped
scalar routing, it is periodic in the loop momentum. This supplies a
well-defined extension to the entire scalar Brillouin zone, including
wraparound pairs. No scalar ultraviolet collar is deleted. After the
complete local Taylor jet is assigned to the counterterm bank, the
renormalized lattice error is $O(a^2\log(1/(am)))$, rather than the divergent
unsubtracted estimate in \eqref{eq:v7-summed-correction}.

\subsection{Scope and nonduplication}

The common matter carrier and its edge-Higgs quadratic form are inherited
from \cite{V6Duality}; the particular sign-averaged Hodge realization is
exactly \eqref{eq:v7-Hodge-writer}. The native phase/scalar symbols, the
unwindowed repaired cubic tensor, and its fixed-order symbol estimates are
inherited from V7. General Gaussian cumulant integration, Schur allocation,
and subtraction before absolute estimates are inherited from V1 and the
monographs. The new objects are the full-relative-momentum extension, the
actual Hodge seagull, the complete second measured coefficient, and a
quantitative convergence and Ward test for their specific loop integral.
Lattice finite-part renormalization and scalar-induced gravitational
polarization are established subjects \cite{V6Reisz,V8ReiszRen,V8ParkWoodard};
no first construction of those subjects is claimed here.

Work in the four-dimensional Euclidean matter reference of V7, with a
fixed $m>0$, $n_H$ real scalar components, and internal links equal to the
identity. For the Higgs packet $n_H=4$. The coefficient has zero Higgs
self-coupling, Yukawa coupling and internal gauge coupling. A positive
reference mass is not identified with all physical Higgs/Goldstone masses.
If a different mass is used in the common action, its compensating quadratic
insertion remains in the ESM state; its resummation and the limit $m\downarrow0$
are not proved below. Gravitational fields are retained geometric sources in
this calculation, not additional integrated Gaussian modes. The scalar
measure is the same flat coefficient measure as in V7. Additional
coframe/connection, local Lorentz and chiral densities are not set to one.
They simply do not belong to this specified scalar integral.

Let $L_\mu/a$ be odd integers, $\mu=0,1,2,3$, and put
\[
 \mathcal B_a=(-\pi/a,\pi/a)^4,\qquad
 \mathcal P_{a,L}=\mathcal B_a\cap\prod_\mu(2\pi/L_\mu)\mathbb Z,
 \qquad \int_p^L=\frac1{\operatorname{Vol}_L}\sum_{p\in\mathcal P_{a,L}}.
\]
The thermodynamic version of $\int_p^L$ is
$\int_{\mathcal B_a}\dd^4p/(2\pi)^4$.
Retained metric/ghost transfers satisfy $|k|\le m/4$ and $am\le1/32$.
A fixed external test window may be placed strictly inside this ball.
All derivatives in $k$ below refer to the stated smooth interpolation of
these finite Fourier coefficients, before restriction to allowed transfers.
A finite-torus subtraction convention is thus specified, not left implicit.

\section{The full scalar carrier and the actual geometric seagull}
\label{sec:v8-seagull}

\subsection{No scalar momentum sector has to be discarded}

For $p\in\mathbb R^4$ and small transfer $k$, define
\begin{equation}
 \mathbb V_a(p,k):=V^\sharp(-p,p+k),\qquad
 \mathbb A_a(p,k;H)=\tfrac12 H_{\mu\nu}\mathbb V_{a,\mu\nu}(p,k),
 \label{eq:v8-full-vertex}
\end{equation}
using the unwrapped arguments on the right side. Here $H$ is a symmetric
metric polarization, normalized with the ordinary Hilbert--Schmidt norm.
For the bare comparison replace $\sharp$ by $b$ and write
$\mathbb A_a^b$. Matrix field insertions $\mathcal A_a(h)$ have precisely
these Fourier kernels. They are first derivatives of the scalar Hessian,
not the tensor $\mathbb V_a$ itself; the factor $1/2$ is important.

\begin{proposition}[Full-relative-momentum extension of the cubic coefficient]
\label[proposition]{prop:v8-full-relative}
For $|ak_\mu|<\pi$, $\mathbb V_a(p,k)$ is periodic in each $p_\mu$ with
period $2\pi/a$. It gives a real symmetric quadratic operator on the full
finite scalar carrier, also when $p+k$ crosses a Brillouin-zone seam. Its
exact Ward contraction is
\begin{equation}
 \kappa_\mu(k)\mathbb V_{a,\mu\nu}(p,k)
 =-v_\nu(p)P_a(p+k)+v_\nu(p+k)P_a(p).
 \label{eq:v8-full-Ward}
\end{equation}
For fixed transfer its Fourier polynomial in $p$ has lattice displacement
length at most two. Only the retained metric-transfer dependence needs the
small chart. In particular there is no inverse soft momentum and no omitted
high-scalar sector in this extension.
\end{proposition}
\begin{proof}
In \eqref{eq:v7-closed-data} set $t=-p-k/2$. Then $w_\mu=\sin(at_\mu)/a$
and $c_\mu=\cos(at_\mu)$ are periodic in $p$, as is
$A=(P_a(p)+P_a(p+k))/2$. All remaining factors depend only on $k$.
Every product contains at most two factors with a nonconstant $p$ frequency,
each of which has one elementary lattice frequency. This proves periodicity
and the displacement claim. The scalar-exchange and reversal identities
proved in V7 give
$\mathbb V_a(p,k)=\mathbb V_a(p+k,-k)$ and
$\mathbb V_a(-p,-k)=\mathbb V_a(p,k)$, with periodic continuation understood.
These are the transpose and reality conditions for the finite operator.
Apply \eqref{eq:v7-exact-repaired-Ward} to $(-p,p+k)$. Its two right-hand
symbols are periodic, so the identity is independent of a change of scalar
representative by a reciprocal lattice vector. The total transfer $k$ is
\emph{not} replaced by the sum of two independently wrapped representatives.
This proves \eqref{eq:v8-full-Ward} at seams as well as away from them.
\end{proof}

This extension removes the two scalar factors of V7's three-leg window,
while restricting the metric/ghost panel as stated. It is an explicitly
specified extension of the unwindowed coefficient, not an assertion that
V7's previously windowed finite action is unchanged. The finite cubic BV
identity extends by \eqref{eq:v8-full-Ward}; its scalar measure divergence
still vanishes by the full paired sum of the odd velocity $v(p)$.
Neither statement closes the nonlinear ghost algebra.

\subsection{Differentiating the same Hodge action twice}

For symmetric matrices $H,G$ define
\begin{align}
 s(H,G)&=\tfrac14\operatorname{tr}H\operatorname{tr}G
                         -\tfrac12\operatorname{tr}(HG),\\*
 T(H,G)&=s(H,G)I-\tfrac12\{(\operatorname{tr}H)G
                    +(\operatorname{tr}G)H\}+HG+GH.
 \label{eq:v8-Hodge-second}
\end{align}
For matrix fields these expressions are pointwise. Let $Q_a^b(g)$ be the
scalar Hessian of \eqref{eq:v7-Hodge-writer}, so that
$S_{H,a}[g,\phi]=\langle\phi,Q_a^b(g)\phi\rangle/2$ and $Q_a^b(I)=P_a$.
Write $\mathcal B_a(h_1,h_2)=D_g^2Q_a^b(I)[h_1,h_2]$.

\begin{proposition}[The native metric--metric--Higgs vertex]
\label[proposition]{prop:v8-seagull}
The complete second variation of the chosen Hodge writer is
\begin{equation}
 \boxed{\mathcal B_a(h_1,h_2)=\frac1{16}
 \sum_{\epsilon\in\{+,-\}^4}(D_\mu^{\epsilon_\mu})^*
 M_{T^{\mu\nu}(h_1,h_2)}D_\nu^{\epsilon_\nu}
                    +m^2 M_{s(h_1,h_2)}.}
 \label{eq:v8-B-operator}
\end{equation}
Thus the action coefficient at order $\varepsilon^2$ for $g=I+\varepsilon h$
is $\langle\phi,\mathcal B_a(h,h)\phi\rangle/4$.
For $h_1=He^{ikx}$ and $h_2=Ge^{-ikx}$ the diagonal scalar symbol is
\begin{equation}
 \boxed{\mathbb B_a(p;H,G)=
  \sum_\mu T_{\mu\mu}(H,G)\kappa_\mu(p)^2
 +\sum_{\mu\ne\nu}T_{\mu\nu}(H,G)v_\mu(p)v_\nu(p)
 +m^2s(H,G).}
 \label{eq:v8-B-symbol}
\end{equation}
It is independent of $k$. In particular this seagull is a local two-metric
contact, not a source-independent vacuum scalar.
\end{proposition}
\begin{proof}
Ordinary matrix differentiation at $I$ gives
\[
 D\sqrt{\det g}[H]=\tfrac12\operatorname{tr}H,\quad
 D^2\sqrt{\det g}[H,G]=s(H,G),\quad
 D^2g^{-1}[H,G]=HG+GH.
\]
The two product-rule cross terms in $D^2(\sqrt{\det g}\,g^{-1})$ give
$-(\operatorname{tr}H)G/2-(\operatorname{tr}G)H/2$, proving
\eqref{eq:v8-Hodge-second} and \eqref{eq:v8-B-operator}.
For opposite metric transfers the multiplier $T$ is constant. The equal-index
sign average of $(D_\mu^\epsilon)^*D_\mu^\epsilon$ is $\kappa_\mu^2$;
for unequal indices the independent averages are $v_\mu v_\nu$.
The mass multiplier contributes $m^2s$. This proves the displayed symbol.
No continuum product rule is used in this calculation.
\end{proof}

\subsection{An actual Gaussian family, not only an unevaluated jet}

Define the full finite scalar quadratic operator
\begin{equation}
 Q_a(h)=Q_a^b(I+h)+\Delta\mathcal A_a(h),\qquad
 \Delta\mathcal A_a=\mathcal A_a^\sharp-\mathcal A_a^b.
 \label{eq:v8-actual-Q}
\end{equation}
The repair is linear in the retained metric. Its first derivative is
$\mathcal A_a^\sharp$ and its second derivative is the native
$\mathcal B_a$. Higher derivatives of the Hodge seed are retained in
\eqref{eq:v8-actual-Q}, not deleted. The absence of an additional quartic
\emph{restoring} term is a stated choice; it will leave the explicit
nonlinear defect in \cref{sec:v8-boundary}.

Use Fourier coefficients $h(x)=\sum_k\widehat h(k)e^{ikx}$ and the norm
$\|h\|_{\mathrm W}=\sum_k\|\widehat h(k)\|$, with support $|k|\le m/4$.

\begin{lemma}[Uniform finite Gaussian domain]
\label[lemma]{lem:v8-positive}
There is a $\rho>0$, independent of $a$, the periods and $m$ on the stated
scaled domain, such that real $\|h\|_{\mathrm W}<\rho$ implies
\begin{equation}
 \tfrac12P_a\preceq Q_a(h)\preceq\tfrac32P_a.
 \label{eq:v8-positive}
\end{equation}
The finite normalized Gaussian integral and its logarithm are analytic in a
possibly smaller complex $\mathrm W$ ball. This is an actual massive scalar
Gaussian family with small retained geometry; it is not positivity of the
Euclidean gravitational action.
\end{lemma}
\begin{proof}
The pointwise metric matrices and their inverses stay in a fixed
neighbourhood of $I$ because $\|h\|_\infty\le\|h\|_{\mathrm W}$.
The Hodge formula gives, by its sixteen positive forms,
\[
 |\langle u,(Q_a^b(I+h)-P_a)u\rangle|
             \le C\|h\|_{\mathrm W}\langle u,P_a u\rangle.
\]
For the correction, $P_a(p)$ and $P_a(p+k)$ are comparable when $|k|\le m/4$.
The full-zone estimates proved in \cref{lem:v8-symbol} give
\[
 \frac{|\Delta\mathbb A_a(p,k;H)|}
       {\sqrt{P_a(p)P_a(p+k)}}
 \le C a^2(m+|p|)^2\|H\|\le C'\|H\|,
\]
since $am\le1/32$ and $a|p|\le2\pi$. The convolution Schur bound in Fourier
space consequently bounds
$\|P_a^{-1/2}\Delta\mathcal A_a(h)P_a^{-1/2}\|$ by
$C'\|h\|_{\mathrm W}$. Choose $\rho$ to make the sum at most $1/2$.
The same operator-norm bound for complex coefficients, together with the
convergent logarithm near the identity, proves analyticity at each finite
regulator on a common radius. The unanchored log determinant itself is
extensive; no volume-independent absolute bound on it is asserted.
\end{proof}

\section{The complete second measured coefficient, including source returns}
\label{sec:v8-measured}

Take any positive finite covariance $C$, including a genuine Higgs innovation
or the full $P_a^{-1}$. Let $\eta\sim N(0,\hbar C)$ and set
$y=\phi+CJ$. For one metric direction abbreviate
$A=\mathcal A_a^\sharp(h)$, $B=\mathcal B_a(h,h)$, and write
$\ell=\mathcal R_v(c)^*\phi^*$ in the fixed ordering for which
$\langle\ell,x\rangle=\langle\phi^*,\mathcal R_v(c)x\rangle$.
The coefficients of $\ell$ are even products of external odd variables.
They commute in the even Grassmann algebra, but their nilpotent terms are
not dropped. Define
\[
 F_1(x)=\tfrac12\langle x,Ax\rangle+\langle\ell,x\rangle,
 \qquad F_2(x)=\tfrac14\langle x,Bx\rangle.
\]
The following formula evaluates the new seagull and return; the general
cumulant rule itself is inherited.

\begin{proposition}[Source-complete quartic-order Higgs pushforward]
\label[proposition]{prop:v8-pushforward}
The expansion of
\[
 -\hbar\log\E\exp\{\hbar^{-1}\langle J,\eta\rangle
 -\hbar^{-1}[\varepsilon F_1(\phi+\eta)
                         +\varepsilon^2F_2(\phi+\eta)+O(\varepsilon^3)]\}
\]
has zeroth coefficient $-\langle J,CJ\rangle/2$, first coefficient as in
\eqref{eq:v7-shell-output}, and second coefficient
\begin{equation}
 \boxed{\begin{aligned}
 W_2={}&\tfrac14\langle y,By\rangle
       -\tfrac12\langle Ay+\ell,C(Ay+\ell)\rangle\\
 &+\frac{\hbar}{4}\operatorname{Tr}(BC)
       -\frac{\hbar}{4}\operatorname{Tr}(ACAC).
 \end{aligned}}
 \label{eq:v8-source-output}
\end{equation}
The trace includes the scalar-component multiplicity. In particular the
new term $-\langle\ell,C\ell\rangle/2$, the two-source contact
$\langle CJ,BCJ\rangle/4-\langle ACJ,CACJ\rangle/2$, and the mixed
antifield--ghost--source terms all remain in the output.
\end{proposition}
\begin{proof}
After the exact source shift, the coefficient is
$\E F_2(y+\eta)-(2\hbar)^{-1}\operatorname{Var}(F_1(y+\eta))$.
The linear centered part of $F_1$ is $\langle Ay+\ell,\eta\rangle$;
its variance is $\hbar\langle Ay+\ell,C(Ay+\ell)\rangle$.
The centered quadratic part has variance
$\hbar^2\operatorname{Tr}(ACAC)/2$. Its covariance with the linear part
vanishes by Gaussian parity. Also
$\E F_2=\langle y,By\rangle/4+\hbar\operatorname{Tr}(BC)/4$.
These identities hold coefficientwise in the finite even Grassmann algebra
and prove \eqref{eq:v8-source-output}. The positivity lemma justifies actual
real differentiation when $C=P_a^{-1}$; other positive shells have their own
finite Gaussian neighbourhood. No complete quartic master identity is used.
\end{proof}

For the full scalar block define the background-normalized one-loop
functional
\[
 \Gamma_a[h]=\frac{\hbar n_H}{2}\log\frac{\det Q_a(h)}{\det P_a}.
\]
Only its source-independent value at $h=0$ has been subtracted. Its mixed
second metric derivative is
\begin{equation}
 D^2\Gamma_a(0)[h_1,h_2]=\frac{\hbar n_H}{2}
 \operatorname{Tr}\{P_a^{-1}\mathcal B_a(h_1,h_2)
       -P_a^{-1}\mathcal A_a^\sharp(h_1)
                    P_a^{-1}\mathcal A_a^\sharp(h_2)\}.
 \label{eq:v8-metric-Hessian}
\end{equation}
For opposite plane-wave transfers its coefficient per physical volume is
\begin{equation}
 \Pi_{a,L}(k)[H,G]=\frac{\hbar n_H}{2}\int_p^L
 \left\{\frac{\mathbb B_a(p;H,G)}{P_a(p)}
 -\frac{\mathbb A_a(p,k;H)\mathbb A_a(p,k;G)}
                      {P_a(p)P_a(p+k)}\right\}.
 \label{eq:v8-polarization}
\end{equation}
The second term is the connected bubble and the first is the full seagull.
The same expression with $\mathbb A^b$ is the unrepaired Hodge calculation.
The scalar covariance is periodic; $p+k$ is routed as in
\cref{prop:v8-full-relative}, not cut off at a seam.

\subsection{One loop cannot be replaced by a sum of independent shell bubbles}

If $C=\sum_{j=0}^J C_j$ is a positive commuting covariance decomposition,
then the bubble in \eqref{eq:v8-metric-Hessian} contains
\[
 \sum_{i,j}\operatorname{Tr}(C_i A_1 C_j A_2),
\]
not only its diagonal terms $i=j$. For the symmetric bilinear metric
response, give the unordered pair $\{i,j\}$ to its larger index. This
partitions the full bubble exactly. Equivalently insert a smooth dyadic
partition of unity in \emph{one routed loop momentum} when estimating
\eqref{eq:v8-polarization}; the other propagator remains the complete one.
Both procedures retain every cross-shell return once. They are applications
of \cref{prop:v4-shell-return}, not new independence assumptions. The
coefficient can be produced by successive measured Higgs shells and the
remaining massive Gaussian block, but this is not an invariant interacting
ESM shell construction.

\section{A finite local bank and a uniform bound on the subtracted integrand}
\label{sec:v8-estimates}

Let $\mathsf T_4$ be the Taylor polynomial in $k$ at zero of total degree at
most four. Reality and the symmetric momentum set make
$\Pi_{a,L}(-k)=\Pi_{a,L}(k)$. Thus only degrees zero, two and four occur.
Define the reference-normalized two-source coefficient by
\begin{equation}
 \Pi^R_{a,L}=(I-\mathsf T_4)\Pi_{a,L}.
 \label{eq:v8-renormalized}
\end{equation}
The Taylor polynomial is stored as an actual local metric-source
counterterm. It is not declared zero in the original integral. In
particular its degree-zero part contains the seagull. A redundant bank of
all even quadratic metric-source monomials through four derivatives has
dimension at most
\[
 \dim\operatorname{Sym}^2(\R^{10})
 \left\{1+\binom52+\binom73\right\}
             =55(1+10+35)=2530.
\]
This is an upper bound for an off-shell \emph{restoring} bank in the fixed
chart, not a count of independent covariant gravitational Wilson
coefficients. Covariant, equation-of-motion and topological identifications
must still be imposed on a common nonlinear action. A common finite local
polynomial may be added after \eqref{eq:v8-renormalized} to impose different
matched quadratic renormalization conditions.

Set
\begin{equation}
 f_a(p,k;H,G)=\frac{\mathbb A_a(p,k;H)\mathbb A_a(p,k;G)}
                         {P_a(p)P_a(p+k)},\qquad
 f_a^+(p,k)=\tfrac12\{f_a(p,k)+f_a(p,-k)\}.
 \label{eq:v8-integrand}
\end{equation}
The polarization arguments in $f_a^+$ are understood. Its integrated value
is that of $f_a$, by $p\mapsto-p$. Write $f_0$ for the expression with
$P_0(p)=m^2+|p|^2$ and the continuum tensor
\eqref{eq:v7-continuum-vertex}.

\begin{lemma}[Full-zone symbol estimates, including differentiated mismatch]
\label[lemma]{lem:v8-symbol}
Put $R_p=m+|p|$. For $p\in\mathcal B_a$, $|k|\le m/4$, $am\le1/32$,
and every fixed multi-indices $\alpha,\beta$,
\begin{align}
 P_a(p+k)&\asymp R_p^2,\\*
 |\partial_p^\beta\partial_k^\alpha f_a(p,k)|
       +|\partial_p^\beta\partial_k^\alpha f_0(p,k)|
   &\le C_{\alpha,\beta}R_p^{-|\alpha|-|\beta|}\|H\|\|G\|,\\*
 |\partial_p^\beta\partial_k^\alpha(f_a-f_0)(p,k)|
   &\le C_{\alpha,\beta}a^2R_p^{2-|\alpha|-|\beta|}\|H\|\|G\|.
 \label{eq:v8-symbol-bounds}
\end{align}
The same bounds hold for the bare vertex. Constants do not depend on the
periods or on $a,m$ in the displayed scaled domain. The last bound is allowed
to be of order one at $|p|\asymp a^{-1}$; it does not assert that ultraviolet
lattice configurations are pointwise close to their continuum counterparts.
\end{lemma}
\begin{proof}
On the enlarged cube containing $p\pm k$,
$\sin(z/2)/z$ is bounded above and below away from zero in absolute value
with its removable value at zero. Since $|k|\le m/4$, this gives the first
comparison. Direct differentiation gives
\[
 |\partial^\gamma P_a|\le C_\gamma R_p^{2-|\gamma|},\quad
 |\partial^\gamma(P_a-P_0)|\le
 C_\gamma a^2R_p^{4-|\gamma|},
\]
and, by differentiating the reciprocal,
\[
 |\partial^\gamma P_a^{-1}|\le C_\gamma R_p^{-2-|\gamma|},\quad
 |\partial^\gamma(P_a^{-1}-P_0^{-1})|
                           \le C_\gamma a^2R_p^{-|\gamma|}.
\]
A bound at derivative orders exceeding the polynomial degree follows from
$aR_p\le C$: each additional trigonometric derivative supplies a factor
$a\le C/R_p$. For the low derivative orders, use
$|\sin z-z|\le |z|^3/6$ and the analogous cosine remainder.

In \eqref{eq:v8-full-vertex}, $w=O(R_p)$, $A=O(R_p^2)$, $r=O(m)$,
$c,C=O(1)$, and $\ell=O(a^2)$. The denominator $1+C_\mu$ has a fixed
positive margin, even for a scalar momentum at the zone boundary. Applying
the same differentiated sine/cosine estimates to
\eqref{eq:v7-closed-vertex} gives
\[
 |\partial_p^\beta\partial_k^\alpha\mathbb A_a|
 \le C R_p^{2-|\alpha|-|\beta|}\|H\|,\qquad
 |\partial_p^\beta\partial_k^\alpha(\mathbb A_a-\mathbb A_0)|
 \le C a^2R_p^{4-|\alpha|-|\beta|}\|H\|.
\]
The bare tensor has the same properties from its explicit sign average.
Use the product rule for the two numerators and two reciprocals, replacing
one factor at a time in the difference. This proves
\eqref{eq:v8-symbol-bounds}. The argument concerns derivatives of the
unwrapped smooth representative near each $p$, which is exactly the
representative used for the finite coefficient.
\end{proof}

\begin{lemma}[Six derivatives, rather than an absolute extensive bound]
\label[lemma]{lem:v8-six-derivatives}
For $r=|\alpha|\le6$ and any fixed $\beta$,
\begin{align}
 |\partial_p^\beta\partial_k^\alpha
       (I-\mathsf T_4)f_0^+(p,k)|
 &\le C_{\alpha,\beta}|k|^{6-r}R_p^{-6-|\beta|}\|H\|\|G\|,\\*
 |\partial_p^\beta\partial_k^\alpha
       (I-\mathsf T_4)(f_a^+-f_0^+)(p,k)|
 &\le C_{\alpha,\beta}a^2|k|^{6-r}R_p^{-4-|\beta|}\|H\|\|G\|.
 \label{eq:v8-subtracted-integrand}
\end{align}
For $r>6$ the respective bounds are $C R_p^{-r-|\beta|}$ and
$C a^2R_p^{2-r-|\beta|}$, with the polarization norms understood.
\end{lemma}
\begin{proof}
The even function $f^+$ has zero homogeneous Taylor components of odd total
degree. Hence its Taylor polynomial through degree five equals the one
through degree four. Apply the multivariate integral Taylor formula through
degree five along $t\mapsto tk$, including differentiated versions for
$r\le6$. All needed derivatives have total order six, and
\cref{lem:v8-symbol} gives the two estimates. For $r>6$ the polynomial has
vanished under differentiation and the undifferentiated symbol estimates
apply. This uses an exact even symmetrization \emph{before} taking absolute
values; a fifth-derivative estimate alone would give a weaker cutoff rate.
\end{proof}

The unsubtracted local jets obey the bounds
\begin{equation}
 |D_k^r\Pi_{a,L}(0)|\le C\hbar n_H
 \begin{cases}
 a^{-4},&r=0,\\
 a^{-2},&r=2,\\
 1+\log(1/(am)),&r=4,
 \end{cases}
 \qquad mL_\mu\ge1.
 \label{eq:v8-local-growth}
\end{equation}
Indeed the bubble integrand of order $r$ is bounded by $R_p^{-r}$;
four-dimensional dyadic counting gives the three lines. The seagull enters
only the first, with $\mathbb B_a/P_a$ bounded. These are actual local
counterterm budgets, not small errors in the nonlocal action.

\section{A full-cutoff theorem for the two-metric Higgs-loop coefficient}
\label{sec:v8-limit}

Define the massive continuum periodic coefficient by the absolutely
convergent sum
\begin{equation}
 \Pi^R_{0,L}(k)[H,G]
 =-\frac{\hbar n_H}{2\operatorname{Vol}_L}
   \sum_{p\in\prod_\mu(2\pi/L_\mu)\mathbb Z}
               (I-\mathsf T_4)f_0^+(p,k;H,G).
 \label{eq:v8-continuum-torus}
\end{equation}
Its infinite-volume counterpart $\Pi^R_{0,\infty}$ uses
$\int_{\mathbb R^4}\dd^4p/(2\pi)^4$. The seagull has not been omitted from
the input: its whole two-source contribution belongs to the degree-zero
term in \eqref{eq:v8-renormalized}.

\begin{theorem}[Massive one-Higgs-loop, two-metric continuum estimate]
\label[theorem]{thm:v8-continuum}
Assume $am\le1/32$, $mL_\mu\ge1$, and $|k|\le m/4$.
For every multi-index $\alpha$ with $r=|\alpha|\le6$,
\begin{equation}
 \boxed{\left|\partial_k^\alpha
 (\Pi^R_{a,L}-\Pi^R_{0,L})(k)[H,G]\right|
 \le C_\alpha\hbar n_H\|H\|\|G\|\,
 |k|^{6-r}a^2\left[1+\log\frac1{am}\right].}
 \label{eq:v8-main-rate}
\end{equation}
The same estimate holds in infinite volume and for the unrepaired bare Hodge
cubic. For any fixed $r>6$ the right side may be replaced by
$C_r\hbar n_H\|H\|\|G\|a^2m^{6-r}[1+\log(1/(am))]$.
Thus the result includes fixed momentum derivatives, not only pointwise
convergence. The constants are uniform in the four periods under the
stated lower-size condition. No mass-uniform limit as $m\downarrow0$ is
asserted.
\end{theorem}
\begin{proof}
The finite seagull is $k$ independent, so applying $I-\mathsf T_4$ to
\eqref{eq:v8-polarization} leaves exactly the subtracted bubble. Both
Taylor differentiation and even symmetrization commute with the finite
sum. Subtract \eqref{eq:v8-continuum-torus}. The difference has two pieces:
the integrand difference on $\mathcal P_{a,L}$, and the continuum tail on
momenta outside the finite Brillouin cube. By
\cref{lem:v8-six-derivatives}, their absolute values at derivative order
$r\le6$ are bounded by $C\hbar n_H\|H\|\|G\||k|^{6-r}$ times
\begin{equation}
 \frac{a^2}{\operatorname{Vol}_L}
       \sum_{p\in\mathcal P_{a,L}}(m+|p|)^{-4},\qquad
 \frac1{\operatorname{Vol}_L}
       \sum_{p\notin\mathcal B_a}(m+|p|)^{-6}.
 \label{eq:v8-two-sums}
\end{equation}
For $R\ge m$ the number of periodic momenta with $|p|\le R$, divided by
$\operatorname{Vol}_L$, is at most
$C\prod_\mu(R+L_\mu^{-1})\le C'R^4$. This includes the zero mode.
Partition the first sum into $|p|\le m$ and dyadic annuli from $m$ to
$2\pi/a$. Each annulus costs a constant, so its bound is
$C a^2[1+\log(1/(am))]$. The second sum starts at $|p|\ge\pi/a$;
its dyadic tail is bounded by $C a^2$. These estimates prove
\eqref{eq:v8-main-rate}. For $r>6$, the corresponding sums have powers
$2-r$ and $-r$ and are bounded by $C a^2m^{6-r}$; the displayed weaker
logarithmic bound is valid as well. Replacing counting by radial integration
proves the infinite-volume version. Bare and repaired vertices satisfy the
same symbol hypotheses, so their limits coincide. All sums are anchored
coefficient sums; no extensive determinant bound is taken.
\end{proof}

\begin{proposition}[The thermodynamic comparison is a separate limit]
\label[proposition]{prop:v8-volume}
Let $L_{\min}=\min_\mu L_\mu$ and $mL_{\min}\ge1$. For every fixed
$s>4$ and $r=|\alpha|\le6$,
\begin{equation}
 \left|\partial_k^\alpha(\Pi^R_{0,L}-\Pi^R_{0,\infty})(k)[H,G]\right|
 \le C_{s,\alpha}\hbar n_H\|H\|\|G\|
 |k|^{6-r}m^{-2}(mL_{\min})^{-s}.
 \label{eq:v8-volume-rate}
\end{equation}
The high-derivative version replaces $|k|^{6-r}$ by $m^{6-r}$.
Consequently the joint lattice/infinite-volume error is the sum of
\eqref{eq:v8-main-rate} and \eqref{eq:v8-volume-rate}; the two limits
have not been conflated.
\end{proposition}
\begin{proof}
The subtracted continuum integrand has integrable $p$ derivatives of every
fixed order. Integrating the first estimate of
\eqref{eq:v8-subtracted-integrand} gives an $L^1_p$ derivative bound
$C|k|^{6-r}m^{-2-s}$. Fourier integration by parts and Poisson summation
therefore bound the difference between the periodic sum and the integral
by this constant times
$\sum_{n\in\mathbb Z^4\setminus\{0\}}|(L_\mu n_\mu)_\mu|^{-s}$.
It is at most $C_sL_{\min}^{-s}$ for $s>4$. Approximation by smooth
compactly supported integrands, followed by the same integrable derivative
bounds, justifies Poisson summation for this polynomially decaying
integrand. This proves the claim without an ultraviolet volume factor.
\end{proof}

\subsection{The corresponding source-kernel topology}

Let $\varpi\in C_c^\infty(\mathbb R^4)$ have support $|z|<1/4$.
The physical source window is $\varpi(k/m)$, chosen radial when rotational
covariance is used. The bounds for all fixed $k$ derivatives imply a
quasilocal kernel statement: if $\mathcal K^\Delta_{a,L}$ is the inverse
Fourier kernel of
$\varpi(k/m)(\Pi^R_{a,L}-\Pi^R_{0,L})(k)$ on the lattice, then
\begin{equation}
 \boxed{a^4\sum_x(1+m d_{\rm per}(x,0))^s
       \|\mathcal K^\Delta_{a,L}(x)\|
 \le C_s\hbar n_H m^6a^2[1+\log(1/(am))].}
 \label{eq:v8-kernel-rate}
\end{equation}
The tensor kernel is measured in fixed finite component norms. To prove this,
extend the smooth momentum symbol by zero beyond the external window,
integrate by parts more than $s+4$ times, and periodize the resulting
infinite-mesh kernel. Its scaled derivative bounds are
$C_q\hbar n_Hm^{6-q}a^2[1+\log(1/(am))]$ by
\cref{thm:v8-continuum}. Exactly the kernel summation argument of
\cref{thm:v7-kernel} then gives \eqref{eq:v8-kernel-rate}.
This is a windowed local-source topology, not a claim about arbitrarily high
external momenta or a massless source space.

For metric profiles in that window, with
$\|h\|_{3,L}^2=\operatorname{Vol}_L\sum_k|k|^6\|\widehat h(k)\|^2$,
\eqref{eq:v8-main-rate} also gives the direct bilinear functional estimate
\begin{equation}
 |D^2(\Gamma^R_a-\Gamma^R_0)[h_1,h_2]|
 \le C\hbar n_Ha^2[1+\log(1/(am))]
                        \|h_1\|_{3,L}\|h_2\|_{3,L}.
 \label{eq:v8-source-rate}
\end{equation}
The degree-zero geometric one-point term is normalized separately by its
actual local linear metric counterterm. For a fixed finite panel of
additional scalar, elementary source, ghost and antifield coefficients with
all external momenta in $|p|<m/32$, every tree term in
\eqref{eq:v8-source-output} converges with $O(a^2)$ coefficient error.
Indeed all its internal momenta are sums of a fixed bounded number of these
external momenta, its massive denominators are bounded below, and its
symbols and source shifts have ordinary $O(a^2)$ Taylor errors there.
There is no loop momentum in these source-return terms. The two trace terms
are precisely the separately renormalized geometric coefficient just
proved. This gives source-complete convergence for the displayed finite
coefficient, not a master identity for all of those sources.

\subsection{Matched schemes and higher finite derivative banks}

In \eqref{eq:v8-renormalized}, both the sharp and bare schemes have their
own \emph{actual} Taylor polynomial subtracted. Subtracting a common
unevaluated bare polynomial in both would not be matching. Once the same
finite physical quadratic polynomial is added back, the two limits agree
by \cref{thm:v8-continuum}. The same proof applies to any other real,
exchange-symmetric full-zone vertex with the bounds of
\cref{lem:v8-symbol} and the same continuum limit. A fixed finite-range
pointwise geometric seagull changes only the local counterterm in this
comparison. No independence of unrelated nonlinear regulators is asserted.

For the thermodynamic sharp-minus-bare local jet one can be more explicit.
At each even $r=|\alpha|\le4$, scaling $u=ap$, $z=ak$, $\mu=am$ gives
\begin{equation}
 \partial_k^\alpha(\Pi_a^\sharp-\Pi_a^b)(0)
 =a^{r-4}\mathcal M_\alpha(am),\qquad
 \mathcal M_\alpha(0)=
 -\frac{\hbar n_H}{2}\int_{(-\pi,\pi)^4}\frac{\dd^4u}{(2\pi)^4}
 \left.\partial_z^\alpha(f_1^\sharp-f_1^b)(u,z;0)\right|_{z=0}.
 \label{eq:v8-finite-matching}
\end{equation}
The last integral is absolutely convergent. Near $u=0$ its integrand is
bounded by $C|u|^{2-r}$, which is locally integrable in dimension four
for $r\le4$; away from zero all denominators have a fixed margin.
Dominated convergence proves $\mathcal M_\alpha(am)\to\mathcal M_\alpha(0)$.
In particular a fourth-derivative matching correction is finite and need not
be declared zero merely because its original vertex was $O(a^2)$ at fixed
momentum. The lower local jets can have the $a^{-4}$ and $a^{-2}$ sizes.
Equation~\eqref{eq:v8-finite-matching} specifies computable matching
integrals, not fitted beta coefficients.

For any fixed integer $q>2$, define $\Pi^{R,2q}_{a,L}$ by subtracting
the actual finite Taylor jet through degree $2q$, and define
$\Pi^{R,2q}_{0,L}$ by the corresponding convergent subtracted continuum sum.
Then the stronger finite-EFT-order comparison is
\begin{equation}
 |\Pi^{R,2q}_{a,L}(k)-\Pi^{R,2q}_{0,L}(k)|
 \le C_q\hbar n_H\|H\|\|G\|\,
 a^2m^{4-2q}|k|^{2q+2},
 \label{eq:v8-higher-subtraction}
\end{equation}
with the corresponding fixed derivative versions. Here the continuum
coefficient is understood through its convergent subtracted sum.
The proof uses $2q+2$ derivatives: the mismatch annulus contributes
$a^2R^{4-2q}$ instead of the constant $a^2$ at $q=2$. Its dyadic sum is
bounded by $C_qa^2m^{4-2q}$, and the omitted continuum tail is smaller.
The local bank remains finite for each $q$, with the redundant bound
$55\sum_{j=0}^q\binom{2j+3}{3}$. This is arbitrary fixed derivative order
\emph{in this one-loop, two-metric coefficient}, not an all-loop ESM theorem.

\section{A derived Ward limit and a nonzero continuum normalization check}
\label{sec:v8-Ward-limit}

The next theorem concerns the thermodynamic continuum coefficient. On a
fixed torus a Taylor subtraction of a chosen continuous momentum
interpolation can alter finite-volume contact terms. We do not infer a
finite-torus covariant subtraction from the continuum argument; its matched
stress/tadpole convention must be treated separately.

\begin{theorem}[Transversality of the limiting nonlocal metric response]
\label[theorem]{thm:v8-Ward-limit}
The tensor represented by $\Pi^R_{0,\infty}$ is rotationally covariant and
satisfies, for the stated external momenta,
\begin{equation}
 k_\mu\Pi^{R\,\mu\nu,\rho\sigma}_{0,\infty}(k)=0.
 \label{eq:v8-limit-Ward}
\end{equation}
The lattice renormalized longitudinal part therefore vanishes at the rate
obtained by multiplying \eqref{eq:v8-main-rate} by $|k|$, and its finite-volume
comparison has the additional term from \eqref{eq:v8-volume-rate}.
This proves the linearized two-source Ward identity of the limiting
nonlocal coefficient. It is not a proof of finite nonlinear BV closure.
\end{theorem}
\begin{proof}
This identity is derived by a convergent loop-momentum argument rather than
imported as an assumption. Regulate the continuum bubble by
$w_M(p,k)=z(p/M)^2z((p+k)/M)^2$, with a smooth even $z$ equal to one near
zero. The degree-four-subtracted coefficient converges to
$\Pi^R_{0,\infty}$, with tail $C\hbar n_H|k|^6M^{-2}$. Indeed, after
even symmetrization, the sixth transfer derivatives of the integrand on
$|p|\ge cM$ are $O(|p|^{-6})$. Terms differentiating the cutoff are
supported on $|p|\asymp M$ and have the same $O(M^{-6})$ bound; their
integrals are $O(M^{-2})$. Taylor's formula proves the stated tail.

For the continuum tensor,
\[
 k_\mu V^{(0)}_{\mu\nu}(-p,p+k)
                 =-p_\nu P_0(p+k)+(p+k)_\nu P_0(p).
\]
Contract the regulated bubble, and change the loop variable by $k$ in the
second term. Up to its common constant factor, the resulting integral is
\begin{equation}
 \int\frac{\dd^4p}{(2\pi)^4}\frac{p_\nu}{P_0(p)}
 \{w_M(p-k,k)V^{(0)}_{\rho\sigma}(-p+k,p)
                  -w_M(p,k)V^{(0)}_{\rho\sigma}(-p,p+k)\}.
 \label{eq:v8-contracted-loop}
\end{equation}
Replace both weights in braces by $w_M(p,0)$. The tensor difference is
exactly
\[
 2(p_\rho k_\sigma+p_\sigma k_\rho
                                      -\delta_{\rho\sigma}p\cdot k),
\]
so this part is a polynomial linear in $k$. The remaining terms are
supported on $|p|\asymp M$ for $M\ge8m$. Their seventh $k$ derivatives
have integral bound $CM^{-2}$: the numerator divided by $P_0(p)$ has
order one, the four-dimensional measure has order four, and each derivative
lowers the scale by one. The contracted response is odd in $k$, as is the
linear polynomial, so its remainder after Taylor degree five starts at
degree seven and is bounded by $C|k|^7M^{-2}$.
Finally
\[
 (I-\mathsf T_5)(k_\mu\Pi^{\mu\nu,\rho\sigma})
              =k_\mu(I-\mathsf T_4)\Pi^{\mu\nu,\rho\sigma}.
\]
The seagull also contracts to a linear polynomial and is included in the
same subtraction. Let $M\to\infty$ to prove \eqref{eq:v8-limit-Ward}.
Rotational covariance follows directly from the absolutely convergent
subtracted continuum integral: its symbols are covariant, the Taylor
operation at zero is covariant, and the full integration domain is
$\mathbb R^4$. No rotational hypothesis on the finite lattice cutoff is
needed after this limit.
\end{proof}

The fixed local counterterms of degrees zero, two and four have not been
identified with a complete nonlinear covariant action by this proof.
It shows that the \emph{nonlocal remainder} has the required transverse
limit. In the massive reference it is analytic near zero and starts at
$|k|^6$. At nonzero $k$ it has the usual two transverse tensor channels:
with $\pi_{\mu\nu}=k^2\delta_{\mu\nu}-k_\mu k_\nu$, they are
\[
 \tfrac12(\pi_{\mu\rho}\pi_{\nu\sigma}+\pi_{\mu\sigma}\pi_{\nu\rho})
                  -\tfrac13\pi_{\mu\nu}\pi_{\rho\sigma},
 \qquad \pi_{\mu\nu}\pi_{\rho\sigma}.
\]
Indeed rotate $k$ to one axis. Transversality leaves the space of symmetric
three-dimensional transverse tensors, whose trace and traceless parts are
the two irreducible orthogonal-group components. Their two scalar
coefficients determine the tensor. Analyticity and the vanishing Taylor
jet remove any artificial soft-pole interpretation of this representation.

\subsection{An evaluated transverse-traceless component}

Take $k$ along the zero direction and a real symmetric $H$ with
$Hk=0$, $\operatorname{tr}H=0$, $\operatorname{Tr}H^2=1$. Put
$z=|k|^2$, $d_x=x(1-x)z$ and $s_x=m^2+d_x$.

\begin{proposition}[A nonzero one-Higgs-loop normalization]
\label[proposition]{prop:v8-TT-value}
In the geometric-source normalization of this installment,
\begin{equation}
 \boxed{\Pi^R_{0,\infty}(k)[H,H]
 =\frac{\hbar n_H}{128\pi^2}\int_0^1
 \left\{s_x^2\log\frac{s_x}{m^2}
                   -m^2d_x-\tfrac32d_x^2\right\}\dd x.}
 \label{eq:v8-TT-formfactor}
\end{equation}
It is strictly positive for real $z>0$, and its first term is
\begin{equation}
 \Pi^R_{0,\infty}(k)[H,H]
   =\frac{\hbar n_H}{53760\pi^2m^2}|k|^6
                         +O(\hbar n_H|k|^8/m^4).
 \label{eq:v8-TT-leading}
\end{equation}
No gravitational coupling has been inserted in these source coefficients;
conversion to a canonically normalized graviton requires the stated
coframe/metric normalization of the common action.
\end{proposition}
\begin{proof}
The continuum vertex contracted with such $H$ is
$\mathbb A_0(p,k;H)=-H_{ij}p_ip_j$, with transverse indices. Use
\[
 \frac1{P_0(p)P_0(p+k)}
       =\int_0^1\frac{\dd x}{[(p+xk)^2+m^2+x(1-x)k^2]^2}.
\]
The numerator is unchanged by the shift because $Hk=0$.
Its angular average in four dimensions is
$\operatorname{Tr}(H^2)|p|^4/12=|p|^4/12$. To avoid shifting a divergent
integral without a justification, first differentiate three times in $z$.
These differentiated integrals are absolutely convergent; equivalently
use the smooth cutoff argument of \cref{thm:v8-Ward-limit}, whose
six-momentum-derivative boundary error vanishes. Integrate back with the
three conditions at $z=0$ supplied by $\mathsf T_4$.

For $s>0$ the radial integral $F(s)=\int \dd^4p/(2\pi)^4\,
|p|^4/(p^2+s)^2$ has a well-defined third derivative
\[
 F'''(s)=-24\int\frac{\dd^4p}{(2\pi)^4}
            \frac{|p|^4}{(p^2+s)^5}=-\frac{3}{8\pi^2s}.
\]
The last equality uses the radial measure
$(16\pi^2)^{-1}t\dd t$ and
$\int_0^\infty t^3(t+s)^{-5}\dd t=1/(4s)$.
The unique primitive with its Taylor polynomial through degree two
subtracted at $s=m^2$ is
\[
 -\frac3{16\pi^2}
 \{s^2\log(s/m^2)-m^2(s-m^2)-\tfrac32(s-m^2)^2\}.
\]
Multiply by the bubble factor $-\hbar n_H/2$, the angular factor $1/12$,
and integrate over $x$. This proves \eqref{eq:v8-TT-formfactor}.
The function inside braces has positive third derivative $2/s$ and zero
first three Taylor data at $m^2$, so it is positive for $s>m^2$.
Its cubic term is $d_x^3/(3m^2)$; since
$\int_0^1x^3(1-x)^3\dd x=1/140$, the coefficient is
$1/(128\cdot420)=1/53760$, proving \eqref{eq:v8-TT-leading}.
This calculation is a normalization check of the limiting integral, not
an independent definition of the lattice answer.
\end{proof}

\section{The quartic BV defect has not been hidden by the loop limit}
\label{sec:v8-boundary}

The native seagull is now explicit, so the first remaining classical
identity can be written without an unspecified action vertex. With the
free metric transformation $R_a c$, the first matter transformation
$R_c=\mathcal R_v(c)$, and no yet-added nonlinear metric transformation,
the antifield-free order-two Noether residual is
\begin{equation}
 \boxed{\mathcal D_{2,0}(h,c,\phi)=\frac12
 \left\langle\phi,
 \{\mathcal B_a(h,R_ac)+R_c^*\mathcal A_a^\sharp(h)
                           +\mathcal A_a^\sharp(h)R_c\}\phi\right\rangle.}
 \label{eq:v8-explicit-defect}
\end{equation}
This follows by varying
$\langle\phi,\mathcal B_a(h,h)\phi\rangle/4$ in $h$ and
$\langle\phi,\mathcal A_a^\sharp(h)\phi\rangle/2$ in $\phi$.
The ghost coefficient is read with a commuting test parameter before its
fixed Grassmann ordering is restored. The antifield component also contains
the commutator of two $R_c$'s, whose nearest-neighbour nonclosure was already
proved in V7. Neither is set to zero here.

To solve the \emph{coupled} quartic master equation one must add the native
pure-gravity cubic and its ghost bracket, its induced nonlinear metric
variation in the mixed vertex, and the appropriate higher matter
transformation and/or restoring seagull. Equation~\eqref{eq:v8-explicit-defect}
is the concrete mixed action term they must cancel on the same measure.
The generated return $-\langle\ell,C\ell\rangle/2$ in
\eqref{eq:v8-source-output} is a further indispensable source term after a
shell; it is not, by itself, a primitive for the full input defect.

The direction change is therefore limited but substantive. Instead of
postulating that this finite nonlinear problem has been solved, the
complete one-Higgs-loop metric two-jet has been shown to possess a matched,
quasilocal continuum limit by a direct comparison. Its nonlocal linearized
Ward identity is proved, while its local counterterms and finite nonlinear
BV defect retain separate ownership. This is a realized coefficient
interface for ESM, not a declaration that the weighted R1--R4 hypotheses
hold for the full interacting theory. In particular no use was made of an
internally supplied running gravitational coupling or of an assumed
all-scale chiral phase.

\subsection{Uniformity and reproducibility}

\begin{center}
\small
\begin{tabular}{>{\raggedright\arraybackslash}p{.22\textwidth}>{\raggedright\arraybackslash}p{.69\textwidth}}
\toprule
Variable & Precise scope in V8\\
\midrule
Scalar ultraviolet modes & Every scalar momentum of the finite Brillouin
zone is retained. Only the declared external geometry/source panel is
restricted to low momentum.\\
Mesh & $am\le1/32$; the subtracted cutoff error is
$a^2[1+\log(1/(am))]$ at fixed $m>0$, with fixed differentiated versions.\\
Volume and time & Cutoff constants are uniform for $mL_\mu\ge1$.
Thermodynamic errors are displayed separately. No finite-torus covariant
subtraction theorem is silently substituted for the infinite-volume Ward
statement.\\
EFT and loop order & One closed Higgs loop with two external metrics;
every prescribed finite derivative subtraction is covered. No all-loop
constant or infinite EFT-bank control is proved.\\
Sources and measure & The exact finite expression retains all scalar,
metric, ghost and antifield contacts through the displayed interaction
order, in the declared flat scalar measure. Fixed windowed source panels
converge; arbitrary massless or relational source algebras are not covered.\\
Gravitational and SM sectors & Metric is retained; pure-graviton loops,
nonlinear coframe/connection measure, local Lorentz, internal gauge/Yukawa
mixing and chiral-line transport are not calculated here.\\
Repeated shells & Every pair of covariance shells has one owner. The summed
subtracted loop estimate is proved for this coefficient, not as an invariant
nonlinear many-shell ESM class.\\
Mass and coupling domain & $m>0$ is fixed. The small-background scalar
Gaussian radius is uniform in the specified norm. The massless limit and
mass-compensation resummation remain separate.\\
\bottomrule
\end{tabular}
\end{center}

The standalone script \srcid{check_ESM_metric_twojet_V8.py} checks the mixed
matrix derivatives of $\sqrt{\det g}\,g^{-1}$, the finite Gaussian second
coefficient and shell-pair allocation, and the continuum contracted-vertex
identity. Exact arithmetic also verifies the rational coefficient
$1/53760$ independently by integrating the six-momentum Taylor coefficient.
Numerical diagnostics test full-zone periodicity at wraparound pairs,
finite Hodge derivatives, actual log-determinant/source differentiation and
the native transverse-traceless sixth coefficient. Numerical tests are
labelled as such and do not prove the cutoff estimates.

\begingroup
\renewcommand{\refname}{Additional methodological references for V8}
\begin{thebibliography}{99}
\bibitem[28]{V8ReiszRen}
T.~Reisz, \emph{Renormalization of Feynman integrals on the lattice},
Communications in Mathematical Physics \textbf{117} (1988), 79--108,
doi:10.1007/BF01228412. Established lattice finite-part renormalization;
not a supplied estimate for the repaired tensor in this installment.
\bibitem[29]{V8ParkWoodard}
S.~Park and R.~P.~Woodard, \emph{Scalar Contribution to the Graviton
Self-Energy during Inflation}, Physical Review D \textbf{83} (2011), 084049,
arXiv:1101.5804. Comparison for scalar-induced gravitational polarization
and its curvature-squared subtraction; its curved, massless setting is not
identified with the massive flat finite-regulator calculation above.
\end{thebibliography}
\endgroup

\section{Final ledgers for V8}
\label{sec:v8-ledgers}

\subsection*{Proved}
\addcontentsline{toc}{subsection}{V8: Proved}
The exact unwindowed cubic tensor extends periodically over the entire
relative scalar momentum zone at small metric transfer
(\cref{prop:v8-full-relative}). The same finite Hodge action has the explicit
seagull \eqref{eq:v8-B-operator}; together with the linear repair it defines
an actual uniformly small-background positive scalar Gaussian family.
Its second measured coefficient is \eqref{eq:v8-source-output}, with all
contact and mixed-shell terms retained. The native massive two-metric
one-loop coefficient, after its actual finite local subtraction, converges
with \eqref{eq:v8-main-rate}; its windowed kernel, source and thermodynamic
bounds are proved separately. Bare and repaired cubic schemes have the same
matched nonlocal limit and the finite local matching integrals
\eqref{eq:v8-finite-matching}. The thermodynamic limit is transverse by the
explicit loop-shift proof. Its TT component is evaluated in
\eqref{eq:v8-TT-formfactor}. Every statement has the domain in the preceding
table; none assumes the full nonlinear master equation.

\subsection*{Inherited}
\addcontentsline{toc}{subsection}{V8: Inherited}
V1--V7 and all their mathematical labels are preserved. The common finite
Higgs carrier, native Hodge seed, free phase/dispersion, cubic tensor,
source-complete Gaussian calculus and Schur ownership are used at their
stated scope. The weighted two-boundary theorem and its general continuum
consequences remain conditional on the primitive ESM hypotheses. The
positive free gravitational reference of V6, the unclosed compact gauge
collar and the chiral-line tasks are neither repeated nor strengthened by
an unrelated scalar Gaussian positivity argument.

\subsection*{Literature input}
\addcontentsline{toc}{subsection}{V8: Literature input}
Massive lattice finite-part renormalization, the BV deformation problem,
and scalar-induced gravitational polarization are established comparison
frameworks. The present proofs evaluate the chosen Hodge seagull, its
full-relative repaired coefficient, and the displayed cutoff/source errors;
no claim of a first perturbative scalar/gravity continuum theory is made.

\subsection*{Conditional}
\addcontentsline{toc}{subsection}{V8: Conditional}
Use inside a complete ESM trajectory requires common nonlinear covariant
counterterms, symmetry-compatible source transport and the actual
gravitational/internal/fermionic measures. The present finite quadratic
subtractor has not been shown to extend to a bounded projection of the
full combined BV complex. Finite-torus Ward matching requires its own stress
contact convention. A physical massless Higgs component or a reference-mass
change requires its infrared and compensating-insertion analysis. Full ESM
R1--R4 invariance does not follow from this realized one-loop component.

\clearpage
\subsection*{Open}
\addcontentsline{toc}{subsection}{V8: Open}
The smallest unresolved nonlinear mixed coefficient is now the explicit
residual \eqref{eq:v8-explicit-defect}, together with its two-ghost antifield
component and the native pure-gravity cubic contribution. A restoring
primitive must include those pieces on one coframe/connection measure with
scale-local bounds. The scalar-loop result provides a direct alternative
comparison mechanism for some symmetry-breaking terms but not that
primitive. Complete gravitational loops, massless/internal-spin sectors,
local-to-global source matching and line transport, and an invariant
filtered interacting ESM shell remain open. No all-orders ESM continuum,
finite-parameter gravity, asymptotic safety, nonperturbative measure or mass
gap is asserted.


\clearpage
\part*{V9: A complete massive Higgs one-loop source module}
\addcontentsline{toc}{part}{V9: A complete massive Higgs one-loop source module}

\section{From an individual response to a closed finite source panel}
\label{sec:v9-start}

Sections~1--63 are preserved. The next result is not another evaluation of
\eqref{eq:v8-polarization}. We construct and estimate the \emph{entire finite
source panel} of a specified massive Higgs Gaussian block: arbitrarily many
metric insertions at any prescribed finite order, electroweak connection
insertions, and a matrix-valued quadratic composite source. The complete
Hodge derivatives, cyclic contractions, and elementary-source contacts are
kept. Their ultraviolet local coefficients and complementary kernels are
assigned by one functional prescription. The resulting one-loop coefficient
trajectory is generated by the measured block, rather than supplied as an
R1--R4 hypothesis.

This is one component of the ESM programme, not the full ESM action class.
The gravitational variables remain retained backgrounds. The construction
does not integrate pure gravitons, the local Lorentz sector, gauge bosons or
chiral fermions, and it does not solve \eqref{eq:v8-explicit-defect}. In
particular, internal background-gauge symmetry and nonlinear gravitational
BV symmetry have different status below. The former is realized and its
subtraction is controlled; the latter is not assumed.

\subsection{Input, source convention, and nonduplication}

The finite Higgs representation, gauge-covariant edge kinetic form, quartic
potential, and common-action stress provenance are those of
\cite{V6Duality}, \srcid{eq:finite-SM-action}. We use the covariant seed already
specified after \eqref{eq:v7-Hodge-writer}: both incoming and outgoing
covariant differences are transported to the same site. At identity links
and vanishing new composite source this is exactly the \emph{bare Hodge
seed}, not V8's linearly repaired family \eqref{eq:v8-actual-Q}. That
choice is explicit. A metric-only comparison with the repaired family is
proved below, but no unconstructed gauge-covariant extension of its repair
is used.

Normalized Gaussian pushforward, ordered coefficient composition and source
contacts are inherited from V1. The general degree-dependent local bank,
Slavnov-splitting criterion and two-boundary theorem are inherited from
\cite{V6CMP} and V6; they are not proved anew. The new content is verification
of the full-zone vertex estimates for the present common family, their
source-degree-dependent subtraction, and the complete finite-panel cutoff,
shell and internal-Ward estimates. Massive lattice finite-part
renormalization is established literature \cite{V8ReiszRen}. Covariant
one-loop effective actions and composite-field renormalization likewise
have established frameworks \cite{V9Vassilevich,V9HollandsWald}; none is
substituted for the finite operator defined next.

\section{The actual gauge-covariant Hodge block and its sources}
\label{sec:v9-family}

Work on the four-dimensional odd periodic grids of V8. Let $\mathcal R$ be
the real orthogonal Higgs representation of the internal group and let
$d_H=\dim_{\R}\mathcal R$ ($d_H=4$ for the inherited doublet). The colour
factor acts trivially on this block. All matrices below act on
$\ell^2(a\mathbb Z^4/L\mathbb Z^4;\mathcal R)$ with physical cell weight
$a^4$. Let $U_{x\mu}$ be orthogonal parallel transport from $x+a\hat\mu$ to
$x$, and define
\[
 (T_\mu^U f)_x=U_{x\mu}f_{x+a\hat\mu},\qquad
 D_\mu^{U,+}=(T_\mu^U-I)/a,\qquad
 D_\mu^{U,-}=(I-(T_\mu^U)^*)/a.
\]
For $g=I+h$ write $\rho_g=\sqrt{\det g}$ and
$G_g^{\mu\nu}=\rho_g g^{\mu\nu}$. A real symmetric endomorphism source
$E_x\in\operatorname{Sym}(\mathcal R)$ gives the quadratic operator
\begin{equation}
 \boxed{Q_a[g,U,E]=\frac1{16}\sum_{\epsilon\in\{+,-\}^4}
 (D_\mu^{U,\epsilon_\mu})^*M_{G_g^{\mu\nu}}D_\nu^{U,\epsilon_\nu}
       +M_{\rho_g(m^2I+E)}.}
 \label{eq:v9-native-Q}
\end{equation}
Repeated spacetime indices are summed. Thus $Q_a[I,I,0]=P_aI$ with the
same $P_a$ as \eqref{eq:v7-native-symbols}. The scalar integration measure
is the flat real coefficient measure used in V7--V8. Its geometric Hodge
factors occur in $Q_a$; they are not inserted a second time in the measure.

For momentum differentiation we declare a source map, rather than a new
quantum action: $U[A]$ is exact parallel transport of a smooth
$\mathfrak{so}(\mathcal R)$-valued connection $A$ along each straight link,
with convention $\nabla_\mu=\partial_\mu+A_\mu$. Then, on smooth fields,
$T_\mu^{U[A]}$ is the sampled flow $e^{a\nabla_\mu}$. The ordered line
integral defines all link source jets. This is also a specification of the
source convention when a momentum is continuously interpolated; replacing
it by a midpoint exponential would give different finite source vertices.
The finite measured input still consists of links, site metrics and site
endomorphisms in \eqref{eq:v9-native-Q}.

The continuum comparison operator, in the same flat coordinate density, is
\begin{equation}
 Q_0[g,A,E]=-\nabla_\mu M_{G_g^{\mu\nu}}\nabla_\nu
                                  +M_{\rho_g(m^2I+E)}.
 \label{eq:v9-continuum-Q}
\end{equation}
This formula will be the limit of the differentiated finite operator, not
a replacement for its ultraviolet vertices.

\begin{proposition}[Measured covariance and a common massive domain]
\label[proposition]{prop:v9-domain}
For arbitrary sitewise orthogonal $O_x$, transform
$U_{x\mu}\mapsto O_xU_{x\mu}O_{x+a\hat\mu}^{-1}$,
$E_x\mapsto O_xE_xO_x^{-1}$, and $J_x\mapsto O_xJ_x$. Then
$Q_a\mapsto OQ_aO^{-1}$ and the finite scalar integral is invariant.
Moreover, a fixed sufficiently small $\varepsilon>0$ gives
\begin{equation}
 \|h\|_{\mathrm W}+m^{-1}\|A\|_{\mathrm W}
                      +m^{-2}\|E\|_{\mathrm W}<\varepsilon
 \quad\Longrightarrow\quad
 \tfrac12P_a\preceq Q_a\preceq\tfrac32P_a,
 \label{eq:v9-positive}
\end{equation}
with constants independent of the periods and $a$ for $am\le1/64$.
Here $\|B\|_{\mathrm W}=\sum_k\|\widehat B(k)\|$, with fixed finite
component norms. The corresponding complex Taylor family has a common
analytic neighbourhood, with transpose interpreted algebraically.
\end{proposition}
\begin{proof}
Each covariant difference transforms by $O$ at its output site. All Hodge
multipliers are scalar on $\mathcal R$, and the potential conjugates.
This proves covariance of the quadratic form. Lebesgue measure is preserved
because $|\det O|=1$; the source pairing is preserved as well.

The line integral gives $\|U[A]-I\|\le a\|A\|_\infty$ for real
skew matrices, hence $\|D^{U,\pm}-D^{I,\pm}\|\le\|A\|_\infty$.
Cauchy--Schwarz bounds the difference of the kinetic forms by
$C(m^{-1}\|A\|_\infty+m^{-2}\|A\|_\infty^2)
\langle f,P_af\rangle$. The pointwise metric inverse and density are
Lipschitz near $I$ and are bounded by the Wiener norm. The potential error
is at most $C(\|h\|_{\mathrm W}+m^{-2}\|E\|_{\mathrm W})
\langle f,P_af\rangle$. Choosing $\varepsilon$ small proves
\eqref{eq:v9-positive}. The same estimates, with a harmless exponential
factor from complex line transport, bound the relative complex perturbation
in operator norm by less than one. Inversion and the logarithm then have
convergent finite-matrix series on a common neighbourhood. No uniform bound
on the unsubtracted extensive determinant is asserted.
\end{proof}

The normalized measured functional is exactly
\begin{equation}
 \boxed{\mathcal W_a[h,A,E,J]
 =\frac{\hbar}{2}\Tr\log(Q_aP_a^{-1})
                -\frac12\langle J,Q_a^{-1}J\rangle.}
 \label{eq:v9-measured-W}
\end{equation}
It is $-\hbar$ times the logarithm of the Gaussian integral divided by its
value at $h=A=E=J=0$. The trace already includes $\mathcal R$. In particular,
$J$-quadratic contacts are not removed by this normalization. Even nilpotent
source coefficients are permitted coefficientwise. For example $J\mapsto
J-\ell$, with V8's fixed even antifield--ghost insertion, retains every
$\ell Q_a^{-1}\ell$ and mixed $J$ contact. This substitution is source
algebra, not a claim of full gravitational master-equation closure.

\subsection{What the endomorphism source represents}

$E=\sigma I$ generates the local quadratic Higgs composite in the chosen
real normalization; general symmetric $E$ retains its matrix-valued
contacts. For the inherited real potential
$\lambda_H(\Phi^T\Phi-v_H^2)^2$, its fluctuation Hessian is
\begin{equation}
 V''(\Phi)=4\lambda_H(\Phi^T\Phi-v_H^2)I
                                      +8\lambda_H\Phi\Phi^T.
 \label{eq:v9-Higgs-Hessian}
\end{equation}
This follows by two ordinary derivatives and transforms by conjugation.
Consequently a finite-order pullback $m^2I+E=V''(\Phi)+\sigma I$ assigns the
appropriate Higgs-background insertions to the same coefficient bank.
A nonzero reference-mass compensation must have a declared positive
perturbative degree. It is not resummed by this substitution, and need not
lie in the positive domain \eqref{eq:v9-positive} when evaluated at its full
physical value. The result concerns the scalar Hessian block; mixed
hard-gravity/gauge Hessian blocks, a shifted stationary background, and
higher-loop scalar interactions require their own measured analysis.

\section{All labelled source vertices and their ultraviolet degree}
\label{sec:v9-vertices}

Fix $N<\infty$. A labelled panel $I$ has $n\le N$ entries, of which $r$
are metric polarizations, $q$ connection polarizations, and $s$ endomorphism
polarizations. Put
\begin{equation}
 w(I)=q+2s,\qquad \omega(I)=4-w(I),\qquad
 D(I)=\max\{6-w(I),0\}.
 \label{eq:v9-weights}
\end{equation}
The polarization matrices are factored out of all estimates through their
norms; their product is denoted $\mathfrak n_I$. Gauge couplings can be
included in the actual connection polarizations and are not fitted later.
Take $\sum_i|k_i|\le Nm/64$, $\sum_i k_i=0$, and
$Nam\le1/64$. Constants may depend on $N$ but not on the periods or mesh.
Every intermediate loop transfer has size at most $Nm/64$. A fixed window
$|k_i|<m/64$ for each individual mark is contained in this domain. The total-
transfer formulation is preserved when momenta are grouped in a Ward
identity; it avoids relying on a window that is not closed under that grouping.
Use $K=(k_1,\ldots,k_{n-1})$, $k_n=-\sum_{i<n}k_i$; all continuously
interpolated amplitudes are symmetrized in their labelled arguments.

For a nonempty subset $B\subset I$, let $Q_{a,B}$ be the mixed derivative
of \eqref{eq:v9-native-Q} in its labelled sources at the flat reference.
Its Fourier vertex has input momentum $p$ and output $p+k_B$. No Taylor
order of the Hodge density or inverse metric needed by $B$ is discarded.

\begin{lemma}[Primitive full-zone vertex estimates]
\label[lemma]{lem:v9-vertices}
For each fixed derivative order and $R_p=m+|p|$,
\begin{align}
 \|\partial_{p,K}^{\alpha} Q_{a,B}(p;K_B)\|
 +\|\partial_{p,K}^{\alpha} Q_{0,B}(p;K_B)\|
 &\le C_{N,\alpha}R_p^{2-w(B)-|\alpha|}\mathfrak n_B,\\*
 \|\partial_{p,K}^{\alpha}(Q_{a,B}-Q_{0,B})(p;K_B)\|
 &\le C_{N,\alpha}a^2R_p^{4-w(B)-|\alpha|}\mathfrak n_B.
 \label{eq:v9-vertex-bounds}
\end{align}
The finite vertices are periodic in the routed scalar momentum. Vertices
with two $E$ entries, or with both $E$ and $A$, vanish. At $a=0$ a vertex
with more than two $A$ entries also vanishes; its finite counterpart is kept
and obeys the same bounds. The estimates hold over the entire scalar zone,
not only for low scalar momenta.
\end{lemma}
\begin{proof}
The metric derivatives of $\rho_g$ and $G_g$ at $I$ are fixed finite tensors
bounded by constants depending on $|B|$. For line transport, the $j$th
connection derivative is a sum of ordered simplex integrals. Acting on a
Fourier mode it is a product of the $j$ polarization matrices, a phase
$e^{iap_\mu}$, and phases whose arguments are linear combinations of the
$ak_i$. Its size is at most $C_j a^j\mathfrak n_B$; each momentum
derivative adds a factor $a$. The two differences in
\eqref{eq:v9-native-Q} contribute at most $a^{-2}$. An undifferentiated
difference contributes $|e^{iap_\mu}-1|/a\le C R_p$ instead. Distributing
the connection derivatives between the two differences proves the first
bound: for $j\ge2$, $a^{j-2}\le C_jR_p^{2-j}$ since $aR_p$ is bounded.
The potential is linear in $E$ and independent of $A$, proving its vanishing
claims and its degree-zero estimate. Differentiating reciprocals is not
involved at this stage.

To obtain the stronger mismatch bound, the order of the finite difference
is essential. Keeping the same smooth connection, $T_\mu(-a)=T_\mu(a)^{-1}$,
so $D_{-a}^{U,+}=D_a^{U,-}$ and conversely. Independent averaging for distinct
$\mu,\nu$, and the common sign average for $\mu=\nu$, make the entire
operator \eqref{eq:v9-native-Q} even in $a$. This holds before source
differentiation, including noncommuting line matrices. Its value at $a=0$
is \eqref{eq:v9-continuum-Q}. Locally fix the numerical scale $R_p$, rescale
all momenta and $m$ by $R_p$, and put $z=aR_p$. After the factor
$R_p^{2-w(B)}$ is removed, the simplex formula is smooth in $z$ on a fixed
compact interval, with removable values at zero. Its odd first derivative
vanishes. Taylor's integral remainder in $z$ is bounded by $Cz^2$.
The same argument after a fixed number of rescaled momentum derivatives
proves the second bound. The scale is fixed during these differentiations;
no differentiation of $|p|$ is used. Shifted scalar momenta lie in a fixed
slightly enlarged cube. All exponential factors are periodic in $p$ there,
including at a seam. The continuum differential operator is quadratic in
$A$, proving the final vanishing statement.
\end{proof}

A multi-vertex trace is periodic even though a continuously interpolated
external momentum is not a lattice label. Scalar representatives are never
independently wrapped inside the source line integrals. This is the same
routing distinction used in V8.

\begin{proposition}[The complete coefficient, with all contact vertices]
\label[proposition]{prop:v9-cycles}
Let $\operatorname{OP}(I)$ be the ordered set partitions of $I$. The mixed
$n$th derivative of the determinant part of \eqref{eq:v9-measured-W} is
\begin{equation}
 \boxed{\Gamma_{a,I}=\frac{\hbar}{2}
 \sum_{(B_1,\ldots,B_l)\in\operatorname{OP}(I)}
 \frac{(-1)^{l-1}}{l}
 \Tr(P_a^{-1}Q_{a,B_1}\cdots P_a^{-1}Q_{a,B_l}).}
 \label{eq:v9-cycles}
\end{equation}
The trace is read as one routed loop, with all successive transfers retained.
In particular the partitions with $|B_j|>1$ are the higher metric, gauge and
composite contact vertices. They are not reconstructed from products of
first vertices. Dividing a plane-wave coefficient by physical volume gives
an amplitude $\gamma_{a,L,I}(K)$ with one finite momentum sum.
\end{proposition}
\begin{proof}
Use independent commuting nilpotent markers $z_i^2=0$. Then
$Q_a-P_a=\sum_{\varnothing\ne B\subset I}z_BQ_{a,B}$. In the exact
finite logarithm series, the coefficient of $z_I$ in its $l$th term is
precisely the sum over ordered disjoint nonempty blocks covering $I$.
The series terminates at $l=n$. Its coefficient is the ordinary mixed
derivative, with no further factorial. This proves \eqref{eq:v9-cycles}
and fixes the cyclic factor $1/l$. It reuses the inherited logarithm
calculus, now with the actual all-source Hodge vertices. Fourier momentum
conservation leaves exactly one unconstrained scalar momentum.
\end{proof}

For example the three-source expression contains
$\Tr(CQ_{123})-\sum_{\rm cyc}\Tr(CQ_1CQ_{23})$ and both oriented
three-singleton cycles, with the overall factor $\hbar/2$. Matrix ordering
in the two cycles matters for noncommuting $E$ or gauge polarizations.

For each summand write $f_{a,I}(p,K)$ for its trace integrand, suppressing
its fixed rational coefficient. The propagator estimates of V8 extend to
all intermediate transfers, with constants depending on $N$. Hence
\begin{align}
 |\partial_p^\beta\partial_K^\alpha f_{a,I}|
       +|\partial_p^\beta\partial_K^\alpha f_{0,I}|
 &\le C_{N,\alpha,\beta}d_H\mathfrak n_I
                           R_p^{-w(I)-|\alpha|-|\beta|},\\*
 |\partial_p^\beta\partial_K^\alpha(f_{a,I}-f_{0,I})|
 &\le C_{N,\alpha,\beta}d_H\mathfrak n_I a^2
                           R_p^{2-w(I)-|\alpha|-|\beta|}.
 \label{eq:v9-loop-bounds}
\end{align}
To verify this, a cycle of length $l$ has $l$ reciprocal factors of order
$-2$ and vertices of total order $2l-w(I)$. Differentiate the products,
and replace one factor at a time for the mismatch. The simplex bounds also
cover the contact vertices that vanish in the continuum but survive at
finite $a$. There is no assumption that the $l$ source occurrences are
independent or that their internal matrices commute.

Under inversion of all four coordinates, $h$ and $E$ are even tensors and
$A$ is a one-form. Thus the integrated coefficient has parity
\begin{equation}
 \gamma_{a,L,I}(-K)=(-1)^q\gamma_{a,L,I}(K).
 \label{eq:v9-parity}
\end{equation}
This follows directly by interchanging the forward and backward covariant
differences and $p\mapsto-p$ on the symmetric odd grid. At integrand level
replace $f$ by $f^{[q]}(p,K)=\{f(p,K)+(-1)^qf(p,-K)\}/2$; its integral is
unchanged. This signed symmetrization must precede absolute estimates.

\section{One counterterm functional and an all-panel continuum estimate}
\label{sec:v9-continuum}

For $w\le4$ let $\mathsf P_I=\mathsf T_{4-w}$ be the complete Taylor
polynomial in $K$ at zero; for $w>4$ put $\mathsf P_I=0$. Define
\begin{equation}
 \gamma^R_{a,L,I}=(I-\mathsf P_I)\gamma_{a,L,I}.
 \label{eq:v9-subtraction}
\end{equation}
All $\mathsf P_I\gamma_{a,L,I}$, with their permutation and transpose
symmetries, are coefficients of \emph{one} local polynomial functional
$\mathcal C_{a,L,N}[h,A,E]$ through total source degree $N$. They remain
in the local state. A common prescribed finite polynomial can subsequently
be added to impose physical normalization conditions.

The connection count, not the metric count, changes the ultraviolet degree:
\begin{center}\small
\begin{tabular}{lll}
\toprule
Panel weight $w=q+2s$ & Local Taylor degree & First retained parity degree\\
\midrule
$0,1,2,3,4$ & $4-w$ & $6-w$\\
$5$ & No ultraviolet subtraction & $1$ (odd coefficient)\\
$6$ & No ultraviolet subtraction & $0$\\
$w>6$ & No ultraviolet subtraction & No gain needed\\
\bottomrule
\end{tabular}
\end{center}
The number of metric marks can be arbitrarily large at fixed $w$; it is the
prescribed finite $N$ that makes their bank finite. A redundant upper bound
on its dimension, with $d_A=4\dim\rho(\mathfrak g_{\rm SM})$ and
$d_E=d_H(d_H+1)/2$, is
\begin{equation}
 \sum_{n=1}^N\ \sum_{\substack{r+q+s=n\\q+2s\le4}}
 \binom{n}{r,q,s}10^r d_A^q d_E^s
 \binom{4(n-1)+4-q-2s}{4-q-2s}.
 \label{eq:v9-bank}
\end{equation}
This intentionally includes redundant component polynomials and integration-
by-parts representatives. It is not a count of independent covariant ESM
Wilson coefficients. In particular no nonlinear diffeomorphism identification
of the metric polynomials follows from this count.

\begin{theorem}[Source-complete massive one-loop cutoff bound]
\label[theorem]{thm:v9-continuum}
Fix $N$ and the source convention of \cref{sec:v9-vertices}. Assume
$Nam\le1/64$ and $mL_\mu\ge1$. For every panel $I$ of size at most $N$,
$\gamma^R_{0,L,I}$ is defined by the absolutely convergent continuum
momentum sum of the subtracted cyclic integrands. For each fixed multi-index
$\alpha$, $r_\alpha=|\alpha|$, one has
\begin{equation}
 \boxed{m^{r_\alpha}
 |\partial_K^\alpha(\gamma^R_{a,L,I}-\gamma^R_{0,L,I})|
 \le C_{N,\alpha}\hbar d_H\mathfrak n_I\,
 m^{4-w(I)}(am)^2\left[1+\log\frac1{am}\right].}
 \label{eq:v9-main}
\end{equation}
The same statement holds for thermodynamic integrals. Constants do not
depend on the periods or mesh on this domain. For $r_\alpha\le D=D(I)$
the sharper right side is
\begin{equation}
 C_{N,\alpha}\hbar d_H\mathfrak n_I\,
 a^2m^{6-w-D}|K|^{D-r_\alpha}
                      [1+\log(1/(am))]
 \label{eq:v9-sharp-rate}
\end{equation}
before the factor $m^{r_\alpha}$ on the left is inserted. The empty
momentum panel $n=1$ has its literal, constant interpretation. No bound
uniform in $N$, or in $m\downarrow0$, is asserted.
\end{theorem}
\begin{proof}
The parity of $w$ is that of $q$. For $w\le4$, the Taylor component of
order $5-w$ has the wrong parity and is zero after
\eqref{eq:v9-parity}. Thus $I-\mathsf T_{4-w}$ is the Taylor remainder
starting at order $D=6-w$. For $w=5$ the integrand symmetrization is odd
and starts at order one, although no local counterterm is required.
For $w\ge6$ use the unsubtracted bound with $D=0$.

Taylor's integral formula, \eqref{eq:v9-loop-bounds}, and their differentiated
versions give, for $r_\alpha\le D$,
\begin{align*}
 |\partial_p^\beta\partial_K^\alpha f_{0,I}^R|
 &\le C d_H\mathfrak n_I |K|^{D-r_\alpha}
                                      R_p^{-w-D-|\beta|},\\
 |\partial_p^\beta\partial_K^\alpha(f_{a,I}^R-f_{0,I}^R)|
 &\le C d_H\mathfrak n_I a^2|K|^{D-r_\alpha}
                                      R_p^{2-w-D-|\beta|}.
\end{align*}
Always $w+D\ge6$. The first bound is integrable in four dimensions. The
second is at worst $a^2R_p^{-4}$, occurring when $w\le6$; its sum over
dyadic annuli up to $a^{-1}$ is at most $Ca^2[1+\log(1/(am))]$.
If $w>6$ its sum is at most $Ca^2m^{6-w}$. The omitted continuum tail
starts at $\pi/a$ and is bounded by $Ca^{w+D-4}|K|^{D-r_\alpha}$,
which is no larger than the displayed mismatch bound because $am<1$.
All normalized periodic counts obey
$\operatorname{Vol}_L^{-1}\#\{|p|\le R\}\le C R^4$ for $R\ge m$,
as $mL_\mu\ge1$. Thus the argument is uniform in the four periods.

At derivative orders above $D$, differentiate the symbols directly; their
extra powers of $R_p^{-1}$ give the required factors of $m^{-1}$. Finally
$|K|\le C_Nm$ on the declared source panel. Summing the finitely many
ordered partitions gives \eqref{eq:v9-main}. Integral estimates give the
thermodynamic version. The same-label Hodge contacts are included in every
step, not estimated as an extensive action. A single-vertex contact with
no relative external transfer is assigned wholly to its Taylor polynomial
when its weight requires subtraction.
\end{proof}

All numerical factors in \eqref{eq:v9-cycles} are fixed by the input
Gaussian law. The estimate is not an assertion that the higher-curvature
or composite-source coefficients are zero. For example, at constant
$E=eI$, $h=A=0$, the coefficient of $e^3$ is already ultraviolet convergent
and has thermodynamic continuum value
\begin{equation}
 \frac{\hbar d_H}{6}\int\frac{\dd^4p}{(2\pi)^4}(p^2+m^2)^{-3}
                  =\frac{\hbar d_H}{192\pi^2m^2}.
 \label{eq:v9-E3}
\end{equation}
The radial integral equals $1/(32\pi^2m^2)$. This is a nonzero triple
composite-source contact coefficient, not the previously calculated
metric two-jet. The $E^2$ contact, by contrast, has weight four and its
zero-momentum local coefficient must be renormalized.

\subsection{Kernel topology, elementary sources, and matching}

Choose a fixed smooth external window in $K/m$ with each $|k_i|<m/64$.
For the inverse Fourier kernel of the difference in
\eqref{eq:v9-main}, on the $n-1$ relative sites, one obtains
\begin{equation}
 \boxed{a^{4(n-1)}\sum_{x_1,\ldots,x_{n-1}}
 (1+m\sum_i d_{\rm per}(x_i,0))^b\|\mathcal K^\Delta_{a,L,I}\|
 \le C_{N,b}\hbar d_H\mathfrak n_I
 m^{4-w}(am)^2[1+\log(1/(am))].}
 \label{eq:v9-kernel}
\end{equation}
To prove it, differentiate the smooth windowed symbol more than
$4(n-1)+b$ times, use \eqref{eq:v9-main} at those fixed orders, integrate
by parts on the full external momentum space, and periodize. The same
weighted lattice summation as in V7 applies in $4(n-1)$ dimensions.
An integrated source functional is then bounded with one source in $L^1$
and the others in $L^\infty$. There is no total-volume factor in this
rooted bound. Normalizing $A$ by $m$ and $E$ by $m^2$ makes the common
coefficient scale $m^4$ for every panel.

Every derivative of the second term of \eqref{eq:v9-measured-W} has two
$J$ ends and a finite ordered resolvent chain. For prescribed total degree,
all its internal momenta are sums of the fixed low external momenta; its
massive denominators have a common lower bound depending only on $N,m$.
The vertex estimates therefore give $O(a^2)$ convergence, with the
corresponding fixed-window kernel bounds. Derivatives with more than two
independent elementary $J$ marks vanish identically in this Gaussian
module. Arbitrarily high \emph{composite} $E$ and metric source derivatives
need not vanish and are instead covered by \cref{thm:v9-continuum}.
The same finite argument applies to $J-\ell$ in a fixed Grassmann panel.
It does not prove a new symmetry identity for the gravitational $\ell$.

For two schemes satisfying \eqref{eq:v9-vertex-bounds} with the same
$Q_0$, subtract their own actual local functionals and impose the same
finite normalization polynomial. Their finite-panel nonlocal limits then
agree by \cref{thm:v9-continuum}. This comparison includes the metric-only
restriction of V8's repaired family: its only altered primitive vertex is
the first one, whose full-zone mismatch is already proved in
\cref{lem:v8-symbol}; every higher vertex is the common Hodge derivative.
All cyclic mixed occurrences of that correction are included by
\eqref{eq:v9-cycles}. No gauge-covariant repaired vertex is inferred from
this metric-only comparison.

At fixed $N$, the local coefficients at momentum Taylor degree $d\le4-w$
have the generated bounds
\begin{equation}
 |u_{a,I,d}|\le C_N\hbar d_H\mathfrak n_I
 \begin{cases}
 a^{-(4-w-d)},&w+d<4,\\
 1+\log(1/(am)),&w+d=4.
 \end{cases}
 \label{eq:v9-local-budgets}
\end{equation}
These follow by summing $R_p^{-w-d}$. They are local normalization
stocks, not small complementary errors. In particular differences between
schemes can leave finite local matching coefficients even when their
fixed-momentum primitive mismatch is $O(a^2)$.

\section{A measured shell trajectory with internally assigned local coordinates}
\label{sec:v9-shells}

The previous result sums all scalar ultraviolet modes. It also gives a
coefficientwise many-shell construction, not merely an existence limit.
Let $\vartheta$ be smooth nonincreasing, equal to one on $[0,1]$ and zero
on $[4,\infty)$. Set $\Lambda_j=2^jm$ and
\begin{equation}
 C_{a,\le j}(p)=\frac{\vartheta((P_a(p)-m^2)/\Lambda_j^2)}{P_a(p)},
 \qquad C_{a,j}=C_{a,\le j}-C_{a,\le j-1}\succeq0,
 \label{eq:v9-covariance-shells}
\end{equation}
with $C_{a,0}=C_{a,\le0}$. These are genuine covariance innovations.
At finite $a$ a finite last scale makes $C_{a,\le j}=P_a^{-1}$ on every
mode. At $a=0$ their sum is the full massive covariance. The low part
includes the scalar zero mode; it is not removed.

Write $V=Q-P$ and let $\Gamma_{\le j}$ be the determinant term from
integrating a Gaussian of covariance $C_{\le j}$ against
$\exp(-\langle\eta,V\eta\rangle/(2\hbar))$. On its support its exact
value is $\hbar\Tr\log(I+C_{\le j}V)/2$. Successive convolution of the
independent innovations gives this same measured coefficient when the
retained quadratic form and all sources are kept by the inherited Schur
rule. Put $\Delta_j=\Gamma_{\le j}-\Gamma_{\le j-1}$.

\begin{proposition}[Exact ownership and a summable high-shell bound]
\label[proposition]{prop:v9-shell-bound}
Every length-$l$ cycle in $\Delta_j$ consists of all internal shell-label
assignments whose maximum label is $j$, with the factor and ordering of
\eqref{eq:v9-cycles}. There is no restriction that its $l$ labels agree.
For $j\ge j_N$, with $\Lambda_j\ge C_Nm$, its subtracted coefficients
obey, for every fixed external derivative or windowed rooted kernel norm,
\begin{equation}
 \boxed{\| (I-\mathsf P_I)\Delta_{j,I}\|_{\rm src,m}
       \le K_N\hbar d_H\mathfrak n_I\,
                     m^{4-w(I)}\,2^{-2j}.}
 \label{eq:v9-shell-decay}
\end{equation}
The bound is uniform in volume and mesh on the stated domain, through the
last finite-lattice shell. The finitely many lower shells have bounds
$K_N\hbar d_H\mathfrak n_I m^{4-w}$. In a thermodynamic shell the local
Taylor degree-$d$ increment obeys
$|\mathsf T_d\Delta_{j,I}|\le K_N\hbar d_H\mathfrak n_I
\Lambda_j^{4-w-d}$ for $d\le4-w$.
\end{proposition}
\begin{proof}
Expand each propagator in a cycle as a sum of innovations. Subtracting the
all-labels-at-most-$j-1$ expansion leaves exactly the stated maximum-label
assignments. Alternatively telescope the product of the $l$ cumulative
covariances. The resulting expression has at least one difference
$C_j$ and the other factors are complete cumulative covariances. No
independence of loop factors is asserted.

For $\Lambda_j\ge C_Nm$, any nonzero high difference forces the single
routed momentum and all its bounded translates to have size comparable
to $\Lambda_j$. Derivatives of the smooth cumulative cutoffs gain
$\Lambda_j^{-1}$. The symbol proof of \eqref{eq:v9-loop-bounds} therefore
applies to the complete telescoped expression. The parity proof still
applies because the covariance cutoffs are even. After the same
$D=\max(6-w,0)$ gain, the annulus costs at most
$C|K|^D\Lambda_j^{4-w-D}$. Since $w+D\ge6$, it is bounded by
$C_Nm^{4-w}(m/\Lambda_j)^2$. Fixed external derivatives, with their
powers of $m$, and the windowed kernel proof give
\eqref{eq:v9-shell-decay}. The low annuli are bounded by the mass.
Differentiating only to degree $d$ without subtraction gives the local
increment estimate. Near the top lattice shell the same bounds use
$P_a-m^2\asymp |p|^2$ in the fundamental cube and smooth periodic
covariance cutoffs. No sharp spatial shell is substituted.
\end{proof}

\begin{theorem}[Realized finite-panel one-loop coefficient trajectory]
\label[theorem]{thm:v9-trajectory}
Fix a finite physical local polynomial $u_{\rm phys}$ in
\eqref{eq:v9-bank}. For this measured Gaussian block there is a uniquely
specified normalized coefficient trajectory
\begin{align}
 u_j&=u_{\rm phys}-\mathsf P\Gamma_{\le j},&
 R_j&=(I-\mathsf P)\Gamma_{\le j},\\*
 u_j-u_{j-1}&=-\mathsf P\Delta_j,&
 R_j-R_{j-1}&=(I-\mathsf P)\Delta_j.
 \label{eq:v9-generated-flow}
\end{align}
Thus $\Gamma_{\le j}+u_j=u_{\rm phys}+R_j$ at every step. In the continuum
reference, at every fixed source degree and source norm of
\cref{prop:v9-shell-bound},
\begin{equation}
 \|R_\infty-R_j\|_{\rm src,m}
       \le K_N\hbar d_H\mathfrak n_I m^{4-w}\,2^{-2j}.
 \label{eq:v9-trajectory-tail}
\end{equation}
The elementary-source resolvent coefficients stabilize once the covariance
plateau contains their finitely many internal low momenta. The last finite
lattice trajectory converges to this continuum source panel with
\eqref{eq:v9-main}. The local sequence has power growth for $w+d<4$ and
at most linear growth in $j$ for $w+d=4$, with its increments determined by
the actual measured traces.
\end{theorem}
\begin{proof}
The exact Gaussian convolution determines each $\Gamma_{\le j}$, hence
each increment. Specifying the physical Taylor normalization then determines
$u_j$ and $R_j$ uniquely; subtraction is an assignment, not deletion.
Summing the geometric series in \eqref{eq:v9-shell-decay} proves
\eqref{eq:v9-trajectory-tail}. Summing its unsubtracted local increment
bound gives the stated power or linear growth. A tree chain has no free
loop momentum; when all cumulative scalar cutoffs equal one at its finitely
many internal momenta it equals the full chain exactly. The finite-lattice
last scale is the complete coefficient in \cref{thm:v9-continuum}, so its
comparison follows from that theorem. None of these arguments uses a
supplied marginal running law.
\end{proof}

\begin{corollary}[An evaluated native composite marginal]
\label[corollary]{cor:v9-composite-flow}
For a constant symmetric endomorphism $E$, at $h=A=0$, the quadratic
composite-source part of the generated counterterm increment is exactly
\begin{equation}
 \delta u_{a,L,j}^{E^2}
 =\frac{\hbar}{4}\operatorname{Tr}_{\mathcal R}(E^2)
 \frac1{\operatorname{Vol}_L}\sum_p
 \left\{C_{a,\le j}(p)^2-C_{a,\le j-1}(p)^2\right\}.
 \label{eq:v9-E2-exact}
\end{equation}
In the continuum thermodynamic reference, for $\Lambda_{j-1}\ge2m$,
\begin{equation}
 \boxed{\delta u_{0,\infty,j}^{E^2}
 =\frac{\hbar\log2}{32\pi^2}\operatorname{Tr}_{\mathcal R}(E^2)
 +O\!\left(\hbar\operatorname{Tr}_{\mathcal R}(E^2)
                         \frac{m^2}{\Lambda_{j-1}^2}\right).}
 \label{eq:v9-E2-flow}
\end{equation}
For each fixed $b>4$, its finite-lattice, finite-volume error has the
additional bound
\begin{equation}
 C_b\hbar\operatorname{Tr}_{\mathcal R}(E^2)
 \left\{(a\Lambda_j)^2+(\Lambda_{j-1}L_{\min})^{-b}\right\},
 \qquad a\Lambda_j\le1/64.
 \label{eq:v9-E2-lattice}
\end{equation}
This is an internally calculated marginal of the measured source block,
not an assumed full ESM beta coefficient.
\end{corollary}
\begin{proof}
The determinant coefficient quadratic in a constant $E$ is
$-\hbar\operatorname{Tr}_{\mathcal R}(E^2)\sum_p C_{\le j}(p)^2/4$.
The counterterm has the opposite sign, giving \eqref{eq:v9-E2-exact}.
In particular the numerator is a \emph{difference of squares}, not the
square of $C_j$; the cross-shell return is present.

For the leading thermodynamic coefficient replace $(p^2+m^2)^{-2}$
by $|p|^{-4}$ on the difference's annulus. Its error integrates to
$C m^2/\Lambda_{j-1}^2$. Four-dimensional radial integration gives
\[
 \frac1{8\pi^2}\int_0^\infty\frac{\dd r}{r}
 \{\vartheta(r^2)^2-\vartheta(4r^2)^2\}
                         =\frac{\log2}{8\pi^2}.
\]
To verify the last equality, integrate from $\epsilon$ to $R$ and set
$t=2r$ in the second term. The remaining lower interval tends to
$\log2$ because $\vartheta=1$ near zero, and the upper interval is zero
for large $R$. Multiplication by $\hbar\operatorname{Tr}_{\mathcal R}(E^2)/4$
proves \eqref{eq:v9-E2-flow}. Thus its leading coefficient is independent
of the chosen smooth cutoff profile.

On the annulus $|p|\asymp\Lambda_j$, $P_a-P_0=O(a^2\Lambda_j^4)$
and the covariance-cutoff change has the same relative
$O((a\Lambda_j)^2)$ size. The integrand has order $\Lambda_j^{-4}$;
integrating its difference proves the first error in
\eqref{eq:v9-E2-lattice}. The continuum annular integrand is smooth and
compactly supported, and its $b$th momentum derivative has $L^1$ norm
$C_b\Lambda_{j-1}^{-b}$. Poisson summation proves the second error.
No estimate for a continuum Ward identity is needed in this calculation.
\end{proof}

The physical quartic-background pullback is also fixed rather than guessed:
\begin{equation}
 \operatorname{Tr}_{\mathcal R}(V''(\Phi))^2
 =16\lambda_H^2\left\{(d_H+8)(\Phi^T\Phi)^2
       -(2d_H+4)v_H^2\Phi^T\Phi+d_Hv_H^4\right\}.
 \label{eq:v9-Higgs-quartic-pullback}
\end{equation}
It follows by squaring \eqref{eq:v9-Higgs-Hessian} and using
$(\Phi\Phi^T)^2=(\Phi^T\Phi)\Phi\Phi^T$. Consequently the scalar-block
contribution to the quartic local counterterm increment, in this real
potential normalization, is
$\hbar(d_H+8)\lambda_H^2\log2\,(\Phi^T\Phi)^2/(2\pi^2)$
at leading high-shell order. A reference-mass compensation affects its
lower polynomial terms, not this quartic coefficient. This is a finite-order
pullback of the displayed Hessian source term; stationary-background,
mixed hard-field and other ESM loop contributions are not thereby supplied.

Uniqueness here is of the generated, normalized \emph{one-scalar-loop
coefficient orbit}. It is not uniqueness of an interacting ESM RG flow.
The theorem realizes a finite local bank and a summable complementary
source sector for this block, but does not show that arbitrary interactions
near it map into an invariant analytic ball. Its local increments are not
identified with the complete ESM beta vector. In particular their logarithmic
growth is retained rather than called a uniformly small analytic coupling.

\section{A concrete internal Ward-compatible subtraction}
\label{sec:v9-Ward}

The source map by exact line transport is important at this interface. For
a smooth commuting test parameter $c(x)\in\rho(\mathfrak g_{\rm SM})$,
\begin{equation}
 \delta_c A_\mu=-\partial_\mu c+[c,A_\mu],\qquad
 \delta_cE=[c,E],\qquad \delta_cJ=cJ,\qquad\delta_ch=0.
 \label{eq:v9-internal-Ward}
\end{equation}
Line transport differentiates to the exact endpoint transformation of the
finite links. Hence \cref{prop:v9-domain} gives the finite measured Ward
identity before subtraction. No lattice pointwise Leibniz rule for the
phase derivative is required. This is an internal symmetry; it says nothing
about the nonexistent naive finite diffeomorphism Leibniz identity of V6.

For a thermodynamic amplitude, define the operator $\mathscr P_{\le4}$
on the full formal source expansion by retaining the Taylor components of
weight
\begin{equation}
 \#A+2\#E+\text{number of external derivatives}\le4.
 \label{eq:v9-Ward-weight}
\end{equation}
This is exactly the family of $\mathsf P_I$ in
\eqref{eq:v9-subtraction}; $J$ has weight three, and its ultraviolet-finite
two-endpoint terms need no subtraction. Give $c$ weight zero.

\begin{theorem}[Intertwining in the internal background complex]
\label[theorem]{thm:v9-Ward}
In the thermodynamic, continuously interpolated source convention,
$\mathscr P_{\le4}$ commutes with the internal Ward operator
\eqref{eq:v9-internal-Ward}. Consequently the thermodynamic finite-lattice
renormalized source coefficients, and their continuum limits from
\cref{thm:v9-continuum}, obey the complete internal Ward identities,
including all metric--current and endomorphism-source contact terms.
The assertion at background degree at most $N-1$ uses the panel through
$N$, since the inhomogeneous $\partial c$ variation can lower background
degree by one. Differentiating with respect to metric marks does not break
the identity. The same commutation holds for the associated internal
Chevalley--Eilenberg source complex with a consistent ghost convention.
For each fixed $N$ the projector also has a mesh- and volume-independent
bound in the scaled finite-panel $C^4$ norm: each amplitude of weight $w$
is measured by the maximum of
$m^{w+|\alpha|-4}\sup_K|\partial_K^\alpha F_I|$, $|\alpha|\le4$,
on a fixed smaller scaled transfer domain, with a sum over its finite
component panels. Its operator norm is at most a constant $C_N$, not a
cutoff-dependent pseudoinverse bound.
\end{theorem}
\begin{proof}
First the unrenormalized thermodynamic coefficients satisfy the Ward
identity at every fixed $a,m$. At finite volume it is the change-of-variable
identity proved above. For fixed $a,m$, all propagators are bounded and
all relevant vertices are smooth periodic functions of the routed scalar
momentum. Riemann sums tend to their thermodynamic integrals. Approximate
any fixed low external panel, including the ghost momentum, by compatible
periodic momenta on successively larger odd tori. Continuity gives the
same coefficient identity for the interpolated thermodynamic amplitudes.
This argument takes the thermodynamic limit at fixed $a$, before any
ultraviolet estimate is used.

Every part of \eqref{eq:v9-internal-Ward} preserves
\eqref{eq:v9-Ward-weight}. Replacing one $A$ by $-\partial c$ removes
one weight-one field and adds one derivative. Replacing it by $[c,A]$
changes no weight; replacing $E$ by $[c,E]$ preserves weight two. The
source variation preserves the weight of $J$. Equivalently, use formal
scaling $k\mapsto tk$, $A\mapsto tA$, $E\mapsto t^2E$, $J\mapsto t^3J$,
$c\mapsto c$. The Ward operation commutes with this scaling, so it
commutes with extraction of the terms of degree at most four in $t$.
This proves the intertwining statement, including the terms that relate
different numbers of connection marks. For the norm assertion, rescale
$K=m\widetilde K$. On the fixed bounded domain, each derivative of the
Taylor polynomial of degree at most four is bounded by a finite sum of
its derivatives at zero, with coefficients depending only on the degree
and the number $4(n-1)$ of variables. Multiplying by $m^{w-4}$ gives
the stated $C_N$ bound. This controls the projection operation, not the
size of the unrenormalized input, whose local coefficients can diverge. Source permutations and total
momentum conservation are retained, so this is a projection of a single
functional, not unrelated projections on its observables.

Applying the projector to the unrenormalized Ward identity proves that both
the local assignment and the remainder obey it. Passage $a\to0$ uses the
fixed derivative and kernel bounds of \cref{thm:v9-continuum}; the Ward
operation involves only finitely many coefficients and one derivative at
any fixed panel. The elementary-source part is separately invariant by
$Q\mapsto OQO^{-1}$ and $J\mapsto OJ$. In the internal
Chevalley--Eilenberg differential the additional ghost-bracket operation
also preserves weight, proving the final assertion without changing a
sign convention from another sector.
\end{proof}

There is a finite-volume qualification. As in V8, a Taylor polynomial of
a particular off-grid momentum interpolation on \emph{one fixed torus}
need not inherit all its on-grid Ward identities. The preceding exact
intertwining conclusion concerns thermodynamic amplitudes and their common
local assignment. The uniform finite-torus cutoff theorem uses its declared
finite-torus interpolation. One cannot silently replace these two
normalization conventions. For a fixed source panel the difference
$\gamma^R_{0,L,I}-\gamma^R_{0,\infty,I}$ tends to zero faster than any
fixed inverse power of $mL_{\min}$: differentiate the subtracted integrand
in the routed $p$ more than four times and apply Poisson summation as in V8.
Explicitly, for each fixed $b>4$,
\begin{equation}
 m^{|\alpha|}|\partial_K^\alpha
  (\gamma^R_{0,L,I}-\gamma^R_{0,\infty,I})|
 \le C_{N,\alpha,b}\hbar d_H\mathfrak n_I
                m^{4-w}(mL_{\min})^{-b}.
 \label{eq:v9-volume}
\end{equation}
All $p$ derivatives are integrable by \eqref{eq:v9-loop-bounds} after
subtraction. Their $L^1$ norms have the extra factor $m^{-b}$, proving the
claim. Thus the two limits and their symmetry statements are explicitly
separated.

Individual smooth covariance shells in \eqref{eq:v9-covariance-shells}
are not gauge covariant under a nonconstant transformation. Their Ward
defect is therefore not set to zero. In the continuum thermodynamic
reference, however, \cref{thm:v9-Ward} and
\eqref{eq:v9-trajectory-tail}, using one additional source degree, prove
\begin{equation}
 \|\mathcal W_c R_j\|_{\rm scaled\,src,N-1}
 \le C_N\hbar d_H\,m^4\|c\|_{\mathrm W,m,1}\,2^{-2j}
 \label{eq:v9-shell-Ward-defect}
\end{equation}
where the background sources are normalized as $h,A/m,E/m^2$ and
$\|c\|_{\mathrm W,m,1}=\|c\|_{\mathrm W}+m^{-1}\|\partial c\|_{\mathrm W}$.
The estimate follows since
$\mathcal W_cR_\infty=0$ and
$\mathcal W_cR_j=-\mathcal W_c(R_\infty-R_j)$. A physical local polynomial
must itself obey the internal identity to preserve it after being added
back. Equation~\eqref{eq:v9-shell-Ward-defect} is a \emph{derived summable
internal source defect}, not a proof of exact invariance of each free
covariance split. Its endpoint identity was derived from the actual finite
measured law, not assumed to prove its own cutoff convergence.

\subsection{A check on why the source map cannot be changed silently}

For one abelian charged component on a flat diagonal background, exact
line sampling gives the one-connection vertex in direction $\mu$
\begin{equation}
 V_{a,\mu}(p,k)=\frac2a\sin(a(p_\mu+k_\mu/2))
                  \operatorname{sinc}(ak_\mu/2),\qquad
 k_\mu V_{a,\mu}\ \text{summed in }\mu=P_a(p+k)-P_a(p).
 \label{eq:v9-abelian-vertex}
\end{equation}
The identity is $2\sin u\sin v=\cos(u-v)-\cos(u+v)$.
At two opposite marks in the same direction the seagull is
$2\cos(ap_\mu)\operatorname{sinc}^2(ak_\mu/2)$. It cancels the contracted
bubble after shifting the complete periodic loop sum. Omitting that
seagull, or replacing the source map by a midpoint value and retaining the
same continuous Ward operator, destroys this check. This example verifies
the conventions; the all-panel theorem does not rely on an abelian
reduction.

\section{ESM interpretation, finite-order consistency, and remaining gates}
\label{sec:v9-scope}

The principal advance is a \emph{whole one-loop source module}, not a
further decimal coefficient. For each finite panel, the actual massive
Higgs Hessian produces a finite local assignment, a generated local
coefficient trajectory, a summable quasilocal complement, elementary and
composite contacts, and internally compatible thermodynamic gauge identities.
Increasing $N$ leaves every previously specified lower coefficient and
its Taylor assignment unchanged. The banks grow by inclusion because
metric degree is not bounded by ultraviolet weight. There is no single
finite-dimensional metric bank closing all source orders.

Finite polynomial substitution of background Higgs fields through
\eqref{eq:v9-Higgs-Hessian}, or of a fixed polynomial composite source
panel, consumes finitely many of these coefficients at a prescribed
perturbative/EFT order, provided every substituted nonconstant or
mass-compensation term carries its recorded degree. The chain rule must
be applied to the complete local functional and the complete remainder,
not only to a renormalized action evaluated at zero sources. This gives the
metric dependence of the Higgs current/composite sector in the declared
coordinate normalization. It does not identify an uncomputed fermion or
gravitational determinant with this block.

For the full ESM endpoint, several distinct questions remain. The metric
counterterms here have not been assembled into the nonlinear gravitational
BV complex or matched to one coframe--connection measure. The native
quartic residual \eqref{eq:v8-explicit-defect}, the pure-gravity cubic,
local Lorentz factors and higher ghost transformations remain the first
unclosed nonlinear gravitational interface. Massless/reference-mass
transport, chiral-line coherence, and the mixed hard-Hessian Schur terms
are not consequences of positivity of a scalar Hessian. Nor does the
one-loop generated orbit above verify V6's weighted R1--R4 hypotheses for
arbitrary interacting ESM inputs.

A useful next calculation can now consume, rather than repeat, this
module: couple its all-source scalar block to the \emph{actual} native
coframe/connection quadratic and cubic data, form the mixed hard-Hessian
Schur return and complete gravitational symmetry defect, and test their
local restoring jets with the same source grading. A separate remaining
choice is to construct a common nonlinear gravitational restoring
projection; the internal weight argument cannot supply it because the
finite gravitational differential is not yet available.

\begin{center}\small
\begin{tabular}{>{\raggedright\arraybackslash}p{.23\textwidth}>{\raggedright\arraybackslash}p{.68\textwidth}}
\toprule
Variable & Scope proved in V9\\
\midrule
Volume & Rooted kernel and cutoff bounds are uniform for $mL_\mu\ge1$.
Unsubtracted determinants are extensive. Exact renormalized internal Ward
statements use thermodynamic amplitudes; the finite-volume normalization
qualification is retained.\\
Mesh & $Nam\le1/64$, with fixed $N$. Every scalar ultraviolet mode is
included. The comparison rate is $(am)^2[1+\log(1/(am))]$ in scaled norms.\\
Loop and EFT order & One scalar loop, arbitrary prescribed finite
background/composite-source degree. Local weight at most four suffices for
this Hessian block, but metric degree grows with $N$. This is not all-loop
ESM power counting.\\
Sources & Metric, represented electroweak connection, symmetric endomorphism,
and elementary scalar sources in the declared low Fourier windows. All
required finite contact vertices are included. No bound uniform over
unbounded source order is proved.\\
Shell count & Complementary scalar-loop increments are summable uniformly
in the number of smooth scalar shells. Local marginal coefficients can
grow linearly with the logarithmic scale. Arbitrary interacting shell
self-mapping is not proved.\\
Mass and backgrounds & $m>0$ fixed; common small real/complex background
domain for the scalar Gaussian. Physical reference-mass compensation is
only a counted insertion unless a separate resummation is supplied.\\
Symmetry and measure & Exact finite internal covariance before subtraction,
and an intertwining thermodynamic internal source projector. Flat scalar
coefficient measure. Full diffeomorphism/local Lorentz BV, chiral measure,
and nonlinear gravitational contours remain separate.\\
\bottomrule
\end{tabular}
\end{center}

The standalone script \srcid{check_ESM_source_module_V9.py} separates exact
finite-algebra checks from numerical diagnostics. It checks ordered
partitions through four marks against a direct determinant coefficient,
noncommuting three-cycle shell ownership, the line-source abelian Ward
vertex and its seagull, the native composite marginal and Higgs-Hessian
pullback, and the weight/subtraction table. It also tests
finite gauge covariance of the actual covariant Hodge matrix with a
nontrivial metric and matrix source, the complete Gaussian source functional,
and full-zone triple/fourfold potential-insertion integrals. Numerical
quadrature, convergence ratios, and floating-point matrix comparisons are
not proofs of the uniform estimates above.

\begingroup
\renewcommand{\refname}{Additional methodological references for V9}
\begin{thebibliography}{99}
\bibitem[30]{V9Vassilevich}
D.~V.~Vassilevich, \emph{Heat kernel expansion: user's manual}, Physics
Reports \textbf{388} (2003), 279--360, arXiv:hep-th/0306138,
doi:10.1016/j.physrep.2003.09.002. Context for one-loop geometric and
background-field coefficients, not an imported native shell estimate.
\bibitem[31]{V9HollandsWald}
S.~Hollands and R.~M.~Wald, \emph{Existence of local covariant time ordered
products of quantum fields in curved spacetime}, Communications in
Mathematical Physics \textbf{231} (2002), 309--345, arXiv:gr-qc/0111108,
doi:10.1007/s00220-002-0719-y. Established curved-spacetime
composite-field framework; no claim of first existence is made here.
\end{thebibliography}
\endgroup

\section{Final ledgers for V9}
\label{sec:v9-ledgers}

\subsection*{Proved}
\addcontentsline{toc}{subsection}{V9: Proved}
The actual covariant Hodge family \eqref{eq:v9-native-Q} has the stated
finite measured covariance and common massive Gaussian domain. Its full
labelled vertices have the verified degree-dependent whole-zone estimates
\eqref{eq:v9-vertex-bounds}, including higher finite-link contact vertices.
The one-loop panel with arbitrary prescribed finite numbers of metric,
connection and quadratic-composite insertions has the cutoff and rooted
source bounds \eqref{eq:v9-main} and \eqref{eq:v9-kernel}, after assignment
to the explicit finite local bank. Its measured covariance shells generate
\eqref{eq:v9-generated-flow} with a summable source tail; no diagonal-shell
replacement is made. Their composite $E^2$ marginal is evaluated in
\eqref{eq:v9-E2-flow}, with native cutoff errors and the explicit Higgs
quartic-background pullback. The thermodynamic internal background Ward projector
intertwines by its homogeneous source weight, and its continuum and
partial-shell consequences have the stated qualification and error bounds.
All elementary and matrix-composite contacts are retained. These are results
for the declared massive Higgs Hessian block, not full interacting ESM.

\subsection*{Inherited}
\addcontentsline{toc}{subsection}{V9: Inherited}
V1--V8 and their labels are unchanged. The common finite Higgs carrier,
internal covariance and quartic potential come from the structural common
action. V7 provides the sign-averaged seed; V8 provides the specific repaired
metric vertex and its two-metric comparison. General measured Gaussian
composition, source contacts, Schur ownership, degree-dependent local banks
and abstract continuum/splitting theorems are inherited at their original
scope. None of those principles is renamed a new ESM theorem here.

\subsection*{Literature input}
\addcontentsline{toc}{subsection}{V9: Literature input}
Massive lattice finite-part renormalization, one-loop background-field and
heat-kernel methods, and curved-spacetime composite-field renormalization
are established subjects. They provide context and comparison; the
all-panel estimates above are obtained directly from the specified finite
Hodge and line-source vertices. A generic EFT existence claim is not the
novelty of this installment.

\subsection*{Conditional}
\addcontentsline{toc}{subsection}{V9: Conditional}
Consumption by a coupled ESM shell requires the native gravitational and
fermionic measures, actual mixed hard-Hessian blocks, common retained source
normalization and nonlinear restoring terms. Replacing the Hessian source
by a physical Higgs background requires correct perturbative grading,
stationary-section bookkeeping and any needed mass-compensation control.
The internal homogeneous projector is not a bounded projection of the
unconstructed full diffeomorphism/Lorentz/SM BV complex. Exact finite-torus
renormalized Ward identities require a compatible local normalization
convention beyond an arbitrary off-grid Taylor interpolation.

\subsection*{Open}
\addcontentsline{toc}{subsection}{V9: Open}
The next load-bearing gravitational calculation is the native coupled
coframe/connection cubic and mixed-Hessian Schur contribution, with its
complete quartic symmetry defect and source-dependent measure. The scalar
module now supplies all its fixed-order metric/current/composite derivatives;
recomputing another isolated scalar polarization is not required. Uniform
nonlinear symmetry restoration, massless and chiral sectors, and a full
invariant filtered interacting ESM state remain open. No infinite-series
convergence, nonperturbative positive gravity measure, asymptotic safety,
finite-parameter gravitational ultraviolet completion or mass-gap conclusion
is drawn.


\clearpage
\part*{V10: Native connection elimination, Lorentz measure, and a local quantum handoff}
\addcontentsline{toc}{part}{V10: Native connection elimination and measured handoff}

\section{Returning to the actual gravitational action}
\label{sec:v10-start}

Sections 1--71 are preserved. The scalar module of V9 is not recomputed.
The first missing gravitational ingredient is upstream of a graviton loop:
the independent Lorentz connection must be eliminated with its actual
nonquadratic link terms and its measure. The present installment computes
its cubic and quartic stationary action and its one-loop logarithmic
amplitude. It also proves that, at every fixed number of coframe and
connection-source insertions about the flat reference, this auxiliary
one-loop contribution has only a local continuum matching part. This does
not assert the same statement for propagating metric, ghost, gauge or
fermion loops.

\subsection{A newly located input, not a new classical inverse theorem}

The Library manuscript \cite{V10GT},
\srcid{sec:GT-shared-link-Cartan}, specifies the completed shared-link
Palatini action that underlies the phase-compatible flat constraints used
in V6. Its exact labels are
\srcid{eq:GT-compensated-shared-link-action},
\srcid{eq:GT-shared-link-normal-form},
\srcid{thm:GT-shared-link-Cartan-recovery}, and
\srcid{eq:GT-phase-compatible-action}.
The construction, its bounded nonlinear Cartan inverse and its exact
quadratic metric identities are inherited. This is a later source than the
attached V073 monograph. It is used explicitly, not silently substituted
for that earlier manuscript.

The same source has a different cellwise polynomial-connection model,
\srcid{thm:GT-finite-Palatini-elimination}. We do not identify that model
with the shared-link action. In particular, the latter contains a nonzero
cubic remainder. The ordinary Gaussian square, Schur identity, stationary
coefficient rule and source differentiation are inherited tools. The new
calculation below identifies their \emph{native operands}, their local
Lorentz/Haar density, and the fixed-order cutoff implication for this
specific auxiliary block. First-order quantum gravity and the dependence
of its metric measure on elimination have established literature
\cite{V10Aros,V10Benedetti}; no general first-order/second-order equivalence
or first construction of perturbative gravity is claimed.

\subsection{Regulator, action and the extra quantum measure choice}

Write $a$ for the spatial mesh, $\eta=\operatorname{diag}(-1,1,1,1)$,
and $g=e^{\mathsf T}\eta e$. Work in one oriented, time-oriented compact
coframe chart with $\det e>0$. A fixed Palatini coefficient $\chi\ne0$ is
retained; the independent Holst coefficient is set to zero in this
calculation. Spatial links are $U_i=\exp(aA_i)$ and $A_0$ is independent.
The invariant Lorentz-algebra pairing is denoted by $b$. The coframe
coefficients $\Pi_i(e),\Sigma_{ij}(e)$ are the ones obtained by inserting
basis curvatures in the source's Palatini contraction; they are quadratic
polynomials in $e$.

After the \emph{exact} kinetic compensation and temporal boundary pairing
of the source, no derivative of $A$ remains in the action. It is therefore
legitimate to take a specified finite temporal quadrature of that normal
form. Its positive weights are $\tau_t$ and $v_{x,t}=a^3\tau_t$.
The temporal load uses the derivative of the declared coframe interpolant.
This is a finite quadrature realization of the same completed normal form,
not an assertion that a different uncompleted time-link action has been
integrated. For the cutoff result in \cref{sec:v10-cutoff}, take periodic
quadrature with $\tau_t=a$ and the trigonometric derivative of the time
record. All finite source panels there are nonaliased. Other temporal
normalizations must keep their actual cell weights.

On a fixed time row the source normal form, with an external Cartan source
$j$, is
\begin{equation}
 S_a(e,A;j)=\sum_x v_x\left\{V(e_x)
 +(f_a^\circ(e)_x-j_x)\cdot A_x
 +\tfrac12 A_x^{\mathsf T}\mathsf C(e_x)A_x\right\}
 +\sum_xv_x r_{a,x}(e,A).
 \label{eq:v10-action}
\end{equation}
There are 24 real connection coordinates per vertex. $V$ is the recorded
cosmological term; $V=0$ in the flat cubic calculation. With the invariant
pairing matrix $\mathsf H$, the load is exactly
\[
 f_{a,0}^\circ=\mathsf H\sum_i\delta_i\Pi_i,\qquad
 f_{a,i}^\circ=\mathsf H\left(-\partial_t\Pi_i
                         -\sum_{j\ne i}\delta_j\Sigma_{ji}\right),
 \qquad \widehat{\delta_i u}(p)=\frac{2i}{a}\sin(ap_i/2)\widehat u(p).
\]
The derivative is dense on the full collocated grid. Nothing below changes
it into a local difference or assumes its Leibniz rule.

For clarity, the full remainder in \eqref{eq:v10-action} is
\begin{equation}
\begin{split}
 r_{a,x}={}&-\sum_i b\left(\Pi_i,
 a^{-1}(e^{a\operatorname{ad}A_i}-I-a\operatorname{ad}A_i)A_{0,+i}\right)\\
 &+\sum_{i<j} b\left(\Sigma_{ij},a^{-2}
 \left[\log\prod_{r=1}^{4}e^{aX_r}
       -a\sum_rX_r-\frac{a^2}{2}\sum_{r<s}[X_r,X_s]\right]\right),
 \label{eq:v10-full-remainder}
\end{split}
\end{equation}
where $(X_1,X_2,X_3,X_4)=(A_i,A_{j,+i},-A_{i,+j},-A_j)$ and the
logarithm has its fixed small chart. All translated occurrences are retained.
Formula \eqref{eq:v10-full-remainder} follows directly from the two
remainders displayed in the source proof; it is not a replacement by its
cubic Taylor term.

The classical action does not determine a quantum reference measure.
For this installment we \emph{declare} product Lebesgue measure on the 16
coframe entries and on $A_0$, product Lorentz Haar measure on $U_i$, and
local quotient by the six Lorentz frame coordinates. Only its local
perturbative germ is used. No normalized noncompact Haar probability,
positive real Euclidean Palatini integral, global gauge slice, or contour
for the complete ESM integral is asserted. A different declared measure
changes the local matching terms derived below.

\section{The missing native cubic and quartic are finite words}
\label{sec:v10-cubic}

For noncommuting matrices $X_1,\ldots,X_4$, put
\[
 E_n(X)=\sum_{k_1+\cdots+k_4=n}
       \frac{X_1^{k_1}\cdots X_4^{k_4}}{k_1!\cdots k_4!},\qquad
 B_n(X)=\sum_{r=1}^{n}\frac{(-1)^{r+1}}r
       \sum_{\substack{n_1+\cdots+n_r=n\\n_i\ge1}}E_{n_1}\cdots E_{n_r}.
\]
Thus $B_n=[z^n]\log\prod_r e^{zX_r}$ is specified by a finite rational
matrix polynomial. In particular
\begin{equation}
\begin{split}
 B_3={}&\frac1{12}\sum_{r<s}
  \bigl([X_r,[X_r,X_s]]+[X_s,[X_s,X_r]]\bigr)\\
 &+\frac16\sum_{r<s<t}
  \bigl([X_r,[X_s,X_t]]+[X_t,[X_s,X_r]]\bigr).
 \label{eq:v10-BCH3}
\end{split}
\end{equation}
This is used as an exact coefficient identity, not as a new BCH principle.

\begin{proposition}[Evaluated native connection jets]
\label[proposition]{prop:v10-native-jets}
The exact remainder has, at $A=0$,
\begin{equation}
 r_a(e,A)=a\mathscr R_3(e,A)+a^2\mathscr R_4(e,A)+O(a^3|A|^5),
 \label{eq:v10-r34}
\end{equation}
where the degree-$n$ local density, $n=3,4$, is
\begin{equation}
 \mathscr R_{n,x}=
 -\sum_i\frac1{(n-1)!}
 b\bigl(\Pi_i,(\operatorname{ad}A_i)^{n-1}A_{0,+i}\bigr)
 +\sum_{i<j}b(\Sigma_{ij},B_n(X)).
 \label{eq:v10-Rn}
\end{equation}
$\mathscr R_n$ includes the same finite set of neighboring link occurrences
as the original plaquette. On a fixed coframe chart and $\|A\|_\infty\le R$,
$\|D_A^2r_a\|$ has row and column bounds $C aR$; all fixed further local
jets have bounds obtained by differentiating the convergent series.
The displayed cubic is not identically zero.
\end{proposition}
\begin{proof}
Expand the exponential in the first line of
\eqref{eq:v10-full-remainder}; its degree $n$ term has coefficient
$a^{n-2}/(n-1)!$. Expand the product exponential and the ordinary
$\log(I+Y)$ power series in the second line. The subtracted terms remove
exactly degrees one and two. This gives the finite formula for $B_n$ and
\eqref{eq:v10-Rn}. The series converge uniformly with all prescribed finite
jets on $aR<c$ because only finitely many matrix entries and neighbors occur.
Differentiating twice gives the stated bound, which is the local mechanism
behind the inherited Cartan inverse.

For a nonzero witness use constant $A_0=A_3=0$, $A_1=J_{12}$,
$A_2=J_{13}$, where $(J_{ij})_{ij}=1$ and $(J_{ij})_{ji}=-1$ are spatial
rotation matrices. The commutator plaquette gives
$B_3=\tfrac12[A_1+A_2,[A_1,A_2]]
=\tfrac12(J_{12}-J_{13})$.
In the flat normalization $b(\Sigma_{12},Z)=2\chi Z^{12}$ its contraction
is $\chi$, while all other plaquettes and the electric remainder vanish.
Thus $\mathscr R_3=\chi$ per site. Translated variables coincide in this
witness only because the chosen record is constant.
\end{proof}

\subsection{Stationary elimination with the cubic retained}

Let $e=e_0+\varepsilon u$, $j=\varepsilon j_1+\varepsilon^2j_2+
\varepsilon^3j_3$, and $e_0=I$, with $V=0$. Write coefficients, not
Taylor derivatives, as
\[
 f_a^\circ(e)-j=\varepsilon b_1+\varepsilon^2 b_2+
 \varepsilon^3 b_3+\cdots,\qquad
 \mathsf C(e)=C_0+\varepsilon C_1+\varepsilon^2 C_2.
\]
In this additive coframe chart $f$ is quadratic, so $b_3=-j_3$.
A nonlinear metric/coframe chart has its actual higher $b_r$ instead.
Expand $\mathscr R_3(e,A)=\mathscr R_3^{(0)}(A)+
\varepsilon\mathscr R_3^{(1)}(u;A)+\cdots$. Pairings in what follows
include cell weights and every translated occurrence; functional gradients
use the same pairing.

\begin{theorem}[Native stationary gravitational action through degree four]
\label[theorem]{thm:v10-reduced-jets}
Put
\begin{align}
 A_1&=-C_0^{-1}b_1,\\
 G_2&=b_2+C_1A_1+aD_A\mathscr R_3^{(0)}(A_1),\\
 A_2&=-C_0^{-1}G_2.
 \label{eq:v10-stationary-jets}
\end{align}
The unique local stationary section has $A_*=\varepsilon A_1+
\varepsilon^2 A_2+O(\varepsilon^3)$. The coefficients of
$S_a^*(e,j)=S_a(e,A_*;j)$ are
\begin{align}
 S_2^*&=-\tfrac12\langle b_1,C_0^{-1}b_1\rangle,\\
 S_3^*&=\langle b_2,A_1\rangle+
       \tfrac12\langle A_1,C_1A_1\rangle
       +a\mathscr R_3^{(0)}(A_1),\\
 S_4^*&=\langle b_3,A_1\rangle+
       \tfrac12\langle A_1,C_2A_1\rangle
       +a\mathscr R_3^{(1)}(u;A_1)\notag\\
 &\quad+a^2\mathscr R_4^{(0)}(A_1)
       -\tfrac12\langle G_2,C_0^{-1}G_2\rangle.
 \label{eq:v10-reduced-jets}
\end{align}
The $a\mathscr R_3$ term and its contribution inside $G_2$ belong to the
actual action. Deleting them changes both its cubic and its quartic
coefficient. Every further fixed degree is determined by the finite word
formula for $B_n$ and the same stationary recursion.
\end{theorem}
\begin{proof}
The degree-one and degree-two equations of
$f-j+\mathsf C A+\partial_A r_a=0$ give
\eqref{eq:v10-stationary-jets}. The derivative at zero is the invertible
block $C_0$, so the finite analytic implicit function theorem identifies
these coefficients with the inherited stationary branch. Substitute into
\eqref{eq:v10-action}. Terms containing $A_2$ at degree three cancel by
$b_1+C_0A_1=0$. At degree four the terms containing $A_3$ cancel for the
same reason, while the remaining $A_2$ terms are
$\langle G_2,A_2\rangle+\langle A_2,C_0A_2\rangle/2
=-\langle G_2,C_0^{-1}G_2\rangle/2$.
The other terms are precisely those in \eqref{eq:v10-reduced-jets}.
This proof works with an indefinite $C_0$ and uses no minimum principle.
At any fixed higher degree the unknown next stationary coefficient occurs
linearly through $C_0$ and all other terms use earlier coefficients, proving
the final assertion. This is a use of V1's compiler with newly evaluated
native vertices, not another abstract compiler theorem.
\end{proof}

At $u=0$, choose $j_1=C_0\mathcal A$ with the constant rotation record
$\mathcal A$ in \cref{prop:v10-native-jets}. Then $A_1=\mathcal A$ and
$S_3^*=a\chi\sum_xv_x$. This is an explicit cubic \emph{Cartan-source
contact} even though no coframe fluctuation is retained. It is not a
source-independent normalization. On fixed smooth, low-momentum coframe
panels the corresponding extra cubic is of order $a\Lambda^3$; the
pointwise Cartan part has two derivatives. No smallness on cutoff-scale
coframes follows from that smooth-panel observation.

\section{The local Lorentz quotient and the spatial-link Haar density}
\label{sec:v10-measure}

A common local Lorentz frame is quotiented through the coframe, not by
regarding all connection components as gauge zero modes. The source's
coframe-adapted representative writes $e=L\widehat e(g)$. Once this is
fixed, the 24-dimensional Cartan block is nondegenerate and is to be
integrated, not canceled by a fictitious Lorentz ghost determinant.

\begin{proposition}[Coframe quotient density in the declared measure]
\label[proposition]{prop:v10-Lorentz-density}
On one oriented Lorentzian metric chart, product Lebesgue coframe measure
factorizes locally as
\begin{equation}
 d^{16}e=c\,|\det g|^{-1/2}\,d^{10}g\,d\mu_{\rm Lor}(L).
 \label{eq:v10-coframe-density}
\end{equation}
Here $d\mu_{\rm Lor}$ is a fixed Haar convention and $c$ is independent of
$g$. With independent symmetric entries for $g$ and the standard six
infinitesimal Lorentz generators, its tangent normalization at $e=I$ is
$c=2^{-4}$ when Haar has unit density at the identity. The local gauge
condition $\log L=0$ has Faddeev--Popov determinant one in these Haar and
gauge coordinates. Thus quotienting this six-dimensional frame factor
leaves the density in \eqref{eq:v10-coframe-density}; it does not remove it.
\end{proposition}
\begin{proof}
The based local action $e\mapsto Le$ on an invertible coframe is free.
Local section coordinates therefore exist. Under right multiplication
$e\mapsto eB$, Lebesgue measure has Jacobian $|\det B|^4$. The induced
map on symmetric metrics is $g\mapsto B^{\mathsf T}gB$, whose determinant
on $\operatorname{Sym}_4$ is $|\det B|^5$: for a diagonal $B$ this is
obtained by multiplying its four diagonal and six off-diagonal weights,
and the general identity follows by polynomial continuation. Changing the
section can only right-translate the Lorentz factor; Haar is unimodular.
Consequently its remaining density obeys
$J(B^{\mathsf T}gB)|\det B|^5=J(g)|\det B|^4$.
Transitivity on the fixed-signature chart gives $J(g)=c|\det g|^{-1/2}$.
At $e=I$, the four diagonal metric variations contribute factors $1/2$;
each metric off-diagonal/Lorentz pair has absolute two-by-two determinant
one. This gives $c=1/16$. Finally a left Lorentz variation at $L=I$ changes
$\log L$ by its own Lie-algebra parameter, so its six-by-six gauge Jacobian
is identity. This calculation is local and does not divide a convergent
integral by an infinite group volume.
\end{proof}

Connection dressing is triangular in $(e,A)$. For a fixed $e$, $A_0$ is
changed by an adjoint and a translation, and the spatial links by endpoint
left/right multiplication. $\det\operatorname{Ad}L=1$ for the connected
Lorentz group; product Lebesgue/Haar measure is therefore unchanged in
these variables. The time derivative of the frame may depend on other
coframe records, but this lies in the off-diagonal part of the full
triangular Jacobian. It creates no extra determinant in the connection
block. This statement uses the actual dressed records, not a discretized
continuum product rule.

\begin{lemma}[Haar density in the actual logarithmic links]
\label[lemma]{lem:v10-Haar}
Up to the source-independent factor $a^6$, Lorentz Haar measure pulled back
by $U=\exp(aA)$ has analytic density
\begin{equation}
 j(aA)=\det_{\mathfrak{so}(1,3)}
       \frac{1-e^{-a\operatorname{ad}A}}{a\operatorname{ad}A}.
 \label{eq:v10-Haar}
\end{equation}
It is positive near zero, and
\begin{equation}
 \log j(aA)=\frac{a^2}{24}\Tr_{\rm ad}(\operatorname{ad}A)^2+O(a^4|A|^4).
 \label{eq:v10-Haar-jet}
\end{equation}
The apparent quotient in \eqref{eq:v10-Haar} denotes its entire power
series, also on the kernel of $\operatorname{ad}A$.
\end{lemma}
\begin{proof}
The left-trivialized differential of the exponential is
$\int_0^1e^{-sa\operatorname{ad}A}\,ds$; its determinant is the pulled-back
left-invariant volume. At zero it is identity. Since
$\Tr_{\rm ad}\operatorname{ad}A=0$, the determinant of the exponential
prefactor is one. Thus its logarithm is the trace of
$\log(\sinh(z/2)/(z/2))=z^2/24+O(z^4)$ at
$z=a\operatorname{ad}A$. Positivity and the remainder follow on the fixed
small chart by analytic continuation from zero.
\end{proof}

Relative to flat metric coordinates and logarithmic connection coordinates,
the complete declared local density is therefore
\begin{equation}
 \rho_a(g,A)=\prod_x|\det g_x|^{-1/2}
                     \prod_{x,i=1}^3j(aA_{x,i}),
 \label{eq:v10-full-density}
\end{equation}
with the identical convention at every time row. This is additional quantum
input, now explicit. It is not recovered from the classical action alone.

\section{A source-local inverse and connected determinant for the actual auxiliary block}
\label{sec:v10-local}

Write $b=j-f_a^\circ(e)$, initially as an independent load, and let
$A_*(e,b)$ solve $\mathsf C(e)A+\partial_A r_a(e,A)=b$ in the inherited
bounded chart. Set
\begin{equation}
 H_a=\mathsf C(e)+R_a,\qquad
 R_a=D_A^2r_a(e,A_*),\qquad K_a=\mathsf C(e)^{-1}R_a.
 \label{eq:v10-H}
\end{equation}
The diagonal matrix $\mathsf C$ is understood on all sites. At one time
row give the spatial lattice its nearest-neighbor graph distance $d$.
There are no cross-time connection entries. For a block kernel define
\[
 \|M\|_\mu=\max\left\{
 \sup_x\sum_y e^{\mu d(x,y)}\|M_{xy}\|,
 \sup_y\sum_x e^{\mu d(x,y)}\|M_{xy}\|\right\}.
\]
This is a submultiplicative norm. All components have fixed size 24.

\begin{theorem}[Uniform auxiliary inverse and rooted loop density]
\label[theorem]{thm:v10-local-inverse}
Fix a coframe compact chart, a connection bound $R$ and $\mu>0$.
There are $c_*,c_R>0$, independent of the number of sites and time rows,
such that for sufficiently small $aR$ the actual stationary Hessian obeys
\[
 \|\mathsf C^{-1}\|_\mu\le c_*^{-1},\qquad
 \|K_a\|_\mu\le\epsilon_a:=c_R aR<1/2,
\]
and hence
\begin{equation}
 \|H_a^{-1}\|_\mu\le\frac{c_*^{-1}}{1-\epsilon_a}.
 \label{eq:v10-inverse-bound}
\end{equation}
Its inertia is the one of $\mathsf C$, namely 12 positive and 12 negative
directions per site. The real relative log determinant has the absolutely
convergent density
\begin{equation}
 \frac12\log\frac{\det H_a}{\det\mathsf C}
 =\sum_x\mathcal L_{a,x},\qquad
 \mathcal L_{a,x}=\frac12\sum_{n\ge1}\frac{(-1)^{n+1}}n
                         \operatorname{tr}_{24}(K_a^n)_{xx},
 \label{eq:v10-loop-density}
\end{equation}
with
\begin{equation}
 |\mathcal L_{a,x}|\le12[-\log(1-\epsilon_a)].
 \label{eq:v10-loop-bound}
\end{equation}
The unrooted sum may be extensive; the bound on each root is not.
\end{theorem}
\begin{proof}
Each primitive remainder Hessian has support in one fixed plaquette
neighborhood of diameter $\ell$ and unweighted row/column bounds $CaR$ by
\eqref{eq:v10-full-remainder}. Multiplication by the diagonal inverse and
by $e^{\mu\ell}$ gives the displayed weighted bound. The Neumann series
$H_a^{-1}=\sum_{n\ge0}(-K_a)^n\mathsf C^{-1}$ proves
\eqref{eq:v10-inverse-bound}. Along
$\mathsf C+tR_a$, $0\le t\le1$, the same argument excludes zero
eigenvalues. Real symmetry therefore preserves inertia. Its determinant
ratio remains positive, so the logarithm continued from $t=0$ is real.
The absolutely convergent matrix logarithm is
$\sum_{n\ge1}(-1)^{n+1}K_a^n/n$. Taking its diagonal trace gives
\eqref{eq:v10-loop-density}; $|\operatorname{tr}_{24}M|\le24\|M\|$
and the weighted norm bound give \eqref{eq:v10-loop-bound}.
No positive-definite Cartan covariance was assumed.
\end{proof}

\begin{proposition}[Fixed-order source derivatives without an extensive loss]
\label[proposition]{prop:v10-marked}
Use the independent site coordinates $\theta=(e,b)$ on a fixed smaller
bounded chart. For every fixed $r$ the stationary map and the rooted
densities above have summable differentiated kernels. If $\tau$ is the
minimum graph-tree length connecting its arguments, then
\begin{align}
 \sup_x\sum_{y_1,\ldots,y_r}e^{\mu'\tau(x,y_1,\ldots,y_r)}
  \|D_{\theta_{y_1}}\cdots D_{\theta_{y_r}}A_{*,x}\|&\le C_r,\\
 \sup_x\sum_{y_1,\ldots,y_r}e^{\mu'\tau(x,y_1,\ldots,y_r)}
  |D_{\theta_{y_1}}\cdots D_{\theta_{y_r}}\mathcal L_{a,x}|&\le C_r a,
 \label{eq:v10-marked}
\end{align}
for any fixed $0<\mu'<\mu$, after choosing $aR$ below a fixed margin.
The corresponding Haar log density has bound $C_r a^2$. Constants depend
on the fixed chart, $R,r,\mu,\mu'$, but not on volume, the number of time
rows or the refinement index.
\end{proposition}
\begin{proof}
At first order, differentiating the stationary equation gives $H_a^{-1}$
applied either to a local unit load or to a finite-support derivative of
$\mathsf C A_*+\partial_A r_a$. Both have bounded weighted rows.
At order $r$, isolate $H_aD^rA_*$; every other term is a finite sum,
indexed by partitions of $r$ marks, of primitive derivatives applied to
stationary derivatives of lower order. Each primitive has fixed support
and uniform jets. The union of the inverse-kernel path and the rooted
trees of these factors connects all marks. Its length bounds $\tau$.
The weighted convolution inequality and induction therefore prove the
first estimate, with constants depending on the finite partition number,
not on sites. This also proves analytic differentiability on the smaller
chart, or alternatively follows by differentiating the uniformly
contracting fixed-point equation.

For the determinant, first differentiate a finite partial sum of
\eqref{eq:v10-loop-density}. Every differentiated $K_a$ has a factor $a$
and bounded marked kernels by the preceding induction and the explicit
remainder. The number of placements of $r$ marks on an $n$-cycle is
bounded by $C_r n^r$. The sums of the undifferentiated factors are controlled
by a fixed geometric margin $\epsilon_a<1/2$; terms with one differentiated
factor, including higher derivatives on that factor, still retain $a$.
The series $\sum_n n^r\epsilon_a^{\max(n-r,0)}$ times its finite
placement factors is convergent uniformly after the fixed margin is chosen.
A closed walk plus its derivative trees connects the root and every mark.
This proves the second estimate and permits differentiation of the limit.
The harmless loss from $\mu$ to $\mu'$ also absorbs path-length factors.
The analytic series for $\log j(aA_*)$ starts with $a^2A_*^2$; the first
estimate gives its last asserted bound. Repeated marks at one site are
included; no contact derivative is excluded.
\end{proof}

These are genuinely source-local estimates for the auxiliary block, but
\emph{in its declared load coordinates}. Substituting the native
$b=j-f_a^\circ(e)$ differentiates the dense phase derivative as well.
The preceding exponential graph norm is not then an automatic metric-source
norm. The finite-panel comparison below pays for that substitution
explicitly; no all-frequency nonlinear diffeomorphism estimate is hidden
in \cref{prop:v10-marked}.

\section{The complete one-loop connection amplitude and its measure contacts}
\label{sec:v10-amplitude}

The inherited Cartan congruence is
\begin{equation}
 \mathsf C(e)=\chi\det(e)\,\mathsf T_e^{\mathsf T}H_0\mathsf T_e,
 \qquad (\mathsf T_e A)_a=(e^{-1})^\mu{}_a A_\mu.
 \label{eq:v10-C-congruence}
\end{equation}
The spectrum of $H_0$ is the inherited
$\{-4\ (4),-2\ (8),2\ (8),4\ (4)\}$.
Its determinant is $2^{32}$, not zero.

\begin{lemma}[Exact coframe power of the Cartan determinant]
\label[lemma]{lem:v10-C-determinant}
In the same positive coordinate convention,
\begin{equation}
 \det\mathsf C(e)=\chi^{24}2^{32}(\det e)^{12}.
 \label{eq:v10-C-det}
\end{equation}
Consequently its half-logarithmic variation relative to $e_0=I$ is
$6\log\det e=3\log|\det g|$.
\end{lemma}
\begin{proof}
$\mathsf T_e$ acts by the inverse coframe on the four one-form slots and
by identity on each of the six Lorentz pairs. Thus
$\det\mathsf T_e=(\det e)^{-6}$.
Taking the determinant of \eqref{eq:v10-C-congruence} gives
\[
 \chi^{24}(\det e)^{24}(\det e)^{-12}\det H_0.
\]
The stated spectrum gives $\det H_0=2^{32}$ and
$|\det g|=(\det e)^2$.
\end{proof}

\begin{theorem}[Measured native one-loop auxiliary amplitude]
\label[theorem]{thm:v10-amplitude}
Fix any finite quadrature regulator and a smaller source chart as above.
A smooth local amplitude cutoff equal to one near $A_*$ and supported
inside its unique stationary chart defines a finite oscillatory connection
integral. Its stationary-phase logarithmic \emph{amplitude} coefficient,
normalized at $e_0=I$, $j_0=0$, is
\begin{equation}
\boxed{\begin{aligned}
 \mathcal M_{a,1}(g,j)={}&\frac72\sum_x\log|\det g_x|
     +\frac12\Tr\log(I+K_a)\\
 &-\sum_{x,i=1}^3\log j(aA_{*,x,i}),\qquad
 b=j-f_a^\circ(\widehat e(g)).
\end{aligned}}
 \label{eq:v10-main-amplitude}
\end{equation}
The same formula is the one-loop coefficient in the formal Laplace/Wick
convention with its continuously chosen square-root branch. In the
Lorentzian convention $-i\hbar\log Z$, it is multiplied by $i\hbar$,
not by $\hbar$; in the formal Laplace convention it is multiplied by
$\hbar$. No Wick-rotation theorem between complete measures is inferred.

For constant coframe and zero Cartan source, $A_*=0$, and the actual
quotiented leading (one-loop) stationary amplitude density relative to
$dg$ is, up to its constant normalization,
\begin{equation}
 \boxed{\rho_{\rm eff}(g)=|\det g|^{-7/2}.}
 \label{eq:v10-constant-density}
\end{equation}
Along spatially and temporally constant metric variations, its first and
second geometric contacts are
\begin{align}
 D\mathcal M_{a,1}[H]&=\tfrac72\sum_x\operatorname{tr}(g^{-1}H),\\
 D^2\mathcal M_{a,1}[H,G]&=-\tfrac72\sum_x
                           \operatorname{tr}(g^{-1}Hg^{-1}G).
 \label{eq:v10-density-contacts}
\end{align}
For general nonconstant variations the displayed terms are the local
measure contribution, and the differentiated loop and Haar terms must be
added. They are not removable source-independent vacuum constants.
\end{theorem}
\begin{proof}
At fixed finite regulator the phase is analytic near its unique critical
point and has a nondegenerate real symmetric Hessian. Diagonalize that
Hessian, use the signed one-dimensional Gaussian stationary-phase factor
in every coordinate, and expand the smooth amplitude. Equivalently the
finite stationary coefficient rule of V1 gives leading modulus
$\rho_a(g,A_*)/\sqrt{|\det(VH_a)|}$, where $V$ is the diagonal matrix
of cell weights, repeated 24 times. The signature has 12 signs of each
kind per site by \cref{thm:v10-local-inverse}, so the signature phase is
constant and cancels in the normalization. The powers of $2\pi\hbar$,
$\det V$, $a^6$ from each spatial link and $\chi$ cancel as well. Hence
the negative logarithm of the relative leading modulus is
\[
 \tfrac12\log\left|\frac{\det H_a}{\det C_0}\right|
                 -\log\frac{\rho_a(g,A_*)}{\rho_a(\eta,0)}.
\]
Insert \eqref{eq:v10-C-det}, \eqref{eq:v10-loop-density} and
\eqref{eq:v10-full-density}. The determinant contributes
$3\sum\log|\det g|$ and the Lorentz quotient contributes
$\tfrac12\sum\log|\det g|$, proving
\eqref{eq:v10-main-amplitude}. The finite stationary-phase error is
$O(\hbar)$ in the normalized leading expansion on a fixed compact source
subchart; \emph{its constant is not claimed uniform in volume or cutoff}.
Only its leading coefficient has the separate uniform bounds proved here.

For constant $e$, the native load is zero and
$\partial_A r_a(e,0)=D_A^2r_a(e,0)=0$. Thus $A_*=0$, $K_a=0$ and
$j(0)=1$. This proves \eqref{eq:v10-constant-density}; ordinary matrix
differentiation along homogeneous backgrounds gives \eqref{eq:v10-density-contacts}. Finally
$Z/Z_0=e^{i(S^*-S_0^*)/\hbar}e^{-\mathcal M_{a,1}}(1+O(\hbar))$
explains the factor $i$ in the Lorentzian effective phase.
\end{proof}

The factor $7/2$ is a calculated \emph{measure convention}, not a universal
gravity observable. Starting with flat metric measure instead of Lebesgue
coframe measure removes its $1/2$ piece. Multiplying the entrance density by
$|\det g|^\alpha$ shifts it to $7/2-\alpha$.
It cannot be called a cosmological constant: $\sum_x\log|\det g_x|$ is
not the covariant density $\int\sqrt{-g}$. It belongs to the coordinate and
symmetry-restoring local bank until a common physical measure and nonlinear
Ward normalization have been established. With isotropic cell volume $a^4$
its coefficient per physical volume is of order $a^{-4}$.

Under any actual smooth connection coordinate change $A=\Psi(y,e)$, the
stationary Hessian becomes $D\Psi^{\mathsf T}H_aD\Psi$, while the amplitude
acquires $|\det D\Psi|$. Their factors cancel in
$\rho_*/\sqrt{|\det H|}$. Thus \eqref{eq:v10-main-amplitude} is not
altered by reexpressing the same measured germ in valid holonomy
coordinates. Replacing the pulled-back measure by flat $dy$ is a change of
measure, not coordinate invariance.

\section{A finite-order local continuum matching theorem for the connection loop}
\label{sec:v10-cutoff}

The bounds in \cref{sec:v10-local} concern bounded classical backgrounds.
They do not by themselves imply small quantum fluctuations in a shrinking
cell: the Gaussian covariance in unweighted connection coordinates has
size $\hbar a^{-4}C_0^{-1}$. The following argument is instead strictly
coefficientwise. It uses the native action's special scaling and pays its
local divergent jets before taking a cutoff limit.

Take $\tau=a$, $\Lambda=0$, and expand around $g=\eta$, $j=0$ to any
fixed total source degree $N$. Metric/coframe conversion uses the source's
analytic section $\widehat e(g)$. Connection sources are normalized with
the fixed $C_0$ convention; no additional physical scaling of a spin current
is inferred. Let $I$ be a labelled panel with $q$ Cartan-source marks and
$n-q$ metric marks, $n\le N$. Denote their norm product by $\mathfrak n_I$.
Independent external momenta form $K\in\R^{4(n-1)}$; take a fixed smooth
window $|K|\le M$. Require $aM<c_N$ and enough lattice sites to exclude
wrapping of a closed graph of at most $N$ elementary plaquette vertices.
For example $L_i/a>8N+8$ is sufficient. These restrictions hold on every
fixed positive physical torus for small enough $a$.

\begin{lemma}[No unbounded loop momentum in a fixed auxiliary coefficient]
\label[lemma]{lem:v10-analytic-symbol}
The coefficient of \eqref{eq:v10-main-amplitude} per physical volume is
exactly of the form
\begin{equation}
 \mathfrak m_{a,I}(K)=a^{q-4}\,F_I(aK),
 \label{eq:v10-scaling}
\end{equation}
where $F_I$ is analytic on a complex neighborhood of $K=0$ after its
argument has been made dimensionless. Every prescribed finite derivative
has a bound $C_{N,r}\mathfrak n_I$ on a smaller neighborhood, independent
of $a$ and the physical periods. All source contacts, repeated marks and
translated link vertices are included.
\end{lemma}
\begin{proof}
Set $X=aA$. Formula \eqref{eq:v10-full-remainder} is exactly
$r_a=a^{-2}\mathfrak r(e,X)$, where $\mathfrak r$ is analytic, has degree
at least three in $X$, and depends on a fixed set of neighboring entries.
The stationary equation becomes
\[
 \mathsf C(e)X+\partial_X\mathfrak r(e,X)
                 =a j-a f_a^\circ(e).
\]
The flat linear inverse is the \emph{pointwise} $C_0^{-1}$. At any fixed
source degree the stationary coefficients are therefore finite trees of
local vertices and this constant inverse. A geometric load vertex contains
either $2i\sin(ak_i/2)$ or $iak_0$, applied to the actual quadratic
coframe coefficient. Each Cartan-source mark contributes exactly one
factor $a$. Products, translations and the analytic section introduce
finite products of $e^{iak_i}$ and functions analytic in $aK$, not soft
momentum inverses. Partial external momenta stay on the nonaliased chart.

Similarly $K_a=\mathsf C^{-1}D_X^2\mathfrak r(e,X_*)$ starts at source
degree one. Its logarithm at degree $n$ has at most $n$ factors. Each closed
walk is made from a fixed number of local vertices and the pointwise flat
inverse; there is no free propagating momentum denominator to sum. The
same holds for the finite stationary trees attached to that walk. The
nonwrapping condition makes their list independent of the lattice size.
The Haar term is $\log j(X_{*,i})$ and the metric term is local. All have
analytic finite coefficients in $aK$, with a factor $a^q$ for $q$ Cartan
marks. Translation invariance supplies one root sum; division by physical
volume changes it to $a^{-4}$. This proves \eqref{eq:v10-scaling}.
Finitely many tree/loop patterns at fixed $N$, the nonzero fixed Cartan
inverse and ordinary Cauchy bounds give the stated uniform finite jets.
Dense phase differentiation has been used only through its \emph{actual}
analytic symbol on each declared external sum, not through an assumed
full-grid local calculus.
\end{proof}

For $q\le4$ let $\mathsf T_{4-q}$ be the complete Taylor polynomial in
$K$ of total degree at most $4-q$. Store it in the common local source
functional. For $q>4$ set $\mathsf T_{4-q}=0$. Denote the complementary
coefficient by
$\mathfrak m^R_{a,I}=(I-\mathsf T_{4-q})\mathfrak m_{a,I}$.

\begin{theorem}[The auxiliary one-loop continuum contribution is local]
\label[theorem]{thm:v10-local-continuum}
For $q\le4$, $r=|\alpha|\le5-q$ and $|K|\le M$,
\begin{equation}
 \boxed{|\partial_K^\alpha\mathfrak m^R_{a,I}(K)|
 \le C_{N,r}\mathfrak n_I\,a\,|K|^{5-q-r}.}
 \label{eq:v10-cutoff-rate}
\end{equation}
For every other fixed derivative order the scaled bound
$M^r|\partial_K^\alpha\mathfrak m^R_{a,I}|
\le C_{N,r}\mathfrak n_I aM^{5-q}$ holds on the same fixed window.
For $q>4$ the coefficient, with no subtraction, is
$O(\mathfrak n_I a^{q-4})$ and all its fixed scaled derivatives are bounded
by the same order. Hence all these nonlocal coefficients tend to zero.
The stored local jets can have precisely the growth
\begin{equation}
 \partial_K^\alpha\mathfrak m_{a,I}(0)
   =a^{q+|\alpha|-4}\partial^\alpha F_I(0),
 \qquad q+|\alpha|\le4.
 \label{eq:v10-local-growth}
\end{equation}
At each fixed $N$ this is a finite local bank. Once its coefficients are
matched to the same physical local source conditions, the connection
one-loop module has a unique, purely local continuum matching contribution.
It is not discarded: the local matching data remain in the ESM state.
\end{theorem}
\begin{proof}
Apply the integral Taylor remainder of degree $D=4-q$ to $F_I$ in
\eqref{eq:v10-scaling}. A differentiated remainder is bounded by
$C a^{q-4+r}|aK|^{D+1-r}$, which is exactly
\eqref{eq:v10-cutoff-rate}. When $r>D+1$, the polynomial has vanished;
the analytic derivative bound is $C a^{q-4+r}$, and $aM<c_N$ implies the
stated weaker scaled bound. If $q>4$, the prefactor in
\eqref{eq:v10-scaling} already vanishes and the same derivative estimate
applies. Evaluating the Taylor derivatives at zero proves
\eqref{eq:v10-local-growth}. A finite number of labelled source types and
Taylor monomials at $n\le N$ gives a finite bank, with no claim of a fixed
finite bank for unbounded metric degree. Adding identical finite local
normalization data to the vanishing remainders proves the matching claim.
No covariance positivity, running Planck parameter or small full quantum
remainder is used in this coefficient proof.
\end{proof}

A smooth external window turns these symbol estimates into rooted
source-kernel estimates. In $4(n-1)$ relative variables, integration by
parts more than $s+4(n-1)$ times, followed by periodization, gives for
$q\le4$
\begin{equation}
 a^{4(n-1)}\sum_{y_2,\ldots,y_n}
 \left(1+M\sum_{i=2}^n d_{\rm per}(y_i,0)\right)^s
 \|\mathcal K^R_{a,I}(y_2,\ldots,y_n)\|
 \le C_{N,s}\mathfrak n_I aM^{5-q}.
 \label{eq:v10-kernel-rate}
\end{equation}
For $q>4$ use $C_{N,s}\mathfrak n_I a^{q-4}$.
The proof needs only the finite derivative bounds just established; the
periodized sum is bounded by the full-space lattice sum. This controls the
multilinear functional with one source in $L^1$ and the rest in $L^\infty$,
with no total-volume factor. A further finite EFT subtraction through degree
$D\ge4-q$ improves the rate to
$a^{q+D-3}|K|^{D+1-r}$ at $r\le D+1$. No parity cancellation is presumed;
for example one additional derivative in the bank gives an $a^2$ rate.

This argument is special to the nonpropagating connection block at the flat
reference. Adding a connection kinetic operator, changing the native
quadratic block, or integrating metric fluctuations replaces its pointwise
inverse and invalidates this proof. It does not show that a propagating
connection or graviton determinant is local. The local polynomial is also
not automatically a nonlinear diffeomorphism-invariant curvature polynomial;
that is precisely part of the remaining restoring problem.

\section{The actual coframe--Higgs Schur interface and the remaining symmetry test}
\label{sec:v10-mixed}

The Higgs field is a spacetime Lorentz scalar. In the supplied common
Hodge matter action it has no direct dependence on the independent
spacetime connection $A_\mu^{IJ}$. The internal electroweak connection of
V9 is a \emph{different variable}; it is not eliminated by the preceding
Cartan integral. At a common declared reference chart write the scalar
quadratic action as $\langle\phi,Q(g)\phi\rangle/2$, with its actual
endomorphism and elementary sources. Then
\[
 \partial_A S_H=0,\qquad \partial_A\partial_\phi S_H=0.
\]
Thus the connection stationary section just constructed is unchanged when
this scalar block is attached at fixed $g$.

\begin{proposition}[Native mixed Hessian after connection elimination]
\label[proposition]{prop:v10-mixed}
Let $q$ be any retained coframe or metric chart coordinate, with all chart
chain rules included, and let $H_a$ be the actual connection Hessian.
At a connection-stationary background $(q,A_*,\phi)$ put
\[
 T=D_qD_A S_a,\qquad
 (L_\phi\psi)[h]=\langle\psi,D_qQ[h]\phi\rangle.
\]
The Hessian in $(q,A,\phi)$ has exactly the block form
\begin{equation}
 \begin{pmatrix}
 S_{qq}^{\rm grav}+\tfrac12\langle\phi,Q_{qq}\phi\rangle&T&L_\phi\cr
 T^*&H_a&0\cr L_\phi^*&0&Q
 \end{pmatrix}.
 \label{eq:v10-full-mixed-Hessian}
\end{equation}
Eliminating $A$ gives the actual metric block
\begin{equation}
 D_q^2S_a^*+\tfrac12\langle\phi,Q_{qq}\phi\rangle,
 \qquad D_q^2S_a^*=S_{qq}^{\rm grav}-TH_a^{-1}T^*.
 \label{eq:v10-metric-Schur}
\end{equation}
If an invertible scalar block is eliminated next, its remaining metric
return is
\begin{equation}
 \boxed{\mathbb H_q^{\rm eff}=D_q^2S_a^*
 +\tfrac12\langle\phi,Q_{qq}\phi\rangle
                              -L_\phi Q^{-1}L_\phi^*.}
 \label{eq:v10-Higgs-return}
\end{equation}
Each determinant and each source return is assigned exactly once. The
native cubic and quartic parts entering $D_q^2S_a^*$ are given by
\cref{thm:v10-reduced-jets}; the scalar operands are all derivatives of
V9's same $Q$, not independent fitted stress vertices.
\end{proposition}
\begin{proof}
Differentiate the common action. The two zero blocks follow from the
absence of $A$ in $Q$. Differentiating $S_A(q,A_*)=0$ gives
$D_qA_*=-H_a^{-1}T^*$; the chain rule then yields
\eqref{eq:v10-metric-Schur}. Complete the scalar quadratic square to obtain
\eqref{eq:v10-Higgs-return}. These are inherited Schur operations, now
applied to the identified native blocks. In particular the mixed Higgs
return is generally nonzero at a nonzero, nonconstant Higgs background.
At $\phi=0$ it vanishes; no nonzero mixed term is manufactured there.
\end{proof}

For example the bare sign-averaged Hodge derivative gives the explicit
mixed row
\begin{equation}
\begin{split}
 (L_\phi\psi)[h]={}&\frac1{16}\sum_\epsilon
 \langle D_\mu^{U,\epsilon_\mu}\psi,
       D_q(\rho_g g^{\mu\nu})[h]D_\nu^{U,\epsilon_\nu}\phi\rangle\\
 &+\langle\psi,D_q\{\rho_g(m^2I+E)\}[h]\phi\rangle.
 \label{eq:v10-Higgs-L}
\end{split}
\end{equation}
An $E$ or internal-link dependence on the geometric coordinate must also
be differentiated if it is part of the chosen physical source map.
Formula \eqref{eq:v10-Higgs-L} uses the source-independent versions held
fixed in V9. The first scalar-loop effective functional then uses precisely
V9's complete contacts; it is not recomputed in this installment.

There is an essential reference qualification. V9's quantitative scalar
bounds use a positive Euclidean matter metric and $m>0$. The native
Palatini record is Lorentzian and its connection Hessian is indefinite.
The algebraic block formulas are valid in either declared invertible
complexified reference, but the two numerical integrals have not thereby
been identified. A common Wick/contour prescription and its source map
must be supplied before their estimates are asserted for a single physical
ESM integral. A formal Laplace coefficient convention can use
\eqref{eq:v10-main-amplitude} and V9 consistently after such reference data
are fixed; this is not a proof that arbitrary Lorentzian and Euclidean
regulators are equivalent.

\subsection{What the gravitational cubic now permits one to test}

The first gravitational deformation equation is
\[
 (\mathcal S_0^{\rm grav},\mathcal S_1^{\rm grav})=0,
\]
with the action part of $\mathcal S_1^{\rm grav}$ now fixed by
\eqref{eq:v10-reduced-jets}, after the actual metric/coframe conversion.
Its remaining operand is the nonlinear geometric and ghost transformation,
not an unspecified action cubic. At quartic order the required equation is
\[
 (\mathcal S_0,\mathcal S_2)=-\tfrac12(\mathcal S_1,\mathcal S_1).
\]
The action component contains the explicit $S_4^*$ and the identified
mixed blocks, as well as the source's genuine fermionic terms when added.
Neither equation is claimed here. The six algebraic Lorentz frame ghosts
just evaluated are not the four diffeomorphism ghosts. The original
phase-compatible derivative still fails Leibniz at cutoff frequencies.

One cannot use the local matching theorem to delete the classical
$a\mathscr R_3$ term before checking this identity. Nor does the fact that
the auxiliary \emph{one-loop remainder} vanishes prove invariance of its
local bank. It isolates the remaining task: evaluate the complete native
cubic BV defect with these operands, construct its local restoring
primitive and the nonlinear source map, and then quantize the actual
propagating metric/ghost Schur block. V6's free TT Gaussian is not a
replacement for that nonlinear reduced measure.

\section{Final ledgers for V10}
\label{sec:v10-ledgers}

\subsection*{Proved}
\addcontentsline{toc}{subsection}{V10: Proved}
The completed shared-link remainder has the explicit finite cubic and
quartic word coefficients \eqref{eq:v10-Rn}; the cubic is nonzero, including
as a pure Cartan-source contact. Its stationary cubic and quartic action
are \eqref{eq:v10-reduced-jets}. In the declared Lebesgue-coframe and
Lorentz-Haar link measure, the local Lorentz quotient, the Haar factor and
the Cartan determinant give the full coefficient
\eqref{eq:v10-main-amplitude}, with the measure contact power $7/2$.
The actual auxiliary stationary Hessian and its fixed-order source
kernels have the uniform bounds \eqref{eq:v10-inverse-bound}--
\eqref{eq:v10-marked}. About the flat reference, every prescribed finite
coframe/Cartan-source one-loop coefficient has the exact scaling
\eqref{eq:v10-scaling} and a vanishing rooted nonlocal remainder after its
actual finite local Taylor bank is assigned. The mixed Higgs and metric
blocks are identified in \eqref{eq:v10-full-mixed-Hessian}--
\eqref{eq:v10-Higgs-L}. No full nonlinear gravitational master equation
is included among these results.

\subsection*{Inherited}
\addcontentsline{toc}{subsection}{V10: Inherited}
V1--V9 remain unchanged. V077 supplies the completed native action, the
actual load and Cartan congruence, the bounded stationary inverse and
quadratic constraints. These are not newly proved here. Its polynomial
cell-connection model is not used as a substitute for shared links.
V1 supplies the measured stationary compiler; the monographs supply Schur
and source-complete elimination. V9 supplies the scalar all-source module
in its own massive positive reference. The classical spacetime and
structural matter theorems retain their original hypotheses and do not
supply a quantum gravitational measure.

\subsection*{Literature input}
\addcontentsline{toc}{subsection}{V10: Literature input}
First-order path-integral measures, tetrad gravity, auxiliary connection
elimination and perturbative Holst calculations are established subjects
\cite{V10Aros,V10Benedetti}. Gaussian stationary phase, Haar volume through
the exponential and matrix-log identities are standard tools used with
proofs of the finite instances needed here. The claimed contribution is the
specified shared-link coefficient and matching calculation, not a new
general equivalence of first- and second-order quantum gravity.

\subsection*{Conditional}
\addcontentsline{toc}{subsection}{V10: Conditional}
A physical ESM interpretation requires a common gravitational/matter
reference and contour, the actual microscopic measure or its stated
matching, nonlinear diffeomorphism-compatible local terms, and the
propagating metric/ghost block. Fermionic spin sources are not identified
with an arbitrary external $j$ without differentiating their actual finite
Dirac operator. Consuming V9's scalar cutoff estimates requires its
positive-mass and reference hypotheses. The independent-load exponential
kernel bound does not by itself control the dense native metric load at all
frequencies. Stationary-phase remainders are not uniform in the full
refinement/volume family.

\subsection*{Open}
\addcontentsline{toc}{subsection}{V10: Open}
The next smallest gravitational test is the native cubic BV defect using
\eqref{eq:v10-reduced-jets}, with a compatible nonlinear transformation and
its source-dependent measure; the quartic equation then consumes the
explicit $S_4^*$ and mixed Schur block. Propagating graviton and
four-diffeomorphism-ghost loops, their joint EFT local projection,
coframe--fermion spin/line transport, massless/reference changes and an
invariant interacting ESM shell remain open. The auxiliary local matching
result is not a proof of an all-loop ESM continuum, a positive quantum
gravity measure, asymptotic safety or a finite-parameter ultraviolet theory.

\begin{center}\small
\begin{tabular}{>{\raggedright\arraybackslash}p{.23\textwidth}>{\raggedright\arraybackslash}p{.68\textwidth}}
\toprule
Variable & Exact scope in V10\\
\midrule
Volume and time rows & Weighted auxiliary inverses and rooted fixed-order
load kernels have no site/time-count factor. Unanchored log determinants
and local measure terms remain extensive.\\
Mesh & Actual bounded-background inverse needs small $aR$ and a common
coframe margin. The coefficient cutoff theorem is about the flat branch,
fixed external momentum window, and an isotropic temporal quadrature.\\
Perturbative/source order & Cubic and quartic classical stationary action;
one auxiliary connection loop at each prescribed finite source degree $N$.
Constants depend on $N$ and derivative order, not uniformly on all orders.\\
Local bank & Metric degree grows with $N$; local source weight
$q+|\alpha|\le4$ is stored. This component bank is not already a covariant
basis for the full nonlinear BV complex.\\
Integration and measure & Local stationary phase or a declared formal Wick
convention. Lebesgue coframes, local Lorentz quotient, actual spatial Haar
factors. No global noncompact normalization or full contour is inferred.\\
Iteration & One auxiliary block and its fixed-order local matching.
No repeated interacting graviton-shell invariant class is claimed.\\
Matter & The structural Higgs field has no spacetime-connection dependence.
The identified scalar Schur interface does not establish the missing
fermionic or Lorentzian/Euclidean reference comparison.\\
\bottomrule
\end{tabular}
\end{center}

The accompanying script \srcid{check_ESM_native_connection_V10.py}
checks the finite BCH coefficients through degree four, the nonzero native
cubic example, the flat Cartan determinant, the coframe quotient Jacobian,
the Haar two-jet and the cubic/quartic stationary substitution in exact
rational algebra. A separate floating-point check tests the signed
Neumann/log-determinant identities and their bounds. These diagnostics are
not substitutes for the source-kernel and cutoff proofs, and no new Lean
formalization is asserted.

\begingroup
\renewcommand{\refname}{Additional sources and comparison for V10}
\begin{thebibliography}{99}
\bibitem[32]{V10GT}
A.~P\'elissier, \emph{Renewal Grand Tensor Monograph}, V077,
Library manuscript \srcid{Renewal_Grand_Tensor_Monograph_V077.tex}.
The newly consumed shared-link action and its scope are in
\srcid{sec:GT-shared-link-Cartan}, source lines 45270--45507;
the distinct polynomial model is in
\srcid{sec:GT-independent-Palatini-holonomy}.
\bibitem[33]{V10Aros}
R.~Aros, M.~Contreras and J.~Zanelli,
\emph{Path integral measure for first order and metric gravities},
arXiv:gr-qc/0303113 (2003). Literature comparison for measure ownership,
not an imported native finite-regulator theorem.
\bibitem[34]{V10Benedetti}
D.~Benedetti and S.~Speziale,
\emph{Perturbative quantum gravity with the Immirzi parameter},
Journal of High Energy Physics \textbf{2011}, 107 (June 2011),
arXiv:1104.4028, doi:10.1007/JHEP06(2011)107.
The present computation fixes the Palatini branch and does not import the
paper's running-coupling conclusions.
\end{thebibliography}
\endgroup



\clearpage
\part*{V11: A native cubic symmetry obstruction and explicit local EFT restoration}
\addcontentsline{toc}{part}{V11: Native cubic test and local EFT restoration}

\section{The unchanged cubic is not the correct object to complete by ghosts alone}
\label{sec:v11-start}

Sections 1--79 and their labels are preserved. The question left in V10 was
whether its actual gravitational cubic admits a nonlinear extension of the
free diffeomorphism generator. A further Library audit finds that the
classical open-basin lineage already answers the general unchanged-action
question negatively: \cite{V11Basin}, Theorems 30.2 and 31.2--31.3 and
Corollary 31.4, evaluates a noncollinear radiative Ward coefficient and
excludes its cancellation by a nonlinear transformation or a near-identity
field redefinition. Its later positive population theorem changes the
field writer; it is not a symmetry theorem for the unchanged signed action.
Those conclusions are inherited, not presented below as newly discovered
no-go principles.

The new calculation makes the obstruction useful for the present quantum
programme. A static, nonaliased three-gauge test isolates the actual
shared-link cubic: its pointwise Cartan contribution vanishes exactly,
yet its ordered plaquette contribution does not. A second test isolates a
separate phase-derivative defect after that link term is compensated.
Both tests work at fixed physical momentum and need no radiative
representative, inverse wave operator, or limiting high-frequency argument.
They determine what a restoring action must change. We then compute the
first two local mesh coefficients of the \emph{same native action},
including the free term and the ghost/antifield partners, and use them to
restore the first gravitational master coefficient in a precise finite
EFT jet. The numerical mesh and the quantum loop parameter are different
expansion parameters throughout.

This is a change of the declared local action, with its counterterms
recorded. It is not a claim that a different gauge condition repairs the
original action. The coframe--connection measure of V10 is retained, and
its connection equation and auxiliary determinant are not changed by the
new coframe-only counterterms. Exact finite-cutoff diffeomorphism symmetry,
the quartic master coefficient, and ultraviolet estimates for propagating
gravitons and their four ghosts remain separate tasks. The general BV
deformation method is inherited from \cite{V7BarnichHenneaux}; broken
finite gravitational symmetry is an established issue
\cite{V11BahrDittrich}. The action-specific coefficients and local
restoring polynomials are the objects evaluated here.

\subsection{A fixed chart and a fixed action}

Use precisely \eqref{eq:v10-action}--\eqref{eq:v10-full-remainder}, with
$j=0$, $\Lambda=0$, fixed $\chi\ne0$, and the zero independent Holst
coefficient. All matter fields are at their zero background for the
classical gravitational test. Coordinates in the local Lorentz quotient
can be chosen as the ten independent entries of a symmetric matrix $q$,
through the \emph{affine ADM coframe}
\begin{equation}
 e(q)=I+\mathsf E(q),\qquad
 \mathsf E(q)^0{}_0=-\tfrac12q_{00},\quad
 \mathsf E(q)^0{}_i=0,\quad
 \mathsf E(q)^i{}_0=q_{0i},\quad
 \mathsf E(q)^i{}_j=\tfrac12q_{ij}.
 \label{eq:v11-chart}
\end{equation}
Its lapse is positive and its spatial coframe is symmetric positive on a
fixed small ball. Thus it is the same ADM representative as in V10,
parametrized differently, not an additional Lorentz gauge. Its metric is
\begin{equation}
 g(q)=\eta+q+\mathsf Q(q,q),\qquad
 \mathsf Q(q,q)=\mathsf E(q)^T\eta\mathsf E(q),\qquad
 \eta=\operatorname{diag}(-1,1,1,1).
 \label{eq:v11-metric-chart}
\end{equation}
In particular $q$ is the metric tangent, not the full metric difference
away from zero. Every nonlinear chart contribution is accounted for by
this formula.

The time record is periodic, with the trigonometric derivative used in
V10; the static witnesses also work on a slab with a long constant
interior and appropriately recorded boundary terms. We use the periodic
coefficient convention below, so there is no boundary work. Set
\[
 d_a(p)=(p_0,\kappa_1(p),\kappa_2(p),\kappa_3(p)),\qquad
 \kappa_i(p)=2a^{-1}\sin(ap_i/2).
\]
The free generator is
\begin{equation}
 (R_ac)_{\mu\nu}=\partial^\circ_\mu c_\nu+
                         \partial^\circ_\nu c_\mu,
 \qquad \widehat{\partial^\circ_\mu}=i d_{a,\mu}.
 \label{eq:v11-free-generator}
\end{equation}
The inherited free Hessian $H_a$ obeys $H_aR_a=0$. No nonlinear Leibniz
identity is inferred. In this part $H_a$ denotes the \emph{metric} free
Hessian; the auxiliary connection Hessian of V10 is always identified as
such when mentioned.

\section{A static test that isolates the ordered shared-link contribution}
\label{sec:v11-static}

\subsection{Evaluating the coframe coefficients without a continuum substitution}

Write $A_\mu^{IJ}$ with $I<J$ for the six upper Lorentz-pair coordinates,
and represent the Lie matrix by $A_\mu=(A_\mu^{IJ})\eta$. The coefficient
of $F_{\mu\nu}^{IJ}$, for $\mu<\nu$, in the actual Palatini contraction is
\begin{equation}
 \ell_{\mu\nu,IJ}(e)=2\chi\det(e)
 \bigl((e^{-1})^\mu{}_I(e^{-1})^\nu{}_J
       -(e^{-1})^\mu{}_J(e^{-1})^\nu{}_I\bigr).
 \label{eq:v11-ell}
\end{equation}
This is a quadratic cofactor polynomial. More explicitly, if
$(r,s)$ and $(K,L)$ are the increasing complements of $(\mu,\nu)$ and
$(I,J)$, respectively, then
\[
 \ell_{\mu\nu,IJ}(e)=2\chi\,
 \epsilon_{\mu\nu rs}\epsilon_{IJKL}
 (e^K{}_r e^L{}_s-e^K{}_s e^L{}_r).
\]
The load for any coefficient of $\ell$, at its own total momentum, is
obtained by adding $+\partial^\circ_\nu\ell_{\mu\nu,IJ}$ to
$f_{\mu,IJ}$ and $-\partial^\circ_\mu\ell_{\mu\nu,IJ}$ to
$f_{\nu,IJ}$. This is just the native load in V10 in explicit pair
coordinates.

Let $U=\mathsf E(q)$ and $\mathsf T_U=U^T\otimes I_6$. The linear Cartan
coefficient is
\begin{equation}
 C_1(q)=(\operatorname{tr}U)C_0-
            \mathsf T_U^TC_0-C_0\mathsf T_U.
 \label{eq:v11-C1}
\end{equation}
It follows by differentiating \eqref{eq:v10-C-congruence}. Equations
\eqref{eq:v11-ell} and \eqref{eq:v11-C1}, together with the finite $B_3$
polynomial of V10, suffice for each computation below.

For a static spatial gauge parameter $c\exp(ip\cdot x)$, $c_0=p_0=0$,
put $k=(0,\boldsymbol\kappa(p))$. Its metric amplitude is
$q=i(kc^T+ck^T)$. The linear stationary connection in this chart is
\begin{equation}
 (A_1)^{IJ}_\mu=\tfrac12 k_\mu(k_Ic_J-k_Jc_I),
 \qquad I,J\in\{1,2,3\},
 \label{eq:v11-gauge-connection}
\end{equation}
and all boost and temporal entries vanish. Substitution of the first
cofactor derivative into $A_1=-C_0^{-1}b_1$ proves the formula; equivalently
it follows by inserting the actual spatial coframe tangent into the
inherited linear Cartan solution. In particular it is a discrete gradient
of a rotation parameter. In a nonaliased three-gauge coefficient the entire
$\langle b_2,A_1\rangle$ contribution is zero: at the momentum of its
$A_1$ leg, its contraction is a symmetric product $k_\mu k_\nu$ against
the antisymmetric curvature slots of $\ell_2$. This cancellation uses the
\emph{total} momentum in the load and no product rule.

\subsection{An exact finite-grid coefficient}

Take a spatial cube of side $L$, $a=L/N$ with $N$ odd, and
$p=2\pi m/L$, $0<m<N/2$. Set $\theta=ap$, $z=e^{i\theta}$ and
$\rho=2a^{-1}\sin(\theta/2)$. Use the three static momenta and lower
parameter polarizations
\begin{equation}
 p_1=(0,p,0,0),\quad p_2=(0,0,p,0),\quad
 p_3=(0,-p,-p,0),\qquad
 c_1=e_2,\quad c_2=c_3=e_3.
 \label{eq:v11-witness}
\end{equation}
Their sum is zero without wrapping; also
$d_a(p_1)+d_a(p_2)+d_a(p_3)=0$. Let $q_r=R_a(c_re^{ip_r x})$ and use
$\mathcal V_a(q_1,q_2,q_3)$ for the coefficient of $t_1t_2t_3$ in
$S_{3,a}^*(\sum_rt_rq_r)$ divided by the physical spacetime volume.

\begin{proposition}[A pure-gauge shared-link coefficient]
\label[proposition]{prop:v11-static-value}
For the actual V10 cubic and the data \eqref{eq:v11-witness},
\begin{equation}
 \boxed{\mathcal V_a(q_1,q_2,q_3)
       =\frac{\chi a\rho^6}{8}(2+z^{-1}).}
 \label{eq:v11-static-value}
\end{equation}
Its geometric Cartan part is zero. Its electric remainder and all
magnetic plaquettes except the $12$ plaquette are zero. Thus the
coefficient is nonzero on every stated grid, including when $p$ is fixed
as $a\downarrow0$.
\end{proposition}
\begin{proof}
Write $J_{ij}$ for the spatial rotation with upper $ij$ entry one, and
set
\[
 X=\tfrac12\rho^2J_{12},\qquad
 Y=\tfrac12\rho^2J_{23},\qquad
 Z=\tfrac12\rho^2(J_{13}+J_{23}).
\]
Formula \eqref{eq:v11-gauge-connection} gives respectively $A_1=X$,
$A_2=Y$, and $A_1=A_2=Z$ for the three labels. The load contribution
vanishes as just proved. In
$\frac12\langle A_1,C_1A_1\rangle$, the three choices of the coframe
label give $0,-i\chi\rho^5/4,+i\chi\rho^5/4$, respectively, by
\eqref{eq:v11-C1}; they sum to zero.

On the $12$ plaquette the three labelled four-slot inputs are
\[
 (X,0,-X,0),\qquad (0,Y,0,-Y),\qquad
 (Z,z^{-1}Z,-z^{-1}Z,-Z).
\]
Polarize the explicit rational polynomial \eqref{eq:v10-BCH3}, meaning
sum its six distinct assignments of the three labels. Direct use of
$[J_{12},J_{23}]=J_{13}$ and the other rotation commutators gives
\[
 [t_1t_2t_3] B_3
   =\frac{\rho^6(2z+1)}{16z}J_{12}.
\]
The flat curvature coefficient is $2\chi$ in this slot. Multiplication
by the remaining $a$ in \eqref{eq:v10-r34} proves
\eqref{eq:v11-static-value}. No temporal connection or $A_3$ occurs.
Plaquettes with direction 3 have unchanged opposite links and are exactly
identity, proving the remaining claims. Each operation is an identity
of finite matrix polynomials, with the shifted $z^{-1}$ factors kept.
\end{proof}

The real test has no hidden complexification issue. For the three real
parameters $c_r\cos(p_r\cdot x)$, the only zero-total-momentum sign
choices are all plus and all minus. Real polarization therefore gives
\begin{equation}
 \mathcal V_a(\operatorname{Re}q_1,\operatorname{Re}q_2,
                         \operatorname{Re}q_3)
       =\frac{\chi a\rho^6}{32}(2+\cos\theta)\ne0.
 \label{eq:v11-real-value}
\end{equation}
A sufficiently small real multiple lies in the declared coframe chart.
There is no use of a complex saddle or of a propagating equation here.

\subsection{What the inherited no-completion result means for a restoring bank}

For completeness, the applicable necessary condition is the elementary
one used in \cite{V11Basin}, Corollary 31.4. If $S_2$ has Hessian $H$,
$HR=0$, and a smooth nonlinear transformation extending $R$ leaves
$S_2+S_3+\cdots$ invariant, then
\[
 DS_3(q)[Rc]+\langle Hq,R_1(q,c)\rangle=0.
\]
On $q\in\operatorname{Ran}R$, the second summand vanishes. Polarization
requires $D^3S_3(Rc_1,Rc_2,Rc_3)=0$. A quadratic field redefinition changes
$S_3$ only by $\langle Hq,F_2(q,q)\rangle$, whose restriction to this
triple kernel is zero. This is an application of the inherited criterion,
not a new general theorem. The static value above supplies a simpler,
all-grid certificate for the present coefficient conventions.

In particular, arbitrary smooth spatially nonlocal $R_1$ on the same
finite carrier cannot fix the unchanged action, and changing only ghost
or antifield coefficients cannot change this antifield-free test.
Neither the six Lorentz-frame ghosts nor the one-loop measure of V10 can
cancel an order-$\hbar^0$ classical defect. Additional fields changing the
free carrier are outside this assertion.

There is also a quantitative constraint on any action repair $\Delta S_3$.
Give its Fourier trilinear coefficient the norm obtained by taking the
supremum over unit Hilbert--Schmidt polarizations. The three complex
metric amplitudes above have norms $\sqrt2\rho,\sqrt2\rho,2\rho$.
Cancellation of their test therefore requires
\begin{equation}
 \boxed{\|\Delta\mathcal V_a\|
   \ge\frac{|\chi|a\rho^3}{32}|2+e^{-i\theta}|
   \ge\frac{|\chi|a\rho^3}{32}.}
 \label{eq:v11-repair-lower}
\end{equation}
This lower bound is independent of the number of sites. At fixed physical
$p$ it has three-derivative size $ap^3$. At cutoff momenta with $ap$
fixed it is not a small relative correction to a two-derivative vertex.
Thus low-momentum consistency and a uniform ultraviolet repair are
different requirements, already in this nonaliased example.

\section{The phase artifact is independent of the shared-link artifact}
\label{sec:v11-phase}

One diagnostic is to compensate the \emph{entire} reduced cubic term
$a\mathscr R_3^{(0)}(A_{1,a})$, keeping the native phase load and $C_1$.
This is not the proposed local counterterm below; it is an exact way to
isolate what remains. It changes the action and must be so recorded.
The resulting geometric cubic still fails the gauge test.

\begin{proposition}[A distinct, nonaliased pointwise-Cartan test]
\label[proposition]{prop:v11-phase-value}
Take $0<\theta=ap<\pi/2$, and
\[
 p_1=(0,p,0,0),\quad p_2=(0,p,p,0),\quad
 p_3=(0,-2p,-p,0),\qquad c_1=c_3=e_3,\quad c_2=e_1.
\]
Put $\rho=2a^{-1}\sin(\theta/2)$ and
$\sigma=2a^{-1}\sin\theta$. The geometric cubic restricted to the three
free gauge waves is exactly
\begin{equation}
 \boxed{\mathcal V_a^{\rm geom}
       =\frac{i\chi}{4}\rho^3\sigma(2\rho-\sigma)
       =i\chi\rho^5\cos(\theta/2)[1-\cos(\theta/2)].}
 \label{eq:v11-phase-value}
\end{equation}
It is nonzero. At fixed $p$ its leading term is
$i\chi a^2p^7/8$.
\end{proposition}
\begin{proof}
The load pairing again vanishes exactly. Insert
$k_1=(0,\rho,0,0)$, $k_2=(0,\rho,\rho,0)$,
$k_3=(0,-\sigma,-\rho,0)$ into
\eqref{eq:v11-gauge-connection}. The three terms with coframe labels
$1,2,3$ in \eqref{eq:v11-C1} are
\[
 i\chi\rho^3\sigma(\rho-\sigma)/4,\qquad 0,\qquad
 i\chi\rho^4\sigma/4.
\]
Their sum is the first expression. The double-angle identity gives the
second. Finally $\rho=p+O(a^2p^3)$ and
$1-\cos(\theta/2)=a^2p^2/8+O(a^4p^4)$.
Replacing the first cosine parameter by a sine extracts a nonzero real
coefficient as in \eqref{eq:v11-real-value}. No momentum wraps when the
displayed condition on $\theta$ holds.
\end{proof}

The native phase satisfies neither $\sigma=2\rho$ nor a Leibniz rule.
Consequently subtracting the shared-link cubic alone is not a symmetry
completion. Conversely, correcting only a phase-addition defect misses
\cref{prop:v11-static-value}, where the three phase covectors already
sum to zero. These two tests prescribe two different entries in the
restoring local bank. They do not require repeating the high--low
radiative obstruction or its Sobolev endpoint estimates from
\cite{V11Basin}.

\section{Explicit local counterterms through the second mesh coefficient}
\label{sec:v11-local-repair}

We now construct a positive, limited result: the first gravitational
master coefficient can be restored in the formal local mesh jet through
$a^2$ by explicit finite-derivative counterterms of the same native
coframe action. The comparison variable $a$ is the \emph{mesh}, not
$\hbar$ or the interaction-order variable. Every estimate below is on a
fixed Fourier window, with no aliasing of the cubic panel. It will not be
used as an estimate for all modes of a graviton loop.

For ordinary derivatives define
\begin{equation}
 \mathfrak d_0=0,\qquad \mathfrak d_i=\tfrac1{24}\partial_i^3,
 \qquad \partial^\circ_\mu=\partial_\mu+a^2\mathfrak d_\mu
                  +O(a^4\partial^5).
 \label{eq:v11-phase-expansion}
\end{equation}
The sign is fixed by $\sin z=z-z^3/6+\cdots$; $\partial_i^3$ has
symbol $-ip_i^3$.
In the affine coframe chart, let $b_{r,0}$ be the homogeneous degree-$r$
load, $r=1,2$, with \emph{ordinary} derivatives in the explicit cofactor
formula \eqref{eq:v11-ell}. Let $b_{r,2}$ be the same expression with
its single derivative replaced by $\mathfrak d$. Then, on every fixed
cubic source panel,
\begin{equation}
 b_{r,a}=b_{r,0}+a^2b_{r,2}+O(a^4),\qquad
 A_{1,0}=-C_0^{-1}b_{1,0},\qquad
 A_{1,2}=-C_0^{-1}b_{1,2}.
 \label{eq:v11-load-jets}
\end{equation}
This notation specifies actual finite differential polynomials; it is
not an assumption about an unspecified effective vertex.

Let $\mathscr R_{3,\rm pt}$ be the unshifted cubic at the flat coframe:
\begin{equation}
\begin{split}
 \mathscr R_{3,\rm pt}(A)={}&
 -\tfrac12\sum_i b(\Pi_i^0,[A_i,[A_i,A_0]])\\
 &+\tfrac12\sum_{i<j}b(\Sigma_{ij}^0,
                       [A_i+A_j,[A_i,A_j]]).
 \label{eq:v11-Rpt}
\end{split}
\end{equation}
Its first spatial-shift coefficient is equally explicit:
\begin{equation}
\begin{split}
 \mathscr R_{3,\rm sh}(A)={}&
 -\tfrac12\sum_i b(\Pi_i^0,[A_i,[A_i,\partial_iA_0]])\\
 &+\sum_{i<j}b\bigl(\Sigma_{ij}^0,
 DB_3(A_i,A_j,-A_i,-A_j)
       [0,\partial_i A_j,-\partial_j A_i,0]\bigr).
 \label{eq:v11-Rsh}
\end{split}
\end{equation}
Here $DB_3[X][Y]$ means replace one occurrence in each of the explicitly
listed triple commutators \eqref{eq:v10-BCH3} by its corresponding $Y$
and sum the three replacements. It is a finite rational four-derivative
polynomial after $A=A_{1,0}$ is inserted. This definition leaves neither
a matrix inversion at nonzero field nor an unevaluated action derivative.

\begin{theorem}[The evaluated native local artifact bank]
\label[theorem]{thm:v11-artifacts}
Define
\begin{align}
 \mathcal O_{2,2}&=\langle b_{1,2},A_{1,0}\rangle,\\
 \mathcal O_{3,1}&=\mathscr R_{3,\rm pt}(A_{1,0}),\\
 \mathcal O_{3,2}&=\langle b_{2,2},A_{1,0}\rangle
             +\langle b_{2,0},A_{1,2}\rangle
             +\langle A_{1,0},C_1 A_{1,2}\rangle
             +\mathscr R_{3,\rm sh}(A_{1,0}).
 \label{eq:v11-artifact-bank}
\end{align}
They have respectively field/derivative degrees $(2,4)$, $(3,3)$,
and $(3,4)$. If $S_{2,0}$ and $S_{3,0}$ denote the zero-mesh coefficients
of the same stationary Palatini action, then
\begin{align}
 S^*_{2,a}&=S_{2,0}+a^2\mathcal O_{2,2}+\mathcal E_{2,a},\\
 S^*_{3,a}&=S_{3,0}+a\mathcal O_{3,1}
                    +a^2\mathcal O_{3,2}+\mathcal E_{3,a}.
 \label{eq:v11-action-expansion}
\end{align}
For source momenta bounded by $\Lambda_*>0$ with $a\Lambda_*\le1/64$,
the polarized remainder symbols and every fixed momentum derivative obey
\begin{equation}
 \left|\partial^\alpha\mathcal E_{2,a}\right|
       \le C_\alpha a^4\Lambda_*^{6-|\alpha|},\qquad
 \left|\partial^\alpha\mathcal E_{3,a}\right|
       \le C_\alpha a^3\Lambda_*^{5-|\alpha|},
 \label{eq:v11-artifact-remainders}
\end{equation}
with the source-polarization norms understood. Constants depend on the
fixed normalization and derivative order, not on the periods or mesh.
\end{theorem}
\begin{proof}
The coefficients $C_0,C_1$ are independent of the mesh. Substituting
\eqref{eq:v11-load-jets} into the exact native formula
\eqref{eq:v10-reduced-jets} first gives the quadratic correction
$-\langle b_{1,2},C_0^{-1}b_{1,0}\rangle$ and the first three terms of
$\mathcal O_{3,2}$. Their signs are exactly those in
\eqref{eq:v11-artifact-bank}. The phase load has no linear mesh term.

In the link term, $A_{1,a}=A_{1,0}+a^2A_{1,2}+O(a^4)$, while each
translation is $T_i=I+a\partial_i+O(a^2\partial_i^2)$ on these panels.
Thus
\[
 a\mathscr R_3^{(0)}(A_{1,a})=
 a\mathscr R_{3,\rm pt}(A_{1,0})
 +a^2\mathscr R_{3,\rm sh}(A_{1,0})+O(a^3).
\]
Inserting the correction $a^2A_{1,2}$ into this cubic first contributes
at order $a^3$, not $a^2$. The electric endpoint and the two shifted
magnetic slots yield exactly \eqref{eq:v11-Rsh}. This proves every
coefficient in \eqref{eq:v11-action-expansion}.

For the bounds, set all momenta equal to $\Lambda_*$ times bounded
dimensionless momenta. The exact cubic consists of a finite sum of
trigonometric entire symbols, phase derivatives at the \emph{total}
subset momentum, and the fixed inverse $C_0^{-1}$. The nonaliased source
window bounds every argument by a fixed multiple of $a\Lambda_*$. Taylor's
integral formula for sine and each shifted exponential, including any
fixed number of differentiated versions, is uniform on that compact set.
The free coefficient has two derivatives and its first omitted factor
is $(a\Lambda_*)^4$. The geometric cubic has the same omitted factor;
the shifted-link cubic has two derivatives times one explicit
$a\Lambda_*$ and its first omitted factor is $(a\Lambda_*)^2$.
Their bounds are therefore $a^4\Lambda_*^6$ and $a^3\Lambda_*^5$.
Each momentum derivative supplies $\Lambda_*^{-1}$. This argument also
covers orders greater than the displayed polynomial degrees, using the
same differentiated integral remainder. It proves
\eqref{eq:v11-artifact-remainders} without taking a norm of a full
extensive action.
\end{proof}

The local change we actually propose is
\begin{equation}
 \boxed{\mathcal C_a(q)=
 -a^2\mathcal O_{2,2}(q)-a\mathcal O_{3,1}(q)
                         -a^2\mathcal O_{3,2}(q).}
 \label{eq:v11-counterterm}
\end{equation}
These are restoring operators in the declared component chart. They are
not asserted to be independent diffeomorphism-invariant curvature
operators. They do not fix any physical higher-curvature Wilson
coefficient. Covariant finite additions remain independent matching data.
The two-derivative Einstein normalization is unchanged. The correction
includes a quadratic higher-derivative term; it does not claim to preserve
the native free action exactly at finite $a$.

\section{The first gravitational BV coefficient in the finite local EFT jet}
\label{sec:v11-BV}

Let $c^\mu$ be the four odd diffeomorphism ghosts, $q^{*\mu\nu}$ the
metric-coordinate antifields, and $c^*_\mu$ the ghost antifields. They are not the
six Lorentz-frame variables already quotiented in V10. We use the
antibracket convention in which the antifield terms below generate
$sq=R(c)$ and $sc=c\cdot\partial c$. All equations can equivalently be
read as Noether and graded-commutator identities, independent of this
notation.

At zero mesh, use the \emph{actual metric map}
\eqref{eq:v11-metric-chart}. Its first nonlinear transformation is
\begin{equation}
 \boxed{R_1^{(0)}(q,c)=\mathcal L_c q
 -\mathsf E(R_0c)^T\eta\mathsf E(q)
 -\mathsf E(q)^T\eta\mathsf E(R_0c),}
 \label{eq:v11-R1}
\end{equation}
where
\[
 (R_0c)_{\mu\nu}=\partial_\mu c_\nu+\partial_\nu c_\mu,
 \quad(\mathcal L_cq)_{\mu\nu}=c^\rho\partial_\rho q_{\mu\nu}
   +q_{\rho\nu}\partial_\mu c^\rho
   +q_{\mu\rho}\partial_\nu c^\rho.
\]
Thus the usual Lie variation of a metric is \emph{not} silently used as
the variation of the affine coframe coordinate $q$.
Define
\[
 (R_{0,2}c)_{\mu\nu}=\mathfrak d_\mu c_\nu+
                                     \mathfrak d_\nu c_\mu,
 \qquad \widehat R_{0,a}=R_a-a^2R_{0,2}.
\]
Set $\widehat S_{2,a}=S^*_{2,a}-a^2\mathcal O_{2,2}$ and
$\widehat S_{3,a}=S^*_{3,a}-a\mathcal O_{3,1}-a^2\mathcal O_{3,2}$.
The specified free and first interaction master coefficients are
\begin{align}
 \widehat{\mathbb S}_{0,a}
 &=\widehat S_{2,a}+\langle q^*,\widehat R_{0,a}c\rangle,\\
 \widehat{\mathbb S}_{1,a}
 &=\widehat S_{3,a}+\langle q^*,R_1^{(0)}(q,c)\rangle
                    +\langle c^*,c\cdot\partial c\rangle.
 \label{eq:v11-master-coefficients}
\end{align}
Pairings and all source derivatives retain their physical quadrature
weights. Products in the following local EFT identity are ordinary
products before a finite coefficient panel is sampled. On the stated
nonaliased panel their integral equals the native discrete pairing.

\begin{theorem}[Explicit first-interaction restoration through mesh degree two]
\label[theorem]{thm:v11-BV-jet}
In the local differential-polynomial algebra with coefficients in
$\R[a]/(a^3)$, and retaining interaction degree at most one,
\begin{equation}
 \boxed{
 (\widehat{\mathbb S}_{0,a},\widehat{\mathbb S}_{0,a})=0,
 \qquad
 (\widehat{\mathbb S}_{0,a},\widehat{\mathbb S}_{1,a})=0.}
 \label{eq:v11-master-jet}
\end{equation}
At finite nonzero $a$ these are consistency estimates, not exact
identities on the entire finite lattice. Specifically, the action-side
cubic Noether residual has, on the same scaled window,
\begin{equation}
 \left|\partial^\alpha
 \{D\widehat S_{3,a}(q)[\widehat R_{0,a}c]
    +D\widehat S_{2,a}(q)[R_1^{(0)}(q,c)]\}\right|
 \le C_\alpha a^3\Lambda_*^{6-|\alpha|}
                              \|q_1\|\|q_2\|\|c\|,
 \label{eq:v11-Noether-rate}
\end{equation}
read as the polarized coefficient in the two metric arguments. The
free action--ghost residual is bounded by
$C_\alpha a^4\Lambda_*^{7-|\alpha|}$, and the first-interaction
metric-antifield/two-ghost residual by
$C_\alpha a^4\Lambda_*^{6-|\alpha|}$, with the corresponding argument
norms. Constants are fixed-order and period independent.
\end{theorem}
\begin{proof}
First identify the zero-mesh action, rather than assign it by a desired
Ward conclusion. In \eqref{eq:v10-action} replace the spatial phase
load by its zero-mesh derivative and take the zero-mesh link coefficient.
The latter is zero. The remaining boundary-completed pointwise Palatini
action is quadratic in its independent connection. Its unique stationary
connection is the ordinary torsion-free connection of $e(q)$, by the
inherited Cartan equation. Consequently its first two nonzero reduced
coefficients are precisely $S_{2,0}$ and $S_{3,0}$ of
\cref{thm:v11-artifacts}. This is the zero-mesh local coefficient of the
native action, not its identification at finite mesh.

For ordinary derivatives, the Palatini integrand is a differential
four-form. Its simultaneous pullback variation is an exact differential;
the periodic convention makes its integral zero. Connection stationarity
removes the connection-variation term in the reduced action. Differentiating
$g(q)$ shows that the induced metric-coordinate generator is
$R_0+R_1^{(0)}+O(q^2)$, with \eqref{eq:v11-R1}. Expanding this invariant
functional therefore proves directly
\[
 DS_{2,0}[R_0c]=0,\qquad
 DS_{3,0}[R_0c]+DS_{2,0}[R_1^{(0)}(q,c)]=0.
\]
Ordinary vector fields satisfy the Lie commutator identity. Conjugating
that identity by the local coordinate map $g(q)$ and retaining its lowest
field degree gives the corresponding two-ghost identity for $R_0,R_1^{(0)}$
and $c\cdot\partial c$. For example, on $g$ itself,
$s(\mathcal L_cg)=\mathcal L_{sc}g-\mathcal L_c\mathcal L_cg=0$
when $sc=c\cdot\partial c$; the invertible chart transfers it to $q$.
These action and transformation identities are exactly the two antibracket
identities at $a=0$.

By \eqref{eq:v11-action-expansion},
$\widehat S_{2,a}=S_{2,0}+O(a^4)$,
$\widehat S_{3,a}=S_{3,0}+O(a^3)$, and
$\widehat R_{0,a}=R_0+O(a^4)$. Hence their classes in $\R[a]/(a^3)$
are those just verified, proving \eqref{eq:v11-master-jet}. This step uses
the evaluated counterterms, not an assumption that every local defect
admits a primitive.

For the finite-window bounds, each $R_0$ or $R_1^{(0)}$ has one derivative.
The cubic remainder in \eqref{eq:v11-artifact-remainders} therefore gives
$a^3\Lambda_*^6$ in the Noether row; the quadratic remainder gives
$a^4\Lambda_*^7$, and the free-generator difference has the same or
smaller size. Since $a\Lambda_*\le1/64$, these are bounded by the stated
cubic rate. The free Noether row has the product of its two-derivative
action with its one-derivative generator and an omitted factor
$(a\Lambda_*)^4$, giving $a^4\Lambda_*^7$. The two-ghost row compares
one first-order generator with another and has the same omitted factor,
giving $a^4\Lambda_*^6$. Product differentiation yields all fixed
momentum-derivative estimates. No inverse external momentum occurs.
\end{proof}

The first corrected free generator is not asserted to exponentiate to an
exact lattice diffeomorphism group. Nor is
$(\widehat{\mathbb S}_{1,a},\widehat{\mathbb S}_{1,a})$ zero: its quartic
completion is the next interaction degree and is not included in
\eqref{eq:v11-master-jet}. In particular this theorem says nothing about
the quantum BV Laplacian of the full Lorentz-Haar density. Its target is
the explicitly stated \emph{classical first-interaction local EFT jet}.

\section{Measured insertion, source locality, and the ultraviolet boundary of the repair}
\label{sec:v11-measured}

The correction \eqref{eq:v11-counterterm} is a function only of the
retained coframe/metric coordinate. It can be inserted into the native
\emph{pre-elimination} action. With exactly V10's finite local contour
and amplitude convention, the identity is
\begin{equation}
 Z_a^{\rm new}[q,j]=e^{i\mathcal C_a(q)/\hbar}Z_a[q,j].
 \label{eq:v11-measured-insertion}
\end{equation}
This is the inherited measured pushforward rule applied to these actual
counterterms. The Cartan stationary section, its Hessian, and every
$j$ derivative of the connected logarithm are unchanged by this insertion. In
particular V10's leading auxiliary amplitude
\eqref{eq:v10-main-amplitude} is retained, not recomputed or suppressed.
Metric differentiation adds the derivatives of \emph{the same}
$\mathcal C_a$, including the quadratic and cubic contact terms.
The scalar block of V9 remains attached to the same metric $g(q)$ with
its actual chain rule. No Euclidean scalar integral is equated to a
Lorentzian gravitational integral by \eqref{eq:v11-measured-insertion}.

The polynomials in \eqref{eq:v11-counterterm} are local on a smooth
coframe interpolation. When sampled as spectral finite derivatives they
are dense finite stencils, just as the native phase load is dense. Their
fixed-window source kernels nevertheless have uniform weighted bounds.
To state one, multiply the cubic remainder symbol by a fixed smooth
window on all three external legs supported in $|p_r|<\Lambda_*$. Let
$\mathcal K_{3,a}$ be the resulting two-relative-coordinate lattice
kernel, with physical Fourier normalization. For every fixed $s\ge0$,
\begin{equation}
 \boxed{a^8\sum_{x,y}
 [1+\Lambda_*(d_{\rm per}(x,0)+d_{\rm per}(y,0))]^s
 \|\mathcal K_{3,a}(x,y)\|
       \le C_s a^3\Lambda_*^5.}
 \label{eq:v11-kernel}
\end{equation}
The quadratic remainder has the analogous one-relative-coordinate bound
$C_s a^4\Lambda_*^6$. The cubic Noether kernel has bound
$C_s a^3\Lambda_*^6$, with ghost argument retained. All field-source
contacts in these finite polynomials are included.

For clarity, these kernel estimates have an elementary proof: rescale the
eight momentum variables by $\Lambda_*$, use
\eqref{eq:v11-artifact-remainders} with more than $s+8$ derivatives,
and integrate by parts in the inverse Fourier transform. The resulting
kernel is bounded by $C a^3\Lambda_*^{13}$ times a function decaying
faster than $(1+\Lambda_*(|x|+|y|))^{-s-9}$. Summation with weights $a^8$
is bounded uniformly for $a\Lambda_*\le1/64$. Periodize, use
$d_{\rm per}\le d$ for each lift, and apply the triangle inequality.
The proof for the other kernels changes only the number of relative
coordinates and the symbol scale. Consequently the corresponding
multilinear functional is bounded with one source in $L^1$ and the
others in $L^\infty$, without an extensive volume multiplier.

There is an important negative conclusion about iteration. If the upper
scale is $\Lambda_J\asymp a^{-1}$, bounds of relative size
$(a\Lambda_j)^3$ alone give
\[
 \sum_{j\le J}(a\Lambda_j)^3\asymp1
 \quad\text{for}\quad\Lambda_j=2^j\Lambda_0,
\]
not a quantity vanishing with the cutoff. Thus
\eqref{eq:v11-Noether-rate} cannot be advertised as a uniform complete
shell Ward theorem. It is a controlled tree-level source/mesh statement.
The full ultraviolet problem needs a restoring map on the actual
graviton/ghost shell family, or a matching theorem paying for its entire
high-frequency defect. The explicit lower bound
\eqref{eq:v11-repair-lower} shows why that extra step is necessary.

\section{Next interaction degree and the status of the changed direction}
\label{sec:v11-next}

The source audit changes the immediate question. It is not useful to keep
asking for ghosts that complete the unchanged V10 cubic: the general
obstruction is inherited, and the new static coefficient displays it
without a radiation or aliasing argument. The constructive question is
whether the local restoration just evaluated can be extended to an
admissible full-source shell map. The first required interaction equation
is now
\[
 (\widehat{\mathbb S}_0,\widehat{\mathbb S}_2)
       =-\tfrac12(\widehat{\mathbb S}_1,
                         \widehat{\mathbb S}_1),
\]
with a \emph{changed, explicitly recorded} input cubic and its native
connection measure. Its action part must use V10's actual quartic,
including the stationary $G_2C_0^{-1}G_2$ return, together with the
quartic restoring terms and the mixed scalar contacts already evaluated.
A cubic coordinate counterterm alone supplies no missing quartic identity.

No invariant higher-derivative physical Wilson coefficient has been fixed
by the artifact subtraction. No new physical mode or compensator has been
added. A microscopic action whose acceptance row matched the uncorrected
finite action does not automatically match this restored action. The
comparison and any additional quantum measure changes remain measurable
or separately stated matching data, exactly as for the original
shared-link compensation in the monograph.

\begingroup
\renewcommand{\refname}{Additional sources for V11}
\begin{thebibliography}{99}
\bibitem[35]{V11Basin}
A.~P\'elissier, \emph{Renewal CMP Open-Basin $3+1$ Research Notes},
V16, Library PDF
\srcid{Renewal_CMP_Open_Basin_3plus1_Research_Notes_V16.pdf}.
The prior native cubic obstruction, field-redefinition criterion and
cokernel test are in its V4 Sections 28--35, especially Theorems 30.2,
31.2--31.3 and Corollary 31.4 (PDF pages 50--61). The shared-representative
analysis is in its V5 Sections 36--42. These classical results are
inherited; its later changed-writer population is not used as a quantum
symmetry theorem.
\bibitem[36]{V11BahrDittrich}
B.~Bahr and B.~Dittrich, \emph{(Broken) gauge symmetries and constraints
in Regge calculus}, Classical and Quantum Gravity \textbf{26} (2009),
225011, arXiv:0905.1670, doi:10.1088/0264-9381/26/22/225011.
Comparison with broken finite diffeomorphism symmetry; not an imported
coefficient or restoring theorem for the present regulator.
\end{thebibliography}
\endgroup

\section{Final ledgers for V11}
\label{sec:v11-ledgers}

\subsection*{Proved}
\addcontentsline{toc}{subsection}{V11: Proved}
The nonaliased, static three-gauge restriction of the actual native cubic
is \eqref{eq:v11-static-value}; its geometric part cancels, so it isolates
a nonzero shared-link term on every stated grid. It gives the explicit
lower bound \eqref{eq:v11-repair-lower} for an action repair. The independent
geometric test \eqref{eq:v11-phase-value} survives complete cubic-link
compensation. The local native mesh coefficients are the evaluated
polynomials \eqref{eq:v11-artifact-bank}. Their coframe-only counterterms,
free-generator correction and chart-correct first nonlinear ghost map
give \eqref{eq:v11-master-jet} in the local mesh jet through degree two,
with the finite-window rates and rooted kernels displayed above.
Their insertion leaves the actual auxiliary connection integral and
leading measure coefficient unchanged apart from the recorded action
factor. These are fixed tree-interaction and mesh-order statements,
not an exact all-frequency master equation.

\subsection*{Inherited}
\addcontentsline{toc}{subsection}{V11: Inherited}
V1--V10 and all labels are unchanged. V10 supplies the native completed
shared-link action, its reduced cubic and quartic and its declared
Lorentz-Haar auxiliary measure. The newly located classical open-basin
source already excludes exact unchanged-action cubic completion and
near-identity field-redefinition repair; that general obstruction is not
new here. Its radiative endpoint bounds and its changed-writer Cauchy
population are not reproved. V9 supplies its scalar source module at its
original positive-mass, reference and symmetry scope. Measured insertion,
BRST deformation language and stationary elimination are inherited tools.

\subsection*{Literature input}
\addcontentsline{toc}{subsection}{V11: Literature input}
BV deformation equations and the distinction between local and unrestricted
primitives are established in \cite{V7BarnichHenneaux}. Finite gravitational
gauge breaking is studied in \cite{V11BahrDittrich}. Neither reference
supplies the static shared-link coefficient, the phase-only test or the
explicit artifact bank computed here. No first general construction of
perturbative gravity or of symmetry-restoring counterterms is claimed.

\subsection*{Conditional}
\addcontentsline{toc}{subsection}{V11: Conditional}
An ESM shell application requires a common matter/gravity reference,
full source normalization, a compatible quantum BV measure and a bounded
restoring map through the ultraviolet shell. The local mesh-jet identity
does not supply those hypotheses. The new action must be matched to the
microscopic finite action when such a matching is part of the programme.
The quartic correction must retain the native stationary return and all
matter/antifield contacts. A vanishing tree source defect at fixed physical
momentum does not imply a summable all-mode quantum defect.

\subsection*{Open}
\addcontentsline{toc}{subsection}{V11: Open}
The next smallest constructive task is the quartic local master coefficient
for the explicitly restored action, and its extension to one complete
propagating gravitational/ghost shell with an actual ultraviolet source
bound. The original unchanged-action cubic-completion route is excluded
on the inherited branch, not left as a repeatedly posed assumption.
Quantum measure restoration, chiral spin/line transport, massless
reference changes and the invariant filtered interacting ESM trajectory
remain open. No nonperturbative or finite-parameter quantum-gravity
conclusion follows.

\begin{center}\small
\begin{tabular}{>{\raggedright\arraybackslash}p{.22\textwidth}>{\raggedright\arraybackslash}p{.69\textwidth}}
\toprule
Variable & Scope in V11\\
\midrule
Volume & Static coefficients are per physical volume. Fixed-window rooted
kernel bounds have no total-volume factor. Periodic or recorded boundary
conventions are essential.\\
Mesh and momentum & Exact witnesses hold at the displayed nonaliased
momenta. Positive restoring bounds require $a\Lambda_*\le1/64$;
no ultraviolet smallness at $\Lambda_*\asymp a^{-1}$ is claimed.\\
Interaction and loop order & Tree-level quadratic and cubic action, first
interaction master coefficient. The old auxiliary one-loop amplitude is
retained, not upgraded to a propagating quantum result.\\
EFT and source order & Explicit mesh degrees 1 and 2, at most cubic
field/ghost panels and every prescribed fixed derivative/moment norm.
No constants uniform in unbounded source or derivative order.\\
Measure and modes & Same native connection block and declared local
Lorentz-Haar measure. Four diffeomorphism ghosts are retained as sources;
they are not integrated or replaced by Lorentz-frame ghosts.\\
Iteration & No all-shell summability follows from the fixed-window
relative defect. Full invariant-class and quantum BV estimates remain
separate.\\
\bottomrule
\end{tabular}
\end{center}

The self-contained verifier \srcid{check_ESM_gravitational_BV_V11.py}
checks the two exact static coefficients, the complete cofactor/load
cancellations and the ordered logarithm in rational symbolic algebra.
It separately tests the native low-momentum artifact expansion and the
assembled Noether residual with finite numerical coefficients. The latter
checks are diagnostics of signs and normalizations, not proofs of the
uniform estimates. No new proof-assistant formalization is asserted.



\clearpage
\part*{V12: Quartic gravitational restoration, mixed Higgs contacts, and the coframe-coordinate measure}
\addcontentsline{toc}{part}{V12: Quartic restoration and the coframe-coordinate measure}

\section{The next deformation equation on the already restored branch}
\label{sec:v12-start}

Sections 1--87 are preserved. The present problem is the second classical
interaction equation
\begin{equation}
 (\widehat{\mathbb S}_{0,a},\widehat{\mathbb S}_{2,a})
 +\tfrac12(\widehat{\mathbb S}_{1,a},\widehat{\mathbb S}_{1,a})=0.
 \label{eq:v12-target}
\end{equation}
The free and cubic coefficients in this equation are the explicitly
\emph{changed} ones in \eqref{eq:v11-master-coefficients}. They are not the
unchanged native cubic excluded by the inherited obstruction. The two
expansion parameters must remain separate: $a$ is spatial mesh, whereas
$\varepsilon$ below counts classical interaction degree. Neither is the
quantum parameter $\hbar$.

The source audit fixes three inputs. V10 already supplies the exact native
stationary quartic \eqref{eq:v10-reduced-jets}; its stationary-elimination
principle is not new here. V11 supplies the quadratic/cubic artifact bank
and the first chart-correct transformation. V9 supplies the actual
sign-averaged Higgs Hodge family and its Euclidean one-loop source module.
The new calculations are the \emph{quartic mesh coefficients of the native
stationary expression}, the second nonlinear coframe transformation, the
mixed signed metric--Higgs restoring contact, and the coordinate Jacobian
needed to express V10's measure in V11's variable $q$. In particular,
a scalar Euclidean determinant is not substituted for a Lorentzian quantum
ESM integral. The signed matter calculation below is a classical coefficient
calculation, in the counterpart explicitly distinguished in
\cite{V6Duality}, \srcid{sec:classical-compatibility}.

All gravity calculations use \eqref{eq:v10-action} with $j=0$, zero
cosmological term and zero independent Holst coefficient. The Palatini
normalization $\chi\ne0$ and V11's affine ADM coframe are fixed. No additional
field or compensator is introduced. Periodic physical pairings are used;
all Fourier panels considered below and all their subset sums are
nonaliased. For quantitative statements every individual momentum has
Euclidean norm at most $\Lambda_*>0$, and $a\Lambda_*\le1/64$.
The Euclidean norm here is a positive norm on labels, not a change in the
Lorentzian contractions. Isotropic temporal sampling and the exact
trigonometric time derivative have V10's declared meaning.

The main positive conclusion is \eqref{eq:v12-target} in the local
coefficient algebra $\R[a]/(a^3)$, with its action and antifield parts both
specified. Fixed-window error bounds are also proved. No quantum BV
Laplacian, ultraviolet graviton loop, or all-frequency finite lattice
symmetry is inferred from that conclusion. The distinction is the one
already required in \cref{sec:v11-measured} and by the ESM programme
\cite{V6Programme}. BV deformation and lattice improvement provide the
established methodological context \cite{V7BarnichHenneaux,V12Symanzik};
the present question concerns their realization for these native operands.

\section{The native quartic artifact bank, including the squared cubic force}
\label{sec:v12-artifacts}

\subsection{Cofactor coefficients and the unshifted fourth word}

Let $E=\mathsf E(q)$ be exactly \eqref{eq:v11-chart},
$T=E^T\otimes I_6$, $d_1=\operatorname{tr}E$, and
$d_2=((\operatorname{tr}E)^2-\operatorname{tr}E^2)/2$.
Besides \eqref{eq:v11-C1}, the second Cartan coefficient is
\begin{equation}
 \boxed{C_2=d_2C_0-d_1(T^TC_0+C_0T)
       +(T^2)^TC_0+C_0T^2+T^TC_0T.}
 \label{eq:v12-C2}
\end{equation}
This is a coefficient of $t^2$, not the second Taylor derivative.
It follows by expanding the actual congruence
\eqref{eq:v10-C-congruence}: $\det(I+tE)=1+td_1+t^2d_2+\cdots$
and $\mathsf T_{I+tE}=I-tT+t^2T^2+\cdots$.
The resulting Cartan form is quadratic in $e$ because its curvature
contraction is the quadratic cofactor \eqref{eq:v11-ell}. Thus there is
no third coframe load at $j=0$: $b_3=0$. In particular, the $\langle b_3,A_1\rangle$
term retained for more general sources in V10 is zero here for a reason,
not by omission.

Write $\Pi_i^{(0)},\Sigma_{ij}^{(0)}$ for the flat cofactor coefficients
and $\Pi_i^{(1)}(q),\Sigma_{ij}^{(1)}(q)$ for their coefficients linear in
$E$. For $r=0,1$ let $\mathscr R_{3,\mathrm{pt}}^{(r)}$ and
$\mathscr R_{3,\mathrm{sh}}^{(r)}$ be the explicit expressions
\eqref{eq:v11-Rpt}--\eqref{eq:v11-Rsh}, with their cofactor coefficients
replaced by these $r$th coefficients. The $r=1$ expressions have a
separate argument $q$. No derivative acts on that argument except those
produced by an explicitly performed integration by parts.

For matrices $X,Y$, put $Z=[X,Y]$ and define
\begin{equation}
 B_{4,\mathrm{com}}(X,Y)=
 \tfrac16[X,[X,Z]]+\tfrac14[Y,[X,Z]]+\tfrac16[Y,[Y,Z]].
 \label{eq:v12-B4-com}
\end{equation}
The required unshifted fourth connection coefficient is
\begin{equation}
 \begin{split}
 \mathscr R_{4,\mathrm{pt}}^{(0)}(A)={}&
 -\tfrac16\sum_i b(\Pi_i^{(0)},[A_i,[A_i,[A_i,A_0]]])\\
 &+\sum_{i<j}b(\Sigma_{ij}^{(0)},
                         B_{4,\mathrm{com}}(A_i,A_j)).
 \end{split}
 \label{eq:v12-R4}
\end{equation}
All these expressions include the physical integral or quadrature when
used as functionals.

\begin{lemma}[Fourth commutator word]
\label[lemma]{lem:v12-B4}
The coefficient \eqref{eq:v12-B4-com} is exactly
\[ [t^4]\log(e^{tX}e^{tY}e^{-tX}e^{-tY}). \]
Consequently \eqref{eq:v12-R4} is the zero-shift coefficient of V10's
actual $\mathscr R_4^{(0)}$, not a replacement curvature invariant.
\end{lemma}
\begin{proof}
Conjugation gives $e^{tX}e^{tY}e^{-tX}=e^{U}$, where
$U=tY+t^2Z+t^3[X,Z]/2+t^4[X,[X,Z]]/6+O(t^5)$.
In $\log(e^Ue^{-tY})$, the term $U-tY$ supplies the first fourth-order
summand. The half commutator supplies $[Y,[X,Z]]/4$.
The two twelfth double commutators each supply $[Y,[Y,Z]]/12$;
all further terms have degree at least five. This proves
\eqref{eq:v12-B4-com}. The electric exponential in
\eqref{eq:v10-full-remainder} has fourth coefficient
$-(\operatorname{ad}A_i)^3A_0/6$. Setting each translated argument to
its unshifted value in the magnetic fourth word proves \eqref{eq:v12-R4}.
These statements are also identities in the free associative algebra,
so they do not depend on a special Lorentz representation.
\end{proof}

\subsection{Three actual force coefficients}

Use the fixed inverse $M=C_0^{-1}$, and abbreviate the previously evaluated
linear stationary coefficients by
\[
 X=A_{1,0}=-M b_{1,0},\qquad Y=A_{1,2}=-M b_{1,2}.
\]
Define the following quadratic-in-$q$ connection covectors:
\begin{equation}
 \begin{aligned}
 G&=b_{2,0}+C_1X,\\
 U&=D_A\mathscr R_{3,\mathrm{pt}}^{(0)}(X),\\
 V&=b_{2,2}+C_1Y+D_A\mathscr R_{3,\mathrm{sh}}^{(0)}(X).
 \end{aligned}
 \label{eq:v12-GUV}
\end{equation}
The symbol $U$ in \eqref{eq:v12-GUV} is a force, not a group link.
These are fully specified local differential polynomials. More explicitly,
$D_A\mathscr R_{3,\mathrm{pt}}^{(0)}$ is obtained by replacing each of
the three occurrences of $A$ in each displayed commutator by a test
coordinate in turn and collecting its coefficient. For the shift term,
if its written density is $P(A,\partial_iA)$, the coefficient in the
same physical pairing is
\begin{equation}
 (D_A\mathscr R_{3,\mathrm{sh}}^{(0)})_\alpha
 =\frac{\partial P}{\partial A_\alpha}
       -\sum_i\partial_i
                  \frac{\partial P}{\partial(\partial_iA_\alpha)}.
 \label{eq:v12-force-rule}
\end{equation}
This replacement-and-adjoint rule is a finite polynomial specification;
there is no inverse external momentum or unspecified nonlinear inverse.
In particular the backward translation in the exact functional gradient
is expanded too. It is not permissible to differentiate just the forward
occurrence and forget its adjoint.

The actual force $G_{2,a}$ of \eqref{eq:v10-stationary-jets} has the expansion
\begin{equation}
 G_{2,a}=G+aU+a^2V+O(a^3).
 \label{eq:v12-force-expansion}
\end{equation}
Indeed $A_{1,a}=X+a^2Y+O(a^4)$, and
$aD_A\mathscr R_3^{(0)}(A_{1,a})
=aD_A\mathscr R_{3,\mathrm{pt}}^{(0)}(X)
 +a^2D_A\mathscr R_{3,\mathrm{sh}}^{(0)}(X)+O(a^3)$.
The $a^2Y$ insertion into this last cubic force first contributes at
$a^3$. This observation fixes the order of every return below.

\begin{theorem}[Explicit native quartic mesh coefficients]
\label[theorem]{thm:v12-quartic-artifacts}
In the declared affine coframe chart define
\begin{align}
 S_{4,0}&=\tfrac12\langle X,C_2X\rangle-\tfrac12\langle G,MG\rangle,
                                                        \label{eq:v12-S40}\\
 \mathcal O_{4,1}&=\mathscr R_{3,\mathrm{pt}}^{(1)}(q;X)
                                      -\langle G,MU\rangle,\nonumber\\
 \mathcal O_{4,2}&=\langle X,C_2Y\rangle
           +\mathscr R_{3,\mathrm{sh}}^{(1)}(q;X)
           +\mathscr R_{4,\mathrm{pt}}^{(0)}(X)\nonumber\\
       &\hspace{9mm}-\langle G,MV\rangle-\tfrac12\langle U,MU\rangle.
                                                        \label{eq:v12-O4}
\end{align}
Then V10's native stationary quartic obeys
\begin{equation}
 \boxed{S_{4,a}^*=S_{4,0}+a\mathcal O_{4,1}
                       +a^2\mathcal O_{4,2}+\mathcal E_{4,a}.}
 \label{eq:v12-quartic-expansion}
\end{equation}
The two artifacts have field/derivative degrees $(4,3)$ and $(4,4)$.
For every fixed momentum multi-index $\alpha$, their polarized remainder
satisfies
\begin{equation}
 |\partial_K^\alpha\mathcal E_{4,a}|
 \le C_\alpha a^3\Lambda_*^{5-|\alpha|}
                              \prod_{r=1}^4\|q_r\|.
 \label{eq:v12-quartic-remainder}
\end{equation}
The constant depends on $\chi$, the fixed tensor norms and derivative
order, not on the mesh, periods or number of sites in a nonaliased panel.
\end{theorem}
\begin{proof}
Use the exact expression \eqref{eq:v10-reduced-jets}, with $b_3=0$.
The first Cartan term gives
$\langle X,C_2X\rangle/2+a^2\langle X,C_2Y\rangle+O(a^4)$.
The coframe-linear cubic word gives
$a\mathscr R_{3,\mathrm{pt}}^{(1)}(q;X)
 +a^2\mathscr R_{3,\mathrm{sh}}^{(1)}(q;X)+O(a^3)$,
and the fourth word gives $a^2\mathscr R_{4,\mathrm{pt}}^{(0)}(X)+O(a^3)$.
Finally, symmetry of the fixed indefinite matrix $M$ gives
\begin{align*}
 -\tfrac12\langle G_{2,a},MG_{2,a}\rangle
 ={}&-\tfrac12\langle G,MG\rangle-a\langle G,MU\rangle\\
 &-a^2\langle G,MV\rangle-\tfrac{a^2}{2}\langle U,MU\rangle+O(a^3).
\end{align*}
These four equalities yield every term in \eqref{eq:v12-O4}.
The stationary return has not been charged twice: it appears once in
V10's exact reduced action and is expanded once here.

The load polynomials $b_{r,0}$ have one derivative, $b_{r,2}$ three,
$X$ one and $Y$ three. The unshifted cubic word has three connection
arguments; its first spatial shift adds one derivative. Thus $G,U,V$
have respectively one, two and three derivatives after substitution.
The stated derivative degrees follow term by term.

For the quantitative remainder, all subset momenta on a four-label panel
are bounded by $4\Lambda_*$. The exact finite coefficient uses only
fixed cofactor polynomials, the constant inverse $M$, the finite fourth
word, finitely many translations, and phase symbols at those subset
momenta. Its dimensionless expression is analytic on a common compact
neighbourhood of $a\Lambda_*=0$. In the geometric part the first omitted
phase power is $a^4\Lambda_*^4$, multiplying a two-derivative vertex.
The cubic and fourth words and their stationary return have their first
omitted powers at $a^3\Lambda_*^3$, again multiplying a two-derivative
vertex. Differentiated Taylor remainders on that compact set prove
\eqref{eq:v12-quartic-remainder}; each momentum derivative removes one
power of $\Lambda_*$. The comparison $a\Lambda_*\le1/64$ absorbs the
smaller $a^4\Lambda_*^6$ term. This is a fixed coefficient estimate,
not a bound on arbitrary high-frequency nonlinear fields.
\end{proof}

The last term of \eqref{eq:v12-O4} is essential even after the cubic
\emph{reduced} action has been corrected. A coframe-only cubic
counterterm leaves the connection equation unchanged and therefore does
not delete $U$ from its force. Erasing the connection cubic before
stationary elimination would instead be a different finite action. The
factor $M$ is indefinite, so $\langle U,MU\rangle$ has no assumed sign.
The checker evaluates a nonzero pure-metric panel of this term; the proof
above does not rely on its sign or on that example.

\section{The second nonlinear coframe transformation is not optional}
\label{sec:v12-chart}

Define the symmetric bilinear map
\begin{equation}
 \mathsf B(q,v)=\mathsf E(q)^T\eta\mathsf E(v)
                    +\mathsf E(v)^T\eta\mathsf E(q).
 \label{eq:v12-B}
\end{equation}
Thus $\mathsf Q(q,q)=\mathsf B(q,q)/2$ and
$Dg(q)[v]=v+\mathsf B(q,v)$, with $g(q)$ fixed by V11.
Write $B_qv=\mathsf B(q,v)$. The operator $I+B_q$ is an invertible
pointwise ten-dimensional matrix on a fixed small coframe ball.

\begin{proposition}[Evaluated second nonlinear gauge coefficient]
\label[proposition]{prop:v12-R2}
For the ordinary Lie derivative and $R_1^{(0)}$ of \eqref{eq:v11-R1}, put
\begin{equation}
 \boxed{R_2^{(0)}(q,c)=\mathcal L_c\mathsf Q(q,q)
                               -\mathsf B(q,R_1^{(0)}(q,c)).}
 \label{eq:v12-R2}
\end{equation}
This local polynomial has degree two in $q$, one in $c$, and one
derivative. The coordinate pullback of the ordinary metric transformation is
\begin{equation}
 (I+B_q)^{-1}\mathcal L_c g(q)
    =R_0c+R_1^{(0)}(q,c)+R_2^{(0)}(q,c)+O(q^3c).
 \label{eq:v12-pullback}
\end{equation}
In particular, the quartic ghost/antifield coefficient is fixed without
an inverse wave operator.
\end{proposition}
\begin{proof}
Multiply $R_0+R_1+R_2$ by $I+B_q$. At field degrees zero, one and two,
the equations are respectively
$R_0=\mathcal L_c\eta$,
$R_1+B_qR_0=\mathcal L_cq$, and
$R_2+B_qR_1=\mathcal L_c\mathsf Q(q,q)$.
The middle one is exactly V11's $R_1$, and the last one is
\eqref{eq:v12-R2}. Every operation is multiplication or an ordinary
first derivative. Local invertibility follows by a finite-dimensional
Neumann series for $B_q$, with radius independent of the lattice volume.
The expansion is therefore the Taylor coefficient of an actual local
coordinate map, not a guessed transformation.
\end{proof}

For example take a constant spatial metric tangent
$q_{11}=u$, $q_{22}=v$, all other entries zero, and a shear
$c^2=e^{ikx^1}$. The exact coefficient of this shear in the off-diagonal
transformation is
\[
 R(q,c)_{12}=ik e^{ikx^1}
          \frac{(1+v/2)^2}{1+(u+v)/4}.
\]
Consequently
\begin{equation}
 R_2^{(0)}(q,c)_{12}
         =\tfrac1{16}(u-v)^2 ik e^{ikx^1}.
 \label{eq:v12-shear}
\end{equation}
It is generally nonzero, even though the source coframe is constant.
The coefficient follows directly by differentiating the spatial matrix
square $g_{\rm sp}=(I+q_{\rm sp}/2)^2$; it is not a curvature term.

\begin{lemma}[The next two-ghost identity]
\label[lemma]{lem:v12-closure}
For commuting vector parameters define $[c,d]^\mu=c^\rho\partial_\rho d^\mu
-d^\rho\partial_\rho c^\mu$. With the convention that the transformation
commutator is $D_qR_c[R_d]-D_qR_d[R_c]$, the degree-one-in-$q$ identity is
\begin{align}
 &R_1^{(0)}(R_1^{(0)}(q,d),c)-R_1^{(0)}(R_1^{(0)}(q,c),d)\nonumber\\
 &\quad+D_qR_2^{(0)}(q,c)[R_0d]
              -D_qR_2^{(0)}(q,d)[R_0c]
            =R_1^{(0)}(q,[c,d]).
 \label{eq:v12-two-ghost}
\end{align}
The derivative in this expression is explicitly
\begin{equation}
 D_qR_2^{(0)}(q,c)[v]
  =\mathcal L_c\mathsf B(q,v)-\mathsf B(v,R_1^{(0)}(q,c))
                            -\mathsf B(q,R_1^{(0)}(v,c)).
 \label{eq:v12-DR2}
\end{equation}
\end{lemma}
\begin{proof}
The ordinary tensor Lie derivatives obey their commutator identity.
The invertible pointwise map $g(q)$ transports it exactly; the second
derivatives of $g$ cancel in the antisymmetric difference. Expand that
identity through field degree one. The degree-zero equation is V11's
first two-ghost relation, and the next equation is
\eqref{eq:v12-two-ghost}. Differentiation of \eqref{eq:v12-R2} gives
\eqref{eq:v12-DR2}. This proof uses ordinary derivations on the local
polynomial algebra; it does not assign a product rule to the native
finite phase derivative. On the nonaliased finite panels the displayed
polynomials have their stated Fourier values.
\end{proof}

\section{The restored quartic master equation and its finite-window error}
\label{sec:v12-master}

Define the additional \emph{action} correction
\begin{equation}
 \mathcal C_{4,a}=-a\mathcal O_{4,1}-a^2\mathcal O_{4,2},\qquad
 \widehat S_{4,a}=S_{4,a}^*+\mathcal C_{4,a},
 \label{eq:v12-counterterm}
\end{equation}
and retain V11's $\widehat S_{2,a},\widehat S_{3,a},\widehat R_{0,a}$
unchanged. The new interaction coefficient in the minimal classical
master functional is
\begin{equation}
 \boxed{\widehat{\mathbb S}_{2,a}
       =\widehat S_{4,a}+\langle q^*,R_2^{(0)}(q,c)\rangle.}
 \label{eq:v12-master2}
\end{equation}
The ghost coefficient $\langle c^*,c\cdot\partial c\rangle$ remains in
$\widehat{\mathbb S}_{1,a}$. No new field-dependent vector-field bracket
is needed for this pointwise reparametrization. This is a statement about
the ordinary local metric symmetry transported to $q$, not about a
nonlinear finite lattice group.

\begin{theorem}[Quartic classical restoration through second mesh degree]
\label[theorem]{thm:v12-master}
In the local differential-polynomial functional algebra modulo total
derivatives, with coefficients in $\R[a]/(a^3)$,
\begin{equation}
 \boxed{(\widehat{\mathbb S}_{0,a},\widehat{\mathbb S}_{2,a})
       +\tfrac12(\widehat{\mathbb S}_{1,a},\widehat{\mathbb S}_{1,a})=0.}
 \label{eq:v12-master-result}
\end{equation}
Together with V11 this is the classical master equation for
$\widehat{\mathbb S}_{0,a}+\varepsilon\widehat{\mathbb S}_{1,a}
+\varepsilon^2\widehat{\mathbb S}_{2,a}$ through interaction degree two.
At finite $a$, the polarized action-side quartic residual obeys
\begin{align}
 \mathcal N_{4,a}={}&D\widehat S_{4,a}[\widehat R_{0,a}c]
       +D\widehat S_{3,a}[R_1^{(0)}(q,c)]
       +D\widehat S_{2,a}[R_2^{(0)}(q,c)],\nonumber\\
 |\partial_K^\alpha\mathcal N_{4,a}|
 &\le C_\alpha a^3\Lambda_*^{6-|\alpha|}
                           \|q_1\|\|q_2\|\|q_3\|\|c\|.
 \label{eq:v12-noether-bound}
\end{align}
The new metric-antifield/two-ghost residual is bounded by
$C_\alpha a^4\Lambda_*^{6-|\alpha|}$ times its argument norms.
All constants have the fixed-order, fixed-window and period-independent
scope of \cref{thm:v12-quartic-artifacts}.
\end{theorem}
\begin{proof}
The zero-mesh action is identified by the native normal form as in V11:
it is the boundary-completed Palatini action after stationary elimination
of the same independent connection. Its quartic coefficient is precisely
\eqref{eq:v12-S40}, by V10's elimination formula with zero mesh. Ordinary
pullback of its differential four-form changes it by a total derivative.
The connection-variation term vanishes at its stationary section, and the
periodic pairing kills the total derivative. Expand this exact reduced
invariance to third field degree in its variation. Using the coordinate
calculation in \cref{prop:v12-R2} gives
\begin{equation}
 DS_{4,0}[R_0c]+DS_{3,0}[R_1^{(0)}(q,c)]
                         +DS_{2,0}[R_2^{(0)}(q,c)]=0.
 \label{eq:v12-zero-noether}
\end{equation}
This derivation identifies the action before using its symmetry and hence
does not assume a finite-cutoff Ward identity to prove itself.

The antifield component is \cref{lem:v12-closure}, with odd parameters
and V11's antibracket ordering. The ghost-antifield component is the
Jacobi identity for ordinary vector fields. These components exhaust the
minimal classical bracket at this interaction degree. Thus
\eqref{eq:v12-master-result} holds at zero mesh. The evaluated expansions
give
$\widehat S_{2,a}=S_{2,0}+O(a^4)$,
$\widehat S_{3,a}=S_{3,0}+O(a^3)$,
$\widehat S_{4,a}=S_{4,0}+O(a^3)$, and
$\widehat R_{0,a}=R_0+O(a^4)$.
Their classes in $\R[a]/(a^3)$ are exactly the zero-mesh coefficients,
proving the stated formal identity. This uses the explicit return terms
in \eqref{eq:v12-O4}, not a declaration that a closed local defect must
have a primitive.

For the finite estimate subtract \eqref{eq:v12-zero-noether} before
bounding. The quartic and cubic action remainders have size
$a^3\Lambda_*^5$; each transformation has one derivative and hence
costs at most $C\Lambda_*$. The quadratic remainder has size
$a^4\Lambda_*^6$ and contributes $C a^4\Lambda_*^7$.
The difference $\widehat R_{0,a}-R_0$ has size $C a^4\Lambda_*^5$;
multiplication by the two-derivative quartic coefficient gives the same
smaller $a^4\Lambda_*^7$ bound. Absorb it using $a\Lambda_*\le1/64$.
Leibniz differentiation proves all fixed momentum derivative estimates.
For the two-ghost row \eqref{eq:v12-two-ghost} cancels exactly except
that $R_0$ is replaced by $\widehat R_{0,a}$ in its two occurrences;
$D_qR_2$ has one derivative, giving $a^4\Lambda_*^6$.
The ghost Jacobi row has no mesh change in these local polynomials.
\end{proof}

As an exact normalization test, the rational four-panel calculation in
the accompanying script gives the three terms of
\eqref{eq:v12-zero-noether} as
$-689i/8$, $773i/8$, and $-21i/2$, at $\chi=1$.
Their sum is zero and the last term is nonzero. The script specifies all
momenta and matrix polarizations. This checks the actual Cartan/cofactor
normalization and detects dropping $R_2$; the proof of the general identity
is the preceding variational and coordinate argument.

\subsection{The missing ghost seagull is an explicit operand}

For the ordinary harmonic functional
\[
 F_\mu(q)=\partial^\nu(q_{\mu\nu}-\tfrac12\eta_{\mu\nu}
                                      \operatorname{tr}_\eta q),
\]
introduce the four antighosts $\bar c^\mu$ and auxiliary fields $b^\mu$,
with $s\bar c=b$, $sb=0$. The gauge fermion
$\Psi=\langle\bar c,F(q)+\alpha b/2\rangle$ gives the ghost row
\begin{equation}
 -\langle\bar c,F(\widehat R_{0,a}c)
       +\varepsilon F(R_1^{(0)}(q,c))
       +\varepsilon^2F(R_2^{(0)}(q,c))\rangle.
 \label{eq:v12-ghost-seagull}
\end{equation}
Its free zero-mesh operator is $F R_0c=\Box_\eta c$.
The last term in \eqref{eq:v12-ghost-seagull} is the two-metric ghost
contact needed beside the quartic graviton action in a two-background
one-loop calculation. It has two derivatives and is local, since
\eqref{eq:v12-R2} is explicit. Its sign follows from
$s(\bar c F)=bF-\bar c\,sF$.
The nonminimal pair is contractible, and canonical gauge fixing preserves
\eqref{eq:v12-master-result} at its declared orders. These statements do
not choose an inverse of $\Box_\eta$, fix the homogeneous modes, or prove
positivity of a gravitational reference. They supply vertices, not an
already constructed Lorentzian quantum integral.

\section{A coupled quartic metric--Higgs repair on the same coframe chart}
\label{sec:v12-Higgs}

The structural common-action manuscript explicitly distinguishes the
Lorentzian signed matter action from its positive Euclidean representative
\cite{V6Duality}. For the following \emph{classical} mixed calculation,
use that signed counterpart of V9's sign-averaged Hodge writer, with
internal links $I$, zero internal gauge/Yukawa couplings, and a fixed
mass coefficient $m^2$. Thus
\begin{equation}
 \begin{split}
 S_{H,a}^{L}(q,\phi)={}&\frac1{32}\sum_{\epsilon\in\{+,-\}^4}
 \int_a W^{\mu\nu}(q)
       D_\mu^{\epsilon_\mu}\phi\cdot D_\nu^{\epsilon_\nu}\phi\\
 &+\frac{m^2}{2}\int_a\rho(q)|\phi|^2,\qquad
 W(q)=\rho(q)g(q)^{-1},\quad \rho(q)=\det e(q)>0.
 \end{split}
 \label{eq:v12-signed-Hodge}
\end{equation}
Here $\phi\in\R^4$ is the real Higgs packet and $\int_a$ includes the
physical weights. The four scalar differences, including its temporal
differences, are exactly the isotropic Hodge choice of V9; the gravity
load still has V11's spectral time convention. Both are specified on one
periodic record. Their zero-mesh coefficients use ordinary derivatives.
No positivity, contour continuation, or use of V9's Euclidean loop theorem
for \eqref{eq:v12-signed-Hodge} is asserted. Overall Lorentzian action-sign
conventions can be transferred consistently to all terms; the identity
below uses the displayed one.

At successive powers of $q$, the actual density and inverse metric give
\begin{align}
 \rho_0&=1,\quad \rho_1=d_1,\quad \rho_2=d_2,\nonumber\\
 W_0&=\eta,\qquad W_1=d_1\eta-\eta q\eta,\nonumber\\
 W_2&=d_2\eta-d_1\eta q\eta+\eta q\eta q\eta
                            -\eta\mathsf Q(q,q)\eta.
 \label{eq:v12-Hodge-qjets}
\end{align}
The final term is the metric/coframe chart contact. Omitting it would
mix the exact metric coordinate of V8 with the affine coordinate used
in the present gravitational action.

For a symmetric matrix coefficient $W$ define the explicit bilinear
scalar functional
\begin{equation}
 \begin{split}
 \mathcal H_2(W;\phi)=\frac12\int\sum_{\mu,\nu}W^{\mu\nu}
 \bigl\{&\tfrac16(\partial_\mu^3\phi\cdot\partial_\nu\phi
                    +\partial_\mu\phi\cdot\partial_\nu^3\phi)\\
       &+\tfrac14\delta_{\mu\nu}|\partial_\mu^2\phi|^2\bigr\}.
 \end{split}
 \label{eq:v12-H2}
\end{equation}
No integration by parts has moved derivatives off the coefficient $W$ in
this definition. It therefore fixes every mixed geometric contact.

\begin{proposition}[Evaluated Higgs artifacts through two metric marks]
\label[proposition]{prop:v12-Higgs-artifacts}
Let $S^H_{r+2,a}$ denote the part of \eqref{eq:v12-signed-Hodge} of
metric degree $r$ and scalar degree two, $r=0,1,2$. Then
\begin{equation}
 S^H_{r+2,a}=S^H_{r+2,0}+a^2\mathcal H_2(W_r;\phi)
                                        +\mathcal E^H_{r+2,a},
 \label{eq:v12-Higgs-artifacts}
\end{equation}
where the polarized remainder symbols obey
$|\partial_K^\alpha\mathcal E^H_{r+2,a}|
 \le C_\alpha a^4\Lambda_*^{6-|\alpha|}$ times their argument norms.
The bound is independent of $m$ because the mass term has no difference
operator and no lattice artifact on these nonaliased panels.
\end{proposition}
\begin{proof}
Expand the \emph{actual} differences:
$D_\mu^\pm=\partial_\mu\pm a\partial_\mu^2/2
                    +a^2\partial_\mu^3/6+O(a^3)$.
For distinct indices the signs are independent, so the average product
is the product of the centered averages. For equal indices the same sign
is used twice. Its variance adds the term
$a^2|\partial_\mu^2\phi|^2/4$, in addition to the two cross terms with
$\partial_\mu^3/6$. This gives exactly \eqref{eq:v12-H2}.
Each sign-averaged product is even in $a$, so the first omitted term has
four mesh powers and six derivatives. The metric coefficient remains at
its root site; ordinary expansion of $\det e$ and $g^{-1}$ gives
\eqref{eq:v12-Hodge-qjets}. Differentiated analytic Taylor estimates on
the nonaliased Fourier window prove the bound. The mass density is the
same finite polynomial coefficient as its zero-mesh expression, which
explains its absence from the error.
\end{proof}

Set $\widehat S^H_{r+2,a}=S^H_{r+2,a}-a^2\mathcal H_2(W_r;\phi)$.
Use the scalar transformation $s\phi=\varepsilon\mathcal L_c\phi$
with $\mathcal L_c\phi=c^\mu\partial_\mu\phi$; it has no $q$-dependent
second coefficient in these variables. The total-field grading can be
obtained by evaluating the unscaled action at $e(\varepsilon q)$ and
$\phi_{\rm phys}=\varepsilon\phi$, dividing by $\varepsilon^2$, and
using the physical vector parameter $\varepsilon c$. It gives exactly
the stated scalar transformation order. This is a specified corrected
transformation, not the unchanged finite difference transformation.

\begin{theorem}[Coupled gravity--Higgs quartic local master coefficient]
\label[theorem]{thm:v12-coupled}
In \eqref{eq:v11-master-coefficients} and \eqref{eq:v12-master2}, add
$\widehat S^H_{2,a}$ to $\widehat{\mathbb S}_{0,a}$,
$\widehat S^H_{3,a}+\langle\phi^*,\mathcal L_c\phi\rangle$ to
$\widehat{\mathbb S}_{1,a}$, and $\widehat S^H_{4,a}$ to
$\widehat{\mathbb S}_{2,a}$. The three classical master coefficients
through interaction degree two vanish in $\R[a]/(a^3)$.
The mixed quartic antifield-free equation at zero mesh is explicitly
\begin{equation}
 \boxed{D_q S^H_{4,0}[R_0c]+D_q S^H_{3,0}[R_1^{(0)}(q,c)]
                   +D_\phi S^H_{3,0}[\mathcal L_c\phi]=0.}
 \label{eq:v12-mixed-noether}
\end{equation}
Its finite-window counterpart, with the hats and $\widehat R_{0,a}$,
has bound $C_{\alpha,M}a^4\Lambda_*^{7-|\alpha|}$ times its four argument
norms when $|m|\le M\Lambda_*$. The pure gravitational quartic bound
remains \eqref{eq:v12-noether-bound}. No internal-gauge or spin/fermion
quartic identity is added to this statement.
\end{theorem}
\begin{proof}
The zero-mesh coefficient of \eqref{eq:v12-signed-Hodge} is the scalar
density $\rho(g)(g^{\mu\nu}\partial_\mu\phi\cdot\partial_\nu\phi
+m^2|\phi|^2)/2$ evaluated at the \emph{same} $g(q)$.
Under the ordinary tensor Lie derivative and $\delta\phi=\mathcal L_c\phi$
its variation is the divergence of $c$ times this density, including the
usual transport term. The periodic integral is zero. Expanding the
identity to one and two metric marks gives the cubic matter equation
and \eqref{eq:v12-mixed-noether}. The scalar Lie transformations close
with the same vector-field bracket. Thus the matter-antifield and ghost
rows obey the corresponding minimal BV identities; the metric rows are
those already proved. There is no missing $D_qS^H_{2,0}[R_2]$ term,
since $S^H_{2,0}$ is independent of $q$.

The two-density coefficients in \eqref{eq:v12-Hodge-qjets} retain all
metric contacts in that expansion, in particular
$-\eta\mathsf Q\eta$. By \cref{prop:v12-Higgs-artifacts}, the hatted
coefficients differ from them only at $a^4$. The corrected free metric
generator differs from $R_0$ at $a^4$, too. This proves the formal
statement through $a^2$. For the finite estimate, each scalar or metric
Lie transformation supplies one momentum derivative. The action
remainders therefore cost $a^4\Lambda_*^7$; the generator discrepancy
acts on a vertex of size $\Lambda_*^2+m^2$, bounded by
$(1+M^2)\Lambda_*^2$. Differentiation yields the asserted bound.
The argument concerns the declared signed classical coefficients and
makes no identification of quantum measures across signatures.
\end{proof}

A quartic Higgs potential may be retained with its own coupling degree;
its zero-metric quartic term is annihilated by the free differential at
this interaction order. Its higher-metric contacts and its nonlinear
coupling count must be included at the next required degrees. It is not
silently used to strengthen \cref{thm:v12-coupled} into the full SM
master equation. Similarly, V9's current/composite-source module remains
available with its original internal-symmetry and massive-reference scope;
the new derivative counterterms have not been shown to preserve that
entire module after ultraviolet integration.

\section{The coframe-coordinate Jacobian must accompany the nonlinear ghosts}
\label{sec:v12-measure}

V10's leading metric-coordinate density was expressed with respect to
$d^{10}g$; V11 introduced the nonlinear coordinate $q$ but did not claim
a quantum BV identity in that coordinate. Holding $q$ as an external
source needs only substitution. Integrating $q$, or evaluating a measured
divergence on its antifields, additionally needs the actual Jacobian.
It is now evaluated; this is transport of the old measure, not a new
choice of reference measure.

Write the affine coframe in blocks as
\begin{equation}
 e(q)=\begin{pmatrix} n&0\\b&S\end{pmatrix},\qquad
 n=1-q_{00}/2>0,\quad b_i=q_{0i},\quad S=I_3+q_{\rm sp}/2\succ0.
 \label{eq:v12-block-chart}
\end{equation}
Let $s_1,s_2,s_3>0$ be the eigenvalues of $S$.

\begin{theorem}[Exact pointwise metric-to-coframe-coordinate density]
\label[theorem]{thm:v12-Jacobian}
In independent symmetric matrix coordinates, normalized so that
$Dg(0)=I$, one has
\begin{equation}
 \boxed{J_g(q):=\det Dg(q)
       =n(\det S)^2\prod_{i<j}\frac{s_i+s_j}{2}>0.}
 \label{eq:v12-Jacobian}
\end{equation}
Consequently V10's ultralocal factor in the $q$ coordinates is
\begin{equation}
 \boxed{|\det g(q)|^{-7/2}J_g(q)
       =n^{-6}(\det S)^{-5}\prod_{i<j}\frac{s_i+s_j}{2}.}
 \label{eq:v12-q-density}
\end{equation}
This factor is independent of the shift coordinate $b$. The complete
leading auxiliary negative-log-modulus in $q$ is therefore
\begin{equation}
 \begin{aligned}
 \mathcal M^{q}_{a,1}={}&
 \sum_x\left\{6\log n_x+5\log\det S_x
                -\sum_{i<j}\log\frac{s_{i,x}+s_{j,x}}2\right\}\\
 &+\tfrac12\operatorname{Tr}\log(I+K_a)
             -\sum_{x,i}\log j(aA_{*,x,i}),
 \end{aligned}
 \label{eq:v12-measure-q}
\end{equation}
with precisely the same $K_a$, stationary section and Haar factors as V10.
\end{theorem}
\begin{proof}
The blocks of $g=e^T\eta e$ are
$g_{\rm sp}=S^2$, $g_{0\rm sp}=Sb$, and $g_{00}=-n^2+b^Tb$.
Order the input variables as $(q_{\rm sp},b,q_{00})$ and the outputs as
$(g_{\rm sp},g_{0\rm sp},g_{00})$. The Jacobian is block triangular.
The spatial block maps a symmetric variation $v$ to $(Sv+vS)/2$.
Orthogonally diagonalizing $S$ gives its eigenvalues $s_i$ on the three
diagonal variations and $(s_i+s_j)/2$ on the three off-diagonal variations.
The shift block has determinant $\det S$, and
$\partial g_{00}/\partial q_{00}=n$. Multiplication proves
\eqref{eq:v12-Jacobian}. Since $|\det g|=n^2(\det S)^2$,
\eqref{eq:v12-q-density} follows. Multiplying V10's complete leading
amplitude by $\prod_x J_g(q_x)$ subtracts $\sum_x\log J_g$ from its
negative logarithm; combining this with \eqref{eq:v10-main-amplitude}
gives \eqref{eq:v12-measure-q}. The calculation introduces no determinant
of a differential operator and no additional Lorentz ghost.
\end{proof}

This density has a finite pointwise source expansion on a common ball.
For $h=q_{\rm sp}$, its Jacobian alone starts as
\begin{equation}
 \log J_g(q)=-\tfrac12q_{00}+\tfrac32\operatorname{tr}h
    -\tfrac18q_{00}^2-\tfrac9{32}\operatorname{tr}h^2
                         -\tfrac1{32}(\operatorname{tr}h)^2+O(q^3).
 \label{eq:v12-Jacobian-twojet}
\end{equation}
At every prescribed higher degree it is specified by the finite trace
formula $\operatorname{Tr}_{\operatorname{Sym}_4}\log(I+B_q)$.
There is no singular eigenvector choice at repeated $s_i$: both this
matrix formula and the symmetric product in \eqref{eq:v12-Jacobian}
are analytic on the positive chart. The sum over sites is a genuine
local one-loop matching contact. Its coefficient per physical volume
can scale as $a^{-4}$; no smallness or cancellation is claimed.

\begin{proposition}[Measured divergence under the actual coordinate change]
\label[proposition]{prop:v12-divergence}
Let $\rho_g(g)$ be any nonvanishing smooth finite reference density on
these even metric variables, including the selected factors when they
are assigned to the reference. For any smooth finite vector field
$R_g$, put $R_q=(Dg)^{-1}R_g\circ g$ and
$\rho_q(q)=\rho_g(g(q))\prod_xJ_g(q_x)$. Then
\begin{equation}
 \boxed{\operatorname{div}_{\rho_q}R_q
          =(\operatorname{div}_{\rho_g}R_g)\circ g.}
 \label{eq:v12-divergence}
\end{equation}
Omitting $J_g$ would change the left side by
$-R_q\sum_x\log J_g(q_x)$.
\end{proposition}
\begin{proof}
The ordinary divergence transforms as
$\operatorname{div}_qR_q=(\operatorname{div}_gR_g)\circ g-R_q\log\det Dg$.
This follows by differentiating the change-of-variables identity for the
flow, or directly by the determinant derivative. Add
$R_q\log\rho_q$ to both sides. The two Jacobian derivatives cancel and
the density chain rule gives \eqref{eq:v12-divergence}. This is a
finite-dimensional identity and requires no lattice Leibniz rule.
For antifield-linear functions, this divergence is the even-metric
operand of the measured BV Laplacian in the corresponding convention.
It does not prove that the operand vanishes, or that a complete quantum
master equation has been solved.
\end{proof}

If density factors are instead moved into the action, their logarithms
must be counted once there and not again in the reference BV operator.
The Lorentzian amplitude convention of V10 is unchanged: the real
quantity \eqref{eq:v12-measure-q} enters $-i\hbar\log Z$ multiplied by
$i\hbar$. Neither the chart determinant nor the classical ghost identity
cancels this amplitude automatically.

\section{Measured contacts, uniform kernels, and the remaining ultraviolet task}
\label{sec:v12-transfer}

The new pure-gravity counterterm \eqref{eq:v12-counterterm}, like V11's,
is independent of the integrated auxiliary connection. Thus at retained
$q$ and Cartan source $j$ the exact insertion is
\begin{equation}
 Z_a^{(12)}[q,j]=\exp\!\left\{\frac{i}{\hbar}
                   (\mathcal C_a(q)+\mathcal C_{4,a}(q))\right\}
                   Z_a[q,j].
 \label{eq:v12-measured-insertion}
\end{equation}
Every $j$ derivative of the connected logarithm and V10's stationary
section and auxiliary Hessian are unchanged. Its $q$ derivatives acquire
the actual derivatives of the single correcting functional, including
all four-metric contacts. When $q$ itself is integrated, the coordinate
factor \eqref{eq:v12-Jacobian} is additionally present. This is precisely
the distinction between an external source reparametrization and a
change of an integrated variable.

The Higgs corrections in \cref{prop:v12-Higgs-artifacts} are also
independent of that spacetime connection. They do change the scalar
quadratic vertices when the scalar is integrated, and must be kept as
insertions in its measured calculus. Their effect on loop counterterms
has not been declared negligible. Mixed elementary-source returns, such
as those of \eqref{eq:v8-source-output}, must be differentiated with the
changed vertices. The coordinate chain rule alone does not renormalize
those new loop coefficients.

For the new quartic remainder, multiply all four legs by a common smooth
window supported inside the declared momentum region, and let
$\mathcal K_{4,a}$ be the inverse Fourier kernel in its three relative
four-dimensional coordinates. Then for every fixed $s\ge0$,
\begin{equation}
 \boxed{a^{12}\sum_{x_1,x_2,x_3}
  [1+\Lambda_*\textstyle\sum_{r=1}^3d_{\rm per}(x_r,0)]^s
  \|\mathcal K_{4,a}(x_1,x_2,x_3)\|
       \le C_s a^3\Lambda_*^5.}
 \label{eq:v12-kernel}
\end{equation}
The quartic Noether kernel has the corresponding bound
$C_s a^3\Lambda_*^6$. The mixed Higgs remainder has bound
$C_s a^4\Lambda_*^6$, and its Noether row the scaled bound stated in
\cref{thm:v12-coupled}. To prove these assertions, rescale the twelve
independent momenta, use more than $s+12$ of the differentiated symbol
bounds, integrate by parts in the inverse Fourier transform, then
periodize. The infinite-space inverse transform is dominated by a
constant times $(1+\Lambda_*\sum_r|x_r|)^{-s-13}$ with the appropriate
normalization factor $\Lambda_*^{12}$. Its weighted integral and sampled
sum are bounded on $a\Lambda_*\le1/64$; torus distance is bounded by the
distance of each lift. This proves \eqref{eq:v12-kernel} without a total
volume factor. It also gives a multilinear functional bound with one
argument in $L^1$ and the other three in $L^\infty$.

The vertices now identified for a propagating gravitational calculation
are the corrected cubic/quartic metric actions, the one- and two-metric
ghost vertices in \eqref{eq:v12-ghost-seagull}, the signed mixed Higgs
contacts, and the chart-correct auxiliary amplitude. Constructing a
common gauge-fixed reference covariance and performing that complete
loop calculation remain separate. In particular,
\[
 \sum_{j\le J}(a\Lambda_j)^3=O(1)
       \quad\hbox{when}\quad \Lambda_j=2^j\Lambda_0,
                              \ \Lambda_J\asymp a^{-1},
\]
so the proved window estimate is not a vanishing all-shell Ward budget.
No number of restatements of fixed-window consistency removes this issue.
The next task is the complete propagating metric/ghost shell, retaining
its quartic action and ghost contacts and the local density assignment,
or a full-zone matching estimate for that same signed family. It is not
another unchanged-cubic ghost ansatz or another isolated scalar determinant.

\begingroup
\renewcommand{\refname}{Additional methodological reference for V12}
\begin{thebibliography}{99}
\bibitem[37]{V12Symanzik}
K.~Symanzik, \emph{Continuum limit and improved action in lattice theories.
I. Principles and $\varphi^4$ theory}, Nuclear Physics B \textbf{226}
(1983), 187--204, doi:10.1016/0550-3213(83)90468-6, DESY-83-016.
Context for finite-mesh irrelevant-operator improvement; no native
ESM symmetry or quantum shell theorem is imported from this reference.
\end{thebibliography}
\endgroup


\section{Final ledgers for V12}
\label{sec:v12-ledgers}

\subsection*{Proved}
\addcontentsline{toc}{subsection}{V12: Proved}
The actual native quartic has the explicit mesh coefficients
\eqref{eq:v12-O4}, including the squared cubic force and all first
translation contributions. The second chart transformation
\eqref{eq:v12-R2} satisfies the field-linear two-ghost identity and
completes the quartic classical master coefficient in $\R[a]/(a^3)$.
The corresponding finite-window residual and rooted kernel are bounded
in \eqref{eq:v12-noether-bound} and \eqref{eq:v12-kernel}.
The signed Hodge counterpart has the explicit mixed metric--Higgs
restoring contact \eqref{eq:v12-H2} and its coupled quartic identity.
The actual pointwise metric-to-$q$ Jacobian is
\eqref{eq:v12-Jacobian}; its transport of V10's leading amplitude is
\eqref{eq:v12-measure-q}, and its measured-divergence contribution is
retained exactly. The harmonic ghost seagull is an explicitly identified
vertex, not an assumed ghost determinant.

\subsection*{Inherited}
\addcontentsline{toc}{subsection}{V12: Inherited}
V1--V11 and all their mathematical labels are unchanged. V10's native
stationary quartic, its classical Cartan inverse and declared leading
Lorentz-Haar amplitude are inputs. V11's cubic obstruction, evaluated
artifact bank and first nonlinear chart transformation are not reproved
as new results. The sign-averaged Hodge writer and source contacts are
inherited from V7--V9; their positive Euclidean loop theorems retain
that reference, mass and symmetry scope. The structural manuscript
supplies its common classical matter carrier and its explicit distinction
between Euclidean and Lorentzian signed representatives. The general
stationary, Schur, measured insertion and BV deformation principles are
inherited tools.

\subsection*{Literature input}
\addcontentsline{toc}{subsection}{V12: Literature input}
Local BV consistency conditions and lattice irrelevant-operator
improvement are established frameworks
\cite{V7BarnichHenneaux,V12Symanzik}. They do not supply the present
native stationary-return coefficients, the affine-ADM chart contacts or
the selected measure's Jacobian. No first construction of perturbative
gravity, scalar gravity, symmetry-restoring counterterms or BV geometry
is claimed.

\subsection*{Conditional}
\addcontentsline{toc}{subsection}{V12: Conditional}
A complete ESM quantum use requires a common graviton/ghost/matter
reference and source convention, an actual measure-compatible quantum
BV operator, and ultraviolet estimates for the newly changed action.
Classical zero-mesh matching does not determine invariant physical EFT
couplings, prove microscopic matching of the action changes, or fix
chiral phases. The internal gauge and spin sectors are not added to
\cref{thm:v12-coupled} without their own transformed vertices. Density
assignment to action versus reference must avoid double counting. The
coordinate-divergence identity is exact, but its value has not been shown
to cancel the full one-loop gravitational anomaly or defect.

\subsection*{Open}
\addcontentsline{toc}{subsection}{V12: Open}
The quartic local classical step has now been carried out at the stated
mesh order. The next load-bearing problem is one complete propagating
graviton plus diffeomorphism-ghost shell on a declared reference,
including the quartic ghost contact, coframe density, mixed source
contacts and full-zone local subtraction. An alternative is a uniform
full-zone symmetry-restoring/matching theorem for that same family.
Higher interaction degrees of the formal tree expansion, a combined
internal/Lorentz/diffeomorphism quantum BV closure, massless and chiral
transport, and the full invariant filtered ESM trajectory remain open.
No infinite-series convergence, positive nonperturbative gravity measure,
asymptotic safety or finite-parameter ultraviolet completion follows.

\begin{center}\small
\begin{tabular}{>{\raggedright\arraybackslash}p{.22\textwidth}>{\raggedright\arraybackslash}p{.69\textwidth}}
\toprule
Variable & Scope in V12\\
\midrule
Volume & Fourier coefficients are per physical volume; rooted bounds have
no total-volume factor. The pointwise Jacobian is exact, but its extensive
logarithm is not uniformly bounded.\\
Mesh & $a\Lambda_*\le1/64$; all subset momenta are nonaliased.
$\R[a]/(a^3)$ is a local coefficient quotient, not an ultraviolet loop
estimate.\\
Orders & Classical action through field degree four and master functional
through interaction degree two. Each fixed momentum-derivative and moment
order is controlled; no unbounded-order uniformity.\\
Matter & Four real Higgs components, internal links $I$, no internal gauge
or Yukawa interaction in this calculation. Signed classical reference;
the mixed Noether bound uses $|m|\le M\Lambda_*$.\\
Measure & V10's auxiliary reference, with the actual $g$-to-$q$ Jacobian
when $q$ is integrated. No new propagating one-loop or quantum BV identity.\\
Iteration & Fixed-window tree error only. A common propagating reference,
full-zone subtraction and the complete ESM invariant class remain open.\\
\bottomrule
\end{tabular}
\end{center}

\subsection*{Verification record}
The standalone \srcid{check_ESM_quartic_BV_V12.py} checks exact associative
fourth-word algebra, generic stationary mesh coefficients, the native
quartic Noether panel, the next two-ghost relation, a nonzero shear $R_2$,
the mixed Higgs/Hodge identities, and the ten-coordinate Jacobian.
Separately labelled numerical tests retain all translated occurrences
and all $2+2$ stationary partitions; their mesh-halving ratios test the
quartic remainder. These checks audit signs and contacts, not the uniform
bounds proved above. No new proof-assistant formalization is asserted.



\clearpage
\part*{V13: A measured propagating metric--ghost shell and its ultraviolet boundary}
\addcontentsline{toc}{part}{V13: Propagating metric--ghost shell and ultraviolet boundary}

\section{The next operand is a reference, not another classical identity}
\label{sec:v13-start}

Sections 1--95 and their labels are preserved. V12 has supplied the cubic
and quartic action, the one- and two-metric ghost vertices, and the
metric-to-coframe Jacobian. Its last open problem is not solved by
repeating the classical equation modulo $a^3$. This installment constructs
a finite propagating reference for those particular coefficients and
computes its complete one-loop, two-retained-metric coefficient. It proves
a subtracted shell bound and a cutoff limit for an explicitly specified
interior covariance-cutoff realization. An exact momentum-boundary test
then identifies what prevents that result from being promoted, without
another theorem, to the unchanged full-zone regulator.

The result integrates ten coframe-coordinate fluctuations and four
complex Grassmann ghost pairs. It is not the two-TT-polarization reference
of V6, and its time symbol is the spectral one retained in V10--V12,
not the different nearest-time-neighbour symbol of V6. The four algebraic
nonminimal fields are integrated with a stated Gaussian convention. The
six local Lorentz directions and the independent auxiliary connection
have already been treated in V10; their measure is not replaced by the
four diffeomorphism ghosts.

\subsection{An additional, explicit coefficient-reference choice}

There is no convergent real Euclidean coframe Gaussian in the supplied
Lorentzian action. We therefore declare a \emph{holomorphic coefficient
reference}, rather than infer a contour theorem from classical BV
closure. Write $D=\operatorname{diag}(i,1,1,1)$. On every finite Fourier
coefficient set
\begin{equation}
 q_L=Dq_ED,\qquad p_L=Dp_E,\qquad
 e_L=D^{-1}e_E D,\qquad
 e_E(q_E)=I+\mathsf E_E(q_E),
 \label{eq:v13-Wick-map}
\end{equation}
where $\mathsf E_E(q)_{00}=q_{00}/2$, $\mathsf E_E(q)_{0i}=0$,
$\mathsf E_E(q)_{i0}=q_{0i}$, and
$\mathsf E_E(q)_{ij}=q_{ij}/2$. Thus
$e_L^T\eta e_L=D e_E^Te_E D$ exactly. Define $S^E_{r,a}$, $r=2,3,4$,
to be \emph{minus} the holomorphic substitution of the corrected
$\widehat S_{r,a}$ of V11--V12. The overall sign is fixed so that the
traceless flat kinetic coefficient below is positive for $\chi>0$.
We restrict to that Palatini-sign branch. This is an explicit linear
operation on the native coefficients, including all their link-word,
stationary-return and restoring terms. It does not replace them by
continuum Feynman rules.

For the contravariant diffeomorphism ghost set $c_L=D^{-1}c_E$;
then $c_L\cdot p_L=c_E\cdot p_E$. Antighost coordinates are transported
with a constant phase chosen so that their quadratic operator is
$-F\widehat R_{0,a}$. Constant contour and ghost phases cancel in normalized finite
integrals. The analytic germ of the V12 negative-log auxiliary amplitude
is transported by the same field substitution, without the minus sign
used for the classical action. Denote it by $\mathcal M^E_{a,1}$.
Here $\log|\det g_L|$ means its real-analytic germ
$\log(-\det g_L)$ at $\eta$, and the determinant branches are anchored
there. No continuation of a Lorentzian distribution, positive
nonperturbative measure, or equality of complete real-time and Euclidean
functional integrals is asserted. Those are additional comparison
problems. The present statements concern the declared finite coefficient
reference.

Work on four odd periodic grids, physical pairing $a^4\sum_x$, momenta
$p\in\mathcal B_a=(-\pi/a,\pi/a)^4$, and fixed $\chi>0$.
The zero Fourier modes of both the metric and the four ghosts are retained,
not inverted. The ordinary harmonic functional in these coordinates is
\[
 F(q)_\mu=\partial_\nu(q_{\mu\nu}-\tfrac12\delta_{\mu\nu}\operatorname{tr}q).
\]
Its spectral derivative acts at the actual total Fourier label. Products
and derivative order on an interior panel are those of the displayed
native polynomials; no lattice Leibniz identity is assumed. For a
literal whole-grid extension one must fix cyclic multiplication and the
ordered derivative placement. The interior results below do not depend
on how two equivalent continuum-polynomial representatives are extended
to aliased high-frequency panels.

The source audit uses the already evaluated operands in
\cref{thm:v10-reduced-jets,thm:v11-artifacts,thm:v12-quartic-artifacts,prop:v12-R2,thm:v12-Jacobian}. Their classical computations are inherited.
Gaussian/Berezin determinant identities and measured covariance
composition are inherited too. The new claims are the simultaneous
reference bound for the actual corrected free operators, its particular
four-term loop tensor, quantitative shell/cutoff estimates, and the
explicit boundary diagnostic. Perturbative gravity as EFT and cutoff
Ward identities have established frameworks
\cite{V13BFR,V13DM}; neither is claimed to begin with this calculation.

\section{A uniformly conditioned reference for all ten metric components}
\label{sec:v13-reference}

Equip $\operatorname{Sym}_4$ with the Frobenius inner product. Introduce
\begin{equation}
 \begin{aligned}
 Jq&=q-\tfrac12I\operatorname{tr}q,&
 F_vq&=qv-\tfrac12v\operatorname{tr}q,\\
 X(v)&=|v|^2J-2F_v^*F_v,&
 u(p)&=(0,-p_1^3/24,-p_2^3/24,-p_3^3/24),\\
 d_a(p)&=(p_0,2a^{-1}\sin(ap_1/2),\ldots,2a^{-1}\sin(ap_3/2)),&
 \widehat d_a(p)&=d_a(p)-a^2u(p).
 \end{aligned}
 \label{eq:v13-symbol-data}
\end{equation}
$F_v$ in this display has no factor $i$; the Fourier symbol of $F$ is
$iF_p$. A star here is the Euclidean matrix adjoint. In real coefficient
identities it is transpose.

\begin{proposition}[The actual corrected free matrices]
\label[proposition]{prop:v13-free}
The Hessian of $S^E_{2,a}+\chi\|Fq\|^2/2$ and the free ghost operator are
\begin{align}
 K_a(p)&=\frac\chi2\{X(d_a(p))-a^2DX(p)[u(p)]
                                      +2F_p^*F_p\},\label{eq:v13-K}\\
 M_a(p)&=(p\cdot\widehat d_a)I_4+
                 \widehat d_a p^T-p\widehat d_a^T.\label{eq:v13-M}
\end{align}
In particular $K_0(p)=\chi|p|^2J/2$ and $M_0(p)=|p|^2I_4$.
No TT projection has been made. The metric numerator has trace-reversal
inertia $(9,1)$ at zero mesh; it is not a positive scalar covariance.
\end{proposition}
\begin{proof}
Substitute the native linear stationary connection in the quadratic
cofactor action. Equivalently, expand the already identified
Fierz--Pauli quadratic form with the same phase in every slot. In the
substitution \eqref{eq:v13-Wick-map} its Hessian is $\chi X(d_a)/2$:
its bilinear form is $\chi/2$ times
\[
 |d|^2\{\operatorname{tr}(qr)-\tfrac12\operatorname{tr}q\operatorname{tr}r\}
       -2(F_dq)\cdot(F_dr).
\]
This normalization also follows directly from
$-b_1^TC_0^{-1}b_1$ with the cofactor \eqref{eq:v11-ell}; it is not
chosen by the desired propagator. The coefficient of $a^2$ in this
quadratic form is $\chi DX(p)[u]/2$. V11 subtracts exactly that
coefficient. The gauge-fixing Hessian is $\chi F_p^*F_p$, proving
\eqref{eq:v13-K}. Apply $-iF_p$ to
$i(\widehat d c^T+c\widehat d^T)$ to obtain \eqref{eq:v13-M}.
This also records the antisymmetric matrix part of the ghost operator;
replacing $M_a$ by a scalar symbol at nonzero mesh is not exact.
\end{proof}

The algebraic nonminimal integral is unambiguous in this reference.
For a dual source $\beta$ its action is
$b^2/(2\chi)+ib\cdot Fq-\beta\cdot b$. Its normalized integration gives
\begin{equation}
 \frac\chi2\|Fq\|^2+i\chi\langle\beta,Fq\rangle
                           -\frac\chi2\|\beta\|^2.
 \label{eq:v13-b-contact}
\end{equation}
Thus its determinant is a background-independent normalization, whereas
its source contact is not zero. With $\beta=0$ it gives exactly the
quadratic gauge term in \eqref{eq:v13-K}. This choice is a contour/phase
realization of the nonminimal Gaussian, not an additional physical field.

\begin{theorem}[A fixed conformal contour with a full-zone free bound]
\label{thm:v13-accretive}
Let $P_0q=I\operatorname{tr}q/4$, $P_T=I-P_0$, and
$T=P_T+iP_0$. For every nonzero $p\in\mathcal B_a$,
\begin{equation}
 \operatorname{Re}(T^TK_a(p)T)\succeq\frac{9\chi}{40}|p|^2I_{10},
 \qquad
 \operatorname{Re}M_a(p)\succeq\frac9{10}|p|^2I_4.
 \label{eq:v13-global-floor}
\end{equation}
Here $\operatorname{Re}A=(A+A^*)/2$. Consequently
\[
 \|K_a(p)^{-1}\|\le\frac{40}{9\chi}|p|^{-2},\qquad
 \|M_a(p)^{-1}\|\le\frac{10}{9}|p|^{-2}.
\]
The finite metric Gaussian converges on the fixed contour $q=Ty$,
$y$ real, and the finite ghost determinant is nonzero. On
$a|p|\le1$, for every fixed multi-index $\alpha$,
\begin{align}
 \|\partial_p^\alpha K_a^{-1}\|&\le C_\alpha\chi^{-1}|p|^{-2-|\alpha|},
 &\|\partial_p^\alpha(K_a^{-1}-K_0^{-1})\|
       &\le C_\alpha\chi^{-1}a^4|p|^{2-|\alpha|},\nonumber\\
 \|\partial_p^\alpha M_a^{-1}\|&\le C_\alpha|p|^{-2-|\alpha|},
 &\|\partial_p^\alpha(M_a^{-1}-M_0^{-1})\|
       &\le C_\alpha a^4|p|^{2-|\alpha|}.
 \label{eq:v13-inverse-symbols}
\end{align}
The global floor does not assert periodic smoothness of these symbols.
\end{theorem}
\begin{proof}
The quadratic operator $X(v)$ has norm at most $2|v|^2$. To check the
constant, rotate $v$ to an axis and diagonalize its scalar block; its
nonzero eigenvalues are $|v|^2$ and $-2|v|^2$. Polarization and scaling
therefore give $\|DX(v)[w]\|\le4|v||w|$.
Set $e=d_a-p-a^2u$ and $r=p+a^2u$. Taylor's theorem gives, componentwise,
\[
 |e_i|\le a^4|p_i|^5/1920,\quad e_0=0.
\]
On the whole cube $|a^2u|\le U|p|$, $|e|\le E|p|$, and
$|r|\le|p|$, where $U=\pi^2/24$, $E=\pi^4/1920$. Since $X$ is
quadratic,
\[
 \Delta_a:=X(d_a)-X(p)-a^2DX(p)[u]
              =X(a^2u)+DX(r)[e]+X(e).
\]
It follows that
\[
 \|\Delta_a\|\le (2U^2+4E+2E^2)|p|^2<\frac{11}{20}|p|^2.
\]
The last inequality follows already from $\pi<22/7$, which gives
\[
 \frac{22712305121}{41506567200}<\frac{11}{20}.
\]
Now $T^TJT=I$ and $T$ is unitary, so
$T^TK_aT=\chi(|p|^2I+T^T\Delta_aT)/2$. This proves the first floor.
The symmetric part of $M_a$ is $(p\cdot(p+e))I$ and
$E<1/10$, which proves the second. The inverse bounds follow by
Cauchy--Schwarz from accretivity and by unitary congruence for $K_a$.

For $a|p|\le1$, the same calculation gives the stronger estimate
$\|\Delta_a\|\le (a|p|)^4|p|^2/128$; use
$|a^2u|\le a^2|p|^3/24$ and the displayed bound for $e$.
Also $\|M_a-|p|^2I\|\le3a^4|p|^6/1920$.
Sine and its differentiated Taylor remainders give the corresponding
bounds with each derivative removing a power of $|p|$.
Differentiate the inverse identity, or its convergent relative Neumann
series, to prove \eqref{eq:v13-inverse-symbols}. The finite Gaussian
converges because its real quadratic part is positive on every integrated
mode; no size-independent infrared gap is claimed as $p\to0$.
\end{proof}

The bound is for a \emph{specified} complex contour and coefficient
reference. It is not an evasion of the conformal-factor issue by calling
the real quadratic form positive. The native zero modes are still
retained. It also does not prove convergence of the exponential of the
full interacting polynomial. Higher interactions are evaluated
coefficientwise by finite Gaussian/Berezin moments, as in the inherited
perturbative calculus.

\section{An actual interior shell and its complete geometric two-jet}
\label{sec:v13-shell}

Fix a smooth radial $w\ge0$, $w\le1$, supported in
$1/2<|z|<2$. Write $w_\Lambda(p)=w(p/\Lambda)$ and define
\begin{equation}
 C^g_{a,\Lambda}(p)=w_\Lambda(p)K_a(p)^{-1},\qquad
 C^c_{a,\Lambda}(p)=w_\Lambda(p)M_a(p)^{-1}.
 \label{eq:v13-covariances}
\end{equation}
Zero weights mean unintegrated coordinates. These are genuine Gaussian
and Berezin innovations: on their support the inverse metric quadratic
form is $w_\Lambda^{-1}T^TK_aT$ in $y$ coordinates. Its real part is
positive. The complementary covariance $(1-w_\Lambda)K_a^{-1}$ has the
same property. Normalized complex Gaussian convolution and finite
Berezin convolution give exact composition. The result is not a sharp
spectral deletion.

For this section impose
\begin{equation}
 a\Lambda\le1/512,\qquad \Lambda L_\nu\ge1,\qquad
 \operatorname{supp}\widehat h\subset\{|k|\le\mu\},\quad
 0<\mu\le\Lambda/32.
 \label{eq:v13-interior-domain}
\end{equation}
Every shifted hard momentum in a two-background loop is inside
$|p|<3\Lambda$, and every required subset momentum is nonaliased.
The metric background is the same affine coframe coordinate, now in the
reference \eqref{eq:v13-Wick-map}. Take the integrated ghosts to have
zero retained background in the coefficient computed below.

The matrices entering the calculation are explicitly specified as follows:
\begin{align}
 A^g_a[h]&=D^3 S^E_{3,a}(0)[h,\cdot,\cdot],&
 B^g_a[h_1,h_2]&=D^4S^E_{4,a}(0)[h_1,h_2,\cdot,\cdot],\nonumber\\
 A^c[h]&=-F R_1^E(h,\cdot),&
 B^c[h_1,h_2]&=-F D_q^2R_2^E(0,\cdot)[h_1,h_2].
 \label{eq:v13-vertices}
\end{align}
The first two derivatives mean coefficient extraction from
\eqref{eq:v10-reduced-jets}, with precisely
\eqref{eq:v11-counterterm} and \eqref{eq:v12-counterterm}, followed by
\eqref{eq:v13-Wick-map}. The polynomial $B_n$ of V10 and the constant
$C_0^{-1}$ specify every matrix entry; no action vertex is left as an
uncomputed independent input. The last two use the explicit
\eqref{eq:v11-R1} and \eqref{eq:v12-R2}, with $\eta$ replaced by $I$
in the transported chart. In particular $B^c$ is not set to zero.
There is no nonlinear gauge-fixing contribution to $A^g,B^g$, because
$F$ was chosen linear in $q$.

\begin{lemma}[Primitive bounds for these hard Hessian vertices]
\label[lemma]{lem:v13-vertex-bound}
On \eqref{eq:v13-interior-domain}, each fixed mixed derivative of the
Fourier kernels of $A^g,B^g$ has bound
$C_\alpha\chi\Lambda^{2-|\alpha|}$ times the background-polarization
norms. Their differences from their zero-mesh kernels have bound
$C_\alpha\chi a^3\Lambda^{5-|\alpha|}$.
The kernels of $A^c,B^c$ have bound $C_\alpha\Lambda^{2-|\alpha|}$ and
are exactly their zero-mesh local-polynomial expressions in this
reference. All constants are independent of the periods and shell scale.
\end{lemma}
\begin{proof}
Polarize the two already evaluated native polynomials. Each geometric
load has one phase derivative; the cofactor and Cartan factors are
finite polynomials; each additional connection word has its recorded
power of $a$. On the interior annulus $a\Lambda\le1/512$, the
trigonometric symbols and their derivatives are uniformly controlled.
V11's cubic and V12's quartic remainder bounds, with the individual
momenta bounded by $3\Lambda$, therefore give exactly the stated
mismatch. Differentiating in hard or external momenta is permitted by
the same finite analytic expressions; the fixed change
\eqref{eq:v13-Wick-map} is bounded on these finite tensors. Each ghost
vertex contains one derivative in $R_1$ or $R_2$ and one in $F$.
This proves its order-two bound. Neither ghost vertex depends on $a$;
the free ghost matrix does, and its difference is accounted for in
\eqref{eq:v13-inverse-symbols}.
\end{proof}

\begin{theorem}[Complete one-loop two-metric coefficient of the interior pushforward]
\label{thm:v13-twojet}
For sufficiently small $\|h\|_{\mathrm W}=\sum_k\|\widehat h(k)\|$,
with a radius independent of the mesh, volume and scale in
\eqref{eq:v13-interior-domain}, the quadratic-hard-field Gaussian and
ghost integrals have a common analytic background germ. Their
background-dependent one-loop functional is
\begin{equation}
 \Gamma^{\mathrm{prop}}_{a,\Lambda}[h]
 =\frac\hbar2\Tr\log\{I+C^g(A^g[h]+\tfrac12 B^g[h,h])\}
 -\hbar\Tr\log\{I+C^c(A^c[h]+\tfrac12 B^c[h,h])\}.
 \label{eq:v13-logdet}
\end{equation}
Its complete polarized Hessian is
\begin{equation}
 \boxed{\begin{aligned}
 D^2\Gamma^{\mathrm{prop}}_{a,\Lambda}[h_1,h_2]
 ={}&\frac\hbar2\Tr\{C^g B^g[h_1,h_2]
                           -C^g A^g[h_1]C^g A^g[h_2]\}\\
 &-\hbar\Tr\{C^c B^c[h_1,h_2]
                           -C^c A^c[h_1]C^c A^c[h_2]\}.
 \end{aligned}}
 \label{eq:v13-four-terms}
\end{equation}
This is also the one-loop, two-retained-metric coefficient of the
full finite perturbative pushforward of the supplied corrected action,
not only a Hessian ansatz. The auxiliary-density contribution at this
loop order is $\hbar D^2\mathcal M^E_{a,1}$ and is assigned once,
not once per propagating shell. The nonminimal contact is
\eqref{eq:v13-b-contact}. No source-dependent seagull or Jacobian has
been included in a removable vacuum constant.
\end{theorem}
\begin{proof}
The free inverse and primitive vertex bounds give a relative hard
operator bound $C(\|h\|_{\mathrm W}+\|h\|_{\mathrm W}^2)$ by the
Fourier convolution Schur estimate; hard input/output momenta remain
comparable. Choose a common radius making it less than one quarter
of the reference margin. On the fixed contour the metric Gaussian is
then convergent and analytic. The ghost matrix is an invertible
perturbation of its free one. Normalized Gaussian and Berezin integration
give \eqref{eq:v13-logdet}, with the opposite determinant powers shown.
Differentiating twice gives \eqref{eq:v13-four-terms}; cyclicity makes
the two polarized bubble orderings equal inside the trace. In
particular, differentiating $R_2$ gives the ghost contact with the full
factor two on a repeated background.

To identify the coefficient with the full perturbative pushforward,
retain hard linear terms before applying loop counting. The free
linear term from $h$ has low momentum. The cubic hard force quadratic
in $h$ has total momentum at most $2\mu$, and the quartic hard force
cubic in $h$ has total momentum at most $3\mu$. All vanish on the hard
support, whose lower radius is $\Lambda/2$. The stationary hard shift
therefore has no term through these background degrees. A connected
one-loop diagram with two retained metric legs is consequently either
the quartic metric contact, two cubic metric vertices, the quartic
ghost contact, or two cubic ghost vertices. A one-loop diagram with a
hard one-legged branch would contain one of the zero hard forces just
identified. Pure hard cubic/quartic vertices without such a branch
first contribute additional loops. This proves the identification and
also explains why an arbitrary off-shell linear hard force could not
have been omitted without the support check.

The old auxiliary amplitude is loop-graded: the action contains
$\hbar\mathcal M^E_{a,1}(h+\eta)$. At one loop it is evaluated at the
classical stationary section, equal to $h$ at the degrees in question.
A fluctuation contraction in this amplitude supplies another $\hbar$.
Thus its contribution is the stated Hessian, once. This argument
preserves its quantum grading instead of inserting that density as a
new ungraded kinetic Hessian. Only the source-independent normalization
at $h=0$ is removed.
\end{proof}

The theorem is source-complete for its declared geometric two-jet and
nonminimal source contact. It is not a theorem for arbitrary elementary
hard-field sources: such sources can move the stationary section and
produce further diagrams. Nor is it a full antifield-dependent Slavnov
functional. Those larger panels must be specified and computed before
a quantum BV conclusion is drawn from this block.

\section{Local assignment before the ultraviolet sum}
\label{sec:v13-subtraction}

For $h_1=He^{ikx}$, $h_2=Ge^{-ikx}$, let
$\Pi_{a,\Lambda,L}(k)[H,G]$ be \eqref{eq:v13-four-terms} divided by
physical volume. Every trace is a finite routed loop sum with physical
normalization $\operatorname{Vol}_L^{-1}\sum_p$. The second propagator,
when present, has momentum $p+k$ and its own weight $w_\Lambda(p+k)$.
Neither the weight nor any translated occurrence is frozen before
differentiation. Define
\[
 \Pi^R_{a,\Lambda,L}=(I-\mathsf T_4)\Pi_{a,\Lambda,L},
\]
where $\mathsf T_4$ is the complete Taylor jet in $k$ through degree four.
It is stored in the local geometric-source bank.

There is no need to assume the tensor is even componentwise in $k$.
Its bilinear symmetry is $\Pi_{AB}(k)=\Pi_{BA}(-k)$. The oriented native
vertices have not been proved reflection-even. A safe redundant local
bank therefore allows symmetric component matrices at even derivative
degree and antisymmetric matrices at odd degree. Its dimension at two
metric marks is at most
\begin{equation}
 55(1+10+35)+45(4+20)=3610.
 \label{eq:v13-bank}
\end{equation}
Integration by parts gives the indicated matrix symmetries. This bank is
a finite component restoring bank, not a count of independent covariant
curvature invariants. It neither deletes its complementary kernel nor
imposes a gravitational Ward identity by projection.

\begin{theorem}[A volume-uniform subtracted propagating shell estimate]
\label{thm:v13-shell-bound}
On \eqref{eq:v13-interior-domain}, for every fixed $r=|\alpha|\le5$,
\begin{align}
 |\partial_k^\alpha\Pi^R_{a,\Lambda,L}(k)[H,G]|
    &\le C_\alpha\hbar\|H\|\|G\|
                  |k|^{5-r}\Lambda^{-1},\label{eq:v13-shell-remainder}\\
 |\partial_k^\alpha(\Pi^R_{a,\Lambda,L}-\Pi^R_{0,\Lambda,L})(k)[H,G]|
    &\le C_\alpha\hbar\|H\|\|G\|
                   a^3|k|^{5-r}\Lambda^2.\label{eq:v13-shell-mismatch}
\end{align}
The unsubtracted local jets of derivative degree $r\le4$ have bounds
$C_r\hbar\Lambda^{4-r}\|H\|\|G\|$. For $r>5$, the first bound becomes
$C_r\hbar\Lambda^{4-r}$ and the second
$C_r\hbar a^3\Lambda^{7-r}$, with the polarization norms understood.
The same statements hold in infinite volume. For a fixed smooth source
window $\varpi(k/\mu)$, the inverse Fourier kernels satisfy, for every
fixed $s\ge0$,
\begin{align}
 a^4\sum_x(1+\mu d_{\rm per}(x,0))^s
       \|\mathcal K^R_{a,\Lambda,L}(x)\|
      &\le C_s\hbar\mu^5\Lambda^{-1},\nonumber\\
 a^4\sum_x(1+\mu d_{\rm per}(x,0))^s
       \|\mathcal K^\Delta_{a,\Lambda,L}(x)\|
      &\le C_s\hbar a^3\mu^5\Lambda^2.
 \label{eq:v13-shell-kernels}
\end{align}
\end{theorem}
\begin{proof}
A metric bubble has two vertices of size $\chi\Lambda^2$ and two
propagators of size $\chi^{-1}\Lambda^{-2}$. A metric contact has
one vertex of size $\chi\Lambda^2$ and one inverse. The corresponding
ghost orders are the same without $\chi$. Thus every full integrand in
\eqref{eq:v13-four-terms} has order zero; each differentiated momentum
lowers its bound by one power of $\Lambda$. All differentiated cutoff
weights have that same scaling. The support has normalized lattice count
at most $C\Lambda^4$, since $\Lambda L_\nu\ge1$. This proves
$C_r\hbar\Lambda^{4-r}$ for the complete differentiated tensor.
The estimate is taken after the signed four terms have been specified;
no identity between their uncomputed individual values is assumed.

Use the exact resolvent difference and
\cref{lem:v13-vertex-bound,thm:v13-accretive}, replacing one factor
at a time. A metric vertex mismatch supplies $(a\Lambda)^3$ and a
propagator or free-ghost mismatch supplies $(a\Lambda)^4$.
Hence each fixed derivative of the integrated difference is bounded by
$C_r\hbar a^3\Lambda^{7-r}$. The lesser power absorbs the latter
terms on the stated domain. Apply the differentiated integral Taylor
remainder through degree four. Its first required derivative is five,
not six: no unproved parity improvement is used. This yields
\eqref{eq:v13-shell-remainder}--\eqref{eq:v13-shell-mismatch}.
For higher derivatives the subtracted polynomial is already zero.

For the kernel statement, the scaled derivatives of the windowed symbol
are bounded by the corresponding right side with $|k|$ replaced by
$\mu$. Integrate by parts in four external momentum variables more
than $s+4$ times, then sample and periodize. The source window lies
strictly inside the lattice zone and $a\mu$ is uniformly small; the
weighted sampled sum of the resulting rapidly decreasing kernel is
bounded by its scaled symbol norm. Periodic distance is no larger than
the distance of a lift. This proves both inequalities without a volume
multiplier. In particular the bilinear source functional is bounded with
one source in $L^1$ and the other in $L^\infty$.
\end{proof}

\section{Measured shell composition and a cutoff limit for the interior family}
\label{sec:v13-iterate}

Fix a lower Wilsonian scale $\Lambda_0\ge32\mu$ and a nonincreasing
smooth radial $\theta$, equal to one on $|z|\le1$ and zero on
$|z|\ge2$. Set $\Lambda_j=2^j\Lambda_0$ and
\begin{equation}
 u_j(p)=\theta(p/\Lambda_j)-\theta(p/\Lambda_0),\qquad
 w_j(p)=u_j(p)-u_{j-1}(p),\quad u_0=0.
 \label{eq:v13-dyadic}
\end{equation}
Both are nonnegative. Use $C^g_{\le j}=u_jK_a^{-1}$,
$C^c_{\le j}=u_jM_a^{-1}$ in \eqref{eq:v13-logdet}; call the result
$\Gamma^{\rm prop}_{a,\le j}$. This is the cumulative measured Gaussian
elimination, not a sum of independent shell determinants. Modes below
$\Lambda_0$, including zero modes, and the upper complementary covariance
remain unintegrated.

\begin{proposition}[No missing mixed-shell pairs]
\label[proposition]{prop:v13-owner}
For $\Delta_j=\Gamma^{\rm prop}_{a,\le j}
-\Gamma^{\rm prop}_{a,\le j-1}$, the two-background bubble includes every
pair of innovations whose larger label is $j$, exactly once. Its
subtracted two-jet obeys
\begin{equation}
 \|D^2(I-\mathsf T_4)\Delta_j\|_{\mathrm{src},\mu}
                  \le C\hbar\mu^5\Lambda_j^{-1},
 \label{eq:v13-step-tail}
\end{equation}
with the corresponding lattice--continuum error
$C\hbar a^3\mu^5\Lambda_j^2$.
\end{proposition}
\begin{proof}
Expand the two cumulative covariances in each bubble. The difference
retains the pairs $(j,j)$, $(j,i)$ and $(i,j)$ for $i<j$.
This is the inherited measured Schur allocation applied to the actual
metric and ghost vertices. In a term with one $w_j(p)$, momentum transfer
$|k|\le\mu$ forces the other propagator momentum to be comparable to
$\Lambda_j$. For sufficiently large $j$, the low cutoff
$\theta((p+k)/\Lambda_0)$ is identically zero on that support; the first
bounded number of steps obey the same estimates with fixed constants.
The cumulative weight is at most one, and its nonconstant derivatives
on this support have order $\Lambda_j^{-r}$. The proof of
\cref{thm:v13-shell-bound} therefore applies, without a factor equal to
the number of preceding shells.
\end{proof}

Let $J(a)$ be the largest integer with $a\Lambda_{J(a)}\le1/512$;
for sufficiently small $a$ it is positive, and
$1/(1024a)<\Lambda_{J(a)}\le1/(512a)$.
At zero mesh define $C^g_\infty=(1-\theta(p/\Lambda_0))K_0^{-1}$ and
$C^c_\infty=(1-\theta(p/\Lambda_0))M_0^{-1}$.
Their degree-four-subtracted one-loop two-jet is denoted
$\Pi^R_{0,\infty,L}$; this symbol $\infty$ refers to the ultraviolet
endpoint, not to physical volume. It has the fixed lower cutoff
$\Lambda_0$, and is absolutely convergent after the subtraction.

\begin{theorem}[A realized interior covariance-cutoff continuum coefficient]
\label{thm:v13-interior-limit}
For $\Lambda_0L_\nu\ge1$, fixed $\mu\le\Lambda_0/32$, and
$r=|\alpha|\le5$,
\begin{equation}
 \boxed{|\partial_k^\alpha
 (\Pi^R_{a,\le J(a),L}-\Pi^R_{0,\infty,L})(k)[H,G]|
 \le C_\alpha\hbar\|H\|\|G\|\,a|k|^{5-r}.}
 \label{eq:v13-band-limit}
\end{equation}
The auxiliary-density two-jet, with its \emph{own actual} local Taylor
jet through degree four assigned once, obeys the same bound and has zero
nonlocal limit. Thus \eqref{eq:v13-band-limit} also holds for the combined
propagating and auxiliary two-jet in this reference. Every fixed
windowed kernel moment has error $C_s\hbar a\mu^5$.
The local coefficients are generated by the exact increments, not supplied
as an external trajectory: for a prescribed physical local polynomial $v$,
\begin{equation}
 v_j=v-\mathsf T_4D^2\Gamma^{\rm prop}_{a,\le j},\qquad
 R_j=(I-\mathsf T_4)D^2\Gamma^{\rm prop}_{a,\le j}.
 \label{eq:v13-local-flow}
\end{equation}
The auxiliary local assignment is made once at the start or end, with
the same result. This is a two-geometric-source, one-loop coefficient
trajectory at a fixed lower scale. It is not an invariant neighbourhood
of arbitrary interacting ESM actions.
\end{theorem}
\begin{proof}
Sum \eqref{eq:v13-step-tail}. Its continuum tail above $J$ is bounded
by $C\hbar|k|^{5-r}\Lambda_J^{-1}$. Compare the native and zero-mesh
increments before summing; their difference is at most
\[
 C\hbar|k|^{5-r} a^3\sum_{j\le J}\Lambda_j^2
 \le C'\hbar|k|^{5-r}a^3\Lambda_J^2.
\]
At the specified $J(a)$ both bounds are $C\hbar a|k|^{5-r}$.
This uses the fifth-derivative local subtraction in addition to the
classical small-mesh comparison. It is not the invalid estimate
$\sum_j(a\Lambda_j)^3\to0$.

V10's fixed-degree auxiliary matching theorem applies to two metric
marks, including their actual stationary-section and Haar contributions.
V12's coordinate Jacobian is ultralocal and hence adds only to its local
jet. Its finite coefficient is analytic at small external $aK$; the
fixed complex substitution \eqref{eq:v13-Wick-map} preserves this
property and its fixed-order bounds. Its subtracted two-mark remainder
is therefore $Ca|k|^{5-r}$, with the same qualification about the
declared reference. This proves the auxiliary statement without
reintegrating the connection or counting its determinant once per shell.
For $r>5$ the summed derivative estimates from
\cref{thm:v13-shell-bound} are bounded by
$C_r\hbar a\Lambda_0^{5-r}$ (a harmless logarithm at $r=7$ is
absorbed using $a\Lambda_0\le1/512$). Rescaling by the smaller
$\mu$ and the Fourier argument of \cref{thm:v13-shell-bound} gives
the kernel error. Finally \eqref{eq:v13-local-flow} follows by exact
finite subtraction, and its differences are the actual measured
$-\mathsf T_4D^2\Delta_j$ and $(I-\mathsf T_4)D^2\Delta_j$.
\end{proof}

A separate thermodynamic limit follows for this fixed lower cutoff:
the subtracted integrand has integrable loop-momentum derivatives of
every fixed order. Fourier integration by parts and Poisson summation
therefore give, for any fixed $b>4$, a bound
$C_b\hbar\mu^5\Lambda_0^{-1}(\Lambda_0L_{\min})^{-b}$
in the same scaled finite-panel norms. The ultraviolet and volume limits
have different error terms. Removing $\Lambda_0$ or identifying the
real-time continuation is not proved. Neither is the equality of this
interior family with the \emph{entire} original Brillouin-zone integral;
that remaining comparison is made explicit in \cref{sec:v13-seam}.

\section{Two evaluated checks, including a retained Higgs contribution}
\label{sec:v13-evaluated}

\subsection{A constant coframe source detects the ghost determinant and measure}

Take $h=tI_4$. Then $e_E=(1+t/2)I_4$. The completed native action is
homogeneous of degree two under constant coframe rescaling, before and
after connection stationarity. Consequently its hard quadratic form
about this background is independent of $t$: the factors
$r^2S(I+\eta/r)$ cancel at hard degree two. The statement holds for
every $a$, including the shared-link terms; their flat connection starts
at hard degree one and their cubic remainder cannot contribute two hard
legs at constant coframe. The local artifact coefficients inherit this
identity. Thus
\begin{equation}
 A^g[tI]=0,\qquad B^g[tI,tI]=0.
 \label{eq:v13-conformal-grav}
\end{equation}
The gauge-fixing action is quadratic and background independent.

In the transported affine chart the exact ordinary conformal transformation
is $R(tI,c)=(1+t/2)R_0c+(c\cdot\partial t)I$.
For constant $t$, $R_1(tI,c)=tR_0c/2$ and $R_2(tI,c)=0$.
The ghost action therefore has $A^c[tI]=t|p|^2I_4/2$ and zero repeated
conformal seagull, while its free operator is the actual $M_a$.

\begin{proposition}[A nonzero measured local two-jet]
\label[proposition]{prop:v13-conformal}
Put $\sigma=p\cdot\widehat d_a$ and
$b_p^2=|p|^2|\widehat d_a|^2-\sigma^2$. For this direction,
\begin{equation}
 \frac1{\operatorname{Vol}_L}
 \frac{\dd^2}{\dd t^2}\Gamma^{\rm prop}_{a,\Lambda}[tI]\big|_{t=0}
 =\frac\hbar4\int_p^L w_\Lambda(p)^2|p|^4
 \left\{\frac2{\sigma^2}
       +\frac{2(\sigma^2-b_p^2)}{(\sigma^2+b_p^2)^2}\right\}>0
 \label{eq:v13-conformal-loop}
\end{equation}
when the shell is nonempty. At zero mesh this is
$\hbar\int_p^Lw_\Lambda(p)^2$.
The repeated conformal derivative of the complete auxiliary measure is
exactly $-9\hbar/(2a^4)$ per volume. It is a local matching term and is
not re-added at each shell. In infinite volume the propagating zero-mesh
coefficient is
\[
 \frac{\hbar\Lambda^4}{8\pi^2}\int_0^\infty r^3w(r)^2\dd r.
\]
This is a coordinate- and gauge-dependent geometric-source coefficient,
not a Newton-coupling or physical cosmological beta function.
\end{proposition}
\begin{proof}
Use \eqref{eq:v13-conformal-grav} in the four-term expression.
The remaining ghost bubble is
$\hbar\Tr(w_\Lambda^2|p|^4M_a^{-2})/4$.
The antisymmetric matrix $\widehat d p^T-p\widehat d^T$ has two zero
eigenvalues and the pair $\pm ib_p$. Thus its inverse-square trace is
the braces in \eqref{eq:v13-conformal-loop}. Its positivity follows by
putting the two summands over a common denominator: the numerator is
$2(2\sigma^4+\sigma^2b_p^2+b_p^4)$.
At $a=0$ the braces equal $4/|p|^4$. The radial formula uses the
four-dimensional sphere area $2\pi^2$ and the physical $(2\pi)^{-4}$
Fourier normalization. At constant coframe $A_*=0$ and $K_a=0$ in
V10's \emph{auxiliary} determinant notation. V12's density is
$r^{-18}$, $r=1+t/2$. Its negative logarithm is $18\sum_x\log r$;
two derivatives give $-9\sum_x/2$, proving the last measure assertion.
\end{proof}

\subsection{What this shell actually supplies to the Higgs sector}

The Higgs is retained in the following calculation, not integrated in
V9's different scalar reference. Impose
$\operatorname{supp}\widehat\phi\subset\{|k|\le\mu\}$ with the same
$\mu$ as in \eqref{eq:v13-interior-domain}. To specify its coefficient continuation
set additionally $\phi_L=i\phi_E$ in minus the signed action of V12.
Its zero-mesh scalar action is then
$\int(W^{\mu\nu}\partial_\mu\phi\cdot\partial_\nu\phi
+m^2\rho|\phi|^2)/2$ with $g=e_E^Te_E$.
This is part of the declared holomorphic comparison, not a physical
reality condition. The finite temporal Hodge symbol continues to a
hyperbolic sine, not V9's sine; no equality to that reference is used.
Since $\phi$ is retained, no new scalar contour or inverse is needed.

Let $V^H_a[\phi,\phi]$ be the metric-hard Hessian of the continued,
corrected $\widehat S^H_{4,a}$ from V12. It is explicitly the second
coframe coefficient $W_2,\rho_2$, with the actual scalar differences and
$-a^2\mathcal H_2(W_2;\phi)$ correction. At background degree two in
$\phi$, the metric/ghost shell generates precisely
\begin{equation}
 \Gamma^{g\to H}_{a,\Lambda}[\phi]
                   =\frac\hbar2\Tr(C^g_{a,\Lambda}V^H_a[\phi,\phi]).
 \label{eq:v13-Higgs-exact}
\end{equation}
The hard metric force from $\widehat S^H_{3,a}$ has momentum at most
$2\mu$ and vanishes on the shell. There is no ghost--Higgs vertex in
the chosen $\phi$-independent gauge.

\begin{proposition}[Evaluated mixed metric-shell contact]
\label[proposition]{prop:v13-Higgs}
Set $I_{\Lambda,L}=\int_p^L w_\Lambda(p)/|p|^2$.
At zero mesh \eqref{eq:v13-Higgs-exact} is the local functional
\begin{equation}
 \boxed{\Gamma^{g\to H}_{0,\Lambda}[\phi]
 =\frac{\hbar I_{\Lambda,L}}{2\chi}\int
 \left\{-\frac12|\partial_0\phi|^2
       +2\sum_{i=1}^3|\partial_i\phi|^2
       -\frac94m^2|\phi|^2\right\}.}
 \label{eq:v13-Higgs-value}
\end{equation}
In infinite volume
$I_{\Lambda}=\Lambda^2(8\pi^2)^{-1}\int_0^\infty r w(r)\dd r$.
The anisotropic local coefficients are in the affine ADM coordinate and
fixed linear gauge. They must be matched as such; they are not asserted
to be a covariant scalar self-energy. A scalar--graviton mixed loop would
also be required if the scalar were integrated, and is not contained in
\eqref{eq:v13-Higgs-value}.
\end{proposition}
\begin{proof}
The zero-mesh metric covariance is $2wJ/(\chi|p|^2)$.
For independent symmetric coordinates its covariance contraction is
\[
 J_{ij,kl}=\tfrac12(\delta_{ik}\delta_{jl}
                  +\delta_{il}\delta_{jk}-\delta_{ij}\delta_{kl}),
 \qquad i\le j,\ k\le l.
\]
Use $\rho_2=((\operatorname{tr}\mathsf E_E)^2
-\operatorname{tr}\mathsf E_E^2)/2$ and
$W_2=\rho_2I-\rho_1q+q^2-\mathsf E_E^T\mathsf E_E$.
Termwise contraction of their ordinary second coordinate derivatives gives
\[
 \sum J_{AB}D_A D_B\rho_2=-9/4,\qquad
 \sum J_{AB}D_A D_B W_2
                   =\operatorname{diag}(-1/2,2,2,2).
\]
For example, the diagonal-coordinate cross derivatives of $\rho_2$
contribute $-3/2$, and the three off-diagonal spatial coordinates contribute
$-3/4$; the shift coordinates do not enter $\det e_E$.
Expanding the four displayed terms of $W_2$ gives the second matrix.
Multiplication by $\hbar C^g/2$ and by the scalar-action factor $1/2$
proves \eqref{eq:v13-Higgs-value}. This is a finite tensor contraction,
not a replacement by a known continuum gravitational mass correction.
\end{proof}

At finite mesh the exact expression \eqref{eq:v13-Higgs-exact} is
retained, including every scalar difference and its restoring insertion.
Its hard covariance integral differs from its zero-mesh value by
$C\chi^{-1}a^4\Lambda^6$, using \eqref{eq:v13-inverse-symbols}.
On a scalar source window $|k|\le\mu$ the corrected scalar vertex has
error $Ca^4\mu^6$, by the holomorphic version of V12's estimate.
These statements quantify this contact but do not turn it into the
uncomputed combined metric--scalar determinant.

\section{The ultraviolet boundary is not controlled by the Gaussian floor}
\label{sec:v13-seam}

The full-zone bound in \cref{thm:v13-accretive} is useful: failure of
invertibility is not the remaining free-reference problem. Smoothness
across the boundary of the principal momentum cube is different. The
ordinary spectral derivative $p_i$ and the half-sine phase are not smooth
periodic functions of that coordinate. The following explicit test shows
why the full-zone scalar proof of V9 cannot simply be renamed a graviton
proof.

\begin{proposition}[An actual ghost-inverse boundary jump]
\label[proposition]{prop:v13-jump}
Use the free ghost matrix \eqref{eq:v13-M} with principal momentum
representatives. At the two sides of the $p_1$ boundary take
\[
 p_\pm=a^{-1}(0,\pm\pi,\pi/2,0).
\]
Let
\begin{align*}
 b_*&=\pi(1-\sqrt2)+\pi^4/64>0,\\
 \sigma_*&=\pi(2+\pi^3/24)
               +(\pi/2)(\sqrt2+\pi^3/192)>0.
\end{align*}
Then
\begin{equation}
 \boxed{(M_a(p_+)^{-1}-M_a(p_-)^{-1})_{12}
            =-\frac{2a^2b_*}{\sigma_*^2+b_*^2}\ne0.}
 \label{eq:v13-jump}
\end{equation}
Consequently this free ghost inverse, periodically extended in $p_1$,
does not possess an absolutely summable Fourier series in that variable.
A uniform full-zone rooted locality estimate cannot be deduced from the
inverse norm bound alone. This is not a proof that the \emph{assembled}
metric/ghost loop has an uncancelled boundary contribution.
\end{proposition}
\begin{proof}
At $p_+$ the two nonzero entries of $a\widehat d$ are
$2+\pi^3/24$ and $\sqrt2+\pi^3/192$.
The $12$ entry of $a^2(\widehat d p^T-p\widehat d^T)$ is $b_*$.
At $p_-$ it changes sign, whereas $a^2p\cdot\widehat d=\sigma_*$
is unchanged. Inverting the resulting two-by-two block proves
\eqref{eq:v13-jump}. Positivity of $b_*$ follows, for example, from
$\pi>157/50$ and $\sqrt2<283/200$ in
$b_*/\pi=1-\sqrt2+\pi^3/64$.

For the Fourier assertion, the inverse is smooth on the closed open-side
interval at this nonzero transverse momentum. Integration by parts in
$t=ap_1$ gives its Fourier coefficients as the nonzero endpoint jump
times $(-1)^n/(2\pi i n)+O(n^{-2})$. Thus absolute summability fails.
Alternatively absolute Fourier convergence would yield a continuous
periodic function, contradicting the jump. Finite odd-grid transverse
momenta approximate this test; it does not assume an even-grid Nyquist
mode belongs to the carrier. A smooth transverse localization near the
test keeps the same obstruction. It is an obstruction to automatically
using the smooth full-zone kernel class, not to a separate signed-loop
cancellation or an explicitly changed high-frequency reference.
\end{proof}

For any fully specified cyclic high-frequency realization, the remaining
one-loop two-source comparison is a concrete object:
\begin{equation}
 \mathfrak C_{a,L}(k)
 =D^2\Gamma^{\rm prop}_{a,\rm full}(k)
                    -D^2\Gamma^{\rm prop}_{a,\le J(a)}(k).
 \label{eq:v13-cap}
\end{equation}
It includes the pure upper-covariance terms \emph{and both mixed upper/interior
pairs}. It is not the sum of only the upper-shell diagonal bubbles.
The theorem still needed to identify the full native endpoint is a
local assignment and source estimate for
$(I-\mathsf T_4)\mathfrak C_{a,L}$, or an explicit high-frequency
restoring prescription with a matched comparison. The boundary jump
prevents obtaining that estimate from the scalar smooth-symbol lemma
without further work. No upper-covariance term has been declared zero.

The same separation applies to symmetry. The free contour and the
classical identities modulo $a^3$ do not imply a quantum Slavnov identity.
The finite propagating cutoff is not invariant under a nonconstant
transformation, the harmonic gauge is fixed rather than background
covariant, and the required antifield source panels have not been
integrated in this theorem. In particular, transversality of just
\eqref{eq:v13-four-terms} is not imposed or claimed. The actual quantum
Ward defect must include the gauge/antifield returns, the inherited
measure divergence and the cap contribution. Standard cutoff Ward
analysis \cite{V13DM} explains the type of additional equation required;
it does not evaluate these native terms.

\section{Ledgers and the next load-bearing comparison}
\label{sec:v13-ledgers}

\subsection*{Proved}
\addcontentsline{toc}{subsection}{V13: Proved}
For the explicitly declared holomorphic coefficient reference,
\cref{thm:v13-accretive} constructs a common fixed conformal contour and
invertible ghost reference for all nonzero finite modes, with explicit
full-zone lower bounds. \Cref{thm:v13-twojet} evaluates the complete
one-loop two-geometric-source interior coefficient, including both metric
and ghost contacts and the separately once-assigned auxiliary density.
\Cref{thm:v13-shell-bound,thm:v13-interior-limit} give a rooted subtracted
shell bound, mixed-shell composition, generated local coefficient
increments and an $O(a)$ cutoff comparison for the specified interior
covariance-cutoff family at a fixed lower Wilsonian scale.
\Cref{prop:v13-conformal,prop:v13-Higgs} evaluate a nonzero conformal
source coefficient and a mixed retained-Higgs contact. Finally
\cref{prop:v13-jump} exhibits the actual principal-symbol boundary jump;
the full-zone missing comparison is \eqref{eq:v13-cap}.

\subsection*{Inherited}
\addcontentsline{toc}{subsection}{V13: Inherited}
All prior mathematical bodies and labels are unchanged. V10 supplies the
native reduced action and auxiliary Lorentz/Haar amplitude; V11--V12
supply their explicit restoring coefficients and nonlinear chart ghost
vertices. V12's $g$-to-$q$ Jacobian is retained. V6's TT construction,
V9's massive internally covariant scalar module, the general measured
Gaussian/Berezin rules and abstract filtered recursion are not reproved
as new results. No full source or symmetry conclusion is imported from a
narrower module.

\subsection*{Literature input}
\addcontentsline{toc}{subsection}{V13: Literature input}
Perturbative effective gravity and its BV treatment are established
\cite{V13BFR}; cutoff-dependent Ward equations are also established
\cite{V13DM}. They supply methodological context, not a native
reference bound, a contour-equivalence theorem, or the unproved cap
estimate. No priority claim is made for Gaussian determinants,
perturbative gravity, or Taylor subtraction.

\subsection*{Conditional}
\addcontentsline{toc}{subsection}{V13: Conditional}
Identification with the original complete Lorentzian ESM source
functional still requires the declared coefficient-reference comparison,
a full-zone high-frequency realization and matching of
\eqref{eq:v13-cap}, and a measure-compatible quantum BV/source theorem.
The scalar is retained in \cref{prop:v13-Higgs}; a scalar--graviton mixed
loop, internal gauge fields, spin/chiral phases and their common contour
have not been supplied by that contact. A finite local bank and a
one-loop generated two-jet do not prove an invariant nonlinear analytic
RG class. Covariant physical Wilson coefficients remain independent
matching data.

\subsection*{Open}
\addcontentsline{toc}{subsection}{V13: Open}
The smallest ultraviolet task is now the \emph{assembled} upper-covariance
comparison \eqref{eq:v13-cap}, including its mixed pairs and source
normalizer. First determine whether the discontinuous principal-symbol
pieces cancel in the complete signed metric/ghost numerator; if not,
construct and record the high-frequency action/reference change and
its matching jets. The next symmetry task is the corresponding complete
quantum Slavnov coefficient with antifield and measure operands, not
another assertion of the classical quartic identity. Full mixed ESM
shell closure, low-scale/infrared transport, chiral-line coherence and
the filtered invariant trajectory remain open. No nonperturbative
positive gravity measure, finite-parameter ultraviolet completion or
mass gap is concluded.

\begin{center}\small
\begin{tabular}{>{\raggedright\arraybackslash}p{.22\textwidth}>{\raggedright\arraybackslash}p{.69\textwidth}}
\toprule
Variable & Scope in V13\\
\midrule
Volume & Free bounds are modewise. Interior loop and rooted bounds require
$\Lambda L_\nu\ge1$ and have no extensive multiplier; raw local density
coefficients can grow as $a^{-4}$.\\
Mesh & The free reference bound covers the whole principal cube. Smooth
vertex/shell bounds use $a\Lambda\le1/512$ and nonaliased panels. The
full-zone cap is explicitly unproved.\\
Orders & One propagating loop and two retained geometric marks. Every fixed
momentum derivative and kernel moment is controlled, not unbounded source
order or the full perturbation series.\\
Iteration & Interior subtracted increments sum as $\Lambda_j^{-1}$.
The upper scale tends to infinity as $a^{-1}$, but its exterior covariance
and mixed pairs remain a distinct comparison.\\
Reference & New explicitly defined holomorphic coefficient substitution
and conformal contour; no real-time Wick-rotation theorem. Zero and low
modes remain retained at a fixed lower Wilsonian scale.\\
Symmetry & The four ghost determinant and its seagull are included. No
finite-shell quantum BV identity or transverse projection is assumed.\\
Matter & A two-Higgs retained-source contact is evaluated. The Higgs,
internal gauge and fermion propagating blocks are not all integrated here.\\
\bottomrule
\end{tabular}
\end{center}

\begingroup
\renewcommand{\refname}{Additional methodological references for V13}
\begin{thebibliography}{99}
\bibitem[38]{V13BFR}
R.~Brunetti, K.~Fredenhagen and K.~Rejzner,
\emph{Locally covariant approach to effective quantum gravity},
arXiv:2212.07800v2 (2022). Established perturbative/BV framework;
used here for context, not as a proof of native lattice estimates.
\bibitem[39]{V13DM}
M.~D'Attanasio and T.~R.~Morris,
\emph{Gauge invariance, the quantum action principle, and the
renormalization group}, Physics Letters B \textbf{378} (1996), 213--221,
arXiv:hep-th/9602156, doi:10.1016/0370-2693(96)00411-X.
\end{thebibliography}
\endgroup

\subsection*{Verification record}
The accompanying \srcid{check_ESM_propagating_shell_V13.py} distinguishes
exact finite algebra from numerical diagnostics. It verifies the native
Cartan normalization of the free Hessian, the ghost matrix and seagull,
the rational global margin, the conformal two-jet, the retained-Higgs
tensor contraction, the upper-boundary inverse jump, and mixed-shell
ownership. Numerical contour floors, Gaussian trace derivatives and
mesh-scaling diagnostics are checks of implementation, not substitutes
for the estimates proved above. No new proof-assistant formalization is
asserted.



\clearpage
\providecommand{\tr}{\operatorname{tr}}
\part*{V14: A signed boundary obstruction and an explicit periodic ultraviolet completion}
\addcontentsline{toc}{part}{V14: Signed boundary test and periodic ultraviolet completion}

\section{The upper covariance must be specified at the action level}
\label{sec:v14-start}

Sections 1--103 and their labels are preserved. V13 established a
propagating metric--ghost coefficient for interior covariances and left
\eqref{eq:v13-cap} open. Its free ghost jump alone did not prove a jump
in the signed loop. This installment first tests the complete signed
loop, including the metric seagull, at that boundary. For an explicitly
specified cyclic realization and a legitimate localized upper
innovation it has a nonzero cusp. Thus a uniform class containing all
such native innovations cannot be justified by local Taylor subtraction.
The test does not assert that integration over an entire unlocalized
face fails to cancel, nor that every possible realization of aliased
products has the same defect.

The constructive response is a \emph{recorded change of the ultraviolet
regulator}, not deletion of its upper covariance. We specify smooth
periodic free matrices and cubic/quartic interactions on the same finite
carrier. They agree exactly with V13 on its interior panels, retain all
high modes, and give those modes nonzero metric and ghost interactions.
For this completed family the full upper contribution, including both
mixed upper/interior pairs, has an explicit local matching jet and a
vanishing source remainder. The resulting full-zone one-loop metric
two-jet has the same nonlocal continuum limit as V13's interior family.
Comparison with the original unmodified cyclic regulator is not inferred
from that theorem.

\subsection*{Inputs and limits}
The finite common gravity--matter carrier, the native connection words,
V11--V12's restored classical jets and coframe chart, and V13's fixed
holomorphic coefficient reference are inherited. The local Lorentz/Haar
amplitude and the metric-coordinate Jacobian retain their original
ownership. No new real-time continuation, physical graviton
normalization, internal-gauge determinant or fermionic phase is introduced.
The calculation below concerns one propagating loop and two geometric
marks at a fixed lower Wilsonian scale. Locality is measured by every
\emph{prescribed finite} polynomial kernel moment; exponential locality
or uniformity in unbounded source order is not asserted.

The audit of the available monographs found no prior signed conformal
boundary calculation or the periodic action completion below. General
Gaussian/Berezin composition, mixed-shell ownership and subtracted
interior estimates are reused. Lattice power counting and the hazards of
nonlocal spectral derivatives are established subjects
\cite{V6Reisz,V14Karsten}; those general subjects are not the novelty.
The assertions identified as proved below concern the displayed native
operands and the explicitly changed finite action.

\section{The signed conformal boundary coefficient}
\label{sec:v14-face}

\subsection{A literal cyclic convention}
For this diagnostic use the principal momentum representative in every
spectral derivative, cyclic multiplication of sampled fields, and the
ordered operator expressions in \eqref{eq:v10-action},
\eqref{eq:v11-artifact-bank} and \eqref{eq:v12-counterterm}. In particular,
a derivative on a cofactor product acts on its \emph{total cyclic label};
no continuum product rule is inserted before sampling. Stationary
elimination uses the same finite pairing. This fixes a realization at
aliased momenta, which V13 had deliberately not identified across all
possible polynomial representatives.

The test is the repeated conformal source
$h(x)=t(x)I_4$, with opposite transfers $k$ and $-k$. The actual coframe is
$e_E=(1+t/2)I_4$ on this background. A smooth nonnegative weight $w(ap)$,
periodic on the momentum torus and supported away from zero, multiplies
both propagating covariances. All ten metric components and four ghost
pairs are included. The metric contour remains V13's $T$.

Here is a basis-independent finite recipe for the mixed geometric
quadratic form needed at a face. Let $\mathsf L_v$ be the $24\times10$
linear load matrix obtained from \eqref{eq:v11-ell} by the substitution
\eqref{eq:v13-Wick-map}, with ordinary covector $v$ in the load slot.
The target quadratic pairing is bilinear: a superscript $T$, not complex
conjugation, occurs in this recipe. Put
\begin{equation}
\begin{split}
 \mathsf H_a(p,q)={}&-\mathsf L_{d_a(p)}^TC_0^{-1}\mathsf L_{d_a(q)}\\
 &+a^2\left(\mathsf L_{u(p)}^TC_0^{-1}\mathsf L_q
                   +\mathsf L_p^TC_0^{-1}\mathsf L_{u(q)}\right).
 \label{eq:v14-cross-form}
\end{split}
\end{equation}
All normalizations, including $\chi$, are those of V13. Thus
$\mathsf H_a(p,q)^T=\mathsf H_a(q,p)$ and
$K_a(p)=\mathsf H_a(p,p)+\chi F_p^*F_p$.
One may use raw independent symmetric coordinates or a Frobenius
orthonormal basis, provided both Hessians and inverses use the same
basis; the traces below are unchanged.

\begin{lemma}[Boundary-leading hard conformal Hessian]
\label{lem:v14-conformal-operator}
Let $r=1+t/2$. Up to terms carrying at least one ordinary derivative of
an external $t$, the corrected geometric hard Hessian about $rI_4$ is
the operator obtained from
\[
 M_r\mathsf L^T M_{r^{-2}}C_0^{-1}\mathsf L M_r
\]
with precisely the signs and the two phase-cross terms in
\eqref{eq:v14-cross-form}. Its gauge-fixing addition is independent of
$t$. Write $P=\mathsf H_a(p,p)$, $Q=\mathsf H_a(q,q)$ and
$X=\mathsf H_a(p,q)$ for two limiting hard representatives at opposite
sides of one face. The limiting one-source vertex and the two limiting
seagulls are
\begin{align}
 A_{pq}&=\tfrac12(P+Q)-X,\\
 B_p&=\tfrac34P+\tfrac14Q-\tfrac12(X+X^T),\\
 B_q&=\tfrac34Q+\tfrac14P-\tfrac12(X+X^T).
 \label{eq:v14-face-vertices}
\end{align}
The one-source ghost vertex in the same limit is
\begin{equation}
 N_{pq}=\tfrac12\{(p\cdot q)I+qp^T-pq^T\},
 \qquad N_{qp}=N_{pq}^T,
 \label{eq:v14-ghost-face}
\end{equation}
and its repeated conformal seagull is zero.
\end{lemma}
\begin{proof}
The cofactor is quadratic in $e$. Its derivative at $rI$ in a hard
coframe direction is $r$ times its derivative at $I$. Therefore the
linearized load is $\mathsf L M_r$, while $\mathsf C(rI)=r^2C_0$.
Terms involving the background stationary connection have an external
derivative: its load vanishes at constant $r$, and its inverse is the
pointwise $C_0^{-1}$ at the displayed degree. Every cubic-or-higher
connection word in a two-hard-leg coefficient also has a soft
connection factor. The stationary returns of those words have the same
property. They vanish in the zero-transfer face limit. This statement
uses the full native quartic return, not an omitted return.

The geometric part of the mesh correction is obtained by replacing one
of the two phase loads by $\mathsf L_u$. Consequently it has exactly the
same three multiplication placements; this proves the operator claim.
For independent scalar source amplitudes, expand
$r=1+t/2$ and $r^{-2}=1-t+3t^2/4+\cdots$ without moving a derivative
through a product. The two outer first variations give $(P+Q)/2$;
the middle variation gives $-X$. At opposite transfers the middle
second variation gives $3P/2$, the two outer variations give
$[\mathsf H(p+k,p+k)+\mathsf H(p-k,p-k)]/4$, and the mixed
outer/middle terms give minus one half of the four corresponding cross
forms. On one crossing strip one shifted representative tends to $q$
and the other to $p$. This gives $B_p$; the opposite strip gives $B_q$.

The chart formula is exact on a variable conformal background:
$R(tI,c)=(1+t/2)R_0c+(c\cdot\partial t)I$.
Applying $-F$ and then setting the \emph{small physical transfer} to zero
while retaining the two different hard representatives gives
\eqref{eq:v14-ghost-face}. There is no quadratic $t$ term in this
formula. No replacement of the free $M_a$ by $|p|^2I$ is made.
\end{proof}

Use dimensionless momenta for the following traces: set $a=1$, $\chi=1$
in the finite recipe and take
$p=(y_0,\pi,y_2,y_3)$, $q=(y_0,-\pi,y_2,y_3)$. Let
$C_p=K_1(p)^{-1}$, $C_q=K_1(q)^{-1}$,
$D_p=M_1(p)^{-1}$ and $D_q=M_1(q)^{-1}$. Define the signed face data
\begin{align}
 L(y)&=\tfrac12\tr(C_pB_p+C_qB_q),\\
 G(y)&=L(y)-\tfrac12\tr(C_pA_{pq}C_qA_{pq}^T),\\
 C(y)&=\tr(D_pN_{pq}D_qN_{pq}^T)
                         -\tfrac14|p|^4\tr(D_p^2),\\
 \Sigma(y)&=G(y)+C(y).
 \label{eq:v14-signed-face}
\end{align}
The last subtraction is the same-side ghost bubble. Its two endpoint
values agree by reflection. $L$ is the metric-seagull contribution;
$G-L$ and $C$ include both signed bubbles. The factors of $\chi$ cancel
between metric vertices and propagators. The same trace result is
obtained after the fixed conformal contour transformation.

\begin{proposition}[An evaluated nonzero signed face]
\label{prop:v14-signed-point}
At $y_*=(\pi/2,0,0)$ the quantities in \eqref{eq:v14-signed-face} obey
\begin{equation}
 \boxed{\frac{199}{20}<L(y_*)<\frac{997}{100},\qquad
              \frac{28}{25}<\Sigma(y_*)<\frac{113}{100}.}
 \label{eq:v14-exact-interval}
\end{equation}
In particular the metric seagull, metric bubble and ghost bubble do not
cancel pointwise at this face. Approximate values, not used as proof,
are $L=9.9573423127$, $G=1.4200568344$,
$C=-0.2965224691$, and $\Sigma=1.1235343652$.
\end{proposition}
\begin{proof}
The following rational expressions provide an exact, reproducible
certificate of the finite traces. In them $x$ is an algebraic placeholder
for $\pi$, $p=(x/2,x,0,0)$ and the limiting half-sine in direction 1
is $2$, not $2\sin(x/2)$ away from the final value $x=\pi$. Set
\begin{align*}
 d_1&=x^4+3x^2+48, &d_2&=x^6-288x^4+96x^3-144x^2-2304x-9216,\\
 d_3&=x^3+6x+48, &d_4&=x^6+96x^3+144x^2+2304.
\end{align*}
Direct cofactor-load multiplication and inversion give
\begin{equation}
 L=\frac{2n_1}{d_1d_2},\qquad
 G=-\frac{24x n_2}{d_1^2d_2^2},\qquad
 C=-\frac{144x^2 n_3}{d_3^2d_4^2},
 \label{eq:v14-rational-face}
\end{equation}
where
\begin{align*}
 n_1={}&5x^{10}-1653x^8+480x^7-1848x^6-10080x^5
       -137088x^4\\
      &+19584x^3-89856x^2-497664x-2764800,\\
 n_2={}&x^{17}+132x^{15}+96x^{14}-91278x^{13}+40896x^{12}
       -138888x^{11}\\
      &-1416672x^{10}-9892512x^9+919296x^8-16588800x^7
       -116909568x^6\\
      &-359313408x^5-91570176x^4-467140608x^3
       -2325086208x^2\\
      &-4586471424x-2038431744,\\
 n_3={}&x^{12}+36x^{10}+192x^9+540x^8+5184x^7+8640x^6
       +51840x^5\\
      &+269568x^4+193536x^3+1244160x^2+3981312x+5308416.
\end{align*}
These formulas result from the $24\times24$ Cartan matrix with the
six upper Lorentz-pair coordinates; raw symmetric metric coordinates
avoid square-root basis factors. They can also be verified by multiplying
the proposed inverses by the finite matrices before taking traces.
To verify the strict intervals, use
\[
 \frac{3141592653589793}{10^{15}}<\pi<
 \frac{3141592653589794}{10^{15}}.
\]
This enclosure follows from Machin's identity
$\pi=16\arctan(1/5)-4\arctan(1/239)$ and respectively twenty and
eight terms of the alternating arctangent series, including their next
term remainders. Rational interval Horner evaluation of
\eqref{eq:v14-rational-face} gives \eqref{eq:v14-exact-interval}.
The denominator intervals exclude zero. All arithmetic in this
certificate is rational; floating-point quadrature is unnecessary.
\end{proof}

\section{A legitimate upper innovation has a nonlocal cusp}
\label{sec:v14-cusp}

\begin{theorem}[Signed localized upper-shell obstruction]
\label{thm:v14-cusp}
There is a smooth real even periodic weight $0\le w\le1$, supported
in fixed dimensionless neighborhoods of the $p_1$ face and away from
zero and the other faces, for which the literal cyclic realization above
has the thermodynamic conformal two-source coefficient
\begin{equation}
 \boxed{\Pi^w_a(ke_1)-\Pi^w_a(0)
      =\hbar a^{-3}\mathfrak c_w |k|
                    +O(\hbar a^{-2}k^2),\qquad \mathfrak c_w>0.}
 \label{eq:v14-cusp}
\end{equation}
The metric and ghost weights are the same $w$, both metric seagulls and
bubbles are retained, and the inherited auxiliary density has no cusp.
No polynomial of fixed degree in $k$, even with cutoff-dependent
coefficients, can make this family bounded on a fixed real momentum
interval as $a\downarrow0$. In particular its thermodynamic Taylor
subtractor through degree four is not defined at zero.
\end{theorem}
\begin{proof}
By \cref{prop:v14-signed-point} and continuity of the finite inverses,
choose a small transverse neighborhood of $y_*$, and its reversal, on
which $L>0$ and $\Sigma>0$. Let $w$ have nonzero trace $v(y)$ on
that face, $0\le v\le1$, supported in these neighborhoods. Take its
normal factor to equal one near the identified face and vanish away
from a slightly larger neighborhood; then extend periodically and
ensure $w(-p)=w(p)$. This is a genuine covariance innovation in V13's
accretive references, with a nonzero high-frequency complement.

For $\varepsilon=ak>0$, partition the dimensionless momentum integral
into the crossing strip of thickness $\varepsilon$, the opposite strip
needed by the seagull, and the noncrossing region. On the strips each
rational symbol has a smooth one-sided limit and an $O(\varepsilon)$
remainder, uniformly on the chosen compact support. The metric seagull
is linear in $v$ and the two bubbles are quadratic in $v$. Therefore
the crossing coefficient is
\begin{equation}
 \mathfrak c_w=\frac1{(2\pi)^4}\int
       \{v(y)(1-v(y))L(y)+v(y)^2\Sigma(y)\}\,\dd^3y>0.
 \label{eq:v14-face-integral}
\end{equation}
There is one bubble crossing strip and two seagull strips. These give
exactly the factors in \eqref{eq:v14-signed-face}, with no extra factor
of two. For the same-side ghost bubble, symmetry of its two endpoints
makes the ordinary first derivative integral one half of the endpoint
difference; that difference is zero. Its crossing strip therefore
subtracts the same-side value displayed in \eqref{eq:v14-signed-face}.
The metric same-side cubic is zero for a constant conformal source.

The terms omitted in \cref{lem:v14-conformal-operator} have an explicit
factor $k$. Their coefficients and one-sided derivatives are bounded on
the compact upper support. Integrating across a strip changes such a
term by $O(k|ak|)$ in scaled variables. The remaining bulk coefficient
has an ordinary derivative at zero, by dominated convergence. The
complete repeated-source Hessian is even in $k$ by symmetry of its
second variation, so its smooth linear part is zero. The total error
is $O(\varepsilon^2)$ before restoring the overall $a^{-4}$ scaling.
The same reasoning at negative transfer gives $|\varepsilon|$ and
proves \eqref{eq:v14-cusp}. Auxiliary graphs at this source degree
contain only low external/subset momenta and a pointwise flat Cartan
inverse; their finite coefficients are analytic in $ak$. They cannot
alter the cusp.

For the polynomial assertion fix a compact interval containing zero.
If $\Pi^w_a-P_a$ were uniformly bounded there for polynomials of fixed
degree, then $a^3(P_a-\Pi^w_a(0))$ would converge uniformly to
$\hbar\mathfrak c_w|k|$, by \eqref{eq:v14-cusp}. The space of
polynomials of that fixed degree is closed in the uniform norm, whereas
$|k|$ does not belong to it. This is a contradiction.
\end{proof}

The theorem is a \emph{localized native-shell} statement with an explicit
cyclic convention. It proves that pointwise signed cancellation and a
uniform local class allowing arbitrary such angular innovations are
false. It does not evaluate the integral of the face coefficient over
an entire unweighted boundary, and hence does not prove that the
unlocalized full native cap has no special cancellation. Different
aliased derivative placements also require their own test. These
qualifications matter: the conclusion licenses a recorded regulator
repair, not a blanket no-go result for the Renewal carrier.

\section{A smooth periodic completion with interacting upper modes}
\label{sec:v14-completion}

Fix smooth even radial profiles on the small central chart of the
four-dimensional momentum torus, with $0\le s,\zeta\le1$, such that
\begin{align}
 s(z)&=1\quad(|z|\le1/64),&s(z)&=0\quad(|z|\ge1/32),\\
 \zeta(z)&=1\quad(|z|\le1/16),&\zeta(z)&=0\quad(|z|\ge1/8).
 \label{eq:v14-profiles}
\end{align}
They are extended by zero periodically; all derivatives vanish near
the faces. Let $S=s(aD)$, $H=I-S$, and
\begin{equation}
 \lambda_a(p)=\frac4{a^2}\sum_{\nu=0}^3\sin^2(ap_\nu/2).
 \label{eq:v14-lambda}
\end{equation}
This \emph{upper reference} includes a nearest-neighbor temporal symbol,
but it does not replace V13's temporal symbol in the interior.
Define
\begin{align}
 \widetilde K_a(p)&=\zeta(ap)K_a(p)
              +(1-\zeta(ap))\frac\chi2\lambda_a(p)J,\\
 \widetilde M_a(p)&=\zeta(ap)M_a(p)
              +(1-\zeta(ap))\lambda_a(p)I_4.
 \label{eq:v14-completed-free}
\end{align}
These equations are finite cutoff-dependent quadratic additions to the
old gauge-fixed action and ghost action. They are not field redefinitions.
The original nonminimal source contact \eqref{eq:v13-b-contact} is retained.

Let $\sigma(q)=\operatorname{tr}(q)/4$ and choose four real bounded
parameters $\gamma_g,\beta_g,\gamma_c,\beta_c$. Put
\begin{equation}
 \begin{split}
 V^g_{\mathrm{up},a}(q)=\frac\chi4 a^4\sum_{x,\nu}
 \big\{\gamma_g\sigma(Sq)_x+\beta_g\sigma(Sq)_x^2\big\}
 (D_\nu^+Hq)_x^T J(D_\nu^+Hq)_x,
 \end{split}
 \label{eq:v14-upper-action}
\end{equation}
with tensor contraction in the Frobenius metric. Define the completed
cubic and quartic metric action by
\begin{equation}
 \widetilde S^E_{3,a}+\widetilde S^E_{4,a}
       =S^E_{3,a}(Sq)+S^E_{4,a}(Sq)+V^g_{\mathrm{up},a}(q).
 \label{eq:v14-completed-interactions}
\end{equation}
The native pieces on the right are evaluated on their ordinary low
Fourier chart. All four-leg subset sums are strictly inside the zone;
there is no choice of aliased representative left in these pieces.
The additional ghost interaction is
\begin{equation}
 a^4\sum_{x,\nu}(D_\nu^+H\bar c)_x^T
 \{\gamma_c\sigma(Sq)_x+\beta_c\sigma(Sq)_x^2\}
                         (D_\nu^+Hc)_x.
 \label{eq:v14-upper-ghost}
\end{equation}
The native ghost cubic and quartic are evaluated on $Sq,Sc,S\bar c$.
No extra fields have been introduced. None of
\eqref{eq:v14-upper-action}--\eqref{eq:v14-upper-ghost} is asserted to
be a diffeomorphism-invariant physical operator. They are explicit
ultraviolet regulator data; physical local normalization is assigned
below.

\begin{theorem}[A periodic, nondegenerate finite reference on the same carrier]
\label{thm:v14-completed-reference}
The free matrices \eqref{eq:v14-completed-free} are smooth periodic
symbols off their common zero mode. For every nonzero principal
momentum,
\begin{equation}
 \boxed{\operatorname{Re}(T^T\widetilde K_a T)
             \succeq\frac\chi5|p|^2I_{10},\qquad
        \operatorname{Re}\widetilde M_a
             \succeq\frac25|p|^2I_4.}
 \label{eq:v14-completed-floor}
\end{equation}
The zero and prescribed low modes remain retained. Every compact upper
region $|ap|\ge\delta>0$ has inverse derivative bounds
\begin{equation}
 \|\partial_p^\alpha\widetilde K_a^{-1}\|
       \le C_{\alpha,\delta}\chi^{-1}a^{2+|\alpha|},\qquad
 \|\partial_p^\alpha\widetilde M_a^{-1}\|
       \le C_{\alpha,\delta}a^{2+|\alpha|}.
 \label{eq:v14-upper-inverse}
\end{equation}
The completed interaction vertices and all fixed derivatives are smooth
periodic in every hard momentum. On upper regions their first and
second metric derivatives have bounds
$C_\alpha\chi a^{|\alpha|-2}$ for metric vertices and
$C_\alpha a^{|\alpha|-2}$ for ghost vertices.
All coefficients coincide \emph{exactly} with V13 when every field leg
has $|ap|\le1/64$. In particular the entire V13 interior calculation
through $J(a)$ is unchanged.
\end{theorem}
\begin{proof}
The lattice inequality $4|p|^2/\pi^2\le\lambda_a(p)\le|p|^2$
holds on the principal cube. Since $T^TJT=I$, convex combination of
V13's accretive matrix with the scalar upper reference gives the lower
constants $\min(9/40,2/\pi^2)>1/5$ and
$\min(9/10,4/\pi^2)>2/5$. The cutoffs remove every nonperiodic
matrix entry in a full neighborhood of each face. In dimensionless
variables the matrices are $a^{-2}$ times smooth matrices on the
compact punctured torus. The inverse identity and the lower bounds
prove \eqref{eq:v14-upper-inverse}.

Each native cubic or quartic Fourier coefficient is a fixed finite
expression in phase loads, translated connection words, the pointwise
Cartan inverse and the stated restoring polynomials. On the support of
all its $s$ factors the subset labels have no aliasing and its
dimensionless derivatives are bounded. Extension by zero is smooth.
The upper vertices are products of smooth periodic $s$, the forward
symbols $(e^{iap_\nu}-1)/a$, and a low source multiplier. Differentiating
therefore gives the asserted bounds. If all field legs are in the
smaller ball, $s=\zeta=1$ and every upper term has a factor $H=0$.
V13's hard momenta through $J(a)$ have size at most $3/(512a)$, which
is smaller than $1/(64a)$, proving the last assertion.
\end{proof}

These additions are quasilocal regulator modifications, with rapid
polynomial decay in lattice distance at every fixed derivative order;
they are not claimed to be finite-range stencils or a finite local EFT
bank by themselves. They do not go to zero in relative norm on the
uppermost modes. The low classical source jets nevertheless agree
exactly with the inherited action for sufficiently fine mesh. The
six Lorentz directions, auxiliary connection, density and coframe chart
are unchanged. A pre-elimination implementation adds the displayed
coframe-only difference and the specified ghost changes before the same
connection pushforward. This preserves the auxiliary stationary section
and its previously assigned amplitude. Equality with a microscopic
acceptance row for the old action does not follow.

\section{The complete upper contribution and its local assignment}
\label{sec:v14-cap}

Fix $\Lambda_0\ge32\mu>0$ and the same low profile $\theta$ as in
V13. Let
\begin{equation}
 \widetilde C^g_a=(1-\theta(p/\Lambda_0))\widetilde K_a^{-1},\qquad
 \widetilde C^c_a=(1-\theta(p/\Lambda_0))\widetilde M_a^{-1}.
 \label{eq:v14-full-cov}
\end{equation}
For sufficiently small $a$ the low profile has support inside the
native region. These are full-zone covariances on the retained-complement
carrier, not just interior bands. Their source vertices are the
Hessians of \eqref{eq:v14-completed-interactions} and
\eqref{eq:v14-upper-ghost}, with all derivatives of the common action
included.

For example, on the source window $Sh=h$, $Hh=0$, the extra vertices are
\begin{align}
 A^g_{\rm up}[h]&=\frac{\chi\gamma_g}{2}
      H\sum_\nu(D_\nu^+)^* M_{\sigma(h)}J D_\nu^+H,\\
 B^g_{\rm up}[h_1,h_2]&=\chi\beta_g
      H\sum_\nu(D_\nu^+)^* M_{\sigma(h_1)\sigma(h_2)}J D_\nu^+H,\\
 A^c_{\rm up}[h]&=\gamma_c
      H\sum_\nu(D_\nu^+)^* M_{\sigma(h)}D_\nu^+H,\\
 B^c_{\rm up}[h_1,h_2]&=2\beta_c
      H\sum_\nu(D_\nu^+)^* M_{\sigma(h_1)\sigma(h_2)}D_\nu^+H.
 \label{eq:v14-upper-vertices}
\end{align}
Their factors are fixed by differentiating the single displayed action.
The native contributions have one $s$ on each hard leg and the source
factor $s=1$. Hence this formula does not discard an upper seagull.

\begin{proposition}[Measured full-zone two-jet and analytic background germ]
\label{prop:v14-measured}
For real sources with $\|h\|_{\mathrm W}<\rho$, where $\rho>0$
is independent of mesh and volume in the stated domain, the completed
quadratic-hard-field Gaussian on $T$ and the ghost integral have a
common analytic source germ. The one-loop Hessian is exactly the
four-term expression \eqref{eq:v13-four-terms}, using the completed
covariances and vertices. The auxiliary contribution is still
$\hbar D^2\mathcal M^E_{a,1}$, assigned once. The same formula is the
one-loop, two-retained-metric coefficient of the full polynomial
pushforward through the specified interaction degree.
\end{proposition}
\begin{proof}
On the support of the integrated covariance a source transfer at most
$\mu$ leaves the hard input/output momenta comparable. In the native
region use V13's relative estimates; on the transition and upper regions
use \eqref{eq:v14-upper-inverse} and the order-two vertex bounds.
The Fourier convolution Schur estimate gives a relative perturbation
$C(\|h\|_{\mathrm W}+\|h\|_{\mathrm W}^2)$ uniformly in the
number of modes. Choose $\rho$ below the free contour and inverse
margins. Gaussian/Berezin integration then yields the signed logarithms
and their four-term Hessian. For the full perturbative coefficient,
linear hard forces with two or three soft fields have Fourier support
below $3\mu<\Lambda_0$; translationally invariant filters do not change
that support. They cannot couple to an integrated covariance. Thus the
possible hard-linear one-loop returns vanish for the same explicit
support reason as in V13. Metric and ghost seagulls are not removed in
this argument. The auxiliary amplitude belongs to loop order one;
contracting any of its hard fields adds a propagating loop and is a
higher-order term, explaining its once-only assignment here.
\end{proof}

Let $C_<^{g,c}$ be V13's cumulative covariance through $J(a)$ and put
$C_>^{g,c}=\widetilde C_a^{g,c}-C_<^{g,c}$. On the interior support
the matrices agree, so these differences are the corresponding
nonnegative scalar weights times the same sector inverses there, and
the full completed inverses elsewhere. They are genuine complementary
innovations. At the two-source level define
\begin{equation}
 \widetilde{\mathfrak C}_{a,L}
      =D^2\widetilde\Gamma^{\rm prop}_{a,L}
                         -D^2\Gamma^{\rm prop}_{a,\le J(a),L}.
 \label{eq:v14-actual-cap}
\end{equation}
If $A,B$ denote the \emph{same completed vertices} in both terms,
each sector's cap is the appropriate determinant sign times
\begin{equation}
 \Tr\{C_>B-C_>A_1C_>A_2-C_>A_1C_<A_2-C_<A_1C_>A_2\}.
 \label{eq:v14-cap-ownership}
\end{equation}
The smaller determinant computed with the completed vertices equals
V13's determinant exactly. Thus \eqref{eq:v14-cap-ownership} proves,
rather than assumes, the ownership of both mixed pairs.

\begin{theorem}[Smooth full-zone cap estimate]
\label{thm:v14-cap}
Assume the stated profiles have fixed finite derivative bounds, the
four upper parameters are bounded, and $a\mu$ and $a\Lambda_0$ are
sufficiently small. For every fixed multi-index $\alpha$,
\begin{equation}
 \left|\partial_k^\alpha\widetilde{\mathfrak C}_{a,L}(k)[H_1,H_2]\right|
 \le C_\alpha\hbar\|H_1\|\|H_2\|a^{|\alpha|-4}.
 \label{eq:v14-cap-raw}
\end{equation}
Store its \emph{actual} local Taylor polynomial through degree four.
For $r=|\alpha|\le5$ the remainder satisfies
\begin{equation}
 \boxed{\left|\partial_k^\alpha(I-\mathsf T_4)
       \widetilde{\mathfrak C}_{a,L}(k)[H_1,H_2]\right|
 \le C_\alpha\hbar\|H_1\|\|H_2\|a|k|^{5-r}.}
 \label{eq:v14-cap-ren}
\end{equation}
For $r>5$ the bound is $C_r\hbar a^{r-4}\|H_1\|\|H_2\|$.
For every fixed windowed kernel moment $b$, the two-source rooted kernel
has bound $C_b\hbar a\mu^5$, with no volume factor. Constants are
uniform in all four periods; the separate interior theorem retains its
condition $\Lambda_0L_\nu\ge1$.
\end{theorem}
\begin{proof}
At least one cap weight vanishes unless $|p|\ge\Lambda_{J(a)}$;
on a bubble the other hard momentum differs by at most $\mu$.
For small $a$ both momenta therefore satisfy $|ap|\ge c>0$, with
$c$ independent of $a$ because
$1/1024<a\Lambda_{J(a)}\le1/512$. The low cutoff and all its
derivatives vanish on this cap support. The only further scale parameter
is $a\Lambda_{J(a)}$, confined to that fixed compact interval.

A propagator contributes $a^2$, a hard vertex $a^{-2}$, and every
momentum derivative contributes $a$. Both the seagull and bubble
integrands in \eqref{eq:v14-cap-ownership} thus have derivative bounds
$C_\alpha a^{|\alpha|}$. These are bounds for smooth periodic
integrands; unlike in the preceding diagnostic they remain valid when
$p+k$ crosses a face. A finite momentum sum divided by physical volume
has exactly $a^{-4}$ times the normalized average over its dimensionless
grid. This proves \eqref{eq:v14-cap-raw}, including derivatives that
fall on the cap weights, and equally proves the integral version.

Apply the multivariate integral Taylor formula through degree four.
Its differentiated versions use fifth total derivatives for $r\le5$,
giving \eqref{eq:v14-cap-ren}; for $r>5$ the polynomial is absent.
Multiply by a fixed smooth source window supported in $|k|<\mu$.
The scaled symbol derivatives have bounds
$C_s\hbar a\mu^{5-s}$, using $a\mu<1$ also for $s>5$.
Fourier integration by parts more than $b+4$ times, followed by sampling
and periodization, proves the rooted kernel bound. Periodic distance is
bounded by the distance of each lift. This last step yields a bilinear
functional bound with one source in $L^1$ and one in $L^\infty$, not
an extensive determinant estimate.
\end{proof}

The polynomial assigned here can have the sizes $a^{r-4}$ at derivative
order $r=0,1,2,3,4$. It is retained in V13's redundant component bank
of dimension $3610$; no parity or covariant-basis reduction is assumed.
The physical local polynomial is imposed only after this actual matching.
Neither a small bare vertex at fixed momentum nor an upper inverse floor
would have implied this conclusion without periodic smoothness.

\section{A full-zone continuum coefficient and matched completion independence}
\label{sec:v14-limit}

\begin{theorem}[Full-zone continuum for the declared periodic completion]
\label{thm:v14-limit}
With a common physical local two-metric polynomial, retain the same
once-assigned auxiliary term in every scheme. Let
$\widetilde\Pi^R_{a,L}$ denote the completed full-zone propagating plus
auxiliary coefficient after their actual local Taylor assignments.
For fixed $\Lambda_0\ge32\mu$, $\Lambda_0L_\nu\ge1$, and
$r=|\alpha|\le5$,
\begin{equation}
 \boxed{\left|\partial_k^\alpha
  (\widetilde\Pi^R_{a,L}-\Pi^R_{0,\infty,L})(k)[H_1,H_2]\right|
 \le C_\alpha\hbar\|H_1\|\|H_2\|a|k|^{5-r}.}
 \label{eq:v14-full-limit}
\end{equation}
Here $\Pi^R_{0,\infty,L}$ is exactly V13's continuum reference with
its lower Wilsonian cutoff, not an infrared limit. Every fixed
windowed source-kernel moment has error $C_b\hbar a\mu^5$.
The conclusion is uniform over profile families with the required
fixed-order seminorm and margin bounds and bounded upper parameters.
\end{theorem}
\begin{proof}
The completed coefficient is the sum of the \emph{unchanged} V13
interior coefficient and \eqref{eq:v14-actual-cap}. Taylor subtraction
is linear and all three summands use their actual source jets. Apply
\cref{thm:v13-interior-limit} to the first summand and
\cref{thm:v14-cap} to the second; the inherited auxiliary remainder is
assigned once. This proves \eqref{eq:v14-full-limit}. The same
addition in the fixed source-kernel topology proves the moment bound.
No assertion about the discontinuous unmodified cap is used.
\end{proof}

\begin{proposition}[Actual local matching between two periodic completions]
\label{prop:v14-matching}
Take two completions of the displayed kind. Their propagating two-jet
difference $\mathcal D_{a,L}$ has, for sufficiently small $a$, the form
\begin{equation}
 \mathcal D_{a,L}(k)=\hbar a^{-4}\mathcal F_{\boldsymbol n}(ak),
 \qquad n_\nu=L_\nu/a,
 \label{eq:v14-matching-scale}
\end{equation}
where $\mathcal F_{\boldsymbol n}$ is the normalized dimensionless
momentum sum of the \emph{difference of their complete four-term
integrands}. It is smooth with fixed-order bounds uniform in
$\boldsymbol n$. Thus the local matching is precisely
\begin{equation}
 \partial_k^\alpha\mathcal D_{a,L}(0)
     =\hbar a^{|\alpha|-4}
                     \partial_z^\alpha\mathcal F_{\boldsymbol n}(0),
 \qquad |\alpha|\le4.
 \label{eq:v14-local-matching}
\end{equation}
After those jets are matched, their nonlocal difference is bounded by
\eqref{eq:v14-cap-ren}. For every prescribed $b$, the dimensionless
local coefficients converge to their torus integrals with error
$C_b\sum_\nu n_\nu^{-b}$. The fourth-order matching can therefore
have a finite nonzero limit; it is not set to zero by definition.
\end{proposition}
\begin{proof}
Both completions agree with the native matrices and vertices on a fixed
central ball. At small transfer the difference of complete loop
integrands consequently vanishes near zero loop momentum. On its
support the common lower profile is identically zero, so no
$a\Lambda_0$ dependence remains. Rescale $p=z/a$ in the finite sum.
All entries are dimensionless smooth periodic functions with two
propagator/vertex powers canceling as in \cref{thm:v14-cap}. This
proves \eqref{eq:v14-matching-scale} and its differentiated formula.
The Taylor remainder proof is the same fifth-derivative argument,
now on an even larger fixed upper support. Fourier coefficients of a
smooth periodic integrand decay faster than any prescribed power.
The product-grid trapezoidal identity retains only reciprocal multiples
of $\boldsymbol n$; summing those coefficients proves the displayed
finite-volume error. This comparison uses the actual integrand
difference, not independently fitted local polynomials.
\end{proof}

The upper modes are not spectators. Where $s=\zeta=0$, for a constant
conformal source $h=tI_4$ they have exactly
\[
 K(t)=(1+\gamma_gt+\beta_gt^2)\frac\chi2\lambda_aJ,
 \qquad
 M(t)=(1+\gamma_ct+\beta_ct^2)\lambda_a I_4.
\]
Their per-mode two-source contribution is
\begin{equation}
 \boxed{\hbar\{5(2\beta_g-\gamma_g^2)
                        -4(2\beta_c-\gamma_c^2)\}.}
 \label{eq:v14-upper-value}
\end{equation}
For example $\gamma_g=1$, $\beta_g=\gamma_c=\beta_c=0$ gives
$-5\hbar$ per such mode. With any additional annular covariance
weight, the seagull terms carry one weight and the bubbles two; replacing
this by a sum of diagonal shell values is again incorrect. The corresponding one-point contribution is $\hbar(5\gamma_g-4\gamma_c)$ per
mode and is also stored; it is not a source-independent vacuum scalar.
The quoted $3610$-dimensional bank is the degree-two bank, supplemented
by the actual degree-one tadpole assignment. The
transition region and the inherited auxiliary density have their own
actual contributions. Formula \eqref{eq:v14-upper-value} is a
coordinate-dependent local matching check, not a physical beta function.
It follows immediately by differentiating the two finite determinant
powers, and proves that the full-zone result was not obtained by setting
the upper interactions to zero.

\section{Source and symmetry consequences, with their boundaries}
\label{sec:v14-scope}

At this loop and source order the completed measured step lands in the
finite local bank plus the rooted remainder space with the explicit cap
bound. Interior local increments remain V13's generated increments;
the final upper step adds $-\mathsf T_4\widetilde{\mathfrak C}_{a,L}$
to the local matching and $(I-\mathsf T_4)\widetilde{\mathfrak C}_{a,L}$
to the remainder. These are generated by the modified finite measured
law. This is a two-source coefficient statement, not a self-map on a
neighborhood of arbitrary interacting ESM actions.

The high completion agrees with V11--V12's classical mesh identities
on every fixed sufficiently low source panel. It does not supply an
exact finite BV complex at high momentum. The new free ghost term,
filtered nonlinear ghost vertices and upper terms must all be varied in
a quantum Slavnov calculation. The measure divergence includes the
unchanged auxiliary density and the coframe Jacobian. A two-metric
functional by itself does not determine the required ghost/antifield
panels or their local restoring projection. In particular, neither
transversality nor vanishing of a quantum anomaly is inferred from
\cref{thm:v14-limit}.

There are now two different possible comparisons, rather than one
repeated undefined cap assumption. The periodic family has a proved
full-zone coefficient theorem and matched independence within its
declared smooth class. For the original literal cyclic family, the
signed localized obstruction \eqref{eq:v14-cusp} rules out a blanket
uniform theorem over angular innovations. A possible special cancellation
in the unlocalized native sum remains an independent question. It would
not turn the localized cusp into a local polynomial, and it is not
needed in the positive theorem for the recorded completion.

For the ESM programme, the next constructive symmetry task is the
one-loop Slavnov coefficient with the actual antifield/ghost and measure
insertions on this now periodic reference. The remaining mixed
scalar--graviton, internal-gauge and chiral propagating sectors must
still be assembled on a common reference. The lower Wilsonian cutoff,
Lorentzian continuation, invariant all-order filtered class and physical
local curvature-basis reduction remain separate gates. The programme's
finite-order EFT endpoint is not weakened or promoted to a
nonperturbative ultraviolet claim by changing its regulator here.

\begin{center}\small
\begin{tabular}{>{\raggedright\arraybackslash}p{.23\textwidth}>{\raggedright\arraybackslash}p{.68\textwidth}}
\toprule
Quantity & Scope of V14\\
\midrule
Boundary obstruction & A specified cyclic realization and a smooth localized
upper innovation; thermodynamic conformal two-jet. Not a theorem about
all unlocalized face sums or all aliased conventions.\\
Full-zone positive result & The explicitly completed metric--ghost family,
all high modes retained, fixed lower Wilsonian scale.\\
Mesh and volume & Cap estimates are period-uniform. The inherited interior
limit requires $\Lambda_0L_\nu\ge1$ and sufficiently small $a\Lambda_0$.
No massless infrared removal is included.\\
Orders & One propagating loop and two geometric sources, all prescribed
finite momentum derivatives/kernel moments. No unbounded-order constants.\\
Regulator class & Uniform bounds on fixed-order profile seminorms, fixed
transition radii, bounded upper parameters and the proved inverse margins.
Not arbitrary discontinuous principal symbols.\\
Measure & Same fixed conformal contour and holomorphic reference. Auxiliary
Lorentz/Haar density and coordinate Jacobian assigned once.\\
Symmetry & Low-panel classical jets retained; full quantum BV/Slavnov and
antifield restoration remain unproved.\\
Matching & Actual finite local source jets retained. Equality is proved
between smooth completions, not with the unmodified microscopic action.\\
\bottomrule
\end{tabular}
\end{center}

\section{Final ledgers for V14}
\label{sec:v14-ledgers}

\subsection*{Proved}
\addcontentsline{toc}{subsection}{V14: Proved}
The native cyclic conformal face operands are
\eqref{eq:v14-face-vertices}--\eqref{eq:v14-signed-face}; the exact
rational point certificate \eqref{eq:v14-exact-interval} proves their
signed noncancellation. It yields the positive localized cusp
\eqref{eq:v14-cusp}, including the metric seagull, and excludes a
fixed-degree local-polynomial cure for that family. The new periodic
free matrices and interacting upper action give
\eqref{eq:v14-completed-floor} on the same finite carrier. Their
complete mixed-pair cap has the actual local jet and source estimate
\eqref{eq:v14-cap-ren}. The completed full-zone two-metric coefficient
converges by \eqref{eq:v14-full-limit}, and two declared smooth
completions agree after their actual local matching
\eqref{eq:v14-local-matching}. The nonzero upper-mode check is
\eqref{eq:v14-upper-value}.

\subsection*{Inherited}
\addcontentsline{toc}{subsection}{V14: Inherited}
V1--V13 and all their mathematical labels are unchanged. V10 supplies
the actual Cartan/cofactor matrices, stationary action and auxiliary
measure; V11--V12 supply the restored action/ghost jets and coordinate
Jacobian. V13 supplies the fixed reference, the signed determinant
convention, full-zone free floors and the interior continuum estimate.
General Gaussian/Berezin integration, finite source contacts,
mixed-shell allocation and Fourier kernel estimates are inherited
tools, not new ESM existence results.

\subsection*{Literature input}
\addcontentsline{toc}{subsection}{V14: Literature input}
Reisz's lattice power-counting framework \cite{V6Reisz} and the
noncovariant continuum problems of nonlocal lattice-fermion
prescriptions discussed in \cite{V14Karsten} are methodological
comparisons. Neither supplies the native metric--ghost face traces,
our chosen periodic completion, or its source estimates. No priority
claim is made for ultraviolet improvement, smooth cutoffs, or
perturbative effective gravity.

\subsection*{Conditional}
\addcontentsline{toc}{subsection}{V14: Conditional}
Using the new family as the quantum realization of a particular
microscopic Renewal action requires its own matched reference/action
comparison; this is not the small local Taylor matching within the
smooth class. Full ESM consumption also requires a quantum BV measure
and the uncomputed antifield, mixed matter and chiral sectors. The
fixed-source two-jet theorem alone does not verify all filtered R1--R4
hypotheses or determine physical Wilson coefficients.

\subsection*{Open}
\addcontentsline{toc}{subsection}{V14: Open}
On the constructive periodic branch the smallest new symmetry target
is the complete one-loop Slavnov source coefficient, including the
cutoff/reference variation and coordinate-correct measure. The
unlocalized original cap may still be tested for a special signed
cancellation, but automatic pointwise cancellation has been ruled out
and is not posed again as an assumption. Full interacting ESM
self-mapping, all-loop filtered closure, infrared transport, chiral-line
coherence and the Lorentzian comparison remain open. No positive
nonperturbative gravitational measure or finite-parameter ultraviolet
completion is claimed.

\begingroup
\renewcommand{\refname}{Additional methodological reference for V14}
\begin{thebibliography}{99}
\bibitem[40]{V14Karsten}
L.~H.~Karsten, \emph{The Lattice Fermion Problem and Weak Coupling
Perturbation Theory}, in W.~R\"uhl (ed.), \emph{Field Theoretical Methods
in Particle Physics}, NATO Advanced Study Institutes Series B,
vol.~55, Springer/Plenum (1980), pp.~235--247,
doi:10.1007/978-1-4684-3722-5\_10. Comparison with nonlocal-derivative
lattice prescriptions, not an imported graviton or ghost result.
\end{thebibliography}
\endgroup

\clearpage

\subsection*{Verification record}
The accompanying \srcid{check_ESM_periodic_UV_V14.py} constructs the
Cartan load matrices independently, recomputes the three signed face
traces in exact algebra, and verifies the rational interval enclosure
using an alternating-series certificate for $\pi$. It checks the
multiplication placements producing the two seagulls, the full
mixed-pair identity and the nonzero upper determinant coefficient.
Numerical periodic-reference floors and matrix determinant diagnostics
are explicitly separate from these exact checks. The new analytic
proofs are the written strip, symbol and source estimates; no proof
assistant formalization or interval quadrature of an unlocalized face
integral is claimed.


\clearpage
\part*{V15: Coordinate-compatible antifield transport and the first quantum diffeomorphism source}
\addcontentsline{toc}{part}{V15: Coordinate-compatible antifields and the first quantum source}

\section{The missing object is a transformation source, not another metric determinant}
\label{sec:v15-start}

Sections 1--111, including their qualifications, are preserved. V14 proves a
full-zone cutoff comparison for a two-metric coefficient of a \emph{changed}
periodic regulator. It does not determine a quantum transformation law. The
first missing coefficient already has one metric antifield and one retained
diffeomorphism ghost. Its evaluation requires the quadratic transformation
$R_2$, a mixed metric--ghost contraction, and a definite prescription for
transporting the cutoff through the nonlinear coframe chart. This installment
constructs that prescription and evaluates the coefficient. No additional
zero-antifield action or propagator is substituted for V14's completed family.

There are two positive results. First, the coordinate Jacobian is compatible
with an explicitly transported finite transformation source \emph{exactly},
not just to a mesh jet. Second, the complete one-loop coefficient with one
metric antifield and one ghost has a matched continuum limit and a rooted
$O(a^2)$ source error. Neither assertion says that the full finite action is
invariant under the new source. Its classical master defect is retained, so
an ordinary Wess--Zumino cocycle is not inferred from an unproved nilpotence.

The supplied ESM programme requires source completeness, a common measure,
and separate cohomological and norm estimates \cite{V6Programme}. The general
BV and measured-pushforward identities are inherited. Cutoff-dependent BRST
transformations and quantum master equations have established treatments
\cite{V15IIM}; the present construction concerns the actual V12 chart and
V14 propagators. It is not a first construction of perturbative gravity.

\subsection{Fixed data and one additional source prescription}

We use V13's holomorphic Euclidean coefficient reference, $\chi>0$, four odd
periodic grids, and the metric contour $T$. The internal matter fields are not
integrated here. Let $S=s(aD)$ be exactly V14's smooth filter
\eqref{eq:v14-profiles}. The derivative $D_\mu$ is the ordinary spectral
multiplier $ip_\mu$ on the finite grid. Every occurrence of this derivative
in the new metric-coordinate source acts on an $S$-filtered argument. Thus
$D_\mu S$ is smooth periodic; no derivative is continued through a momentum
face by an unspecified representative.

Take $\Lambda_0\ge32\mu>0$, $a\Lambda_0\le1/4096$, and
$\Lambda_0L_\nu\ge1$. The retained ghost and antifield have momenta
$|k|\le\mu$. Put
\[
 \nu_0(p)=1-\theta(p/\Lambda_0),\qquad
 C^g_a=\nu_0\widetilde K_a^{-1},\qquad
 C^c_a=\nu_0\widetilde M_a^{-1},
\]
with the same $\theta,\widetilde K_a,\widetilde M_a$ as V14. The zero
modes and the lower Wilsonian modes are not inverted. Source coefficients are
normalized per physical volume $V_L=\prod_\nu L_\nu$; write
$\int_p^L=V_L^{-1}\sum_{p\in\mathcal P_{a,L}}$.
Literal finite Fourier coefficients have allowed external transfers. Momentum
derivatives and Taylor jets below use the smooth interpolation specified by
the same routed symbols in each summand; no exact renormalized finite-torus
Slavnov identity is inferred from that interpolation.

Write $\mathsf E$ for the Euclidean affine ADM map of
\eqref{eq:v13-Wick-map}, and define
\begin{equation}
 \begin{split}
 G(q)&=(I+\mathsf E(q))^T(I+\mathsf E(q))
       =I+q+\mathsf Q(q,q),\\
 \mathsf B(q,v)&=\mathsf E(q)^T\mathsf E(v)+\mathsf E(v)^T\mathsf E(q),
 \qquad A(q)=DG(q)=I+\mathsf B_q.
 \end{split}
 \label{eq:v15-chart}
\end{equation}
Here $\mathsf B_qv=\mathsf B(q,v)$ is pointwise. The metric antifield is
$\eta=q^*$, of odd parity. Its sign convention is fixed by adding
$+\langle\eta,\mathscr R_a(q,c)\rangle$ to the action; coefficients below are
read with $\eta$ before the external ghost. This additional source prescription
is not an assertion of a completed finite BRST symmetry.
The full transformation is an analytic source germ. All Gaussian claims
below use its polynomial Taylor coefficients at the declared interaction
order; no inverse coordinate chart is integrated over an unbounded Gaussian
configuration space.

\section{Transport the filter through the actual coframe chart}
\label{sec:v15-coordinate-source}

For a symmetric metric record $g$, let $\widehat c=Sc$ and set
\begin{equation}
 \begin{split}
 \mathscr V_a(g,c)&=S\{R_0\widehat c+
                   \mathcal L_{\widehat c}S(g-I)\},\\
 (\mathcal L_c h)_{\alpha\beta}
   &=c^\rho D_\rho h_{\alpha\beta}
     +h_{\rho\beta}D_\alpha c^\rho
     +h_{\alpha\rho}D_\beta c^\rho,\\
 \boxed{\mathscr R_a(q,c)=A(q)^{-1}\mathscr V_a(G(q),c).}
 \end{split}
 \label{eq:v15-canonical-source}
\end{equation}
Products are the literal finite site products. The inverse is a pointwise
$10\times10$ coordinate inverse on the nondegenerate chart, not an inverse
Laplacian or a ghost propagator. It is important that $S$ is applied in the
\emph{metric} record before pulling the vector field back. Filtering $q$
inside the nonlinear map $G$ is a different operation.

\begin{theorem}[Exact coordinate-measure compatibility of the finite source]
\label{thm:v15-coordinate}
On a fixed sufficiently small real or complex Wiener ball in $q$,
\eqref{eq:v15-canonical-source} is analytic, with a radius independent of
mesh and volume. Its fixed-order coefficient kernels are quasilocal. Let
$J_G(q)=\prod_x\det DG(q_x)$, with the anchored branch from V12. Then
\begin{equation}
 \boxed{\operatorname{div}_{J_G(q)\,dq}\mathscr R_a(q,c)=0}
 \label{eq:v15-J-divergence}
\end{equation}
exactly on each finite grid. More generally, for an even nonvanishing density
$\rho_g(g)$ and its actual pullback $\rho_q(q)=\rho_g(G(q))J_G(q)$,
\begin{equation}
 \boxed{\operatorname{div}_{\rho_q}\mathscr R_a
     =\bigl(\mathscr V_a\log\rho_g\bigr)\circ G.}
 \label{eq:v15-weighted-divergence}
\end{equation}
No invariance of V14's action, Lie-algebra closure of $\mathscr V_a$, or
quantum master equation is assumed in this statement.
\end{theorem}
\begin{proof}
The map $\mathsf B_q$ has pointwise norm at most $C|q|$. The Wiener norm
controls the pointwise norm and is a multiplication algebra. Its Neumann
series therefore inverts $I+\mathsf B_q$ on one common ball. Every finite
coefficient is a finite composition of pointwise multilinear maps, $S$, and
$D S$. Their Fourier symbols are smooth periodic. At any prescribed source
order their physical derivative bound is $C_Na^{-1}$ times fixed filter
seminorms; integrating by parts gives rapid polynomial lattice-distance
bounds. These are fixed-order statements, not bounds uniform in $N$.

The field derivative of $\mathscr V_a$ is the affine, $g$-independent
operator $S\mathcal L_{Sc}S$. Its finite trace is zero. For the transport
term, cyclicity reduces its scalar trace to
$\Tr(M_{\widehat c^\rho}D_\rho S^2)$. The diagonal of the circulant
$D_\rho S^2$ vanishes: its Fourier trace is the paired sum of
$ip_\rho s(ap)^2$. Each tensor-index term is a fixed component matrix
multiplying $S M_{D_\alpha\widehat c^\rho}S$; its trace is the constant
diagonal of $S^2$ times $\sum_x D_\alpha\widehat c^\rho(x)=0$.
The inhomogeneous $S R_0Sc$ has no field derivative. Hence
$\operatorname{div}_{dg}\mathscr V_a=0$ on the full finite metric carrier.

The ordinary finite-dimensional change-of-variables formula for divergence
now applies to the actual diffeomorphism $G$. It gives
$\operatorname{div}_{J_Gdq}(A^{-1}\mathscr V_a\circ G)
=(\operatorname{div}_{dg}\mathscr V_a)\circ G=0$.
Multiplying by $\rho_g\circ G$ adds precisely the logarithmic derivative
in \eqref{eq:v15-weighted-divergence}. These formulas are algebraic and
extend to the declared analytic germ and contour. They also hold
coefficientwise when $c$ is odd. No continuum product rule was used in the
finite trace or coordinate argument.
\end{proof}

\subsection{Explicit source vertices, including the high--high contact}

Let $\mathscr R_a=\mathscr R_{0,a}+\mathscr R_{1,a}+\mathscr R_{2,a}
+O(q^3c)$, with the subscript counting metric degree. Direct multiplication
of the Neumann series gives
\begin{align}
 \mathscr R_{0,a}(c)&=S R_0Sc,\\
 \mathscr R_{1,a}(q,c)&=S\mathcal L_{Sc}(Sq)
                    -\mathsf B(q,\mathscr R_{0,a}c),\\
 \mathscr R_{2,a}(q,c)&=\tfrac12S\mathcal L_{Sc}(S\mathsf B(q,q))
          -\mathsf B(q,S\mathcal L_{Sc}(Sq))
          +\mathsf B(q,\mathsf B(q,\mathscr R_{0,a}c)).
 \label{eq:v15-source-jets}
\end{align}
In particular the first term of $\mathscr R_{2,a}$ does not filter its two
individual $q$ factors before multiplying them. Two high momenta whose sum
is low give a retained source contact. Deleting that contact does not define
the pullback in \eqref{eq:v15-canonical-source} and invalidates
\eqref{eq:v15-J-divergence}.

On a fixed nonaliased panel for which all required subset labels lie where
$s=1$, the nonlinear coefficients in \eqref{eq:v15-source-jets} are exactly
$R_1^E,R_2^E$ of V12--V13. The free coefficient is $R_0$, not
$\widehat R_{0,a}$. This is an \emph{explicit} additional source change:
\begin{equation}
 \|\partial_p^\alpha(\mathscr R_{0,a}-\widehat R_{0,a})(p)\|
 \le C_\alpha a^4\Lambda_*^{5-|\alpha|},\qquad |p|\le\Lambda_*.
 \label{eq:v15-free-source-change}
\end{equation}
It leaves the established classical mesh identities through $a^2$ unchanged,
but is not exact finite-cutoff equality with the old free source. V14's
zero-antifield metric and ghost actions, $\widetilde K_a$ and
$\widetilde M_a$, are not changed. In particular the new antifield map is
not retroactively declared to generate all their gauge-fixing terms.

The corresponding filter on tangent $q$ records is
\begin{equation}
 S_q=A(q)^{-1}SA(q),\qquad
 \boxed{DS_q(0)[h]=[S,\mathsf B_h].}
 \label{eq:v15-filter-transport}
\end{equation}
For an $h$-mark of momentum $k$, this commutator has kernel
$[s(a(p+k))-s(ap)]\mathsf B(h,\cdot)$. The fundamental theorem of calculus
therefore bounds its $r$ hard-momentum derivatives by
$C_r a^{r+1}|k|\|h\|$. After a fixed soft window $|k|\le\mu$, its
rooted weighted kernel has the corresponding $C_ba\mu$ bound. This is a
computed reference/filter-variation term. An unvaried $S$ in both coordinate
systems omits it; its small factor does not by itself pay the ultraviolet
mode count of a loop.

\section{The complete first antifield coefficient is a tadpole minus a mixed loop}
\label{sec:v15-current}

Use a Frobenius-orthonormal basis $E_A$ of $\operatorname{Sym}_4$, $A=1,\ldots,10$,
and the standard ghost basis $e_\alpha$, $\alpha=0,\ldots,3$. In this section
$h$ denotes a polarization, not a background to be integrated. Let
\begin{align}
 L_{a,A\alpha}(p,t)&=\mathscr R_{1,a}(E_Ae^{ipx},e_\alpha e^{itx})
                   \big|_{p+t},\\
 T_{a,AB}(p,t;c,k)&=D_q^2\mathscr R_{2,a}(0,ce^{ikx})
                       [E_Ae^{ipx},E_Be^{itx}]\big|_{p+t+k},\\
 V_{a,\beta B}(p,k;c)&=-s(a(p+k))s(ap)s(ak)
     \left[F_{p+k}^{\rm Four}R_1^E(E_B,p;c,k)\right]_\beta,
 \quad F_p^{\rm Four}=iF_p.
 \label{eq:v15-LTV}
\end{align}
The notation $|_r$ means its Fourier amplitude at total momentum $r$.
These are explicit finite tensors: $L,T$ are the three displayed
polynomials \eqref{eq:v15-source-jets}; $R_1^E=\mathcal L-\mathsf B R_0$
is the already evaluated chart polynomial. The $V$ vertex is from
V14's \emph{actual} ghost action, not a newly gauge-fixed version of the
antifield source. The upper ghost terms contain $Hc$ and vanish for the
external ghost used here, since $Sc=c$.

Let $\mathcal Z_a(k):\C^4\to\operatorname{Sym}_4(\C)$ be defined by the
coefficient
$\hbar\langle\eta,\mathcal Z_a c\rangle$ in the negative logarithm of
the finite normalized pushforward, at zero retained metric and through
interaction degree two. Thus $\hbar$ is \emph{not} included in $\mathcal Z_a$.

\begin{theorem}[Evaluated one-loop transformation source]
\label{thm:v15-current}
On the declared reference and source panel,
\begin{equation}
 \boxed{\begin{aligned}
 \mathcal Z_a(k)c={}&\frac12\int_p^L
           C^g_{a,AB}(p)T_{a,AB}(p,-p;c,k)\\
 &-\int_p^L C^g_{a,AB}(p)C^c_{a,\alpha\beta}(k-p)
       L_{a,A\alpha}(p,k-p)V_{a,\beta B}(-p,k;c).
 \end{aligned}}
 \label{eq:v15-current-formula}
\end{equation}
All repeated component labels are summed. The first term is the $R_2$
tadpole; the second is the complete mixed metric--ghost bubble. The formula
includes all ten metric components and the non-scalar finite ghost inverse.
It is independent of $\gamma_g,\beta_g,\gamma_c,\beta_c$ at this particular
source order. It does not describe the geometric two-jet, which does depend
on those parameters.
\end{theorem}
\begin{proof}
Let $y,\gamma,\bar\gamma$ be the hard metric and ghost variables. The only
one-loop vertices with these two external marks are
\[
 \tfrac12\eta_i T^i_{AB}(c)y^Ay^B,\qquad
 \eta_i L^i_{A\alpha}y^A\gamma^\alpha,\qquad
 \bar\gamma_\beta V^\beta_B(c)y^B.
\]
The first has Gaussian expectation $\hbar\eta_iT^i_{AB}C^g_{AB}/2$.
For the other two, complete the finite Grassmann square with action
$\bar\gamma M\gamma+\eta Ly\gamma+\bar\gamma V(c)y$.
Its additional effective metric term is
$-\eta Ly M^{-1}V(c)y$. One further metric expectation gives the minus
sign and both covariances in \eqref{eq:v15-current-formula}. This argument
fixes the Grassmann ordering without a convention-dependent diagram sign.
It works on the fixed complex metric contour by Gaussian algebra.

A source vertex linear in a hard metric has momentum equal to the sum of
the two external source momenta, of magnitude at most $2\mu<\Lambda_0$;
it cannot attach to an integrated covariance. Thus the possible
source-force/pure-gravity tadpole graphs vanish by support, not by a
stationarity assumption. The free antifield vertex cannot create a hard
ghost because its source momentum is low. A connected graph with additional
hard cubic or quartic vertices and these two external marks either contains
one of those forbidden hard-linear branches or has at least two loops.
These observations exhaust the one-loop possibilities.

An insertion of the old one-loop auxiliary amplitude could contribute at
this order only through a hard-linear branch of the transformation source;
that branch is absent. Its field-independent normalization also cancels.
The auxiliary density nevertheless remains in the measured Slavnov
\emph{divergence}, evaluated below. It has not been removed from the action.

In the bubble, $V$ forces both $p$ and $k-p$ into the $s$ chart. The completed
free matrices there equal the native ones because $\zeta=1$. The upper
ghost interaction has $Hc=0$, and no interacting metric vertex enters this
coefficient. The tadpole uses only the free matrices and the new source
jets. None of the four upper interaction parameters therefore occurs.
\end{proof}

\subsection{The ordinary zero-mesh tadpole is explicitly nonzero}

For the annular continuum covariance
$C^g_{0,\Lambda}(p)=2w_\Lambda(p)J/(\chi|p|^2)$ define
$I_\Lambda=\int_p w_\Lambda(p)/|p|^2$. The tadpole contribution in
\eqref{eq:v15-current-formula}, with the ordinary $R_2^E$, equals
\begin{equation}
 \boxed{\mathcal Z^{\rm tad}_{0,\Lambda}(k)c
                     =\frac{I_\Lambda}{\chi}\,\mathcal T(k,c),}
 \label{eq:v15-tadpole}
\end{equation}
where spatial indices range from 1 to 3 and
\begin{align}
 \mathcal T_{00}&=0,&\mathcal T_{0i}&=\frac{5i}{2}c^0k_i,\\
 \mathcal T_{ij}&=\frac{3i}{8}(k_i c^j+k_j c^i)
                 -\frac{i}{4}\delta_{ij}\sum_{\ell=1}^3k_\ell c^\ell.
 \label{eq:v15-tadpole-components}
\end{align}
The $i$ in these formulas is the imaginary unit. These are
coordinate-dependent local transformation-source coefficients, not a
physical graviton normalization or a beta function.

\begin{proof}
Use
$\sum_{AB}J_{AB}E_A\otimes E_B
 =\sum_AE_A\otimes E_A-\tfrac12I\otimes I$.
Polarization of $R_2^E=\mathcal L_c\mathsf Q-\mathsf B(q,R_1^E)$ gives
\[
 T_0(h,p;g,t;c,k)=\mathcal L_c\mathsf B(h,g)
 -\mathsf B(h,R_1^E(g,t;c,k))-\mathsf B(g,R_1^E(h,p;c,k)).
\]
For $h=g$ and $t=-p$, the terms linear in $p$ cancel. Insert the four
diagonal basis matrices and the six off-diagonal ones divided by $\sqrt2$.
Their sum, minus half the $I,I$ value, gives precisely
\eqref{eq:v15-tadpole-components}. The covariance factor $2/\chi$ cancels
the tadpole factor $1/2$, proving \eqref{eq:v15-tadpole}. This is a finite
component computation; no angular averaging or Ward identity is used.
\end{proof}

The mixed loop is not an optional replacement of that tadpole. For example,
at $\chi=1$, continuum unweighted propagators, $p=(2,3,1,4)$,
$k=(0,1,0,0)$ and $c=e_0$, its integrand \emph{before} the overall minus
sign has $01$ entry $-29i/600$. This exact rational test is obtained by
substituting the same $\mathsf E,\mathsf B,F$ into \eqref{eq:v15-LTV}; it
is not a value of the integrated loop. The accompanying verifier checks it
independently of the cutoff estimate.

\section{A full-zone cutoff theorem for the first antifield source}
\label{sec:v15-limit}

The coefficient \eqref{eq:v15-current-formula} is odd in $k$. Indeed $S$,
$\nu_0$, $\widetilde K_a^{-1}$ and $\widetilde M_a^{-1}$ are even under
simultaneous momentum reversal. The tensors $L,T$ have one derivative and
$V$ has two. Sending $(p,k)$ to $(-p,-k)$ therefore changes the sign of
both integrands. The paired finite grid proves
$\mathcal Z_a(-k)=-\mathcal Z_a(k)$, including the coordinate-transport
contacts. This parity assertion concerns this source coefficient; it does
not assert an evenness of V14's general geometric tensor.

Let $\mathsf T_3$ be the Taylor polynomial at zero in the external momentum
through total degree three, and set
\begin{equation}
 \mathcal Z^R_a=(I-\mathsf T_3)\mathcal Z_a.
 \label{eq:v15-ZR}
\end{equation}
The linear and cubic Taylor terms are stored as actual local antifield
counterterms $\hbar\langle\eta,(\mathsf T_3\mathcal Z_a)c\rangle$ with the
chosen subtraction sign. A redundant component bank has dimension
$10\cdot4\{4+20\}=960$. It is a source-renormalization bank, not a count of
physical gravitational Wilson coefficients.

\begin{lemma}[Bounds for the explicit source integrands]
\label{lem:v15-symbol}
Remove from the tadpole the part which is exactly linear in $k$ and comes
from the first and last terms of $\mathscr R_{2,a}$ in
\eqref{eq:v15-source-jets}. This removal means assignment to the linear
local bank. Denote the remaining tadpole-minus-bubble integrand by
$f_a(p,k)c$, and its unfiltered zero-mesh version by $f_0$. For all fixed
multi-indices $\alpha,\beta$, with $R_p=\Lambda_0+|p|$,
\begin{align}
 \|\partial_p^\beta\partial_k^\alpha f_a\|
 +\|\partial_p^\beta\partial_k^\alpha f_0\|
 &\le C_{\alpha,\beta}\chi^{-1}R_p^{-1-|\alpha|-|\beta|},\\
 \|\partial_p^\beta\partial_k^\alpha(f_a-f_0)\|
 &\le C_{\alpha,\beta}\chi^{-1}a^4R_p^{3-|\alpha|-|\beta|}
 \quad(p\in\mathcal B_a).
 \label{eq:v15-symbol-bounds}
\end{align}
All source tensor norms are ordinary fixed finite component norms.
\end{lemma}
\begin{proof}
For the first tadpole term, the two hard momenta add to zero, so $S$ on their
product equals one and the Lie derivative supplies only $k$. The last term
has only $R_0c$, also linear in $k$. Their full covariance trace, including
all upper modes, is therefore a genuine local linear polynomial. It is not
discarded before being assigned to $\mathsf T_3\mathcal Z_a$.

Every remaining term has filtered hard arguments. For the middle tadpole
term these give $s(ap)s(a(p\pm k))$; the bubble's $V$ gives the same hard
support. All relevant momenta are a fixed distance from the faces. On that
support the source tensors have bounds $C R_p$ for $L,T$ and $C R_p^2$ for
$V$, and each derivative lowers the degree by one. Both propagators have
order $-2$, one with factor $\chi^{-1}$. The tadpole has one propagator;
the bubble has two. Both resulting integrands have order $-1$.
Derivatives of $\nu_0$ are supported at scale $\Lambda_0$ and obey the
same degree bounds there.

If $|ap|\le1/128$ and $a\mu$ is as stated, every needed filter equals one.
The source vertices are exactly the ordinary ones. Only the free inverses
differ, and \eqref{eq:v13-inverse-symbols} gives relative mismatch
$C(aR_p)^4$, with all fixed differentiated versions. On the remaining
support $aR_p$ is bounded below by a fixed positive constant. The separate
order $-1$ estimates therefore imply the weaker mismatch bound in
\eqref{eq:v15-symbol-bounds}. The same argument covers the portion where
$f_a=0$ but $f_0\ne0$. Smooth periodic extension of the filtered factors
justifies the estimates at seams; the unassigned local tadpole is not used
in a differentiated nonlocal estimate.
\end{proof}

Define $f_a^-(p,k)=(f_a(p,k)-f_a(p,-k))/2$. Its integral equals that of
$f_a$. Oddness implies $\mathsf T_4f_a^-=\mathsf T_3f_a^-$. Integral
Taylor remainders and \cref{lem:v15-symbol} give, for $r=|\alpha|\le5$,
\begin{align}
 \|\partial_k^\alpha(I-\mathsf T_3)f_0^-\|
 &\le C_\alpha\chi^{-1}|k|^{5-r}R_p^{-6},\\
 \|\partial_k^\alpha(I-\mathsf T_3)(f_a^--f_0^-)\|
 &\le C_\alpha\chi^{-1}a^4|k|^{5-r}R_p^{-2}.
 \label{eq:v15-five-derivatives}
\end{align}
Thus the continuum periodic and thermodynamic coefficients are defined by
absolutely convergent subtracted sums and integrals, not by an unsubtracted
massless determinant.

\begin{theorem}[First quantum transformation source: full-cutoff convergence]
\label{thm:v15-limit}
On the stated fixed-lower-scale domain, for $r=|\alpha|\le5$,
\begin{equation}
 \boxed{\|\partial_k^\alpha(\mathcal Z^R_{a,L}-\mathcal Z^R_{0,L})(k)\|
 \le C_\alpha\chi^{-1}a^2|k|^{5-r}.}
 \label{eq:v15-main-rate}
\end{equation}
Moreover
$\|\partial_k^\alpha\mathcal Z^R_{0,L}(k)\|
\le C_\alpha\chi^{-1}|k|^{5-r}\Lambda_0^{-2}$.
For each fixed $r>5$, the cutoff error is bounded by
$C_r\chi^{-1}a^2\Lambda_0^{5-r}$. The same conclusions hold in infinite
volume. Constants are independent of the mesh and periods when
$\Lambda_0L_\nu\ge1$, but not of unbounded source or derivative order.
The actual one-loop action error has the additional factor $\hbar$.
\end{theorem}
\begin{proof}
Subtract the absolutely convergent continuum sum from the finite one.
By \eqref{eq:v15-five-derivatives}, the mismatch in the cube and the
omitted continuum tail are bounded, respectively, by
\[
 C_\alpha\chi^{-1}|k|^{5-r}\,
 \frac{a^4}{V_L}\sum_{p\in\mathcal P_{a,L}}R_p^{-2},
 \qquad
 C_\alpha\chi^{-1}|k|^{5-r}\,
 \frac1{V_L}\sum_{p\notin\mathcal B_a}R_p^{-6}.
\]
For $R\ge\Lambda_0$, the number of periodic momenta of size at most $R$
divided by $V_L$ is bounded by $C R^4$. The first dyadic sum is at most
$C a^4 a^{-2}=C a^2$; the second starts at order $a^{-1}$ and is at most
$C a^2$. Summing the continuum estimate from $\Lambda_0$ gives the
$\Lambda_0^{-2}$ bound. There is no logarithm in this nonlocal comparison.

For $r>5$, the subtracted polynomial is absent. The mismatch annular bound
is $a^4R^{7-r}$ and the continuum one is $R^{3-r}$. For $r=6$ their
ultraviolet sums cost $O(a^3)$, for $r=7$ the mismatch is
$O(a^4[1+\log(1/(a\Lambda_0))])$, and for $r\ge8$ it is bounded by
$C_ra^4\Lambda_0^{7-r}$. Each is at most the displayed weaker
$C_ra^2\Lambda_0^{5-r}$ when $a\Lambda_0$ is small. This proves the
fixed higher-derivative statement. Replacing counting by integration
proves the thermodynamic version.
\end{proof}

The raw local budgets are
\begin{equation}
 \|D_k\mathcal Z_a(0)\|\le C\chi^{-1}a^{-2},\qquad
 \|D_k^3\mathcal Z_a(0)\|\le
 C\chi^{-1}[1+\log(1/(a\Lambda_0))].
 \label{eq:v15-local-budgets}
\end{equation}
Even jets vanish. The local high--high tadpole contributes to the first
bound. These terms are retained source normalization; smallness of the
nonlocal error does not make the local source correction zero.

\section{Rooted source topology, shell allocation, and matched source choices}
\label{sec:v15-shells}

Multiply the coefficient difference by a fixed smooth external window
$\varpi(k/\mu)$ supported in $|k|<\mu$. The fixed derivative estimates of
\cref{thm:v15-limit}, with more than $b+4$ derivatives, give
\begin{equation}
 \boxed{a^4\sum_x(1+\mu d_{\rm per}(x,0))^b
 \|\mathcal K^\Delta_{a,L}(x)\|
 \le C_b\chi^{-1}a^2\mu^5.}
 \label{eq:v15-kernel}
\end{equation}
Indeed the scaled derivative bounds are
$C_r\chi^{-1}a^2\mu^{5-r}$; Fourier integration by parts, then
periodization, proves the result. It implies the corresponding bilinear
source estimate with one coefficient function in $L^1$ and the other in
$L^\infty$, with no total-volume factor. Odd sources are treated
coefficientwise in their finite exterior algebra.

The volume comparison is separate. The same $p$-derivative estimates and
Poisson summation show, for every fixed $b>4$ and $r\le5$,
\begin{equation}
 \|\partial_k^\alpha(\mathcal Z^R_{0,L}-\mathcal Z^R_{0,\infty})(k)\|
 \le C_{b,\alpha}\chi^{-1}|k|^{5-r}\Lambda_0^{-2}
                        (\Lambda_0L_{\min})^{-b}.
 \label{eq:v15-volume}
\end{equation}
There is no infrared removal in either limit.

For a cumulative covariance tower, the bubble in
\eqref{eq:v15-current-formula} contains every pair
$C_i^g C_j^c$. Its increment at level $n$ contains
\[
 C_n^g C_{<n}^c+C_{<n}^g C_n^c+C_n^g C_n^c,
\]
with their actual noncommuting tensors between the factors. The tadpole
has its single $C_n^g$. This is the inherited mixed-shell allocation
applied to a new source, not a diagonal-shell approximation. An interior
annulus of size $\Lambda\ge\Lambda_0$ obeys
\begin{equation}
 \|\partial_k^\alpha\mathcal Z^R_{a,\Lambda}\|
 \le C_\alpha\chi^{-1}|k|^{5-r}\Lambda^{-2},\qquad
 \|\partial_k^\alpha(\mathcal Z^R_{a,\Lambda}-\mathcal Z^R_{0,\Lambda})\|
 \le C_\alpha\chi^{-1}a^4|k|^{5-r}\Lambda^{2}.
 \label{eq:v15-shell}
\end{equation}
These follow from the same five-derivative argument with one routed
annular weight. Summing to a fixed fraction of $a^{-1}$ costs $O(a^2)$;
the remaining smooth source cap and the continuum tail cost $O(a^2)$ as
well. Local linear and cubic source increments are generated by the actual
Taylor jets of these diagrams.

Two profile choices obeying the declared fixed-radius and derivative
bounds, with their \emph{own} coordinate-compatible source
\eqref{eq:v15-canonical-source}, have the same continuum $\mathcal Z^R_0$.
Matching their actual local antifield jets therefore gives a nonlocal
difference bounded by the sum of \eqref{eq:v15-main-rate}. This is
independence within the specified source-completion class. A different
untransported high-momentum antifield prescription need not have the same
local Jacobian identity, and no equivalence with an unmodified cyclic
microscopic regulator is inferred.

\section{An evaluated measure contact and the curved Slavnov consistency equation}
\label{sec:v15-measure}

Let $\rho_g$ be V10's auxiliary leading density relative to $dg$, with its
link/Haar factors and branch conventions, and let
$\rho_q=(\rho_g\circ G)J_G=e^{-\mathcal M^E_{a,1}}$ be V12's actual density
relative to $dq$. Equation \eqref{eq:v15-weighted-divergence} shows precisely
what the coordinate Jacobian cancels and what it does not cancel. In
particular it cancels the divergence caused purely by changing $g$ to $q$;
it does not make the physical weight $\rho_g$ invariant under the new source.

\begin{proposition}[A nonzero complete auxiliary-density contact on the lapse line]
\label{prop:v15-measure-contact}
Take $q_{00}(t)=\varepsilon e^{-ikt}$, all other $q$ entries zero, and
$c^0(t)=\vartheta e^{ikt}$, all other ghost entries zero, with $|k|\le\mu$.
Here $\vartheta$ is a single odd coefficient. In the holomorphic reference,
the coefficient of $\varepsilon\vartheta$ per physical volume is
\begin{equation}
 \boxed{\left[\operatorname{div}_{\rho_q}\mathscr R_a\right]_{q_{00}c^0}
                   =\frac{7i}{2a^4}k.}
 \label{eq:v15-measure-contact}
\end{equation}
This is a local measured-divergence operand. It is not the full quantum
Slavnov defect and is not a demonstration of a diffeomorphism anomaly.
\end{proposition}
\begin{proof}
On this line $e=\operatorname{diag}(n,1,1,1)$ with $n=1+q_{00}/2$.
The native cofactor load vanishes even for time-dependent $n$: the temporal
load differentiates the spatial cofactor, which is constant, and the other
load terms are spatial derivatives. Hence $A_*=0$, $K_a$ of the auxiliary
closed-walk formula is zero, and every Haar logarithm is zero. The metric
weight is therefore exactly $\rho_g=n^{-7}$ on this line. Its actual
coordinate factor is $J_G=n$, so $\rho_q=n^{-6}$, consistently with V12.

The low panel is unfiltered. Its zeroth transformation is
$R_0q_{00}=2\partial_0c^0$ and its first one is
$R_1q_{00}=c^0\partial_0q_{00}+q_{00}\partial_0c^0$.
The latter is zero for the opposite plane waves above. Expanding
$\log n=q_{00}/2-q_{00}^2/8+\cdots$ gives
$[\mathscr R_a\log n]_{q_{00}c^0}=-ik/2$.
Equation \eqref{eq:v15-weighted-divergence} multiplies this by $-7$,
not by $-6$: the divergence of the transformed field itself supplies
the coordinate-Jacobian part. The physical site count per volume is
$a^{-4}$, proving \eqref{eq:v15-measure-contact}. Differentiations remain
within the same lapse line, so the zero link/Haar terms have no hidden
variation in this coefficient.
\end{proof}

One must also keep track of where the density is represented. On the finite
odd cotangent algebra use the convention
$\Delta_\rho F=\Delta_0F+(\log\rho,F)$ for a field density $\rho$.
Equivalently the odd cotangent volume has base density $\rho^2$.
For an even classical master candidate $\mathcal S$ and a one-loop
propagating coefficient $\Gamma_1^{\rm prop}$,
\begin{equation}
 \boxed{(\mathcal S,\Gamma_1^{\rm prop})-\Delta_\rho\mathcal S
  =(\mathcal S,\Gamma_1^{\rm prop}+\mathcal M)-\Delta_0\mathcal S,
  \qquad \rho=e^{-\mathcal M}.}
 \label{eq:v15-measure-once}
\end{equation}
This follows immediately from graded antisymmetry for the two even
functionals. Including $\mathcal M$ in the action \emph{and} using
$\Delta_\rho$ without compensating it would double-count the same contact.

\subsection{Consistency before, not after, a master equation is proved}

V14 has not supplied an exact high-frequency classical master solution.
The antifield source above does not do so either. On any fixed finite BV
extension of these candidates, put
\begin{equation}
 B_0=\tfrac12(\mathcal S,\mathcal S),\quad
 d_{\mathcal S}=(\mathcal S,\cdot),\quad
 A_1=(\mathcal S,\Gamma_1)-\Delta_\rho\mathcal S.
 \label{eq:v15-curved-data}
\end{equation}
The actual identities are
\begin{equation}
 \boxed{d_{\mathcal S}^{\,2}F=(B_0,F),\qquad
 d_{\mathcal S}A_1=(B_0,\Gamma_1)+\Delta_\rho B_0.}
 \label{eq:v15-curved-consistency}
\end{equation}
For completeness, the first is graded Jacobi. The BV derivation rule gives
$\Delta_\rho B_0=-(\mathcal S,\Delta_\rho\mathcal S)$; applying
$d_{\mathcal S}$ to $A_1$ proves the second. These are general algebraic
identities, included as an audit of their application, not as newly
invented BV theorems. They show why a bounded inverse on a putative
$cq$ anomaly space cannot be invoked before the tree-breaking and its
insertions have been retained. If $B_0=0$ on an actual filtered complex,
the usual closedness follows; V12's local identity modulo $a^3$ is not
that all-frequency premise.

The present source map is affine in $g$ before coordinate pullback, but the
fixed filters do not make its commutator the Lie bracket of vector fields.
This is exactly the remaining classical breaking, not a missing
coordinate Jacobian. A ghost-antifield extension and reference variation
must be specified together before $A_1$ is called a calculated anomaly.

\section{The remaining first Slavnov row is now a finite measured problem}
\label{sec:v15-next}

An exact finite diagnostic is available without assuming symmetry. Let $x$
be retained coordinates, $y$ the integrated innovation coordinates,
$\mathscr R=(\mathscr R_x,\mathscr R_y)$ an actual finite vector field,
and $S_{\rm all}$ the action including the \emph{actual} Gaussian/Berezin
reference. Then integration by parts gives
\begin{equation}
 \boxed{\operatorname{Div}_x
   \left(Z\,\langle\mathscr R_x\rangle\right)
 =-\hbar^{-1}Z\left\langle
       \mathscr R S_{\rm all}
               -\hbar\operatorname{div}_\rho\mathscr R
                         \right\rangle.}
 \label{eq:v15-measured-breaking}
\end{equation}
The super-divergence signs on odd retained variables are understood in the
same fixed ordering. To prove it, differentiate
$\rho\exp(-S_{\rm all}/\hbar)\mathscr R$ in all variables, and integrate
the $y$ divergence. Its polynomial-Gaussian coefficients vanish at the
bosonic contour ends; Berezin integration has the same exact integration
by parts. The remaining $x$ divergence is the left side. The proof works
coefficientwise without convergence of the full interacting exponential.
Normalized expectations on the right use the same $Z$ as on the left.

In particular $S_{\rm all}$ is not just V14's interaction polynomial.
For a supported covariance $C$, its hard reference includes
$y^TC^{-1}y/2$ and $\bar\gamma(C^c)^{-1}\gamma$. Their variations, including
changes of the retained/innovation maps, belong to the insertion in
\eqref{eq:v15-measured-breaking}. The commutator
\eqref{eq:v15-filter-transport} is one explicit term of that transport.
It cannot be set to zero because the unvaried determinant has a continuum
limit. Equation \eqref{eq:v15-measured-breaking} is the ordinary inherited
measured-divergence principle applied without dropping the breaking; it
is not a claim that its right side has already been shown to vanish.

For orientation, after reintroducing the retained algebraic nonminimal
field, the zero-nonminimal, one-ghost/one-metric part of the
\emph{action--antifield} pairing has the form
\begin{equation}
 \langle \Pi_1 q,R_0c\rangle+
 \langle H_{\rm cl}q,\mathcal Z_1c\rangle+
 \langle\tau_1,R_1(q,c)\rangle.
 \label{eq:v15-first-row}
\end{equation}
Here $\Pi_1$ and $\tau_1$ are the loop two-point coefficient and tadpole,
$H_{\rm cl}$ is the ungauge-fixed quadratic action on that retained slice,
and $\mathcal Z_1$ is the source just calculated (with its loop factor
assigned consistently). The retained nonminimal source contact, the
measured divergence, and the reference/breaking insertion in
\eqref{eq:v15-measured-breaking} must accompany this expression. It is not
legitimate to test only $\Pi_1R_0$ for transversality and discard the
second and third terms. Nor is it legitimate to replace $H_{\rm cl}$ by
the gauge-fixed Gaussian without restoring the nonminimal terms.

Thus one exact coordinate-measure interface and one full-cutoff quantum
source operand have been closed. The next calculation is the \emph{signed
sum} in the first measured Slavnov row, with a common local assignment for
$\Pi_1,\mathcal Z_1,\tau_1$, the auxiliary contact, and the reference
insertion. The present $960$-component source projector and V14's metric
projector have not been proved to intertwine that combined row; their
separate bounds are not that theorem.

There is a useful order audit for subsequent work. At one loop, the
$cq^m$ action--antifield identity consumes a geometric coefficient with
$m+1$ backgrounds, hence action vertices through degree $m+3$, and an
antifield coefficient with at most $m-1$ backgrounds, whose tadpole uses
$R_{m+1}$. Thus the row $cq$ uses $S_4$ and $R_2$ and is within the
present native bank. The next row $cq^2$ needs $S_5$ and $R_3$. Merely
renaming the two-source result a complete Slavnov functional would
conceal that additional work.

\begingroup
\renewcommand{\refname}{Additional methodological reference for V15}
\begin{thebibliography}{99}
\bibitem[41]{V15IIM}
Y.~Igarashi, K.~Itoh and T.~R.~Morris,
\emph{BRST in the Exact RG}, Progress of Theoretical and Experimental
Physics (2019), 103B01, arXiv:1904.08231,
doi:10.1093/ptep/ptz099. Comparison framework for cutoff-dependent BRST
and modified Slavnov identities; not an imported estimate for the
Renewal regulator or its coframe measure.
\end{thebibliography}
\endgroup

\section{Final ledgers for V15}
\label{sec:v15-ledgers}

\subsection*{Proved}
\addcontentsline{toc}{subsection}{V15: Proved}
The metric-coordinate filtered source and its exact pullback
\eqref{eq:v15-canonical-source} obey the coordinate-Jacobian identity
\eqref{eq:v15-J-divergence} on each finite grid. Its explicit $R_1,R_2$
vertices retain the high--high source contact, and its filter variation is
\eqref{eq:v15-filter-transport}. The first quantum antifield coefficient
is the complete tadpole-minus-mixed-loop expression
\eqref{eq:v15-current-formula}. Its native continuum tadpole is evaluated
in \eqref{eq:v15-tadpole-components}; its matched full-cutoff error is
\eqref{eq:v15-main-rate}, with rooted source, volume and shell bounds.
The actual auxiliary measure gives the local lapse--ghost contact
\eqref{eq:v15-measure-contact}. These results supply a transformation-source
module, not a full quantum symmetry theorem.

\subsection*{Inherited}
\addcontentsline{toc}{subsection}{V15: Inherited}
V1--V14 and their labels are unchanged. V12 provides the chart, its
Jacobian and the ordinary nonlinear ghost polynomials. V13 provides the
holomorphic reference and native inverse estimates. V14 provides the
periodic completed action, full-zone covariance bounds, geometric two-jet
and matched regulator class. Gaussian/Berezin integration, mixed-shell
allocation, coordinate divergence, and BV algebra identities are inherited
methods; they are not renamed new abstract continuum theorems.

\subsection*{Literature input}
\addcontentsline{toc}{subsection}{V15: Literature input}
Cutoff-dependent BRST/quantum-master frameworks \cite{V15IIM} and
perturbative effective gravity \cite{V13BFR} are established. The new
calculations evaluate the source lift, finite contacts and cutoff estimates
for the declared native chart and reference. No priority claim is made for
BV calculus, Ward identities or perturbative gravity as such.

\subsection*{Conditional}
\addcontentsline{toc}{subsection}{V15: Conditional}
A common renormalized first Slavnov row still requires the actual reference
variation and breaking insertion, and a single compatible local assignment
for its metric, antifield and measure components. The source completion
changes the high antifield prescription and the low free source at $O(a^4)$;
it does not prove exact invariance of the unchanged V14 action. Any
microscopic matching remains separate from matching within the specified
smooth source family. The odd cotangent measure convention must be fixed
when \eqref{eq:v15-measure-once} is used.

\subsection*{Open}
\addcontentsline{toc}{subsection}{V15: Open}
The smallest remaining quantum target is the complete $cq$ measured
Slavnov row, not another isolated geometric or antifield determinant.
Its source and coordinate-measure operands are now explicit; its reference
and tree-breaking insertion must be evaluated in the same signed sum.
A bounded common restoring projector, the larger antifield bank, chiral
and mixed ESM sectors, all-loop filtered self-mapping, infrared removal and
Lorentzian comparison remain open. No positive nonperturbative gravitational
measure, asymptotic safety or finite-parameter ultraviolet completion is
asserted.

\begin{center}\small
\begin{tabular}{>{\raggedright\arraybackslash}p{.24\textwidth}>{\raggedright\arraybackslash}p{.67\textwidth}}
\toprule
Variable & Scope in V15\\
\midrule
Volume & Source and cutoff bounds uniform for $\Lambda_0L_\nu\ge1$.
Raw local contacts can be extensive; their per-volume budgets are explicit.\\
Mesh & $a\Lambda_0\le1/4096$, fixed profile radii; all metric modes
are integrated according to V14. High--high source contacts are retained.\\
Loop/source degree & One propagating loop, one metric antifield and one
ghost, at zero retained metric. The coordinate-density identity is finite
and exact on its chart.\\
Derivative order & Every prescribed finite momentum derivative and kernel
moment; constants may depend on that order.\\
Lower scale/coupling & $\Lambda_0\ge32\mu>0$, $\chi>0$ fixed. The
source error displays $\chi^{-1}$; no infrared or $\chi\downarrow0$ limit.\\
Shell count & Nonlocal increments of this coefficient are summable.
No invariant neighborhood of arbitrary interacting ESM actions is proved.\\
Symmetry & Exact coordinate-Jacobian compatibility, not exact classical
or quantum BV closure of the completed action.\\
Measure & The V10--V12 auxiliary density and the fixed metric contour;
measure counted once. No signature-continuation theorem.\\
\bottomrule
\end{tabular}
\end{center}

\subsection*{Verification record}
The standalone \srcid{check_ESM_antifield_source_V15.py} checks the
coordinate-filter divergence and source-jet expansion in an exact finite
coordinate model, the full four-dimensional low tadpole, the Grassmann sign of the mixed loop, and a rational
bubble operand. It separately checks numerical parity and symbol scaling.
The cutoff and rooted estimates are proved above, not inferred from those
finite diagnostics. No new proof-assistant formalization is asserted.



\clearpage
\part*{V16: The first measured Slavnov row and a common local assignment}
\addcontentsline{toc}{part}{V16: First measured Slavnov row and common local assignment}

\section{From separate source limits to an assembled identity}
\label{sec:v16-start}

Sections 1--119 are preserved. V15 provides an evaluated antifield
coefficient, not a Slavnov identity. A fixed lower Wilsonian scale is
still present. The right hand side of the identity at that scale need
not vanish: it contains the variation of the actual Gaussian reference.
Treating this term as an ultraviolet anomaly, or setting it to zero
before the lower scale is removed, would change the problem.

This installment uses the \emph{same} V14--V15 completed branch. The
nonminimal field is restored before extracting the one-ghost/one-metric
row. A finite Hessian formula evaluates the reference insertion and the
insertion of the actual classical breaking. The coordinate and
ultralocal measure contacts are calculated in all ten metric directions.
Taylor orders four, three and five then give a common bounded local
assignment for the geometric, transformation and breaking coefficients.
The renormalized row converges to a normalized modified continuum
identity, with its lower-scale reference trace retained.

The general Legendre and BV change-of-variables principles are standard
and already used in the lineage; see also \cite{V15IIM}. Their
rederivation below fixes the signs and the specific operands. The new
application is the assembled row of the completed native action,
including the off-shell nonminimal cancellation, its 40-component
measure contact, finite local-assignment mismatch and cutoff implication.
This is not a splitting theorem for the full nonlinear ESM BV complex.
A result for one row does not imply identities with two retained metric
fields, which consume further native vertices.

\subsection{Data, normalization, and distinct conclusions}

Use V13's holomorphic Euclidean coefficient convention and conformal
contour, V14's fixed smooth profiles and bounded upper parameters, and
V15's source \eqref{eq:v15-canonical-source}. The Palatini branch has
$\chi>0$, zero independent Holst coefficient and zero cosmological term.
No matter field is integrated here. Put
\[
 \Theta(p)=\theta(p/\Lambda_0),\qquad \nu(p)=1-\Theta(p),\qquad
 \Lambda_0\ge32\mu>0,\qquad a\Lambda_0\le 1/4096.
\]
Evaluate the row at $h e^{-ikx}$ and the contravariant ghost $c e^{ikx}$,
$|k|\le\mu$. Coefficients are per physical volume, with $\hbar$
factored out of the one-loop term. Thus $\Pi$ below is the coefficient
of $\hbar$ in the metric Hessian; V14's $\Pi$ included that factor.
The transformation coefficient $\mathcal Z$ has V15's normalization.

Finite matrix identities hold on allowed grid momenta. Smooth Taylor
identities and the exact renormalized modified identity are asserted
in the thermodynamic continuous-momentum convention. Cutoff and rooted
estimates use fixed source windows; finite-volume comparisons are
separate. An arbitrary off-grid interpolation on one fixed torus is not
declared to preserve its on-grid Ward identities.

\section{Restore the nonminimal field before testing the row}
\label{sec:v16-offshell}

Let $F$ be V13's harmonic functional and let
$\widetilde K_a,\widetilde M_a$ be V14's completed free matrices. Set
$H_a^{\rm cl}=\widetilde K_a-\chi F^*F$. This off-shell assignment
reproduces \emph{exactly} the existing gauge-fixed matrix after the
nonminimal Gaussian; it does not assert full-zone ungauge-fixed
symmetry. Through the needed degree define
\begin{equation}
 \begin{split}
 \mathscr S_a(q,b,c,\bar c)={}&
 \tfrac12\langle q,H_a^{\rm cl}q\rangle
 +\tfrac1{2\chi}\langle b,b\rangle+i\langle b,Fq\rangle
 +\langle\bar c,\widetilde M_a c\rangle\\
 &+\widetilde S^E_{3,a}(q)+\widetilde S^E_{4,a}(q)
 +\widetilde S^{\rm gh}_{3,a}(q,c,\bar c)
 +\widetilde S^{\rm gh}_{4,a}(q,c,\bar c).
 \label{eq:v16-offshell-action}
 \end{split}
\end{equation}
The last four terms are literally
\eqref{eq:v14-completed-interactions} and \eqref{eq:v14-upper-ghost},
with the native ghost terms of V13. Retain the auxiliary-connection
amplitude as the order-$\hbar$ action insertion
$\mathcal M_a(q)=\mathcal M^E_{a,1}(q)$, using
\eqref{eq:v12-measure-q} after the declared holomorphic substitution.
We use the \emph{flat} coefficient divergence; the density is not
also put into that operator.

Complete the transformation source by
\begin{equation}
 \begin{aligned}
 r_a^q&=\mathscr R_a(q,c),&
 r_a^c&=S\{(Sc)^\rho D_\rho(Sc)\},\\
 r_a^{\bar c}&=ib,& r_a^b&=0.
 \label{eq:v16-full-source}
 \end{aligned}
\end{equation}
This fixes an additional ghost-source prescription with the same $S$
as V15. It does not change the zero-antifield action, and nilpotence of
this finite filtered source is not assumed. Its zero-mesh coefficients
are the ordinary diffeomorphism BRST transformations pulled back by the
actual chart. Its $q$ component is exactly V15's source, and the added
components do not change the coefficient $\mathcal Z_a$ calculated there.

\begin{proposition}[Off-shell representation of the existing measured block]
\label{prop:v16-offshell}
Integrating $b$, with the source $\beta$ of
\eqref{eq:v13-b-contact}, reproduces the V14 metric and ghost integral
and its nonminimal contact without an additional determinant. The even
free inverse on nonzero modes, in $(q,b)$ coordinates, is
\begin{equation}
 \begin{pmatrix}
 \widetilde K_a^{-1}&-i\chi\widetilde K_a^{-1}F^*\\
 -i\chi F\widetilde K_a^{-1}&
          \chi I-\chi^2F\widetilde K_a^{-1}F^*
 \end{pmatrix}.
 \label{eq:v16-even-inverse}
\end{equation}
The hard covariance is $\nu$ times this matrix, with ghost covariance
$\nu\widetilde M_a^{-1}$. At $b=\bar c=0$ the retained first
Slavnov row uses $H_a^{\rm cl}$, not $\widetilde K_a$.
\end{proposition}
\begin{proof}
Complete the square:
\[
 \frac{b^2}{2\chi}+ibFq
 =\frac{(b+i\chi Fq)^2}{2\chi}+\frac\chi2(Fq)^2.
\]
The inherited Gaussian calculation proves the action and source
statement; block inversion proves \eqref{eq:v16-even-inverse}.
The antighost transformation is $ib$, so its retained contribution
vanishes at $b=0$ \emph{before} that field is integrated. The classical
metric derivative on that slice is $H_a^{\rm cl}q$. Replacing it by
the gauge-fixed matrix would add $\chi F^*Fq$ without its nonminimal
partner. The triangular square completion has constant unit Jacobian
and is the contour convention already used in V13.
\end{proof}

For algebraic purposes use a regulated Legendre representation of these
same coefficients. Write $\Phi=(q,b,c,\bar c)$ and let
$\mathbb H_{a,0}$ be the free graded Hessian. Replace zeros of $\nu$
by $\nu_\epsilon=\nu+\epsilon\Theta$. On its invertible blocks put
\begin{equation}
 \mathbb C_{a,\epsilon}=\nu_\epsilon\mathbb H_{a,0}^{-1},\qquad
 \mathbb R_{a,\epsilon}=(\nu_\epsilon^{-1}-1)\mathbb H_{a,0}.
 \label{eq:v16-IR-regulator}
\end{equation}
At a physical zero mode one may instead take any invertible temporary
regulator with covariance tending to zero. Its final data are
$\mathbb C(0)=0$ and $\Theta(0)=I$; no inverse of a physical zero-mode
kinetic operator occurs. The modified Legendre transform subtracts
$\Phi\mathbb R\Phi/2$. Only its fixed perturbative coefficients are
used, not a nonperturbative effective action.

\begin{lemma}[The low source coefficients do not change]
\label{lem:v16-Legendre}
After $\epsilon\downarrow0$, the one-loop tadpole, metric Hessian
and $q^*c$ coefficient in this Legendre representation equal
$\tau_a,\Pi_a,\mathcal Z_a$ of V14--V15, with $\mathcal M_a$
assigned once. The potentially different tree-attached one-loop terms
vanish on the stated panel.
\end{lemma}
\begin{proof}
The Legendre transform fixes the fluctuation mean. To one loop its
insertion coefficient is the Wick contraction of the insertion's second
derivative; the linear mean correction is zero. The difference from a
Wilsonian expression is a one-particle-reducible graph with a hard
propagator leaving a branch containing only the stated external sources.
At these degrees that line has momentum at most $3\mu<\Lambda_0$
and covariance zero. The same support argument removes a pure-gravity
tadpole attached to a hard-linear transformation vertex. These are the
conditions used in \cref{prop:v14-measured,thm:v15-current}. Neither
the source nor the ghost matrix is changed. The $b$ square introduces
no interaction. In every remaining finite graph the $\epsilon$ limit
replaces its covariance by the prescribed supported covariance.
\end{proof}

\section{An explicit reference trace and the complete finite first row}
\label{sec:v16-finite-row}

Left Grassmann derivatives are used, the metric antifield precedes the
retained ghost, and the Wick supertrace convention is
\[
 \tfrac12\operatorname{STr}\log\mathbb H
 =\tfrac12\Tr\log K-\Tr\log M+\text{constant}
\]
for $qKq/2+\bar\gamma M\gamma$. Equivalently, all contractions below
are defined by the Gaussian and Berezin square completions with that
ordering. Multiplying the odd transformation by an independent odd test
parameter gives an even change of variables and fixes its Berezinian
and super-divergence signs.

Write
\[
 \mathbb H_a(\Phi)=D^2\mathscr S_a(\Phi),\quad
 \mathbb C_a(\Phi)=(\mathbb H_a(\Phi)+\mathbb R_a)^{-1},\quad
 \mathbb T_a(\Phi)=Dr_a(\Phi),\quad
 \mathcal B_a(\Phi)=D\mathscr S_a(\Phi)[r_a(\Phi)].
\]
The last quantity is the \emph{actual classical breaking}, formed
from the displayed input action and source before any loop is evaluated.
For $\Phi=(t h e^{-ikx},0,\zeta c e^{ikx},0)$, with $t$ even and
$\zeta$ odd, define
\begin{align}
 \mathcal I_a(h,c;k)&=\tfrac12[t\zeta]\operatorname{STr}
       \{\mathbb C_a(\Phi)D^2\mathcal B_a(\Phi)\},
       \label{eq:v16-I}\\
 \mathcal X_a(h,c;k)&=[t\zeta]\operatorname{STr}
       \{(I-\mathbb C_a(\Phi)\mathbb H_a(\Phi))\mathbb T_a(\Phi)\},
       \label{eq:v16-X}\\
 \mathcal Y_a(h,c;k)&=[t\zeta]D\mathcal M_a(q)[r_a^q(q,c)],\qquad
 \mathfrak d_a(h,c;k)=[t\zeta]\operatorname{sdiv}r_a(\Phi).
       \label{eq:v16-Yd}
\end{align}
Supertraces are divided by physical volume. The $\epsilon$ prescription
is understood, with the following cancellation performed before its
limit.

These formulas have a finite explicit vertex implementation. Put
$V=\mathbb H_a(\Phi)-\mathbb H_{a,0}$ and
$\mathbb C=\mathbb C_a(0)$. Keep only $1,t,\zeta,t\zeta$.
Then
\begin{align}
 \mathcal I_a&=\tfrac12[t\zeta]\operatorname{STr}
 \{(\mathbb C-\mathbb C V\mathbb C+
       \mathbb C V\mathbb C V\mathbb C)D^2\mathcal B_a\},
       \label{eq:v16-I-compiler}\\
 \mathcal X_a&=[t\zeta]\operatorname{STr}
 \{(\Theta-\mathbb C V\Theta+
       \mathbb C V\mathbb C V\Theta)\mathbb T_a\}.
       \label{eq:v16-X-compiler}
\end{align}
Order is retained, including the position of $\Theta$; it does not
commute with a source vertex. The matrices $V$ are derivatives of the
explicit $S_3,S_4$ and ghost polynomials, and $\mathbb T$ uses
$R_0,R_1,R_2$. Form $\mathcal B_a$ first, then differentiate, keeping
the gauge-fixing mismatch and the classical breaking of the upper
completion. No fifth action vertex or third nonlinear metric
transformation is needed for this row. Expanding the coefficient
extraction yields finitely many tadpole and two-vertex traces.

\begin{theorem}[Assembled native one-ghost/one-metric identity]
\label{thm:v16-finite-row}
For the declared action, source and measure, with the auxiliary
$\mathcal M_a$ included in $\Pi_a,\tau_a$ once, one has
\begin{equation}
 \boxed{\begin{aligned}
 &\Pi_a(k)[h,R_0(k)c]+\langle H_a^{\rm cl}(k)h,\mathcal Z_a(k)c\rangle\\
 &\quad+\tau_a[R_1^{(0),E}(h,-k;c,k)]+\mathfrak d_a(h,c;k)\\
 &\hspace{20mm}=\mathcal I_a(h,c;k)+\mathcal X_a(h,c;k)
                                      +\mathcal Y_a(h,c;k).
 \end{aligned}}
 \label{eq:v16-native-row}
\end{equation}
This is exact for finite coefficients on allowed grid momenta and for
thermodynamic loop integrals on the continuous source window. The
reference insertion is \eqref{eq:v16-X-compiler}, with no inverse
cutoff factor. Finite-cutoff BRST invariance is not assumed.
\end{theorem}
\begin{proof}
Fix the mean $\phi$ by ordinary linear sources in the regulated Legendre
representation. Finite super integration by parts with density
$e^{-\mathcal M}$ gives
\[
 \begin{split}
 D\Gamma[\langle r\rangle]={}&\langle D\mathscr S[r]\rangle
 -\hbar\langle\operatorname{sdiv}r\rangle
 +\hbar\langle D\mathcal M[r]\rangle\\
 &+\langle(\mathbb R(\Phi-\phi))\cdot r(\Phi)\rangle.
 \end{split}
\]
At fixed mean the first insertion correction is
$\hbar\operatorname{STr}(\mathbb C D^2f)/2$.
Also $\Gamma=\mathscr S+\hbar\Gamma_1+O(\hbar^2)$, where
$\Gamma_1=\operatorname{STr}\log(\mathbb H+\mathbb R)/2+
\mathcal M$ up to a source-independent constant. Hence
\begin{equation}
 \begin{split}
 D\Gamma_1[r]+D\mathscr S[r_1]+\operatorname{sdiv}r
 ={}&\tfrac12\operatorname{STr}(\mathbb C D^2\mathcal B)\\
 &+\operatorname{STr}(\mathbb C\mathbb R Dr)+D\mathcal M[r].
 \label{eq:v16-one-loop-chain}
 \end{split}
\end{equation}
Wick-contracting the last quadratic fluctuation fixes the plus sign of
the reference trace. A bosonic check of the algebra is
\[
 \tfrac12\Tr(\mathbb C D^2(D\mathscr S[r]))
 =D\Gamma_1^{\rm prop}[r]+D\mathscr S[r_1]
                                  +\Tr(\mathbb C\mathbb H Dr).
\]
The graded product rule gives the same identity with supertraces and
reproduces the negative ghost determinant and V15 mixed-loop sign.
Neither $r^2=0$ nor $D\mathscr S[r]=0$ was used.

Since $\mathbb C\mathbb R=I-\mathbb C\mathbb H$ and
$\mathbb C(0)\mathbb R\to\Theta$, its resolvent expansion gives
\eqref{eq:v16-X-compiler}. The same expansion gives
\eqref{eq:v16-I-compiler}. No $\nu^{-1}$ remains. Extract $t\zeta$
at $b=\bar c=0$. The classical metric derivative is $H_a^{\rm cl}h$;
the source window gives exactly $R_0c$ and $R_1^{(0),E}$. The loop
derivative contributes $\Pi_aR_0+\tau_aR_1$, and the quantum source
is $H_a^{\rm cl}\mathcal Z_a$ by \cref{lem:v16-Legendre}.
The action's $c$ derivative has no term at zero antighost and $b=0$;
$r^{\bar c}=ib$ and $r^b=0$ remove the remaining nonminimal terms
only after this off-shell extraction.

The potentially nonperiodic high-frequency $F$ entries in the $b$
block are not estimated separately. Complete that constant Gaussian
square first; the assembled row then uses the periodic metric/ghost
operands of V14--V15. The finite proof is matrix algebra. In
thermodynamic coefficients cyclicity is justified first with compact
smooth momentum regulators and subsequently in the subtracted
integrable expressions below.
\end{proof}

Every nonpolynomial word in \eqref{eq:v16-X-compiler} contains
$\Theta$. The other loop momenta differ by at most $2\mu$, so the
reference loop is confined to $|p|\le C\Lambda_0$, with $C$ fixed
by the lower profile. It is not an ultraviolet anomaly. The insertion
$\mathcal I_a$, by contrast, is not assumed small term by term;
its cutoff statement below is proved only after signed assembly.

\section{The full coordinate and ultralocal measure contact}
\label{sec:v16-measure}

Let $j(q)=\log\det DG(q)$ in V15's Euclidean chart, and write $u_s$
for the spatial block of a symmetric $u$. Its first two derivatives are
\begin{align}
 j_1(u)&=\tfrac12u_{00}+\tfrac32\operatorname{tr}u_s,\\
 j_2(u,v)&=-\tfrac14u_{00}v_{00}
 -\tfrac9{16}\operatorname{tr}(u_sv_s)
 -\tfrac1{16}\operatorname{tr}u_s\operatorname{tr}v_s.
 \label{eq:v16-j12}
\end{align}
The ultralocal part $m(q)$ of the actual negative-log auxiliary density,
relative to $dq$, has
\begin{align}
 m_1(u)&=3u_{00}+2\operatorname{tr}u_s,\\
 m_2(u,v)&=-\tfrac32u_{00}v_{00}
 -\tfrac{19}{16}\operatorname{tr}(u_sv_s)
 +\tfrac1{16}\operatorname{tr}u_s\operatorname{tr}v_s.
 \label{eq:v16-m12}
\end{align}
Here $j_2,m_2$ are second derivatives, not homogeneous quadratic
coefficients. Both are shift-independent before the transformation is
inserted. The remainder $\mathcal M_a-\sum_xm(q_x)$ is V10--V12's
actual closed-walk/Haar term and is not set to zero.

\begin{proposition}[All metric polarizations of the local measure row]
\label{prop:v16-measure}
On the low source panel,
\begin{align}
 \mathfrak d_a(h,c;k)&=-a^{-4}
 \{j_2(h,R_0c)+j_1(R_1^{(0),E}(h,-k;c,k))\},\\
 \mathcal Y_a^{\rm ul}(h,c;k)&=a^{-4}
 \{m_2(h,R_0c)+m_1(R_1^{(0),E}(h,-k;c,k))\}.
 \label{eq:v16-all-contacts}
\end{align}
In particular
\begin{equation}
 \boxed{\mathfrak d_a-\mathcal Y_a^{\rm ul}
     =\frac{7i}{2a^4}(c\cdot k)\operatorname{tr}h.}
 \label{eq:v16-full-measure-contact}
\end{equation}
This specifies all 40 metric--ghost components of the ultralocal
weighted divergence, not the full auxiliary density on arbitrary sources.
\end{proposition}
\begin{proof}
Put $n=1+q_{00}/2$ and $S_3=I+q_s/2$. V12's exact formulas give
\[
 \begin{split}
 j&=\log n+2\log\det S_3+
                    \sum_{i<j}\log((s_i+s_j)/2),\\
 m&=6\log n+5\log\det S_3-
                    \sum_{i<j}\log((s_i+s_j)/2).
 \end{split}
\]
For the eigenvalues $\lambda_i$ of $q_s$,
\[
 \sum_{i<j}(\lambda_i+\lambda_j)=2\operatorname{tr}q_s,\qquad
 \sum_{i<j}(\lambda_i+\lambda_j)^2
 =\operatorname{tr}q_s^2+(\operatorname{tr}q_s)^2.
\]
Taylor expansion and polarization prove
\eqref{eq:v16-j12}--\eqref{eq:v16-m12}; no eigenvector choice is
needed at repeated eigenvalues.

V15's exact coordinate identity implies
$\operatorname{div}_{dq}r^q=-r^q\sum_xj(q_x)$.
The ghost-source divergence is independent of $q$ and the antighost
source has no antighost derivative, so their $hc$ coefficients vanish.
This proves the first line of \eqref{eq:v16-all-contacts}; the second
is the chain rule for the actual $m$.

Finally $m+j=7\log\det e=(7/2)\log\det G$. On this low panel the
metric source is the ordinary Lie variation. Differentiating the matrix
determinant along it gives
$c\cdot D\log\det g+2D\cdot c$. The metric-coordinate source has
zero flat divergence by V15. Therefore its weighted divergence for
$(\det g)^{-7/2}dg$ is minus $7/2$ times that expression. At opposite
source momenta its $hc$ coefficient is
$+(7i/2)(c\cdot k)\operatorname{tr}h$. The site density is $a^{-4}$.
Direct insertion of $R_1^E=\mathcal L_ch-\mathsf B(h,R_0c)$ in
\eqref{eq:v16-all-contacts} gives the identical answer.
\end{proof}

For $h_{00}=1$, $c=e_0$, $k=ke_0$, the two separate terms are
$ik/(2a^4)$ and $-3ik/a^4$. Their difference is V15's
$7ik/(2a^4)$. The functions $j,m$ are explicit local primitives for
these particular variation terms. This does not infer cohomology of an
unconstructed full finite differential, or cancel a nonlocal anomaly.

\section{A common Taylor assignment, with its finite-mesh spill retained}
\label{sec:v16-projector}

On the low source window set
$H_0^{\rm cl}(k)=\chi X(k)/2$ and
$E_a(k)=H_a^{\rm cl}(k)-H_0^{\rm cl}(k)$.
The evaluated quadratic correction of V11--V13 implies
\begin{equation}
 D^r E_a(0)=0\quad(0\le r\le5),\qquad
 \|\partial_k^\alpha E_a(k)\|
 \le C_\alpha\chi a^4|k|^{6-|\alpha|}\quad(|\alpha|\le6).
 \label{eq:v16-free-spill}
\end{equation}
The first assertion is about the exact analytic Taylor jet, whose first
possible term has degree six. Higher derivatives have the usual fixed
window versions.

Let $\mathsf T_d$ be Taylor projection at zero through momentum degree
$d$. Define
\begin{equation}
 \mathbb P(\tau,\Pi,\mathcal Z)
   =(\tau,\mathsf T_4\Pi,\mathsf T_3\mathcal Z),\qquad
 \mathbb P_{\rm row}=\mathsf T_5.
 \label{eq:v16-projector}
\end{equation}
This retains V14's full component symmetry convention, without assuming
componentwise evenness of $\Pi$. Oddness of $\mathcal Z$ is V15's
specific result. A local action/source counterterm implementing the
assignment through this panel is
\begin{equation}
 \mathcal C_{1,a}=-\tau_a[q]
  -\tfrac12\langle q,(\mathsf T_4\Pi_a)q\rangle
  -\langle q^*,(\mathsf T_3\mathcal Z_a)c\rangle.
 \label{eq:v16-counterterm}
\end{equation}
All three are actual output jets, including auxiliary contacts. Fixed
finite physical normalizations may be added subject to their same
first-row relation. Zero finite matching is used in the comparison,
not a change of the tree Palatini coefficient.

\begin{theorem}[Bounded common local assignment for the first row]
\label{thm:v16-projector}
For
\[
 \mathcal L_a(\tau,\Pi,\mathcal Z)
 =\Pi[h,R_0c]+\langle H_a^{\rm cl}h,\mathcal Zc\rangle
                                +\tau[R_1^{(0),E}(h,c)],
\]
one has
\begin{equation}
 \boxed{\mathsf T_5\mathcal L_a(\tau,\Pi,\mathcal Z)
          =\mathcal L_0(\mathbb P(\tau,\Pi,\mathcal Z)).}
 \label{eq:v16-intertwine}
\end{equation}
The assignment is bounded independently of mesh and physical volume
in scaled finite-panel $C^s$ norms for every fixed $s\ge5$. It is
not asserted to commute with the full nonhomogeneous finite Hessian
off its local jet.

Set $\Pi_a^R=(I-\mathsf T_4)\Pi_a$ and
$\mathcal Z_a^R=(I-\mathsf T_3)\mathcal Z_a$; on row terms use
$\mathcal I_a^R,\mathcal X_a^R,\mathcal Y_a^R
=(I-\mathsf T_5)(\mathcal I_a,\mathcal X_a,\mathcal Y_a)$.
The exact thermodynamic renormalized balance is
\begin{equation}
 \boxed{\begin{aligned}
 &\Pi_a^R[h,R_0c]+\langle H_a^{\rm cl}h,\mathcal Z_a^Rc\rangle
       -\mathcal X_a^R\\
 &\qquad=\mathcal I_a^R+\mathcal Y_a^R
        -\langle E_a h,(\mathsf T_3\mathcal Z_a)c\rangle.
 \end{aligned}}
 \label{eq:v16-renormalized-balance}
\end{equation}
For $r=|\alpha|\le6$, the last term obeys
\begin{equation}
 |\partial_k^\alpha\langle E_a h,(\mathsf T_3\mathcal Z_a)c\rangle|
    \le C_\alpha a^2|k|^{7-r}\|h\|\|c\|.
 \label{eq:v16-spill-bound}
\end{equation}
All local terms, including $\mathfrak d_a$, are assigned, not silently
removed from the measured state.
\end{theorem}
\begin{proof}
$R_0$ is homogeneous of degree one, $H_0^{\rm cl}$ of degree two,
and the tadpole row is of degree one in the opposite source momenta.
Multiplication gives Taylor orders $4,3,0$ inside row degree five.
$E_a$ has zero jet through degree five. This proves
\eqref{eq:v16-intertwine}. Use on each component the scaled norm
\[
 \sum_{|\alpha|\le s}\mu^{|\alpha|}
       \sup_{|k|\le\mu}|\partial^\alpha f(k)|.
\]
For example, a compatible packet norm is
\[
 \|(\tau,\Pi,Z)\|_{\mu,s}^{\rm packet}
 =\mu\|\tau\|+\mu\|\Pi\|_{C^s_\mu}
                  +\chi\mu^2\|Z\|_{C^s_\mu},
\]
where $C^s_\mu$ is the preceding scaled norm, and the row uses
$C^s_\mu$ after unit polarization of $h,c$. Taylor evaluation and
multiplication by the fixed homogeneous symbols are bounded by constants
depending on $s$, finite dimensions and fixed $\chi$, not the mesh
or volume. This does not assert that raw loop jets are small.

Apply $I-\mathsf T_5$ to \eqref{eq:v16-native-row}.
The divergence contact and the tadpole row are linear in $k$, and
\[
 (I-\mathsf T_5)(H_a^{\rm cl}\mathcal Z_a)
 =H_a^{\rm cl}\mathcal Z_a^R+E_a\mathsf T_3\mathcal Z_a.
\]
This proves the balance with the displayed sign. V15 gives
\[
 \|\mathsf T_3\mathcal Z_a(k)\|
 \le C\chi^{-1}\{a^{-2}|k|
 +(1+\log(1/(a\Lambda_0)))|k|^3\}.
\]
Multiply by \eqref{eq:v16-free-spill}. Boundedness of
$(a\Lambda_0)^2[1+\log(1/(a\Lambda_0))]$ and
$|k|\le\Lambda_0/32$ proves \eqref{eq:v16-spill-bound}.
Product differentiation proves the stated derivative bounds. Thus the
large local-source coefficients have been included when the finite
Hessian is compared with its homogeneous part.
\end{proof}

The local part of the same measured identity is retained as the ledger
\begin{equation}
 \mathcal L_0(\mathbb P(\tau_a,\Pi_a,\mathcal Z_a))+
 \mathfrak d_a=\mathsf T_5(\mathcal I_a+\mathcal X_a+\mathcal Y_a).
 \label{eq:v16-local-ledger}
\end{equation}
Thus $\mathsf T_5\mathcal X_a$ is stored as the local normalization of
the reference insertion, not assumed zero or removed from the measured
state. An identity with the \emph{unsubtracted} reference trace additionally
requires finite matching data satisfying this ledger. The theorem below
uses the stated normalized reference $\mathcal X_a^R$; it does not claim
that independently imposed physical finite normalizations automatically
satisfy the unnormalized identity.

This is an intertwining statement for an explicitly displayed linear
\emph{row}, not the full nonlinear BV operator. It supplies neither
an all-order cohomological splitting nor an inverse on arbitrary local
anomaly candidates. The finite component banks may grow at subsequent
orders, as required by the ESM EFT filtration.

\section{The continuum reference row and signed localization of the breaking}
\label{sec:v16-continuum}

Here subscript $0$ denotes zero mesh with ordinary derivatives, the
actual affine coframe chart, the V12 restored vertices, and the same
strictly positive lower scale $\Lambda_0$. It does not remove that
infrared scale.

\begin{proposition}[The reference trace is a lower-scale source]
\label{prop:v16-reference-limit}
Nonpolynomial terms of \eqref{eq:v16-X-compiler} are supported at
$|p|\le C\Lambda_0$. For every fixed derivative,
\begin{align}
 |\partial_k^\alpha\mathcal X_0(h,c;k)|
 &\le C_\alpha\Lambda_0^{5-|\alpha|}\|h\|\|c\|,\\
 |\partial_k^\alpha(\mathcal X_a^R-\mathcal X_0^R)(h,c;k)|
 &\le C_\alpha a^3\Lambda_0^2|k|^{6-r}\|h\|\|c\|,
                      \quad r=|\alpha|\le6.
 \label{eq:v16-X-bounds}
\end{align}
Also
\begin{equation}
 |\partial_k^\alpha\mathcal X_0^R(h,c;k)|
 \le C_\alpha\Lambda_0^{-1}|k|^{6-r}\|h\|\|c\|\quad(r\le6).
 \label{eq:v16-XR}
\end{equation}
Vanishing of $\mathcal X_0^R$ is not asserted.
\end{proposition}
\begin{proof}
Every word has the root $\Theta$; all other labels differ by sums of
the two soft transfers. Terms with a covariance contain the smooth
factor $\nu$, zero near the physical zero modes. Contributions with
no covariance are finite polynomial contacts. After the constant $b$
square, the remaining matrices are the interior propagators of V13
and source polynomials of V15, multiplied by fixed lower-scale
profiles. Their derivatives are bounded on a compact rescaled domain.

The mode count supplies four powers and the transformation one,
giving $\Lambda_0^5$; each momentum derivative removes one power.
In each word of \eqref{eq:v16-X-compiler}, a hard vertex of order
two is paired with a covariance of order minus two. The nonminimal
blocks cancel to this expression before estimation.
For sufficiently small $a$, $s=\zeta=1$ on this fixed ball.
The corrected geometric vertices differ by
$O(a^3\Lambda_0^5)$, the free inverses relatively by
$O((a\Lambda_0)^4)$, and the nonlinear source/ghost vertices have
the ordinary chart values. The raw difference therefore has bound
$C_\alpha a^3\Lambda_0^{8-|\alpha|}$. Taylor's integral formula
through degree five proves the second line of
\eqref{eq:v16-X-bounds}; the first line gives
\eqref{eq:v16-XR}. Mode counting proves the analogous interpolated
finite-sum bounds for $\Lambda_0L_\nu\ge1$.
\end{proof}

\begin{lemma}[Zero-mesh first modified identity]
\label{lem:v16-continuum-identity}
Let $\Pi_0^R,\mathcal Z_0^R$ be the thermodynamic continuum
coefficients of V14--V15 with common zero finite matching. Then
\begin{equation}
 \boxed{\Pi_0^R(k)[h,R_0(k)c]
 +\langle H_0^{\rm cl}(k)h,\mathcal Z_0^R(k)c\rangle
              =\mathcal X_0^R(h,c;k).}
 \label{eq:v16-continuum-identity}
\end{equation}
This is derived from the zero-mesh action and trace algebra, not
assumed in a finite-scale estimate.
\end{lemma}
\begin{proof}
The zero-mesh independent-connection action is the boundary-completed
Palatini action, whose stationary reduction is the Einstein action in
$G(q)$. V11--V12 prove its first two deformation equations with
$R_1,R_2$. The nonminimal part here is $s\Psi$, with
\[
 \Psi=\langle\bar c,Fq+b/(2i\chi)\rangle,\qquad
 sc=c\cdot\partial c,\quad s\bar c=ib,\quad sb=0.
\]
The ordinary Lie identity transferred through $G$ gives $s^2=0$
at the required degree. Consequently all classical-breaking
coefficients feeding this loop vanish: $q^3c$, $bq^2c$,
$\bar c q c^2$, and the necessary lower coefficients. These use
$S_4,R_2$, not an all-order finite-lattice identity.

Justify the trace first with a smooth ordinary ultraviolet covariance
factor $u(p/M)$, one on a central ball and zero beyond twice it.
Apply the graded Hessian product identity in
\eqref{eq:v16-one-loop-chain} to the continuum vertex polynomials.
The local action identity just proved eliminates their breaking
coefficient. The reference factor is now
$1-\nu u=\Theta+\nu(1-u)$.
Its $\Theta$ term is \eqref{eq:v16-X-compiler} at zero mesh.
In the upper reference terms a covariance or cutoff derivative pins
the nonpolynomial expression to $|p|\gtrsim M$. A term with no
covariance is a polynomial divergence contact, assigned before tracing.
The assembled row has ultraviolet degree five; six external
momentum derivatives bound an annulus of radius $R$ by $C/R$.
The dyadic tail is $C/M$, and differentiated cutoff factors have the
same bound. Taylor subtraction through degree five therefore leaves
$C|k|^6/M$. This supplies the finite-part justification of cyclic
loop shifts, rather than assuming an unsubtracted infinite trace.

After common orders $(4,3,5)$ the ultralocal coordinate and measure
contacts have no remainder. The auxiliary nonlocal part has zero
continuum limit by V10 at this fixed order. Let $M\to\infty$.
The result is \eqref{eq:v16-continuum-identity}. The lower profile
was not removed, so the theorem gives a modified, not an ordinary
zero-right-hand-side, identity.
\end{proof}

\begin{theorem}[Cutoff limit of the first renormalized measured row]
\label{thm:v16-row-limit}
In the thermodynamic convention define
\begin{equation}
 \mathcal W_a^R(h,c;k)=\Pi_a^R(k)[h,R_0(k)c]
 +\langle H_a^{\rm cl}(k)h,\mathcal Z_a^R(k)c\rangle
 -\mathcal X_a^R(h,c;k).
 \label{eq:v16-WR}
\end{equation}
For every fixed $r=|\alpha|\le6$,
\begin{equation}
 \boxed{|\partial_k^\alpha\mathcal W_a^R(h,c;k)|
 \le C_\alpha a|k|^{6-r}\|h\|\|c\|.}
 \label{eq:v16-main-limit}
\end{equation}
Higher fixed derivatives have bound $C_r a\mu^{6-r}$ on a fixed
window, with constants allowed to depend on $\mu/\Lambda_0$,
profiles, $\chi$ and derivative order. The same cutoff rate applies
to the signed renormalized insertion $\mathcal I_a^R$, with the
auxiliary and spill terms owned as in
\eqref{eq:v16-renormalized-balance}. Thus the ultraviolet classical
breaking has no surviving nonlocal $hc$ remainder in this limit;
its local matching coefficients have not been set to zero beforehand.
\end{theorem}
\begin{proof}
V14 gives $\Pi_a^R-\Pi_0^R=O(a|k|^5)$ with fixed higher
derivatives and the auxiliary contribution assigned once. Multiplying
by $R_0$ gives $O(a|k|^6)$. V15 gives
$\mathcal Z_a^R-\mathcal Z_0^R=O(\chi^{-1}a^2|k|^5)$;
multiplication by $H_0^{\rm cl}$ gives $O(a^2|k|^7)$.
The continuum subtracted source bound is
$\mathcal Z_0^R=O(\chi^{-1}\Lambda_0^{-2}|k|^5)$.
Together with \eqref{eq:v16-free-spill}, it bounds the remaining
free-matrix change by a smaller quantity on this window.
Proposition~\ref{prop:v16-reference-limit} bounds the reference change
by $Ca^3\Lambda_0^2|k|^6$. Subtract the independently proved
\eqref{eq:v16-continuum-identity}; $a\Lambda_0$ and $a\mu$ are
bounded. This proves \eqref{eq:v16-main-limit}. Product
differentiation and the inherited higher derivative estimates give
all stated versions, including $r=6$.

Do not bound individual ultraviolet terms of
$D\mathscr S_a[r_a]$ first. The exact signed balance gives
\[
 \mathcal I_a^R=\mathcal W_a^R-\mathcal Y_a^R
        +\langle E_a h,(\mathsf T_3\mathcal Z_a)c\rangle.
\]
The ultralocal $\mathcal Y_a$ is linear in $k$; its remaining part
is the variation of V10's actual auxiliary two-jet. Its subtracted
bound is $Ca|k|^6$ by that theorem and the one-derivative source.
Use \eqref{eq:v16-spill-bound} for the last term. This proves the
insertion statement. The logical order is component cutoff estimates,
finite algebra, and an independent zero-mesh local action identity,
followed by the conclusion about the assembled breaking. No finite
Ward identity is assumed to obtain the estimates it is used with.
\end{proof}

An added finite local packet $(\tau_f,\Pi_f,Z_f)$ must satisfy its
matching first-row relation; arbitrary independent normalizations need
not preserve the identity. Nor does this result identify the completed
regulator with the discontinuous unmodified one of V14's negative test.

\section{Rooted topology, finite volume, and an independent algebra check}
\label{sec:v16-checks}

Multiply \eqref{eq:v16-WR} by a fixed smooth window supported strictly
inside $|k|<\mu$. Its thermodynamic kernel, sampled and periodized,
satisfies for every prescribed $b$
\begin{equation}
 \boxed{a^4\sum_x(1+\mu d_{\rm per}(x,0))^b
       \|\mathcal K_{a,L}^{\mathcal W}(x)\|
                     \le C_b a\mu^6.}
 \label{eq:v16-kernel}
\end{equation}
The action row has the additional factor $\hbar$. To prove this,
rescale $k=\mu z$, integrate the Fourier transform by parts more than
$b+4$ times using the fixed higher-derivative estimates, sample and
periodize. Periodic distance is no larger than the distance of each
lift. The associated bilinear functional is bounded with one source
in $L^1$ and the other in $L^\infty$, not by an extensive determinant
norm.

The original finite-volume coefficients have a separate volume error.
For every fixed $N>4$, differentiation of the subtracted continuum
integrands $N$ times in the loop label is integrable. Poisson
summation, as in V8--V15, gives the conservative comparison
\begin{equation}
 \begin{split}
 |\partial_k^\alpha\mathcal W^R_{a,L}(h,c;k)|\le{}&
 C_\alpha a|k|^{6-r}\|h\|\|c\|\\
 &+C_{N,\alpha}\Lambda_0^{5-r}
       (\Lambda_0L_{\min})^{-N}\|h\|\|c\|.
 \label{eq:v16-volume}
 \end{split}
\end{equation}
Here the finite-volume expression uses the declared smooth
interpolations for comparison. No exact renormalized identity for
arbitrary off-grid torus interpolations is asserted. The second term
follows from the compact support of the reference trace and the
integrable differentiated propagating/source integrands. Its lower-scale
power is fixed by the row dimension five. This deliberately weaker
bound distinguishes ultraviolet and thermodynamic limits; neither
removes $\Lambda_0$.

\subsection{An exact finite check of reference and nonminimal signs}

The following calculation tests the contraction compiler; it is not a
surrogate gravity model or another continuum theorem. Let an even
coordinate $z$, an auxiliary $b$ and an odd pair $(c,\bar c)$ have
\[
 S=V(z)+b^2/(2\chi)+ibfz+\bar c M(z)c,\qquad
 rz=r(z)c,\quad r\bar c=ib,\quad rc=rb=0.
\]
Use arbitrary independent parameters in
\[
 V'(z)=\alpha z+gz^2/2+\lambda z^3/6,\quad
 M(z)=m_0+m_1z+m_2z^2/2,\quad
 r(z)=r_0+r_1z+r_2z^2.
\]
For covariance fraction $0<\nu\le1$ and mean $z=h$, set
\[
 C_g=\frac{\nu}{\alpha+\chi f^2+\nu(gh+\lambda h^2/2)},\quad
 C_c=\frac{\nu}{m_0+\nu(m_1h+m_2h^2/2)},\quad
 Z=\tfrac12C_gr''-C_gC_cM'r'.
\]
The explicit one-loop contraction of $B(z)c+ibD(z)c$ is
\[
 \mathfrak i(B,D)=\tfrac12C_gB''-C_gC_cM'B'
                     +\chi f C_gD'-\chi f C_gC_cM'D.
\]
The actual classical-breaking insertion is
$\mathfrak i(V'r,fr+M)$ and the reference insertion is
\begin{equation}
 X=(\nu^{-1}-1)
       \{\mathfrak i(\alpha z r,fr+m_0)-\alpha hZ\}.
 \label{eq:v16-check-reference}
\end{equation}
Exact finite Gaussian/Berezin algebra yields
\begin{equation}
 (\tfrac12C_gV'''-C_cM')r+V'Z+r'
                   =\mathfrak i(V'r,fr+M)+X.
 \label{eq:v16-check-identity}
\end{equation}
It is an identity of rational functions on their nonsingular domain.
The checker verifies its constant and linear mean coefficients with
independent symbolic parameters, the two needed for the present row.
It also checks cancellation of the apparent $\nu^{-1}$ singularity.
The $\chi f$ terms come from the nonminimal square and cannot be
omitted. This check fixes signs and trace order; it is not a numerical
proof of the cutoff theorem.

\subsection{Verification scope}

The self-contained script
\srcid{check_ESM_first_Slavnov_row_V16.py} verifies the complete four-dimensional contact
\eqref{eq:v16-full-measure-contact} for arbitrary symmetric $h$,
ghost $c$ and momentum $k$, the common Taylor assignment on explicit
polynomial matrices, and the ordered resolvent compiler with
noncommuting matrices. It separately compares numerical values of
\eqref{eq:v16-check-identity} on nonsingular inputs. The finite algebra
and the cutoff proof have different evidentiary roles. No new Lean
formalization or independent peer review is asserted. A separate
SHA-256 comparison verifies that the complete V15 mathematical body
is retained.

\section{Final ledgers for V16}
\label{sec:v16-ledgers}

\subsection*{Proved}
\addcontentsline{toc}{subsection}{V16: Proved}
The actual completed action and V15 source have the explicit off-shell
representation \eqref{eq:v16-offshell-action}. The first measured row
is \eqref{eq:v16-native-row}, with independently defined breaking
and reference insertions given by
\eqref{eq:v16-I-compiler}--\eqref{eq:v16-X-compiler}. The reference
has no inverse-cutoff singularity. All ultralocal coordinate/auxiliary
contacts are \eqref{eq:v16-all-contacts}, with the full tensor
simplification \eqref{eq:v16-full-measure-contact}. The common
row assignment obeys \eqref{eq:v16-intertwine}, with an explicitly
bounded finite-mesh spill. The component cutoff theorems and the
independent zero-mesh action identity give the normalized modified
continuum row \eqref{eq:v16-continuum-identity}, cutoff error
\eqref{eq:v16-main-limit}, and rooted bound \eqref{eq:v16-kernel}.
The signed nonlocal breaking insertion tends to zero at this source
order; its complete local matching ledger is retained.

\subsection*{Inherited}
\addcontentsline{toc}{subsection}{V16: Inherited}
V1--V15 and their labels are unchanged. The completed cubic/quartic
and zero-mesh BRST coefficients are V11--V14 inputs. V15 supplies
the source and $q^*c$ theorem. V10--V12 supply the auxiliary measure,
coframe Jacobian and fixed-order matching. Generic finite Gaussian,
Berezin, Legendre and integration-by-parts identities are tools, not
newly claimed general principles. The structural SM and classical
Einstein results retain their original scope and are not used as
quantum estimates.

\subsection*{Literature input}
\addcontentsline{toc}{subsection}{V16: Literature input}
The relationship of Wilsonian quantum-master identities to modified
Legendre Slavnov identities is established in \cite{V15IIM};
\cite{V13BFR} provides broader perturbative gravity context. The
explicit trace, contact and matching statements here concern the
chosen Renewal completion. No first general construction of
perturbative gravitational Ward identities or all-orders ESM is
claimed.

\subsection*{Conditional}
\addcontentsline{toc}{subsection}{V16: Conditional}
The common assignment is a bounded projector for this row, not a
splitting of the full combined diffeomorphism/Lorentz/SM BV complex.
The fixed lower scale retains its reference trace; removing it needs
massless/infrared source estimates. Matching within V14's smooth
completion class is not matching to the unmodified microscopic action.
A larger source panel or higher loop order requires its native
vertices, measure and norm estimates. Physical finite normalizations
must respect \eqref{eq:v16-local-ledger}; independent choices are
not automatically Slavnov compatible. The reference remains
holomorphic Euclidean without a Lorentzian reconstruction theorem.

\subsection*{Open}
\addcontentsline{toc}{subsection}{V16: Open}
The next smallest quantum row has one ghost and two metrics. It needs
the native fifth action coefficient, third nonlinear coframe
transformation, and the one-loop three-metric and $q^*qc$
coefficients. Its lower-scale reference and source contacts must be
retained. A useful next attack is that joint bank and its common
assignment, not a relabeling of the present result as the full Slavnov
functional. Internal gauge, chiral fermion, determinant-line and mixed
ESM extensions, filtered interacting R1--R4 realization and arbitrary
loop-order continuum remain open in this lineage. No nonperturbative
gravity, asymptotic safety, finite-parameter UV completion or YM mass
gap conclusion follows.

\begin{center}\small
\begin{tabular}{>{\raggedright\arraybackslash}p{.23\textwidth}>{\raggedright\arraybackslash}p{.68\textwidth}}
\toprule
Variable & Scope in V16\\
\midrule
Loop/source order & One propagating loop, one retained metric and one
retained diffeomorphism ghost; action/source inputs through $S_4,R_2$.\\
Mesh & Fixed V14 smooth completion, $a\Lambda_0\le1/4096$;
all ultraviolet modes included with their actual weights.\\
Momentum/lower scale & $|k|\le\mu$, $\Lambda_0\ge32\mu>0$;
no infrared removal or bound as $\chi\downarrow0$.\\
Volume & Raw finite-grid identity on allowed transfers. Exact normalized
modified identity in thermodynamic normalization; separate finite-volume
comparison \eqref{eq:v16-volume}. No extensive rooted bound.\\
Derivative order & Every prescribed finite momentum derivative and
kernel moment, with order-dependent constants.\\
Reference/measure & Existing holomorphic contour and nonminimal,
coordinate, auxiliary and Haar factors, represented once.\\
Iteration & Summable nonlocal coefficient tails and a common first-row
assignment, not an invariant interacting ESM ball.\\
\bottomrule
\end{tabular}
\end{center}


\clearpage
\part*{V17: A nonlinear native counterterm generator and the second measured Slavnov row}
\addcontentsline{toc}{part}{V17: Native generator and second measured Slavnov row}

\section{Move upstream once, then assemble the next quantum row}
\label{sec:v17-start}

Sections 1--127 are preserved. V16 leaves the coefficient with one
retained diffeomorphism ghost and two retained metric-coordinate fields
unclosed. Its required new quantities are a three-metric one-loop
coefficient, a field-dependent quantum transformation insertion, the fifth
native action coefficient, and the third nonlinear transformation. An
independent subtraction of these quantities would not establish their
common measured identity. Nor does the first row determine the next row.

There is a useful upstream reduction. The first two mesh corrections of
the stationary native action are coefficients of one explicit local
functional, analytic in the coframe. It generates the earlier quadratic,
cubic and quartic counterterms and the required fifth-order counterterms
without selecting them by the desired Ward identity. The affine-coframe
map likewise gives a recursion for every nonlinear transformation
coefficient. We establish these two statements first, then use their
fifth-order output in a complete second-row calculation.

The previous coefficient compiler, Cartan inverse, Gaussian/Berezin
identities, first-row proof and earlier ultraviolet estimates are
inherited. The new claims below concern the displayed native generator,
the new vertices and loop coefficients, and the next common source
assignment. General perturbative effective gravity and cutoff-modified
Slavnov identities are established frameworks
\cite{V13BFR,V15IIM}; they are comparison literature, not proofs of this
regulator's estimates. The target remains the finite-order ESM programme
\cite{V6Programme}. In particular, this installment does not supply the
chiral or mixed spin/gauge sector, an all-loop invariant class, or removal
of the retained infrared scale.

\subsection{Conventions and the exact extension being made}

For the classical native calculation use V10's completed Palatini action,
V11's affine ADM coframe, $j=0$, zero cosmological and independent Holst
coefficients, and the stated periodic boundary convention. For propagating
loops use exactly V13's holomorphic Euclidean coefficient reference and
fixed trace contour, V14's periodic free completion, V15's pulled-back
source, and V16's off-shell nonminimal field. This is the declared complex
perturbative reference, not a newly asserted Lorentzian continuation.

Only interaction degree five is added to V16's action. All lower action
and source coefficients remain unchanged. The auxiliary connection density
is inserted once at order $\hbar$, relative to the actual $q$ coordinate;
the coefficient divergence remains flat. The extra upper fifth-order
vertices will be specified explicitly in \cref{sec:v17-measured}.

For loop estimates fix $\chi>0$ and smooth profiles in the classes already
used, and impose the slightly smaller common domain
\begin{equation}
 \Lambda_0\ge64\mu>0,\qquad a\Lambda_0\le1/8192,\qquad
 \Lambda_0L_\nu\ge1.
 \label{eq:v17-domain}
\end{equation}
All three external momenta have norm at most $\mu$ and sum to zero.
Use $K=(k_1,k_2)\in\R^8$, $k_c=-k_1-k_2$, and $|K|$ for the
ordinary eight-dimensional norm. Linear changes of this routing change
only fixed constants. Fourier coefficients and traces are divided by
physical volume. A superscript $R$ denotes assignment of the actual local
Taylor jet, not deletion of a source coefficient. The common zero finite
matching convention is used in continuum comparisons.

\section{A local generating functional for the native mesh artifacts}
\label{sec:v17-generator}

Write $C(q)=\mathsf C(e(q))$ and $b_a(q)=f_a^\circ(e(q))$.
The inverse $C(q)^{-1}$ is the pointwise $24$-component Cartan inverse on
the inherited nondegenerate chart. Explicitly, the native congruence
\eqref{eq:v10-C-congruence} gives
\begin{equation}
 C(q)^{-1}=\frac{1}{\chi\det e(q)}
 (e(q)^T\otimes I_6)H_0^{-1}(e(q)\otimes I_6).
 \label{eq:v17-local-inverse}
\end{equation}
Here $H_0$ is the fixed flat Cartan matrix of that congruence. This is
not a propagator in spacetime. Let $b_0$ be the same cofactor load with
ordinary derivatives, and let
$b_2$ replace its single spatial derivative by
$\mathfrak d_i=\partial_i^3/24$, with $\mathfrak d_0=0$. Thus
$b_a=b_0+a^2b_2+O(a^4)$ on each fixed nonaliased coefficient panel.

In this section superscripts on the connection functionals denote
coframe degree only when explicitly stated. Define
\begin{align}
 P(q,A)&=\mathscr R_{3,\mathrm{pt}}(e(q),A),\\
 Q(q,A)&=\mathscr R_{3,\mathrm{sh}}(e(q),A)
                         +\mathscr R_{4,\mathrm{pt}}(e(q),A).
 \label{eq:v17-PQ}
\end{align}
Here the first two functionals are exactly
\eqref{eq:v11-Rpt}--\eqref{eq:v11-Rsh}, with their flat cofactor
coefficients replaced by the full quadratic cofactor of $e(q)$.
The fourth-word functional is \eqref{eq:v12-R4}, with the same full
cofactor. Its magnetic word is
\[
 B_{4,\mathrm{com}}(X,Y)=\tfrac16[X,[X,[X,Y]]]
 +\tfrac14[Y,[X,[X,Y]]]+\tfrac16[Y,[Y,[X,Y]]].
\]
The functional gradient in the shift term uses its complete Euler
adjoint, as in \eqref{eq:v12-force-rule}. In particular all backward
occurrences produced by variation are retained. These instructions
specify finite differential expressions; they leave no unknown
connection vertex.

Put
\begin{equation}
 X(q)=-C(q)^{-1}b_0(q),\qquad U(q)=D_AP(q,X(q)),
 \label{eq:v17-XU}
\end{equation}
and define
\begin{align}
 \mathcal A_0(q)&=-\tfrac12\langle b_0,C^{-1}b_0\rangle,\\
 \mathcal A_1(q)&=P(q,X(q)),\\
 \mathcal A_2(q)&=\langle b_2,X\rangle+Q(q,X)
                      -\tfrac12\langle U,C^{-1}U\rangle.
 \label{eq:v17-generators}
\end{align}
Every pairing here is the native physical pairing. $P$ has no derivatives
of its connection arguments before $X$ is substituted. $Q$ has the
explicit first-shift derivative. Thus the three functionals have total
spacetime derivative orders two, three, and four. Their coefficients are
analytic local functions of the coframe on the fixed chart.

\begin{theorem}[Native nonlinear artifact generator]
\label{thm:v17-generator}
In the local mesh expansion of the actual stationary action,
\begin{equation}
 \boxed{S_a^*(q)=\mathcal A_0(q)+a\mathcal A_1(q)
                              +a^2\mathcal A_2(q)+O(a^3).}
 \label{eq:v17-generated-action}
\end{equation}
The statement holds at every prescribed finite field degree. If $[\cdot]_n$
denotes the homogeneous field coefficient and all of its external momenta
have norm at most $\Lambda_*$, then for each fixed $n\ge3$ and fixed
multi-index $\alpha$, on $na\Lambda_*\le1/64$,
\begin{equation}
 \left|\partial_K^\alpha
 [S_a^*-\mathcal A_0-a\mathcal A_1-a^2\mathcal A_2]_n\right|
 \le C_{n,\alpha}a^3\Lambda_*^{5-|\alpha|}
                                      \prod_{j=1}^n\|q_j\|.
 \label{eq:v17-generator-rate}
\end{equation}
At field degree two the stronger $C_\alpha a^4\Lambda_*^{6-|\alpha|}$
bound applies. Constants are independent of periods and mesh on this
panel, but not of $n$ or derivative order. There is no claim of convergence
of an all-field series in a cutoff-independent ultraviolet norm.
\end{theorem}
\begin{proof}
At fixed coframe the actual pre-elimination action through mesh degree two
is
\[
 \langle b_0,A\rangle+\tfrac12\langle A,CA\rangle
 +aP(q,A)+a^2\{\langle b_2,A\rangle+Q(q,A)\}.
\]
This follows by expanding the precise electric exponential and ordered
magnetic logarithm of \eqref{eq:v10-full-remainder}, including the first
translation in the cubic word. Degrees at least five in the connection
begin with $a^3$. The phase load has no term linear in $a$.

Write the stationary connection as $A_*=X+aZ+a^2W+O(a^3)$.
The zeroth and first mesh equations give
$CX+b_0=0$ and $CZ+U=0$. Hence $Z=-C^{-1}U$.
Inserting this expression in the action cancels every term containing
$W$ by the zeroth stationary equation. The coefficient at $a^2$ involving
$Z$ is
\[
 \tfrac12\langle Z,CZ\rangle+\langle U,Z\rangle
                         =-\tfrac12\langle U,C^{-1}U\rangle.
\]
The remaining terms are precisely \eqref{eq:v17-generators}. The argument
uses an invertible symmetric Cartan form, not positivity or a minimum.

To justify the stated finite coefficient estimates, expand the analytic
pointwise inverse in coframe degree around $C_0$. At any fixed degree
there are finitely many coefficient trees, each with fixed $C_0^{-1}$,
quadratic cofactor vertices and the actual finite connection words.
This is the inherited finite coefficient compiler applied to the displayed
native operands. Each subset momentum has norm at most $n\Lambda_*$.
The phase factors and translations have uniformly bounded differentiated
Taylor remainders on the declared compact chart. A tree-level coefficient
has two physical derivatives; a mesh factor $a^j$ has $j$ additional
momentum powers. The first unassigned power is $j=3$, except at quadratic
field degree where odd powers are absent and $j=4$ is first. Rescale
momenta by $\Lambda_*$ and apply Taylor's integral remainder to each of
the finitely many trees. This proves \eqref{eq:v17-generator-rate} and
all its fixed derivative versions. No estimate for an extensive action is
used.
\end{proof}

\begin{corollary}[One compatible counterterm bank rather than independent repairs]
\label[corollary]{cor:v17-bank}
For each $N$ the coframe-only insertion
\begin{equation}
 \mathcal C_{a,\le N}=-[a\mathcal A_1+a^2\mathcal A_2]_{\le N}
 \label{eq:v17-counterterm-generator}
\end{equation}
restores the classical native action to its zero-mesh coefficients through
mesh degree two and field degree $N$. At $N\le4$ it is exactly the
counterterm bank of V11--V12. Increasing $N$ does not change any earlier
coefficient. Every bank is finite; the analytic generating functional is
not a claim that finitely many constant couplings determine all gravity
EFT orders.
\end{corollary}
\begin{proof}
Coefficient extraction in \eqref{eq:v17-generators} gives at degree two
$\langle b_{1,2},A_{1,0}\rangle$, at degree three the three terms of
\eqref{eq:v11-artifact-bank}, and at degree four the two expressions
\eqref{eq:v12-O4}. In particular the squared cubic force is included
once. Uniqueness of coefficients in the native expansion proves equality
with those earlier banks. The remaining conclusions follow from
\cref{thm:v17-generator}.
\end{proof}

The new fifth-order insertion is only the degree-five part of
\eqref{eq:v17-counterterm-generator}; the earlier insertions are not
added a second time. It changes the reduced and pre-elimination actions
by the same coframe-only factor. It leaves the actual stationary Cartan
section and the auxiliary determinant unchanged. Local covariance under
a full nonlinear finite diffeomorphism group is not asserted for the
individual restoring operators.

\section{The fifth native action and the third transformation are explicit}
\label{sec:v17-fifth}

Use the native homogeneous notation of \eqref{eq:v10-stationary-jets}:
$X_a=A_{1,a}$ and $Y_a=A_{2,a}=-C_0^{-1}G_{2,a}$.
The quantities $\mathscr R_{n,a}^{(r)}$ below are the coefficients of
coframe degree $r$ in \eqref{eq:v10-Rn}, including their actual shifts.
In the affine chart, $b$ and $C$ have no terms above degree two.

\begin{proposition}[Fifth coefficient without a third stationary solve]
\label[proposition]{prop:v17-S5}
The exact native homogeneous coefficient is
\begin{equation}
\boxed{\begin{aligned}
 S_{5,a}^*={}&\langle X_a,C_2Y_a\rangle
                  +\tfrac12\langle Y_a,C_1Y_a\rangle\\
 &+a\{\mathscr R_{3,a}^{(2)}(X_a)
        +D\mathscr R_{3,a}^{(1)}(X_a)[Y_a]
        +\tfrac12D^2\mathscr R_{3,a}^{(0)}(X_a)[Y_a,Y_a]\}\\
 &+a^2\{\mathscr R_{4,a}^{(1)}(X_a)
                      +D\mathscr R_{4,a}^{(0)}(X_a)[Y_a]\}
       +a^3\mathscr R_{5,a}^{(0)}(X_a).
\end{aligned}}
\label{eq:v17-S5}
\end{equation}
Here $\mathscr R_{5,a}$ is the explicit finite word $B_5$ and electric
factor of \eqref{eq:v10-Rn}, not an unspecified new vertex. The corrected
coefficient
\begin{equation}
 \widehat S_{5,a}=S_{5,a}^*-[a\mathcal A_1+a^2\mathcal A_2]_5
                  =S_{5,0}+\mathcal E_{5,a}
 \label{eq:v17-corrected-S5}
\end{equation}
has the bound \eqref{eq:v17-generator-rate} at $n=5$.
\end{proposition}
\begin{proof}
Substitute $A_* = tX_a+t^2Y_a+t^3A_3+t^4A_4+\cdots$ in the native
action and take $[t^5]$. The $A_4$ terms multiply $b_1+C_0X_a=0$.
The $A_3$ terms multiply
$b_2+C_1X_a+C_0Y_a+aD\mathscr R_{3,a}^{(0)}(X_a)=0$.
The remaining Cartan terms are the first line of \eqref{eq:v17-S5}.
Taylor expansion of the cubic connection polynomial gives its three
terms on the second line; the fourth and fifth words give the last line.
These are all partitions of total degree five, and none remains involving
an unknown stationary coefficient. The generating theorem supplies the
last assertion. This is a new evaluation of the native coefficient, not a
new general stationary-elimination principle.
\end{proof}

Let $\mathsf B_qv=\mathsf B(q,v)$ in the actual coframe chart. In the
Lorentzian calculation its contractions use $\eta$ as in V12; for the
holomorphic coefficient reference use the Euclidean map of V15. Since
$DG=I+\mathsf B_q$ and $G-I=q+\mathsf Q(q,q)$, multiplication of the
coordinate equation gives
\begin{equation}
 \boxed{R_3^{(0)}(q,c)=-\mathsf B_qR_2^{(0)}(q,c),\qquad
 R_n^{(0)}=(-\mathsf B_q)^{n-2}R_2^{(0)}\quad(n\ge2).}
 \label{eq:v17-R3}
\end{equation}
The same recursion holds \emph{exactly} for V15's finite filtered source
coefficients once its first three coefficients have been fixed:
\begin{equation}
 \mathscr R_{n,a}=-\mathsf B_q\mathscr R_{n-1,a}\quad(n\ge3).
 \label{eq:v17-finite-R3}
\end{equation}
No individual high-momentum $q$ factor is filtered inside $\mathsf B_q$.
The high--high contacts already required by V15 are therefore preserved.

\begin{proposition}[Local third transformation and classical degree-five closure]
\label[proposition]{prop:v17-classical}
The coefficients in \eqref{eq:v17-R3} are local, of first differential
order, and the expression at $n=3$ is cubic in $q$. With the corrected
actions through degree five and their chart transformations,
\begin{equation}
 \boxed{DS_{5,0}[R_0c]+DS_{4,0}[R_1^{(0)}]
              +DS_{3,0}[R_2^{(0)}]+DS_{2,0}[R_3^{(0)}]=0.}
 \label{eq:v17-classical-row}
\end{equation}
The corresponding classical master equation at interaction degree three
holds modulo $a^3$. On a fixed source window the action-side defect has
bound $C_\alpha a^3\Lambda_*^{6-|\alpha|}$ times the four metric and one
ghost polarization norms. The new ghost-action contact is
$-\langle\bar c,F R_3^{(0),E}(q,c)\rangle$.
\end{proposition}
\begin{proof}
The equation
$(I+\mathsf B_q)R=R_0c+\mathcal L_cq+\mathcal L_c\mathsf Q(q,q)$
has a right side of field degree at most two. Equating degrees proves
\eqref{eq:v17-R3}; replacing its right side by V15's two filtered terms
proves \eqref{eq:v17-finite-R3}. Pointwise inversion and multiplication
introduce no derivative inverse.

At zero mesh the native stationary action is the boundary-completed
Palatini action evaluated at its unique torsion-free stationary
connection, as established in V11. Its pullback variation is zero.
Extract its coefficient with four $q$ fields and one ghost to get
\eqref{eq:v17-classical-row}. The same invertible coordinate map transports
the ordinary Lie algebra identity. Its next field coefficient gives the
metric-antifield/two-ghost part of
$(\mathbb S_0,\mathbb S_3)+(\mathbb S_1,\mathbb S_2)=0$.
The ghost bracket has no metric dependence, so no new ghost-antifield
structure constant is inserted. The generating theorem and the current
low-window source $R_0$ identify the corrected classes with these
zero-mesh coefficients modulo $a^3$. Each transformation contributes one
derivative; the stated finite-window defect follows from
\eqref{eq:v17-generator-rate}. Acting with $F$ in the declared gauge
fermion yields the displayed ghost contact with its inherited sign.
\end{proof}

For a normalization check, on a constant diagonal spatial background with
$q_{11}=u,q_{22}=v$ and shear $c^2=e^{ikx^1}$, V12's nonzero $R_2$ gives
\begin{equation}
 R_{3,12}^{(0)}=-\frac{(u+v)(u-v)^2}{64}\,ik e^{ikx^1}.
 \label{eq:v17-shear}
\end{equation}
This follows because $\mathsf B_q$ multiplies that off-diagonal component
by $(u+v)/4$. No constant-background argument licenses setting $R_3$ to
zero.

\section{The enlarged measured family and its two new loop coefficients}
\label{sec:v17-measured}

Retain V14's $S=s(aD)$ and write $H=I-S$, $\sigma(q)=\operatorname{tr}q/4$.
Add precisely
\begin{align}
 \widetilde S^E_{5,a}(q)&=\widehat S^E_{5,a}(Sq)
 +\frac{\chi\delta_g}{4}a^4\sum_{x,\nu}
     \sigma(Sq)_x^3(D_\nu^+Hq)_x^TJ(D_\nu^+Hq)_x,\\
 \widetilde S^{\rm gh}_{5,a}&=-\langle S\bar c,F
                      R_3^{(0),E}(Sq,Sc)\rangle
 +\delta_c a^4\sum_{x,\nu}(D_\nu^+H\bar c)_x^T
                      \sigma(Sq)_x^3(D_\nu^+Hc)_x.
 \label{eq:v17-completed-5}
\end{align}
The two new parameters are fixed in a bounded real interval; zero is
allowed. Upper modes still have V14's nonzero lower interactions.
Equations \eqref{eq:v17-completed-5} define a recorded degree-five
extension on the same carrier, not a claim that its upper interactions
are dictated by diffeomorphism invariance. They leave V14--V16's lower
coefficients unchanged. All five filtered native legs and their subset
momenta lie inside the nonaliased chart.

Let $\mathbb H_a(\Phi)$ be the graded Hessian of this extended off-shell
action and $\mathbb C_a(\Phi)=(\mathbb H_a(\Phi)+\mathbb R_a)^{-1}$.
The reference and its zero-mode prescription are those of V16. Through
one loop define, with $\hbar$ removed from coefficient notation,
\begin{equation}
 \Gamma_{a,1}(\Phi)=\tfrac12\operatorname{STr}
       \log(\mathbb H_a(\Phi)+\mathbb R_a)+\mathcal M_a(q),\qquad
 z_a^q(\Phi)=\tfrac12\operatorname{STr}
           \{\mathbb C_a(\Phi)D^2r_a^q(\Phi)\}.
 \label{eq:v17-one-loop-functionals}
\end{equation}
The additive background-free normalization is understood. The trace
in the second formula leaves the $q$-output index uncontracted. Equivalently
it is the specified Gaussian/Berezin contraction of each component.
Write
\[
 \Gamma_{a,1}(q,0)=\tau_a(q)+\tfrac12\Pi_a(q,q)
                      +\tfrac16\Upsilon_a(q,q,q)+\cdots,
 \qquad z_a^q(q,c)=Z_a c+W_a(q,c)+\cdots.
\]
Here $Z_a=\mathcal Z_a$ is V15's old coefficient. The new quantities are
$\Upsilon_a$ and $W_a$, not a new evaluation of $\Pi_a,Z_a$.

For three even background directions set
$V_I=D_I\mathbb H_a(0)$, with no factorial in a labelled derivative, and
$\mathbb C=\mathbb C_a(0)$. Direct differentiation gives
\begin{equation}
\boxed{\begin{aligned}
 \Upsilon_a(h_1,h_2,h_3)={}&\tfrac12\operatorname{STr}
 \{\mathbb C V_{123}-\mathbb C V_{12}\mathbb C V_3
 -\mathbb C V_{13}\mathbb C V_2-\mathbb C V_{23}\mathbb C V_1\\
 &\hspace{8mm}+\mathbb C V_1\mathbb C V_2\mathbb C V_3
             +\mathbb C V_1\mathbb C V_3\mathbb C V_2\}
             +D^3\mathcal M_a(0)[h_1,h_2,h_3].
\end{aligned}}
\label{eq:v17-three-metric}
\end{equation}
In particular, $V_{123}$ uses the new fifth action and ghost contact.
The two oriented three-vertex traces are not replaced by one trace for
commuting matrices.

For one even metric mark $h$ and one odd ghost mark $c$, put
$P(\Phi)=D^2r_a^q(\Phi)$. Subscripts $h,c,hc$ denote ordered labelled
coefficients. Then
\begin{equation}
\boxed{\begin{aligned}
 W_a(h,c)=\tfrac12\operatorname{STr}\{&
 \mathbb C P_{hc}-\mathbb C V_h\mathbb C P_c
                    -\mathbb C V_c\mathbb C P_h
                    -\mathbb C V_{hc}\mathbb C P_0\\
 &+\mathbb C V_h\mathbb C V_c\mathbb C P_0
                    +\mathbb C V_c\mathbb C V_h\mathbb C P_0\}.
\end{aligned}}
\label{eq:v17-field-source}
\end{equation}
Odd coefficients are ordered with the antifield before the retained
ghost; the supertrace is V16's Berezin convention. $P_{hc}$ contains
the actual $R_3$ contraction. $V_h$ now includes the gravitational cubic,
unlike the older coefficient $Z_a$.

\begin{proposition}[Actual measured coefficients and source contacts]
\label[proposition]{prop:v17-measured}
Equations \eqref{eq:v17-three-metric}--\eqref{eq:v17-field-source} are the
complete specified one-loop coefficients of the extended measured family.
They have no missing hard-linear stationary branch on
\eqref{eq:v17-domain}. Their covariance decompositions retain every
mixed-shell pair and triple. The background quadratic metric Gaussian and
ghost determinant have a common small analytic domain at each finite
regulator, with radius independent of mesh and volume for the fixed
source order and bounded completion parameters.
\end{proposition}
\begin{proof}
Expand $\log(\mathbb H_0+\mathbb R+V)$ through three labelled even
marks; cyclicity gives exactly \eqref{eq:v17-three-metric}. Expand
$(\mathbb H_0+\mathbb R+V)^{-1}$ through $hc$ and multiply by $P$;
this gives the six ordered terms of \eqref{eq:v17-field-source}.
These are coefficient evaluations of the inherited measured calculus,
including the ghost minus sign through the supertrace. A hard propagator
leaving a branch with only the external marks has momentum at most
$4\mu<\Lambda_0$, so it vanishes. Thus no mean-shift tree changes
these coefficients between the specified Legendre representation and the
hard pushforward. An insertion of the order-$\hbar$ density in that
absent tree does not modify $W_a$. The density is already present once in
$\Upsilon_a$, and remains separately present in the measured row.

Substitution of $\mathbb C=\sum_j\mathbb C_j$ in these finite traces
keeps all shell words; assign a word to its largest shell label, without
requiring its labels to agree. This is exact ownership, not independence.
For the analytic Gaussian assertion, use V14's accretive free margins.
The extra fifth-order Hessian is cubic in the background and has relative
form bound $C\|q\|_{\mathrm W}^3$ on the soft source chart. The lower
vertices have their inherited relative bounds. A sufficiently small
common ball makes their sum less than one half. The complex Gaussian on
the fixed contour and the local branch of the ghost determinant are then
analytic. This is a background quadratic reference statement, not an
absolutely convergent unbounded interacting gravity integral.
\end{proof}

\section{Whole-zone bounds for the three-metric and field-dependent source coefficients}
\label{sec:v17-bounds}

The covariance has its fixed lower factor $\nu=1-\theta(p/\Lambda_0)$.
After the constant nonminimal square is completed, every propagating
coefficient in \eqref{eq:v17-three-metric} is a finite one-loop word with
metric or ghost propagators of order $-2$ and vertices of order $2$.
Every word in \eqref{eq:v17-field-source} has one transformation insertion
of order $1$ and the corresponding two additional contracted legs. Its
loop-symbol order is $-1$. The higher metric degree does not increase
these ultraviolet orders.

\begin{lemma}[Differentiated routed integrands]
\label[lemma]{lem:v17-symbols}
Use a smooth partition of one routed loop momentum, with its inner
part supported in $|ap|<1/1024$ and equal to one on $|ap|\le1/2048$,
and put $R=\Lambda_0+|p|$. On this inner region all profiles $s,\zeta$
are one. Together with \eqref{eq:v17-domain}, its soft routed shifts
satisfy the degree-five window of \cref{thm:v17-generator}.
The assembled integrands $f_a$ of $\Upsilon_a^{\rm prop}$ and
$g_a$ of $W_a$ satisfy, for each fixed total derivative order $r$,
\begin{align}
 |\partial^r f_a|+|\partial^r f_0|&\le C_rR^{-r},&
 |\partial^r(f_a-f_0)|&\le C_ra^3R^{3-r},\\
 |\partial^r g_a|+|\partial^r g_0|&\le C_r\chi^{-1}R^{-1-r},&
 |\partial^r(g_a-g_0)|&\le C_r\chi^{-1}a^3R^{2-r}.
 \label{eq:v17-interior-symbols}
\end{align}
Derivatives may be in the external or routed momentum, with appropriate
fixed profile constants. On any complementary upper region
$|ap|\ge\delta>0$ the corresponding bounds are
\begin{equation}
 |\partial^r f_a|\le C_{r,\delta}a^r,\qquad
 |\partial^r g_a|\le C_{r,\delta}\chi^{-1}a^{1+r}.
 \label{eq:v17-upper-symbols}
\end{equation}
All hard-momentum symbols there are smooth and periodic. The low-mode
cutoff removes all zero-denominator singularities.
\end{lemma}
\begin{proof}
On an interior annulus, inverse differentiation and V13's floors give
$C^g=O(\chi^{-1}R^{-2})$, $C^c=O(R^{-2})$, with one inverse power
per derivative. The relative free mismatch is $O((aR)^4)$.
\cref{thm:v17-generator} supplies the corrected fifth vertex with
mismatch $O(\chi a^3R^5)$; V11--V12 supply the same bound for the
lower metric vertices. Interior nonlinear ghost and transformation
vertices are the exact ordinary chart polynomials, of orders two and
one respectively. Replacing one factor at a time in each finite trace
therefore gives \eqref{eq:v17-interior-symbols}. In a source word its
one extra metric covariance leaves the factor $\chi^{-1}$.

For the upper estimates rescale $z=ap$. V14's free inverse is $a^2$
times a smooth periodic matrix, with its stated fixed margins. A completed
metric or ghost vertex is $a^{-2}$ times a smooth periodic symbol, and
any fixed coefficient of the pulled-back source is $a^{-1}$ times
such a symbol. The recursion \eqref{eq:v17-finite-R3} adds only
pointwise products and no new derivative. Each source word consequently
has one overall power $a$, and each momentum derivative adds $a$.
Filtering the total momentum of a product, instead of each factor, does
not invalidate periodicity; every differentiated spectral derivative
still acts on an $S$-filtered expression. Some contractions are local
polynomials in the external momenta; they are retained until subtraction.
There are only finitely many words at these source degrees. This proves
\eqref{eq:v17-upper-symbols} and all mixed derivative versions.
\end{proof}

Set
\begin{equation}
 \Upsilon_a^R=(I-\mathsf T_4)\Upsilon_a,\qquad
 W_a^R=(I-\mathsf T_3)W_a,
 \label{eq:v17-new-subtractions}
\end{equation}
where Taylor degree is the \emph{total} degree of the two independent
external momenta. The jets are stored in the same local action/source
functional as the earlier $\tau,\Pi,Z$ jets. No componentwise parity of
$\Upsilon_a$ or $W_a$ is assumed. V15's stronger oddness statement
applies to $Z_a$, whose graphs did not contain the gravitational cubic;
it is not transferred to $W_a$.

\begin{theorem}[The two new continuum coefficients]
\label{thm:v17-coefficients}
In the thermodynamic convention, and also for the specified finite-volume
coefficient comparison with its own continuum sums, one has
\begin{align}
 |\partial_K^\alpha(\Upsilon_a^R-\Upsilon_0^R)(K)|
       &\le C_\alpha a|K|^{5-r},&&r=|\alpha|\le5,\\
 \|\partial_K^\alpha(W_a^R-W_0^R)(K)\|
       &\le C_\alpha\chi^{-1}a|K|^{4-r},&&r\le4.
 \label{eq:v17-coefficient-rates}
\end{align}
Unit source polarizations are understood. The corresponding continuum
bounds have respective right-hand sides
\[
 C_\alpha\Lambda_0^{-1}|K|^{5-r},\qquad
 C_\alpha\chi^{-1}\Lambda_0^{-1}|K|^{4-r}.
\]
All fixed higher derivatives are controlled on a fixed window by replacing
the displayed momentum powers by the corresponding powers of $\mu$;
constants may depend on $\mu/\Lambda_0$, the profiles and derivative order.
There is no infinite-source-order uniformity.
\end{theorem}
\begin{proof}
Apply Taylor's integral remainder through degree four to $f_a$, and
through degree three to $g_a$, before estimating the routed integrals.
The differentiated interior mismatches are bounded by
$Ca^3|K|^{5-r}R^{-2}$ and
$C\chi^{-1}a^3|K|^{4-r}R^{-2}$ respectively. A four-dimensional
annulus contributes at most $CR^4$, so summing up to a fixed multiple of
$a^{-1}$ gives $Ca^3a^{-2}=Ca$. The omitted continuum tails use
$R^{-5}$ after the corresponding Taylor derivatives and contribute
$C\sum_{R\gtrsim a^{-1}}R^{-1}=Ca$ in dyadic notation.

On the compact upper region the number of momenta per physical volume
is bounded by $Ca^{-4}$. The fifth derivative of $f_a$ and fourth
derivative of $g_a$ are both $O(a^5)$, with $\chi^{-1}$ for the
latter. Thus the entire completed upper contribution is $O(a)$ after
its actual subtraction, including every mixed upper/interior word.
Routed labels differ by bounded soft transfers, so if one label is in
this upper region all inverse estimates used there have a common margin.
Smooth partition derivatives obey the same powers. This accounts for
all loop momenta, including crossings of the periodic faces.

The auxiliary third metric derivative is assigned once. V10's finite
source-degree cutoff estimate, transported through V12's local coordinate
map, gives its subtracted $Ca|K|^5$ bound. The ultralocal part belongs
entirely to its degree-zero jet. The continuum estimates follow by
summing the convergent dyadic tails from $\Lambda_0$. Momentum derivatives
above the subtraction degree annihilate the Taylor polynomial and are
estimated by the same symbols. Finite-volume counting is uniform for
$\Lambda_0L_\nu\ge1$; no Ward assertion for an arbitrary off-grid
interpolation is required in this comparison.
\end{proof}

For later use the actual local jets satisfy the conservative budgets
\begin{align}
 \|D_K^r\Upsilon_a(0)\|&\le C_r
 \begin{cases}a^{r-4},&r<4,\\1+\log(1/(a\Lambda_0)),&r=4,\end{cases}\\
 \|D_K^rW_a(0)\|&\le C_r\chi^{-1}
 \begin{cases}a^{r-3},&r<3,\\1+\log(1/(a\Lambda_0)),&r=3.\end{cases}
 \label{eq:v17-local-budgets}
\end{align}
They follow by summing the unsubtracted annular bounds. In particular,
we do not assume the constant or quadratic local jets of $W_a$ vanish.
The complementary annular increments have bounds
$C\mu^5\Lambda^{-1}$ and $C\chi^{-1}\mu^4\Lambda^{-1}$ and are
summable. This is a fixed-order coefficient flow, not a nonlinear
invariant ball of arbitrary ESM interactions.

\section{The complete second row and its measured contact}
\label{sec:v17-row}

Let $\mathcal V_{3,a}=D^3\widehat S^E_{3,a}(0)$ on the low source panel
and let $H_a^{\rm cl}$ be V16's off-shell free Hessian. Set
$R_2^{12}(c)=D_q^2R_2^{(0),E}(0,c)[h_1,h_2]$.
For the packet $\mathbf u_a=(\tau_a,\Pi_a,\Upsilon_a,Z_a,W_a)$ define
\begin{equation}
\begin{aligned}
 \mathcal L_{2,a}(\mathbf u_a)={}&
 \Upsilon_a[h_1,h_2,R_0c]
 +\Pi_a[h_1,R_1^{(0),E}(h_2,c)]
 +\Pi_a[h_2,R_1^{(0),E}(h_1,c)]\\
 &+\tau_a[R_2^{12}(c)]
 +\langle H_a^{\rm cl}h_1,W_a(h_2,c)\rangle
 +\langle H_a^{\rm cl}h_2,W_a(h_1,c)\rangle\\
 &+\mathcal V_{3,a}[h_1,h_2,Z_ac].
 \label{eq:v17-L2}
\end{aligned}
\end{equation}
The momenta are those dictated by the displayed source products and
conservation. No output of $R_1$ or $W$ is assigned the momentum of one
input alone.

Use $\Phi=(t_1h_1+t_2h_2,0,\zeta c,0)$, with one odd coefficient $\zeta$.
For coefficient extraction take $t_1^2=t_2^2=\zeta^2=0$.
As in V16 put
\[
 V=\mathbb H_a(\Phi)-\mathbb H_{a,0},\qquad
 B_a(\Phi)=D\mathscr S_a(\Phi)[r_a(\Phi)],\qquad
 \mathbb T_a=Dr_a.
\]
Define the truncated resolvent polynomials
\begin{align}
 \mathcal I_{2,a}&=\tfrac12[t_1t_2\zeta]\operatorname{STr}
 \left\{\sum_{j=0}^3(-\mathbb C V)^j\mathbb C\,D^2B_a\right\},\\
 \mathcal X_{2,a}&=[t_1t_2\zeta]\operatorname{STr}
 \left\{\sum_{j=0}^3(-\mathbb C V)^j\Theta\,\mathbb T_a\right\},\\
 \mathcal Y_{2,a}&=[t_1t_2\zeta]D\mathcal M_a(q)[r_a^q(q,c)],\qquad
 d_{2,a}=[t_1t_2\zeta]\operatorname{sdiv}r_a(\Phi).
 \label{eq:v17-insertions}
\end{align}
There are at most three factors $V$ in one trace; every expression is
a one-loop coefficient. $\Theta$ stays on the indicated side of each vertex. The
limit removing the temporary zero-covariance regulator is taken only
after the inherited cancellation $\mathbb C\mathbb R=\Theta$.

\begin{theorem}[Finite one-ghost/two-metric measured identity]
\label{thm:v17-finite-row}
For the extended action and current source prescription,
\begin{equation}
 \boxed{\mathcal L_{2,a}(\mathbf u_a)+d_{2,a}
             =\mathcal I_{2,a}+\mathcal X_{2,a}+\mathcal Y_{2,a}.}
 \label{eq:v17-finite-row}
\end{equation}
This is an exact finite coefficient identity on allowed momenta, and an
identity of the regulated thermodynamic traces. Neither finite
nilpotence nor vanishing of $B_a$ is assumed. Its required primitive
orders are $S_5,R_3$, now supplied; no sixth action coefficient is used.
\end{theorem}
\begin{proof}
Apply the established chain identity \eqref{eq:v16-one-loop-chain} to
the enlarged action. Extract $t_1t_2\zeta$ at $b=\bar c=0$.
The derivative of $\Gamma_1$ produces the first four terms of
\eqref{eq:v17-L2}; the derivative of the classical action acting on
$z_a^q$ gives the next two terms and $\mathcal V_{3,a}Z_a$.
The factors follow by differentiating $\Pi(q,q)/2$, $\Upsilon(q,q,q)/6$
and $S_3(q)=\mathcal V_3(q,q,q)/6$. The other retained field components
vanish on this slice only after off-shell differentiation, as in V16.

The resolvent series through total external degree three is exactly the
finite polynomial in \eqref{eq:v17-insertions}. $D^2B_a$ at this row
uses the classical action identity with four metric fields and one
ghost, together with its nonminimal partners. Those coefficients require
$S_5$ and $R_3$ and no higher terms. The Gaussian change-of-variables
identity is valid for the actual, possibly breaking, input. It therefore
gives \eqref{eq:v17-finite-row} without using the desired continuum
symmetry as a hypothesis. Complete the constant $b$ square before
estimating any apparently nonperiodic entry in its mixed inverse.
\end{proof}

\subsection{All tensor components of the quadratic metric measure contact}

Use the Euclidean chart, $n=1+q_{00}/2$ and $S_q^{\rm sp}=I+q_{\rm sp}/2$.
The letter $S_q^{\rm sp}$ in this paragraph is a $3\times3$ spatial
coframe block, not the Fourier filter. The pure ultralocal auxiliary
density relative to $dg$ is $(\det G)^{-7/2}$. Its pullback is the
inherited density of V12. On the present low Fourier panel its full
flat-divergence minus density-action variation is
\begin{equation}
 \boxed{d_{2,a}-\mathcal Y_{2,a}^{\rm ul}
 =-\frac{7i}{4a^4}(k_c\cdot c)
       \{(h_1)_{00}(h_2)_{00}
                    +\operatorname{tr}((h_1)_{\rm sp}(h_2)_{\rm sp})\}.}
 \label{eq:v17-measure-contact}
\end{equation}
All shift-polarization contributions cancel in this combined contact.
The remaining native Haar and closed-walk variations are still in
$\mathcal Y_{2,a}$.

\begin{proposition}[Derivation of the measured quadratic contact]
\label[proposition]{prop:v17-measure}
Equation \eqref{eq:v17-measure-contact} is an exact coefficient of the
same density and coordinate source used in V16. In particular it is not
obtained by differentiating the $q$-density alone.
\end{proposition}
\begin{proof}
V15's exact change-of-variables formula gives
$\operatorname{div}_{\rho_q}r^q=(\mathscr V_a\log\rho_g)\circ G$.
The ghost super-divergence has no quadratic metric component. All
external subset momenta in this coefficient are inside the filter
plateau. Ordinary differentiation is consequently valid on these
finite polynomial coefficients, so
\[
 \operatorname{div}_{\rho_q}r^q
 =-\tfrac72\sum_x\{c\cdot D\log\det G+2D\cdot c\}
 =\tfrac72\sum_x(D\cdot c)\log\det G.
\]
The second equality is exact periodic coefficient integration by parts;
no product rule on unrestricted high lattice fields is assumed.
Since
\[
 \log\det G=2\log n+2\operatorname{tr}\log(I+q_{\rm sp}/2),
\]
its mixed second derivative at zero is
$-\{(h_1)_{00}(h_2)_{00}+
\operatorname{tr}((h_1)_{\rm sp}(h_2)_{\rm sp})\}/2$.
The Fourier coefficient of the divergence of the plane-wave ghost is
$i k_c\cdot c$. Divide site counting by physical volume, giving
$a^{-4}$. This proves \eqref{eq:v17-measure-contact}.
The equality with $d_{2,a}-\mathcal Y_{2,a}^{\rm ul}$ uses
$\rho_q=e^{-\mathcal M^{\rm ul}}$ in the flat-divergence convention.
The coordinate Jacobian has thus been counted exactly once.
\end{proof}

The contact is a local linear momentum polynomial. It is not evidence
of a nonlocal anomaly. On the pure lapse line its sign is the derivative
of V16's $7ik/(2a^4)$ contact, with the two independent metric marks
and their conserved ghost momentum retained.

\section{A common nonlinear source assignment and its two spill terms}
\label{sec:v17-projector}

The row requires a joint extension of V16's local packet, not a new
independent choice at every coefficient. Define
\begin{equation}
 \mathbb P\mathbf u=(\tau,\mathsf T_4\Pi,\mathsf T_4\Upsilon,
                           \mathsf T_3Z,\mathsf T_3W),\qquad
 \mathbb P_{\rm row}=\mathsf T_5.
 \label{eq:v17-projector}
\end{equation}
A single one-loop counterterm functional implementing it adds, to the
previous lower terms,
\begin{equation}
 -\tfrac16(\mathsf T_4\Upsilon_a)[q,q,q]
                     -\langle q^*,(\mathsf T_3W_a)(q,c)\rangle.
 \label{eq:v17-new-local-action}
\end{equation}
Its coefficients are the actual calculated output jets. The corresponding
local jets of the reference and breaking insertions remain in their
matching ledger.

Write
\[
 E_{2,a}=H_a^{\rm cl}-H_0^{\rm cl},\qquad
 E_{3,a}=\mathcal V_{3,a}-\mathcal V_{3,0}.
\]
The inherited native expansions imply exact Taylor vanishing through
degrees five and four respectively, with
\begin{equation}
 E_{2,a}=O(\chi a^4|K|^6),\qquad E_{3,a}=O(\chi a^3|K|^5).
 \label{eq:v17-tree-errors}
\end{equation}
These are low-source symbols, with all fixed differentiated versions.

\begin{theorem}[Bounded second-row intertwiner]
\label{thm:v17-projector}
For the actual packet, using only V15's established oddness $Z_a(-k)=-Z_a(k)$,
\begin{equation}
 \boxed{\mathsf T_5\mathcal L_{2,a}(\mathbf u_a)
                    =\mathcal L_{2,0}(\mathbb P\mathbf u_a).}
 \label{eq:v17-intertwine}
\end{equation}
The assignment has a mesh- and volume-independent bound in scaled
finite-panel $C^s$ norms for every fixed $s\ge5$. No parity is imposed
on $W_a$. Define
\begin{align}
 \mathcal F_{2,a}={}&
 \langle E_{2,a}h_1,(\mathsf T_3W_a)(h_2,c)\rangle
 +\langle E_{2,a}h_2,(\mathsf T_3W_a)(h_1,c)\rangle\\
 &+E_{3,a}[h_1,h_2,(\mathsf T_3Z_a)c].
 \label{eq:v17-spill}
\end{align}
Then for $r=|\alpha|\le6$,
\begin{equation}
 |\partial_K^\alpha\mathcal F_{2,a}|
              \le C_\alpha a|K|^{6-r}\|h_1\|\|h_2\|\|c\|.
 \label{eq:v17-spill-bound}
\end{equation}
Thus the finite-mesh error multiplied by a divergent local coefficient
is paid explicitly, rather than discarded by consistency.
\end{theorem}
\begin{proof}
The zero-mesh tensors $R_0,R_1,R_2$ are homogeneous of momentum degree
one and $H_0^{\rm cl},\mathcal V_{3,0}$ have degree two.
The substitutions of summed external momenta are linear maps. Taylor
projection therefore gives degrees $(4,4,0,3,3)$ for the five types of
terms in \eqref{eq:v17-L2}. $E_{2,a}$ starts at degree six, so its
product with any $W_a$ has zero Taylor jet through degree five.
$E_{3,a}$ starts at degree five. Its product with $Z_a$ also starts
at degree at least six because $Z_a(0)=0$ by the inherited exact
oddness. This last fact is necessary; one cannot deduce the intertwiner
for an arbitrary constant part of $Z$.

On each coefficient use the norm
$\sum_{|\alpha|\le s}\mu^{|\alpha|}
\sup_{K\in\mathcal K_\mu}\|\partial_K^\alpha f\|$.
A compatible packet norm multiplies $\tau,\Pi,\Upsilon$ by $\mu$
and $Z,W$ by $\chi\mu^2$. Fixed-dimensional Taylor evaluation,
linear routing and multiplication by the displayed homogeneous tensors
are bounded in these norms independently of mesh and periods. This
proves boundedness and \eqref{eq:v17-intertwine}.

For the size of the spill, \eqref{eq:v17-local-budgets} yields
\[
 \|\mathsf T_3W_a\|\le C\chi^{-1}
 \{a^{-3}+a^{-2}|K|+a^{-1}|K|^2+
                 (1+\log(1/(a\Lambda_0)))|K|^3\}.
\]
V15 gives
$\|\mathsf T_3 Z_a\|\le C\chi^{-1}
\{a^{-2}|k_c|+(1+\log(1/(a\Lambda_0)))|k_c|^3\}$.
Multiply by \eqref{eq:v17-tree-errors}. The leading products are
$a^4|K|^6a^{-3}=a|K|^6$ and
$a^3|K|^5a^{-2}|K|=a|K|^6$.
The others have extra bounded factors $a|K|$ or
$(a|K|)^j[1+\log(1/(a\Lambda_0))]$, $j\ge2$.
The domain \eqref{eq:v17-domain} bounds them uniformly. Differentiating
these finite products gives \eqref{eq:v17-spill-bound}, including all
orders up to six. No raw local jet was assumed small.
\end{proof}

Let $\mathcal L_{2,a}^R$ be \eqref{eq:v17-L2} with the old and new
renormalized coefficients substituted and with its local tadpole term
removed from the remainder (and stored). Use $\mathsf T_5$ on every
row insertion. The exact renormalized balance is
\begin{equation}
 \boxed{\mathcal L_{2,a}^R-\mathcal X_{2,a}^R
           =\mathcal I_{2,a}^R+\mathcal Y_{2,a}^R-\mathcal F_{2,a}.}
 \label{eq:v17-renormalized-row}
\end{equation}
Indeed $(I-\mathsf T_5)\mathcal L_{2,a}=
\mathcal L_{2,a}^R+\mathcal F_{2,a}$ by the proof above, and $d_{2,a}$
is an assigned local polynomial. The local ledger is simultaneously
\begin{equation}
 \mathcal L_{2,0}(\mathbb P\mathbf u_a)+d_{2,a}
       =\mathsf T_5(\mathcal I_{2,a}+\mathcal X_{2,a}+\mathcal Y_{2,a}).
 \label{eq:v17-local-ledger}
\end{equation}
Finite physical matching additions must satisfy this relation and the
previous first-row relation. The result is a bounded assignment on two
nested rows, not a full cohomological splitting of the ESM BV complex.

\section{The second modified continuum identity and a cutoff bound}
\label{sec:v17-continuum}

The reference insertion is confined to a fixed lower-scale region for
the same explicit reason as in V16: each nonpolynomial word in
\eqref{eq:v17-insertions} contains $\Theta$, and all routed labels differ
by sums of the three soft transfers. There is no inverse-cutoff factor.
The fifth-action extension is consumed here only through the finite
background derivatives that were specified above.

\begin{lemma}[Second-row reference and zero-mesh identity]
\label[lemma]{lem:v17-continuum}
With thermodynamic normalization,
\begin{align}
 |\partial_K^\alpha\mathcal X_{2,0}|&\le C_\alpha\Lambda_0^{5-r}
                                        \|h_1\|\|h_2\|\|c\|,\\
 |\partial_K^\alpha(\mathcal X_{2,a}^R-\mathcal X_{2,0}^R)|
       &\le C_\alpha a^3\Lambda_0^2|K|^{6-r}
                                        \|h_1\|\|h_2\|\|c\|\quad(r\le6).
 \label{eq:v17-reference-bound}
\end{align}
The continuum coefficients obtained by the preceding cutoff comparisons
satisfy
\begin{equation}
 \boxed{\mathcal L_{2,0}^R=\mathcal X_{2,0}^R.}
 \label{eq:v17-zero-mesh-row}
\end{equation}
The right side is not asserted to vanish at fixed $\Lambda_0$.
\end{lemma}
\begin{proof}
Each reference trace is supported at $|p|\le C\Lambda_0$ unless it
is a polynomial contact. After completing the nonminimal square it has
one transformation derivative and propagators paired with two-derivative
vertices. Four-dimensional integration gives degree five. Ordinary
symbol differentiation gives the first bound. The actual corrected
metric vertices through degree five differ by
$O(a^3\Lambda_0^5)$ on this compact region; the free difference is
relatively $O((a\Lambda_0)^4)$, and the nonlinear source/ghost
polynomials have their ordinary values there. Thus the raw reference
mismatch has degree $a^3\Lambda_0^8$. Taylor's integral remainder
through degree five proves the second bound, with the stated derivative
versions.

For \eqref{eq:v17-zero-mesh-row}, first use the zero-mesh native action
and its ordinary chart transformation, not a guessed continuum Ward
identity. \cref{prop:v17-classical} establishes the missing $q^4c$
Noether coefficient and its two-ghost partner. The zero-mesh nonminimal
action is the same $s\Psi$ as in V16, now including
$-\bar c F R_3$. Consequently all classical-breaking coefficients
entering $[q^2c]\operatorname{STr}(\mathbb C D^2B)$ vanish.

Insert a smooth ordinary ultraviolet covariance cutoff $u(p/M)$.
The finite graded Hessian identity of V16 now applies with the reference
factor $1-\nu u=\Theta+\nu(1-u)$. The first term is exactly the
second-row reference trace; the second is the ultraviolet boundary term.
After subtracting the common degrees $(4,3,5)$ and forming the complete
signed row, six external derivatives bound an annulus at radius $R$
by $C/R$. Differentiated cutoff factors obey the same bound. The
ultraviolet tail is therefore at most $C|K|^6/M$ and tends to zero.
This justifies cyclic traces and routing changes without using undefined
unsubtracted infinite traces. The coordinate and ultralocal measure
contacts are assigned local polynomials; the auxiliary nonlocal remainder
tends to zero by V10 at this fixed source degree. Removing $M$ proves
\eqref{eq:v17-zero-mesh-row}. This proof is independent of any desired
vanishing of the finite classical breaking.
\end{proof}

\begin{theorem}[Second-row cutoff convergence]
\label{thm:v17-row-limit}
Define the normalized residual
\begin{equation}
 \mathcal W_{2,a}^R=\mathcal L_{2,a}^R-\mathcal X_{2,a}^R.
 \label{eq:v17-residual}
\end{equation}
For every fixed $r=|\alpha|\le6$ on \eqref{eq:v17-domain},
\begin{equation}
 \boxed{|\partial_K^\alpha\mathcal W_{2,a}^R|
 \le C_\alpha a|K|^{6-r}\|h_1\|\|h_2\|\|c\|.}
 \label{eq:v17-main}
\end{equation}
Every fixed higher derivative has the corresponding fixed-window bound
$C_r a\mu^{6-r}$. The signed renormalized insertion
$\mathcal I_{2,a}^R$ has the same bound. The action-level row has an
additional factor $\hbar$.
\end{theorem}
\begin{proof}
Subtract the independently established \eqref{eq:v17-zero-mesh-row}.
The new geometric mismatch is $a|K|^5$ multiplied by the one-derivative
$R_0$, giving $a|K|^6$. The two older geometric terms use V14's
$a|k|^5$ mismatch and the one-derivative $R_1$ with its correct linear
routing. The new source mismatch is $\chi^{-1}a|K|^4$ multiplied
by the two-derivative $H_0^{\rm cl}$, again giving $a|K|^6$.
The old source mismatch in $\mathcal V_{3,0}Z_a^R$ is at most
$Ca^2|K|^7$ by V15.

The differences $E_{2,a}W_0^R$ and $E_{3,a}Z_0^R$ are smaller, using
\eqref{eq:v17-tree-errors},
$W_0^R=O(\chi^{-1}\Lambda_0^{-1}|K|^4)$ and V15's
$Z_0^R=O(\chi^{-1}\Lambda_0^{-2}|k_c|^5)$.
The reference difference is bounded by \eqref{eq:v17-reference-bound}.
The fixed domain bounds all extra ratios. Leibniz differentiation gives
the result through $r=6$ and the higher-window versions.

Finally use the exact balance \eqref{eq:v17-renormalized-row}:
$\mathcal I_{2,a}^R=\mathcal W_{2,a}^R-\mathcal Y_{2,a}^R
+\mathcal F_{2,a}$. The spill is bounded in
\cref{thm:v17-projector}. The nonlocal part of $\mathcal Y_{2,a}$
is the variation of V10's complete auxiliary three-source coefficient
and its lower coefficients under the one-derivative local chart
transformation. Its assigned remainder is $O(a|K|^6)$; the
ultralocal part has no remainder. This proves the insertion estimate
\emph{after} signed assembly. No separate ultraviolet smallness of
$D\mathscr S_a[r_a]$ was assumed.
\end{proof}

For a fixed smooth window in the eight external momentum variables, the
thermodynamic kernel, sampled and periodized, obeys for every fixed $b$
\begin{equation}
 \boxed{a^8\sum_{x,y}
 (1+\mu[d_{\rm per}(x,0)+d_{\rm per}(y,0)])^b
 \|\mathcal K^{\mathcal W}_{2,a,L}(x,y)\|
                     \le C_b a\mu^6.}
 \label{eq:v17-kernel}
\end{equation}
Rescale by $\mu$, integrate by parts more than $b+8$ times, and
periodize to prove this inequality from the fixed derivative bounds.
The resulting trilinear functional is bounded with one source in $L^1$
and two in $L^\infty$, without an extensive volume multiplier.

The raw finite-volume interpolation has, separately, the conservative
comparison
\begin{equation}
 |\partial_K^\alpha\mathcal W_{2,a,L}^R|
 \le C_\alpha a|K|^{6-r}
 +C_{N,\alpha}\Lambda_0^{5-r}(\Lambda_0L_{\min})^{-N},\quad N>4,
 \label{eq:v17-volume}
\end{equation}
with unit polarizations. The proof is Poisson summation after $N$
integrable derivatives in the single routed loop momentum, as in V16.
The upper comparison and finite-source counting are uniform on the stated
domain. This does not upgrade an arbitrary finite-torus off-grid Taylor
convention into an exact renormalized Ward convention. Ultraviolet,
thermodynamic, and infrared limits remain distinct.

\section{Checks, remaining scope, and the V17 ledgers}
\label{sec:v17-ledgers}

A useful exact normalization is available in the zero-mesh annular
reference. On the constant conformal background $q=tI_4$, the
hard metric Hessian is independent of $t$ by the native constant-coframe
homogeneity. The ghost Hessian is $(1+t/2)M_0$ and its higher contact
vertices in this repeated direction vanish. For an actual covariance
weight $w$ the third metric coefficient per volume is consequently
\begin{equation}
 \boxed{\Upsilon^{\rm prop}_{0,w}(I,I,I)=-\int_p w(p)^3.}
 \label{eq:v17-conformal}
\end{equation}
Indeed the ghost logarithm is $-4\operatorname{Tr}\log(1+tw/2)$;
three derivatives give the displayed answer. The third ultralocal
auxiliary contribution is separately $9/(2a^4)$, because the
coordinate-correct density is $(1+t/2)^{-18}$. The physical action
coefficients carry $\hbar$. These are local, reference-dependent
matching checks, not a gravitational beta function. Replacing $w^3$ by
a diagonal-shell sum would fail the exact mixed-shell allocation.

The standalone checker \srcid{check_ESM_second_Slavnov_row_V17.py}
separates exact algebra from numerical diagnostics. It checks the fifth
stationary coefficient including cancellation of unknown $A_3,A_4$,
the nonlinear mesh generator, the explicit $R_3$ recursion and shear,
the third log-determinant and source-resolvent coefficients with ordered
matrices, the joint Taylor assignment and its two spill terms, and the
quadratic metric measure contact. It also reports finite numerical
matrix and remainder checks separately. None of these finite tests is
presented as a proof of the uniform Fourier estimates, and no new Lean
formalization is asserted.

\subsection*{Proved}
\addcontentsline{toc}{subsection}{V17: Proved}
The nonlinear native artifact functional
\eqref{eq:v17-generators} generates the first two mesh counterterms at
every prescribed finite coframe degree, with explicit fixed-panel bounds.
Its fifth coefficient, including the fifth connection word and all
stationary returns, is \eqref{eq:v17-S5}. The exact chart recursion gives
$R_3$ and the required classical degree-five identity. The specified
periodic degree-five extension gives the complete three-metric and
field-dependent transformation coefficients
\eqref{eq:v17-three-metric}--\eqref{eq:v17-field-source}, their local
jets, source bounds and cutoff comparisons. The second measured row,
reference polynomial, full quadratic metric measure contact, and bounded
common assignment are explicit. The normalized modified continuum
identity has the rate \eqref{eq:v17-main}, with the two finite-mesh
spill products retained. These are one-loop pure-gravity/ghost and
auxiliary-source results, not a complete interacting ESM theorem.

\subsection*{Inherited}
\addcontentsline{toc}{subsection}{V17: Inherited}
V1--V16 and all their labels are preserved. The actual completed
shared-link action and Cartan inverse, the earlier stationary cubic and
quartic and their mesh artifacts, the affine coframe measure, the
periodic free completion, the exact pulled-back source and V16's
finite integration-by-parts identity retain their original scopes.
In particular $\Pi_a,Z_a$, the lower-scale convention and the absence
of a proven matching to the discontinuous unmodified regulator are
unchanged. General finite coefficient compilation and Schur or
Gaussian identities are not relabelled as new results.

\subsection*{Literature input}
\addcontentsline{toc}{subsection}{V17: Literature input}
Perturbative effective gravity and cutoff-modified BV/Slavnov identities
are established in other frameworks \cite{V13BFR,V15IIM}.
Their use here is methodological comparison. No first existence of
perturbative gravity, universal anomaly-cancellation theorem, or
finite-parameter renormalizability claim is made.

\subsection*{Conditional}
\addcontentsline{toc}{subsection}{V17: Conditional}
Physical local normalizations must satisfy both row matching ledgers.
The two-row projector is not yet a bounded splitting of the combined
diffeomorphism/local-Lorentz/SM complex. The native action modification
and its upper completion require their own microscopic matching when
that comparison is part of the programme. The gravity coefficient
reference is still holomorphic/contour-based; no complete Lorentzian
measure is obtained by multiplying it by the positive Euclidean Higgs
reference. Coupled gauge, chiral, Yukawa, massless and higher-loop
estimates remain additional inputs.

\subsection*{Open}
\addcontentsline{toc}{subsection}{V17: Open}
The next obstruction is not an unknown fifth vertex or another isolated
scalar loop. A substantive next step is a finite-source-degree induction
for the \emph{full} quantum source complex, including the two-ghost and
antifield-squared sectors and their contact subtractions. The native
mesh generator and chart recursion now remove repeated classical
vertex discovery from that induction. Their existence alone does not
bound multi-loop forests or restore the missing chiral line and mixed
ESM blocks. An invariant filtered interacting ESM shell, all-orders
coefficient continuum, lower-scale removal and Lorentzian reconstruction
remain open. The result does not imply asymptotic safety, convergence of
the perturbative series, a nonperturbative gravitational measure, or a
Yang--Mills mass gap.

\begin{center}\small
\begin{tabular}{>{\raggedright\arraybackslash}p{.22\textwidth}>{\raggedright\arraybackslash}p{.69\textwidth}}
\toprule
Variable & Uniformity actually used in V17\\
\midrule
Volume & Rooted kernels have no total-volume factor. Finite-sum comparisons
assume $\Lambda_0L_\nu\ge1$; exact normalized identities use the
thermodynamic convention.\\
Mesh and momenta & Whole-zone completed loop estimates on
\eqref{eq:v17-domain}. Native local artifact jets use
$na\Lambda_*\le1/64$; their constants depend on fixed degree $n$.\\
Loop and sources & One quantum loop; up to three metric sources, the
$q^*qc$ source coefficient, and the $q^2c$ row. The classical local mesh
generator holds at each fixed field degree. No unbounded-source norm.\\
Coupling and profiles & Fixed $\chi>0$ and bounded upper parameters;
constants depend on fixed profile seminorms, inverse margins and
$\mu/\Lambda_0$.\\
Shell count & Renormalized one-loop increments at these source orders
are summable with all mixed labels. No interacting all-loop self-map.\\
Infrared and measure & $\Lambda_0>0$ is retained. Same coordinate-correct
auxiliary measure, flat divergence and fixed complex metric contour.\\
\bottomrule
\end{tabular}
\end{center}


\clearpage
\part*{V18: A two-ghost source obstruction and the minimal-antifield loop module}
\addcontentsline{toc}{part}{V18: Two-ghost audit and minimal-antifield loop module}

\section{The next source sector requires a separate algebra test}
\label{sec:v18-start}

Sections 1--136 and their labels are preserved. The two one-ghost rows
proved in V16--V17 are not a theorem about the square of the finite
transformation source. Indeed, \eqref{eq:v16-full-source} explicitly did
not assume nilpotence. Before extending those rows to the full source
complex, we test that missing property. The test below gives a transverse
obstruction already at the flat background. It concerns the \emph{filtered
source}, not the native cubic action obstruction of V11, and is not a
contradiction of the earlier integration-by-parts identities.

A recorded change of filter placement removes this first obstruction,
without changing the zero-antifield action, reference covariance or
auxiliary measure. It is deliberately not called a complete BRST
restoration: its next field-dependent curvature and its three-ghost
residual are displayed explicitly. In parallel, a special feature of the
actual ghost action allows the entire zero-metric-background
\emph{minimal-antifield} one-loop module to be evaluated. This includes
$c^*cc$, $q^*q^*cc$, and their higher mixed contacts at every fixed
antifield degree. Its finite local assignments and cutoff estimates are
proved without supposing that the new finite source is a symmetry.

The scope of this module is narrower than all source completeness.
Retained ordinary metric, antighost and nonminimal backgrounds are zero;
retained ghosts and the two minimal antifields are arbitrary in a fixed
finite low-momentum panel. Matter fields are not integrated. The
nonminimal field is integrated with exactly V16's Gaussian convention.
The zero modes and a positive lower Wilsonian scale are retained. The
metric contour and the \emph{completed}, rather than unmodified,
regulator of V14 remain in use. No result below removes an infrared
cutoff, supplies a Lorentzian reconstruction, or proves a chiral-line
statement.

\subsection{Source audit and conventions}

The direct inputs are \eqref{eq:v14-completed-free},
\eqref{eq:v14-full-cov}, the cubic ghost action in V14, the exact
coframe map \eqref{eq:v15-chart}, the old source
\eqref{eq:v15-canonical-source}--\eqref{eq:v16-full-source}, and the
continuation ledger of V17. The native mesh-artifact generator,
stationary fifth coefficient, geometric loop coefficients and the two
Slavnov-row theorems are inherited, not recomputed. The supplied programme
\cite{V6Programme} distinguishes finite source compilation from a
bounded splitting of the combined ESM complex. The present construction
respects that distinction. BV/renormalization-group and effective-source
methods are established comparison frameworks \cite{V15IIM}; the new
claims are the following specified filter test, source modification and
minimal-antifield coefficient estimates.

We use Euclidean coefficient indices $0,1,2,3$, $\chi>0$, and V14's
smooth even radial filter $S=s(aD)$. The ordinary finite spectral
multiplier $D_\mu$ is $ip_\mu$. Put
\[
 P=S^2,\qquad G(q)=I+q+\mathsf Q(q,q),\qquad
 A(q)=DG(q)=I+\mathsf B_q.
\]
Here $\mathsf Q(q,q)=\mathsf E(q)^T\mathsf E(q)$ and
$\mathsf B(q,v)=\mathsf E(q)^T\mathsf E(v)+\mathsf E(v)^T\mathsf E(q)$
are precisely V15's maps. Odd parameters are placed on the left. Thus
$\mathfrak s(c\,f)=(\mathfrak sc)f-c\,\mathfrak sf$ for odd $c$.
For two ordinary vector waves $u e^{ipx},v e^{itx}$, the vector-field
bracket has amplitude
\[
 [u_p,v_t]=i\{(u\cdot t)v-(v\cdot p)u\}
 \quad\hbox{at momentum }p+t.
\]
All source amplitudes and traces below keep the physical finite cell
weights. No Hilbert-space positivity is assigned to ghost contractions.

\section{An exact two-ghost obstruction in the old finite source}
\label{sec:v18-obstruction}

In metric coordinates $g=I+h$ the old transformation is
\[
 \mathfrak s_{\rm old}h=R_0S^2c+S\mathcal L_{Sc}Sh,
 \qquad
 \mathfrak s_{\rm old}c=S\{(Sc)^\rho D_\rho(Sc)\}.
\]
Write $c=\vartheta_1u_p+\vartheta_2v_t$, with the plane
waves understood and $\vartheta_1,\vartheta_2$ independent odd
parameters. Take nonaliased momenta $p,t,k=p+t$ and abbreviate
$\alpha=s(ap)$, $\beta=s(at)$, $\gamma=s(ak)$.

\begin{proposition}[The flat two-ghost defect]
\label{prop:v18-old-square}
At $h=0$, the coefficient of $\vartheta_1\vartheta_2$ in the square of
the old transformation is
\begin{equation}
 \boxed{\mathcal N_{\rm old}(u_p,v_t)
 =\gamma\alpha\beta\bigl[
 (\gamma^2-\beta^2)\mathcal L_{u_p}R_0v_t
 +(\alpha^2-\gamma^2)\mathcal L_{v_t}R_0u_p\bigr].}
 \label{eq:v18-old-square}
\end{equation}
The same tensor is the flat $q$-component of its square after pullback
by $G$.
\end{proposition}
\begin{proof}
The ghost variation in the inhomogeneous metric term contributes
$\gamma^3\alpha\beta R_0[u_p,v_t]$. The minus sign from varying the odd
coefficient of $S\mathcal L_{Sc}Sh$ contributes
\[
 -\gamma\alpha\beta^3\mathcal L_{u_p}R_0v_t
 +\gamma\beta\alpha^3\mathcal L_{v_t}R_0u_p.
\]
On this nonaliased two-wave panel the spectral derivative is the ordinary
derivative and
$R_0[u,v]=\mathcal L_uR_0v-\mathcal L_vR_0u$.
Collecting terms proves the formula. The square of an odd derivation
transforms as a vector field under an ordinary coordinate change: the
second derivative of $G$ cancels against the antisymmetry of its two odd
arguments. Since $DG(0)=I$, the flat coefficient is unchanged.
\end{proof}

Choose
\begin{equation}
 p=(0,r,0,0),\qquad t=(0,0,3r,0),\qquad u=v=e_3,
 \qquad z=\frac{(0,3,-1,0)}{\sqrt{10}},\quad H=zz^T.
 \label{eq:v18-waves}
\end{equation}
The underlying vector fields commute, $[u_p,v_t]=0$, and $Hk=0$.
Choose allowed grid momenta for which
$\alpha=1$, $0<\beta<1$, and $\gamma>0$. Such momenta exist on fine
odd grids for, for example, a radial profile strictly decreasing between
$1/64$ and $1/32$: it suffices that
$1/192<ar<1/(32\sqrt{10})$.

\begin{corollary}[A transverse obstruction, not a ghost-bracket normalization]
\label{cor:v18-transverse}
For \eqref{eq:v18-waves},
\begin{equation}
 \boxed{\langle H,\mathcal N_{\rm old}\rangle
       =\frac95\gamma\beta(1-\beta^2)r^2>0.}
 \label{eq:v18-transverse}
\end{equation}
No change of the field-independent ghost bracket alone cancels this
coefficient while retaining the old metric source. At the flat background, adding only regular higher powers of the
\emph{minimal} antifields, while fixing the nonminimal source sector,
also cannot cancel this transverse coefficient. Along a sequence with
$ar$ in a fixed compact subinterval of the displayed range, its size is
bounded below by $c a^{-2}$.
\end{corollary}
\begin{proof}
Both Lie tensors in \eqref{eq:v18-old-square} equal
$-(p t^T+t p^T)$ for these polarizations. Their contraction with $H$ is
$9r^2/5$. This proves \eqref{eq:v18-transverse}. Changing the ghost
bracket changes the flat square only by $R_0S^2v$ for some vector $v$;
its contraction with $H$ is zero because $Hk=0$.
In the coefficient of one metric antifield and two ghosts in the classical
master equation, a term quadratic in the minimal antifields is multiplied
by a minimal zero-antifield Euler row. At $q=\bar c=b=0$ those rows
vanish, also with the ghosts retained. The nonminimal source sector is
fixed in this statement; terms involving its antifields are not covered.
Terms of still higher antifield degree do not contribute to this coefficient. These changes therefore cannot cancel
the displayed transverse term. The last assertion follows from continuity
of the positive profile factor and $r\asymp a^{-1}$. This is not a
no-go theorem for changes of the metric source or of the field carrier.
\end{proof}

This test has no Brillouin-face crossing. All three effective labels are
in a very small central chart, although their physical size grows as the
mesh is removed. At any fixed physical panel the filter is eventually
one, and the defect is exactly zero. Thus the result does not invalidate
the previously proved low-panel continuum comparisons. It prevents
promoting them into all-frequency nilpotence of the old source.

\section{A recorded source repair and its exact remaining curvature}
\label{sec:v18-source}

The zero-antifield action and both covariances are left unchanged. Modify
only the minimal transformation sources by
\begin{equation}
 \boxed{\begin{aligned}
 u&=Pc,\\
 \breve r_a^g(g,c)&=R_0u+P\mathcal L_u(g-I),\\
 \breve r_a^q(q,c)&=A(q)^{-1}\breve r_a^g(G(q),c),\\
 \breve r_a^c(c)&=u^\rho D_\rho u,\\
 \breve r_a^{\bar c}&=ib,\qquad \breve r_a^b=0.
 \end{aligned}}
 \label{eq:v18-repaired-source}
\end{equation}
In particular there is no outer $S$ on $\breve r^c$. The metric
inhomogeneous term $R_0Pc$ is exactly the old one. No $P^{-1}$ occurs.
Products use their finite cyclic labels; the derivative of a product
with a low cyclic label is evaluated at that label, not at a sum of two
independently unwrapped high representatives.

\begin{theorem}[Flat closure and coordinate-measure compatibility]
\label{thm:v18-flat-repair}
On the same small coframe chart as V15, \eqref{eq:v18-repaired-source}
is an analytic source germ with a mesh- and volume-independent radius.
Its coefficients through every fixed field degree are quasilocal. On
every stated finite grid,
\begin{equation}
 \boxed{(\breve{\mathfrak s}_a^2q)|_{q=0}=0,\qquad
 \operatorname{div}_{J_Gdq}\breve r_a^q=0,
 \qquad \operatorname{STr}D_c\breve r_a^c=0.}
 \label{eq:v18-flat-measure}
\end{equation}
For any actual base density $\rho_g$, its pullback gives
\begin{equation}
 \operatorname{div}_{(\rho_g\circ G)J_G}\breve r_a^q
     =(\breve r_a^g\log\rho_g)\circ G.
 \label{eq:v18-density}
\end{equation}
No invariance of the completed action or full nilpotence is asserted.
\end{theorem}
\begin{proof}
At $h=g-I=0$ the square is
\[
 R_0P(u\cdot Du)-P\mathcal L_uR_0u=0.
\]
All modes of $u=Pc$ have $|ap|<1/32$. Products of two such modes
are nonaliased, so the ordinary identity
$\mathcal L_u\mathcal L_u I=\mathcal L_{u\cdot Du}I$
applies exactly. A further $P$ does not alter the equality. Pullback by
$G$ gives the assertion in $q$ coordinates.

The metric field derivative is $P\mathcal L_u$. The transport trace
is a sum of $\operatorname{Tr}(P M_{u^\rho}D_\rho)$.
Cyclicity puts $D_\rho P$ on the diagonal, which is zero by paired
momenta. The index terms are the constant diagonal of $P$ times a sum
of $D_\mu u^\nu$, which is zero. Thus the metric-coordinate divergence
is zero. In the ghost trace, differentiating the first $u$ produces
$P_{xx}D_\mu u^\mu$, whose site sum is zero. Differentiating the second
produces $(D_\rho P)_{xx}u^\rho$, also zero. The overall supertrace
sign is immaterial to either zero. The ordinary coordinate-divergence
identity inherited from V12 gives the last two statements; this proof
verifies its new finite-source premise, rather than assuming it.

The inverse of $I+\mathsf B_q$ has a Neumann series on a common
Wiener ball. In $P\mathcal L_{Pc}h$, a contributing output and ghost
momentum both have size $O(a^{-1})$ in the central chart; their
difference is also away from a principal face. Every derivative
multiplier therefore occurs where it is a smooth periodic symbol.
The same observation applies when $h$ is the product $\mathsf Q(q,q)$:
its total cyclic label, not each factor, is differentiated. Each fixed
coefficient is a finite composition of pointwise products and such
multipliers. It has symbol derivatives of order $j$ bounded by
$C_{N,j}a^{j-1}$. Finite Fourier integration by parts gives the stated
quasilocality, including local delta kernels. The constants may depend
on the fixed degree and filter seminorms.
\end{proof}

The first two metric-dependent coefficients are explicitly
\begin{align}
 \breve R_0(c)&=R_0Pc,\\
 \breve R_1(q,c)&=P\mathcal L_{Pc}q-\mathsf B(q,R_0Pc),\\
 \breve R_2(q,c)&=P\mathcal L_{Pc}\mathsf Q(q,q)
                 -\mathsf B(q,\breve R_1(q,c)).
 \label{eq:v18-source-jets}
\end{align}
The product in the first term of $\breve R_2$ is formed before the
output filter is applied. Its high--high/low-total contacts are retained.
On panels whose relevant labels lie in the plateau these are exactly
V17's ordinary chart coefficients. They are not identical to the old
finite source on high modes, and old finite identities continue to refer
to their old source unless explicitly recalculated.

\subsection{The repaired source is still not a full differential}

Set $b_c=u\cdot Du$. On any nonaliased panel supporting the ordinary
Lie identities, the exact remaining squares in metric coordinates are
\begin{align}
 \breve{\mathfrak s}^{\,2}h
   &=P\{\mathcal L_{(P-I)b_c}h
                +\mathcal L_u(I-P)\mathcal L_u h\},\\
 \breve{\mathfrak s}^{\,2}c
   &=((P-I)b_c)\cdot Du-u\cdot D((P-I)b_c).
 \label{eq:v18-next-curvatures}
\end{align}
To prove these formulas, differentiate \eqref{eq:v18-repaired-source}
with the odd Leibniz sign. Subtract
$\mathcal L_u^2=\mathcal L_{b_c}$ in the first line and the ordinary
three-ghost Jacobi identity in the second. These subtractions prove the
formulas on the declared panel; arbitrary aliased products would retain
their separate spectral product defect.

The first line is generally nonzero even on a constant conformal
$h=t_0I$. For the same two waves as \eqref{eq:v18-waves}, its
transverse coefficient is
\begin{equation}
 \boxed{\langle H,[\breve{\mathfrak s}^{\,2}h]_{12}\rangle
       =\frac95 t_0\gamma^2\beta^2(1-\beta^2)r^2.}
 \label{eq:v18-next-witness}
\end{equation}
Indeed $[u_p,v_t]=0$ and the two ordered double Lie derivatives of $I$
coincide; their intermediate filters have values $\alpha^2$ and
$\beta^2$. This is a new field-dependent source-repair task, not a
term to be declared zero. The constructive gain of
\cref{thm:v18-flat-repair} is exact cancellation of the \emph{first}
flat two-ghost obstruction with regular filters, while preserving the
coordinate measure. It is not closure of the full combined BV complex.

\section{The actual zero-background minimal-antifield integral}
\label{sec:v18-module}

Let $\eta=q^*$ be odd, and let $\sigma=c^*$ be even; $\sigma$ is not
the conformal scalar $\operatorname{tr}q/4$ used in V14. Add exactly
\[
 \langle\eta,\breve r_a^q(q,c)\rangle
             +\langle\sigma,\breve r_a^c(c)\rangle
\]
to the source action. No term quadratic in antifields is added at input.
Such terms will be retained in the measured output. Set the retained
ordinary metric and antighost to zero. Denote the hard metric by $x$ and
hard ghost/antighost by $\beta,\bar\beta$. Covariances are
$C_g=\nu\widetilde K_a^{-1}$ and $C_c=\nu\widetilde M_a^{-1}$,
with the old positive lower scale and the same metric contour.

The following four operands are definitions by \emph{explicit native
vertices}, not free tensors:
\begin{align}
 L_\eta(x)\beta&=\langle\eta,\breve R_1(x,\beta)\rangle,\\
 \tfrac12\langle x,T_{\eta c}x\rangle
   &=\langle\eta,\breve R_2(x,c)\rangle,\\
 \langle\bar\beta,B_cx\rangle
   &=\widetilde S^{\rm gh}_{3,a}(x,c,\bar\beta),\\
 Q_\sigma(\beta)&=\langle\sigma,(P\beta)^\rho
                                            D_\rho(P\beta)\rangle.
 \label{eq:v18-operands}
\end{align}
The first two use \eqref{eq:v18-source-jets}; the third is exactly the
cubic ghost action of V14, not the source's new generator. Its upper
added term vanishes here because the retained ghost obeys $(I-S)c=0$.
For example, at momenta $p$ for $x$ and $k$ for $c$, this operand is
\[
 B_c(p,k)=-s(a(p+k))s(ap)s(ak)
             F_{p+k}R_1^{(0),E}(x_p,c_k).
\]
All transposes in the next definition refer to the even metric slots;
the ordering of external odd coefficients is unchanged.

Define two even metric quadratic forms by
\begin{align}
 \tfrac12\langle x,D_\eta x\rangle
  &=\tfrac12\langle x,T_{\eta c}x\rangle
                       -L_\eta(x)C_cB_cx,\\
 \tfrac12\langle x,D_\sigma x\rangle
  &=Q_\sigma(C_cB_cx),\\
 D&=D_\eta+D_\sigma.
 \label{eq:v18-dressed-D}
\end{align}
In particular $D_\eta$ has one $\eta$ and one retained $c$; $D_\sigma$
has one $\sigma$ and two retained $c$'s. The definition by quadratic
forms fixes the Grassmann ordering without an ambiguous ``graded
transpose'' abbreviation. Polarizing the even $x$ variables gives the
unique symmetric coefficient kernel.

\begin{theorem}[Measured minimal-antifield generator at one loop]
\label{thm:v18-generator}
Fix $N<\infty$. Take any minimal-antifield panel with at most $N$
occurrences of $\eta,\sigma$ and at most $2N$ retained ghosts, each
of momentum at most $\mu$, and assume $\Lambda_0\ge128N\mu$.
At zero ordinary backgrounds the complete one-loop source coefficient,
with its common factor $\hbar$ removed, is
\begin{equation}
 \boxed{\Gamma^{\rm min}_{a,1}[\eta,\sigma,c]
       =\frac12\operatorname{Tr}\log(I+C_gD)
       =\sum_{n\ge1}\frac{(-1)^{n+1}}{2n}
                               \operatorname{Tr}(C_gD)^n.}
 \label{eq:v18-master-generator}
\end{equation}
At each specified exterior-algebra panel this sum is finite. Its
coefficient with $r$ metric antifields and $s$ ghost antifields has
exactly $r+2s$ ghosts and comes from $n=r+s$. No claim that this
functional itself solves a master equation is part of the theorem.
\end{theorem}
\begin{proof}
With a temporary nonsingular covariance on the hard support, the
quadratic fluctuation action is exactly
\begin{equation}
 \tfrac12 x C_g^{-1}x+\bar\beta C_c^{-1}\beta
 +\bar\beta B_cx+L_\eta(x)\beta
 +\tfrac12 xT_{\eta c}x+Q_\sigma(\beta).
 \label{eq:v18-quadratic-action}
\end{equation}
Finite Berezin integration over $\bar\beta$ first is an algebraic
substitution of $\beta=-C_cB_cx$, with determinant
$\det C_c^{-1}$. This can equivalently be proved by translating $\beta$
and extracting the top coefficient of each $\bar\beta$. A polynomial
left in $\beta$ is evaluated at that value, since the other $\beta$
coefficients have no available $\bar\beta$ partners. The linear term
therefore gives $-L_\eta C_cB_cx$ and the quadratic term gives
$Q_\sigma(C_cB_cx)$. The determinant has no minimal-source dependence.
Gaussian integration of $x$ proves \eqref{eq:v18-master-generator}.
Remove the temporary nonsingularity parameter on the integrated support;
all formulas use $C_g,C_c$, not an inverse on retained zero modes.

It remains to justify using the quadratic fluctuation action for the
\emph{complete} stated perturbative coefficient. Every hard-linear
vertex at this external order has momentum a sum of at most $3N$
soft labels. It is below the hard covariance support and cannot occur
in a nonzero graph. A connected one-loop graph has the same number of
internal lines and vertices. If every vertex has at least two hard
legs, this equality forces every vertex to have exactly two. Thus no
unlisted higher-hard-degree action or source vertex contributes. The
native cubic ghost vertex and precisely the source vertices in
\eqref{eq:v18-operands} are all the quadratic possibilities at the
specified backgrounds. Integrating the nonminimal Gaussian only gives
the already prescribed $q$ covariance. The auxiliary insertion has
loop degree one and no antifields; contracting it into this source
coefficient would either add a loop or require the forbidden hard-linear
return. It does not appear a second time in this module.

All coefficients of $D$ contain external ghosts. On a finite panel they
are nilpotent, while their even coefficients commute in the Grassmann
sense. The determinant formula is therefore coefficientwise exact with
no analytic smallness hypothesis on antifields. Counting the two types
of $D$ gives the last assertion.
\end{proof}

This calculation uses the general Gaussian/Berezin principle already
proved earlier, but identifies its new operands and the exact
source-sector coverage. It avoids constructing higher gravitational
vertices just to compute the zero-background antifield contacts.
It is a one-loop statement at every fixed source degree, not an
all-loop statement.

\subsection{The first new coefficients and their signs}

The ghost-antifield coefficient is
\begin{equation}
 [\Gamma^{\rm min}_{a,1}]_{\sigma cc}
                   =\tfrac12\operatorname{Tr}(C_gD_\sigma).
 \label{eq:v18-ghost-source}
\end{equation}
It includes two internal ghost propagators inside $D_\sigma$ and the
metric propagator closing the loop. Setting it to zero because there
is no elementary $\bar c$ in the ghost transformation would miss
these native mixed ghost-action vertices.

For the metric-antifield square choose independent odd variables
$\xi_1,\xi_2,\vartheta_1,\vartheta_2$ and write
$D_\eta=\sum_{i,j=1}^2\xi_i\vartheta_jD_{ij}$.
With that fixed ordering,
\begin{equation}
 \boxed{[\Gamma^{\rm min}_{a,1}]_{\xi_1\xi_2\vartheta_1\vartheta_2}
  =\frac12\operatorname{Tr}
       (C_gD_{11}C_gD_{22}-C_gD_{12}C_gD_{21}).}
 \label{eq:v18-antifield-square}
\end{equation}
This follows from the $-\operatorname{Tr}(C_gD)^2/4$ term:
$\xi_1\vartheta_1\xi_2\vartheta_2=-\xi_1\xi_2\vartheta_1\vartheta_2$,
whereas the crossed pairing has the opposite sign. Each $D_{ij}$
contains both the transformation seagull and its ghost-mediated return.
It is not legitimate to replace the displayed formula by only a
seagull--seagull contraction.

These are identities of full finite kernels, including momentum routing.
They do not infer a sign or nonvanishing of a full integrated
renormalized coefficient from a single loop-momentum summand. For an
independent check of the native tensor operands, in the unfiltered
zero-mesh reference with $\chi=1$, choose
\[
 p=(2,3,1,4),\quad k=(0,1/256,0,0),\quad
 \eta_1=E_{00},\ \eta_2=E_{01}+E_{10},\quad
 c_1=e_0e^{ikx},\ c_2=e_1e^{-ikx},
\]
with zero antifield momenta. Take $\mu=1/256$ and $\Lambda_0=1$;
the displayed hard modes are in the region where the lower covariance
weight is one. The fixed routing of
\eqref{eq:v18-antifield-square} has rational integrand
\[
 \frac{8656370045527516036831382943360737}
 {20792664351909365772859143135155076465000}.
\]
This is a nonzero \emph{summand test},
not a value or sign theorem for its integrated amplitude.

\section{Whole-zone power counting for the specified source module}
\label{sec:v18-symbols}

Fix a panel $I$ with $r$ copies of $\eta$, $s$ copies of $\sigma$,
$n=r+s\ge1$, and exactly $r+2s$ copies of $c$. Its number of external
marks is $m_I=2r+3s$. Let $K$ denote $m_I-1$ independent four-momenta,
with the last fixed by conservation; $|K|$ may be their maximum norm,
with equivalent norms changing only fixed-panel constants. Let
$\mathfrak n_I$ be the product of all polarization norms. Write
$R_p=\Lambda_0+|p|$ and use the normalized loop sum $\int_p^L$.

\begin{lemma}[Native dressed-vertex bounds]
\label{lem:v18-D-bounds}
On the full completed momentum torus and the stated soft panel,
all prescribed fixed hard and external derivatives satisfy
\begin{align}
 \|\partial_p^\beta\partial_K^\alpha C_g\|
  &\le C_{\alpha,\beta}\chi^{-1}R_p^{-2-|\alpha|-|\beta|},\\
 \|\partial_p^\beta\partial_K^\alpha C_c\|
  &\le C_{\alpha,\beta}R_p^{-2-|\alpha|-|\beta|},\\
 \|\partial_p^\beta\partial_K^\alpha D_\eta\|
       +\|\partial_p^\beta\partial_K^\alpha D_\sigma\|
  &\le C_{I,\alpha,\beta}R_p^{1-|\alpha|-|\beta|}
                                                 \mathfrak n_D.
 \label{eq:v18-D-bounds}
\end{align}
In an interior region where every routed label is in the plateau,
\begin{align}
 \|\partial^{\alpha,\beta}(C_{g,a}-C_{g,0})\|
       &\le C\chi^{-1}a^4R_p^{2-|\alpha|-|\beta|},\\
 \|\partial^{\alpha,\beta}(C_{c,a}-C_{c,0})\|
       &\le Ca^4R_p^{2-|\alpha|-|\beta|},\\
 \|\partial^{\alpha,\beta}(D_a-D_0)\|
       &\le C_Ia^4R_p^{5-|\alpha|-|\beta|}\mathfrak n_D.
 \label{eq:v18-D-difference}
\end{align}
The completed kernels and these source vertices are smooth periodic in
the routed hard momentum. The constants are independent of periods and
mesh, for fixed panel, profile seminorms and $\chi$.
\end{lemma}
\begin{proof}
The first two estimates are V14's actual inverse bounds, combined with
smoothness of $\nu$ and its vanishing below the fixed lower scale.
Soft shifts change $R_p$ by bounded factors. In the central region,
V13's quadratic and ghost repairs give
$K_a-K_0=O(\chi a^4R_p^6)$ and $M_a-M_0=O(a^4R_p^6)$ with
differentiated bounds. The inverse identity gives the first two
comparisons, including the common lower cutoff.

Each tensor in $\breve R_1,\breve R_2$ has one derivative. Its hard
order is therefore at most one. The apparent unfiltered differentiated
$\mathsf Q(x,x)$ has a total soft momentum at this coefficient; it
introduces no derivative of an individual hard factor at a principal
face. Every other differentiated hard momentum is cut off before a
face. Thus $T,L$, and $Q_\sigma$ have order one and differentiated
bounds lowering that order. The actual ghost operand $B_c=-SFR_1$
has order two. Consequently the return $LC_cB_c$ has order one,
and $Q_\sigma(C_cB_cx)$ has order one as a quadratic form in $x$:
its two factors $C_cB_c$ have order zero and $Q_\sigma$ has order
one. This proves \eqref{eq:v18-D-bounds} term by term, before any
cancellation is needed.

On the interior plateau the source and $B_c$ are exactly their
ordinary zero-mesh polynomials. The only differences inside $D$ are
then the one or two ghost covariances. Replacing one at a time gives
$a^4R_p^5$; for example $L(C_{c,a}-C_{c,0})B_c$ has orders
$1+2+2=5$. The other return is identical by the product rule. In the
transition and upper regions rescale $z=ap$: propagators have scale
$a^2$, the primitive source has scale $a^{-1}$ and the ghost-action
vertex scale $a^{-2}$. Every differentiated smooth profile contributes
$a$. These are exactly the remaining whole-zone bounds. They apply
also to high--high source contacts formed at low total momentum.
\end{proof}

Expanding \eqref{eq:v18-master-generator} with its matrix order and
external Grassmann order retained defines a finite sum of routed
integrands $f_{a,I}(p,K)$. There is no cancellation assumption in the
following consequence:
\begin{equation}
 \boxed{\begin{aligned}
 |\partial_p^\beta\partial_K^\alpha f_{a,I}|
       +|\partial_p^\beta\partial_K^\alpha f_{0,I}|
   &\le C_{I,\alpha,\beta}\chi^{-n}\mathfrak n_I
                       R_p^{-n-|\alpha|-|\beta|},\\
 |\partial_p^\beta\partial_K^\alpha(f_{a,I}-f_{0,I})|
   &\le C_{I,\alpha,\beta}\chi^{-n}\mathfrak n_I a^4
                       R_p^{4-n-|\alpha|-|\beta|}.
 \end{aligned}}
 \label{eq:v18-loop-bounds}
\end{equation}
The second line is for the plateau region. The first is whole-zone
with periodic representatives for the lattice expression. Indeed each
$C_gD$ has order $-1$, and replacing one factor introduces four
powers of $aR_p$. The number of replacements and label orders is
finite at fixed $I$. A loop with $n$ minimal-antifield insertions thus
has superficial degree $4-n$ in this module. Counting ghost number
alone would not give this result; the two ghost propagators inside the
$c^*$ insertion are essential to its degree.

\section{Finite local assignment and an arbitrary finite antifield-order cutoff theorem}
\label{sec:v18-cutoff}

For $1\le n\le4$ set $d_I=4-n$ and store the \emph{actual} Taylor
polynomial $\mathsf T_{d_I}\gamma_{a,L,I}$ at zero external momentum,
where $\gamma_{a,L,I}=\int_p^Lf_{a,I}$. Define
\begin{equation}
 \gamma^R_{a,L,I}=(I-\mathsf T_{d_I})\gamma_{a,L,I}
 \quad(n\le4),\qquad
 \gamma^R_{a,L,I}=\gamma_{a,L,I}\quad(n>4).
 \label{eq:v18-subtraction}
\end{equation}
Use the same integrand subtraction to define the absolutely convergent
continuum sum over all periodic momenta, or its thermodynamic integral.
For every prescribed $N$ the component source bank is finite, and
increasing $N$ leaves old coefficients and their assignments unchanged.
It is a counterterm bank, not a deletion of source-only contacts. Its
entries are not independent physical gravitational couplings. In
particular, for $q^*q^*cc$ the Taylor degree is two, not the degree
three used for $q^*c$.

\begin{theorem}[Cutoff convergence of the minimal-antifield module]
\label{thm:v18-continuum}
Fix $N<\infty$, $\chi>0$, bounded completion parameters and profile
seminorms, and assume
\begin{equation}
 \Lambda_0\ge128N\mu>0,\qquad
 a\Lambda_0\le(8192N)^{-1},\qquad
 \Lambda_0L_\nu\ge1.
 \label{eq:v18-domain}
\end{equation}
All external source labels have size at most $\mu$ and $n\le N$.
For $1\le n\le4$, put $M_I=5-n$. For every multi-index with
$j=|\alpha|\le M_I$,
\begin{equation}
 \boxed{|\partial_K^\alpha
   (\gamma^R_{a,L,I}-\gamma^R_{0,L,I})(K)|
 \le C_{I,\alpha}\chi^{-n}\mathfrak n_I
                             a|K|^{M_I-j}.}
 \label{eq:v18-main}
\end{equation}
The same estimate holds thermodynamically. If $n>4$, no subtraction is
needed and, for every fixed $j$, the bound is
\begin{equation}
 |\partial_K^\alpha(\gamma_{a,L,I}-\gamma_{0,L,I})|
 \le C_{I,\alpha}\chi^{-n}\mathfrak n_I
       \Lambda_0^{4-n-j}\epsilon_{n+j}(a\Lambda_0),
 \label{eq:v18-convergent-orders}
\end{equation}
where for integers $t\ge5$,
\[
 \epsilon_t(x)=
 \begin{cases}
 x^{t-4},&5\le t\le7,\\
 x^4[1+\log(1/x)],&t=8,\\
 x^4,&t\ge9.
 \end{cases}
\]
For $n\le4$ and $j>M_I$, the same expression as
\eqref{eq:v18-convergent-orders} bounds derivatives of the subtracted
difference. These are coefficientwise one-loop statements; action
coefficients carry an additional $\hbar$.
\end{theorem}
\begin{proof}
For $n\le4$ use integral Taylor's formula through degree $d_I$ on
the fully ordered integrand. All required remainder derivatives have
order $M_I=5-n$. Equation \eqref{eq:v18-loop-bounds} gives
\[
 |\partial_K^\alpha(I-\mathsf T_{d_I})f_{0,I}|
       \le C\chi^{-n}\mathfrak n_I|K|^{M_I-j}R_p^{-5},
\]
and, in the interior plateau,
\[
 |\partial_K^\alpha(I-\mathsf T_{d_I})(f_{a,I}-f_{0,I})|
       \le C\chi^{-n}\mathfrak n_I a^4
                         |K|^{M_I-j}R_p^{-1}.
\]
Introduce a smooth partition between a fixed small dimensionless
interior region and its complement. The separating radius is
$c_N/a$, with $c_N>0$ independent of $a,L$. Below it the last bound
integrates, or sums per volume, to
$C a^4\int^{c_N/a} R^2\,\dd R\le Ca$.
The transition derivatives of this partition have the same scaling.
Above it, the subtracted continuum tail costs
$C\int_{c_N/a}^{\infty}R^{-2}\,\dd R\le Ca$.
For the lattice upper region, $M_I$ external derivatives of the
integrand have size $a^{n+M_I}=a^5$; there are at most $Ca^{-4}$
momenta per physical volume. Its Taylor remainder also costs $Ca$.
These three bounds prove \eqref{eq:v18-main} with every indicated
external derivative. No parity improvement is used.

For a finite periodic box the needed counting estimate is
\[
 V_L^{-1}\#\{p:|p|\le R\}
        \le C\prod_\nu(R+L_\nu^{-1})\le C'R^4,
        \qquad R\ge\Lambda_0.
\]
Dyadic annuli then give exactly the same three bounds. The lower cutoff
extends the covariance smoothly by zero through $p=0$, so no omitted
zero-mode inverse appears.

For $n>4$, or for $j>M_I$, no Taylor remainder is needed after
$j$ derivatives. The interior mismatch is bounded by
$a^4\int_{\Lambda_0}^{c_N/a}R^{7-n-j}\,\dd R$.
It is $O(a^{n+j-4})$ for $5\le n+j\le7$, logarithmic at $n+j=8$,
and $O(a^4\Lambda_0^{8-n-j})$ above eight. The lattice upper term
and continuum tail are $O(a^{n+j-4})$. This proves the remaining
formulas and also gives all higher derivatives needed for source kernels.
\end{proof}

The corresponding local coefficient budgets, for $j\le d_I$, are
\begin{equation}
 |D_K^j\gamma_{a,L,I}(0)|\le C_I\chi^{-n}\mathfrak n_I
 \begin{cases}
 a^{n+j-4},&n+j<4,\\
 1+\log(1/(a\Lambda_0)),&n+j=4.
 \end{cases}
 \label{eq:v18-local-budgets}
\end{equation}
They follow by summing $R_p^{-n-j}$. Some coefficients vanish by
additional symmetries or Grassmann antisymmetry, but these bounds do not
assume such a cancellation. A finite list of Taylor assignments exists
at each $N$; a bounded projection \emph{commuting with the entire BV
operator} is not proved by this power count.

\subsection{Rooted source norms and mixed-shell ownership}

For $n\le4$, multiply the difference in \eqref{eq:v18-main} by a fixed
smooth external window, sampled and periodized. With $m_I-1$ relative
four-dimensional coordinates, every fixed moment $b$ obeys
\begin{equation}
 \boxed{a^{4(m_I-1)}\!\sum_{x_1,\ldots,x_{m_I-1}}
 (1+\mu\!\sum_\ell d_{\rm per}(x_\ell,0))^b
 \|\mathcal K_{a,L,I}^{\Delta}(x_1,\ldots,x_{m_I-1})\|
 \le C_{I,b}\chi^{-n}a\mu^{5-n}.}
 \label{eq:v18-kernel}
\end{equation}
The proof rescales each external momentum by $\mu$, integrates by
parts more than $b+4(m_I-1)$ times, and uses the higher-derivative
bounds just proved. Periodization can only decrease the distance
weight of a lift. This gives a multilinear bound with one source in
$L^1$ and the others in $L^\infty$, with their fixed component norms,
without a total-volume factor. Formula \eqref{eq:v18-convergent-orders}
gives the corresponding bounds for $n>4$.

Each $D$ includes its actual one or two ghost covariance factors.
Expand those factors and every $C_g$ into cumulative innovations before
assigning a shell. Give a word to the largest shell index among
\emph{all} its internal lines. This partitions the exact trace once;
it never replaces $C_gD(C_c)C_gD(C_c)$ by equal-shell products.
Equivalently take differences of the cumulative generator
\eqref{eq:v18-master-generator}, with all covariances changed together.
Since adjacent routed momenta differ by at most $3N\mu$, a word with
maximum internal scale $\Lambda\ge\Lambda_0$ lies in a comparable
loop annulus. The subtracted integrand estimates prove
\[
 \|\Delta_\Lambda\gamma_I^R\|_{\mathrm{src},\mu}
 \le C_I\chi^{-n}
 \begin{cases}
 \mu^{5-n}\Lambda^{-1},&n\le4,\\
 \Lambda^{4-n},&n>4.
 \end{cases}
\]
Both are summable over geometric scales. This proves shell-count
uniformity for the specified one-loop module; it is not an interacting
all-loop self-map.

\section{Matching of source prescriptions and what is not a master theorem}
\label{sec:v18-matching}

The change \eqref{eq:v18-repaired-source} is an explicit antifield
source change. It leaves every zero-antifield coefficient of V14--V17
unchanged, but changes their finite high-frequency source insertions.
One must not retain an old finite reference or breaking insertion while
silently substituting the new source on the other side of its identity.
Finite integration by parts applies to the new source with all operands
changed together, irrespective of nilpotence.

\begin{proposition}[Matched minimal-source comparison]
\label{prop:v18-matching}
Apply \eqref{eq:v18-master-generator} to the old V15--V16 minimal
source as well, using its actual $L,T,Q_\sigma$ and the same ghost
action. Both prescriptions have the same zero-mesh minimal-antifield
coefficients after their respective actual local Taylor assignments.
Their difference satisfies the sum of the bounds of
\cref{thm:v18-continuum}. For $q^*c$ the stronger comparison inherited
from V15 remains valid for the new prescription:
\[
 \|\partial_k^\alpha(\breve Z_a^R-Z_a^R)\|
 \le C_\alpha\chi^{-1}a^2|k|^{5-|\alpha|},\qquad |\alpha|\le5.
\]
For $q^*qc$ the same primitive estimates give the conservative
$C_\alpha\chi^{-1}a|K|^{4-|\alpha|}$ comparison through four
external derivatives, after the actual degree-three assignments.
\end{proposition}
\begin{proof}
The Berezin substitution is unchanged for the old source: its ghost
source is also quadratic, independent of $q$, and has no antighost.
Its vertices obey the same one-derivative bounds. Old and new sources
coincide on an interior plateau where all their needed labels lie.
Their difference is therefore in the smooth upper/transition
contribution, whose Taylor remainder was bounded above, while both
interior limits are identical. Applying the triangle inequality proves
the first assertion, including the locally integrable $n>4$ cases.

For one $q^*c$ pair, simultaneous reversal of routed and external
momenta makes the complete integrand odd in the external momentum
after integration: each term has one odd source derivative, an even
ghost-action vertex and even covariances. Both profiles are even.
The degree-four Taylor coefficient is therefore zero. Five rather
than four external derivatives of its smooth upper difference have
size $a^6$ before the four-dimensional count, giving $a^2|k|^5$.
The same property holds for the continuum tail. The field-dependent
coefficient has no required odd parity. Differentiating the new source
once in a retained metric introduces only pointwise coframe maps and
one-derivative source vertices; the extra action Hessian has order two.
The proof of V17's fixed $q^*qc$ comparison consequently has the same
orders, with upper differences again beginning outside the common
plateau. Four external derivatives and the four-dimensional count
give its $a|K|^4$ remainder. This transfers a \emph{coefficient
comparison}, not the old finite Slavnov insertion unchanged.
\end{proof}

One can test the next classical tasks without calculating another
metric determinant. The two formulas \eqref{eq:v18-next-curvatures}
are the required field-dependent two-ghost and three-ghost source
curvatures. A full quantum algebra row also contains the now-defined
$c^*cc$ insertion \eqref{eq:v18-ghost-source}, the old $Z,W$ with
the appropriate source matching, and the multi-antifield contacts
\eqref{eq:v18-antifield-square}. At nonzero metric background the
Euler covector multiplies the quadratic-antifield kernel; omitting
that kernel changes the source complex. The counterterm projector
\eqref{eq:v18-subtraction} has not been proved to intertwine all
these terms. Nor has the nonlinear finite source curvature been
shown to have a summable renormalized insertion in every such row.
These are the next algebra and norm questions, rather than another
unknown fifth gravitational vertex.

The previous row theorems remain at their stated scope. The new
transverse test does not contradict their fixed-soft-source continuum
limits: all filter values in that test vary only when the ghost
momenta grow with the ultraviolet cutoff. The programme needs control
of just such internal source configurations for a full master theorem.
A fixed-panel loop continuum theorem and a nilpotent finite source are
logically independent assertions.

\section{Verification and closing ledgers for V18}
\label{sec:v18-ledgers}

The self-contained verifier
\srcid{check_ESM_minimal_antifields_V18.py}
checks the old two-ghost tensor, the transverse rational coefficient,
flat closure of the rearranged source, its next conformal curvature,
and a finite coordinate-measure trace test. It also performs an exact
Berezin elimination with both minimal antifields and checks the
quadratic-antifield ordering sign, the ghost-antifield contact, the
native rational summand displayed above, and the cutoff-degree table.
Finite numerical checks, when reported, are diagnostics and are
identified separately. Neither the script nor its exact finite
identities is a proof assistant verification of the continuum estimates.

\subsection*{Proved}
\addcontentsline{toc}{subsection}{V18: Proved}
The old filtered source has the transverse flat two-ghost obstruction
\eqref{eq:v18-transverse}, which cannot be corrected by a ghost
bracket change alone on its fixed flat stationary carrier. The explicit
source modification \eqref{eq:v18-repaired-source} cancels that flat
component exactly and preserves the verified coordinate-Jacobian
measure identity. Its remaining curvatures are explicit, not zero.
The actual ghost-action vertex, the first two metric-source vertices
and the ghost-bracket source compile the zero-background
minimal-antifield one-loop module by \eqref{eq:v18-master-generator}.
This gives every prescribed finite $\eta^r\sigma^sc^{r+2s}$ coefficient,
including both missing lowest contacts. Their degree-dependent local
assignments, differentiated cutoff bounds, rooted kernels, mixed-shell
summability and matched-source limits are proved on
\eqref{eq:v18-domain}.

\subsection*{Inherited}
\addcontentsline{toc}{subsection}{V18: Inherited}
V1--V17 and their labels are unchanged. The finite carrier, coframe
chart, complex contour, completed inverse bounds, lower Wilsonian
reference, native cubic ghost action, auxiliary connection measure,
and earlier geometric and linear-antifield cutoff coefficients are
used at their original scopes. General Gaussian/Berezin elimination,
BV notation, coordinate divergence transport and maximum-shell
allocation are inherited tools. The present work supplies their new
minimal-antifield operands and the missing source-algebra audit.

\subsection*{Literature input}
\addcontentsline{toc}{subsection}{V18: Literature input}
Cutoff-modified BV identities and higher effective antifield dependence
are established subjects, with \cite{V15IIM} providing the methodological
comparison already used in this lineage. No first construction of
perturbative gravity, a universal anomaly theorem, or a new general
Berezin identity is claimed. The displayed old-filter obstruction,
source placement and native graded trace are derived here.

\subsection*{Conditional}
\addcontentsline{toc}{subsection}{V18: Conditional}
A complete ESM use still requires a common source prescription with
controlled nonlinear curvature, a bounded compatible local assignment
for the full quantum algebra, and the actual mixed internal-gauge and
fermionic/chiral sectors. Matching of the completed regulator to the
unmodified microscopic action is not obtained. The new minimal-source
comparison does not by itself transfer an entire old finite Slavnov
identity: its breaking and reference operands must be recomputed.
No ordinary-background induction or arbitrary source-order uniform
radius is proved by the zero-background module.

\subsection*{Open}
\addcontentsline{toc}{subsection}{V18: Open}
The immediate upstream target is the field-dependent two-ghost and
three-ghost defect in \eqref{eq:v18-next-curvatures}, including its
renormalized insertions and compatible higher source terms. The
quadratic-antifield and $c^*cc$ loop operands are now available and
should be consumed, not omitted. A full nonlinear quantum master
identity, a combined diffeomorphism/local-Lorentz/SM projection, an
invariant interacting ESM class and arbitrary-loop EFT continuum
remain open. Lower-scale removal, a full Lorentzian measure, global
chiral-line transport and a nonperturbative gravitational continuum
are not conclusions of this installment.

\begin{center}\small
\begin{tabular}{>{\raggedright\arraybackslash}p{.22\textwidth}>{\raggedright\arraybackslash}p{.69\textwidth}}
\toprule
Variable & Uniformity and limitation in V18\\
\midrule
Volume & Rooted bounds have no extensive factor. Finite cutoff
comparisons assume $\Lambda_0L_\nu\ge1$ and the stated Fourier
interpolation; no exact finite-torus renormalized BV identity is claimed.\\
Mesh and momenta & Whole completed zone in \eqref{eq:v18-domain}.
Flat algebra repair is exact on each grid; the obstruction uses
cutoff-size but nonaliased ghosts. Fixed source momenta remain low.\\
Loop and sources & One loop; any prescribed finite minimal-antifield
degree at zero retained metric/antighost/nonminimal backgrounds.
Constants depend on the degree, not uniformly on its unbounded limit.\\
Couplings & $\chi>0$ fixed; bounds retain $\chi^{-n}$ and fixed
profile/inverse margins. No matter-coupling or gravitational UV flow is
inferred.\\
Shell count & Summable one-loop complement with all internal metric
and ghost shell labels retained. No all-loop invariant state space.\\
Measure and infrared & Old coordinate-correct auxiliary measure and
complex reference retained. $\Lambda_0>0$ is not removed. The new
source is not a complete finite nilpotent symmetry.\\
\bottomrule
\end{tabular}
\end{center}


\clearpage
\part*{V19: The remaining source curvature is ultraviolet-local at one loop}
\addcontentsline{toc}{part}{V19: Measured source-curvature insertions and local matching}

\section{Consume the source curvature rather than change the source again}
\label{sec:v19-start}

Sections 1--144 and their labels are preserved. The outstanding operand
in V18 is not another minimal-antifield determinant: it is the insertion
of the two explicit components of the square of the same transformation.
The present installment keeps \emph{exactly}
\eqref{eq:v18-repaired-source}, the completed zero-antifield action,
its propagators, the coordinate chart and the auxiliary measure. No
further filter rearrangement, action replacement, or nilpotence assumption
is made.

The result is a locality theorem for those actual curvature insertions.
Their marked one-loop integrals are supported in a fixed dimensionless
ultraviolet annulus. After their own finite local assignments, their
nonlocal parts tend to zero at every prescribed finite minimal-antifield
degree. The theorem includes zero or one retained metric insertion, so it
also tests a first ordinary-background derivative of the V18 module.
It does not identify an insertion subtraction with an action-level
BV primitive. The latter is a separate finite local consistency problem,
spelled out at the end.

The completed regulator and complex metric contour are those of V14,
not the discontinuous unmodified ultraviolet regulator. The positive
lower scale is retained. Matter fields, chiral transport and unrestricted
ordinary metric backgrounds are not added by this calculation. The
structural ESM programme \cite{V6Programme} and its preceding classical
target remain the surrounding framework; neither the classical
Einstein theorem nor internal anomaly cancellation is used as a quantum
symmetry assumption. The general language of modified BV identities is
inherited from \cite{V15IIM}. The new input is the \emph{marked native
curvature vertex and its proved annular support}, not a general claim
that all regulator breakings are local.

\subsection{Fixed conventions and the domain}

Use the notation of V18. In particular
\[
 P=S^2,\quad S=s(aD),\quad G(q)=I+q+\mathsf Q(q,q),
 \quad A(q)=DG(q)=I+\mathsf B_q,
\]
with the ordinary finite spectral derivative $D_\mu\leftrightarrow ip_\mu$.
Let the fixed smooth even radial profile obey
\begin{equation}
 s(z)=1\quad(|z|\le r_-),\qquad s(z)=0\quad(|z|\ge r_+),
 \qquad r_-=1/64,\quad r_+=1/32.
 \label{eq:v19-profile}
\end{equation}
It is enough that these plateau and support conditions hold; no strict
monotonicity is used. Profiles and the V14 completions range over a
fixed class with bounded prescribed derivative seminorms and fixed
inverse margins. Physical Fourier pairings and cell weights are unchanged.

Fix $N\ge1$ and a finite exterior-algebra panel. There are $n\le N$
minimal-antifield occurrences, $r$ of $\eta=q^*$ and $s_*$ of
$\sigma=c^*$, so $n=r+s_*$. There is one source-curvature insertion
and $r+2s_*+1$ retained ghosts. There are $d=0$ or $1$ retained metric
marks. Thus the number of ordinary external marks is
\begin{equation}
 m_I=2r+3s_*+1+d.
 \label{eq:v19-mark-count}
\end{equation}
The subscript $s_*$ is used here to avoid confusion with the profile $s$.
Each external momentum has norm at most $\mu$. The last label is fixed
by conservation; $K$ denotes the $m_I-1$ independent labels. Fix
\begin{equation}
 \boxed{\Lambda_0\ge256(N+1)\mu>0,\qquad
 a\Lambda_0\le\frac1{16384(N+1)},\qquad
 \Lambda_0 L_\nu\ge1.}
 \label{eq:v19-domain}
\end{equation}
The periodic component sizes $L_\nu/a$ are odd integers. These
conservative constants separate all soft subset sums from the hard
support, and all momenta in the marked annulus from the lower cutoff.
The statements with momentum derivatives use the same smooth
interpolation convention as V18. There is no claim of exact
renormalized finite-torus BV covariance for that interpolation.

\section{Exact curvature operands on the finite carrier}
\label{sec:v19-curvature}

Put $u=Pc$, $b_c=u\cdot Du$, and $w_c=(P-I)b_c$. Introduce
\begin{align}
 \mathcal K_c H
  &=P\{\mathcal L_{w_c}H+\mathcal L_u(I-P)\mathcal L_u H\},\\
 \mathcal J(c)
  &=w_c\cdot Du-u\cdot Dw_c.
 \label{eq:v19-KJ}
\end{align}
Here $H$ is a symmetric metric array, not an antifield. Compositions
have their displayed order; Lie derivatives carry the external
Grassmann coefficients on the left. These are the expressions already
identified on nonaliased panels in \eqref{eq:v18-next-curvatures}.

\begin{lemma}[The displayed curvature formulas are sufficient on the full grid]
\label[lemma]{lem:v19-global-curvature}
For the supported profile \eqref{eq:v19-profile}, the square of the
actual finite source \eqref{eq:v18-repaired-source} satisfies
\begin{equation}
 \boxed{\breve{\mathfrak s}^{\,2}g=\mathcal K_c(g-I),\qquad
 \breve{\mathfrak s}^{\,2}c=\mathcal J(c),\qquad
 \breve{\mathfrak s}^{\,2}q
       =A(q)^{-1}\mathcal K_c(G(q)-I).}
 \label{eq:v19-exact-curvature}
\end{equation}
These are identities of finite arrays, with cyclic multiplication.
No derivative product rule on unfiltered factors is assumed.
The coefficients of the last expression of metric degrees one, two
and three are respectively
\begin{align}
 N_q^{[1]}&=\mathcal K_c q,\\
 N_q^{[2]}&=\mathcal K_c\mathsf Q(q,q)
                      -\mathsf B_q\mathcal K_cq,\\
 N_q^{[3]}&=\mathsf B_q^2\mathcal K_cq
                      -\mathsf B_q\mathcal K_c\mathsf Q(q,q).
 \label{eq:v19-N123}
\end{align}
\end{lemma}
\begin{proof}
Before applying any continuum identity, direct differentiation of the
finite metric source gives
\[
 R_0P b_c-P\mathcal L_uR_0u
       +P\mathcal L_{P b_c}H-P\mathcal L_uP\mathcal L_uH.
\]
The first two terms cancel by the already verified flat identity in
V18. The remaining rearrangement requires
$P(\mathcal L_u^2-\mathcal L_{b_c})H=0$.
Every Fourier label of $u$ has $|ap|<r_+$. If a summand in this
identity survives the outer $P$, its final label has size at most
$r_+/a$, its input $H$ label has size at most $3r_+/a$, and all
intermediate labels have size at most $4r_+/a$. Cyclic wraparound
cannot convert a principal-face $H$ mode into a central output by
adding two labels of size $r_+/a$. Since $4r_+<\pi$, these particular
products are ordinary nonaliased products and the ordinary graded
Lie identity is valid. It yields the first formula in
\eqref{eq:v19-KJ}. This argument also covers an arbitrary $H$ that was
itself formed as a cyclic product of two high modes: it uses its total
label only and never differentiates its two factors separately.

All labels in the cubic ghost calculation are sums of three labels
of $u$, of size at most $3r_+/a<\pi/a$. The ordinary three-ghost
Jacobi cancellation therefore applies there exactly. Differentiating
$b_c$ and subtracting that cancellation leaves $\mathcal J(c)$.
Finally the square of an odd derivation is an even vector field.
Under an ordinary finite-dimensional coordinate map, the second
coordinate derivative contracts two odd tangent vectors and cancels
by antisymmetry. Thus its square pulls back by $A^{-1}$, proving the
last identity. Expanding $(I+\mathsf B_q)^{-1}$ through degree two
then gives \eqref{eq:v19-N123}. These are finite pointwise matrix
operations, not a nonlocal coordinate inverse.
\end{proof}

Define the odd, ghost-number-one insertion
\begin{equation}
 \boxed{\Omega_a(q,c;\eta,\sigma)
 =\langle\eta,N_q(q,c)\rangle+\langle\sigma,\mathcal J(c)\rangle.}
 \label{eq:v19-Omega}
\end{equation}
With the antibracket convention of V15, it is the curvature component
$\tfrac12(\mathcal H_r,\mathcal H_r)$ of
$\mathcal H_r=\langle\eta,\breve r^q\rangle+
\langle\sigma,\breve r^c\rangle$ (the nonminimal doublet has zero
square). It is \emph{not} the full master defect of the action.
An independent odd tag $\zeta$ is placed on the left so that
$\zeta\Omega_a$ is even in the measured exponent. The tag is a
constant bookkeeping source, not an additional integrated field.

\subsection{The lowest dressed curvature vertices}

The following formulas make explicit what is consumed from V18.
For a hard ghost test $\beta$, let $v=P\beta$ and set
\[
 b_1(c,\beta)=c\cdot Dv+v\cdot Dc,\qquad
 w_1=(P-I)b_1.
\]
When $c$ and the contracted antifield labels are soft, the coefficient
of hard degree two in \eqref{eq:v19-Omega}, at zero retained metric,
is the quadratic form
\begin{align}
 \tfrac12\langle x,\mathscr E_0 x\rangle
 ={}&\Big\langle\eta,
       \mathcal K_c\mathsf Q(x,x)-\mathsf B(x,\mathcal K_cx)
                  +\mathcal K_{c,\beta}x\Big\rangle\notag\\
 &+\Big\langle\sigma,[t^2]\mathcal J(c+t\beta)\Big\rangle,
 \qquad \beta=-C_cB_cx,\quad
 \mathcal K_{c,\beta}=[t]\mathcal K_{c+t\beta}.
 \label{eq:v19-E0}
\end{align}
The auxiliary parameter $t$ is even. The substitution is performed
\emph{after} extracting the hard-degree coefficients. The old native
ghost-action operand $B_c$ and $C_c$ are exactly those of V18.

For clarity, after pairing with these soft labels the same expression
can be simplified to
\begin{align}
 \tfrac12\langle x,\mathscr E_{0,\eta}x\rangle
 ={}&\Big\langle\eta,
 -\mathsf B(x,P\mathcal L_c(I-P)\mathcal L_cx)
 +P\{\mathcal L_{w_1}x+\mathcal L_v(I-P)\mathcal L_cx\}
 \Big\rangle,\label{eq:v19-Eeta}\\
 \tfrac12\langle x,\mathscr E_{0,\sigma}x\rangle
 ={}&\langle\sigma,w_1\cdot Dv-v\cdot Dw_1\rangle,
 \qquad v=-PC_cB_cx.
 \label{eq:v19-Esigma}
\end{align}
These are odd quadratic forms in the even variable $x$; their odd
coefficients keep the order in which they occur in these formulas.
In particular \eqref{eq:v19-Esigma} contains \emph{two} ghost
covariances. It is not a one-propagator contact.

To verify the simplifications, $Pc=c$ and $(P-I)(c\cdot Dc)=0$.
The total label of $\mathsf Q(x,x)$ in its first pairing in
\eqref{eq:v19-E0} is a sum of soft labels, so $\mathcal K_c$ annihilates
it exactly. In $P\mathcal L_c(I-P)\mathcal L_vx$, the intermediate
label selected by the soft output and soft outer ghost is also in the
plateau, so this term vanishes after the pairing. Finally
$(P-I)(v\cdot Dv)$ in the cubic ghost expression is paired only with
soft $\sigma$ and $c$ and vanishes for the same total-label reason.
These are support cancellations made \emph{before} bounds. They are
not permission to remove high--high contacts from other source panels.

\section{A marked measured integral, including one retained metric derivative}
\label{sec:v19-marked}

Use a formal even parameter $\tau$ with $\tau^2=0$, and set the
retained coframe coordinate to $q_\tau=\tau h$. This requests a first
background derivative, not a finite-background approximation theorem.
Let $A^g[h]$ and $A^c[h]$ be the first hard Hessian derivatives of the
\emph{same} completed action of V14. Define the actual inverse jets
\begin{align}
 C_g(\tau)&=C_g-\tau C_gA^g[h]C_g,\\
 C_c(\tau)&=C_c-\tau C_cA^c[h]C_c.
 \label{eq:v19-background-cov}
\end{align}
Inverses here are only notation for the series on the integrated
support; no inverse is applied to a retained zero mode.

The ghost-action mixed block $B_c(\tau)x$ is the coefficient of
$\bar\beta x$ at the retained background $(\tau h,c)$, through degree
one in $\tau$. Its two terms use the existing cubic and quartic ghost
actions. The latter contains the actual $R_2$ ghost contact of V12.
Define
\begin{align}
 L_{\eta,\tau}(x)\beta
   &=\langle\eta,D_q\breve r^q(\tau h,\beta)[x]\rangle,\\
 \tfrac12\langle x,T_{\eta c,\tau}x\rangle
   &=\tfrac12\langle\eta,D_q^2\breve r^q(\tau h,c)[x,x]\rangle,\\
 \tfrac12\langle x,D_\tau x\rangle
   &=\tfrac12\langle x,T_{\eta c,\tau}x\rangle
       -L_{\eta,\tau}(x)C_c(\tau)B_c(\tau)x\notag\\
   &\quad+Q_\sigma(C_c(\tau)B_c(\tau)x).
 \label{eq:v19-background-D}
\end{align}
$Q_\sigma$ is V18's ghost-source quadratic form. All derivatives here
are of the closed explicit source \eqref{eq:v18-repaired-source}.
The source coefficients required are at most $\breve R_3$; the
curvature coefficients needed separately are exactly
\eqref{eq:v19-N123} and $\mathcal J$. No new unspecified higher
action vertex is used.

The marked quadratic operand at this background is
\begin{equation}
 \boxed{\tfrac12\langle x,\mathscr E_\tau x\rangle
 =[t^2]\Omega_a(\tau h+tx,c+t\beta;\eta,\sigma)
       \big|_{\beta=-C_c(\tau)B_c(\tau)x}.}
 \label{eq:v19-background-E}
\end{equation}
Equations \eqref{eq:v19-KJ}, \eqref{eq:v19-N123} and
\eqref{eq:v19-background-E} are a finite evaluation rule for this
operand, with a fixed Grassmann ordering. In particular, differentiating
$\tau$ acts on \emph{both} ghost inverses and all source and ghost-action
vertices. It does not merely attach $h$ to an external leg.

\begin{theorem}[The complete one-loop curvature insertion]
\label{thm:v19-marked-generator}
On \eqref{eq:v19-domain}, at the declared backgrounds and every fixed
minimal-antifield panel, the coefficient of $\zeta$ in the one-loop
connected logarithm after adding $\zeta\Omega_a$ to the action is
\begin{equation}
 \boxed{\mathfrak I_{a,1}(\tau;\eta,\sigma,c)
 =\frac12\operatorname{Tr}
  \left[(I+C_g(\tau)D_\tau)^{-1}C_g(\tau)\mathscr E_\tau\right].}
 \label{eq:v19-marked-generator}
\end{equation}
The common factor $\hbar$ is removed. The coefficient with $n$ total
minimal-antifield occurrences comes from exactly $n-1$ factors $D$:
\begin{equation}
 \mathfrak I_{a,1}
 =\frac12\sum_{m\ge0}(-1)^m\operatorname{Tr}
      [(C_g(\tau)D_\tau)^m C_g(\tau)\mathscr E_\tau].
 \label{eq:v19-marked-words}
\end{equation}
Every word keeps all internal ghost propagators inside $D$ and
$\mathscr E$. The two coefficients $[\tau^d]$, $d=0,1$, are the complete
stated one-loop insertions, not only selected diagrams.
\end{theorem}
\begin{proof}
On the hard quadratic sector, Berezin integration over the antighost
makes exactly the substitution
$\beta=-C_c(\tau)B_c(\tau)x$, as in V18. The insertion contains no
antighost, so it does not change this constraint or its source-independent
determinant. The remaining quadratic bosonic exponent is
\[
 \tfrac12\langle x,
   (C_g(\tau)^{-1}+D_\tau+\zeta\mathscr E_\tau)x\rangle.
\]
Extracting the leftmost odd tag from its finite determinant gives
\eqref{eq:v19-marked-generator}. Each $D$ is even; moving the tag past
it makes no additional sign. $\mathscr E$ is odd, and its internal
odd factors retain their written order. Cyclicity with the even
factors proves the word expansion. There is no extra $1/n$ in a
marked word: the $n$ possible positions of the distinguished insertion
cancel the $1/n$ of the logarithm. V18's scalar determinant is thereby
\emph{used}, not reproved as a new abstract result.

For completeness of the quadratic-hard reduction, first
$\Omega_a(\tau h,c)=0$ on this soft panel: all labels involved in
its fixed jet are in the plateau, and the ordinary source square
vanishes there. Every hard-linear vertex, marked or unmarked, has
momentum a sum of at most $3N+2$ soft labels. It lies below the
integrated support by \eqref{eq:v19-domain}. A connected one-loop
graph has $E=V$; if every participating vertex has at least two hard
legs, $2E=\sum_v\deg(v)$ forces every degree to be two. The coefficient
in the theorem therefore uses exactly the listed quadratic vertices.
Source-independent hard vertices of degree at least three cannot
contribute. An order-$\hbar$ auxiliary-measure insertion would require
another contraction loop or a forbidden hard-linear return; its
soft tree product cancels from the normalized marked expectation.
It is not reassigned to this coefficient. These arguments also hold
with one retained metric mark.
\end{proof}

Writing $\mathcal C=(C_g^{-1}+D_0)^{-1}$, the first background
coefficient has the particularly useful form
\begin{equation}
 [\tau]\mathfrak I_{a,1}
 =\tfrac12\Tr\{\mathcal C\mathscr E'_0
       -\mathcal C(A^g[h]+D'_0)\mathcal C\mathscr E_0\}.
 \label{eq:v19-marked-first-jet}
\end{equation}
The prime means the full derivative in
\eqref{eq:v19-background-cov}--\eqref{eq:v19-background-E}.
This displays the metric covariance return and the ghost/source
returns separately. Omitting $D'_0$, or differentiating $C_g$ while
holding the ghost inverses fixed, changes the measured insertion.

\section{The marked loop has an exact ultraviolet-annulus support}
\label{sec:v19-annulus}

For each panel $I$ and $d=0,1$, let $f_{a,I}^{(d)}(p,K)$ be the
finite sum of routed integrands obtained from
\eqref{eq:v19-marked-words}; traces are divided by physical volume.
We choose the routed variable on a metric line incident to the
marked vertex. Up to the orientation of each internal line, all other momenta differ
by fixed soft subset sums, also after expansion of the ghost-mediated
returns.

\begin{theorem}[Support and exact scaling of the assembled marked integrand]
\label{thm:v19-annulus}
There are fixed $\delta_->0,\delta_+<\pi$, which may be taken as
$\delta_-=r_-/4$ and $\delta_+=4r_+$, such that
\begin{equation}
 \boxed{f_{a,I}^{(d)}(p,K)=0
 \quad\hbox{unless}\quad \delta_-\le |ap|\le\delta_+.}
 \label{eq:v19-annulus}
\end{equation}
This statement applies to the signed, support-simplified integrand,
not to separately estimated summands before the cancellations in
\eqref{eq:v19-E0}. On its support every lower covariance factor
$\nu$ is exactly one. The integrand extends smoothly and periodically
in $p$ and has the exact scaling
\begin{equation}
 \boxed{f_{a,I}^{(d)}(p,K)
  =\chi^{-n}a^{n-1}\mathfrak f_I^{(d)}(ap,aK).}
 \label{eq:v19-integrand-scaling}
\end{equation}
For every fixed multi-index, $\mathfrak f_I^{(d)}$ and its
 derivatives have compact annular support and uniform bounds depending
only on the panel, the fixed profile class and inverse margins.
Polarization norms multiply these bounds.
\end{theorem}
\begin{proof}
First work at $d=0$. If every hard label has dimensionless magnitude
less than $r_-$ with a margin exceeding all soft subset sums, every
$P$ encountered in the curvature operand equals one. Then
$\mathcal K_c$ and $\mathcal J$ vanish by the exact identities
of \cref{lem:v19-global-curvature}. Their coordinate pullback vanishes
as well. This is an \emph{algebraic} cancellation; it uses no equation
of motion, reference covariance, or integrated Ward identity.

For the upper support use the explicit quadratic formulas. In
\eqref{eq:v19-Eeta}, the term with two hard metric legs has an outer
$P$ on one of the hard labels, shifted only by soft ghosts. Every
mixed metric/ghost term has the hard ghost factor $v=P\beta$;
after substitution its momentum differs from a hard metric momentum
only by the retained ghost label in $B_c$. The ghost curvature term
\eqref{eq:v19-Esigma} has two such factors. If the hard labels lie
above $r_+$ plus the soft margin, all these terms vanish. The term
$\mathcal K_c\mathsf Q(x,x)$ that might have carried arbitrarily
large paired hard labels was annihilated at its \emph{soft total}
label before this upper estimate. This is why the coordinate contact
must be kept until that cancellation is made.

At $d=1$, use \eqref{eq:v19-N123}.
A pointwise $\mathsf B_h$ changes a hard label only by the soft
background momentum. The possible $\mathcal K_c$ acting on a product
whose total label is soft still vanishes. Each remaining term either
has a hard ghost under $P$, or has a hard metric input/output under
one of the $P$ factors in $\mathcal K_c$. The additional covariance
and ghost-action derivatives carry only soft momentum differences.
The same lower and upper argument therefore holds with the fixed
larger soft margin $ (3N+2)\mu$. The domain
\eqref{eq:v19-domain} makes this margin much smaller than $r_-/a$;
the conservative annulus in \eqref{eq:v19-annulus} covers it.
It also places every internal line above the support of $\Theta$, so
$\nu=1$ on every nonzero term. There is no lower-scale denominator
in the ensuing dimensionless integral.

On the marked annulus the completed covariances have scales
$\chi^{-1}a^2$ and $a^2$. Each ordinary dressed $D$ has scale $a^{-1}$
by its actual formula in V18. Each curvature operand has two
spectral derivatives: $\mathscr E$ has scale $a^{-2}$, including
both substitutions $C_cB_c$ of scale zero. Thus $n$ metric
covariances, $n-1$ ordinary $D$ vertices and one marked $\mathscr E$
give $\chi^{-n}a^{n-1}$. Differentiating a background covariance
adds one covariance and an action vertex of order two; their scales
cancel. In particular the $A^g$ contribution contains one factor
$\chi$ which cancels the extra $\chi^{-1}$ covariance. This proves
\eqref{eq:v19-integrand-scaling} also at $d=1$.

All symbols on this support are the specified smooth periodic
symbols of the completed reference. Each differentiated spectral
factor inside a source is confined to a central chart by its filters,
or acts on a soft total label. The support argument in
\cref{lem:v19-global-curvature} excludes a hidden seam at such a
factor. Smooth extension by zero is valid at the profile boundaries.
Fixed differentiated products of the explicit matrices and their
uniformly invertible annular inverses now prove the asserted bounds.
\end{proof}

This support result is stronger than a small low-momentum estimate.
The nonzero source square at cutoff-size ghosts, established in V18,
has not been set to zero. Its measured one-loop insertion is confined
to the region \eqref{eq:v19-annulus}, where a direct local Taylor
assignment can be paid for. The proof makes no claim that the whole
classical action breaking has the same support.

\section{A finite local curvature bank and a vanishing nonlocal limit}
\label{sec:v19-limit}

Write $\iota_{a,L,I}^{(d)}(K)=\int_p^L f_{a,I}^{(d)}(p,K)$ for the
one-loop coefficient of the marked insertion per physical volume.
Let $\boldsymbol\ell=(L_0/a,\ldots,L_3/a)$ be the vector of odd site
counts. \Cref{thm:v19-annulus} gives an exact finite-sum formula
\begin{equation}
 \boxed{\iota_{a,L,I}^{(d)}(K)
   =\chi^{-n}a^{n-5}\mathcal F_{I,\boldsymbol\ell}^{(d)}(aK),
 \quad
 \mathcal F_{I,\boldsymbol\ell}^{(d)}(z)
   =\frac1{\prod_\nu\ell_\nu}
      \sum_{v\in\mathcal B_1\cap\prod_\nu(2\pi/\ell_\nu)\mathbb Z}
                 \mathfrak f_I^{(d)}(v,z).}
 \label{eq:v19-scaling}
\end{equation}
The summand is zero outside the fixed annulus. This is not a
continuum integral inserted in place of the measured finite trace.
The latter is exactly the displayed sum. Its derivative bounds follow
from a normalized sum with at most $\prod\ell_\nu$ terms and contain
no extensive determinant norm.

For $n\le5$ define
\begin{equation}
 \mathsf P_{\Omega,I}=\mathsf T_{5-n},\qquad
 \iota_{a,L,I}^{R,(d)}=(I-\mathsf P_{\Omega,I})\iota_{a,L,I}^{(d)}.
 \label{eq:v19-projector}
\end{equation}
For $n>5$ take $\mathsf P_{\Omega,I}=0$.
The local Taylor terms are stored as \emph{ghost-number-one insertion
coefficients}; an independent derivation is needed before identifying
them with the variation of ghost-number-zero action counterterms.
Their field degree is finite at each declared $N$ and their derivative
degree is at most $5-n$. Reality, external permutation and Grassmann
relations may reduce a redundant component bank. No covariant or
cohomological identification is assumed here.

\begin{theorem}[Measured source-curvature locality at every fixed antifield degree]
\label{thm:v19-vanishing}
For \eqref{eq:v19-domain}, $d=0,1$, every fixed panel with
$n\le5$, and $j=|\alpha|\le6-n$, one has
\begin{equation}
 \boxed{\left|\partial_K^\alpha
        \iota_{a,L,I}^{R,(d)}(K)\right|
 \le C_{I,\alpha}\chi^{-n}\mathfrak n_I\,
                    a|K|^{6-n-j}.}
 \label{eq:v19-main}
\end{equation}
For $j>6-n$ a bound is
$C_{I,\alpha}\chi^{-n}\mathfrak n_I a^{n-5+j}$.
For $n>5$, without subtraction,
\begin{equation}
 \boxed{|\partial_K^\alpha\iota_{a,L,I}^{(d)}(K)|
 \le C_{I,\alpha}\chi^{-n}\mathfrak n_I a^{n-5+j}.}
 \label{eq:v19-high-degree}
\end{equation}
Thus the limiting nonlocal coefficient of this insertion is zero.
The actual local coefficients are
\begin{equation}
 \boxed{\partial_K^\alpha\iota_{a,L,I}^{(d)}(0)
 =\chi^{-n}a^{n-5+j}
       \partial_z^\alpha\mathcal F_{I,\boldsymbol\ell}^{(d)}(0),
       \qquad j\le5-n.}
 \label{eq:v19-local-data}
\end{equation}
They can diverge or remain finite. The top allowed derivative
coefficient is in general a finite matching datum, not an error.
\end{theorem}
\begin{proof}
The dimensionless derivatives of $\mathcal F$ in
\eqref{eq:v19-scaling} are bounded uniformly in all site counts on
the declared domain. For $n\le5$ use the multivariate integral Taylor
formula through degree $5-n$, with $q=6-n$ derivatives. The prefactor
is $a^{n-5}$, while those $q$ derivatives in $K$ supply $a^q$.
Their product is $a$. Differentiating this Taylor formula $j\le q$
times gives \eqref{eq:v19-main}. If $j>q$, the subtraction polynomial
has disappeared, and direct differentiation of \eqref{eq:v19-scaling}
gives the stated bound. The same direct differentiation proves
\eqref{eq:v19-high-degree} and \eqref{eq:v19-local-data}.
There is no logarithm from an infrared-to-ultraviolet annulus sum:
the entire marked loop lies in one fixed dimensionless annulus.
The assertion that its zero-mesh nonlocal coefficient is zero follows
from these estimates and not from an assumed continuum Ward identity.
\end{proof}

\begin{center}\small
\begin{tabular}{ccc}
\toprule
Minimal-antifield count $n$ & Stored Taylor degree & Nonlocal bound at $j=0$\\
\midrule
1 & 4 & $C\chi^{-1}a|K|^5$\\
2 & 3 & $C\chi^{-2}a|K|^4$\\
3 & 2 & $C\chi^{-3}a|K|^3$\\
4 & 1 & $C\chi^{-4}a|K|^2$\\
5 & 0 & $C\chi^{-5}a|K|$\\
$n>5$ & none & $C\chi^{-n}a^{n-5}$\\
\bottomrule
\end{tabular}
\end{center}
The same table applies with one retained metric mark. No unproved
odd/even parity improvement is used. Compared with V18's unmarked
module, taking the square of the source adds one derivative and
raises the subtraction degree by one.

\subsection{Rooted kernels and the separate thermodynamic comparison}

Take a fixed smooth compact external window $\varpi(K/\mu)$ on the
declared panel and let $\mathscr K_{a,L,I}$ be the inverse Fourier
kernel of its product with $\iota^{R,(d)}$. There are
$M=4(m_I-1)$ relative coordinates. For every fixed $b\ge0$,
\begin{equation}
 \boxed{a^M\sum_{x_1,\ldots,x_{m_I-1}}
 (1+\mu\textstyle\sum_i d_{\rm per}(x_i,0))^b
       \|\mathscr K_{a,L,I}(\boldsymbol x)\|
 \le C_{I,b}\chi^{-n}
 \begin{cases}
 a\mu^{6-n},&n\le5,\\
 a^{n-5},&n>5.
 \end{cases}}
 \label{eq:v19-kernel}
\end{equation}
To prove this, every scaled external derivative of the windowed
symbol is bounded by its displayed right-hand-side scale times
$\mu^{-j}$; for high derivatives use $a\mu\ll1$ and the high-derivative
part of \cref{thm:v19-vanishing}. Fourier integration by parts more
than $M+b$ times gives an integrable rapidly decreasing kernel.
Periodization and $d_{\rm per}\le d$ for any lift give
\eqref{eq:v19-kernel}. As usual, one external source can then be
measured in $L^1$ and the others in $L^\infty$. This is a genuine
rooted multilinear estimate and has no total-volume multiplier.

For fixed source tensors replace the normalized dimensionless sum
in \eqref{eq:v19-scaling} by
\[
 \mathcal F_{I,\infty}^{(d)}(z)
   =\int_{\mathcal B_1}\frac{\dd^4v}{(2\pi)^4}
                       \mathfrak f_I^{(d)}(v,z).
\]
For every $b>4$ and fixed derivative order, smooth periodic Poisson
summation gives
\begin{equation}
 \|\partial_z^\alpha(\mathcal F_{I,\boldsymbol\ell}^{(d)}
                   -\mathcal F_{I,\infty}^{(d)})\|
 \le C_{I,\alpha,b}\ell_{\min}^{-b}.
 \label{eq:v19-volume}
\end{equation}
Indeed integration by parts gives Fourier coefficients bounded by
$C_b(1+|m|)^{-b}$, and the rectangular sampling error is the sum at
nonzero multiples of the site counts. This proof uses smoothness
on the actual momentum torus, not a shift across a discontinuous face.
Multiplying the Taylor-remainder proof by this bound gives an extra
factor $(a/L_{\min})^b$ in the finite-volume comparison of
\eqref{eq:v19-main}. The local matching integrals also have the
comparison \eqref{eq:v19-volume}, with their actual powers of $a$
retained. This separates a finite-volume estimate from a symmetry claim.

\section{Shell ownership and stability of the marked local assignment}
\label{sec:v19-shells}

Expand every metric and ghost covariance in
\eqref{eq:v19-marked-words} into the same smooth dyadic innovations.
This includes the ghost covariances inside $D$, $\mathscr E$, and
$D'_0,\mathscr E'_0$. A word is assigned to the maximum of \emph{all}
its internal shell indices. The partition is exact. In particular a
marked $\sigma$ operand has two internal ghost labels, not one.
This uses the inherited mixed-shell ownership rule with the new
operands; it is not a new independence assumption.

\begin{proposition}[Only a bounded number of ultraviolet levels can contribute]
\label[proposition]{prop:v19-shells}
For dyadic annular innovations and fixed $N$, all nonzero marked words
on \eqref{eq:v19-domain} have internal scales comparable to $a^{-1}$.
The number of possible maximum shell levels is bounded independently
of $a$, the number of sites and the total number of preceding shells.
After the actual shellwise local assignments, their rooted remainder
sum obeys \eqref{eq:v19-kernel}. The local coefficient sums agree
with \eqref{eq:v19-local-data}.
\end{proposition}
\begin{proof}
By \eqref{eq:v19-annulus}, the routed line is in
$[\delta_-/a,\delta_+/a]$. Every other line differs by at most
$(3N+2)\mu$, which is smaller than a fixed fraction of
$\delta_-/a$. Hence all its annular shell indices range over an
interval of length bounded by a constant depending on
$\delta_+/\delta_-$ and the shell support ratio only. For fixed $N$
there are finitely many words and line assignments per such interval.
Their smooth dimensionless symbols have the same derivative bounds as
in \cref{thm:v19-annulus}; repeat the Taylor proof on each word.
Summing a bounded number of levels gives the assertion. Taylor
projection is linear, so the local sums reproduce the actual full
local coefficient. This does not apply to arbitrary multiloop graphs,
whose independent loop momenta need not all be in this annulus.
\end{proof}

For two supported profile choices, the two constructions define
\emph{different finite insertions}. Each has a local jet specified by
\eqref{eq:v19-local-data} and a remainder bounded by
\eqref{eq:v19-main}. Thus, after their own local assignments, their
nonlocal difference tends to zero with the sum of those bounds.
This conclusion concerns matched insertion prescriptions. It does
not transfer a finite Slavnov identity without also transforming its
reference, action-breaking and measure operands.

The fixed-panel Taylor map is bounded, for example in the scaled
$C^{5-n}$ norm on a fixed external momentum ball with the Taylor
polynomial equipped with its coefficient norm. Its bound depends on
the panel and degree but not on mesh or volume. This is an analytic
projection of the specified insertion space. It has not been proved
to be a projection commuting with the full ESM BV differential.

\section{What the result removes from the quantum algebra problem}
\label{sec:v19-BV-scope}

Let $S_{\rm cl}$ be the same zero-antifield action and
$\mathcal H_r$ its recorded minimal transformation source. In the
convention of this lineage the classical master expression separates as
\begin{equation}
 \tfrac12(S_{\rm cl}+\mathcal H_r,S_{\rm cl}+\mathcal H_r)
       =(S_{\rm cl},\mathcal H_r)+\Omega_a,
 \label{eq:v19-master-separation}
\end{equation}
with the nonminimal action and doublet included consistently when that
sector is restored. The first term is action breaking; the second is
source-algebra curvature. Neither is silently replaced by the other.
There is also an exact consistency check on this particular curvature.
For the coefficient density $J_G\,dq\,dc$ of V18, the verified total
source divergence is zero. The graded divergence identity and graded
Jacobi identity consequently give
\[
 \operatorname{div}_{J_G\,dq\,dc}(\breve{\mathfrak s}^{\,2})=0,
 \qquad [\breve{\mathfrak s},\breve{\mathfrak s}^{\,2}]=0.
\]
These are applications of inherited vector-field calculus to the actual
source, not an assertion that its square vanishes. Its Hamiltonian
operator still has square given by the bracket with $\Omega_a$.
Thus one cannot treat it as a differential on a complex before paying
for that curvature and the separate action breaking.
The general identity \eqref{eq:v19-master-separation} is inherited BV
algebra. What is proved here is the measured insertion theorem for its
\emph{second} operand, including one ordinary metric derivative and
arbitrary prescribed minimal-antifield degree.

This is precisely where the old minimal-antifield contacts are needed.
Already the $n=2$ marked term is
\[
 -\tfrac12\Tr(C_gD C_g\mathscr E),
\]
not a seagull-only approximation to the $\eta\eta ccc$ coefficient. The ghost-mediated $D_\eta$
return and $D_\sigma$ with two ghost covariances are part of it.
The one-metric coefficient additionally contains
\eqref{eq:v19-marked-first-jet}. Thus the theorem consumes rather
than bypasses the contacts established in V18.

There are two distinct implications:
\begin{equation}
 \boxed{\text{proved: }\ (I-\mathsf P_\Omega)
     [\text{one-loop measured }\Omega_a]_I\longrightarrow0,}
 \label{eq:v19-proved-implication}
\end{equation}
whereas the existence of a local ghost-number-zero counterterm $C_I$
whose full BV variation cancels all remaining local terms is
\emph{not inferred}. A subtraction of a ghost-number-one insertion
is not a proof of that second statement. The finite data
\eqref{eq:v19-local-data} must be combined with the locally assigned
action-breaking, reference and measure terms on the same source
prescription. Their combined consistency conditions, a restoring
primitive if it exists, and a cutoff-uniform norm for its choice are
still required. Finite local cohomological existence alone would not
supply the last estimate.

Accordingly, the nonzero finite witnesses in V18 remain true. This
installment does not make the source nilpotent at arbitrary lattice
momenta. It proves that their explicitly identified curvature mechanism
has no surviving \emph{nonlocal one-loop insertion} in the specified
soft-source topology. This was not a consequence of the two earlier
one-ghost Slavnov rows, and it was not assumed in the proof.

The next load-bearing object is now the \emph{combined local}
quantum algebra row with one minimal antifield and two or three ghosts,
including its ordinary-background derivative. Its source-curvature
piece is the finite computable jet \eqref{eq:v19-local-data}; its
nonlocal remainder has been paid for. The actual action-breaking and
reference insertions must be matched to that piece before a common
BV-restoring counterterm can be claimed. Extending to arbitrary ordinary
background degree, mixed SM/chiral fields and higher loops remains a
separate induction, not a corollary of the present marked determinant.

\section{Verification and closing ledgers for V19}
\label{sec:v19-ledgers}

The standalone verifier
\srcid{check_ESM_source_curvature_V19.py}
separates finite exact algebra from analytic estimates. It checks the
supported finite Fourier curvature identities including high cyclic
inputs, coordinate transport of the squared source, the marked
Berezin/Gaussian sign and marked-word multiplicity, the complete first
background derivative, and the degree/scaling and shell-ownership
bookkeeping. Such finite tests do not certify the continuum proof in
a proof assistant. The proofs of support, symbol estimates, Taylor
bounds and kernel summation are given above.

\subsection*{Proved}
\addcontentsline{toc}{subsection}{V19: Proved}
The supported finite source has the exact global curvature operands
\eqref{eq:v19-exact-curvature}; the necessary quadratic and cubic
coframe jets are \eqref{eq:v19-N123}. Their actual one-loop insertion
is \eqref{eq:v19-marked-generator}, at zero or one retained metric
mark and every prescribed finite minimal-antifield degree. Its
assembled integrand has proved compact dimensionless ultraviolet
support, exact scaling and period-independent derivative bounds.
The resulting finite local assignment gives
\eqref{eq:v19-main}--\eqref{eq:v19-high-degree}, the rooted kernel
bound \eqref{eq:v19-kernel}, and a bounded number of active ultraviolet
shell levels. All internal ghost returns and the full background
covariance derivative are retained. This is locality and vanishing
of a specified renormalized breaking insertion, not full quantum BV
closure.

\subsection*{Inherited}
\addcontentsline{toc}{subsection}{V19: Inherited}
V1--V18 and their labels are unchanged. The completed action and
covariances, metric contour, positive lower reference, coframe chart,
coordinate Jacobian, auxiliary measure, native cubic and quartic ghost
actions, source recursion and unmarked minimal-antifield module are
used at their stated scopes. V18's nonzero higher-square witnesses
remain valid. General Berezin elimination, Gaussian marking,
coordinate vector-field transport and mixed-shell allocation are
inherited tools; the marked curvature operand and its support theorem
are the new application needed here.

\subsection*{Literature input}
\addcontentsline{toc}{subsection}{V19: Literature input}
The modified quantum-master/Slavnov framework and composite insertions
are established, with \cite{V15IIM} providing the comparison already
used in this lineage. No general anomaly-removal theorem or first
construction of perturbative gravity is claimed. No external continuum
Ward identity is used to establish the annular support or the marked
cutoff estimate.

\subsection*{Conditional}
\addcontentsline{toc}{subsection}{V19: Conditional}
A full ESM application requires the action-breaking, reference,
nonminimal and measure insertions to be assembled with these local
curvature jets. A common bounded BV-restoring primitive and compatible
physical normalization remain to be constructed. Matching to the
unmodified microscopic regulator, ordinary-background induction and
mixed internal-gauge/fermionic source estimates are not supplied by
the present completed gravitational reference. Neither a vanishing
subtracted insertion nor the bounded Taylor map alone implies a
nilpotent finite differential or a cohomological splitting.

\subsection*{Open}
\addcontentsline{toc}{subsection}{V19: Open}
The next concrete target is the combined local quantum algebra
consistency/restoration problem, using
\eqref{eq:v19-local-data} with the other actual local breaking and
reference coefficients. The nonlocal contribution from the two
explicit V18 source curvatures is now controlled in the stated module.
Full source-algebra closure at nonzero unrestricted backgrounds,
chiral-line transport, mixed ESM loops, an invariant filtered
interacting class and arbitrary-loop continuum control remain open.
No infrared removal, full Lorentzian reconstruction, nonperturbative
quantum measure or finite-parameter gravitational ultraviolet completion
is asserted.

\begin{center}\small
\begin{tabular}{>{\raggedright\arraybackslash}p{.22\textwidth}>{\raggedright\arraybackslash}p{.69\textwidth}}
\toprule
Variable & Uniformity and limitation in V19\\
\midrule
Volume & Coefficients per physical volume and rooted kernels have no
extensive factor. All estimates hold for $\Lambda_0L_\nu\ge1$;
thermodynamic local jets have the separate Poisson comparison.\\
Mesh and momenta & The complete smooth regulator is integrated on
\eqref{eq:v19-domain}. The assembled marked coefficient is supported
in \eqref{eq:v19-annulus}; no arbitrary hard mode is deleted from the
underlying theory.\\
Sources and backgrounds & Every prescribed finite minimal-antifield
degree $n\le N$, one curvature mark, zero or one retained metric
mark, and the ghost count \eqref{eq:v19-mark-count}. No bound uniform
as $N\to\infty$ or on an unrestricted background ball.\\
Loops and shells & One propagating loop. All internal metric/ghost
labels are included. The number of contributing ultraviolet levels
is bounded, not the number of independent multiloop momenta.\\
Couplings & Explicit $\chi^{-n}$ for $\chi>0$, fixed completion
parameters and inverse margins. No gravitational running trajectory
is inferred.\\
Measure and infrared & The old coordinate-correct auxiliary measure
is retained once; it is not a second insertion at this loop order.
The lower scale remains positive. No positive gravity measure or
Lorentzian continuation is claimed.\\
\bottomrule
\end{tabular}
\end{center}

\clearpage
\part*{V20: Explicit local primitives for the one-antifield matching terms}
\addcontentsline{toc}{part}{V20: Local primitives and their BV descendant contacts}

\section{From an insertion counterterm to a counterterm in the source action}
\label{sec:v20-start}

Sections 1--152 and their labels are preserved. V19 proves that the
measured source-curvature insertion has a vanishing nonlocal remainder,
not that its local part is the variation of one action counterterm. The
next operation should therefore act on its actual local coefficients,
rather than alter the filtered transformation again. This installment
constructs such an operation on the zero-background, one-minimal-antifield
sector. It removes a specified part of that sector and gives every
additional contact which that removal creates in the free BV complex.
It does not assert that these contacts already agree with the other
uncomputed components of the combined quantum breaking.

There is a useful exact reduction. In a rooted local polynomial with one
metric antifield and two ghosts, every term of total derivative degree
three or four is the linearized gauge variation of an explicit local
source counterterm. For one ghost antifield and three ghosts the same is
true at derivative degree four. More generally a fixed rational map removes
all terms containing a symmetrized first ghost derivative or a ghost
derivative of order at least two. The residual is obtained by testing only
an undifferentiated ghost and its antisymmetric first derivative. Its
finite dimension is computed below. In particular the fourth-derivative,
finite $a^0$ matching terms supplied by V19 have no obstruction at this
\emph{first longitudinal step}. Their lower-antifield descendants and the
remaining local components still have to be matched together.

The construction of contractible jet pairs and the distinction between
longitudinal and full BRST cohomology are established methods
\cite{V20BBH1994,V20BBH1995}. They are not claimed as new methods here.
The contribution is their explicit, norm-controlled realization for the
native $R_0$ and V19's finite annular coefficient arrays, including their
physical derivative degrees, Grassmann signs, the finite quotient size,
and the resulting shifts of the already retained matching equations.
The abstract finite-dimensional splitting criterion of \cite{V6CMP} is
inherited and is not reproved or assumed to provide a local inverse.
No classification theorem from \cite{V20BBH1995} is used to declare the
present regulator anomaly free; in particular its discussion of abelian
factors cannot be bypassed in the eventual Standard Model application.

\subsection{Scope and terminology}

The mesh is $a$, the loop parameter is $\hbar$, and $\chi>0$ is the
inherited Palatini normalization. Use the completed regulator, density,
and \emph{unchanged classical source} of V18--V19. This calculation concerns
one propagating loop, no integrated Standard Model fields, and a fixed
positive lower scale. Its algebraic identities are identities of local
polynomial densities, or of their integrals on a periodic smooth core.
Their samples give literal finite-panel identities whenever all relevant
Fourier sums lie in the central plateau. They are not identities on all
arbitrary high-frequency lattice configurations.

A \emph{primitive} below means an explicit ghost-number-zero local
functional whose stated variation is the specified part of the breaking.
A primitive for the longitudinal operator alone is distinguished from a
primitive for the full BV operator. In particular a nonzero remainder in
the quotient constructed below is not called a physical anomaly.

\section{The actual free gravitational jet map has a rational local section}
\label{sec:v20-jet-section}

Use the Euclidean coefficient chart $G(q)=I+q+\mathsf Q(q,q)$ already
fixed in V15--V19. Its tangent is the ordinary symmetric metric variation.
Let $\eta^{\mu\nu}=q^{*\mu\nu}$ and $\sigma_\mu=c^*_\mu$. The parities
of $q,c,\eta,\sigma$ are respectively $0,1,1,0$, and their ghost numbers
are $0,1,-1,-2$. All products are graded products in one fixed exterior
source algebra. Define the longitudinal free derivation by
\begin{equation}
 \gamma_0q_{\mu\nu}=\partial_\mu c_\nu+\partial_\nu c_\mu,
 \qquad \gamma_0c=\gamma_0\eta=\gamma_0\sigma=0,
 \qquad [\gamma_0,\partial_\mu]=0.
 \label{eq:v20-gamma}
\end{equation}
This is the zero-mesh free generator of the actual chart, not an assigned
nonlinear finite diffeomorphism algebra. In the plateau it is also the
free part of V18's source exactly.

For a multi-index $I\in\mathbb N^4$ use raw jets
$c_{\nu,I}=\partial^Ic_\nu$. Split the first jets into
\begin{equation}
 y_{\mu\nu}=\tfrac12(\partial_\mu c_\nu+\partial_\nu c_\mu),
 \quad y_{\mu\nu}=y_{\nu\mu},\qquad
 z_{\mu\nu}=\tfrac12(\partial_\mu c_\nu-\partial_\nu c_\mu),
 \quad z_{\mu\nu}=-z_{\nu\mu}.
 \label{eq:v20-first-jets}
\end{equation}
Their partners are $x_{\mu\nu}=q_{\mu\nu}/2$. For $|I|\ge2$, choose
once and for all the first two entries $\alpha,\beta$ of the sorted
multi-index $I$, with multiplicity. Put
\begin{equation}
 \begin{split}
 \Gamma_{\nu;\alpha\beta}(q)
 &=\tfrac12(\partial_\alpha q_{\beta\nu}
             +\partial_\beta q_{\alpha\nu}
             -\partial_\nu q_{\alpha\beta}),\\
 x_{\nu,I}
 &=\partial^{I-e_\alpha-e_\beta}
                         \Gamma_{\nu;\alpha\beta}(q),
 \qquad y_{\nu,I}=c_{\nu,I}.
 \end{split}
 \label{eq:v20-Christoffel-section}
\end{equation}
There is no summation in this choice of $\alpha,\beta$. This is the
linear Christoffel polynomial in the tangent coordinate; it is not a
replacement of the native nonlinear auxiliary connection.

\begin{lemma}[Explicit local section of every contractible ghost jet]
\label[lemma]{lem:v20-section}
All the partners in \eqref{eq:v20-Christoffel-section} and
\eqref{eq:v20-first-jets} satisfy
\begin{equation}
 \boxed{\gamma_0x_A=y_A,\qquad \gamma_0y_A=0.}
 \label{eq:v20-pairs}
\end{equation}
The remaining independent ghost coordinates are just $c_\mu$ and
$z_{\mu\nu}$, of dimensions four and six. The partner of a ghost jet
of derivative degree $r$ has one metric field and derivative degree
$r-1$. Its raw coefficient $\ell^1$ norm is at most $3/2$.
For ghost jets through order four there are 270 independent pairs. The
corresponding map from the 350 metric jets through order three has
rank 270 and kernel dimension 80.
\end{lemma}
\begin{proof}
The first assertion for $x_{\mu\nu}$ is immediate. For the second type,
substitute \eqref{eq:v20-gamma} into the three terms of $\Gamma$.
The terms $\partial_\alpha\partial_\nu c_\beta$ and
$\partial_\beta\partial_\nu c_\alpha$ cancel their negative partners,
leaving $\partial_\alpha\partial_\beta c_\nu$. Commuting the remaining
ordinary derivatives proves \eqref{eq:v20-pairs}. There are no momentum
denominators. Each partner has at most three terms of coefficient $1/2$;
coincident indices can only improve the asserted upper bound.

The first ghost jets have the invertible decomposition into ten symmetric
and six antisymmetric entries. The higher symmetric multi-index jets are
independent, with $4\binom{r+3}{3}$ entries at order $r$. Thus the ranks
by metric-jet order $0,1,2,3$ are $10,40,80,140$. Their sum is 270.
The explicit partners prove surjectivity onto those entries. They are
independent because a linear relation among partners would, after
$\gamma_0$, be a relation among independent $y_A$. There are
$10\sum_{r=0}^3\binom{r+3}{3}=350$ metric jets. Rank--nullity gives 80.
This kernel contains the familiar linearized curvature combinations and
their first derivatives, but no curvature-cohomology classification is
needed for this count.
\end{proof}

\section{A bounded primitive map and the exact affine-ghost quotient}
\label{sec:v20-homotopy}

Consider a local density linear in exactly one undifferentiated root
$\eta^{\mu\nu}$ or $\sigma_\mu$, independent of ordinary $q$, and
polynomial in the ghost jets. Move any derivative on the root onto the
ghost coefficient by integration by parts first. This fixes a faithful
rooted representative: its Euler derivative with respect to that root is
its ghost coefficient. Thus a nonzero rooted coefficient is not erased
by quotienting integrated total derivatives.

Let $\pi_{\rm aff}$ set every $y_A$ to zero and retain $c_\mu,z_{\mu\nu}$.
It is the algebraic test of the ghost jet of an affine Killing field,
$c_\nu(x)=u_\nu+\omega_{\mu\nu}x^\mu$ with $\omega$ antisymmetric,
at a point. The test does not assert that such a nonperiodic polynomial
is itself an admissible periodic ghost history.

On a monomial with $k>0$ occurrences of paired odd variables $y_A$,
define
\begin{equation}
 \mathfrak h F=\frac1k\,\mathfrak k F,
 \qquad \mathfrak k=\sum_A x_A\frac{\partial^L}{\partial y_A},
 \qquad \mathfrak hF=0\quad(k=0).
 \label{eq:v20-homotopy}
\end{equation}
The derivative is left graded differentiation, including the parity of
the root preceding it. The $x_A$ in this formula are subsequently replaced
by their explicit polynomials in \eqref{eq:v20-Christoffel-section}.
The divisor $k$ counts paired ghost factors, not derivatives.

\begin{theorem}[Constructive longitudinal primitive on the rooted source bank]
\label[theorem]{thm:v20-homotopy}
For every rooted ghost polynomial $B$ just specified,
\begin{equation}
 \boxed{\gamma_0\mathfrak hB=B-\pi_{\rm aff}B.}
 \label{eq:v20-homotopy-identity}
\end{equation}
The primitive has one ordinary metric field, the same one antifield, one
fewer ghost, and one fewer derivative. It is local and rational in its
input coefficients. In the adapted coefficient $\ell^1$ norm,
\begin{equation}
 \|\mathfrak hB\|_1\le\tfrac32\|B\|_1.
 \label{eq:v20-homotopy-norm}
\end{equation}
For a polynomial with $m$ ghosts expressed in raw ghost jets a bound is
$\tfrac32 2^m\|B\|_1$, after returning the answer to raw metric/ghost
jets. On derivative-homogeneous components, the scaled norm obeys
$\|\mathfrak hB\|_\mu\le C_m\mu^{-1}\|B\|_\mu$.

Moreover $\pi_{\rm aff}B=0$ is necessary and sufficient for a primitive
linear in the ordinary metric, with the root held fixed under
$\gamma_0$. The statement remains true for integrated densities in the
rooted convention. It is a statement about $\gamma_0$, not about full
local BRST cohomology.
\end{theorem}
\begin{proof}
On the freely generated pair algebra the graded anticommutator
$\gamma_0\mathfrak k+\mathfrak k\gamma_0$ counts $x$ and $y$ factors.
The input $B$ contains no $x$ and is annihilated by $\gamma_0$. On each
monomial with $k$ paired factors, dividing by $k$ gives the original
monomial. On a monomial with no paired factor the result is zero. The
explicit section preserves the same identity after substitution, proving
\eqref{eq:v20-homotopy-identity} for the actual metric jets.

There are $k$ replacement terms, each with raw partner norm at most
$3/2$, and the divisor is $k$. This proves
\eqref{eq:v20-homotopy-norm}. Rewriting a first raw ghost derivative as
$y+z$ or $y-z$ costs at most two per ghost. Rewriting $y,z$ back into
raw first derivatives has norm at most one. This proves the raw bound.
Exactly one derivative is lost in each replacement, proving the scaled
version. These constants do not involve a momentum, a cutoff, or a
spacetime volume.

Conversely $\gamma_0$ of any metric jet is in the span of the $y_A$.
The zero-metric part of the variation of a metric-linear polynomial
therefore lies in the ideal generated by them. Its affine projection
vanishes. This proves necessity, and the constructed map proves
sufficiency. Integration by parts can move derivatives of the root
but cannot change this conclusion: the root Euler derivative of a total
derivative is zero, and the rooted coefficient determines the functional.
\end{proof}

\begin{corollary}[The finite one-antifield matching degree has a primitive]
\label[corollary]{cor:v20-top-degree}
A ghost monomial of ghost count $m$ and derivative degree $d>m$ has
zero affine projection. Hence every $\eta cc$ coefficient of derivative
degree three or four, and every $\sigma ccc$ coefficient of derivative
degree four, admits the explicit primitive \eqref{eq:v20-homotopy}.
\end{corollary}
\begin{proof}
Every ghost surviving $\pi_{\rm aff}$ has derivative degree zero or
one. A product of $m$ such ghosts has degree at most $m$.
\end{proof}

\subsection{Size and limitations of the quotient}

The following dimensions impose only the declared tensor components,
commuting derivative multi-indices, and Grassmann antisymmetry. They
are not counts of covariant Wilson coefficients or of possible anomalies.
\begin{center}\small
\begin{tabular}{c|rr|rr}
\toprule
Derivative degree & $\eta cc$: full & affine residual
 & $\sigma ccc$: full & affine residual\\
\midrule
0&60&60&16&16\\
1&640&240&384&144\\
2&2800&150&2880&240\\
3&9600&0&14400&80\\
4&26200&0&55520&0\\
\midrule
Total&39300&450&73200&480\\
\bottomrule
\end{tabular}
\end{center}

\begin{proposition}[Exact component count for this longitudinal quotient]
\label[proposition]{prop:v20-count}
Through derivative degree four the rooted one-antifield bank has
dimension 112500. Its quotient by the image of $\gamma_0$ on
metric-linear primitives has dimension 930; the image dimension is
111570. Actual symmetry or source restrictions can make the populated
subspace smaller.
\end{proposition}
\begin{proof}
There are $N_r=4\binom{r+3}{3}$ raw ghost jets of derivative degree
$r$. Count exterior monomials using
$\prod_{r=0}^4(1+t z^r)^{N_r}$, retaining the coefficient of $t^m z^d$.
Multiply by ten roots for $m=2$ and four roots for $m=3$.
For the affine projection replace the product by
$(1+t)^4(1+tz)^6$. This gives the displayed table. The image equals
the kernel of the projection by \cref{thm:v20-homotopy}, so subtracting
dimensions proves the assertion. No numerical rank test is required.
\end{proof}

For example $\eta^{11}z_{01}z_{23}$ and
$\sigma_0 z_{01}z_{02}z_{03}$ survive this projection. They show why
the derivative thresholds cannot be lowered for arbitrary rooted
polynomials. They are not asserted to occur in the native measured row.
The 930 residual coordinates are a remaining \emph{test space}, not
930 nonzero anomalies. A full differential can further identify or
exclude them, but that requires the coupled equations.

\section{Applying the primitive to the actual annular matching arrays}
\label{sec:v20-native-counterterm}

Specialize V19 to one minimal antifield and zero ordinary metric
background. Let $B_{\eta,a,L}$ and $B_{\sigma,a,L}$ be the rooted
local densities corresponding to its actual Taylor assignments
$\mathsf T_4\iota^{(0)}_{a,L,I}$ in the $\eta cc$ and $\sigma ccc$
slots. They are not new arbitrary parameters. At this antifield degree
V19's marked expression is precisely $\frac12\Tr(C_g\mathscr E_0)$,
with \eqref{eq:v19-Eeta} and \eqref{eq:v19-Esigma}, including the two
ghost covariances in the latter. Passing from a Fourier Taylor
polynomial to $B$ replaces a momentum on each ghost leg by
$\partial/i$ and roots all derivatives on the ghosts. Use the same
Grassmann order as that trace.

Writing $B=B_\eta+B_\sigma$ and $B^{[d]}$ for derivative degree $d$,
\eqref{eq:v19-local-data} gives
\begin{equation}
 B_{a,L}^{[d]}=\chi^{-1}a^{d-4}\,b_{\boldsymbol\ell}^{[d]},
 \qquad 0\le d\le4,
 \qquad \boldsymbol\ell=(L_\nu/a)_\nu.
 \label{eq:v20-native-b}
\end{equation}
Every coefficient of $b_{\boldsymbol\ell}^{[d]}$ is the fixed
linear Fourier-to-jet conversion of
$\partial_z^\alpha\mathcal F_{I,\boldsymbol\ell}^{(0)}(0)/\alpha!$
from \eqref{eq:v19-scaling}. Thus it is a normalized finite sum of
\emph{the original measured annular integrand}. None of these
coefficient values is fitted by the primitive construction.

\begin{theorem}[A realized one-loop local source-action counterterm]
\label[theorem]{thm:v20-native-primitive}
Define
\begin{equation}
 \boxed{H_{a,L}=\mathfrak h B_{a,L},\qquad
        C_{a,L}=-H_{a,L},\qquad
        S^{\rm new}_{a}=S_a+\hbar C_{a,L}.}
 \label{eq:v20-counterterm}
\end{equation}
Then $C$ is a specified local ghost-number-zero, even source functional
of the form
\begin{equation}
 C_{a,L}=\int\bigl\{\eta^{\mu\nu}U_{\mu\nu,a,L}(q,c)
                 +\sigma_\mu V^\mu_{a,L}(q,c,c)\bigr\},
 \label{eq:v20-counterterm-shape}
\end{equation}
where $U,V$ are linear in $q$, and have at most three derivatives.
There is an exact identity
\begin{equation}
 \boxed{B_{a,L}+\gamma_0 C_{a,L}=\pi_{\rm aff}B_{a,L}.}
 \label{eq:v20-native-cancel}
\end{equation}
In particular both fourth-derivative source-curvature matching slots
are cancelled, and the third-derivative $\eta cc$ slot is cancelled.
These include \emph{all} finite $a^0$ matching terms of V19 in the
two specified zero-background slots. Other finite contacts generated
by the variation of $C$ are not claimed to vanish.

For $d\ge1$ the derivative-$d-1$ part of $C$ satisfies
\begin{equation}
 \|C_{a,L}^{[d-1]}\|_\mu
 \le C_d\chi^{-1}a^{d-4}\mu^{d-1}.
 \label{eq:v20-counterterm-bound}
\end{equation}
The degree-zero part of $B$ has no primitive supplied by this map.
For fixed profiles and completion parameters, the fourth-derivative
input gives a bounded three-derivative counterterm with a limit
determined by the corresponding annular integral. For every fixed $b>4$ its finite-volume coefficient error is
bounded by $C_b\chi^{-1}\mu^3(a/L_{\min})^b$.
\end{theorem}
\begin{proof}
Apply \cref{thm:v20-homotopy} to the actual rooted coefficient array.
Replacing one ghost factor by its even metric partner reduces ghost
number by one and reverses parity, so $C$ has ghost number zero and
is even. It has the asserted field content and derivative degree.
Equation \eqref{eq:v20-native-cancel} follows without any consistency
hypothesis on the other breaking components. The high-derivative
claims follow from \cref{cor:v20-top-degree} and
\eqref{eq:v20-native-b}. The coefficient bound follows from the
finite pair norm and V19's bounded dimensionless arrays. Finally
$\mathfrak h$ is a fixed bounded linear map on those arrays. It
commutes with their normalized finite sums and with the limiting
annular integral. Applying \eqref{eq:v19-volume} proves the stated
coefficient comparison.
\end{proof}

The $a^{-3}$, $a^{-2}$ and $a^{-1}$ coefficients permitted by
\eqref{eq:v20-counterterm-bound} are local source renormalizations,
not uniformly small errors. The fourth-derivative input is finite,
but its numerical value and that of its primitive can depend on the
completion profile and matching convention. The theorem specifies
these values by the inherited finite sums; no numerical quadrature of
them is claimed in this installment.

If only the certainly removable top part is desired, set
\begin{equation}
 B^{\rm top}=B_\eta^{[3]}+B_\eta^{[4]}+B_\sigma^{[4]},
 \qquad C^{\rm top}=-\mathfrak hB^{\rm top}.
 \label{eq:v20-top-primitive}
\end{equation}
This restricted counterterm has coefficients growing at most as
$a^{-1}$ and cancels those three slots exactly. It is sometimes useful
for retaining a smaller leading-power matching budget. The full map
\eqref{eq:v20-counterterm} also removes the paired parts at lower
derivative degrees, while keeping their affine remainders visible.


\begin{proposition}[Rooted norm of the actual local counterterm]
\label[proposition]{prop:v20-kernel}
Put a fixed smooth window on every external leg, at scale $\mu$ strictly
inside the inherited soft domain. The derivative-$d-1$ primitive has
$M=8$ relative dimensions in the $\eta qc$ sector and $M=12$ in the
$\sigma qcc$ sector. Its Fourier kernel obeys, for every fixed $b$,
\begin{equation}
 a^M\sum_{\boldsymbol x}
 \left(1+\mu\sum_i d_{\rm per}(x_i,0)\right)^b
 \|K_{a,L}^{C,[d-1]}(\boldsymbol x)\|
 \le C_{d,b}\chi^{-1}a^{d-4}\mu^{d-1}.
 \label{eq:v20-counterterm-kernel}
\end{equation}
The norm is rooted: it gives one source in $L^1$ and the others in
$L^\infty$, not a total-volume bound on an unanchored functional.
\end{proposition}
\begin{proof}
The primitive is a finite polynomial in the external momenta. On the
window its scaled derivatives are bounded by
\eqref{eq:v20-counterterm-bound}; differentiation of the fixed window
has the same scaled bounds. Fourier integration by parts more than
$M+b$ times gives an integrable kernel bounded by the right-hand-side
scale times $\mu^M(1+\mu|\boldsymbol x|)^{-M-b-1}$. Sampling and
periodization, with $d_{\rm per}$ no larger than the distance of any
lift, prove the result for $a\mu\ll1$. V19 bounds the input coefficient
array independently of the periods, so no additional volume factor
enters this argument.
\end{proof}

\section{The free BV descendants cannot be suppressed}
\label{sec:v20-descendants}

Let $s_0=\gamma_0+\delta_0$ be the inherited free minimal BV operator.
The signs of its dual coordinates are fixed by the inherited
antibracket. To avoid imposing an independent sign convention, write
\begin{equation}
 \delta_0\eta=\mathcal E_0(q),\qquad
 \delta_0\sigma=\mathcal B_0(\eta),\qquad
 \delta_0q=\delta_0c=0.
 \label{eq:v20-KT}
\end{equation}
Here $\mathcal E_0$ is the signed free Euler operator, of order two
and size proportional to $\chi$, and $\mathcal B_0$ is the signed
adjoint of $R_0$, of order one. Both are defined by the same free
master action, not separately chosen operators. The native free
Noether identities give
$\mathcal E_0R_0=0$ and $\mathcal B_0\mathcal E_0=0$.
Thus $s_0^2=0$ in the ordinary local coefficient algebra. Nothing here
asserts $s_a^2=0$ for the full finite filtered classical action.

\begin{theorem}[An exact primitive packet, including its lower-antifield rows]
\label[theorem]{thm:v20-packet}
Write $H=\mathfrak hB=H_\eta+H_\sigma$ and put
$\mathcal D_B=\delta_0H$. Then
\begin{equation}
 \boxed{s_0H=B-\pi_{\rm aff}B+\mathcal D_B,
 \qquad s_0\bigl(B-\pi_{\rm aff}B+\mathcal D_B\bigr)=0.}
 \label{eq:v20-packet}
\end{equation}
The descendant of $H_\eta$ has no antifield, two metric fields, one
ghost, and derivative degree $d+1$ for an input of degree $d$.
The descendant of $H_\sigma$, after rooting by integration by parts,
has one $\eta$, one metric, two ghosts, and derivative degree $d$.
In particular the first fits the five-derivative local row of V17,
and the second fits the four-derivative, one-background row addressed
by V19. These are actual induced contacts, not optional choices of
normalization.
\end{theorem}
\begin{proof}
The first equation is \eqref{eq:v20-homotopy-identity} plus the
explicit complementary part of the differential. Applying $s_0$ and
using its nilpotence gives the second equation. To compute the
contacts, apply \eqref{eq:v20-KT} to the single root of each term of
$H$; its other factors are $q$ and ghosts and have zero $\delta_0$.
Replacing $\eta$ by $\mathcal E_0(q)$ introduces a second metric and
two derivatives. Replacing $\sigma$ by $\mathcal B_0(\eta)$ introduces
one derivative and changes its antifield type. Rooting that derivative
changes no total derivative degree. As $H$ already lost one derivative,
the asserted degrees follow. This computation retains the signed Euler
and adjoint operators rather than using a gauge-fixed inverse in their
place.
\end{proof}

If the full one-loop local breaking is denoted by $\mathcal A_{a,\rm loc}$,
adding \eqref{eq:v20-counterterm} changes its free-coefficient part by
\begin{equation}
 \boxed{\mathcal A_{a,\rm loc}\longmapsto
 \mathcal A_{a,\rm loc}
 -(B-\pi_{\rm aff}B)-\mathcal D_B.}
 \label{eq:v20-ledger-update}
\end{equation}
Consequently cancelling $B-\pi_{\rm aff}B$ \emph{alone} while leaving
the other local matching equations untouched is not the operation of
adding this source action. In particular a packet consisting only of
$B-\pi_{\rm aff}B$ need not be $s_0$-closed at all:
$ s_0(B-\pi_{\rm aff}B)=-s_0\mathcal D_B$.
This gives a direct consistency test, rather than assuming the local
curvature component is an anomaly cocycle by itself.

\subsection{Two explicit sign and contact examples}

Set $y=\partial_0^2c_1$, $v=\partial_1^2c_2$,
$X=\Gamma_{1;00}$, and $Y=\Gamma_{2;11}$. For
$B=\eta^{12}yv$, the exact primitive is
\begin{equation}
 H=-\tfrac12\eta^{12}(Xv-Yy),\qquad
 \gamma_0H=\eta^{12}yv,\qquad
 \delta_0H=-\tfrac12\mathcal E_0^{12}(q)(Xv-Yy).
 \label{eq:v20-eta-example}
\end{equation}
The minus sign in $H$ comes from crossing the odd root. There are two
paired factors, so the divisor is two even though the input has four
derivatives. The displayed descendant is a five-derivative
$q^2c$ contact.

For $B=\sigma_0 y\,y_{13}z_{23}$ put $X_s=q_{13}/2$. Then
\begin{equation}
 H=\tfrac12\sigma_0(Xy_{13}-X_s y)z_{23},\qquad
 \gamma_0H=B,\qquad
 \delta_0H=\tfrac12\mathcal B_{0,0}(\eta)
                       (Xy_{13}-X_s y)z_{23}.
 \label{eq:v20-sigma-example}
\end{equation}
The last expression is subsequently integrated by parts onto the same
root convention. Both examples are exact polynomial checks of the
primitive map. Their coefficients are not asserted to be the numerical
native annular coefficients; the map applies linearly to those actual
coefficients without changing them.

\section{Nonlinear transport, finite-mesh spill, and measured ownership}
\label{sec:v20-transport}

The preceding free packet is not the whole nonlinear variation. The
ordinary zero-mesh source in the inherited chart is known explicitly:
\begin{equation}
 r_\infty^q=(I+\mathsf B_q)^{-1}
       \{R_0c+\mathcal L_cq+\mathcal L_c\mathsf Q(q,q)\},
 \qquad r_\infty^c=c\cdot\partial c.
 \label{eq:v20-full-source}
\end{equation}
At every prescribed finite field degree its coefficients and their
adjoints are finite polynomials by V17's recursion. Define $s_\infty$
by this source and the same zero-mesh stationary action. Then
\begin{equation}
 s_\infty H=(B-\pi_{\rm aff}B)+\mathcal D_B+
                    \mathcal T_B,
 \qquad \mathcal T_B=(s_\infty-s_0)H.
 \label{eq:v20-nonlinear-contacts}
\end{equation}
This formula is an explicit finite coefficient prescription: act on
$H$ with \eqref{eq:v20-full-source}, its antifield adjoints, and the
native stationary-action derivatives. It does not solve a new inverse
problem. Every monomial of $\mathcal T_B$ has at least one ordinary
metric field. The antifield-free terms of $\mathcal T_B$ have at least
three ordinary metric fields. Thus neither $\mathcal T_B$ nor
$\mathcal D_B$ changes the zero-background one-antifield cancellation.
At one background their nonzero source contacts must be retained.

\begin{proposition}[Precise finite-panel transfer to the existing regulator]
\label[proposition]{prop:v20-spill}
On the V19 soft domain, the zero-background, one-antifield projection of
the actual finite variation of $C$ equals
$-(B-\pi_{\rm aff}B)$ exactly. In the antifield-free $q^2c$ descendant
row, the difference between the finite free operator and
\eqref{eq:v20-KT} has the bound
\begin{equation}
 \left|\partial_K^\alpha
 [ (s_a-s_\infty)H_\eta]_{q^2c}\right|
 \le C_\alpha a\mu^{6-|\alpha|}
 \quad(|K|\le\mu),
 \label{eq:v20-spill}
\end{equation}
for every fixed derivative order, with the source norms understood.
The corresponding bound for $H^{\rm top}$ is
$C_\alpha a^3\mu^{8-|\alpha|}$. This is not an assertion that every
higher nonlinear row has the same small bound.
\end{proposition}
\begin{proof}
All external subset momenta entering these finite local densities lie
inside the plateau by \eqref{eq:v19-domain}. V18's source therefore
agrees with the ordinary coordinate source there, including the required
product contacts. At zero ordinary background the only part of the
variation which removes the single $q$ from $H$ is $\gamma_0$.
Action/Euler and antifield-adjoint terms have an ordinary metric factor
left. This proves the first assertion without finite high-mode nilpotence.

For the $q^2c$ row the sole dispersion difference is the signed free
Euler matrix $E_{2,a}=H_a^{\rm cl}-H_0^{\rm cl}$, of order
$\chi a^4|K|^6$ with fixed differentiated versions by V16. The degree
$d$ input contributes to $H_\eta$ at order
$\chi^{-1}a^{d-4}\mu^{d-1}$. Their product is bounded by
$a^d\mu^{d+5}$. Only $d\ge1$ is paired. Summing $1\le d\le4$ and
using $a\mu\ll1$ proves \eqref{eq:v20-spill}. Fixed momentum
derivatives follow by the same product estimates on a scaled window.
For the top packet $d\ge3$, the first power is instead
$a^3\mu^8$. No inverse lower momentum enters either estimate.
\end{proof}

There is a useful warning for the next interaction degree. The native
cubic artifact is $O(\chi a^3\mu^5)$, whereas a possible $d=1$
primitive has coefficient $O(\chi^{-1}a^{-3})$. Their product is only
$O(\mu^5)$, not necessarily small. Its limiting piece must be kept as
an additional local matching term in the $q^3c$ row. Restricting to
\eqref{eq:v20-top-primitive} avoids this particular order-one product,
but cannot replace a calculation of the complete remaining local row.
This explains why bounded inversion of a local coefficient map and
arbitrary interacting RG control are different questions.

At one loop the measured ownership is unambiguous. In the formal
Laplace convention, insert $\hbar C$ in the \emph{same} source action.
The one-loop BV expression changes by $(S_a,C)$; the term
$\hbar^2\Delta C$ belongs to the next loop order. At zero antifields
$C=0$, so the zero-antifield action, free covariances, auxiliary connection
stationary section and its density are unchanged. This does not make
the transformation insertions unchanged. Their finite local
normalizations are precisely the new coefficients of $C$.
The original reference insertion and density remain counted once.
Any exact measured identity is recomputed with the complete changed
source action, not with a change on one side alone.

\section{A concrete remaining compatibility problem, not an anomaly claim}
\label{sec:v20-remainder}

This installment supplies an actual local counterterm for the
longitudinally exact part of V19's one-antifield, zero-background
matching array. The residual test is now smaller and explicit:
\begin{equation}
 \boxed{\pi_{\rm aff}(\mathcal A_{\rm combined})\quad\hbox{and}\quad
 \mathcal A_{\rm combined}-s_\infty\mathfrak hB,}
 \label{eq:v20-next-test}
\end{equation}
where $\mathcal A_{\rm combined}$ must contain the action-breaking,
reference, nonminimal and measure contributions in the \emph{same}
source convention, together with the curvature part used here.
The projection in the first expression is taken only on the specified
rooted zero-background slots. It must not be applied to the full
functional while discarding its other components.

In particular, the 930-dimensional affine-ghost quotient is neither
$H^1(s\mid d)$ nor a proof that its populated coefficient is nonzero.
The one-background consistency row also has linearized curvature and
antifield-adjoint contributions not present in the simple quotient.
Its ghost-number-zero primitive, if required, must reproduce the
contacts \eqref{eq:v20-packet} and
\eqref{eq:v20-nonlinear-contacts}, not cancel a selection of rows
independently. The general BRST classification and the modified
Slavnov framework \cite{V20BBH1995,V15IIM} do not establish these
regulator-specific equalities by themselves.

The construction respects the programme's separation between local
restoration and a full filtered ESM continuum. It does not modify the
classical spacetime theorem, prove a chiral line theorem, remove the
positive lower scale, or construct an invariant analytic RG class.
It does show explicitly which finite matching terms already have a
bounded pole-free primitive, where that primitive lives in the action,
and which previously retained rows it necessarily changes.

\begingroup
\renewcommand{\refname}{Additional methodological references for V20}
\begin{thebibliography}{99}
\bibitem[42]{V20BBH1994}
G.~Barnich, F.~Brandt and M.~Henneaux,
\emph{Local BRST cohomology in the antifield formalism: I. General theorems},
Communications in Mathematical Physics \textbf{174} (1995), 57--92,
arXiv:hep-th/9405109. General local cohomology and antifield framework;
not an imported native matching or norm theorem.
\bibitem[43]{V20BBH1995}
G.~Barnich, F.~Brandt and M.~Henneaux,
\emph{Local BRST cohomology in Einstein--Yang--Mills theory},
Nuclear Physics B \textbf{455} (1995), 357--408,
arXiv:hep-th/9505173, doi:10.1016/0550-3213(95)00471-4.
The treatment of abelian factors is relevant to the eventual ESM
extension; no anomaly conclusion for the present finite regulator is
imported from its classification.
\end{thebibliography}
\endgroup

\section{Verification and closing ledgers for V20}
\label{sec:v20-ledgers}

The standalone script \srcid{check_ESM_local_primitive_V20.py} verifies
all 270 native ghost-jet sections using exact rational coefficients.
It tests the primitive identity exhaustively on every independent
rooted ghost monomial type through degree four for the two sectors,
including the odd-root sign, and reproduces the component counts.
It also checks the two explicit primitive examples, the free
Euler/Noether symbol, coefficient scaling and the finite-mesh spill
powers. These are finite algebraic tests of the construction, not
numerical evaluations of the V19 annular integrals, a certification
of their analytic proofs, or a new Lean formalization.

\subsection*{Proved}
\addcontentsline{toc}{subsection}{V20: Proved}
The actual free gravitational jet map has the explicit local section
\eqref{eq:v20-Christoffel-section} and the bounded primitive
\eqref{eq:v20-homotopy}. Its rooted one-antifield quotient through
four derivatives has the exact dimensions in \cref{prop:v20-count}.
Applied to V19's actual finite matching arrays, it gives the
one-loop ghost-number-zero source counterterm
\eqref{eq:v20-counterterm} and the slotwise cancellation
\eqref{eq:v20-native-cancel}. All fourth-derivative local terms in
these two slots, and the third-derivative metric-antifield terms,
are included. The full free-BV descendant packet is explicit, as
are its nonlinear source transport and the stated lower-row mesh
spill. Counterterm norms, their normalized volume bounds and the
finite top-degree matching limits have been derived.

\subsection*{Inherited}
\addcontentsline{toc}{subsection}{V20: Inherited}
V1--V19 and their labels are unchanged. V19 supplies the marked
annular integrands, exact finite local coefficients and their source
and volume bounds. V14--V18 supply the completed regulator, ordinary
coframe chart, contour, filtered source, stationary action, minimal
antifield module, reference and auxiliary density. The native free
Noether identities and V17's higher classical recursion are used at
their original scopes. The prior abstract splitting criterion and
measured insertion calculus are not renamed as new results.

\subsection*{Literature input}
\addcontentsline{toc}{subsection}{V20: Literature input}
Contractible jet pairs, longitudinal/Koszul--Tate separation, and
local BRST cohomology are established. Their use here is constructive
and coefficient specific. No external cohomology theorem is used
to set the native combined anomaly to zero or to infer a uniform
inverse for an unconstructed full ESM complex.

\subsection*{Conditional}
\addcontentsline{toc}{subsection}{V20: Conditional}
A complete local quantum restoration requires the other components
of the combined breaking to match the displayed descendants and
the surviving affine-ghost coefficients. It also requires the
nonlinear and nonminimal consistency equations and common physical
normalizations. Those coefficients have not been evaluated by this
primitive operation. It is not a counterterm which is already proved
to preserve every earlier physical normalization or to restore the
entire nonlinear master equation.

\subsection*{Open}
\addcontentsline{toc}{subsection}{V20: Open}
Evaluate the actual affine-ghost projection of the \emph{combined}
local one-antifield row, and its one-background descendant equation,
using \eqref{eq:v20-ledger-update} rather than repeating an insertion
subtraction. The curvature contribution to the finite top derivative
matching slots now has an explicit primitive. The remaining coupled
local primitive, higher-antifield/background closure, mixed SM and
chiral transport, higher-loop estimates and the invariant filtered
ESM class remain open. No nonperturbative or finite-parameter gravity
conclusion is drawn.

\begin{center}\small
\begin{tabular}{>{\raggedright\arraybackslash}p{.23\textwidth}>{\raggedright\arraybackslash}p{.68\textwidth}}
\toprule
Variable & Uniformity and limitation in V20\\
\midrule
Mesh and volume & The primitive map is a fixed rational jet map. Input
coefficients have the explicit powers $\chi^{-1}a^{d-4}$. Its operator
norm, not every unscaled counterterm, is uniform in mesh and volume.
V19's size and plateau conditions remain in force.\\
Source degree & One minimal antifield, zero ordinary background in the
input, two or three ghosts, derivatives through four. The counterterm
has one ordinary metric field. No arbitrary-source-order bound.\\
Symmetry & Exact longitudinal cancellation and full free-BV descendant
identity. This is not a full nonlinear cohomological splitting or
an all-mode finite quantum master equation.\\
Loops and couplings & One-loop source counterterm; $\chi>0$ fixed.
The map uses no inverse momentum or inverse cutoff filter. The
zero-antifield measure and action remain the same.\\
Locality & Finite differential polynomials on the smooth core. On a
fixed soft window, Fourier sampling/periodization has rooted bounds.
This does not turn spectral derivatives into fixed-range lattice stencils.\\
Iteration & No accumulated shell-count factor is introduced by applying
one bounded linear map to V19's assembled local array. Preservation of
an interacting RG neighbourhood is not proved.\\
\bottomrule
\end{tabular}
\end{center}


\clearpage
\part*{V21: The affine residual differential and a measured six-coefficient test}
\addcontentsline{toc}{part}{V21: Affine residual differential and measured compatibility}

\section{Replace the raw affine quotient by its actual consistency complex}
\label{sec:v21-start}

Sections 1--160, including the qualifications in their ledgers, are
preserved. V20 constructed a longitudinal primitive and left a
930-dimensional affine-ghost coefficient space. That space was explicitly
not identified with anomaly cohomology. The next step is to determine the
\emph{induced differential} on it, keeping the coupling between the metric
and ghost antifields. This is smaller and more informative than evaluating
930 unrelated integrals or adding another source prescription.

The result below is an explicit finite compatibility reduction. In the
rooted, one-minimal-antifield affine quotient, the actual zero-mesh
source induces integer matrices
\[
 \mathscr C^0\xrightarrow{\mathscr D_0}\mathscr C^1
             \xrightarrow{\mathscr D_1}\mathscr C^2,
 \qquad (\dim\mathscr C^0,\dim\mathscr C^1,\dim\mathscr C^2)
             =(280,930,2040).
\]
Their ranks are 231 and 693. There are six, rather than 930, residual
classes after imposing this consistency equation. Fixed rational maps
satisfy a contraction identity with coefficient norm bounds 49 and 33.
The six classes are evaluated by constant translation ghosts. For the
one-loop minimal-source module and its measured reference/breaking
balance, these six coefficients vanish by an exact zero-momentum
Grassmann cancellation, before integration. The remaining term is an
explicit bounded image of the \emph{consistency defect}. It is not
silently set to zero.

No higher ghost source is chosen in this installment. The completed
regulator, the source of V18--V19, the actual ghost action of V14, the
positive lower scale, and the once-counted auxiliary measure remain
unchanged. The finite computation below concerns the ordinary zero-mesh
local differential, which is also the external coefficient differential
on the soft plateau. It does not manufacture a nilpotent finite
high-frequency source. Standard local BRST methods distinguish precisely
such a residual complex from full local cohomology
\cite{V20BBH1994,V20BBH1995}; those references do not supply the particular
matrices, norm bounds or measured zero-momentum test calculated here.

\subsection{What will and will not be concluded}

Let $B$ be a rooted ghost-number-one local coefficient in the two
zero-background slots $q^*cc$ and $c^*ccc$, with derivatives through four.
The new construction gives a source-action counterterm and proves
\[
 [B+s_\infty C_B]_{\text{one antifield},\,q=0}
 =\text{six-coefficient term}
       +\text{a specified image of }\mathscr D_1 B_{\rm aff}.
\]
Here $s_\infty$ is the ordinary local differential of the inherited
stationary action and its chart-correct source. A closed input with zero
six-coefficient term has a primitive in these slots. A nonclosed input
has the displayed residual. The fact that a curvature insertion is
local, or that its original classical expression is a square of a
source, does not prove that its \emph{isolated quantum local coefficient}
is a closed input. The reference and action-breaking terms must still
be assembled consistently. This limitation is part of the theorem, not
an omitted final step.

\section{Derivation of the coupled affine differential}
\label{sec:v21-affine-differential}

Work in the Euclidean coefficient convention of V18--V20. At a point
write an affine Killing ghost as
\begin{equation}
 c^\mu(x)=u^\mu+Z^\mu{}_{\nu}x^\nu,
 \qquad Z_{\mu\nu}=-Z_{\nu\mu}.
 \label{eq:v21-affine-ghost}
\end{equation}
The four $u$ and six independent $Z$ entries are odd. This is a test of
local jets, not a claim that a nonperiodic polynomial is a periodic
history. Relative to \eqref{eq:v20-first-jets},
$Z^\mu{}_{\nu}=z_{\nu\mu}=-z_{\mu\nu}$ on this quotient. This fixed
sign conversion is kept in the checker.

Define even formal total derivatives $\partial^K_\nu$ on this exterior
algebra by
\begin{equation}
 \partial^K_\nu u^\mu=Z^\mu{}_{\nu},\qquad
 \partial^K_\nu Z^\mu{}_{\rho}=0.
 \label{eq:v21-affine-total}
\end{equation}
The ordinary ghost law $sc=c\cdot\partial c$ restricts to
\[
 d_Ku=-Zu,\qquad d_KZ=-Z^2.
\]
The order in these formulas matters: $u^\nu Z^\mu{}_{\nu}
=-Z^\mu{}_{\nu}u^\nu$. Subtracting transport gives an odd derivation
\begin{equation}
 \mathfrak d:=d_K-u^\nu\partial^K_\nu,
 \qquad \mathfrak d u=0,\quad \mathfrak d Z=-Z^2.
 \label{eq:v21-bar-d}
\end{equation}
Here $\mathfrak d$ is only notation for this finite algebra map; it is
not a spacetime differential.

For symmetric-matrix and vector coefficients, respectively, put
\begin{align}
 (d_TF)_{\mu\nu}
 &=\mathfrak d F_{\mu\nu}
   +Z_{\mu\rho}F_{\rho\nu}+Z_{\nu\rho}F_{\mu\rho},\
 (d_VG)^\mu&=\mathfrak d G^\mu+Z^\mu{}_{\rho}G^\rho,\
 (\mathcal PG)_{\mu\nu}
 &=\partial^K_\mu G_\nu+\partial^K_\nu G_\mu.
 \label{eq:v21-tensor-maps}
\end{align}
All $Z$ coefficients are on the left; this fixes the signs even when
$F$ is odd. Set
\begin{equation}
 \mathscr C^g=
 \bigl(\operatorname{Sym}^2\R^4\otimes\Lambda^{g+1}\R^{10}\bigr)
 \oplus\bigl(\R^4\otimes\Lambda^{g+2}\R^{10}\bigr).
 \label{eq:v21-spaces}
\end{equation}
An exterior power with degree outside $0,\ldots,10$ is zero.

We use vector-symbol coordinates for the rooted Hamiltonians. With
independent symmetric entries and their canonical dual antifields, the
conversion is
\begin{equation}
 \mathcal A_g(F,G)=(-1)^g\langle\eta,F\rangle
                         +\langle\sigma,G\rangle.
 \label{eq:v21-Hamiltonian-sign}
\end{equation}
This is the fixed Koszul conversion to the evolutionary field generated
by a degree-$g$ Hamiltonian. It gives the inherited plus signs in the
classical minimal source at $g=0$. Equivalently all the following
identities may be checked as commutators of evolutionary fields, and
then converted back by \eqref{eq:v21-Hamiltonian-sign}. In particular,
an input rooted coefficient $+\eta T$ has $F=-T$ at ghost number one;
no physical source sign is changed by this convention.

\begin{theorem}[The induced affine consistency operator]
\label{thm:v21-induced}
Let $\mathscr Q$ mean: take the coefficient linear in one minimal
antifield, put ordinary metric and nonminimal fields to zero, root by
integration by parts, and restrict the ghost jets to
\eqref{eq:v21-affine-ghost}. The ordinary local minimal differential
induces
\begin{equation}
 \boxed{\mathscr D_g(F,G)=
       \bigl(d_TF+(-1)^g\mathcal PG,\ d_VG\bigr).}
 \label{eq:v21-D}
\end{equation}
Precisely, $\mathscr Q(s_\infty A)=\mathscr D_g\mathscr Q(A)$ in the
vector-symbol convention, whenever the projection has the stated
single-root meaning. Components with more metric fields cannot enter
this projection; components with more antifields cannot enter it
through a stationary zero-background Euler row. The same coefficient
identity applies to the finite external soft plateau. Moreover
$\mathscr D_{g+1}\mathscr D_g=0$.
\end{theorem}
\begin{proof}
For the ordinary chart source, $R_0c=0$ on an affine Killing jet.
At $q=0$ its first metric derivative is therefore $\mathcal L_c$:
the coframe-chart correction contains $R_0c$ and vanishes. Higher
ordinary metric factors and their derivatives have zero variation in
this restriction. The zero-background action Euler row is zero on the
flat branch with its recorded vanishing cosmological coefficient.
Thus the remaining calculation is the cotangent lift of the actual
source, not the gauge-fixed Hessian.

Let $X$ be the evolutionary vector symbol with components $(F,G)$
and parity $g+1$. In $[s,X]^q$, variation of the metric through $F$
subtracts its Lie transport and its two index transformations;
variation of the ghost through $G$ supplies
$-(-1)^{g+1}R_0G$. Rooting the transport derivative moves it onto $F$.
The coefficient differential is consequently
$d_KF-u\cdot\partial^KF-Z^TF-FZ$ with $Z$ moved to the left, which
is $d_TF$ in components, and the additional term is
$(-1)^g\mathcal PG$. Similarly $[s,X]^c$ is
$d_KG-u\cdot\partial^KG+ZG=d_VG$. This proves
\eqref{eq:v21-D}, including its off-diagonal term. The latter is the
rooted signed adjoint-generator contribution from the ghost antifield;
it cannot be omitted. The Euler and higher-antifield observations
above prove the projection statement.

For a direct check of nilpotence, $\mathfrak d Z=-Z^2$ gives
$d_T^2=d_V^2=0$ by the vector and symmetric-tensor representations.
Furthermore
$[\mathfrak d,\partial^K_\mu]=Z^\rho{}_{\mu}\partial^K_\rho$;
substitution in \eqref{eq:v21-tensor-maps} gives
$d_T\mathcal P=\mathcal P d_V$. The signs $(-1)^g$ in consecutive
degrees are opposite, so the two cross terms cancel. These are finite
polynomial identities and use no nonlocal inverse or high-frequency
product rule. On a fixed external plateau all the filter derivatives
vanish and the same local coefficient calculation is literal.
\end{proof}

The subtraction of the transport term in \eqref{eq:v21-bar-d} is
essential. Using $d_K$ in its place would calculate a different complex.
Nor may the 450 metric-antifield and 480 ghost-antifield coefficients
be solved independently. For example take, in degree zero,
$G^\mu=\delta^\mu_2u^0u^1$, $F=0$. Its vector image alone is
$d_VG$. Although $d_V^2G=0$, the degree-one pair $(0,d_VG)$ has a
nonzero metric image: its $00$ component contains
\begin{equation}
 2u^0Z_{01}Z_{02}.
 \label{eq:v21-coupling-witness}
\end{equation}
The omitted $\mathcal PG$ is exactly what cancels it in
$\mathscr D_1\mathscr D_0(0,G)$. The checker verifies this integer
coefficient separately.

\section{An exact six-dimensional reduction with bounded rational maps}
\label{sec:v21-ranks}

Order the odd generators as
\[
 u^0,u^1,u^2,u^3,Z_{01},Z_{02},Z_{03},Z_{12},Z_{13},Z_{23}.
\]
Use increasing exterior monomials, lexicographic independent symmetric
entries, and then vector entries. This completely specifies the matrices
in \eqref{eq:v21-D}. Their entries are integers. Write
$F^{(0)}(u)$ for the part of $F$ containing no $Z$. Define
\begin{equation}
 \boxed{\Pi(F,G)=
  \left(\frac14I_4\operatorname{tr}F^{(0)}(u),\ 0\right).}
 \label{eq:v21-Pi}
\end{equation}
On $\mathscr C^1$ its range has basis
\begin{equation}
 U_{ij}=(I_4u^iu^j,0),\qquad 0\le i<j\le3.
 \label{eq:v21-six}
\end{equation}
Their root-left action representatives have the additional minus sign
in \eqref{eq:v21-Hamiltonian-sign}. It has no effect on a vanishing test.

\begin{theorem}[Exact finite affine reduction]
\label{thm:v21-contraction}
In the stated component norms,
\begin{equation}
 \boxed{\operatorname{rank}\mathscr D_0=231,\qquad
        \operatorname{rank}\mathscr D_1=693,\qquad
        \dim\frac{\ker\mathscr D_1}{\operatorname{im}\mathscr D_0}=6.}
 \label{eq:v21-ranks}
\end{equation}
There are explicitly constructed fixed rational matrices
$\mathsf H:\mathscr C^1\to\mathscr C^0$ and
$\mathsf Q_1:\mathscr C^2\to\mathscr C^1$ with
\begin{equation}
 \boxed{I_{\mathscr C^1}=
       \mathscr D_0\mathsf H+\Pi+\mathsf Q_1\mathscr D_1,
       \qquad \|\mathsf H\|_{1\to1}=49,\quad
              \|\mathsf Q_1\|_{1\to1}=33.}
 \label{eq:v21-contraction}
\end{equation}
Thus for an arbitrary, not necessarily closed, coefficient $B_{\rm aff}$,
\begin{equation}
 \|B_{\rm aff}-\mathscr D_0\mathsf H B_{\rm aff}-\Pi B_{\rm aff}\|_1
 \le33\|\mathscr D_1B_{\rm aff}\|_1.
 \label{eq:v21-defect-bound}
\end{equation}
A closed coefficient has an affine primitive if and only if its six
coefficients \eqref{eq:v21-Pi} vanish. Both $\mathsf H$ and
$\mathsf Q_1$ lower the number of $Z$ factors, hence the total derivative
degree, by one on their nonzero entries.
\end{theorem}
\begin{proof}
Here is the finite rational construction and certificate used for the
rank assertion and the norm estimate. It is also a reproducible
specification of the two matrices, rather than an appeal to a
pseudoinverse at varying cutoff.

Form the columns of \eqref{eq:v21-D} by exterior multiplication in the
displayed order. To reduce a column, use its least occupied row as pivot,
normalize the pivot to one, and subtract previously stored pivot
columns in increasing pivot-row order. Carry along the domain column
under every rational operation. For $\mathscr D_0$ this procedure
produces 231 pivots. Appending the six columns $U_{ij}$ produces 237
pivots. Their independence also follows directly because every image
of $\mathscr D_0$ has a $Z$ factor whereas the six $U_{ij}$ have none.
The 693 coordinate vectors at the nonpivot rows are a complement to
these 237 columns. Their images under $\mathscr D_1$ have 693 pivots.
Every remaining image column reduces in this span. All operations are
rational, not numerical rank tests.

Reducing a vector of $\mathscr C^1$ in the 237-column basis gives its
boundary preimage, six-class coefficients and a complement remainder.
The boundary preimage defines $\mathsf H$. Reduce a vector of
$\mathscr C^2$ against the 693 image pivots, carry its complement lift,
and apply $I-\Pi$ to this lift; this defines $\mathsf Q_1$. For a
vector already in $\operatorname{im}\mathscr D_1$, this is the unique
chosen complement lift followed by that projection. The identity
\eqref{eq:v21-contraction} follows from these two reductions.

For completeness the exact certificate supplied with this installment
records the pivot rows, input columns and rational pivots, and verifies
all 930 columns of \eqref{eq:v21-contraction}. It also verifies all
280 columns of $\mathscr D_1\mathscr D_0=0$ and all 930 columns of
$\mathscr D_2\mathscr D_1=0$. The resulting $\mathsf H$ has 1292
nonzero rational entries and $\mathsf Q_1$ has 7697. Taking maximum
absolute column sums gives exactly 49 and 33. These equalities are
checked using fractions; no floating-point singular values enter.
The complete column generator and reduction algorithm are contained in
\srcid{check_ESM_affine_compatibility_V21.py}, with no nonstandard
runtime dependency.

Since $\Pi\mathscr D_0=0$ and $\mathscr D_1\Pi=0$, the contraction
identity proves that the six representatives exhaust the cohomology.
A boundary cannot contain a nonzero combination of them. The derivative
grading follows either from the same columns or directly: every entry
of $\mathscr D$ increases the number of $Z$ factors by one, and the
reductions are block homogeneous in that number. Finally take norms in
\eqref{eq:v21-contraction}. This proves all the assertions, including
the nonclosed-input estimate.
\end{proof}

The number six is a property of \emph{this quotient complex}. It is
not the number of Einstein--SM anomalies and not a computation of
$H^1(s\mid\dd)$. In particular, possible obstructions to lifting a
quotient cocycle through ordinary-background or higher-antifield rows
are not decided by \eqref{eq:v21-ranks}. The map is also not inferred
from a claim of rotational invariance of the regulator: its tensor
operator is derived from the local source in
\cref{thm:v21-induced}, while the regulator may break that symmetry
at finite frequency.

\section{The native zero-momentum test eliminates the six curvature coefficients}
\label{sec:v21-measured-zero}

The six coordinates are the trace of the coefficient with two
\emph{undifferentiated constant ghosts}. They can therefore be tested
without evaluating a differentiated annular integral. Set the retained
ghost to a single formal constant odd vector $u$. All its four
components remain independent. Every source carries zero external
momentum, and the two hard legs at a quadratic vertex have momenta
$p,-p$. Define
\begin{equation}
 \alpha_p=i\,u\cdot p,
 \qquad \alpha_{-p}=-\alpha_p,\qquad \alpha_p^2=0.
 \label{eq:v21-alpha}
\end{equation}
The last equation is an exterior-algebra identity; it is not an
angular-average or parity argument.

\begin{lemma}[Actual ghost-action return for a constant retained ghost]
\label{lem:v21-constant-return}
For the unchanged ghost-action operand of \eqref{eq:v18-operands},
\begin{equation}
 B_u(p)x_p=-s(ap)^2\alpha_p F_p x_p,
 \qquad \beta_p=-C_c(p)B_u(p)x_p=\alpha_p V_p x_p,
 \label{eq:v21-beta-factor}
\end{equation}
with $V_p$ an even matrix independent of $u$. The optional upper ghost
interaction contributes zero to this expression. Every return through
$P\beta$ has the same factor $\alpha_p$.
\end{lemma}
\begin{proof}
The explicit native formula in V18 is
$B_c(p,k)=-s(a(p+k))s(ap)s(ak)F_{p+k}R_1^{(0),E}(x_p,c_k)$.
At $k=0$, $R_0u=0$, so the chart correction in $R_1$ vanishes and
$R_1(x_p,u)=\mathcal L_u x_p=\alpha_p x_p$. The upper added action
has a factor $(I-S)c$ and hence vanishes for $c=u$. Multiplication
by the actual, possibly non-scalar, $C_c(p)$ changes only the even
matrix. No inverse is replaced by a scalar here.
\end{proof}

\begin{theorem}[Zero constant-ghost coefficients of the measured curvature]
\label{thm:v21-curvature-zero}
Let $B^\Omega_{a,L}$ denote exactly V19's local one-loop coefficient of
\eqref{eq:v19-Omega}, at one minimal antifield and zero retained
ordinary metric. Then
\begin{equation}
 \boxed{[B^\Omega_{a,L}]_{\eta uu,\,K=0}=0,
 \qquad [B^\Omega_{a,L}]_{\sigma uuu,\,K=0}=0,
 \qquad \Pi\mathscr Q B^\Omega_{a,L}=0.}
 \label{eq:v21-zero-curvature}
\end{equation}
The cancellation holds for each loop momentum and every finite regulator
in the declared family, before integration or Taylor assignment. It is
uniform in volume and does not depend on the numerical values of the
annular matching integrals.
\end{theorem}
\begin{proof}
For constant $u$, $u\cdot Du=0$, and $\mathcal L_u=u\cdot D$
commutes with $P$. Consequently
$\mathcal K_u=P\mathcal L_u(I-P)\mathcal L_u=0$ on every metric
array, since its square contains $u^\mu u^\nu D_\mu D_\nu=0$.
Both direct coordinate terms in \eqref{eq:v19-E0} therefore vanish.

In its ghost-mediated term put $v=P\beta$. By
\eqref{eq:v21-beta-factor}, $v_p$ contains $\alpha_p$. Since
$Du=0$, $w_1=(P-I)(u\cdot Dv)$ contains a second $\alpha_p$ and
is zero. The remaining term in \eqref{eq:v19-Eeta} is
$P\mathcal L_v(I-P)\mathcal L_u x$. In a zero-external-momentum
quadratic pairing its two hard labels are $p,-p$; its two odd scalar
factors are $\alpha_{-p}$ and $\alpha_p$, whose product is zero.
Derivatives and matrix components do not alter that scalar factor.
The ghost-antifield operand \eqref{eq:v19-Esigma} vanishes because
$w_1=0$. Thus $\mathscr E_0$ is zero in both indicated sectors.
At one antifield the complete marked trace is
$\operatorname{Tr}(C_g\mathscr E_0)/2$ by V19, proving the result.
The six-class map reads only a subset of these zero coefficients.
\end{proof}

\subsection{The canonical row and the reference must remain distinguished}

There is a second exact use of the constant-ghost test. In the
zero-background minimal sector a retained constant ghost enters a
quadratic-hard transformation or ghost-action vertex through
$\alpha_p$; differentiating the transformation in a hard direction
can remove a ghost but cannot supply a new independent constant-ghost
momentum factor. The zero-antifield vertices have no retained ghost.
All hard lines of a zero-external-momentum one-loop word carry $p$ or
$-p$. Therefore a word with one $\eta$ and two constant retained ghosts
has two proportional odd factors and vanishes. The same argument
applies to the three-constant-ghost $\sigma$ word. For the reference
trace, the factor $C\mathbb R=\Theta$ is even and does not change
this counting. There is no singular inverse-cutoff step. This proves
\begin{equation}
 [\mathcal X_{a,\eta cc}]_{u u,\,K=0}=0
 \label{eq:v21-reference-zero}
\end{equation}
for the actual reference insertion in this minimal-source sector.
It does not say that its first or higher external derivatives vanish.

One can also read the canonical part without an integral. On these
soft local jets, projection of its one-antifield part is
$\mathscr D_0\mathscr Q(\Gamma_{a,1}^{\rm loc})$; this is simply
\cref{thm:v21-induced} applied to the actual local quantum source
coefficients of V18. Its six-class projection is zero since
$\Pi\mathscr D_0=0$. The classical master Hamiltonian is linear in
minimal antifields, so its BV Laplacian has no one-minimal-antifield
coefficient. The auxiliary density similarly has no such direct
coefficient, and a one-loop connection-density insertion cannot be
contracted into it without an extra loop or a forbidden hard-linear
branch, as in V18. Restoring the nonminimal pair introduces no
constant-ghost alternative to \eqref{eq:v21-alpha}.

Consequently the six coefficients of the \emph{complete measured
balance} in these specified slots vanish as well, when action,
reference and source are included with the same convention. The
corresponding statement for the action-breaking side follows from the
inherited finite integration-by-parts identity and
\eqref{eq:v21-reference-zero}; it is not inferred from an action Ward
identity. This is a zero-momentum statement only. At nonzero external
derivative degree the reference trace is retained, and the
reference-subtracted local row need not equal a boundary or be closed.
In particular this argument does not set the reference term of
V16--V17 to zero at fixed lower Wilsonian scale.

\section{Lift the affine calculation to a source-action counterterm}
\label{sec:v21-lift}

The affine primitive alone is not a complete local functional identity.
Its lift can produce paired ghost jets, and its ordinary-field
descendants must be retained. The following construction makes this
accounting explicit.

Let $\mathcal E$ extract the rooted, zero-ordinary-background,
one-minimal-antifield slots with unrestricted ghost jets through four
derivatives. Let $\mathfrak h$ be exactly V20's paired-jet primitive,
and let $\iota$ insert an affine polynomial by replacing $u=c$ and
$Z^\mu{}_{\nu}=(\partial_\nu c^\mu-\partial_\mu c^\nu)/2$.
It is a fixed finite differential-polynomial lift. Sign conversion is
always \eqref{eq:v21-Hamiltonian-sign}. For an actual input $B$ put
\begin{align}
 R&=\mathscr Q B,\
 A&=\iota\mathsf H R,\
 V_B&=\mathcal E(s_\infty A)
                  -\iota\mathscr D_0\mathsf H R,\
 H_B&=\mathfrak hB+A-\mathfrak hV_B.
 \label{eq:v21-lift-construction}
\end{align}
Here $\mathfrak hB$ is understood in V20's rooted sign convention.
The operation defining $V_B$ uses the \emph{actual} local source
\eqref{eq:v20-full-source}, its antifield adjoint, and the native
Euler operator. At this projection it is a finite operation; there
is no unknown inverse or fitted counterterm.

\begin{theorem}[Source-action reduction with a visible consistency residual]
\label{thm:v21-lift}
The polynomial $V_B$ has zero affine projection. The ghost-number-zero
counterterm $C_B=-H_B$, inserted as $\hbar C_B$ in the same source
action, satisfies
\begin{equation}
 \boxed{B+\mathcal E(s_\infty C_B)
       =\iota\bigl(\Pi R+\mathsf Q_1\mathscr D_1R\bigr).}
 \label{eq:v21-lift-identity}
\end{equation}
In particular, for V19's curvature coefficient the six term is zero.
For a genuinely consistent combined coefficient, $\mathscr D_1R=0$
then gives complete cancellation in these zero-background slots.
Without that consistency input, the residual in
\eqref{eq:v21-lift-identity} is retained with the bound
\eqref{eq:v21-defect-bound}. All counterterm coefficients are fixed
rational linear combinations of the actual input coefficients.
\end{theorem}
\begin{proof}
By \cref{thm:v21-induced},
$\mathscr Q\mathcal E(s_\infty A)=\mathscr D_0\mathsf H R$, so
$\mathscr QV_B=0$. V20 therefore gives
$\gamma_0\mathfrak hV_B=V_B$. In the zero-background projection the
only contribution of $s_\infty$ that removes the one ordinary metric
in $\mathfrak hB$ or $\mathfrak hV_B$ is $\gamma_0$. Hence
\[
 \mathcal E(s_\infty H_B)
   =B-\iota R+\iota\mathscr D_0\mathsf HR+V_B-V_B.
\]
Apply \eqref{eq:v21-contraction}. This proves
\eqref{eq:v21-lift-identity} for arbitrary inputs, including nonclosed
ones. The zero-momentum result supplies the assertion for the
curvature coefficient. Each factor in the construction is a fixed
finite coefficient operation; the inherited input coefficients, not
independent freely chosen numbers, enter the answer.
\end{proof}

This theorem does not infer the needed consistency equation from
$\Pi R=0$. The two tests are separate. If $B$ is extracted from an
entire combined local breaking $\mathcal A$, then
\begin{equation}
 \mathscr D_1R=\mathscr Q(s_\infty\mathcal A).
 \label{eq:v21-combined-consistency}
\end{equation}
It is this coefficient, with the actual reference and finite-source
breaking included, that must be evaluated or bounded. The classical
identity $[r,r^2]=0$ by itself is not a statement that an isolated
renormalized insertion of $r^2$ is closed in the unmodified ordinary
local differential. V19's nonlocal convergence theorem cannot replace
\eqref{eq:v21-combined-consistency}.

\section{New descendants enter the first as well as the second Slavnov row}
\label{sec:v21-descendants}

V20's primitive had one ordinary metric field. The affine part $A$ in
\eqref{eq:v21-lift-construction} has none. This changes which earlier
rows it consumes. In root-left coordinates write
\[
 A=\langle\eta,F_A(c,\partial c)\rangle
                       +\langle\sigma,G_A(c,\partial c)\rangle.
\]
At fixed input degree its coefficients contain only the specified
antisymmetric first ghost derivatives; every monomial has one or two
ghosts as appropriate. The complementary free-BV derivative is
\begin{equation}
 \delta_0 A
 =\langle\mathcal E_0(q),F_A\rangle
                     +\langle\mathcal B_0(\eta),G_A\rangle,
 \label{eq:v21-descendants}
\end{equation}
with V20's signed native operators. The second term is rooted by
integration by parts. It is precisely the off-diagonal contact
$\mathcal P$ in the residual complex, with the conversion
\eqref{eq:v21-Hamiltonian-sign}. The first term is an antifield-free
\emph{one-metric/one-ghost} contact, not a two-metric contact.

Thus this normalization can change the local row of V16, not merely
the row of V17. At input derivative degree $d$, its derivative order
is $d+1$. The two paired-jet primitives in
\eqref{eq:v21-lift-construction} generate the two-metric descendants
already described by V20. The nonlinear adjoint of the source supplies
the additional one-background algebra contacts. A complete ledger is
obtained by retaining the whole finite polynomial
\begin{equation}
 \mathcal T_B:=s_\infty H_B-\mathcal E(s_\infty H_B),
 \qquad
 \mathcal A\longmapsto\mathcal A-s_\infty H_B.
 \label{eq:v21-whole-ledger}
\end{equation}
Nothing in the construction authorizes holding those local equations
fixed while changing only the zero-background algebra row.

At zero antifields $C_B=0$. Consequently the zero-antifield action,
free propagators, auxiliary stationary section, Lorentz/Haar factors
and coordinate Jacobian are unchanged. The quantum transformation
and ghost-bracket source normalizations are changed. In the formal
one-loop convention the change in the quantum master expression is
$(S_a,C_B)$; the $\Delta C_B$ term carries the next loop power. The
reference continues to be counted once and is recomputed in the same
source convention. No positive quantum-gravity measure or Lorentzian
continuation is supplied by this local operation.

\section{Uniformity, finite-mesh transfer and the remaining measured equation}
\label{sec:v21-uniformity}

The fixed matrix norms in \eqref{eq:v21-contraction} are actual
cutoff-independent constants. They do not mean that all their inputs
are bounded without weights. For the known curvature input,
\eqref{eq:v19-local-data} has at one antifield
\[
 B^{[d]}_{a,L}=\chi^{-1}a^{d-4}b_{\boldsymbol\ell}^{[d]},
 \qquad0\le d\le4.
\]
The affine and paired primitives lower derivative degree by one.
The new affine part is absent at degree zero, and the measured
constant-ghost test removes that degree from the curvature input.
For each nonzero degree, therefore, its windowed coefficient has the
bound
\begin{equation}
 \|\mathsf H R^{[d]}\|_\mu
 \le49\chi^{-1}a^{d-4}\mu^{d-1}
                       \|b_{\boldsymbol\ell}^{[d]}\|_1,
 \qquad d\ge1.
 \label{eq:v21-scaled}
\end{equation}
Here the norm assigns one factor $\mu$ to each derivative. The input
and output root signs have unit norm. The complete lift has the same
power of $a$ and $\mu$, with an additional fixed coefficient constant
from V20's section and the finite source variation in
\eqref{eq:v21-lift-construction}. More explicitly it is bounded by
\[
 \|\mathfrak hB\|_1+49\|R\|_1
       +\tfrac32\|\mathcal E s_\infty\iota\mathsf HR
                          -\iota\mathscr D_0\mathsf HR\|_1
\]
in adapted coordinates. Each operator here is explicitly specified on
at most four ghost derivatives; its finitely many rational column
sums are independent of mesh, volume and shell count. This conclusion
comes from fixed coefficient maps, not an unproved uniform inverse
for an increasing-dimensional continuum complex.

Multiplying the output symbols by a fixed smooth external window and
integrating by parts in their Fourier transforms gives the same
scaled bound for every prescribed rooted kernel moment. The bound
has one $L^1$ source and the other source norms in $L^\infty$, with
no total-volume factor. It inherits exactly V19's domain
\eqref{eq:v19-domain}; no new infrared limit or unbounded-source-order
uniformity is claimed.

On that external plateau the source part of $s_a$ agrees with the
ordinary local source. Its free Euler difference is
$O(\chi a^4|K|^6)$ as in V20. Applied to an affine primitive of input
degree $d\ge1$, its first-row product is bounded by
\begin{equation}
 O(\chi a^4\mu^6)\,
 O(\chi^{-1}a^{d-4}\mu^{d-1})
       =O(a^d\mu^{d+5})=O(a\mu^6).
 \label{eq:v21-spill}
\end{equation}
The same estimate, with differentiated scaled norms, holds for each
fixed momentum derivative. At the next field degree the cubic
artifact can again multiply a degree-one primitive of size $a^{-3}$
and leave a finite local term. This term is retained in
\eqref{eq:v21-whole-ledger}; its vanishing is not concluded from
\eqref{eq:v21-spill}.

For a combined local input, the useful residual bound is now the
specific one
\begin{equation}
 \boxed{\|\text{unremoved affine coefficient}\|_1
       \le33\|\mathscr Q(s_\infty\mathcal A_{\rm combined})\|_1,}
 \label{eq:v21-final-test}
\end{equation}
when the six constant-ghost coefficients have been assembled in the
specified measured sector. If the consistency coefficient is zero,
there is a primitive in the stated slots. If it has an independently
proved vanishing cutoff bound, that bound transfers through 33. If
neither is known, \eqref{eq:v21-final-test} is a reduction of the
remaining problem, not its solution. In particular a fixed positive
lower Wilsonian scale still requires its modified reference identity.
There is no legitimate shortcut from the zero constant-ghost test to
the complete quantum Slavnov identity.

\section{Verification and closing ledgers for V21}
\label{sec:v21-ledgers}

The companion \srcid{check_ESM_affine_compatibility_V21.py} uses only
Python's standard library and exact rational arithmetic. It constructs
the matrices from the displayed ghost and tensor operations, verifies
both adjacent square-zero identities, checks all 930 columns of the
contraction, the six representatives, the norm bounds and derivative
grading. It also verifies the coupled-antifield witness and the
universal constant-ghost exterior factor. Its optional coefficient
compiler accepts an already affine-projected rooted array and returns
the primitive, six-class projection, consistency defect and unremoved
term separately. It does not invent the uncomputed combined coefficient
array or numerically evaluate an annular integral. The certificate file
records the rational elimination pivots so the rank computation can be
reproduced independently. No new Lean formalization is asserted.

\subsection*{Proved}
\addcontentsline{toc}{subsection}{V21: Proved}
The coupled affine differential \eqref{eq:v21-D} is derived from the
native ordinary source and its antifield adjoint. Its fixed matrices
have the exact ranks \eqref{eq:v21-ranks} and contraction
\eqref{eq:v21-contraction}, with coefficient norms 49 and 33. The
six representatives are undifferentiated translation-ghost trace
coefficients. The actual marked curvature coefficients vanish at
this zero-momentum test before integration; the same constant-ghost
factor and the canonical projection identify the zero six-class
coefficient of the specified measured minimal balance. The lifted
source-action construction \eqref{eq:v21-lift-identity} retains a
precise consistency residual and all induced first- and second-row
contacts. Its fixed-degree scaled bounds and first-row mesh spill
are explicit.

\subsection*{Inherited}
\addcontentsline{toc}{subsection}{V21: Inherited}
V1--V20 and their mathematical labels are unchanged. V18 supplies the
completed one-loop minimal-source integral and the actual cubic
ghost-action return. V19 supplies the marked curvature operands,
finite local coefficients and annular/source estimates. V20 supplies
the paired-jet primitive, its root signs and full-BV descendant
bookkeeping. The stationary action, coframe chart, physical reference,
nonminimal pair and once-counted measure retain their previous scopes.
No classical or source module is relabelled a full Einstein--SM shell.

\subsection*{Literature input}
\addcontentsline{toc}{subsection}{V21: Literature input}
Contractible pairs, local BRST cohomology and quotient consistency
complexes are established methods \cite{V20BBH1994,V20BBH1995}.
Modified reference identities are inherited with the role described
in \cite{V15IIM}. None of these inputs sets the native combined local
breaking to zero. The six-dimensional computation is for the
specified affine quotient, not a new classification of gravitational
anomalies.

\subsection*{Conditional}
\addcontentsline{toc}{subsection}{V21: Conditional}
Complete cancellation in the displayed zero-background slots follows
when the actual combined coefficient obeys
$\mathscr D_1\mathscr Q\mathcal A_{\rm combined}=0$, after the
reference, action-breaking and nonminimal conventions are kept the
same. Compatibility of the resulting counterterm with the ordinary
background and higher-antifield equations still requires the full
ledger \eqref{eq:v21-whole-ledger}. Those consistency equations are
not proved merely by this finite coefficient contraction. Physical
finite normalizations cannot be selected independently in each row.

\subsection*{Open}
\addcontentsline{toc}{subsection}{V21: Open}
The next smallest measured object is
$\mathscr Q(s_\infty\mathcal A_{\rm combined})$ with its actual
reference and finite-source curvature insertions, followed by the
first-background descendant equation for the lifted primitive. The
raw 930-dimensional coefficient space no longer needs to be treated
as 930 possible anomaly classes, and the six translation-ghost tests
are zero in the specified measured module. The remaining task is
consistency and lifting, not another source redesign or another
metric-only loop. Mixed internal-gauge and chiral sectors, a common
full ESM local projection, higher-loop control, infrared transport
and the invariant filtered interacting class remain open. No
nonperturbative quantum-gravity or finite-parameter ultraviolet
completion is concluded.

\begin{center}\small
\begin{tabular}{>{\raggedright\arraybackslash}p{.22\textwidth}>{\raggedright\arraybackslash}p{.69\textwidth}}
\toprule
Variable & Uniformity and limitation in V21\\
\midrule
Mesh and volume & The matrices and 49/33 operator bounds are independent
of both. Actual coefficients retain their stated powers of $a$;
rooted windowed norms have no extensive factor.\\
Derivative and source order & One minimal antifield, zero ordinary
background for the quotient, at most four derivatives in the input.
The lift creates the explicitly retained background contacts.\\
Symmetry & An exact ordinary local quotient differential and finite
coefficient contraction. No nilpotent high-frequency source or full
quantum master identity is inferred.\\
Loop and coupling & One-loop matching application; fixed $\chi>0$
and positive lower scale. The algebraic matrices contain no inverse
coupling or momentum.\\
Regulator & The existing smooth periodic completion and ghost action.
No upper-mode deletion, profile averaging, or new source prescription.\\
Iteration & Applying fixed coefficient maps introduces no shell-count
factor. An invariant interacting RG neighborhood is not proved.\\
\bottomrule
\end{tabular}
\end{center}



\clearpage
\part*{V22: The combined affine consistency coefficient is a lower-reference effect}
\addcontentsline{toc}{part}{V22: Combined consistency, lower-reference transport and local restoration}

\section{Assemble the input before solving another local equation}
\label{sec:v22-start}

Sections 1--168 are retained, with their original hypotheses and limitations.
The unfinished object in \eqref{eq:v21-final-test} was the consistency
coefficient of the \emph{combined} local breaking. It was not the rank of
another complex. The distinction between that coefficient and the isolated
curvature insertion is essential here. No action, cutoff profile, source
prescription, contour or measure is changed in this installment.

We use the V14 completed action, the V18 transformation source and the
V16 off-shell nonminimal convention. Work at one loop, with its common
factor $\hbar$ removed, at zero ordinary background and with one minimal
antifield. The relevant ghost-number-one slots are $\eta cc$ and
$\sigma ccc$, where $\eta=q^*$ and $\sigma=c^*$. The lower scale is denoted
by $\lambda=\Lambda_0>0$. A convenient common domain, slightly smaller
than the inherited domains, is
\begin{equation}
 \lambda\ge512\mu>0,\qquad a\lambda\le 1/32768,
 \qquad \lambda L_\nu\ge1.
 \label{eq:v22-domain}
\end{equation}
The lower profile $\Theta(p)=\theta(p/\lambda)$ equals one on
$|p|\le\lambda$ and zero on $|p|\ge2\lambda$, with fixed smooth
seminorms. Set $\nu=1-\Theta$. Small changes of these radii only change
fixed constants. Exact local consistency below uses thermodynamic
continuous-momentum coefficients. Rectangular finite-volume errors are
kept separately.

Let $\Gamma_{a,1}$ be the actual one-loop Legendre coefficient, with all
minimal sources retained and with the auxiliary density assigned once as
an action insertion. Only its already constructed minimal-source slots
are consumed. Let $\mathcal X_a$ be the actual reference insertion, not
its zero-mesh value. We use the measured one-loop convention of V16,
so the canonical expression has the form
\begin{equation}
 \mathcal A^{\rm comb}_a
   =(S_a,\Gamma_{a,1})+\operatorname{sdiv}r_a-\mathcal X_a.
 \label{eq:v22-combined-definition}
\end{equation}
On allowed finite grid momenta the change-of-variables identity identifies
this with the complete classical-breaking insertion and the recorded measure-variation term,
in the same source prescription. In particular, it is not a second copy
of the curvature insertion $\Omega_a$. A convention using the weighted
BV Laplacian moves the once-counted density term between the two sides;
it gives the same minimal-antifield extraction used below.

Denote the V21 rooted affine extraction by $\mathscr Q$. Taylor degrees
refer to the total degree in all independent external momenta, and are
taken \emph{after} the complete signed coefficient is formed. Define
\begin{equation}
 \begin{split}
 g_{a,L}&=\mathscr Q\,[\Gamma_{a,1}^{\rm loc}]_{\eta c,\,\sigma cc}
                 \in\mathscr C^0,\qquad
       \Gamma_{a,1}^{\rm loc}:=\mathsf T_3\Gamma_{a,1}^{\rm min},\\
 x_{a,L}&=\mathscr Q\mathsf T_4
                   [\mathcal X_a]_{\eta cc,\,\sigma ccc}
                 \in\mathscr C^1,\\
 R_{a,L}&=\mathscr Q\mathsf T_4
                   [\mathcal A^{\rm comb}_a]_{\eta cc,\,\sigma ccc}.
 \end{split}
 \label{eq:v22-three-arrays}
\end{equation}
The sign conversion is exactly \eqref{eq:v21-Hamiltonian-sign}.
The arrays in \eqref{eq:v22-three-arrays} are fixed by the measured
functional, not independent renormalization parameters. On a finite torus,
\eqref{eq:v22-combined-definition} also specifies the common interpolation
of the canonical expression: canonical coefficient operations are applied
to the same interpolants before extraction. It does not identify an
independently chosen off-grid interpolation of a breaking insertion with
that expression. The on-grid measured identity and thermodynamic local
normalization remain distinct.

\begin{proposition}[The ultraviolet local coefficients cancel from consistency]
\label{prop:v22-canonical-cancellation}
For the declared affine extraction and the actual common local assignment,
\begin{equation}
 \boxed{R_{a,L}=\mathscr D_0g_{a,L}-x_{a,L},\qquad
        \mathscr D_1R_{a,L}=-\mathscr D_1x_{a,L}.}
 \label{eq:v22-canonical-cancellation}
\end{equation}
The formulas also hold in thermodynamic normalization. They do not assert
$\mathscr D_1x_{a,L}=0$ at finite mesh.
\end{proposition}
\begin{proof}
On the external plateau, the V18 source is the ordinary source pulled
back by the same coframe chart. On an affine Killing ghost, $R_0c$ and
all its derivatives vanish. The metric Euler row is zero on the flat branch. The remaining
source-free ghost Euler operator is second order in spacetime derivatives and is zero
on an affine Killing jet; the nonminimal ordinary coordinates are zero.
Thus source-free Euler terms cannot lower a higher-antifield term into
this extraction without leaving an ordinary field factor. Thus the
rooted canonical coefficient is exactly the operator of
\cref{thm:v21-induced}, acting on the one-antifield part of
$\Gamma_{a,1}$. The antifield-free density and the flat divergence have
no one-minimal-antifield coefficient here. The nonminimal doublet is
restored before setting its ordinary coordinates to zero; it supplies
no extra term in this extraction.

Every nonzero entry of $\mathscr D$ increases total ghost-derivative
degree by one. Lie transport on a Fourier coefficient acts by a linear
momentum vector field and preserves its homogeneous Taylor degree before
this additional ghost derivative. Consequently Taylor degree four of
the canonical row consumes precisely degree three of the source action.
There is no contribution from a discarded higher Taylor jet. This proves
the first equality. Apply the exact finite identity
$\mathscr D_1\mathscr D_0=0$ to obtain the second. No absolute bound is
taken on the potentially divergent $g_{a,L}$ before this cancellation.
\end{proof}

This calculation goes upstream of V19's annular matching integrals. Their
large local coefficients occur in $R_{a,L}$, but the canonical part of
those coefficients has identically zero affine consistency. What remains
to control is a trace at the fixed \emph{lower} reference scale.

\section{The required reference coefficient is a finite lower-annulus trace}
\label{sec:v22-reference}

Use the full graded field list $\Phi=(q,b,c,\bar c)$ and the actual free
Hessian $\mathbb H_{a,0}$ of \eqref{eq:v16-offshell-action}. Set
\begin{equation}
 \mathbb C_a=\nu\mathbb H_{a,0}^{-1},\qquad
 V= D_\Phi^2(S_a-\tfrac12\Phi\mathbb H_{a,0}\Phi),\qquad
 J=D_\Phi r_a.
 \label{eq:v22-reference-objects}
\end{equation}
As in V16, a temporary $\nu_\epsilon>0$ can be used to define the
Legendre regulator. Its inverse-cutoff factor is removed algebraically:
$\mathbb C_{a,\epsilon}\mathbb R_{a,\epsilon}=(1-\epsilon)\Theta$.
The limit is taken only after this cancellation. The physical zero mode
has covariance zero and is never inverted.

At zero ordinary background, the only interaction Hessian blocks that
can enter the indicated source slots are obtained from
\begin{equation}
 \langle\bar c,B_cq\rangle,
 \quad\langle\eta,R_1(q,c)\rangle,
 \quad\langle\eta,R_2(q,c)\rangle,
 \quad\langle\sigma,c\cdot Dc\rangle.
 \label{eq:v22-vertex-list}
\end{equation}
Their respective retained-label types after two hard derivatives are
$c$, $\eta$, $\eta c$ and $\sigma$. The $R_i$ in this display mean
the actual V18 coefficients on the active momentum labels. The first
operand is the \emph{action's} ghost vertex, as in
\eqref{eq:v18-operands}. The Jacobian $J$ has a constant part $J_0$
and a part $J_c$ linear in the retained ghost. The constant nonminimal
block is included. No ordinary metric source is present in this section.

\begin{proposition}[Explicit compiler and support]
\label{prop:v22-reference-compiler}
With $\mathcal E_1$ extracting the two ghost-number-one minimal slots,
\begin{equation}
 \boxed{x_{a,L}=\mathscr Q\mathsf T_4\mathcal E_1
 \operatorname{STr}\!\left[
       \sum_{j=0}^{4}(-\mathbb C_aV)^j\Theta(J_0+J_c)
                         \right].}
 \label{eq:v22-reference-compiler}
\end{equation}
Every surviving word contains at least one $\mathbb C_a$. With the
momentum through $\Theta$ used as routing variable, its integrand and
all its external Taylor derivatives are supported in
\begin{equation}
             \lambda/2\le |p|\le3\lambda.
 \label{eq:v22-lower-annulus}
\end{equation}
On \eqref{eq:v22-domain}, all source filters and all low/upper completion
profiles in these words equal their interior values. The operands
$V,J_0,J_c$ are therefore exactly their zero-mesh polynomials. Only
$\mathbb C_a$ carries mesh dependence.
\end{proposition}
\begin{proof}
The full reference factor is
$(I+\mathbb C_aV)^{-1}\Theta J$. Each $V$ has at least one retained
label. There are three labels in $\eta cc$ and four in $\sigma ccc$;
$J$ supplies zero or one of them. Hence no power greater than four can
contribute. The term with $j=0$ has no antifield and vanishes under
$\mathcal E_1$. The Hessians in \eqref{eq:v22-vertex-list} exhaust the
remaining possibilities: higher metric action vertices leave an ordinary
metric factor, and higher ghost or source vertices do so as well. The
nonminimal coupling is quadratic and is already in $\mathbb H_{a,0}$.
This proves both completeness and the existence of a covariance in every
surviving word.

The line carrying $\Theta$ has momentum of norm at most $2\lambda$.
Every other line differs by a sum of at most four external labels. At
least one covariance contains $\nu$ and is zero below $\lambda$.
Since $4\mu\le\lambda/128$, these facts imply
\eqref{eq:v22-lower-annulus}; derivatives of $\Theta$ and $\nu$ do
not enlarge their closed supports. All internal labels and all partial
sums used at a vertex are bounded by a fixed multiple of $\lambda$.
The strict smallness of $a\lambda$ places them inside the common
plateau. Thus $B_c=-FR_1$ there, all completed upper interactions
vanish, and the transformation coefficients are precisely the ordinary
ones. The ghost Hessian at zero background and the metric free Hessian
are still the native corrected matrices, not their zero-mesh values.
\end{proof}

The support statement is stronger than a generic ultraviolet estimate.
It removes the Brillouin-face problem, the high--high upper-source
contacts, and an unbounded shell sum from \emph{this} coefficient.
It does not remove those phenomena from other coefficients. The matrix
$\Theta$ remains at its displayed side of every ordered source word.

\section{Why the zero-mesh reference coefficient is a cocycle}
\label{sec:v22-cocycle}

We record the finite trace argument being used, since declaring a
reference term to be closed merely because its integral is finite would
be circular. The underlying BV and modified Slavnov identities are
standard \cite{V15IIM,V20BBH1994}; the application concerns the actual
lower-annulus coefficient in \eqref{eq:v22-reference-compiler}.

\begin{lemma}[Finite one-loop reference consistency]
\label{lem:v22-finite-reference}
Let $S$ be an even, antifield-linear finite BV action satisfying the
classical master equation, with its nonminimal doublets included.
Let $s=(S,\cdot)$ in the convention of V16, and let $r=s\Phi$.
For a constant graded-symmetric reference matrix $R$, assume that
$H+R$, $H=D_\Phi^2S$, is invertible on the integrated coordinates.
Put
\[
 G_R=(H+R)^{-1},\qquad
 \Gamma_1^R=\tfrac12\operatorname{STr}\log(H+R),\qquad
 X_R=\operatorname{STr}(G_R R D_\Phi r).
\]
Then
\begin{equation}
       s\Gamma_1^R+\operatorname{sdiv}r=X_R,
             \qquad sX_R=0.
 \label{eq:v22-finite-cocycle}
\end{equation}
The same statement holds coefficientwise in nilpotent external source
parameters. A field-dependent physical density must be assigned once;
it is not absorbed into the constant $R$.
\end{lemma}
\begin{proof}
Differentiate the classical master identity twice in the integrated
field coordinates, using the same left/right convention as the graded
Hessian. Contract with $G_R/2$ and take the supertrace. The two
terms differentiating the source generator give
$\operatorname{STr}(G_R H D_\Phi r)$, whereas the remaining term is
$s\Gamma_1^R$. The terms in which an Euler derivative multiplies a
second derivative of $r$ are exactly the antifield-derivative part of
$s\Gamma_1^R$; omitting those terms would give only an on-shell formula.
The result is
\[
 s\Gamma_1^R=-\operatorname{STr}(G_R H D_\Phi r).
\]
Use $G_RH=I-G_RR$ to obtain the first equation of
\eqref{eq:v22-finite-cocycle}. Equivalently this is the finite
Gaussian/Berezin chain rule of \eqref{eq:v16-one-loop-chain}, now
including the ghost-antifield variations and using the classical master
identity rather than retaining its breaking insertion. This derivation
uses the complete ghost matrix and the nonminimal inverse.

The master equation gives $s^2=0$. The graded divergence identity
$\operatorname{sdiv}[r,r]=2r(\operatorname{sdiv}r)$ gives
$s\operatorname{sdiv}r=0$, since $r^2=0$. Apply $s$ to the first
identity. This proves the second without setting the reference trace
itself to zero. All operations are finite graded algebra.
\end{proof}

\begin{theorem}[Continuum lower-reference cocycle in the required quotient]
\label{thm:v22-reference-cocycle}
Let $x_{0,\infty}$ be the zero-mesh thermodynamic coefficient of
\eqref{eq:v22-reference-compiler}, with the same lower profile and
ordinary source coordinate. Then
\begin{equation}
               \boxed{\mathscr D_1x_{0,\infty}=0.}
 \label{eq:v22-reference-closed}
\end{equation}
Its six constant-ghost coefficients vanish. Thus it has an affine
primitive, but it need not be the zero coefficient.
\end{theorem}
\begin{proof}
The zero-mesh action used here is not a replacement gravity theory. It
is $\mathcal A_0$ of \cref{thm:v17-generator}, in the same coframe
coordinate, with its ordinary cotangent source and the nonminimal
harmonic gauge sector. The local source is the pullback of the ordinary
Lie action and the ordinary ghost bracket. The inherited stationary
Palatini argument proves invariance of its action; the chart map
transports the Lie bracket. Consequently the classical master identity
holds at the finite field/source degrees used in this calculation.
The all-degree local generating description allows the extra degrees
needed when differentiating an identity before extraction. No identity
of the nonnilpotent finite ultraviolet source is used.

Apply the trace-algebra calculation of
\cref{lem:v22-finite-reference} to these ordinary differential
vertices and the cancelled reference representation. There is no
unregulated infinite determinant in this application. After the two
source slots have been extracted, every term of $sX$ is a finite
ordered trace containing $\Theta$ and at least one supported covariance.
By \cref{prop:v22-reference-compiler}, its loop momenta are confined
to a fixed compact annulus. The extra external differentiation or
rooting in $\mathscr D_1$ neither introduces an independent loop
momentum nor changes this support property. All cyclic interchanges
and all momentum integrations by parts are therefore operations on
smooth compactly supported integrands.

More explicitly, insert an auxiliary loop cutoff equal to one on
$|p|\le 8\lambda$. At the displayed source degree every operator
identity used in the finite proof involves labels within this region.
A contraction or differentiation in the proof cannot reach the
auxiliary cutoff boundary. The differentiated ordinary master identities
are literal polynomial identities on these labels; replacing an
unrestricted finite carrier by this auxiliary truncation cannot change
their extracted terms. The remaining cyclic changes of routing are
justified by translation of the full continuum loop integral. Hence the
finite trace cancellation passes to this coefficient. The temporary
$\nu_\epsilon$ may equivalently be removed first at the finite-word
level: the extracted words contain only $\mathbb C$, $\Theta$ and
polynomial vertices, and \eqref{eq:v22-lower-annulus} supplies the
limiting nonsingular representation. Thus $sX=0$ holds in this
coefficient without an infrared inverse or ultraviolet boundary term.

Rooting and the affine projection now give
$\mathscr Q(sX)=\mathscr D_1\mathscr QX$, by the same local
coefficient computation as \cref{thm:v21-induced}. Homogeneity makes
Taylor degree four commute with this calculation, the target having
degree five. This proves \eqref{eq:v22-reference-closed}.
Finally \eqref{eq:v21-reference-zero} applies to this same reference:
two constant retained ghosts produce proportional odd momentum factors
in each word. The six projection is zero. Its vanishing is an inherited
measured test, not a new rotational average.
\end{proof}

The theorem is a consistency identity, not an assertion $X_R=0$.
The checker gives a nonzero finite reference example in a nonlinear
coordinate with an anisotropic quadratic regulator. Its reference
coefficient is closed and exact but not zero. A fixed physical lower
scale must therefore keep its actual reference normalization.

\section{The native lower-reference error begins at fourth mesh order}
\label{sec:v22-mesh}

The stronger power in this installment is due to its particular source
slots. No pure-gravity cubic is used in
\eqref{eq:v22-vertex-list}; the $O(a^3)$ cubic artifact of V17 is
therefore absent from this comparison. On the active annulus, only the
free matrices differ from their zero-mesh values.

Use $X(v),F_v,u(p)$ from \eqref{eq:v13-symbol-data}, and set
\begin{equation}
 v(p)=\frac1{1920}(0,p_1^5,p_2^5,p_3^5).
 \label{eq:v22-v}
\end{equation}
The sine series and the actual quadratic correction give
\begin{align}
 K_a&=K_0+a^4K_4+O(\chi a^6|p|^8),&
 K_4&=\frac\chi2\{DX(p)[v(p)]+X(u(p))\},\\
 M_a&=M_0+a^4M_4+O(a^6|p|^8),&
 M_4&=(p\cdot v)I+vp^T-pv^T.
 \label{eq:v22-free-four}
\end{align}
The full off-shell graded coefficient $\mathbb H_4$ has $K_4$ in
its $qq$ block, the corresponding $M_4$ ghost blocks, and zero
correction in the $qb$ and $bb$ blocks. This keeps the nonminimal
inverse rather than replacing it by the gauge-fixed classical Hessian.

\begin{lemma}[Reference covariance coefficient with the correct cutoff factor]
\label{lem:v22-covariance-four}
On the annulus, with every prescribed finite momentum derivative,
\begin{equation}
 \mathbb C_a=\mathbb C_0+a^4\mathbb C_4+O(a^6),\qquad
 \boxed{\mathbb C_4=-\nu\mathbb H_0^{-1}
                           \mathbb H_4\mathbb H_0^{-1}.}
 \label{eq:v22-C4}
\end{equation}
The remainders have the dimensions obtained from
\eqref{eq:v22-free-four}. In particular the metric and ghost inverse
remainders are respectively $O(\chi^{-1}a^6|p|^4)$ and
$O(a^6|p|^4)$ there. The coefficient is \emph{not}
$-\mathbb C_0\mathbb H_4\mathbb C_0$: that expression would
incorrectly introduce a second factor $\nu$.
\end{lemma}
\begin{proof}
The quadratic map $X$ obeys
$X(p+t u+t^2v)=X(p)+tDX(p)[u]+t^2(DX(p)[v]+X(u))+O(t^3)$.
Put $t=a^2$ and use the explicit subtraction in \eqref{eq:v13-K}.
Similarly $\widehat d_a=p+a^4v+O(a^6|p|^7)$ proves the ghost
formula. The uniform annular inverse bounds and the resolvent identity
give \eqref{eq:v22-C4}. Differentiating the same finite identities
gives all fixed derivative bounds. Multiply by the single original
factor $\nu$ only after inverting. The reference is varied with the
native free Hessian, so $\mathbb C_a\mathbb R_a=\Theta$ remains
fixed; holding $\mathbb R_a$ constant would describe another split.
\end{proof}

\begin{theorem}[Evaluated first coefficient and uniform comparison]
\label{thm:v22-reference-rate}
There is an explicitly specified array $x_{4,L}$ such that
\begin{equation}
 x_{a,L}=x_{0,L}+a^4x_{4,L}+e_{6,a,L}.
 \label{eq:v22-x-expansion}
\end{equation}
It is obtained by replacing each covariance, once at a time, in the
finite word \eqref{eq:v22-reference-compiler} by
\eqref{eq:v22-C4}. In formulas,
\begin{equation}
 \boxed{x_{4,L}=\mathscr Q\mathsf T_4\mathcal E_1
 \sum_{j=1}^{4}(-1)^j\sum_{\ell=0}^{j-1}
 \operatorname{STr}\!\left[
 (\mathbb C_0V)^\ell\mathbb C_4V
 (\mathbb C_0V)^{j-1-\ell}\Theta J\right].}
 \label{eq:v22-x4}
\end{equation}
For each total Taylor degree $d\le4$,
\begin{align}
 \|x_{4,L}^{[d]}\|_1&\le C_d\chi^{-1}\lambda^{8-d},&
 \|e_{6,a,L}^{[d]}\|_1&\le C_d\chi^{-1}a^6\lambda^{10-d},\\
 \|(x_{a,L}-x_{0,L})^{[d]}\|_1
       &\le C_d\chi^{-1}a^4\lambda^{8-d}.&&
 \label{eq:v22-x-bounds}
\end{align}
Constants are independent of mesh and periods on
\eqref{eq:v22-domain}. For every fixed $b>4$,
\begin{equation}
 \|(x_{0,L}-x_{0,\infty})^{[d]}\|_1
 \le C_{d,b}\chi^{-1}\lambda^{4-d}
                  (\lambda L_{\min})^{-b}.
 \label{eq:v22-x-volume}
\end{equation}
The formulas specify the actual finite normalized sums and their
thermodynamic integrals; no numerical value of $x_4$ is asserted.
\end{theorem}
\begin{proof}
There are at most four covariance factors and finitely many fixed source
vertices. Insert \eqref{eq:v22-C4} into each word and collect $a^4$;
this gives \eqref{eq:v22-x4}, with every matrix order unchanged.
Products of two fourth-order replacements start at $a^8$, so the
next possible order is the single sixth-order free replacement.
There is no odd power of $a$ on this fixed lower annulus.

Here are the scale bounds in physical normalization. A metric covariance
costs $\chi^{-1}\lambda^{-2}$, a ghost covariance costs
$\lambda^{-2}$, an ordinary transformation vertex costs one derivative,
and the native ghost-action vertex costs two. Expanding the four block
possibilities in \eqref{eq:v22-vertex-list}, the indicated single-antifield
reference words have total integrand order zero and one net factor
$\chi^{-1}$. This is equivalently the single-antifield degree count of
V18 with the additional first-order $Dr$ in place of the marked
curvature operand. Keeping the $b$ blocks gives the same count after
the exact constant square of V16. A fourth-order covariance replacement
increases this integrand order by four. A sixth-order replacement
increases it by six. Each external derivative lowers it by one.
All smooth profile derivatives have the same scaled bounds.

The number of periodic momenta in a ball of radius $3\lambda$, divided
by physical volume, is bounded by
$C\prod_\nu(\lambda+L_\nu^{-1})\le C'\lambda^4$.
Thus the full normalized annular sum gives
\eqref{eq:v22-x-bounds}, including the zero modes through their zero
covariance convention. No extensive determinant norm is used.
The integrands are smooth and supported away from zero. Integrating by
parts $b$ or more times in their loop Fourier transform and applying
Poisson summation gives \eqref{eq:v22-x-volume}, exactly with the
additional factor $\lambda^{-b}$ per physical length. The finite
number of coefficient and affine projections only changes constants.
\end{proof}

The expression \eqref{eq:v22-x4} is a more precise result than declaring
an unspecified small breaking. It tells which native free insertions
produce the first possible consistency error. It does not claim that
$\mathscr D_1x_4$ is nonzero; additional signed cancellation is allowed.

\section{The V21 combined consistency gate now has a vanishing bound}
\label{sec:v22-consistency}

For a degree-$g$ affine array $T$, use the dimensionless norm
\begin{equation}
 \|T\|_{g,\lambda}
       :=\chi\sum_d\lambda^{d-(g+3)}\|T^{[d]}\|_1,
 \label{eq:v22-scaled-norm}
\end{equation}
where $d$ is its number of rotation-ghost jets. It is the same coefficient
basis as V21, with the physical powers displayed. Since $\mathscr D$
raises $d$ by one and the contraction maps lower it by one, their
unscaled column-sum norms are unchanged in these norms. Direct addition
of the inherited integer columns gives
\begin{equation}
 \|\mathscr D_0\|_{0\to1}=11,\qquad
 \|\mathscr D_1\|_{1\to2}=15.
 \label{eq:v22-D-norms}
\end{equation}
These two small finite sums are independently verified by the checker;
the ranks and the norms 49 and 33 are inherited.

\begin{theorem}[Combined local affine consistency]
\label{thm:v22-consistency}
The actual combined coefficient of \eqref{eq:v22-three-arrays} obeys
\begin{equation}
 \boxed{\|\mathscr D_1R_{a,L}\|_{2,\lambda}
  \le15C_b\left[(a\lambda)^4+
                      (\lambda L_{\min})^{-b}\right],\quad b>4.}
 \label{eq:v22-consistency-bound}
\end{equation}
In thermodynamic normalization its expansion is
\begin{equation}
 \boxed{\mathscr D_1R_{a,\infty}
        =-a^4\mathscr D_1x_{4,\infty}+O_{2,\lambda}((a\lambda)^6).}
 \label{eq:v22-consistency-expansion}
\end{equation}
The six-class projection of $R_{a,L}$ is zero under the inherited
constant-ghost convention.
\end{theorem}
\begin{proof}
Insert \eqref{eq:v22-reference-closed} in
\eqref{eq:v22-canonical-cancellation}:
\[
 \mathscr D_1R_{a,L}
     =-\mathscr D_1(x_{a,L}-x_{0,\infty}).
\]
Use \eqref{eq:v22-D-norms}, \eqref{eq:v22-x-bounds} and
\eqref{eq:v22-x-volume}. The expansion follows in the same way from
\eqref{eq:v22-x-expansion}. The six term is zero because
$\Pi\mathscr D_0=0$ and the actual reference has the zero
constant-ghost coefficient established in V21. No cancellation between
separately bounded divergent ultraviolet integrals is being assumed:
their entire canonical array disappeared algebraically first.
\end{proof}

This evaluates the \emph{combined} consistency problem by a controlled
reference comparison. It does not prove that V19's isolated curvature
array is closed. Equating that isolated array with $R_{a,L}$ would omit
the action-breaking and reference balance which made
\eqref{eq:v22-canonical-cancellation} possible.

\begin{corollary}[Bounded source-action restoration of the affine slots]
\label{cor:v22-restoration}
Apply V21's complete local lift to the actual combined local input
$B_{a,L}=\mathsf T_4[\mathcal A_a^{\rm comb}]_{\eta cc,\sigma ccc}$,
and insert its counterterm $\hbar C_{B_{a,L}}$. Its remaining affine
coefficient $\varepsilon_{a,L}$ satisfies
\begin{equation}
 \boxed{\|\varepsilon_{a,L}\|_{1,\lambda}
 \le495C_b\left[(a\lambda)^4+
                         (\lambda L_{\min})^{-b}\right].}
 \label{eq:v22-restored-bound}
\end{equation}
In particular it tends to zero as $a\downarrow0$ in thermodynamic
normalization, or jointly with $\lambda L_{\min}\to\infty$ at fixed
positive $\lambda$. This is a local one-antifield conclusion, not a
claim of vanishing of every Slavnov row.
\end{corollary}
\begin{proof}
The lift identity \eqref{eq:v21-lift-identity} gives
$\varepsilon=\iota\mathsf Q_1\mathscr D_1R$ because $\Pi R=0$.
Use $\|\mathsf Q_1\|=33$ and \cref{thm:v22-consistency}.
The factor $495=33\cdot15$ is explicit. The lift's fixed component
conversion is kept separate when raw differential-polynomial norms,
rather than the stated affine coefficient norm, are used.
\end{proof}

At a fixed finite torus the last term in
\eqref{eq:v22-restored-bound} must not be discarded. The affine Killing
test is a local continuous-momentum normalization, not a global rotation
symmetry of an arbitrary rectangular torus.

\section{A reference-matched primitive and the unchanged descendant obligations}
\label{sec:v22-matching}

There is a second useful normalization, which makes the remaining
reference mismatch especially transparent. By the zero-mesh cocycle
and the six-term test,
\[
  x_{0,\infty}=\mathscr D_0\mathsf H x_{0,\infty}.
\]
Hence the explicit affine counterterm
\begin{equation}
 t_{a,L}=-g_{a,L}+\mathsf H x_{0,\infty}
              \quad\in\mathscr C^0
 \label{eq:v22-reference-matched-counterterm}
\end{equation}
satisfies
\begin{equation}
 \boxed{R_{a,L}+\mathscr D_0t_{a,L}
                 =x_{0,\infty}-x_{a,L}.}
 \label{eq:v22-direct-matching}
\end{equation}
All entries of $\mathsf Hx_{0,\infty}$ are fixed rational combinations
of the finite lower-annulus integrals specified above. This is not an
assertion that those entries vanish. The constant-source normalizations
left in $\ker\mathscr D_0$ remain freely matchable subject to the
other source equations.

For completeness, the actual local lift of this alternative is also
explicit. Let $\iota$ and $\mathcal E$ be those of V21, and put
\begin{equation}
 \begin{split}
 A&=\iota t_{a,L},\\
 V_A&=\mathcal E(s_\infty A)-\iota\mathscr D_0t_{a,L},\\
 C^{\rm ref}_{a,L}&=A-\mathfrak h B_{a,L}-\mathfrak h V_A.
 \end{split}
 \label{eq:v22-direct-lift}
\end{equation}
Since $\mathscr QV_A=0$, V20's paired-jet section applies. The same
zero-background calculation as in the V21 lift gives
\begin{equation}
 B_{a,L}+\mathcal E(s_\infty C^{\rm ref}_{a,L})
                    =\iota(x_{0,\infty}-x_{a,L}).
 \label{eq:v22-direct-lift-identity}
\end{equation}
This is a choice of source normalization, not an additional physical
coupling prescription. In particular it does not justify choosing
incompatible finite normalizations in the earlier geometric rows.

Both counterterms are ghost-number-zero and vanish at zero antifields.
The underlying zero-antifield action, propagators, auxiliary stationary
section, coordinate Jacobian and Lorentz/Haar density are unchanged.
Nevertheless the \emph{whole} variation $s_\infty C$ must be recorded.
An affine metric-antifield primitive gives an antifield-free $qc$
contact; the ghost-antifield primitive contributes the coupled
$\eta cc$ adjoint-generator term. Paired-jet lifts generate the
$q^2c$ and $\eta qcc$ contacts of V20. Nonlinear variations can reach
further background rows. Thus restoration here updates V16 and V17's
local matching equations; it cannot be superposed while freezing those
equations' old source normalizations.

The inherited finite free-Euler difference still multiplies potentially
divergent counterterm coefficients. Its first-row contribution has the
$O(a\mu^6)$ budget recorded in V21. A cubic artifact multiplying a
leading $a^{-3}$ primitive can still leave a finite local term at the
next ordinary-background degree. The stronger reference estimate in
this installment does not erase that descendant. No complete
ordinary-background induction is claimed.

The residual polynomial itself has controlled local source topology.
For any fixed smooth external window of scale $\mu\le\lambda/512$
and any fixed kernel moment, \eqref{eq:v22-restored-bound} implies
\begin{equation}
 \|\mathcal K^{\rm aff}_{a,L}\|_{\rm rooted,\mu,b}
 \le C_b'\chi^{-1}\lambda^4
       \left[(a\lambda)^4+(\lambda L_{\min})^{-b'}\right],
               \qquad b'>4.
 \label{eq:v22-affine-kernel}
\end{equation}
Indeed a degree-$d$ coefficient is bounded by
$C\chi^{-1}\lambda^{4-d}$ times the bracketed error. Its windowed
inverse Fourier transform has weighted $L^1$ norm at most the same
bound times $\mu^d$, by integration by parts. Summing the finite
polynomial and periodizing proves the claim. One source can be placed
in $L^1$, all others in $L^\infty$, without an extensive multiplier.
This statement concerns the residual polynomial, not the norm of the
large local counterterm before its prescribed scaling.

\section{What the new gate closes and what it leaves visible}
\label{sec:v22-scope}

The earlier requirement that the 930-component input satisfy an
uncomputed consistency equation is no longer left as an independent
assumption for the \emph{combined} coefficient
\eqref{eq:v22-combined-definition}. Its ultraviolet local part is an
exact affine boundary; its consistency error is determined by the
lower-reference array. The ordinary reference array is a cocycle, and
the native-to-ordinary difference has a fourth-order mesh estimate on
a fixed compact momentum annulus. This gives a source-action
restoration with a vanishing affine remainder.

This is a result at one loop, one minimal antifield and zero ordinary
background. It is not the consistency theorem of the entire finite
filtered source, which still has the higher-square witnesses of V18.
It is not a proof that all higher-antifield or ordinary-background
components form a closed subcomplex. The exact finite identities,
the thermodynamic cocycle and the finite-volume comparison have
different roles and have been kept distinct.

In particular, the next missing object is not another computation of
$\mathscr D_1\mathscr D_0$. It is the \emph{first ordinary-background
coefficient of the combined local equation after the full descendant
update}, including the finite cubic-artifact product, the reference
variation and the two-antifield contacts. The primitive just constructed
must be inserted into that equation before claiming it vanishes.
The matter/chiral part of the Einstein--SM programme, an invariant
interacting shell class, higher loops and infrared removal remain
separate. The positive lower scale and the declared complex metric
reference are not removed by any assertion here.

\section{Verification and closing ledgers for V22}
\label{sec:v22-ledgers}

The companion \srcid{check_ESM_combined_consistency_V22.py} reuses the
specified V21 integer column construction as inherited data. It tests
the new canonical/reference reduction on exact rational arrays, records
the column norms 11 and 15, checks the reference-matched counterterm and
its residual, and verifies the $a^4$ native free coefficients and the
single-$\nu$ covariance replacement. A nonlinear finite-coordinate
example has a nonzero reference coefficient; its exact Ward and
consistency identities are checked with exterior arithmetic. These tests
are not numerical evaluations of the native $x_4$ loop integral and not
a proof-assistant formalization of the analytic estimates.

\subsection*{Proved}
\addcontentsline{toc}{subsection}{V22: Proved}
The combined affine coefficient is
$\mathscr D_0g_{a,L}-x_{a,L}$, so all its ultraviolet canonical local
coefficients cancel from consistency. The required reference array is
the complete finite ordered trace \eqref{eq:v22-reference-compiler},
with support \eqref{eq:v22-lower-annulus}. Its zero-mesh thermodynamic
value is a cocycle. Its first possible mesh correction is the explicit
fourth-order insertion \eqref{eq:v22-x4}, with the stated coefficient
and volume bounds. Equations \eqref{eq:v22-consistency-bound} and
\eqref{eq:v22-restored-bound} discharge the V21 consistency estimate
for this combined one-antifield module. The alternative source-action
normalization \eqref{eq:v22-direct-lift} retains the exact reference
mismatch. All conclusions are restricted to the stated source,
background, measure and reference conventions.

\subsection*{Inherited}
\addcontentsline{toc}{subsection}{V22: Inherited}
V1--V21 remain unchanged. The completed metric/ghost reference and free
inverse margins, the ordinary zero-mesh action and chart transformations,
minimal-antifield coefficient construction, finite measured reference
identity, constant-ghost zero test, paired-jet section, affine complex
and rational contraction are used at their original scopes. The ranks,
six-class quotient and the norms 49 and 33 are not new results here.
The classical Lorentzian, spectral and structural SM sources remain
classical or structural interfaces, not a quantum conclusion.

\subsection*{Literature input}
\addcontentsline{toc}{subsection}{V22: Literature input}
The BV consistency mechanism and modified Slavnov reference terms are
standard \cite{V15IIM,V20BBH1994,V20BBH1995}. Their role here is
methodological. The source-specific content is the combined-row
cancellation, the finite lower-annulus compiler and its native
fourth-order comparison. No claim of first construction of perturbative
gravity, Einstein--SM EFT or local BRST cohomology is made.

\subsection*{Conditional}
\addcontentsline{toc}{subsection}{V22: Conditional}
Restoration of the complete local quantum functional requires matching
all descendants of the displayed source-action counterterm and the
remaining background/higher-antifield equations. A stronger finite-torus
symmetry statement requires its own compatible normalization convention.
No estimate is uniform when the positive lower scale is removed, when
$\chi\to0$, or when the source/derivative order becomes unbounded.
Reference matching to the discontinuous uncompleted microscopic
regulator is not supplied by this calculation.

\subsection*{Open}
\addcontentsline{toc}{subsection}{V22: Open}
Evaluate the updated one-background combined local BV coefficient with
its reference and measure terms and its finite cubic-artifact product.
Then test the coupling to the higher-antifield rows using the same
counterterm, rather than independent normalizations. Mixed SM/chiral
transport, higher-loop source closure and a full filtered interacting
ESM trajectory remain open. Neither analytic ultraviolet gravity nor
a nonperturbative positive gravitational measure is asserted.

\begin{center}\small
\begin{tabular}{>{\raggedright\arraybackslash}p{.23\textwidth}>{\raggedright\arraybackslash}p{.68\textwidth}}
\toprule
Variable & Scope in V22\\
\midrule
Volume & Coefficients use normalized sums; constants are period independent
for $\lambda L_\nu\ge1$. Exact affine cocycle uses thermodynamic
normalization; finite-volume error is explicit.\\
Mesh & $a\lambda\le1/32768$. The new reference error is
$O((a\lambda)^4)$ in scaled local norms, with an evaluated first
possible coefficient and $O((a\lambda)^6)$ remainder.\\
Loop/source degree & One loop; the combined $\eta cc$ and $\sigma ccc$
local slots at zero ordinary background.\\
Derivative order & Taylor degree four for the row, three for its source
primitive, and every prescribed fixed derivative or kernel moment in
the comparison.\\
Lower scale/coupling & $\lambda\ge512\mu>0$ and $\chi>0$ fixed;
no infrared or zero-coupling normalization limit.\\
Shell count & The remaining reference loop has only a bounded number of
levels near the fixed lower scale. Ultraviolet local arrays cancel
algebraically in consistency.\\
Symmetry & Vanishing combined affine consistency error and local source
restoration in the specified quotient; not full nonlinear quantum BV.\\
Measure & Same off-shell nonminimal block, coordinate Jacobian,
Lorentz/Haar density and complex contour, all counted once.\\
\bottomrule
\end{tabular}
\end{center}


\clearpage
\part*{V23: A derivative floor removes the putative finite descendant return}
\addcontentsline{toc}{part}{V23: Derivative floor, descendant transport and the first background reference}

\section{The dangerous product must be tested on the populated coefficients}
\label{sec:v23-start}

Sections 1--176 and their labels are retained. The question left by V22 is
not whether its affine matrices can be inverted again. It is whether the
\emph{actual} source counterterm can be transported into the next ordinary-
background equations without an unaccounted finite term. The conservative
estimate in \cref{sec:v22-matching} permitted an $a^{-3}$ coefficient in a
zero-derivative source primitive. Multiplying that coefficient by the
$a^3$ cubic artifact would leave a finite local return. This installment
checks whether that coefficient is populated, rather than assigning another
counterterm to a possibly absent term.

The answer on the specified gravitational module is negative. At zero external
momentum, its one-loop source functional and its first \emph{constant} metric
variation vanish. The complete zero-background local breaking consequently
starts at derivative degree two. Both particular counterterms constructed in
V21--V22 start at derivative degree one and have worst power $a^{-2}$, not
$a^{-3}$. Their cubic-artifact descendant is therefore $O(a)$ in a fixed
source window. This strengthens a conservative earlier bound; it changes
neither an earlier action nor an earlier local matching coefficient.

The second calculation identifies the first-background reference error. Its
first possible mesh term is the insertion of an explicitly evaluated cubic
artifact of order $a^3$. That insertion vanishes for a constant metric mark.
For a nonconstant mark it is a fixed lower-annulus integral. These statements
are not an independent proof that the whole one-background BV complex closes.
The higher-antifield and higher-background entries in its consistency equation
are retained in \cref{sec:v23-remaining}.

\subsection{Fixed family, notation and scope}

Use the smooth completed action of V14, its degree-five extension in V17,
the V18 minimal source, V16's off-shell nonminimal convention, and V10--V12's
once-counted auxiliary density. Work in the Euclidean coefficient/complex
contour convention already declared in V13, with $\chi>0$ fixed. No
Lorentzian-continuation statement is added. The ordinary background is the
affine coframe coordinate $q$, not the exact metric difference. The common
factor $\hbar$ is omitted from one-loop coefficients. Write
\[
 \eta=q^*,\qquad \sigma=c^*,\qquad
 G(q)=I+q+\mathsf Q(q,q),\qquad DG(q)=I+\mathsf B_q.
\]
The particular counterterms below are those in
\eqref{eq:v21-lift-construction} and \eqref{eq:v22-direct-lift}, without an
additional independently chosen element of their homogeneous solution
spaces. Adding an arbitrary derivative-free source normalization would not
preserve the conclusion being proved.

We take the cofinal odd product grids, so their principal momentum sets are
invariant under $p\mapsto-p$ without an unpaired Nyquist label. All profiles
$s,\zeta,\Theta$ are the inherited real even profiles, and the lower covariance
vanishes on the retained zero mode. A convenient common domain is
\begin{equation}
 \lambda=\Lambda_0\ge1024\mu>0,\qquad
 a\lambda\le 1/65536,\qquad \lambda L_\nu\ge1.
 \label{eq:v23-domain}
\end{equation}
The extra restriction to odd coordinate counts is used for the exact finite
momentum-pair cancellation. No assertion for an unpaired spectral Nyquist
convention is needed. All estimates are at fixed source order, profile
seminorms and coupling normalization. The matter/chiral sectors are not
included in this gravitational source calculation.

For one even mark $h$, use $q=\tau h$, $\tau^2=0$, and precisely the inverse
and source jets of \eqref{eq:v19-background-cov}--\eqref{eq:v19-background-D}.
In particular, $D_\tau=D_{\eta,\tau}+D_{\sigma,\tau}$ differentiates both
ghost returns, not just the metric propagator. The unmarked counterpart of
V19's measured computation is
\begin{equation}
 \Gamma^{\min}_{a,1}(\tau)
 =\frac12\Tr\log\bigl(I+C_g(\tau)D_\tau\bigr).
 \label{eq:v23-background-generator}
\end{equation}
This is an application of the inherited quadratic-hard reduction. All
external labels in the coefficient panel are soft, so a hard-linear vertex
has no covariance support. At one loop every remaining vertex has exactly
two hard legs. The ghost determinant without minimal antifields cancels
from \eqref{eq:v23-background-generator}; the auxiliary density has no direct
minimal-antifield term at this order. Thus the displayed formula includes the
whole required minimal-source coefficient, not a selected subset of graphs.

\section{The native constant-background Hessian has a special inversion symmetry}
\label{sec:v23-constant}

The full oriented cubic does not have a general simultaneous-inversion parity.
The argument uses a more restricted fact, which is enough to determine its
zero-transfer source coefficient.

\begin{lemma}[Even hard Hessians at one constant metric mark]
\label{lem:v23-even-hessian}
Let $h$ be a constant coframe-coordinate polarization. For the completed
metric and ghost covariances and their first background derivatives,
\begin{equation}
 \begin{aligned}
 C_g(-p)&=C_g(p),&\qquad C_c(-p)&=C_c(p),\\
 A^g[h](-p)&=A^g[h](p),& A^c[h](-p)&=A^c[h](p).
 \end{aligned}
 \label{eq:v23-even-hessian}
\end{equation}
Here $A^g[h](p)$ and $A^c[h](p)$ mean the two-hard-leg kernels with zero
momentum on $h$. Consequently $C_g(\tau)$ and $C_c(\tau)$ are even through
$\tau^1$. The equalities are matrix equalities in the fixed component bases,
not statements obtained by taking an angular average.
\end{lemma}
\begin{proof}
The free matrices in \eqref{eq:v14-completed-free} are even: each component
of $d_a$, $\widehat d_a$ and $p$ is odd, the free matrices are quadratic in
those components, and all profiles and $\lambda_a$ are even. Their inverses
and the even lower weight $\nu$ preserve this property.

For the metric derivative, use the actual reduced cubic
\eqref{eq:v10-reduced-jets} and the V11 counterterms. The load $b_{1,a}(h)$
is zero because $h$ is constant, and hence $A_{1,a}(h)=0$. In
$\langle b_{2,a},A_{1,a}\rangle$, each surviving polarized term has one
phase derivative on the hard cofactor load and one on the hard $A_1$.
Both change sign when $p$ is reversed. The term
$\langle A_1,C_1(h)A_1\rangle/2$ has the same two odd factors. The native
shared-link term $a\mathscr R_3^{(0)}(A_1)$ is trilinear in $A_1$; with one
constant metric label it is exactly zero, including every translated
occurrence. Its V11 point and shift counterterms vanish for the same reason.
The remaining geometric $a^2$ counterterm has one plus three phase-jet
orders and is even. This proves evenness for the complete native corrected
cubic Hessian at this mark, without assigning a parity to a nonconstant mark.

The upper metric interaction has $Hh=0$, $Sh=h$, so only its source factor
can carry this constant mark. The product of the two forward symbols is
$|e^{iap_\nu}-1|^2/a^2$, which is even. The ghost-action vertices are
$-FR_1$ and $-FR_2$ on the filtered chart. Each has two ordinary derivatives,
and their constant-mark hard symbols are even. The upper ghost term again
has the product of opposite forward symbols. This proves the two remaining
identities. Finally $C'(\tau)|_0=-CA[h]C$ is even as a product of the
established even matrices; no commutation of these matrices is required.
\end{proof}

\begin{theorem}[Vanishing zero-transfer minimal-source coefficients]
\label{thm:v23-zero-transfer}
At every prescribed finite nonzero minimal-antifield degree, set all retained
antifields and ghosts to constant formal fields. Then
\begin{equation}
 \boxed{\Gamma^{\min}_{a,1}(0)\big|_{K=0}=0,
 \qquad [\tau]\Gamma^{\min}_{a,1}(\tau h)\big|_{K=0}=0}
 \label{eq:v23-zero-transfer}
\end{equation}
for every constant $h$. These are identities of the exterior-algebra
coefficients at each finite odd regulator. For one metric antifield the
cancellation uses the pair $p,-p$. For a ghost antifield, or at least two
minimal antifields, the constant-ghost factors already vanish in each
closed-momentum word.
\end{theorem}
\begin{proof}
Let $u$ be the constant retained odd vector ghost and
$\alpha_p=i\,u\cdot p$. The native mixed ghost-action block has the inherited
factorization
$B_u(p)x=-s(ap)^2\alpha_p F_px$. Its first constant-background derivative
is zero. Indeed the ordinary coordinate formula gives
\[
 R_2^{(0)}(h,x_p,u)
 =\mathcal L_u\mathsf B(h,x_p)
   -\mathsf B(h,\mathcal L_u x_p)-\mathsf B(x_p,\mathcal L_u h)=0,
\]
and the upper ghost interaction contains $(I-S)u=0$. The derivative of the
ghost inverse does not remove the factor $\alpha_p$. Thus every returned
hard ghost at order $\tau^0$ or $\tau^1$ has the form
$\beta_p=\alpha_p V_p(\tau)x_p$, where $V_p(\tau)$ is even in the
Grassmann grading and independent of $u$.

For the direct transformation contact, the coefficient of two opposite hard
metric labels in $\breve R_2(x,u)$ is zero. The term
$P\mathcal L_u\mathsf Q(x,x)$ has zero total momentum; the two terms in
$-\mathsf B(x,P\mathcal L_u x)$ cancel because $P(p)=P(-p)$ and
$\mathsf B$ is symmetric. The same statement after one constant background
mark follows from the source recursion: its two-field terms with one
constant mark are
$P\mathcal L_u\mathsf B(h,x)-\mathsf B(h,P\mathcal L_u x)=0$.
Thus the direct $T_{\eta u,\tau}$ vanishes in this zero-transfer
quadratic form.

It remains to trace the ghost-mediated $D_{\eta,\tau}$. The kernel
$L_{\eta,\tau}$ has exactly one derivative and is odd under simultaneous
reversal of its two opposite hard labels. The block $B_u$ has two derivatives
and is even, as are $C_c(\tau)$ and $C_g(\tau)$ by
\cref{lem:v23-even-hessian}. Therefore the one-$\eta$ trace integrand is odd
in $p$. The symmetric finite sum is zero. This reasoning holds coefficientwise
in $\tau$, and keeps the actual nonscalar ghost inverse.

The form $Q_\sigma(C_c(\tau)B_ux)$ has two returned ghost factors. At
zero total transfer they contain $\alpha_p\alpha_{-p}=0$, since
$\alpha_{-p}=-\alpha_p$ and $\alpha_p^2=0$ in the exterior algebra.
Hence $D_{\sigma,\tau}=0$. Each remaining $D_{\eta,\tau}$ contains one
$\alpha_p$. A closed word with at least two such insertions has only the
momenta $p,-p$ and therefore contains $\alpha_p^2$, after its signs are
retained. It vanishes. The constant-background covariance insertion carries
no ghost and does not change this reasoning. Expand
\eqref{eq:v23-background-generator} at the required finite source degree.
This proves both identities, including every metric and ghost return.
\end{proof}

For clarity, denote the one-antifield quantum coefficients by
\begin{equation}
 [\Gamma_{a,1}]_{\text{one antifield},\,q^{\le1}}
 =\langle\eta,Z_a(c)+W_a(q,c)\rangle
  +\langle\sigma,Y_a(c,c)+U_a(q,c,c)\rangle.
 \label{eq:v23-source-panels}
\end{equation}
These symbols refer to the V18 source, not to an unchanged older source
prescription. The theorem states, in particular,
\begin{equation}
 \boxed{Z_a^{[0]}=W_a^{[0]}=Y_a^{[0]}=U_a^{[0]}=0.}
 \label{eq:v23-four-zero-jets}
\end{equation}
The superscript is total external derivative degree. It does not mean
$W_a$ or $U_a$ vanishes at nonzero external momentum. In particular no
general oddness of $W_a(K)$ has been used. The theorem also removes the
zero-derivative coefficient of the quadratic-antifield module needed later;
its higher local jets remain present.

\section{The actual V22 counterterm has at least one derivative}
\label{sec:v23-floor}

Let
\begin{equation}
 B_{a,L}=\mathsf T_4
 [\mathcal A_a^{\rm comb}]_{\eta cc,\sigma ccc,q=0}
 \label{eq:v23-B}
\end{equation}
with exactly \eqref{eq:v22-combined-definition}, before any new source
normalization. All densities are rooted on their single antifield by the
same integration-by-parts convention as V20.

\begin{lemma}[The reference has no first-derivative coefficient in these slots]
\label{lem:v23-reference-parity}
At zero ordinary background the two indicated coefficients of
$\mathcal X_a$ are even under reversal of all external momenta.
Their derivative-degree-zero and derivative-degree-one jets vanish.
\end{lemma}
\begin{proof}
Use the complete finite words of \eqref{eq:v22-reference-compiler}.
There is no pure-gravity interaction Hessian in that extraction. On the
full off-shell field list assign reversal sign $-1$ to $b$ and $+1$
to $q,c,\bar c$, and let $E_b$ be the corresponding constant diagonal
matrix. The free Hessian and its covariance transform by conjugation with
$E_b$: the $bFq$ entries are odd and every other nonzero free entry is
even. This remains true for the native corrected matrices and the reference
profile.

A ghost-action vertex has two derivatives and transforms by that conjugation
with no extra sign. A minimal-antifield Hessian has one derivative and has
one extra minus sign. The source Jacobian $Dr$ also has one extra minus
sign; its derivative-free $\bar c\mapsto ib$ entry has precisely the same
sign because of $E_b$. A contributing word has exactly one minimal-antifield
vertex. The two additional minus signs consequently cancel. The $E_b$
matrices telescope inside the supertrace. The term $\Theta$ is even and
is not moved through a vertex. This proves evenness of the assembled
coefficient and the vanishing first Taylor jet. The constant jet is the
constant-ghost cancellation already established in V21: the required two
or three retained ghosts supply proportional odd momentum factors.
Equivalently apply that finite-word argument to this same compiler.
There is no assertion of parity for a reference coefficient with a
nonconstant metric mark.
\end{proof}

\begin{theorem}[Derivative floor for the combined input and its particular primitive]
\label{thm:v23-derivative-floor}
The actual combined input satisfies
\begin{equation}
                    \boxed{B_{a,L}^{[0]}=B_{a,L}^{[1]}=0.}
 \label{eq:v23-B-floor}
\end{equation}
Let $C_{a,L}$ be either the particular V21 complete lift applied in
\cref{cor:v22-restoration}, or the reference-matched lift
\eqref{eq:v22-direct-lift}. Without adding a homogeneous source normalization,
\begin{equation}
 \boxed{C_{a,L}=C_{a,L}^{[1]}+C_{a,L}^{[2]}+C_{a,L}^{[3]},
 \qquad C_{a,L}^{[0]}=0.}
 \label{eq:v23-C-floor}
\end{equation}
It is linear in the minimal antifields and has at most one ordinary metric
factor, as in the existing lift.
\end{theorem}
\begin{proof}
Extract the canonical bracket in \eqref{eq:v22-combined-definition} before
Taylor truncation. At zero ordinary background a nonzero contribution to
one minimal antifield comes from the four panels in
\eqref{eq:v23-source-panels}. A variation by $R_0c$ consumes their one-metric
part; the remaining terms are the linear index/transport action on their
zero-metric part, the ordinary ghost bracket, and the rooted
metric--ghost-antifield coupling. Each such operation introduces exactly
one external derivative. In particular the degree-one coefficient depends
only on the four zero-degree coefficients in
\eqref{eq:v23-four-zero-jets}, and is zero.

There is no hidden source-free Euler contribution at this extraction. The
metric Euler row vanishes at the flat vacuum. A higher-minimal-antifield
coefficient can enter through it only with an ordinary metric factor. The
source-free ghost Euler row can be nonzero on a nonconstant retained ghost,
but it multiplies a nonminimal-antifield derivative of $\Gamma_1$. That
derivative is zero here: the nonminimal master term $i\bar c^*b$ is linear
in the integrated fields and has zero Hessian, while $b^*$ does not occur.
Hard-linear returns have no support. The direct super-divergence and the
once-counted auxiliary density have no one-minimal-antifield coefficient.
Thus the listed canonical terms exhaust the extraction. Its constant
Taylor jet is zero since each has a derivative. Apply
\cref{lem:v23-reference-parity} to the reference part to prove
\eqref{eq:v23-B-floor}.

The paired-jet map $\mathfrak h$ of V20 replaces one paired ghost derivative
by its metric partner and lowers total derivative degree by one. The affine
map $\mathsf H$ of V21 also lowers degree by one; this is a proved grading
property of its fixed rational matrix, not a new rank computation. The lift
$\iota$ preserves degree. The zero-background extraction of
$s_\infty\iota\mathsf H\mathscr Q B$ raises it by one; its paired-jet
correction lowers it back. Starting from degrees two through four therefore
gives only degrees one through three in every term of
\eqref{eq:v21-lift-construction}. For the alternative lift,
$g_{a,L}^{[0]}=0$ by \eqref{eq:v23-four-zero-jets}, and the reference array
starts at degree two. The same reasoning applies to
\eqref{eq:v22-direct-lift}. Neither construction introduces a product of two
antifields or more than one ordinary metric factor. This proves all claims.
\end{proof}

\begin{proposition}[Improved native coefficient budget]
\label{prop:v23-budget}
For the fixed panels above, in the original component coefficient norms,
\begin{equation}
 \|C^{[1]}_{a,L}\|_1\le C\chi^{-1}a^{-2},\quad
 \|C^{[2]}_{a,L}\|_1\le C\chi^{-1}a^{-1},\quad
 \|C^{[3]}_{a,L}\|_1
       \le C\chi^{-1}\{1+\log(1/(a\lambda))\}.
 \label{eq:v23-coefficient-budget}
\end{equation}
In a source window $|K_i|\le\mu$, every fixed scaled momentum seminorm of
$C_{a,L}$ is bounded by $C_r\chi^{-1}a^{-2}\mu$. The corresponding windowed
rooted kernel norm has the same bound, with a constant depending on its
fixed moment order, not on total physical volume.
\end{proposition}
\begin{proof}
One factor $C_gD$ in \eqref{eq:v23-background-generator} has symbol order
$-1$. A first background derivative of a covariance replaces it by
$-CA[h]C$ and preserves that order: $A[h]$ has order two. Differentiating
$D$ retains the ghost inverse and the source/ghost-action vertices and
also preserves order. Thus, on an annulus $|p|\asymp R\ge\lambda$, each
$r$th external derivative of these four one-antifield panels is bounded by
$C_r\chi^{-1}R^{-1-r}$. On the completed upper region the same estimate is
read as $C_r\chi^{-1}a^{1+r}$. The inherited inverse margins and finite
number of source vertices justify the constants even where oriented cubic
terms occur; there $aR$ is bounded, not taken as a small relative error.

Counting four-dimensional momenta per physical volume gives
$a^{r-3}$ for $r<3$ and $1+\log(1/(a\lambda))$ for $r=3$. The zero-degree
coefficient is exactly absent by \cref{thm:v23-zero-transfer}. The canonical
operation producing $B$ raises degree by one. The reference coefficient is
a lower-annulus trace with degree-$d$ size $C_d\chi^{-1}\lambda^{4-d}$,
by V22; it has no degrees zero or one. These bounds are dominated by the
stated ones because $a\lambda\ll1$. Apply the fixed norms of the paired
and affine lifts to obtain \eqref{eq:v23-coefficient-budget}. No growing
matrix dimension is used in this step.

Multiplication of the degree-$d$ coefficient by $\mu^d$ proves the scaled
symbol claim. The logarithmic term is harmless in that bound since
$(a\mu)^2[1+\log(1/(a\lambda))]$ is bounded on
\eqref{eq:v23-domain}. Multiply by a smooth fixed external window, integrate
by parts more times than the number of relative momentum coordinates plus
the requested moment order, and periodize. The same scaled polynomial
seminorms bound the weighted $L^1$ kernel. This gives a multilinear estimate
with one source in $L^1$ and all other sources in $L^\infty$.
\end{proof}

\section{The full descendant update has no finite cubic-artifact return at this order}
\label{sec:v23-descendants}

Write the metric-antifield part of the selected counterterm as
\begin{equation}
 C_{a,L}^{\eta}
   =\langle\eta,U_{0,a}(c)+U_{1,a}(q,c)\rangle.
 \label{eq:v23-U}
\end{equation}
Both $U_0$ and $U_1$ start at derivative degree one; either can be zero.
Keep the full ghost-antifield part as well. Let $s_a$ and $s_\infty$ be the
canonical differentials of the finite and zero-mesh source actions in the
same off-shell convention. They need not both square to zero on arbitrary
cutoff fields. On the stated soft panels their transformation-source
vertices agree exactly. Set $b=\bar c=0$ only after forming the canonical
bracket.

Let $E_{2,a}$ and $E_{3,a}$ be the signed quadratic-Euler and polarized
cubic-action differences, respectively. Their inherited fixed-window
bounds are
\begin{equation}
 \|\partial^\alpha E_{2,a}\|\le C_\alpha\chi a^4\mu^{6-|\alpha|},
 \qquad
 \|\partial^\alpha E_{3,a}\|\le C_\alpha\chi a^3\mu^{5-|\alpha|}.
 \label{eq:v23-E-bounds}
\end{equation}
The word ``signed'' fixes the same convention as $s\eta=\mathcal E(q)$
in V20; no extra minus sign is hidden in a norm.

\begin{theorem}[Transport of the particular counterterm into the first two action rows]
\label{thm:v23-descendant-bound}
For the actual $C_{a,L}$ of \cref{thm:v23-derivative-floor}, the canonical
finite-to-ordinary differences in the antifield-free rows are exactly
\begin{align}
 [(s_a-s_\infty)C]_{qc}
   &=\langle E_{2,a}q,U_{0,a}(c)\rangle,\notag\\
 [(s_a-s_\infty)C]_{q^2c}^{\rm pol}
   &=\langle E_{2,a}h_1,U_{1,a}(h_2,c)\rangle\notag\\
   &\quad+\langle E_{2,a}h_2,U_{1,a}(h_1,c)\rangle
         +E_{3,a}(h_1,h_2,U_{0,a}(c)).
 \label{eq:v23-exact-descendants}
\end{align}
For every fixed momentum-derivative order $r$ their scaled-window norms
satisfy
\begin{equation}
 \boxed{\|[(s_a-s_\infty)C]_{qc}\|_{r,\mu}
          \le C_r a^2\mu^7,\qquad
 \|[(s_a-s_\infty)C]_{q^2c}\|_{r,\mu}
          \le C_r a\mu^6.}
 \label{eq:v23-descendant-bounds}
\end{equation}
The same bounds hold in every fixed windowed rooted kernel moment norm.
The minimal-source algebra descendants of this linear-antifield counterterm
have no finite-to-ordinary difference on the same soft panels.
\end{theorem}
\begin{proof}
The bracket with the source-free action removes $\eta$ from
\eqref{eq:v23-U}. Its ordinary degree one and two terms are exactly
\eqref{eq:v23-exact-descendants}. A derivative of the source-free ghost
action in the $\sigma$ part still has an antighost factor, and vanishes on
the displayed extraction. The nonminimal terms are unchanged. The remaining
brackets involve source generators only. Each input label and the finitely
many subset labels lie in their common plateau, where the V18 source is the
ordinary source in the same chart. Those terms cancel in $s_a-s_\infty$,
without asserting a symmetry of the full finite action.

Apply \eqref{eq:v23-coefficient-budget}. The free term is bounded by
$C\chi a^4\mu^6\,\chi^{-1}a^{-2}\mu=C a^2\mu^7$.
The cubic term is bounded by
$C\chi a^3\mu^5\,\chi^{-1}a^{-2}\mu=C a\mu^6$.
The remaining free term is smaller since $a\mu\ll1$. Leibniz's rule for
external momentum derivatives gives the fixed scaled seminorms, and the
windowed Fourier argument of \cref{prop:v23-budget} gives the kernel bounds.
All source-polarization norms are understood multiplicatively.
\end{proof}

The absence of an $a^{-3}$ coefficient is the load-bearing point. The
previously allowed product $a^3a^{-3}$ is not populated by this particular
counterterm. The replacement $a^3a^{-2}=a$ closes that finite-artifact
budget. It does \emph{not} set $s_\infty C$ to zero. Its ordinary-background,
ghost-antifield and nonminimal descendants must still be added to the
matching equations. Choosing an independent derivative-free homogeneous
source counterterm would reinstate the older, weaker estimate and would
require its own matching calculation.

\section{The first background reference consumes an explicit third-mesh cubic}
\label{sec:v23-cubic}

To distinguish a constant background from a general background, calculate the
next native cubic coefficient. This is a Taylor coefficient of the unchanged
V11-restored action, not a further action modification. Use
$A=A_{1,0}=-C_0^{-1}b_{1,0}$ and
$A^{(2)}=A_{1,2}=-C_0^{-1}b_{1,2}$ from V11. For a plaquette $ij$ set
\[
 X=(A_i,A_j,-A_i,-A_j),\quad
 Y=(0,\partial_iA_j,-\partial_jA_i,0),\quad
 Z=(0,\tfrac12\partial_i^2A_j,-\tfrac12\partial_j^2A_i,0).
\]
These $X,Y,Z$ denote the four matrix slots of the inherited rational $B_3$
polynomial, not new connection variables. Define
\begin{equation}
 \begin{split}
 P_{\rm sh2}(A)={}&
 -\frac14\sum_i b(\Pi_i^0,[A_i,[A_i,\partial_i^2A_0]])\\
 &+\sum_{i<j}b\!\left(\Sigma_{ij}^0,
       DB_3(X)[Z]+\tfrac12D^2B_3(X)[Y,Y]\right),\\
 \mathcal O_{3,3}(q)={}&
 D\mathscr R_{3,\rm pt}(A)[A^{(2)}]+P_{\rm sh2}(A).
 \end{split}
 \label{eq:v23-O33}
\end{equation}
The derivatives of $B_3$ mean finite replacement of its matrix occurrences.
Thus \eqref{eq:v23-O33} is an explicit local differential polynomial, with
three fields and five derivatives. Its Euclidean coefficient is obtained by
the same fixed holomorphic substitution as the earlier cubic.

\begin{proposition}[Third-mesh coefficient and zero constant-metric slot]
\label{prop:v23-O33}
On each fixed nonaliased cubic panel,
\begin{equation}
 \widehat S^E_{3,a}=S^E_{3,0}+a^3\mathcal O^E_{3,3}
                         +O(a^4\Lambda_*^6).
 \label{eq:v23-cubic-expansion}
\end{equation}
Every fixed momentum derivative has the corresponding differentiated bound.
Moreover,
\begin{equation}
 \boxed{D^3\mathcal O^E_{3,3}(0)[h,x,y]=0
                   \quad\text{if }h\text{ is constant}.}
 \label{eq:v23-constant-O33}
\end{equation}
For a soft metric momentum $k_h$ and hard momenta in a fixed annulus of size
$\lambda$, its Hessian insertion $A_3[h]$ therefore obeys
\begin{equation}
 \|A_3[h](p,p+k_h)\|
       \le C\chi |k_h|\lambda^4\|h\|,
 \qquad \|\partial^\alpha A_3[h]\|
       \le C_\alpha\chi\lambda^{5-|\alpha|}\|h\|.
 \label{eq:v23-A3-bounds}
\end{equation}
\end{proposition}
\begin{proof}
The geometric cubic depends on the phase derivative only through an even
mesh expansion, and hence has no $a^3$ coefficient. In the remaining
term $a\mathscr R_3^{(0)}(A_{1,a})$, expand
$A_{1,a}=A+a^2A^{(2)}+O(a^4)$ and
$T_i=I+a\partial_i+a^2\partial_i^2/2+O(a^3)$.
The $a^3$ coefficient from the first expansion is
$D\mathscr R_{3,\rm pt}(A)[A^{(2)}]$. The second-shift coefficient of the
electric endpoint is the first line of $P_{\rm sh2}$, while the ordinary
second Taylor coefficient of the magnetic polynomial is
$DB_3(X)[Z]+D^2B_3(X)[Y,Y]/2$. V11 subtracted precisely the $a$ and $a^2$
coefficients, leaving \eqref{eq:v23-cubic-expansion}. Taylor's integral
remainder for the finite sine and exponential symbols gives the stated
fixed-order estimates, including the even geometric $a^4$ remainder.

Every field occurrence in \eqref{eq:v23-O33} is in $A$, $A^{(2)}$ or a
further spatial derivative of one of them. Each is zero on a constant
metric label. Polarization proves \eqref{eq:v23-constant-O33}.
As a polynomial in $(k_h,p)$ the insertion vanishes at $k_h=0$.
The integral identity
$A_3(k_h,p)-A_3(0,p)=\int_0^1 k_h\cdot\partial_{k_h}A_3(tk_h,p)\,dt$
and its total five-derivative degree prove \eqref{eq:v23-A3-bounds}.
No division by a soft momentum is used.
\end{proof}

\section{A controlled first-background reference coefficient}
\label{sec:v23-reference}

Keep the full off-shell graded field list and define the \emph{unprojected}
rooted local coefficient
\begin{equation}
 y_{a,L}[h]=\mathsf T_4
  [D_q\mathcal X_a(0)[h]]_{\eta cc,\sigma ccc}.
 \label{eq:v23-y}
\end{equation}
No one-background affine complex is assumed in this definition. Let $V_0$ be
the zero-background source/ghost Hessian used in V22 and
$J^{(0)}=J_0+J_c$. The derivative $V_h$ includes every source-marked derivative
as well as the source-free metric and ghost Hessian insertions. Similarly
$J_h=D_qDr_a(0)[h]$ retains the actual chart derivatives. Then
\begin{equation}
 \begin{split}
 y_{a,L}[h]=\mathsf T_4\mathcal E_1\operatorname{STr}\bigg\{&
 \sum_{j=1}^{5}(-1)^j\sum_{r=0}^{j-1}
 (\mathbb C_aV_0)^r\mathbb C_aV_h
 (\mathbb C_aV_0)^{j-1-r}\Theta J^{(0)}\\
 &+\sum_{j=0}^{4}(-\mathbb C_aV_0)^j\Theta J_h\bigg\}.
 \end{split}
 \label{eq:v23-y-compiler}
\end{equation}
Terms outside the indicated exterior labels are zero. The apparent $j=0$
term on the second line contains no antifield and is annihilated by
$\mathcal E_1$.

\begin{theorem}[First-background lower-reference comparison]
\label{thm:v23-reference}
Every contributing integrand in \eqref{eq:v23-y-compiler}, with each fixed
external Taylor derivative, is supported in
$\lambda/2\le|p|\le3\lambda$. On \eqref{eq:v23-domain},
\begin{equation}
 y_{a,L}[h]=y_{0,L}[h]+a^3 y_{3,L}[h]+e_{4,a,L}[h],
 \label{eq:v23-y-expansion}
\end{equation}
where the fully specified first possible coefficient is
\begin{equation}
 \begin{split}
 y_{3,L}[h]=\mathsf T_4\mathcal E_1\operatorname{STr}
 \sum_{j=1}^{5}(-1)^j\sum_{r=0}^{j-1}
 (\mathbb C_0V_0)^r\mathbb C_0 A_3[h]
 (\mathbb C_0V_0)^{j-1-r}\Theta J^{(0)}.
 \end{split}
 \label{eq:v23-y3}
\end{equation}
Here $A_3[h]$ occupies only the $qq$ block and is defined in
\cref{prop:v23-O33}. For total external derivative degree $d\le4$,
\begin{equation}
 \|y_{3,L}^{[d]}[h]\|_1\le C_d\chi^{-1}\lambda^{7-d}\|h\|,
 \qquad
 \|e_{4,a,L}^{[d]}[h]\|_1
       \le C_d\chi^{-1}a^4\lambda^{8-d}\|h\|.
 \label{eq:v23-y-bounds}
\end{equation}
If $h$ is constant, $y_{3,L}[h]=0$. The comparison on that constant-mark
restriction therefore retains fourth mesh order. For nonconstant $h$ the
third-order coefficient is retained, not asserted to be zero.
\end{theorem}
\begin{proof}
Differentiate the finite ordered resolvent in V22 once in $h$.
This gives \eqref{eq:v23-y-compiler}. There are at most five retained labels,
so no longer word contributes. Every surviving word has both a covariance
factor $\nu$ and the displayed $\Theta$. The same routing proof as in V22,
now allowing one additional soft transfer, gives the stated annulus. The
strict margins in \eqref{eq:v23-domain} keep every internal label and every
subset label inside the common source/completion plateau.

On this annulus the source and ghost interaction derivatives are precisely
the ordinary polynomial ones. The only additional tree operand with mesh
dependence is the source-free metric insertion:
$A^g_a[h]=A^g_0[h]+a^3A_3[h]+O(\chi a^4\lambda^6)$, by
\cref{prop:v23-O33}. The free covariance error starts at $a^4$, with the
single-$\nu$ rule of \cref{lem:v22-covariance-four}. There can be only one
$h$-carrying insertion. Replacing it by $A_3[h]$ gives
\eqref{eq:v23-y3}; no derivative of $J$ has an $a^3$ correction. This proves
the expansion and the matrix order of its coefficient.

Rescale $p=\lambda z$ in each finite annular sum. The unperturbed
one-antifield reference coefficient has degree-$d$ size
$C_d\chi^{-1}\lambda^{4-d}$. Replacing an order-two Hessian by the
order-five $A_3$ adds three powers of $\lambda$. The remaining insertions
have at least four powers of $a\lambda$, whether from a covariance or
from the next metric vertex. The number of annular momenta per physical
volume is at most $C\lambda^4$ for $\lambda L_\nu\ge1$. This gives
\eqref{eq:v23-y-bounds}. Momentum derivatives of the smooth profiles obey
the same scaled estimates. Finally \eqref{eq:v23-constant-O33} makes every
word of \eqref{eq:v23-y3} zero for a constant $h$.
\end{proof}

The proof also gives, for every fixed $b>4$,
\begin{equation}
 \|(y_{0,L}-y_{0,\infty})^{[d]}[h]\|_1
       \le C_{d,b}\chi^{-1}\lambda^{4-d}
              (\lambda L_{\min})^{-b}\|h\|.
 \label{eq:v23-y-volume}
\end{equation}
Indeed the integrands and all needed derivatives are smooth on a fixed compact
annulus, so Fourier integration by parts followed by Poisson summation gives
this estimate. The same argument with degree $7-d$ applies to $y_3$.
Windowed rooted kernel bounds follow from these fixed finite Taylor arrays.
Only finitely many lower-shell levels occur. None of these estimates removes
the lower Wilsonian scale or proves a Ward identity for a chosen finite-torus
interpolation.

At one-loop order, adding $\hbar C_{a,L}$ changes the canonical row by
$s_aC_{a,L}$. It does not change the already order-$\hbar$ reference trace
or the once-counted density at that order; their variations induced by the
order-$\hbar$ source change would be of order $\hbar^2$. The ordinary
one-loop descendants \emph{are} changed and remain those in the full bracket.

\section{What the first-background equation still has to retain}
\label{sec:v23-remaining}

The results above do not identify a constant-background test with the full
local one-background complex. At an arbitrary soft mark $q$, the free Euler
row is nonzero. In particular, let
\[
 G^{(2)}_{a,\loc}
   =\mathsf T_2[\Gamma_{a,1}]_{\eta\eta cc,q=0}
\]
be the actual quadratic-metric-antifield coefficient of V18. The canonical
one-background equation contains the prescribed return
\begin{equation}
 \left[
 \left\langle\mathcal E_0(q),
          \frac{\delta G^{(2)}_{a,\loc}}{\delta\eta}\right\rangle
 \right]_{\eta qcc}.
 \label{eq:v23-two-antifield-return}
\end{equation}
It is evaluated by the coefficient of
$-\Tr(C_gD)^2/4$, including both ghost returns inside $D$.
\Cref{thm:v23-zero-transfer} makes its zero-derivative input zero; its
first- and second-derivative local coefficients are not zeroed. The latter
can still supply a finite term in this equation. This is different from
the now-absent $a^3a^{-3}$ counterterm artifact.

The first-background consistency extraction also receives the free
$R_0c$ variation of coefficients with two ordinary metric marks. Therefore
a truncation to only zero and one ordinary background is not automatically
a subcomplex. Differentiating V22's $930$-component identity while freezing
these additional entries would miss actual canonical terms. Neither the
six-class computation nor a tensor product of its matrices with the ten
metric components is asserted to solve that enlarged equation.

The exact one-loop update to be used is
\begin{equation}
 \mathcal A_{a}^{\rm new}
       =\mathcal A_a^{\rm comb}+s_aC_{a,L},
 \label{eq:v23-full-update}
\end{equation}
with all its field and antifield components. The new results supply two
verified inputs to its first-background analysis: the native-to-ordinary
counterterm update in \eqref{eq:v23-descendant-bounds}, and the explicit
reference transport in \eqref{eq:v23-y-expansion}. They do not assert that
$[\mathcal A_a^{\rm new}]_{\eta qcc,\sigma qccc}$ vanishes. Its local primitive
must still be solved with \eqref{eq:v23-two-antifield-return}, the required
two-background coefficient, and the original measure/nonminimal terms.

No integral for those remaining local coefficients is evaluated numerically
in this installment. Their absence would not follow from a general anomaly
statement: local BRST cohomology and modified reference identities provide
established comparison frameworks \cite{V20BBH1994,V15IIM}, but the native
bounded projection and its populated source equations still have to be
verified. The new derivative floor narrows that task and removes a specific
putative finite return without assuming its conclusion.

\section{Verification and closing ledgers for V23}
\label{sec:v23-ledgers}

The companion \srcid{check_ESM_derivative_floor_V23.py} uses exact symbolic
and rational arithmetic. It checks the constant-slot cofactor/Cartan
parity, the actual filtered chart vertices and their constant-mark
cancellations, the source-trace reversal and ghost-factor identities,
the second-shift cubic coefficient, and the ordered first-background
reference replacement. It separately records the purely algebraic derivative
budgets. These checks are not a numerical evaluation of a native loop integral,
a machine-checked proof of the estimates, or a proof-assistant formalization.
The analytic proofs are the preceding arguments with their stated norms.

\subsection*{Proved}
\addcontentsline{toc}{subsection}{V23: Proved}
The native completed hard metric and ghost Hessians are even at one constant
metric mark, despite the absence of a general parity for the oriented cubic.
The measured minimal-source functional has zero zero-transfer coefficients
at every prescribed finite antifield degree, through one constant metric
mark. The actual combined V22 input has no derivative-degree-zero or
-degree-one coefficient. Its two specified lifts consequently have at least
one derivative and worst power $a^{-2}$. The exact descendant difference in
the first two antifield-free rows is bounded by $O(a^2\mu^7)$ and
$O(a\mu^6)$, respectively. The first possible one-background reference error
is the displayed $a^3$ insertion of the evaluated native cubic artifact; it
vanishes for a constant metric mark and has fixed lower-annulus bounds in
general.

\subsection*{Inherited}
\addcontentsline{toc}{subsection}{V23: Inherited}
The completed propagators and interactions, measured trace/Berezin calculus,
nonminimal sector, source/coordinate maps, auxiliary density, local source
banks, V20's paired-jet primitive, V21's rational affine contraction and V22's
zero-background consistency theorem are used at their original scope. Their
ranks, dimensions and abstract existence statements are not recomputed as
new results. The earlier $a^{-3}$ budget was conservative, not a proved
nonzero matching coefficient.

\subsection*{Literature input}
\addcontentsline{toc}{subsection}{V23: Literature input}
Local antifield cohomology and cutoff-modified BRST/Slavnov identities are
established methods \cite{V20BBH1994,V15IIM}. They supply the framework, not
the constant-background native cancellation or the derivative-floor estimate
proved here. No novelty claim for perturbative Einstein--matter EFT itself
is made.

\subsection*{Conditional}
\addcontentsline{toc}{subsection}{V23: Conditional}
A complete first-background local restoring theorem must combine the
available estimates with the actual higher-antifield and two-background
coefficients and a compatible finite local primitive. That theorem is not
inferred from the constant-source test. The improved counterterm budget
applies to the particular matching prescriptions used here, without an extra
derivative-free homogeneous source normalization. Different source choices,
noneven profiles or a different Nyquist convention require a separate test.

\subsection*{Open}
\addcontentsline{toc}{subsection}{V23: Open}
Assemble \eqref{eq:v23-full-update} at the next ordinary-background degree,
retaining \eqref{eq:v23-two-antifield-return} and the free-generator descent
from the two-background coefficient, and solve its remaining local
compatibility equation. Do not reintroduce a finite $a^3a^{-3}$ artifact
for the counterterm whose derivative-free part has just been proved zero.
Mixed matter/chiral quantum identities, higher loops, an invariant full ESM
shell class, and removal of the positive lower scale remain separate. No
nonperturbative measure, asymptotic safety, finite-parameter gravitational
ultraviolet completion or Yang--Mills mass-gap conclusion is drawn.

\begin{center}\small
\begin{tabular}{>{\raggedright\arraybackslash}p{.23\textwidth}>{\raggedright\arraybackslash}p{.68\textwidth}}
\toprule
Variable & Scope in V23\\
\midrule
Volume & Exact zero-transfer tests use paired finite momenta on odd grids.
Coefficient and rooted bounds are period independent for $\lambda L_\nu\ge1$;
thermodynamic comparison is separate.\\
Mesh & $a\lambda\le1/65536$. Native counterterm transport is $O(a)$ in the
first cubic-artifact descendant, not an all-mode smallness theorem.\\
Loop/source degree & One loop. Constant-mark zeros hold at each fixed finite
minimal-antifield degree; the counterterm and reference application uses one
minimal antifield and zero or one ordinary background.\\
Derivative order & Counterterm degrees one through three; local breaking and
reference degree at most four; every separately fixed symbol/kernel moment.\\
Lower scale/coupling & $\lambda\ge1024\mu>0$ and $\chi>0$ fixed. No infrared
removal or uniform estimate as $\chi\to0$.\\
Shell count & The first-background reference is confined to a fixed lower
annulus. No unrestricted interacting RG iteration follows.\\
Measure & The same complex metric reference, nonminimal fields, coframe
Jacobian and Lorentz/Haar auxiliary density are retained and counted once.\\
Symmetry & An actual derivative floor and a descendant-transport bound;
not a new nilpotence theorem for the finite ultraviolet source or a completed
one-background BV subcomplex.\\
\bottomrule
\end{tabular}
\end{center}


\clearpage
\part*{V24: A degree-closed local BV bank and explicit subcritical restoration}
\addcontentsline{toc}{part}{V24: Degree-closed local BV bank and scale-transgression restoration}

\section{Close the source degree before solving another background row}
\label{sec:v24-start}

Sections 1--184, including their proofs and qualifications, are preserved.
V23 identified a genuine obstruction to a background-only truncation:
free variation of a two-background coefficient and Euler variation of a
quadratic-antifield coefficient both enter the one-background equation.
The present installment closes that bookkeeping problem at its source.
The organizing cutoff is the \emph{total polynomial arity}, counting
ordinary fields, ghosts and antifields, together with a homogeneous
momentum weight. It is not a cutoff on the number of ordinary metrics.

There are three new conclusions. First, the ordinary zero-mesh gravitational
BV differential has a finite, degree-closed local bank and an explicit
bounded Taylor retraction on that bank. Second, the actual completed
regulator has no derivative-free minimal-source coefficients through the
constant-background orders needed by this bank. This gives an exact local
canonical cancellation, not merely an estimate of each ultraviolet term.
Third, differentiating the \emph{actual lower-reference} trace supplies a
local primitive for every subcritical weight. The remaining critical-weight
coefficient is retained; its primitive is not inferred from consistency.

The application below uses total arity at most five. This includes the
coupled one-background slots $\eta qcc$ and $\sigma qccc$, but also their
required two-background and quadratic-antifield parents. The calculation is
one-loop, gravitational, and at fixed positive lower scale. It is not a
combined internal-gauge/chiral Einstein--SM closure theorem. The supplied
programme permits an order-dependent finite local bank and requires the
local projection and its norm to be separate from anomaly cancellation
\cite{V6Programme}. The BRST and modified-reference identities used below
are established methods \cite{V20BBH1994,V20BBH1995,V15IIM}; their application
to this degree-closed native bank is the result here.

\subsection{The higher input jets are specified, not silently assumed}

Keep V14's reference, profiles, contour and upper interactions, V17's
degree-five extension, V18's source, and the once-counted auxiliary density.
Let $S=s(aD)$ be the existing filter. To determine all one-loop coefficients
of arity at most five, add just the following degree-six and -seven input
jets:
\begin{align}
 \widetilde S^g_{a,n}(q)
  &=\bigl[S_a^*-a\mathcal A_1-a^2\mathcal A_2\bigr]_n(Sq),
       &&n=6,7,\notag\\
 \widetilde S^{\rm gh}_{a,n}(q,c,\bar c)
  &=-\langle S\bar c,F R_{n-2}^{(0),E}(Sq,Sc)\rangle,
       &&n=6,7.
 \label{eq:v24-higher-input}
\end{align}
The functions $\mathcal A_1,\mathcal A_2$ are precisely
\eqref{eq:v17-generators}. The ordinary chart recursion supplies
$R_4^{(0),E},R_5^{(0),E}$. No new upper polynomial of degree six or seven
is added. Upper modes remain interacting through the already specified
lower-degree terms. Equation~\eqref{eq:v24-higher-input} is a recorded
extension of the completed regulator on its existing carrier. It does not
claim that the original microscopic action uniquely selected these upper
jets. Every action coefficient of degree at most five is unchanged.

The finite stationary recursion of V10 determines $[S_a^*]_6,[S_a^*]_7$:
at order $j$ solve the actual equation
$C_0A_j=-[b+(C-C_0)A+D_Ar_a]_j$, whose right side uses the preceding
$A_i$, and substitute into the original action. This is an explicit finite
coefficient prescription with the fixed $24\times24$ inverse, not an
assumption about unknown continuum vertices. V17 already proves the
fixed-degree mesh estimate for these coefficients. All six- and seven-leg
filtered subset sums are inside the nonaliased chart.

Use odd product grids and thermodynamic local normalization as before. A
single convenient domain for the comparisons is
\begin{equation}
 \lambda=\Lambda_0\ge4096\mu>0,\qquad
 a\lambda\le2^{-18},\qquad \lambda L_\nu\ge1.
 \label{eq:v24-domain}
\end{equation}
All statements are at fixed $\chi>0$, fixed profile seminorms, and fixed
source degree. The common factor $\hbar$ is suppressed in one-loop
coefficients. It is restored when a counterterm is inserted into the action.

\section{A finite complex and a common nonlinear Taylor projection}
\label{sec:v24-complex}

Work first in the ordinary zero-mesh local theory obtained from the same
stationary Palatini action and the same affine coframe chart. Its minimal
variables, ghost numbers, and auxiliary momentum weights are
\begin{center}\small
\begin{tabular}{c|rrrr}
 & $q$ & $c$ & $\eta=q^*$ & $\sigma=c^*$\\\hline
 ghost number & $0$&$1$&$-1$&$-2$\\
 momentum weight & $0$&$-1$&$2$&$3$
\end{tabular}
\end{center}
An ordinary derivative has weight one. These weights are bookkeeping
degrees, not additional physical scaling dimensions. In particular $\chi$
is held fixed, not assigned a new running law.

For a monomial with $b$ metric factors, $r$ metric antifields, $s$ ghost
antifields, $t$ ghosts, and $d$ derivatives, define
\begin{equation}
 A=b+r+s+t,\qquad g=t-r-2s,\qquad
 w=d+2r+3s-t=d+r+s-g.
 \label{eq:v24-gradings}
\end{equation}
$A$ counts a differentiated field once, not once per derivative. Here $N$
will denote an arity cap, not a perturbative loop order. Put
$n_*=r+s$. At one loop the common local bank is $w\le4$; equivalently its
Taylor degree in a ghost-number-$g$ slot is
\begin{equation}
 d_I=4+g-n_*.
 \label{eq:v24-degree-cap}
\end{equation}
A negative $d_I$ means that this slot has no assigned local Taylor term.
There is no bound on the number of metrics arising from $w$ alone. The
separate arity cutoff makes the bank finite.

Let $\mathfrak L^g_{N,4}$ be the space of translation-invariant local
polynomial functionals with $A\le N$ and $w\le4$, modulo total derivatives
and the actual tensor/Grassmann symmetries. Equivalently one may use labelled
Fourier coefficient polynomials on momentum conservation, with the induced
exchange symmetries. This formulation keeps source contacts as components
of a functional rather than independent normalizations of each row.

\begin{theorem}[Degree-closed ordinary gravitational source complex]
\label{thm:v24-complex}
Let $s_\infty=(\mathbb S_\infty,\cdot)$ be the ordinary classical BV
differential of the zero-mesh native Palatini action in this chart. Then
\begin{equation}
 d_N=\operatorname{pr}_{A\le N}s_\infty,
 \qquad d_N^2=0,\qquad
 d_N:\mathfrak L^g_{N,4}\longrightarrow\mathfrak L^{g+1}_{N,4}.
 \label{eq:v24-complex}
\end{equation}
The ideal of arity greater than $N$ is invariant. Consequently no coefficient
with more than $N$ total arguments can enter any of these classical local
rows, although a coefficient with more ordinary metric arguments can enter.
The statement also holds after adjoining the off-shell nonminimal doublet
and its antifields and after the fixed harmonic gauge-fermion transformation.
\end{theorem}
\begin{proof}
The native zero-mesh action begins at quadratic order on the flat,
zero-cosmological-term branch. Its action part has two derivatives. Its
minimal source terms are $\eta R(q,c)$ and $\sigma(c\cdot\partial c)$,
with one derivative; $R$ is the actual pointwise coframe-coordinate pullback
of the Lie variation. Every term of the full classical master action
therefore has weight two and arity at least two. The two canonical pairs
have total weights $0+2=2$ and $-1+3=2$. An antibracket reduces weight by
two and total arity by two. Hence bracketing with this action preserves
$w$ and cannot lower $A$. Ordinary diffeomorphism invariance and the
cotangent chart transformation give $(\mathbb S_\infty,\mathbb S_\infty)=0$;
this is the same zero-mesh action identification used in V11--V12, not an
identity asserted for the filtered finite source. Jacobi proves square zero
before quotienting, and invariance of the arity ideal proves it afterward.
Total derivatives form an invariant subspace because the evolutionary
variation commutes with ordinary differentiation.

For the nonminimal pair take weights $1,1,1,1$ on
$\bar c,b,\bar c^*,b^*$ respectively. The term $\bar c^*b$ has weight two
and arity two. The harmonic gauge fermion has weight two and begins at
arity two. Its canonical shift is linear in these fields and preserves
both filtrations. It conjugates, rather than deletes, their complex. At
fixed $N$ the derivative order is at most $4+N$, so all spaces involved are
finite dimensional. No cohomological vanishing follows just from this
finiteness.
\end{proof}

\subsection{The projection on coefficient germs}

Let $\mathfrak V^g_N$ consist of finite collections of smooth Fourier
coefficient germs of arity at most $N$ at zero external momenta. The germs
are taken on the common thermodynamic momentum convention; arbitrary
independent off-grid interpolations of different rows are not identified.
For each slot $I$ set
\begin{equation}
 (\mathsf P_{N,4}F)_I=\mathsf T_{d_I}F_I,
 \qquad \mathsf T_d=0\quad(d<0).
 \label{eq:v24-projector}
\end{equation}
The ordinary differential on germs uses the same finite local vertices and
canonical operations as on local polynomials.

\begin{theorem}[A bounded retraction for all coupled rows at once]
\label{thm:v24-projection}
The map in \eqref{eq:v24-projector} is a retraction onto the local bank and
\begin{equation}
 \boxed{\mathsf P_{N,4}d_N=d_N\mathsf P_{N,4}.}
 \label{eq:v24-chain-projection}
\end{equation}
It therefore splits the short exact sequence of these \emph{ordinary
coefficient complexes}. In scaled coefficient-jet norms its norm is at most
one; inclusion of the local polynomial in the corresponding scaled smooth
norm has norm at most $2^{N+4}$. These constants are independent of the
mesh, periods and number of integrated shells. They are not uniform in $N$.
\end{theorem}
\begin{proof}
At fixed output arity, only finitely many local vertices of
$\mathbb S_\infty$ act. Each contribution is a finite linear combination
of a homogeneous polynomial in the external momenta times an input germ
whose arguments have been replaced by linear subset sums. Its output and
input weights agree. Homogeneous Taylor expansion commutes with a linear
substitution of momenta, and multiplication by a degree-$j$ polynomial
changes the allowed Taylor order by exactly $j$. Extracting $w\le4$ before
or after this operation gives the same result. This also proves descent
under exchange symmetries and integration by parts. The arity truncation
commutes for the reason in \cref{thm:v24-complex}.

For explicit norms, write $K$ for independent labelled momenta and use
$\sum_{|\alpha|\le d}\mu^{|\alpha|}\sup|\partial^\alpha F(K)|/\alpha!$
on a cube of radius $\mu$, summed over the retained component slots.
Use $\sum_\alpha\mu^{|\alpha|}|F_\alpha|$ on Taylor coefficients.
Evaluation at zero proves the retraction bound. On a monomial
$F_\alpha K^\alpha$, the sum of all scaled derivative norms is bounded
by $2^{|\alpha|}\mu^{|\alpha|}|F_\alpha|$, by the binomial identity.
Since $d_I\le N+4$, this proves the inclusion bound. Quotient norms may
be used for tensor symmetries and total derivatives; the same estimates
descend. Polynomial coordinate changes and their cotangent lifts preserve
weight and do not lower arity, so they give the same intertwining statement
with all induced contacts. This is not a claim of a uniform projection on
an unbounded-order or high-momentum function space.
\end{proof}

This projection is onto the \emph{full component counterterm complex}, not
onto its invariant physical operators. It constructs no primitive for a
nontrivial cocycle. It also does not turn $s_a^2\ne0$ for the finite filtered
source into $s_a^2=0$.

\section{The actual higher constant-background source coefficients have a derivative floor}
\label{sec:v24-floor}

The new bank consumes the formerly unestimated zero-transfer slots
$\eta q^2c$, $\eta q^3c$ and $\sigma q^2cc$. Merely bounding them by their
superficial ultraviolet degrees would permit an artificial derivative-free
return. A constant-background identity decides this question.

Let $\bar q$ be a constant formal coframe background, retaining its degree
only as far as the chosen finite input jets allow. At the needed orders
$e(\bar q)$ is constant and the actual Cartan load is zero. The stationary
connection is therefore zero. For a hard metric perturbation $x$ the
stationary connection starts at order $x$.

\begin{lemma}[Even hard Hessians through the required constant-background jet]
\label{lem:v24-constant-hessian}
For the action \eqref{eq:v24-higher-input}, the hard metric and ghost
Hessians, their inverses with the lower reference, and all their
constant-background derivatives needed at external arity at most five are
even under $p\mapsto-p$. In the same jet the actual ghost-action return at
constant retained odd ghost $u$ is
\begin{equation}
 B_u(\bar q;p)x_p=-s(ap)^2\alpha_pF_px_p,
 \qquad \alpha_p=i\,u\cdot p.
 \label{eq:v24-constant-return}
\end{equation}
In particular this return has no constant-background derivative.
\end{lemma}
\begin{proof}
At fixed constant coframe the hard Hessian of the reduced native action is
\[
 -b'_a(\bar q)^TC(\bar q)^{-1}b'_a(\bar q),
\]
with the fixed normalization understood. Each $b'_a$ has one odd phase/time derivative;
the Cartan inverse is pointwise and independent of $p$. The complete
shared-link remainder has at least three stationary connection factors,
and hence at least three hard fields at this background. It contributes
nothing to this second hard variation, even after any number of constant
background differentiations. This is stronger than assigning a parity to
the general oriented interaction.

In V17's restoring functionals, $X(\bar q)=0$, so the terms containing
the cubic and fourth connection words and the squared cubic-force return
again have hard degree at least three. The surviving hard quadratic
correction is the geometric $\langle b_2,X\rangle$ term; its derivatives
have orders three and one and have even simultaneous inversion parity.
The even profile on each hard leg preserves the result. The explicit
upper terms in V14 and V17 have two opposite forward symbols and constant
coefficient insertions, so are even as well. This proves the metric claim
coefficient by coefficient, including the degree-six and -seven extensions.

For ordinary constant $u$, pointwise chain differentiation gives the exact
identity $R^{(0),E}(q,u)=u\cdot Dq$ at every field degree, because it is the
pullback of translation by the same local coordinate map. Every nonlinear
coefficient beyond the linear one vanishes. Filtering the hard legs leaves
exactly \eqref{eq:v24-constant-return}; all optional upper ghost terms have
$(I-S)u=0$. The constant-background ghost Hessian contains two ordinary
momentum derivatives, so is even. Inversion and repeated differentiation
of an even matrix symbol preserve evenness without commuting its factors.
All arguments are identities of finite formal background jets.
\end{proof}

\begin{theorem}[Zero-transfer minimal sources on the closed arity-five bank]
\label{thm:v24-floor}
For every nonzero minimal-antifield sector of total external arity at most
five, all its constant-source coefficients vanish in the completed finite
one-loop functional:
\begin{equation}
 \boxed{[\Gamma_{a,1}]_{I,K=0}=0\qquad(n_*(I)\ge1,\ A(I)\le5).}
 \label{eq:v24-source-floor}
\end{equation}
In particular the zero-derivative parts of $\eta q^2c$, $\eta q^3c$,
$\sigma q^2cc$, $\eta^2qcc$ and $\eta\sigma ccc$ are zero. No assertion
about their first or higher momentum derivatives is made.
\end{theorem}
\begin{proof}
At constant $\bar q$ the hard-quadratic Berezin reduction is the same one
as V18, with the actual background covariances. The returned hard ghost
has the form $\beta_p=\alpha_pV_p(\bar q)x_p$, with $V_p$ even in the
Grassmann grading and independent of $u$. Since
$\alpha_{-p}=-\alpha_p$ and $\alpha_p^2=0$, the dressed ghost-antifield
quadratic form and every closed word with at least two minimal antifields
vanish at zero transfer.

For one metric antifield, the direct quadratic contact vanishes before
tracing. To see this without a background expansion, write
$A_\circ=DG(\bar q)$. At opposite hard momenta the quadratic part of
$DG(\bar q+x)^{-1}P\mathcal L_u(G(\bar q+x)-I)$ contains a term with
zero total derivative and the two terms
$-A_\circ^{-1}\mathsf B(x,A_\circ^{-1}P\mathcal L_uA_\circ x)$.
The latter cancel upon polarizing the two hard fields, since $A_\circ$
is constant, $P(p)=P(-p)$ and $\mathsf B$ is symmetric. The remaining
metric-antifield contribution has one-derivative source kernel $L_\eta$,
two-derivative return $B_u$, and even metric and ghost covariances. Its
trace is odd in $p$ and sums to zero on the paired odd grids. These
statements hold coefficientwise in the constant background. Extracting
precisely the powers available in \eqref{eq:v24-higher-input} proves
\eqref{eq:v24-source-floor}. No covariance or high-frequency mode has
been replaced by a scalar surrogate.
\end{proof}

\section{The finite measured bank and its mixed-antifield parent}
\label{sec:v24-parents}

The structural component types at exact arities four and five are shown
below. Tensor components, derivative distributions and Grassmann exchange
symmetries are still retained inside each entry.
\begin{center}\small
\begin{tabular}{c|p{.35\textwidth}|p{.40\textwidth}}
 & $A=4$ & $A=5$\\\hline
 $g=0$ & $q^4,\ \eta q^2c,\ \sigma qcc,\ \eta^2cc$
 & $q^5,\ \eta q^3c,\ \sigma q^2cc,\ \eta^2qcc,\ \eta\sigma ccc$\\[2mm]
 $g=1$ & $q^3c,\ \eta qcc,\ \sigma ccc$
 & $q^4c,\ \eta q^2cc,\ \sigma qccc,\ \eta^2ccc$\\[2mm]
 $g=2$ & $q^2cc,\ \eta ccc$
 & $q^3cc,\ \eta qccc,\ \sigma cccc$
\end{tabular}
\end{center}
For a ghost-number-zero slot the local Taylor orders are four, three and
two at $n_*=0,1,2$ respectively. The corresponding breaking orders are
five, four and three. The lower arities are included, not replaced by this
table. At all ghost numbers the finite bank is determined by
\eqref{eq:v24-gradings}, not by a guessed list.

\begin{proposition}[Finite input budget for the complete one-loop bank]
\label{prop:v24-input-budget}
On \eqref{eq:v24-domain}, all one-loop minimal coefficients of arity at
most five are determined by the ordinary action and source jets of arity
at most seven, together with the inherited covariances and auxiliary
measure. A finite prescription for them is
\begin{equation}
 \Gamma_{a,1}^{\le5}
 =\operatorname{pr}_{A\le5}\left\{
   \frac12\sum_{r=1}^{5}\frac{(-1)^{r+1}}r
       \operatorname{STr}(\mathbb C_aV_a)^r+\mathcal M_a(q)\right\}.
 \label{eq:v24-loop-compiler}
\end{equation}
Here $V_a$ is the actual background-dependent graded Hessian minus its
zero-background value; the hard nonminimal Gaussian is eliminated exactly
and its off-shell contact is restored as in V16. Every ghost, metric and
antifield Hessian entry is differentiated from the specified common action.
\end{proposition}
\begin{proof}
A hard-linear vertex at this external order has a momentum equal to a sum
of at most five soft momenta and therefore has no hard covariance support.
After the exact algebraic nonminimal integration, every remaining vertex
has at least two hard legs. For a connected one-loop graph, $I=V$, so
summing hard valences forces every vertex to have exactly two hard legs.
A vertex with at most five external legs consequently needs at most seven
input legs. The Hessian perturbation has positive external arity, so the
logarithm stops at five at the required order. The supertrace implements
the already fixed metric and ghost signs. This is a use of the inherited
measured compiler, now on a degree-closed bank. It is not an additional
all-source nonlocal cutoff theorem.
\end{proof}

A previously unextracted parent is particularly important. V18's actual
zero-background quadratic forms give
\begin{equation}
 \mathcal F_{\eta\sigma,a}
   =-\frac12\operatorname{Tr}
           (C_gD_\eta C_gD_\sigma)\big|_{\eta\sigma ccc},
 \qquad \mathcal F_{\eta\sigma,a}^{\loc}=\mathsf T_2\mathcal F_{\eta\sigma,a}.
 \label{eq:v24-mixed-parent}
\end{equation}
This is the coefficient of the inherited trace, including the two ghost
covariances within $D_\sigma$ and the ghost return within $D_\eta$.
It is not an independently adjustable tensor.

To display its descent, root its density as
$\eta_A\sigma_\alpha U^{A\alpha}(c,\partial c,\ldots)$, with the ordering
shown. The free complementary BV operator obeys
$\delta_0\eta_A=\mathcal E_{0,A}(q)$ and
$\delta_0\sigma_\alpha=\mathcal B_{0,\alpha}(\eta)$ in the inherited sign
convention. Thus
\begin{equation}
 \boxed{\delta_0(\eta_A\sigma_\alpha U^{A\alpha})
  =\mathcal E_{0,A}(q)\sigma_\alpha U^{A\alpha}
      -\eta_A\mathcal B_{0,\alpha}(\eta)U^{A\alpha}.}
 \label{eq:v24-mixed-descent}
\end{equation}
Derivatives on an antifield are transported by the same integrations by
parts. The two outputs are respectively $\sigma qccc$ and $\eta^2ccc$.
This explains why retaining only the $\eta^2cc$ Euler return from V23
cannot close the requested pair of one-background equations. The
$\eta^2qcc$ coefficient and the $\eta q^3c,\sigma q^2cc$ coefficients are
also determined by \eqref{eq:v24-loop-compiler} and are not frozen during
variation. Equation~\eqref{eq:v24-mixed-descent} is a new extraction of an
inherited coefficient, not a claim to have evaluated its momentum integral.

For example, for any independent odd $c_0,c_1,c_2$ and even component
coefficient $u$, the ordered monomial $\eta_A\sigma_\alpha u c_0c_1c_2$
has both displayed descendants. Choosing components where the free Euler
and adjoint-generator maps are nonzero gives two nonzero polynomial
summands; neither is removed by Grassmann antisymmetry. The verifier checks
their exact relative sign and square-zero relation separately.

\section{The ultraviolet canonical coefficients cancel in the complete local bank}
\label{sec:v24-cancellation}

Write $d=d_5$, $P_4=\mathsf P_{5,4}$, and use the minimal restriction after
restoring the off-shell nonminimal fields. Let
\begin{equation}
 g_a=P_4\Gamma_{a,1}^{\le5},\qquad
 j=P_4\sum_x\log\det DG(q_x),\qquad
 x_a=P_4\mathcal X_a,\qquad
 B_a=P_4\mathcal A_a^{\rm comb}.
 \label{eq:v24-local-arrays}
\end{equation}
The definition of $\mathcal A_a^{\rm comb}$ is exactly
\eqref{eq:v22-combined-definition}. In particular $\mathcal M_a$ remains
once in $\Gamma_{a,1}$ as its recorded action insertion. The term $j$
in this equation records the \emph{flat divergence} of the transformed
vector field; it is not a second choice of physical density.

\begin{theorem}[Native local canonical cancellation at arity five]
\label{thm:v24-canonical}
On the common interpolation and normalization convention,
\begin{equation}
 \boxed{P_4(s_a\Gamma_{a,1})=d g_a,\qquad
 B_a=d(g_a-j)-x_a,\qquad dB_a=-d x_a.}
 \label{eq:v24-canonical-cancellation}
\end{equation}
These equations concern the specified minimal local bank. They do not
assert nilpotence of $s_a$ or closure of its full finite-cutoff functional.
\end{theorem}
\begin{proof}
On sufficiently small external momenta all source profiles and their
subset arguments equal one. The source part of $s_a$ is therefore exactly
the ordinary chart source in each Taylor jet being extracted. The
nonminimal doublet can be restored before restriction: the loop coefficient
has no $b^*$ dependence; a $\bar c^*$ insertion is linear in the auxiliary
field and cannot occur in a one-loop minimal coefficient. Terms involving
$s\bar c=b$ vanish on the restriction. The fixed gauge-fermion shift is
undone on both sides, not on just the metric Hessian.

The only possible difference from $s_\infty$ on the minimal local jets is
an action-Euler insertion. Its quadratic difference begins at momentum
degree six. Every higher corrected native vertex in use begins its
ordinary-momentum difference at degree five, by V17 and
\eqref{eq:v24-higher-input}. Suppose a loop coefficient has $n_*$
antifields. Euler variation removes one, so the output breaking has Taylor
cap $6-n_*$. For $n_*\ge2$, a degree-five or -six difference already
exceeds this cap. For $n_*=1$, the only possible contribution at the cap is
the degree-five cubic-or-higher artifact times the zero-derivative loop
source. That coefficient is zero by \cref{thm:v24-floor}, including the
needed two- and three-background slots. The degree-six free difference
never contributes. Terms with more than five total arguments cannot descend
back by \cref{thm:v24-complex}. This proves the first identity, using the
measured derivative floor rather than a bound on divergent local terms.

V18's exact coordinate divergence is
$\operatorname{sdiv}r_a=-r_a\log J_G$, since the metric-coordinate and
ghost-coordinate flat divergences vanish. Its low external jet is
therefore $-dj$. Inserting this and the first identity into the fixed
measured definition proves the second. The last follows from $d^2=0$.
The density $\mathcal M_a$ can contain further local and nonlocal pieces;
none was discarded in $g_a$. No ultraviolet estimate for those pieces is
needed to take the algebraic square.
\end{proof}

The particular source counterterm $C_a^{\rm old}$ in V21--V23 has weight
at most four, arity at most four, and no zero-derivative part. The same
degree test proves
\begin{equation}
 P_4s_a C_a^{\rm old}=d C_a^{\rm old},\qquad
 B_a^{\rm old}=d(g_a-j+C_a^{\rm old})-x_a.
 \label{eq:v24-old-update}
\end{equation}
All descendants, including both terms of \eqref{eq:v24-mixed-descent},
are included in $d$. A zero- or one-background truncation would not permit
this cancellation.

\section{A lower-scale derivative supplies an explicit local primitive}
\label{sec:v24-transgression}

Let $x_{0,\lambda}=P_4\mathcal X_{0,\lambda}$ use the ordinary zero-mesh
vertices of the native action and source, and the fixed profile
$\Theta_\lambda(p)=\theta(p/\lambda)$. This is a reference insertion, not
a full ultraviolet determinant. Its ordered expansion at this arity is
\begin{equation}
 x_{a,\lambda,L}
 =P_4\operatorname{pr}_{A\le5}
  \operatorname{STr}\left[
       \sum_{r=0}^{5}(-\mathbb C_a V_a)^r\Theta_\lambda Dr_a\right].
 \label{eq:v24-reference-compiler}
\end{equation}
The $r=0$ term has no antifield-dependent coefficient. Its source-free
part is retained as a regular trace of $\Theta Dr$, supported in the lower
ball and containing no inverse propagator. Every remaining nontrivial word
contains a covariance and is routed through $\Theta$ to a lower annulus.
The same formula defines the zero-mesh thermodynamic coefficient.

Set $\dot C_0=\lambda\partial_\lambda C_0
=-(\lambda\partial_\lambda\Theta_\lambda)\mathbb H_0^{-1}$, with its
value zero near the retained origin. Define the ghost-number-zero local
functional
\begin{equation}
 \begin{split}
 Y_\lambda=P_4\operatorname{pr}_{A\le5}
 \frac12\sum_{r=1}^{5}\frac{(-1)^{r+1}}r
 \sum_{i=0}^{r-1}\operatorname{STr}
 \left[(C_0V_0)^i(\dot C_0V_0)
                         (C_0V_0)^{r-1-i}\right].
 \end{split}
 \label{eq:v24-Y}
\end{equation}
This is a finite lower-annulus integral, with every matrix in its original
order. No inverse of $1-\Theta$ occurs. The auxiliary density is independent
of $\lambda$ and is not differentiated a second time.

\begin{theorem}[Native reference-scale transgression]
\label{thm:v24-transgression}
In thermodynamic normalization,
\begin{equation}
 \boxed{d x_{0,\lambda}=0,\qquad
       dY_\lambda=\lambda\partial_\lambda x_{0,\lambda}.}
 \label{eq:v24-transgression}
\end{equation}
For its homogeneous weight-$w$ component the actual integral has the exact
scaling
\begin{equation}
 x_{0,\lambda}^{[w]}=\lambda^{4-w}x_{0,1}^{[w]},\qquad
 Y_\lambda^{[w]}=\lambda^{4-w}Y_1^{[w]}.
 \label{eq:v24-reference-scaling}
\end{equation}
The rescaling holds at fixed $\chi$ and fixed dimensionless profile, with
component-dependent powers of $\chi$ retained. Consequently, for every
integer $w<4$,
\begin{equation}
 \boxed{H_\lambda^{[w]}=\frac{Y_\lambda^{[w]}}{4-w},\qquad
                         dH_\lambda^{[w]}=x_{0,\lambda}^{[w]}.}
 \label{eq:v24-primitive}
\end{equation}
This is a ghost-number-zero local primitive for the entire coupled
subcritical coefficient, not separate primitives for independently chosen
rows. At $w=4$ the conclusion is only $dY_\lambda^{[4]}=0$; division by
$4-w$ is not allowed.
\end{theorem}
\begin{proof}
The zero-mesh action is the same ordinary, antifield-linear master action
used in V22's trace lemma, with its nonminimal doublet restored. Its
classical square is zero in the formal local theory. The differentiated
master equation gives the regulated trace identity
$s\Gamma_1^{\mathbb R}+\operatorname{sdiv}r=\mathcal X_{\mathbb R}$.
Applying $s$ and using the superdivergence identity for a square-zero
vector field gives $s\mathcal X_{\mathbb R}=0$. The proof uses all
antifield variations of the Hessian; dropping them would not give this
identity. This is the inherited finite trace lemma, now applied before
extracting any background or antifield subspace.

To justify its present continuum use, form the ordered words first and
cancel $C\mathbb R=\Theta$. Every reference word containing a covariance has support at zero
external momentum in $\lambda\le|p|\le2\lambda$. A small external
neighborhood enlarges it only to a fixed compact lower annulus. The
no-covariance source-free trace is regular on its compact lower ball and
is kept separately.
All derivatives needed here are integrable, and changes of routed momentum
are legitimate with a smooth compact guard. No ultraviolet boundary term
survives the guard. The same statements hold for the scale derivative of
the determinant expansion because each such word contains $\dot C_0$,
which is supported in that annulus. Thus one may differentiate the trace
identity in $\lambda$ without giving a value to an unsubtracted infinite
ultraviolet determinant. The scale derivative of the source-dependent
logarithm is exactly \eqref{eq:v24-Y}; the flat divergence does not depend
on $\lambda$. Apply the common chain projection of
\cref{thm:v24-projection} to obtain both equations in
\eqref{eq:v24-transgression}.

For scaling, each propagator has momentum degree minus two, each ordinary
metric or ghost-action Hessian vertex degree two, and each minimal-antifield
vertex degree one. A loop with $n_*$ minimal antifields therefore has degree
$4-n_*$. A reference insertion has one additional derivative, giving degree
$5-n_*$. At external Taylor degree $d$, these are respectively
$4-n_*-d$ and $5-n_*-d$. By \eqref{eq:v24-gradings} each equals $4-w$
in its own ghost-number sector. The algebraic nonminimal elimination
preserves this count. The change of variables $p=\lambda z$ in the
compact, thermodynamic integrals proves \eqref{eq:v24-reference-scaling}.
There is no mass, external momentum or second scale left in their Taylor
coefficients. Holding a fixed torus instead would invalidate this exact
rescaling and is not done here. Finally $d$ preserves $w$, so
\eqref{eq:v24-transgression} gives
$dY^{[w]}=(4-w)x^{[w]}$, proving the displayed primitive.
\end{proof}

This construction uses the actual ordinary reference integrals as finite
matching data. It does not determine their numerical values, and it does
not replace them by zero. The denominators are nonzero positive integers
for the subcritical bank. The map has norm at most one on the direct sum
of its weight components before the fixed trace-compiler norm is applied.
For a slot with $n_*$ antifields and derivative degree $d$, the coefficients
of $Y$ are bounded by
$C_I\chi^{-n_*}\lambda^{4-n_*-d}$; those of $x$ are bounded by
$C_I\chi^{-n_*}\lambda^{5-n_*-d}$. Rescaling the compact annulus and using
the inherited free inverse margin proves these bounds directly. Their
constants depend on the finite jet order and the profile, not on the mesh,
volume or ultraviolet shell count.

\section{Subcritical coupled matching and the remaining finite critical block}
\label{sec:v24-matching}

\subsection{The reference comparison uses only finitely many lower-scale words}

\begin{proposition}[Full-bank lower-reference comparison]
\label{prop:v24-reference-rate}
For the degree-five bank and \eqref{eq:v24-domain}, a minimal component
with $n_*$ antifields and derivative degree $d\le5-n_*$ satisfies
\begin{equation}
 \left|(x_{a,\lambda,L}-x_{0,\lambda,\infty})_I^{[d]}\right|
 \le C_{I,b}\chi^{-n_*}\lambda^{5-n_*-d}
       \left[(a\lambda)^3+(\lambda L_{\min})^{-b}
              +\frac{\mathbf1_{n_*=0}}{\lambda^4V_L}\right],\qquad b>4.
 \label{eq:v24-reference-rate}
\end{equation}
The estimate includes the needed two-background and two-antifield slots.
It does not assert that their critical local values are zero. The last
term conservatively retains the omitted zero-mode contact in a source-free
trace; it is absent for every antifield-containing word.
\end{proposition}
\begin{proof}
Expand \eqref{eq:v24-reference-compiler} at total arity at most five. A
word has at most five background vertices. In covariance-bearing words,
routing through $\Theta$ and using the lower-cutoff covariance confines
every momentum to
$\lambda/2\le|p|\le3\lambda$, up to harmless fixed margins. All field
and source profiles equal one on these arguments by
\eqref{eq:v24-domain}. The difference of free inverses has relative size
$O((a\lambda)^4)$, as proved in V22. Each restored interaction Hessian
has relative discrepancy $O((a\lambda)^3)$ by V17 and
\eqref{eq:v24-higher-input}; source vertices on this annulus are exactly
the ordinary coordinate polynomials. Telescoping the finite ordered product
therefore bounds its difference by $C(a\lambda)^3$ times its ordinary
scaling. At degree $d$ this scaling is
$\chi^{-n_*}\lambda^{5-n_*-d}$, by the same count used in
\cref{thm:v24-transgression}. Differentiated inverse identities and compact
profile derivatives give the bounds for every prescribed Taylor derivative.
The normalized finite sum has $O(\lambda^4)$ points per physical volume
under $\lambda L_\nu\ge1$. For the remaining volume comparison apply
Poisson summation to the compactly supported smooth zero-mesh integrand
and differentiate it more than $b+4$ times. This gives the displayed
superalgebraic period error. The no-covariance term is a smooth polynomial trace on the lower ball;
its comparison has the same scaling. Omitting its retained zero mode from
the finite trace changes its normalized sum by at most
$C\chi^{-n_*}\lambda^{5-n_*-d}/(\lambda^4V_L)$, where here $n_*=0$.
Covariance-bearing words vanish in a neighborhood of that mode and have no
such correction. No estimate of a divergent ultraviolet local coefficient
is used.
\end{proof}

For clarity a dimensionless coefficient norm on weight $w$, ghost number
$g$, is obtained by multiplying each component by
$\chi^{n_*}\lambda^{w-4}$, then taking its finite coefficient $\ell^1$
norm. In these norms $d$ has the fixed norm of a finite matrix determined
by the ordinary coframe source and the native quadratic/cubic vertices.
It has no mesh or volume dependence. This statement does not assign the
49/33 affine norms of V21 to the larger complex.

\begin{theorem}[An explicit coupled subcritical local counterterm]
\label{thm:v24-matching}
Let $C_a^{\rm old}$ denote the previously chosen source counterterm, with
its full descendants retained. Define an \emph{additional}, ghost-number-zero
local insertion by
\begin{equation}
 \boxed{\Delta C_{a,<4}
   =-\bigl(g_a-j+C_a^{\rm old}\bigr)_{w<4}
        +\sum_{w<4}\frac{Y_\lambda^{[w]}}{4-w}.}
 \label{eq:v24-counterterm}
\end{equation}
Then the updated complete local breaking satisfies
\begin{equation}
 \boxed{(B_a^{\rm new})_{w<4}
    =(x_{0,\lambda,\infty}-x_{a,\lambda,L})_{w<4}.}
 \label{eq:v24-restoration}
\end{equation}
In the scaled finite coefficient norm this is bounded by
$C_b[(a\lambda)^3+(\lambda L_{\min})^{-b}+(\lambda^4V_L)^{-1}]$.
The last term is absent in the antifield-containing rows. In particular the coupled
$\eta qcc$ and $\sigma qccc$ equations are restored through external
derivative degree three, with this explicit error. Their degree-four
critical coefficients are not canceled by this theorem.
\end{theorem}
\begin{proof}
Put $U_a=g_a-j+C_a^{\rm old}$. Equations
\eqref{eq:v24-canonical-cancellation}--\eqref{eq:v24-old-update} give
$B_a^{\rm old}=dU_a-x_a$. Since $d$ preserves weight, adding the ordinary
variation of \eqref{eq:v24-counterterm} leaves exactly
$-x_a^{<4}+x_0^{<4}$. The native projected variation equals this ordinary
variation as well: the one-antifield part of $g_a$ has no degree-zero
coefficient by \cref{thm:v24-floor}; the same holds for $Y_\lambda$ by
differentiating that zero-transfer identity in $\lambda$. The old source
counterterm has the V23 floor. For two or more antifields the degree-five
Euler discrepancy exceeds the allowed output order without using a floor.
A source-free counterterm varies only through the common soft ordinary
transformation. Thus the degree argument in \cref{thm:v24-canonical}
applies to $\Delta C$ and proves \eqref{eq:v24-restoration} for the actual
measured local update. The quantitative estimate follows from
\cref{prop:v24-reference-rate}.
\end{proof}

Equation~\eqref{eq:v24-counterterm} is a new \emph{common one-loop
normalization}, not a source-only change. Its source-free components assign
the actual power-divergent local action coefficients too. It leaves the
tree action, reference propagators and declared contour unchanged but
changes the order-$\hbar$ local action and transformation data together.
The subtraction of $C_a^{\rm old}$ inside the increment prevents adding
its contribution twice. Any desired subcritical homogeneous physical
normalization can be added as one $d$-closed local functional. Arbitrarily
independent constants in each displayed row are not licensed.

The primitive map is bounded; the raw counterterm coefficients need not
be small. Source-free coefficients can contain $a^{-4}$ terms, and the
one-antifield coefficients retain their allowed $a^{-2}$ and $a^{-1}$
powers. These are stored renormalization data, not an invariant analytic
neighborhood of microscopic actions. At fixed source window, every
polynomial insertion has a finite rooted distributional/source-jet norm.
No summation of an extensive vacuum density is used to bound a local
observable.

\subsection{What has become the genuinely remaining finite problem}

All positive-power scale factors of the ordinary reference breaking have
been inverted by \eqref{eq:v24-primitive}. The critical block obeys
\begin{equation}
 d x_{0,\lambda}^{[4]}=0,\qquad
 \lambda\partial_\lambda x_{0,\lambda}^{[4]}=0,
 \qquad dY_\lambda^{[4]}=0.
 \label{eq:v24-critical}
\end{equation}
None implies that $x_{0,\lambda}^{[4]}$ is exact. The concrete next
calculation is the finite equation
\begin{equation}
 d C_{\rm ref}^{[4]}=x_{0,\lambda}^{[4]}
 \quad\text{in }\mathfrak L^1_{5,4}\text{ at weight four},
 \label{eq:v24-next}
\end{equation}
with the already established lower-arity normalizations and the actual
coefficients of \eqref{eq:v24-reference-compiler}. Its one-background
slots include both descendants of \eqref{eq:v24-mixed-parent}. If a
primitive fails, a left-null functional of this \emph{complete} finite
matrix evaluated on the actual reference coefficient is the relevant
obstruction; the six-class affine quotient alone is not that test.

Arity above five cannot enter \eqref{eq:v24-next}. This is the precise
sense in which the upstream change of grading stops the repeated demand
for another background row. It does not remove the need to solve the
remaining finite critical system. Subcritical restoration cannot change
that system because $d$ preserves $w$. Higher-loop BV Laplacians may
consume additional arities and must be budgeted separately; no all-loop
closure is inferred from this one-loop bank.

\section{Verification and closing ledgers for V24}
\label{sec:v24-ledgers}

The accompanying script \srcid{check_ESM_degree_closed_BV_V24.py}
verifies the weight and arity rules, enumerates the component-type tables,
checks the exact relative signs of the mixed-antifield descent, tests the
constant-background coordinate identities on nontrivial rational matrices,
and verifies the scale-transgression cancellation on an explicit finite
reference example. It also checks that every possible native Euler/Taylor
spill at arity five is either beyond the cap or uses precisely a
zero-derivative one-antifield coefficient. The script does not compute the
rank of the full critical local complex, evaluate the native annular
integrals, or prove their analytic bounds. Those distinctions are included
in its machine-readable report. No proof-assistant formalization is claimed.

\subsection*{Proved}
\addcontentsline{toc}{subsection}{V24: Proved}
The ordinary native gravitational BV complex admits the arity/weight cutoff
\eqref{eq:v24-gradings}, the finite bank \eqref{eq:v24-degree-cap}, and the
bounded common Taylor retraction \eqref{eq:v24-chain-projection}. The
declared six/seven input jets determine the complete one-loop arity-five
bank. Its actual constant-background minimal coefficients have the required
derivative floor. The local canonical ultraviolet coefficients then cancel
in consistency before estimation. The lower-reference scale derivative
\eqref{eq:v24-Y} is a finite, pole-free primitive for every weight below
four, and the additional common normalization
\eqref{eq:v24-counterterm} restores those coupled local rows with the
stated third-order lower-scale cutoff error. The mixed-antifield parent
has both the one-background and higher-antifield descendants displayed.

\subsection*{Inherited}
\addcontentsline{toc}{subsection}{V24: Inherited}
V1--V23 remain unchanged. V17 provides the all-finite-degree native mesh
functionals and stationary recursion; V14 and V17 provide the lower-degree
periodic completion; V18 provides the finite source and Berezin operands;
V20--V23 provide the particular source counterterms and their lower-arity
properties. The pointwise coframe map, auxiliary measure, off-shell
doublet, zero-mesh master action and regulated trace lemma are inherited
with their existing hypotheses. Their general Gaussian, cohomological and
source-transport principles are not presented as new discoveries.

\subsection*{Literature input}
\addcontentsline{toc}{subsection}{V24: Literature input}
Local BRST cohomology, antifield gradings, and modified cutoff Slavnov
identities are established frameworks \cite{V20BBH1994,V20BBH1995,V15IIM}.
No general novelty claim for perturbative gravity, homogeneity arguments or
one-loop renormalization is made. The new application is the stated native
finite bank, its all-constant-background floor, and the explicit
reference-matched subcritical source-action insertion.

\subsection*{Conditional}
\addcontentsline{toc}{subsection}{V24: Conditional}
Critical local restoration depends on the populated finite equation
\eqref{eq:v24-next}, not merely its consistency. Matching the enlarged
smooth regulator to the unmodified microscopic action, joint internal-gauge
and chiral compatibility, arbitrary-order source estimates and full
interacting R1--R4 invariance retain their earlier requirements. The
one-loop extension \eqref{eq:v24-higher-input} is additional recorded
regulator data. No nonlocal cutoff theorem for all new metric-only
four/five-source coefficients is asserted here.

\subsection*{Open}
\addcontentsline{toc}{subsection}{V24: Open}
Compute the critical weight-four reference array and its complete
arity-five local primitive, with the lower-arity choices respected. The
required higher-antifield parents and two-background entries now belong to
one finite complex; they must not be suppressed or treated as new external
assumptions. The larger ESM programme still requires mixed matter/chiral
quantum control, higher loops, infrared transport and invariant filtered
interacting shells. Nothing here proves asymptotic safety, a positive
nonperturbative gravity measure, finite-parameter ultraviolet gravity, or a
Yang--Mills mass gap.

\begin{center}\small
\begin{tabular}{>{\raggedright\arraybackslash}p{.22\textwidth}>{\raggedright\arraybackslash}p{.69\textwidth}}
\toprule
Variable & Scope of uniformity in V24\\\midrule
Volume & Constant-source zeros use paired odd-grid momenta. Algebraic
projection bounds are volume independent. Exact scale transgression is
thermodynamic; finite periods have the separate displayed error.\\
Mesh & The input and reference comparison use \eqref{eq:v24-domain}.
The local residual is $O((a\lambda)^3)$ at fixed $\lambda$, not an
all-mode smallness statement.\\
Arity and derivatives & Application at arity five, action/source input
through seven. The structural projection exists at each finite $N$;
its constants need not be uniform in $N$ or derivative order.\\
Loop and EFT & One loop and local momentum weight at most four. The
critical weight-four primitive remains open. An arity cutoff is not a
claim of finite-dimensional all-orders gravity couplings.\\
Couplings and lower scale & $\chi>0$, $\lambda>0$ fixed. No infrared
removal and no estimate uniform as $\chi\to0$.\\
Reference and measure & Same completed covariances, source, contour and
once-counted auxiliary density; only the recorded higher input jets and
order-$\hbar$ local normalization are added.\\
Iterations & Reference words have fixed lower-annulus support. Neither
analytic many-shell self-mapping nor higher-loop closure is concluded.\\
\bottomrule
\end{tabular}
\end{center}


\clearpage
\part*{V25: Critical local restoration by a full-germ argument}
\addcontentsline{toc}{part}{V25: Critical exactness and a bounded finite prescription}

\section{The critical question requires an extendable local cocycle}
\label{sec:v25-start}

Sections 1--192 are preserved. The last unresolved local equation in V24
was
\[
 d_5 H^{[4]}=x_{0,\lambda}^{[4]}.
\]
The scale derivative cannot solve this equation: at weight four it gives
\(d_5Y^{[4]}=0\), not a primitive of \(x^{[4]}\). This installment uses a
different input, stated explicitly: the established local BRST classification
of \emph{ordinary pure four-dimensional Einstein gravity}. It is not a
classification of the full Einstein--Standard Model or of the finite
filtered transformation. Its application requires showing that the actual
reference coefficient belongs to the ordinary local germ algebra, before
any truncation in the number of fields.

The resulting claim is stronger than another conditional splitting
criterion but narrower than a numerical counterterm evaluation. We prove
that the specified critical coefficient is in the image of the actual
finite local differential. We give a finite rational prescription for a
primitive, with an explicit matrix-height bound independent of the lattice
cutoff. The large native coefficient matrix and its native loop-integral
right-hand side are \emph{not numerically evaluated in this installment}.
No rank, small singular value, numerical anomaly coefficient or practical
conditioning estimate for that matrix is asserted.

The endpoint here is the one-loop, pure-gravity, local bank of V24 at total
arity five and weight four. The propagating reference, its infrared scale,
the coframe chart, the nonminimal sector and the auxiliary measure are
unchanged. We add no compensating field, no new high-frequency source
prescription, and no new finite input vertex. An all-arity \emph{zero-mesh
comparison germ} will be used to prove membership in a local differential
image. Its terms of arity beyond seven are not consumed by the arity-five
one-loop coefficient and are not an extension of the finite regulated
quantum action.

\subsection{Inputs and scope of the outside theorem}

The native ordinary action is the stationary Palatini action
\(\mathcal A_0\) in \eqref{eq:v17-generators}, with zero cosmological
term, zero independent Holst coefficient, and fixed \(\chi>0\). The source
is its ordinary coframe-coordinate Lie pullback. The lower reference is
\(\Theta_\lambda(p)=\theta(p/\lambda)\), where \(\theta\) is smooth,
real and even, equals one for \(|p|\le1\), and vanishes for \(|p|\ge2\).
All local coefficients use thermodynamic momentum normalization. We use
V24's weights, tensor symmetries, integrated functional convention and
ordinary differential \(d_N\).

The external classification used below is the local theorem of Barnich,
Brandt and Henneaux for pure Einstein gravity, including its antifields
\cite{V25BBHEinstein}. We use its local, contractible coframe-chart version;
global frame topology, anomalies of internal gauge factors and chiral
matter are not covered by this use. The corresponding Einstein--Yang--Mills
classification explicitly distinguishes additional abelian and matter
possibilities \cite{V20BBH1995}. V25 therefore cannot be transferred to the
full Standard Model merely by replacing the field list in its conclusion.

\section{An all-arity local lift of the actual critical reference}
\label{sec:v25-germ}

Let \(\mathscr A_{\rm jet}\) denote local densities with no explicit
coordinate dependence. They are analytic in the undifferentiated coframe
on a smaller contractible chart, polynomial in finitely many differentiated
fields, ghosts and antifields, and have the tensor and Grassmann symmetries
of the native ordinary theory. Local functionals are taken modulo total
ordinary derivatives. A formal germ version, expanded at the flat coframe,
is sufficient for every finite coefficient statement below.

Write \(s=(\mathbb S_\infty,\cdot)\) for the ordinary zero-mesh BV
differential. The harmonic gauge-fermion transformation is undone before
restriction to minimal variables. The off-shell nonminimal doublet is
restored first, as in V24; its removal after this canonical transformation
is a chain map. Let \(\mathscr X_{\theta,\lambda}\) be the local
weight-four component of the full ordinary reference trace. Its Taylor
coefficients at each finite arity are defined by the same ordered compiler
as \eqref{eq:v24-reference-compiler}, not by a new amplitude.

\begin{theorem}[Local, extendable critical reference germ]
\label{thm:v25-germ}
The reference has an all-arity local germ
\(\mathscr X_{\theta}^{[4]}\in\mathscr A_{\rm jet}\) of ghost number one
and weight four such that
\begin{equation}
 \boxed{s\mathscr X_{\theta}^{[4]}=0\quad\text{modulo total derivatives},
 \qquad
 \operatorname{pr}_{A\le N}\mathscr X_{\theta}^{[4]}
       =x_{0,\lambda,\theta}^{[4],\le N}.}
 \label{eq:v25-extendable}
\end{equation}
In particular the \(N=5\) coefficient is exactly V24's critical
right-hand side. The germ is independent of \(\lambda>0\), is analytic in
undifferentiated coframe variables on a sufficiently small fixed chart, and
has finite total derivative degree in each antifield sector. At ghost
number one and weight four that degree is \(5-n_*\), so \(n_*\le5\).
Its coefficients at each fixed arity are bounded uniformly in \(\lambda\)
for a fixed bounded class of profiles and a fixed \(\chi>0\).
\end{theorem}
\begin{proof}
The ordinary stationary functional and source are
\[
 -\tfrac12\langle b_0,C(q)^{-1}b_0\rangle,
 \qquad DG(q)^{-1}\mathcal L_cG(q).
\]
They are analytic local functions of the undifferentiated coframe, with
respectively two and one ordinary derivatives. The native pointwise inverse \(C(q)^{-1}\), not an inverse
wave operator, occurs in the first expression. These are the same
all-finite-degree functions already constructed in V17. Their classical
master equation is the ordinary one, before an arity quotient.

One can construct the required reference germ without assigning a value to
an unregulated ultraviolet determinant. After eliminating the algebraic
nonminimal Gaussian, put \(C_0=(1-\Theta_\lambda)\mathbb H_0^{-1}\) on
the propagating blocks and form
\[
 \operatorname{STr}\bigl[(I+C_0V_0)^{-1}\Theta_\lambda Dr\bigr].
\]
This notation denotes its ordered local symbol expansion. Freeze the
undifferentiated background \(\bar q\). On the reference annulus the
zero-derivative background perturbation obeys
\(\|C_0V_0(\bar q)\|\le C\|\bar q\|\); by decreasing the fixed
coframe chart it is less than \(1/2\). Its inverse is analytic there. The
nilpotent variables and positive derivative jets are then expanded only as
far as the requested total weight requires. For example, resumming the
undifferentiated background first changes a geometric series into its
analytic inverse; differentiating it a fixed number \(D\) of times is
bounded by a convergent sum \(\sum_{r\ge0}(1+r)^D2^{-r}\).

At weight four there are at most five ordinary derivatives in the
minimal reference density, and at most five antifields. There are thus
only finitely many differentiated or nilpotent jet factors at each
undifferentiated-background coefficient. Ordinary composition of symbols,
or equivalently the Taylor expansion in all routed external momenta,
produces a local polynomial in those jets. The equality of these two
constructions follows by expanding the frozen inverse back into its
convergent geometric series; at any finite arity only finitely many words
remain. Those words are precisely the inherited compiler.

Every nontrivial word contains both a covariance factor and the final
\(\Theta_\lambda\). At zero external momenta, their common support is
\(\lambda\le|p|\le2\lambda\); a fixed finite number of external
derivatives does not enlarge that closed support. The no-covariance trace
has only one derivative and no minimal antifield, so has no weight-four
component. The coefficient integrals in question are therefore compact
annular integrals with uniformly bounded inverses. They are analytic in
\(\bar q\) by dominated differentiation on that annulus.

The inherited ordinary regulated trace identity is applied before taking
weight or arity components. Its differentiated form gives
\(s\mathscr X=0\) with all antifield and nonminimal variations retained.
All operations on the displayed reference words are justified by the same
compact annular support; no cyclic rearrangement of an unregulated
ultraviolet trace is involved. Since \(s\) preserves weight, the
weight-four component is closed. Since it cannot lower arity, its
restriction is the stated finite cocycle. Finally the change of variables
\(p=\lambda z\), including the external Taylor factors, gives the
factor \(\lambda^{4-w}=1\). Annular inverse bounds and the finite
Taylor compiler bound each fixed coefficient. These arguments prove the
claims about locality, analyticity, consistency and scaling separately.
\end{proof}

\begin{remark}[Why extending the cocycle matters]
\label{rem:v25-false-truncation}
Vanishing of full local cohomology does not imply vanishing of cohomology
after arbitrary arity truncation. For example, take a degree-one element
\(b\) of arity five, a degree-two element \(z\) of arity six, and
\(db=z\). The full degree-one cohomology is zero, while the quotient at
arity five contains the nontrivial cycle \(b\). The theorem above rules
out this difficulty for the \emph{populated reference coefficient}; it does
not calculate the cohomology of every arity-five component array.
\end{remark}

\section{Critical exactness in the translation-invariant local algebra}
\label{sec:v25-cohomology}

We state the external input at the point where it is used. In the local
coframe chart, the ordinary pure-Einstein BRST classification removes the
antifields cohomologically and represents its positive-ghost-number
zero-form classes by
\[
 \widehat\Theta\,L(\mathcal T,\vartheta_1,\vartheta_2),
 \qquad \widehat\Theta=\widehat c^0\widehat c^1
                         \widehat c^2\widehat c^3.
\]
Here \(L\) is a Lorentz-invariant polynomial in curvature and covariant
curvature derivatives; each \(\vartheta_i\) is a degree-three Lorentz
ghost generator. In particular positive ghost numbers below seven can only
be four. These are classification results from
\cite[Theorems 1--3]{V25BBHEinstein}, not new V25 cohomology computations.
The statements hold in the Euclidean coefficient version as well as the
Lorentzian one. We do not infer an equality of their quantum measures.

There are two additional interfaces to check: the local metric/coframe
quotient, and the exclusion of coordinate-dependent counterterms.

\begin{lemma}[The native quotient and nonminimal pairs do not change local exactness]
\label{lem:v25-quotient}
On the contractible chart, exactness of an ordinary local integrated
coframe cocycle descends to the native metric-coordinate minimal complex.
The descent preserves analyticity, translation invariance and momentum
weight. It commutes with any fixed arity truncation after the resulting
local functional has been expressed in \(q,c,\eta,\sigma\).
\end{lemma}
\begin{proof}
Write a coframe locally as \(e=L\widehat e(q)\) and use six analytic
coordinates \(\ell\) on \(L\). Replace the six Lorentz ghosts by
\(v=s\ell\). The coefficient of the original Lorentz ghost in this
change is invertible on the small chart; it is the infinitesimal free
Lorentz action on the frame coordinate. This is a local invertible ghost
change, with the diffeomorphism ghost retained. The ordinary identity
\(s^2=0\) gives \(s\ell=v\), \(sv=0\). Its cotangent lift splits
the master action into the metric minimal action and the corresponding
frame doublets, up to the same local canonical change of coordinates.
Contracting a doublet and its antifield pair gives an acyclic factor;
explicitly the usual count-and-replace homotopy divides a monomial by
its positive number of doublet entries. Restriction to zero doublet
coordinates is a chain map. The off-shell nonminimal pair is removed by
the identical argument after undoing its local gauge fermion.

These coordinate changes use pointwise analytic inverses and ordinary
finite derivatives only. Frame coordinates and Lorentz ghosts have weight
zero; the canonical pairs have the same total weight two as the original
ones. There is no explicit coordinate dependence. The metric germ is
finally expanded at \(q=0\), where \(s\) cannot lower arity; the
restriction therefore commutes with the quotient. The proof is local in
the chart and makes no global gauge-slice claim.
\end{proof}

\begin{theorem}[No critical ghost-number-one class for an extendable pure-gravity germ]
\label{thm:v25-critical-exact}
Let \(a\) be a translation-invariant local pure-gravity density of ghost
number one and weight four on the native ordinary chart, analytic in the
undifferentiated coframe and polynomial in finite remaining jets. If
\(sa\) is a total derivative, then there are translation-invariant local
\(h\) and \(k^\mu\) with
\begin{equation}
 \boxed{a=sh+\partial_\mu k^\mu,
 \qquad \operatorname{gh}(h)=0,\quad w(h)=4.}
 \label{eq:v25-critical-exact}
\end{equation}
The statement is literature-assisted through the classification just
specified. It applies in particular to
\(\mathscr X_\theta^{[4]}\), and yields
\begin{equation}
 \boxed{x_{0,\lambda,\theta}^{[4],\le5}\in
 \operatorname{im}\left(d_5:
   \mathfrak L^0_{5,4}[w=4]\longrightarrow
   \mathfrak L^1_{5,4}[w=4]\right).}
 \label{eq:v25-range}
\end{equation}
No assertion that the entire truncated first cohomology vanishes is needed.
\end{theorem}
\begin{proof}
Use \cref{lem:v25-quotient} to work with the ordinary tetrad complex.
Denote the horizontal exterior derivative by \(d_H\), to distinguish it
from the finite differential \(d_5\). Give \(dx^\mu\) weight \(-1\).
Then both \(s\) and \(d_H\) preserve the total weight of a local form.
The four-form \(a\,d^4x\) has total weight zero.

The translation-invariant algebraic Poincare lemma in form degrees below
four has only constant horizontal forms as additional classes. Such forms
have ghost number zero. It therefore permits the descent of the given
positive-ghost-number consistency equation without introducing explicit
\(x\)-dependence. The successive bottom candidates have form/ghost
degrees \((0,5),(1,4),(2,3),(3,2),(4,1)\). By the imported zero-form
classification the candidates with ghost numbers five, three, two and one
are exact. The only possible exception while lifting the descent is a
constant one-form times
\(\widehat\Theta L(\mathcal T)\) at form/ghost degree \((1,4)\).

This exception is absent at the stated weight. The four ghosts in
\(\widehat\Theta\) have total weight \(-4\), and the one-form has
weight \(-1\). Total weight zero would force \(L\) to have derivative
weight five. A pure-gravity Lorentz scalar made of \(r\) curvature
tensors and \(m\) covariant derivatives has weight \(2r+m\) and
\(4r+m\) tensor indices. Complete contraction uses invariant metrics and,
on the oriented branch, invariant four-index alternating tensors. It
therefore requires an even number of indices. Weight five would instead
make \(m=5-2r\) odd and \(4r+m=5+2r\) odd. No such scalar exists.
This also excludes linear combinations and integrations by parts at that
homogeneous weight. The degree-three Lorentz ghost generators cannot enter
this ghost-four exception. The local chart excludes independent global
frame-topology terms.

All descent components can consequently be removed by translation-invariant
trivial terms. Lifting them back gives \(a=sh+\partial_\mu k^\mu\).
The construction takes place in the ordinary local jet algebra; the
imported classification, not a numerical matrix test, supplies its
cohomological step. Homogeneous weight extraction gives \(w(h)=4\),
since \(s\) preserves weight. Analytic germ expansion and the local
quotient transfer in \cref{lem:v25-quotient} give the result on \(q\).
Apply it to \cref{thm:v25-germ} and truncate at arity five. Both \(s\)
and ordinary total differentiation cannot lower arity, so higher terms of
\(h\) and \(k\) cannot contribute to that truncation. This proves
\eqref{eq:v25-range} for the actual coefficient.
\end{proof}

\begin{remark}[What has and has not been imported]
The absence of a four-dimensional pure gravitational anomaly is not claimed
as a new discovery. The new bridge proves that the actual compact-reference
coefficient belongs to the relevant all-arity local algebra and uses the
homogeneous weight to keep the primitive translation invariant. Neither the
non-nilpotent lattice source nor the Standard Model's abelian and chiral
sectors satisfy this pure-gravity input merely by notation.
\end{remark}

\section{A finite prescription with a cutoff-independent bound}
\label{sec:v25-finite}

Existence in a local cohomology group alone does not give a useful
refinement statement. The following construction separates that issue from
cohomological existence. Normalize \(\eta=\chi\widehat\eta\) and
\(\sigma=\chi\widehat\sigma\). The ordinary minimal differential then
has rational Taylor coefficients independent of \(\chi\); undoing this
rescaling restores the component powers of \(\chi\).

To avoid a hidden choice of representatives modulo total derivatives, let
\(E_0\) be the raw density space with arity at most five, ghost number
zero and weight four. Let \(J_1\) be the raw density space with ghost
number one, weight three and the same arity cap; its four copies represent
a current. Let \(E_1\) be the corresponding ghost-one, weight-four
density space. Use independent symmetric derivative multi-indices, the ten
symmetric metric components and the four ghost components, with the
corresponding antifield components and their actual Grassmann parities.
Define the rational matrix
\begin{equation}
 \mathsf A:E_0\oplus J_1^{\oplus4}\longrightarrow E_1,
 \qquad
 \mathsf A(h,k)=\operatorname{pr}_{A\le5}(sh)+\partial_\mu k^\mu.
 \label{eq:v25-matrix}
\end{equation}
Every column is produced by differentiating the native ordinary master
coefficients through arity six. These are already supplied by the V24
input budget. This is a finite local coefficient matrix, not a propagator.

For orientation, the \emph{raw} dimensions before integrations by parts are
\begin{equation}
 \dim E_0=27\,236\,680,\quad
 \dim J_1=26\,669\,300,\quad
 \dim E_1=130\,001\,064.
 \label{eq:v25-raw-counts}
\end{equation}
They are monomial counts, not ranks or numbers of physical operators. The
verifier computes them by a generating-function count of the independent
jet variables. They explain why building an uncompressed dense native
matrix is not a reasonable numerical strategy. No such matrix is built
or inverted in this work.

\begin{theorem}[Finite rational primitive and an arithmetic norm bound]
\label{thm:v25-finite}
Choose a raw density representative \(b_\theta\in E_1\) of
\(x_{0,\lambda,\theta}^{[4],\le5}\). Put
\(B=\mathsf A\mathsf A^T\) and let \(e_j\) be the coefficients of
\(\det(I+tB)=\sum_j e_jt^j\). Let \(r\) be the largest index with
\(e_r\ne0\). For \(r>0\) set
\begin{equation}
 \boxed{\mathsf A^+
 =\frac1{e_r}\mathsf A^T
       \sum_{j=0}^{r-1}(-1)^j e_{r-1-j}B^j.}
 \label{eq:v25-pseudoinverse}
\end{equation}
For \(r=0\) set \(\mathsf A^+=0\). Then
\begin{equation}
 \boxed{(h_\theta,k_\theta)=\mathsf A^+b_\theta,
       \qquad d_5[h_\theta]=x_{0,\lambda,\theta}^{[4],\le5}.}
 \label{eq:v25-primitive}
\end{equation}
This is a prescribed finite rational linear operation on the actual
annular coefficient array. It does not set that array to zero or replace it
by an invariantly guessed polynomial.

Let \(m,n\) be the row and column counts of \(\mathsf A\), let
\(Q\) be any positive common denominator of its entries, and let
\(H\ge1\) bound the absolute entries of the integer matrix \(Q\mathsf A\).
Then
\begin{equation}
 \boxed{\|\mathsf A^+\|_{2\to2}
 \le Q\,\max\{1,\sqrt{mn}H\}^{\min(m,n)-1}.}
 \label{eq:v25-height-bound}
\end{equation}
This bound has no lattice spacing, period, shell count or lower-scale
parameter. It is extremely conservative and is not a claim of practical
conditioning. Fixed finite coefficient-norm and windowed source-norm bounds
follow by the stated finite-dimensional conversions.
\end{theorem}
\begin{proof}
The matrix is rational because the coframe chart, its Taylor inverse, the
native ordinary action coefficients, the ghost recursion, the Leibniz rule
and integration by parts all have rational coefficients after the displayed
\(\chi\) normalization. \Cref{thm:v25-critical-exact} proves
\(b_\theta\in\operatorname{im}\mathsf A\), including any change of
raw representative by a total derivative.

The nonzero eigenvalues of \(B\) are positive. Their characteristic
polynomial is
\(z^r-e_1z^{r-1}+\cdots+(-1)^re_r\).
At a nonzero eigenvalue \(z\), this polynomial gives
\[
 z^{-1}=e_r^{-1}\sum_{j=0}^{r-1}(-1)^j e_{r-1-j}z^j.
\]
On \(\ker B\), multiplication by \(\mathsf A^T\) gives zero.
Thus \eqref{eq:v25-pseudoinverse} is the Moore--Penrose inverse, by the
spectral theorem. It is rational and \(\mathsf A\mathsf A^+\) is the
orthogonal projection onto \(\operatorname{im}\mathsf A\). This proves
\eqref{eq:v25-primitive}. The current part records an ordinary total
derivative; the functional counterterm is its \(h\) component.

For the bound write \(Z=Q\mathsf A\). A nonzero rank-\(r\) minor
of this integer matrix has absolute value at least one. By Cauchy--Binet,
the squared product of its nonzero singular values is the sum of squared
rank-\(r\) minors, hence is at least one. Every singular value is at most
\(\|Z\|_F\le\sqrt{mn}H\). Therefore
\(\sigma_{\min}^+(Z)\ge
\max(1,\sqrt{mn}H)^{-(r-1)}\).
Division by \(Q\) and \(r\le\min(m,n)\) prove
\eqref{eq:v25-height-bound}. For example
\(\|\mathsf A^+b\|_1\le\sqrt n\|\mathsf A^+\|_2\|b\|_1\).
Fixed powers of a source-window scale are assigned using the homogeneous
component weights, as in V24. All conversion constants depend on the finite
bank, not on refinement. The estimate supplies an explicit bound rather
than inferring analytic uniformity from the existence of a primitive.
\end{proof}

\begin{remark}[Status of the prescription]
Equations \eqref{eq:v25-matrix}--\eqref{eq:v25-primitive} determine a
mathematical counterterm, but its expanded native coefficient vector is not
reported. The checker verifies the inverse formula and its certificates on
finite test matrices and supplies a solver for a matrix explicitly provided
to it. It does not certify an unconstructed 130-million-row matrix. A
compressed tensor or jet implementation is still needed for a practical
coefficient evaluation. The range proof for the native right-hand side is
\cref{thm:v25-critical-exact}, not those small matrix tests.
\end{remark}

\section{Explicit transport between reference profiles at the critical weight}
\label{sec:v25-profile}

There is a useful part of the primitive which does not require any local
matrix solve. Take two profiles \(\theta_A,\theta_B\) with the common
plateaux stated above. They change only the lower reference, not the
ordinary native action, source or auxiliary density. Put
\(C_i=(1-\Theta_i)\mathbb H_0^{-1}\). Define
\begin{equation}
 \begin{split}
 T_{AB}^{[4]}=
 \operatorname{pr}_{w=4}\operatorname{pr}_{A\le5}
 \frac12\sum_{r=1}^5\frac{(-1)^{r+1}}r
 \operatorname{STr}\left[(C_BV_0)^r-(C_AV_0)^r\right],
 \end{split}
 \label{eq:v25-profile-primitive}
\end{equation}
where the local weight extraction is the common external Taylor extraction.
The difference is formed before integration.

\begin{theorem}[Critical profile transport with actual matrix order]
\label{thm:v25-profile}
The expression \eqref{eq:v25-profile-primitive} is a convergent finite
annular coefficient integral, and
\begin{equation}
 \boxed{d_5T_{AB}^{[4]}=x_B^{[4]}-x_A^{[4]},\qquad
 T_{AC}^{[4]}=T_{AB}^{[4]}+T_{BC}^{[4]}.}
 \label{eq:v25-profile}
\end{equation}
It is independent of the common lower scale \(\lambda\). For a bounded
class of profile derivatives through a fixed sufficient order \(M\),
\begin{equation}
 \|T_{AB}^{[4]}\|_{\rm coeff}
 \le C_{5,M,\chi}\|\theta_B-\theta_A\|_{C^M}.
 \label{eq:v25-profile-bound}
\end{equation}
Thus one may choose a single reference primitive \(h_A\) from
\cref{thm:v25-finite} and set \(h_B=h_A+T_{AB}^{[4]}\). Its profile
transport then composes exactly, rather than differing by unspecified local
terms. This is reference-profile matching, not matching to the discontinuous
unmodified ultraviolet regulator.
\end{theorem}
\begin{proof}
For noncommuting \(X_A,X_B\),
\[
 X_B^r-X_A^r=\sum_{j=0}^{r-1}X_B^j(X_B-X_A)X_A^{r-1-j}.
\]
Take \(X_i=C_iV_0\). Every difference word has one factor
\(C_B-C_A=-(\Theta_B-\Theta_A)\mathbb H_0^{-1}\), supported in the
common lower transition region. At a fixed external Taylor order all other
routed momenta coincide there up to the soft shifts used to define the
derivative. Hence the integrals are compact and the inverse margins bound
all differentiated words. Keeping their actual order is essential; no
commutation of a profile through a source vertex is used.

Subtract the two ordinary regulated trace identities before extracting
arity or weight. Their divergence terms, measure conventions and native
action are identical and cancel. The determinant difference is the finite
word expression above. Compact support makes its differentiation and
cyclicity legitimate, so \(d_5T_{AB}=x_B-x_A\). The composition identity
is the telescoping difference of the same determinant coefficients. Critical
homogeneity gives scale independence. The preceding telescoping expression,
the finite number of words and external derivatives, and rescaling of the
annulus give \eqref{eq:v25-profile-bound}. This estimate is for the
critical local coefficients, not for the unsubtracted determinant.
\end{proof}

\section{The common local update now includes the critical block}
\label{sec:v25-matching}

Let \(U_a=g_a-j+C_a^{\rm old}\) be exactly V24's complete local input.
Continue its already constructed subcritical update and add
\begin{equation}
 \boxed{\Delta C_{a,4}=-(U_a)_{w=4}+h_\theta^{[4]}.}
 \label{eq:v25-critical-counterterm}
\end{equation}
This is an additional common order-\(\hbar\) local action-and-source
normalization. It is not an extra copy of the old source counterterm and
not a modification of the tree-level propagators.

\begin{theorem}[Critical local restoration for the measured arity-five bank]
\label{thm:v25-matching}
With the finite prescription of \cref{thm:v25-finite}, the actual projected
finite variation obeys
\begin{equation}
 \boxed{\left(B_a^{\rm new}\right)_{w=4}
       =x_{0,\lambda,\infty,\theta}^{[4]}
                       -x_{a,\lambda,L,\theta}^{[4]}.}
 \label{eq:v25-matching}
\end{equation}
Consequently the full local pure-gravity bank \(A\le5,w\le4\) has one
common normalization with the vanishing residual proved in V24 for the
reference comparison. In the critical components the regular
no-covariance trace contributes no Taylor coefficient, so the separate
source-free omitted-zero-mode term is absent. At fixed profile bounds and
on \eqref{eq:v24-domain}, a critical slot with \(n_*\) antifields and
\(d=5-n_*\) external derivatives satisfies
\begin{equation}
 \boxed{\left|(B_a^{\rm new})_I^{[d]}\right|
 \le C_{I,b}\chi^{-n_*}
       \left[(a\lambda)^3+(\lambda L_{\min})^{-b}\right],\quad b>4.}
 \label{eq:v25-critical-rate}
\end{equation}
The common factor \(\hbar\) is restored in the action-level statement.
\end{theorem}
\begin{proof}
V24 gives \(B_a^{\rm old}=d_5U_a-x_a\) on the complete local bank.
The ordinary variation of \eqref{eq:v25-critical-counterterm} leaves
\(x_0^{[4]}-x_a^{[4]}\). One must also check that this is the
\emph{native projected} variation, not just its continuum substitute.

In a ghost-zero weight-four counterterm a component with \(n_*\)
antifields has exactly \(4-n_*\) derivatives. At arity at most five,
antifield-containing components have \(n_*=1\) or \(2\), hence at least
two derivatives. The native higher-action Euler discrepancy starts at
five derivatives and the free discrepancy at six. Their product with any
of these components has derivative order strictly above the allowed output
Taylor cap \(6-n_*\). Source-free components vary through the ordinary
soft source, since all external subset momenta are on the plateau. Thus
\(\mathsf P_{5,4}s_a\Delta C_{a,4}=d_5\Delta C_{a,4}\) exactly in
this local extraction. Unlike the subcritical proof this conclusion does
not require an extra zero-derivative hypothesis on the new primitive.

The critical right-hand side is the same reference coefficient already
estimated in \cref{prop:v24-reference-rate}. Its factor
\(\lambda^{5-n_*-d}\) equals one. The no-covariance source-free trace
has at most one derivative and no weight-four coefficient, as in
\cref{thm:v25-germ}; thus its retained-zero-mode contact is not present in
\eqref{eq:v25-critical-rate}. Covariance-bearing terms have the compact
lower-annulus support and the independent finite-volume comparison of V24.
This proves the estimate with its stated scope. The old Jacobian and
auxiliary measure are included once in \(U_a\) and are not reset by this
calculation.
\end{proof}

A polynomial local coefficient is a finite derivative contact. With a
fixed smooth source window, its Fourier kernel has the usual rooted
weighted bound obtained by differentiating the compact momentum window
more than the number of relative coordinates plus the required moment.
Applying that construction to \eqref{eq:v25-critical-rate} gives the same
cutoff/volume factor times the appropriate fixed source-window power. The
number of relative variables is at most sixteen at arity five. There is no
extensive volume multiplier. The primitive coefficient bound itself is
\eqref{eq:v25-height-bound} times the native annular coefficient bound;
it is not an invariant-neighbourhood bound for an interacting RG map.

\section{Prescribed lower-arity normalizations and the unresolved larger gates}
\label{sec:v25-relative}

The absolute critical equation and an equation with arbitrarily prescribed
finite constants are different problems. To make this distinction explicit,
let \(\mathsf L(h,k)=\ell\) encode any requested linear lower-arity or
physical normalization conditions; it acts only on the functional
counterterm, not on irrelevant choices of current representatives. Form
\begin{equation}
 \mathsf A_{\rm rel}=\begin{pmatrix}\mathsf A\\\mathsf L\end{pmatrix},
 \qquad b_{\rm rel}=\binom{b_\theta}{\ell}.
 \label{eq:v25-relative}
\end{equation}
Then a primitive respecting those conditions exists exactly when
\begin{equation}
 \boxed{(I-\mathsf A_{\rm rel}\mathsf A_{\rm rel}^+)b_{\rm rel}=0.}
 \label{eq:v25-relative-test}
\end{equation}
This is an ordinary finite linear-algebra equivalence. The same arithmetic
bound applies after denominators are cleared for the enlarged matrix.
Without an actual prescribed array \(\ell\), this note does not report a
numerical evaluation of the relative test.

The subcritical choices of V24 are preserved automatically because the
critical correction has a different weight. The earlier measured
identities and every retained source contact also remain; their local
normalizations are updated by the same functional, not separately by row.
A chosen critical normalization is certainly admissible if its difference
from \(h_\theta\) is the restriction of one ordinary \(s\)-closed
local functional. This includes physical finite additions presented with
all their symmetry partners. Merely solving a lower-arity truncated
equation does not prove extendability of an arbitrary list of constants.
We therefore do not silently identify the particular matrix prescription
with every previously conceivable finite normalization.

The new result discharges the absolute critical \emph{local existence}
question for the native pure-gravity reference and gives a cutoff-uniform
finite prescription for it. A practical compressed evaluation of the
counterterm, or of normalization-constrained coefficients, is still a
separate computation. More importantly, this local result does not prove
the missing nonlocal estimates for every new four- and five-metric panel,
the full mixed Einstein--SM source complex, the chiral measure interface,
higher-loop locality, infrared removal, or invariant filtered interacting
R1--R4 hypotheses. Absence of a pure gravitational local obstruction is
not a theorem about the combined Standard-Model anomaly and line data.

The next analytic target need not be another local rank hunt. It is the
full one-loop nonlocal source estimate on the already closed arity-five
bank, using this \emph{single} critical/subcritical normalization and the
same measured reference. Any later mixed matter theorem must explicitly
check its enlarged local cohomology and measure rather than inherit the
pure-gravity theorem by analogy.

\section{Verification and closing ledgers for V25}
\label{sec:v25-ledgers}

The standalone checker verifies the raw density counts, the homogeneous
weight obstruction in the translation-invariant descent, the finite
pseudoinverse formula and integer-minor bound on supplied finite examples,
a noncommuting reference-profile trace identity, the native-versus-ordinary
critical derivative caps, and the common incremental normalization rule.
It also accepts a separately provided finite rational coefficient matrix
and target and returns either a certified solution or its explicit
orthogonal range residual. This utility is not a populated native matrix
calculation. No new native loop integral is numerically evaluated. The
external local-cohomology theorem is identified as literature input and is
not claimed to have been formalized or re-proved by the checker.

\begingroup
\renewcommand{\refname}{Additional source for V25}
\begin{thebibliography}{99}
\bibitem[44]{V25BBHEinstein}
G.~Barnich, F.~Brandt and M.~Henneaux,
\emph{General solution of the Wess--Zumino consistency condition for
Einstein gravity}, Physical Review D \textbf{51} (1995), 1435--1439,
arXiv:hep-th/9409104, doi:10.1103/PhysRevD.51.1435.
Theorems 1--3 supply the local pure-gravity BRST classification used in
\cref{sec:v25-cohomology}. This source does not supply the native
reference integral, its lattice matching bound, or the finite numerical
counterterm coefficients.
\end{thebibliography}
\endgroup

\subsection*{Proved}
\addcontentsline{toc}{subsection}{V25: Proved}
The actual critical reference array admits the all-arity local germ of
\cref{thm:v25-germ}. Applying the explicitly stated external classification
with the translation-invariant weight argument proves that its arity-five
array lies in the image of the native ordinary local differential. The
finite rational prescription \eqref{eq:v25-primitive} and the independent
arithmetic norm bound \eqref{eq:v25-height-bound} then give a local
critical primitive at this fixed order. Reference-profile differences have
the explicit ordered primitive \eqref{eq:v25-profile-primitive}. The
common measured critical update has residual
\eqref{eq:v25-critical-rate}. These are existence and coefficient-map
theorems, not a reported large native matrix evaluation.

\subsection*{Inherited}
\addcontentsline{toc}{subsection}{V25: Inherited}
V1--V24 and their labels are unchanged. Native ordinary stationary action
and chart-source germs, the finite compiler, the off-shell trace identity,
the degree-closed ordinary differential, the subcritical primitive, the
reference cutoff/volume estimate and the once-counted measure convention
are used at their established scopes. No new proof of the finite
Yang--Mills or Standard-Model perturbative programme is asserted.

\subsection*{Literature input}
\addcontentsline{toc}{subsection}{V25: Literature input}
The local pure-Einstein BRST classification, including antifield
triviality and its zero-form classes, is imported from
\cite{V25BBHEinstein}. The modified-reference framework remains
\cite{V15IIM}. Its use is confined to a contractible ordinary pure-gravity
chart. It is not a theorem about the finite filtered source, global
measure-line holonomy or the Standard Model's chiral/internal sectors.

\subsection*{Conditional}
\addcontentsline{toc}{subsection}{V25: Conditional}
Specified additional finite normalization constraints require the actual
relative test \eqref{eq:v25-relative-test}, unless supplied as a common
closed functional. Efficient evaluation requires a compressed
representation of the native local matrix and actual annular coefficient
array. The all-EFT quantum endpoint requires the unclosed mixed, nonlocal,
infrared, higher-loop and interacting-shell estimates listed above.

\subsection*{Open}
\addcontentsline{toc}{subsection}{V25: Open}
Evaluate selected native critical matching coefficients in a tractable
symmetry/jet basis and establish the full nonlocal one-loop source bound
for the completed arity-five bank with its single local normalization.
The large-matrix rank and coefficients are not reported here. The next
mixed Einstein--SM step must separately address its combined local
cohomology and chiral measure. No nonperturbative positivity, asymptotic
safety, finite-parameter gravity completion, or mass-gap conclusion follows.

\begin{center}\small
\begin{tabular}{>{\raggedright\arraybackslash}p{.22\textwidth}>{\raggedright\arraybackslash}p{.69\textwidth}}
\toprule
Variable & Uniformity and limitation in V25\\\midrule
Volume & Local primitive coefficients use thermodynamic normalization.
The finite-period residual is separately bounded; no total-volume factor
enters a rooted local coefficient.\\
Mesh and scale & The primitive matrix has no mesh dependence. Critical
reference coefficients are scale independent at fixed profile. The finite
comparison uses V24's domain and fixed positive lower scale.\\
Field/source degree & Application at total arity five and weight four.
The raw matrix is finite but very large; its rank is not computed. Bounds
are not uniform in unbounded arity or derivative order.\\
Loop and couplings & One loop, pure gravity, fixed \(\chi>0\). The
normalization rescaling is undone componentwise; no bound uniform at
\(\chi=0\).\\
Profiles & Fixed bounded smooth even profile class with common plateaux.
Transport constants depend on finitely many profile seminorms.\\
Measure and topology & Same finite completion and once-counted auxiliary
measure. Local contractible coframe chart; no global quantum-measure claim.\\
Iteration & Local symmetry restoration does not prove analytic multiscale
self-mapping or arbitrary-loop EFT closure.\\
\bottomrule
\end{tabular}
\end{center}

\clearpage
\part*{V26: Nonlocal closure of the one-loop arity-five gravitational source module}
\addcontentsline{toc}{part}{V26: Nonlocal source closure with one common measured normalization}

\section{From local solvability to a cutoff-independent source module}
\label{sec:v26-start}

Sections 1--200 are preserved. V25 establishes solvability and a bounded
coefficient prescription for the critical \emph{local} reference equation;
it explicitly leaves the enlarged nonlocal source estimate open. This
installment supplies that estimate for the completed pure-gravity regulator
already specified in V14, V17 and V24. It does not repeat the critical
cohomology calculation or evaluate its very large matrix. The application
of that calculation below retains its external pure-gravity classification
and translation-invariant-germ hypotheses.

The new analytic content is the simultaneous treatment of all ordinary
metric and minimal-antifield backgrounds of total arity at most five. The
previous zero-background antifield reduction extends to this bank only if
both background propagators, the ghost-action return, and the nonlinear
coordinate contacts are differentiated. A dressed quadratic formula does
that without enumerating a new determinant separately for each panel. Its
symbol estimates prove a complete Brillouin-zone comparison. The resulting
functional uses the \emph{single} local normalization of V24--V25; it is not
a collection of separately imposed Slavnov rows.

The scope remains one loop, pure gravity, a fixed positive lower Wilsonian
scale, and the declared holomorphic metric reference. The six Lorentz
variables and auxiliary connection retain their earlier measured treatment.
The four diffeomorphism-ghost pairs are integrated. This is not a theorem
about the discontinuous uncompleted regulator, a global Lorentzian measure,
an invariant interacting analytic ball, or the full Einstein--Standard
Model field list. The general relation of regulated determinants to modified
Slavnov identities is inherited from \cite{V15IIM}. No first existence claim
for perturbative gravity is made.

\subsection{Precise finite panel and normalization convention}

Let \(\eta=q^*\), \(\sigma=c^*\), and use the actual V18 source
\eqref{eq:v18-repaired-source}. Ordinary nonminimal backgrounds are set to
zero \emph{after} the fixed off-shell restoration and harmonic canonical
shift used in V16 and V24. The minimal ghost-number-zero panels are
\begin{equation}
 I=(b,r,s,t),\qquad t=r+2s,\qquad
 A(I)=b+2r+3s\le5,\qquad n(I)=r+s.
 \label{eq:v26-panel}
\end{equation}
Here \(b\) counts metrics, \(r\) metric antifields, \(s\) ghost
antifields, and \(t\) ghosts. Vacuum constants are normalized away, but no
field-independent \emph{source} contact is removed. Thus \(n=0,1,2\).
The entire structural bank is
\begin{center}\small
\begin{tabular}{c|c|c|c}
\toprule
Slot & allowed metric degree & local Taylor degree & first remainder degree\\\midrule
\(q^b\)&\(1\le b\le5\)&4&5\\
\(\eta q^bc\)&\(0\le b\le3\)&3&4\\
\(\sigma q^bcc\)&\(0\le b\le2\)&3&4\\
\(\eta^2q^bcc\)&\(0\le b\le1\)&2&3\\
\(\eta\sigma ccc\)&0&2&3\\\bottomrule
\end{tabular}
\end{center}
These are structural families, not a count of their tensor components or
independent physical couplings. Antifield and ghost exchange signs are
retained. The input budget through arity seven is precisely
\eqref{eq:v24-higher-input}; no further finite action vertex is added here.

Use odd periodic product grids with physical periods \(L_\nu\) and common
mesh \(a\). Set \(\lambda=\Lambda_0\), and retain the convenient domain
\begin{equation}
 \boxed{\lambda\ge4096\mu>0,\qquad a\lambda\le2^{-18},\qquad
                   \lambda L_\nu\ge1.}
 \label{eq:v26-domain}
\end{equation}
Every external momentum has size at most \(\mu\); the dependent last
momentum is also kept on that window. In estimates on independent momentum
cubes one may enlarge these windows by a fixed factor at most five. The
strong separation in \eqref{eq:v26-domain} allows that enlargement. Write
\(K\) for all independent external momenta, \(\mathfrak n_I\) for the
product of polarization norms, and \(V_L=\prod_\nu L_\nu\). The normalized
loop sum is \(\int_p^L=V_L^{-1}\sum_p\). Source coefficients below have
\(\hbar\) removed. Constants depend on the stated finite panel, derivative
order, fixed \(\chi>0\), completion parameters and finitely many smooth
profile seminorms, not on mesh, periods or the number of integrated shells.

\section{Background-complete Berezin reduction}
\label{sec:v26-reduction}

Let \(C_g,C_c\) be the full completed covariances
\eqref{eq:v14-full-cov}, with the retained low modes excluded. Let
\(V_g(q)\) and \(V_c(q)\) be the ordinary-background Hessian perturbations
of the \emph{actual} extended metric and ghost actions. Define
\begin{equation}
 \mathcal G_g(q)=(I+C_gV_g(q))^{-1}C_g,\qquad
 \mathcal G_c(q)=(I+C_cV_c(q))^{-1}C_c.
 \label{eq:v26-bg-cov}
\end{equation}
These expressions mean finite formal background jets of the resolvents.
They do not invert a zero covariance on retained modes. They also agree
with the analytic hard-quadratic Gaussian germ on its sufficiently small
source chart. The latter assertion does not claim convergence of the full
nonquadratic gravitational integral.

For hard variables \((x,\beta,\bar\beta)\), define the primitive quadratic
operands by their coefficients in the existing action and source:
\begin{align}
 \tfrac12\langle x,T_{\eta c}(q)x\rangle
    &=[x^2]\langle\eta,\breve r_a^q(q+x,c)\rangle,\notag\\
 L_\eta(q;x)\beta
    &=[x\beta]\langle\eta,\breve r_a^q(q+x,\beta)\rangle,\notag\\
 \langle\bar\beta,B(q,c)x\rangle
    &=[\bar\beta x]\widetilde S^{\rm gh}_a(q+x,c+\beta,\bar\beta),\notag\\
 Q_\sigma(\beta)&=\langle\sigma,(P\beta)\cdot D(P\beta)\rangle,
       \qquad P=S^2.
 \label{eq:v26-operands}
\end{align}
Brackets denote the displayed hard degree, not a projection on external
source degree. In particular \(B\) comes from the V14/V17/V24 ghost
\emph{action}, which is not replaced by the V18 transformation source.
With the inherited odd-source ordering put
\begin{align}
 \tfrac12\langle x,D_\eta(q)x\rangle
   &=\tfrac12\langle x,T_{\eta c}(q)x\rangle
        -L_\eta(q;x)\mathcal G_c(q)B(q,c)x,\notag\\
 \tfrac12\langle x,D_\sigma(q)x\rangle
   &=Q_\sigma\bigl(\mathcal G_c(q)B(q,c)x\bigr).
 \label{eq:v26-dressed}
\end{align}
All coefficients of \(D=D_\eta+D_\sigma\) are even in total Grassmann
parity; their external odd parameters still anticommute in the stipulated
order. There are two ghost resolvents inside \(D_\sigma\).

\begin{proposition}[One measured expression on the full minimal panel]
\label{prop:v26-factorization}
Through the external arity in \eqref{eq:v26-panel}, the inherited finite
one-loop compiler is exactly
\begin{equation}
\begin{split}
 \Gamma_{a,1}={}&\mathcal M_a(q)
 +\tfrac12\Tr\log(I+C_gV_g(q))
 -\Tr\log(I+C_cV_c(q))\\
 &+\tfrac12\Tr\log\bigl(I+\mathcal G_g(q)D(q)\bigr).
\end{split}
 \label{eq:v26-factorization}
\end{equation}
The auxiliary coefficient \(\mathcal M_a\), including the coordinate
Jacobian, occurs once. Every background derivative of every resolvent and
of every source and ghost-action operand in this formula is retained.
\end{proposition}
\begin{proof}
At a fixed coefficient of external arity at most five, a hard-linear
vertex has a momentum equal to a sum of at most five external momenta.
It lies strictly in the zero region of the hard covariance by
\eqref{eq:v26-domain}. The algebraic nonminimal Gaussian is first
eliminated with its exact contact, as in V16. In a connected one-loop graph
of the remaining fields, the number of internal lines equals the number
of vertices. Every surviving vertex has at least two hard legs, so each
has exactly two. This is the same finite input-budget mechanism as
\cref{prop:v24-input-budget}, now used with the whole background bank.

The quadratic ghost dependence is
\(\bar\beta M(q)\beta+\bar\beta B(q,c)x\), together with the source
terms \(L_\eta\beta+Q_\sigma(\beta)\). Berezin integration over
\(\bar\beta\) makes the exact substitution
\(\beta=-\mathcal G_c(q)B(q,c)x\) and contributes its determinant.
It leaves precisely \eqref{eq:v26-dressed}; the minus sign in the first
line and the two ghost factors in the second line follow before the
metric integration. The latter integration supplies its determinant and
then the last trace in \eqref{eq:v26-factorization}. Square-root and
contour conventions are inherited, with the ratio normalized at zero
sources. Each remaining trace is a finite formal series at the prescribed
arity. A further contraction of \(\mathcal M_a\) would add a loop;
a hard-linear return containing it has already been excluded. Thus it
is an insertion at the retained background, not a second loop determinant.
\end{proof}

This proposition extends the zero-background reduction of V18; it is not
a new general Gaussian-elimination theorem. To make its background content
unambiguous, for a nonempty labelled set \(J\) of metric marks,
\begin{equation}
 D_J\mathcal G_g(0)=
 \sum_{\pi=(J_1,\ldots,J_\ell)}(-1)^\ell
 C_g(V_g)_{J_1}C_g\cdots(V_g)_{J_\ell}C_g,
 \label{eq:v26-inverse-jets}
\end{equation}
where the sum is over ordered set partitions of \(J\), and the analogous
formula holds for ghosts. No factorial is present for distinct labelled
derivatives. For example the \(\eta^2qcc\) coefficient includes all four
terms obtained by differentiating
\(-\Tr(\mathcal G_gD_\eta\mathcal G_gD_\eta)/4\).
In \(D'_\eta\) one must also differentiate
\(L_\eta\mathcal G_cB\) in all three positions. These formulas specify
the formerly missing background contacts without leaving a propagator
frozen or introducing a new fitted vertex.

\section{Whole-zone symbol bounds with ordinary backgrounds retained}
\label{sec:v26-symbols}

Put \(R_p=\lambda+|p|\), using periodic representatives on the lattice.
A soft shift by any external subset changes \(R_p\) by bounded factors.
An interior plateau can be fixed, for example, by \(|ap|\le1/1024\).
All routed labels and differentiated subset labels then lie on the native
plateau after a harmless smaller support margin is inserted.

\begin{lemma}[Marked resolvents and dressed vertices]
\label{lem:v26-symbols}
At every background degree used in \eqref{eq:v26-panel}, and for every
fixed pair of hard/external multi-indices \((\beta,\alpha)\),
\begin{align}
 \|\partial_p^\beta\partial_K^\alpha(\mathcal G_g)_J\|
   &\le C\chi^{-1}R_p^{-2-|\alpha|-|\beta|},\notag\\
 \|\partial_p^\beta\partial_K^\alpha(\mathcal G_c)_J\|
   &\le CR_p^{-2-|\alpha|-|\beta|},\notag\\
 \|\partial_p^\beta\partial_K^\alpha(D_\eta)_J\|
 +\|\partial_p^\beta\partial_K^\alpha(D_\sigma)_J\|
   &\le CR_p^{1-|\alpha|-|\beta|}\mathfrak n_D.
 \label{eq:v26-dressed-bounds}
\end{align}
For nonempty \(J\), the interior comparison satisfies
\begin{align}
 \|\partial^{\alpha,\beta}[(\mathcal G_{g,a})_J-(\mathcal G_{g,0})_J]\|
   &\le C\chi^{-1}a^3R_p^{1-|\alpha|-|\beta|},\notag\\
 \|\partial^{\alpha,\beta}[(\mathcal G_{c,a})_J-(\mathcal G_{c,0})_J]\|
   &\le Ca^4R_p^{2-|\alpha|-|\beta|},\notag\\
 \|\partial^{\alpha,\beta}[(D_a)_J-(D_0)_J]\|
   &\le Ca^4R_p^{5-|\alpha|-|\beta|}\mathfrak n_D.
 \label{eq:v26-dressed-comparison}
\end{align}
For empty \(J\) the free metric comparison is the stronger
\(C\chi^{-1}a^4R_p^{2-|\alpha|-|\beta|}\). The dressed kernels are
smooth and periodic in the hard loop momentum. Polarization factors for
ordinary metric marks are understood on the right.
\end{lemma}
\begin{proof}
V14 supplies the actual completed free inverse bounds, and multiplication
by the smooth lower cutoff makes them regular near the retained origin.
The native degree-six and -seven action inputs in V24 have Hessian order
two: their vertices are \(O(\chi R_p^2)\), and the ghost-action vertices
are \(O(R_p^2)\). For the metric comparison V17's evaluated mesh generator
and V24's extension give \(O(\chi a^3R_p^5)\). Ghost interactions are
exactly ordinary on the plateau; their only comparison is the free inverse,
which is \(O(a^4R_p^2)\). Formula \eqref{eq:v26-inverse-jets} therefore
proves the first two bounds and their differences. In particular a metric
vertex difference between two inverses has order
\(R_p^{-2}a^3R_p^5R_p^{-2}=a^3R_p\). This explains why the stronger
zero-background \(a^4\) estimate must not be used for this whole bank.

Every differentiated transformation vertex has one derivative. The
pointwise inverse \((I+\mathsf B_q)^{-1}\) contributes no momentum
inverse or extra derivative. Its differentiated quadratic metric product
must be examined before assigning hard degree. If both hard metric legs
are in \(\mathsf Q\), their total momentum, including any metric marks
at that location, is soft. The derivative of that product is therefore a
soft derivative. If only one hard leg enters the differentiated product,
its momentum is cut off by the outer \(P\) before reaching a principal
face. If the differentiated leg is a hard ghost, it has the inner factor
\(P\). Undifferentiated hard legs in \((I+\mathsf B_q)^{-1}\) cause no
symbol discontinuity. This exhausts the two-hard-leg possibilities at
every retained background degree. It proves periodic smoothness without
filtering the individual high--high legs in a contact.

Thus \(T,L,Q_\sigma\) have order at most one and the actual return
\(B\) has order at most two. The products \(L\mathcal G_cB\) and
\(Q_\sigma(\mathcal G_cB x)\) have order one, because
\(\mathcal G_cB\) has order zero. Derivatives of ordinary backgrounds
insert only the order-zero chains in \eqref{eq:v26-inverse-jets} or
further finite coefficients of the same source polynomials. On the
plateau the source and ghost-action vertices coincide with their
ordinary counterparts, so each difference in \(D\) replaces a ghost
inverse and gives \(a^4R_p^5\). In the remaining zone, put \(z=ap\):
propagators scale as \(a^2\), action Hessian vertices as \(a^{-2}\),
and source vertices as \(a^{-1}\). All profiles are smooth periodic
there. Each momentum derivative contributes \(a\), giving exactly the
stated global symbol bounds. All expressions use finitely many terms at
this arity, so differentiated product constants are uniform as stated.
\end{proof}

Expand \eqref{eq:v26-factorization} in labelled sources, retaining cyclic
order and all Grassmann signs. Let \(f_{a,I}(p,K)\) be the resulting sum
of propagating integrands. The purely auxiliary coefficient is treated
separately by \cref{thm:v10-local-continuum}.

\begin{corollary}[Uniform superficial degree on the enlarged bank]
\label{cor:v26-loop-bounds}
For each panel in \eqref{eq:v26-panel},
\begin{align}
 |\partial_p^\beta\partial_K^\alpha f_{a,I}|
 +|\partial_p^\beta\partial_K^\alpha f_{0,I}|
 &\le C_{I,\alpha,\beta}\chi^{-n}\mathfrak n_I
                     R_p^{-n-|\alpha|-|\beta|},\notag\\
 |\partial_p^\beta\partial_K^\alpha(f_{a,I}-f_{0,I})|
 &\le C_{I,\alpha,\beta}\chi^{-n}\mathfrak n_I
                     a^3R_p^{3-n-|\alpha|-|\beta|}.
 \label{eq:v26-loop-bounds}
\end{align}
The first line is whole-zone; the second is on the common interior
plateau. The ultraviolet degree is \(4-n\), independently of the number
of retained metric marks. No inversion-parity improvement is assumed.
\end{corollary}
\begin{proof}
A metric or ghost background Hessian insertion between the appropriate
propagators adds order zero. Each \(\mathcal G_gD\) adds order minus
one and one factor \(\chi^{-1}\). Differentiating backgrounds leaves
these counts unchanged. The first two logarithms of
\eqref{eq:v26-factorization} have \(n=0\). Telescoping each finite
product between its lattice and ordinary values adds at worst a relative
factor \(a^3R_p^3\); the \(a^4\) differences are smaller on the
plateau. Differentiated product rules prove every stated seminorm.
\end{proof}

\section{The full arity-five cutoff theorem}
\label{sec:v26-cutoff}

Let \(\gamma_{a,L,I}\) denote the complete coefficient per volume,
including \(\mathcal M_a\) in source-free slots. Set
\begin{equation}
 d_I=4-n,\qquad m_I=5-n,\qquad
 \gamma^R_{a,L,I}=(I-\mathsf T_{d_I})\gamma_{a,L,I}.
 \label{eq:v26-subtraction}
\end{equation}
This is precisely the ghost-zero component of the common V24 projection.
The continuum coefficient \(\gamma^R_{0,L,I}\) is defined by subtraction
\emph{inside} its loop sum, not by the difference of divergent integrals.
The auxiliary nonlocal continuum coefficient is zero, with its local
matching data retained. For a single metric source there is no independent
external momentum; its full coefficient is local, so its remainder is zero.

\begin{theorem}[Nonlocal convergence on every retained minimal source slot]
\label{thm:v26-cutoff}
On \eqref{eq:v26-domain}, for \(j=|\alpha|\le m_I\),
\begin{equation}
 \boxed{|\partial_K^\alpha(\gamma^R_{a,L,I}-\gamma^R_{0,L,I})(K)|
 \le C_{I,\alpha}\chi^{-n}\mathfrak n_I
                         a|K|^{m_I-j}.}
 \label{eq:v26-main}
\end{equation}
The same estimate holds in the thermodynamic convention. For every fixed
integer \(r\), all scaled seminorms with \(|\alpha|\le r\) satisfy
\begin{equation}
 \sum_{|\alpha|\le r}\frac{\mu^{|\alpha|}}{\alpha!}
  \sup_{|K|\le\mu}|\partial_K^\alpha
                 (\gamma^R_{a,L,I}-\gamma^R_{0,L,I})|
 \le C_{I,r}\chi^{-n}\mathfrak n_I a\mu^{m_I}.
 \label{eq:v26-all-derivatives}
\end{equation}
This theorem covers, in particular, the four- and five-metric panels,
\(\eta q^2c,\eta q^3c,\sigma qcc,\sigma q^2cc,\eta^2qcc\), and
\(\eta\sigma ccc\) with one simultaneous local assignment.
\end{theorem}
\begin{proof}
For the propagating terms, Taylor's integral remainder to order
\(d_I=m_I-1\), applied also after up to \(m_I\) external derivatives,
uses the \(m_I\)-th derivative of the integrand. By
\eqref{eq:v26-loop-bounds}, its continuum majorant is
\(C\chi^{-n}\mathfrak n_I R_p^{-5}\), while the interior mismatch
majorant is \(C\chi^{-n}\mathfrak n_I a^3R_p^{-2}\).
The factor outside each majorant is \(|K|^{m_I-j}\).

Here are the elementary sum bounds needed uniformly in volume. With
\(\lambda L_\nu\ge1\), comparison with cells of the momentum grid,
or dyadic annular counting, gives
\begin{equation}
 \int_{|p|\le c/a}^{L}R_p^{-2}\le Ca^{-2},\qquad
 \int_{|p|\ge c/a}^{L}R_p^{-5}\le Ca,
 \label{eq:v26-sums}
\end{equation}
where the second sum is on the continuum torus momentum lattice. Each
annulus of radius \(R\ge\lambda\) has normalized mode count at most
\(CR^4\); summing \(R^{4-2}\) up to \(a^{-1}\) and
\(R^{4-5}\) above it proves \eqref{eq:v26-sums}. No total-site count
is used as a bound on a local coefficient.

Choose a smooth partition in \(ap\) supported strictly within the
common plateau. In its inner part the mismatch integral is bounded by
\(a^3\cdot a^{-2}=a\). In the complementary lattice region, all
routed lines have momentum of order at least \(a^{-1}\). The first
line of \eqref{eq:v26-loop-bounds} gives the same \(a\) bound after
\(m_I\) derivatives, including face crossings because the completed
symbols are periodic and smooth. The continuum complementary tail has
that bound by \eqref{eq:v26-sums}. This comparison includes all native
and upper vertices; it does not discard the upper region. Taylor
subtraction commutes with the finite lattice sum and with the now
absolutely convergent continuum integral. It proves
\eqref{eq:v26-main} for the propagating part.

When \(j>m_I\), the Taylor polynomial is annihilated. Integration of
the differentiated mismatch gives respectively
\(O(a^2)\), \(O(a^3[1+\log(1/(a\lambda))])\), or
\(O(a^3\lambda^{7-n-j})\), according as
\(j=m_I+1\), \(j=m_I+2\), or \(j>m_I+2\).
The upper and continuum tails are \(O(a^{n+j-4})\). Multiplication
by \(\mu^j\), and \(a\mu\le a\lambda\ll1\), bound all these
terms by \(Ca\mu^{m_I}\). The first part of the theorem handles
lower derivatives and proves \eqref{eq:v26-all-derivatives}.

For \(n=0\), the auxiliary amplitude obeys the already proved fixed
source-degree scaling \(a^{-4}F_I(aK)\). Its subtracted remainder and
all fixed scaled derivatives are bounded by \(Ca|K|^5\) and
\(Ca\mu^5\), respectively, by \cref{thm:v10-local-continuum}.
The pointwise coordinate Jacobian is local and has zero complementary
remainder. The nonwrapping conditions of that theorem follow from
\eqref{eq:v26-domain}. For \(n>0\) an auxiliary insertion inside
a loop would be of the next loop order. This completes the measured
coefficient comparison.
\end{proof}

The local jets are not estimated by \eqref{eq:v26-main}. They have the
usual possible sizes \(a^{n+d-4}\) for \(d<4-n\), and
\(1+\log(1/(a\lambda))\) at \(d=4-n\), with normalization factors
understood. The known derivative-free source cancellations sharpen some
of these upper bounds, but are not needed for this nonlocal theorem.
Finite critical matching coefficients and divergent lower coefficients
remain in the single local state. No term is removed on grounds that it
has only source legs.

\section{Rooted source norms, finite-volume comparison and shell allocation}
\label{sec:v26-kernels}

Multiply each remainder by fixed smooth windows on all its external legs.
For \(A(I)\ge2\), let \(\mathcal K_{a,L,I}\) be the inverse Fourier
kernel of the lattice--continuum difference, in \(A(I)-1\) relative
coordinates. The physical normalization gives the following estimate.

\begin{theorem}[Source topology and independent volume limit]
\label{thm:v26-kernels}
For every fixed kernel moment \(b\ge0\),
\begin{equation}
\begin{split}
 &a^{4(A(I)-1)}\sum_{\boldsymbol x}
 \left[1+\mu\sum_{j=1}^{A(I)-1}d_{\rm per}(x_j,0)\right]^b
                       \|\mathcal K_{a,L,I}(\boldsymbol x)\|\\
 &\hspace{22mm}\le C_{I,b}\chi^{-n}a\mu^{5-n}.
\end{split}
 \label{eq:v26-kernel}
\end{equation}
Polarization norms are incorporated in the operator norm of the kernel.
Consequently the corresponding functional has a bound with one argument
in physical \(L^1\) and all others in \(L^\infty\), with the same
constant and no extensive volume multiplier. Separately, for \(b>4\)
and \(j\le5-n\),
\begin{equation}
\begin{split}
 |\partial_K^\alpha(\gamma^R_{0,L,I}-\gamma^R_{0,\infty,I})(K)|
 \le C_{I,\alpha,b}\chi^{-n}\mathfrak n_I
 |K|^{5-n-j}\lambda^{-1}(\lambda L_{\min})^{-b}.
\end{split}
 \label{eq:v26-volume}
\end{equation}
The analogous fixed scaled derivative and windowed kernel bounds hold.
\end{theorem}
\begin{proof}
Rescale all independent momenta by \(\mu\). Use
\eqref{eq:v26-all-derivatives} with more than
\(4(A(I)-1)+b+1\) derivatives and integrate by parts in the inverse
Fourier transform. Its absolute moment is bounded by
\(C\chi^{-n}a\mu^{5-n}\). Periodize this kernel; on each lift
\(d_{\rm per}\le |x|\), so the triangle inequality preserves the
bound. This proves \eqref{eq:v26-kernel} also when a source-window
length exceeds a period. Pairing with one root argument and bounding the
others proves the functional statement.

For \eqref{eq:v26-volume}, the subtracted continuum integrand has
\(b\) hard-momentum derivatives in \(L^1\) with norm at most
\(C\chi^{-n}\mathfrak n_I|K|^{5-n-j}\lambda^{-1-b}\), by
\eqref{eq:v26-loop-bounds}. Poisson summation followed by integration
by parts bounds a nonzero dual mode by this norm times its physical
length to power \(-b\). Summing the dual modes requires \(b>4\)
and proves the stated estimate. The low covariance cutoff removes any
singularity at zero. There is no omitted-zero-mode contact in these
subtracted propagating amplitudes. Finite-torus off-grid coefficients
and thermodynamic coefficients remain distinct normalization conventions;
this estimate, not an asserted torus Ward identity, relates them.
\end{proof}

Let \(C_{\bullet,\le j}\) be common smooth cumulative covariance
splits of the \emph{primitive} completed metric and ghost covariances,
with \(\Lambda_j=2^j\lambda\). Expand every dressed inverse before
assigning a shell. In particular both ghost lines in \(D_\sigma\),
the line inside \(D_\eta\), and every background-resolvent return are
included. A word belongs to its largest internal shell label, not to a
selected metric line alone.

\begin{proposition}[A summable complement for the entire module]
\label{prop:v26-shells}
For a smooth annular increment at scale \(\Lambda_j\ge\lambda\),
\begin{equation}
 \|(I-\mathsf T_{4-n})\Delta_j\Gamma_{a,1,I}\|_{r,\mu}
 \le C_{I,r}\chi^{-n}\mathfrak n_I
                         \mu^{5-n}\Lambda_j^{-1}.
 \label{eq:v26-shells}
\end{equation}
The same holds for the ordinary coefficient and its rooted kernel. On
interior increments the comparison is bounded by
\(C\chi^{-n}\mathfrak n_Ia^3\mu^{5-n}\Lambda_j^2\).
A final smooth upper-cap grouping obeys the full upper bound already used
in \cref{thm:v26-cutoff}. All mixed-shell terms are retained exactly once.
\end{proposition}
\begin{proof}
For a product of primitive factors use the ordered telescoping identity
\[
 A_1\cdots A_m-B_1\cdots B_m
 =\sum_{i=1}^m B_1\cdots B_{i-1}(A_i-B_i)A_{i+1}\cdots A_m.
\]
Take \(A_i,B_i\) to be the cumulative covariances times their following
vertices at successive scales. At least one internal covariance is an
annular increment. All other routed momenta differ by at most five soft
momenta, so every nonzero line is comparable to \(\Lambda_j\).
After \(5-n\) external derivatives, the integrand has order \(-5\);
normalized four-dimensional annular counting gives \(\Lambda_j^{-1}\).
The comparison has order \(a^3\Lambda_j^{-2}\) before integration,
and hence size \(a^3\Lambda_j^2\). The argument applies after expanding
each background inverse into its finite primitive chains. It proves the
ownership assertion even for words containing several types of covariance.
The auxiliary coefficient is assigned once as a local/density insertion,
not at each propagating increment. Geometric summation proves uniformity
in the number of these propagating shells. It proves no self-map of an
arbitrary nonlinear action neighborhood.
\end{proof}

\section{The convergent action--density representative and matched schemes}
\label{sec:v26-normalization}

A normalization detail matters before stating a continuum functional.
Let \(\mathsf P=\mathsf P_{5,4}\), \(g_a=\mathsf P\Gamma_{a,1}\),
and \(j_a=\mathsf P\log J_G\), with \(J_G\) the actual coordinate
Jacobian of V12. V24's measured identity is
\(B_a=d_5(g_a-j_a)-x_a\). Let
\begin{equation}
 H_{\lambda,\theta}
 =\sum_{w<4}\frac{Y_{\lambda,\theta}^{[w]}}{4-w}
                          +h_\theta^{[4]}+H_{\rm phys},
 \qquad d_5H_{\rm phys}=0,
 \label{eq:v26-H}
\end{equation}
where the first term is V24's explicit transgression and the second is
V25's particular critical primitive. Optional finite physical terms must
be one compatible closed functional; a source-independent constant is
anchored to zero. The critical primitive retains the
literature input of \cref{thm:v25-critical-exact}; its large rational
matrix is not evaluated anew. Then
\(d_5H_{\lambda,\theta}=\mathsf P\mathcal X_{0,\lambda,\theta}\).

After accumulating, rather than double-counting, the previous increments,
the total order-\(\hbar\) counterterm is
\begin{equation}
 C_a^{\rm tot}=-g_a+j_a+H_{\lambda,\theta}.
 \label{eq:v26-total-C}
\end{equation}
Define the Jacobian-compensated coefficient representative
\begin{equation}
 \boxed{\mathcal F_{a,L}
  =\Gamma_{a,1}+C_a^{\rm tot}-j_a
  =(I-\mathsf P)\Gamma_{a,1}+H_{\lambda,\theta}.}
 \label{eq:v26-functional}
\end{equation}
The theorem below concerns this Jacobian-compensated representative.
It does not establish a limit of the unpaired coordinate coefficient
\(\Gamma_{a,1}+C_a^{\rm tot}\). At the retained coefficient order the
corresponding action--density bookkeeping is the algebraic identity
\[
 e^{-(\Gamma_{a,1}+C_a^{\rm tot})}J_G
                  =e^{-\mathcal F_{a,L}},
\]
where \(J_G\) and its logarithm are consistently truncated in external
arity. This reorganizes the additional coordinate term in the total
normalization; it does not add a second physical auxiliary determinant. In a convention keeping flat \(dq\) and displaying the Jacobian
term in the coefficient action instead, \(\Gamma_{a,1}+C_a^{\rm tot}
=\mathcal F_{a,L}+j_a\). That coordinate contact need not converge
separately. No new base measure is inferred from the existence of the
coefficient representative.

\begin{theorem}[One common normalized continuum source functional]
\label{thm:v26-functional}
On the bank \eqref{eq:v26-panel}, and retaining the exact scope of V25's
local primitive, \(\mathcal F_{a,L}\) converges to
\begin{equation}
 \mathcal F_{0,L}=(I-\mathsf P)\Gamma_{0,1,L}
                                      +H_{\lambda,\theta}
 \label{eq:v26-limit}
\end{equation}
in every stated finite windowed multilinear source norm. Its component
errors are \eqref{eq:v26-main} and \eqref{eq:v26-kernel}. The independent
thermodynamic comparison is \eqref{eq:v26-volume}. These statements use
one functional counterterm, including all its background and antifield
contacts, not independent subtractions of different observations.

Two bounded smooth ultraviolet completions in the declared class, with
the same ordinary action/source jets, lower reference profile and local
functional \(H_{\lambda,\theta}\), have the same limit. On a common
source core their complementary difference is bounded by
\(C_I\chi^{-n}(a+a')\mu^{5-n}\) in each fixed rooted source norm.
\end{theorem}
\begin{proof}
The same local \(H_{\lambda,\theta}\) appears at every cutoff, so it
cancels from the comparison. Apply \cref{thm:v26-cutoff,thm:v26-kernels}
to every one of the finitely many retained panels. Their exchange-symmetric
coefficients define one polynomial source functional, and the common
projection \(\mathsf P\) preserves those symmetries. This proves the
first assertion with all derivative contacts. For two completions, compare
each to the same ordinary subtracted integral and add the estimates.
Every completion uses its \emph{own actual} local jet \(g_a\) in
\eqref{eq:v26-total-C}; equality before that matching is not asserted.
Changing the lower reference profile instead changes a finite Wilsonian
functional, not just an ultraviolet scheme. Its local comparison remains
V25's ordered profile transgression, and no equality of the untransported
lower-reference functionals is asserted here.
\end{proof}

\section{The full minimal modified identity at this arity}
\label{sec:v26-identity}

It remains to specify what symmetry statement follows, rather than calling
convergence alone a BV theorem. Let \(\mathcal X_{0,\lambda}\) be the
\emph{unsubtracted} lower-reference insertion of the inherited off-shell
trace identity, now restricted only after assembling the nonminimal terms.
It is a smooth finite-coefficient functional because its nonpolynomial
words are confined to the lower annulus. This is not set equal to zero.

\begin{lemma}[Ordinary nonlocal identity with the reference retained]
\label{lem:v26-ordinary-identity}
In thermodynamic normalization and at total arity at most five,
\begin{equation}
 \boxed{d_5\bigl[(I-\mathsf P)\Gamma_{0,1}\bigr]
                       =(I-\mathsf P)\mathcal X_{0,\lambda}.}
 \label{eq:v26-ordinary-identity}
\end{equation}
\end{lemma}
\begin{proof}
Use the ordinary classical master action and its fixed harmonic canonical
shift, not the finite filtered source. Its off-shell one-loop trace identity
is the one used in V22--V25. Before removing an auxiliary smooth ultraviolet
trace cutoff, retain both the lower and upper reference terms and the flat
coefficient divergence. The latter is a local coordinate contact and is
included in \(j\). No nilpotence of the lattice source enters this step.

At the required arity each coefficient is a finite ordered word. The
nontrivial upper-reference words contain an upper transition factor and at
least one covariance; they can be routed to a fixed annulus at that upper
scale \(M\). A no-covariance term is a local trace polynomial and is
assigned to \(\mathsf P\). If an output slot has \(n\) antifields,
the marked word has superficial degree \(5-n\): its one transformation
variation adds one derivative to the degree \(4-n\) proved above.
After its common subtraction through degree \(5-n\), its
\(6-n\) derivatives have integrand order \(-5\). Annular counting
bounds the upper remainder by \(C M^{-1}|K|^{6-n}\). The same bound
controls momentum shifts made under this finite regulated trace. It tends
to zero. This is the convergence justification needed to apply the
ordinary trace identity, rather than an identity for a divergent loop.

The lower-reference words remain integrable and unchanged. Passing to the
limit is justified by \eqref{eq:v26-loop-bounds} and the same differentiated
majorants. Finally the common ordinary coefficient projection commutes
with \(d_5\) by \cref{thm:v24-projection}. This gives
\eqref{eq:v26-ordinary-identity}. The argument treats all source slots
in the bank simultaneously; no conclusion about the next arity or the
\(\hbar^2\) antibracket of two quantum coefficients follows.
\end{proof}

\begin{theorem}[Normalized minimal one-loop modified Slavnov identity]
\label{thm:v26-identity}
With the local input \eqref{eq:v26-H}, the thermodynamic continuum
functional satisfies
\begin{equation}
 \boxed{d_5\mathcal F_{0,\infty}
                       =\mathcal X_{0,\lambda}^{\le5}.}
 \label{eq:v26-modified-identity}
\end{equation}
This is the ghost-number-one identity on the complete minimal arity-five
bank, including its metric-, ghost-, and multiple-antifield coefficients.
It uses V25's local exactness with its stated literature input, not a new
calculation of that local cohomology.

For a finite-cutoff diagnostic define
\(\mathcal W_{a,L}=\operatorname{pr}_{A\le5}
(s_a\mathcal F_{a,L})-\mathcal X_{a,\lambda,L}\), using the same
source and reference conventions on both terms. For every fixed scaled
momentum seminorm on a window bounded by \(\mu\), an output slot with
\(n\) minimal antifields obeys the conservative estimate
\begin{equation}
\begin{split}
 \|\mathcal W_{a,L,I}\|_{r,\mu}
 \le C_{I,r,b,H}\chi^{-n}\mathfrak n_I\lambda^{5-n}
 \left[a\lambda+(\lambda L_{\min})^{-b}
                  +\frac{\mathbf1_{n=0}}{\lambda^4V_L}\right],
 \qquad b>4.
\end{split}
 \label{eq:v26-W-rate}
\end{equation}
Thus the entire stated minimal row has a vanishing cutoff residual in
thermodynamic normalization, or in a joint cutoff/volume limit. Exact
finite-grid nilpotence is not asserted.
\end{theorem}
\begin{proof}
Add \(d_5H=\mathsf P\mathcal X_{0,\lambda}\) to
\eqref{eq:v26-ordinary-identity}. This proves
\eqref{eq:v26-modified-identity} and explains why the full lower-reference
term, including its local part, appears on its right.

For the quantitative comparison, \(d_5\) is a finite sum of products of
local ordinary vertices and input coefficients with linearly substituted
momenta. A source variation contributes one derivative. An Euler insertion
removes one antifield and contributes two derivatives and its factor
\(\chi\). Therefore \eqref{eq:v26-all-derivatives} contributes at
most \(C\chi^{-n}a\mu^{6-n}\) to an output slot, and the independent
volume comparison contributes
\(C\chi^{-n}\mu^{6-n}\lambda^{-1}(\lambda L_{\min})^{-b}\).
These are bounded by the corresponding terms in
\eqref{eq:v26-W-rate}.

The unprojected reference words obey the same compact-support comparison
as V24's local words: before taking a Taylor jet, at each fixed retained
arity their momenta lie in \(\lambda/2\le|p|\le3\lambda\), apart
from the regular source-free no-covariance trace in the lower ball. On
that fixed set the action-Hessian difference is
\(O(\chi a^3\lambda^5)\), the free inverse difference starts at
fourth relative mesh order, and the source agrees with its ordinary
polynomial. Differentiating a finite ordered product gives
\(C\chi^{-n}\lambda^{5-n}(a\lambda)^3\) in scaled norms. Poisson
summation gives its displayed volume error. Exclusion of the single zero
mode in the source-free no-covariance trace contributes at most the last
term in \eqref{eq:v26-W-rate}; no such trace occurs when \(n>0\).

Finally, replacing the ordinary differential by \(s_a\) acts only on
external soft labels in this diagnostic. Its source part agrees on the
common plateau. Its free action discrepancy starts at degree six and its
higher action discrepancy at degree five with factor \(a^3\). The
renormalized representative \(\mathcal F_{a,L}\), unlike the raw
counterterm, is uniformly bounded in every fixed source seminorm by
\cref{thm:v26-functional}. Hence these insertions are
\(O((a\lambda)^3\lambda^{5-n})\), with constants also depending on
the chosen finite local \(H\). The possibly divergent \(j_a\) was
not estimated this way: its source variation cancels its coordinate
measure contribution exactly and is separately retained in the pair
\eqref{eq:v26-functional}. Adding the preceding bounds, and using
\((a\lambda)^3\le a\lambda\), proves \eqref{eq:v26-W-rate}.
No finite lattice symmetry was assumed in proving its continuum bound.
\end{proof}

The identity just proved is for the reduced minimal coefficient complex
and its specified nonminimal restoration. It does not assert a theorem
for arbitrary independently retained nonminimal antifields, or for the
internal Standard-Model ghosts. It is also a modified identity at fixed
\(\lambda>0\), not an unmodified zero-reference master equation. Its
local matching part is inherited from V25, while its nonlocal convergence
and the continuity needed to assemble every retained row are the new
analytic conclusions.

\section{Verification, uniformity and the next Einstein--SM interface}
\label{sec:v26-ledgers}

The accompanying checker verifies exact labelled inverse and determinant
coefficients through five marks on noncommuting rational matrices,
Berezin elimination with a ghost-antifield quadratic term, complete
mixed-shell ownership, structural arity counts, ultraviolet exponents,
Taylor remainders, and action--density bookkeeping. These checks are finite
algebraic diagnostics of the displayed identities. They neither evaluate
the native loop integrals nor verify V25's external cohomology theorem.
The native large local matrix remains unevaluated. The appended source
contains the analytic proofs and is not presented as proof-assistant output.

\subsection*{Proved}
\addcontentsline{toc}{subsection}{V26: Proved}
The background-complete factorization \eqref{eq:v26-factorization}
identifies every measured one-loop coefficient on the minimal arity-five
bank. The marked whole-zone symbol estimates and full-cutoff comparison
\eqref{eq:v26-main} hold also for the previously unbounded multi-metric
and mixed-antifield panels. Their rooted source kernels and independent
volume limits obey \eqref{eq:v26-kernel}--\eqref{eq:v26-volume}.
Every mixed primitive shell is assigned once and the complementary shell
increments are summable. Consuming the inherited particular local
primitive gives the common coefficient representative
\eqref{eq:v26-functional}, its scheme-matched limit, and the minimal
one-loop modified identity \eqref{eq:v26-modified-identity} with the
finite diagnostic bound \eqref{eq:v26-W-rate}.

\subsection*{Inherited}
\addcontentsline{toc}{subsection}{V26: Inherited}
V1--V25 are unchanged. In particular the native finite input, complex metric
reference, periodic upper completion, auxiliary amplitude, source lift,
stationary mesh generator, finite arity-seven input budget, ordinary
component projection and the particular critical/subcritical local
primitive keep all their original qualifications. The new result does not
rename the already proved zero-background source compiler or a general
Gaussian identity as a new theorem about ESM. The classical Einstein--SM
interface is not used as a quantum premise.

\subsection*{Literature input}
\addcontentsline{toc}{subsection}{V26: Literature input}
The general BRST/exact-RG and modified-reference setting is
\cite{V15IIM}. The local critical existence consumed here is precisely
V25's use of \cite{V25BBHEinstein} on its extendable ordinary pure-gravity
germ. Neither literature input supplies a finite-regulator nonlocal bound
or a Standard-Model determinant-line theorem. The cutoff estimates above
are obtained from the specified finite operands, not an invocation of a
generic EFT existence assertion.

\subsection*{Conditional}
\addcontentsline{toc}{subsection}{V26: Conditional}
The local finite normalization remains the particular solution provided by
V24--V25. Additional incompatible linear normalization conditions are not
licensed; they require V25's relative range test. A microscopic matching
to the discontinuous pre-completion regulator remains separate. Reading
the result as a full ESM theorem would require the mixed gravity--gauge,
Higgs, spin/chiral, and combined-measure estimates that have not been
established for this full source bank. No all-loop or infinite-series
convergence follows.

\subsection*{Open}
\addcontentsline{toc}{subsection}{V26: Open}
The one-loop minimal pure-gravity arity-five analytic/local interface is
now assembled, at fixed lower scale and with the declared completion.
The next task should consume it in a \emph{mixed} common-action Hessian
rather than revisit an isolated pure-gravity coefficient. A useful first
step is the simultaneous hard metric--Higgs determinant with its mixed
Schur terms and geometry-dependent scalar sources, using V9's scalar
module and the actual common measure. Internal gauge and chiral fields
then require their own covariant source completion and local cohomology
and line tests. Independently, higher perturbative orders require
subdivergence/forest bounds; the absence of one-loop subgraphs does not
extend automatically to two loops. The lower scale, contour comparison,
full interacting EFT self-map and numerical critical matching remain
unresolved where earlier ledgers place them.

\begin{center}\small
\begin{tabular}{>{\raggedright\arraybackslash}p{.23\textwidth}>{\raggedright\arraybackslash}p{.68\textwidth}}
\toprule
Variable & Scope in V26\\\midrule
Volume & Same-torus cutoff and rooted norms are uniform for
\(\lambda L_\nu\ge1\). Thermodynamic normalization of the modified
identity and the separate volume/zero-mode errors are explicit.\\
Mesh & Full completed Brillouin zone, \(a\lambda\le2^{-18}\).
No matching to an uncompleted discontinuous symbol is asserted.\\
Sources & Ghost-zero minimal bank \(A\le5\), and its ghost-one rows.
All finite momentum derivatives and rooted moments; constants depend on
those orders. No independent full nonminimal source theorem.\\
Loop and coupling & One loop, fixed \(\chi>0\), lower scale
\(\lambda\ge4096\mu>0\). No infrared removal or uniform
\(\chi\downarrow0\) bound.\\
Measure & Same complex Gaussian coefficient reference and once-counted
auxiliary measure; coordinate Jacobian explicitly paired with its action
representative. No nonperturbative positivity or Lorentzian continuation.\\
Shell count & Summable complementary increments of this finite module.
Not an invariant neighborhood of arbitrary interacting effective actions.\\
Local computation & Critical existence and its finite norm prescription
inherited; native large matrix and matching integrals remain unevaluated.\\
\bottomrule
\end{tabular}
\end{center}

\clearpage
\part*{V27: A simultaneous hard metric--Higgs block and mixed source continuum}
\addcontentsline{toc}{part}{V27: Simultaneous metric--Higgs block and mixed source continuum}

\section{The mixed determinant is the next common-action interface}
\label{sec:v27-start}

Sections 1--208 and their labels are preserved. The pure-gravity module in
V26 and the scalar-background module in V9 cannot be added and called a
simultaneous gravity--matter elimination. The off-diagonal hard Hessian
is nonzero at a retained Higgs background. Its contraction contributes
at the same loop order as the metric contact in
\cref{prop:v13-Higgs}. That contact was explicitly calculated with the
scalar retained, and its mixed loop was explicitly left open. The present
part computes the missing coupled coefficient and its geometric
variations, using one finite Gaussian/Berezin reference.

The claimed component has two ordinary Higgs fields and at most three
ordinary metric marks, at one loop and at zeroth order in internal gauge,
Yukawa and Higgs self-interaction couplings. The internal links are the
identity and the real Higgs carrier has four components. It is the
\(\chi^{-1}\phi^2\) component of the programme, not a replacement of the
entire Higgs potential by a free physical theory. A mass parameter
\(0\le m\le M\lambda\) is retained, where the positive lower Wilsonian
scale \(\lambda\) is never removed. The allowance \(m=0\) concerns only
this hard-covariance problem; all modes below \(\lambda\) stay retained.
No assertion about a massless infrared functional follows.

The one-loop coefficient is defined by the finite Gaussian/Berezin
moment expansion of the displayed action jets. At fixed external degree,
the one-loop part has finitely many connected graphs. A convergent
quadratic-hard integral is established separately; no uniform
\(O(\hbar^2)\) remainder for the nonlinear theory is asserted.

The zero-background pure-gravity action, ghost action, contour, and
auxiliary measure remain those of V26. Minimal matter antifields,
internal gauge fields and chiral fields are not integrated into a new
Slavnov theorem here. In particular, the pure-gravity cohomology theorem
consumed in V25 is not applied to an enlarged matter complex.

\subsection{Source audit and the additional coefficient-reference choice}

The native scalar carrier and Hodge action are inherited from
\eqref{eq:v9-native-Q}. The affine coframe chart and classical scalar
mesh correction are inherited from
\eqref{eq:v12-signed-Hodge}--\eqref{eq:v12-H2}.
The actual holomorphic continuation in V13 has a hyperbolic temporal
scalar symbol; it is not V9's Euclidean sine symbol. This matters once
that scalar is integrated. We therefore record a smooth periodic
completion of this continued scalar family. It agrees exactly with the
continued, corrected native coefficients on the common interior
plateau and keeps the periodic Hodge interactions in the upper region.
This is additional ultraviolet regulator data, on the same field
carrier. It is not a proof of equivalence between complete Lorentzian
and Euclidean measures, nor a match to an uncompleted nonperiodic
regulator.

One-loop gravity coupled to scalars, background Hessians and their
coordinate dependence are established subjects; see
\cite{V27Moss} for a distinct covariant background-field construction.
No result or coefficient from that construction is imported below.
The elementary Schur identity is also inherited. The new claims are the
specified simultaneous operands, their full-zone mixed estimates,
and the signed tensor contraction controlling the two-Higgs response.

\section{One scalar completion consistent with the inherited interior}
\label{sec:v27-scalar}

Use four odd periodic grids, physical pairing \(a^4\sum_x\), and the
Euclidean coefficient chart of \eqref{eq:v13-Wick-map}:
\[
 e(q)=I+\mathsf E_E(q),\qquad G(q)=e(q)^Te(q),\qquad
 \rho(q)=\det e(q),\qquad W(q)=\rho(q)G(q)^{-1}.
\]
The metric hard variable lies on the fixed contour \(T\R^{10}\), with
\(T=(I-P_{\rm tr})+iP_{\rm tr}\). Scalar hard variables are real in this
coefficient reference. This uses the recorded substitution
\(\phi_L=i\phi_E\) from V13, not a new physical scalar reality condition.
All pairings in the action and all transposes below are algebraic; complex
conjugation is used only for norm estimates.

For two scalar momenta \(k,r\), define
\[
 \epsilon_0=-1,\quad \epsilon_i=1\ (i=1,2,3),\qquad
 v^n_0(p)=a^{-1}\sinh(ap_0),\quad
 v^n_i(p)=a^{-1}\sin(ap_i).
\]
Also put \(\kappa^n_0(p)=2a^{-1}\sinh(ap_0/2)\),
\(\kappa^n_i(p)=2a^{-1}\sin(ap_i/2)\), and
\(c^n_0(t)=\cosh(at/2)\), \(c^n_i(t)=\cos(at/2)\).
The continued sign-averaged kinetic tensor is
\begin{equation}
 D^n_{\mu\nu}(k,r)=
 \begin{cases}
 v^n_\mu(k)v^n_\nu(r),&\mu\ne\nu,\\
 \kappa^n_\mu(k)\kappa^n_\mu(r)c^n_\mu(k_\mu+r_\mu),&\mu=\nu.
 \end{cases}
 \label{eq:v27-Dnative}
\end{equation}
Only its symmetric part contracts with \(W\). Define the polynomial
\begin{equation}
 E_{2,\mu\nu}(k,r)=
 \frac{\epsilon_\mu k_\mu^3r_\nu+
             \epsilon_\nu k_\mu r_\nu^3}{6}
       +\frac{\delta_{\mu\nu}\epsilon_\mu}{4}k_\mu^2r_\mu^2,
 \qquad \widehat D^n=D^n+a^2E_2.
 \label{eq:v27-Dcorrected}
\end{equation}
These definitions are obtained by continuing the actual difference
symbols and the V12 correction. The minus sign in \(\epsilon_0\)
records the hyperbolic temporal symbol. Let \(S^n_{H,a}\) be the scalar
quadratic action with Fourier coefficient
\begin{equation}
 \frac12\,\phi(k)\cdot\phi(r)
 \left\{m^2\widehat\rho(-k-r)
       -\widehat W^{\mu\nu}(-k-r)\widehat D^n_{\mu\nu}(k,r)\right\}.
 \label{eq:v27-native-scalar-action}
\end{equation}
The physical sums in this formula are understood. Hats on \(\rho,W\)
mean Fourier coefficients, whereas the hat on \(D^n\) means the explicit
mesh correction. Formula \eqref{eq:v27-native-scalar-action} is used
only on its nonaliased interior chart.

Let \(S^b_{H,a}(q,\phi)\) be V9's uncorrected, periodic Euclidean
Hodge action at \(g=G(q),U=I,E=0\). Its tensor \(D^b\) is
\eqref{eq:v27-Dnative} with sine and cosine in all four directions.
Write \(I^n=S^n_{H,a}(q,\phi)-S^n_{H,a}(0,\phi)\), and define
\(I^b\) analogously. Use exactly the profiles \(s,\zeta\) of
\eqref{eq:v14-profiles} and \(S=s(aD)\). The free scalar symbols are
\begin{align}
 P^n_a(p)&=m^2+\sum_\mu(\kappa^n_\mu(p))^2
                +\frac{a^2}{12}\sum_\mu\epsilon_\mu p_\mu^4,\nonumber\\
 P^b_a(p)&=m^2+\lambda_a(p),\qquad
 \widetilde P_a=\zeta(ap)P^n_a+(1-\zeta(ap))P^b_a.
 \label{eq:v27-P}
\end{align}
The complete scalar reference and interaction are specified by one action:
\begin{equation}
 \boxed{\widetilde S_{H,a}(q,\phi)
 =\frac12\langle\phi,\widetilde P_a\phi\rangle
   +I^b(q,\phi)+I^n(Sq,S\phi)-I^b(Sq,S\phi).}
 \label{eq:v27-completed-Higgs}
\end{equation}
For the present coefficients, only metric degree at most five in this
analytic expression is consumed. This is total action degree at most
seven, matching V24's existing gravitational input budget. All scalar
Hodge factors occur in the action; the scalar coefficient measure is
not multiplied by a second \(\rho\).

\begin{proposition}[Interior matching, scalar floor and retained upper interactions]
\label{prop:v27-scalar-reference}
The action \eqref{eq:v27-completed-Higgs} agrees with the continued,
corrected native scalar action whenever every field leg has
\(|ap|\le1/64\). It is smooth periodic in every hard momentum.
For all principal momenta,
\begin{equation}
 m^2+\frac4{\pi^2}|p|^2\le\widetilde P_a(p)
                         \le m^2+C|p|^2.
 \label{eq:v27-scalar-floor}
\end{equation}
Its interior free and fixed-degree interaction symbols differ from
their ordinary counterparts by \(O(a^4R^6)\), with a factor
\(R^{-1}\) for each differentiated momentum, when all labels have
size at most \(R\), \(aR<c\), and \(m\le MR\).
On compact upper regions their order-two derivative bounds are
\(C_\alpha a^{|\alpha|-2}\), with the mass contribution retained.
Upper scalar and metric modes are not spectators.
\end{proposition}
\begin{proof}
Expanding sine and hyperbolic sine, including the same-sign diagonal
average, gives
\(D^n_{\mu\nu}=k_\mu r_\nu-a^2E_{2,\mu\nu}+O(a^4R^6)\).
This verifies both the correction and its differentiated bound. At
\(r=-k\) it gives exactly \(P^n_a\). In the temporal direction,
\(4a^{-2}\sinh^2(ap_0/2)-a^2p_0^4/12\ge p_0^2\), by the
power series with nonnegative remaining terms. Spatially,
\(\kappa_i^2\ge4p_i^2/\pi^2\), and the quartic addition is nonnegative.
The same bounds hold for \(P^b_a\); convex combination proves the
lower bound. The upper bound follows on the compact principal cube.

The only nonperiodic tensors in \eqref{eq:v27-completed-Higgs} occur
inside the central profiles. At a fixed degree every field occurrence
in the difference \(I^n(Sq,S\phi)-I^b(Sq,S\phi)\) has its own
profile. Thus no branch choice occurs at a momentum face. The bare
piece is periodic already. On the plateau the two bare interactions
cancel and \(\widetilde P=P^n\); this proves exact matching.
Rescaling by \(a\) proves the upper derivative estimates. The pointwise
coefficients \(W,\rho\) have bounded fixed jets on the coframe chart.

If a field leg is above the support of \(s\), the native correction
term vanishes but the corresponding bare Hodge interaction generally
does not. For example, at nonzero mass the one-metric vertex with a
constant scalar leg is \(m^2I_4/2\), independently of the other scalar
momentum. At zero mass a small nonzero scalar momentum and a hard
momentum in the same spatial direction give a nonzero sine kinetic
vertex. The completion has therefore not removed high interactions.
\end{proof}

\section{The simultaneous Hessian and its exact mixed coefficient}
\label{sec:v27-hessian}

Write \(\widetilde S_{H,a}=\langle\phi,Q_a(q)\phi\rangle/2\).
The symbol \(Q_a\) in this part refers to
\eqref{eq:v27-completed-Higgs}, not silently to V9's unmodified operator.
Let \(\nu_\lambda=1-\Theta_\lambda\), and set
\[
 C_g=\nu_\lambda\widetilde K_a^{-1},\qquad
 C_H=\nu_\lambda\widetilde P_a^{-1}I_4.
\]
Both are zero on retained modes; inverses below are taken only on their
hard supports. The metric covariance has its fixed contour meaning.
Define the background inverses without ever inverting a zero covariance:
\begin{equation}
 G_g(q)=(I+C_gV_g(q))^{-1}C_g,\qquad
 G_H(q)=(I+C_H[Q_a(q)-\widetilde P_a])^{-1}C_H.
 \label{eq:v27-background-inverses}
\end{equation}
Here \(V_g\) is V26's actual pure-gravity hard Hessian perturbation.

For hard metric directions \(x,y\) and a hard scalar \(z\), define
\begin{equation}
 \begin{split}
 U(q;\phi)[x,y]&=\tfrac12\langle\phi,D_q^2Q_a(q)[x,y]\phi\rangle,\\
 \langle x,L(q;\phi)z\rangle
       &=\langle D_qQ_a(q)[x]\phi,z\rangle.
 \end{split}
 \label{eq:v27-UL}
\end{equation}
The coupled bosonic hard Hessian is, exactly,
\begin{equation}
 \boxed{\mathbb H(q,\phi)=
 \begin{pmatrix}G_g(q)^{-1}+U(q;\phi)&L(q;\phi)\\
                  L(q;\phi)^T&G_H(q)^{-1}\end{pmatrix}.}
 \label{eq:v27-hard-Hessian}
\end{equation}
The gravitational ghost block is unchanged and has no \(\phi\)
dependence in this \(\phi\)-independent harmonic gauge. The Higgs
has no spacetime Lorentz-connection charge. Thus its direct Hessian
with the auxiliary Lorentz connection vanishes; the Cartan stationary
section and its once-counted density remain functions of \(q\), not
of \(\phi\). None of these statements applies to the internal
weak connection or to fermions, which are absent in this coefficient.

\begin{proposition}[A common joint Gaussian germ]
\label{prop:v27-joint-germ}
For fixed \(M\), sufficiently small
\(\|q\|_{\mathrm W}+\chi^{-1/2}\|\phi\|_{\mathrm W}\), low retained
Fourier support, and \(0\le m\le M\lambda\), the Hessian
\eqref{eq:v27-hard-Hessian} has a convergent Gaussian integral on
\((T\R^{10})\oplus\R^4\) and a common analytic coefficient germ.
The smallness threshold is independent of mesh and periods under
\eqref{eq:v27-domain} below. This is a statement about the joint
quadratic hard integral, not convergence of the full nonlinear
coframe path integral.
\end{proposition}
\begin{proof}
Use V14's metric accretivity and \eqref{eq:v27-scalar-floor} as the
reference energy. On the hard support every nonzero covariance label
has size at least \(\lambda\). The fixed-degree metric/background
perturbations obey the inherited relative energy bound. The scalar
Hodge form and its filtered correction give a relative bound
\(C_M\|q\|_{\mathrm W}\): each difference has order one, and each
correction has two additional derivatives compensated by \(a^2\)
on its central support. The upper bare Hodge term satisfies the same
bound directly. In the off-diagonal block the two derivatives are split
between the soft \(\phi\) and the hard scalar, or replaced by \(m^2\).
Cauchy--Schwarz bounds its normalized operator by
\(C_M\chi^{-1/2}\|\phi\|_{\mathrm W}\). The normalized \(U\) bound
is \(C_M\chi^{-1}\|\phi\|_{\mathrm W}^2\). Translations of hard labels
by the soft momenta preserve the energy comparison. Convolution is
bounded by the Wiener norm, with no volume factor. Choosing their sum
less than half the free accretivity margin proves the assertion.
The same relative norm estimate gives a holomorphic inverse and a
logarithm near the reference. Small values of \(\nu\) improve, rather
than worsen, these covariance-weighted estimates.
\end{proof}

\begin{theorem}[The complete two-Higgs component of the simultaneous pushforward]
\label{thm:v27-mixed-trace}
Remove the common factor \(\hbar\) from one-loop coefficients. At
scalar degree two, the difference between the simultaneous coefficient
and its value at \(\phi=0\) is
\begin{equation}
 \boxed{\Gamma^{\rm mix}_{a,1}(q,\phi)
 =\frac12\Tr(G_g(q)U(q;\phi))
 -\frac12\Tr(G_g(q)L(q;\phi)G_H(q)L(q;\phi)^T).}
 \label{eq:v27-main-trace}
\end{equation}
All geometric derivatives are derivatives of this one expression.
It is the measured one-loop coefficient on the external bank
\(q^b\phi^2\), \(0\le b\le3\), not only a formal background-Hessian
ansatz. The scalar determinant at \(\phi=0\), the pure-gravity/ghost
determinants, and the auxiliary density are present in the full
functional, but disappear from this difference.
\end{theorem}
\begin{proof}
The ordinary block determinant gives
\[
 \det\mathbb H=\det G_H^{-1}\det G_g^{-1}
 \det\{I+G_g(U-LG_HL^T)\}.
\]
Both \(U\) and \(LG_HL^T\) have scalar degree two. The first term
of the last logarithm is therefore the whole coefficient at that
degree. This uses the inherited Schur identity on the newly identified
operands. Integrating the metric first gives the same determinant and
hence the same signs and factors.

To justify its use for the nonlinear pushforward, a hard-linear vertex
in a graph with at most five external arguments carries the sum of a
subset of those momenta, of size at most \(5\mu\). It cannot meet a
covariance supported above \(\lambda\). With such vertices excluded,
a connected one-loop graph has as many internal edges as vertices,
and hard valence at least two at each vertex. Every hard valence is
therefore exactly two. These are precisely the Hessian cycles above.
The ghost action is \(\phi\)-independent and its determinant cancels
in the stated difference. An order-\(\hbar\) auxiliary density insertion
would acquire another loop factor if a hard field were contracted;
its stationary evaluation is also \(\phi\)-independent here, since
the possible hard force has just been excluded. This proves completeness
at the declared coefficient order.
\end{proof}

For example the first metric derivative of the mixed bubble is
\begin{equation}
 \begin{split}
 D_h\Tr(G_gLG_HL^T)={}&
 -\Tr(G_gA^g[h]G_gLG_HL^T)
 +\Tr(G_gL_hG_HL^T)\\
 &-\Tr(G_gLG_HQ_hG_HL^T)
 +\Tr(G_gLG_HL_h^T).
 \end{split}
 \label{eq:v27-background-four}
\end{equation}
There is also the derivative of the contact term, including
\(-G_gA^g[h]G_g\). Holding either propagator fixed would lose an
actual geometric-source coefficient.

A quadratic-hard source check is useful but has a narrower scope.
For external hard-linear probes \(J_g,J_H\), the same Gaussian has
\(-\langle(J_g,J_H),\mathbb H^{-1}(J_g,J_H)\rangle/2\) in its
negative logarithm. Its off-diagonal inverse is
\(-G_gL G_H+O(\phi^3)\). Thus mixed probe contacts are retained.
This identity for the Hessian Gaussian is not asserted to account for
all nonlinear composite-source coefficients of the complete theory.

\section{The actual mixed vertex and its signed continuum contraction}
\label{sec:v27-vertex}

At \(q=0\), both scalar families have
\(\rho_1=\operatorname{tr}q/2\),
\(W_1=(\operatorname{tr}q)I/2-q\).
For any kinetic tensor \(D\), put
\[
 \mathcal V[D](k,r)=\tfrac12(D+D^T)
                      +\tfrac12(m^2-\operatorname{tr}D)I_4.
\]
The actual one-metric/two-scalar vertex in \eqref{eq:v27-UL} is
\begin{equation}
 \boxed{\mathcal V_a(k,r)=\mathcal V[D^b](k,r)
 +s(a(k+r))s(ak)s(ar)
       \{\mathcal V[\widehat D^n](k,r)-\mathcal V[D^b](k,r)\}.}
 \label{eq:v27-actual-V}
\end{equation}
The mass parts cancel in the braces. The flat free adjustment does not
produce a metric vertex. The two-metric contact is obtained from the
same formula for the action, with the product of profiles on both
metric legs; its coframe coefficient is
\begin{equation}
 W_2=d_2I-d_1q+q^2-\mathsf E_E(q)^T\mathsf E_E(q),\qquad
 \rho_2=d_2.
 \label{eq:v27-W2}
\end{equation}
This retains the chart contact already identified in V12. It is not the
seagull in exact metric-difference coordinates.

At zero mesh define
\begin{equation}
 V_0(k,r)=\tfrac12(kr^T+rk^T)
                       +\tfrac12(m^2-k\cdot r)I_4,
 \qquad JX=X-\tfrac12\operatorname{tr}(X)I_4.
 \label{eq:v27-V0}
\end{equation}
Let \(\Pi^{\rm mix}\) be the second scalar derivative, so that
\(\Gamma^{\rm mix}=\frac12\int_k\phi(-k)\cdot\phi(k)\Pi^{\rm mix}(k)\)
at \(q=0\).

\begin{proposition}[Evaluated numerator of the mixed metric--scalar loop]
\label{prop:v27-numerator}
For all real \(k,r\) and \(m\ge0\),
\begin{equation}
 \boxed{V_0(k,r):JV_0(k,r)
       =\frac12 k^2r^2+m^2k\cdot r-m^4=:N(k,r).}
 \label{eq:v27-numerator}
\end{equation}
Consequently the continuum mixed bubble is
\begin{equation}
 \boxed{\Pi^{\rm bub}_0(k)
 =-\frac2\chi\int_p
 \frac{\nu_\lambda(p)\nu_\lambda(p+k)
                     N(k,-p-k)}{p^2((p+k)^2+m^2)}.}
 \label{eq:v27-cont-bubble}
\end{equation}
Divergent expressions in this formula mean the common regulated
integral before its local Taylor assignment; subtracted expressions
below are absolutely convergent. The coefficient is diagonal in the
four real Higgs components. There is no extra factor of four for a
fixed external pair.
\end{proposition}
\begin{proof}
Set \(s=k\cdot r\) and \(A=(kr^T+rk^T)/2\).
Then \(\operatorname{tr}A=s\),
\(\operatorname{tr}A^2=(k^2r^2+s^2)/2\), and
\(\operatorname{tr}V_0=2m^2-s\). Substitution into
\(V_0:V_0-(\operatorname{tr}V_0)^2/2\) gives
\eqref{eq:v27-numerator}. This signed contraction uses the inherited
flat metric covariance \(2\nu J/(\chi p^2)\), not a positive
replacement for its trace component. The factor two in
\eqref{eq:v27-cont-bubble} is the second scalar derivative of
\eqref{eq:v27-main-trace}. The scalar line connects the two external
Higgs indices by a Kronecker delta; it is not a closed independent
four-species trace. Contracting on the rotated contour gives the same
answer since \(T(T^TKT)^{-1}T^T=K^{-1}\).
\end{proof}

The contact contribution is inherited, not recomputed as a new result:
\begin{equation}
 \Pi^{\rm con}_0(k)=\frac{I_\lambda}{\chi}
 \left[-\frac12 k_0^2+2\sum_{i=1}^3k_i^2-\frac94m^2\right],
 \qquad I_\lambda=\int_p\frac{\nu_\lambda(p)}{p^2}.
 \label{eq:v27-contact}
\end{equation}
It is exactly the second scalar derivative of
\eqref{eq:v13-Higgs-value}. It is local and needs its actual ultraviolet
assignment. It is not to be added a second time after
\eqref{eq:v27-main-trace} has been used. Its anisotropy is a coordinate
and gauge dependent local coefficient, not an invariant physical mass.

For later checks, with \(D=(p+k)^2+m^2\), the signed numerator obeys
\begin{equation}
 \frac{N(k,-p-k)}{p^2D}
 =\frac{k^2-m^2}{2p^2}+\frac{m^2}{2D}
               -\frac{m^2(k^2+m^2/2)}{p^2D}.
 \label{eq:v27-partial-fraction}
\end{equation}
This is an integrand identity; divergent summands are never assigned
independent momentum routings by it. Averaging at \(p=R\omega\),
\(\omega\in S^3\), gives
\begin{equation}
 \frac1{2\pi^2}\int_{S^3}\frac{N(k,-R\omega-k)}{R^2((R\omega+k)^2+m^2)}\dd\omega
 =\frac{k^2}{2R^2}-\frac{m^2(k^2+m^2)}{R^4}+O(R^{-6}).
 \label{eq:v27-angular}
\end{equation}
The remainder is uniform for bounded \(k,m\); the linear angular
moment vanishes and
\(\operatorname{av}_{S^3}(\omega\cdot k)^2=k^2/4\). In this normalization there is
no \(k^4\) logarithmic coefficient in this specific flat mixed
bubble. The logarithmic coefficient under a radial loop-momentum
cutoff is \(m^2(k^2+m^2)/(4\pi^2\chi)\). This is a normalization
check on this one component, not the physical Higgs beta function.

\section{Whole-zone bounds for two scalar sources and their geometric derivatives}
\label{sec:v27-bounds}

An external panel \(I\) has exactly two Higgs arguments and \(b\)
metric arguments, \(0\le b\le3\). Thus it has \(b+1\) independent
four-momenta \(K\). Let \(\gamma_{a,L,I}\) be the coefficient per
physical volume obtained by differentiating
\eqref{eq:v27-main-trace}. The fixed component norms are those of V26;
\(\mathfrak n_I\) is the product of argument norms. Use total Taylor
degree four in all independent external momenta:
\begin{equation}
 \gamma^R_{a,L,I}=(I-\mathsf T_4)\gamma_{a,L,I}.
 \label{eq:v27-subtraction}
\end{equation}
Every Taylor jet is stored in one local \(q^b\phi^2\) functional.
This redundant component bank is not identified with a covariant
curvature--matter basis before the coupled symmetry equations are solved.

The common domain is
\begin{equation}
 \boxed{\lambda\ge4096\mu>0,\qquad a\lambda\le2^{-18},\qquad
 \lambda L_\nu\ge1,\qquad 0\le m\le M\lambda.}
 \label{eq:v27-domain}
\end{equation}
All external legs have magnitude at most \(\mu\). The mass parameter
and \(\lambda\) are fixed during a mesh limit. Constants can depend
on \(M\), fixed profile seminorms, derivative order and the stated
coupling/reference margins, not on mesh, periods or shell count.

\begin{lemma}[Native mixed integrand estimates]
\label{lem:v27-symbols}
With \(R_p=\lambda+|p|\), the complete finite sums of loop integrands
and their zero-mesh comparisons obey, for each fixed pair of
multi-indices,
\begin{equation}
 |\partial_p^\beta\partial_K^\alpha f_{a,I}|+
 |\partial_p^\beta\partial_K^\alpha f_{0,I}|
 \le C_{I,\alpha,\beta}\chi^{-1}\mathfrak n_I
                     R_p^{-|\alpha|-|\beta|}.
 \label{eq:v27-f-bound}
\end{equation}
On a common interior plateau the difference is bounded by
\begin{equation}
 |\partial_p^\beta\partial_K^\alpha(f_{a,I}-f_{0,I})|
 \le C_{I,\alpha,\beta}\chi^{-1}\mathfrak n_I
                         a^3R_p^{3-|\alpha|-|\beta|}.
 \label{eq:v27-f-difference}
\end{equation}
For \(b=0\), the right side improves to
\(C_{\alpha,\beta}\chi^{-1}\mathfrak n_I
 a^4R_p^{4-|\alpha|-|\beta|}\).
On an upper region \(|ap|\ge\delta>0\), the corresponding bound is
\(C\chi^{-1}\mathfrak n_I a^{|\alpha|+|\beta|}\).
\end{lemma}
\begin{proof}
Every metric or scalar covariance has order minus two; a metric
covariance contributes \(\chi^{-1}\). Pure-gravity Hessian insertions
have order two and factor \(\chi\). Each scalar Hessian vertex has
order at most two, by its explicitly displayed sign averages and
filtered correction. A mixed \(L\) has two scalar derivative factors
or one \(m^2\), so it also has order at most two. The contact term has
one metric line and an order-two vertex; the bubble has one metric and
one scalar line and two order-two vertices. Both therefore have order
zero with exactly one net factor \(\chi^{-1}\). A metric derivative
of either background inverse inserts one order-two vertex and one
additional inverse, preserving that degree and the net coupling power.
Leibniz differentiation covers every ordered term through three marks.

In the plateau, the scalar covariance and every scalar vertex differ
from the ordinary ones by fourth relative mesh order, by
\cref{prop:v27-scalar-reference}. The free metric covariance has the
same improvement; a differentiated gravitational background Hessian
has the inherited third relative mesh order. Telescoping a product
proves \eqref{eq:v27-f-difference}. With no metric mark there is no
pure-gravity interaction insertion, so every difference starts at
fourth relative order. On upper regions use the smooth periodic
symbols and rescale all hard labels by \(a\). Each momentum
derivative supplies a factor \(a\), proving the upper bound.

Cutoff derivatives are supported at \(R_p\asymp\lambda\) and obey
the same scaled inequalities. On the support of any complete word,
all hard momenta differ by at most \(5\mu\); hence they have comparable
\(R_p\). Kernels with zero covariance below the lower scale are
extended smoothly by zero. The scalar mass never creates an infrared
singularity on this support. The estimates take no norm of an extensive
action and use no Ward identity.
\end{proof}

\begin{theorem}[Full-cutoff mixed source theorem]
\label{thm:v27-continuum}
For \(0\le b\le3\), \(j=|\alpha|\le5\), and
\eqref{eq:v27-domain},
\begin{equation}
 \boxed{|\partial_K^\alpha(\gamma^R_{a,L,I}-\gamma^R_{0,L,I})(K)|
 \le C_{I,\alpha}\chi^{-1}\mathfrak n_I\,a|K|^{5-j}.}
 \label{eq:v27-main-rate}
\end{equation}
Every prescribed higher derivative is controlled on a fixed external
window. In the unmarked two-Higgs sector, write \(K=k\). Then
\begin{equation}
 \boxed{|\partial_k^\alpha(\Pi^R_{a,L}-\Pi^R_{0,L})(k)|
 \le C_\alpha\chi^{-1}a^2|k|^{6-j},\qquad j\le6.}
 \label{eq:v27-two-point-rate}
\end{equation}
These estimates include the entire completed Brillouin zone. They
concern the coefficients with their common \(\hbar\) removed.
\end{theorem}
\begin{proof}
Taylor's integral remainder through degree four uses five external
derivatives. The continuum integrand is then bounded by
\(C\chi^{-1}|K|^5R_p^{-5}\), and the interior difference by
\(C\chi^{-1}a^3|K|^5R_p^{-2}\). In four dimensions the continuum
tail above \(c/a\) is \(O(a)\), while
\(a^3\int_\lambda^{c/a}R\,\dd R=O(a)\).
The upper region has \(O(a^{-4})\) normalized mode count and an
\(O(a^5)\) fifth derivative, also giving \(O(a)\). The same
estimates hold for normalized finite sums: covering momentum space by
cubes of side \(2\pi/L_\nu\), and using \(\lambda L_\nu\ge1\),
bounds the sums of the displayed radial majorants by fixed multiples
of their integrals. This proves \eqref{eq:v27-main-rate}, including
differentiated remainders.

For \(b=0\), scalar flavour symmetry and the symmetry of the second
variation imply \(\Pi(k)=\Pi(-k)\). The complete integrand can be
replaced by its even external average without changing the coefficient.
Consequently its fifth Taylor jet is its fourth. Six derivatives give
\(R_p^{-6}\) for the continuum and \(a^4R_p^{-2}\) for the interior
difference. Their tail and sum are both \(O(a^2)\). The sixth upper
derivative has size \(a^6\) times the \(a^{-4}\) mode count. This
proves \eqref{eq:v27-two-point-rate}. No parity assumption is made for
the multi-metric panels. Higher differentiated estimates follow from
the same inverse and symbol bounds, or from the scaled smooth Taylor
remainder on the fixed window. The constants need not be uniform in
that derivative order.
\end{proof}

The subtraction precedes the limiting integral. It does not set a
divergent trace to zero. The local coefficients have the conservative
sizes \(\chi^{-1}a^{r-4}\) at derivative order \(r<4\) and
\(\chi^{-1}[1+\log(1/(a\lambda))]\) at order four, with their actual
mass factors and cancellations retained. No universality of those raw
coefficients is claimed. Their finite matching specifies the selected
renormalized component functional.

\section{Source kernels, simultaneous shell transport, and a massless check}
\label{sec:v27-kernels}

\begin{corollary}[Rooted estimates and independent volume comparison]
\label{cor:v27-kernels}
Window the external symbols smoothly at scale \(\mu\). For each fixed
moment order \(s\), their difference kernels obey
\begin{equation}
 \boxed{a^{4(b+1)}\sum_{\boldsymbol x}
 \left(1+\mu\sum_{i=1}^{b+1}d_{\rm per}(x_i,0)\right)^s
 \|\mathcal K_{a,L,I}^{\Delta}(\boldsymbol x)\|
 \le C_{I,s}\chi^{-1}a\mu^5.}
 \label{eq:v27-kernel}
\end{equation}
For \(b=0\) this improves to \(C_s\chi^{-1}a^2\mu^6\).
For every fixed \(t>4\), the separate same-reference volume comparison is
\begin{equation}
 |\partial_K^\alpha(\gamma^R_{0,L,I}-\gamma^R_{0,\infty,I})|
 \le C_{I,\alpha,t}\chi^{-1}\mathfrak n_I
 |K|^{5-j}\lambda^{-1}(\lambda L_{\min})^{-t},\qquad j\le5.
 \label{eq:v27-volume}
\end{equation}
The unmarked even sector has instead
\(|k|^{6-j}\lambda^{-2}(\lambda L_{\min})^{-t}\), \(j\le6\).
\end{corollary}
\begin{proof}
Rescale the \(4(b+1)\) independent external variables by \(\mu\).
The proof of \cref{thm:v27-continuum} gives bounds for more than
\(s+4(b+1)\) derivatives. Repeated integration by parts in the inverse
Fourier transform gives an integrable weighted majorant. Periodizing
and using \(d_{\rm per}\le d\) for every lift proves
\eqref{eq:v27-kernel}. This yields a multilinear functional estimate
with one argument in \(L^1\) and all others in \(L^\infty\), without
multiplication by the volume. For \eqref{eq:v27-volume}, differentiate
the subtracted continuum integrand in the loop momentum more than four
times and apply Poisson summation. Each derivative supplies
\(\lambda^{-1}\), and the basic convergent integral has scale
\(\lambda^{-1}\), or \(\lambda^{-2}\) in the even sector. Smoothness
at the lower reference boundary justifies these derivatives.
\end{proof}

\begin{proposition}[Mixed shell ownership and matched completions]
\label{prop:v27-shells}
Expand every occurrence of \(C_g\) and \(C_H\), including those in
background inverses, into the actual smooth positive-weight innovations
of its own reference. Assign a complete loop word to its largest
primitive covariance index. Then every mixed word occurs once, and
\begin{equation}
 \|(I-\mathsf T_4)\Delta_j\Gamma_{a,1,I}^{\rm mix}\|_{\rm src}
       \le C_I\chi^{-1}\mathfrak n_I\mu^5\Lambda_j^{-1}.
 \label{eq:v27-shell-rate}
\end{equation}
At \(b=0\), the complete even increment is bounded by
\(C\chi^{-1}\mu^6\Lambda_j^{-2}\). The complements are summable
uniformly in the number of shells. Two smooth completions with the
same ordinary limit and the stated interior/upper bounds have the same
nonlocal limiting mixed coefficients after their actual local jets
are matched. This does not establish matching to an uncompleted
nonperiodic theory.
\end{proposition}
\begin{proof}
A word of primitive covariances has a unique maximal shell label;
expanding the cumulative product and subtracting the preceding one
is the sum over exactly those labels. In particular, for the unmarked
bubble the increment of \(C_gLC_HL^T\) is not just its diagonal-shell
term: both cross-species mixed scale terms remain. An annular normalized
sum costs \(\Lambda_j^4\); five differentiated order-zero factors give
\(\Lambda_j^{-5}\), proving \eqref{eq:v27-shell-rate}. In the even
sector six derivatives give \(\Lambda_j^{-6}\). All shifted internal
labels are comparable on the stated external window. Finally subtract
the two common local jets and compare both completions to their shared
ordinary coefficient by \cref{thm:v27-continuum}. Local coefficients
are matched as one background functional, not independently discarded.
\end{proof}

\begin{proposition}[Massless two-Higgs remainder is a lower-reference term]
\label{prop:v27-massless}
In thermodynamic normalization at \(m=0\), define the finite integral
\begin{equation}
 J_\lambda(k)=\int_p\frac{\nu_\lambda(p)\Theta_\lambda(p+k)}{p^2}.
 \label{eq:v27-Jlambda}
\end{equation}
Then the continuum coefficient obtained above satisfies
\begin{equation}
 \boxed{\Pi_0^R(k)=\frac{k^2}{\chi}
             \{J_\lambda(k)-\mathsf T_2J_\lambda(k)\}.}
 \label{eq:v27-massless}
\end{equation}
Its defining nonlocal integral is supported near the fixed lower
reference scale. This statement is specific to the declared gauge,
coordinate and coupling component; it is not absence of quantum
matter effects or an infrared-removal theorem.
\end{proposition}
\begin{proof}
At \(m=0\), \eqref{eq:v27-numerator} gives
\(N=k^2(p+k)^2/2\). Cancelling this factor before integration reduces
the bubble to
\(-k^2\int_p\nu(p)\nu(p+k)/\chi p^2\). Write
\(\nu(p+k)=1-\Theta(p+k)\) with a common ultraviolet regulator.
The first term is a local quadratic polynomial. The contact
\eqref{eq:v27-contact} is also quadratic. Both disappear from the
complement, but remain in the local assignment. The second term gives
\eqref{eq:v27-massless}; inversion symmetry makes \(J\) even, so
multiplication by \(k^2\) intertwines these two Taylor assignments.
On its support \(|p|\ge\lambda\) and \(|p+k|\le2\lambda\), proving
localization of the loop momentum. Nothing in this argument removes
\(\lambda\), and the small-external-window assumption depends on it.
\end{proof}

\section{What the coupled coefficient supplies, and its remaining symmetry equation}
\label{sec:v27-ledgers}

The result is a mixed common-reference calculation, not an addition of
two already renormalized sector determinants. It supplies the
\(q^b\phi^2\), \(b\le3\), connected coefficients and their metric
source derivatives from \eqref{eq:v27-main-trace}, with the same scalar
coframe dependence in its contact and mixed blocks. The continuum
coefficient's zero-loop action is the zero-mesh Hodge coupling of the
common finite carrier. This is a tree-level compatibility statement,
not a deduction of a quantum theorem from the upstream classical
Einstein--SM interface.

The new component local bank includes, before symmetry and redundancy
reduction, all \(q^b\phi^2\) polynomials through four derivatives.
Terms such as the metric variations of \(R\phi^2\), the scalar kinetic
and mass terms, and higher-derivative redundant terms may participate
in a coupled normalization. Their invariant decomposition and physical
coefficients are not determined by convergence alone. In particular,
the pure-gravity local primitive from V25 does not renormalize the
scalar transformation source automatically.

The next smallest symmetry test has one diffeomorphism ghost and two
Higgs fields. Besides the mixed coefficient above it consumes the
scalar-antifield insertion \(\phi^*c\cdot\partial\phi\), its quantum
correction, the metric source variation of the scalar two-point
coefficient, and the actual lower-reference and measure variations.
In schematic form its action--antifield terms include
\[
 D_q\Gamma_1^{(\phi^2)}[R_0c]
 +D_\phi\Gamma_1^{(\phi^2)}[\mathcal L_c\phi]
 +D_\phi S_{H,0}[z_1^\phi]
 +D_qS_{H,0}[z_1^q].
\]
The required arities and background terms must be extracted together;
this display is not an asserted complete finite identity. The source
completion of internal gauge, spin and chiral sectors remains separate.
No inference of the mixed critical cohomology from the pure-gravity
classification is made.

\subsection*{Proved}
\addcontentsline{toc}{subsection}{V27: Proved}
The recorded scalar completion preserves the continued, corrected native
interior and has a periodic positive scalar free symbol. It combines
with the inherited metric contour in one small joint Hessian germ.
The actual simultaneous two-Higgs coefficient is the contact-minus-mixed
trace \eqref{eq:v27-main-trace}; its first three metric variations,
all inverse variations and all cross-species shell pairs are included.
The explicit finite vertex has the zero-mesh contraction
\eqref{eq:v27-numerator}. The complete coefficient has the full-zone
\(O(a)\) source comparison, improved to \(O(a^2)\) for the unmarked
even two-point sector, with independent volume and rooted-kernel bounds.
The massless fixed-lower-scale remainder is
\eqref{eq:v27-massless}. These are one-loop coefficient results on the
declared gravitational--Higgs component.

\subsection*{Inherited}
\addcontentsline{toc}{subsection}{V27: Inherited}
V1--V26 remain unchanged. The metric/ghost periodic completion and its
contour estimates, the native coframe map, the V12 scalar correction,
V9's finite Higgs carrier and sign-averaged Hodge prescription, V13's
metric contact, and the general measured Schur and shell calculus are
used at their original scope. The interior scalar continuation is not
identified with V9's different temporal reference. The pure-gravity
critical primitive remains a pure-gravity input.

\subsection*{Literature input}
\addcontentsline{toc}{subsection}{V27: Literature input}
Gravity--scalar background quantization is established, including
constructions different from the present fixed gauge and coframe chart
\cite{V27Moss}. No external self-energy, beta coefficient, or coupled
anomaly theorem supplies a step in the native cutoff estimates here.
The proof is a finite-reference calculation; its Gaussian and Schur
identities are not claimed as novel general principles.

\subsection*{Conditional}
\addcontentsline{toc}{subsection}{V27: Conditional}
Interpretation as a complete mixed ESM source functional requires the
combined diffeomorphism, internal-gauge and local-Lorentz BV normalization, matter
antifields and composite-source variations. The component counterterms
here need not separately satisfy those equations. Equivalence to a
Lorentzian distribution or a microscopic acceptance law requires its
own matching and contour theorem. V9's full scalar determinant can be
compared with the new scalar completion only with its actual local
matching; it is not set equal to it at finite mesh.

\subsection*{Open}
\addcontentsline{toc}{subsection}{V27: Open}
Construct and renormalize the mixed one-ghost/two-Higgs source row using
this simultaneous determinant, rather than recomputing an isolated
scalar or metric determinant. Then add internal gauge and fermionic
hard blocks and their phase data. Higgs self-interactions, nonzero Higgs
vacuum backgrounds, infrared removal, higher-loop forests and a full
invariant interacting EFT class are not consequences of this result.
No nonperturbative or finite-parameter ultraviolet-gravity claim is made.

\begin{center}\small
\begin{tabular}{>{\raggedright\arraybackslash}p{.23\textwidth}>{\raggedright\arraybackslash}p{.68\textwidth}}
\toprule
Variable & Scope in V27\\\midrule
Cutoff and volume & Full completed zone; fixed profile margins and
\(a\lambda\le2^{-18}\). Rooted and same-torus estimates are uniform
for \(\lambda L_\nu\ge1\); thermodynamic comparison is separate.\\
Sources & Two ordinary Higgs fields and zero through three metric marks;
all prescribed finite derivative and kernel-moment orders. No all-source-
order constant and no matter-antifield theorem.\\
Couplings and mass & One loop, fixed \(\chi>0\), coefficient of zero
internal gauge, Yukawa and scalar self-couplings; \(0\le m\le M\lambda\).
The positive lower scale is retained, including at \(m=0\).\\
Measure & Same metric contour, four real hard Higgs coordinates,
\(\phi\)-independent gravitational gauge; auxiliary and coordinate
densities counted once. Only a small Hessian Gaussian germ is asserted.\\
Iteration & Summable complementary mixed coefficients and exact
cross-species shell allocation, not an analytic interacting self-map.\\
\bottomrule
\end{tabular}
\end{center}

The accompanying checker uses exact polynomial and rational algebra for
the temporal/spatial correction, joint determinant and source signs,
background inverse derivatives, the native continuum numerator, the
large-momentum coefficient and the complete cross-species shell sum.
A separate floating-point finite-Hessian diagnostic is labelled as
such. Neither these checks nor successful typesetting independently
certifies the inherited proofs or the analytic continuum estimates.
No new Lean formalization or numerical value for a native lattice
matching integral is reported.

\begingroup
\renewcommand{\refname}{Additional comparison reference for V27}
\begin{thebibliography}{99}
\bibitem[45]{V27Moss}
I.~G.~Moss, \emph{Covariant one-loop quantum gravity and Higgs inflation},
arXiv:1409.2108 (2014). A comparison on scalar--gravity background quantization
and coordinate dependence; no native coefficient or bound is imported.
\end{thebibliography}
\endgroup


\clearpage
\part*{V28: The scalar transformation source and a mass-graded mixed identity}
\addcontentsline{toc}{part}{V28: Scalar antifield, mass grading, and the mixed identity}

\section{The missing source and the mass-dependent subtraction interface}
\label{sec:v28-start}

Sections 1--215, including their qualifications and references, are retained.
V27 computes the simultaneous metric--Higgs determinant and explicitly leaves
its coupled symmetry equation open. A determinant alone does not determine
how the integrated Higgs transforms. This installment specifies that source,
computes its quantum correction in the same simultaneous Gaussian, and
assembles the one-ghost/two-Higgs row. No internal gauge, Yukawa or Higgs
self-coupling is switched on. There are still four real Higgs components,
with a common mass, on the fixed coefficient contour and periodic completion
of \eqref{eq:v27-completed-Higgs}.

A second interface appears before a symmetry-compatible subtraction can be
claimed. At positive mass the scalar Euler operator contains a term of
momentum degree zero. Consequently the momentum-only Taylor orders used for
separate coefficients in V27 are not, without extra finite matching, a
chain assignment for the coupled row. We resolve this by giving the squared
mass weight two in a common local Taylor assignment. The propagators keep
their actual mass. This is an explicitly changed local normalization, not
a reinterpretation of V27's fixed-mass subtraction or an infrared limit.

The new results are the scalar-source loop and its cutoff estimate, the
mass-graded row and its finite breaking/reference identity, and convergence
of its nonlocal modified identity. Construction of a ghost-number-zero
primitive for the \emph{local coupled} row remains open. In particular,
V25's pure-gravity cohomology input is not applied to gravity with matter.
The exact-RG/modified-Slavnov formalism is established methodology
\cite{V28IIM}; the results below identify its additional native operands.

\subsection{A precise obstruction to reusing the momentum-only assignment}

Write \(\mathfrak m=m^2\). For a function of momenta, \(\mathsf T_d^K\)
denotes its Taylor polynomial at zero momentum, with \(\mathfrak m\)
held fixed. The following calculation is about the assignment, not a claim
that this particular test function is the native loop.

\begin{proposition}[The massive Euler factor changes the Taylor interface]
\label{prop:v28-mass-obstruction}
At fixed \(\mathfrak m>0\), the identity
\[
 \mathsf T_5^K\{(k^2+\mathfrak m)Z\}
       =(k^2+\mathfrak m)\mathsf T_3^KZ
\]
does not hold for general analytic odd source coefficients. More precisely,
\[
 \mathsf T_5^K\{(k^2+\mathfrak m)Z\}
 =k^2\mathsf T_3^KZ+\mathfrak m\mathsf T_5^KZ.
\]
\end{proposition}
\begin{proof}
Multiplication by the quadratic homogeneous polynomial shifts Taylor degree
by two; multiplication by the fixed constant \(\mathfrak m\) does not.
For an explicit one-variable test take
\(Z(t)=t^5/(\lambda^2+t^2)\), \(\lambda>0\). Then
\(\mathsf T_3^t Z=0\), whereas the left side has coefficient
\(\mathfrak m t^5/\lambda^2\). This proves the assertion even with odd
momentum parity. No ultraviolet or infinite-volume argument is involved.
\end{proof}

V27 did not assert the failed identity: its counterterms were explicitly
componentwise. The correction here is to the \emph{new use} of those
components in a coupled equation. One could instead retain additional
finite momentum jets at each positive mass. We use the following finite
weighted bank because it also records the mass insertions themselves.

\section{A scalar source on the existing coframe and ghost carrier}
\label{sec:v28-source}

Keep \(S=s(aD)\), \(P=S^2\), the source
\eqref{eq:v18-repaired-source}, and the same zero-antifield action.
To avoid confusing the profile \(P\) with the scalar Euler operator, the
latter is always denoted \(\widetilde P_a\). Add the scalar antifield
\(\upsilon_A=\phi_A^*\), \(A=1,\ldots,4\), with odd parity, and set
\begin{equation}
 \boxed{r_a^H(\phi,c)=P\bigl((Pc)^\mu D_\mu\phi\bigr),\qquad
 S_{\mathrm{src}}^H=\langle\upsilon,r_a^H(\phi,c)\rangle.}
 \label{eq:v28-scalar-source}
\end{equation}
The antifield is placed before the ghost in the coefficient convention.
This is additional \emph{source} data. It changes neither V27's scalar
quadratic form nor V14's ghost action. The latter is not replaced by a
gauge fixing of the new source. The source is independent of the coframe
coordinate, as appropriate to a spacetime scalar rather than a scalar
half-density. Its Hodge volume remains inside the action; there is no
additional scalar coefficient-density factor.

For a scalar wave of momentum \(r\) and ghost wave of momentum \(t\),
its exact vertex is
\begin{equation}
 \boxed{\mathscr S_a(r,t)c
   = iP(a(r+t))P(at)\,r_\mu c^\mu.}
 \label{eq:v28-source-symbol}
\end{equation}
Here \(P(az)=s(az)^2\), and all sums use the finite cyclic convention.
A contributing output and ghost are in the central chart, so their
scalar difference is away from any momentum face. Thus the derivative
in this vertex has a smooth periodic realization on its support.

\begin{proposition}[Finite scalar divergence and the remaining source square]
\label{prop:v28-divergence}
The scalar source has exactly zero coefficient divergence on every odd
finite grid:
\begin{equation}
 \boxed{\operatorname{Tr}_{\phi}D_\phi r_a^H=0.}
 \label{eq:v28-divergence}
\end{equation}
With \(u=Pc\), \(b_c=u\cdot Du\), and the V18 ghost source, its square
on nonaliased panels is
\begin{equation}
 \boxed{\breve s^{\,2}\phi
 =P\{\mathcal L_{(P-I)b_c}\phi
             +\mathcal L_u(I-P)\mathcal L_u\phi\}.}
 \label{eq:v28-source-square}
\end{equation}
It is zero on an interior panel where every required filter equals one;
it is not asserted to vanish on all finite modes.
\end{proposition}
\begin{proof}
The field derivative is \(P M_{u^\mu}D_\mu\). Cyclicity puts
\(D_\mu P\) on the diagonal of the trace. This diagonal is constant
and equals the paired Fourier sum of \(ip_\mu s(ap)^2\), which is
zero. Each of the four scalar components has this same zero. No product
rule is used for this finite trace.

The odd Leibniz rule gives
\(\breve s^{\,2}\phi=P\mathcal L_{P b_c}\phi
-P\mathcal L_uP\mathcal L_u\phi\). On the declared panels
\(\mathcal L_u\mathcal L_u=\mathcal L_{b_c}\). Adding and subtracting
\(P\mathcal L_{b_c}\phi\) proves \eqref{eq:v28-source-square}.
The formula retains the same filter-curvature mechanism as V18; adding a
scalar representation does not turn that filtered source into a full
finite differential.
\end{proof}

The ordinary source at zero mesh is \(r_0^H=c\cdot\partial\phi\).
Together with the ordinary chart-correct gravitational source it preserves
the zero-mesh scalar action
\(\frac12\int\sqrt{\det g}(g^{\mu\nu}\partial_\mu\phi\partial_\nu\phi
+\mathfrak m\phi^2)\): its variation is the integral of a Lie derivative
of a scalar density. This is a check on the ordinary scalar representation,
not an invariance assertion for the completed finite action.

\section{The quantum scalar transformation is a three-propagator return}
\label{sec:v28-loop}

Use the background inverses \(G_g(q),G_H(q)\) of
\eqref{eq:v27-background-inverses} and the actual ghost inverse \(G_c(q)\)
of V26. Let \(x,z,\beta,\bar\beta\) be respectively hard metric,
scalar, ghost and antighost variables. Define \(B(q,c)\) by the actual
action term \(\langle\bar\beta,B(q,c)x\rangle\). In particular at
\(q=0\) it is the ghost-action vertex in \eqref{eq:v15-LTV}; it is
not a derivative of \eqref{eq:v28-scalar-source}.

The mixed action block \(L(q;\phi)\) is exactly \eqref{eq:v27-UL}.
Define an even, antifield-linear rectangular matrix \(N_a(q;\upsilon,c)\)
by
\begin{equation}
 \boxed{\langle x,N_a(q;\upsilon,c)z\rangle
  =\left\langle\upsilon,
    P\left([P G_c(q)B(q,c)x]^\mu D_\mu z\right)\right\rangle.}
 \label{eq:v28-N}
\end{equation}
The product of the odd antifield and the odd retained ghost is read in
that order. Every occurrence of \(G_c\) in this formula is part of the
loop; it is not absorbed into an unmarked numerical vertex.

\begin{theorem}[Complete scalar-antifield coefficient in the simultaneous law]
\label{thm:v28-source-loop}
Remove the common factor \(\hbar\) from one-loop coefficients. On the
source bank \(\upsilon q^b\phi c\), \(0\le b\le2\), the simultaneous
one-loop source functional is
\begin{equation}
 \boxed{\langle\upsilon,\mathcal Z_a^H(q;\phi,c)\rangle
       =\Tr\bigl(G_g(q)N_a(q;\upsilon,c)G_H(q)L(q;\phi)^T\bigr).}
 \label{eq:v28-source-loop}
\end{equation}
There is no direct scalar transformation tadpole. At zero metric background
this is one cycle with one metric, one ghost and one scalar propagator.
It is diagonal in the scalar species for a fixed external pair; it has no
extra factor of four.
\end{theorem}
\begin{proof}
The only hard-degree-two part of the new source is
\(\langle\upsilon,P((P\beta)\cdot Dz)\rangle\).
The finite antighost integral substitutes
\(\beta=-G_c(q)B(q,c)x\). Consequently the bosonic off-diagonal block is
\(L(q;\phi)-N_a(q;\upsilon,c)\). The term of the joint determinant
linear in each of \(\phi\) and \(\upsilon\) is
\[
 \frac12\Tr\{G_g L G_H N_a^T+G_g N_a G_H L^T\}.
\]
Both coefficients are even. Algebraic symmetry of the bosonic covariances
and the transpose/cyclic trace identity make the two summands equal.
This proves the sign and the factor in \eqref{eq:v28-source-loop}.
The finite Berezin calculation, rather than a chosen diagram sign, fixes
the minus sign in the off-diagonal block before the determinant expansion.

The source is linear in \(\phi\), independent of \(q\), and contains
only one ghost. It has no two-metric or two-scalar hard contact with an
external ghost. Hence no source tadpole is available. At most five external
arguments occur in the stated bank. A vertex with one hard leg would have
momentum at most the sum of their soft momenta, below \(\lambda\); its
covariance is zero. A connected one-loop graph with all hard valences at
least two has exactly two hard legs at every vertex, by
\(\sum_v(\deg v-2)=2(\ell-1)=0\). Thus the quadratic-hard reduction
exhausts the measured coefficient. All additional background marks are
obtained by differentiating its inverses and vertices. An order-\(\hbar\)
auxiliary-amplitude insertion could enter this source coefficient at the
same loop order only through a hard-linear branch, which is absent.

The scalar propagator connects the external scalar and antifield indices.
No independent closed species trace remains. The fixed metric contour
implements the algebraic inverse, just as in V27.
\end{proof}

For example, \(D_hG_c=-G_cA^c[h]G_c\), so \(D_hN_a\) contains both
\(D_hB\) and \(-G_cA^c[h]G_cB\). The first geometric derivative of
\eqref{eq:v28-source-loop} has the four terms obtained by differentiating
\(G_g,N_a,G_H,L^T\), with
\(D_hG_g=-G_gA^g[h]G_g\) and
\(D_hG_H=-G_HQ_hG_H\). The second derivative acts on all these
terms again. Holding any species inverse fixed changes the source module.

\subsection{An explicit ordinary momentum integrand}

Let the external scalar have momentum \(k\), the external ghost momentum
\(t\), and use the hard metric momentum \(p\). Put
\(r=k-p\), \(w=p+t\). With the tensor \(V_0\) from
\eqref{eq:v27-V0}, define
\[
 \mathcal R_1(H,p;c,t)=i\{(c\cdot p)H+t c^TH+Hc t^T
                    -\mathsf B(H,tc^T+ct^T)\},
\]
\[
 F_w^{\rm Four}H=i\{H w-\tfrac12 w\operatorname{tr}H\},\qquad
 B_{p,t}(H,c)=-F_w^{\rm Four}\mathcal R_1(H,p;c,t).
\]
All matrix multiplications use the existing Euclidean coframe map in
\(\mathsf B\). In the ordinary reference the exact source integrand is
\begin{equation}
 \boxed{z_0(p;k,t)c
 =i\,r^T C_c(w)
       B_{p,t}\bigl(C_g(p)V_0(k,-r),c\bigr)\,C_H(r).}
 \label{eq:v28-routed-integrand}
\end{equation}
Here \(C_H(r)=\nu_\lambda(r)/(r^2+\mathfrak m)\),
\(C_g(p)=2\nu_\lambda(p)J/(\chi p^2)\), and
\(C_c(w)=\nu_\lambda(w)I/w^2\).
It is simply \eqref{eq:v28-N} and \eqref{eq:v28-source-loop} with a
fixed routing. At finite mesh the actual scalar vertex and ghost matrix
replace these ordinary tensors, and the source factors in
\eqref{eq:v28-source-symbol} are retained. No loop integral is evaluated
by naming this integrand. For unweighted ordinary propagators, \(\chi=1\), \(\mathfrak m=2\),
\(p=(2,3,1,4)\), \(k=(1,0,0,0)\), \(t=(0,1,0,0)\), and
\(c=(1,0,0,0)\), its value is \(-421i/32190\). This is one
exact rational integrand evaluation, not a value of the integrated loop;
the accompanying checker verifies it independently.

\section{Mass-graded subtraction and cutoff control of the new source}
\label{sec:v28-subtraction}

For a smooth hard coefficient \(f(K,\mathfrak m)\), define
\begin{equation}
 \boxed{\mathbf T_D f
 =\sum_{|\alpha|+2j\le D}
       \frac{K^\alpha\mathfrak m^j}{\alpha!j!}
       \partial_K^\alpha\partial_{\mathfrak m}^j f(0,0).}
 \label{eq:v28-weighted-T}
\end{equation}
The squared mass is a constant, ghost-neutral parameter. Its use in this
Taylor assignment does not change the propagator to its massless value.
All finite remainders use the unexpanded massive propagator.
Here ``nonlocal remainder'' means the complement of this declared finite
local bank. It may contain higher-weight local terms; the terminology
does not assert separated spatial support for every retained coefficient.

On the hard support the mass derivatives are regular even at zero mass:
\begin{equation}
 \boxed{\partial_{\mathfrak m}^j C_H(p;\mathfrak m)
  =(-1)^j j!\,\nu_\lambda(p)
                 \widetilde P_a(p;\mathfrak m)^{-j-1}.}
 \label{eq:v28-mass-derivative}
\end{equation}
In particular \(\partial_{\mathfrak m}C_H\ne-C_H^2\) on the cutoff
transition: the derivative has one \(\nu\), not two. Mass derivatives
of background inverses are taken by the corresponding actual inverse
rule with their free and interaction mass insertions included.

\begin{lemma}[Weighted Taylor calculus on the massive hard domain]
\label{lem:v28-T}
If \(P_l(K,\mathfrak m)\) is a homogeneous polynomial of weight \(l\),
with weights one for \(K\) and two for \(\mathfrak m\), then
\begin{equation}
 \boxed{\mathbf T_{D+l}(P_l f)=P_l\mathbf T_Df.}
 \label{eq:v28-homogeneous-T}
\end{equation}
The assignment is bounded by one in the finite weighted coefficient-jet
norm. On a compact scaled hard domain its polynomial inclusion into any
fixed smooth norm is bounded by a constant depending only on the degree,
number of variables and that norm. No mesh or volume constant occurs.
\end{lemma}
\begin{proof}
Expand both factors into monomials and add their weights. This proves
\eqref{eq:v28-homogeneous-T}, including matrix-valued polynomials with
their order fixed. The coefficient-jet statement is coordinate projection.
For the smooth-norm statement rescale \(K=\lambda\widehat K\),
\(\mathfrak m=\lambda^2\widehat{\mathfrak m}\); differentiation of
a finite polynomial is a finite coefficient map on the compact scaled
set. To obtain Taylor remainder bounds one can apply ordinary Taylor's
integral formula to \(f(K,z^2)\). This is even in \(z\), and its
ordinary degree-\(D\) Taylor polynomial is exactly
\(\mathbf T_Df(K,z^2)\). Derivatives with respect to \(\mathfrak m\)
commute with truncation after lowering the degree by two; their integral
bounds follow instead from the smooth hard inverse
\eqref{eq:v28-mass-derivative}. There is no division by \(z\) at zero.
\end{proof}

For each coefficient \(\mathcal Z^H_{a,I}\) in
\eqref{eq:v28-source-loop}, assign its actual local jet
\(\mathbf T_3\mathcal Z^H_{a,I}\) to the source action and write
\(\mathcal Z^{H,R}_{a,I}=(I-\mathbf T_3)\mathcal Z^H_{a,I}\).
For the V27 metric--Higgs coefficients use \(I-\mathbf T_4\) instead
of their earlier fixed-mass momentum projector. At each fixed mass the
difference
\begin{equation}
 (\mathsf T_D^K-\mathbf T_D)f
 \label{eq:v28-matching-change}
\end{equation}
is a local polynomial in the external momenta. Its actual coefficients
are part of the change of matching convention, not deleted. This makes
the use of V27's component theorem in a new normalization explicit.

\begin{lemma}[Mixed source symbol estimates with mass derivatives]
\label{lem:v28-symbols}
Let \(R_p=\lambda+|p|\), \(0\le m\le M\lambda\), and
\(\omega=|\alpha|+2j\). The finite sum of source-loop integrands obeys,
for all fixed multi-indices, on its smooth support,
\[
 |\partial_p^\beta\partial_K^\alpha
       \partial_{\mathfrak m}^j f_{a,I}|+
 |\partial_p^\beta\partial_K^\alpha
       \partial_{\mathfrak m}^j f_{0,I}|
 \le C\chi^{-1}\mathfrak n_I R_p^{-1-|\beta|-\omega}.
\]
On a common interior plateau their difference is at most
\[
 C\chi^{-1}\mathfrak n_I a^3R_p^{2-|\beta|-\omega}.
\]
When \(I=\upsilon\phi c\), with no metric mark, this improves to
\(C\chi^{-1}\mathfrak n_I a^4R_p^{3-|\beta|-\omega}\).
That unmarked coefficient is odd under simultaneous reversal of its
external momenta at fixed \(\mathfrak m\).
\end{lemma}
\begin{proof}
The cycle contains one order-one scalar source, one order-two ghost-action
vertex, one order-two scalar--metric vertex, and three order-minus-two
propagators. Exactly one metric propagator contributes \(\chi^{-1}\).
This gives order minus one. The scalar propagator derivative
\eqref{eq:v28-mass-derivative} lowers order by two. The scalar vertices
are affine in \(\mathfrak m\), and differentiating their mass term
also lowers the bound by two. Fixed momentum derivatives lower order by
one. V14's and V27's periodic upper estimates apply to the remaining
factors. The source output and ghost factors keep every source derivative
away from a principal momentum face.

A metric-background derivative inserts an order-two vertex together with
its extra order-minus-two inverse, or differentiates an existing vertex.
It does not increase the net degree. At most two such marks are retained,
so there are finitely many explicitly ordered words. Their interior
comparison uses V26's \(a^3R^3\) relative vertex estimate and V27's
stronger scalar estimate. At zero metric background no gravitational
interaction vertex is consumed; the ghost-action vertex is the exact
filtered ordinary \(R_1\) vertex and the scalar/free discrepancies
start at relative order \(a^4R^4\). This proves the improvement.

For the unmarked coefficient both bosonic inverses and the ghost inverse
are even under full reversal. The scalar source has one derivative, the
ghost-action vertex has two, and the scalar--metric vertex has even parity
(the mass term is included). The profiles are even. Reversing both the
loop and external momenta changes the sign of the integrand. Paired odd-grid
summation proves the assertion, including its mass derivatives.
\end{proof}

Set \(\epsilon=|K|+m\). Below \(\mathfrak n_I\) denotes the product of
all fixed source/field polarization norms, and all errors have their common
one-loop factor \(\hbar\) removed. Use the common domain
\begin{equation}
 \boxed{\lambda\ge8192\mu>0,\qquad a\lambda\le2^{-20},\qquad
 \lambda L_\nu\ge1,\qquad 0\le m\le M\lambda,}
 \label{eq:v28-domain}
\end{equation}
with every external momentum at most \(\mu\). The constants can depend
on \(M\), fixed profile seminorms and fixed differentiation order.

\begin{theorem}[Cutoff convergence of the scalar transformation source]
\label{thm:v28-cutoff}
For \(I=\upsilon q^b\phi c\), \(0\le b\le2\), and
\(\omega=|\alpha|+2j\le4\),
\begin{equation}
 \boxed{\left|\partial_K^\alpha\partial_{\mathfrak m}^j
 (\mathcal Z^{H,R}_{a,L,I}-\mathcal Z^{H,R}_{0,L,I})\right|
 \le C\chi^{-1}\mathfrak n_I a\epsilon^{4-\omega}.}
 \label{eq:v28-general-source-rate}
\end{equation}
With no metric mark, oddness improves the estimate to
\begin{equation}
 \boxed{\left|\partial_K^\alpha\partial_{\mathfrak m}^j
 (\mathcal Z^{H,R}_{a,L}-\mathcal Z^{H,R}_{0,L})\right|
 \le C\chi^{-1}\mathfrak n_I a^2\epsilon^{5-\omega},
 \qquad\omega\le5.}
 \label{eq:v28-source-rate}
\end{equation}
Higher fixed derivatives have the corresponding bounds with
\(\lambda^{4-\omega}\) and \(\lambda^{5-\omega}\), respectively.
All completed modes are included; zero and lower Wilsonian modes remain
retained. The limiting integrals are the weighted-subtracted ordinary
ones with the actual mass.
\end{theorem}
\begin{proof}
For the general panel, four external weighted derivatives yield a
continuum majorant \(C\chi^{-1}R_p^{-5}\) and an interior mismatch
\(C\chi^{-1}a^3R_p^{-2}\). Physical four-dimensional summation through
\(c/a\) and continuum integration beyond it give, respectively,
\[
 a^3\int_\lambda^{c/a}R\,dR=O(a),\qquad
 \int_{c/a}^{\infty}R^{-2}\,dR=O(a).
\]
On the upper transition region, rescaling by \(a\) gives the same bound.
There is no missing edge strip because the completed symbols and supported
source vertices are smooth periodic. The finite-torus sums obey the same
estimates for \(\lambda L_\nu\ge1\), by covering momentum space by
physical lattice cells and using a fixed number of differentiated symbol
bounds. Constants do not grow with the total volume.

Apply the weighted Taylor remainder of \cref{lem:v28-T}. The parameter
path \((tK,t^2\mathfrak m)\) stays in the massive hard domain. Its
coefficients supply \(\epsilon^{4-\omega}\). For the unmarked source,
odd momentum parity and evenness in the auxiliary mass variable imply
\(\mathbf T_4\mathcal Z^H=\mathbf T_3\mathcal Z^H\).
Five weighted derivatives now give \(R_p^{-6}\) and \(a^4R_p^{-2}\).
Their two integrals are \(O(a^2)\), proving
\eqref{eq:v28-source-rate}. Repeating the differentiated estimates proves
the stated higher-order bounds without a small-momentum inverse.
\end{proof}

The same mass differentiation of V27's explicit determinant words gives
\begin{equation}
 \left|\partial_K^\alpha\partial_{\mathfrak m}^j
 (\gamma^{R}_{a,L,I}-\gamma^{R}_{0,L,I})\right|
 \le C\chi^{-1}\mathfrak n_I a\epsilon^{5-\omega},\qquad\omega\le5,
 \label{eq:v28-action-rate}
\end{equation}
for \(q^b\phi^2\), \(b\le3\), in the \emph{weighted} assignment.
Indeed V27's explicit scalar denominators and mass vertices have precisely
the derivatives in \eqref{eq:v28-mass-derivative}; five weighted
derivatives give its same \(R_p^{-5},a^3R_p^{-2}\) majorants.
Thus \eqref{eq:v28-action-rate} is an extension to the stated mass jets,
not an assumption that the previous fixed-mass projector already intertwined.
At zero mass it reduces to the earlier momentum assignment.

\section{The six terms of the mixed one-ghost row}
\label{sec:v28-row}

At zero ordinary background let
\(\Pi_a=D_\phi^2\Gamma_{a,1}^{\rm mix}\),
\(\mathcal T_a=D_qD_\phi^2\Gamma_{a,1}^{\rm mix}\), and
\(V_a=D_qD_\phi^2\widetilde S_{H,a}\).
Let \(Z_a^g(c)\) be the inherited quantum metric transformation for the
V18 source, in the same reference convention, and let
\(Z_a^H(\phi,c)=\mathcal Z_a^H(0;\phi,c)\) be
\eqref{eq:v28-source-loop}. Scalar flavour indices are retained in the
pairings. Define
\begin{equation}
 \boxed{\begin{aligned}
 \mathfrak L_a(c;\phi_1,\phi_2)={}&
 \mathcal T_a(R_0c;\phi_1,\phi_2)\\
 &+\Pi_a(\mathcal L_c\phi_1,\phi_2)
       +\Pi_a(\phi_1,\mathcal L_c\phi_2)\\
 &+V_a(Z_a^gc;\phi_1,\phi_2)\\
 &+\langle\widetilde P_a\phi_1,Z_a^H(\phi_2,c)\rangle
       +\langle\widetilde P_a\phi_2,Z_a^H(\phi_1,c)\rangle.
 \end{aligned}}
 \label{eq:v28-six-row}
\end{equation}
There are six summands after the two polarizations are counted. In
particular the fourth term is the \emph{tree matter stress contracted
with the quantum metric transformation}. It would be missing from a Ward
test using only the mixed determinant and a scalar antifield source.

No \(\phi\)-dependent gauge fixing is introduced. The off-shell
nonminimal fields are restored before this extraction. The pure-gravity
Euler derivative vanishes at the flat background, but its matter stress
at scalar degree two does not. The quantum metric source at \(q=\phi=0\)
is unchanged by the scalar: its two-hard-leg source Hessian has no scalar
block, and the classical hard inverse is block diagonal there.

\begin{theorem}[Finite measured identity for the mixed row]
\label{thm:v28-finite-row}
Use the full graded field list \((q,\phi,c,\bar c,b)\), the actual
zero-antifield action \(S_a\), and the source
\((r_a^q,r_a^H,r_a^c,r_a^{\bar c},r_a^b)\). At the coefficient
\(c\phi^2\), finite Gaussian/Berezin integration by parts gives
\begin{equation}
 \boxed{\mathfrak L_a=\mathcal I_a^{H}+\mathcal X_a^{H}.}
 \label{eq:v28-finite-identity}
\end{equation}
Here \(\mathcal I_a^{H}\) is the one-loop insertion of the actual
classical breaking \(DS_a[r_a]\), and \(\mathcal X_a^{H}\) is its
reference trace. Neither is set to zero at finite mesh. No independent
auxiliary-density or coefficient-divergence term survives in this particular
slot.
\end{theorem}
\begin{proof}
The finite coefficient identity can be derived without nilpotence. Let
\(\mathbb C(\Phi)=(\mathbb H(\Phi)+\mathbb R)^{-1}\), with the
reference matrix temporarily invertible on its hard support. Differentiating
\(B=DS[r]\) twice and contracting gives, in the graded convention,
\[
 \tfrac12\operatorname{STr}(\mathbb C D^2B)
 =D\Gamma_1[r]+DS\!\bigl[\tfrac12\operatorname{STr}
       (\mathbb C D^2r)\bigr]
       +\operatorname{sdiv}r
       -\operatorname{STr}(\mathbb C\mathbb R Dr).
\]
The two middle Hessian terms in \(D^2B\) combine into
\(\operatorname{STr}(\mathbb C\mathbb H Dr)\);
\(\mathbb C\mathbb H=I-\mathbb C\mathbb R\) gives the displayed
formula. It is equivalently the Gaussian/Berezin integration-by-parts
identity. Background expansions are finite at the requested order, so
zero covariance subspaces are handled by support restriction, not by
inverting their zeros. Hard-linear branches have already been excluded
by the external support argument. This identifies the insertion with
the actual measured one-loop coefficient.

In the extraction, \(D\Gamma_1[r]\) gives the first three summands
of \eqref{eq:v28-six-row}; \(DS[z_1]\) gives the matter-stress term
and the two scalar-Euler terms. The scalar divergence is zero by
\cref{prop:v28-divergence}; every other flat source divergence is independent
of \(\phi\). The once-counted auxiliary logarithmic density is a function
of \(q\) alone, and \(r_a^q\) has no scalar dependence. Its direct
variation therefore has no \(c\phi^2\) coefficient. This last fact does
not discard the density in any other row. The retained nonminimal contact
implements the existing \(\phi\)-independent gauge and adds no separate
term in this slot. The identity now has precisely the two right-hand
operands in \eqref{eq:v28-finite-identity}.
\end{proof}

\subsection{The explicit lower-reference operand}

Write \(\mathbb C_0=\nu_\lambda\mathbb H_0^{-1}\),
\(V=\mathbb H(\Phi)-\mathbb H_0\), and
\(\Theta=1-\nu_\lambda\), acting on every primitive species.
The inverse-reference cancellation is
\(\mathbb C_0\mathbb R=\Theta\), including the massive scalar block.
It gives the finite ordered expression
\begin{equation}
 \boxed{\mathcal X_a^{H}
 =[c\phi^2]\operatorname{STr}
 \left[\sum_{j=0}^{3}(-\mathbb C_0V)^j\Theta\,Dr_a\right].}
 \label{eq:v28-reference}
\end{equation}
The \(j=0\) term has no \(\phi^2\) coefficient. Every remaining
term contains both a hard covariance and \(\Theta\). All routed
momenta therefore lie in a fixed lower annulus, which can be taken to be
\(\lambda/2\le|p|\le3\lambda\) on
\eqref{eq:v28-domain}. The matrix order in \eqref{eq:v28-reference}
matters; \(\Theta\) is not commuted through a source vertex.

At this background order the action vertices are the scalar mixed and
contact Hessians and the ghost-action vertex with one retained ghost.
There is no pure-gravity interaction with a retained metric mark.
Their interior discrepancy, and the free covariance discrepancy, start
at \(a^4\). Source vertices are exactly ordinary on this annulus.
Mass derivatives act on their actual scalar factors as above. Thus
\begin{equation}
 \left|\partial_K^\alpha\partial_{\mathfrak m}^j
 (\mathcal X_a^{H,R}-\mathcal X_0^{H,R})\right|
 \le C\chi^{-1}\mathfrak n_I a^4\lambda^3
                        \epsilon^{6-\omega},\qquad\omega\le6,
 \label{eq:v28-reference-rate}
\end{equation}
where \(\mathcal X^{H,R}=(I-\mathbf T_5)\mathcal X^H\).
To verify the powers, the unprojected row has ultraviolet degree five:
its integrand has order one. Six weighted derivatives on the fixed
annulus yield \(\lambda^{-5}\), the relative error contributes
\(a^4\lambda^4\), and physical integration contributes
\(\lambda^4\). The result is \(a^4\lambda^3\), followed by the
weighted Taylor factor. Direct product differentiation of the finitely
many words proves the same estimate at all fixed derivative orders.
This estimate is for the actual lower reference, not a ultraviolet bound
obtained by assuming a Ward identity.

\section{A common massive assignment and the nonlocal modified identity}
\label{sec:v28-compatibility}

The ordinary scalar Euler multiplier \(P_0(k)=k^2+\mathfrak m\) and
the ordinary matter-stress tensor \(V_0\) have weight two. The ordinary
scalar and metric transformation generators have weight one. Consequently
\cref{lem:v28-T} proves, at the row under consideration,
\begin{equation}
 \boxed{\mathbf T_5\mathfrak L_0
 =\mathfrak L_0\bigl(\mathbf T_4\mathcal T,
 \mathbf T_4\Pi,\mathbf T_3 Z^g,\mathbf T_3 Z^H\bigr).}
 \label{eq:v28-intertwining}
\end{equation}
This is an explicit bounded intertwining assignment for the displayed
row. It is not a splitting theorem for the full matter BV complex.

The same projected identity holds with the finite action on the left
and its ordinary multipliers on the right. Indeed,
\(\widetilde P_a-P_0\) and \(V_a-V_0\) have weighted Taylor order
at least six on soft legs. They cannot contribute to \(\mathbf T_5\)
when multiplied by a regular source coefficient. Put
\(E_H=\widetilde P_a-P_0\) and \(E_V=V_a-V_0\). After subtracting the
common local assignment, the finite remainder still contains
\begin{equation}
 \begin{split}
 \mathfrak F_a={}&E_V(\mathbf T_3 Z_a^gc;\phi_1,\phi_2)\\
 &+\langle E_H\phi_1,(\mathbf T_3Z_a^H)(\phi_2,c)\rangle
 +\langle E_H\phi_2,(\mathbf T_3Z_a^H)(\phi_1,c)\rangle.
 \end{split}
 \label{eq:v28-finite-local-product}
\end{equation}
This is a product of an omitted classical artifact with a retained local
quantum coefficient, not a term that can be dropped because the artifact
is small. Both unmarked source coefficients are odd. Their weight-one
local coefficients are at most \(C\chi^{-1}a^{-2}\), and their
weight-three coefficients at most
\(C\chi^{-1}[1+\log(1/(a\lambda))]\), by their order-minus-one
integrands. Since \(E_H,E_V=O(a^4\epsilon^6)\),
\begin{equation}
 |\partial_K^\alpha\mathfrak F_a|
 \le C\chi^{-1}\mathfrak n_I a^2(\mu+m)^{7-|\alpha|}
 \label{eq:v28-finite-local-bound}
\end{equation}
on fixed external windows through order seven, with the usual scaled
higher-derivative interpretation. The fixed \(M\) domain bounds the
possible higher \(a\epsilon\) factors. The weight-three contribution
is no larger because \((a\lambda)^2\log(1/(a\lambda))\) is bounded.

\begin{theorem}[Nonlocal modified identity for one ghost and two Higgs fields]
\label{thm:v28-nonlocal-row}
Let a superscript \(R\) in \eqref{eq:v28-six-row} mean the common
weighted complement of each quantum coefficient; retain the actual tree
multipliers. In thermodynamic normalization at the fixed positive scale
\(\lambda\),
\begin{equation}
 \boxed{\mathfrak L_0^R=\mathcal X_0^{H,R}.}
 \label{eq:v28-continuum-row}
\end{equation}
For the finite diagnostic
\(\mathfrak W_{a,L}^R=\mathfrak L_{a,L}^R-\mathcal X_{a,L}^{H,R}\),
the mesh comparison at fixed volume obeys
\begin{equation}
 \boxed{|\partial_K^\alpha\partial_{\mathfrak m}^j
 (\mathfrak W_{a,L}^R-\mathfrak W_{0,L}^R)|
 \le C\chi^{-1}\mathfrak n_I a\epsilon^{6-\omega},
 \qquad\omega\le6.}
 \label{eq:v28-row-rate}
\end{equation}
The finite-volume comparison with \eqref{eq:v28-continuum-row} is
separate; in scaled external-window norms it is bounded by
\(C_b\chi^{-1}\lambda^5(\lambda L_{\min})^{-b}\), \(b>4\).
No exact normalized off-grid finite-torus identity is inferred.
\end{theorem}
\begin{proof}
The ordinary source extends the scalar action by its actual Lie variation;
the ordinary gravity, ghost and nonminimal sectors have the inherited
classical identity. Thus \(DS_0[r_0]=0\) before loop integration.
Apply the finite trace identity used in \cref{thm:v28-finite-row} with
a common smooth upper regulator on the ordinary loop. It has a lower
reference term and an upper-reference term. No no-covariance trace
contributes at \(c\phi^2\). Every upper-reference word consequently
lies on its upper transition annulus. After \(I-\mathbf T_5\), six
weighted derivatives give order minus five. Integration on an annulus
of radius \(\Lambda\) is bounded by
\(C\chi^{-1}\epsilon^6\Lambda^{-1}\), which vanishes as
\(\Lambda\to\infty\). All differentiated subtracted integrals are
absolutely convergent. Their loop-momentum changes of variables have
no boundary term in that limit. The common weighted assignment commutes
with the ordinary row by \eqref{eq:v28-intertwining}. This proves
\eqref{eq:v28-continuum-row}; no finite-cutoff invariance was assumed to
prove a cutoff estimate.

For the finite comparison, the first term in the row uses
\eqref{eq:v28-action-rate} with one metric mark and acquires one
external derivative from \(R_0c\), giving \(a\epsilon^6\).
The two self-energy terms have the same or smaller bound. The stress
term uses the inherited unmarked metric-source bound, at worst
\(C\chi^{-1}a|K|^4\), times a weight-two polynomial. The two scalar
Euler terms use \eqref{eq:v28-source-rate} times \(P_0\), giving
\(C\chi^{-1}a^2\epsilon^7\). Replacing the ordinary tree multipliers
by the finite ones gives still smaller soft-window terms by their
\(a^4\) comparison and the bounded subtracted limiting source
coefficients. All these terms are bounded by the right side of
\eqref{eq:v28-row-rate} on \eqref{eq:v28-domain}. The reference
comparison \eqref{eq:v28-reference-rate} is
\(a(a\lambda)^3\epsilon^6\) and fits the same bound.
Leibniz differentiation proves the stated fixed orders. Beyond weighted
order six the scaled bound is \(C\chi^{-1}a\lambda^{6-\omega}\).

The continuum subtracted words and the lower-reference words have
integrable loop-momentum derivatives of every fixed order, because the
lower covariance vanishes near zero. Poisson summation gives the displayed
separate volume estimate. It need not preserve the precise small-\(K\)
zero of the thermodynamic Ward identity for an arbitrary chosen off-grid
interpolation, so no such extra finite-volume factor is asserted.
\end{proof}

Equation \eqref{eq:v28-finite-identity} also gives the exact normalized
balance
\begin{equation}
 \boxed{\mathfrak W_a^R=(I-\mathbf T_5)\mathcal I_a^H-\mathfrak F_a.}
 \label{eq:v28-breaking-limit}
\end{equation}
Thus the subtracted nonlocal breaking insertion tends to zero in the
thermodynamic or declared joint limit, \emph{after} the signed coefficient
comparison. The local equation is not thereby solved:
\(\mathbf T_5\mathcal I_a^H\), the local reference trace, and the
permitted quantum action/source counterterms still have to be matched in
the combined scalar--gravity complex.

\section{Source topology, shell ownership, and what the local problem still requires}
\label{sec:v28-topology}

For the two independent momenta of \(\upsilon\phi c\) or
\(c\phi^2\), a fixed smooth external window produces a kernel on two
relative four-dimensional coordinates. Fixed differentiated symbol bounds,
rescaling by \(\mu\), and Fourier integration by parts give
\begin{align}
 a^8\sum_{x,y}[1+\mu(d_{\rm per}(x,0)+d_{\rm per}(y,0))]^b
          \|\mathscr K_{Z^H}^{\Delta}(x,y)\|
 &\le C_b\chi^{-1}a^2(\mu+m)^5,\\
 a^8\sum_{x,y}[1+\mu(d_{\rm per}(x,0)+d_{\rm per}(y,0))]^b
          \|\mathscr K_{\mathfrak W}^{\Delta}(x,y)\|
 &\le C_b\chi^{-1}a(\mu+m)^6.
 \label{eq:v28-kernels}
\end{align}
The second inequality is the same-volume mesh difference from
\eqref{eq:v28-row-rate}; the separate volume comparison is added when
comparing directly with the zero thermodynamic residual. For
\(\upsilon q^b\phi c\), \(b=1,2\), the analogous bound is
\(C\chi^{-1}a(\mu+m)^4\) in its \(b+2\) relative coordinates.
To verify these inequalities, take more momentum derivatives than the
relative dimension plus the requested moment, use a compact smooth window,
and sum the rapidly decaying Fourier majorant with the physical powers of
\(a\). Periodization can only decrease the distance in each lift.
There is no total-volume multiplier. The resulting multilinear form is
bounded with one argument in \(L^1\) and all others in \(L^\infty\).

Every primitive covariance in \eqref{eq:v28-source-loop} has its own
shell decomposition. Expanding \(G_c\) and all background inverses first,
then assigning a word to its largest primitive shell index, includes all
metric--ghost--scalar mixed levels exactly once. The subtracted general
source-shell increments have bound
\(C\chi^{-1}(\mu+m)^4\Lambda^{-1}\); in the unmarked odd sector
this improves to \(C\chi^{-1}(\mu+m)^5\Lambda^{-2}\).
These bounds follow from the same differentiated integrands in the proof
of \cref{thm:v28-cutoff} restricted to an annulus of radius \(\Lambda\).
They are summable for fixed lower \(\lambda\). Keeping only equal
metric, ghost and scalar shell labels would not give this measured trace.

The common Taylor projection supplies \emph{compatibility of a row}, not
existence of every desired finite normalization. In particular the scalar
mass, wave-function, mixed geometric and scalar-antifield finite terms
must solve one local equation. A schematic remaining target is
\[
 [\,(S_{\rm grav+H,0},C_{\rm loc}^{(1)})\,]_{c\phi^2}
 =-[\,\mathfrak L_{\rm loops}-\mathcal X^H\,]_{c\phi^2,\loc},
\]
with the inherited pure-gravity normalization included, not reset.
The enlarged equation has scalar-antifield and scalar-Euler components
that are not present in the pure-gravity classification used in V25.
This installment does not invoke that classification to solve it.

A comparison of two admissible smooth completions, with the same interior
and scalar source and their own actual weighted local jets assigned, obeys
the same nonlocal estimates: the difference of their loop integrands has
only the controlled upper-region and native-artifact contributions.
Independent finite changes of the scalar transformation source would require
matching its own contacts and breaking/reference insertions as well.
Neither an infrared limit nor a full interacting shell self-map follows
from these fixed-order statements.

\section{Final ledgers for V28}
\label{sec:v28-ledgers}

\subsection*{Proved}
\addcontentsline{toc}{subsection}{V28: Proved}
The scalar source \eqref{eq:v28-scalar-source} has the explicit vertex,
finite divergence and recorded source square. The simultaneous law gives
the complete \(\upsilon q^b\phi c\), \(b\le2\), one-loop coefficient
\eqref{eq:v28-source-loop}, with its three species inverses and correct
Berezin/determinant sign. Its mass-graded source coefficients have the
full-cutoff and rooted bounds above. At zero metric background their odd
parity gives the stronger \(a^2\) comparison. The actual massive Euler
factor exposes the failure of a naive momentum-only Taylor interface;
the weighted assignment resolves this precisely without changing the
massive propagator. The one-ghost/two-Higgs row has all six canonical
terms, its finite breaking/reference identity, and a nonlocal modified
continuum identity with the stated mesh and volume qualifications.

\subsection*{Inherited}
\addcontentsline{toc}{subsection}{V28: Inherited}
V1--V27 and their labels are unchanged. In particular V27 supplies the
periodically completed, temporally continued scalar action, the simultaneous
mixed determinant and its native symbol bounds. V14--V26 supply the metric
contour, ghost action, source chart, reference conventions and zero-background
metric-antifield estimates. The current source uses V18's actual filtered
source and does not restore V15's older prescription. General Gaussian,
Berezin and stationary pushforward identities are inherited tools. The
ordinary scalar density action and Lie transformation are read from the
native zero-mesh action; no unrelated gravitational model is substituted.

\subsection*{Literature input}
\addcontentsline{toc}{subsection}{V28: Literature input}
Modified Slavnov identities at fixed Wilsonian scale and their BV relation
are established methods \cite{V28IIM}. They provide context for the lower
reference term. No external gravity--Higgs beta function, anomaly theorem,
or matching coefficient supplies the source triangle or its cutoff bounds.
The weighted Taylor algebra is proved here directly. There is no claim of
a new general construction of perturbative gravity with scalar matter.

\subsection*{Conditional}
\addcontentsline{toc}{subsection}{V28: Conditional}
Full local quantum symmetry still requires the combined scalar--gravity
counterterm primitive and compatible physical normalization conditions.
The component Taylor assignment and the nonlocal modified row do not
supply that primitive. All analytic conclusions consume the earlier
contour, periodic-symbol and source estimates at their stated scopes;
this installment is not an independent audit of all 27 preceding versions.
A complete ESM theorem additionally requires the internal gauge and chiral
sectors, their common measure and source transport, and higher-loop bounds.

\subsection*{Open}
\addcontentsline{toc}{subsection}{V28: Open}
The next smallest target is the mass-graded \emph{local} one-ghost/two-Higgs
normalization equation, including its scalar-antifield terms and the
already chosen pure-gravity reference normalization. Its embedding in a
finite combined field/antifield bank should retain the scalar mass as a
weighted parameter rather than applying a pure-gravity primitive unchanged.
Higher antifield/matter source squares, nonzero internal couplings, the
physical chiral line, infrared removal and an invariant interacting EFT
class remain open. No infinite-series convergence, asymptotic safety,
finite-parameter ultraviolet gravity or nonperturbative continuum is claimed.

\begin{center}\small
\begin{tabular}{>{\raggedright\arraybackslash}p{.24\textwidth}>{\raggedright\arraybackslash}p{.67\textwidth}}
\toprule
Variable & Scope in V28\\\midrule
Cutoff and volume & Smooth completed full zone, odd grids and
\eqref{eq:v28-domain}. Same-volume mesh and rooted bounds are uniform;
thermodynamic Ward normalization and its volume error are separate.\\
Sources and mass & \(\upsilon q^b\phi c\), \(b\le2\), and the
\(c\phi^2\) row. Fixed external derivative/moment orders; constant mass
parameter \(0\le m\le M\lambda\). The mass is not sent to zero in
propagators to define the local assignment.\\
Couplings and loops & One loop, fixed \(\chi>0\); zero internal gauge,
Yukawa and Higgs self-couplings. Four scalar components, no extra closed
species factor in the displayed external-pair source.\\
Measure & Same metric contour and coefficient/Haar/Jacobian convention.
The scalar-source divergence is zero, but other measured rows retain their
actual density contributions.\\
Iteration & All primitive mixed-shell indices are retained; subtracted
increments are summable at fixed lower scale. This is not analytic
invariance of arbitrary interacting inputs.\\
Symmetry & Exact finite change-of-variables identity with breaking retained;
ordinary nonlocal modified row after cutoff removal. No all-frequency
nilpotence or complete coupled local BV primitive.\\
\bottomrule
\end{tabular}
\end{center}

The accompanying checker tests the scalar-source sign by finite Berezin
algebra, a nonzero ordinary source integrand, the joint determinant and its
background derivatives, mass-weighted Taylor identities and their
momentum-only counterexample, the single-cutoff-factor mass derivatives,
finite divergence, and all three-species shell labels. Its certificate
separates exact finite checks from the analytic proofs and inherited inputs.
No numerical native loop integral or new Lean formalization is claimed.

\begingroup
\renewcommand{\refname}{Additional comparison reference for V28}
\begin{thebibliography}{99}
\bibitem[46]{V28IIM}
Y.~Igarashi, K.~Itoh and T.~R.~Morris,
\emph{BRST in the Exact RG}, Progress of Theoretical and Experimental
Physics \textbf{2019} (2019), 103B01, arXiv:1904.08231.
Comparison for modified reference identities, not an imported native
mixed coefficient or local coupled anomaly cancellation.
\end{thebibliography}
\endgroup

\clearpage
\part*{V29: A mass-regular relative primitive for the mixed scalar--gravity row}
\addcontentsline{toc}{part}{V29: Relative scalar--gravity local normalization}

\section{The relative local equation, rather than a new matter anomaly assumption}
\label{sec:v29-start}

Sections 1--223 and their labels are unchanged. V28 leaves the local
\(c\phi^2\) equation open after constructing its nonlocal coefficient and
mass-weighted Taylor assignment. The task here is precisely that local
row. We retain the four real equal-mass scalars, the coefficient contour,
the smooth periodic completion, the V18 gravitational source and V28 scalar
source, and zero internal gauge, Yukawa and scalar self-couplings. The
one-loop factor \(\hbar\) is omitted from coefficient formulas. In
particular this is not a claim about the chiral Standard Model.

The new method does not import the pure-gravity cohomology theorem into a
matter complex. It solves a finite polynomial row explicitly. Its scalar
Euler polynomial is divided monically in the \emph{local mass parameter};
the remaining antisymmetric ghost-momentum jet is controlled by rotations
which fix a scalar momentum. A Koszul homotopy and a symmetric-divergence
homotopy supply the actual counterterm. The full rotation-cocycle test is
\emph{not} imposed: it would incorrectly reject some perfectly valid
scalar-antifield variations. This distinction is checked below by an exact
counterexample.

\subsection{The fixed gravity normalization is part of the right-hand side}

Let \(H_g\) denote the one-loop local pure-gravity normalization already
chosen in V24--V26, in the paired coordinate-density convention. Its source
part satisfies the ordinary local reference identity at the required
arities. The existence and finite coefficient prescription for its critical
part are inherited from V25, with exactly the qualifications stated there;
we do not independently certify that external input here. In this section
\(s\) is the ordinary, zero-mesh \emph{coupled} BV differential, not the
filtered finite source square.

Write \(z=m^2\). The ordinary scalar density is the zero-mesh coefficient
of the same action used in V27,
\[
 S_H=\frac12\sum_{A=1}^{4}\int \sqrt{\det g(q)}
       \bigl(g^{\mu\nu}(q)\partial_\mu\phi_A\partial_\nu\phi_A
                  +z\phi_A^2\bigr),\qquad s\phi_A=c^\mu\partial_\mu\phi_A.
\]
Its metric is the existing coframe-coordinate map \(g(q)\), with
\(Dg(0)=I\). No scalar half-density is substituted. Define the scalar-degree-two
relative reference row
\begin{equation}
 \boxed{B=\mathbf T_5[\mathcal X_{0,\lambda}-sH_g]_{c\phi^2}.}
 \label{eq:v29-relative-B}
\end{equation}
Here \(\mathcal X_{0,\lambda}\) is the ordinary coupled lower-reference
trace; its displayed component is exactly \eqref{eq:v28-reference} at
zero mesh and in thermodynamic normalization. In this slot
\([sH_g]_{c\phi^2}\) is the tree scalar stress contracted with the already
chosen local metric transformation. It is not zero just because \(H_g\)
has no scalar fields. Subtracting it is what makes the construction
\emph{relative} to the previous gravity normalization.

All formulas below act on the actual coefficient array of
\eqref{eq:v29-relative-B}. That array remains specified by finite source
derivatives of the native annular trace. No numerical loop integral is
reported. The polynomial maps neither fit a new reference coefficient nor
set an unevaluated coefficient to zero.

\subsection{A small coupled consistency envelope}

For clarity, the ordinary consistency identity used below has enough
fields to include its Euler terms. Give \(q,\phi\) weight zero,
\(c\) weight \(-1\), \(\eta=q^*,\upsilon=\phi^*\) weight two,
\(\sigma=c^*\) weight three, each derivative weight one and \(z\) weight
two. Every term of the ordinary minimal master action has weight two.
The scalar mass term and the scalar-antifield term have respectively
\(2+0+0=2\) and \(2-1+1=2\). The antibracket removes weight two.
The action has no term of field/ghost/antifield arity below two.
Consequently the extension of V24's arity-five, weight-four local
projection is a complex for the \emph{ordinary coupled} differential.
For a component of ghost number \(g\) and \(n_*\) antifields its allowed
momentum--mass Taylor degree is \(4+g-n_*\), just as required by V28's
scalar Euler multiplier. The proof is the weight and arity addition in
each bracket, including \(s\eta\)'s scalar stress and \(s\upsilon\)'s
massive Euler term. Neither term is deleted. This observation licenses
the finite consistency extraction; it is not a new all-source cutoff
estimate or an assertion that a single matter degree is a subcomplex.

\section{The complete polynomial normalization map at two scalar legs}
\label{sec:v29-row-map}

It suffices to keep a fixed external Higgs component: the equal-mass,
zero-internal-coupling action and all stated sources preserve the common
orthogonal flavour action, so the two-scalar row is proportional to
\(\delta_{AB}\). There is no extra factor four for a specified external
pair. Let all three momenta be incoming, and write
\[
 k+r+t=0,\qquad r=p-k/2,\qquad t=-p-k/2,\qquad
 P(p,z)=p^2+z.
\]
The scalar exchange is \(p\mapsto-p\). Strip the common Fourier factor
\(i\) from the row and write its density as
\(\frac{i}{2}\int c^\nu(k)\phi(r)\phi(t)B_\nu(k,p,z)\).
The polynomial vector \(B\) is even in \(p\) and has weighted degree at
most five. Complex coefficients are permitted during Fourier calculation;
real densities are recovered by the usual conjugation involution.

The three permitted local counterterm pieces have Fourier kernels
\[
 C_f=\frac12\int\phi(-p)f(p,z)\phi(p),\qquad
 C_T=\frac14\int q_{\mu\nu}(k)\phi(r)\phi(t)T_{\mu\nu}(k,p,z),
\]
\[
 C_U=\int\upsilon(-k-t)c^\nu(k)\,iU_\nu(k,t,z)\phi(t).
\]
Here \(f\) is even and has weight at most four, \(T=T^{\mathsf T}\) is
exchange even with weight at most four, and \(U\) has weight at most
three. The inherited antibracket convention is used so that the Euler
contribution of \(C_U\) has the sign in the following formula. Reversing
that convention simply reverses the displayed source kernel.

\begin{proposition}[The exact action/source row map]
\label{prop:v29-map}
The coefficient of \(c\phi^2\) in \(s(C_f+C_T+C_U)\), at zero metric
and zero antifields, is
\begin{equation}
 \boxed{\begin{aligned}
 \mathcal L(T,f,U)_\nu={}&k^\mu T_{\mu\nu}
                 +r_\nu f(t,z)+t_\nu f(r,z)\\
       &+P(r,z)U_\nu(k,t,z)+P(t,z)U_\nu(k,r,z).
 \end{aligned}}
 \label{eq:v29-map}
\end{equation}
No equation of motion has been imposed in this identity.
\end{proposition}
\begin{proof}
The free metric variation is
\(i(k_\mu c_\nu+k_\nu c_\mu)\); symmetry of \(T\) and the factor
\(1/4\) give its first term with the stated \(i/2\) row convention.
Vary each scalar of \(C_f\) by \(c\cdot\partial\phi\) and exchange the
two scalar labels to obtain the next two terms. Replace the scalar
antifield in \(C_U\) by the \emph{actual ordinary} Euler multiplier
\(P\phi\), then polarize, to obtain the last two terms. At this
background the other variations of \(C_U\) retain an antifield and
therefore belong to other rows, not to \eqref{eq:v29-map}.
\end{proof}

Define the two low ghost-momentum jets
\begin{equation}
 b_\nu(p,z)=B_\nu(0,p,z),\qquad
 A_{\mu\nu}(p,z)=
 [\partial_{k_\mu}B_\nu-\partial_{k_\nu}B_\mu]_{k=0}.
 \label{eq:v29-jets}
\end{equation}
Evaluation of a \emph{polynomial} at \(z=-p^2\) is denoted by a bar.
It is the quotient by the ideal \((P)\). It does not evaluate a propagator
at negative squared mass, alter a contour, or replace a massive covariance
by a massless one. Monic division in \(z\) is valid including at \(m=0\).

\begin{lemma}[Two necessary polynomial tests]
\label{lem:v29-tests}
Every row in the image of \eqref{eq:v29-map} satisfies
\begin{equation}
 \boxed{\bar b=0,\qquad
 p_\alpha\bar A_{\mu\nu}-p_\mu\bar A_{\alpha\nu}
                   +p_\nu\bar A_{\alpha\mu}=0.}
 \label{eq:v29-tests}
\end{equation}
Conversely, these tests are sufficient. They involve finite coefficient
identities; they are not a choice to discard a failed component.
\end{lemma}
\begin{proof}[Proof of necessity]
The metric and kinetic terms vanish at \(k=0\). The scalar-source term
there is \(P(p,z)(U(0,p,z)+U(0,-p,z))\), proving the first test. Split
\(U(0,p,z)\) into its even and odd parts. The even part contributes no
antisymmetric first \(k\) jet. If \(U_o\) is the odd part, its
antisymmetric jet modulo \(P\) is
\(2(p_\mu\bar U_{o,\nu}-p_\nu\bar U_{o,\mu})\).
Terms explicitly linear in \(k\) in \(U\) give a multiple of \(P\).
The kinetic term contributes
\(p_\nu\partial_{p_\mu}\bar f-p_\mu\partial_{p_\nu}\bar f\),
and the metric term contributes zero. All these antisymmetric tensors
are a wedge product with \(p\), so wedging once more with \(p\) gives
zero. Sufficiency is constructive in \cref{thm:v29-primitive}.
\end{proof}

\section{The stabilizer quotient and two pole-free homotopies}
\label{sec:v29-homotopies}

\subsection{Why a full rotation-cocycle test would be too strong}

Let \(L_{ij}=p_j\partial_{p_i}-p_i\partial_{p_j}\). An initial idea
would be to require \(\bar A\) to be a Chevalley--Eilenberg cocycle for
these six rotations on scalar polynomials. That condition is not necessary
for \eqref{eq:v29-map}: it has omitted the scalar-antifield Euler channel.
For example take \(U_1(0,p)=p_0^3\), all other components zero, and
\(B=\mathcal L(0,0,U)\). Then
\[
 (\bar A_{01},\bar A_{02},\bar A_{03},\bar A_{12},\bar A_{13},\bar A_{23})
 =(2p_0^4,0,0,-2p_0^3p_2,-2p_0^3p_3,0).
\]
Its rotation-cochain differential has, among other nonzero entries,
\(-6p_0^3p_2\), although \(B\) is already the variation of the displayed
local source action. It satisfies \(p\wedge\bar A=0\) exactly. Thus no
full-rotation cocycle condition is silently assumed in the proof below.

\subsection{The Koszul homotopy}

 Regard \(\bar A\) as a polynomial two-form on momentum space. On an
exterior \(r\)-form of homogeneous polynomial degree \(n\), define
\[
 d_p\omega=p_i\,e^i\wedge\omega,\qquad
 h_p\omega=\sum_i\partial_{p_i}\,\iota_i\omega.
\]
\begin{lemma}[A finite polynomial section of the stabilizer test]
\label{lem:v29-koszul}
In four dimensions,
\begin{equation}
 h_pd_p+d_ph_p=(4+n-r)I.
 \label{eq:v29-koszul}
\end{equation}
In particular, if a degree-\(n\) polynomial two-form satisfies
\(p\wedge\bar A=0\), then
\begin{equation}
 \boxed{\bar A_{\mu\nu}=p_\mu V_\nu-p_\nu V_\mu,
 \qquad V_\nu=\frac1{n+2}\sum_\mu\partial_{p_\mu}\bar A_{\mu\nu}.}
 \label{eq:v29-V}
\end{equation}
The only denominators for this problem are the nonzero integers
\(n+2=2,4,6\), with \(n=0,2,4\). The degree-zero closed two-form is zero.
\end{lemma}
\begin{proof}
Use \(\iota_i(e^j\wedge\omega)=\delta_i^j\omega-e^j\wedge\iota_i\omega\)
and differentiate \(p_j\). The resulting diagonal terms are the dimension
four, the polynomial Euler operator of eigenvalue \(n\), and minus the
exterior-degree operator of eigenvalue \(r\). The remaining terms cancel
between the two compositions. Set \(r=2\) and \(d_p\bar A=0\) to obtain
\eqref{eq:v29-V}. For \(n=0\), \(h_p\bar A=0\), so the same identity
forces \(\bar A=0\).
\end{proof}

\subsection{The symmetric metric section}

\begin{lemma}[A symmetric divergence section with no soft pole]
\label{lem:v29-symmetric}
Let a vector polynomial \(R_\nu(k)\), with any spectator variables, obey
\(R(0)=0\) and have symmetric first derivative at zero. Decompose it into
homogeneous \(k\)-degrees, \(R=\sum_{n\ge1}R^{(n)}\). Then
\(R_\nu=k^\mu T_{\mu\nu}\) with a symmetric polynomial tensor. Explicitly,
\begin{equation}
 \begin{split}
 T^{(0)}_{\mu\nu}&=\partial_{k_\mu}R^{(1)}_\nu,\\
 T^{(n-1)}_{\mu\nu}
 &=\frac{\partial_{k_\mu}R^{(n)}_\nu+
          \partial_{k_\nu}R^{(n)}_\mu}{n-1}
   -\frac{\partial_{k_\mu}\partial_{k_\nu}
                      (k^\rho R^{(n)}_\rho)}{n(n-1)},\quad n\ge2.
 \end{split}
 \label{eq:v29-symmetric}
\end{equation}
It lowers total momentum--mass degree by one and preserves the parity
of every spectator variable.
\end{lemma}
\begin{proof}
The first line uses the assumed symmetry. For \(n\ge2\), put
\(F=k\cdot R^{(n)}\), of degree \(n+1\). Euler's identity gives
\[
 k^\mu(\partial_\mu R_\nu+\partial_\nu R_\mu)
      =(n-1)R_\nu+\partial_\nu F,\qquad
 k^\mu\partial_\mu\partial_\nu F=n\partial_\nu F.
\]
Substitution proves the assertion. These are polynomial identities also
at \(k=0\); no division by \(k^2\) or \(p^2\) is made.
\end{proof}

\section{A complete finite right inverse for the mixed local row}
\label{sec:v29-primitive}

\begin{theorem}[Mass-regular local action/source primitive]
\label{thm:v29-primitive}
Let \(B\) be exchange even and of weighted degree at most five. It belongs
to the image of \eqref{eq:v29-map} if and only if
\eqref{eq:v29-tests} holds. One can choose \(f=0\), so the particular
right inverse need not change a prescribed scalar kinetic normalization.
The following finite operations construct \(U,T\).

First divide the constant jet monically and set
\[
 U_e(p,z)=\frac{b(p,z)}{2P(p,z)},\qquad
 B^t=B-\mathcal L(0,0,U_e).
\]
The quotient is a polynomial of weight at most two. Next compute \(\bar A\)
from \(B^t\), decompose it into homogeneous \(p\)-degrees, and set
\begin{equation}
 \boxed{U_{o,\nu}(p)=
   \sum_{n=0,2,4}\frac1{2(n+2)}
       \partial_{p_\mu}\bar A^{(n)}_{\mu\nu}(p).}
 \label{eq:v29-Uo}
\end{equation}
Put \(B^1=B^t-\mathcal L(0,0,U_o)\) and form its antisymmetric first
jet \(A^1\). The polynomial tensor
\[
 D_{\mu\nu}(p,z)=A^1_{\mu\nu}(p,z)/P(p,z)
\]
is antisymmetric, exchange even, and has weight at most two. Set
\begin{equation}
 \boxed{U_{1,\nu}(k,p,z)=\frac14 k^\mu D_{\mu\nu}(p,z),\qquad
 U=U_e+U_o+U_1.}
 \label{eq:v29-U1}
\end{equation}
Finally apply \eqref{eq:v29-symmetric} to
\(B-\mathcal L(0,0,U)\). Its result is \(T\). Exactly,
\begin{equation}
 \boxed{\mathcal L(T,0,U)=B.}
 \label{eq:v29-right-inverse}
\end{equation}
\end{theorem}
\begin{proof}
The first test gives the first monic division and removes \(b\). Its even
source has no antisymmetric linear \(k\) jet. Formula \eqref{eq:v29-Uo}
and \cref{lem:v29-koszul} give
\(\bar A=2p\wedge U_o\), exactly the modulo-\(P\) first jet of this
odd source. Thus \(A^1\) has zero remainder upon division by \(P\).
For \(U_1\), differentiation of \eqref{eq:v29-map} at \(k=0\) gives
\(\partial_{k_\mu}\mathcal L_\nu=P D_{\mu\nu}/2\). Antisymmetrizing
therefore gives \(PD=A^1\). This source has zero constant jet. The
remaining vector has zero constant jet and symmetric linear jet, so
\cref{lem:v29-symmetric} applies. All degree bounds follow from subtracting
one degree for a ghost-momentum derivative and two for division by the
monic Euler polynomial. Necessity was proved in \cref{lem:v29-tests}.
\end{proof}

An arbitrary prescribed even polynomial \(f_{\rm phys}\) of weight at
most four can be accommodated by applying the same construction to
\(B-\mathcal L(0,f_{\rm phys},0)\). Its tests are unchanged by
\cref{lem:v29-tests}. This is freedom in the displayed row, not a claim
that every independent finite normalization extends to all BV equations.
The source \(U_o\) can be nonzero for a constant ghost. It is an explicit
\emph{quantum scalar-transformation normalization}; it is not advertised
as preserving the old unnormalized scalar-antifield coefficient.

\subsection{An exact rank calculation, not an anomaly count}

The row space contains
\[
 4\sum_{i+j+2\ell\le5,\;j\ {
m even}}
      \binom{i+3}{3}\binom{j+3}{3}=2964
\]
coefficients. At zero ghost momentum, the even scalar polynomials of
weight at most four have dimension 58. Their quotient by \((P)\) has
basis all even polynomials in \(p\) of degrees 0, 2, 4 and dimension
\(1+10+35=46\). Thus the first test has \(4\cdot46=184\) independent
components. The remaining first-jet quotient has dimension \(6\cdot46=276\).
The ranks of \(p\wedge\) on its degrees 0, 2, 4 are respectively
\(6,45,140\). Indeed its kernels have dimensions 0, 15, 70: the latter
are the images of polynomial one-forms of degrees 1 and 3, of dimensions
16 and 80, modulo \(p\) times scalar polynomials of dimensions 1 and 10.
The homotopy proves exactness of these counts. All symmetric first jets
and all higher \(k\) jets are in the metric image. Hence
\begin{equation}
 \boxed{\operatorname{rank}\mathcal L=2964-184-191=2589.}
 \label{eq:v29-rank}
\end{equation}
These are counts of this particular row map, not of the full scalar--gravity
BV complex or of nonzero anomalies. The tests are triangular and
independent: even \(U_e\) fills the Euler-divisible constant jets;
odd \(U_o\) fills the wedge part; \(U_1\) fills the Euler-divisible
antisymmetric jets; \(T\) fills the remaining components.

\subsection{A fixed, explicit coefficient bound}

Use the sum of absolute monomial coefficients, with a matrix tensor counted
in its full symmetric or antisymmetric array, and the corresponding scaled
norms at momentum scale \(\lambda\). The constructed linear map obeys
\begin{equation}
 \boxed{\|(T,U)\|_{\rm out,\lambda}
        \le 2^{30}\|B\|_{\rm row,\lambda}.}
 \label{eq:v29-norm}
\end{equation}
The output degrees are four for \(T\), three for \(U\), and five for
\(B\); the norm includes the compensating powers of \(\lambda\).
The numerical constant is intentionally conservative, not an optimized
conditioning result. To verify it directly, substitution
\(z=-p^2\) costs at most 16; monic division in \(z\) costs at most 5;
first-jet antisymmetrization at most 2; the odd-source homotopy at most
\(1/2\); scalar-leg shifts cost at most \((3/2)^d\) at degree \(d\).
Consequently the scalar Euler row has norms at most 46 and 68 on sources
of scalar degree at most two and three, respectively. Taking the original
row norm as one, bounds for the three successive remainders are
\[
 C_1=1+46(5/2)=116,\quad
 C_2=(1+68\cdot16)C_1=126324,\quad
 C_3=(1+46\cdot5)C_2=29180844.
\]
The final symmetric section costs at most 8. The sum of its bound and the
three source bounds is less than \(234080231<2^{30}\). This proof is
unchanged at zero physical mass; none of its denominators involves mass,
momentum, the cutoff or a volume. Reality is retained by averaging the output with its coefficient
conjugation involution. The maps act componentwise on the common flavour
Kronecker tensor, so they preserve the stated flavour symmetry.

\section{Why the actual relative reference passes the polynomial tests}
\label{sec:v29-native}

The sufficient conditions in \cref{thm:v29-primitive} are not imposed on an
arbitrary breaking. This section obtains them for
\eqref{eq:v29-relative-B}. The argument uses the ordinary coupled reference
identity, but no finite-regulator BRST invariance.

\begin{lemma}[The reference consistency extraction with scalar Euler terms]
\label{lem:v29-consistency}
The ordinary coupled lower-reference functional has a local, mass-weighted
jet \(X\) satisfying \(sX=0\) at the arities used here. At \(\phi=\upsilon=0\),
its coefficients containing a gravitational minimal antifield equal the
pure-gravity ones. The difference can include a scalar vacuum
\emph{metric} reference coefficient, but has no such antifield coefficient.
Consequently \(\mathcal B=X-sH_g\) has zero purely gravitational
minimal-antifield component through the required arity and weight, and
satisfies the coupled consistency identity.
\end{lemma}
\begin{proof}
Use the ordinary master action of the zero-mesh stationary Palatini
functional plus \(S_H\), its actual coordinate-induced gravitational
transformations and the scalar Lie transformation. The scalar density
changes by a total divergence; the scalar Lie representation and its
semidirect bracket obey the graded commutator identity. Together with the
inherited gravitational and nonminimal identities this proves the ordinary
classical master equation. The mass is an inert parameter, so the same
identity holds coefficientwise in \(z\).

Apply the differentiated ordinary reference trace identity used in V22,
now with the scalar Hessian and scalar source included. Its proof only uses
the classical master identity, graded cyclicity and differentiation of
\((S''+\mathbb R)^{-1}\). Cancel \(C\mathbb R=\Theta\) before expanding.
Every covariance-bearing word needed here has a factor \(\Theta\) and a
hard weight \(\nu\), and is supported between fixed multiples of
\(\lambda\). Thus its derivatives and trace rearrangements are justified
on a compact annulus. The source-free no-covariance term is treated as its
finite coefficient trace. Its scalar \(c\phi^2\) extraction is zero.
The same identity gives \(sX=0\), including the scalar-antifield Euler
terms, before the mass-weighted local projection. The weight/arity
argument in \cref{sec:v29-start} transfers it to the stated local bank.

At zero scalar and scalar-antifield background the ordinary Hessian is
block diagonal between scalar and gravitational/ghost variables. Neither
V18's metric source nor the ghost action has a scalar block in its
metric-antifield variation; this is also the zero-background identification
in V28 preceding \cref{thm:v28-finite-row}. The scalar block of \(Dr\)
has no gravitational antifield. Differentiating the inverse therefore
proves the asserted equality of gravitational-antifield coefficients.
Subtracting \(sH_g\) removes precisely those local source coefficients,
by the inherited pure-gravity normalization. The scalar vacuum reference
has no antifields and so cannot supply a scalar-stress return through
\(s\eta\). Finally \(s^2=0\) is used only for this ordinary master action,
giving \(s(X-sH_g)=0\).
\end{proof}

\begin{theorem}[Relative Killing test for the native two-scalar row]
\label{thm:v29-native-tests}
For \(B\) of \eqref{eq:v29-relative-B},
\begin{equation}
 \boxed{B(0,p,z)=0,\qquad p\wedge\bar A(p)=0.}
 \label{eq:v29-native-tests}
\end{equation}
The statements are polynomial identities in the ordinary thermodynamic
coefficient convention. They do not assert exact finite-torus rotational
symmetry of the microscopic grid.
\end{theorem}
\begin{proof}
A constant ordinary ghost acts as translation on both the metric coordinate
and the scalar: the pointwise coordinate chain rule gives
\(r^q=c\cdot\partial q\). Constant-ghost vanishing can also be checked
directly on the full reference words after restoring the nonminimal
contact. The ordinary action and ghost Hessian vertices have two derivatives,
the transformation has one, and the paired propagators and profiles are
even under simultaneous reversal of loop and external momenta. The resulting
\(c\phi^2\) row is odd under \((k,p)\mapsto(-k,-p)\).
At \(k=0\) scalar exchange makes it even in \(p\), so its value is zero.
This check includes the ghost return, rather than treating only a bosonic
translation determinant. The fixed pure-gravity normalization has no zero-background,
derivative-free metric-transformation coefficient: this is the preserved
constant-ghost source normalization of V23--V26. Its scalar-stress return
therefore has none either. This proves the first identity. This argument
uses the fixed normalization, not arbitrary additional homogeneous
source terms which would need their own constant-ghost matching.

To prove the second identity, fix a nonzero real momentum \(p\) and work
\emph{only with local polynomial coefficients} modulo \(z+p^2\).
Take two scalar plane waves of momenta \(p,-p\). Their ordinary Euler
polynomials vanish in this quotient. Let
\(\mathfrak k_p=\{A\in\mathfrak{so}(4):Ap=0\}\).
An affine ghost \(c(x)=Ax\), \(A\in\mathfrak k_p\), fixes both waves
and the flat metric. It is used here as a local jet test, not a periodic
physical ghost. The scalar Euler terms of the coupled consistency equation
vanish on these waves, including derivatives of the Euler equations.
The metric Euler term can have scalar stress, but it multiplies the
pure-gravity antifield coefficient which was removed in
\cref{lem:v29-consistency}. Nonminimal equations vanish on these ordinary
Killing jets after their off-shell contact has been included. No scalar
Euler term is omitted before making this restriction.

It follows that the only surviving term of consistency on two such
rotations is the row evaluated on their commutator. Hence its functional
on \(\mathfrak k_p\) annihilates \([\mathfrak k_p,\mathfrak k_p]\).
Since \(\mathfrak k_p\cong\mathfrak{so}(3)\) and the three elementary
rotation generators are each commutators of two others, this commutator
space is the entire stabilizer. The row vanishes on it. Because its
constant-ghost coefficient is zero, its affine-ghost value is exactly the
antisymmetric first \(k\) jet, with no contribution from differentiating
the external scalar factors. Thus \(\bar A\) annihilates all rotations
on \(p^\perp\). Elementary orthogonal decomposition of two-forms then
gives \(\bar A=p\wedge v\), or equivalently
\(p\wedge\bar A=0\), for this \(p\). The polynomial vanishes for every
nonzero real \(p\), so it vanishes identically. No division by \(p^2\)
is used in the subsequent construction.

The formal evaluation \(z=-p^2\) in this proof concerns a finite polynomial
identity following from the ordinary master equation. No annular trace is
evaluated there. Its physical massive coefficients were differentiated at
\(z=0\) on the positive hard domain as in V28. This distinction is what
makes the argument regular in the physical mass.
\end{proof}

The stabilizer is used rather than the full rotation algebra because only
these rotations fix the scalar waves. Away from the stabilizer, scalar
Euler/source terms can contribute. The counterexample in
\cref{sec:v29-homotopies} demonstrates concretely why discarding them would
produce a false test. This argument proves a row-specific primitive; it
is not a replacement for a classification of the full matter BV cohomology.

\section{The relative counterterm and the finite-cutoff local remainder}
\label{sec:v29-update}

\begin{theorem}[Local mixed normalization preserving the gravity boundary]
\label{thm:v29-relative-primitive}
Apply \cref{thm:v29-primitive} to \eqref{eq:v29-relative-B} and form
\(h_H=C_T+C_U\), or include a specified compatible \(C_f\) as described
above. Then
\begin{equation}
 \boxed{[s(H_g+h_H)]_{c\phi^2,\loc}
                 =[\mathcal X_{0,\lambda}]_{c\phi^2,\loc},
 \qquad h_H|_{\phi=\upsilon=0}=0.}
 \label{eq:v29-relative-primitive}
\end{equation}
The map from the actual reference coefficient to \(h_H\) is finite,
linear, mass-polynomial and bounded as in \eqref{eq:v29-norm}.
Its output has the local types
\(q\phi^2\) through weight four and
\(\upsilon\phi c\) through weight three, with an optional scalar quadratic
normalization through weight four. No gravitational counterterm is reset.
\end{theorem}
\begin{proof}
The native tests in \cref{thm:v29-native-tests} give
\([sh_H]_{c\phi^2}=B\) by \eqref{eq:v29-right-inverse}.
Substitute \eqref{eq:v29-relative-B}. Every term of \(h_H\) contains either
two scalar fields or one scalar and its antifield, proving the relative
restriction. The bounds and polynomial dependence follow from the two
homotopies and the two monic divisions. They apply to the coefficient
array itself and require no numerical value of an annular integral.
\end{proof}

Let \(g_{a,H}\) be the actual weighted local jets of the two-scalar action,
one-metric/two-scalar action, and unmarked scalar-antifield coefficient
from V27--V28. Add the order-\(\hbar\) local insertion
\begin{equation}
 \boxed{C_{a,H}^{\rm new}=-g_{a,H}+h_H.}
 \label{eq:v29-counterterm}
\end{equation}
The already accumulated gravity normalization is not included in
\(-g_{a,H}\) and is not counted a second time. The scalar mass and
wave-function coefficients of the particular reference choice \(f=0\)
are assigned at the physical normalization point; this is a normalization,
not a prediction that their loop coefficients vanish. Nonzero physical
finite choices can be inserted through the prescribed \(f\) channel or
through full ordinary closed functionals.

\begin{proposition}[The new local reference comparison]
\label{prop:v29-local-rate}
Under \eqref{eq:v28-domain}, let \(d=|\alpha|+2j\le5\). In a common
coefficient convention,
\begin{equation}
 \begin{split}
 &\left|\partial_K^\alpha\partial_z^j\,
 \mathbf T_5(\mathcal X^H_{a,\lambda,L}
                -\mathcal X^H_{0,\lambda,\infty})(0,0)\right|\\
 &\quad\le C_{\alpha,j,b}\chi^{-1}\lambda^{5-d}
       \left[(a\lambda)^4+(\lambda L_{\min})^{-b}\right],\qquad b>4.
 \end{split}
 \label{eq:v29-local-rate}
\end{equation}
The same bound is the local modified-identity residual after
\eqref{eq:v29-counterterm}. Its coefficient map has no cutoff or volume
loss beyond the displayed reference error.
\end{proposition}
\begin{proof}
The same ordered words as in \eqref{eq:v28-reference} are confined to
\(\lambda/2\le|p|\le3\lambda\). At this background order their vertices
and free inverses differ from their ordinary counterparts at fourth
relative mesh order, as established immediately before
\eqref{eq:v28-reference-rate}. The unprojected row integrand has order one;
\(d\) weighted external derivatives lower it to order \(1-d\).
Four-dimensional integration therefore gives
\(\chi^{-1}a^4\lambda^{9-d}\). Unlike V28's complementary estimate,
no sixth-order Taylor remainder is taken here: these are the actual local
coefficients. Smooth compact annular support and Poisson summation give
the separate period error. There is no source-free zero-mode contact in
this scalar row.

After the local update the ordinary local row is exactly
\(\mathbf T_5\mathcal X^H_{0,\lambda,\infty}\), by
\cref{thm:v29-relative-primitive}. The finite ordinary multiplier differences
\(\widetilde P_a-P_0\) and \(V_a-V_0\) have weight at least six.
They cannot contribute to the local weight-five row when applied to these
regular local polynomials. The same applies to the fixed gravity-source
coefficient. Thus the local finite residual is exactly the reference
comparison displayed in \eqref{eq:v29-local-rate}.
\end{proof}

The common norm of \(h_H\), at scale \(\lambda\), is bounded by a constant
depending on the fixed gravity reference normalization and the native
profile seminorms, not on \(a\) or \(L\). The divergent coefficients in
\(g_{a,H}\) are retained in the counterterm. They are not bounded by this
statement and do not define an invariant analytic ball of microscopic
actions.

\begin{corollary}[The first full local-plus-nonlocal mixed identity]
\label{cor:v29-full-row}
Combine the update with the paired pure-gravity normalization of V26 and
the nonlocal limits of V28. On the displayed row, the ordinary normalized
functional satisfies
\begin{equation}
 \boxed{\mathfrak L^{\rm ren}_{0,\lambda}
                       =\mathcal X^H_{0,\lambda},}
 \label{eq:v29-full-row}
\end{equation}
with the \emph{unsubtracted} lower-reference term. In fixed scaled source
norms a conservative finite comparison is
\begin{equation}
 \|\mathfrak L^{\rm ren}_{a,\lambda,L}
                    -\mathcal X^H_{a,\lambda,L}\|_{\rm src}
 \le C_H\chi^{-1}\lambda^5
       [a\lambda+(\lambda L_{\min})^{-b}],\qquad b>4.
 \label{eq:v29-full-rate}
\end{equation}
The action-level bound has the additional factor \(\hbar\).
\end{corollary}
\begin{proof}
V28 supplies \(\mathfrak L_0^R=\mathcal X_0^{H,R}\).
Equation \eqref{eq:v29-relative-primitive} supplies exactly its missing
local part with the fixed gravity normalization included. Add the two
identities. For the finite comparison use V28's nonlocal source estimate
and \cref{prop:v29-local-rate}. Finite multiplier differences acting on
\(h_H,H_g\) have fixed coefficients and higher weight and fit the same
conservative bound. The previously divergent local products of V28 are
already assigned consistently by subtracting \(g_{a,H}\); they are not
counted once again in the limiting normalized coefficients. Apply the
inherited windowed Fourier estimates to obtain the corresponding rooted
multilinear bound with one source in \(L^1\) and the others in
\(L^\infty\). The period error is separate. No exact off-grid finite-torus
Ward identity or removal of \(\lambda\) is inferred.
\end{proof}

\section{Explicit examples, descendants, and the next coupled interface}
\label{sec:v29-descendants}

A useful massive example is an antisymmetric tensor
\(D_{01}=z+p_2^2\), \(D_{10}=-D_{01}\), with all other entries zero.
The source
\[
 U_1=\tfrac14k_0(z+p_2^2),\qquad
 U_0=-\tfrac14k_1(z+p_2^2)
\]
produces an antisymmetric first ghost jet \(P D\). The compiler removes
it by the monic Euler division and recovers the remaining symmetric metric
vertex, with \(f=0\). The example is a nonzero local normalization row,
not a claimed value of the native reference integral.

The earlier example \(U_1(0,p)=p_0^3\) is also recovered by the complete
compiler even though it fails the unnecessarily strong rotation-cocycle
test. As a distinct metric-only check, every homogeneous vector
\(B_\nu=k_\mu^n\delta_{\nu\rho}\), \(2\le n\le5\), is recovered by
\eqref{eq:v29-symmetric}. These checks independently test the two different
homotopies and the scalar Euler source rather than treating successful
matrix rank computation as a proof of the native reference identity.

\subsection{The displayed row is not the full matter complex}

The new insertion is a genuine ghost-number-zero functional. Its entire
BV variation, not only \eqref{eq:v29-map}, must therefore be used at the
next step. In particular, \(s C_U\) includes the scalar-antifield
coefficient obtained by varying its scalar and ghost arguments. It has
type \(\upsilon\phi cc\) and need not vanish. The \(q\phi^2\) term also
produces a \(q\phi^2c\) row through the nonlinear metric transformation,
and the metric variation of the scalar Euler operator in \(C_U\)
contributes there. The scalar Euler return of higher-antifield parents
must be retained when those rows are assembled. The primitive is not
claimed to solve them merely because its first action row is exact.

The gravity boundary is preserved strictly, since
\(h_H|_{\phi=\upsilon=0}=0\). This is stronger than solving an absolute
mixed equation and then choosing independent gravity constants afterward.
It is weaker than preserving every unnormalized scalar source coefficient:
\(U_o\) is explicitly allowed to change the constant-ghost translation
source at one loop. Its associated higher-source contacts must change with
it. The mass, wave-function and covariant higher-derivative physical
normalizations remain free compatible data, not predictions of the row
compiler.

The next load-bearing calculation is the local and nonlocal
\(\upsilon\phi cc\) equation, together with the one-metric extension of
\(c\phi^2\), after inserting \emph{this same} \(C_{a,H}^{\rm new}\).
It consumes the actual scalar source square of V28, the field-dependent
metric/ghost source defects, the mixed higher-antifield returns, and their
reference terms. No internal gauge or fermion/phase theorem has been
established by the present scalar calculation. General matter-inclusive
local BRST classification is available in the literature
\cite{V29BBHMatter}; it is not used here to replace the row calculation.

\section{Final ledgers for V29}
\label{sec:v29-ledgers}

\subsection*{Proved}
\addcontentsline{toc}{subsection}{V29: Proved}
The complete mass-polynomial \(c\phi^2\) row map is
\eqref{eq:v29-map}. Its image is characterized by the two polynomial tests
\eqref{eq:v29-tests}, has rank 2589 in the 2964-dimensional declared row
space, and has the explicit right inverse
\eqref{eq:v29-Uo}--\eqref{eq:v29-right-inverse}, bounded by
\eqref{eq:v29-norm}. The scalar-antifield Euler channel invalidates the
stronger full-rotation-cocycle shortcut. The native ordinary relative
reference passes the correct translation and stabilizer tests by its
coupled consistency identity with scalar Euler and gravity stress terms
retained. This supplies a local mixed counterterm without changing the
previous gravity normalization. The local reference comparison and the
first full local-plus-nonlocal modified row have the stated cutoff and
period bounds. These are results for one displayed mixed row.

\subsection*{Inherited}
\addcontentsline{toc}{subsection}{V29: Inherited}
All earlier mathematical bodies and labels are retained. V27--V28 supply
the common massive hard integral, its source triangle, weighted Taylor
calculus and nonlocal estimates. V22 supplies the ordinary reference
trace mechanism. V24--V26 supply the fixed pure-gravity normalization,
including V25's external local-cohomology input and the qualifications of
its unevaluated coefficient prescription. V23--V26 supply its chosen
constant-ghost source normalization. The present application is not an
independent audit of those inherited proofs. The paired action--density
convention, actual Haar/Jacobian contributions and lower Wilsonian scale
remain unchanged.

\subsection*{Literature input}
\addcontentsline{toc}{subsection}{V29: Literature input}
The exact-RG/modified-reference framework of \cite{V28IIM} remains the
comparison for the nonzero lower-reference term. The broader
Einstein--Yang--Mills analysis \cite{V29BBHMatter} includes scalar matter
and stresses the role of antifields and abelian factors. No conclusion
about the full Standard Model or its chiral measure is imported here.
The two polynomial homotopies used in this installment are proved directly.

\subsection*{Conditional}
\addcontentsline{toc}{subsection}{V29: Conditional}
The native application consumes the established ordinary coupled source
identity and the fixed gravity normalization with its exact source
components. A different gravity normalization with derivative-free
constant-ghost sources would require its actual boundary matching before
\cref{thm:v29-native-tests} is applied. Numerical counterterm values require
evaluating the prescribed native annular coefficients; this installment
provides their finite linear compiler, not numerical integrals. Further
physical normalization constraints must be tested together with the
higher-source rows, not independently imposed on every coefficient.

\subsection*{Open}
\addcontentsline{toc}{subsection}{V29: Open}
The \(\upsilon\phi cc\) and \(q\phi^2c\) equations after the complete
counterterm update are not solved here. The full simultaneous
matter-antifield bank, internal gauge and chiral source/measure sectors,
higher-loop EFT closure, and an invariant interacting measured RG class
remain open. The result is not an all-orders ESM theorem, an infrared
limit, a positive Lorentzian measure, or a finite-parameter ultraviolet
completion of gravity.

\begin{center}\small
\begin{tabular}{>{\raggedright\arraybackslash}p{.24\textwidth}>{\raggedright\arraybackslash}p{.67\textwidth}}
\toprule
Variable & Scope in V29\\\midrule
Cutoff and volume & Coefficient primitive has fixed finite operator norm.
Native mesh estimates use \eqref{eq:v28-domain}; thermodynamic symmetry
and finite-period comparison are distinguished.\\
Sources & One ghost and two equal-species scalar fields; the counterterm
retains its metric and scalar-antifield contacts. No full matter BV rank
or all-source-order constant is asserted.\\
Mass & Polynomial weighted jets in \(z=m^2\), weight two. No inverse
\(m\), inverse \(m^2\), or replacement of a massive propagator. Physical
hard domain \(0\le m\le M\lambda\) with fixed \(\lambda>0\).\\
Couplings and loops & One loop, fixed \(\chi>0\); zero internal gauge,
Yukawa and scalar self-couplings. Higgs flavour symmetry is used only in
this equal-mass component.\\
Measure & Same completed contour, ghost action and paired
Haar/coordinate-density convention. No additional scalar density chosen.\\
Iteration & Inherited summability for the displayed coefficients;
no analytic interacting self-map or higher-loop claim.\\
\bottomrule
\end{tabular}
\end{center}

The self-contained checker \srcid{check_ESM_relative_scalar_local_V29.py}
verifies the finite polynomial sections, exact stabilizer ranks, mass
quotient identities, complete example row recoveries, and a counterexample
to the stronger rotation test. It also retains an independent exact
rotation-matrix diagnostic solely to explain that rejected shortcut.
It does not numerically integrate a native loop or independently prove
the inherited analytic estimates. The certificate records the preserved
V1--V28 body hash. No new Lean formalization is claimed.

\begingroup
\renewcommand{\refname}{Additional comparison source for V29}
\begin{thebibliography}{99}
\bibitem[47]{V29BBHMatter}
G.~Barnich, F.~Brandt and M.~Henneaux,
\emph{Local BRST cohomology in Einstein--Yang--Mills theory},
Nuclear Physics B \textbf{455} (1995), 357--408,
arXiv:hep-th/9505173. Sections 2 and 10 discuss the matter carrier,
antifields and anomaly classification. Comparison only; no new
matter-anomaly vanishing theorem is assumed in the polynomial row proof.
\end{thebibliography}
\endgroup


\clearpage
\part*{V30: The scalar algebra insertion and its missing quadratic-antifield parent}
\addcontentsline{toc}{part}{V30: Scalar algebra insertion and quadratic-antifield return}

\section{The next scalar equation requires a new parent, not another one-row primitive}
\label{sec:v30-start}

Sections 1--231 and their labels are retained unchanged. V29 supplies a
particular local counterterm for the first mixed action row. It expressly
does not establish the scalar-antifield/two-ghost equation after that
counterterm is inserted. Two different operands of the latter equation
must be distinguished: the measured insertion of the finite scalar
source square, and the Euler return of the quantum coefficient with two
scalar antifields. Neither is V28's one-antifield triangle or V19's
metric-antifield curvature insertion. We evaluate both here. We also
calculate the entire zero-background source-algebra descendant of V29's
scalar-antifield counterterm. No new regulator or reference is introduced.

The conclusions below concern four equal-mass real Higgs components,
with internal gauge, Yukawa and scalar self-couplings set to zero, and
one quantum loop. The existing metric contour, nonminimal ghost action,
Haar/coordinate amplitude and smooth ultraviolet completion remain in
force. The factor \(\hbar\) is removed from one-loop coefficients. The
auxiliary amplitude is not independently inserted again into the new
source coefficient. The ordinary background metric is zero. This is a
calculation of specific mixed source operands, not a claim that the
complete matter quantum BV complex now closes.

\subsection{Source audit and domains}

We inherit the actual scalar action
\eqref{eq:v27-completed-Higgs}, the metric--scalar block of V27, the
scalar source \eqref{eq:v28-scalar-source}, V18's ghost source, and the
unmodified ghost-action vertex \eqref{eq:v18-operands}. The Berezin and
Gaussian elimination rules themselves are inherited. V19's source
curvature concerned \(q^*,c^*\), not \(\phi^*\); its loop had no
propagating scalar inverse. V28's scalar-antifield coefficient had only
one retained ghost and was linear in \(\phi^*\). V29's local polynomial
compiler is not repeated. Its input coefficients are still specified by
their native reference integrals and are not numerically evaluated here.

Use \(z=m^2\), \(\upsilon=\phi^*\), \(S=s(aD)\), and \(P=S^2\).
The profile is real, even and radial, equal to one below \(r_-=1/64\)
and zero above \(r_+=1/32\). Write \(\pi_a(p)=s(ap)^2\) for its
symbol, to distinguish the filter from a scalar Euler multiplier. The
ordinary finite spectral derivative has symbol \(ip_\mu\); products
are the actual cyclic products. The covariances are
\[
 C_g=\nu_\lambda\widetilde K_a^{-1},\qquad
 C_c=\nu_\lambda\widetilde M_a^{-1},\qquad
 C_H=\nu_\lambda\widetilde P_a^{-1},\qquad
 \nu_\lambda=1-\Theta_\lambda.
\]
The zero-mode/low-mode sectors remain retained. Throughout the estimates,
\begin{equation}
 \boxed{\lambda\ge8192\mu>0,\qquad
 a\lambda\le2^{-20},\qquad \lambda L_\nu\ge1,\qquad
 0\le m\le M\lambda.}
 \label{eq:v30-domain}
\end{equation}
The numbers \(L_\nu/a\) are odd integers, every external momentum has
norm at most \(\mu\), and there are four external arguments in each
new coefficient. Constants depend on the fixed profile seminorms,
\(M\), contour/inverse margins, \(\chi>0\), and the prescribed
finite derivative or moment order. None depends on the total number of
sites. We use the inherited off-grid smooth interpolation only to define
momentum derivatives; no exact off-grid torus Slavnov identity is inferred.

\section{The scalar square has an exact finite momentum symbol}
\label{sec:v30-square}

For an odd ghost \(c\), set
\[
 u=Pc,\qquad b_c=u\cdot Du,\qquad w_c=(P-I)b_c,
 \qquad
 \mathcal N_a(c)\varphi
 =P\{\mathcal L_{w_c}\varphi+
               \mathcal L_u(I-P)\mathcal L_u\varphi\}.
\]
Here \(\mathcal L_v\varphi=v\cdot D\varphi\) is the scalar Lie
operation, not a tensor or half-density operation.

\begin{proposition}[Finite square and its evaluated polarized symbol]
\label{prop:v30-square}
On the full finite carrier,
\begin{equation}
 \boxed{\breve s^{\,2}\phi=\mathcal N_a(c)\phi.}
 \label{eq:v30-square}
\end{equation}
For even polarizations \(v,w\), ghost waves \(v e^{ikx},w e^{ilx}\),
and scalar wave \(e^{irx}\), the coefficient of the ordered exterior
product of their two ghost labels is
\begin{equation}
\begin{split}
 \mathfrak n_a(k,l,r;v,w)
 ={}&\pi_a(k+l+r)\pi_a(k)\pi_a(l)\bigl\{
 \pi_a(r+l)(w\cdot r)(v\cdot(r+l))\\
 &\quad-\pi_a(r+k)(v\cdot r)(w\cdot(r+k))\\
 &\quad-\pi_a(k+l)
       [(v\cdot l)(w\cdot r)-(w\cdot k)(v\cdot r)]\bigr\}.
 \label{eq:v30-square-symbol}
\end{split}
\end{equation}
Whenever a factor contributes, its differentiated intermediate labels
are nonaliased; thus the formula uses their ordinary representatives.
It is antisymmetric under \((k,v)\leftrightarrow(l,w)\), even under
simultaneous reversal of all momenta, and zero when every required
filter equals one. It also annihilates a constant scalar exactly.
\end{proposition}
\begin{proof}
Differentiate the actual source rather than impose a representation
identity. The odd Leibniz rule gives
\[
 \breve s^{\,2}\phi=P\mathcal L_{Pb_c}\phi
                         -P\mathcal L_uP\mathcal L_u\phi.
\]
If a term survives the outer \(P\), its output and both filtered ghost
momenta have dimensionless norms below \(r_+\). Its scalar input has
norm at most \(3r_+\), and all differentiated intermediate labels have
norm at most \(4r_+<\pi\). A principal-face input cannot be moved to
the central output by these two small labels. For these terms the
ordinary identity \(\mathcal L_u^2=\mathcal L_{b_c}\) is valid.
Adding and subtracting \(P\mathcal L_{b_c}\) proves
\eqref{eq:v30-square}, without a product rule for unrestricted grid
fields.

Directly, \(r_a^H(v_k)\varphi_r
=i\pi_a(k+r)\pi_a(k)(v\cdot r)\varphi_{k+r}\).
The bracket of the two filtered ghosts is
\(i\pi_a(k)\pi_a(l)((v\cdot l)w-(w\cdot k)v)\).
Subtract the two ordered compositions of the preceding scalar operators
from the source applied to this bracket. This gives
\eqref{eq:v30-square-symbol}, including its signs. Interchanging the
two ghost labels changes the sign, and reversing all momenta changes
neither product of two momentum contractions. For unit filters the
braces cancel by expanding \(v\cdot(r+l)\) and
\(w\cdot(r+k)\). Every term contains the scalar momentum \(r\).
\end{proof}

The finite square is generally nonzero. Its vanishing on an interior
panel does not license dropping its loop insertion: in that insertion
some of the ghost and scalar arguments are hard. The scalar equation
\eqref{eq:v30-square} is the missing representation component of the
source curvature, not the entire classical master defect.

Introduce an independent odd tag on the left of
\begin{equation}
 \boxed{\Omega_H=\langle\upsilon,\mathcal N_a(c)\phi\rangle.}
 \label{eq:v30-insertion}
\end{equation}
The tag makes the perturbation of the exponent even. It is a source
bookkeeping device, not an additional integrated field. The coefficient
of the tag is defined by differentiating the same normalized finite
integral as in V27--V28.

\section{Two measured traces and the quadratic-antifield parent}
\label{sec:v30-measured}

Let \(x,\zeta,\beta,\bar\beta\) denote hard metric, scalar, ghost
and antighost variables. In this section \(\zeta\) is even; it is not
the odd tag just described. The mixed action block is
\(\langle x,L_\phi\zeta\rangle\), with \(L_\phi=L(0;\phi)\)
from V27. The unchanged ghost action gives
\(\langle\bar\beta,B_cx\rangle\). Its flat momentum vertex is
\[
 B_c(p,k)x_p=-s(a(p+k))s(ap)s(ak)
                  F_{p+k}R_1^{(0),E}(x_p,c_k).
\]
In particular it has the displayed hard profiles even though the
scalar action retains interacting upper modes. The optional upper
pure-ghost vertex vanishes here because its retained ghost factor is
\((I-S)c=0\).

Put \(\beta_x=-C_cB_cx\), and define odd coefficient kernels by
\begin{align}
 \tfrac12\langle x,A_\Omega(\phi)x\rangle
   &=\langle\upsilon,\mathcal N_a(\beta_x)\phi\rangle,\\
 \langle x,B_\Omega\zeta\rangle
   &=\langle\upsilon,
       [t]\mathcal N_a(c+t\beta_x)\zeta\rangle.
 \label{eq:v30-AB}
\end{align}
The extraction parameter \(t\) is even, the exterior coefficients
keep their original order, and both definitions follow coefficient
extraction in independent hard variables. In particular \(A_\Omega\)
contains two ghost propagators, while \(B_\Omega\) contains one.

\begin{theorem}[Complete scalar-curvature insertion at one loop]
\label{thm:v30-marked}
At zero retained metric, scalar degree one, one scalar antifield and
two retained ghosts, the coefficient of \(\Omega_H\) is
\begin{equation}
 \boxed{\mathfrak I_{a,1}^H
 =\tfrac12\operatorname{Tr}(C_g A_\Omega(\phi))
       -\operatorname{Tr}(C_gB_\Omega C_HL_\phi^T).}
 \label{eq:v30-marked}
\end{equation}
The trace is over the specified native operators, not freely chosen
vertices. There is no additional Higgs multiplicity factor for a fixed
external scalar/antifield species. The formula contains all one-loop
contributions to this insertion in the declared panel.
\end{theorem}
\begin{proof}
The hard-degree-two coefficient before eliminating ghosts is exactly
\[
 \langle\upsilon,\mathcal N_a(\beta)\phi
             +[t]\mathcal N_a(c+t\beta)\zeta\rangle.
\]
There is no direct hard-metric derivative of the scalar source.
Berezin integration over the hard antighost substitutes
\(\beta=-C_cB_cx\); applying it to the displayed expression gives
\eqref{eq:v30-AB}. At the required scalar degree the bosonic inverse is
\[
 \begin{pmatrix}
 C_g&-C_gL_\phi C_H\\
 -C_HL_\phi^TC_g&C_H
 \end{pmatrix}+O(\phi^2)
\]
in its diagonal blocks and \(O(\phi^3)\) in its off-diagonal blocks.
The first variation of its half-logarithm against the marked Hessian
\(\left(\begin{smallmatrix}A_\Omega&B_\Omega\\B_\Omega^T&0\end{smallmatrix}\right)\)
is \eqref{eq:v30-marked}. The tag makes this use of ordinary even
matrix differentiation legitimate; it is removed without reordering the
remaining odd coefficients.

Every hard-linear vertex would carry a sum of at most four soft external
labels, below the support of every hard covariance on
\eqref{eq:v30-domain}. Such diagrams are zero. In a connected one-loop
graph without a hard-linear vertex,
\(\sum_v(\deg v-2)=2(1-1)=0\), so every hard degree is exactly two.
The quadratic-hard calculation therefore exhausts the insertion. The
auxiliary amplitude could contribute at this quantum order only by a
tree hard-return diagram; its potentially relevant hard-linear branch
is absent. The scalar index follows the one scalar line and is fixed
by the external pair.
\end{proof}

The same joint Hessian determines a second, different operand. Let
\(N_\upsilon\) be exactly V28's source block at zero metric:
\[
 \langle x,N_\upsilon\zeta\rangle
 =\left\langle\upsilon,
 P\bigl([PC_cB_cx]^\mu D_\mu\zeta\bigr)\right\rangle.
\]
It is even, with one \(\upsilon\) and one retained ghost. At zero
retained scalar the off-diagonal Hessian is \(-N_\upsilon\).

\begin{proposition}[The native two-scalar-antifield coefficient]
\label{prop:v30-parent}
The complete one-loop coefficient with two scalar antifields and two
ghosts, at zero ordinary backgrounds, is
\begin{equation}
 \boxed{\mathfrak G_{a,1}^{\upsilon\upsilon}
       =-\tfrac12\operatorname{Tr}
              (C_gN_\upsilon C_HN_\upsilon^T).}
 \label{eq:v30-parent}
\end{equation}
The exterior product of independent antifield/ghost source labels must
be extracted after forming this signed trace. Its two ghost inverses
are part of the loop. This coefficient is not determined by the
one-antifield triangle alone.
\end{proposition}
\begin{proof}
After the same antighost integration, the joint determinant at zero
retained scalar is
\(\frac12\operatorname{Tr}\log(I-C_gN_\upsilon C_HN_\upsilon^T)\).
The displayed term is its entire coefficient of antifield degree two.
The hard-valence argument in \cref{thm:v30-marked} applies unchanged.
Since each \(N_\upsilon\) is even, its matrix transpose acts on the
even hard slots; no further exchange of external odd labels is made.
For example, for a scalar hard block and
\(N_\upsilon=a_1\upsilon_1c_1+a_2\upsilon_2c_2\), the coefficient
in \(-C_gC_HN_\upsilon^2/2\) of
\(\upsilon_1\upsilon_2c_1c_2\) is
\(+C_gC_Ha_1a_2\). This fixes the exterior sign without a diagrammatic
sign convention. It is an algebraic example, not a numerical native
loop evaluation.
\end{proof}

\section{Annular localization and a genuine parity gain}
\label{sec:v30-support}

The following support statements are for the assembled marked
coefficient. They do not remove upper modes from the underlying action.
Let \(v=P\beta_x\) and
\(w_1=(P-I)(c\cdot Dv+v\cdot Dc)\). Pairing with the soft source
labels simplifies \eqref{eq:v30-AB} to
\begin{align}
 \tfrac12\langle x,A_\Omega(\phi)x\rangle
  &=\langle\upsilon,P\mathcal L_v(I-P)\mathcal L_v\phi\rangle,\\
 \langle x,B_\Omega\zeta\rangle
  &=\langle\upsilon,
       P\{\mathcal L_{w_1}\zeta+
                      \mathcal L_v(I-P)\mathcal L_c\zeta\}\rangle.
 \label{eq:v30-simplified}
\end{align}
In the first line the term with \((P-I)(v\cdot Dv)\) has soft total
momentum and vanishes. In the second line the intermediate momentum in
\(P\mathcal L_c(I-P)\mathcal L_v\zeta\) is the soft final output
minus a soft ghost momentum; that term vanishes too. These cancellations
are made before taking absolute values.

\begin{theorem}[Support and scaling of the scalar-square insertion]
\label{thm:v30-annulus}
Choose the routed loop label on its hard metric line. The integrand in
\eqref{eq:v30-marked} vanishes unless
\begin{equation}
 \boxed{1/256\le|ap|\le1/8.}
 \label{eq:v30-annulus}
\end{equation}
Every lower covariance factor is one on this support. The integrand
is smooth periodic in the routed label and has the exact form
\begin{equation}
 f_a^\Omega(p,K,z)=\chi^{-1}
                   F^\Omega(ap,aK,a^2z).
 \label{eq:v30-scaling}
\end{equation}
All prescribed derivatives of the dimensionless function are bounded
on its fixed annulus, uniformly in the domain
\eqref{eq:v30-domain}. Consequently its per-volume coefficient has the
exact finite-sum representation
\begin{equation}
 \boxed{\iota_{a,L}(K,z)=\chi^{-1}a^{-4}
                  \mathcal F_{\boldsymbol\ell}(aK,a^2z),
       \qquad \ell_\nu=L_\nu/a.}
 \label{eq:v30-finite-scaling}
\end{equation}
The finite normalized sum is retained in \(\mathcal F_{\boldsymbol\ell}\).
\end{theorem}
\begin{proof}
Below the inner bound all routed labels differ from \(p\) or \(-p\)
by at most four soft momenta. They lie strictly on the unit plateau.
Every \((I-P)\) or \((P-I)\) in
\eqref{eq:v30-simplified} is then zero. Above the outer bound the hard
profiles in \(B_c\) vanish. At least one such factor occurs in every
term and two occur in the contact. The domain margins are far larger
than four soft transfers. They also place all remaining lines above
\(2\lambda\). The source support stays away from principal momentum
faces, so the inherited smooth completed symbols give periodic
regularity with no boundary discontinuity.

The source square has two derivatives. A ghost return \(C_cB_c\)
has net hard order zero. Thus \(A_\Omega\) and \(B_\Omega\) have
hard order two, each metric covariance contributes \(\chi^{-1}a^2\),
the scalar covariance contributes \(a^2\), and \(L_\phi\) has order
two with its mass term retained as \(a^{-2}(a^2z)\). Both traces have
net integrand order zero and one metric normalization \(\chi^{-1}\).
The free symbols and finite vertices are exactly functions of the
scaled momenta and mass. Rescaling proves \eqref{eq:v30-scaling}.
The physical mode sum contributes \(a^{-4}\), proving
\eqref{eq:v30-finite-scaling}. Uniform finite-sum bounds follow from
the fixed compact support and normalized mode counting, including all
periods with \(\lambda L_\nu\ge1\).
\end{proof}

\begin{lemma}[Parity and the derivative-free coefficient]
\label{lem:v30-parity}
Both \(\iota_{a,L}\) and the coefficient
\(\gamma^{\upsilon\upsilon}_{a,L}\) of
\eqref{eq:v30-parent} are even under simultaneous reversal of all
external momenta. Moreover, for every allowed squared mass,
\begin{equation}
 \boxed{\iota_{a,L}(0,z)=0,\qquad
        \gamma^{\upsilon\upsilon}_{a,L}(0,z)=0.}
 \label{eq:v30-constant-zero}
\end{equation}
No arbitrary source coefficient or metric-marked parity is covered by
this statement.
\end{lemma}
\begin{proof}
The free metric, scalar and matrix ghost inverses are even in their
momentum. The flat \(B_c\) has two derivatives; the scalar square has
two; and the actual scalar vertex \(L_\phi\) is even under simultaneous
reversal of both scalar momenta. For the latter, each term of the
sign-averaged Hodge symbol, its continued native correction and its
upper completion is a product of either two odd sine/sinh factors or
two even diagonal factors. The mass term is even. Thus the two traces
in \eqref{eq:v30-marked} are even before reversing the routed loop
label. The block \(N_\upsilon\) is odd, so its product in
\eqref{eq:v30-parent} is even. Odd finite grids admit the required
pairing \(p\leftrightarrow-p\) exactly. No oriented gravitational
interaction vertex enters either calculation.

At zero external momenta the contact in \eqref{eq:v30-simplified}
annihilates the constant scalar. For the cross term take one retained
ghost of zero momentum and a hard ghost of momentum \(l\); the hard
scalar has momentum \(-l\). In \eqref{eq:v30-square-symbol}, set
\(k=0,r=-l\). The two remaining contractions cancel, for arbitrary
hard ghost polarization and arbitrary \(\pi_a(l)\). The cancellation
holds for each mode and is independent of the massive scalar inverse.

For \eqref{eq:v30-parent}, the inherited constant-ghost action return
is \(\beta_p=\alpha_pV_px_p\), with \(V_p\) even and independent
of the constant odd ghost \(u\), and \(\alpha_p=i\,u\cdot p\).
The two source factors contain \(\alpha_p\alpha_{-p}
=-\alpha_p^2=0\), with the scalar-antifield labels moved only with
their Grassmann signs. This argument works species by species and for
every scalar mass. The no-extra-species-trace convention is unchanged.
\end{proof}

\section{Mass-regular local data and quantitative nonlocal limits}
\label{sec:v30-limits}

Use V28's mass-weighted Taylor operator \(\mathbf T_D\), with
\(\deg K=1\), \(\deg z=2\). The scalar algebra insertion has
subtraction degree four and the quadratic-antifield parent degree two:
\[
 \iota^R_{a,L}=(I-\mathbf T_4)\iota_{a,L},\qquad
 \gamma^{\upsilon\upsilon,R}_{a,L}
            =(I-\mathbf T_2)\gamma^{\upsilon\upsilon}_{a,L}.
\]
Both are even in \(K\), so respectively
\(\mathbf T_5=\mathbf T_4\) and
\(\mathbf T_3=\mathbf T_2\) on these coefficients.
Taylor derivatives in \(z\) are taken at zero on a hard domain with a
fixed positive lower scale. They do not replace the massive covariance.
In particular, the differentiated propagator is
\[
 \partial_z^j C_H=(-1)^j j!\,\nu_\lambda
                         \widetilde P_a^{-j-1},
\]
with one factor of \(\nu_\lambda\), not a power of that cutoff.

\begin{theorem}[Vanishing nonlocal scalar-curvature insertion]
\label{thm:v30-insertion-limit}
Let \(\epsilon=|K|+m\), \(d=|\alpha|+2j\), and let
\(\mathfrak n_I\) denote the product of external coefficient norms.
For \(d\le6\),
\begin{equation}
 \boxed{|\partial_K^\alpha\partial_z^j\iota^R_{a,L}(K,z)|
 \le C_{I,\alpha,j}\chi^{-1}\mathfrak n_I
                    a^2\epsilon^{6-d}.}
 \label{eq:v30-insertion-limit}
\end{equation}
Its continuum nonlocal limit is zero. Every fixed derivative with
\(d>6\) is bounded by
\(C\chi^{-1}\mathfrak n_I a^2\lambda^{6-d}\).
The local data are the actual finite sums
\begin{equation}
 \partial_K^\alpha\partial_z^j\iota_{a,L}(0,0)
  =\chi^{-1}a^{d-4}
      \partial_\kappa^\alpha\partial_\zeta^j
                       \mathcal F_{\boldsymbol\ell}(0,0).
 \label{eq:v30-local-data}
\end{equation}
They have only momentum degrees at least two. The retained weight-two
terms may grow as \(a^{-2}\); weight-four terms may have a finite
limit. No logarithmic scale sum is created by this fixed annulus.
\end{theorem}
\begin{proof}
Apply weighted Taylor expansion to the smooth dimensionless function
in \eqref{eq:v30-finite-scaling}. Evenness in \(K\) removes all odd
weighted degrees. After degree four, the first remaining degree is six.
All mass derivatives have uniform bounds because the dimensionless
scalar denominator is bounded away from zero on the annulus for
\(0\le a^2z\le M^2(a\lambda)^2\), including zero. The integral
Taylor remainder along \((K,z)\mapsto(tK,t^2z)\) therefore gives
\(a^{-4}(a\epsilon)^6=a^2\epsilon^6\), with the differentiated
versions stated. Above degree six, direct differentiated scaling gives
\(a^{d-4}\le a^2\lambda^{6-d}\).
The same scaling without subtraction proves \eqref{eq:v30-local-data}.
Equation \eqref{eq:v30-constant-zero} eliminates every pure mass
monomial; evenness eliminates linear momentum terms. Thus the only
possible local structures at these weights are \(K^2\), \(K^4\)
and \(zK^2\), with their full tensor/source index coefficients.
No coefficient integral has been evaluated or set to zero beyond these
proved restrictions.
\end{proof}

\begin{theorem}[Continuum control of the quadratic scalar-antifield parent]
\label{thm:v30-parent-limit}
Let \(\gamma^{\upsilon\upsilon,R}_{0,L}\) be the subtracted ordinary
zero-mesh trace on the same periodic volume and at the same lower scale.
For \(d=|\alpha|+2j\le4\),
\begin{equation}
 \boxed{|\partial_K^\alpha\partial_z^j
 (\gamma^{\upsilon\upsilon,R}_{a,L}
            -\gamma^{\upsilon\upsilon,R}_{0,L})|
 \le C_{I,\alpha,j}\chi^{-1}\mathfrak n_I
                    a^2\epsilon^{4-d}.}
 \label{eq:v30-parent-limit}
\end{equation}
There is no assertion that the limiting parent is zero. Higher fixed
derivatives have bound \(C\chi^{-1}a^2\lambda^{4-d}\).
Its entire local assignment consists of momentum-quadratic terms, with
coefficient bound
\begin{equation}
 \boxed{\|\mathbf T_2\gamma^{\upsilon\upsilon}_{a,L}\|_{\mathrm{coeff}}
       \le C\chi^{-1}\left(1+\log\frac1{a\lambda}\right).}
 \label{eq:v30-parent-local}
\end{equation}
\end{theorem}
\begin{proof}
Each \(N_\upsilon\) has hard order one, since its scalar source has
one derivative and its ghost return has net order zero. The parent
integrand has order \(-2\). A momentum derivative lowers order by one,
and a squared-mass derivative lowers it by two. In the central plateau,
only the improved free inverses distinguish it from the ordinary
integrand; the source and the cubic ghost-action operands there agree
exactly. The metric, ghost and scalar free comparisons are fourth-order
relative comparisons. Thus, with \(R=\lambda+|p|\),
\[
 |\partial^{(d)}f_0|\le C\chi^{-1}R^{-2-d},\qquad
 |\partial^{(d)}(f_a-f_0)|\le C\chi^{-1}a^4R^{2-d}.
\]
Here \(\partial^{(d)}\) denotes any stated weighted combination.
At \(d=4\) the first is \(R^{-6}\), and the second
\(a^4R^{-2}\). Their four-dimensional tail/interior sums are
respectively \(O(a^2)\) and \(O(a^2)\). The hard ghost profiles
make the remaining upper terms smooth periodic; rescaling them gives
the same \(a^2\) bound after four derivatives. The inherited full
scalar completion remains in the scalar inverse throughout.

The physical finite-volume sums obey the same integral majorants by
partitioning momentum space into boxes of mesh \(2\pi/L_\nu\) and
using the smooth hard floor \(\lambda L_\nu\ge1\). No spatial
volume factor appears in the per-volume coefficient. Evenness permits
using four derivatives after degree-two subtraction. Weighted integral
Taylor remainders prove \eqref{eq:v30-parent-limit}. Further derivatives
are integrable with the stated powers of \(\lambda\).
Finally, the derivative-free coefficients vanish for every \(z\) by
\cref{lem:v30-parity}; first momentum derivatives vanish by evenness.
Only second momentum derivatives remain locally, and their unsummed
bound is \(R^{-4}\), giving \eqref{eq:v30-parent-local}. This is an
upper bound, not a computed logarithmic coefficient.
\end{proof}

Both new panels have four external arguments. Multiplication by a fixed
smooth external window and Fourier inversion therefore give, for every
fixed moment \(b\),
\begin{align}
 a^{12}\sum_{x_1,x_2,x_3}
 (1+\mu\textstyle\sum_i d_{\rm per}(x_i,0))^b
      \|\mathscr K_{\iota^R}\|
 &\le C_b\chi^{-1}a^2(\mu+m)^6,\\
 a^{12}\sum_{x_1,x_2,x_3}
 (1+\mu\textstyle\sum_i d_{\rm per}(x_i,0))^b
      \|\mathscr K_{\gamma_a^R-\gamma_0^R}\|
 &\le C_b\chi^{-1}a^2(\mu+m)^4.
 \label{eq:v30-kernels}
\end{align}
To verify these estimates, rescale all twelve independent momenta by
\(\mu\), use the preceding differentiated bounds beyond order
\(b+12\), and integrate by parts in each Fourier variable. Derivatives
of the window are included. Periodize the resulting rapidly decreasing
kernel, using the periodic distance no larger than the distance of each
lift. Sampling with physical weights is uniformly summable on
\(a\mu\ll1\). These are multilinear estimates with one source in
\(L^1\) and three in \(L^\infty\), not extensive supremum norms.

For the marked insertion only a bounded number of ultraviolet shell
levels contribute, by \eqref{eq:v30-annulus}. For the unmarked parent,
assign every word to the largest label among its metric, scalar and both
ghost covariances. The subtracted increment is bounded by
\(C\chi^{-1}(\mu+m)^4\Lambda_j^{-2}\). This is summable. Omitting
the two ghost labels inside \(N_\upsilon\) changes this assignment.
Thermodynamic comparison is independent of mesh comparison: smooth
lower-cutoff integrands give an error
\(C_b\chi^{-1}\epsilon^{4-d}\lambda^{-2}
 (\lambda L_{\min})^{-b}\), \(b>4\), for the parent. The marked
annular local data have their usual finite-sum to annular-integral
comparison in \(L_\nu/a\). No lower-scale limit is taken.

\section{The exact two-ghost descendant of V29's counterterm}
\label{sec:v30-descendant}

The particular scalar source counterterm of V29 has the form
\[
 C_U=\langle\upsilon,Q_c\phi\rangle,\qquad
 Q_{v_k}\phi_p=i\,U_v(k,p,z)\phi_{k+p},\qquad
 U_v=v^\nu U_\nu,
\]
where \(U\) is the actual output of its polynomial compiler and has
weighted degree at most three. No value for an unevaluated native
coefficient is supplied in this notation. Use the ordinary source
convention \(s\phi=\mathcal L_c\phi\), \(sc=c\cdot\partial c\).

\begin{theorem}[A finite polynomial compiler for the induced algebra row]
\label{thm:v30-descendant}
The scalar-antifield/two-ghost part of the variation of \(C_U\), at
zero metric, is \(\langle\upsilon,\mathcal N_U(c)\phi\rangle\),
where
\begin{equation}
 \boxed{\mathcal N_U(c)
       =Q_{c\cdot\partial c}-\mathcal L_cQ_c-Q_c\mathcal L_c.}
 \label{eq:v30-NU}
\end{equation}
Its coefficient on two even ghost tests \((k,v),(l,w)\) and scalar
momentum \(p\) is
\begin{equation}
\begin{split}
 \mathfrak d_U={}&-(v\cdot l)U_w(k+l,p,z)
              +(w\cdot k)U_v(k+l,p,z)\\
 &+(v\cdot(p+l))U_w(l,p,z)
                 -(v\cdot p)U_w(l,p+k,z)\\
 &-(w\cdot(p+k))U_v(k,p,z)
                 +(w\cdot p)U_v(k,p+l,z).
 \label{eq:v30-NU-symbol}
\end{split}
\end{equation}
It is antisymmetric in the two ghost tests, polynomial in \(z\), and
has weighted degree at most four. For fixed coefficient \(\ell^1\)
norms and unit ghost polarizations it obeys the conservative bound
\begin{equation}
 \boxed{\|\mathfrak d_U\|_{\mu,1}
           \le384\mu\|U\|_{\mu,1}.}
 \label{eq:v30-NU-bound}
\end{equation}
It vanishes on two constant ghost tests, but not in general.
\end{theorem}
\begin{proof}
Deform the scalar source by \(\varepsilon Q_c\), keeping the ordinary
ghost source fixed. The coefficient linear in \(\varepsilon\) of its
square is precisely \eqref{eq:v30-NU}. Equivalently, the source part
of the master antibracket is obtained by polarization of the source
Hamiltonian's square, in the convention used for
\eqref{eq:v30-insertion}. Substitution of the Fourier formula for
\(Q\), and of
\([v_k,w_l]=i((v\cdot l)w-(w\cdot k)v)_{k+l}\), gives
\eqref{eq:v30-NU-symbol}.

Every shifted polynomial has degree at most three. Substitution of a
sum for any group of momentum variables increases its monomial
coefficient norm by at most \(2^3\). A contracted momentum sum costs
at most \(8\mu\) in four dimensions with the component norms used
here. The six terms therefore give the safe bound
\(6\cdot8\cdot8\mu=384\mu\). Mass variables are not substituted
or divided. Antisymmetry follows term by term, or from the ordered
exterior extraction. Setting \(k=l=0\) makes the two shifted pairs
cancel. The formula is the negative ordinary Lie-algebra cochain
differential of \(Q\); its next cochain differential is zero by the
Jacobi identity and the representation identity of \(\mathcal L\).
This last statement is about ordinary sources, not the filtered finite
source's nilpotence.
\end{proof}

For an explicit example already admissible in V29's input family, take
\(U_1(k,p,z)=p_0^3\) and all other components zero. With
\(v=e_0\), \(k=te_0\), \(w=e_1\), \(l=0\), formula
\eqref{eq:v30-NU-symbol} gives
\begin{equation}
 \boxed{\mathfrak d_U
        =-3t p_0^3-3t^2p_0^2-t^3p_0\ne0.}
 \label{eq:v30-descendant-example}
\end{equation}
This is a counterterm diagnostic, not a computed native loop value.
It shows why solving the action row does not make the accompanying
algebra row vanish. It does not constitute an anomaly or an obstruction
to a further coupled completion.

For the actual finite source, the same derivative before setting the
filters to one is
\[
 \mathcal N_{a,U}
       =Q_{b_c}-R_a^H(c)Q_c-Q_cR_a^H(c),\qquad u=Pc.
\]
On the soft four-leg panel all composite labels in this local
counterterm remain on the common plateau, so
\(\mathcal N_{a,U}=\mathcal N_U\) there exactly. This statement
is not used on hard loop labels. An order-\(\hbar\) counterterm
changes the one-loop defect by this tree insertion; its insertion into
one of the new loops would be of order \(\hbar^2\). Hence the
source-square part of the updated one-loop local coefficient is
\begin{equation}
 \boxed{B_{\Omega,H}^{\rm new}
       =\mathbf T_4\iota_{a,L}+\mathfrak d_U,\qquad
       (I-\mathbf T_4)(\iota_{a,L}+\mathfrak d_U)=\iota^R_{a,L}.}
 \label{eq:v30-updated-insertion}
\end{equation}
The sign of \(U\) here is its sign in the inserted action; for V29's
particular normalized update it includes its actual subtraction of the
old quantum local source. Formula \eqref{eq:v30-updated-insertion}
does not assert that the remaining local sum is zero.

\section{The parent return and the precise remaining scalar algebra equation}
\label{sec:v30-row}

The new parent supplies a second contribution that no adjustment of
\(U\) may discard. Its free scalar-Euler return is
\begin{equation}
 \boxed{\mathcal E_a^{\upsilon\upsilon}
 =\sum_A\int(\widetilde P_a\phi_A)
       \frac{\delta^L\mathfrak G_{a,1}^{\upsilon\upsilon}}
                       {\delta\upsilon_A}.}
 \label{eq:v30-Euler-return}
\end{equation}
The signed scalar Euler convention is the one in V29's exact row map.
Left antifield differentiation fixes the sign for each ordered source
panel. In particular \(\mathcal E_a^{\upsilon\upsilon}\) has type
\(\upsilon\phi cc\); it is not removed by setting the scalar
on shell inside an off-shell identity.

\begin{proposition}[Subtraction and transport of the parent return]
\label{prop:v30-parent-return}
For the ordinary Euler multiplier \(P_0=k^2+z\),
\begin{equation}
 \boxed{\mathbf T_4(P_0\gamma^{\upsilon\upsilon})
              =P_0\mathbf T_2\gamma^{\upsilon\upsilon}.}
 \label{eq:v30-Euler-projector}
\end{equation}
The same local identity holds with the finite Euler multiplier in the
left side, since its difference from \(P_0\) starts at weighted degree
six. The nonlocal same-volume return comparison has bound
\(C\chi^{-1}a^2\epsilon^{6-d}\) for weighted derivative degree
\(d\le6\). In addition the exact finite subtraction retains the
product
\[
 (\widetilde P_a-P_0)\mathbf T_2
                 \gamma^{\upsilon\upsilon}_{a,L},
\]
whose fixed-window norm is bounded by
\begin{equation}
 C\chi^{-1}a^4\epsilon^8
             \left(1+\log\frac1{a\lambda}\right).
 \label{eq:v30-parent-small-product}
\end{equation}
\end{proposition}
\begin{proof}
The ordinary multiplier is homogeneous of weight two, so the identity
is V28's proven weighted multiplication rule, applied to the new
parent. The actual free discrepancy is \(O(a^4|k|^6)\); no degree
at most four can result from it. Multiplying
\eqref{eq:v30-parent-limit} by the ordinary Euler polynomial gives
the stated return estimate. Multiplication of the continuum parent
remainder by the finite discrepancy is smaller on the domain and has
the same fixed-derivative bound. The remaining exact product is bounded
using \eqref{eq:v30-parent-local}. Since the local parent has no
constant or pure-mass term, its first possible momentum degree is two,
giving \eqref{eq:v30-parent-small-product}. This product is retained;
its smallness is proved rather than inferred from classical consistency
alone.
\end{proof}

For orientation of the next calculation, let \(Z_H(c)\phi\) be
V28's scalar quantum transformation, let \(Z_{H,1}(h;c)\phi\) be its
first metric derivative, and let \(Y_g(c,c)\) be the retained quantum
ghost transformation from the pure-gravity minimal source module. In
ordinary zero-mesh coordinates, the minimal canonical coefficient of
\(\upsilon\phi cc\) is assembled from
\begin{equation}
 \boxed{\mathcal A_H^{\rm can}
   =\mathcal N_{Z_H}(c)\phi
     +\mathcal P_{R_0}Z_{H,1}(c,c)\phi
     +\mathcal L_{Y_g(c,c)}\phi
     +\mathcal E_0^{\upsilon\upsilon}.}
 \label{eq:v30-canonical-row}
\end{equation}
Here \(\mathcal P_{R_0}Z_{H,1}\) is the ordered two-ghost
polarization obtained by varying the one retained metric argument by
\(R_0c\). Formula \eqref{eq:v30-NU} defines
\(\mathcal N_{Z_H}\) equally for a nonlocal kernel. The four terms
are fixed by differentiating the ordinary master antibracket: scalar
representation transport, free metric variation of the scalar source,
the ghost-bracket return into the tree scalar source, and the parent
Euler return. At this scalar degree the stress replacement of a metric
antifield would require a further scalar pair and does not contribute.
The uneliminated nonminimal convention must be restored before making
this minimal extraction, as in V16 and V28.

All canonical operands of \eqref{eq:v30-canonical-row} are now specified
by actual measured coefficients. This statement does not replace its
right-hand side by \(\mathfrak I_{a,1}^H\). The latter is only the
scalar-source-square part of the breaking insertion. The complete
measured equation also contains action-breaking insertions differentiated
with respect to the scalar antifield, the actual lower-reference trace,
and the indicated nonminimal/measure bookkeeping. In particular,
source differentiation of an antifield-free action-breaking insertion
need not be zero. No estimate for that complete combination is inferred
from annular localization of one of its operands.

After the V29 counterterm is inserted, the \(Z_H\) part changes by
\(Q_c\); the resulting algebra descendant is exactly
\eqref{eq:v30-NU}. The counterterm is at most linear in
\(\upsilon\), so it does not itself alter the parent
\eqref{eq:v30-parent} at order \(\hbar\). The local first action
row stays fixed only if the subsequent algebra counterterm and its
\(q\phi^2c\) descendants are solved together. The remaining local
problem is therefore the common cancellation of
\(\mathbf T_4\mathcal A_H^{\rm can}\) against its complete
measured reference/breaking terms, with
\(\mathbf T_4\iota+\mathfrak d_U\) and
\(P_0\mathbf T_2\gamma^{\upsilon\upsilon}\) included at their
proper positions. It is not another application of V29's
\(c\phi^2\) right inverse to an unrelated row.

\section{Final ledgers for V30}
\label{sec:v30-ledgers}

\subsection*{Proved}
\addcontentsline{toc}{subsection}{V30: Proved}
The scalar square has the exact finite symbol
\eqref{eq:v30-square-symbol}. Its complete marked one-loop coefficient
is the contact-minus-mixed-return expression
\eqref{eq:v30-marked}, with its two distinct ghost-propagator counts.
Its assembled integrand has fixed ultraviolet-annulus support, parity
and an exactly vanishing constant-momentum coefficient for all masses
in the declared domain. Its degree-four-subtracted remainder tends to
zero with \eqref{eq:v30-insertion-limit}. The new quadratic
scalar-antifield parent is \eqref{eq:v30-parent}; its degree-two local
assignment, logarithmic upper bound, full-cutoff comparison and Euler
return are specified with the displayed source norms. The induced
algebra variation of V29's actual source counterterm is compiled by
\eqref{eq:v30-NU-symbol}, with a bounded pole-free polynomial map and
an explicit nonzero example. These results identify operands of the
next equation and do not claim its complete local cancellation.

\subsection*{Inherited}
\addcontentsline{toc}{subsection}{V30: Inherited}
V1--V29 are retained, not independently re-certified. The continued,
periodically completed scalar action and simultaneous reference are V27;
the source and triangle are V28; the first mixed local primitive is
V29. The metric/ghost inverses, actual ghost action, source ordering,
finite lower-cutoff convention, contour, and coordinate/Haar amplitude
are unchanged. The Gaussian, Berezin, annular Fourier and weighted
Taylor mechanisms are inherited tools; the scalar marked operands,
parent, parity restrictions and counterterm descendant are the new
applications established here. No new classical Einstein or structural
Standard-Model reconstruction is asserted.

\subsection*{Literature input}
\addcontentsline{toc}{subsection}{V30: Literature input}
The distinction between action and algebra consistency, antifield
returns, and full local BRST cohomology is established in
\cite{V29BBHMatter}. The distinction between a fixed-scale modified
reference identity and an unmodified master equation is retained from
\cite{V28IIM}. Neither is imported as a proof that the actual combined
scalar row is restored. The matrix and polynomial formulas used here
are proved directly; no general first result on scalar--gravity
quantization or anomaly cancellation is claimed.

\subsection*{Conditional}
\addcontentsline{toc}{subsection}{V30: Conditional}
The native realization consumes the prior completed-symbol and contour
bounds at their stated scope. The row display
\eqref{eq:v30-canonical-row} becomes a normalized modified identity only
after its action-breaking, reference and nonminimal terms are assembled
and their common local assignment is proved compatible. A local
counterterm cancelling the remaining scalar algebra coefficient must
preserve V29's action row together with all its descendants. Native
local constants are still defined by finite sums or convergent
subtracted integrals; no numerical integral value is reported.

\subsection*{Open}
\addcontentsline{toc}{subsection}{V30: Open}
The full updated \(\upsilon\phi cc\) equation and its coupled
\(q\phi^2c\) normalization are not solved here. The scalar square's
nonlocal suppression is not a proof of vanishing of the full quantum
breaking. A general mixed-source invariant shell class, internal-gauge
and chiral determinant-line transport in the common gravitational
reference, arbitrary-loop EFT control, and removal of the lower scale
remain separate tasks. No positive Lorentzian measure or finite-parameter
ultraviolet completion follows.

\begin{center}\small
\begin{tabular}{>{\raggedright\arraybackslash}p{.24\textwidth}>{\raggedright\arraybackslash}p{.67\textwidth}}
\toprule
Variable & Scope in V30\\\midrule
Volume & Exact finite identities on odd periodic grids. Per-volume symbols
and twelve-relative-coordinate rooted kernels have no total-volume
factor. Thermodynamic comparison is separate.\\
Cutoff & All completed modes retained. Marked support is exactly
\(1/256\le|ap|\le1/8\). Rates require \eqref{eq:v30-domain}.\\
Mass & Weighted jets in \(z=m^2\); \(0\le m\le M\lambda\), including
zero, at fixed \(\lambda>0\). No inverse mass.\\
Sources & \(\upsilon\phi cc\) insertion and \(\upsilon^2cc\)
parent, at zero retained metric. No metric-marked parity extension.\\
Couplings/order & One loop, fixed \(\chi>0\), zero internal gauge,
Yukawa and scalar self-couplings. No higher-loop uniformity.\\
Transport & Bounded marked ultraviolet shell count and summable parent
complements; this is not an analytic interacting RG self-map.\\
\bottomrule
\end{tabular}
\end{center}

The companion checker \srcid{check_ESM_scalar_algebra_V30.py} verifies
finite symbol, exterior-sign, marked-block and polynomial-descendant
identities, weighted Taylor bookkeeping, and preservation of the prior
mathematical body. Its certificate separates these exact checks from
the analytic estimates above and the unevaluated native loop constants.
No new proof-assistant formalization is asserted.


\clearpage
\part*{V31: The scalar Euler compatibility complex and a mass-regular algebra counterterm}
\addcontentsline{toc}{part}{V31: Scalar Euler compatibility and a relative algebra counterterm}

\section{Go upstream from the scalar square to its coupled consistency equation}
\label{sec:v31-start}

V1--V30 and their mathematical labels are retained. V30 has already
computed the scalar source-square insertion, its nonlocal estimate, the
quadratic scalar-antifield parent, and the descendant of V29's source
counterterm. None is recalculated here as a new loop theorem. The remaining
question is whether a \emph{combined} local scalar-algebra coefficient can
be removed without undoing the normalized one-ghost/two-scalar action row.
There is a useful algebraic answer that is stronger than another isolated
component calculation, but weaker than an assertion that the entire mixed
Slavnov identity has been restored.

The answer uses the scalar Euler operator that is already in the native
zero-mesh action. A local scalar-algebra coefficient compatible with its
off-shell quadratic action has a unique factorization through that Euler
operator. At the weight retained in V28--V30, its factor is obtained simply
by taking the coefficient of the squared mass and antisymmetrizing the two
scalar momenta. That factor gives a quadratic scalar-antifield counterterm.
The latter has no antifield-free variation, so it leaves the previously
normalized action row unchanged. For an arbitrary, not necessarily
compatible input, the same map removes exactly the Euler-compatible part
and returns an explicit remainder, together with a bound in terms of the
actual consistency defect.

This does not prove that the still-unassembled native remainder is zero.
Instead it identifies its dependence on the one-metric action row and the
already retained gravity-algebra row. It changes the order of attack: once
those action-side data supply the necessary compatibility, the scalar
algebra normalization is a fixed bounded coefficient operation, not another
independent cohomological assumption.

\subsection{Inherited branch and new scope}

Keep four equal-mass real scalar components, zero internal gauge, Yukawa
and scalar self-couplings, and the gravitational reference and source
conventions of V27--V30. Work on the flavour-diagonal coefficient supplied
by that branch; a fixed external scalar species is enough, and summing
identical diagonal components does not introduce a new closed scalar loop.
All quantum coefficients discussed below have their common factor
\(\hbar\) removed. No new integration contour, action vertex, filter,
propagator, or reference scale is introduced.

The new proofs are finite polynomial statements in \(z=m^2\) and ordinary
external momenta. Their inputs are the actual local Taylor coefficients of
the previous measured calculation, whenever that complete array is
assembled. In particular, this is \emph{not} a statement that the isolated
coefficient \(\mathbf T_4\iota+\mathfrak d_U\) in
\eqref{eq:v30-updated-insertion} is the whole breaking.

The use of the equations of motion through the Koszul--Tate part of BV is
standard \cite{V20BBH1994}; the full longitudinal differential alone is not
the relevant object. Here the finite scalar-pair kernel, its explicit
coefficient section, its norm and its relation to the existing mixed row
are proved directly. No general scalar--gravity anomaly classification is
imported, and the pure-gravity classification used in V25 is not reused as
a theorem about this enlarged complex.

\section{The exact scalar-pair Euler compatibility map}
\label{sec:v31-syzygy}

Let \(k,l\) be the two ghost momenta, \(p\) the scalar momentum and
\(r=-p-k-l\) the scalar-antifield momentum. Put
\begin{equation}
 \begin{gathered}
 \mathscr R=\R[k_0,\ldots,k_3,l_0,\ldots,l_3,p_0,\ldots,p_3,z],\\
 P=p^2+z,\qquad Q=(p+k+l)^2+z,\\
 \iota F(k,l,p,z)=F(k,l,-p-k-l,z).
 \end{gathered}
 \label{eq:v31-ring}
\end{equation}
Thus \(\iota^2=I\) and \(\iota P=Q\). We assign weight one to
momentum and weight two to \(z\). The same statements hold over \(\C\)
or over a fixed characteristic-zero coefficient field containing the
renormalized parameters. Ghost polarization indices are retained as
coefficient indices. All divisibility arguments below are applied to
their ordinary polynomial components, never to a Grassmann variable.

Let \(D(k,l,p,z;v,w)\) be the coefficient of the local functional
\begin{equation}
 \mathcal B_D=\int_{k+l+p+r=0}
       \upsilon(r)c(k)c(l)D(k,l,p,z)\phi(p).
 \label{eq:v31-row-functional}
\end{equation}
This notation means the ordered, polarized two-ghost coefficient. It is
antisymmetric under exchange of the two ghost tests. Antifield derivatives
have been moved onto the other factors by integration by parts, with
momentum conservation imposed only by eliminating \(r\). The two
transfers \(k,l\) remain independent variables.

The free scalar Euler differential is normalized by
\(\delta_H\upsilon(p)=(p^2+z)\phi(p)\),
\(\delta_H\phi=\delta_Hc=0\). Then
\begin{equation}
 \delta_H\mathcal B_D
 =\frac12\int_{k+l+p+r=0}
       c(k)c(l)\phi(r)\phi(p)\,\mathscr J D,
 \qquad
 \boxed{\mathscr J D=QD+P\iota D.}
 \label{eq:v31-J}
\end{equation}
The factor one half records symmetrization of the two identical scalar
legs. This is an off-shell identity, not evaluation of a propagator on
an energy shell.

\begin{theorem}[Complete polynomial Euler kernel]
\label{thm:v31-syzygy}
For a polynomial \(D\), the following are equivalent:
\begin{enumerate}
\item \(\mathscr J D=0\) as a polynomial in all the variables in
\eqref{eq:v31-ring};
\item there is a polynomial \(E\) with
\begin{equation}
 \boxed{D=PE,\qquad \iota E=-E.}
 \label{eq:v31-factor}
\end{equation}
\end{enumerate}
The polynomial \(E\) is unique. If \(D\) has weight at most four,
then \(E\) has weight at most two, is independent of \(z\), and has
no constant term in the centred scalar momentum
\(x=p+(k+l)/2\). The construction commutes with ghost exchange and
with simultaneous reversal of all momenta.
\end{theorem}
\begin{proof}
The quotient \(\mathscr R/(P)\) is the polynomial domain
\(\R[k,l,p]\), obtained by the substitution \(z=-p^2\).
Consequently \(P\) is prime. In that quotient,
\[
 Q-P=2p\cdot(k+l)+(k+l)^2
\]
is a nonzero polynomial. Reducing \(QD+P\iota D=0\) modulo \(P\)
therefore gives \(D=0\) in the quotient: \(P\) divides every ordinary
component of \(D\). Write \(D=PE\). Applying \(\iota\) gives
\(\iota D=Q\iota E\), so the equation is
\(PQ(E+\iota E)=0\). This polynomial domain has no zero divisors,
hence \(\iota E=-E\). The converse follows by substitution. Uniqueness
uses the same domain property.

The Euler factor is homogeneous of weight two, so its nonzero product
raises the largest homogeneous weight by exactly two. Thus
\(\operatorname{wt}E\le2\). In the variables \((k,l,x,z)\), the
involution fixes \(k,l,z\) and sends \(x\) to \(-x\). Its odd
polynomials of weight at most two contain exactly one \(x\), multiplied
by a polynomial of momentum degree at most one. A term containing \(z\)
would have weight at least three, and is impossible. Both exchange of the
ghost arguments and total momentum reversal commute with the two Euler
factors and with \(\iota\), giving the last assertion.
\end{proof}

The substitution \(z=-p^2\) was used only in a polynomial quotient. No
native loop is evaluated at negative mass squared, and no factor
\(1/(p^2+z)\) is introduced in a counterterm. In particular the proof is
regular at the physical massless point. It does require the mass Taylor
coefficient as part of the recorded local input; a massless numerical
value alone does not supply that coefficient.

\begin{remark}[The full transfer variables are essential]
Take \(D=p_0\) times an antisymmetric constant two-ghost tensor.
On the subspace \(k+l=0\), the two Euler factors coincide and
\(\mathscr J D=0\) there. Globally \(\mathscr J D\ne0\), and
\(D\) is not divisible by \(P\). Thus constant-total-transfer tests
or a single on-shell momentum do not establish
\cref{thm:v31-syzygy}. Its norm estimates below are coefficient-jet
estimates on the full independent panel, not pointwise estimates near
this degenerate subspace.
\end{remark}

\section{A mass-coefficient projector and the exact finite component count}
\label{sec:v31-projector}

For arbitrary \(D\) of weight at most four, define
\begin{equation}
 \begin{gathered}
 D_1=[z]D=\partial_zD\big|_{z=0},\qquad
 \boxed{\mathfrak h_ED=\tfrac12(D_1-\iota D_1),}\\
 \Pi_ED=P\mathfrak h_ED,\qquad
 \mathfrak r_ED=D-P\mathfrak h_ED.
 \end{gathered}
 \label{eq:v31-projector}
\end{equation}
All three are explicit linear operations on the actual local coefficients.
The subscript distinguishes this map from V20's longitudinal primitive;
the latter resolves a different differential.

\begin{theorem}[A bounded projection onto the removable scalar-algebra sector]
\label{thm:v31-projector}
The image of \(\Pi_E\) is precisely \(\ker\mathscr J\), and
\begin{equation}
 \boxed{\Pi_E^2=\Pi_E,\qquad
 \mathfrak h_E\mathfrak r_E=0,\qquad
 \mathscr J\mathfrak r_E=\mathscr J.}
 \label{eq:v31-projector-identities}
\end{equation}
In particular \(\mathfrak r_ED=0\) if and only if
\(\mathscr J D=0\). In the expanded coefficient norm
\[
 \|F\|_{\lambda,1}
   =\sum_{\alpha,j,\text{polarizations}}
      |F_{\alpha j}|\lambda^{|\alpha|+2j},\qquad\lambda>0,
\]
one has
\begin{equation}
 \boxed{\begin{aligned}
 \|\mathfrak h_ED\|_{\lambda,1}
       &\le5\lambda^{-2}\|D\|_{\lambda,1},\\
 \|\Pi_ED\|_{\lambda,1}&\le25\|D\|_{\lambda,1},\\
 \|\mathfrak r_ED\|_{\lambda,1}&\le26\|D\|_{\lambda,1}.
 \end{aligned}}
 \label{eq:v31-projector-bounds}
\end{equation}
No parity restriction or consistency assumption is required for these
three maps.
\end{theorem}
\begin{proof}
The mass coefficient \(D_1\) has ordinary momentum degree at most two.
Thus \(E=\mathfrak h_ED\) is mass independent, has weight at most two
and is odd under \(\iota\). Theorem~\ref{thm:v31-syzygy} gives
\(\mathscr J(PE)=0\). Conversely, on the kernel its unique factor
is mass independent, so \([z](PE)=E\) and
\(\mathfrak h_E(PE)=E\). This proves the image, idempotence and all
identities in \eqref{eq:v31-projector-identities}.

Extracting \([z]\) loses exactly two powers of \(\lambda\).
Replacing each \(p_i\) by \(-p_i-k_i-l_i\) increases a monomial's
coefficient norm by at most \(3^2=9\), because \(D_1\) has degree
at most two. Hence half the identity-minus-involution has norm at most
five on that array. The coefficient norm of \(P\) is
\(5\lambda^2\), and the triangle inequality gives the remaining
bounds. Fixed real/conjugation conditions can be imposed afterward;
the maps commute with their momentum involution.
\end{proof}

\subsection{What the quadratic-antifield kernel can contain}

In position space every scalar, formally skew-adjoint \(E\) in this
weight range has the normal form
\begin{equation}
 \begin{split}
 E_c&=a^\rho(c)\partial_\rho
                 +\tfrac12\partial_\rho a^\rho(c),\\
 a^\rho(c)&=
 \sum_{\mu<\nu}b^{\rho\mu\nu}c^\mu c^\nu
 +\sum_{\mu,\nu,\sigma}t^{\rho\mu\nu\sigma}
                    c^\mu\partial_\sigma c^\nu .
 \end{split}
 \label{eq:v31-normal-form}
\end{equation}
The coefficients may depend on the cutoff, lower scale and declared
couplings, but not on a negative power of the mass introduced by this
map. The divergence term in \eqref{eq:v31-normal-form} is required by
formal adjointness; this formula is not a change of the physical scalar
into a half-density. Fourier factors of \(i\) are assigned by the
inherited reality convention.

Indeed, centred scalar momentum is odd and appears exactly once. The
remaining factor has either zero or one ghost derivative. In an operator
presentation a scalar even-order principal symbol cannot be formally
skew; its apparent second-order terms reduce to the displayed first-order
operator with a differentiated ghost coefficient. There is no independent
zeroth-order scalar skew multiplication.

\begin{proposition}[Exact sizes of this scalar Euler kernel]
\label{prop:v31-counts}
For one flavour-diagonal scalar and four diffeomorphism ghosts, the rooted
\(\upsilon\phi cc\) polynomial bank of weight at most four has
\(14976\) components. The kernel of \(\mathscr J\) has dimension
\(280=24+256\), and its rank is \(14696\). If the input is even
under simultaneous momentum reversal, the corresponding numbers are
\begin{equation}
 \boxed{\begin{gathered}
 \dim D_{\mathrm{even}}=11960,\qquad
 \dim\ker\mathscr J_{\mathrm{even}}=256,\\
 \operatorname{rank}\mathscr J_{\mathrm{even}}=11704.
 \end{gathered}}
 \label{eq:v31-counts}
\end{equation}
Every nonzero even compatible coefficient has weight exactly four and is of the
form \((p^2+z)E^{[2]}\). These numbers are ranks of the displayed
Euler map, not local BRST cohomology or anomaly counts.
\end{proposition}
\begin{proof}
There are \(a_d=4\binom{d+3}{3}\) ghost jets of derivative order
\(d\). If \(a(t)=4(1-t)^{-4}\), the two-ghost exterior-pair series is
\[
 g_2(t)=\tfrac12\{a(t)^2-a(t^2)\}.
\]
Multiplication by \((1-t)^{-4}\) for scalar derivatives and by
\((1-t^2)^{-1}\) for mass powers gives the component counts
\[
 (6,88,602,2928,11352)
\]
at weights zero through four. Their sum is \(14976\); the even sum
is \(11960\). For the kernel use the invertible centred coordinate
\(x=p+(k+l)/2\). An odd polynomial of weight at most two contains one
of its four components. At degree one it multiplies one of the six
undifferentiated ghost pairs, giving \(4\cdot6=24\). At degree two
it multiplies one of the \(4\cdot16=64\) pairs consisting of one
undifferentiated ghost and one first ghost jet, giving \(256\).
These centred monomials are linearly independent. Multiplication by
\(P\) is injective and raises weight by two. Thus only the latter
256 survive total evenness. Rank-nullity completes the argument; no
floating-point rank or uncomputed cohomology is used.
\end{proof}

The even specialization is available for V30's isolated marked insertion
and its quadratic-antifield parent, whose parity was proved there. It is
\emph{not} silently applied to the complete finite mixed residual: the
one-metric quantum transformation can contain oriented native artifacts.
The general 280-dimensional kernel and the bounds in
\cref{thm:v31-projector} require no such extra claim.

\section{A quantitative reconstruction of the unremoved consistency defect}
\label{sec:v31-recovery}

The remainder \(\mathfrak r_ED\) is not just bounded by the original
array. It is determined by \(\mathscr J D\). A direct polynomial
reconstruction gives a uniform bound without forming a large matrix or
assuming the coefficient is closed.

Write \(A=p^2\), \(B=(p+k+l)^2\), and expand
\(D=D_0+zD_1+z^2D_2\). The last coefficient is momentum independent.
For \(W=\mathscr J D=\sum_{j=0}^3z^jW_j\), define
\begin{align}
 d_2&=W_3/2,\\
 s_1&=\tfrac12\{W_2-(A+B)d_2\},\\
 v_0&=\tfrac12\{W_1-(A+B)s_1\},\\
 u_0&=\operatorname{Div}_{k_0,\,B-A}
                    \{W_0-(A+B)v_0\}.
 \label{eq:v31-recover-ingredients}
\end{align}
The last line denotes \emph{monic polynomial division} in \(k_0\),
not evaluation of the quotient of two momentum functions. The divisor
\(B-A\) is monic of degree two in \(k_0\). On the indicated image
its polynomial remainder is zero.

\begin{theorem}[Consistency controls the entire projected residual]
\label{thm:v31-recovery}
For every \(D\) of weight at most four,
\begin{equation}
 \boxed{\mathfrak r_ED=v_0+u_0+zs_1+z^2d_2.}
 \label{eq:v31-recovery}
\end{equation}
The quotient in \eqref{eq:v31-recover-ingredients} is exact, and
\begin{equation}
 \boxed{\|\mathfrak r_ED\|_{\lambda,1}
       \le2^{40}\lambda^{-2}
                       \|\mathscr J D\|_{\lambda,1}.}
 \label{eq:v31-consistency-bound}
\end{equation}
The constant is a deliberately conservative finite coefficient bound,
independent of mass, mesh, volume and the numerical loop coefficients.
It is not asserted to be a sharp condition number.
\end{theorem}
\begin{proof}
Set \(e=(D_1-\iota D_1)/2\), \(s=(D_1+\iota D_1)/2\), and
\(R_0=D_0-Ae\). The coefficients of \(W\) give
\[
 W_3=2D_2,\qquad
 W_2=2s+(A+B)D_2,\qquad
 W_1=R_0+\iota R_0+(A+B)s.
\]
Thus \(d_2=D_2\), \(s_1=s\), and
\(v_0=(R_0+\iota R_0)/2\). At degree zero in \(z\),
\[
 W_0=BR_0+A\iota R_0
     =(A+B)v_0+(B-A)(R_0-\iota R_0)/2.
\]
Hence the monic division returns
\(u_0=(R_0-\iota R_0)/2\) with zero remainder. This proves
\eqref{eq:v31-recovery}. In particular the computation never has to
invert \(B-A\) at a physical zero of that polynomial.

For the bound rescale momenta by \(\lambda\) and \(z\) by
\(\lambda^2\). The coefficient norms of \(A+B\) and \(B-A\)
are at most 40 and 32. Mass extraction successively gives, writing
\(w=\|W\|_{\lambda,1}\),
\[
 \|d_2\|\le\tfrac12\lambda^{-6}w,\quad
 \|s_1\|\le10.5\lambda^{-4}w,\quad
 \|v_0\|\le210.5\lambda^{-2}w.
\]
The numerator defining \(u_0\) therefore has norm at most \(8421w\).
After its monic \(k_0^2\) term is removed, the dimensionless divisor
has coefficient norm at most 31. The numerator has \(k_0\)-degree at
most six; at most five descending coefficient eliminations suffice.
Using the safe growth factor 32 at each elimination bounds the quotient
norm by \(32^5\lambda^{-2}\) times the numerator norm. The sum of
all four output norms is less than
\((8421\cdot32^5+222)\lambda^{-2}w<2^{40}\lambda^{-2}w\).
This establishes the displayed bound componentwise and hence after
summing polarization indices.
\end{proof}

This reconstruction also defines a verification test for supplied exact
arrays. A nonzero polynomial remainder in the division means that the
supplied \(W\) was not the asserted \(\mathscr J D\), or that the
coefficient conventions do not match; it must not be rounded away. The
much smaller bounds \(5,25,26\) suffice for applying the counterterm
itself. The larger bound is used only to transfer a separately proved
consistency estimate to the unremoved sector.

\section{Which mixed equations actually imply scalar Euler compatibility?}
\label{sec:v31-coupled}

Let \(\mathcal A\) be a complete ghost-number-one ordinary residual
functional in the inherited scalar--gravity BV convention. It can be the
continuum modified residual after its reference and coordinate-measure
terms have been assigned, or a finite-cutoff comparison record tested by
the ordinary differential. Do not assume that \(s_\infty\mathcal A=0\)
unless that consistency identity has been established for the same record.

A useful bookkeeping fact follows directly from the common action. Count
each scalar and scalar antifield as one matter argument and all other
fields and antifields as zero. On the zero-self-coupling branch,
\begin{equation}
 \boxed{s_\infty=d_0+d_2,\qquad
 d_0^2=0,\quad d_0d_2+d_2d_0=0,\quad d_2^2=0.}
 \label{eq:v31-matter-filtration}
\end{equation}
Here \(d_0\) preserves matter count and \(d_2\) raises it by two.
To verify this, split the classical master action into its pure-gravity
part and its quadratic scalar part, the latter including
\(\upsilon c\partial\phi\). A scalar antibracket with the quadratic
part preserves count; its gravitational antibracket raises count by two.
The pure-gravity bracket preserves it. Separating the three matter-count
parts of the ordinary identity \(s_\infty^2=0\) gives
\eqref{eq:v31-matter-filtration}. This is a triangular filtration, not
a claim that an arbitrary fixed matter-count slice is a subcomplex of the
full coupled differential. In particular \(d_2\) retains the scalar
stress in \(s\eta\) and the scalar contribution in \(s\sigma\).

Extract from \(\mathcal A\) the following local functionals:
\begin{align*}
 a_0&=[\mathcal A]_{c\phi^2,q=0},&
 a_1&=[\mathcal A]_{q c\phi^2},\\
 \mathcal B_D&=[\mathcal A]_{\upsilon\phi cc,q=0},&
 n_g&=[\mathcal A]_{\eta cc,\phi=\upsilon=q=0}.
\end{align*}
Retain the nonminimal sector until the ordinary off-shell extraction is
made. Let \(\mathscr C a_0\) denote the scalar/ghost Lie variation of
\(a_0\); let \(\mathscr R a_1\) be its metric-linear slot varied by
\(R_0c\); and let \(\mathscr T n_g\) be the scalar-stress replacement
of its metric antifield. These are literal coefficient extractions with
the inherited signed Euler maps, not independently fitted right sides.

\begin{theorem}[The exact action-to-algebra consistency relation]
\label{thm:v31-consistency}
At zero ordinary metric and antifields, with two scalars and two ghosts,
\begin{equation}
 \boxed{
 [s_\infty\mathcal A]_{\phi^2cc}
 =\delta_H\mathcal B_D+
       \mathscr C a_0+\mathscr R a_1+\mathscr T n_g.}
 \label{eq:v31-coupled-consistency}
\end{equation}
Consequently, if this consistency coefficient and the three displayed
other contributions vanish, then \(\mathscr J D=0\) and the scalar
algebra coefficient is removed by the unique quadratic-antifield kernel
of \cref{thm:v31-syzygy}.

If only \(a_0\) and \(n_g\) have been normalized, compatibility of the
scalar-algebra coefficient is determined by
\([s_\infty\mathcal A]_{\phi^2cc}-\mathscr R a_1\), not by the
isolated source square. With the factor-two polarization in
\eqref{eq:v31-J}, \cref{thm:v31-recovery} supplies a quantitative bound
on its remaining scalar coefficient from that exact difference.
\end{theorem}
\begin{proof}
Enumerate the ways an ordinary BV differential can produce the specified
output. Varying the ghost or either scalar in \(a_0\) gives
\(\mathscr C a_0\). The free metric variation of the single metric
in \(a_1\) gives \(\mathscr R a_1\). Replacing the scalar
antifield in \(\mathcal B_D\) by its free scalar Euler expression
gives \eqref{eq:v31-J}. Finally the matter-quadratic term of
\(s\eta\) acting on the pure-gravity coefficient \(n_g\) supplies
\(\mathscr T n_g\).

A term with two or more ordinary metrics still has one metric after a
single variation. A term with two or more antifields still has an
antifield. Replacing a metric antifield already accompanied by a scalar
pair supplies either a metric from its gravitational Euler term or a
second scalar pair from the stress. Neither survives this extraction.
The scalar term in \(s\sigma\) retains a scalar antifield and does
not contribute here. Terms from the nonminimal doublet either retain a
nonminimal field or belong to its already restored off-shell contacts.
This proves exhaustion of the listed contributions. It is equally the
matter-count-two part of \eqref{eq:v31-matter-filtration}.

A quadratic scalar functional vanishes for arbitrary scalar tests exactly
when its symmetrized momentum coefficient vanishes. Using
\eqref{eq:v31-J} therefore gives \(\mathscr J D=0\) under the
stated hypotheses. The kernel and bound follow from the preceding
polynomial theorems. No finite-filter nilpotence or new measured
consistency estimate has entered the argument.
\end{proof}

V29 normalizes the first mixed action row in its declared continuum
reference convention, but does not yet normalize \(a_1\) with the new
algebra and parent contributions in one functional. V30 explicitly
retains that task. Equation~\eqref{eq:v31-coupled-consistency} is why
closing \(a_1\) and its consistency record is the upstream problem;
checking \(\Omega_H\) again is not a substitute. At a finite cutoff
\([s_\infty\mathcal A_a]_{\phi^2cc}\) must be kept as an actual
comparison coefficient, not set to zero because its ordinary limit
has a nilpotent differential.

\section{The source-action insertion, its descendants, and cutoff bookkeeping}
\label{sec:v31-insertion}

Take the local scalar-algebra array \(D_{a,L}\) of the \emph{complete}
post-V29 measured residual. When using it, include the canonical terms
in \eqref{eq:v30-canonical-row}, the quadratic-antifield Euler return,
the actual reference and action-breaking insertions, the counterterm
\(\mathfrak d_U\), and the declared nonminimal/measure convention.
This specifies the input to the map; it is not a claim to have numerically
evaluated that array or to have assembled its missing action row in this
installment. The same construction applies unconditionally to any
already supplied exact array.

Put \(E_{a,L}=\mathfrak h_ED_{a,L}\), and insert
\begin{equation}
 \boxed{C_{EE,a,L}=\frac12\sum_A\int
       \upsilon_A E_{a,L;c}\upsilon_A,\qquad
       S_a^{\mathrm{new}}=S_a+\hbar C_{EE,a,L}.}
 \label{eq:v31-counterterm}
\end{equation}
The index \(A\) is the diagonal scalar species. This coefficient is
real in the inherited reality convention, even in Grassmann parity, and
has ghost number zero. Its formal scalar transpose is skew, as required
for a nonzero quadratic expression in the odd \(\upsilon\)'s.

\begin{theorem}[A relative quadratic-antifield update]
\label{thm:v31-counterterm}
At zero metric the free Euler variation is
\begin{equation}
 \boxed{\delta_H C_{EE,a,L}
        =-\sum_A\langle\upsilon_A,E_{a,L;c}P_0\phi_A\rangle.}
 \label{eq:v31-counterterm-sign}
\end{equation}
Thus the local ordinary scalar-algebra array becomes
\(\mathfrak r_ED_{a,L}\). The antifield-free action rows, in
particular V29's normalized \(c\phi^2\) row and its one-metric
extension, are unchanged by this counterterm. The update has no effect
on the tree propagators, stationary Cartan section or the once-counted
auxiliary density.

Its remaining full-BV variation is retained: the replacement of
\(P_0\) by \(Q_0(q)\) gives the metric-marked
\(\upsilon q\phi cc\) contacts, and the scalar-density/ghost
variation gives quadratic-antifield three-ghost contacts. No claim that
these higher equations vanish is made.
\end{theorem}
\begin{proof}
Using the left odd derivation and formal adjointness gives
\[
 \delta_H(\tfrac12\upsilon E\upsilon)
 =\tfrac12\{(P_0\phi)E\upsilon-\upsilon E(P_0\phi)\}
 =-\upsilon E(P_0\phi).
\]
This agrees with the Fourier scalar exchange in
\eqref{eq:v31-factor}; the operator order is \(EP_0\), not
\(P_0E\). The change of the coefficient is therefore
\(-P\mathfrak h_ED\).

Every monomial of \(C_{EE}\) contains two scalar antifields. One
classical antibracket can remove at most one of them. Its variation
has no antifield-free term, so it cannot alter either action row just
named. The ordinary scalar Euler map at a background is \(Q_0(q)\phi\),
which yields the stated first descendants. The remaining source
variations preserve the two scalar antifields and add one ghost.
Keeping those terms is precisely use of the full \(sC_{EE}\), rather
than only its selected zero-background component. An order-\(\hbar\)
quadratic-antifield insertion affects the one-loop coefficient by this
tree variation; reinsertion in a one-loop graph starts at
\(\hbar^2\). At zero antifields the counterterm itself vanishes.
\end{proof}

\subsection{Finite Euler operator and the already computed parent}

On the inherited soft panel,
\(\widetilde P_a-P_0=O(a^4|p|^6)\). For a parent of degree at most
two the local rule is exactly
\begin{equation}
 \boxed{\mathbf T_4\{\widetilde P_aE\}=P_0E.}
 \label{eq:v31-native-local}
\end{equation}
There is no finite local mismatch at this order. Its unprojected remainder
is \(E(\widetilde P_a-P_0)\) and must still be retained. More
precisely, if the degree-one and degree-two coefficient norms of \(E\)
are bounded by \(B_1,B_2\), then on a window of scale \(\mu\)
\begin{equation}
 \|E(\widetilde P_a-P_0)\|_{r,\mu}
 \le C_r a^4\{B_1\mu^7+B_2\mu^8\}.
 \label{eq:v31-native-product}
\end{equation}
Here \(\|\cdot\|_{r,\mu}\) includes derivatives rescaled by
\(\mu\). For \(a\mu\ll1\), the fixed differentiated native symbol
bounds used in V30 prove this by Leibniz differentiation. No estimate
for \(B_1,B_2\) of an unassembled input is presumed.

For V30's \emph{known} parent, parity excludes \(B_1\) and its
local bound is
\(B_2\le C\chi^{-1}(1+\log(1/(a\lambda)))\). Equation
\eqref{eq:v31-native-product} then recovers, as an application rather
than a new loop estimate, the bound retained in
\eqref{eq:v30-parent-small-product}. Its Euler-compatible local image
is extracted and cancelled by \(\Pi_E\) exactly. The nonlocal parent
is not deleted and retains V30's nonzero possible continuum limit.

For a compatible even \(D\), the primitive is especially simple:
\begin{equation}
 \boxed{E=\partial_zD\big|_{z=0},\qquad
 D(k,l,p,z)=(p^2+z)E^{[2]}(k,l,p).}
 \label{eq:v31-even-mass}
\end{equation}
Thus the finite-mass and zero-mass coefficients of a compatible even
algebra row are linked, not independent normalizations. A weight-two
even breaking cannot lie in this Euler kernel unless it is zero. This is
a test to apply to the \emph{combined} compatible row; it does not erase
V30's potentially nonzero weight-two isolated insertion.

\subsection{Source topology of the new maps}

All maps above act on finite jets. They do not assert an operator-norm
bound on unwindowed ultraviolet fields. Multiplying the output by a
fixed smooth window on its four external legs and rescaling the twelve
relative momentum variables gives, for each fixed moment \(b\),
\begin{equation}
 \begin{split}
 a^{12}\sum_{x_1,x_2,x_3}
 (1+\mu\textstyle\sum_i d_{\mathrm{per}}(x_i,0))^b
       \|\mathcal K_E(x_1,x_2,x_3)\|
 &\le C_b\|E\|_{\mu,1},\\
 &\le5C_b\mu^{-2}\|D\|_{\mu,1}.
 \end{split}
 \label{eq:v31-kernel}
\end{equation}
Indeed, on the rescaled window a polynomial of degree at most two has
all derivatives bounded by its coefficient norm times fixed window
seminorms. Integration by parts in more than \(b+12\) derivatives
gives integrable decay in twelve dimensions. Sampling with physical
weights is uniformly summable when \(a\mu\) is small. Periodizing
and bounding periodic distance by the distance of each lift proves the
stated estimate. The corresponding four-linear functional is bounded
with one argument in \(L^1\) and the others in \(L^\infty\), without
an extensive multiplier.

These norm bounds preserve whatever quantitative cutoff estimate is
actually proved for the relevant input coefficient or for its
consistency defect via \eqref{eq:v31-consistency-bound}. They do not
supply a new loop convergence rate for an unknown combined array, nor a
bounded analytic RG self-map on arbitrary interacting ESM actions.

\section{Exact examples and the remaining native action-side target}
\label{sec:v31-examples}

\paragraph{An admissible Euler-trivial algebra term.}
Take
\(a^0(c)=c^0\partial_1c^1\), all other components zero, and put
\(E_c=a^0\partial_0+\frac12\partial_0a^0\).
It is formally skew, has two derivatives in total and two ghosts.
The coefficient \(D=EP_0\) satisfies \(\mathscr J D=0\), and
\eqref{eq:v31-even-mass} returns exactly \(E\). The parent
\(\frac12\int\upsilon E_c\upsilon\) cancels this scalar-algebra
term with the sign in \eqref{eq:v31-counterterm-sign}, without changing
an antifield-free action normalization. This is a polynomial diagnostic,
not a native integral evaluation.

\paragraph{Why V29's diagnostic still needs its metric-marked action row.}
For V30's test \(U_1=p_0^3\), with ghost polarizations
\(v=e_0,w=e_1\), the component of \(\mathfrak d_U\) is
\[
 D_U=p_0^4-p_0(p_0+k_0)^3
     =-3k_0p_0^3-3k_0^2p_0^2-k_0^3p_0.
\]
It has no mass coefficient. Hence \(\mathfrak h_ED_U=0\) and
\(\mathfrak r_ED_U=D_U\). At
\(p=e_0,k=e_0,l=0,z=0\), its compatibility coefficient is
\begin{equation}
 \boxed{\mathscr J D_U=-14.}
 \label{eq:v31-witness}
\end{equation}
A quadratic-antifield correction alone cannot cancel this component.
There is no obstruction to the programme in that statement: the same
source counterterm has an action descendant in \(a_1\), and
\eqref{eq:v31-coupled-consistency} requires it there. Forgetting that
descendant would produce exactly the false compatibility claim detected
by this example.

\paragraph{The next task is not another calculation of the parent.}
The native remaining coefficient is
\[
 [s_\infty\mathcal A^{\mathrm{post29}}]_{\phi^2cc}
 -\mathscr C a_0-\mathscr R a_1-\mathscr T n_g.
\]
Its scalar polarization is \(\mathscr J D\). The already solved
first action row and pure-gravity normalization remove their own
contributions only in their stated common reference convention. The
one-metric action coefficient \(a_1\), its full local assignment and
its reference terms must still be assembled with the higher-antifield
parents retained. A consistency or convergence assertion about the
isolated square is insufficient.

V31 removes an independent choice from that problem: once the combined
compatibility is established, the zero-background scalar-algebra parent
is uniquely \(\mathfrak h_ED\) in this scalar Euler channel. For a
nonzero finite compatibility residual, its remaining coefficient is
explicitly \(\mathfrak r_ED\) with the stated norm. These conclusions
are compatible with the modified-reference framework
\cite{V28IIM}; the lower reference is not replaced by zero and the finite
filtered transformation is not declared nilpotent.

\section{Final ledgers for V31}
\label{sec:v31-ledgers}

\subsection*{Proved}
\addcontentsline{toc}{subsection}{V31: Proved}
The exact scalar Euler compatibility kernel is
\(D=(p^2+z)E\) with the scalar-exchange-odd factor \(E\), and at
weight four or below that factor is mass independent. The coefficient map
\(\mathfrak h_E\), its projection and remainder satisfy
\eqref{eq:v31-projector-identities} and the bounds
\eqref{eq:v31-projector-bounds}. The rooted component counts, kernel
dimensions and ranks are proved by generating series and an explicit
centred-coordinate basis. The remainder is recovered from the exact
compatibility defect by monic polynomial division, with
\eqref{eq:v31-consistency-bound}. The coupled consistency extraction
identifies all action, gravity-stress and scalar-Euler contributions.
The quadratic-antifield source-action insertion cancels the projected
part without changing an antifield-free action row, and its higher
source descendants, finite Euler product and windowed kernel norm are
retained. None of these claims requires an unproved vanishing of the
native combined residual.

\subsection*{Inherited}
\addcontentsline{toc}{subsection}{V31: Inherited}
The complete V1--V30 mathematical body is retained, not independently
re-certified. V27 supplies the common scalar--gravity branch; V28 supplies
the mass-graded assignment and scalar source; V29 supplies the first
relative action-row normalization; V30 supplies the scalar-square
insertion, quadratic-antifield parent, local bounds and explicit source
counterterm descendant. The ordinary zero-mesh master action,
nonminimal conventions, reference insertion and coordinate/Haar
bookkeeping retain their original scope. The newly supplied classical
spacetime, spectral and structural manuscripts do not supply a missing
quantum mixed Ward coefficient by themselves.

\subsection*{Literature input}
\addcontentsline{toc}{subsection}{V31: Literature input}
The role of equations of motion and the Koszul--Tate differential in
local BV cohomology is standard \cite{V20BBH1994}. The distinction
between a modified fixed-reference identity and an unmodified quantum
master equation remains as in \cite{V28IIM}. No external theorem is used
to declare the mixed source block anomaly-free. The polynomial kernel
proof, concrete mass-coefficient map and coupled extraction are supplied
above rather than attributed to an unevaluated classification.

\subsection*{Conditional}
\addcontentsline{toc}{subsection}{V31: Conditional}
Exact scalar-algebra restoration of the complete native coefficient
follows if its coupled action and consistency data make
\(\mathscr J D=0\); a bound on those data gives an explicit bound
on the unremoved coefficient. V31 has not established that native input
condition. Applying the compiler to a numerical approximation requires
its coefficient error to be retained. Finite product and kernel bounds
consume the displayed native soft-symbol estimates and whatever
coefficient bounds have actually been proved for the input; they are not
a replacement for those estimates.

\subsection*{Open}
\addcontentsline{toc}{subsection}{V31: Open}
Assemble and normalize the complete post-V29 one-metric/two-scalar action
row and its reference/consistency coefficient, keeping the already
computed scalar and mixed antifield returns. That is the remaining
input to the explicit Euler correction, not a new arbitrary choice of
quadratic-antifield counterterm. The resulting higher-metric and
three-ghost descendants must then enter their own coupled equations.
A complete mixed source complex, internal gauge and chiral-line
interfaces, all-loop filtered estimates, and an invariant interacting
ESM shell class are still outside the new result. No all-order
perturbative closure, positive quantum-gravity measure or
nonperturbative continuum is claimed.

\begin{center}\small
\begin{tabular}{>{\raggedright\arraybackslash}p{.24\textwidth}>{\raggedright\arraybackslash}p{.67\textwidth}}
\toprule
Variable & Scope in V31\\\midrule
Cutoff and volume & New polynomial maps and constants are independent of
mesh, periods and loop values. Native insertions use local soft panels;
no all-frequency field norm follows.\\
Mass & Polynomial jets in \(z=m^2\), regular at zero. Extraction of a mass
coefficient requires that coefficient to have been retained; it is not
inferred from a single massless value.\\
Source order & One scalar species in the flavour-diagonal branch,
\(\upsilon\phi cc\) through weight four, and its quadratic-antifield
parent through weight two. No general flavour-matrix count is claimed.\\
Parity & The general 280-dimensional kernel needs no parity. The
256-dimensional even specialization is used only when evenness is proved
for the actual input.\\
Norms & Fixed coefficient-jet norms and fixed smooth source-window
moments. No uniformity in unbounded source or derivative order.\\
Quantum order & Counterterm inserted at order \(\hbar\). Loop
reinsertion would be a higher-order calculation. No new loop integral
has been evaluated.\\
Coupled closure & Action/consistency conditions remain explicit. A
quadratic-antifield correction preserves action rows but does not solve
them.\\
\bottomrule
\end{tabular}
\end{center}

The self-contained checker \srcid{check_ESM_Euler_syzygy_V31.py}
checks every scalar momentum--mass monomial through weight four,
verifies the projector and defect reconstruction exactly, checks the
finite counts, Grassmann sign and the V30 diagnostic, and provides a
rational single-component coefficient compiler. Its certificate separates
these exact algebraic checks from the analytic proofs, inherited inputs,
unevaluated native loop integrals and the unresolved coupled coefficient.
No proof-assistant formalization is asserted.


\clearpage
\part*{V32: The simultaneous matter--antifield block and the missing metric-source return}
\addcontentsline{toc}{part}{V32: Simultaneous matter--antifield block and metric-source return}

\section{An upstream operand of the metric-marked scalar equation}
\label{sec:v32-start}

Sections 1--247 and their labels are preserved. V31 identified the
remaining input to its Euler section as the complete post-V29
one-metric/two-scalar action row and its consistency coefficient. There
is an additional measured operand in that row which cannot be obtained
by merely differentiating the six displayed terms of the flat row while
holding the quantum metric transformation independent of the scalar.
At zero retained metric the pure-gravity Euler covector is zero; its
first metric derivative is not. Consequently a scalar-dependent quantum
metric transformation contributes for the first time in this row.

This installment calculates that transformation from the simultaneous
hard integral, together with the mixed-antifield parents which the same
integral generates. The previous metric--scalar action, scalar source,
and quadratic scalar-antifield formulas are recovered as restrictions,
not presented again as new sector theorems. The new source estimates
complete the one-loop matter-count-two minimal bank through total arity
five. They give the full finite measured metric-marked action row and
its subtracted, nonlocal continuum comparison. They do not supply a
local primitive for that whole row or certify V31's remaining local
Euler compatibility condition.

\subsection{Fixed inputs and the meaning of the new calculation}

Use exactly the completed metric/ghost reference of V14--V18, the scalar
completion and common contour of V27, the scalar source of V28, and
V28's polynomial grading in \(z=m^2\). There are four equal-mass real
scalar components. Internal gauge, Yukawa and scalar self-interaction
couplings remain zero in this coefficient. A fixed external scalar
species is used in formulas; the Kronecker contraction connecting its
two matter arguments is retained. There is no additional closed flavour
trace. No new finite action, cutoff profile, source, or contour is chosen.

The inputs are \eqref{eq:v27-UL}, \eqref{eq:v28-N}, the gravitational
source kernels of V18/V26, and their actual background inverses. The
symbol estimates of V14, V26--V28 are inherited at their stated scopes.
All proofs below distinguish these inputs from their new products and
coefficient extractions. The preceding body is not independently
recertified by appending this installment. The classical, spectral and
structural manuscripts provide the programme's carrier and classical
interfaces; they do not evaluate the missing quantum source coefficient.

The coefficient \(\Gamma_{a,1}\) below has its factor \(\hbar\) removed.
Previously chosen order-\(\hbar\) counterterms are added as local tree
insertions after this raw one-loop calculation. Putting such a counterterm
inside one of the displayed loops would be an order-\(\hbar^2\) operation.
The auxiliary amplitude and coordinate Jacobian have their inherited
once-counted action--density convention. Neither depends on an external
scalar or scalar antifield before the hard metric is integrated.

General source-dependent Gaussian/Berezin reduction and modified
Slavnov identities are established tools; compare \cite{V28IIM}.
The present contribution is the simultaneous native operand, the new
mixed source sectors, their estimates, and the resulting extra Euler
term. No new general existence theorem for perturbative gravity is
claimed.

\section{One joint Hessian for all matter-count-two minimal sources}
\label{sec:v32-joint}

Denote the retained minimal antifields by
\(\eta=q^*,\ \sigma=c^*,\ \upsilon=\phi^*\).
Count each \(\phi\) and each \(\upsilon\) as one matter argument.
Ordinary metric and ghost backgrounds are \(q,c\); ordinary antighost
and nonminimal backgrounds are set to zero only after their off-shell
Gaussian convention has been retained. Let
\(G_g(q),G_H(q),G_c(q)\) be the inherited background hard inverses,
read as finite formal background jets. On zero covariance directions
these mean the resolvents \((I+CV)^{-1}C\), not an inverse of zero.

The metric action contact \(U\) and the mixed action block \(L\) are
\begin{align}
 U(q;\phi)[x,y]
  &=\tfrac12\langle\phi,D_q^2Q_a(q)[x,y]\phi\rangle,\\
 \langle x,L(q;\phi)\zeta\rangle
  &=\langle D_qQ_a(q)[x]\phi,\zeta\rangle.
 \label{eq:v32-UL}
\end{align}
Here \(x,y\) are hard metric fields and \(\zeta\) is a hard scalar.
The ghost action, which is not replaced by a gauge fixing of the filtered
source, has a hard term \(\bar\beta B(q,c)x\).
The exact hard antighost integral substitutes
\begin{equation}
 \beta_x=-G_c(q)B(q,c)x.
 \label{eq:v32-beta}
\end{equation}

Let \(D_\eta(q),D_\sigma(q)\) be the actual even metric quadratic
kernels after that substitution:
\begin{align}
 \tfrac12\langle x,D_\eta(q)x\rangle
  &=\tfrac12\langle x,T_{\eta c}(q)x\rangle
       -L_\eta(q;x)G_c(q)B(q,c)x,\\
 \tfrac12\langle x,D_\sigma(q)x\rangle
  &=Q_\sigma\bigl(G_c(q)B(q,c)x\bigr),
 \qquad D=D_\eta+D_\sigma.
 \label{eq:v32-D}
\end{align}
The first source kernel includes its transformation contact and its
one-ghost-propagator return; the second contains two ghost propagators.
Their definitions agree with V18 and V26. Each background differentiation
acts on those ghost inverses as well as on the other factors.

The scalar source gives the inherited rectangular block
\begin{equation}
 \langle x,N(q;\upsilon,c)\zeta\rangle
 =\left\langle\upsilon,
   P\left([PG_c(q)B(q,c)x]^\mu D_\mu\zeta\right)\right\rangle .
 \label{eq:v32-N}
\end{equation}
Thus the source-substituted even bosonic quadratic matrix is
\begin{equation}
 \boxed{\mathbb H_{\mathrm{even}}=
 \begin{pmatrix}
 G_g^{-1}+U+D&L-N\\
 (L-N)^T&G_H^{-1}
 \end{pmatrix}.}
 \label{eq:v32-Hessian}
\end{equation}
The notation \(G^{-1}\) in this display is restricted to the hard
support. All subsequent formulas use products of the covariances and
remain defined on their full finite carriers. Transpose is the bosonic
bilinear transpose, including reversal of the routed momentum, not
Hermitian conjugation. Entries involving retained odd parameters are
placed in the inherited exterior order; the complete entries of this
matrix are even.

Put
\begin{align}
 W_\phi&=U-LG_HL^T,\\
 W_{\upsilon\phi}&=LG_HN^T+NG_HL^T,\\
 W_{\upsilon\upsilon}&=-NG_HN^T,
 \qquad W_{(2)}=W_\phi+W_{\upsilon\phi}+W_{\upsilon\upsilon}.
 \label{eq:v32-W}
\end{align}
The subscript \((2)\) counts matter arguments, not loop number.

\begin{theorem}[Simultaneous matter--antifield coefficient]
\label{thm:v32-joint}
On total external arity at most five, the complete one-loop coefficient
of matter count two is
\begin{equation}
 \boxed{\Gamma_{a,1}^{(2)}
 =\tfrac12\operatorname{pr}_{A\le5}
       \Tr\left[(I+G_gD)^{-1}G_gW_{(2)}\right].}
 \label{eq:v32-master-trace}
\end{equation}
It uses the actual joint scalar, metric and ghost integral, and all
retained geometric derivatives act on every inverse and source factor.
No purely scalar, purely gravitational, or auxiliary determinant is added
again to this matter-count-two coefficient.
\end{theorem}
\begin{proof}
The only hard antighost terms are its existing action terms, so its
Berezin integral enforces \eqref{eq:v32-beta}, independently of the
minimal antifields. The gravitational source then gives exactly
\eqref{eq:v32-D}. The hard-quadratic scalar source is
\(\langle\upsilon,P((P\beta)\cdot D\zeta)\rangle\);
substitution makes the off-diagonal block \(L-N\), with that sign.
The ghost determinant is independent of \(\phi,\upsilon\).

The block determinant of \eqref{eq:v32-Hessian} has metric Schur
factor \(G_g^{-1}+D+W_{(2)}\). Every term of \(W_{(2)}\) has
matter count two and \(D\) has count zero. The coefficient linear in
\(W_{(2)}\) in its half logarithm is consequently
\[
 \tfrac12\Tr\bigl[(G_g^{-1}+D)^{-1}W_{(2)}\bigr].
\]
Writing the inverse as \((I+G_gD)^{-1}G_g\) proves the formula.
This identity is an identity of finite formal jets even when a source
matrix contains nilpotent even exterior coefficients. Cyclic trace is
used only for even complete factors.

To verify that the hard-quadratic calculation exhausts the measured
one-loop coefficient, retain a labelled panel with at most five
external arguments. A hard-linear vertex has momentum bounded by the
sum of these soft momenta, strictly below the support of every hard
covariance. It therefore contributes zero. A connected one-loop graph
with all hard valences at least two satisfies
\(\sum_v(\deg_hv-2)=0\); hence every hard valence is exactly two.
The Hessian cycles are precisely these graphs. Hard-linear returns from
an order-\(\hbar\) auxiliary amplitude are absent for the same support
reason, and its direct tree value is independent of matter. The scalar
Hessian itself has no retained \(\phi\) dependence on the zero
self-coupling branch. These facts also explain the absence of an
additional flavour multiplicity for a fixed external pair.
\end{proof}

For example, setting \(\eta=\sigma=\upsilon=0\) recovers V27's
\(\tfrac12\Tr(G_gU-G_gLG_HL^T)\). Keeping only one \(\upsilon\)
recovers V28's triangle, and keeping two recovers V30's parent.
These restrictions are inherited results, not the new assertion.
The new factors arise by expanding the same metric resolvent in
\(D_\eta,D_\sigma\), rather than evaluating disconnected source
systems on the separate diagonal species blocks.

\section{The new metric and mixed-antifield coefficients}
\label{sec:v32-new}

Let \(b,r,s,v,f,t\) count metrics, metric antifields, ghost antifields,
scalar antifields, scalars and ghosts, respectively. Matter count two
and ghost number zero give
\begin{equation}
 f+v=2,\qquad t=r+2s+v,\qquad
 A=b+2+2r+3s+v\le5.
 \label{eq:v32-bank}
\end{equation}
There are the following six structural types (thirteen choices of the
listed metric degree):
\begin{center}\small
\begin{tabular}{c c c c}
\toprule
Type & Metric degree & Antifield count \(n=r+s+v\) & Status\\
\midrule
\(q^b\phi^2\)&\(0\le b\le3\)&0&V27/V28\\
\(\upsilon q^b\phi c\)&\(0\le b\le2\)&1&V28\\
\(\upsilon^2q^bcc\)&\(0\le b\le1\)&2&new at \(b=1\)\\
\(\eta q^b\phi^2c\)&\(0\le b\le1\)&1&new\\
\(\sigma\phi^2cc\)&0&1&new\\
\(\eta\upsilon\phi cc\)&0&2&new\\
\bottomrule
\end{tabular}
\end{center}
A source superscript here labels a coefficient type; two occurrences of
an odd source are distinct test arguments. Their exterior signs are not
suppressed by this table.

\begin{proposition}[Explicit additional source words]
\label{prop:v32-new-words}
At zero ordinary metric, the additional gravitational-antifield
coefficients are obtained from
\begin{align}
 [\Gamma_{a,1}^{(2)}]_{\eta\phi^2c}
  &=-\tfrac12\Tr(G_gD_\eta G_gW_\phi),\\
 [\Gamma_{a,1}^{(2)}]_{\sigma\phi^2cc}
  &=-\tfrac12\Tr(G_gD_\sigma G_gW_\phi),\\
 [\Gamma_{a,1}^{(2)}]_{\eta\upsilon\phi cc}
  &=-\tfrac12\Tr(G_gD_\eta G_gW_{\upsilon\phi}).
 \label{eq:v32-new-words}
\end{align}
The \(\eta q\phi^2c\) coefficient is the full metric derivative of
the first expression. The \(\upsilon^2qcc\) coefficient is the full
metric derivative of \(-\tfrac12\Tr(G_gNG_HN^T)\).
\end{proposition}
\begin{proof}
The factors \(D_\eta,D_\sigma\) each contain one gravitational
minimal antifield. Expanding \((I+G_gD)^{-1}\) once produces all the
new allowed terms; two such factors would have arity at least six in
matter count two. The two source terms inside \(W_{\upsilon\phi}\)
must both be kept before exterior labels are extracted. Differentiation
in \(q\) commutes with extraction of these finite jets and acts on
every product factor. This proves all the statements, including the
signs. It does not set any of these integrated coefficients to zero.
\end{proof}

In particular, the scalar-dependent metric transformation is defined by
\begin{equation}
 [\Gamma_{a,1}]_{\eta\phi^2c}
   =\tfrac12\langle\eta,\mathcal M_a^H(\phi,\phi;c)\rangle.
 \label{eq:v32-MH}
\end{equation}
It is the sum of a contact word and a further mixed-species word:
\begin{equation}
 -\tfrac12\Tr(G_gD_\eta G_gU)
 +\tfrac12\Tr(G_gD_\eta G_gLG_HL^T).
 \label{eq:v32-MH-two}
\end{equation}
Both metric covariances are required. Consequently this source coefficient
has order \(\chi^{-2}\), whereas the action row in which it is
contracted with the gravitational Euler operator has order
\(\chi^{-1}\). This is not a new loop order.

For clarity, even the first background derivative of the new trace
contains more than a single vertex derivative. Use
\(A^g[h]=D_hG_g^{-1}\), so \(D_hG_g=-G_gA^g[h]G_g\).
Then
\begin{align}
 D_h\{-\tfrac12\Tr(G_gD_\eta G_gW_\phi)\}
 ={}&\tfrac12\Tr(G_gA^g[h]G_gD_\eta G_gW_\phi)\notag\\
 &-\tfrac12\Tr(G_g(D_hD_\eta)G_gW_\phi)\notag\\
 &+\tfrac12\Tr(G_gD_\eta G_gA^g[h]G_gW_\phi)\notag\\
 &-\tfrac12\Tr(G_gD_\eta G_gD_hW_\phi).
 \label{eq:v32-first-background}
\end{align}
Inside the second and fourth lines,
\(D_hG_c=-G_cA^c[h]G_c\) and
\(D_hG_H=-G_H(D_hQ_a)G_H\) are part of the derivative.
No species inverse may be held fixed in this operation.

\subsection{A finite signed-matrix diagnostic}

The extra term is not eliminated by the Schur algebra. In a two-metric,
one-scalar diagnostic take
\[
 C_g=\begin{pmatrix}1&0\\0&-1\end{pmatrix},\quad
 C_H=\tfrac12,\quad L=\binom12,\quad
 U=\begin{pmatrix}2&1\\1&3\end{pmatrix},\quad
 D_\eta=\begin{pmatrix}1&2\\2&0\end{pmatrix}.
\]
Suppress a single even product of external antifield and ghost labels.
Then
\[
 W_\phi=\begin{pmatrix}3/2&0\\0&1\end{pmatrix},\qquad
 -\tfrac12\Tr(C_gD_\eta C_gW_\phi)=-\tfrac34.
\]
The contact and mixed words separately give \(1\) and \(-7/4\).
This is an exact matrix check of normalization and noncancellation, not
a native loop integral or an integrated physical coefficient. The
native formula is \eqref{eq:v32-MH-two}; its value is not replaced by
this diagnostic.

\section{Source bounds for the closed matter-count-two one-loop bank}
\label{sec:v32-bounds}

Use the mass assignment \(\deg K=1,\ \deg z=2\), and the
operator \(\mathbf T_D\) of \eqref{eq:v28-weighted-T}. Define
\[
 n=r+s+v,\qquad \kappa_I=\chi^{-(1+r+s)},\qquad
 R_p=\lambda+|p|,\qquad \epsilon=|K|+m.
\]
The factor \(\kappa_I\), rather than \(\chi^{-n}\), records the
extra metric inverse in a matter-quadratic loop. For example
\(\kappa_I=\chi^{-2}\) for \(\eta\phi^2c\) and
\(\eta\upsilon\phi cc\), but \(\kappa_I=\chi^{-1}\) for
\(\upsilon^2qcc\).

We use the common domain
\begin{equation}
 \boxed{\lambda\ge16384\mu>0,\qquad a\lambda\le2^{-22},\qquad
  \lambda L_\nu\ge1,\qquad 0\le m\le M\lambda.}
 \label{eq:v32-domain}
\end{equation}
Every external momentum is at most \(\mu\); \(M\), the profile
seminorms, contour/inverse margins, source order and prescribed
differentiation orders are fixed. The slight strengthening of the
previous domain only ensures all five-leg routings fit their established
plateaux. No lower-reference scale is removed.

\begin{lemma}[Actual mixed source symbol degrees]
\label{lem:v32-symbols}
After expansion into primitive inverse/vertex words, every coefficient
in \eqref{eq:v32-bank} has ordinary and completed integrands with
\begin{align}
 |\partial_p^\beta\partial_K^\alpha\partial_z^j f_{a,I}|
 +|\partial_p^\beta\partial_K^\alpha\partial_z^j f_{0,I}|
 &\le C\kappa_I\mathfrak n_I
     R_p^{-n-|\beta|-|\alpha|-2j},\\
 |\partial_p^\beta\partial_K^\alpha\partial_z^j(f_{a,I}-f_{0,I})|
 &\le C\kappa_I\mathfrak n_I a^3
     R_p^{3-n-|\beta|-|\alpha|-2j}
 \label{eq:v32-symbols}
\end{align}
on the common interior plateau. Here \(\mathfrak n_I\) is the
product of the external polarization norms. The completed upper words
have the corresponding smooth periodic rescaling bounds. The estimates
hold for each prescribed finite differentiation order, including all
background derivatives required by \eqref{eq:v32-bank}.
\end{lemma}
\begin{proof}
The inherited background covariances have symbol degree \(-2\).
Each background derivative replaces an inverse by a product containing
a degree-two Hessian vertex and another inverse, so it leaves that degree
unchanged. Every further derivative of such a product is expanded by the
ordinary inverse rule; this includes the scalar mass insertions.
The actual action operands have degrees
\[
 \deg U=2,\quad\deg L=2,\quad
 \deg D_\eta=\deg D_\sigma=1,\quad\deg N=1.
\]
For \(D_\eta\) these are the one-derivative source contact and
\(1-2+2=1\) for its ghost return. For \(D_\sigma\), each
\(G_cB_c\) has degree zero, followed by the one derivative of the
ghost source. For \(N\), the source derivative of the hard scalar
has degree one and its \(G_cB_c\) factor has degree zero. Thus
\[
 \deg W_\phi=2,\quad \deg W_{\upsilon\phi}=1,
 \quad\deg W_{\upsilon\upsilon}=0.
\]
The initial metric covariance contributes \(-2\); each additional
\(G_gD\) contributes \(-1\). The result is precisely \(-n\).
Each metric inverse carries one \(\chi^{-1}\); background gravitational
vertices cancel the additional \(\chi\) factors in differentiated
inverses. The scalar and source vertices add no such factor. This proves
\(\kappa_I=\chi^{-(1+r+s)}\).

The free discrepancies are fourth relative order. The background metric
action vertices have the already established third relative order, while
the scalar and supported source vertices have at least this much control.
Expand the difference of each complete product by ordered telescoping,
placing its one difference in every possible position. With fixed arity
there are finitely many words. This gives the second line of
\eqref{eq:v32-symbols}; no unproved parity improvement is used.
Mass derivatives obey the inherited rule
\(\partial_z^jC_H=(-1)^jj!\nu_\lambda\widetilde P_a^{-j-1}\),
with one cutoff factor. Each derivative therefore lowers degree by two,
and the same remains true after background differentiation. On the upper
region the inherited action completion and supported sources are smooth
periodic functions of their dimensionless labels. The pointwise coframe
inverse adds coefficients, not momentum poles. Primitive counting and
rescaling yield the asserted upper bounds, including routed crossings.
\end{proof}

\begin{theorem}[Full-cutoff mixed minimal-source estimate]
\label{thm:v32-cutoff}
Let \(\gamma_{a,L,I}\) be any coefficient of
\eqref{eq:v32-master-trace}, per physical volume, and put
\(\gamma^R_I=(I-\mathbf T_{4-n})\gamma_I\).
On \eqref{eq:v32-domain}, for \(d=|\alpha|+2j\le5-n\),
\begin{equation}
 \boxed{\left|\partial_K^\alpha\partial_z^j
 (\gamma^R_{a,L,I}-\gamma^R_{0,L,I})\right|
 \le C_{I,\alpha,j}\kappa_I\mathfrak n_I
       a\epsilon^{5-n-d}.}
 \label{eq:v32-cutoff}
\end{equation}
Every higher fixed weighted derivative has the corresponding scaled
bound \(C\kappa_I a\lambda^{5-n-d}\) on the fixed window.
The limit uses the massive ordinary propagators with the same lower
reference. The stronger special estimates of V27--V30 remain available
on their original restrictions and are not replaced by a false improved
parity assertion for the whole bank.
\end{theorem}
\begin{proof}
The superficial degree in four dimensions is \(4-n\). After the
\(5-n\) weighted derivatives in the Taylor remainder, the continuum
majorant is \(R_p^{-5}\) and the interior discrepancy is
\(a^3R_p^{-2}\). Four-dimensional radial integration gives
\[
 \int_{c/a}^{\infty}R^{-2}\,dR=O(a),\qquad
 a^3\int_\lambda^{c/a}R\,dR=O(a).
\]
These are respectively the omitted continuum tail and the interior
comparison, not two replacements for the same region. On the completed
upper region, \(p=a^{-1}u\) and the corresponding source rescaling
give \(a\) after those derivatives. Smooth periodicity permits the
same argument through a face; no discontinuous native extension is
being used. Normalized finite sums obey these bounds for
\(\lambda L_\nu\ge1\) by subdividing momentum space into physical
cells and using finitely many differentiated symbol bounds.

Apply V28's mass-weighted Taylor remainder along
\((tK,t^2z)\), which stays in the massive hard domain. Its scalar
coefficient bound gives \(\epsilon^{5-n-d}\). For mass derivatives
one differentiates the actual inverse first, rather than dividing by
\(m\) after an auxiliary square-root substitution. This proves
regularity at zero mass and the claimed derivative estimates. The new
products consume only the stated inherited primitive bounds.
\end{proof}

The local assignment remains genuine data. For a weighted coefficient
of degree \(d\le4-n\), the same symbol proof gives the bound
\begin{equation}
 |\gamma_{a,L,I}^{[d]}|
 \le C_I\kappa_I\mathfrak n_I
 \begin{cases}
 a^{n+d-4},&n+d<4,\\
 1+\log(1/(a\lambda)),&n+d=4.
 \end{cases}
 \label{eq:v32-local-bound}
\end{equation}
No value of a native local annular integral, and no vanishing of its
finite part, is asserted by this estimate.

\begin{corollary}[Rooted source topology, thermodynamic comparison and shells]
\label{cor:v32-kernel}
For a panel of arity \(A\), multiply the remainder difference by a fixed
smooth external window. For every fixed kernel moment \(b_0\),
\begin{equation}
 \boxed{a^{4(A-1)}\sum_{\boldsymbol x}
 \left(1+\mu\sum_i d_{\rm per}(x_i,0)\right)^{b_0}
 \|\mathcal K^\Delta_{a,L,I}(\boldsymbol x)\|
 \le C_{I,b_0}\kappa_I a(\mu+m)^{5-n}.}
 \label{eq:v32-kernel}
\end{equation}
The separate thermodynamic error is bounded, for \(b_0>4\), by
\(C\kappa_I\mathfrak n_I\epsilon^{5-n-d}\lambda^{-1}
 (\lambda L_{\min})^{-b_0}\).
Each complementary shell increment, with all primitive species lines
assigned to its largest shell label, has source norm at most
\(C_I\kappa_I(\mu+m)^{5-n}\Lambda_j^{-1}\).
\end{corollary}
\begin{proof}
Rescale all \(4(A-1)\) external momentum coordinates by \(\mu\),
use more than \(b_0+4(A-1)\) of the differentiated bounds, and
integrate by parts in the inverse Fourier transform. Periodization and
\(d_{\rm per}\le d\) transfer the absolute weighted sum without an
extensive factor. Poisson summation of the smooth ordinary subtracted
integrand and its differentiated integrable tails proves the independent
volume estimate. For a shell increment expand every background inverse,
and also every \(G_c\) inside \(D\) and \(N\), into its primitive
covariances. Ordered telescoping assigns each word once. At least one
annular line confines all other routed labels to the same scale up to
five soft transfers. The degree \(-5\) remainder integrates to
\(\Lambda_j^{-1}\). The sum over scales is geometric. Neither this
nor the kernel bound asserts invariance of an unrestricted interacting
analytic class.
\end{proof}

\section{The eleven terms of the metric-marked scalar action row}
\label{sec:v32-row}

The new metric source has a necessary place in the equation. To make all
normalizations unambiguous define the raw quantum tensors by actual
labelled derivatives at zero retained fields:
\[
 \Pi=D_\phi^2\Gamma_{a,1}^{(2)},\quad
 T=D_qD_\phi^2\Gamma_{a,1}^{(2)},\quad
 V=D_q^2D_\phi^2\Gamma_{a,1}^{(2)}.
\]
The letter \(V\) in this section denotes this quantum two-metric tensor,
not the generic hard Hessian perturbation in a reference trace.
Let \(Z_{g,0},Z_{g,1}(h)\) be the existing scalar-independent quantum
metric transformation and its first metric derivative. Let
\(Z_{H,0},Z_{H,1}(h)\) be V28's scalar transformation and derivative.
Use \(\mathcal M_a^H\) from \eqref{eq:v32-MH}; its definition by two
scalar derivatives fixes its polarization factors. Put
\[
 P_a=Q_a(0),\quad Q_{a,h}=D_qQ_a(0)[h],\quad
 Q_{a,h_1h_2}=D_q^2Q_a(0)[h_1,h_2].
\]
Let \(H_a^{\rm cl}\) be the off-shell, not gauge-fixed, pure-gravity
Euler Hessian of V16. The nonminimal variable is restored before taking
this derivative. The soft classical metric transformations on this
panel are the inherited ordinary \(R_0,R_1\).

\begin{theorem}[Complete first-metric measured action row]
\label{thm:v32-row}
At one loop, with its factor \(\hbar\) removed, the coefficient with
one metric \(h\), one ghost \(c\), and two scalars is
\begin{equation}
\begin{aligned}
 \mathfrak L_{1,a}(h,c;\phi_1,\phi_2)={}&
 V(h,R_0c;\phi_1,\phi_2)
 +T(R_1(h,c);\phi_1,\phi_2)\\
 &+T(h;\mathcal L_c\phi_1,\phi_2)
  +T(h;\phi_1,\mathcal L_c\phi_2)\\
 &+\langle\phi_1,Q_{a,h,Z_{g,0}c}\phi_2\rangle
  +\langle\phi_1,Q_{a,Z_{g,1}(h)c}\phi_2\rangle\\
 &+\langle P_a\phi_1,Z_{H,1}(h;\phi_2,c)\rangle
  +\langle P_a\phi_2,Z_{H,1}(h;\phi_1,c)\rangle\\
 &+\langle Q_{a,h}\phi_1,Z_{H,0}(\phi_2,c)\rangle
  +\langle Q_{a,h}\phi_2,Z_{H,0}(\phi_1,c)\rangle\\
 &+\langle H_a^{\rm cl}h,
                  \mathcal M_a^H(\phi_1,\phi_2;c)\rangle.
\end{aligned}
\label{eq:v32-eleven}
\end{equation}
The actual finite measured identity is
\begin{equation}
 \boxed{\mathfrak L_{1,a}=\mathcal I_{1,a}^H+\mathcal X_{1,a}^H.}
 \label{eq:v32-finite}
\end{equation}
Here \(\mathcal I_{1,a}^H\) is the coefficient of the actual classical
breaking insertion, not just of the scalar source square. The reference
insertion has its complete graded field list \((q,\phi,c,\bar c,b)\).
There is no independent direct auxiliary-density or flat-divergence
term in this particular scalar-quadratic slot.
\end{theorem}
\begin{proof}
Before extracting any ordinary background, finite Gaussian/Berezin
integration by parts gives the identity proved in V28:
\[
 \tfrac12\operatorname{STr}(\mathbb C D^2(DS[r]))
 =D\Gamma_1[r]+DS[z_1]+\operatorname{sdiv}r
       -\operatorname{STr}(\mathbb C\mathbb R Dr),
 \quad z_1=\tfrac12\operatorname{STr}(\mathbb C D^2r).
\]
It follows by expanding the two Hessian derivatives, not by assuming
\(r^2=0\). In \(D\Gamma_1[r]\), differentiation by \(h\) gives
the first four terms of \eqref{eq:v32-eleven}. The scalar stress in
\(DS[z_1^q]\) contributes the two \(Q\)-terms involving
\(Z_{g,0},Z_{g,1}\). The scalar Euler row contributes the four terms
involving \(Z_{H,0},Z_{H,1}\).

There is one further term: the pure-gravity Euler row times the
scalar-dependent part of \(z_1^q\). Its Euler row is zero at \(q=0\)
but its first derivative is \(H_a^{\rm cl}h\). Its transformation
coefficient is precisely \eqref{eq:v32-MH}, so it gives the last line.
No scalar-dependent ghost source contributes to this antifield-free
row: the minimal ghost Euler term is zero at the retained nonminimal
background used for this extraction. Its source coefficient is still
retained for the subsequent antifield equations.

The metric/ghost source divergence is independent of \(\phi\),
and the scalar source divergence is zero. The auxiliary logarithmic
density depends only on \(q\), while the classical metric source
has no scalar dependence. Its \(h c\phi^2\) derivative is zero.
Multiplication by a quantum transformation correction would already
carry a second \(\hbar\). The nonminimal convention changes the
metric Euler multiplier to \(H_a^{\rm cl}\), rather than licensing
use of its gauge-fixed replacement. These facts give
\eqref{eq:v32-finite}, with all eleven terms and no extra measure count.
\end{proof}

The last line cannot be generated by differentiating a flat identity
in which the full scalar-dependent quantum metric source was first set
to zero. That setting was harmless for the flat action row and is not
harmless for its derivative. The factor \(\chi\) in
\(H_a^{\rm cl}\) makes its product with \(\mathcal M_a^H\)
of the same \(\chi^{-1}\) order as the other mixed terms.

The reference word can be written without an inverse cutoff weight:
\begin{equation}
 \mathcal X_{1,a}^H
 =[hc\phi^2]\operatorname{STr}
    \left[\sum_{j=0}^{4}(-\mathbb C_0\mathbb V)^j
                     \Theta\,Dr_a\right],
 \label{eq:v32-reference}
\end{equation}
where now \(\mathbb V=\mathbb H(\Phi)-\mathbb H_0\).
The cancellation \(\mathbb C_0\mathbb R=\Theta\) is performed
before expansion. All nonzero terms in this extraction have at least
one covariance. Their routed labels are between fixed multiples of
\(\lambda\), since a \(\Theta\) factor and a hard weight occur
in the same connected word. On \eqref{eq:v32-domain} this is contained
in \(\lambda/2\le|p|\le3\lambda\). There is no upper-cutoff
singularity in this reference coefficient.

\section{One subtraction and the nonlocal continuum equation}
\label{sec:v32-comparison}

For the ghost-number-zero coefficients use
\(\mathbf P_I=\mathbf T_{4-n}\), including the new metric source.
For the action row use \(\mathbf T_5\). The ordinary scalar Euler
and stress vertices have weight two, the ordinary gravitational Euler
operator has weight two, and each classical transformation has weight
one. Linear subset routing preserves weighted Taylor degree. Therefore
\begin{equation}
 \mathbf T_5\mathfrak L_{1,0}(\Gamma_1)
       =\mathfrak L_{1,0}(\mathbf P\Gamma_1).
 \label{eq:v32-intertwine}
\end{equation}
This is a row-compatible component assignment. It is not a projector
onto invariant physical terms or a primitive in the mixed BV complex.

Let \(\mathfrak L_{1,a}^R\) mean \eqref{eq:v32-eleven} with
complementary quantum coefficients; retain the actual classical finite
multipliers. Let \(\mathcal X_{1,a}^{H,R}=(I-\mathbf T_5)
\mathcal X_{1,a}^H\). In thermodynamic normalization the following
statement concerns the complement of the local bank, not yet its
unsolved normalization.

\begin{theorem}[Nonlocal metric-marked modified identity and cutoff comparison]
\label{thm:v32-nonlocal}
For the declared ordinary zero-mesh action/source and fixed lower
reference,
\begin{equation}
 \boxed{\mathfrak L_{1,0}^R=\mathcal X_{1,0}^{H,R}.}
 \label{eq:v32-limit-row}
\end{equation}
Set \(\mathfrak W_{1,a,L}^R=\mathfrak L_{1,a,L}^R-
\mathcal X_{1,a,L}^{H,R}\). For \(d=|\alpha|+2j\le6\),
\begin{equation}
 \boxed{\left|\partial_K^\alpha\partial_z^j
 (\mathfrak W_{1,a,L}^R-\mathfrak W_{1,0,L}^R)\right|
 \le C\chi^{-1}\mathfrak n_I a\epsilon^{6-d}.}
 \label{eq:v32-row-rate}
\end{equation}
There are analogous fixed higher derivative and rooted kernel bounds;
for the four-argument row its three-relative-coordinate difference has
norm at most \(C_b\chi^{-1}a(\mu+m)^6\).
The separate thermodynamic error has the inherited lower-scale decay;
no exact Ward identity is assigned to arbitrary off-grid finite-volume
interpolations. In particular the lower-reference term in
\eqref{eq:v32-limit-row} is not set to zero.
\end{theorem}
\begin{proof}
The ordinary simultaneous classical action and its coframe-pulled metric
source, scalar Lie source and ghost/nonminimal terms satisfy the
ordinary master identity. This is the ordinary source identification
used in V29, not the finite filtered identity. Use its regulated trace
identity with a smooth upper reference and the same lower reference.
After extraction of \(hc\phi^2\), its canonical left side is now
all of \eqref{eq:v32-eleven}. The upper-reference terms have the
same symbol counting as these complete words. Six weighted external
derivatives give order \(-6\); the upper-reference trace has one
additional derivative from \(Dr\), leaving an annular integral
bounded by a constant times the inverse upper scale. Thus its
\(I-\mathbf T_5\) coefficient tends to zero. The remaining lower
trace is \eqref{eq:v32-reference} at zero mesh. The ordinary
intertwining \eqref{eq:v32-intertwine} yields
\eqref{eq:v32-limit-row}. This proof uses the ordinary identity only
after the independent finite-symbol comparisons have been established.
It supplies no finite-cutoff nilpotence assertion.

For the quantitative comparison, the action coefficients \(T,V\)
have the inherited \(a\epsilon^5\) error, and their one-derivative
classical transformations give \(a\epsilon^6\). The scalar-source
and scalar-independent metric-source comparisons give the same bound
after multiplication by the two-derivative Euler/stress polynomials.
The new term uses \eqref{eq:v32-cutoff} with \(n=1\): its
\(\chi^{-2}a\epsilon^4\) error multiplied by
\(H_0^{\rm cl}=O(\chi\epsilon^2)\) gives exactly
\(\chi^{-1}a\epsilon^6\). Finite classical multipliers differ
from the ordinary ones by their displayed fourth-order free/scalar
errors on this panel; when multiplying a complementary coefficient
they are smaller on \eqref{eq:v32-domain}.

The reference difference is compactly supported at scale \(\lambda\).
A first metric mark can use a third-relative-order gravitational vertex,
so its six-weighted-derivative comparison is bounded by
\(C\chi^{-1}a^3\lambda^2\epsilon^{6-d}\), smaller than the
stated \(a\)-bound. The inverse-cutoff cancellation has already been
made, so no reciprocal \(\nu\) is differentiated. Together these
bounds prove \eqref{eq:v32-row-rate}. Fourier integration by parts
in twelve relative coordinates, using the higher derivative versions,
gives the asserted source-kernel bound. Poisson summation of the smooth
ordinary remainder yields the separate volume comparison.
\end{proof}

\subsection{Local finite-mesh products are not silently removed}

Equation \eqref{eq:v32-intertwine} uses ordinary multipliers. At finite
mesh, subtracting the whole row also produces products of a classical
artifact and an actual local quantum coefficient. A new example is
\begin{equation}
 \left\langle(H_a^{\rm cl}-H_0^{\rm cl})h,
               (\mathbf T_3\mathcal M_a^H)(\phi_1,\phi_2;c)
 \right\rangle.
 \label{eq:v32-local-product}
\end{equation}
The crude bounds already suffice:
\(H_a^{\rm cl}-H_0^{\rm cl}=O(\chi a^4|K|^6)\),
whereas \eqref{eq:v32-local-bound} allows
\(\mathbf T_3\mathcal M_a^H=O(\chi^{-2}a^{-3})\)
in its lowest coefficient. Their product is \(O(\chi^{-1}a|K|^6)\),
not zero. Higher local degrees are bounded after their compensating
powers of \(a\mu\) are retained. The scalar Euler/stress artifacts
and the known local transformation coefficients obey the analogous
bound. All such products have weighted degree at least six, so they
have zero \(\mathbf T_5\) projection, but contribute to the finite
complementary balance. They are included in its \(O(a)\) budget.
No vanishing constant-source coefficient is assumed to justify this
new product estimate.

\section{The remaining local equation now has its missing operands}
\label{sec:v32-local}

The new source coefficients are raw measured coefficients plus their
own recorded local jets. Assigning \(\mathbf T_3\mathcal M_a^H\)
or \(\mathbf T_2[\Gamma]_{\eta\upsilon\phi cc}\) componentwise
does not choose their common physical or Slavnov-compatible finite
normalization. Likewise, \eqref{eq:v32-limit-row} is not the identity
with an unsubtracted reference on its right.

Let \(H_g\) and \(h_H\) be the previously selected pure-gravity
and V29 scalar-row local functionals. The next ordinary local target
retains
\begin{equation}
 B_{1}^{\rm post29}
 =\mathbf T_5\left[
       \mathcal X_{0,\lambda}-s_\infty(H_g+h_H)
                  \right]_{q c\phi^2},
 \label{eq:v32-next-local}
\end{equation}
together with its scalar-algebra and higher-antifield rows. This is the
relative equation, not permission to reset the earlier normalizations.
The eleven-term compiler and \eqref{eq:v32-master-trace} identify all
raw local quantum operands entering the finite representative of this
equation. Their native annular coefficients have not been numerically
evaluated here.

In particular, the new mixed parent gives a term which must be kept in
the scalar-algebra equation. If
\(C_{\eta\upsilon}=\eta_A\upsilon_B F^{AB}(\phi,c,c)\)
is written with that exterior order, its free gravitational/scalar Euler
variation is
\begin{equation}
 \delta_0C_{\eta\upsilon}
  =\mathcal E_{g,A}(q)\upsilon_B F^{AB}
       -\eta_A(P_0\phi)_B F^{AB}.
 \label{eq:v32-parent-descendants}
\end{equation}
It supplies both a one-background scalar-algebra contact and a
metric-antifield matter-stress contact. Their signs come from the odd
\(\eta\); they are not independently adjustable. The
\(\sigma\phi^2cc\) coefficient likewise contributes to the
metric-antifield algebra through the rooted ghost Euler map. These are
reasons to retain the whole bank rather than just
\(\mathcal M_a^H\).

V31's exact identity remains the local compatibility criterion:
\[
 [s_\infty\mathcal A]_{\phi^2cc}
   =\delta_H\mathcal B_D+\mathscr C a_0+
                      \mathscr R a_1+\mathscr T n_g.
\]
This installment supplies the complete measured first-metric action
row and the additional parents needed to evaluate its sources. It does
not prove that the post-counterterm local coefficient
\(B_1^{\rm post29}\) is zero or that the combined consistency
coefficient makes \(\mathscr JD=0\). The distinction is sharp:
V31's mass-coefficient section is immediately applicable to a compatible
array, but no input to that section is replaced by the isolated scalar
source square. A complete local relative primitive, preserving V29's
first row and including \eqref{eq:v32-parent-descendants}, is still the
next task. No scalar-inclusive anomaly classification is imported here.

\section{Final ledgers for V32}
\label{sec:v32-ledgers}

\subsection*{Proved}
\addcontentsline{toc}{subsection}{V32: Proved}
The simultaneous source-substituted metric--scalar Hessian gives
\eqref{eq:v32-master-trace}. It generates the scalar-dependent metric
and ghost transformation coefficients, the mixed metric/scalar-antifield
parent, and the first metric derivative of the quadratic scalar-antifield
parent. Their new words and all background inverse derivatives are
retained. On the stated inherited primitive domain the resulting
matter-count-two arity-five bank has the cutoff, local-coefficient,
rooted-kernel and shell bounds proved above. The first-metric action row
has all eleven canonical terms, including the gravitational Euler times
scalar-dependent quantum metric transformation. Its finite measured
identity and subtracted ordinary continuum comparison retain the actual
reference term. The finite artifact times local-source products and
mixed-parent descendants are explicitly assigned.

\subsection*{Inherited}
\addcontentsline{toc}{subsection}{V32: Inherited}
The V1--V31 body and labels are unchanged. The completed propagators,
scalar continuation, primitive symbol bounds, contour, source filters,
nonminimal convention and auxiliary action--density record are inherited
from V14--V28. V27, V28 and V30 already give the action, one-scalar-
antifield and zero-background quadratic-scalar-antifield restrictions.
V29's first local mixed normalization and V31's polynomial Euler section
retain their original scope. This iteration does not independently
recertify the preceding 422-page proof body.

\subsection*{Literature input}
\addcontentsline{toc}{subsection}{V32: Literature input}
Gaussian/Berezin elimination, Schur reduction and modified-reference
Slavnov methods are established tools; \cite{V28IIM} supplies comparison
context, not the native mixed coefficient or a missing matter-inclusive
primitive. No external cohomology theorem is used to claim that the
remaining local scalar--gravity coefficient is exact.

\subsection*{Conditional}
\addcontentsline{toc}{subsection}{V32: Conditional}
All native cutoff conclusions consume the inherited primitive symbol
and measured-reference results on their stated completed-regulator
branch. Convergence after a common full local normalization also
requires the unresolved relative local primitive at this background
order. V31 removes the scalar-algebra compatible part once its complete
local input is assembled and satisfies its criterion; that criterion
has not been discharged by nonlocal convergence alone. Numerical
coefficient errors would have to be transported through both maps.

\subsection*{Open}
\addcontentsline{toc}{subsection}{V32: Open}
Solve \eqref{eq:v32-next-local} and its coupled local consistency
problem using the newly supplied source parents, preserving the first
mixed action row and the pure-gravity normalization. The counterterm's
higher-metric and three-ghost descendants must remain in their rows.
The full internal-gauge/chiral source complex, determinant-line
transport, arbitrary-loop filtered estimates and an invariant
interacting ESM shell class remain outside this result. No mass-gap,
nonperturbative gravitational continuum, infrared removal, or
finite-parameter ultraviolet-gravity conclusion is drawn.

\begin{center}\small
\begin{tabular}{>{\raggedright\arraybackslash}p{.24\textwidth}>{\raggedright\arraybackslash}p{.67\textwidth}}
\toprule
Variable & Scope in V32\\\midrule
Cutoff and volume & Full smooth completed ultraviolet zone; fixed lower
scale. Constants are independent of mesh and periods on
\eqref{eq:v32-domain}; thermodynamic comparison is separate.\\
Sources & Matter count two, minimal ghost-number-zero arity at most
five; the four-argument first-metric action row. Arbitrary external
antighost/nonminimal panels are not covered.\\
Couplings and mass & Four equal real scalar components; zero internal,
Yukawa and self-couplings. \(\chi>0\) and inverse margins fixed;
\(0\le m\le M\lambda\), with mass jets regular at zero.\\
Quantum order & One loop; old and new local counterterms are separate
order-\(\hbar\) tree insertions. No higher-loop reinsertion theorem.\\
Norms & Fixed coefficient and differentiated/windowed source norms;
up to sixteen relative coordinates. No total-volume factor or
uniformity in unbounded arity or derivative order.\\
Iteration & Largest-primitive-shell ownership and summable complementary
coefficients. No analytic action-class self-map has been inferred.\\
Symmetry & Exact finite measured balance with its breaking insertion;
ordinary subtracted modified identity. Full local restoration and
finite filtered nilpotence are not asserted.\\
\bottomrule
\end{tabular}
\end{center}

The accompanying checker \srcid{check_ESM_joint_matter_sources_V32.py}
compares the full joint logarithm with the simultaneous Schur formula
in an exact finite exterior/matrix algebra. It tests the newly allowed
source families, noncommuting background inverse derivatives, the
mass-weighted degree rule, the eleven-term chain rule and primitive
shell ownership. Its certificate distinguishes these finite identities
from the analytic bounds and from the native loop integrals which have
not been numerically evaluated. No proof-assistant formalization is
claimed.


\clearpage
\part*{V33: Relative scalar--gravity transgression and the boundary of rowwise normalization}
\addcontentsline{toc}{part}{V33: Relative matter transgression and normalization compatibility}

\section{A simultaneous primitive and a genuine relative-normalization question}
\label{sec:v33-start}

Sections 1--255 and all their labels are retained. V32 supplies the
simultaneous matter-count-two source coefficients and the complete
metric-marked action row. It does not solve
\eqref{eq:v32-next-local}. V31 already supplies the matter-count
filtration \eqref{eq:v31-matter-filtration}; V24 supplies the ordinary
arity/weight projection, and V29 supplies its mass-graded extension.
None of those algebraic observations is presented again as a new
classification.

The present step applies the lower-reference derivative to the
\emph{simultaneous} metric--scalar--ghost functional. Its matter-count-two
component is not a cocycle on its own: its differential contains the
scalar stress and scalar ghost-current returns from the gravitational
reference. Keeping those terms gives a common subcritical local primitive
relative to the paired gravitational scale normalization. This is an
actual finite annular integral, not a choice of a desired Ward tensor.
A native scalar-shift test also removes the zero-derivative source
coefficients which would otherwise obstruct the finite local projection.

There is a separate issue of relative normalization. A solution of V29's
first action row need not extend with all its coefficients fixed. An
explicit homogeneous first-row solution below fails the next scalar
algebra equation on affine Killing jets, even on the scalar stationary
surface. Thus an arbitrary rowwise choice cannot be declared extendible.
The new simultaneous normalization is recorded as a \emph{comparison
scheme}; the particular V29 coefficients remain the current prescription
until the stated relative comparison has been resolved. The first
normalized identity and literal equality of its counterterm coefficients
are different requirements.

Throughout this part the internal gauge, Yukawa and scalar self-couplings
are zero, the real scalar carrier has four equal components, and
\(z=m^2\) is assigned weight two. The common factor \(\hbar\) is omitted
from one-loop coefficients. The contour, actual ghost action, scalar
completion, filters and auxiliary action--density convention are exactly
those of V27--V32. There is no new propagator or finite action vertex.
The ordinary differential used for local identities is the one derived
from the native zero-mesh action, not the nonnilpotent filtered finite
transformation. The fixed lower Wilsonian scale is denoted by
\(\lambda>0\).

\section{The relative differential retains the scalar stress and ghost current}
\label{sec:v33-relative}

Let \(\mathbb S_g\) be the ordinary gravitational minimal master action,
and set
\begin{equation}
 \mathbb S_H=\frac12\sum_A\int
 \left[W^{\mu\nu}(q)\partial_\mu\phi_A\partial_\nu\phi_A
                  +z\rho(q)\phi_A^2\right]
      +\sum_A\int\upsilon_A c^\mu\partial_\mu\phi_A .
 \label{eq:v33-SH}
\end{equation}
Here \(W=\rho G^{-1}\), \(\rho=\det e\), \(G=e^Te\), and \(e(q)\)
is the same affine coframe chart as before. Signs of Euler replacements
are those of the inherited antibracket. In particular the scalar stress
is the signed \(q\)-derivative of the first integral and the matter
contribution to the ghost-antifield row is the signed derivative of the
second integral with respect to \(c\).

Count each \(\phi\) and \(\upsilon\) once, writing this count as \(M\).
Split the antibracket into gravitational/ghost and scalar parts. On
matter-count-zero coefficients write \(d_g=(\mathbb S_g,\cdot)_g\).
On matter-count two write
\begin{equation}
 d_H=(\mathbb S_g,\cdot)_g+(\mathbb S_H,\cdot)_H,
 \qquad t=(\mathbb S_H,\cdot)_g.
 \label{eq:v33-differentials}
\end{equation}
The letter \(t\) is an operator here, not a time coordinate. It acts
on a matter-count-zero functional by its metric and ghost antifields:
\begin{equation}
 tF=\sum_a\mathcal T_a(\phi)\frac{\delta^L F}{\delta\eta_a}
       +\sum_\mu\mathcal J_\mu(\upsilon,\phi)
                                \frac{\delta^L F}{\delta\sigma_\mu},
 \label{eq:v33-stress-map}
\end{equation}
with the signs incorporated in the signed coefficients
\(\mathcal T_a,\mathcal J_\mu\). Before that common sign convention,
\(\mathcal T_a\) is the explicit quadratic polynomial
\[
 \frac12\sum_A\left[
 (\partial_{q_a}W^{\mu\nu})\partial_\mu\phi_A\partial_\nu\phi_A
                 +z(\partial_{q_a}\rho)\phi_A^2\right],
 \qquad \mathcal J_\mu\ \text{comes from}\
                   \sum_A\upsilon_A\partial_\mu\phi_A .
\]
Equation \eqref{eq:v33-stress-map} is understood as a functional
coefficient identity, with integration by parts and the ordered odd
roots retained. It is not a new choice of stress tensor.

The inherited matter filtration implies
\begin{equation}
 d_g^2=d_H^2=0,\qquad d_Ht+td_g=0.
 \label{eq:v33-anticommutation}
\end{equation}
Modulo matter count at least four the full operator is consequently
\begin{equation}
 d(F_0,F_2)=(d_gF_0,\ tF_0+d_HF_2).
 \label{eq:v33-triangular}
\end{equation}
This is a quotient of the full coupled complex. Discarding \(tF_0\)
would instead give a different complex.

All statements below are truncated at total arity \(A\le5\), after
forming this operator. For ghost number \(g\), antifield count \(n_*\)
and momentum/mass degree \(D=|\alpha|+2j\), the local weight is
\begin{equation}
 w=D+n_*-g,\qquad \mathbf P_I=\mathbf T_{4+g-n_*}.
 \label{eq:v33-weight}
\end{equation}
The ordinary differential preserves \(w\) and cannot lower arity.
The massive Euler map preserves this assertion since \(z\) has weight
two. These are applications of the already established degree argument.
They give a finite quotient complex without assuming closure of a
zero- or one-metric slice.

\begin{proposition}[The coupled local target]
\label{prop:v33-target}
Let \((x_g,x_2)\) be the ordinary local reference coefficient in this
quotient. Suppose the fixed gravitational normalization obeys
\(d_gH_g=x_g\). Then the relative coefficient
\begin{equation}
 \Xi_2=x_2-tH_g
 \label{eq:v33-Xi}
\end{equation}
is a \(d_H\)-cocycle. A matter-count-two normalization \(H_2\) solves
all its retained local rows at once precisely when
\begin{equation}
 d_HH_2=\Xi_2.
 \label{eq:v33-relative-equation}
\end{equation}
The reference coefficient and the gravitational normalization must use
the same lower-reference convention.
\end{proposition}
\begin{proof}
Ordinary reference consistency gives
\(d_gx_g=0\) and \(d_Hx_2+tx_g=0\). Applying \(d_H\) to
\eqref{eq:v33-Xi} and using \eqref{eq:v33-anticommutation} yields
\(-tx_g+td_gH_g=0\). Equation \eqref{eq:v33-triangular}
then proves the second assertion. In particular the term \(tH_g\)
contains V29's matter stress times the gravitational transformation,
and also the ghost-current return. It is not set to zero by evaluating
the ordinary metric at zero.
\end{proof}

There is a harmless background-scalar determinant at \(M=0\) in the
full simultaneous integral. It has no gravitational antifield and is
annihilated by \(t\). Its scale derivative and reference variation also
have no such antifield. It can therefore be retained in the gravitational
background block or subtracted together with its own reference when
forming \eqref{eq:v33-relative-equation}; neither choice changes that
equation. This observation is not a new local normalization theorem for
that determinant. The nontrivial \(tH_g\) comes from the already
retained gravitational source coefficients.

\section{A native massless-shift test licenses the local finite comparison}
\label{sec:v33-floor}

The zero-mass Taylor origin is useful even when the physical scalar mass
is nonzero. At \(z=0\), the full V27 scalar action obeys
\begin{equation}
 Q_a(q)\mathbf 1=0
 \quad\text{for every admissible }q.
 \label{eq:v33-shift}
\end{equation}
Every kinetic difference annihilates a constant; the continued
counterterm \(E_{2,\mu\nu}\) contains a derivative on each scalar
leg; and both filters leave the constant field unchanged. The upper
periodic Hodge term has the same property. This proves
\eqref{eq:v33-shift} on the finite array, not only on a continuum
interpolation. Differentiating it in any retained metric direction
preserves the zero.

\begin{theorem}[No derivative-free matter source at the mass Taylor origin]
\label{thm:v33-floor}
For the actual ghost-number-zero matter-count-two bank
\eqref{eq:v32-bank}, every coefficient containing a minimal antifield
satisfies
\begin{equation}
 [\Gamma_{a,1}]_{I,K=0,z=0}=0.
 \label{eq:v33-source-floor}
\end{equation}
The assertion includes all constant metric marks permitted by arity five.
It is a statement about the constant Taylor coefficient only. It does
not remove a coefficient containing \(z\) or a nonzero momentum jet.
\end{theorem}
\begin{proof}
Use V32's joint formula, not separate species determinants. If an external
ordinary scalar is constant, \eqref{eq:v33-shift} gives
\(L_\phi=U_\phi=0\) at \(z=0\), including every retained metric
derivative. Consequently \(W_\phi=W_{\upsilon\phi}=0\).
This proves the assertion for the \(\eta\phi^2c\),
\(\sigma\phi^2cc\), \(\upsilon\phi c\) and
\(\eta\upsilon\phi cc\) families and their allowed constant
metric marks.

The remaining family is \(\upsilon^2q^bcc\), \(b\le1\).
At constant retained ghost \(u\), the actual ghost-action return
on a hard momentum \(p\) has the factor
\(\alpha_p=i\,u\cdot p\), also at a constant coframe background.
The two factors in \(W_{\upsilon\upsilon}=-NG_HN^T\) therefore
contain \(\alpha_p\alpha_{-p}=-\alpha_p^2=0\) in the
exterior algebra. Constant metric differentiation changes the even
matrices and not this common factor. There is no additional
\(D_\eta\) or \(D_\sigma\) factor at this arity. This proves
\eqref{eq:v33-source-floor}. All primitive ghost inverses and the
upper ghost-action return are those of V32; the latter vanishes on a
constant retained ghost when it contains \((I-S)u\).
\end{proof}

Let \(g_{a,2}=\mathbf P[\Gamma_{a,1}]_{M=2,A\le5}\), and let
\(g_{a,g}\) be the retained gravitational local coefficient, including
its established coordinate-measure convention. The purely metric
coordinate logarithm has \(t j=0\) and no direct matter coefficient.

\begin{proposition}[Exact projected native differential on the populated bank]
\label{prop:v33-projection}
On a common nonaliased soft coefficient panel,
\begin{equation}
 \mathbf P[s_a\Gamma_{a,1}]_{M=2,A\le5}
       =d_Hg_{a,2}+tg_{a,g}.
 \label{eq:v33-native-projection}
\end{equation}
Thus its combined canonical/reference local breaking, after the fixed
gravity normalization and an already inserted matter counterterm, is
\begin{equation}
 B_{a,2}=d_HU_{a,2}+tH_g-x_{a,2},
 \qquad U_{a,2}=g_{a,2}+C_{a,2}^{\rm old}.
 \label{eq:v33-finite-breaking}
\end{equation}
Here the old matter counterterm is the actual V29 insertion; any further
previously chosen parent is included only if it has actually been
inserted. Equation \eqref{eq:v33-finite-breaking} is not an assertion
of finite-filter nilpotence.
\end{proposition}
\begin{proof}
The ordinary chain projection is inherited. It remains to bound the
\emph{degree}, not the size, of a possible finite artifact contribution.
On the soft plateau all source maps equal their ordinary coordinate/Lie
polynomials. Free metric and scalar Euler discrepancies start at
weighted derivative degree six. The corrected nonlinear gravitational
Euler discrepancy starts at degree five; the continued corrected scalar
vertices have degree at least six. Multiplication by such a discrepancy
raises local weight by at least four or three, respectively.

An antifield-containing ghost-number-zero coefficient has
\(w=D+n_*\ge n_*\). With \(n_*\ge2\), even the increase by
three exceeds the retained weight four. With \(n_*=1\), the only
possible boundary term would require \(D=0\), and it vanishes by
\cref{thm:v33-floor}. The corresponding gravitational coefficient has
the inherited V23--V24 derivative floor. A functional with no antifield
can only be varied through the soft source maps, which agree exactly.
This proves \eqref{eq:v33-native-projection}. The same degree argument
applies to V29's counterterm: its source part has no derivative-free
coefficient, as retained in its actual compiler. In fact projection to
\(w<4\) does not need that last floor, since any antifield coefficient
has weight at least one. The signed measured identity then gives
\eqref{eq:v33-finite-breaking}. The auxiliary logarithmic amplitude is
already included in the gravitational coefficient; its direct matter
extraction and its \(t\)-image are zero.
\end{proof}

Finite artifact products outside \(\mathbf P\) have not disappeared.
For example the gravitational Hessian discrepancy times a local
\(\eta\phi^2c\) coefficient remains in V32's complementary balance.
The proposition is an exact local degree statement, not an improvement
of that nonlocal error by deletion.

\section{The simultaneous lower-reference derivative and its relative term}
\label{sec:v33-transgression}

Use the full ordinary hard field list \((q,\phi,c,\bar c,b)\), with
nonminimal variables retained until its algebraic elimination. Let
\(\mathbb H_0(z)\) be its actual free Hessian and
\(V_0=\mathbb S_\infty''-\mathbb H_0\). Write
\begin{equation}
 C_0=(1-\Theta_\lambda)\mathbb H_0(z)^{-1},\qquad
 \dot C_0=\lambda\partial_\lambda C_0
      =-(\lambda\partial_\lambda\Theta_\lambda)
                                      \mathbb H_0(z)^{-1}.
 \label{eq:v33-dotC}
\end{equation}
All inverse expressions mean the operative hard-support covariance;
retained zero modes are not inverted. The scalar block is the actual
massive inverse. For example
\[
 \partial_z^jC_H=(-1)^jj!(1-\Theta_\lambda)(p^2+z)^{-j-1}.
\]
It has one cutoff factor. At the mass Taylor origin all these derivatives
are regular on the support of \(\dot C_0\).

Define the ordinary local scale coefficient
\begin{equation}
 \begin{split}
 Y_{\lambda}&=\mathbf P\operatorname{pr}_{A\le5,M\le2}
 \frac12\sum_{r=1}^{5}\frac{(-1)^{r+1}}r
 \sum_{i=0}^{r-1}\operatorname{STr}
 \left[(C_0V_0)^i(\dot C_0V_0)(C_0V_0)^{r-1-i}\right],\\
 Y_{g,\lambda}&=[Y_\lambda]_{M=0},\qquad
 Y_{2,\lambda}=[Y_\lambda]_{M=2}.
 \end{split}
 \label{eq:v33-Y}
\end{equation}
The notation \(Y_g\) suppresses the passive source-free scalar
background determinant described after \cref{prop:v33-target}; its
\(t\)-image is zero. Every primitive propagator is differentiated
once in the ordered sum. In particular a ghost inverse inside a Schur
source return is not held fixed. This formula is the scalar-coupled
operand, not a repeat of the gravitational-only integral of V24.

\begin{theorem}[Relative matter scale transgression]
\label{thm:v33-transgression}
In ordinary thermodynamic normalization the local reference components
and \eqref{eq:v33-Y} obey
\begin{align}
 d_gx_g&=0,& d_Hx_2+tx_g&=0,\
 d_gY_g&=\lambda\partial_\lambda x_g,&
 d_HY_2+tY_g&=\lambda\partial_\lambda x_2.
 \label{eq:v33-transgression}
\end{align}
For each homogeneous local weight, with momentum and mass jets taken
before rescaling,
\begin{equation}
 x_2^{[w]}=\lambda^{4-w}x_{2,1}^{[w]},\qquad
 Y_2^{[w]}=\lambda^{4-w}Y_{2,1}^{[w]}.
 \label{eq:v33-homogeneity}
\end{equation}
In particular the term \(tY_g\) is part of the measured transgression;
\(d_HY_2=\lambda\partial_\lambda x_2\) is generally false.
\end{theorem}
\begin{proof}
The zero-mesh common action \(\mathbb S_g+\mathbb S_H\) obeys
the ordinary classical master identity. Use the regulated trace identity
already proved in V22--V24, with this same full field list:
\(s_\infty\Gamma_1^{\mathbb R}+\operatorname{sdiv}r
=\mathcal X_{\mathbb R}\).
Its nilpotence and divergence identity give
\(s_\infty\mathcal X_{\mathbb R}=0\). Extract matter counts zero
and two, using \eqref{eq:v33-triangular}. This gives the first line.

For the scale derivative, first cancel \(C\mathbb R=\Theta\) in
reference words. A scale-differentiated determinant word has a factor
\(\dot C_0\), supported on the lower transition annulus. At the
retained arity, its other routed momenta differ by a sum of at most five
soft momenta. They therefore lie in a fixed enlarged lower annulus.
Every prescribed momentum and mass derivative is integrable there.
A compact upper guard can be removed from these words without changing
them. The auxiliary density and coordinate logarithm have no
\(\lambda\)-dependence. Differentiating the trace identity consequently
gives \(s_\infty Y=\lambda\partial_\lambda x\) without assigning
a value to an unsubtracted ultraviolet determinant. Its two components
are the second line of \eqref{eq:v33-transgression}. This is an
application of the inherited trace identity after verifying its mixed
operands and support, not an assumption of finite-cutoff BRST symmetry.

For scaling, expand the retained polynomial in \(z\) at zero and
then put \(p=\lambda y\) in each compact integral. A one-loop
matter-count-two ghost-number-zero component with \(n_*\) antifields
has superficial degree \(4-n_*\), as proved by the primitive counts
in V32. A reference component has one additional derivative and degree
\(5-n_*\). An external coefficient with \(|\alpha|\) momentum
derivatives and \(j\) mass derivatives lowers that degree by
\(|\alpha|+2j\). In its respective ghost-number sector, the result
is \(4-w\). This proves \eqref{eq:v33-homogeneity}. The local
counterterms are polynomial in \(z\); the physical massive propagator
has not been replaced by a massless one in the unsubtracted functional.
Exact rescaling is asserted in thermodynamic normalization, not on a
fixed finite torus.
\end{proof}

For a fixed labelled component \(I\), let \(r,s\) count its
metric and ghost antifields and set \(\kappa_I=\chi^{-(1+r+s)}\),
as in V32. On the fixed inverse-margin and profile class,
\begin{equation}
 \left|(Y_2)_I^{[D]}\right|
 \le C_I\kappa_I\lambda^{4-n_*-D},\qquad
 \left|(x_2)_I^{[D]}\right|
 \le C_I\kappa_I\lambda^{5-n_*-D}.
 \label{eq:v33-local-norms}
\end{equation}
These follow directly by rescaling the annulus, bounding its finite
volume and each prescribed primitive vertex/inverse derivative.
Their constants depend on the finite coefficient order, not the lattice
cutoff, volume or number of ultraviolet shells. They are not a bound on
\(Y\) in terms of \(x\) alone: closed local additions to a primitive
are real normalization freedom.

\section{A common subcritical mixed normalization and its cutoff error}
\label{sec:v33-subcritical}

First take the gravitational subcritical scale prescription already
constructed in V24:
\begin{equation}
 H_g^{[w]}=\frac{Y_g^{[w]}}{4-w},\qquad w<4.
 \label{eq:v33-paired-gravity}
\end{equation}
Its higher-weight and physical terms are discussed explicitly below.

\begin{theorem}[Simultaneous subcritical primitive relative to the paired reference]
\label{thm:v33-primitive}
For \(w<4\), define
\begin{equation}
 H_{2,\mathrm{can}}^{[w]}=\frac{Y_2^{[w]}}{4-w}.
 \label{eq:v33-primitive}
\end{equation}
Then
\begin{equation}
 d_HH_{2,\mathrm{can}}^{[w]}
           =x_2^{[w]}-tH_g^{[w]}.
 \label{eq:v33-canonical-solution}
\end{equation}
This one local functional includes action, transformation, mixed-antifield
and quadratic-scalar-antifield components. It therefore solves the
subcritical metric-marked action and scalar-algebra rows together.
The coefficient division in \eqref{eq:v33-primitive} has norm at most
one in the direct sum of the homogeneous coefficient norms. Combined
with \eqref{eq:v33-local-norms}, it is bounded independently of cutoff
and volume at each prescribed finite source order, and contains no
inverse mass or external momentum.

The same conclusion preserves a gravitational normalization of the
explicitly decomposed form
\begin{equation}
 H_g^{[w]}=\frac{Y_g^{[w]}}{4-w}+I_g^{[w]}(q)+d_gF_g^{[w]},
 \qquad d_gI_g=0,
 \label{eq:v33-physical-addition}
\end{equation}
by replacing \(H_{2,\mathrm{can}}^{[w]}\) by
\(Y_2^{[w]}/(4-w)+tF_g^{[w]}\).
An arbitrary closed gravitational source normalization is not assumed
here to come with such a decomposition if it has not been supplied.
\end{theorem}
\begin{proof}
At weight \(w\), \cref{thm:v33-transgression} gives
\(d_HY_2+tY_g=(4-w)x_2\). Divide by the positive integer
\(4-w\). This proves \eqref{eq:v33-canonical-solution} and
shows why the stress term must be paired with the same gravitational
reference primitive. In \eqref{eq:v33-physical-addition}, \(tI_g=0\)
since \(I_g\) has no gravitational antifield, while
\(td_gF_g=-d_HtF_g\). The stated extra term therefore gives exactly
the new relative right side. Locality, polynomial mass dependence and
the norm assertion follow from the finite annular compiler and the
nonzero integer division. There is no cohomology classification in
this argument.
\end{proof}

On the actual degree-five bank the canonical mixed contribution below
weight four has a still smaller support:
\begin{equation}
 H_{2,\mathrm{can}}^{[w]}=0\quad(w=0,1,3),\qquad
 H_{2,\mathrm{can},<4}=\tfrac12Y_2^{[2]}.
 \label{eq:v33-even-weights}
\end{equation}
To see this, ordinary two-derivative action vertices and one-derivative
source vertices, together with even propagators and the even reference
profile, give the ghost-number-zero component with \(n_*\) antifields
parity \((-1)^{n_*}\) in all external momenta. Thus \(D+n_*\)
is even; mass jets have even weight. A zero-weight matter action has
no antifield, derivative or mass factor, and its constant-scalar
coefficient vanishes by \eqref{eq:v33-shift}. These two observations
prove \eqref{eq:v33-even-weights} for the ordinary scale compiler.
They are not a parity assertion for arbitrary finite oriented vertices.
The denominators \(4,3,2,1\) in the general theorem are retained because
physical and explicitly chosen scheme additions need not use this
particular parity-reduced representative.

\begin{proposition}[Finite reference error in the complete matter-count-two bank]
\label{prop:v33-reference-rate}
Fix
\begin{equation}
 \lambda\ge 32768\mu>0,\qquad
 a\lambda\le2^{-24},\qquad
 \lambda L_\nu\ge1,\qquad 0\le m\le M\lambda.
 \label{eq:v33-domain}
\end{equation}
For every retained ghost-number-one component and
\(D=|\alpha|+2j\le5-n_*\), its local reference coefficient obeys
\begin{equation}
 \left|(x_{a,2}-x_{0,\lambda,\infty,2})_I^{[D]}\right|
 \le C_{I,b}\kappa_I\lambda^{5-n_*-D}
       \left[(a\lambda)^3+(\lambda L_{\min})^{-b}\right],
 \quad b>4.
 \label{eq:v33-reference-rate}
\end{equation}
All mass jets are taken in the convention of \(\mathbf P\).
There is no omitted-zero-mode contact in this matter-count-two
reference estimate.
\end{proposition}
\begin{proof}
Use the full-field ordered reference expansion through five vertices.
A no-covariance diagonal trace of \(\Theta Dr\) has matter count
zero: the scalar--scalar derivative of the scalar Lie source is
independent of \(\phi\), its scalar--ghost derivative is off diagonal,
and the geometric source has no scalar dependence. Thus every contributing
matter-count-two word has a covariance as well as \(\Theta\).
At the small external transfers in \eqref{eq:v33-domain}, all its
routed labels lie in \(\lambda/2\le|p|\le3\lambda\). The
field/source completion profiles equal one there.

On that annulus, the free inverse differences and the continued scalar
vertex differences start at fourth relative mesh order. A retained
metric background can use a third-relative-order gravitational vertex.
An ordered telescope of the finitely many primitive factors bounds the
whole difference by \(C(a\lambda)^3\) times its ordinary scale.
Momentum differentiation lowers scale by one and mass differentiation
by two, with the single-cutoff-factor rule of \eqref{eq:v33-dotC}.
This proves the mesh term. Poisson summation after more than \(b+4\)
fixed derivatives of the smooth compactly supported integrand proves
the stated separate volume term. The annular support excludes the
retained origin, explaining the final assertion.
\end{proof}

If the canonical subcritical comparison scheme is chosen, its incremental
finite counterterm is
\begin{equation}
 \Delta C_{a,2,<4}=-(U_{a,2})_{<4}+H_{2,\mathrm{can},<4}.
 \label{eq:v33-finite-update}
\end{equation}
The old mixed counterterm is included in \(U_{a,2}\), not added
again. Combining \eqref{eq:v33-finite-breaking} and
\eqref{eq:v33-canonical-solution} gives
\begin{equation}
 (B_{a,2}^{\mathrm{can}})_{<4}
          =(x_{0,\lambda,\infty,2}-x_{a,2})_{<4}.
 \label{eq:v33-finite-result}
\end{equation}
This restores both \(q c\phi^2\) and \(\upsilon\phi cc\)
through their subcritical degrees, with \eqref{eq:v33-reference-rate}.
For the former these are degrees at most four, and for the latter at
most three. Critical weight four is not included.

For a fixed smooth external window, the residual polynomial in each
component has rooted kernel norm at most
\[
 C_{I,s,b}\kappa_I
 \sum_{D\le5-n_*}\lambda^{5-n_*-D}(\mu+m)^D
       \left[(a\lambda)^3+(\lambda L_{\min})^{-b}\right].
\]
Fourier integration by parts in its finite number of relative coordinates
proves this bound; periodization does not enlarge the physical weighted
\(L^1\) norm. It yields a multilinear bound with one source in \(L^1\)
and all others in \(L^\infty\), without a total-volume multiplier.
The counterterm coefficients themselves are retained normalization data;
this residual estimate is not a bound on a full nonlinear action.

\section{Why literal preservation of the first-row coefficients is an extra condition}
\label{sec:v33-relative-test}

The canonical comparison just constructed preserves the chosen pure-gravity
normalization and solves the earlier mixed \emph{equation} on its
subcritical projection. It need not preserve the particular coefficients
of V29's \(h_H\), including its choice of scalar kinetic normalization.
There is no automatic implication from one assertion to the other.

Let \(L\) restrict a matter-count-two ghost-number-zero functional to
the precise V29 coefficients which are to be held fixed, for example
\(\phi^2,q\phi^2,\upsilon\phi c\) through the declared weights.
In the finite subcritical quotient let \(Z\) synthesize a basis of
\(\ker d_H\). Reusing the ordinary finite relative range test already
established in V25 gives
\begin{equation}
 \boxed{L(h_H-H_{2,\mathrm{can}})\in\operatorname{Ran}(LZ).}
 \label{eq:v33-relative-test}
\end{equation}
This is necessary and sufficient for a simultaneous solution with exactly
those coefficients. No rank or native right side of this particular
relative system has been computed here. In particular the comparison
scheme \eqref{eq:v33-finite-update} is not silently substituted for
V29's current scheme.

The need for \eqref{eq:v33-relative-test} is substantive even below
critical weight. It is not merely a numerical-conditioning caveat.

\begin{theorem}[A first-row homogeneous freedom which has no fixed-gravity extension]
\label{thm:v33-nonextension}
At zero mass consider the homogeneous change of the cubic matter
normalization
\begin{equation}
 C_\alpha=\alpha\left(S_{H,3}
               +\langle\upsilon,\mathcal L_c\phi\rangle\right),
 \qquad \alpha\ne0,
 \label{eq:v33-test-counterterm}
\end{equation}
where \(S_{H,3}\) is the actual zero-mesh \(q\phi^2\) kinetic
vertex in the affine coframe chart. It belongs to weight two. Its first
antifield-free \(c\phi^2\) row is zero. Nevertheless it cannot be
extended by terms of arity at least four, with the pure-gravity
normalization and its existing scalar coefficients fixed, to a
\(d_H\)-closed matter-count-two functional through arity four.
This statement permits quadratic-antifield terms.
\end{theorem}
\begin{proof}
The free scalar Noether identity is
\(D_qS_{H,3}[R_0c]+\langle P_0\phi,\mathcal L_c\phi\rangle=0\).
It follows either by differentiating the ordinary scalar action or by
contracting its explicit vertex. In the all-incoming convention,
\(k+r+t=0\), set
\[
 T_{\mu\nu}=r_\mu t_\nu+r_\nu t_\mu
                     -\delta_{\mu\nu}(r\cdot t-z),
 \qquad U_\nu(k,p)=p_\nu .
\]
Then
\(k^\mu T_{\mu\nu}+P(r)U_\nu(k,t)+P(t)U_\nu(k,r)=0\)
identically, even before setting \(z=0\). This is a homogeneous
solution of V29's first row with \(f=0\).

Its scalar transformation perturbation is \(Q_c=\alpha\mathcal L_c\).
V30's exact variation formula gives at zero metric
\[
 Q_{c\cdot\partial c}-\mathcal L_cQ_c-Q_c\mathcal L_c
                  =-\alpha\mathcal L_{c\cdot\partial c}.
\]
Test the local jets of two affine Killing fields
\(X=\partial_0\) and
\(Y=x^0\partial_1-x^1\partial_0\), for which
\([X,Y]=\partial_1\). Put \(c=\theta_1X+\theta_2Y\), with
independent odd labels. Both \(R_0X\) and \(R_0Y\) vanish.
On the scalar stationary jet \(\phi=\varepsilon x^1\),
\(P_0\phi=-\partial^2\phi=0\), whereas the coefficient of
\(\theta_1\theta_2\) in the displayed defect is
\(-\alpha\varepsilon\ne0\).

An arity-four term with a metric could contribute to this zero-metric
row only through \(R_0c\), which vanishes on these Killing jets.
An additional scalar antifield could contribute only through its free
Euler replacement, which vanishes on this stationary jet. A metric
antifield contributes the free gravitational Euler row, zero at the
flat metric; its scalar stress has at least two scalar factors and
cannot cancel a term linear in \(\phi\). A ghost-antifield correction
to the pure gravitational bracket would change the prescribed gravity
normalization and is excluded. Higher arity cannot return to the tested
arity-four coefficient. These are all possibilities under the degree
rules. Thus the nonzero coefficient cannot be cancelled by an allowed
extension.

The test concerns local jets and can be placed inside a region where
smooth test functions agree with the displayed polynomials. It does not
postulate a periodic affine ghost or integrate an unbounded scalar
configuration. The factor \(\varepsilon\) isolates the matter-linear
coefficient before any scalar-stress term.
\end{proof}

The theorem is not a claim that V29's particular native coefficient
contains \eqref{eq:v33-test-counterterm}. Its coefficient there has not
been evaluated. It proves that testing only the first action-row image
cannot establish extendibility, even for a mass-regular, two-derivative
homogeneous choice. The new comparison primitive and the old one must be
compared in the full complex or by an equally sufficient relative test.
The counterexample would be invisible to a calculation which set the
scalar Euler terms or the ghost bracket aside at the wrong stage.

\section{What is left for the critical block and the scalar Euler section}
\label{sec:v33-next}

At weight four the transgression gives only
\begin{equation}
 d_HY_2^{[4]}+tY_g^{[4]}=0.
 \label{eq:v33-critical}
\end{equation}
The unsolved reference equation is still
\begin{equation}
 d_HH_2^{[4]}=x_2^{[4]}-tH_g^{[4]},
 \label{eq:v33-critical-target}
\end{equation}
with the actual critical pure-gravity primitive, not a rescaled
subcritical substitute. Neither division by \(4-w\) nor the
pure-gravity classification of V25 supplies its relative scalar solution.
A scalar-inclusive existence theorem, if used, would additionally have
to preserve the prescribed lower coefficients and provide a bound for
the relevant projection; none is imported in this part.

In the canonical subcritical comparison scheme,
\eqref{eq:v33-finite-result} supplies both the action and scalar-algebra
residuals as one reference discrepancy. Applying the already proved
\eqref{eq:v31-coupled-consistency} therefore bounds its scalar Euler
compatibility residual by a fixed finite coefficient-map norm times that
same discrepancy and the fixed gravity comparison. V31's quadratic-
antifield section can then remove its Euler-compatible component without
altering the action rows. This is a consequence of a common normalization;
it is not a new proof of V31's polynomial section.

For the current V29 coefficient prescription, the same conclusion needs
\eqref{eq:v33-relative-test}. Otherwise the unremoved first-background
action coefficient reappears in \(\mathscr JD\), as V31 explicitly
requires. The nonextension example makes it inappropriate to assume that
this is merely an unrecorded higher-metric contact which can always be
chosen.

There are thus two sharply specified remaining local questions, not a
new background-by-background loop task. First compare the particular
V29 subcritical coefficients with the native simultaneous comparison
primitive in the relative kernel image. Second solve the weight-four
relative equation with the same lower normalizations. The primitive
integrals in this installment and the V29 coefficient compiler give their
right sides. Their numerical values and the relative matrix range have
not been evaluated. The internal-gauge/chiral sectors, higher-loop
filtered shell control and infrared removal remain separate.

\section{Verification and closing ledgers for V33}
\label{sec:v33-ledgers}

The accompanying verifier
\srcid{check_ESM_relative_transgression_V33.py} checks the finite
source-family enumeration and weights, the two-scalar cubic homogeneous
identity, the affine-Killing nonextension witness, the massive covariance
scale/mass derivatives, and the ordered trace derivative. It also tests
the triangular relative identities on exact finite complexes and the
source-floor factorization in a finite exterior algebra. Such model
checks verify algebraic formulas and signs; they are not evaluations of
the native annular tensor integrals or independent verification of the
inherited continuum estimates. Its relative-range utility accepts an
explicit finite coefficient matrix; the full native matrix is not
constructed in this installment.

\subsection*{Proved}
\addcontentsline{toc}{subsection}{V33: Proved}
The massless-shift identity of the actual scalar completion removes the
constant Taylor coefficients of all minimal matter-count-two sources at
arity five. It licenses the exact local finite-to-ordinary projection on
the populated bank. The simultaneous scale derivative has the relative
term \(tY_g\) and supplies a bounded common subcritical primitive when
paired with the same gravitational reference normalization. Its reference
comparison has the displayed annular cutoff and volume bound. A precise
first-row homogeneous freedom is shown not to extend with fixed gravity,
so literal preservation of the old row coefficients is a real additional
relative requirement. No value is assigned to an unevaluated native
matching integral.

\subsection*{Inherited}
\addcontentsline{toc}{subsection}{V33: Inherited}
V1--V32 remain unchanged. In particular V31 already proves the triangular
matter filtration and scalar Euler section; V24--V29 supply the ordinary
arity/weight chain projection and trace identity; V32 supplies the actual
simultaneous Hessian, ghost returns and symbol estimates. The current
V29 local coefficients are not silently replaced. The new comparison
scheme has been specified separately, including its possible change of
mixed finite normalization. Pure-gravity numerical/source choices beyond
the explicitly stated decomposition retain their own relative tests.

\subsection*{Literature input}
\addcontentsline{toc}{subsection}{V33: Literature input}
The relation between exact-RG reference identities and quantum BV is
established methodology, for example \cite{V28IIM}. Local BRST
cohomology and the distinction between consistency and a local primitive
are comparison context, not an imported proof that
\eqref{eq:v33-critical-target} or \eqref{eq:v33-relative-test} has a
solution. The proofs above use explicit graded identities, finite native
vertices and the identified inherited reference estimates.

\subsection*{Conditional}
\addcontentsline{toc}{subsection}{V33: Conditional}
Adopting the simultaneous canonical subcritical scheme gives the common
restored rows with the stated error. Preserving \emph{every} specified
V29 coefficient additionally requires the uncomputed relative range test.
The theorem accommodates the explicit gravitational physical additions
in \eqref{eq:v33-physical-addition}; an unspecified antifield-dependent
closed addition is not assumed to have that decomposition. No arbitrary
all-loop or analytic action-class result follows from these one-loop,
finite-arity local statements.

\subsection*{Open}
\addcontentsline{toc}{subsection}{V33: Open}
Evaluate the relative comparison for V29's populated coefficient array
and the critical scalar--gravity equation
\eqref{eq:v33-critical-target}. Keep their scalar, metric, ghost and
higher-antifield descendants in the same degree-closed bank. Numerical
native integrals and a compressed relative matrix remain to be computed
if a numerical counterterm list is desired. Combined internal-gauge and
chiral-line transport, higher loops, the filtered R1--R4 invariant-class
realization and the full ESM source continuum remain unproved here.

\begin{center}\small
\begin{tabular}{>{\raggedright\arraybackslash}p{.24\textwidth}>{\raggedright\arraybackslash}p{.67\textwidth}}
\toprule
Variable & Scope in V33\\\midrule
Quantum/EFT order & One loop, local weight at most four and total arity
at most five. The new constructive primitive is subcritical.\\
Matter and couplings & Matter count two, four equal real scalars, no
internal gauge, Yukawa or scalar self-coupling. \(\chi>0\) fixed.\\
Mass & Physical \(0\le m\le M\lambda\); local polynomial jets in
\(z=m^2\) at zero. No inverse mass or infrared limit.\\
Cutoff and volume & Native coefficient tests are finite. Quantitative
reference comparison is uniform on \eqref{eq:v33-domain}; exact scaling
uses thermodynamic normalization.\\
Sources and norms & Fixed labelled finite panels and prescribed derivative
and rooted-kernel orders. No total-volume factor or unbounded-order
uniformity is asserted.\\
Normalization & Canonical paired-reference comparison constructed;
literal V29 coefficient matching remains separately testable, not assumed.\\
Iteration and measure & Only finitely many lower-annulus reference levels;
old auxiliary measure counted once. No nonlinear RG self-map or positive
continuum gravitational measure is concluded.\\
\bottomrule
\end{tabular}
\end{center}


\clearpage
\part*{V34: A populated relative-normalization obstruction and a compatible scalar kinetic prescription}
\addcontentsline{toc}{part}{V34: Populated kinetic obstruction and compatible normalization}

\section{The particular relative test can be decided before the large matrix is built}
\label{sec:v34-start}

Sections 1--263 and their labels are preserved. V33 left the relative
range test \eqref{eq:v33-relative-test} unevaluated for V29's particular
coefficient array. The issue was not whether the first action row has a
primitive: V29 already constructed one. It was whether that particular
primitive, including its exactly zero flat scalar kinetic coefficient,
can be extended with the gravitational normalization held fixed.

This part evaluates a separating component of that actual range test.
The flat weight-two part of V33's simultaneous primitive has a nonzero
time--space kinetic difference. We prove that every translation-invariant,
matter-count-two homogeneous BV cocycle at the same weight has an
\emph{isotropic} flat kinetic tensor. Thus the difference cannot be
removed by an allowed homogeneous addition. One explicit left-null test
of the relative system already decides failure; the full relative matrix
and the many other annular coefficients need not be computed to reach
that conclusion.

The positive continuation is recorded, not hidden: from this part onward
we use a simultaneous subcritical normalization whose flat scalar mass
coefficient and mean kinetic coefficient are zero, while retaining the
uniquely required trace-free kinetic tensor. The full functional, including
its antifield and higher-metric coefficients, is specified below. It is
not obtained by changing just one component of V29's counterterm. V29's
first-row result remains true in its original scheme, but that scheme is
not the active subcritical prescription for the coupled continuation on
the branch considered here. Its historical body is unaltered.

The branch is precisely the paired reference of
\eqref{eq:v33-paired-gravity}, at weight two:
\begin{equation}
 H_g^{[2]}=\tfrac12Y_g^{[2]}+I_g^{[2]}(q),\qquad
 d_gI_g^{[2]}=0,\qquad tI_g^{[2]}=0.
 \label{eq:v34-gravity-branch}
\end{equation}
The last two conditions allow the already specified antifield-free
physical additions. An independently chosen exact gravitational
source addition \(d_gF_g\) can change the test through \(tF_g\); its
effect is stated separately and is not silently set to zero.

All other conventions are inherited: the four equal real scalars,
zero internal gauge, Yukawa and self-couplings, the affine Euclidean
coframe coefficient chart, the fixed metric contour, the massive hard
scalar reference, the once-counted auxiliary density, the V18/V28 source
prescriptions, and the ordinary local differential \(d_H\). Work is at
one loop, matter count two, total arity at most five. The common factor
\(\hbar\) is omitted. The squared mass \(z=m^2\) has weight two.
No new field, propagator, ultraviolet completion or cohomology
classification is introduced.

\section{The flat kinetic coefficient of the simultaneous primitive}
\label{sec:v34-kinetic}

Write \(\Theta_\lambda(p)=\theta(p/\lambda)\),
\(\nu_\lambda=1-\Theta_\lambda\). We use the inherited nonnegative
profile class
\begin{equation}
 0\le\theta\le1,\qquad
 \theta(y)=1\ (|y|\le1),\qquad
 \theta(y)=0\ (|y|\ge2),
 \label{eq:v34-profile}
\end{equation}
with the stated smoothness and fixed seminorms. Radial symmetry is not
needed for the kinetic separation proved here. In thermodynamic
normalization put
\begin{equation}
 \mathfrak a_\theta=\int\frac{\dd^4y}{(2\pi)^4}
                         \frac{\theta(y)}{|y|^2},\qquad
 \mathfrak b_\theta=\int\frac{\dd^4y}{(2\pi)^4}
                         \frac{\theta(y)^2}{|y|^2}.
 \label{eq:v34-moments}
\end{equation}
These are finite positive profile moments. The origin is integrable in
four dimensions; no inverse is being applied to a retained zero mode.
They obey
\begin{equation}
 \frac1{16\pi^2}\le\mathfrak b_\theta\le\mathfrak a_\theta
                                  \le\frac1{4\pi^2}.
 \label{eq:v34-moment-bounds}
\end{equation}
For a radial profile they equal
\((8\pi^2)^{-1}\int_0^2r\theta(r)\dd r\) and its
\(\theta^2\) counterpart. The inequalities follow by comparison with
the unit and radius-two balls, not by evaluating a particular smooth
transition numerically.

Define the flat coefficient by
\begin{equation}
 [H_{2,\mathrm{can}}^{[2]}]_{q=\eta=\sigma=\upsilon=0}
  =\frac12\sum_A\int_k\phi_A(-k)
          \bigl(k^TF_{\lambda,\theta}k+d_{\lambda,\theta}z\bigr)
                                                  \phi_A(k).
 \label{eq:v34-flat-definition}
\end{equation}
Here \(H_{2,\mathrm{can}}^{[2]}=Y_2^{[2]}/2\) is the actual
V33 simultaneous primitive, not an arbitrary solution of its equation.
Let \(P_t=e_0e_0^T\) in the inherited affine coframe coordinates.

\begin{theorem}[Evaluated flat weight-two normalization]
\label{thm:v34-flat}
On \eqref{eq:v34-gravity-branch},
\begin{equation}
 \boxed{
 F_{\lambda,\theta}
   =\frac{\lambda^2}{\chi}
        \left(\frac52\mathfrak a_\theta P_t
                           -\mathfrak b_\theta I_4\right),\qquad
 d_{\lambda,\theta}
   =\frac{9\lambda^2}{4\chi}\mathfrak a_\theta .}
 \label{eq:v34-flat-value}
\end{equation}
In particular, for every spatial \(i\),
\begin{equation}
 \boxed{
 (F_{\lambda,\theta})_{00}-(F_{\lambda,\theta})_{ii}
       =\frac{5\lambda^2}{2\chi}\mathfrak a_\theta>0.}
 \label{eq:v34-gap}
\end{equation}
The coefficient is per physical volume and per fixed external scalar
pair; there is no additional factor of four. The scalar mass in the
unsubtracted functional is not set to zero: \eqref{eq:v34-flat-value}
is its weighted Taylor coefficient in \((k,z)\).
\end{theorem}
\begin{proof}
The two inputs are already evaluated in V27: the native coframe contact
\eqref{eq:v27-contact} and the complete signed mixed bubble
\eqref{eq:v27-cont-bubble}. We use them, rather than recomputing a new
matter determinant. Write
\[
 K_{\rm con}=\operatorname{diag}(-1/2,2,2,2).
\]
The contact is
\(\chi^{-1}I_\nu(k^TK_{\rm con}k-9z/4)\), with
\(I_\nu=\int_p\nu_\lambda(p)/p^2\).
At \(z=0\), the exact mixed numerator cancels the scalar denominator,
so its weight-two coefficient is
\(-\chi^{-1}k^2I_{\nu^2}\), where
\(I_{\nu^2}=\int_p\nu_\lambda(p)^2/p^2\).
Its coefficient linear in \(z\) at \(k=0\) is zero: the numerator
there is \(-z^2\). Thus, before taking the lower-scale derivative, the
only weight-two terms are
\begin{equation}
 \frac1\chi\left[
 I_\nu(k^TK_{\rm con}k-9z/4)-I_{\nu^2}k^2\right].
 \label{eq:v34-guarded-input}
\end{equation}
The undifferentiated integrals in this display are ultraviolet divergent.
Introduce one common smooth upper guard independent of \(\lambda\),
equal to one throughout the support of every scale-differentiated
factor, and differentiate before removing that guard. No independently
routed value is assigned to either divergent integral.

Both derivatives are compactly supported and satisfy exactly
\begin{align}
 \dot I_\nu&=-\int_p\frac{\dot\Theta_\lambda(p)}{p^2}
                 =-2\lambda^2\mathfrak a_\theta,\nonumber\\
 \dot I_{\nu^2}
   &=-2\int_p\frac{\nu_\lambda(p)\dot\Theta_\lambda(p)}{p^2}
                 =-4\lambda^2\mathfrak a_\theta
                              +2\lambda^2\mathfrak b_\theta.
 \label{eq:v34-dot-moments}
\end{align}
For example \(\int_p\Theta_\lambda/p^2
=\lambda^2\mathfrak a_\theta\); differentiating this finite integral
proves the first identity. Applying the same argument to
\(\Theta_\lambda^2\) proves the second. These identities require
scaling of the profile, not spherical symmetry.

V33's \(Y_2\) is the logarithmic lower-scale derivative of the
simultaneous determinant, including every differentiated covariance.
Its weight-two primitive is half that derivative. Applying
\eqref{eq:v34-dot-moments} to \eqref{eq:v34-guarded-input} gives
\[
 F_{\lambda,\theta}
 =\frac{\lambda^2}{\chi}
       [-\mathfrak a_\theta K_{\rm con}
                    +(2\mathfrak a_\theta-\mathfrak b_\theta)I_4],
\]
and \(d_{\lambda,\theta}=9\lambda^2\mathfrak a_\theta/(4\chi)\).
Since \(2I_4-K_{\rm con}=5P_t/2\), this is
\eqref{eq:v34-flat-value}. The bubble contributes only an isotropic
kinetic coefficient at this weight; it cannot cancel the contact's
anisotropic part. The ghost determinant has no ordinary scalar
background, and the auxiliary density has no \(\lambda\)-derivative
or scalar dependence in this extraction. No omitted term in the common
integral supplies an additional contribution to this component.
\end{proof}

The inherited contact can also be checked without a loop calculation:
its coframe polynomial obeys
\[
 \sum_{A,B}J_{AB}\,\partial_A\partial_BW_2
       =K_{\rm con},\qquad
 \sum_{A,B}J_{AB}\,\partial_A\partial_B\rho_2=-9/4,
\]
with the actual ten symmetric coordinate covariances of V13. The new
checker independently verifies these finite rational contractions. They
are not presented as new seagull theorems.

\section{A Killing-jet restriction on homogeneous subcritical counterterms}
\label{sec:v34-killing}

A homogeneous difference between two simultaneous solutions on the
same gravity/reference branch satisfies
\(d_H\Delta=0\), modulo total derivatives, through arity five.
It is translation invariant and has matter count two. This section
restricts its weight-two component only. Coefficients may depend on the
fixed reference parameters; they have no explicit spacetime-coordinate
or inverse-mass dependence. The ordinary coframe chart and scalar
transformation are unchanged.

At zero metric and zero antifields write
\begin{equation}
 \Delta_{\rm flat}=\frac12\sum_A\int
       \left(F^{\mu\nu}\partial_\mu\phi_A\partial_\nu\phi_A
                                  +b\,z\phi_A^2\right),\qquad F=F^T.
 \label{eq:v34-hom-flat}
\end{equation}
The zero-metric scalar-antifield part has one derivative. Its most
general flavour-diagonal scalar transformation is consequently
\begin{equation}
 Q_X\phi=(MX)^\rho\partial_\rho\phi
                         +N_{\nu\rho}\partial_\rho X^\nu\,\phi,
 \label{eq:v34-Q}
\end{equation}
for constant real matrices \(M,N\). In Fourier notation these are the
coefficients linear in the scalar momentum and in the ghost momentum.
Terms with higher metric degree and the other minimal antifields have
not been excluded from \(\Delta\); their vanishing on the tests below
is proved separately.

\begin{lemma}[The homogeneous scalar source is trivial on flat Killing jets]
\label{lem:v34-source-killing}
If \(d_H\Delta=0\) at the stated arity and weight, then
\begin{equation}
 M=0,\qquad N=N^T.
 \label{eq:v34-MN}
\end{equation}
Hence \(Q_X=0\) for every flat translation or rotation \(X\).
\end{lemma}
\begin{proof}
Test the scalar-antifield/two-ghost coefficient at \(q=0,z=0\) on
flat Killing vector fields and on scalar stationary jets. The free metric
variation of any Killing field is zero. Thus terms with an additional
metric can contribute nothing by a free metric replacement. Terms with
a further scalar antifield replace it by the scalar Euler expression,
which is zero on the stationary jet. A metric antifield replaces to the
pure-gravity Euler expression, also zero at the flat metric. Matter
stress from such a replacement would raise matter count and is absent
in \(d_H\) on this quotient. A matter-dependent ghost-bracket change
would first change the scalar transformation at matter count greater
than the tested one. Independent nonminimal doublets are put to zero
only after restoring their off-shell equations; their restriction is a
chain map. Therefore the remaining identity is
\begin{equation}
 Q_{[X,Y]}-[\mathcal L_X,Q_Y]-[Q_X,\mathcal L_Y]=0
 \quad\text{on scalar stationary jets}.
 \label{eq:v34-cocycle-Killing}
\end{equation}

Choose \(X=t^\mu\partial_\mu\) and
\(Y=(\Omega x)^\mu\partial_\mu\), where \(\Omega^T=-\Omega\).
Substitution of \eqref{eq:v34-Q} into
\eqref{eq:v34-cocycle-Killing} gives
\(-\Omega Mt\cdot\partial\phi\). Affine scalars are massless
stationary jets with arbitrary gradient. Thus \(\Omega Mt=0\) for
all \(t\) and all rotations. The intersection of the kernels of the
six rotation generators is zero, so \(M=0\).

Next take two rotations. With \(M=0\), both \(Q_X\) and \(Q_Y\)
are constant multiplication operators. Their commutators with the Lie
operators vanish. Their bracket condition therefore states that \(N\)
annihilates every rotation commutator. Such commutators span
\(\mathfrak{so}(4)\): for distinct indices each elementary rotation
is a bracket of the two rotations through a third index. Hence the
antisymmetric part of \(N\) is zero. A symmetric \(N\) annihilates
the antisymmetric derivative of a rotation and also the zero derivative
of a translation, proving the last assertion.

These are tests of local differential-polynomial identities. No affine
field is asserted to be periodic. They can equivalently be taken in a
smooth region where the test jets have the indicated values. The use
of \(z=0\) introduces no inverse mass: \(Q_X\) at weight two cannot
contain a factor \(z\), so the conclusion concerns its entire retained
polynomial.
\end{proof}

\begin{theorem}[Exact flat-action projection of the homogeneous kernel]
\label{thm:v34-kinetic-kernel}
For translation-invariant, flavour-diagonal, weight-two matter-count-two
local functionals through arity five,
\begin{equation}
 \boxed{
 \operatorname{pr}_{\phi^2}\ker d_H
    =\left\{\tfrac12\int\phi(-k)(\alpha k^2+\beta z)\phi(k):
                                         \alpha,\beta\in\R\right\}.}
 \label{eq:v34-kernel-projection}
\end{equation}
Thus all nine trace-free kinetic coefficient functionals annihilate the
homogeneous kernel. This is a projection statement, not the cohomology
of the complete matter BV complex.
\end{theorem}
\begin{proof}
The first action row on a flat Killing rotation has no metric variation.
By \cref{lem:v34-source-killing}, its scalar transformation correction
also vanishes. The other antifield replacements vanish at the stated
zero-background restriction. Hence the quadratic action
\eqref{eq:v34-hom-flat} is invariant under every rotation, on arbitrary
scalar tests, not just stationary ones. Since a rotation Lie derivative
is skew-adjoint in the flat pairing, this says that its commutator with
the constant-coefficient scalar quadratic operator is zero. Equivalently
\([F,\Omega]=0\) for every antisymmetric \(\Omega\). A real symmetric
matrix commuting with all elementary rotations is a scalar multiple of
\(I_4\). The coefficient of \(z\) is already a scalar, proving the
inclusion into the right-hand side.

For the reverse inclusion use the explicit ordinary covariant
functionals
\begin{equation}
 \mathcal K_H=\frac12\sum_A\int \rho(q)G(q)^{\mu\nu}
                         \partial_\mu\phi_A\partial_\nu\phi_A,
 \qquad
 \mathcal V_H=\frac z2\sum_A\int\rho(q)\phi_A^2.
 \label{eq:v34-closed-functionals}
\end{equation}
Their Lie variations are total derivatives. They are antifield-free,
so \(d_H\mathcal K_H=d_H\mathcal V_H=0\) in this matter quotient.
Their expansions truncated after arity five stay closed because the
ordinary differential cannot lower arity. Their flat quadratic kernels
are \(k^2\) and \(z\), respectively. Every pair \((\alpha,\beta)\)
is therefore attained.
\end{proof}

The new finite certificate has small size. In component bases, the
translation--rotation constraints on \(M\) have rank 16, the rotation
commutator constraints on \(N\) have rank 6, and the six commutators of
a symmetric kinetic matrix have rank 9 with kernel \(\R I_4\).
The checker verifies these ranks in exact rational arithmetic. None is
an estimate of the large relative matrix considered in V33.

\section{The populated V29 relative range test fails on the paired branch}
\label{sec:v34-obstruction}

\begin{theorem}[A nonzero separator for the actual first-row prescription]
\label{thm:v34-no-extension}
On \eqref{eq:v34-gravity-branch}, no simultaneous weight-two
matter-count-two solution of
\(d_HH_2=x_2-tH_g\) through arity five can retain V29's prescribed
flat scalar coefficient \(f=0\). In particular it cannot retain all
of V29's \(\phi^2,q\phi^2,\upsilon\phi c\) coefficients.
For each spatial \(i\), the exact separator of the relative test is
\begin{equation}
 \boxed{
 \ell_i\bigl(h_H^{[2]}-H_{2,\mathrm{can}}^{[2]}\bigr)
 =-\frac{5\lambda^2}{2\chi}\mathfrak a_\theta\ne0,
 \qquad \ell_i(F)=F_{00}-F_{ii}.}
 \label{eq:v34-native-separator}
\end{equation}
Its magnitude has the profile-uniform lower bound
\begin{equation}
 |\ell_i|\ge\frac{5\lambda^2}{32\pi^2\chi}.
 \label{eq:v34-separator-lower}
\end{equation}
\end{theorem}
\begin{proof}
The particular V29 compiler in \cref{thm:v29-primitive} sets \(f=0\)
exactly. V33's simultaneous solution has the flat tensor evaluated in
\cref{thm:v34-flat}. If a solution with the prescribed V29 coefficient
existed, its difference from \(H_{2,\mathrm{can}}^{[2]}\) would be a
homogeneous \(d_H\)-cocycle with flat tensor \(-F_{\lambda,\theta}\).
By \cref{thm:v34-kinetic-kernel}, every \(\ell_i\) must annihilate
that difference. Equation \eqref{eq:v34-gap} contradicts this. The lower
bound follows from \eqref{eq:v34-moment-bounds}.

More explicitly, let \(L\) and \(Z\) be V33's restriction and kernel
synthesis. The row extracting \(F_{00}-F_{ii}\) satisfies
\(\ell_iLZ=0\), but its pairing with the actual right-hand side is
\eqref{eq:v34-native-separator}. This is an evaluated left-null
certificate for nonmembership in \(\operatorname{Ran}(LZ)\). No rank
of the full matrix, no value of the other source coefficients and no
numerical annular quadrature is needed.
\end{proof}

This answers the subcritical relative question negatively for the
\emph{actual} V29 choice, rather than just exhibiting a possible bad
homogeneous freedom as V33 did. It does not contradict the proof of
V29's first row: a right inverse for one row need not be a section of
the coupled complex. The proof also shows that changing only the
scalar-antifield or higher-metric coefficients cannot repair the
prescription while retaining \(f=0\).

There are precise boundaries to this conclusion. A different
antifield-dependent gravity normalization can alter \(tH_g\).
For the explicitly permitted change \(H_g\mapsto H_g+d_gF_g\), the
simultaneous primitive changes by \(tF_g\). Let \(F_{tF_g}\) be its
flat kinetic tensor. The necessary separator then becomes
\begin{equation}
 \operatorname{tf}\!\bigl(F_{\lambda,\theta}+F_{tF_g}\bigr),
 \qquad \operatorname{tf}F=F-\tfrac14\operatorname{tr}(F)I_4.
 \label{eq:v34-other-gravity}
\end{equation}
An unspecified \(F_g\) is not assumed either to cancel or to preserve
this tensor. Changing that gravity/source normalization, transforming a
new background spurion, or allowing explicit coordinate-dependent
counterterms is a different relative problem. None is used here.

\section{A compatible replacement with zero mass and mean kinetic coefficient}
\label{sec:v34-replacement}

The attainable flat kinetic tensors on the paired branch are now
completely known, even though the full homogeneous kernel is not:
\begin{equation}
 \{F_{\rm attainable}\}=F_{\lambda,\theta}+\R I_4,
 \qquad d_{\rm attainable}\in\R.
 \label{eq:v34-attainable}
\end{equation}
The inclusion follows from \cref{thm:v34-kinetic-kernel}; the converse
uses the two explicit closed functionals
\eqref{eq:v34-closed-functionals}.

Define
\begin{align}
 \bar f_{\lambda,\theta}
   &=\tfrac14\operatorname{tr}F_{\lambda,\theta}
     =\frac{\lambda^2}{\chi}
             \left(\frac58\mathfrak a_\theta-\mathfrak b_\theta\right),
             \nonumber\\
 H_{2,\min}^{[2]}
   &=H_{2,\mathrm{can}}^{[2]}
           -\bar f_{\lambda,\theta}\mathcal K_H
           -d_{\lambda,\theta}\mathcal V_H.
 \label{eq:v34-Hmin}
\end{align}
All functionals in this equation are retained through the same arity-five
bank. In particular the subtracted kinetic and mass actions include all
their metric contacts; they are not just flat two-point subtractions.

\begin{theorem}[An explicit compatible subcritical continuation]
\label{thm:v34-compatible}
The functional \eqref{eq:v34-Hmin} obeys
\begin{equation}
 d_HH_{2,\min}^{[2]}=x_2^{[2]}-tH_g^{[2]}.
 \label{eq:v34-min-equation}
\end{equation}
It preserves the chosen pure-gravity normalization and has zero flat
scalar mass coefficient and zero mean kinetic coefficient. Its remaining
flat kinetic tensor is
\begin{equation}
 \boxed{
 F_{\min}=\frac{5\lambda^2\mathfrak a_\theta}{2\chi}
                         \left(P_t-\frac14I_4\right).}
 \label{eq:v34-Fmin}
\end{equation}
This is the unique flat tensor of minimum Hilbert--Schmidt norm among
all simultaneous solutions with fixed gravity. Its norm and the minimum
distance of that flat action projection from V29's zero tensor are
\begin{equation}
 \boxed{
 \|F_{\min}\|_{\mathrm{HS}}
   =\frac{5\sqrt3\lambda^2}{4\chi}\mathfrak a_\theta.}
 \label{eq:v34-min-distance}
\end{equation}
No uniqueness of the \emph{whole} counterterm functional is asserted.
\end{theorem}
\begin{proof}
V33 supplies \(d_HH_{2,\mathrm{can}}^{[2]}=x_2^{[2]}-tH_g^{[2]}\).
The two added terms are closed by
\eqref{eq:v34-closed-functionals}, proving
\eqref{eq:v34-min-equation}. Their flat coefficients subtract precisely
\(\bar f I_4\) and \(d\), leaving \eqref{eq:v34-Fmin}.
The affine line \eqref{eq:v34-attainable} is orthogonally decomposed
into its trace-free part and a scalar multiple of \(I_4\). Its unique
least-norm point is therefore \(\operatorname{tf}F_{\lambda,\theta}\).
The squared norm of \(P_t-I_4/4\) is \(3/4\), giving
\eqref{eq:v34-min-distance}. Closed additions with zero flat projection
remain possible, so this is only a minimality statement for the
explicitly stated projection.
\end{proof}

\textbf{Recorded normalization change.} For subsequent subcritical
calculations in this lineage, \(H_{2,\min}^{[2]}\), with the inherited
zero weight-zero/one/three coefficients, is the active mixed local
normalization. The gravitational normalization stays fixed. V29's
particular \(f=0\) primitive is retained historically, not silently
relabelled as this simultaneous functional. The scalar-antifield and
higher-metric coefficients are those of the full functional
\eqref{eq:v34-Hmin}; they are not held at their V29 values.

At finite cutoff use V33's actual accumulated coefficient
\(U_{a,2}=g_{a,2}+C_{a,2}^{\rm old}\). The incremental counterterm is
\begin{equation}
 \Delta C_{a,2,<4}=-(U_{a,2})_{<4}+H_{2,\min}^{[2]}.
 \label{eq:v34-finite-update}
\end{equation}
The old insertion is included in \(U_{a,2}\), so it is not counted
twice. The exact local projection already proved in V33 gives
\begin{equation}
 \boxed{
 (B_{a,2}^{\rm new})_{<4}
     =(x_{0,\lambda,\infty,2}-x_{a,2})_{<4}.}
 \label{eq:v34-new-breaking}
\end{equation}
All action, scalar-algebra and mixed-antifield descendants are included
in this one equation. The shift from the canonical comparison is
antifield-free and closed, so it changes none of its reference equations.

The tensor \eqref{eq:v34-Fmin} is a coordinate- and reference-dependent
local normalization in a modified identity at fixed \(\lambda\). It
is not a physical Lorentz-violating coupling prediction. The reference
is itself part of that identity; a noninvariant local coefficient need
not be removable while every other source normalization is frozen.
The distinctions are familiar in cutoff-modified BRST treatments
\cite{V28IIM}; the coefficient and nonextension proof here concern the
specific native regulator. No infrared-removal assertion is made merely
because this particular tensor is proportional to \(\lambda^2\).

\section{Finite-cutoff stability of the separating test and source bounds}
\label{sec:v34-bounds}

The separator is an ordinary reference coefficient; it also has a
controlled finite approximation. Let \(F_{a,L}^{\rm sc}\) be the flat
kinetic tensor extracted from \(Y_{a,2}^{[2]}/2\), where \(Y_{a,2}\)
is obtained by logarithmically differentiating the actual V27 simultaneous
finite determinant at fixed upper regulator. Its definition keeps every
differentiated metric and scalar covariance. It is a diagnostic of the
finite coefficient, not an assertion that the finite scale derivative
already solves the exact finite master equation.

\begin{proposition}[Stability of the populated kinetic separation]
\label{prop:v34-finite-gap}
On the inherited domain
\begin{equation}
 \lambda\ge32768\mu>0,\qquad
 a\lambda\le2^{-24},\qquad \lambda L_\nu\ge1,
 \qquad 0\le m\le M\lambda,
 \label{eq:v34-domain}
\end{equation}
for every prescribed \(b>4\),
\begin{equation}
 \|F_{a,L}^{\rm sc}-F_{\lambda,\theta}\|_{\mathrm{HS}}
 \le C_b\frac{\lambda^2}{\chi}
       \left[(a\lambda)^4+(\lambda L_{\min})^{-b}\right].
 \label{eq:v34-F-rate}
\end{equation}
The same comparison holds for the mass coefficient with the corresponding
absolute norm. In particular the diagonal separator stays nonzero once
its displayed error is smaller than half the lower bound
\eqref{eq:v34-separator-lower}. No numerical value of that
profile-dependent threshold is asserted.
\end{proposition}
\begin{proof}
A scale derivative puts \(\dot\Theta_\lambda\) on a primitive
covariance. At the zero external Taylor origin the other routed momentum
is identical or differs only by a differentiated soft transfer. Every
nonzero differentiated word is confined to a fixed lower annulus, away
from the retained origin. The profile of the periodic ultraviolet
completion is identically one there. The flat metric inverse and both
continued, corrected scalar vertices differ from their ordinary values
by fourth relative mesh order, as established in V27. No retained metric
mark is present, so the third-relative-order gravitational cubic is not
an operand of this calculation. Differentiating the cutoff factors once
does not change that order. An ordered telescope of the finite number
of inverse and vertex factors is bounded by \(C(a\lambda)^4\)
times the ordinary weight-two scale \(\lambda^2/\chi\).

After rescaling the annulus, all fixed momentum derivatives are bounded.
Poisson summation with more than \(b+4\) such derivatives gives the
separate volume error. The support is separated from zero, so no
omitted-zero-mode term occurs. This proves \eqref{eq:v34-F-rate}.
The kinetic difference is a bounded coefficient functional and therefore
inherits the estimate; \eqref{eq:v34-separator-lower} gives the final
claim. The mass statement uses its actual derivative at \(z=0\),
regular on the same annulus, with one cutoff factor per scalar covariance.
\end{proof}

For a smooth external window of scale \(\mu\), the flat replacement
and each of its prescribed kernel moments obey
\begin{equation}
 \|\mathcal K_{\min}^{\rm flat}\|_{\rm rooted,s}
               \le C_s\frac{\lambda^2}{\chi}\mu^2.
 \label{eq:v34-flat-kernel}
\end{equation}
This follows by Fourier integration by parts for the degree-two symbol
\eqref{eq:v34-Fmin}. Its finite-comparison difference carries the
additional bracket in \eqref{eq:v34-F-rate}. More generally the
finite arity-five coefficients of \eqref{eq:v34-Hmin} inherit V33's
annular bounds; the two closed additions have bounded analytic coframe
jets on the same fixed chart. They retain the existing derivative floor
and all mass jets. Rooted bounds have no total-volume factor, with one
source in \(L^1\) and the others in \(L^\infty\).

The restored local residual itself uses the previously proved V33 bound:
\begin{equation}
 |(B_{a,2}^{\rm new})_I^{[D]}|
 \le C_{I,b}\chi^{-(1+r+s)}\lambda^{5-n_*-D}
       \left[(a\lambda)^3+(\lambda L_{\min})^{-b}\right],
 \label{eq:v34-restored-rate}
\end{equation}
for the stated subcritical slots. This is inherited through
\eqref{eq:v34-new-breaking}, not presented as a new multiscale proof.
Counterterms may still contain divergent local coefficients; the
bounded remainder and bounded coefficient maps do not prove an invariant
analytic ball of microscopic actions.

\section{Consequences for the critical relative problem}
\label{sec:v34-next}

The former subcritical ambiguity is now resolved on the paired branch:
V29's zero scalar kinetic prescription fails a populated left-null test,
while \eqref{eq:v34-Hmin} supplies a simultaneous compatible replacement.
It is therefore no longer useful to seek a higher-background extension
which leaves that zero tensor fixed. The failure is a relative
normalization obstruction, not a nontrivial anomaly of the scalar--gravity
theory: the simultaneous solution explicitly exists at the same weight.

The critical question remains
\begin{equation}
 d_HH_2^{[4]}=x_2^{[4]}-tH_g^{[4]}.
 \label{eq:v34-critical}
\end{equation}
The ordinary differential preserves the weight, so the recorded
weight-two change does not alter this right-hand side. The strict
preservation of arbitrary V29 weight-four row coefficients has not
been proved possible; it must not be assumed merely because the
subcritical obstruction is now understood. Any such conditions remain
an explicit relative restriction rather than part of an unnamed choice
of primitive.

A useful next computation is the critical relative source coefficient
on a compressed degree-closed bank, including the mixed metric/scalar
antifields. Alternatively a scalar-inclusive local existence theorem
would have to be applied to the \emph{actual extendable relative germ}
with the fixed gravity normalization and with explicit control of the
finite primitive map. V25's pure-gravity classification by itself is
not that theorem. The present proof neither redoes that classification
nor claims that the critical scalar equation has been solved.

The paired-reference choice in \eqref{eq:v34-gravity-branch} is explicit.
For a different supplied \(d_gF_g\), the extra flat coefficient in
\eqref{eq:v34-other-gravity} must be evaluated first. This provides an
upstream diagnostic rather than an invitation to change the gravity
source merely to make a failed scalar condition disappear.

The new results concern local one-loop normalization. The nonlocal
matter-count-two estimates of V27--V32 are inherited and are not
reproved. Matched source functionals must use the new local data
consistently. No arbitrary-loop theorem, chiral determinant-line
construction, full internal-gauge restoration, or analytic interacting
ESM trajectory follows from this finite step.

\section{Verification and closing ledgers for V34}
\label{sec:v34-ledgers}

The self-contained verifier
\srcid{check_ESM_kinetic_compatibility_V34.py} independently checks the
ten-coordinate contact contractions, the signed mixed numerator,
the scale-moment algebra, the Killing-jet constraints and their exact
ranks, the trace-free separator and least-norm projection, the constant
mass coefficient, and the append-only preservation test. Its certificate
keeps these finite rational identities distinct from the integration
by parts, local-jet arguments and analytic reference estimates proved
above. The separating coefficient is evaluated analytically in terms
of a strictly positive profile moment; no decimal value of the complete
native counterterm array or full relative matrix is claimed.

\subsection*{Proved}
\addcontentsline{toc}{subsection}{V34: Proved}
The flat weight-two coefficient of the actual simultaneous scale
primitive is \eqref{eq:v34-flat-value}. Every homogeneous local
matter-count-two cocycle with fixed gravity has the isotropic flat
projection \eqref{eq:v34-kernel-projection}. Their combination gives
the populated nonextension certificate \eqref{eq:v34-native-separator}
for V29's particular \(f=0\), without a large matrix computation.
The compatible replacement \eqref{eq:v34-Hmin} has zero mass and mean
kinetic coefficient and the unique minimum-norm flat tensor
\eqref{eq:v34-Fmin}. Its finite scale diagnostic is stable under the
displayed cutoff/volume comparison. The change of active subcritical
normalization is explicitly recorded, with all source and metric
contacts carried by one functional.

\subsection*{Inherited}
\addcontentsline{toc}{subsection}{V34: Inherited}
V1--V33 and all labels are unchanged. V13/V27 provide the actual contact,
signed mixed bubble and continued scalar reference; V24--V33 provide
the ordinary degree-closed complex, common projection, simultaneous
scale primitive and reference estimates. V29 still supplies its valid
first-row right inverse, not a simultaneously extendible scheme.
The old auxiliary density and geometric source conventions remain
unchanged. None of the previous analytic proofs is certified merely
by rerunning a finite symbolic identity.

\subsection*{Literature input}
\addcontentsline{toc}{subsection}{V34: Literature input}
The relation between local consistent deformations and antifield
terms is established BRST methodology, for example
\cite{V29BBHMatter}. Cutoff-dependent wave-function normalizations
must be organized with their source identities, as in
\cite{V28IIM}. No general scalar-inclusive cohomology theorem is
imported in this part. The particular separator is proved from local
Killing tests and the actual native loop coefficient.

\subsection*{Conditional}
\addcontentsline{toc}{subsection}{V34: Conditional}
The no-extension certificate assumes the declared paired gravity branch;
an extra exact gravity source normalization changes it by the explicit
\(tF_g\) term. The simultaneous continuation imports V33's already
proved subcritical primitive and reference theorem at their original
scope. Matching to a different microscopic acceptance law, to an
uncompleted ultraviolet regulator, or to extra physical normalization
conditions is not supplied by this comparison. The exact finite
threshold for stability depends on the retained profile seminorms.

\subsection*{Open}
\addcontentsline{toc}{subsection}{V34: Open}
Solve the critical relative equation \eqref{eq:v34-critical}, with
its own explicitly declared lower-arity conditions, on the same
matter/gravity carrier. The particular zero subcritical kinetic target
of V29 is not left as a repeatedly untested extension assumption: it
fails on the paired branch. The new subcritical choice is
\eqref{eq:v34-Hmin}. Coupled internal gauge/chiral sources, higher
loops, filtered R1--R4 realization, composite source transport beyond
the proved panels and the full ESM continuum remain open.

\begin{center}\small
\begin{tabular}{>{\raggedright\arraybackslash}p{.24\textwidth}>{\raggedright\arraybackslash}p{.67\textwidth}}
\toprule
Variable & Scope in V34\\\midrule
Order and fields & One loop; matter count two; arity at most five.
New normalization result is weight two, hence subcritical.\\
Couplings and mass & Four equal real scalars; zero internal gauge,
Yukawa and self-couplings; fixed \(\chi>0\). Mass-polynomial jets in
\(z=m^2\); no inverse mass.\\
Reference & Same coframe coordinates, contour and paired gravity
normalization; nonnegative smooth lower profile and fixed \(\lambda>0\).
No new transforming background spurion.\\
Cutoff and volume & Separator is an exact thermodynamic coefficient;
finite comparison is \eqref{eq:v34-F-rate} on
\eqref{eq:v34-domain}. No exact finite-torus scaling is assumed.\\
Uniformity & Fixed profile seminorms, source order and derivative/kernel
moments. Rooted estimates have no total-volume factor; constants need
not be uniform in unbounded degree or \(\chi\downarrow0\).\\
Normalization & V29 body retained; active subcritical prescription changed
explicitly to \eqref{eq:v34-Hmin}. Critical coefficients not silently
replaced. No full-functional uniqueness is asserted.\\
Iteration & Only the actual lower-reference coefficient is integrated
in the new calculation. No invariant analytic RG class, all-loop
summability or infrared removal is proved.\\
\bottomrule
\end{tabular}
\end{center}



\clearpage
\part*{V35: A mass-polynomial critical complex and an explicit highest-mass solver}
\addcontentsline{toc}{part}{V35: Critical mass recursion and highest-mass coefficient solver}

\section{The critical equation must be solved as a mass-polynomial equation}
\label{sec:v35-start}

Sections 1--271 and all their labels are preserved. The active subcritical
normalization remains exactly \eqref{eq:v34-Hmin}, including its nonzero
trace-free scalar kinetic tensor. The fixed gravitational critical
normalization is not changed. This part attacks
\eqref{eq:v34-critical} by separating its three possible powers of
\(z=m^2\) \emph{inside the differential}, not by solving independently at
several numerical masses.

The main positive results are an exact triangular coefficient equation,
a signed ordered formula for each of its native reference inputs, and a
complete finite solver for the last, highest-mass equation. The latter is
an actual 286-dimensional local potential problem: its action map has
rank 285, and its primitive is given by radial integration in the ten
\emph{field coordinates}, not by an inverse spacetime derivative. We also
exhibit an eleven-dimensional family of critical exact counterterms. It
shows concretely why V34's obstruction to a two-derivative anisotropic
kinetic change must not be copied to the four-derivative critical bank.

These results do not evaluate the two earlier native range problems. In
particular, a primitive at \(z=0\), chosen with an arbitrary homogeneous
part, need not lift to a polynomial primitive at nonzero \(z\). The
compatibility terms governing that lift are retained below. There is no
claim in V35 that the complete critical scalar--gravity equation is solved.

The fields, ordinary master action, completed finite regulator, metric
contour and reference measure are those of V27--V34. Internal gauge,
Yukawa and scalar self-couplings remain zero. The scalar carrier consists
of four equal real components; we keep its flavour-singlet quadratic
sector. Work is at one loop, matter count two, and total arity at most
five. The common factor \(\hbar\) is suppressed in one-loop coefficient
arrays. It is restored when a counterterm is inserted. All local
functionals below are modulo total ordinary derivatives.

\section{The native mass term supplies a square-zero second operator}
\label{sec:v35-bicomplex}

Separate the mass term in the \emph{same} scalar master action
\eqref{eq:v33-SH}:
\begin{equation}
 \mathbb S_H(z)=\mathbb S_H(0)+z\mathcal V,
 \qquad
 \mathcal V=\frac12\sum_A\int\rho(q)\phi_A^2,
 \qquad \rho(q)=\det e(q).
 \label{eq:v35-mass-action}
\end{equation}
On the matter-count-two complex set
\begin{equation}
 D=d_H(0),\qquad M=(\mathcal V,\,\cdot\,)_H,
 \qquad t(z)=t_0+zt_1,\quad t_1=(\mathcal V,\,\cdot\,)_g.
 \label{eq:v35-DM}
\end{equation}
The subscripts denote the inherited scalar and gravity/ghost parts of the
antibracket. In the Euler sign convention of V31,
\(M\upsilon_A=\rho\phi_A\), with \(M\) zero on the other minimal
coordinate generators. It commutes with ordinary differentiation. Thus
\(M\) replaces a scalar antifield by the actual density-weighted scalar;
it does not discard a scalar equation of motion. The operator \(t_1\)
replaces a metric antifield by the derivative of \(\mathcal V\); it is the
mass part of the scalar stress.

Use the residual weight that excludes the explicit power of \(z\).
Then \(D\) preserves this weight and \(M\) lowers it by two. Both raise
ghost number by one, preserve matter count two, and cannot lower total
arity. Let \(\mathscr C^g_{v}\) denote the corresponding mass-free local
bank of ghost number \(g\), residual weight \(v\), and arity at most five.

\begin{proposition}[Native mass bicomplex]
\label{prop:v35-bicomplex}
On the declared quotient,
\begin{equation}
 \boxed{d_H(z)=D+zM,\qquad D^2=M^2=0,\qquad DM+MD=0.}
 \label{eq:v35-bicomplex}
\end{equation}
No inverse mass is needed to define any of these operators. Their
coefficient matrices are fixed rational matrices after the inherited
normalizations and the finite expansion of \(e(q)\).
\end{proposition}
\begin{proof}
The first identity follows directly from \eqref{eq:v33-differentials}
and \eqref{eq:v35-mass-action}; the ordinary transformations do not depend
on the mass. The mass potential has zero scalar antibracket with itself,
which proves \(M^2=0\). Alternatively compare coefficients of \(z^2\)
in the inherited identity \(d_H(z)^2=0\). Its constant and linear
coefficients give \(D^2=0\) and \(DM+MD=0\). The identity is a polynomial
identity before choosing a numerical mass. The arity quotient is legitimate
because all these operations are nondecreasing in arity. The replacement
\(\upsilon\mapsto\rho\phi\) is local and has no derivatives, and
\(\rho\) is a finite coframe polynomial. The other coefficient operations
are those already used to construct \(d_H\). This proves the stated
matrix and locality properties.
\end{proof}

Let \(H_g^{[4]}\) denote the fixed gravitational \emph{source}
normalization. Its source part is independent of \(z\). Any separately
retained antifield-free scalar-background determinant is annihilated by
\(t\), as already specified in V33; it does not affect the following
relative coefficient. Define
\begin{equation}
 \Xi(z)=x_2^{[4]}(z)-t(z)H_g^{[4]}
       =\xi_0+z\xi_1+z^2\xi_2.
 \label{eq:v35-Xi}
\end{equation}
Here \(\xi_j\in\mathscr C^1_{4-2j}\). The expansion stops at \(z^2\):
a ghost-one critical component with \(n_*\) antifields has
\(D_K+2j=5-n_*\), so \(j\ge3\) is impossible. The source-dependent
part of \(tH_g\) is at most affine in \(z\).

Relative consistency from V33 now has four, not three, coefficient rows:
\begin{equation}
 \boxed{D\xi_0=0,\quad
 D\xi_1=-M\xi_0,\quad
 D\xi_2=-M\xi_1,\quad M\xi_2=0.}
 \label{eq:v35-input-consistency}
\end{equation}
In particular \(\xi_1\) and \(\xi_2\) need not separately be
\(D\)-cocycles. Their residual weights two and zero do not turn them
into the already solved subcritical reference coefficients of V33.

\section{The three inputs are ordered derivatives of the same reference trace}
\label{sec:v35-native-input}

This section specifies the \(\xi_j\) without substituting a massless
propagator into an unsubtracted massive loop. Work first on the ordinary
thermodynamic reference at the retained scale \(\lambda>0\), and form
all primitive scalar, metric and ghost blocks of the common graded Hessian.
Let
\[
 C(z)=\nu_\lambda\mathbb H_0(z)^{-1},\quad
 V(z)=V_0+zV_1,\quad W(z)=C(z)V(z),\quad E(z)=(I+W(z))^{-1}.
\]
The scalar action is quadratic and affine in \(z\); hence its full
background-Hessian perturbation, including its metric derivatives, is
also affine in \(z\). The ordinary transformation Jacobian \(Dr\) is
mass independent. All inverses here are on hard support; the finite
external-arity expansion of \(E\) is used after the indicated derivatives.

At \(z=0\) the only nonzero entries of \(C_1=\partial_zC|_0\) and
\(C_2=\partial_z^2C|_0\) are in the scalar block:
\begin{equation}
 (C_1)_H=-\frac{\nu_\lambda(p)}{|p|^4}I_4,
 \qquad (C_2)_H=\frac{2\nu_\lambda(p)}{|p|^6}I_4.
 \label{eq:v35-Cjets}
\end{equation}
There is one \(\nu_\lambda\) in each formula. The mass derivative also
hits the vertices, not just the scalar inverse. Put
\[
 W_1=C_1V_0+C_0V_1,\qquad
 W_2=C_2V_0+2C_1V_1,\qquad E_0=(I+C_0V_0)^{-1}.
\]
Denote by \(\mathscr E_j\) the actual extraction of matter count two,
arity at most five, and mass-free residual weight \(4-2j\), after
rooting and retaining all Grassmann signs.

\begin{proposition}[Native ordered mass inputs]
\label{prop:v35-mass-inputs}
With coefficient rather than derivative normalization for \(z^2\),
\begin{align}
 x^{(0)}&=\mathscr E_0\operatorname{STr}(E_0\Theta_\lambda Dr),
 \nonumber\\
 x^{(1)}&=-\mathscr E_1\operatorname{STr}
                   (E_0W_1E_0\Theta_\lambda Dr),
 \nonumber\\
 x^{(2)}&=\mathscr E_2\operatorname{STr}
 \left[
 \left(E_0W_1E_0W_1E_0-\tfrac12E_0W_2E_0\right)
                                      \Theta_\lambda Dr\right],
 \label{eq:v35-reference-jets}\\
 \xi_0&=x^{(0)}-t_0H_g^{[4]},\qquad
 \xi_1=x^{(1)}-t_1H_g^{[4]},\qquad \xi_2=x^{(2)}.
 \label{eq:v35-relative-jets}
\end{align}
These are finite ordered coefficient integrals of the existing measured
reference. No matrix factor is commuted through a source vertex.
\end{proposition}
\begin{proof}
Differentiate \((I+W)E=I\), obtaining
\(E'=-EW'E\) and
\(E''=2EW'EW'E-EW''E\). Substitute the displayed derivatives of
\(C,V\) into the inherited trace
\(\operatorname{STr}(E\Theta_\lambda Dr)\), and divide the second
derivative by two. This proves the formulas, including the mixed
\(C_1V_1\) term and its multiplicity. The factor
\(\Theta_\lambda\) is not differentiated with respect to \(z\).

Every contributing matter-count-two word contains a covariance and the
final lower-reference factor. On a small external panel, their routed
support is in a fixed lower annulus, for example
\(\lambda/2\le|p|\le3\lambda\). Thus all displayed mass derivatives
are dominated there. At the local Taylor origin this is compact support
away from zero. The finite external degree requires only the existing
input jets through degree seven. Expanding each \(E_0\) in external
arity produces exactly the signed finite compiler, so the operation is
not a new determinant model. The relative terms follow from
\eqref{eq:v35-DM}.\end{proof}

In the original, unscaled variables, a component of these critical jets
with \(n_*\) antifields and \(D_K=5-n_*-2j\) momentum derivatives has
lower-scale factor
\(\lambda^{5-n_*-D_K-2j}=1\), apart from its already stated powers of
\(\chi\). This is why differentiating with respect to the lower scale
cannot by itself construct the critical primitive. Formula
\eqref{eq:v35-reference-jets} is a prescription for its actual mass
coefficients; their native annular integrals are not numerically evaluated
in this part.

\section{An exact three-step reconstruction, with the lift defects retained}
\label{sec:v35-recursion}

A ghost-zero critical matter-count-two counterterm has precisely the form
\begin{equation}
 H_2^{[4]}(z)=h_0+zh_1+z^2h_2,
 \qquad h_j\in\mathscr C^0_{4-2j}.
 \label{eq:v35-Hansatz}
\end{equation}
Its last coefficient has no antifields and no spacetime derivatives:
\begin{equation}
 h_2=\frac12\int f(q)\sum_A\phi_A^2,
 \qquad f(q)\text{ a polynomial jet through degree three}.
 \label{eq:v35-h2}
\end{equation}
The degree-three cutoff follows from the two scalar arguments and arity
five. In particular \(Mh_2=0\).

\begin{theorem}[Critical mass reconstruction and residual identities]
\label{thm:v35-recursion}
The critical relative equation is equivalent to
\begin{equation}
 \boxed{
 Dh_0=\xi_0,\qquad
 Dh_1=\xi_1-Mh_0,\qquad
 Dh_2=\xi_2-Mh_1.}
 \label{eq:v35-three-steps}
\end{equation}
For arbitrary candidates define
\(e_0=\xi_0-Dh_0\),
\(e_1=\xi_1-Mh_0-Dh_1\), and
\(e_2=\xi_2-Mh_1-Dh_2\). Then
\begin{equation}
 \boxed{
 \Xi-(D+zM)H=e_0+ze_1+z^2e_2,\quad
 De_0=0,\quad De_1=-Me_0,\quad De_2=-Me_1,\quad Me_2=0.}
 \label{eq:v35-residuals}
\end{equation}
Every operation is polynomial at \(z=0\). There is no missing coefficient
of \(z^3\) in the counterterm variation.
\end{theorem}
\begin{proof}
Expand \((D+zM)(h_0+zh_1+z^2h_2)\); its last term is
\(z^3Mh_2=0\). Comparison gives \eqref{eq:v35-three-steps} and the
first identity of \eqref{eq:v35-residuals}. Apply \(D\), use
\eqref{eq:v35-input-consistency}, \(DM=-MD\), and \(M^2=0\).
For the final relation,
\(Me_2=M\xi_2-MDh_2=DMh_2=0\). This establishes each coefficient
identity even when the proposed counterterms do not solve the equation.
\end{proof}

There is a genuine choice obstruction between these steps. If \(h_0\)
is changed by a \(D\)-cocycle \(k_0\), the second target changes by
\(-Mk_0\). It is liftable without changing \(h_0\) again only if
that change lies in the required image of \(D\), with the imposed finite
normalizations. Likewise a homogeneous change in \(h_1\) changes the
last target. The identities above prove consistency of the targets after
earlier steps have been solved; they do not prove exactness.

For a finite algebraic example, take \(Da=0\), \(Ma=b\ne0\),
\(Db=Mb=0\), with \(a\) of ghost number zero and residual weight four,
and \(b\) of ghost number one and residual weight two. Let the weight-two
ghost-zero space be zero. The target \(\Xi=0\) has the solution zero,
but the chosen massless primitive \(h_0=a\) cannot be extended. This
example is not a native anomaly: it proves that an arbitrary homogeneous
choice at the first step cannot be accepted on consistency alone.

In matrix language the actual finite problem has the block form
\begin{equation}
 \begin{pmatrix}
 D_4&0&0\\ M_4&D_2&0\\0&M_2&D_0
 \end{pmatrix}
 \begin{pmatrix}h_0\\h_1\\h_2\end{pmatrix}
 =\begin{pmatrix}\xi_0\\\xi_1\\\xi_2\end{pmatrix}.
 \label{eq:v35-block-system}
\end{equation}
The subscripts here are residual weights, not cohomological degrees.
This is a single mass-regular coefficient equation. Solving unrelated
systems at several positive masses would not certify it.

\section{The highest-mass equation is an explicitly solvable potential map}
\label{sec:v35-potential}

The final map \(D_0\) in \eqref{eq:v35-block-system} can be completely
specified without a large BV rank calculation. This section works first
in the metric-difference coordinate \(h=G(q)-I\). It is only a local
coefficient conversion: the original variable \(q\), integration measure
and Jacobian are not changed. The maps \(h=q+\mathsf Q(q,q)\) and its
inverse have fixed rational jets, and are tangent to the identity.

Write \(g=I+h\), \(\rho(h)=\sqrt{\det(I+h)}\), and let \(A^\mu{}_\nu\)
be the independent first ghost jet \(\partial_\nu c^\mu\). Denote
\(\mathcal T_h(A)=A^Tg+gA\). Once derivatives are moved off the scalar
square, the variation of the potential is
\begin{equation}
 D\left(\frac12\int f(h)\sum_A\phi_A^2\right)
 =\frac12\int\sum_A\phi_A^2\,\mathcal Bf(h;\partial c),
 \qquad
 \mathcal Bf=\jet_h^2\{Df[\mathcal T_h(A)]-f\operatorname{tr}A\}.
 \label{eq:v35-potential-map}
\end{equation}
Indeed the transport \(c\cdot\partial h\) cancels the derivative of
\(f\) produced by integrating the scalar Lie variation by parts. The
remaining density term is \(-f\operatorname{div}c\). This cancellation
is exact before taking field degree.

Use the ten independent symmetric coordinates \(h_i\), \(1\le i\le10\).
For an arbitrary row \(B(h;A)\), linear in the 16 entries of \(A\) and
polynomial through degree two in \(h\), define
\begin{align}
 J_h(V)&=\tfrac12g^{-1}V,
 &\alpha_h(V)&=\jet_h^2\{\rho^{-1}B(h;J_h(V))\},\nonumber\\
 B^{\mathrm{hor}}(h;A)
 &=\jet_h^2\{\rho\alpha_h(\mathcal T_h(A))\},
 &B^{\mathrm{vert}}&=B-B^{\mathrm{hor}}.
 \label{eq:v35-horizontal}
\end{align}
The only inverse is the finite-dimensional background matrix \(g^{-1}\),
expanded at \(I\); it has nothing to do with an external momentum.
The notation ``horizontal'' refers solely to the kernel of
\(A\mapsto A^Tg+gA\) in this finite coefficient calculation.

Define two radial operators on polynomial field-space forms:
\begin{align}
 (K_1\alpha)(h)&=\int_0^1\alpha_{th}(h)\,\dd t,
 \nonumber\\
 (K_2\beta)_i(h)&=\int_0^1 t\,h^j\beta_{ji}(th)\,\dd t.
 \label{eq:v35-radial}
\end{align}
All integrals in \eqref{eq:v35-radial} are finite coefficient operations:
the denominators are integers from one to three.

\begin{theorem}[Potential range, primitive, and certified residual]
\label{thm:v35-potential}
The row \(B\) is in the image of \eqref{eq:v35-potential-map} if and
only if
\begin{equation}
 \boxed{B^{\mathrm{vert}}=0,\qquad d_h\alpha=0.}
 \label{eq:v35-potential-tests}
\end{equation}
On this image every potential is
\begin{equation}
 \boxed{f=\jet_h^3\{\rho(K_1\alpha+c)\},\qquad c\in\R.}
 \label{eq:v35-potential-primitive}
\end{equation}
For arbitrary input and the particular choice \(c=0\), the exact
unremoved row is
\begin{equation}
 \boxed{
 B-\mathcal Bf
 =B^{\mathrm{vert}}
  +\jet_h^2\{\rho(K_2d_h\alpha)(\mathcal T_h(A))\}.}
 \label{eq:v35-potential-residual}
\end{equation}
At arity five the potential domain has dimension 286 and the matrix of
\(\mathcal B\), into the 1056-component row space just specified, has
rank 285. Its one-dimensional kernel is \(\jet_h^3\rho\).
\end{theorem}
\begin{proof}
First \(\mathcal T_hJ_h(V)=V\) for every symmetric \(V\). Hence
\(B\mapsto B^{\mathrm{hor}}\) is a projection, and the map
\(\alpha\mapsto\jet_h^2\rho\alpha(\mathcal T_h(A))\) is injective.
The determinant identity
\(D\rho[\mathcal T_h(A)]=\rho\operatorname{tr}A\) gives
\(\mathcal Bf=\jet_h^2\rho\,d(f/\rho)(\mathcal T_h(A))\).
Thus image rows have zero vertical part and exact \(\alpha\).

For a polynomial one-form of degree at most two, differentiation under
\eqref{eq:v35-radial} gives
\[
 \alpha=d_hK_1\alpha+K_2d_h\alpha.
\]
For example its \(i\)-component follows by integrating the derivative of
\(t\alpha_i(th)\), then subtracting
\(\partial_i\int_0^1h^j\alpha_j(th)\dd t\). This proves
sufficiency, \eqref{eq:v35-potential-residual}, and the displayed
primitive. The same calculation shows that two potentials differ by
\(\rho\) times a constant to the retained degree.

There are \(\binom{13}{3}=286\) monomials of degree at most three in ten
variables. Differentiation annihilates only the constant polynomial, so
its gradient has rank 285. Multiplication by the invertible jet of
\(\rho\) and the horizontal embedding do not change this rank. The
row space has dimension \(16\binom{12}{2}=1056\). As a separate useful
check, polynomial one-forms have dimension \(10\binom{12}{2}=660\),
while their curl has rank \(660-285=375\). These ranks follow from the
explicit homotopy, without numerical singular values or an unevaluated
cohomology assertion.
\end{proof}

An arbitrary highest-mass target \(\xi_2-Mh_1\) may also have
antifield components or, after rooting on the scalar square,
\(c^\mu\partial_\nu h_i\) terms. Potentials have neither. Such components
must first be zero for \eqref{eq:v35-potential-tests} to be the complete
test. The theorem does not infer those zeros from the isolated reference
coefficient. It applies to the \emph{updated} target after \(h_1\) has
been chosen in the preceding equation.

\begin{proposition}[Explicit coefficient bounds]
\label{prop:v35-potential-norm}
In independent symmetric \(h\) coordinates, let \(\|\cdot\|_1\) sum
absolute monomial coefficients and, for forms, absolute coefficients in
independent ordered components. Then
\begin{equation}
 \|K_1\alpha\|_1\le\|\alpha\|_1,\qquad
 \|K_2\beta\|_1\le\|\beta\|_1,
 \qquad
 \|f\|_1\le\tfrac{39}{2}\|\alpha\|_1\quad(c=0),
 \label{eq:v35-adapted-bound}
\end{equation}
and
\begin{equation}
 \boxed{
 \|f\|_1\le2048\|B\|_1,\qquad
 \|B-\mathcal Bf\|_1
 \le\|B^{\mathrm{vert}}\|_1+80\|d_h\alpha\|_1.}
 \label{eq:v35-raw-bound}
\end{equation}
These bounds have no cutoff, volume, mass or momentum denominator.
\end{proposition}
\begin{proof}
A coefficient of degree \(r\) in \(K_1\) is divided by \(r+1\).
An independent two-form component contributes twice to \(K_2\), with
denominator at least two. This proves the first two bounds. The actual
metric density jet is
\[
 \rho_{\le3}=1+\tfrac12\operatorname{tr}h
 +\tfrac18(\operatorname{tr}h)^2-\tfrac14\operatorname{tr}(h^2)
 +\tfrac1{48}(\operatorname{tr}h)^3
 -\tfrac18\operatorname{tr}h\operatorname{tr}(h^2)
 +\tfrac16\operatorname{tr}(h^3).
\]
Its coefficient norm is \(39/2\); the degree-two density and inverse
density have norms 8 and 9. A row of
\((I+h)^{-1}_{\le2}=I-h+h^2\) has summed coefficient norm at most 21.
The section \(J_h\) has the extra factor \(1/2\), so
\(\|\alpha\|_1\le(189/2)\|B\|_1\). The product
\((39/2)(189/2)=7371/4\) is less than 2048. Finally each component of
\(\mathcal T_h(A)\), summed in its \(A\) entries and \(h\) monomials,
has norm at most 10. Its composition with \(\rho_{\le2}\) is bounded
by 80. Apply \eqref{eq:v35-potential-residual}. All estimates are for
fixed finite polynomial maps. The actual coframe-coordinate conversion
is a further fixed rational jet map and does not change uniformity in
cutoff or volume; the displayed numerical constants refer to metric
coordinates, not to an unspecified coordinate-invariant norm.
\end{proof}

\section{Critical exact freedoms differ from the subcritical kinetic freedom}
\label{sec:v35-exact-family}

The preceding recursion must allow compatible homogeneous changes, but
V34's fixed weight-two tensor cannot be reset. Here is a constructive
family at \emph{weight four} that illustrates the distinction.
Let \(A=A^T\) be a constant real \(4\times4\) matrix and \(b\in\R\).
Define the self-adjoint scalar operator
\[
 V=-A^{\mu\nu}\partial_\mu\partial_\nu+bz,
 \qquad v(p,z)=p^TAp+bz,
\]
and the ghost-number-minus-one functional
\begin{equation}
 K_V=\frac12\sum_A\int\upsilon_A V\phi_A,
 \qquad \Delta_V=d_H(z)K_V.
 \label{eq:v35-exact-family}
\end{equation}
The repeated flavour label in the integral is independent of the matrix
label \(A\). The functional is truncated in total arity only after its
ordinary variation has been taken.

\begin{proposition}[An explicit critical homogeneous family]
\label{prop:v35-exact-family}
Every \(\Delta_V\) is an ordinary matter-count-two, ghost-zero,
weight-four cocycle with zero pure-gravity restriction. Its flat action
coefficient is
\begin{equation}
 \boxed{
 [\Delta_V]_{q=\eta=\sigma=\upsilon=0}
 =\frac12\sum_A\int_p\phi_A(-p)
             (p^2+z)(p^TAp+bz)\phi_A(p).}
 \label{eq:v35-flat-exact}
\end{equation}
The resulting flat family has dimension eleven. Adding the covariant
closed mass functional
\(\frac12\int\rho z^2\sum_A\phi_A^2\) gives a twelve-dimensional
constructively realized family. No claim that this is the entire
homogeneous kernel is made.
\end{proposition}
\begin{proof}
Nilpotence gives \(d_H\Delta_V=0\). The generator \(K_V\) has matter
count two and residual weight four, and vanishes on the pure-gravity
carrier. The Euler part is
\(\frac12\int(P_g\phi_A)V\phi_A\), where
\(P_g=-\partial_\mu W^{\mu\nu}\partial_\nu+z\rho\) is the same
ordinary scalar Hessian. At flat metric it commutes with the constant
operator \(V\), giving \eqref{eq:v35-flat-exact}. The accompanying
source term is
\(\frac12\int\upsilon_A[\mathcal L_c,V]\phi_A\) in this convention.
In particular it contains
\[
 A^{\mu\nu}
 \left[(\partial_\mu\partial_\nu c^\rho)\partial_\rho
       +2(\partial_\mu c^\rho)\partial_\nu\partial_\rho\right],
\]
inside the commutator. These terms are kept; the action polynomial alone
is not the cocycle.

Multiplication by the nonzero polynomial \(p^2+z\) is injective. The
space of \(p^TAp+bz\) has dimension \(10+1\), proving the first
count. The polynomial \(z^2\) is not in its image, since reduction
modulo \(p^2+z\) gives \((p^2)^2\ne0\). The covariant mass functional
is invariant under the ordinary simultaneous metric/scalar transformation,
so it supplies the additional independent flat direction.
\end{proof}

For instance, take \(A=\operatorname{diag}(1,-1,0,0)\) and \(b=0\).
The anisotropic \emph{critical} polynomial
\((p^2+z)(p_0^2-p_1^2)\) has the complete exact source completion
\eqref{eq:v35-exact-family}. This does not contradict V34: that result
concerns a two-derivative weight-two kinetic difference, not this
four-derivative polynomial with its Euler and source partners. Since
\(d_H\) preserves total weight, none of these additions changes
\eqref{eq:v34-Hmin}. They are possible normalization freedoms, not a
counterterm already fitted to \(\Xi\).

\section{What the construction establishes for the native critical programme}
\label{sec:v35-status-interface}

The native right-hand side now has the signed finite prescription
\eqref{eq:v35-reference-jets}, and its critical primitive must obey the
single block system \eqref{eq:v35-block-system}. The last diagonal block
is no longer an unspecified local-cohomology problem: the scalar potential
compiler \eqref{eq:v35-potential-primitive} evaluates it, with the exact
range and residual tests above. The first and second blocks still require
their \emph{populated} relative range calculations, including the choice
of homogeneous terms that can pass into the next block.

The scalar-inclusive classification of local BRST cohomology in
\cite{V29BBHMatter} explicitly allows matter transforming in tensor
representations, and distinguishes the additional possibilities associated
with abelian gauge factors. It is a relevant external comparison. It is
not used here to assert exactness of an arbitrary arity-truncated relative
array, to impose an unverified translation-invariant lift, or to replace
our fixed gravitational normalization. In particular V25's absolute
pure-gravity primitive does not, by itself, decide either of the first
two relative mass equations. Those are distinct maps and constraints.

If native coefficient errors \(\delta\xi_j\) or intermediate residuals
are supplied, \eqref{eq:v35-residuals} and
\eqref{eq:v35-raw-bound} propagate them by fixed finite operators. This
is a bound on the compiler, not a new estimate for an unevaluated native
loop integral. Inserting a potential found by that compiler adds
\(\hbar z^2h_2\); it does not change the tree propagators or generate a
new auxiliary determinant. The same action functional supplies every
metric contact of that potential.

For a fixed smooth external window, a degree-at-most-four local
counterterm has rooted kernel bounds obtained from its finite coefficient
norm and derivatives of the window. For the highest-mass term there are
no spacetime derivatives in \(h_2\), so the bound is
\(C_{N,s}m^4\|f\|_1\), with the corresponding scalar/polarization norms.
There is no total-volume factor; there is also no claim that \(\|f\|_1\)
is small before its native right-hand side has been controlled. The other
critical terms have the appropriate fixed powers of the external window
scale. All statements are mass regular and remain meaningful at \(m=0\).

The unaltered analytic domain for inherited comparisons may be taken as
\eqref{eq:v34-domain}. The finite mass and potential maps themselves do
not require a mesh restriction. They hold before taking a continuum
limit and supply no all-loop estimate, no infrared removal, and no
analytic invariant RG neighbourhood. Internal gauge/chiral and composite
source interfaces beyond the declared panels remain outside this part.

\section{Verification and closing ledgers for V35}
\label{sec:v35-ledgers}

The verifier \srcid{check_ESM_critical_mass_V35.py} uses exact rational
arithmetic. It checks the radial field-space identities on every basis
one-form through degree two, reconstructs every nonconstant scalar
potential monomial through degree three, verifies the density norms,
the mass-resolvent coefficients with noncommuting test matrices, and the
critical exact-family ranks and counterterm commutator. It also tests a
nonclosed potential input and the finite mass-lift example. Its optional
coefficient compiler accepts a ten-component polynomial one-form and
returns the potential primitive and its curl residual. It does not claim
to derive that one-form from an unevaluated native reference array.

\subsection*{Proved}
\addcontentsline{toc}{subsection}{V35: Proved}
The native critical equation is the mass-polynomial system
\eqref{eq:v35-three-steps}, with the exact residual identities
\eqref{eq:v35-residuals}. The input mass jets are the ordered reference
coefficients \eqref{eq:v35-reference-jets}, including all differentiated
mass vertices and one cutoff factor in each scalar covariance derivative.
The highest-mass potential map has the necessary-and-sufficient tests
\eqref{eq:v35-potential-tests}, explicit primitive and residual, rank 285,
and the displayed coefficient bounds. Eleven critical exact normalization
directions and one additional covariant mass direction are constructed.

\subsection*{Inherited}
\addcontentsline{toc}{subsection}{V35: Inherited}
V1--V34 and all their labels are unchanged. The simultaneous hard
metric--scalar functional, relative matter complex, lower-reference
identity, mass-weighted Taylor assignment and prior analytic estimates
are inputs at their stated scopes. V34's active weight-two normalization
and the gravitational normalization remain fixed. V35 does not recalculate
the subcritical kinetic separator or claim an independent audit of all
458 preceding pages.

\subsection*{Literature input}
\addcontentsline{toc}{subsection}{V35: Literature input}
Local BRST classification with tensor matter is discussed in
\cite{V29BBHMatter}; the distinction between antifield/Euler and
longitudinal operations is also inherited from the cited BRST framework.
No external cohomology theorem is invoked as an unproved assertion of
solvability for the populated finite critical target. The elementary radial
homotopy is proved here on the precise finite polynomial space.

\subsection*{Conditional}
\addcontentsline{toc}{subsection}{V35: Conditional}
If the first two actual relative equations admit compatible chosen
solutions, the final input is \(\xi_2-Mh_1\), not \(\xi_2\) alone.
Its excluded components, horizontal test and curl test must pass before
\eqref{eq:v35-potential-primitive} solves the last equation. Error bounds
for this step consume actual errors of that input. No native value or
vanishing range obstruction is inferred from the small exact test systems.

\subsection*{Open}
\addcontentsline{toc}{subsection}{V35: Open}
Solve the massless critical relative equation \(Dh_0=\xi_0\) on the
fixed gravity branch, and determine which homogeneous choices have
\(\xi_1-Mh_0\in\operatorname{Ran}D_2\). Carry the resulting choice
into the explicit highest-mass test rather than changing the active
subcritical normalization. The full critical primitive and the combined
Einstein--Standard Model quantum continuum are not claimed here.

\begin{center}\small
\begin{tabular}{>{\raggedright\arraybackslash}p{.24\textwidth}>{\raggedright\arraybackslash}p{.67\textwidth}}
\toprule
Variable & Scope in V35\\\midrule
Fields and order & One-loop coefficient data; four equal neutral real
scalars, matter count two, total arity at most five.\\
Weight and mass & Critical weight four; exactly three mass powers
\(1,z,z^2\). Polynomial identities are regular at \(z=0\).\\
Cutoff and volume & New finite coefficient maps are independent of both.
Native reference coefficients use the existing thermodynamic convention;
no new finite-volume rate is proved.\\
Norms & The numerical bounds 2048 and 80 use independent symmetric metric
coordinates and specified coefficient \(\ell^1\) norms. Coframe pullback
has its own fixed finite conversion factor.\\
Reference & Same finite carrier, contour, ghost action, sources and density.
No profile is tuned to make an obstruction disappear.\\
Normalization & V34's subcritical tensor and the gravity source remain
unchanged. The critical exact family is described, not silently inserted.\\
Remaining computation & The first two native mass-block images and their
populated right-hand sides are not solved or numerically evaluated.\\
Uniformity & Fixed field/derivative degree and chart. No bound uniform in
unbounded arity, \(\chi\downarrow0\), or infinitely many loop orders.\\
\bottomrule
\end{tabular}
\end{center}


\clearpage
\part*{V36: Critical flat-normalization rigidity and the complete guarded two-point jet}
\addcontentsline{toc}{part}{V36: Critical flat normalization, mass-lift tests, and guarded matching}

\section{An upstream question about the homogeneous choices in the mass lift}
\label{sec:v36-start}

Sections 1--279 and all their labels are preserved. The active subcritical
prescription is still \eqref{eq:v34-Hmin}; the gravitational source
normalization is fixed. V35 supplied the three actual inhomogeneous
critical equations \eqref{eq:v35-relative-jets} and
\eqref{eq:v35-three-steps}, but did not solve their first two range
problems. In particular, the twelve flat homogeneous directions in
\cref{prop:v35-exact-family} were explicitly not proved exhaustive.
This part closes that last classification question and turns it into a
small quantitative test for admissible choices in the mass lift.

The result is a projection of the \emph{homogeneous} ordinary BV kernel,
not a substitute for evaluating the inhomogeneous reference arrays.
The complete flat critical projection has dimension twelve. Its
complement has thirty-four independently testable directions: twenty-five
harmonic quartic directions and nine anisotropic mismatches between the
massless and first-mass coefficients. A coefficient compiler and its
remainder have exact column norms two and nine-halves. These results
constrain changes between solutions of the \emph{same} critical equation;
they do not assert that its populated right-hand side lies in its image.

A second calculation consumes V27's already evaluated signed mixed
numerator. With one common smooth ultraviolet guard on the two ordinary
propagators, we calculate the \emph{entire weight-four flat local jet},
including its finite profile tensor. Its logarithm agrees with V27 and
is not claimed as new. The new finite tensor has an exact image under the
thirty-four-dimensional test. This calculation is on the inherited
ordinary comparison coefficient. The auxiliary ultraviolet guard is not
silently made the microscopic regulator, and no active profile is changed.

The scope is the same flavour-singlet, quadratic sector of four equal
neutral real scalars, at one loop, with internal gauge, Yukawa and
self-interaction couplings zero. The local proofs use the ordinary
zero-mesh differential \(d_H(z)\), \(z=m^2\), modulo total ordinary
derivatives. They hold through total arity five. The common factor
\(\hbar\) is suppressed in coefficient arrays and restored on action
insertion. The free scalar Euler symbol is
\[
 s_p=\sum_{i=0}^3p_i^2,\qquad P(p,z)=s_p+z.
\]
The letter \(s_p\) denotes this polynomial, not a BV differential.
No statement below assumes finite-cutoff nilpotence of the filtered source.

\section{Killing-jet rigidity of the flat critical projection}
\label{sec:v36-killing}

Let \(\Delta\) be a translation-invariant, flavour-singlet, matter-count-two
local functional of ghost number zero and weight four with
\(d_H(z)\Delta=0\) through arity five. Its pure-gravity restriction is zero.
The flat antifield-free scalar coefficient is necessarily
\begin{equation}
 [\Delta]_{\rm flat}
 =\frac12\sum_A\int_p\phi_A(-p)F(p,z)\phi_A(p),
 \qquad
 F=F_4(p)+zF_2(p)+d z^2,
 \label{eq:v36-flat-bank}
\end{equation}
where \(F_j\) is homogeneous of momentum degree \(j\). The coefficient
space has dimension \(35+10+1=46\). Scalar exchange eliminates odd
momentum polynomials. This is not a rotational-invariance assumption.

The zero-metric, one-scalar-antifield part of \(\Delta\) gives a scalar
transformation correction \(Q_X\), linear in the test vector field \(X\).
On affine Killing fields it has the normal-ordered form
\begin{equation}
 Q_X=X^\nu A_\nu(\partial,z)
       +(\partial_\alpha X^\nu)B^\alpha{}_{\nu}(\partial,z).
 \label{eq:v36-Q-affine}
\end{equation}
The operators \(A_\nu\) and \(B^\alpha{}_{\nu}\) have weights three
and two respectively. Terms containing two derivatives of \(X\) vanish
on this test. No restriction has been placed on those terms away from
this test. Work temporarily with real commuting differential symbols
\(y_i\), so the Euler polynomial is \(P_\partial=z-\sum_i y_i^2\).
The physical Fourier substitution \(y_i=ip_i\) returns \(P=s_p+z\).

\begin{lemma}[The translation part is scalar-Euler divisible]
\label{lem:v36-A-divisible}
The homogeneous cocycle equations imply
\begin{equation}
 \boxed{A_\nu(y,z)=P_\partial(y,z)L_\nu(y),}
 \label{eq:v36-A-divisible}
\end{equation}
where every \(L_\nu\) is linear in \(y\). This is polynomial divisibility,
not division of a propagator on a mass shell.
\end{lemma}
\begin{proof}
The scalar-algebra row on a flat metric and to first scalar order is the
linearized representation equation
\[
 [\mathcal L_X,Q_Y]+[Q_X,\mathcal L_Y]-Q_{[X,Y]}=0
       \quad\bmod\{T P_\partial:T\text{ a local operator}\}.
\]
Its right-ideal remainder is the scalar-Euler contribution of the
quadratic scalar-antifield terms. Higher-metric terms have free variation
\(R_0X=R_0Y=0\). A scalar-dependent ghost transformation contributes at
least three scalar arguments here. A mixed metric--scalar-antifield Euler
return contains either a metric field or two additional matter arguments.
Thus none supplies an extra term at this restriction. This is an
extraction of the complete cocycle equation, not an off-shell omission
of antifields.

Take \(X=u\) constant and \(Y=\Omega x\), \(\Omega^T=-\Omega\).
The terms \([\mathcal L_u,Q_{\Omega x}]\) and
\(Q_{[u,\Omega x]}\) cancel exactly. Since
\([A(y),x_j]=\partial_{y_j}A(y)\), the remaining constant-coefficient
operator is
\[
 u^\nu\Omega_{ij}y_i\partial_{y_j}A_\nu(y,z).
\]
Consequently every rotational derivative of each \(A_\nu\) is divisible
by \(P_\partial\). Reduction modulo \(z-\sum y_i^2\) commutes with
rotations and identifies the quotient with \(\R[y_0,\ldots,y_3]\).
The image of \(A_\nu\) is a homogeneous cubic rotational invariant.
Such a polynomial is zero: invariance under all infinitesimal rotations
implies invariance under the connected group \(SO(4)\), whose element
\(-I\) acts by minus one on cubics. Polynomial division by the monic
polynomial \(z-\sum y_i^2\) proves \eqref{eq:v36-A-divisible}.
Weight three leaves a weight-one quotient, hence a linear polynomial in
\(y\). No inverse numerical mass was used.
\end{proof}

\begin{lemma}[The flat action is rotationally invariant modulo its Euler polynomial]
\label{lem:v36-on-shell-action}
For every \(\Omega\in\mathfrak{so}(4)\),
\begin{equation}
 (\Omega p)\cdot\partial_p F\in (P).
 \label{eq:v36-on-shell-invariant}
\end{equation}
Equivalently \(F(p,-s_p)\) is a rotationally invariant homogeneous
quartic in the four independent momenta.
\end{lemma}
\begin{proof}
Evaluate the homogeneous action equation on the flat Killing field
\(\Omega x\), leaving the scalar arbitrary. The metric variation of all
higher-metric vertices is zero. In the flat pairing the rotation Lie
derivative is skew-adjoint. The remaining operator identity is the
commutator of the quadratic action with this Lie derivative plus
\(P_\partial Q_{\Omega x}+Q_{\Omega x}^{\dagger}P_\partial\), with
the harmless common normalization fixed by the scalar-antifield source.

Modulo the right Euler ideal, the second expression equals
\([P_\partial,Q_{\Omega x}]\). The constant-coefficient
\(B\) terms commute with \(P_\partial\). The other commutator is a
constant multiple of
\(\Omega_{ij}y_j A_i(y,z)\), which is divisible by \(P_\partial\)
by \cref{lem:v36-A-divisible}. The action commutator is the rotational
derivative of its Fourier polynomial. This proves
\eqref{eq:v36-on-shell-invariant}. In particular, the proof did not set
\(P\phi=0\) before retaining the scalar-antifield contribution.
\end{proof}

\begin{theorem}[Complete flat critical normalization image]
\label{thm:v36-flat-image}
For the declared local cocycles through arity five,
\begin{equation}
 \boxed{
 \operatorname{pr}_{\rm flat}\ker d_H(z)
 =\left\{(s_p+z)v_2(p,z)+\beta z^2:
        v_2=p^TAp+bz,\ A=A^T,\ b,\beta\in\R\right\}.}
 \label{eq:v36-image}
\end{equation}
Its dimension is exactly twelve. The equality concerns the flat
projection, not the dimension or cohomology of the complete BV kernel.
\end{theorem}
\begin{proof}
By \cref{lem:v36-on-shell-action}, the polynomial
\(r(p)=F(p,-s_p)\) is a homogeneous quartic invariant under all rotations.
Rotations act transitively on each sphere, so homogeneity gives
\(r=\beta s_p^2\). The polynomial \(F-\beta z^2\) therefore has zero
remainder upon division by the monic polynomial \(z+s_p\). Its quotient
has weight two and hence has the stated form \(v_2\).

For the reverse inclusion use the \emph{already constructed} family
\eqref{eq:v35-exact-family}: \(d_H K_V\), with constant self-adjoint
symbol \(v_2\), has flat coefficient \(P v_2\). Add
\(\beta\mathcal V_2\), where
\(\mathcal V_2=\frac12\sum_A\int\rho(q)z^2\phi_A^2\).
This is an ordinary covariant invariant. Both preserve the pure-gravity
restriction; truncating their full ordinary variations after arity five
preserves closedness. Thus every polynomial in the right-hand side is
attained with all of its source and metric partners.

Multiplication by \(P\) is injective on the eleven-dimensional
weight-two scalar polynomial space. The polynomial \(z^2\) is not in
that image, since its reduction modulo \(P\) is \(s_p^2\ne0\).
This proves the dimension. V35 proved the inclusion by construction;
the Killing-jet necessity above is the new exhaustion argument.
\end{proof}

The Killing fields in this proof are tests of local coefficient identities.
No affine ghost is asserted to be a periodic physical history. The
Euclidean rotation test uses the inherited ordinary coefficient branch;
complexifying its finite identities does not change the local ranks.
Independent flavour-breaking matrices and changes to the gravitational
source normalization are outside \eqref{eq:v36-image}.

\section{A norm-two coefficient compiler and the exact first-mass matching constraint}
\label{sec:v36-compiler}

For \(F\) in \eqref{eq:v36-flat-bank} put
\begin{equation}
 r=F_4-s_pF_2+d s_p^2,\qquad
 \beta=\frac{\Delta_p^2r}{192},\qquad
 v_2=F_2+(\beta-d)(s_p-z),
 \label{eq:v36-compiler}
\end{equation}
where \(\Delta_p=\sum_i\partial_{p_i}^2\). Define
\begin{equation}
 O(F)=F-Pv_2-\beta z^2
      =r-\beta s_p^2.
 \label{eq:v36-residual}
\end{equation}
The coefficient 192 is \(\Delta_p^2(s_p^2)\) in four dimensions.
All operations are polynomial coefficient operations.

\begin{theorem}[Explicit projection and exact coefficient norms]
\label{thm:v36-compiler}
The map \(F\mapsto Pv_2+\beta z^2\) is a projection onto
\eqref{eq:v36-image}; its complementary projection \(O\) has rank 34.
In the ordinary monomial coefficient \(\ell^1\) norms,
\begin{equation}
 \boxed{\|(v_2,\beta)\|_1\le2\|F\|_1,
       \qquad \|O(F)\|_1\le\tfrac92\|F\|_1.}
 \label{eq:v36-norms}
\end{equation}
Both constants are the exact induced column norms in these bases.
\end{theorem}
\begin{proof}
The formula for \(v_2\) gives \eqref{eq:v36-residual} by expansion.
If \(F=PV+\beta_0z^2\), then \(r=\beta_0s_p^2\), so the compiler
returns \(V,\beta_0\). Thus it is a left inverse on the stated image,
and \(O^2=O\). The image has dimension twelve inside the
46-dimensional domain, so \(O\) has rank 34.

For the exact norms it suffices to enumerate the eight permutation types
of unit monomials. The entries below are sums of absolute coefficients;
all are obtained by \eqref{eq:v36-compiler}.
\begin{center}\small
\begin{tabular}{lcc}
\toprule
Input & \(\|(v_2,\beta)\|_1\) & \(\|O(F)\|_1\)\\\midrule
\(p_0^4\) & \(3/4\) & \(11/4\)\\
\(p_0^3p_1\) & \(0\) & \(1\)\\
\(p_0^2p_1^2\) & \(1/4\) & \(3/2\)\\
\(p_0^2p_1p_2\), \(p_0p_1p_2p_3\) & \(0\) & \(1\)\\
\(zp_0^2\) & \(2\) & \(9/2\)\\
\(zp_0p_1\) & \(1\) & \(4\)\\
\(z^2\) & \(1\) & \(0\)\\\bottomrule
\end{tabular}
\end{center}
Every basis monomial belongs to one row, and the maxima are attained on
\(zp_0^2\). The column-norm formula for linear maps on coefficient
\(\ell^1\) proves \eqref{eq:v36-norms}.
\end{proof}

The thirty-four residual coordinates have an intrinsic split. Write
\[
 F_4=H_4+s_pH_2+a_0s_p^2,\qquad
 F_2=B_2+b_0s_p,
\]
where \(\Delta_pH_4=\Delta_pH_2=\Delta_pB_2=0\). The explicit
formulas are
\[
 a_0=\Delta_p^2F_4/192,\quad
 H_2=(\Delta_pF_4-24a_0s_p)/16,\quad
 b_0=\Delta_pF_2/8.
\]
Here harmonic means a polynomial in the four momenta, not a spacetime
gauge condition. Direct substitution gives
\begin{equation}
 \boxed{O(F)=H_4+s_p(H_2-B_2).}
 \label{eq:v36-harmonic-test}
\end{equation}
The degree-four harmonic space has dimension \(35-10=25\), and the
traceless quadratic space has dimension \(10-1=9\). These two parts
are independent: applying \(\Delta_p\) to their sum first identifies
the quadratic part, since \(\Delta_p(s_pH_2)=16H_2\).

\begin{corollary}[Flat matching across the critical mass lift]
\label{cor:v36-mass-lift}
A homogeneous change with flat coefficients
\(F_4+zF_2+dz^2\) is extendable to the stated ordinary critical cocycle
if and only if
\begin{equation}
 \boxed{H_4=0,\qquad H_2=B_2.}
 \label{eq:v36-mass-lift-test}
\end{equation}
Equivalently, \(F_4=s_pA_2\) and \(F_2-A_2\) is a scalar multiple
of \(s_p\). The coefficient \(d\) has no additional flat obstruction.
\end{corollary}
\begin{proof}
Combine \eqref{eq:v36-harmonic-test} with \cref{thm:v36-flat-image}.
When \(F_4=s_pA_2\) and \(F_2=A_2+b s_p\), use
\(v_2=A_2+bz\), \(\beta=d-b\) in the inherited full completion.
\end{proof}

For example the massless flat term \(s_pp_0^2\) has the explicit
massless exact completion from V35. Keeping its first-mass flat term
zero is not a massive homogeneous lift: \(H_2\ne B_2\). The
compatible polynomial \((s_p+z)p_0^2\) supplies the missing first-mass
coefficient. Conversely, \(p_0p_1p_2p_3\) is a harmonic quartic and
cannot be a flat homogeneous change on this branch. These examples are
normalization tests, not claimed values of the native reference arrays.

For two actual solutions of the same critical inhomogeneous equation,
\eqref{eq:v36-mass-lift-test} necessarily holds for their flat difference.
A flat difference passing it has a displayed homogeneous completion, but
this does not decide additional requirements on already prescribed
nonflat coefficients. The theorem therefore does not replace the full
relative range test by a flat one.

\section{The complete critical flat jet of the guarded native mixed coefficient}
\label{sec:v36-loop}

This calculation evaluates more local data than the logarithmic check
already contained in \eqref{eq:v27-angular}. It uses exactly the ordinary
mixed coefficient \eqref{eq:v27-cont-bubble} and contact
\eqref{eq:v27-contact}, not a new continuum Feynman rule.
Let \(w(p)\) be a real, even, smooth function of compact support, vanishing
near \(p=0\). Place the \emph{same} factor on the two primitive covariances,
\[
 C_g^w(p)=\frac{2w(p)}{\chi p^2}J,
 \qquad C_H^w(p)=\frac{w(p)}{p^2+z}.
\]
For instance \(w=\nu_\lambda\psi_\Lambda\) is a common finite ultraviolet
guard on the established ordinary reference. All terms are first
integrated with this guard. No divergent term is shifted independently.
The calculation below is thermodynamic; it does not assign an exact
continuous-momentum shift to an interpolated finite torus.

Set
\begin{equation}
 I_{ij}(w)=\int_p\frac{w(p)\partial_i\partial_j w(p)}{p^2},
 \qquad J(w)=\int_p\frac{w(p)^2}{p^4},
 \qquad \int_p=(2\pi)^{-4}\int_{\R^4}d^4p.
 \label{eq:v36-IJ}
\end{equation}
These are finite integrals. The tensor \(I\) is symmetric. Mass
derivatives act on the inverses, with the single unchanged \(w\)
factor, as in V28 and V35.

\begin{theorem}[Evaluated weight-four flat jet]
\label{thm:v36-critical-jet}
Let \(\Pi_w(k,z)\) be the complete contact-minus-mixed-loop two-Higgs
coefficient with the above guard, with the normalization of V27 and
\(\hbar\) removed. Its homogeneous weighted-degree-four part is
\begin{equation}
 \boxed{
 F_w(k,z):=[\Pi_w]_{[4]}
 =-\frac{k^2}{2\chi}k_iI_{ij}(w)k_j
   +\frac{2J(w)}{\chi}\,z(k^2+z).}
 \label{eq:v36-guarded-jet}
\end{equation}
In particular its image under the exact homogeneous-normalization test is
\begin{equation}
 \boxed{
 O(F_w)=-\frac{k^2}{2\chi}
 k^T\left(I(w)-\tfrac14\operatorname{tr}I(w)I_4\right)k.}
 \label{eq:v36-profile-test}
\end{equation}
There is no harmonic-quartic component in this coefficient; its possible
nonzero residual consists of the nine quadratic-anisotropy components.
\end{theorem}
\begin{proof}
The contact \eqref{eq:v27-contact}, with \(\nu\) replaced by \(w\),
is exactly of weighted degree two in the external momentum and mass.
It remains in that local assignment and contributes no degree-four term.
For the bubble use the already proved integrand identity
\eqref{eq:v27-partial-fraction}. Put
\[
 A(k)=\int_p\frac{w(p)w(p+k)}{p^2},\quad
 B(k,z)=\int_p\frac{w(p)w(p+k)}{(p+k)^2+z},
\]
\[
 C(k,z)=\int_p\frac{w(p)w(p+k)}{p^2((p+k)^2+z)},\quad
 D(k,z)=\int_p\frac{w(p)w(p+k)}{p^2(p^2+z)}.
\]
The paired finite integrals obey \(B=A-zD\): change variables
\(p\mapsto-p-k\) in \(B\), using evenness of \(w\), and use the
ordinary reciprocal identity. This operation is legitimate on the
\emph{guarded complete integral}. The signed numerator therefore gives
\begin{equation}
 \Pi_w^{\rm bub}
 =-\chi^{-1}k^2 A+\chi^{-1}z^2D
       +\chi^{-1}(2zk^2+z^2)C.
 \label{eq:v36-bubble-decomp}
\end{equation}
At the weighted Taylor origin,
\(C(0,0)=D(0,0)=J(w)\). Evenness makes \(A\) even, and its quadratic
coefficient is \(k_i k_j I_{ij}/2\). Taking degree four in
\eqref{eq:v36-bubble-decomp} proves \eqref{eq:v36-guarded-jet}.
The second summand is an inherited Euler-exact flat polynomial, so its
\(O\)-projection is zero. Decomposing \(I\) into trace and trace-free
parts in the first summand proves \eqref{eq:v36-profile-test}.
No extra flavour trace multiplies a fixed external Higgs pair.
\end{proof}

The profile tensor can also be written without second derivatives:
\begin{equation}
 I_{ij}
 =-\int_p\frac{\partial_iw\,\partial_jw}{p^2}
  -\delta_{ij}J(w)
  +4\int_p\frac{w^2p_ip_j}{p^6},
 \qquad
 \operatorname{tr}I=-\int_p\frac{|\nabla w|^2}{p^2}.
 \label{eq:v36-profile-energy}
\end{equation}
Both identities follow by integration by parts; boundary terms vanish
at infinity and near zero. This is an evaluated profile-dependent
criterion. Evenness alone need not make its trace-free tensor zero.

The object \(F_w\) is a local coefficient of the guarded \emph{ordinary
quantum functional}. It is not the populated ghost-one target
\(\Xi\), and it is not automatically the flat part of a solution of
\(d_HH=\Xi\). In particular, a nonzero value in
\eqref{eq:v36-profile-test} is a separator for changing this flat
coefficient by a homogeneous cocycle \emph{alone}, not a proof of a
quantum anomaly or a failed inhomogeneous critical equation.

\section{Radial guards: the finite quartic term and the inherited logarithm}
\label{sec:v36-radial}

For a radial \(w(p)=w_0(|p|)\), define
\[
 E_w=\int_0^\infty r\,[w_0'(r)]^2\,dr,\qquad
 L_w=\int_0^\infty w_0(r)^2\,\frac{dr}{r}.
\]
Both integrals are finite on the preceding class. The angular area is
\(2\pi^2\), with \(\int_p\) as in \eqref{eq:v36-IJ}.

\begin{proposition}[Exact radial finite matching data]
\label{prop:v36-radial}
For every such radial guard,
\begin{equation}
 I_{ij}(w)=-\frac{E_w}{32\pi^2}\delta_{ij},\qquad
 J(w)=\frac{L_w}{8\pi^2},
 \label{eq:v36-radial-IJ}
\end{equation}
and hence
\begin{equation}
 \boxed{
 F_w(k,z)=a_w(k^2)^2+b_wz(k^2+z),\qquad
 a_w=\frac{E_w}{64\pi^2\chi},\quad
 b_w=\frac{L_w}{4\pi^2\chi}.}
 \label{eq:v36-radial-jet}
\end{equation}
Its homogeneous-normalization residual is zero. A complete local
homogeneous functional realizing this flat coefficient is
\begin{equation}
 \boxed{
 \mathcal C_w=d_HK_{V_w}+a_w\mathcal V_2,
 \qquad v_w(p,z)=a_ws_p+(b_w-a_w)z.}
 \label{eq:v36-radial-completion}
\end{equation}
\end{proposition}
\begin{proof}
Rotational invariance gives \(I_{ij}=\delta_{ij}\operatorname{tr}I/4\).
Equation \eqref{eq:v36-profile-energy} and the radial integral give
\(\operatorname{tr}I=-E_w/(8\pi^2)\). The formula for \(J\) is
the same radial integration. Inserting these into
\eqref{eq:v36-guarded-jet} proves \eqref{eq:v36-radial-jet}.
Finally
\[
 P\{a_ws_p+(b_w-a_w)z\}+a_wz^2
       =a_ws_p^2+b_wzP.
\]
The full completion is exactly the inherited V35 construction, with the
newly evaluated coefficients substituted. Its BV variation is zero;
this is a permissible homogeneous shift, not a particular primitive for
\(\Xi\).
\end{proof}

There is a useful separation of the two transitions. Take
\[
 w_0(r)=[1-\theta(r/\lambda)]\psi(r/\Lambda),\qquad \Lambda>2\lambda,
\]
where \(\theta,\psi\) are smooth, equal to one for arguments at most
one, zero for arguments at least two, and take values in \([0,1]\).
This is a radial \emph{subcase} for analysis of the ordinary guard; no
radiality or replacement is imposed on the active finite reference.
Set
\[
 E_\theta=\int_1^2t\theta'(t)^2dt,\quad
 E_\psi=\int_1^2t\psi'(t)^2dt,
\]
\[
 c_\theta=\int_1^2[1-\theta(t)]^2\frac{dt}{t},\quad
 c_\psi=\int_1^2\psi(t)^2\frac{dt}{t}.
\]
The disjoint support of the derivatives gives the exact expressions
\begin{equation}
 \boxed{
 E_w=E_\theta+E_\psi,\qquad
 L_w=\log\frac{\Lambda}{\lambda}-\log2+c_\theta+c_\psi.}
 \label{eq:v36-separated-profiles}
\end{equation}
The lower and upper quartic contact contributions are both retained.
By Cauchy--Schwarz,
\(1\le E_\theta\log2\) and \(1\le E_\psi\log2\), so for
\(\chi>0\),
\begin{equation}
 a_w\ge\frac{1}{32\pi^2\chi\log2}.
 \label{eq:v36-quartic-bound}
\end{equation}
This lower bound concerns a finite reference-matching term in the stated
guard, not a physical higher-derivative coupling.

The logarithmic part of \eqref{eq:v36-radial-jet} is
\begin{equation}
 \frac{z(s_p+z)}{4\pi^2\chi}\log\frac{\Lambda}{\lambda}.
 \label{eq:v36-log-check}
\end{equation}
Its value was already established by \eqref{eq:v27-angular}; V36 does
not count it as a new beta-function calculation. The new information
is the complete finite quartic profile term and its image under the
exhaustive flat-normalization map. The logarithm has the inherited
Euler-exact completion with \(v_2\) proportional to \(z\).
Neither \eqref{eq:v36-log-check} nor
\eqref{eq:v36-radial-completion} identifies all metric-marked critical
loop coefficients with that completion.

\section{What the new test does and does not solve in the critical recursion}
\label{sec:v36-use}

The first two equations of V35 remain inhomogeneous:
\[
 Dh_0=\xi_0,\qquad Dh_1=\xi_1-Mh_0.
\]
Their homogeneous ambiguities can no longer be treated as independent
flat coefficients. If two \emph{full} mass-polynomial primitives are
compared, their massless flat difference must have the form \(s_pA_2\),
and their first-mass flat difference must equal \(A_2\) up to a scalar
multiple of \(s_p\). The scalar potential coefficient at the final
step has the additional invariant freedom already identified in V35.
A desired set of flat normalization changes violating these conditions
has an explicit finite left-null witness from
\eqref{eq:v36-harmonic-test}. This can be checked before attempting a
large relative system.

Conversely, when only a flat change is prescribed and it passes the test,
\(d_HK_V+\beta\mathcal V_2\) is a complete local homogeneous extension.
It supplies all metric and scalar-antifield contacts automatically. Thus
flat compatible changes do not require solving the large matrix again.
Prescribing nonflat coefficients simultaneously is a different relative
problem and may constrain the same choice further.

For a native coefficient approximation \(F+\delta F\),
\eqref{eq:v36-norms} gives errors at most \(2\|\delta F\|_1\) in
the compiled homogeneous parameters and
\(\tfrac92\|\delta F\|_1\) in its residual. The finite coframe
Taylor expansion of the inherited completion has a fixed additional
operator norm at arity five. It depends on the chart and degree, not on
mesh, volume or loop values. After a smooth external window is inserted,
the corresponding ordinary action kernels obey bounds of the form
\[
 C_{N,b}(\mu+m)^4\|(v_2,\beta)\|_1,
\]
with analogous degree-three scalar-antifield kernels and the declared
argument normalizations. These are rooted coefficient bounds, with one
argument in \(L^1\) and the others in \(L^\infty\), not extensive
action norms. No uniformity in unbounded arity or derivative order is
asserted.

Every homogeneous correction discussed here is of weight four and is
inserted at order \(\hbar\). It leaves V34's active weight-two
normalization unchanged, and has zero pure-gravity restriction. It does
not change the tree covariance, auxiliary Cartan integral, Haar density,
or the metric-coordinate Jacobian. Its scalar-antifield commutator and
all ordinary metric derivatives are retained by the same functional.
No member of the new homogeneous family is silently added to the active
critical scheme in this installment.

The guarded computation is a finite thermodynamic integral on the
ordinary comparison branch. It is not a new finite-cutoff Ward theorem,
a numerical evaluation of \(\xi_0,\xi_1\), or a proof of their local
exactness. In particular, the two-Higgs projection cannot detect a
curvature--scalar coefficient which vanishes on flat backgrounds.
The unresolved full critical equations contain those terms, as well as
additional antifield equations. The test supplies necessary normalization
data and an exact homogeneous compiler, not permission to drop them.

\section{Verification and comparison with established methods}
\label{sec:v36-verification}

The general distinction between consistent action deformations, Euler
terms and source changes is established in the local BRST framework,
including tensor matter \cite{V29BBHMatter}. The quoted classification
is not used to assert vanishing of a cohomology group in this part.
The necessity proof above is an explicit translation--rotation
calculation in the actual scalar source and Euler equations. The
sufficient family was already proved in V35. The signed loop numerator,
contact, partial fraction and logarithmic check are inherited from V27.
Only their complete guarded weight-four Taylor coefficient and its new
normalization test are evaluated here.

The self-contained script
\srcid{check_ESM_critical_flat_V36.py} generates all 46 weighted
flat monomials, checks the compiler and harmonic identities in exact
rational arithmetic, and computes the ranks of the 12-dimensional image
and 34-dimensional residual. It checks the exact column norms and the
25-plus-9 decomposition. Separate finite differential-operator tests
verify the translation--rotation commutator formula and its polynomial
mass reduction. The signed numerator and guarded partial-fraction
identity are checked without floating-point loop integration.
The radial profile identities are tested by exact integration on
explicit polynomial transition examples, labelled as algebraic tests
rather than new choices of the smooth regulator.

The preservation certificate compares the entire inherited mathematical
body, beginning at Section 1 and ending immediately before the old
continuation marker, byte for byte. The front matter and continuation
marker alone are updated. Exact finite identities do not independently
formalize every analytic or field-theoretic assertion in the cumulative
manuscript. No proof-assistant formalization or numerical native critical
range calculation is reported.

\section{Closing ledgers for V36}
\label{sec:v36-ledgers}

\subsection*{Proved}
\addcontentsline{toc}{subsection}{V36: Proved}
The exact flat projection of the homogeneous critical scalar--gravity
kernel is \eqref{eq:v36-image}; V35's twelve constructed flat directions
are exhaustive. The translation--rotation test proves the scalar-Euler
divisibility of the constant-ghost transformation coefficient. The
coefficient compiler and residual have exact column norms two and
nine-halves. Their rank-34 complement splits into a harmonic-quartic
rank-25 test and a rank-9 mass-lift anisotropy test. The complete guarded
ordinary flat mixed coefficient is \eqref{eq:v36-guarded-jet}, with the
explicit profile residual \eqref{eq:v36-profile-test}. The radial finite
term, its homogeneous completion, and its two-transition formula are
proved at the stated scope.

\subsection*{Inherited}
\addcontentsline{toc}{subsection}{V36: Inherited}
V1--V35 and their labels are unchanged. V35 supplies the actual
mass-polynomial inhomogeneous target, its mass-Euler recursion, the
highest-mass potential solver, and the sufficient exact family. V27
supplies the native signed continuum numerator and contact; its
logarithm is recovered, not presented as a new result. The active V34
subcritical scalar prescription and the fixed gravity source remain
unchanged. None of the earlier finite or analytic proofs is claimed to
have been independently reverified in its entirety.

\subsection*{Literature input}
\addcontentsline{toc}{subsection}{V36: Literature input}
The BRST consistent-deformation and antifield framework is the established
one described in \cite{V29BBHMatter}. No external no-anomaly theorem or
unproved relative splitting is used for the new finite projection.
The ordinary guarded integral is compared with the already derived
native continuum coefficient, not substituted for a new microscopic
Einstein--Standard Model action.

\subsection*{Conditional}
\addcontentsline{toc}{subsection}{V36: Conditional}
If two full critical primitives exist, their flat difference obeys
\eqref{eq:v36-mass-lift-test}. Conversely, a prescribed compatible flat
change has the stated local homogeneous completion; preservation of
additional nonflat conditions remains a separate test. Error bounds on
a native array can be transported through the norm-two compiler, but no
error estimate for an unevaluated array is manufactured here. The
radial-guard formulas apply only to that explicitly described subcase;
no radial profile is adopted for the active scheme.

\subsection*{Open}
\addcontentsline{toc}{subsection}{V36: Open}
Solve the populated massless inhomogeneous equation \(Dh_0=\xi_0\),
and choose its homogeneous part so that the complete next target
\(\xi_1-Mh_0\) lies in the required image. The new rank-34 flat test
can reject incompatible normalization choices and construct every
compatible flat homogeneous change, but it cannot decide the remaining
curved-background and antifield coefficients. The last input to V35's
potential solver is still \(\xi_2-Mh_1\). No complete critical primitive,
new many-shell theorem, mixed internal-gauge/chiral closure, or all-order
Einstein--Standard Model continuum is claimed in this part.

\begin{center}\small
\begin{tabular}{>{\raggedright\arraybackslash}p{.24\textwidth}>{\raggedright\arraybackslash}p{.67\textwidth}}
\toprule
Variable & Scope in V36\\\midrule
Fields and loop order & Four equal neutral real scalars, flavour-singlet
matter count two. New coefficient use is one loop; the local homogeneous
identities are ordinary BV algebra through arity five.\\
Mass and weight & Homogeneous critical weight four, polynomial in
\(z=m^2\), including all three powers \(1,z,z^2\). No inverse mass or
external momentum.\\
Cutoff and volume & The finite maps and their norms do not depend on
mesh or volume. The guarded trace is a thermodynamic comparison integral;
no new finite-volume rate is asserted.\\
Reference & Same action, ghost source, contour and measure as inherited.
The common ultraviolet guard is an explicit comparison operation;
radiality is a subcase, not a change of the active profile.\\
Norms & Exact constants 2 and \(9/2\) use ordinary monomial coefficient
\(\ell^1\) norms. Full coframe/source lift and window norms have fixed
additional degree- and chart-dependent constants.\\
Normalization & V34 weight two and the gravitational source are fixed.
The weight-four homogeneous compiler is described, not silently inserted.\\
Unresolved computation & No numerical native \(\xi_0,\xi_1\) integral,
full inhomogeneous matrix image, or complete critical primitive is
reported.\\
Uniformity & Fixed arity and derivative/window order; no uniformity in
unbounded source order, \(\chi\downarrow0\), infrared removal, or
arbitrarily many loop orders.\\\bottomrule
\end{tabular}
\end{center}

\clearpage
\part*{V37: The four-matter critical completion and the provenance of a relative primitive}
\addcontentsline{toc}{part}{V37: Four-matter completion and relative-normalization provenance}

\section{The quadratic-matter quotient is not the complete critical reference germ}
\label{sec:v37-start}

Sections 1--287 and their labels are preserved. V35's equations
\(Dh_0=\xi_0\) and \(Dh_1=\xi_1-Mh_0\) remain the active
inhomogeneous critical problem. V36 classifies its flat homogeneous
freedom but does not solve those equations. Repeating that flat test
would not determine a curved-background or antifield coefficient.
This part goes upstream to the ordinary \emph{joint} reference germ and
to the provenance of the gravitational normalization held fixed in the
relative problem.

There are two distinct issues. First, the matter-count-two quotient is
not the full matter-dependent local germ. We determine its missing
matter-count-four parent, and calculate a genuinely nonzero four-scalar,
four-derivative coefficient of the same simultaneous integral. Its full
completed-cutoff nonlocal remainder is controlled separately. Second, an
absolute local primitive, even if available by an external classification,
need not agree with a prescribed finite gravitational primitive. We give
the exact relative condition, an explicit native stress lift for compatible
changes, and a finite counterexample showing why a minimum-norm finite
primitive cannot simply be assumed to extend to a full germ.

The scope is one loop, four equal real scalars, vanishing internal gauge,
Yukawa and scalar self-couplings, and the completed finite reference of
V27--V36. The new four-scalar cutoff theorem is stated at \(m=0\)
with the lower Wilsonian scale \(\lambda>0\) retained. Mass derivatives
in the critical-reference construction use \(z=m^2\) and the previous
mass-weighted convention. A common factor \(\hbar\) is suppressed in
one-loop coefficients. The ordinary master differential is nilpotent;
no such assertion is made about the filtered finite source.

V34's active subcritical normalization and the prescribed gravitational
normalization are not changed. Newly calculated local terms are stored
as coefficient data. No critical source counterterm is silently selected
by the arguments below. The only outside classification discussed is
\cite{V29BBHMatter}; its conclusions concern the full specified local
BRST algebra, not an arbitrary normalized finite quotient.

\section{The measured four-scalar coefficient and its completed-cutoff limit}
\label{sec:v37-four-scalar}

Work first at zero retained metric and minimal antifields. Write \(C_g\)
and \(C_H\) for the actual metric and scalar hard covariances. They include
the inherited contour and lower and upper profiles. Use the actual
quadratic forms from V27:
\[
 U_\phi[x,y]=\tfrac12\langle\phi,Q_{qq}[x,y]\phi\rangle,
 \qquad \langle x,L_\phi\zeta\rangle
       =\langle Q_q[x]\phi,\zeta\rangle,
 \qquad W_\phi=U_\phi-L_\phi C_HL_\phi^T.
\]
\(U_\phi\) is quadratic and \(L_\phi\) linear in the retained scalar.
The superscript \(T\) means the inherited signed transpose in the even
metric--scalar Hessian; all functional pairings use the same physical
quadrature. There is no replacement of \(C_g\) by a positive TT covariance.

\Needspace{8\baselineskip}
\begin{proposition}[The complete one-loop four-scalar coefficient]
\label{prop:v37-four-scalar}
The homogeneous scalar-degree-four coefficient of the simultaneous
one-loop functional, with the normalization of V27, is
\begin{equation}
 \boxed{\Gamma_{a,1}^{\phi^4}
       =-\tfrac14\Tr(C_gW_\phi C_gW_\phi).}
 \label{eq:v37-four-scalar}
\end{equation}
It contains, with their actual relative signs,
\begin{align}
 -\tfrac14\Tr(C_gU_\phi C_gU_\phi),\qquad&
 +\tfrac12\Tr(C_gU_\phi C_gL_\phi C_HL_\phi^T),\
 &-\tfrac14\Tr(C_gL_\phi C_HL_\phi^T
                         C_gL_\phi C_HL_\phi^T).
 \label{eq:v37-three-words}
\end{align}
No scalar-flavour multiplicity is to be appended to these traces: the
flavour indices already occur in \(L_\phi\) and \(C_H\).
\end{proposition}
\begin{proof}
Introduce a scalar-background parameter \(\varepsilon\). The even hard
matrix is
\[
 \begin{pmatrix}C_g^{-1}+\varepsilon^2U_\phi&\varepsilon L_\phi\\
          \varepsilon L_\phi^T&C_H^{-1}\end{pmatrix}.
\]
Inverses in this display are restricted to the integrated support; the
resulting formula uses only covariances. Its Schur determinant is the
source-independent scalar determinant times
\(\det(C_g^{-1}+\varepsilon^2W_\phi)\). The coefficient of
\(\varepsilon^4\) in one half its logarithm is
\eqref{eq:v37-four-scalar}. Expanding the two \(W_\phi\)'s and using
cyclicity of the finite trace gives \eqref{eq:v37-three-words}.

Completeness requires more than this matrix identity. Every candidate
hard-linear vertex in a four-source coefficient carries only a sum of
four soft external momenta; it cannot meet the integrated support when
\(\lambda\ge 32\mu\). A connected one-loop graph with every hard valence
at least two has every hard valence equal to two, by
\(\sum_v(n_v-2)=2(L-1)=0\). All such vertices are in the displayed
Hessian. The gravitational ghost action has no scalar dependence in this
gauge at zero retained metric, and the auxiliary amplitude and coordinate
Jacobian are scalar independent. The pure scalar determinant is also
independent of the retained scalar when its self-interaction is zero.
None supplies an additional four-scalar term. This argument is on the
same perturbative hard support as V27; it is not a stationary-phase claim
for a different full nonlinear path integral.
\end{proof}

At zero mass the exact scalar-shift identity permits a representation of
each vertex with a derivative on every external scalar leg. Consequently
the coefficient in \eqref{eq:v37-four-scalar} has an external derivative
floor of four. Its Taylor jet below degree four vanishes. Simultaneous
momentum reversal leaves this source-free coefficient invariant: the
improved free matrices are even, and all scalar vertices entering these
three words are the sign-averaged vertices with at most two metric legs.
No oriented gravitational cubic enters this extraction. Thus its fifth
Taylor coefficient also vanishes.

\begin{theorem}[Massless four-scalar completed-cutoff comparison]
\label{thm:v37-cutoff}
Let \(\gamma_{a,L}\) be any polarized four-scalar coefficient of
\eqref{eq:v37-four-scalar}, and put
\(\gamma^R_{a,L}=(I-\mathsf T_4)\gamma_{a,L}\).
For
\begin{equation}
 \lambda\ge32768\mu>0,\qquad a\lambda\le2^{-24},\qquad
 \lambda L_\nu\ge1,\qquad m=0,
 \label{eq:v37-domain}
\end{equation}
and for four external legs bounded by \(\mu\), the same-volume comparison
satisfies
\begin{equation}
 \boxed{
 |\partial_K^\alpha(\gamma^R_{a,L}-\gamma^R_{0,L})(K)|
 \le C_\alpha\chi^{-2}\mathfrak n_I a^2|K|^{6-r},
 \qquad r=|\alpha|\le6.}
 \label{eq:v37-cutoff}
\end{equation}
Here \(K\) denotes three independent four-momenta and
\(\mathfrak n_I\) is the product of scalar polarization norms.
The result includes all completed Brillouin-zone modes. Constants are
independent of mesh and periods on \eqref{eq:v37-domain}, with the
inherited fixed profile seminorms and inverse margins. Higher prescribed
momentum derivatives have the corresponding fixed-window bounds.
\end{theorem}
\begin{proof}
The primitive words are precisely \eqref{eq:v37-three-words}. On their
common interior chart the V27 scalar symbols, including their two-metric
variation, agree with the ordinary ones through relative mesh degree
two. Sign averaging leaves the next discrepancy at relative degree
four. The metric inverse also has relative discrepancy
\(O((aR_p)^4)\), where \(R_p=\lambda+|p|\). The four-scalar extraction
uses no background derivative of the gravitational inverse and therefore
no gravitational cubic artifact of relative order three. All required
scalar vertices are already among the improved V12/V27 inputs.

A uniform, conservative symbol bound for each completed word is
\(C\chi^{-2}R_p^{-r}\) after \(r\) external derivatives;
its differentiated interior difference is bounded by
\(C\chi^{-2}a^4R_p^{4-r}\). A sharper undifferentiated bound keeps
one external momentum per scalar, but is not needed for the radial count.
Both bounds follow by differentiating the actual finite products of
inverse matrices and scalar vertices. The inverse derivative formula
retains every factor, and external derivatives of each fixed smooth
profile cost one power of \(R_p^{-1}\) on its transition region.

The coefficient is even, so \(\mathsf T_5\gamma=\mathsf T_4\gamma\).
Taylor's integral remainder after five degrees uses six derivatives and
gives the majorants
\[
 C\chi^{-2}|K|^6 R_p^{-6},\qquad
 C\chi^{-2}a^4|K|^6 R_p^{-2}
\]
for the ordinary remainder and its interior mismatch. In four dimensions
\(a^4\int_\lambda^{c/a}R\,dR\le Ca^2\), and
\(\int_{c/a}^\infty R^{-3}\,dR\le Ca^2\).
On the smooth periodic upper region a word has degree zero and its sixth
external derivative is bounded by \(Ca^6\); its normalized mode count
is \(O(a^{-4})\). This gives the same bound without assuming a smooth
extension of an uncompleted principal symbol. The fixed lower scale makes
all massless inverse differentiations regular on their hard support.

Replacing integrals by the normalized finite sums uses the inherited
four-dimensional counting estimate under \(\lambda L_\nu\ge1\).
Taylor remainder differentiation gives the factor \(|K|^{6-r}\) for
\(r\le6\). Higher fixed derivatives are bounded directly by the same
inverse/profile estimates with the chosen window scale. These arguments
prove \eqref{eq:v37-cutoff}; they do not choose the finite part of the
stored local fourth jet.
\end{proof}

\section{An evaluated nonzero critical four-gradient coefficient}
\label{sec:v37-gradient}

The next calculation is a guarded \emph{ordinary} coefficient, analogous
to V36's explicitly declared comparison. It does not replace the finite
regulator or its active profiles. Give both ordinary propagators the same
smooth radial weight \(w(p)=w_0(|p|)\), compactly supported away from zero:
\[
 C_g^w(p)=\frac{2w(p)}{\chi p^2}\mathsf J,
 \qquad C_H^w(p)=\frac{w(p)}{p^2}I_4.
\]
Work at \(m=0\). Take constant local scalar gradients
\(v_{A\mu}=\partial_\mu\phi_A\) and define their Gram matrix
\begin{equation}
 \mathsf V_{\mu\nu}=\sum_{A=1}^4v_{A\mu}v_{A\nu}
   =\begin{pmatrix}t&b^T\\b&S\end{pmatrix},\qquad
 s=\operatorname{tr}S,\quad u=\operatorname{tr}S^2,\quad
 \beta=b^Tb.
 \label{eq:v37-gradient-Gram}
\end{equation}
This is a local jet evaluation, not an affine scalar configuration
assumed to be periodic or bounded at infinity. The derivative floor above
shows that it determines the entire flat, massless, four-scalar
weight-four action polynomial: all four external scalar legs must carry
exactly one derivative.

For clarity the matrix calculation is specified completely. Let
\(q\) be a symmetric four-by-four matrix and use its ten independent
entries as coordinates. The inherited Euclidean coefficient coframe is
\[
 e=I+E,\quad E_{00}=q_{00}/2,\quad E_{0i}=0,\quad
 E_{i0}=q_{0i},\quad E_{ij}=q_{ij}/2\quad(i,j>0).
\]
Put \(d_1=\operatorname{tr}E\),
\(d_2=((\operatorname{tr}E)^2-\operatorname{tr}E^2)/2\). The second
coefficient of \(\det(e)(e^Te)^{-1}\) is
\begin{equation}
 W_2=d_2I-d_1q+q^2-E^TE.
 \label{eq:v37-W2}
\end{equation}
Let \(U\) be the ten-coordinate Hessian of
\(\tfrac12\operatorname{tr}(\mathsf V W_2)\), and define the vector
\(\ell_A(p)\) by differentiating
\[
 \tfrac12(v_A\cdot p)\operatorname{tr}q-v_A^Tqp
\]
with respect to those coordinates. In the same coordinates the constant
covariance matrix is
\begin{equation}
 \mathsf J_{(\mu\nu),(\rho\sigma)}
 =\tfrac12(\delta_{\mu\rho}\delta_{\nu\sigma}
           +\delta_{\mu\sigma}\delta_{\nu\rho})
       -\tfrac12\delta_{\mu\nu}\delta_{\rho\sigma},
 \quad \mu\le\nu,\ \rho\le\sigma.
 \label{eq:v37-J}
\end{equation}
In particular its off-diagonal-coordinate variance is \(1/2\).
Using an identity matrix in those ten coordinates would change the result.

Define three quadratic polynomials in the gradient Gram data:
\begin{align}
 A(\mathsf V)&=\tfrac14t^2-\tfrac3{16}ts-\tfrac3{64}s^2
                    +\tfrac{41}{64}u+\tfrac78\beta,\notag\\
 B(\mathsf V)&=\tfrac1{16}t^2+\tfrac1{32}ts+\tfrac1{64}s^2
                    +\tfrac{11}{64}u+\tfrac14\beta,\notag\\
 C(\mathsf V)&=\tfrac14t^2+\tfrac14u+\tfrac12\beta.
 \label{eq:v37-ABC}
\end{align}
Finally put \(L_j(w)=\int_0^\infty w_0(r)^j\,dr/r\), for \(j=2,3,4\).

\begin{theorem}[Complete guarded four-gradient coefficient]
\label{thm:v37-gradient}
With the normalization of \eqref{eq:v37-four-scalar}, the homogeneous
quartic action density on these scalar jets is
\begin{equation}
 \boxed{
 \mathcal F_w(\mathsf V)
 =-\frac{1}{8\pi^2\chi^2}
       \{L_2(w)A(\mathsf V)-2L_3(w)B(\mathsf V)
                              +L_4(w)C(\mathsf V)\}.}
 \label{eq:v37-finite-gradient}
\end{equation}
This is a homogeneous action coefficient, not a fourth derivative with
an additional factorial. Polarization gives the four-source tensor.
\end{theorem}
\begin{proof}
Let \(S_p=\sum_A\ell_A(p)\ell_A(p)^T\). The two Fourier factors in
the scalar Schur return are \(i\ell_A\) and \(-i\ell_A\), so its
sign is positive before the subtraction from \(U\). Inserting the guarded
covariances into \eqref{eq:v37-four-scalar} gives
\[
 -\frac{1}{\chi^2}\int_p\frac1{p^4}
 \left\{w^2\Tr(\mathsf JU\mathsf JU)
 -\frac{2w^3}{p^2}\Tr(\mathsf JU\mathsf JS_p)
 +\frac{w^4}{p^4}\Tr(\mathsf JS_p\mathsf JS_p)\right\}.
\]
There are respectively two, three, and four primitive covariance weights.
They must not all be replaced by \(w^2\).

The definitions \eqref{eq:v37-W2}--\eqref{eq:v37-J} give
\[
 \Tr(\mathsf JU)=-t/4+s,
 \qquad
 \ell_A(p)^T\mathsf J\ell_B(p)
                 =\tfrac12p^2(v_A\cdot v_B).
\]
The first reproduces the inherited two-scalar contact normalization.
The second gives
\[\Tr(\mathsf JS_p\mathsf JS_p)/(p^2)^2
 =\tfrac14\operatorname{tr}\mathsf V^2=C(\mathsf V).\]
Direct contraction of the remaining finite matrices gives
\[
 \Tr(\mathsf JU\mathsf JU)=A(\mathsf V),\qquad
 \operatorname{av}_{S^3}
 \frac{\Tr(\mathsf JU\mathsf JS_p)}{p^2}=B(\mathsf V).
\]
For an explicit way to check the latter without an angular ansatz, write
\(\ell_A(p)=\sum_\mu p_\mu T_\mu v_A\), where \(T_\mu\) is obtained
by differentiating the displayed linear form. Then the average is
\(\tfrac14\sum_\mu\Tr(\mathsf JU\mathsf JT_\mu
\mathsf V T_\mu^T)\). Substitution gives exactly the five coefficients
of \(B\) in \eqref{eq:v37-ABC}; the same substitution gives those of
\(A\). These are polynomial identities for an arbitrary symmetric
\(\mathsf V\), independently checked in the accompanying script.
The radial factor with \(\int_p=(2\pi)^{-4}\int d^4p\) is
\((8\pi^2)^{-1}\int dr/r\), proving the formula.
\end{proof}

Take the separated radial guard
\(w_0(r)=[1-\theta(r/\lambda)]\psi(r/\Lambda)\),
with \(\Lambda>2\lambda\) and the same endpoint conventions as V36.
Then, exactly,
\begin{equation}
 L_j(w)=\log\frac{\Lambda}{\lambda}-\log2+c_j^-+c_j^+,
 \quad c_j^-=\int_1^2\frac{(1-\theta(t))^j}{t}\,dt,
 \quad c_j^+=\int_1^2\frac{\psi(t)^j}{t}\,dt.
 \label{eq:v37-separated}
\end{equation}
Both transition contributions remain in the finite local coefficient.
The logarithmic polynomial is
\begin{equation}
 \boxed{
 \mathcal K(\mathsf V)=A-2B+C
 =\tfrac38t^2-\tfrac14ts-\tfrac5{64}s^2
                         +\tfrac{35}{64}u+\tfrac78\beta.}
 \label{eq:v37-K}
\end{equation}
Thus the log contribution is
\(-\mathcal K(\mathsf V)\log(\Lambda/\lambda)/(8\pi^2\chi^2)\).

\begin{corollary}[A strictly nonzero populated critical coefficient]
\label{cor:v37-positive}
For every scalar gradient Gram \(\mathsf V\succeq0\),
\begin{equation}
 \boxed{\frac9{280}(\operatorname{tr}\mathsf V)^2
 \le\mathcal K(\mathsf V)
 \le\frac{15}{32}(\operatorname{tr}\mathsf V)^2.}
 \label{eq:v37-positive}
\end{equation}
The lower bound is sharp on the four-scalar carrier. In particular this
critical four-matter coefficient is not identically zero.
\end{corollary}
\begin{proof}
For the block decomposition \eqref{eq:v37-gradient-Gram},
\(u\ge s^2/3\) and \(\beta\ge0\). Hence
\[
 \mathcal K\ge\tfrac38t^2-\tfrac14ts+\tfrac5{48}s^2
 =\tfrac9{280}(t+s)^2+\tfrac{(24t-11s)^2}{1680}.
\]
Equality is attained for \(b=0\), \(S=(s/3)I_3\), and
\(24t=11s\), which is a positive four-by-four Gram and can be realized
by four scalar gradient vectors. Conversely \(u\le s^2\) and
\(\beta\le ts\), the latter following from positive two-by-two
principal minors. Inserting these gives
\(\mathcal K\le3t^2/8+5ts/8+15s^2/32\le15(t+s)^2/32\).
\end{proof}

For one scalar gradient, or parallel flavour gradients,
\(\mathsf V=vv^T\), the logarithmic action density is explicitly
\begin{equation}
 -\frac{12v_0^4+20v_0^2|\boldsymbol v|^2
                    +15|\boldsymbol v|^4}{256\pi^2\chi^2}
                 \log\frac{\Lambda}{\lambda}.
 \label{eq:v37-rankone}
\end{equation}
The time--space distinction reflects the declared coframe-coordinate and
gauge reference. It is not a physical Lorentz-violation prediction, an
on-shell amplitude, or an anomaly. A full source-complete counterterm has
its metric and antifield partners. The calculation determines a populated
local component and does not identify that component with the ghost-one
critical reference target.

\section{The critical reference has only three possible even matter degrees}
\label{sec:v37-reference}

To discuss local cohomology, return to the \emph{ordinary} complete
metric--scalar minimal master action, with its nonminimal restoration and
its original source. Scale only retained matter fields and their
antifields by a parameter \(\varepsilon\). The unscaled hard carrier
and its covariances are kept fixed. Since the scalar master action is
quadratic in matter, its hard Hessian perturbation and source Jacobian have
exact background expansions
\[
 V(\varepsilon)=V_0+\varepsilon V_1+\varepsilon^2V_2,
 \qquad J(\varepsilon)=Dr=J_0+\varepsilon J_1.
\]
The zero-matter block \(V_0\) includes the scalar hard kinetic operator on
the retained geometry; it must not be replaced by a gravity-only Hessian.
The ownership of its scalar determinant follows the same once-counted
convention used in V33's relative reference equation.

Put
\[
 E_0=(I+CV_0)^{-1},\qquad A=E_0CV_1,\qquad B=E_0CV_2.
\]
This notation differs from the gradient polynomials in the preceding
section and is used only in this ordered operator formula.

\begin{proposition}[The missing four-matter reference word]
\label{prop:v37-reference-word}
The matter-four part of the ordinary reference trace is obtained, before
its local weight and arity extraction, from
\begin{equation}
 \boxed{\begin{aligned}
 \mathscr X_4=\operatorname{STr}\bigl[&
 (A^4-A^2B-ABA-BA^2+B^2)E_0\Theta J_0\\
 &+(-A^3+AB+BA)E_0\Theta J_1\bigr].
 \end{aligned}}
 \label{eq:v37-reference-four}
\end{equation}
The analogous matter-two term is
\(\operatorname{STr}[(A^2-B)E_0\Theta J_0-AE_0\Theta J_1]\).
Every product retains its displayed matrix order. These are coefficients
of the same reference functional, not new physical source laws.
\end{proposition}
\begin{proof}
Factor
\((I+CV(\varepsilon))^{-1}
 =(I+\varepsilon A+\varepsilon^2B)^{-1}E_0\).
If its first factor is \(\sum_n\varepsilon^nE_n\), then
\(E_0^{\rm series}=I\), \(E_1=-A\), and
\(E_n=-AE_{n-1}-BE_{n-2}\). This gives
\(E_2=A^2-B\), \(E_3=-A^3+AB+BA\), and
\(E_4=A^4-A^2B-ABA-BA^2+B^2\).
Multiplication by \(E_0\Theta(J_0+\varepsilon J_1)\) yields
\eqref{eq:v37-reference-four}. No commutative binomial formula is used.
\end{proof}

\begin{theorem}[Finite matter completion of the actual critical germ]
\label{thm:v37-germ}
On the inherited ordinary local coframe chart, the critical reference
has a local germ, analytic in the undifferentiated coframe and polynomial
in the remaining finite jets, of the form
\begin{equation}
 \boxed{\mathscr X^{[4]}(z)
   =\mathscr X_0^{[4]}(z)+\mathscr X_2^{[4]}(z)
                         +\mathscr X_4^{[4]}(z),
 \qquad \deg_z\mathscr X^{[4]}\le2.}
 \label{eq:v37-germ}
\end{equation}
The subscripts are total matter counts, including scalar antifields.
There is no matter-six or higher critical reference coefficient. The
full ordinary reference consistency identity implies
\begin{equation}
 \begin{gathered}
 d_g\mathscr X_0=0,\qquad
 d_{(2)}\mathscr X_2+t\mathscr X_0=0,\\
 d_{(4)}\mathscr X_4+t\mathscr X_2=0,\qquad
 t\mathscr X_4=0.
 \end{gathered}
 \label{eq:v37-reference-consistency}
\end{equation}
Here \(d_{(m)}\) preserves matter count \(m\), and \(t\) raises it by
two. At each fixed field degree the coefficients are compact lower-annulus
integrals. The all-arity ordinary germ is a comparison construction;
it is not a claim of convergence of an interacting finite-lattice field
series or a new finite action with infinitely many vertices.
\end{theorem}
\begin{proof}
The local construction uses the same frozen-coframe inverse and ordered
symbol expansion as V25, now on the joint metric--scalar--ghost carrier.
On the lower annulus the ordinary free inverse is uniformly bounded;
resumming the undifferentiated metric perturbation on a sufficiently small
fixed chart gives a convergent analytic inverse. Positive derivative
jets and nilpotent jets are then expanded only to the chosen weight.
The nonconstant reference words contain a covariance and \(\Theta\), so
at zero external transfers their support is
\(\lambda\le|p|\le2\lambda\). The weight-four no-covariance term is
absent by degree. All differentiations of these local coefficients are
therefore under finite, smooth, lower-annulus integrals.

It remains to bound matter degree; analyticity alone would not do so.
At zero mass, the displayed scalar kinetic action and source have a
derivative on each ordinary external scalar. The same is true after hard
Hessian or source-Jacobian differentiation. A mass insertion can supply at
most two undifferentiated external scalars per power of \(z\); derivatives
of a scalar inverse supply none. Thus a term with \(d\) external
spacetime derivatives and mass power \(z^j\) has at most \(d+2j\)
ordinary scalar factors, in this shift-adapted representation.

Let \(r,s,v\) count metric, ghost and scalar antifields. At ghost number
one and weight four,
\(d+2j+r+s+v=5\). Its total matter count is therefore at most
\(d+2j+v=5-r-s\le5\), and \(j\le2\).
The simultaneous sign change of all scalar fields and scalar antifields
leaves the ordinary master action, reference and trace invariant. Odd
matter degrees consequently vanish. This proves the precise bound
\(M\le4\), rather than assuming a quadratic-matter truncation.

The full ordinary regulated trace identity used in V22--V26 gives
\(s_\infty\mathscr X^{[4]}=0\) modulo total derivatives, before matter
or arity truncation. On this action
\(s_\infty=d_\parallel+t\), with \(t^2=0\): the scalar master action
has no gravitational antifields, so its gravitational self-antibracket
vanishes. Extracting its successive even matter degrees gives
\eqref{eq:v37-reference-consistency}. Restricting to any finite arity
commutes only after forming these identities. The bound on differentiated
scalar factors also proves that this germ is polynomial in scalar jets,
not merely an arbitrary analytic scalar dependence. Its native weight
and mass coefficients remain those of the original finite-word compiler.
\end{proof}

The third equation in \eqref{eq:v37-reference-consistency} is essential
when attempting a full-germ cohomological argument. The count-two quotient
used in V33 is valid for its relative equation, but it is not a substitute
for the entire cocycle. At arity five the new source-free family includes
\(c\phi^4\). Some of its higher-arity metric and antifield partners are
needed to certify extension before truncation; they are not inferred from
closure of the arity-five array alone.

\section{What a full local primitive must do to preserve the fixed gravity source}
\label{sec:v37-relative}

The following relative statement separates two requirements that an
absolute cohomology result cannot identify. It is a theorem about the
specified complexes, not an assertion that the remaining native image
test has already been evaluated.

Let \(x=(x_g,x_2)\) be the actual reference in the matter-count-at-most-two
quotient and let \(H_g\) be the prescribed gravity primitive,
\(d_gH_g=x_g\). Suppose an absolute primitive
\((\widetilde H_g,\widetilde H_2)\) has been obtained in the \emph{same}
weight, mass-polynomial, current and arity convention:
\[
 d_g\widetilde H_g=x_g,\qquad
 d_H\widetilde H_2=x_2-t\widetilde H_g.
\]
Set \(k_g=H_g-\widetilde H_g\), so \(d_gk_g=0\).

\begin{theorem}[The exact relative lift and its invariant part]
\label{thm:v37-relative}
A primitive agreeing with \(H_g\) exists if and only if
\begin{equation}
 \boxed{t k_g\in\operatorname{im}d_H.}
 \label{eq:v37-relative-test}
\end{equation}
If, in the same local category, the mismatch has a representation
\begin{equation}
 k_g=I_g+d_gQ_g,
 \qquad tI_g=0,
 \label{eq:v37-mismatch}
\end{equation}
then an explicit relative matter primitive is
\begin{equation}
 \boxed{H_2=\widetilde H_2+tQ_g.}
 \label{eq:v37-relative-lift}
\end{equation}
Here \(Q_g\) has ghost number minus one. The construction preserves local
polynomial mass dependence and introduces no inverse spacetime operator.
An antifield-free gravitational invariant satisfies \(tI_g=0\).
\end{theorem}
\begin{proof}
The difference of the two matter equations is
\(d_H(H_2-\widetilde H_2)=-tk_g\), proving the equivalence.
For \eqref{eq:v37-mismatch}, use the actual relation
\(d_Ht+td_g=0\). It gives
\(d_H(tQ_g)=-td_gQ_g=-tk_g\), which proves
\eqref{eq:v37-relative-lift}. The stress map differentiates a local
functional in its gravitational antifields and inserts the quadratic
scalar Euler/stress packet. It is local and polynomial in \(z\).
On an antifield-free gravity invariant it vanishes directly, without
using equations of motion.
\end{proof}

The formula is useful beyond flat tests. In metric coordinates, take
\(Q_g=\int\eta_g^{\mu\nu}V_{\mu\nu}(g)\), with a local
symmetric metric variation \(V\) of weight two. Then the action part
of its native matter lift is exactly
\begin{equation}
 tQ_g=\frac12\sum_A\int\sqrt g\,
 V_{\mu\nu}
 \left[-\nabla^\mu\phi_A\nabla^\nu\phi_A
  +\tfrac12g^{\mu\nu}
       (|\nabla\phi_A|^2+z\phi_A^2)\right].
 \label{eq:v37-stress-lift}
\end{equation}
The cotangent change back to the affine coframe coordinate is included
when this functional is expressed in the active variables. This is the
variation of the same scalar action, not an independently defined stress.
For \(V_{\mu\nu}=a_0R_{\mu\nu}+b_0g_{\mu\nu}R\), it gives
\begin{equation}
 \begin{split}
 tQ_g=\sum_A\int\sqrt g\,[&-\tfrac{a_0}{2}R_{\mu\nu}
                          \nabla^\mu\phi_A\nabla^\nu\phi_A\\
 &+(\tfrac{a_0}{4}+\tfrac{b_0}{2})R|\nabla\phi_A|^2
 +(\tfrac{a_0}{4}+b_0)zR\phi_A^2].
 \end{split}
 \label{eq:v37-curved-lift}
\end{equation}
All three terms have zero flat quadratic-scalar restriction. They
illustrate concretely why V36's flat test cannot determine the curved
relative lift. For a noncovariant local \(V\), the gravitational
source terms in \(d_gQ_g\) are also retained; the Jacobi proof of the
relative lift is unchanged. Neither \(Q_g\) nor these curvature
coefficients is installed as an active normalization in V37.

\section{A finite minimum-norm primitive need not have the needed full-germ lift}
\label{sec:v37-finite-obstruction}

V25 proved full-germ membership of the \emph{reference target}, and then
prescribed a finite minimum-norm primitive. It did not prove that every
primitive of that finite target is the truncation of a full primitive.
The distinction matters when that very primitive is subsequently held
fixed while coupling matter. The following example isolates this issue;
it is not evidence that the native gravitational primitive fails.

\begin{proposition}[An exact normalized finite-quotient counterexample]
\label{prop:v37-counterexample}
There is a filtered triangular complex with zero full degree-one
cohomology and an extendable target, but whose minimum-norm gravitational
primitive at the retained arity admits no relative matter completion.
\end{proposition}
\begin{proof}
Let the gravity degree-zero generators be \(a,b\), both of arity four;
let its degree-one generators be \(x\) of arity four and \(y\) of
arity six. Define
\[
 d_ga=x,\qquad d_gb=x+y,\qquad d_gx=d_gy=0.
\]
Let the matter space have \(z_1\) in degree one and arity five, and
\(w_2\) in degree two and arity seven, with
\(d_Hz_1=-w_2\). There is no matter degree-zero generator.
Set \(tb=z_1\), \(ty=w_2\), and \(ta=tx=0\).
Then \(d_Ht+td_g=0\) exactly. The full differential is
\[
 sa=x,\qquad sb=x+y+z_1,\qquad sy=w_2,\qquad sz_1=-w_2.
\]
Its degree-one kernel is \(\operatorname{span}\{x,y+z_1\}\), exactly
the image of degree zero. Thus its degree-one cohomology is zero, and the
target \(x\) has the full primitive \(a\).

Truncate at arity five. The gravitational matrix is now \((1\ \ 1)\).
In the orthonormal coefficient basis its minimum-norm primitive of \(x\)
is \(H_g=(a+b)/2\). The total truncated differential sends it to
\(x+z_1/2\). The required matter correction would have differential
\(-z_1/2\), but the matter degree-zero space is zero. No such correction
exists. All differentials are arity nondecreasing. Both the target and an
absolute solution extend to the full complex; it is the \emph{fixed finite
normalization} that fails to extend. This proves the claim.
\end{proof}

The example rules out a shortcut, not a native branch. Applying
\cite{V29BBHMatter} to a full tensor-matter local algebra, where its
hypotheses have been checked, addresses an absolute local cohomology
question. The theorem does not state that an arbitrary finite
minimum-norm choice is the restriction of its primitive. The additional
translation-invariance, mass-polynomial and relative-normalization
requirements must also be respected. V37 neither assumes a native
counterexample nor transfers the above finite example into the actual
coefficient arrays.

A direct finite diagnostic remains available. Choose integrated finite
bases, or include current columns before quotienting total derivatives,
and let \(A_g,A_H,T\) be the actual three matrices. If \(A_gH_g=x_g\),
put
\begin{equation}
 \boxed{r_H=x_2-TH_g,\qquad
 \mathcal O_H=(I-A_HA_H^+)r_H.}
 \label{eq:v37-direct-test}
\end{equation}
This is equivalent to relative solvability and does not require an
absolute primitive first. The pseudoinverse and matrix-height bounds are
the inherited V25 finite-linear-algebra machinery, not new spectral
estimates. The supplied utility evaluates \eqref{eq:v37-direct-test}
for an explicitly supplied rational coefficient system. The native
matrices and their annular right-hand sides are not built or evaluated
by V37. In particular \(\mathcal O_H\) is not reported as zero.

For a mass-block implementation the single operator is the whole
\(D+zM\) system of V35. Solving separate numerical-mass systems or using
independent homogeneous choices in its first two equations is not
\eqref{eq:v37-direct-test}. The necessary flat tests remain those of V36.

\section{Source topology, EFT ownership, and the next noncircular computation}
\label{sec:v37-use}

The new four-scalar nonlocal theorem has the same physical source
interpretation as the inherited two-scalar module. Windowing the three
independent momenta at scale \(\mu\) and differentiating beyond the
chosen moment order gives
\begin{equation}
 \boxed{a^{12}\sum_{x_1,x_2,x_3}
 \left[1+\mu\sum_{i=1}^3d_{\rm per}(x_i,0)\right]^b
 \|\mathcal K_{a,L}^{\Delta}(x_1,x_2,x_3)\|
 \le C_b\chi^{-2}a^2\mu^6.}
 \label{eq:v37-kernel}
\end{equation}
This follows by inverse Fourier integration by parts using the higher
fixed derivatives in \cref{thm:v37-cutoff}, then periodization. It bounds
a polarized functional with one scalar source in \(L^1\) and the other
three in \(L^\infty\), without multiplying by the physical volume.
The separate ordinary finite-volume remainder is bounded by
\(C_b\chi^{-2}\mu^6\lambda^{-2}(\lambda L_{\min})^{-b}\) for fixed
\(b>4\), using Poisson summation on the six-times differentiated loop.

The three words in \eqref{eq:v37-three-words} have two, three and four
primitive covariance factors. Assigning their increments to the largest
internal scale, including every scalar line, gives complementary bounds
\(C\chi^{-2}\mu^6\Lambda_j^{-2}\). They are summable. Their local
fourth jets are \emph{not} summable without renormalization: the evaluated
logarithm in \cref{thm:v37-gradient} is precisely a retained local
critical contribution. Matched completions have the same complementary
limit when their own fourth jets are assigned to the same physical local
data. No proof of a full covariant/source completion of those data is
contained in this scalar-only comparison.

The local four-gradient operator is of order \(\hbar\chi^{-2}\).
It is neither the renormalizable Higgs potential nor a replacement for
its physical coupling. It is a generated EFT interaction, with all four
real components represented by \(\mathsf V\). Adding its source family
to the recorded coefficient bank is required by the programme's filtered
operator inventory. The calculation does not introduce a tree-level
scalar self-coupling or modify V34's weight-two prescription.

The critical relative problem now has three cleanly separated pieces:
the actual full local reference is available, including its four-matter
parent; compatible full homogeneous gravitational changes have the
explicit stress lift \eqref{eq:v37-relative-lift}; and a prescribed finite
gravity normalization still has to pass \eqref{eq:v37-direct-test}, or be
shown to have an appropriate full-germ lift. A literature exactness result
cannot silently provide that last piece. Conversely, a failed flat scalar
test is not a substitute for evaluating it.

A productive next computation is therefore the relative image for the
\emph{actual fixed critical gravity prescription}, or an explicit
full-germ representation of that prescription. Either route must retain
its massless and first-mass coupling and the curved stress terms. The
present four-scalar calculation proves that a full-germ argument cannot
be carried out on the quadratic-matter quotient alone. No conclusion is
made about which of the two native routes will close first.

\section{Verification and closing ledgers for V37}
\label{sec:v37-ledgers}

The checker derives the ten-coordinate scalar contact matrix directly
from \eqref{eq:v37-W2}, verifies the three quartic gradient polynomials
for an arbitrary symmetric gradient Gram, checks the sharp lower-bound
identity and its realizing Gram, and checks the one-loop determinant sign
on an independent rational joint block. It also verifies the ordered
four-matter resolvent coefficient, the exact triangular-complex example,
and the native curvature--stress contraction. Its relative-system utility
returns a solution and residual for a \emph{supplied} rational system.
No native critical \(\xi_0\), \(\xi_1\), or gravitational-normalization
array is numerically evaluated. Polynomial and matrix checks are not
presented as verification of all inherited analytic proofs.

\subsection*{Proved}
\addcontentsline{toc}{subsection}{V37: Proved}
The simultaneous four-scalar coefficient has the three retained Schur
words \eqref{eq:v37-three-words}; its massless full completed-cutoff
nonlocal remainder obeys \eqref{eq:v37-cutoff} and
\eqref{eq:v37-kernel}. The entire guarded flat four-gradient coefficient
is \eqref{eq:v37-finite-gradient}; its logarithm is nonzero with the sharp
lower bound \eqref{eq:v37-positive}. The critical joint reference has only
matter degrees zero, two and four, and its ordered matter-four operand is
\eqref{eq:v37-reference-four}. The exact relative-lift theorem, native
stress formulas, and finite normalized-quotient counterexample are proved.

\subsection*{Inherited}
\addcontentsline{toc}{subsection}{V37: Inherited}
V1--V36, including every label, are unchanged. The completed carrier,
contour, scalar improvement through two metric legs, propagator bounds,
ordinary master action, reference identity, and source-shift property are
used at their stated scope. V34's subcritical normalization is active;
V25's fixed gravity prescription is not changed. The general finite
pseudoinverse test, stationary Schur principle and method of rooted Fourier
bounds are inherited and are not advertised as new general theorems.

\subsection*{Literature input}
\addcontentsline{toc}{subsection}{V37: Literature input}
The tensor-matter classification in \cite{V29BBHMatter} includes
antifields and separates semisimple from abelian internal-gauge cases.
It motivates a full-local-germ approach, but is not used here as an
unverified relative image theorem for the populated finite normalization.
The modified-reference framework is \cite{V15IIM}. No first construction
of scalar--gravity perturbation theory or discovery of general gravitational
nonrenormalizability is claimed.

\subsection*{Conditional}
\addcontentsline{toc}{subsection}{V37: Conditional}
If an absolute critical primitive is established in the same local
category and the fixed gravitational mismatch has the representation
\eqref{eq:v37-mismatch}, the relative matter primitive is explicitly
\eqref{eq:v37-relative-lift}. Neither hypothesis is silently inferred
from a minimum-norm finite solve. A simultaneous covariant normalization
of the newly populated four-scalar sector requires its source and metric
partners; its flat logarithm alone does not provide them.

\subsection*{Open}
\addcontentsline{toc}{subsection}{V37: Open}
The populated equations \(Dh_0=\xi_0\) and
\(Dh_1=\xi_1-Mh_0\), with V35's final potential target, remain open.
The upstream normalization issue has an exact relative test but no
reported native residual value. Internal gauge/chiral fields, arbitrary
loop order, infrared removal and an invariant full interacting ESM class
remain beyond this installment. No analytic or nonperturbative quantum
continuum is inferred from the new one-loop coefficient.

\begin{center}\small
\begin{tabular}{>{\raggedright\arraybackslash}p{.24\textwidth}>{\raggedright\arraybackslash}p{.67\textwidth}}
\toprule
Variable & Scope in V37\\\midrule
Fields and loop order & Four real equal scalars; new source-free
four-scalar coefficient at one loop and order \(\chi^{-2}\).
No internal, Yukawa or scalar self-interaction couplings.\\
Mass & New nonlocal theorem at \(m=0\), with \(\lambda>0\) retained.
The critical reference completion is polynomial through \(z^2\).\\
Cutoff and volume & Completed finite regulator on \eqref{eq:v37-domain};
rooted bounds are uniform in periods. Radial guarded coefficient uses
thermodynamic normalization; it is not the active finite reference.\\
Sources & Four ordinary scalar marks, all prescribed finite momentum
and rooted-moment orders. The full reference germ is used only as an
ordinary local comparison.\\
Normalization & V34 and the fixed gravity coefficients unchanged.
Local fourth jets stored; no new critical primitive installed.\\
Uniformity & Constants depend on fixed profile bounds, chart,
\(\chi>0\), and derivative order. No uniform infinite-arity,
all-loop or zero-infrared-scale claim.\\
Computation & Exact local coefficient contractions and symbolic
identities evaluated. Native inhomogeneous critical matrix and right-hand
side not evaluated.\\\bottomrule
\end{tabular}
\end{center}


\clearpage
\part*{V38: A source-complete stress lift of the native four-gradient logarithm}
\addcontentsline{toc}{part}{V38: Native logarithm, stress lift, and composite-source transport}

\section{Which part of the critical problem can be completed without changing gravity}
\label{sec:v38-start}

Sections 1--295 and their labels are preserved. V37 leaves the populated
critical relative equation unresolved. In particular, the finite
minimum-norm gravitational primitive has not been proved to extend to a
full germ. We neither assume that extension nor repeat its abstract image
test. Instead we complete a \emph{specific measured local coefficient}
that V37 has already evaluated: its four-scalar, four-gradient logarithm.
The completion uses a scalar-dependent metric redefinition, whose
curvature and antifield partners are calculated below. It has zero
pure-gravity restriction, so the prescribed gravitational normalization is
not changed.

There are three logically different statements. First, the actual finite
local fourth jet has the universal logarithmic coefficient of V37, with
a bounded finite remainder. Second, every quadratic polynomial of the
four-scalar gradient Gram has a bounded algebraic preimage under the
\emph{native metric stress}; the logarithmic polynomial is solved
explicitly. Third, that preimage gives a complete ordinary BV-exact local
functional, not just a four-scalar subtraction. This is a homogeneous
critical freedom. It does not solve the inhomogeneous equation
\(d_HH_2=x_2-tH_g\), and cannot change its obstruction class.

The scope remains the same completed Euclidean coefficient reference,
Palatini branch, four real equal scalars, and zero internal gauge, Yukawa
and scalar self-interaction couplings. New ultraviolet statements concern
one loop at \(m=0\), with \(\lambda>0\) fixed. The displayed local
functional also has a polynomial continuation in \(z=m^2\). This
continuation is not a computation of a massive four-scalar logarithm.
V34's subcritical normalization is unchanged. Only the explicitly stated
critical logarithmic subtraction is added.

Field redefinitions and BV-exact deformations are established tools
\cite{V29BBHMatter,V38Criado}. The new objects are the native stress map,
its exact small coefficient matrix, the evaluated preimage, the retained
curvature/source contacts, and their connection to the completed-cutoff
local jet. No cohomology-vanishing theorem is used in the new proofs.

\section{The native completed regulator has the same local logarithm}
\label{sec:v38-log}

Let \(\mathcal P_{a,L}(\mathsf V)\) be the homogeneous action polynomial
specified by the actual massless four-scalar local fourth jet of
\eqref{eq:v37-four-scalar}. All primitive covariances and their finite
normalizations are retained. The gradient Gram is
\[
 \mathsf V_{\mu\nu}=\sum_{A=1}^4
       (\partial_\mu\phi_A)(\partial_\nu\phi_A)
       =\begin{pmatrix}t&b^T\\b&S\end{pmatrix},\qquad
 s=\operatorname{tr}S,\quad u=\operatorname{tr}(S^2),\quad
 \beta=b^Tb.
\]
Flavour contractions in the three Schur words pair scalar legs using only
Kronecker tensors. The shift identity gives one derivative on each
external scalar at this local order. Hence \(\mathcal P_{a,L}\) is a
quadratic polynomial in the ten entries of \(\mathsf V\), even when the
finite ultraviolet completion has no continuous spatial rotation
symmetry. There are 55 such polynomial coefficients. Four scalar flavours
allow every positive definite Gram locally, so equality on physical
Gram data determines the polynomial uniquely.

\begin{theorem}[A bounded finite part after extraction of the native logarithm]
\label{thm:v38-log}
On the domain \eqref{eq:v37-domain}, with fixed admissible profile
seminorms and inverse margins,
\begin{equation}
 \boxed{
 \mathcal P_{a,L}(\mathsf V)
 =-\frac{\log(1/(a\lambda))}{8\pi^2\chi^2}
                  \mathcal K(\mathsf V)
       +\mathcal B_{a,L}(\mathsf V),\qquad
 \|\mathcal B_{a,L}\|_{\mathrm{coeff}}
                  \le C\chi^{-2},}
 \label{eq:v38-local-log}
\end{equation}
where the inherited evaluated polynomial is
\begin{equation}
 \mathcal K=\frac38t^2-\frac14ts-\frac5{64}s^2
                    +\frac{35}{64}u+\frac78\beta.
 \label{eq:v38-K}
\end{equation}
The constant has no total-volume or mesh dependence for
\(\lambda L_\nu\ge1\). Boundedness of \(\mathcal B_{a,L}\) is not
convergence or scheme independence of its coefficients.
\end{theorem}
\begin{proof}
Use the actual three words of \eqref{eq:v37-three-words}, before taking
norms. V37 supplies their differentiated interior mismatch bound
\(C\chi^{-2}a^4R_p^{4-r}\), with \(R_p=\lambda+|p|\).
At the fourth external jet this is \(C\chi^{-2}a^4\).
Integrating or summing through \(c/a\), for one fixed common plateau
constant \(c>0\), costs at most \(Ca^{-4}\); its total contribution
is therefore bounded by \(C\chi^{-2}\). The smooth completed upper
region has fourth-jet integrand bounded by \(C\chi^{-2}a^4\) and
normalized mode count \(O(a^{-4})\), giving the same conclusion.
The lower transition, after \(p=\lambda y\), also has a bounded
fourth-jet integral and finite sum.

In the intermediate region \(2\lambda<|p|<c/a\) the ordinary
covariance weights are one. Its local integrand is homogeneous of degree
minus four. The angular contraction already evaluated in V37 is exactly
\(-\mathcal K/(8\pi^2\chi^2)\); radial integration gives
\(-\mathcal K\log(c/(2a\lambda))/(8\pi^2\chi^2)\).
The constant \(\log(c/2)\) belongs to \(\mathcal B_{a,L}\).
This uses a radial decomposition of the interior integral, not a radial
assumption on either active transition profile.

For the finite sums, partition the intermediate region with a fixed
smooth dyadic partition. The ordinary shell symbol at scale \(R\) is
\(R^{-4}\) times a uniformly smooth compactly supported symbol.
Poisson summation, or integration by parts in its Fourier transform,
bounds sum-minus-integral by \(C_b(RL_{\min})^{-b}\), \(b>4\).
The dyadic sum of these errors is bounded when
\(\lambda L_{\min}\ge1\). The lower and upper normalized counting
bounds are the same as in V37. This proves the asserted volume-uniform
finite-part bound. No individual divergent continuum integral has been
assigned a value, and no native finite matching coefficient is discarded.
\end{proof}

The logarithm in \eqref{eq:v38-local-log} is now identified for the
completed regulator itself. Its coefficient is inherited from V37; the
new assertion is the bounded comparison of the \emph{unsubtracted local
fourth jet}, rather than another proof of the previously established
nonlocal \(O(a^2)\) theorem.

\section{An evaluated stress map with a bounded right inverse}
\label{sec:v38-stress}

Work first at flat metric and zero mass. On
\(\mathcal S=\operatorname{Sym}_4(\R)\) use the trace inner product and
its orthogonal involution
\[
 \mathsf JX=X-\tfrac12\operatorname{tr}(X)I_4,
 \qquad \mathsf J^2=I.
\]
The derivative of the \emph{same} scalar action with respect to the
covariant metric is
\begin{equation}
 \mathcal E_H^{\mu\nu}
       =-\tfrac12(\mathsf J\mathsf V)_{\mu\nu}
       \quad(g=I,z=0).
 \label{eq:v38-flat-stress}
\end{equation}
Thus a linear metric variation \(F:\mathcal S\to\mathcal S\) produces
\(\mathcal E_H:F(\mathsf V)\).

\begin{theorem}[The full 55-coefficient stress lift]
\label{thm:v38-stress}
For every quadratic polynomial
\(p_B(\mathsf V)=\langle\mathsf V,B\mathsf V\rangle\),
\(B=B^*\) on the ten-dimensional space \(\mathcal S\), define
\begin{equation}
 \boxed{F_B(\mathsf V)=-2\mathsf J B\mathsf V
                         =-\mathsf J\nabla p_B(\mathsf V).}
 \label{eq:v38-general-F}
\end{equation}
Then \(\mathcal E_H:F_B=p_B\). The map from all 100 coefficients
of \(F\) to the 55 quadratic coefficients has rank 55 and kernel
\begin{equation}
 \boxed{F=\mathsf J\Omega,\qquad \Omega^*=-\Omega,}
 \label{eq:v38-kernel}
\end{equation}
of dimension 45. Formula \eqref{eq:v38-general-F} is the unique
least-Hilbert--Schmidt-norm lift, and
\(\|F_B\|_{\mathrm{HS}(\mathcal S)}=2\|B\|_{\mathrm{HS}(\mathcal S)}\).
These norms refer to linear maps on \(\mathcal S\), not to an
extensive action.
\end{theorem}
\begin{proof}
The output quadratic operator is
\(-\tfrac14(\mathsf JF+F^*\mathsf J)\).
Taking \(F=-2\mathsf JB\) gives \(B\). The zero-output condition
is precisely that \(\mathsf JF\) be skew-adjoint, proving the kernel
and ranks. Every lift is \(-2\mathsf JB+\mathsf J\Omega\).
The involution \(\mathsf J\) is an isometry, while self-adjoint and
skew-adjoint operators are orthogonal. Hence its squared norm is
\(4\|B\|_{\mathrm{HS}}^2+\|\Omega\|_{\mathrm{HS}}^2\), proving
uniqueness and the exact bound. All operations are at one spacetime point.
\end{proof}

There is a smaller rational compiler for the native logarithm. Write
\(p=k_1t^2+k_2ts+k_3s^2+k_4u+k_5\beta\), and set
\[
 F_{00}=at+bs,\qquad F_{0i}=cb_i,\qquad
 F_{ij}=(dt+es)\delta_{ij}+fS_{ij}.
\]
In coefficient order \((a,b,c,d,e,f)\), the native stress map is
\begin{equation}
 \mathsf D=\begin{pmatrix}
 -1/4&0&0&3/4&0&0\\
 1/4&-1/4&0&1/4&3/4&1/4\\
 0&1/4&0&0&1/4&1/4\\
 0&0&0&0&0&-1/2\\
 0&0&-1&0&0&0
 \end{pmatrix}.
 \label{eq:v38-D}
\end{equation}

\begin{proposition}[Small native coefficient compiler]
\label{prop:v38-small}
An exact right inverse is
\begin{equation}
 \mathsf R=\begin{pmatrix}
 -1&3/2&0&0&0\\
 0&-1/2&3&1&0\\
 0&0&0&0&-1\\
 1&1/2&0&0&0\\
 0&1/2&1&1&0\\
 0&0&0&-2&0
 \end{pmatrix},\qquad
 \mathsf D\mathsf R=I_5.
 \label{eq:v38-R}
\end{equation}
Its coefficient norm is \(\|\mathsf R\|_{1\to1}=4\), and
\(\|\mathsf D\|_{1\to1}=1\). The kernel of \(\mathsf D\) is
\(\operatorname{span}\{(3,1,0,1,-1,0)\}\).
For \eqref{eq:v38-K}, the lift is
\begin{equation}
 \boxed{\begin{aligned}
 F_{00}^{\mathcal K}&=-\tfrac34t+\tfrac7{16}s,\\
 F_{0i}^{\mathcal K}&=-\tfrac78b_i,\\
 F_{ij}^{\mathcal K}&=
       (\tfrac14t+\tfrac{11}{32}s)\delta_{ij}
                    -\tfrac{35}{32}S_{ij}.
 \end{aligned}}
 \label{eq:v38-native-F}
\end{equation}
It satisfies \(\operatorname{tr}F^{\mathcal K}=3s/8\).
\end{proposition}
\begin{proof}
Insert \eqref{eq:v38-flat-stress}. The five output coefficients are
\((-a+3d)/4\), \((a-b+d+3e+f)/4\), \((b+e+f)/4\),
\(-f/2\), and \(-c\). This proves \eqref{eq:v38-D}. Direct rational
multiplication proves \eqref{eq:v38-R}; the absolute column sums give
its norm and that of \(\mathsf D\). Solving the five homogeneous
relations gives the stated kernel. Applying \(\mathsf R\) to
\((3/8,-1/4,-5/64,35/64,7/8)\) gives
\eqref{eq:v38-native-F} and its trace. The formula agrees with
\(-\mathsf J\nabla\mathcal K\), with off-diagonal entries counted
twice in the trace pairing.
\end{proof}

The zero-output kernel is not discarded: its curved and source partners
can be nonzero. The rank statements concern the evaluated flat stress
map, not the unresolved 130-million-row critical system of V25.

\section{The ordinary BV completion includes curvature and antifields}
\label{sec:v38-BV}

Use metric coordinates only to write the following local functional.
The active variable remains \(q\), with
\(g=G(q)=(I+\mathsf E(q))^T(I+\mathsf E(q))\).
Extend \(F^{\mathcal K}\) by the same component polynomial
\eqref{eq:v38-native-F} in \(\partial\phi\), independent of \(g\).
The constant coordinate tensors in that formula are not asserted to
transform as additional dynamical fields. Its noncovariance is paid by
the antifield term below.

Let \(\eta_g\) be the metric antifield, with its cotangent relation
\(\langle\eta_g,\delta g\rangle
=\langle\eta_q,DG(q)^{-1}\delta g\rangle\). Define
\begin{equation}
 Q_F=\int\eta_g^{\mu\nu}F_{\mu\nu}(\partial\phi)
     =\int\langle\eta_q,DG(q)^{-1}F\rangle,
 \qquad \mathcal C_F=s_\infty Q_F.
 \label{eq:v38-generator}
\end{equation}
Here \(s_\infty\) is the same ordinary full scalar--gravity differential
as in V35--V37. Its metric Euler term is defined by the inherited action
variation \(\delta S=\int\mathcal E^{\mu\nu}\delta g_{\mu\nu}\).

\begin{theorem}[Full local completion of the evaluated polynomial]
\label{thm:v38-BV}
The ghost-zero, weight-four functional \(\mathcal C_F\) is
\begin{equation}
 \boxed{
 \mathcal C_F
 =\int\mathcal E_g:F
   +\int\mathcal E_H:F
   +\int\eta_g:\mathcal R_F(c,\phi),\qquad
 \mathcal R_F=\mathcal L_cF-DF[\mathcal L_c\phi].}
 \label{eq:v38-full-completion}
\end{equation}
It obeys \(s_\infty\mathcal C_F=0\) modulo a total derivative and
vanishes when the scalar and scalar-antifield sector is removed. It has
only matter counts two and four. For \(F=F^{\mathcal K}\),
\begin{equation}
 \boxed{
 [\mathcal C_F]_{g=I,\eta_g=0}
   =\int\left[\mathcal K(\mathsf V)
           +\frac{3z}{32}\Big(\sum_A\phi_A^2\Big)s\right].}
 \label{eq:v38-mass-extension}
\end{equation}
The restriction at \(z=0\) is exactly the native polynomial. The
additional massive term is a specified homogeneous completion, not an
evaluation of the native massive logarithm.
\end{theorem}
\begin{proof}
Apply the inherited antibracket to \(Q_F\), using physical quadrature
for all adjoints. Its action part is the full metric Euler covector
contracted with \(F\). The metric-antifield part is the linearized
metric transformation applied to \(F\), minus the variation of \(F\)
under the scalar transformation. This gives the sign in
\eqref{eq:v38-full-completion}, consistent with the scalar
\([\mathcal L_c,V]\) term in V35. No ghost-antifield or scalar-antifield
term is omitted: the generator contains neither, and the scalar
transformation has no metric dependence in these coordinates.
Nilpotence of the ordinary differential gives the closedness assertion.

Explicitly the native scalar metric Euler density is
\[
 \mathcal E_H^{\mu\nu}
 =\frac{\sqrt g}{2}\sum_A
 \left[-\nabla^\mu\phi_A\nabla^\nu\phi_A
 +\frac12g^{\mu\nu}
       (|\nabla\phi_A|^2+z\phi_A^2)\right].
\]
The massless flat contraction is Theorem~\ref{thm:v38-stress}.
The remaining contraction at flat metric is
\(z(\sum_A\phi_A^2)\operatorname{tr}F/4\), proving
\eqref{eq:v38-mass-extension}. The gravitational Euler density contains
no scalar, so its product with \(F\) has matter count two. The last
term has two scalars and a metric antifield. The scalar stress product
has matter count four. These observations prove every stated restriction.
\end{proof}

For a constant tensor map \(F=\mathsf B(\mathsf V)\), the source partner
has the completely explicit component form
\begin{equation}
 \begin{split}
 (\mathcal R_F)_{\mu\nu}
 ={}& (\partial_\mu c^\lambda)F_{\lambda\nu}
       +(\partial_\nu c^\lambda)F_{\mu\lambda}\\
 &-\mathsf B_{\mu\nu}^{\rho\sigma}
    \left[(\partial_\rho c^\lambda)\mathsf V_{\lambda\sigma}
          +(\partial_\sigma c^\lambda)\mathsf V_{\rho\lambda}\right].
 \end{split}
 \label{eq:v38-source-explicit}
\end{equation}
Transport terms cancel because \(\mathsf B\) is constant. This formula
has two scalar derivatives and one ghost derivative. It is genuinely
nonzero. For \(\mathsf V=e_0e_0^T\) and
\(\partial_\nu c^\mu\) equal to the rotation with entries
\(A^0{}_1=1\), \(A^1{}_0=-1\), direct substitution gives
\begin{equation}
 (\mathcal R_{F^{\mathcal K}})_{01}=-\frac18.
 \label{eq:v38-source-witness}
\end{equation}
Thus the logarithmic polynomial is not being relabelled as an invariant
scalar interaction by itself.

Writing \(\mathcal C_F=\mathcal C_{F,2}+\mathcal C_{F,4}\), the
same triangular differential as in V37 gives
\begin{equation}
 \boxed{
 d_{(2)}\mathcal C_{F,2}=0,\qquad
 d_{(4)}\mathcal C_{F,4}+t\mathcal C_{F,2}=0.}
 \label{eq:v38-triangle}
\end{equation}
In particular, adding this packet is a homogeneous matter normalization
with fixed gravity. It cannot alter the class of the inhomogeneous target
\(x_2-tH_g\). Its first curved action contact is
\(\int (H_g^{\mathrm{cl}}h):F(\partial\phi)\), together with the
explicit \(\eta\phi^2c\) source above. These are precisely the
contacts that a flat four-scalar subtraction alone would lose.

\section{Recorded logarithmic subtraction and its finite local scope}
\label{sec:v38-subtraction}

At one loop define the explicit logarithmic counterterm
\begin{equation}
 \boxed{
 C_{a,\log}
 =\frac{\log(1/(a\lambda))}{8\pi^2\chi^2}
       \operatorname{pr}_{A\le5}\mathcal C_{F^{\mathcal K}},
 \qquad S_a\longmapsto S_a+\hbar C_{a,\log}.}
 \label{eq:v38-log-counterterm}
\end{equation}
Its massless flat four-scalar restriction cancels the logarithm in
\eqref{eq:v38-local-log}. The finite coefficient
\(\mathcal B_{a,L}\) is retained unchanged. No unevaluated finite
critical normalization is selected.

\begin{proposition}[Preservation of gravity, subcritical data, and the local identities]
\label{prop:v38-preservation}
The update \eqref{eq:v38-log-counterterm} has zero pure-gravity
restriction and weight four only. It leaves the fixed pure-gravity
normalization, V34's weight-two scalar normalization, and the flat
quadratic-scalar critical coefficient unchanged. In the ordinary
arity-five complex it is exactly closed. On the inherited nonaliased
soft panel, its actual finite-source variation satisfies
\begin{equation}
 \boxed{\mathbf P_{A\le5,w\le4}(s_aC_{a,\log})=0.}
 \label{eq:v38-finite-projection}
\end{equation}
The unprojected finite variation is not asserted to vanish.
\end{proposition}
\begin{proof}
Every monomial of \(Q_F\) has two scalar arguments. Every term of its
variation therefore has matter count at least two. The first action
term needs a nonzero gravitational Euler row, so its flat quadratic-scalar
restriction is zero; the other action term is scalar quartic. Weight
counting gives four throughout. Since the ordinary differential cannot
lower total arity, truncation of its closed full local functional remains
closed through arity five.

On the soft panel all profile arguments, including sums of the finitely
many legs, lie on their plateaux. The finite and ordinary transformation
vertices consequently agree there. The only difference in the variation
of the antifield term of \eqref{eq:v38-full-completion} is the native
metric Euler discrepancy. Its linear term has at least six derivatives;
its nonlinear restored gravitational terms have at least five. The
source coefficient multiplying that Euler row has three derivatives,
including one on the ghost. Hence all these outputs have weight strictly
greater than four. The corrected scalar-stress discrepancy has the same
or a greater excess. The antifield-free part varies only by the
transformation, already equal on the plateau. This proves
\eqref{eq:v38-finite-projection}. It does not use finite all-mode
nilpotence or change the reference measure.
\end{proof}

The local counterterm is the variation of one generator, so its curved
mixed action and antifield coefficients must be changed at the same time.
Their logarithmic factors can have different powers of \(\chi\): the
factor \(\mathcal E_g\) itself includes the gravitational normalization.
We do not assign \(\chi^{-2}\) indiscriminately to every component.

An actual finite fourth jet not invariant under spatial rotations can
also be lifted using Theorem~\ref{thm:v38-stress}, with all 55 coefficient
slots. That is a specified linear operation on a \emph{supplied} array,
not a numerical evaluation of that array. V38 installs only the evaluated
universal logarithmic packet. Its addition neither fixes the unknown
critical reference primitive nor resolves the relative minimum-norm
extension question of V37.

\section{Composite sources and the measured coordinate change}
\label{sec:v38-sources}

The operator \(F\) is local, but it changes geometric observables.
For a physical source coupled to an actual observable \(\mathcal O(g,\phi)\),
the corresponding first contact is
\begin{equation}
 \boxed{\delta_F\mathcal O=D_g\mathcal O[F].}
 \label{eq:v38-observable}
\end{equation}
This is the existing measured field-redefinition rule applied to the
explicit generator, not a new assertion that action normalization alone
renormalizes all insertions.

\begin{proposition}[Evaluated volume and curvature contacts]
\label{prop:v38-curvature}
At flat metric, the native variation \eqref{eq:v38-native-F} gives
\begin{equation}
 \delta_F\sqrt g=\frac{3}{16}s,
 \qquad
 \delta_FR_{\mu\nu}
 =\frac12\left(
 \partial_\rho\partial_\mu F_{\rho\nu}
 +\partial_\rho\partial_\nu F_{\rho\mu}
 -\partial^2F_{\mu\nu}
 -\partial_\mu\partial_\nu\operatorname{tr}F\right),
 \label{eq:v38-Ricci}
\end{equation}
and
\begin{equation}
 \boxed{\begin{aligned}
 \delta_FR={}&
 (-\tfrac34\partial_0^2+\tfrac14\Delta_{\mathrm{sp}})t
 +(\tfrac1{16}\partial_0^2-\tfrac1{32}\Delta_{\mathrm{sp}})s\\
 &-\tfrac74\sum_i\partial_0\partial_i b_i
 -\tfrac{35}{32}\sum_{i,j}\partial_i\partial_j S_{ij}.
 \end{aligned}}
 \label{eq:v38-curvature-R}
\end{equation}
Thus a curvature source acquires a specific four-derivative,
scalar-quadratic contact. It is not a source-independent vacuum term.
\end{proposition}
\begin{proof}
Differentiate the determinant and the ordinary torsion-free curvature
in the same covariant-metric convention. At flat metric these give
\(\delta\sqrt g=\operatorname{tr}F/2\) and
\(\delta R=\partial_\mu\partial_\nu F_{\mu\nu}
-\partial^2\operatorname{tr}F\). Insert
\eqref{eq:v38-native-F}; the off-diagonal time--space entries occur
twice. Their coefficient is \(-7/4\). Using
\(\operatorname{tr}F=3s/8\) gives the remaining coefficients exactly.
\end{proof}

For completeness, the cotangent lift in the active coordinate has a
controlled local domain. The native linear coframe map obeys
\(\|\mathsf E(q)\|_{\mathrm{HS}}
\le2^{-1/2}\|q\|_{\mathrm{HS}}\). Hence
\[
 \|\mathsf B_q\|\le\|q\|_{\mathrm{HS}},\qquad
 \|DG(q)^{-1}\|\le(1-r)^{-1}\quad(\|q\|_{\mathrm{HS}}\le r<1).
\]
Through arity five its generator is simply
\begin{equation}
 [Q_F]_{A\le5}
 =\int\langle\eta_q,(I-\mathsf B_q+\mathsf B_q^2)F\rangle.
 \label{eq:v38-coframe-jet}
\end{equation}
There is no spatial inverse hidden in this coordinate inverse.

An order-\(\hbar\) redefinition
\(g\mapsto g+\hbar c_aF\) changes the tree action at order \(\hbar\)
by its Euler contraction, and the classical sources by
\(\hbar c_aJ\,D_g\mathcal O[F]\). Its ordinary Jacobian has
\(\log\operatorname{Jac}=O(\hbar c_a)\), so its contribution to
\(-\hbar\log Z\) starts at order \(\hbar^2\). The already present
one-loop auxiliary density similarly changes only at the next loop order.
This observation is coefficientwise at each finite regulator; no uniform
control of its ultraviolet-enhanced two-loop trace is inferred. It does
not justify omitting the source contacts \eqref{eq:v38-observable}.
The exact finite density/Jacobian rule of V1 remains the rule at higher
orders.

\section{Uniformity of the local lift and what it does not solve}
\label{sec:v38-bounds}

The native pointwise map in \eqref{eq:v38-native-F} satisfies
\begin{equation}
 \|F^{\mathcal K}(\mathsf V)\|_{\mathrm{HS}}
   \le\frac32\|\mathsf V\|_{\mathrm{HS}},\qquad
 \|\mathcal R_{F^{\mathcal K}}\|_{\mathrm{HS}}
   \le6\|\partial c\|_{\mathrm{op}}\|\mathsf V\|_{\mathrm{HS}}.
 \label{eq:v38-pointwise}
\end{equation}
To verify the first bound, decompose \(\mathsf V\) into its
\(t,s/\sqrt3\), time--space, and traceless-spatial parts. The
squared Frobenius norm of the first two-dimensional block is \(85/64\),
and the other multipliers are \(7/8\) and \(35/32\). All are
smaller than \(3/2\). Formula \eqref{eq:v38-source-explicit} proves
the second bound. For the actual Gram,
\(\|\mathsf V\|_{\mathrm{HS}}\le\sum_A|\partial\phi_A|^2\).
The coframe factor is additionally bounded by \((1-r)^{-1}\).

For every fixed arity \(N\), coefficient derivative order, and smooth
external window \(\mu\), the local lift and its cotangent conversion
therefore have fixed finite bounds. The new action symbols have four
spacetime derivatives; the metric-antifield symbols have three. Their
windowed rooted kernels obey, respectively,
\begin{equation}
 \|\mathcal K_{C_{a,\log}}^{\mathrm{action}}\|_{\mathrm{rooted},b}
 \le C_{N,b,\chi}\log(1/(a\lambda))\mu^4,
 \qquad
 \|\mathcal K_{C_{a,\log}}^{\mathrm{source}}\|_{\mathrm{rooted},b}
 \le C_{N,b,\chi}\log(1/(a\lambda))\mu^3.
 \label{eq:v38-rooted}
\end{equation}
Here the powers of \(\chi\) are included in the explicitly fixed
component-dependent constant as explained above. For the flat
four-scalar action the factor is \(\chi^{-2}\).
The kernel proof multiplies the finite polynomial symbols by the declared
smooth windows, differentiates more than \(b+4(N-1)\) times, integrates
by parts, and periodizes. It has one source in \(L^1\) and the others
in \(L^\infty\), with no total-volume factor. The logarithm is an
actual local divergence, not an error that is claimed uniformly small.

On each fixed soft source panel the unprojected finite breaking of the
new packet is bounded by the native restored Euler discrepancy times
these logarithmic coefficient bounds. Thus, for the massless action rows
in the retained bank, a conservative bound is
\begin{equation}
 \|(s_a-s_\infty)C_{a,\log}\|_{r,\mu}
 \le C_{r,\chi}a^3\mu^8\log(1/(a\lambda))
 \quad(a\mu\ \hbox{within the inherited plateau}).
 \label{eq:v38-finite-error}
\end{equation}
The free-row contribution is actually of order \(a^4\mu^9\).
The source coefficient has three derivatives and the nonlinear
Euler error at least five, giving the conservative exponent eight.
This is a fixed-window estimate, not an all-mode symmetry theorem or an
arbitrary-loop remainder estimate.

The \emph{nonlocal} four-scalar remainder of V37 is unchanged by
\eqref{eq:v38-log-counterterm}: its new flat restriction is a polynomial
of degree four, annihilated by \(I-\mathsf T_4\). Its already proved
\(O(a^2)\) cutoff and mixed-shell estimates therefore remain in force.
The paired geometric, mixed and antifield local terms are not
annihilated from their own source state; they are stored together.

There is no inference that the physical four-scalar amplitude is zero.
Removing the displayed flat contact trades it for curvature, source and
observable terms in a redundant critical direction. The full functional
rather than its flat restriction is the relevant comparison object.
Likewise \(\mathcal C_{F,2}\) is homogeneous, so adding it cannot solve
\(Dh_0=\xi_0\) or change the inhomogeneous class. The lift of the fixed
V25 gravitational primitive, the native populated relative coefficient
system, and the remaining source-complete critical primitive are still
unproved. This concrete completion must not be substituted for those
statements.

\section{Verification and closing ledgers for V38}
\label{sec:v38-ledgers}

The accompanying checker independently forms the native scalar stress
from \(\sqrt g\,g^{-1}\), builds its full 55-by-100 coefficient map,
checks the 55 polynomial basis lifts, evaluates the 5-by-6 restricted
matrix and its right inverse, and verifies the exact native tensor,
mass contact, curvature symbols, and nonzero rotation-source contact.
It also tests the scalar-dependent canonical variation on a finite
first-jet algebra, and the cotangent series through the retained arity.
It checks no new native loop integral or large relative BV matrix.
Analytic shell bounds use the inherited estimates and the proof above;
the symbolic checks are not substituted for those estimates.

\begingroup
\renewcommand{\refname}{Additional source for V38}
\begin{thebibliography}{99}
\bibitem[48]{V38Criado}
J.~C.~Criado and M.~P\'erez-Victoria,
\emph{Field redefinitions in effective theories at higher orders},
JHEP \textbf{03} (2019), 038, arXiv:1811.09413,
doi:10.1007/JHEP03(2019)038.
Established field-redefinition and matching framework; it does not supply
the native stress tensor coefficients evaluated here.
\end{thebibliography}
\endgroup

\subsection*{Proved}
\addcontentsline{toc}{subsection}{V38: Proved}
The completed finite local fourth jet has the inherited universal
logarithm with the bounded finite remainder \eqref{eq:v38-local-log}.
The native stress map has the exact 55-dimensional image, 45-dimensional
kernel, canonical norm-two Hilbert-space lift, and the evaluated
coefficient map \eqref{eq:v38-R}. The tensor
\eqref{eq:v38-native-F} gives the explicit ordinary closed packet
\eqref{eq:v38-full-completion}, its mass-polynomial extension, and the
curvature and physical-source contacts. Its recorded one-loop logarithmic
subtraction preserves the fixed pure-gravity and subcritical
normalizations and the retained local identities. Finite-window and
rooted bounds have the scope stated above.

\subsection*{Inherited}
\addcontentsline{toc}{subsection}{V38: Inherited}
V1--V37 and all labels are unchanged. The actual simultaneous determinant,
computed four-gradient polynomial, improved symbols, local source
conventions, ordinary BV differential and its arity filtration,
nonlocal four-scalar theorem, and measured coordinate rule are inherited.
No new evaluation of the logarithmic angular contraction is claimed.
V34 remains the active subcritical scalar prescription; the fixed gravity
primitive is not modified or declared extendible.

\subsection*{Literature input}
\addcontentsline{toc}{subsection}{V38: Literature input}
Local BV-exact deformations and field redefinitions are established in
\cite{V29BBHMatter,V38Criado}. Their physical interpretation requires
source and measure transport; no general first construction is claimed.
No external cohomological existence theorem is used to settle the
uncomputed native inhomogeneous problem.

\subsection*{Conditional}
\addcontentsline{toc}{subsection}{V38: Conditional}
The 55-coefficient compiler applies to a fully supplied local fourth jet,
but its unspecified finite coefficients have not been evaluated. It can
be combined with any already constructed particular critical primitive
as a homogeneous normalization change. This does not establish that such
a complete particular primitive agreeing with the fixed gravity choice
has been constructed. Higher-loop application must retain the transformed
measure, source contacts and nonlinear coordinate terms.

\subsection*{Open}
\addcontentsline{toc}{subsection}{V38: Open}
The native critical equations \(Dh_0=\xi_0\),
\(Dh_1=\xi_1-Mh_0\), and the full fixed-gravity relative extension
remain open. The new exact packet removes one \emph{evaluated local
logarithmic coefficient} together with its partners; it does not remove
the outstanding ghost-number-one reference class. Internal gauge/chiral
sectors, unbounded source order, higher-loop closure, removal of the lower
reference scale and an invariant analytic ESM class remain separate.

\begin{center}\small
\begin{tabular}{>{\raggedright\arraybackslash}p{.23\textwidth}>{\raggedright\arraybackslash}p{.68\textwidth}}
\toprule
Variable & Scope in V38\\\midrule
Loop and fields & One-loop massless four-scalar logarithm on the existing
completed metric--scalar reference. The added local packet includes its
mixed metric and antifield partners.\\
Mass & Algebraic completion polynomial in \(z\); no new massive logarithm
is asserted. The cutoff logarithm theorem is at \(m=0\).\\
Cutoff and volume & Bounded finite fourth-jet remainder on V37's domain;
fixed-window rooted estimates have no extensive volume factor.\\
Norms & Exact coefficient-map norms refer to the displayed bases; local
counterterm size grows logarithmically. No all-order smallness claim.\\
Normalization & Pure gravity and V34 unchanged. The new critical
logarithmic packet is explicit; finite matching data retained.\\
Computation & Finite rational stress and source calculations performed.
Native critical reference integrals and the full relative matrix not
computed.\\\bottomrule
\end{tabular}
\end{center}

\clearpage
\part*{V39: The critical gravitational tail is removable in the relative problem}
\addcontentsline{toc}{part}{V39: Critical tail exactness and low-arity relative matching}

\section{The fixed gravitational prescription can be tested at lower arity}
\label{sec:v39-start}

Sections 1--303 are preserved. V38 completes a homogeneous logarithmic
counterterm; it does not solve the inhomogeneous equation
\(d_HH_2=x_2-tH_g\). V37 correctly identified the distinction between a
full reference germ and a particular finite gravitational primitive. We
now investigate that distinction for the \emph{native ordinary Einstein
complex}, rather than for a general filtered complex.

The result is a reduction of the fixed-normalization problem, not a new
loop coefficient. At critical weight four, a homogeneous gravitational
cocycle whose first nonzero arity is four or five has a local exact
primitive through the retained arity. Moreover, the scalar stress map
raises arity by at least one. Consequently changing the quartic and
quintic gravitational normalization, while keeping the lower coefficients
fixed, cannot create a new relative matter obstruction at arity five.
The compensating matter functional is explicit in terms of one
quartic, ghost-number-minus-one primitive. It belongs precisely to the
higher source families already constructed in V32.

This removes the need to match all gravitational arity-five coefficients
in an absolute-to-relative comparison. It does not prove that the actual
low-arity coefficients of V25's finite minimum-norm prescription agree
with those of a full primitive. Those low coefficients have not been
computed. No gravitational normalization is changed in the active theory:
the comparison theorem below returns a solution with the original
\(H_g\) if a comparison solution passing the smaller test is supplied.
V34's active subcritical scalar normalization and V38's complete
logarithmic packet are untouched.

The only external cohomological input used here is the local
\emph{linearized} spin-two reduction stated in the next section. The
four-derivative conclusion is not imported from a theorem that assumes
at most two derivatives. No full matter anomaly-vanishing theorem, no
finite-lattice nilpotence statement, and no native relative matrix rank
is assumed.

\section{A four-derivative gap in the free gravitational complex}
\label{sec:v39-free-gap}

Work with translation-invariant local polynomial densities, modulo total
ordinary derivatives, at the flat Euclidean coefficient reference.
The nonminimal pairs are restored and removed in the already declared
ordinary canonical convention. The linear tangent of the coframe map is
the symmetric metric variation. With \(\eta=q^*\) and \(\sigma=c^*\),
write
\[
 s_0=\delta_0+\gamma_0,\qquad
 \gamma_0q_{\mu\nu}=\partial_\mu c_\nu+\partial_\nu c_\mu,
 \qquad \delta_0\eta=\mathcal E_{g,2}(q),\quad
 \delta_0\sigma=R_0^*\eta.
\]
The sign of the last expression is the inherited adjoint-generator sign;
all subsequent identities use that same sign. Both pieces commute with
ordinary differentiation, and all other generator actions vanish. The
normalization \(\eta=\chi\widehat\eta\),
\(\sigma=\chi\widehat\sigma\) removes \(\chi\) from this finite
coefficient differential.

Arity \(A\) counts each field, ghost, or antifield once. The weights are
\[
 w(q)=0,\quad w(c)=-1,\quad w(\eta)=2,\quad w(\sigma)=3,
 \quad w(\partial)=1.
\]
They are V24's weights, not canonical particle dimensions. The operator
\(s_0\) preserves both arity and weight. Antifield number, in contrast,
assigns two to \(\sigma\), and must not be confused with arity or with
the number of antifield occurrences.

\subsection{The precise imported reduction}

We use the following free spin-two local-cohomology facts from
\cite[Sections 3--5 and Appendix A]{V39BDGH}, with the general
Koszul--Tate framework of \cite{V20BBH1994}:
\begin{enumerate}[label=(\roman*)]
\item The longitudinal cohomology is generated by the linearized Riemann
  tensor \(\mathscr R(q)\), its derivatives, the antifields and their
  derivatives, and \(c_\mu,\partial_{[\mu}c_{\nu]}\).
\item A ghost-number-zero cocycle modulo a divergence can be reduced to
  antifield number at most two. For a nontrivial highest antifield-two
  term, its invariant coefficients belong to
  \(H_2(\delta_0\mid\partial)\), represented by undifferentiated
  \(\sigma_\mu\) with constant coefficients. Terms quadratic in
  \(\eta\) do not give extra representatives.
\item After a trivial antifield-two term is removed, the antifield-one
  term may be made \(\gamma_0\)-closed. Its invariant representatives
  are linear in \(\eta\) and its derivatives and polynomial in
  \(\mathscr R\) and the surviving ghost jets.
\end{enumerate}
These algebraic reduction facts allow free tensor indices and do not
require rotation-invariant coefficient tensors. We do not invoke the
paper's Poincare-invariant two-derivative interaction classification.
Its general unresolved higher-derivative current classification is also
not needed: the elementary degree bound below excludes the relevant
current terms at the arities used here. The reductions can be taken in
the translation-invariant polynomial algebra. Since all their
differentials preserve homogeneous weight and arity, projecting any
resulting primitive to the input weight and arity preserves the identity.
These are literature inputs, not ranks computed by the new checker.

\begin{lemma}[High-arity longitudinal remainder has zero Euler derivative]
\label{lem:v39-euler-degree}
Let \(a_0(q)\) be a homogeneous local density of field degree
\(r\ge4\) and total derivative degree four. If
\(\gamma_0a_0\) is a divergence, then \(a_0\) is a divergence.
The assertion permits arbitrary constant component coefficients.
\end{lemma}
\begin{proof}
Set \(E^{\mu\nu}=\delta\int a_0/\delta q_{\mu\nu}\). The free
gauge displacement \(R_0c\) has no field dependence. Therefore Euler
differentiation commutes with its prolonged variation, giving
\(\gamma_0E=0\). Each coefficient of \(E\) is thus a polynomial
in \(\mathscr R(q)\) and its derivatives, by input (i). It has field
degree \(r-1\) and derivative degree four. A product of \(r-1\)
linearized curvatures has at least \(2(r-1)\ge6\) derivatives.
There can be no such coefficient, so \(E=0\).

For completeness, vanishing of this Euler derivative gives a divergence
without explicit coordinates. Repeated integration by parts in
\(\sum_{\mu\nu,I}\partial_Iq_{\mu\nu}
(\partial a_0/\partial(\partial_Iq_{\mu\nu}))\) gives
\(r a_0=q_{\mu\nu}E^{\mu\nu}+\partial_\mu j^\mu\).
The current is a finite polynomial in the same jets. Since \(r>0\)
and \(E=0\), divide by the integer \(r\). This is not an inverse
spacetime differential.
\end{proof}

\begin{theorem}[Native free critical tail exactness]
\label{thm:v39-free-gap}
Let \(k_r\) have ghost number zero, weight four, and homogeneous
arity \(r\ge4\) in the native free gravitational complex. If
\(s_0 k_r=0\) modulo total derivatives, then
\begin{equation}
 \boxed{k_r=s_0Q_r\quad\text{modulo total derivatives},
 \qquad \operatorname{gh}Q_r=-1,\quad w(Q_r)=4,\quad A(Q_r)=r.}
 \label{eq:v39-free-gap}
\end{equation}
The primitive is a translation-invariant finite differential polynomial.
In particular this applies at arities four and five. It is not a
statement about ghost-number-one cohomology or the filtered lattice
transformation.
\end{theorem}
\begin{proof}
Apply the preceding imported reduction and retain homogeneous arity and
weight at every step. A potentially nontrivial antifield-two
representative is \(\sigma\) times two surviving ghost jets.
After integration by parts its only ghost derivatives are curls, with
at most two derivatives in total. Its weight is therefore at most
\(3-2+2=3\), not four. Thus the antifield-two term is trivial and can
be removed. Antifield number greater than two has already been removed
by input (ii).

The remaining antifield-one representative contains one \(\eta\),
one ghost, and, at arity \(r\), exactly \(r-2\) linearized curvature
factors. Derivatives of \(\eta\) can be transferred by parts; if
this produces higher ghost derivatives their classes are trivial by
input (i). Even before those additional derivatives are counted, the
weight is at least
\(2+2(r-2)-1=2r-3\ge5\).
There is no weight-four representative of this type. Hence the
antifield-one term is also trivial. This argument does not classify
all higher-derivative rigid symmetries; it excludes their occurrence at
the present weight and arity.

The remaining antifield-free functional is longitudinally closed modulo
a divergence, has \(r\) metric fields and four derivatives, and
vanishes modulo a divergence by \cref{lem:v39-euler-degree}. Combining
the local trivializing changes proves \eqref{eq:v39-free-gap}. Projection
to ghost number \(-1\), weight four and arity \(r\) gives the
asserted primitive. No momentum or mass division appears.
\end{proof}

\section{The nonlinear quartic and quintic tail is an exact finite jet}
\label{sec:v39-nonlinear-tail}

Let \(s_g\) be the full ordinary gravitational differential. Its
homogeneous expansion at the flat coframe is
\[
 s_g=s_0+s_1+s_2+\cdots,
 \qquad A(s_jF)=A(F)+j.
\]
The operators are supplied by the native stationary action and the
coframe-coordinate source recursion; this is not a replacement by a
new Einstein action. They satisfy
\(s_0^2=0\), \(s_0s_1+s_1s_0=0\), and the higher coefficient
identities. All assertions below are identities of integrated local
functionals; raw coefficient implementations must retain current columns
or work in a fixed integration-by-parts quotient.

\begin{theorem}[Exact tail through arity five]
\label{thm:v39-nonlinear-tail}
Suppose \(k=k_4+k_5\), with each part of ghost number zero and weight
four, and
\(\operatorname{pr}_{A\le5}s_gk=0\). There are local polynomials
\(Q_4,Q_5\) of ghost number minus one and the respective arities such
that
\begin{equation}
 \boxed{k=\operatorname{pr}_{A\le5}s_g(Q_4+Q_5).}
 \label{eq:v39-tail}
\end{equation}
They may be obtained in the ordered steps
\begin{equation}
 s_0Q_4=k_4,\qquad
 s_0Q_5=k_5-s_1Q_4.
 \label{eq:v39-recursion}
\end{equation}
Thus this particular finite tail is the truncation of a full closed
ordinary local germ, namely \(s_g(Q_4+Q_5)\). No claim that a specified
low-arity primitive extends follows from this theorem alone.
\end{theorem}
\begin{proof}
The arity-four closure equation is \(s_0k_4=0\); apply
\cref{thm:v39-free-gap}. At arity five, closure gives
\(s_0k_5+s_1k_4=0\). Hence
\[
 s_0(k_5-s_1Q_4)
 =-s_1k_4+s_1s_0Q_4=0.
\]
Apply the same theorem at arity five. All remaining terms in
\(s_g(Q_4+Q_5)\) have arity at least six. Its full ordinary square
is zero, proving the stated lift. The full germ is used only as a local
identity; higher-arity vertices are not inserted into the finite
regulated action by this argument.
\end{proof}

In particular, V37's arbitrary filtered-complex counterexample cannot be
used to infer a native obstruction carried exclusively by this tail:
its arity-four kernel does not satisfy the free exactness theorem just
proved. The counterexample remains valid as a warning against an
\emph{unqualified} truncation argument.

\subsection{A fixed, finite norm prescription}

The exactness result is not a numerical singular-value assertion. A
refinement-independent bound can nevertheless be specified explicitly.
Choose rational bases in the finite integrated spaces at one fixed
arity \(r\), ghost numbers \(-1,0,1\), and weight four. Let
\(A_r\) and \(B_r\) be the two \(s_0\) matrices. The native free
operator has rational coefficients after the above normalization.
\Cref{thm:v39-free-gap} gives
\(\ker B_r=\operatorname{im}A_r\). Consequently
\[
 L_r=A_rA_r^T+B_r^TB_r>0,
 \quad H_r=A_r^TL_r^{-1},\quad Q_r^{\mathrm{err}}=L_r^{-1}B_r^T
\]
satisfy
\begin{equation}
 \boxed{I=A_rH_r+Q_r^{\mathrm{err}}B_r.}
 \label{eq:v39-hodge}
\end{equation}
Indeed the two orthogonal ranges \(\operatorname{im}A_r\) and
\(\operatorname{im}B_r^T\) exhaust the middle space, and \(L_r\)
preserves each. This proves \eqref{eq:v39-hodge} directly. In
particular \(H_r\) is a right section on the cocycles; for a nonclosed
input its error is displayed, not silently erased.

For an arithmetic bound set \(Z_r=(A_r^T,B_r)^T\), let its shape be
\(m_r\) by \(n_r\), choose a common denominator \(D_r\), and let
\(J_r\ge1\) bound the integer entries of \(D_rZ_r\). Put
\(L=\max\{1,\sqrt{m_rn_r}J_r\}\). A nonzero full-column minor
of \(D_rZ_r\), followed by Cauchy--Binet, gives
\begin{equation}
 \|L_r^{-1}\|_2\le D_r^2L^{2(n_r-1)},\qquad
 \|H_r\|_2\le\|A_r\|_2D_r^2L^{2(n_r-1)}.
 \label{eq:v39-bound}
\end{equation}
This uses the same elementary matrix-height principle as V25, not a new
analytic theorem. The new input is the proved absence of a middle
cohomology class in these native homogeneous spaces. Fixed basis
conversion then bounds the chosen raw local representative. The native
matrices are not assembled or inverted here; these deliberately
conservative constants are specified rather than numerically estimated.
There is no dependence on lattice size or number of shells, and no
claim of uniformity as derivative order or arity grows.

\section{The scalar compensation starts at arity five}
\label{sec:v39-stress-arity}

Return to the same ordinary gravity--scalar complex with four real scalar
components, internal gauge and Yukawa couplings zero, and \(z=m^2\).
On the matter-count-zero/two quotient write
\[
 s(F_g,F_2)=(s_gF_g,\ t(z)F_g+d_H(z)F_2),
 \quad d_H(z)=D+zM,\quad t(z)=t_0+zt_1.
\]
These are the actual triangular operators in V33--V37 and satisfy
\(d_Ht+ts_g=0\).

\begin{lemma}[Native arity gain of the stress insertion]
\label{lem:v39-arity-gain}
The map \(t\) raises total arity by at least one. If \(Q_4,Q_5\)
are as above, then
\begin{equation}
 \boxed{\operatorname{pr}_{A\le5}t(Q_4+Q_5)=t_{[1]}(z)Q_4,}
 \label{eq:v39-leading-t}
\end{equation}
where \(t_{[1]}\) is the leading, one-unit arity-raising part of \(t\).
Its output has matter count two, ghost number zero, and arity five.
\end{lemma}
\begin{proof}
An antifield-free gravitational density has zero scalar-stress
insertion. Acting on a metric antifield, \(t\) replaces one argument
by the scalar stress, whose first coefficient has two scalar arguments.
Acting on a ghost antifield it replaces one argument by the scalar
current \(\sum_A\upsilon_A\partial_\mu\phi_A\), again two
arguments. Derivatives commute with these replacements. Every further
coframe coefficient adds metric arguments. Thus the least arity gain is
one, proving the formula. The statement is independent of whether the
mass is zero; \(z\phi^2\) has the same field count.
\end{proof}

For the integrated scalar action
\(\frac12\int\sqrt g\sum_A(g^{\mu\nu}\partial_\mu\phi_A
\partial_\nu\phi_A+z\phi_A^2)\), the actual tangent stress is
\begin{equation}
 \mathcal E_{H,0}^{\mu\nu}
 =\frac14\delta^{\mu\nu}\sum_A
        (|\partial\phi_A|^2+z\phi_A^2)
  -\frac12\sum_A\partial_\mu\phi_A\partial_\nu\phi_A.
 \label{eq:v39-flat-stress}
\end{equation}
In independent off-diagonal metric coordinates its dual coefficient is
twice the corresponding displayed tensor entry. With the inherited
signed adjoint convention, denote the scalar current replacing
\(\sigma_\mu\) by \(\mathcal J_{H,0}^\mu\); its entries are the
correspondingly signed \(\sum_A\upsilon_A\partial_\mu\phi_A\).
Thus
\(t_{[1]}\eta=\mathcal E_{H,0}\),
\(t_{[1]}\sigma=\mathcal J_{H,0}\). In the normalized hatted
antifields both replacements have the additional factor \(\chi^{-1}\).
The following formulas use the unhatted convention for readability.

\begin{proposition}[All required compensation families are already measured]
\label{prop:v39-slots}
The possible structural types of \(Q_4\), with total derivatives
understood, and their leading stress images are
\begin{center}\small
\begin{tabular}{lll}
\toprule
Generator type & Derivative degree & Matter image types\\\midrule
\(\eta q^3\) & 2 & \(q^3\phi^2\)\\
\(\sigma q^2c\) & 2 & \(\upsilon q^2\phi c\)\\
\(\eta^2qc\) & 1 & \(\eta q\phi^2c\)\\
\(\eta\sigma cc\) & 1 & \(\sigma\phi^2cc\), \(\eta\upsilon\phi cc\)\\
\bottomrule
\end{tabular}
\end{center}
No new arity-five matter family is required. These are the coupled source
families already supplied in V32. In particular \(t_{[1]}Q_4\)
changes no matter coefficient of arity at most four.
\end{proposition}
\begin{proof}
If \(r,s\) count metric and ghost antifields in \(Q_4\), its ghost
count is \(r+2s-1\), so its metric count is
\(5-2r-3s\). Nonnegative counts give exactly the four rows above.
At ghost number minus one, weight four requires derivative degree
\(3-r-s\), giving the stated two or one derivatives. The
replacement rule proves the output types and keeps every graded sign.
\end{proof}

For example, the local generator
\(Q_4=\int\eta^{00}q_{00}(\partial_0q_{00})^2\) has the measured
compensation
\begin{equation}
 t_{[1]}Q_4
 =\frac14\int q_{00}(\partial_0q_{00})^2
   \sum_A\left[-(\partial_0\phi_A)^2+
        \sum_i(\partial_i\phi_A)^2+z\phi_A^2\right].
 \label{eq:v39-example-one}
\end{equation}
It contains both its massless and first-mass terms and vanishes on the
flat scalar two-point restriction. A second ordered generator gives
\begin{equation}
 \begin{split}
 t_{[1]}(\eta^{00}\sigma_0c^1\partial_0c^2)
 ={}&\mathcal E_{H,0}^{00}\sigma_0c^1\partial_0c^2\\
    &-\eta^{00}\mathcal J_{H,0}^{0}c^1\partial_0c^2.
 \end{split}
 \label{eq:v39-example-two}
\end{equation}
The minus sign is from the odd metric antifield. Keeping only the first
output would not be a stress insertion of the generator.

In the raw normalized jet \(\ell^1\) norm, an explicit safe estimate is
\begin{equation}
 \boxed{\|t_{[1]}Q_4\|_1\le40\chi^{-1}\|Q_4\|_1.}
 \label{eq:v39-stress-bound}
\end{equation}
To see this, a flat stress component has coefficient sum at most five
for four real scalars, including the mass term and the independent
symmetric-coordinate normalization; a current component has sum at most
four. Distributing at most two derivatives over a bilinear coefficient
costs at most \(2^2\). There are at most two antifields to replace.
The triangle inequality gives forty. This is a deliberately safe raw
coefficient bound; quotient norms require their fixed representative
maps as in \eqref{eq:v39-bound}.

\section{A low-arity equivalence for the actual fixed-gravity problem}
\label{sec:v39-relative}

All equations in this section are in the arity-five local integrated
complex at weight four. Denote its two-matter reference target by \(x_2\)
and its pure gravitational target by \(x_g\). The value \(H_g\)
is the fixed critical prescription of V25, not a new minimizer.

\begin{theorem}[Quartic/quintic normalization cannot change relative solvability]
\label{thm:v39-relative}
Let \(H_g\) and \(\widetilde H_g\) solve the same gravitational
reference equation, and suppose
\(\operatorname{pr}_{A\le3}(H_g-\widetilde H_g)=0\).
Put \(k=H_g-\widetilde H_g\), and choose \(Q_4,Q_5\) from
\cref{thm:v39-nonlinear-tail}. Then
\begin{equation}
 \boxed{
 d_H\widetilde H_2=x_2-t\widetilde H_g
 \quad\Longleftrightarrow\quad
 d_HH_2=x_2-tH_g,
 \quad H_2=\widetilde H_2+t_{[1]}Q_4.}
 \label{eq:v39-relative-equivalence}
\end{equation}
The map is reversible after the same generator choice. It changes only
arity-five matter coefficients, and changes no pure-gravity coefficient.
It requires no full-germ lift of either individual gravitational
primitive.
\end{theorem}
\begin{proof}
The difference is closed and starts at arity four, so
\(k=s_gQ\) in the quotient, with \(Q=Q_4+Q_5\).
The triangular identity gives \(tk=-d_HtQ\). Therefore
\[
 d_H(\widetilde H_2+tQ)
 =x_2-t\widetilde H_g-tk=x_2-tH_g.
\]
Use \cref{lem:v39-arity-gain} to replace \(tQ\) by
\(t_{[1]}Q_4\). The converse subtracts the same term. All identities
are local and polynomial in mass. The full exact packet
\((s_gQ,tQ)\) explains why both members, not just the scalar action
part, must be retained.
\end{proof}

More generally it is enough that the low-arity difference is exact:
\[
 \operatorname{pr}_{A\le3}(H_g-\widetilde H_g)
 =\operatorname{pr}_{A\le3}s_gQ_{\le3}.
\]
Subtract that full ordinary exact change first and apply the theorem to
the remaining tail. The resulting matter correction includes
\(tQ_{\le3}\), which may change lower-arity matter coefficients;
those changes must be matched to any prescribed lower normalization.
The strict-equality version \eqref{eq:v39-relative-equivalence} avoids
this additional lower-row change.

\begin{corollary}[Reduced boundary data for the native range problem]
\label{cor:v39-reduced-problem}
The fixed-gravity critical equation has a solution if and only if the
simultaneous local system has a solution \((\widetilde H_g,\widetilde H_2)\)
with
\begin{equation}
 \boxed{
 s_g\widetilde H_g=x_g,\quad
 t\widetilde H_g+d_H\widetilde H_2=x_2,\quad
 \operatorname{pr}_{A\le3}\widetilde H_g
 =\operatorname{pr}_{A\le3}H_g.}
 \label{eq:v39-reduced-problem}
\end{equation}
Equality constraints on the gravitational quartic and quintic coefficients
are unnecessary for this existence test. Once a solution of
\eqref{eq:v39-reduced-problem} is found, the theorem restores those
coefficients to their fixed values while compensating in the existing
matter bank.
\end{corollary}
\begin{proof}
A solution with the fully fixed gravity choice is a solution of the
reduced system. Conversely apply \cref{thm:v39-relative} to a solution
of the reduced system. This proves an equivalence, not the assertion that
the reduced system is already populated by a solution.
\end{proof}

For orientation, the inherited raw ghost-zero weight-four density counts
by arity one through five are
\[
 350,\quad21\,350,\quad388\,710,\quad
 3\,698\,400,\quad23\,127\,870.
\]
The lower-arity equality data therefore have at most \(410\,410\)
raw coordinates rather than \(27\,236\,680\). These are ambient
counts before integration by parts, not ranks or physical couplings.
This reduction removes high-arity \emph{normalization constraints}; it
does not delete the equations or unknowns needed to solve the simultaneous
system. No reduced native matrix is evaluated here.

\section{The mass lift and measured bookkeeping are preserved}
\label{sec:v39-mass-measure}

The compensating map \(t_{[1]}(z)\) is affine in \(z\), while the
pure-gravity generator \(Q_4\) has no mass parameter. Consequently the
three coefficients in V35's critical candidate transform as
\begin{equation}
 \boxed{
 h_0=\widetilde h_0+t_{0,[1]}Q_4,\qquad
 h_1=\widetilde h_1+t_{1,[1]}Q_4,\qquad
 h_2=\widetilde h_2.}
 \label{eq:v39-mass-lift}
\end{equation}
The same \(Q_4\) occurs in both of the first two changes. Comparing
coefficients in
\((D+zM)t(z)Q=-t(z)s_gQ\) verifies the coupled V35 equations.
There is no independent choice of a first-mass normalization, no inverse
mass, and no new highest-mass potential obstruction caused by the high
pure-gravity tail. The unresolved low-arity matching can still affect
all stages of the populated relative problem.

The comparison restores the originally fixed gravitational coefficient.
It is therefore not an instruction to update the active \(H_g\).
If it is implemented after a comparison solution is constructed, the
order-\(\hbar\) change is the complete ordinary exact packet. Its
pure-gravity part restores \(H_g\), and its matter part is
\eqref{eq:v39-leading-t}. Ordinary source and composite-observable
contacts are differentiated from that packet, exactly as in V38.
The tree action, hard covariances and auxiliary stationary point do not
change at this loop order. A field-change implementation requires the
same cotangent/source transport; the Jacobian of an order-\(\hbar\)
change belongs to the subsequent action coefficient and is not counted
again as a new one-loop determinant.

The matter compensation has two derivatives plus either two scalar
stress derivatives or the mass factor, or one derivative plus a scalar
current. It belongs to the existing critical weighted local bank. Fixed
windowing and integration by parts therefore give a finite multilinear
source bound with one argument in \(L^1\) and the others in
\(L^\infty\), without a total-volume factor. Combining
\eqref{eq:v39-stress-bound} with a chosen finite section gives, in
normalized coefficient representatives,
\[
 \|H_2-\widetilde H_2\|_1
 \le40\chi^{-1}\|H_4\|_{1\to1}\|k_4\|_1,
\]
where \(H_4\) is the free section in \eqref{eq:v39-hodge}, followed
by its fixed raw-representative conversion. This is a finite-order bound
on a local comparison, not a claim that arbitrary cutoff-dependent
input coefficients stay small or that a many-shell trajectory has been
constructed.

\section{The next native computation is the low-arity boundary, not another tail}
\label{sec:v39-next}

V37's general relative image criterion remains correct. V39 determines
which part of the fixed gravitational boundary can carry the obstruction
at this retained order: it factors through its arity-at-most-three
restriction, with exact changes treated by their scalar stress lifts.
In particular, an obstruction cannot be attributed solely to an
arbitrary quartic or quintic homogeneous gravity choice.

This is not a classification of the low-arity cocycles. Bilinear
four-derivative metric terms and cubic source deformations can carry
nontrivial free classes or nonlinear compatibility constraints. Their
actual coefficients in V25's minimum-norm primitive have not been
computed. No claim that the low restriction extends to a full germ is
made. The existence of a full reference germ or the absence of a full
unconstrained anomaly does not alone impose those lower normalization
values on a chosen primitive.

A direct next task is thus \eqref{eq:v39-reduced-problem}, or a
construction of a full primitive with those specific lower coefficients.
A successful comparison then needs only the one local quartic section to
recover the fully fixed normalization at matter arity five. The massless
and first-mass equations remain coupled as in \eqref{eq:v39-mass-lift}.
V38's four-scalar logarithmic packet is independent of this reduction and
cannot substitute for the missing low-arity inhomogeneous solution.

\section{Verification and closing ledgers for V39}
\label{sec:v39-ledgers}

The self-contained checker verifies the native free metric/gauge symbol
identity, linearized curvature invariance, the finite arity/weight bounds,
raw density counts, the complete quartic-generator source-type table,
the signs of the two displayed stress examples, and the finite ordered
tail and mass compensation identities. It tests the finite Hodge-section
formula on separately declared rational complexes. These tests do not
compute the native cohomology by rank: the reduction inputs in
\cref{sec:v39-free-gap} are imported, and the elementary degree proof is
given above. Neither a native reference integral, V25's low coefficient
array, nor the full or reduced relative coefficient matrix is evaluated.

\begingroup
\renewcommand{\refname}{Additional source for V39}
\begin{thebibliography}{99}
\bibitem[49]{V39BDGH}
N.~Boulanger, T.~Damour, L.~Gualtieri and M.~Henneaux,
\emph{Inconsistency of interacting, multi-graviton theories},
Nuclear Physics B \textbf{597} (2001), 127--171,
arXiv:hep-th/0007220, doi:10.1016/S0550-3213(00)00718-5.
Sections 3--5 and Appendix A supply the free longitudinal and invariant
Koszul--Tate reductions. The two-derivative interaction classification
and its Poincare restrictions are not used as a four-derivative theorem.
\end{thebibliography}
\endgroup

\subsection*{Proved}
\addcontentsline{toc}{subsection}{V39: Proved}
The native ordinary free critical cohomology has no ghost-zero class at
arity four or higher, using the identified linearized reduction input
and the explicit degree/Euler proof. Closed quartic/quintic gravity tails
are exact through arity five. Their matter stress correction is exactly
\(t_{[1]}Q_4\), has the five existing source types, obeys the declared
raw coefficient bound, and preserves all lower-arity matter terms. The
fixed-gravity relative range problem is equivalent to matching gravity
only through arity three. Its massless and first-mass changes are coupled
and its highest-mass coefficient is unchanged by this tail correction.

\subsection*{Inherited}
\addcontentsline{toc}{subsection}{V39: Inherited}
V1--V38 and their labels are unchanged. The ordinary coframe/master
coefficients, fixed V25 primitive, V34 active scalar prescription,
V35 mass bicomplex, V37 relative identity, V32 matter-source bank, and
V38 source-complete logarithmic packet retain their original status.
The matrix-height estimate and source-window conversion are inherited
methods applied to the newly proved free-image statement.

\subsection*{Literature input}
\addcontentsline{toc}{subsection}{V39: Literature input}
Free spin-two longitudinal and invariant Koszul--Tate reductions are
those of \cite{V39BDGH,V20BBH1994}. The normal forms, not an unrestricted
matter anomaly-vanishing assertion, are used. No Lorentz-invariant
coefficient ansatz is imposed on the native finite component arrays.
No general higher-derivative deformation classification is claimed.

\subsection*{Conditional}
\addcontentsline{toc}{subsection}{V39: Conditional}
A comparison primitive passing the actual arity-three gravitational
boundary test can be converted to the fully fixed prescription by the
proved local map. Such a comparison has not been produced here. The
finite Hodge sections have explicit arithmetic bounds but their native
coefficient matrices have not been built. Fixed-order bounds transfer
actual input coefficient bounds; they do not establish analytic RG
invariance or uniformity in unbounded arity.

\subsection*{Open}
\addcontentsline{toc}{subsection}{V39: Open}
The populated low-arity boundary of the fixed gravitational critical
primitive and the reduced simultaneous equations
\eqref{eq:v39-reduced-problem} remain open. Therefore the native
\(Dh_0=\xi_0\), \(Dh_1=\xi_1-Mh_0\) critical system is not
claimed solved. The unresolved part is no longer an arbitrary high-arity
gravity tail. Internal gauge/chiral interfaces, higher-loop closure,
removal of the lower reference, and an invariant analytic full ESM class
remain outside the new result.

\begin{center}\small
\begin{tabular}{>{\raggedright\arraybackslash}p{.23\textwidth}>{\raggedright\arraybackslash}p{.68\textwidth}}
\toprule
Variable & Scope in V39\\\midrule
Loop and order & One-loop local normalization comparison; ordinary critical
weight four, gravity and matter functionals retained through arity five.\\
Mass & Matter correction affine in \(z=m^2\), regular at zero; no new
massive loop integral.\\
Cutoff and volume & Algebraic comparison independent of mesh and periods;
fixed-window local kernels have no extensive volume factor. No new
cutoff convergence rate for an unresolved coefficient.\\
Source order & Exactly the displayed retained bank. Arithmetic section
bounds are fixed-order; no uniformity as the bank grows.\\
Normalization & Fixed gravity, V34 and V38 remain active. The theorem
restores, rather than resets, fixed gravity after a successful comparison.\\
Verification & Symbol and finite-coefficient checks performed; free
cohomological reduction identified as literature input; native lower
normalization and relative matrices not evaluated.\\\bottomrule
\end{tabular}
\end{center}

\clearpage
\part*{V40: A computed native bilinear boundary for critical relative matching}
\addcontentsline{toc}{part}{V40: Native bilinear boundary and cubic lifting freedom}

\section{The next reduction is computed, rather than assumed}
\label{sec:v40-start}

Sections 1--311 are preserved. V39 removes a closed critical gravitational
normalization difference beginning at arity four. It leaves the actual
boundary through arity three uncomputed. The present installment resolves
one finite part of that boundary by constructing and certifying the
\emph{complete native free bilinear complex}. It does not replace the
inhomogeneous reference coefficient by an invariant curvature ansatz.

There are 2,725 integrated bilinear coefficient coordinates, 946 independent
local field/ghost-redefinition directions, and a 43-dimensional quotient of
closed coefficients by those directions. All 43 classes have explicit
curvature-square representatives. A rational section is computed, not only
defined by a pseudoinverse. Its coefficient norm and its response to a
nonclosed input are certified. The ambiguity in a quadratic primitive is
also computed: modulo a free exact generator, only nine parameters can
influence the subsequent cubic equation. These are two different
cohomological degrees and must not be conflated.

The native free metric Euler operator and all ten metric components are
retained. No TT reduction, gauge-fixed inverse, inverse soft momentum,
Lorentz-invariant coefficient restriction, or numerical rank threshold is
used. The calculation concerns the ordinary zero-mesh local complex; it is
not a nilpotence theorem for the filtered finite source. The free spin-two
framework is established \cite{V39BDGH,V20BBH1994}, but the ranks and sections
below are certified directly from the displayed finite matrices and do not
invoke a cohomology classification as a missing matrix-rank input.

The active fixed gravitational prescription, the V34 subcritical scalar
normalization, and the complete V38 logarithmic packet are not changed. In
particular no value of V25's unevaluated critical reference primitive is
invented. The theorem applies to the difference of two primitives only after
they are specified in the same conventions. Its nonclosed extension also
retains the actual inhomogeneous free row.

\section{The complete integrated bilinear polynomial complex}
\label{sec:v40-complex}

Work over real or rational coefficients, with complexification allowed for
polynomial tests. Integration by parts is imposed exactly by using the one
independent Fourier momentum \(p=(p_0,p_1,p_2,p_3)\). The ten independent
symmetric coordinates are ordered
\[
 (00,01,02,03,11,12,13,22,23,33).
\]
Write \(q\) both for their column and for the corresponding symmetric
\(4\times4\) matrix, according to the displayed operation. Define
\begin{equation}
 \begin{split}
 S_f(p,q)=\frac12\bigl\{
 &(p^Tp)\operatorname{tr}(q^2)-2p^Tq^2p
 +2\operatorname{tr}(q)p^Tqp\\
 &-(p^Tp)(\operatorname{tr}q)^2\bigr\},\qquad
 E(p)=\partial_q^2S_f(p,q),\\
 R_{(ij),\mu}(p)&=p_i\delta_{j\mu}+p_j\delta_{i\mu}.
 \end{split}
 \label{eq:v40-native-symbols}
\end{equation}
Here \(E\) is the independent-coordinate Euler matrix, not the gauge-fixed
matrix. In the conventions of \eqref{eq:v13-symbol-data} it is the
independent-coordinate matrix of \(X(p)\); the physical ordinary Euler
matrix is \(\chi E/2\). Thus a physical pair \((K_{\rm phys},T_{\rm phys})\)
is represented below by \((2K_{\rm phys}/\chi,T_{\rm phys})\). The odd Fourier
symbols have their common factor \(i\) removed. This fixes the phase and
normalization of the following coefficient maps; it changes no physical
coupling. All norm constants below refer to this normalized finite basis.

Direct expansion gives
\begin{equation}
 E^T=E,\qquad ER=0.
 \label{eq:v40-Noether}
\end{equation}
A ghost-number-zero critical bilinear coefficient consists of
\[
 \frac12q(-p)^TK(p)q(p)+\widehat\eta(-p)^TT(p)c(p),
\]
with the understood odd Fourier phase. The matrix \(K\) is symmetric and
homogeneous of momentum degree four; \(T\) is \(10\times4\), homogeneous of
degree three. Self-adjointness gives symmetry because the degree of \(K\)
is even. These symbols describe the full integrated bilinear bank, not just
its rotation-invariant part.

For \(\mathscr P_d\), the homogeneous degree-\(d\) polynomials in four
momenta, the relevant segment of the local free complex is
\begin{equation}
 \begin{gathered}
 \mathscr C^{-2}\xrightarrow{\partial_{-2}}
 \mathscr C^{-1}\xrightarrow{\partial_{-1}}
 \mathscr C^0\xrightarrow{\partial_0}\mathscr C^1,\\
 \mathscr C^{-2}=\operatorname{Alt}_{10}(\mathscr P_0)
                   \oplus\operatorname{Mat}_{4\times10}(\mathscr P_1),\\
 \mathscr C^{-1}=\operatorname{Mat}_{10}(\mathscr P_2)
                   \oplus\operatorname{Mat}_{4}(\mathscr P_2),\\
 \mathscr C^0=\operatorname{Sym}_{10}(\mathscr P_4)
                   \oplus\operatorname{Mat}_{10\times4}(\mathscr P_3),
 \qquad
 \mathscr C^1=\operatorname{Mat}_{10\times4}(\mathscr P_5).
 \label{eq:v40-spaces}
 \end{gathered}
\end{equation}
The differential need not end at \(\mathscr C^1\); only this segment is
required for the two cohomology groups computed here. The first space
represents \(\eta\eta\) and \(\sigma q\) generators, and the second represents
\(\eta A(\partial)q+\sigma B(\partial)c\). In the fixed integrated phase
convention the maps are
\begin{equation}
 \boxed{\begin{aligned}
 \partial_{-2}(C,D)&=(-CE+RD,\ DR),\\
 \partial_{-1}(A,B)&=(EA+A^TE,\ -AR+RB),\\
 \partial_0(K,T)&=KR+ET.
 \end{aligned}}
 \label{eq:v40-maps}
\end{equation}
The skew matrix \(C\) is not a covariance. The map \(D\) is a linear
momentum symbol, not a BRST differential.

\begin{proposition}[Exact finite encoding of the native free bilinear row]
\label{prop:v40-encoding}
Equations \eqref{eq:v40-spaces}--\eqref{eq:v40-maps} encode the complete
integrated, translation-invariant, arity-two, weight-four free gravitational
BV segment in the stated normalization. Their dimensions are
\begin{equation}
 \boxed{205\longrightarrow1160\longrightarrow2725\longrightarrow2240.}
 \label{eq:v40-dimensions}
\end{equation}
Both consecutive products vanish exactly. No free tensor indices are
restricted by rotational averaging.
\end{proposition}
\begin{proof}
At arity two and ghost number zero the only structural types are \(q^2\)
and \(\eta c\); weight four requires four and three derivatives,
respectively. At ghost number minus one they are \(\eta q\) and \(\sigma c\),
both with two derivatives. At ghost number minus two they are \(\eta^2\)
without derivatives and \(\sigma q\) with one derivative. Grassmann
antisymmetry makes the first coefficient skew. At ghost number one the row
is \(qc\) with five derivatives. Integration by parts moves all derivatives
to the single independent momentum. This gives
\[
 45+40\cdot4=205,\quad116\cdot10=1160,\quad
 55\cdot35+40\cdot20=2725,\quad40\cdot56=2240.
\]
The prolonged free transformations and their signed Euler replacements give
\eqref{eq:v40-maps}. These formulas can also be checked without any graded
notation: substitute \(ER=0\), \(E^T=E\), \(C^T=-C\) to obtain
\(\partial_0\partial_{-1}=0\) and
\(\partial_{-1}\partial_{-2}=0\). The raw-symbol compiler checks these
identities in every column. A critical arity-one integrated density has all
four derivatives on its single metric field and is a divergence, so it does
not provide an additional integrated boundary coordinate.
\end{proof}

\section{Certified native ranks, representatives, and a bounded section}
\label{sec:v40-ranks}

The following is a finite arithmetic proof. Its complete rational witness is
\srcid{ESM_bilinear_boundary_V40_certificate.json}; the accompanying standalone
script reconstructs \(E,R\) and all differential matrices from
\eqref{eq:v40-native-symbols}--\eqref{eq:v40-maps} before checking the witness.
It does not trust an asserted numerical rank.

For reproducibility, momentum monomials are ordered lexicographically by
their exponent tuple, with the first exponent increasing first. Matrix
components have the order already displayed, and symmetric or skew pairs
have increasing first and then second index. Coordinate reflections split
each matrix into eight independent parity blocks. The parity is the sum,
modulo two, of the momentum exponents and all tensor-index occurrences.
Only even-total parities occur. The exact block results are
\begin{center}\small
\begin{tabular}{lrrrrrrr}
\toprule
Parity & \(\dim C^{-2}\)&\(\dim C^{-1}\)&\(\dim C^0\)&\(\dim C^1\)&
\(\operatorname{rank}\partial_{-1}\)&\(\operatorname{rank}\partial_0\)&\(\dim H^0\)\\
\midrule
0000 &34&188&431&328&154&268&9\\
Each two-odd parity (six)&26&144&338&280&117&216&5\\
1111 &15&108&266&232&90&172&4\\
\midrule
Total &205&1160&2725&2240&946&1736&43\\
\bottomrule
\end{tabular}
\end{center}
The first differential is injective in every block.

\begin{theorem}[Computed free critical bilinear quotient]
\label{thm:v40-ranks}
For the native matrices above,
\begin{equation}
 \boxed{\operatorname{rank}\partial_{-2}=205,\quad
 \operatorname{rank}\partial_{-1}=946,\quad
 \operatorname{rank}\partial_0=1736.}
 \label{eq:v40-ranks}
\end{equation}
Consequently
\begin{equation}
 \boxed{\dim H^{-1}=9,\qquad \dim H^0=43.}
 \label{eq:v40-cohom}
\end{equation}
Every ghost-number-zero class has an antifield-free curvature-square
representative. The certificate gives fixed rational maps
\(\mathsf S:\mathscr C^0\to\mathscr C^{-1}\),
\(\Pi:\mathscr C^0\to\mathbb Q^{43}\),
\(\mathsf Z:\mathbb Q^{43}\to\mathscr C^0\), and
\(\mathsf Q:\mathscr C^1\to\mathscr C^0\) such that
\begin{equation}
 \boxed{
 I=\partial_{-1}\mathsf S+\mathsf Z\Pi+\mathsf Q\partial_0,
 \quad \Pi\partial_{-1}=0,\quad\partial_0\mathsf Z=0,
 \quad\Pi\mathsf Z=I_{43}.}
 \label{eq:v40-contraction}
\end{equation}
In the specified raw monomial coefficient \(\ell^1\) norms,
\begin{equation}
 \boxed{
 \|\mathsf S\|=\frac{145}{2},\qquad
 \|\Pi\|=16,\qquad \|\mathsf Z\|=22,\qquad
 \|\mathsf Q\|=\frac{191}{4}.}
 \label{eq:v40-norms}
\end{equation}
These are exact induced column norms for this section, not invariant norms
of the cohomology and not minimum norms over every possible section.
\end{theorem}
\begin{proof}
Here is the entire certificate protocol. Index curvature components by the
six ordered bivectors \(01,02,03,12,13,23\), retain the symmetric bivector
pairs, and omit \(((03),(12))\) using the first Bianchi identity. This gives
20 components. Use the integer normalization
\[
 \mathscr R_{ij,kl}
 =p_kp_jq_{il}+p_lp_iq_{jk}-p_lp_jq_{ik}-p_kp_iq_{jl},
\]
which is twice the conventional linearized curvature, up to the harmless
Fourier sign. Their 210 symmetric products supply columns with \(T=0\).
Their closure is checked as an exact polynomial identity, not inferred
from a floating-point evaluation.

In each parity block select the independent \(\partial_{-1}\) columns listed
in the certificate, and then the listed curvature-product columns. Append
standard coordinate columns whose \(\partial_0\) images are independent.
Call this square matrix \(\mathsf T\). The certificate contains its rational
inverse and verifies both products with the identity. It also verifies that
every \(\partial_{-1}\) column has zero coordinates outside the selected
image, that the curvature columns are closed, and that a listed square
row minor of \(\partial_0\) on the appended standard columns has the displayed
inverse. These identities give both lower and upper bounds for each stated
rank. In particular, they prove spanning, not merely independence of 43
examples.

The inverse coordinates of \(\mathsf T\) are partitioned into the exact,
curvature, and nonclosed parts. Send the first part to its selected
generator columns to define \(\mathsf S\), use the second for \(\Pi\), and
embed those columns by \(\mathsf Z\). The inverse of the independent boundary
row minor defines \(\mathsf Q\). The certificate verifies that the last
coordinates of \(\mathsf T^{-1}\) are that inverse applied to the actual
\(\partial_0\) row; this proves \eqref{eq:v40-contraction} on \emph{all}
inputs. Exact rational column sums give \eqref{eq:v40-norms}.

For ghost number minus one the witness similarly appends nine closed columns
to all 205 \(\partial_{-2}\) columns, then appends independent standard
columns for \(\partial_{-1}\). The rational inverse of this complete basis
is verified. Thus \(\partial_{-2}\) is injective and its image has codimension
nine in \(\ker\partial_{-1}\). This proves the second cohomology count and
supplies the nine explicit representatives used below. The arithmetic
witness is part of the proof; the script evaluates it over \(\mathbb Q\).
\end{proof}

For an arbitrary coefficient \(b=(K,T)\), not assumed closed, the compiler
therefore returns a primitive of its exact part, its 43 curvature coordinates,
and its actual free-identity error. In particular,
\begin{equation}
 \boxed{
 \|b-\partial_{-1}\mathsf Sb\|_1
 \le22\|\Pi b\|_1+\frac{191}{4}\|\partial_0b\|_1.}
 \label{eq:v40-error}
\end{equation}
For a closed input, \(\Pi b=0\) is necessary and sufficient for a local
bilinear primitive. For a nonclosed reference primitive, \(\Pi b\) alone is
not a cohomology class; the last term in \eqref{eq:v40-contraction} must
remain. No cutoff-dependent local coefficient is replaced by its bound.

\section{The residual directions are anisotropic, but cannot be averaged away}
\label{sec:v40-rotations}

The certificate also evaluates the action of all six infinitesimal ordinary
Euclidean rotations on the 43 quotient coordinates. For a skew matrix
\(\Omega\), let \(U_\Omega\) act on independent metric coordinates by
\(q\mapsto\Omega q-q\Omega\). The actual symbol action is
\begin{equation}
 \begin{aligned}
 \mathcal L_\Omega K&=(\Omega p)\cdot\partial_pK
                         +U_\Omega^TK+KU_\Omega,\\
 \mathcal L_\Omega T&=(\Omega p)\cdot\partial_pT
                         -U_\Omega T+T\Omega.
 \label{eq:v40-rotation-action}
 \end{aligned}
\end{equation}
It is the change of an integrated coefficient under the tensor and momentum
rotation, not a change of the fixed cutoff profile. It preserves the image
and kernel of the free complex. Put
\(J_{ab}=\Pi\mathcal L_{ab}\mathsf Z\) and
\(\mathcal C=-\sum_{a<b}J_{ab}^2\).

\begin{theorem}[Exact rotational test on the bilinear quotient]
\label{thm:v40-Casimir}
The computed rational matrices obey
\begin{equation}
 \boxed{
 (\mathcal C-24I)(\mathcal C-40I)=0,\qquad
 \operatorname{rank}(40I-\mathcal C)=25,\qquad
 \operatorname{rank}(\mathcal C-24I)=18.}
 \label{eq:v40-Casimir}
\end{equation}
In particular,
\begin{equation}
 \boxed{\mathcal C^{-1}=(64I-\mathcal C)/960.}
 \label{eq:v40-Casimir-inverse}
\end{equation}
There is no rotation-invariant class in \(H^0\). If a closed bilinear
coefficient has all six rotational variations in
\(\operatorname{im}\partial_{-1}\), then the coefficient itself is exact.
\end{theorem}
\begin{proof}
The script differentiates the native curvature columns using
\eqref{eq:v40-rotation-action}, verifies their closure, and applies the
already certified \(\Pi\). This reconstructs each displayed \(43\times43\)
rational matrix. Exact matrix multiplication verifies the quadratic
polynomial, and the idempotents
\((40I-\mathcal C)/16\) and \((\mathcal C-24I)/16\) have traces 25 and 18.
This proves the multiplicities and inverse. An invariant vector is killed
by all six \(J_{ab}\), hence by the invertible \(\mathcal C\), and is zero.
For a closed input the induced representation gives
\(J_{ab}\Pi b=\Pi\mathcal L_{ab}b\). The last claim follows.
\end{proof}

Equivalently the class can be reconstructed from its six rotational
variations:
\[
 \Pi b=-\mathcal C^{-1}\sum_{a<b}
                 J_{ab}\Pi\mathcal L_{ab}b.
\]
In a rotation-invariant positive inner product on this finite quotient,
obtained by compact group averaging, the further estimate is
\[
 \|\Pi b\|\le24^{-1/2}
  \left(\sum_{a<b}\|\Pi\mathcal L_{ab}b\|^2\right)^{1/2}.
\]
This norm is not identified with the raw coefficient \(\ell^1\) norm.
Neither estimate asserts that the native reference or the prescribed
minimum-norm primitive has the required rotational invariance modulo
exact terms. That is a test of its actual coefficients.

There are genuine nonexact anisotropic examples. Consider
\(\int (R^{(1)}_{0202})^2\). At the complex polynomial data
\[
 p=(1,i,0,0),\qquad v=(0,0,1,i),\qquad q=vv^T,
\]
one has \(p^2=0\), \(pq=0\), \(\operatorname{tr}q=0\), and \(E(p)q=0\),
but \((R^{(1)}_{0202})^2=1/4\). An exact bilinear action coefficient
\(EA+A^TE\) vanishes on these data. Thus this curvature square is not exact.
This is a complexification test of a real polynomial identity, not a mode
inserted into the Euclidean Gaussian integral. It illustrates why replacing
an arbitrary native tensor by just \(R^2\) and \(R_{\mu\nu}^2\) would erase
possible boundary data.

For comparison, the familiar isotropic squares are indeed exact in this
free bank. In tensor notation the normalized Euler tensor is
\(2G^{(1)}\). Its contraction with
\(-\tfrac12\delta_{\mu\nu}R^{(1)}\) gives \((R^{(1)})^2\), and its contraction
with \(\tfrac12R^{(1)}_{\mu\nu}-\tfrac14\delta_{\mu\nu}R^{(1)}\) gives
\(R^{(1)}_{\mu\nu}R^{(1)\mu\nu}\). Their partners are gauge-invariant linear
curvatures. These elementary examples do not exhaust the 2,725-component
native array and are not used instead of the computed quotient.

\section{Only nine quadratic primitive choices reach the cubic obstruction}
\label{sec:v40-cubic}

Let \(n_1,\ldots,n_9\) be the closed ghost-number-minus-one representatives
specified by the negative-degree certificate. Then
\begin{equation}
 \boxed{
 \ker\partial_{-1}
 =\operatorname{im}\partial_{-2}
       \oplus\operatorname{span}\{n_1,\ldots,n_9\}.}
 \label{eq:v40-negative}
\end{equation}
The full kernel has dimension 214; 205 directions are free exact generators.
These nine parameters are not the 43 ghost-number-zero boundary classes.

Write the ordinary gravitational differential as in V39,
\(s_g=s_0+s_1+s_2+\cdots\), with \(s_j\) raising arity by \(j\).
Suppose \(k=k_2+k_3+\cdots+k_5\) is a closed critical normalization
\emph{difference}. Thus \(s_0k_2=0\) and
\(s_0k_3+s_1k_2=0\). The map \(s_1\) below is the actual native ordinary
cubic differential, not a newly chosen deformation.

\begin{theorem}[Complete two-stage low-arity exactness test]
\label{thm:v40-two-stage}
The restriction of \(k\) through arity three is ordinary exact if and only if
\begin{equation}
 \boxed{\begin{aligned}
 &\Pi k_2=0,\\
 &k_3-s_1\mathsf S k_2
    \in\operatorname{im}(s_0:\mathscr C^{-1}_3\to\mathscr C^0_3)
       +\operatorname{span}\{s_1n_1,\ldots,s_1n_9\}.
 \end{aligned}}
 \label{eq:v40-two-stage}
\end{equation}
The second line is an explicit native cubic range problem; its matrix and
populated right-hand side are not evaluated in this installment.
\end{theorem}
\begin{proof}
A quadratic primitive exists exactly when \(\Pi k_2=0\), by
\eqref{eq:v40-contraction}. Every such primitive is
\[
 Q_2=\mathsf S k_2+\partial_{-2}z+\sum_{\nu=1}^9 a_\nu n_\nu.
\]
The cubic equation is \(k_3=s_1Q_2+s_0Q_3\). Since
\(s_1s_0=-s_0s_1\), the term \(s_1\partial_{-2}z\) is a cubic free exact
term and can be absorbed into \(Q_3\). This gives exactly the second line
of \eqref{eq:v40-two-stage}. Conversely any representation of that line
constructs \(Q_2,Q_3\). Each \(s_1n_\nu\) is closed under \(s_0\), because
\(s_0n_\nu=0\). No claim is made that these nine cubic images are
independent modulo free exact terms; nine is the complete parameter bank
that must be tested, not the rank of a computed nonlinear obstruction.
\end{proof}

Thus a deterministic quadratic section is not permission to ignore its
homogeneous freedom at the next equation. In particular, checking only
whether \(k_3-s_1\mathsf Sk_2\) is free exact would impose an unnecessary
normalization choice. Conversely, freely choosing all 214 kernel
coefficients would retain 205 directions that cannot change this cubic
cohomological test.

\section{Relative matter matching and the boundary still to be populated}
\label{sec:v40-relative}

Take two solutions \(H_g,\widetilde H_g\) of the same gravitational
reference equation through arity five and set \(k=H_g-\widetilde H_g\).
The fixed \(H_g\) is still V25's prescription. No restriction of it has been
numerically evaluated here. In particular the theorem is not applied to
\(H_g\) alone as though its inhomogeneous row were zero.

If \eqref{eq:v40-two-stage} passes, subtract the resulting
\(s_g(Q_2+Q_3)\). The residual begins at arity four. V39's inherited tail
result supplies \(Q_4,Q_5\), giving
\[
 k=\operatorname{pr}_{A\le5}s_g(Q_2+Q_3+Q_4+Q_5).
\]
A comparison matter primitive then returns to the fixed gravitational
prescription by the same triangular identity:
\begin{equation}
 \boxed{
 H_2=\widetilde H_2+\operatorname{pr}_{A\le5}t(Q_2+Q_3+Q_4+Q_5),
 \qquad d_HH_2=x_2-tH_g.}
 \label{eq:v40-matter-compensation}
\end{equation}
This is conditional on a comparison solution and the tested cubic image,
not a claim that either has now been produced.

Let \(t_{[j]}\) raise arity by \(j\). The compensation retains every term
\begin{equation}
 \begin{aligned}
 [tQ]_3&=t_{[1]}Q_2,\\
 [tQ]_4&=t_{[2]}Q_2+t_{[1]}Q_3,\\
 [tQ]_5&=t_{[3]}Q_2+t_{[2]}Q_3+t_{[1]}Q_4.
 \label{eq:v40-matter-arities}
 \end{aligned}
\end{equation}
For \(Q_2=\int(\eta A(\partial)q+\sigma B(\partial)c)\), the leading term
is the actual scalar stress contracted with \(Aq\), together with the scalar
antifield current contracted with \(Bc\). Its families are \(q\phi^2\) and
\(\upsilon\phi c\), with the mass and derivative contacts included. The
higher lines are precisely existing families of V32. There is no flat
quadratic-scalar term of arity two in this compensation. Critical lower
source and curvature coefficients can nevertheless change and must not be
held fixed without their relative tests.

The finite matrices and the maps \(t_{[j]}\) at this fixed order have no
mesh or volume dependence. Any physical coefficient estimate must first
apply the stated \(\chi\)-normalization and the actual input coefficient
bound. Windowing a finite polynomial gives the inherited rooted
multilinear estimates, with one source in \(L^1\) and the others in
\(L^\infty\). This is not a new loop remainder estimate and does not turn a
divergent local normalization coefficient into a uniformly small action.

The advance is therefore specific: the previously uncomputed quadratic
boundary now has a 43-coordinate exact test and a bounded native
coefficient solver; the subsequent cubic ambiguity is reduced to nine
specified directions. The remaining operands are the populated difference
of the fixed and comparison primitives, and the cubic range in
\eqref{eq:v40-two-stage}. A nonzero one of the 43 coordinates obstructs a
\emph{bilinear exact comparison}; it is not automatically an anomaly or a
proof that the relative matter equation has no solution by another route.
Neither the complete \(\xi_0\) nor \(\xi_1\) reference array has been evaluated.

\section{Verification and closing ledgers for V40}
\label{sec:v40-ledgers}

The standalone checker reconstructs the native polynomial matrices,
verifies both consecutive zero products, all 210 curvature-square columns,
the block inverses and independent boundary minors, all rank upper bounds,
the nine primitive-ambiguity representatives, the quotient normalization,
exact coefficient norms, the six induced rotation matrices and their
Casimir identity. It then tests its public solver on every generator column
and every cohomology representative. The local complex is reduced by exact
integration by parts before these matrices are formed. Complex null data
provide the separately labelled nonexact curvature witness.

A supplied bilinear array can be compiled using
\srcid{check_ESM_bilinear_boundary_V40.py} in
\texttt{--solve input.json --output solution.json} mode.
The input format and physical normalization are in the script's help.
It returns the free identity defect, 43 quotient coordinates, and a local
\((A,B)\) primitive of the exact component. It reports a full primitive only
when both the defect and the quotient coordinates vanish. The certificate
contains rational inverse data rather than a floating-point singular-value
threshold. These finite computations are not a proof-assistant
formalization, a recertification of all inherited analyses, or a numerical
evaluation of the critical reference loops.

\subsection*{Proved}
\addcontentsline{toc}{subsection}{V40: Proved}
The complete integrated native bilinear complex, exact ranks
\eqref{eq:v40-ranks}, 43-dimensional ghost-zero quotient, nine-dimensional
ghost-minus-one quotient, explicit curvature representatives, rational
contraction and exact coefficient-map norms are established by the supplied
finite arithmetic witness. The rotational Casimir is invertible with the
displayed spectrum; the invariance test and anisotropic example follow.
The two-stage quadratic/cubic criterion and complete matter-compensation
arity bookkeeping are proved.

\subsection*{Inherited}
\addcontentsline{toc}{subsection}{V40: Inherited}
V1--V39 and their labels are unchanged. The ordinary native Euler and gauge
symbols, its nonlinear differential, the V39 high-arity tail theorem, the
V25 fixed gravitational prescription, V34 active scalar normalization,
V35 coupled mass equations, V32 source bank and V38 logarithmic packet
retain their original status. No previously derived continuum estimate is
claimed as new.

\subsection*{Literature input}
\addcontentsline{toc}{subsection}{V40: Literature input}
The free spin-two BV language and deformation interpretation are the
established framework of \cite{V39BDGH,V20BBH1994}. Unlike V39's
high-arity conclusion, the two native cohomology dimensions in this
installment are obtained from the explicitly certified coefficient complex,
not from an imported vanishing or classification theorem. No new assertion
about the full matter-inclusive anomaly cohomology is made.

\subsection*{Conditional}
\addcontentsline{toc}{subsection}{V40: Conditional}
If an actual closed normalization difference passes its 43-coordinate test
and the stated cubic image test, the finite primitive and inherited tail
restore the fixed gravitational prescription with all matter contacts. The
rotational test applies only after the actual rotational variations are
shown exact. Neither property is assumed for the unevaluated fixed native
primitive. Bounds transfer real input coefficient bounds; they are not
uniform in a growing EFT or source bank.

\subsection*{Open}
\addcontentsline{toc}{subsection}{V40: Open}
The fixed primitive's populated low-arity reference coefficients and the
native cubic system in \eqref{eq:v40-two-stage} remain to be evaluated.
Therefore \(Dh_0=\xi_0\), \(Dh_1=\xi_1-Mh_0\), and the complete relative
critical normalization are not claimed solved. The current computation
replaces the quadratic part of that uncomputed range problem by an exact
small output test; it does not remove the remaining cubic or inhomogeneous
data. Internal gauge/chiral interfaces, higher loops, removal of the lower
reference and an invariant analytic Einstein--SM class remain separate.

\begin{center}\small
\begin{tabular}{>{\raggedright\arraybackslash}p{.23\textwidth}>{\raggedright\arraybackslash}p{.68\textwidth}}
\toprule
Variable & Scope in V40\\\midrule
Order & Ordinary free critical bilinear complex, and an exact reduction of
the cubic comparison. No new loop integral or loop-order extension.\\
Fields & All ten metric components and four minimal ghosts, with both
antifield types. No rotation-invariant coefficient ansatz.\\
Cutoff and volume & Rational maps independent of mesh, periods, lower
reference and shell count at this fixed bank. Input local coefficients can
still depend on all these data.\\
Mass and matter & Gravity matrices have no mass parameter. Their inherited
stress lift retains its polynomial scalar mass dependence; no new massive
coefficient is evaluated.\\
Norms & Exact values in the declared monomial basis. Physical Euler
normalization requires the explicit factor \(\chi/2\). No unbounded-order
uniformity.\\
Normalization & Fixed gravity, V34 and V38 unchanged. No counterterm is
inserted without an actual comparison input.\\
Computation & Native bilinear matrices and inverse certificates computed;
native cubic matrices, inhomogeneous reference arrays, and loop integrals
not computed.\\\bottomrule
\end{tabular}
\end{center}

\clearpage
\part*{V41: A faithful native cubic test for all nine quadratic primitive classes}
\addcontentsline{toc}{part}{V41: Faithful cubic test and unique ambiguity extraction}

\section{The nine-channel cubic question has a finite on-shell answer}
\label{sec:v41-start}

Sections 1--318 and their labels are preserved. V40 supplies nine rational
representatives \(n_\nu\) of the free bilinear group
\(H^{-1}(s_0)\), but explicitly does not evaluate their cubic images
\(s_1n_\nu\). Here \(s_g=s_0+s_1+\cdots\) is the ordinary zero-mesh
native gravitational differential, with \(s_j\) raising arity by \(j\).
The present calculation settles that particular missing question:
\emph{all nine cubic images are independent modulo free exact terms}.

The proof uses nine real linear tests of the actual native cubic on complex
free solutions. It constructs a nonsingular rational \(9\times9\) matrix,
not the full cubic cohomology matrix. The tests annihilate every free-exact
cubic, including all its possible antifield partners. Consequently their
rank proves an injection
\[
 H^{-1}_{A=2,w=4}(s_0)\longrightarrow H^0_{A=3,w=4}(s_0),
 \qquad [n]\longmapsto[s_1n].
\]
A closed bilinear generator representing a nonzero one of these nine
classes therefore cannot be extended to an ordinary variational symmetry
through the first nonlinear interaction. This does not prevent using it
as a field-redefinition generator: its nonzero cubic image is precisely
what must then be included in a normalization comparison.

The calculation does not assume rotational invariance of an unknown
coefficient. Its finitely many rotated tests probe different components.
Nor does it place a complex null momentum in the positive Euclidean
Gaussian integral. Complexification is used solely to separate identities
of local momentum polynomials. No propagator, loop integral, infrared limit,
new cutoff profile, or change of the active counterterms is involved.

The classification of full local generalized symmetries of vacuum Einstein
gravity provides established comparison material \cite{V41AndersonTorre}.
It is not used to identify the first order at which the present free
symmetries fail. That order and the rank-nine statement are proved below
by the native finite calculation. In particular neither that classification
nor the nine tests solves the remaining inhomogeneous critical equations.

\section{The native cubic and the free-exact-annihilating tests}
\label{sec:v41-evaluator}

Use exactly V40's ten-coordinate order, Euler matrix \(E(p)\), generator
\(R(p)\), and action normalization. Thus the normalized action is
\(2/\chi\) times the physical ordinary action and its quadratic Hessian is
\(E\). Let \(\mathscr V_3(h_1,h_2,h_3;p_1,p_2,p_3)\) be the third
polarized coefficient of its native zero-mesh action, with
\(p_1+p_2+p_3=0\). The complete stationary formula
\eqref{eq:v10-reduced-jets} is used; its ordinary limit is not replaced by
an independently chosen vertex.

Here is a finite formula fixing all phases and factors. Put
\(D=\operatorname{diag}(i,1,1,1)\),
\(p_i^L=Dp_i\), and \(u_i=\mathsf E_L(Dh_iD)\), where
\[
 (\mathsf E_L(h))_{00}=-h_{00}/2,\qquad
 (\mathsf E_L(h))_{i0}=h_{0i},\qquad
 (\mathsf E_L(h))_{ij}=h_{ij}/2\quad(i,j>0),
 \qquad (\mathsf E_L(h))_{0i}=0.
\]
These are the inherited V13 continuation and affine coframe coordinates.
Let \(C_0\) be the \(\chi=1\) Lorentz Cartan block and let \(f_1,f_2\)
be the first and polarized second cofactor-load variations of
\eqref{eq:v10-action}. Set
\[
 X_i=-C_0^{-1}f_1(u_i;p_i^L),\qquad
 C_1(u)=\operatorname{tr}(u)C_0-T_u^TC_0-C_0T_u,
 \qquad T_u=u^T\otimes I_6.
\]
Then
\begin{equation}
 \boxed{
 \mathscr V_3=-2\sum_{i=1}^3
 \left\{
 f_2(u_j,u_k;p_j^L+p_k^L)^TX_i
 +X_j^TC_1(u_i)X_k
 \right\},\quad \{j,k\}=\{1,2,3\}\setminus\{i\}.}
 \label{eq:v41-native-cubic}
\end{equation}
The order of \(j,k\) is immaterial because the two pairings are symmetric.
The factor \(-2\) is the continued action sign followed by V40's
\(2/\chi\) normalization. The checker reconstructs the cofactor table
and the 24-variable Cartan quadratic directly and verifies that the
quadratic formula gives all 100 entries of \(E\).

For clarity, the cofactor derivative used in that reconstruction is the
derivative of
\[
 \ell_{\mu\nu,ij}(e)
 =2\epsilon_{\mu\nu rt}\epsilon_{ijab}
 (e_{ar}e_{bt}-e_{at}e_{br}),
\]
where each complementary pair is ordered increasingly. Its load puts
\(ip_\nu\ell_{\mu\nu,ij}\) in slot \((\mu,ij)\) and
\(-ip_\mu\ell_{\mu\nu,ij}\) in slot \((\nu,ij)\).
This specifies every finite operand in \eqref{eq:v41-native-cubic}.
No finite-mesh shared-link remainder is discarded in a finite-mesh claim:
only its already established zero-mesh local coefficient is being tested.

\begin{lemma}[A local cubic test annihilating every free exact coefficient]
\label{lem:v41-annihilator}
Let \(p_i\in\mathbb C^4\) and \(h_i\in\operatorname{Sym}_4(\mathbb C)\)
satisfy
\[
 p_1+p_2+p_3=0,\qquad E(p_i)\operatorname{col}(h_i)=0.
\]
Extract the antifield-free, ghost-free cubic coefficient of a local
functional and evaluate its polarized symbol on these data. This linear
functional annihilates \(s_0Q_3\) for every local ghost-number-minus-one
cubic generator \(Q_3\), as well as every total divergence. Its real and
imaginary parts are real linear annihilators on real coefficient arrays.
\end{lemma}
\begin{proof}
After setting ghosts and antifields to zero, the only term of \(s_0Q_3\)
that can survive is a free metric Euler covector contracted with the
quadratic coefficient multiplying a metric antifield in \(Q_3\).
Every polarized term contains one of the factors
\(E(p_i)\operatorname{col}(h_i)\), and is zero. Ghost-antifield and
multiple-antifield terms retain a ghost or antifield and cannot contribute
to this extraction. A divergence contains \(p_1+p_2+p_3\).
Real polynomial identities complexify, so complex data are legitimate
separators; taking real and imaginary parts gives real coefficient maps.
\end{proof}

Use the following single seed:
\begin{equation}
 \begin{gathered}
 p=(1,i,0,0)^T,\quad t=(0,0,1,i)^T,\quad
 (p_1,p_2,p_3)=(p,t,-p-t),\\
 e_1=(0,0,1,-i)^T,\quad e_2=(1,-i,0,0)^T,\quad e_3=p-t,
 \qquad h_i=e_ie_i^T.
 \label{eq:v41-seed}
 \end{gathered}
\end{equation}
All three momenta are null; each polarization is transverse and traceless.
They are therefore free solutions for V40's unreduced Euler matrix.
The pairings are bilinear, with transpose rather than Hermitian conjugation.
Direct native evaluation gives
\begin{equation}
 \mathscr V_3(h_1,h_2,h_3;p,t,-p-t)=-64.
 \label{eq:v41-seed-value}
\end{equation}
This value is also checked against the ordinary metric
\(\Gamma\Gamma\) expression. The latter is only a normalization check:
all entries of the obstruction matrix below are calculated by
\eqref{eq:v41-native-cubic} in the inherited coframe chart.

If \(O\) is a real rational orthogonal matrix, rotate every momentum and
polarization by \(p_i\mapsto Op_i\), \(h_i\mapsto Oh_iO^T\).
Let \(\mathscr E_O\) denote the evaluation of a cubic coefficient on
that data. It is not rotational averaging. Define nine tests by
\(\ell_a=\frac1{64}\operatorname{Re}\mathscr E_O\) or
\(\frac1{64}\operatorname{Im}\mathscr E_O\), according to this table:
\begin{center}\small
\begin{tabular}{cll}
\toprule
Row & Orthogonal matrix \(O\) & Real or imaginary part\\\midrule
1 & \(I\) & real\\
2 & \(I\) & imaginary\\
3 & \(P_{(0,1,3,2)}\) & real\\
4 & \(P_{(0,1,3,2)}\) & imaginary\\
5 & \(P_{(0,2,1,3)}\) & real\\
6 & \(P_{(0,2,3,1)}\) & imaginary\\
7 & \(O_{02}\) & real\\
8 & \(O_{03}\) & imaginary\\
9 & \(O_{12}\) & imaginary\\\bottomrule
\end{tabular}
\end{center}
Here \(P_\pi e_j=e_{\pi_j}\). The matrix \(O_{ab}\) is identity outside
its \(ab\)-plane and has block
\(\left(\begin{smallmatrix}3/5&4/5\\-4/5&3/5\end{smallmatrix}\right)\)
in that plane. These are seven distinct complex panels and nine real tests.
The permutation matrices need not preserve orientation; each simply gives
another free polynomial test. Write \(\ell=(\ell_1,\ldots,\ell_9)^T\).

\section{The nine cubic classes are independent}
\label{sec:v41-injection}

Use V40's nine representatives in their certified order. Export their
sparse polynomials as \(n_\nu=\int\eta^TA_\nu(\partial)q\), with the
Fourier convention of V40. A newly useful feature of those particular
representatives is
\begin{equation}
 B_\nu=0,\qquad
 A_\nu R=0,\qquad EA_\nu+A_\nu^TE=0.
 \label{eq:v41-representatives}
\end{equation}
The V41 certificate contains all 345 nonzero rational monomial entries
of their \(A\) matrices and no new fitted coefficient. It records the hash
of the V40 source certificate. Their parities are, in order,
\(0011,0101,0110,1001,1010,1100,1111,1111,1111\).
The polynomial identities in \eqref{eq:v41-representatives} are rechecked
from the native matrices. V40's dimension theorem is inherited.

Let
\[
 L_\nu=s_1n_\nu.
\]
Each \(L_\nu\) is a complete local cubic cocycle because
\(s_0s_1+s_1s_0=0\). Its ghost-free action coefficient is precisely
\begin{equation}
 (L_\nu)_{q^3}(h_1,h_2,h_3)
 =\sum_{j=1}^3
 \mathscr V_3(h_1,\ldots,
 \operatorname{mat}(A_\nu(p_j)\operatorname{col}h_j),
 \ldots,h_3;p_1,p_2,p_3).
 \label{eq:v41-cubic-insertion}
\end{equation}
Equation \eqref{eq:v41-representatives} ensures that every replaced leg
is still a free solution. The ghost and antifield pieces of \(s_1n_\nu\)
are not asserted zero; they simply do not enter this particular annihilator.

\Needspace{22\baselineskip}
\begin{theorem}[Faithful native cubic obstruction matrix]
\label{thm:v41-injection}
In the ordered tests and representatives above,
\begingroup\small\setlength{\arraycolsep}{3.5pt}
\begin{equation}
\boxed{
 \mathsf C=(\ell_a(L_\nu))_{a,\nu}=
 \begin{pmatrix}
 0&-1&0&0&-2&0&0&0&0\\
 0&0&2&-1&0&0&0&0&0\\
 0&0&-2&-1&0&0&0&0&0\\
 0&1&0&0&-2&0&0&0&0\\
 1&0&0&0&0&2&0&0&0\\
 -1&0&0&0&0&2&0&0&0\\
 0&7/25&0&0&-2&0&0&24/25&0\\
 0&0&-14/25&-1&0&0&24/25&0&-24/25\\
 0&0&2&7/25&0&0&0&24/25&24/25
 \end{pmatrix}.}
 \label{eq:v41-matrix}
\end{equation}
\endgroup
It has
\begin{equation}
 \boxed{
 \det\mathsf C=\frac{884736}{15625},\qquad
 \|\mathsf C^{-1}\|_{1\to1}=\frac{25}{8},\qquad
 \|\mathsf C^{-1}\|_{\infty\to\infty}=\frac{59}{12}.}
 \label{eq:v41-matrix-norms}
\end{equation}
The map
\(H^{-1}_{A=2,w=4}(s_0)\to H^0_{A=3,w=4}(s_0)\),
\([n]\mapsto[s_1n]\), is injective. In particular all nine images are
nonzero and independent modulo free exact coefficients.
\end{theorem}
\begin{proof}
Substitute the finite cofactor and Cartan data into
\eqref{eq:v41-native-cubic}, then substitute the exported polynomials
\(A_\nu\) into \eqref{eq:v41-cubic-insertion}. All momenta, polarizations,
and rotations lie in \(\mathbb Q(i)\). The real and imaginary parts of the
81 outputs, divided by 64, are exactly \eqref{eq:v41-matrix}.
For example the unrotated complex row before division is
\[
 (0,-64,128i,-64i,-128,0,0,0,0).
\]
The three rational rotations supply the last three independent rows;
permutations alone would not establish rank nine for this seed.

This is a finite arithmetic certificate, not a numerical rank inference.
The companion script reconstructs the 24-variable Cartan matrix, both
cofactor variations, all nine representative matrices, and every test
before checking the output. It also checks an independent metric
\(\Gamma\Gamma\) expression on the seed and three nontrivial altered-leg
panels. Exact multiplication of the displayed rational matrix and the
certificate's inverse gives the identity on both sides; direct determinant
and absolute column/row sums give \eqref{eq:v41-matrix-norms}.

For the cohomological conclusion, first note that the assignment is
well-defined: replacing \(n\) by \(n+s_0z\) changes \(s_1n\) by
\(-s_0s_1z\). If \(\sum a_\nu L_\nu=s_0Q_3\) modulo a divergence,
\cref{lem:v41-annihilator} gives \(\mathsf Ca=0\), hence \(a=0\).
V40 proves that the nine \([n_\nu]\) span the bilinear group. This proves
injectivity on the entire group, without computing the full cubic target
cohomology or its matrix rank.
\end{proof}

\begin{corollary}[No nontrivial bilinear symmetry in this bank extends through cubic order]
\label{cor:v41-no-extension}
Suppose \(N_2\) is a closed weight-four bilinear generator and there is
a local cubic \(N_3\) with
\[
 s_1N_2+s_0N_3=0.
\]
Then \([N_2]=0\) in V40's nine-dimensional quotient.
\end{corollary}
\begin{proof}
Express the class of \(N_2\) in the nine representatives and apply
\cref{thm:v41-injection}. The statement is restricted to this derivative
weight, local field class, and first nonlinear equation. It makes no claim
that a nonzero \(s_gN_2\) cannot be used as an exact counterterm.
\end{proof}

\section{A unique nine-coefficient extraction for the cubic comparison}
\label{sec:v41-cubic-compiler}

For any real local cubic coefficient \(r\), set
\begin{equation}
 a(r)=\mathsf C^{-1}\ell(r),\qquad
 \mathscr P_9r=\sum_{\nu=1}^9a_\nu(r)L_\nu,\qquad
 \mathscr R_9r=r-\mathscr P_9r.
 \label{eq:v41-projector}
\end{equation}
The maps use only fixed evaluations. They do not divide by a variable
momentum or by a mass. Their nine-coefficient part can be written without
matrix inversion. For \(y=\ell(r)\),
\begin{align}
 a_1&=(y_5-y_6)/2,& a_2&=(-y_1+y_4)/2,\notag\\
 a_3&=(y_2-y_3)/4,& a_4&=-(y_2+y_3)/2,\notag\\
 a_5&=-(y_1+y_4)/4,& a_6&=(y_5+y_6)/4,\notag\\
 a_8&=(25y_7-7a_2+50a_5)/24,\notag\\
 a_9&=(25y_9-50a_3-7a_4-24a_8)/24,\notag\\
 a_7&=a_9+(25y_8+14a_3+25a_4)/24.
 \label{eq:v41-explicit-compiler}
\end{align}

\begin{proposition}[Exact splitting of the evaluated cubic channels]
\label{prop:v41-projector}
The map \(\mathscr P_9\) is a projection onto
\(\operatorname{span}\{L_1,\ldots,L_9\}\), and
\[
 \ell\mathscr R_9=0,\qquad
 \mathscr P_9(s_0Q)=0,\qquad
 s_0\mathscr R_9r=s_0r.
\]
For every input, closed or not,
\begin{equation}
 \boxed{
 \|a(r)\|_1\le\frac{25}{8}\|\ell(r)\|_1,
 \qquad
 \|a(r)\|_\infty\le\frac{59}{12}\|\ell(r)\|_\infty.}
 \label{eq:v41-coeff-bound}
\end{equation}
Neither \(\ell\mathscr R_9r=0\) nor \(s_0\mathscr R_9r=0\) by itself
asserts that \(\mathscr R_9r\) is free exact.
\end{proposition}
\begin{proof}
Use \(\ell(L)=\mathsf C\), the annihilator lemma, and \(s_0L=0\).
The matrix norm calculation proves the inequalities. These operations do
not remove a nonclosed input's actual free identity defect.
\end{proof}

For a direct coefficient-norm interpretation, take the \(\ell^1\) norm
of the polarized ghost-free cubic kernel after eliminating its third
momentum. On the seven panels above, every momentum component has modulus
at most \(7/5\), and every independent metric entry has modulus at most
\((7/5)^2\). Its homogeneous derivative degree is four. Therefore
\[
 \|\ell(r)\|_\infty
 \le \frac{7^{10}}{64\,5^{10}}\|r_{q^3}\|_{\mathrm{pol},1}
 <\|r_{q^3}\|_{\mathrm{pol},1}.
\]
The polarized-kernel norm includes its actual symmetrization coefficients;
no unrecorded factorial is suppressed. Also the exported representatives
have total coefficient norms
\[
 (\|n_\nu\|_1)_{\nu=1}^9=(31,24,44,24,44,44,28,28,32),
\]
and therefore
\[
 \Big\|\sum a_\nu n_\nu\Big\|_1\le44\|a\|_1.
\]
These are finite local-map bounds, not estimates on a critical loop input
which has not yet been computed.

\begin{theorem}[V40's cubic ambiguity is fixed by the actual cubic target]
\label{thm:v41-two-stage}
Let \(k_2+k_3\) be a closed critical normalization difference through
arity three. Suppose V40's quadratic test passes:
\(\Pi k_2=0\). With its fixed section \(\mathsf S\), define
\[
 r_3=k_3-s_1\mathsf Sk_2,\qquad
 a=\mathsf C^{-1}\ell(r_3),\qquad r_3^\perp=\mathscr R_9r_3.
\]
Then a local primitive through arity three exists if and only if
\begin{equation}
 \boxed{r_3^\perp\in\operatorname{im}
 (s_0:\mathscr C^{-1}_3\longrightarrow\mathscr C^0_3).}
 \label{eq:v41-remaining-range}
\end{equation}
When this condition holds, it has a primitive with
\begin{equation}
 Q_2=\mathsf Sk_2+\sum_\nu a_\nu n_\nu,\qquad
 s_0Q_3=r_3^\perp.
 \label{eq:v41-primitive}
\end{equation}
The nine coefficients are unique relative to the chosen quadratic section,
modulo bilinear free-exact generators. No search over nine undetermined
coefficients remains in the cubic image test.
\end{theorem}
\begin{proof}
Closure gives \(s_0k_3=-s_1k_2\), while
\(s_0s_1\mathsf Sk_2=-s_1k_2\); hence \(r_3\) is closed.
V40 expresses every quadratic primitive as
\(\mathsf Sk_2+s_0z+\sum b_\nu n_\nu\).
Its cubic equation implies
\(r_3=\sum b_\nu L_\nu+s_0Q'_3\), because
\(s_1s_0z=-s_0s_1z\). Apply \(\ell\):
\(\ell(r_3)=\mathsf Cb\). Invertibility forces \(b=a\), and the
remaining equation is \eqref{eq:v41-remaining-range}. Conversely a
solution of that equation gives \eqref{eq:v41-primitive} directly.
A different quadratic section changes the extracted nine coordinates
by exactly the corresponding change of representative; the class of the
total quadratic primitive is unchanged. Free-exact changes of its
representative do not affect the nine test values.
\end{proof}

This theorem does not discard the nine directions because they fail to
extend as symmetries. In a normalization comparison their nonzero cubic
images are needed. What is removed is their status as \emph{free choices}:
the proposed cubic target fixes them uniquely.

\section{Relative matter transport}
\label{sec:v41-relative}

The application still concerns the difference \(k=H_g-\widetilde H_g\)
of two solutions of the same inhomogeneous gravitational reference equation.
It is not legitimate to apply a closed-coefficient test to \(H_g\) alone.
Neither the actual \(k_2,k_3\) nor the native reference arrays
\(\xi_0,\xi_1\) have been supplied numerically by this calculation.

Once the quadratic and remaining cubic tests pass, V39 supplies the closed
quartic/quintic tail primitive. The complete scalar compensation remains
\[
 H_2=\widetilde H_2+
 \operatorname{pr}_{A\le5}t(Q_2+Q_3+Q_4+Q_5).
\]
The new information is that the part of \(Q_2\) in the nine-class quotient
is fixed by \eqref{eq:v41-explicit-compiler}. Since the representatives
have \(B_\nu=0\), their leading matter correction is exactly
\begin{equation}
 \boxed{
 t_{[1]}\Big(\sum_\nu a_\nu n_\nu\Big)
 =\int\mathcal E_{H,0}^{\mu\nu}
 \left[\operatorname{mat}
 \Big(\sum_\nu a_\nu A_\nu(-i\partial)\operatorname{col}q\Big)
 \right]_{\mu\nu}.}
 \label{eq:v41-matter-leading}
\end{equation}
The notation \(A_\nu(-i\partial)\) means its declared Fourier multiplier;
the symbol is real and of even degree. The native scalar stress is the
one used in V39--V40, including its mass term. There is no leading scalar
current replacement from these nine representatives because their ghost
redefinition matrix is zero. Other parts of \(Q_2\) and higher arities
still have the inherited current and source contacts. No part of them is
removed by the on-shell diagnostic.

The correction in \eqref{eq:v41-matter-leading} has type \(q\phi^2\),
not a flat two-scalar kinetic term. It has four total derivatives or the
corresponding mass-weighted components. Thus it lies in V32's existing
source bank. Its higher-coframe contacts come from the same full stress
map, not by holding \(\mathcal E_H\) fixed at flat metric after the first
variation. With the fixed finite conversion and \(\chi\) normalization,
the map inherits ordinary local rooted-kernel estimates from its polynomial
coefficients. No new shell summability or cutoff rate is asserted.

The exact new numerical data are the 81 native on-shell cubic entries,
their determinant, and the inverse coefficient bounds. The full off-shell
cubic matrices, the dimension of the complementary cubic quotient, and
\eqref{eq:v41-remaining-range} on a populated input have not been evaluated.
The nine tests are faithful on \(\operatorname{span}[s_1n_\nu]\), not on
the entire cubic cohomology. It remains possible for a nonzero free cubic
class to have all nine evaluations zero.

An immediate useful computational sequence is now fixed: form the actual
closed quadratic/cubic difference; use V40's 43-coordinate test and section;
evaluate the nine native panels on \(r_3\); subtract the uniquely determined
\(\sum a_\nu s_1n_\nu\); and solve the remaining free cubic image equation
with the full ghost and antifield coefficients retained. The present step
replaces an unevaluated nine-channel obstruction by an exact invertible
calculation. It does not claim the final relative critical primitive.

\section{Verification and closing ledgers for V41}
\label{sec:v41-ledgers}

The self-contained V41 checker reads its small rational certificate,
reconstructs the native free and Cartan matrices, and rechecks all nine
representative identities. It evaluates the native cubic on seven complex
panels, extracts the nine real rows, and verifies the displayed matrix and
its inverse over exact rational arithmetic. It separately checks the
quadratic normalization and selected metric-coordinate cubic evaluations.
The rotations, conservation equations, null/free-Euler conditions, inverse
norms and coefficient round trips are exact. The optional public
\texttt{--solve} operation consumes nine rational values in the declared
normalized test convention; it returns the nine coefficients together with
an explicit warning that the residual full cubic exactness has not been
tested. No floating-point rank threshold or numerical integration is used.

\begingroup
\renewcommand{\refname}{Additional comparison source for V41}
\begin{thebibliography}{99}
\bibitem[50]{V41AndersonTorre}
I.~M.~Anderson and C.~G.~Torre,
\emph{Classification of generalized symmetries for the vacuum Einstein equations},
Communications in Mathematical Physics \textbf{176} (1996), 479--539,
arXiv:gr-qc/9404030, doi:10.1007/BF02099248.
The full nonlinear classification is comparison material; it is not used
as a proof of the native rank-nine cubic obstruction computed here.
\end{thebibliography}
\endgroup

\Needspace{9\baselineskip}
\subsection*{Proved}
\addcontentsline{toc}{subsection}{V41: Proved}
The native cubic gives a nonsingular rational evaluation matrix on all nine
V40 bilinear primitive classes. Their first nonlinear images are independent
modulo free exact terms. A nontrivial class in that bank cannot extend to an
ordinary variational symmetry through cubic order. The actual cubic target
uniquely fixes all nine ambiguity coordinates, with exact inverse norms
\(25/8\) and \(59/12\). The residual cubic image criterion, its invariance
under free-exact quadratic changes, and the leading scalar-stress
compensation are established. The remaining full cubic exactness is not
inferred from these finite tests.

\subsection*{Inherited}
\addcontentsline{toc}{subsection}{V41: Inherited}
V1--V40, their labels and bibliographies remain unchanged. V40's bilinear
dimension and complete section, V39's high-arity tail theorem, the native
stationary coframe cubic and continuation, the ordinary nonlinear BV
identity, the fixed V25 prescription, V34's active scalar normalization,
V35's mass lift, V32's source bank and V38's logarithmic packet retain their
original scope. Rechecking selected algebraic operands is not a new proof
of every inherited statement.

\subsection*{Literature input}
\addcontentsline{toc}{subsection}{V41: Literature input}
The variational-symmetry and deformation interpretation uses the established
BV framework. The nonlinear vacuum symmetry classification
\cite{V41AndersonTorre} is used only for comparison. No external
cohomology-vanishing or full-symmetry theorem substitutes for any entry or
rank of the displayed native obstruction matrix.

\subsection*{Conditional}
\addcontentsline{toc}{subsection}{V41: Conditional}
A specified quadratic/cubic normalization difference passing V40's
quadratic test and \eqref{eq:v41-remaining-range} can be returned to fixed
gravity with the full inherited matter compensation. Neither the comparison
primitive nor its remaining cubic image is constructed here. Coefficient
map bounds control actual input perturbations; they are not bounds on an
unevaluated critical reference array or on an unbounded-order action class.

\subsection*{Open}
\addcontentsline{toc}{subsection}{V41: Open}
The populated native critical boundary \(k_2,k_3\), the remaining free
cubic exactness equation after the nine-channel subtraction, and the
inhomogeneous scalar--gravity equations \(Dh_0=\xi_0\),
\(Dh_1=\xi_1-Mh_0\) remain open. The nine homogeneous channels are no
longer an unspecified rank problem, but they do not exhaust all cubic
coefficients. Mixed internal-gauge/chiral interfaces, higher loops, removal
of the lower reference, and an invariant analytic full ESM class remain
separate tasks.

\begin{center}\small
\begin{tabular}{>{\raggedright\arraybackslash}p{.23\textwidth}>{\raggedright\arraybackslash}p{.68\textwidth}}
\toprule
Variable & Scope in V41\\\midrule
Order & Ordinary local weight four; bilinear generators and their first
nonlinear cubic images. No new quantum loop coefficient.\\
Fields & All ten metric components, four ghosts and their antifield
complex retained. Only the diagnostic evaluation sets ghosts/antifields
to zero after applying the differential.\\
Cutoff and volume & The rational maps have no mesh, period, lower-scale or
shell-count dependence. They test ordinary local coefficient arrays.\\
Test momenta & Fixed complex free solutions with exact momentum conservation;
not physical lattice integration modes and not an infrared limit.\\
Norms & V40 normalized action and the stated monomial/polarized bases.
Physical coefficients require the recorded \(2/\chi\) conversion. No
uniformity in growing source or EFT order.\\
Matter and mass & Leading stress lift retains the actual scalar mass term;
no new massive or matter-loop integral is evaluated.\\
Active normalization & Fixed gravity, V34 scalar prescription, and V38
packet unchanged; no counterterm is inserted without a successful
populated comparison.\\
Computation & Native cubic evaluations and rank-nine certificate computed.
Full cubic coefficient matrix and inhomogeneous loop arrays not computed.\\\bottomrule
\end{tabular}
\end{center}

\clearpage
\part*{V42: The complete Weyl-square current block and nine further cubic tests}
\addcontentsline{toc}{part}{V42: Conserved Weyl currents and an eighteen-class cubic test}

\section{The remaining cubic test has a genuinely larger obstruction space}
\label{sec:v42-start}

Sections 1--324, including their labels and bibliographies, are preserved.
V41 proves that the first nonlinear images of nine quadratic primitive
classes are independent. It carefully does not identify them with the
whole cubic cohomology. This installment tests that distinction on a
concrete part of the remaining native cubic problem: symmetric currents
quadratic in the linearized Weyl tensor, with arbitrary constant component
coefficients and no extra curvature derivatives.

The complete current space has 550 coefficients. Its conservation matrix
has exact rank 541, and its nine-dimensional kernel consists of contractions
of the continued Bel--Robinson tensor with a constant trace-free symmetric
tensor. Each current is supplied with an explicit local antifield partner,
so it defines a cubic cocycle in the same unreduced free metric complex.
Eighteen native on-shell tests separate these nine cocycles from all nine
V41 images. In particular, a displayed nonexact cubic passes every old V41
test. This gives a stronger necessary condition for the remaining image
equation, rather than declaring that condition solved.

Bel--Robinson tensors and their vacuum conservation are established
structures \cite{V42Senovilla}. The current/deformation question also
appears in the higher-derivative discussion of
\cite[Section 11]{V39BDGH}. Neither source is used as a proof that the
550-component native current matrix has a specified rank or that the
whole cubic local cohomology has been classified. Those finite rank and
separation calculations are carried out below. The construction makes no
rotational-invariance assumption about an unknown coefficient.

All statements concern the ordinary zero-mesh coefficient complex on the
existing Euclidean-continuation chart. The ten metric coordinates and all
four diffeomorphism ghosts remain present. No propagator, contour, measure,
reference profile, or active counterterm is changed. In particular the
fixed gravitational prescription, V34's subcritical scalar prescription,
and V38's complete logarithmic packet remain unchanged.

\section{Native Weyl values and their independently realizable first jets}
\label{sec:v42-weyl}

Use the V40 normalized free action. Its Euler covector is
\(e_A=(E(\partial)q)_A\) in the ten independent coordinates
\(A=(00,01,02,03,11,12,13,22,23,33)\).
As a symmetric tensor the corresponding Euler field is
\(\mathcal E_{ab}=e_{(ab)}/(2-\delta_{ab})=2G^{(1)}_{ab}\).
The Fourier convention is \(\partial_a\mapsto ip_a\).
The linearized curvature is fixed by
\[
 R^{(1)}_{abcd}
 =\tfrac12(\partial_b\partial_cq_{ad}
 +\partial_a\partial_dq_{bc}
 -\partial_a\partial_cq_{bd}
 -\partial_b\partial_dq_{ac}).
\]
With the linearized Schouten tensor
\[
 S_{bd}=\tfrac12\bigl(R^{(1)}_{bd}-\tfrac16\delta_{bd}R^{(1)}\bigr)
       =\tfrac14\mathcal E_{bd}
        -\tfrac1{12}\delta_{bd}\operatorname{tr}\mathcal E,
\]
define \(C=R^{(1)}-\delta\mathbin{\wedge}S\), where
\[
 (\delta\mathbin{\wedge}S)_{abcd}
 =\delta_{ac}S_{bd}-\delta_{ad}S_{bc}
  -\delta_{bc}S_{ad}+\delta_{bd}S_{ac}.
\]
These are linear differential polynomials in the inherited tangent \(q\),
not new independent physical fields. They obey
\begin{equation}
 \gamma_0C=0,\qquad
 \partial^a C_{abcd}=\partial_cS_{bd}-\partial_dS_{bc}.
 \label{eq:v42-cotton}
\end{equation}
The checker verifies the latter identity against the native ten-coordinate
matrix, including the off-diagonal coordinate factors.

\subsection{A rational ten-coordinate Weyl basis}

Represent an algebraic Weyl tensor by two symmetric trace-free
\(3\times3\) matrices \(E_W,B_W\):
\begin{align}
 C_{0i0j}&=(E_W)_{ij},&
 C_{0ijk}&=\epsilon_{jkl}(B_W)_{il},&
 C_{ijkl}&=\epsilon_{ijm}\epsilon_{kln}(E_W)_{mn}.
 \label{eq:v42-EB}
\end{align}
Spatial indices run from one to three; the other entries follow from
curvature symmetries. The independent coordinate vector is
\[
 w=((E_W)_{11},(E_W)_{12},(E_W)_{13},(E_W)_{22},(E_W)_{23},
     (B_W)_{11},(B_W)_{12},(B_W)_{13},(B_W)_{22},(B_W)_{23}).
\]
Let \(C_I\), \(I=0,\ldots,9\), be the resulting integer basis.
For a first derivative \(v_{aI}\), form the \(24\times40\) matrix
\begin{equation}
 S_{(b,cd),(a,I)}=(C_I)_{abcd},\qquad c<d.
 \label{eq:v42-div-table}
\end{equation}
Thus \(Sv\) is the coefficient array of \(\partial^aC_{abcd}\).

\begin{lemma}[The vacuum first-jet space is exactly the tested kernel]
\label{lem:v42-jets}
The matrix \(S\) has rank 16 and kernel dimension 24. Every algebraic
Weyl value \(C\) and every first derivative in \(\ker S\) are
simultaneously realized at the origin by a polynomial metric perturbation
solving the native free Euler equation. The value and derivative can be
chosen independently.
\end{lemma}
\begin{proof}
The rational table \eqref{eq:v42-EB} gives the stated rank and a specified
\(40\times24\) kernel basis \(Z\); both are reconstructed by the
checker. Algebraically, for a Weyl tensor in four dimensions the left and
right Hodge duals agree. Divergence zero is consequently equivalent to the
differential Bianchi identity on its first derivative. This equivalence
also follows by substituting the ten tables in \eqref{eq:v42-EB}.

For prescribed value and derivative, use the normal-coordinate polynomial
\[
 q_{ab}(x)=-\tfrac13C_{acbd}x^cx^d
          -\tfrac16C_{acbd;e}x^cx^dx^e.
\]
Direct differentiation, the algebraic symmetries, and the differential
Bianchi identity recover the prescribed curvature and first curvature
derivative. The Ricci tensor of this polynomial is affine and has zero
value and derivative, so its free Einstein tensor vanishes identically.
The two homogeneous terms realize the value and derivative separately.
The checker verifies the reconstruction on all ten value basis tensors
and all 24 derivative basis vectors. Complexification is unnecessary for
this jet statement.
\end{proof}

\section{All conserved algebraic Weyl-square currents are evaluated}
\label{sec:v42-current}

A current in the declared class is
\begin{equation}
 J_{ab}(C)=J_{ba}(C)
   =\sum_{I\le J}j_{(ab),IJ}w_Iw_J.
 \label{eq:v42-current-class}
\end{equation}
There are \(10\cdot55=550\) independent real coefficients. They are
arbitrary constants: no invariant contraction ansatz is imposed.
No derivatives of \(C\), Ricci factors, or additional metric factors
are included in this class. The last restriction matters in applying the
result to a general cubic BV coefficient.

Substitute \(\partial_aw_I=Z_{(a,I),u}z_u\) in
\(\partial^aJ_{ab}\). Coefficients of \(w_Iz_u\), with four
choices of \(b\), give an integer matrix
\[
 \mathsf M:\R^{550}\longrightarrow\R^{960}.
\]
Explicitly, a column labelled by \(((a,b),(I,J))\) contributes
\(Z_{(a,I),u}\) to row \((b,J,u)\), and
\(Z_{(a,J),u}\) to row \((b,I,u)\); for \(a\ne b\) add the
same two entries with \(a,b\) interchanged. Coincident terms are added.
This is a complete coefficient specification of the matrix.

Let \({}^\star C_{ab cd}=\frac12\epsilon_{ab}{}^{ef}C_{efcd}\).
On the present Euclidean coefficient chart define
\begin{equation}
 \mathbb B_{abcd}(C)
 =C_{aecf}C_{bedf}-{}^\star C_{aecf}\,{}^\star C_{bedf}.
 \label{eq:v42-BR}
\end{equation}
The minus sign belongs to this continuation. It must not be replaced by a
Lorentzian positivity convention; no positive-energy assertion is made.
The polynomial \(\mathbb B\) is completely symmetric and trace-free.
For a constant symmetric trace-free tensor \(T\), put
\begin{equation}
 J^T_{ab}=T^{cd}\mathbb B_{abcd}(C).
 \label{eq:v42-BR-current}
\end{equation}

\begin{theorem}[Complete coefficient classification of the declared current block]
\label{thm:v42-current-classification}
The following conditions on \eqref{eq:v42-current-class} are equivalent:
\begin{enumerate}[label=(\roman*)]
\item \(\partial^aJ_{ab}=0\) on every local solution of the native free
      gravitational Euler equation;
\item \(\mathsf Mj=0\);
\item there is a unique constant trace-free symmetric \(T\) such that
      \(J=J^T\).
\end{enumerate}
In particular,
\begin{equation}
 \boxed{\operatorname{rank}\mathsf M=541,
 \qquad\dim\ker\mathsf M=9.}
 \label{eq:v42-current-rank}
\end{equation}
\end{theorem}
\begin{proof}
The jet realization lemma makes (i) and (ii) equivalent, rather than merely
making (ii) a sufficient formal test. Substitution of the ten Weyl tables
in \eqref{eq:v42-BR} verifies all symmetry and trace identities.
Contraction against the nine basis tensors below produces a rational
\(550\times9\) matrix \(\mathsf B\) with
\(\mathsf M\mathsf B=0\). A displayed rational decoder proves that
these nine columns are independent.

For the converse, the certificate supplies 541 row and column indices of
\(\mathsf M\). Their integer square minor has determinant
\[
 \det\mathsf M_{\mathrm{minor}}\equiv41419\pmod{65521}.
\]
The modulus is prime and the residue is nonzero. Hence the determinant is
nonzero over \(\mathbb Q\), proving rank at least 541. The nine
independent rational null columns give rank at most 541. This is an exact
rank proof; no floating-point threshold or probabilistic rank claim is
used. It proves (ii) implies (iii), including uniqueness.
\end{proof}

\subsection{A small coefficient decoder}

Use parameter order
\[
 (T_{00},T_{01},T_{02},T_{03},T_{11},T_{12},T_{13},T_{22},T_{23}),
 \qquad T_{33}=-T_{00}-T_{11}-T_{22}.
\]
For an off-diagonal parameter the symmetric tensor has both corresponding
entries equal to that parameter. Flatten \(j\) in lexicographic
\((a,b)\), \(a\le b\), followed by lexicographic \((I,J)\),
\(I\le J\), with indices starting at zero. Select entries
\[
 \mathcal A=(0,1,3,4,6,7,9,10,13).
\]
If \(z=(j_i)_{i\in\mathcal A}\), the parameter vector is
\(\tau=\mathsf Dz\), where
\begingroup\small
\begin{equation}
\mathsf D=
\begin{pmatrix}
0&0&1/8&0&0&0&0&1/4&0\\
0&0&0&0&0&0&1/4&0&0\\
0&0&0&0&0&-1/8&0&0&0\\
0&0&0&0&1/4&0&0&0&0\\
1/2&0&-3/8&0&0&0&0&-1/4&0\\
0&1/4&0&0&0&0&0&0&0\\
0&0&0&0&0&0&0&0&1/4\\
-1/2&0&1/8&0&0&0&0&1/4&0\\
0&0&0&-1/4&0&0&0&0&0
\end{pmatrix}.
\label{eq:v42-decoder}
\end{equation}
\endgroup
The equation \(\mathsf D\mathsf B_{\mathcal A,:}=I_9\) is exact.
In the stated coefficient norms,
\begin{equation}
 \boxed{\|\mathsf D\|_{1\to1}=1,
 \qquad\|\mathsf B\|_{1\to1}=224.}
 \label{eq:v42-current-norms}
\end{equation}
For an arbitrary current array, \(j-\mathsf B\mathsf D j_{\mathcal A}\)
is its explicitly retained nonconserved part for this section; it vanishes
if and only if the current is conserved. Its norm is at most
\(225\|j\|_1\). This decoder does not itself test exactness of a
complete cubic action/source coefficient.

\section{Every current has an explicit native local BV completion}
\label{sec:v42-BV}

Conservation on the free stationary surface is not substituted for an
off-shell identity. We now factor its defect through the actual Euler
operator. For a current \(J^T\), let \(N_b(w)\) be the 40-entry
row such that \(\partial^aJ^T_{ab}=N_b(w)v\) for an arbitrary
Weyl derivative \(v\). It is linear in \(w\).

Select the independent rows
\[
 \mathcal I=(0,1,2,3,4,5,6,7,8,9,10,11,13,14,15,16)
\]
of \(S\), with row order \((b,cd)\), \(c<d\), as above.
Let \(S_0=S_{\mathcal I,:}\), and choose its pivot columns
\[
 \mathcal J=(0,1,2,3,4,5,6,7,8,9,10,11,12,15,16,17).
\]
Define the \(40\times16\) matrix \(R\) by putting
\(\bigl((S_0)_{:,\mathcal J}\bigr)^{-1}\) in those rows and zero in the other rows.
Then \(S_0R=I_{16}\). The local coefficient is
\begin{equation}
 Y_{b\alpha}(w)=(N_b(w)R)_\alpha.
 \label{eq:v42-Y}
\end{equation}
For \(\alpha=(b_\alpha,c_\alpha,d_\alpha)\), let
\(\Lambda_\alpha^{aA}\) be the fixed rational coefficients in
\[
 \partial_{c_\alpha}S_{b_\alpha d_\alpha}
 -\partial_{d_\alpha}S_{b_\alpha c_\alpha}
 =\Lambda_\alpha^{aA}\partial_a e_A.
\]
Their normalization includes \(\mathcal E_{ab}=e_{(ab)}/(2-\delta_{ab})\).

\begin{proposition}[Off-shell Cotton factor and coefficient bound]
\label{prop:v42-Cotton-factor}
The actual current obeys
\begin{equation}
 \boxed{\partial^aJ^T_{ab}
 =Y_{b\alpha}(C)\Lambda_\alpha^{dA}\partial_d e_A.}
 \label{eq:v42-div-factor}
\end{equation}
All coefficients are fixed rational linear functions of \(T\).
In the specified coordinate bases,
\(\|Y\|_{\mathrm{coeff},1}\le280\|\tau\|_1\).
\end{proposition}
\begin{proof}
Every \(N_b(w)\) annihilates \(\ker S\), by the complete current
test. Since \(S_0\) has full row rank and the same kernel, it follows
that \(N_b=N_bRS_0\). This identity is also checked coefficient by
coefficient on all nine columns. Substitute \eqref{eq:v42-cotton} to
obtain \eqref{eq:v42-div-factor}. The sum of absolute coefficients in
\(Y\) for the nine unit parameters is, respectively,
\[
 232,\ 280,\ 280,\ 232,\ 240,\ 248,\ 248,\ 216,\ 216.
\]
The maximum gives the stated induced coefficient bound. This is a local
factorization through the Euler equation, not division by that equation.
\end{proof}

Define
\begin{align}
 \mathcal V_T(q)&=\int q_{ab}J^T_{ab}(C(q)),\nonumber\\
 \mathcal R_{T,A}(q,c)&=
 -2\Lambda_\alpha^{dA}
       \partial_d\bigl(c^bY_{b\alpha}(C(q))\bigr),\nonumber\\
 \mathcal B_T&=\mathcal V_T+\int\eta_A\mathcal R_{T,A}.
 \label{eq:v42-cocycle}
\end{align}
Repeated tensor indices in \(\mathcal V_T\) run over all 16 ordered
pairs; the source term uses the ten independent coordinates. The free
antibracket convention is \(\delta_0\eta_A=e_A\).

\begin{theorem}[Local cubic BV cocycle for every conserved current]
\label{thm:v42-cocycle}
For each trace-free symmetric \(T\),
\begin{equation}
 \boxed{s_0\mathcal B_T=0\quad\text{modulo a divergence}.}
 \label{eq:v42-master}
\end{equation}
The action part has three metric fields and four derivatives; the source
part has one metric, one ghost, one metric antifield, and three derivatives.
No ghost-antifield term is required at this free-cocycle stage. The local
source coefficient has the conservative bound
\(\|\mathcal R_T\|_{\mathrm{coeff},1}\le1024\|\tau\|_1\).
\end{theorem}
\begin{proof}
Since \(\gamma_0C=0\), the only longitudinal variation of the action
comes from its explicit \(q\):
\[
 \gamma_0\mathcal V_T
 =-2\int c^bY_{b\alpha}\Lambda_\alpha^{dA}\partial_de_A
 =2\int e_A\Lambda_\alpha^{dA}\partial_d(c^bY_{b\alpha}).
\]
This cancels \(\delta_0\int\eta\mathcal R_T\) by the displayed
choice of sign. The remaining longitudinal source variation is zero:
\(Y\) is linear in the gauge-invariant Weyl tensor and
\(\gamma_0c=0\). There are no other antifield-number terms.

For the bound, each selected Cotton row has coefficient norm at most
\(2/3\) as a polynomial in a derivative of the independent Euler
coordinates. The product derivative in \(\mathcal R_T\) has two
Leibniz terms. Hence its coefficient norm is at most
\(4(2/3)280\|\tau\|_1<1024\|\tau\|_1\).
This counts the complete source polynomial, including derivative-of-ghost
and derivative-of-curvature pieces. The proof uses no inverse spacetime
operator. It proves only the free cubic cocycle, not its extension to a
nonlinear interacting symmetry.
\end{proof}

These formulas provide complete test cocycles in the native tangent
complex. They are not inserted into the active action. Their nonlinear
quartic extensions, if sought, would be separate deformation equations.
Similarly, the classification of \(J(C,C)\) does not show that every
cubic cocycle can first be reduced to this current form: Ricci-dependent
terms, improvements, and other antifield representatives require their
own reduction.

\section{Eighteen tests separate the two cubic blocks}
\label{sec:v42-eighteen}

Let \(L_\nu=s_1n_\nu\) be V41's nine complete cubic cocycles and
\(\mathcal B_\nu\) the nine just constructed. Use V41's seed
\[
 p=(1,i,0,0),\quad t=(0,0,1,i),\quad -p-t,
\]
with its transverse, traceless polarizations
\[
 (0,0,1,-i)(0,0,1,-i)^T,\quad
 (1,-i,0,0)(1,-i,0,0)^T,\quad
 (p-t)(p-t)^T.
\]
All evaluations use the ghost-free trilinear coefficient divided by 64.
Their normalization is the same native free-action convention as V41.
V41's annihilator lemma applies to every panel below: a free-exact cubic
has a free Euler factor on an external leg, while total derivatives vanish
by momentum conservation. The evaluation therefore annihilates all
\(s_0Q_3\), without deleting possible antifield partners beforehand.

Keep the original nine rows \(\ell\) unchanged. Add nine rows
\(m\), with the following matrices applied simultaneously to every
momentum and polarization:
\begin{center}\small
\begin{tabular}{cll}
\toprule
Added row & Matrix & Component\\\midrule
1&\(O_{01}\)&real\\
2&\(O_{01}\)&imaginary\\
3&\(O_{02}\)&imaginary\\
4&\(O_{03}\)&real\\
5&\(O_{12}\)&real\\
6&\(O_{13}\)&real\\
7&\(O_{13}\)&imaginary\\
8&\(P_{(0,2,1,3)}\)&imaginary\\
9&\(O_{03}\operatorname{diag}(-1,1,1,1)\)&imaginary\\\bottomrule
\end{tabular}
\end{center}
The plane matrices \(O_{ab}\) have V41's rational block
\(\left(\begin{smallmatrix}3/5&4/5\\-4/5&3/5\end{smallmatrix}\right)\).
No averaging or assumption of rotational invariance is involved.

Form the block matrix
\begin{equation}
 \mathsf A_{18}
 =\begin{pmatrix}C&B_0\\G&B_1\end{pmatrix}
 =\begin{pmatrix}\ell(L)&\ell(\mathcal B)\\
                         m(L)&m(\mathcal B)\end{pmatrix}.
 \label{eq:v42-augmented}
\end{equation}
Here \(C\) is precisely V41's already specified invertible matrix.
The new current entries are calculated from \eqref{eq:v42-BR}; the old
columns are recalculated from the native stationary Cartan/cofactor cubic,
not an unrelated vertex.

\begin{theorem}[Two independent cubic blocks and a faithful augmented test]
\label{thm:v42-eighteen}
The matrix in \eqref{eq:v42-augmented} has
\begin{equation}
 \boxed{
 \det\mathsf A_{18}=\frac{2^{45}3^5\,11}{5^{17}},
 \qquad \|\mathsf A_{18}^{-1}\|_{1\to1}=\frac{75}{8}.}
 \label{eq:v42-augmented-det}
\end{equation}
Consequently the 18 classes
\([L_1],\ldots,[L_9],[\mathcal B_1],\ldots,[\mathcal B_9]\)
are independent in the full native free cubic group
\(H^0_{A=3,w=4}(s_0\mid\partial)\). In particular,
\begin{equation}
 \boxed{\dim H^0_{A=3,w=4}(s_0\mid\partial)\ge18.}
 \label{eq:v42-lowerbound}
\end{equation}
No assertion that 18 is the full dimension is made.
\end{theorem}
\begin{proof}
Each current cocycle is closed by \cref{thm:v42-cocycle}; the old ones
are closed by V41. In the nine unit-parameter current polynomials, polarize
\(C(q)^2\), contract with the remaining explicit metric leg, and sum
over the three choices of that leg. This gives every new trilinear entry.
Re-evaluate each old column by the same Cartan inverse, cofactor variation,
and stationary-section formula used in V41.

The certificate supplies all 324 rational entries and the complete inverse.
Both inverse products are verified over \(\mathbb Q\). For a smaller
independent check, eliminate the old columns. The resulting matrix is
\(D=B_1-GC^{-1}B_0\), explicitly displayed below. Its determinant is
\(2^{30}3^2\,11/5^{11}\). Multiplication by V41's determinant gives
\eqref{eq:v42-augmented-det}; the inverse column sums give its norm.
Since all 18 rows annihilate \(\operatorname{im}s_0\), a linear
combination of the 18 cocycles that is free-exact has zero coefficient
vector. This proves independence. It does not test an arbitrary cubic
cohomology class outside their span.
\end{proof}

\Needspace{15\baselineskip}
The additional obstruction matrix is
\begingroup\small\setlength{\arraycolsep}{2.6pt}
\begin{equation}
D=\begin{pmatrix}
-12/25&-8/25&0&0&32/25&0&0&4/5&8/5\\
-16/25&-44/25&0&0&-24/25&0&0&-8/5&4/5\\
8/5&0&0&8/5&8/5&-8/5&0&0&0\\
-44/25&0&8/5&-48/25&-32/25&0&8/5&-12/25&0\\
4/5&0&-8/5&0&32/25&48/25&-8/5&-12/25&0\\
0&88/25&0&0&0&0&0&0&88/25\\
8/5&0&0&-8/5&8/5&8/5&0&0&0\\
0&2&2&0&0&0&2&0&2\\
0&-2/5&0&0&-2&0&16/5&-2&2
\end{pmatrix}.
\label{eq:v42-schur-matrix}
\end{equation}
\endgroup
Its inverse has exact norms
\begin{equation}
 \boxed{\|D^{-1}\|_{1\to1}=\frac{575}{192},
 \qquad\|D^{-1}\|_{\infty\to\infty}=\frac{2297}{1056}.}
 \label{eq:v42-schur-norms}
\end{equation}

\begin{corollary}[A nonexact native cubic invisible to every V41 row]
\label{cor:v42-invisible}
Let \(T_{03}=T_{30}=1\) and all other entries vanish. Its complete
cocycle \(\mathcal B_T\) satisfies
\begin{equation}
 \boxed{\ell(\mathcal B_T)=0,
 \qquad m_3(\mathcal B_T)=\frac85\ne0.}
 \label{eq:v42-invisible}
\end{equation}
It therefore is not free-BV exact despite passing all nine old tests.
\end{corollary}
\begin{proof}
It is the fourth current column in the recorded parameter order. Its
first nine entries vanish and its added \(O_{02}\)-imaginary entry is
\(8/5\). The latter annihilates all free-exact cubics. No numerical
integration, limit, or rank tolerance enters this certificate.
\end{proof}

\section{The fixed-normalization comparison gains nine independent necessary tests}
\label{sec:v42-comparison}

Apply the construction only to a correctly formed closed difference of two
solutions of the same reference equation. Let \(k_2\) pass V40's
quadratic test, and let
\[
 r_3=k_3-s_1\mathsf Sk_2
\]
be V41's remaining cubic. Set
\begin{equation}
 a_0=C^{-1}\ell(r_3),\qquad
 b=D^{-1}\bigl(m(r_3)-Ga_0\bigr),\qquad
 a=a_0-C^{-1}B_0b.
 \label{eq:v42-compiler}
\end{equation}
This is the block solution
\((a,b)=\mathsf A_{18}^{-1}(\ell(r_3),m(r_3))\).

\begin{theorem}[A stronger rejection test, with the unsampled remainder retained]
\label{thm:v42-rejection}
If
\[
 r_3\in\operatorname{im}s_0+\operatorname{span}\{L_1,\ldots,L_9\},
\]
then necessarily \(b=0\). When \(b=0\), the old ambiguity
coordinates remain exactly \(a_0=C^{-1}\ell(r_3)\); the final condition
is still
\begin{equation}
 \boxed{r_3-\sum_\nu(a_0)_\nu L_\nu\in\operatorname{im}s_0.}
 \label{eq:v42-still-open}
\end{equation}
The added zero test is necessary, not sufficient, for this equation.
For arbitrary input, the projection onto the 18 displayed cocycles
annihilates free-exact inputs and preserves the actual closure defect of
the unremoved coefficient. Its extracted coordinates satisfy
\[
 \|(a,b)\|_1\le\tfrac{75}{8}\|(\ell(r_3),m(r_3))\|_1.
\]
\end{theorem}
\begin{proof}
All tests vanish on free-exact terms. On an old cocycle combination the
test vector is \(\mathsf A_{18}(a,0)\), whose solution has zero
second block. Conversely, setting \(b=0\) in the block equations
leaves exactly the old coefficients; subtracting them gives
\eqref{eq:v42-still-open}, which cannot be decided by finitely many
rows that have not been proved exhaustive. The 18-cocycle projection is
idempotent because its test matrix is invertible. Its range consists of
closed coefficients, so subtracting it cannot change \(s_0r_3\).
The norm follows from the verified inverse.
\end{proof}

A nonzero \(b\) is not an additional adjustable primitive parameter.
V40 classified the entire quadratic primitive ambiguity: 205 directions
are already free-exact and nine are the \(n_\nu\). Their nonlinear
images cannot remove a current obstruction detected here. Higher-arity
primitives cannot change a cubic coefficient either. Changing the
normalization itself would be a different, explicitly recorded operation;
no such change is made here.

If a full populated comparison eventually passes, its matter compensation
remains the one in V39--V41, including the actual stress, scalar mass,
current, and higher-source contacts. The newly constructed current cocycles
are test directions, not new physical Wilson coefficients. No matter
counterterm is inserted merely because a test direction has been found.

The local coefficient maps have no mesh, volume, lower-reference, or
shell-count dependence. At a fixed external window of scale \(\mu\),
the current action and its source partner have the ordinary rooted-kernel
bounds \(C_b\|\tau\|_1\mu^4\) and
\(C_b\|\tau\|_1\mu^3\), respectively: multiply their polynomial
symbols by the fixed smooth external windows and use integration by parts
in the two relative four-dimensional coordinates. Constants depend on
those windows and the prescribed moment order. These estimates are local
polynomial bounds, not a new shell-summability or continuum theorem.

\section{Verification and closing ledgers for V42}
\label{sec:v42-ledgers}

The checker reconstructs the ten integer Weyl tables, the derivative
constraint, and the full 960-by-550 conservation matrix. Its exact rank
certificate combines a nonzero integer minor modulo the prime 65521 with
nine independent rational null columns; the modular calculation is a
proof over a finite field, not a numerical rank threshold. It separately
checks the normal-coordinate realizations, the complete rational current
map and decoder, the off-shell Cotton factorization, and the native
Euler/Cotton normalization. The action--antifield cancellation is tested
on independent off-shell polynomial panels in addition to the general
factorization proof.

The full run recomputes all 324 on-shell entries using the native cubic
and the new current vertices. It checks the 18-by-18 inverse and the
9-by-9 Schur inverse in rational arithmetic, including the old-test
invisible example. The public current compiler consumes the 550-component
array in the recorded basis. The cubic-test compiler consumes 18 rational
evaluations and reports the additional obstruction coordinates, with an
explicit warning that a zero result is not full cubic exactness.

\begingroup
\renewcommand{\refname}{Additional comparison source for V42}
\begin{thebibliography}{99}
\bibitem[51]{V42Senovilla}
J.~M.~M.~Senovilla,
\emph{Super-energy tensors},
Classical and Quantum Gravity \textbf{17} (2000), 2799--2842,
arXiv:gr-qc/9906087,
doi:10.1088/0264-9381/17/14/313.
The Bel--Robinson and conservation framework is established literature.
The present rational current classification and native cubic separation
are verified independently; no Lorentzian positivity property is asserted
for the Euclidean-continuation polynomial used here.
\end{thebibliography}
\endgroup

\subsection*{Proved}
\addcontentsline{toc}{subsection}{V42: Proved}
The entire algebraic symmetric Weyl-square current space is classified:
its conservation matrix has rank 541 and a nine-dimensional kernel, with
explicit bounded parameter extraction. Every current has a finite,
pole-free off-shell Cotton factorization and a complete native cubic BV
partner. The nine resulting cubic classes are independent of V41's nine
classes, giving an exact lower bound of 18 for the cubic quotient.
An explicit nonexact current cocycle passes every old test. The additional
Schur test, its exact coefficient bounds, and its rejection criterion are
established. Neither this current space nor the 18 evaluations is asserted
to exhaust the full cubic complex.

\subsection*{Inherited}
\addcontentsline{toc}{subsection}{V42: Inherited}
V1--V41, their mathematical bodies, labels, and bibliographies are retained.
The unreduced free Euler/gauge matrices, normalized native stationary cubic,
continuation, nine quadratic primitive representatives, V40's complete
quadratic ambiguity theorem, V41's on-shell annihilator, V39's high-arity
tail result, and the full relative stress transport keep their original
scope. Selected input identities are rechecked, not all inherited proofs.

\subsection*{Literature input}
\addcontentsline{toc}{subsection}{V42: Literature input}
Bel--Robinson tensors and their Lorentzian super-energy interpretation are
comparison material \cite{V42Senovilla}. The relation between conserved
currents and higher-derivative cubic deformations is discussed in
\cite[Section 11]{V39BDGH}. No external completeness claim for all such
deformations is imported. The new current rank, BV factorization, and
18-class independence follow from the explicit native coefficient
constructions and certificates.

\subsection*{Conditional}
\addcontentsline{toc}{subsection}{V42: Conditional}
A supplied current-form cubic can be classified within the proved block.
A supplied correctly formed normalization difference can be rejected if
its new obstruction coordinates are nonzero. If all displayed coordinates
vanish, full free cubic exactness still requires a further calculation.
Transferring a successful comparison to fixed gravity and matter remains
conditional on that full equation and the previously stated input bounds.

\subsection*{Open}
\addcontentsline{toc}{subsection}{V42: Open}
The actual native critical arrays \(k_2,k_3,\xi_0,\xi_1\), the full
remaining cubic image equation, and the complete inhomogeneous relative
scalar--gravity primitive remain unresolved. A reduction of an arbitrary
cubic cocycle to the algebraic Weyl-current form has not been proved.
In particular, the new rank cannot be used to replace the entire cubic
cohomology by 18 parameters. Combined internal-gauge/chiral interfaces,
higher loops, lower-reference removal, and the invariant interacting ESM
class remain separate tasks.

\begin{center}\small
\begin{tabular}{>{\raggedright\arraybackslash}p{.24\textwidth}>{\raggedright\arraybackslash}p{.67\textwidth}}
\toprule
Variable & Scope in V42\\\midrule
Order & Ordinary free complex; ghost-number-zero cubic weight four.
No new quantum loop or shell coefficient.\\
Current restriction & All 550 constant component coefficients in symmetric
algebraic \(J(C,C)\), without extra curvature derivatives. Not the full
cubic action/source bank.\\
Fields and sources & Ten metric coordinates, four ghosts, and the explicit
metric-antifield partners; no gauge-fixed or TT replacement.\\
Cutoff and volume & Finite algebraic maps independent of mesh, periods,
lower scale, and iteration count. Kernel bounds are fixed-window local
polynomial bounds only.\\
Test momenta & Complex free solutions with exact conservation, used only as
polynomial annihilators; not Gaussian integration modes.\\
Norms & The stated monomial and tensor-coordinate norms in V40's normalized
free convention. Physical normalization retains its \(2/\chi\) conversion.
No uniformity in growing EFT, source, or derivative order.\\
Active prescription & Fixed gravity, V34 and V38 unchanged. The new
cocycles are test directions, not inserted interactions.\\
Computation & Current matrix, modular-minor/rational-kernel proof, Cotton
factors, and all 324 augmented entries computed. Full cubic matrix and
native inhomogeneous loop arrays not evaluated.\\\bottomrule
\end{tabular}
\end{center}


\clearpage
\part*{V43: Quartic lifting separates the conserved-current cubic sector}
\addcontentsline{toc}{part}{V43: Native quartic lifting and an extension-compatible cubic projection}

\section{The next equation, rather than another free cubic census}
\label{sec:v43-start}

Sections 1--331, their labels, and their bibliographies are preserved.
V42 constructs nine conserved Weyl-square currents and their complete cubic
BV partners. Those are nontrivial \emph{free} cocycles, independent of
V41's nine images. Their presence does not establish that they are possible
leading terms of a counterterm closed under the ordinary Einstein
differential. The next equation is now evaluated rather than postponed:
\begin{equation}
 s_0b_3=0,\qquad s_0b_4+s_1b_3=0.
 \label{eq:v43-lifting-equation}
\end{equation}
The quadratic coefficient has already been set to zero by an actual exact
subtraction, or is zero in the candidate under examination. A nonzero
quadratic coefficient would also contribute \(s_2b_2\) and may not be
omitted.

The new result is a rank-nine obstruction for the nine V42 currents in
\eqref{eq:v43-lifting-equation}. A constant translation ghost and three
free metric waves suffice to detect it. In contrast, V41's nine images
have explicit quartic lifts. In the eighteen-dimensional cohomological
subspace already constructed, the liftable part is exactly the old
nine-dimensional image. This statement does not assert completeness of
the eighteen-dimensional subspace in the full free cubic cohomology.

All calculations use the same unreduced, ordinary zero-mesh metric complex,
normalized native stationary Palatini cubic, and coframe-coordinate source.
No propagator, contour, reference profile, finite regulator, or physical
normalization is changed. The fixed gravitational prescription, V34's active
subcritical scalar prescription, and V38's full logarithmic packet remain
unchanged. The calculation is a local classical compatibility test, not a
new quantum shell or continuum theorem.

The connection between higher-derivative conserved currents, their source
partners, and subsequent deformation obstructions is part of the established
BV framework; see \cite[Sections 2 and 11]{V39BDGH}. The full generalized
symmetry classification \cite{V41AndersonTorre} is comparison material.
Neither result is used as a substitute for the native quartic matrix
computed below. In particular, a two-derivative uniqueness theorem is not
applied to the present four-derivative counterterm problem.

\section{A constant-ghost annihilator of every free-exact quartic}
\label{sec:v43-annihilator}

Let \(\mathscr C_r^g\) denote translation-invariant local integrated
coefficients of total arity \(r\), ghost number \(g\), and inherited
weight four, modulo total derivatives. All ten independent metric
coordinates, four ghosts, and their antifields are retained. Write
\[
 s_g=s_0+s_1+s_2+\cdots,
 \qquad s_0s_1+s_1s_0=0,
 \qquad s_0s_2+s_1^2+s_2s_0=0.
\]
These are identities of the ordinary native master action. They are not
claims about nilpotence of the completed finite filtered source.

Use the independent metric order \((00,01,02,03,11,12,13,22,23,33)\)
and the free matrices \(E,R\) of V40, with \(ER=0\). For momenta and
polarizations obeying
\begin{equation}
 p_1+p_2+p_3=0,\qquad
 E(p_j)\operatorname{col}h_j=0,
 \label{eq:v43-test-conditions}
\end{equation}
put \(q=\sum_{j=1}^3t_jh_j e^{ip_jx}\) and let the ghost be the
constant odd vector \(c=\vartheta v\). Define \(\mathscr E_{h,p;v}\)
on \(\mathscr C_4^1\) by taking its antifield-free
\([t_1t_2t_3\vartheta]\) zero-total-momentum coefficient. Each evaluation
below is a real or imaginary part of \(\mathscr E/64\).
This is a polynomial test; it does not insert complex waves or a constant
hard ghost into any Gaussian measure.

\begin{lemma}[Free-exact quartic terms vanish on the complete test]
\label{lem:v43-quartic-annihilator}
For every local \(b_4\in\mathscr C_4^0\),
\begin{equation}
 \mathscr E_{h,p;v}(s_0b_4)=0.
 \label{eq:v43-annihilation}
\end{equation}
The evaluation also vanishes on total derivatives. No restriction on the
antifield content of \(b_4\) is needed.
\end{lemma}
\begin{proof}
The longitudinal variation of its antifield-free action replaces a metric
by \(R_0c\), which is zero for a constant ghost. To produce an
antifield-free term from the complementary free variation, the only
possibility is an occurrence of one metric antifield, replaced by the
linear metric Euler covector. After polarization that factor is
\(E(p_j)\operatorname{col}h_j\) on one of the three external legs and
vanishes by \eqref{eq:v43-test-conditions}. A ghost-antifield replacement
still contains a metric antifield; terms containing more than one antifield
also retain an antifield. They vanish under the stated extraction, not by
being discarded from the functional. The constant ghost has zero momentum,
so every total derivative carries \(i(p_1+p_2+p_3)\) and vanishes.
\end{proof}

Choose nine such real tests and collect them in \(\mathscr Q_{\rm tr}\).
Define the cubic-to-test map
\begin{equation}
 \mathscr N b=\mathscr Q_{\rm tr}(s_1b).
 \label{eq:v43-N-definition}
\end{equation}
It annihilates every free-exact cubic: if \(b=s_0q_3\), then
\(s_1b=-s_0s_1q_3\) and the lemma applies. In particular,
\begin{equation}
 s_0b_4+s_1b_3=0\quad\Longrightarrow\quad\mathscr N b_3=0.
 \label{eq:v43-N-necessary}
\end{equation}
This necessary condition acts on the \emph{complete} cubic, including its
metric-antifield coefficient.

For a constant ghost the first nonlinear metric transformation is exactly
\(R_1(q,c)=\mathcal L_cq=c^\mu\partial_\mu q\) on the inherited
coframe chart. The \(R_0c\)-dependent chart terms vanish. The action
variation of any translation-invariant cubic density is a total derivative.
It follows that the part detected in \eqref{eq:v43-N-definition} is
precisely the native cubic Euler covector contracted with that cubic's
\(\eta q c\) source coefficient. This observation will avoid a separate
computation of the full quartic matrix without omitting a term in the test.

\section{The actual current source in three exact complex panels}
\label{sec:v43-native-source}

Recall, as inherited data rather than a new free-cocycle theorem,
\begin{equation}
 \mathcal B_T=\mathcal V_T+\int\eta_A\mathcal R_{T,A},\qquad
 \mathcal R_{T,A}=-2\Lambda_\alpha^{dA}
       \partial_d\bigl(c^bY^T_{b\alpha}(C(q))\bigr).
 \label{eq:v43-current-recall}
\end{equation}
The rational tables \(Y\) and \(\Lambda\) are exactly those of
\eqref{eq:v42-Y}--\eqref{eq:v42-cocycle}. The parameter order remains
\[
 \tau=(T_{00},T_{01},T_{02},T_{03},T_{11},T_{12},T_{13},T_{22},T_{23}),
 \qquad T_{33}=-T_{00}-T_{11}-T_{22}.
\]
Both symmetric off-diagonal entries equal the corresponding parameter.

For one plane wave and a constant ghost vector \(v\), the source symbol is
\begin{equation}
 Z_T(h,p;v)_A
 =-2i\Lambda_\alpha^{dA}p_dv^bY^T_{b\alpha}(C(h,p)).
 \label{eq:v43-Z-symbol}
\end{equation}
Let \(\mathscr V_3\) be V41's polarized \emph{native} stationary cubic
in the normalized free convention. By the preceding extraction,
\begin{equation}
 \mathscr E_{h,p;v}(s_1\mathcal B_T)
 =\sum_{j=1}^3\mathscr V_3
   (h_1,\ldots,\operatorname{mat}Z_T(h_j,p_j;v),\ldots,h_3;p_1,p_2,p_3).
 \label{eq:v43-native-test-formula}
\end{equation}
This includes the action and antifield parts of \(\mathcal B_T\).
The action part vanishes on the constant-translation variation, while its
source part need not vanish.

Use the seed
\begin{align}
 p_1&=(1,i,0,0)^T,&p_2&=(0,0,1,i)^T,&p_3&=-p_1-p_2,\\
 v_1&=(0,0,1,-i)^T,&v_2&=(1,-i,0,0)^T,&v_3&=p_1-p_2,\\
 h_j&=v_jv_j^T.&&&&
 \label{eq:v43-seed}
\end{align}
The three panels are the seed and its simultaneous permutations
\[
 O_0=I,\qquad O_1=P_{(0,1,3,2)},\qquad O_2=P_{(0,2,1,3)},
 \qquad P_\pi e_j=e_{\pi_j}.
\]
They replace \(p_j,h_j\) by \(O_ap_j,O_ah_jO_a^T\).
All nonzero components of the test momenta and polarizations have modulus
one. The ghost vectors used are coordinate vectors.

\begin{lemma}[Finite native reduction on the declared panels]
\label{lem:v43-panel-reduction}
For each of the three panels, each of the four coordinate ghost vectors,
and each of the nine unit current parameters,
\begin{equation}
 E(p_j)Z_T(h_j,p_j;v)=0,
 \qquad
 Z_T(h_j,p_j;v)
 =-i(v\cdot p_j)(p_j^TTp_j)\operatorname{col}h_j+R(p_j)\xi_j.
 \label{eq:v43-gauge-remainder}
\end{equation}
The native cubic is \(-64\) on each panel. Consequently,
\begin{equation}
 \boxed{\frac1{64}\mathscr E_{h,p;v}(s_1\mathcal B_T)
 =i\sum_{j=1}^3(v\cdot p_j)(p_j^TTp_j).}
 \label{eq:v43-scalar-test}
\end{equation}
The gauge remainder is asserted here on these specific test panels,
not as an unproved formula for arbitrary momentum.
\end{lemma}
\begin{proof}
Substitute the ten Weyl tables and the sixteen selected Cotton rows in
\eqref{eq:v43-Z-symbol}. This is a linear calculation in \(T\) and \(v\),
so the stated nine-by-four basis tests suffice on each leg. The following
explicit formula checks the gauge remainder without a null-momentum
inverse. Put
\[
 \Delta=\operatorname{mat}Z_T+i(v\cdot p)(p^TTp)h,
 \qquad
 \xi_i=\frac{\Delta_{ji}}{p_j}
       -\frac{p_i\Delta_{jj}}{2p_j^2},
\]
where \(j\) is the first nonzero coordinate of the fixed test momentum.
For these panels \(p_j\in\{1,-1,i,-i\}\). Direct rational substitution
gives \(p_i\xi_k+p_k\xi_i=\Delta_{ik}\) and \(EZ_T=0\) on every
listed leg. This division concerns finitely specified nonzero constants
in a proof certificate; it is not a momentum-dependent local counterterm.
The checker verifies all 324 vector instances of each identity.

The value \(-64\) is recomputed from the original twenty-four-variable
Cartan inverse and the Palatini cofactor variations, not assumed from a
chosen continuum three-point vertex. The inherited cubic Noether identity
implies that a pure-gauge insertion in \(\mathscr V_3\) vanishes when
its other two legs solve the free Euler equation: its compensating terms
contain those Euler covectors. Substitution of
\eqref{eq:v43-gauge-remainder} into \eqref{eq:v43-native-test-formula}
therefore gives \eqref{eq:v43-scalar-test}. Independently, all 108 complex
contractions are recomputed directly using the ten polarized cubic
gradients on each leg. This supplies a finite arithmetic certificate of
the displayed identities, rather than a numerical inference from a few
approximately null waves.
\end{proof}

\section{All nine current classes have independent quartic obstructions}
\label{sec:v43-rank}

The nine real evaluations \(\mathscr Q_{\rm tr}\) are specified by
\begin{center}\small
\begin{tabular}{ccll}
\toprule
Row & Panel & Constant ghost vector & Part of \(\mathscr E/64\)\\\midrule
1 & \(O_0\) & \(e_0\) & real\\
2 & \(O_0\) & \(e_0\) & imaginary\\
3 & \(O_0\) & \(e_2\) & real\\
4 & \(O_0\) & \(e_2\) & imaginary\\
5 & \(O_1\) & \(e_0\) & real\\
6 & \(O_1\) & \(e_0\) & imaginary\\
7 & \(O_1\) & \(e_2\) & real\\
8 & \(O_2\) & \(e_0\) & real\\
9 & \(O_2\) & \(e_0\) & imaginary\\\bottomrule
\end{tabular}
\end{center}

\begin{theorem}[Faithful first nonlinear obstruction of the current block]
\label{thm:v43-injection}
For the inherited nine current cocycles,
\begin{equation}
 \boxed{\mathscr N\mathcal B_T=\mathsf D_{\rm tr}\tau,}
 \qquad
 \mathsf D_{\rm tr}=
 \begin{pmatrix}
 0&0&0&2&0&2&0&0&2\\
 -1&0&-2&0&-1&0&2&-2&0\\
 0&2&0&2&0&2&0&0&0\\
 -1&0&-2&0&1&0&2&0&0\\
 0&0&2&0&0&0&2&0&2\\
 1&0&0&-2&1&2&0&2&0\\
 1&0&0&2&-1&-2&0&0&0\\
 0&0&0&2&0&2&2&0&0\\
 -1&-2&0&0&-2&0&0&-1&2
 \end{pmatrix}.
 \label{eq:v43-matrix}
\end{equation}
It satisfies
\begin{equation}
 \boxed{\det\mathsf D_{\rm tr}=-3072,\qquad
 \|\mathsf D_{\rm tr}^{-1}\|_{1\to1}=\frac52.}
 \label{eq:v43-inverse-bound}
\end{equation}
The induced map
\[
 \operatorname{span}\{[\mathcal B_T]\}
 \longrightarrow H^1_{A=4,w=4}(s_0\mid\partial),
 \qquad [\mathcal B_T]\longmapsto[s_1\mathcal B_T],
\]
is injective. In particular, if a translation-invariant local \(b_4\)
satisfies \(s_0b_4+s_1\mathcal B_T=0\) modulo a divergence, then \(T=0\).
\end{theorem}
\begin{proof}
Equation \eqref{eq:v43-scalar-test} computes every matrix entry. For
example, the seed with \(v=e_0\) has real coefficient
\(2(T_{03}+T_{12}+T_{23})\) and imaginary coefficient
\(-T_{00}-2T_{02}-T_{11}+2T_{13}-2T_{22}\), giving the first two rows.
The two permutations give the remaining displayed rows. The certificate
contains the rational inverse; direct multiplication on both sides gives
the identity, and its largest absolute column sum is \(5/2\).
Thus the determinant and norm are exact, not floating-point rank tests.

The target cochain is closed because \(s_0s_1\mathcal B_T=0\).
Changing the cubic representative by \(s_0q_3\) changes its image by
\(-s_0s_1q_3\). If the image of a current combination is free-exact,
\cref{lem:v43-quartic-annihilator} makes its nine tests zero. The invertible
matrix forces all nine current parameters to vanish.
\end{proof}

\begin{corollary}[An explicit failure invisible to V41's action tests]
\label{cor:v43-witness}
For \(T_{03}=T_{30}=1\) and all other components zero, V42's old
cubic-action tests vanish, but
\[
 \mathscr N\mathcal B_T=(2,0,2,0,0,-2,2,2,0)^T.
\]
On the seed and constant ghost \(e_0\), the unnormalized quartic
coefficient is exactly \(128\). No local quartic coefficient can cancel
it while keeping the quadratic term zero and this complete cubic fixed.
\end{corollary}
\begin{proof}
This is the fourth column of \eqref{eq:v43-matrix}; the assertion about the
old action tests is the inherited zero column in V42. The annihilator
lemma excludes cancellation by an arbitrary quartic, including its ghost
and antifield parts.
\end{proof}

\begin{proposition}[Quantitative local separation]
\label{prop:v43-separation}
For any local quartic \(b_4\), let
\(\omega=s_1\mathcal B_T+s_0b_4\). Then
\begin{equation}
 \|\mathscr Q_{\rm tr}\omega\|_1
 =\|\mathsf D_{\rm tr}\tau\|_1
 \ge\frac25\|\tau\|_1.
 \label{eq:v43-probe-separation}
\end{equation}
In the raw translation-invariant local monomial norm, with independent
metric coordinates, ordinary derivatives on individual arguments, and no
Taylor-factorial normalization,
\begin{equation}
 \boxed{\|\omega\|_{\mathrm{coeff},1}\ge\frac{64}{135}\|\tau\|_1.}
 \label{eq:v43-raw-separation}
\end{equation}
The estimate holds for every representative of the integrated density.
\end{proposition}
\begin{proof}
The first bound follows from \eqref{eq:v43-inverse-bound} and the
annihilator lemma. In the antifield-free \(q^3c\) slot, each raw
monomial contributes at most \(3!=6\) in modulus to a labelled test:
the test entries and momenta have modulus at most one, and derivatives
of the constant ghost vanish. The sum of nine \(1/64\)-normalized
real evaluations therefore has norm at most \(54/64=27/32\).
Other antifield components do not increase that evaluation. Dividing
\(2/5\) by \(27/32\) gives \(64/135\). Total derivatives have zero
evaluation, so the lower bound persists for every density representative
and for the corresponding quotient seminorm.
\end{proof}

These inequalities are in fixed dimensionless coefficient coordinates.
Rescaling every test momentum by \(\mu\) multiplies this
\(q^3c\) row by \(\mu^5\). They are not estimates in a cutoff-dependent
quantum measure.

\section{An eighteen-class projection compatible with quartic lifting}
\label{sec:v43-projector}

Retain V41's old cubic-action evaluations \(\ell\), their matrix
\(\mathsf C\), and V42's current block
\(B_0=(\ell_a(\mathcal B_\nu))\). Let \(L_\nu=s_1n_\nu\).
The ordinary master identities imply
\begin{equation}
 s_0(s_2n_\nu)+s_1L_\nu=0,
 \qquad \mathscr N L_\nu=0.
 \label{eq:v43-old-lifts}
\end{equation}
Thus the old nine classes have actual quartic lifts. Their entire higher
extension is the exact local germ \(s_gn_\nu\), whose quadratic term
vanishes because \(s_0n_\nu=0\).

\begin{theorem}[Complete lifting result on the displayed eighteen classes]
\label{thm:v43-eighteen-lifts}
If
\[
 b_3=\sum_\nu a_\nu L_\nu+\mathcal B_T+s_0q_3,
\]
then a local solution of \(s_0b_4+s_1b_3=0\) exists if and only if
\(T=0\). Consequently the liftable part of V42's specified
18-dimensional free cubic subspace is exactly the nine-dimensional
image of V40's quadratic primitive classes.
\end{theorem}
\begin{proof}
Applying \(\mathscr N\) gives \(\mathsf D_{\rm tr}\tau\), so necessity
follows from \cref{thm:v43-injection}. If \(T=0\), an explicit lift is
\[
 b_4=\sum_\nu a_\nu s_2n_\nu+s_1q_3.
\]
Its free variation cancels \(s_1b_3\) by the two ordinary master
identities. In fact it is the fourth-degree restriction of the exact germ
\(s_g(\sum a_\nu n_\nu+q_3)\).
\end{proof}

For an arbitrary real cubic coefficient \(r\), define
\begin{equation}
 \begin{aligned}
 b(r)&=\mathsf D_{\rm tr}^{-1}\mathscr N r,\\
 a(r)&=\mathsf C^{-1}\bigl(\ell(r)-B_0b(r)\bigr),\\
 \mathscr P_{18}^{\mathrm{lift}}r
 &=\sum_\nu a_\nu(r)L_\nu+\sum_\nu b_\nu(r)\mathcal B_\nu.
 \end{aligned}
 \label{eq:v43-projector}
\end{equation}
This uses the source-sensitive quartic tests instead of V42's second bank
of action-only cubic evaluations. No prior coefficient is altered by
changing a diagnostic map.

\begin{theorem}[Bounded projection preserving the quartic-lifting domain]
\label{thm:v43-compatible-projection}
The map \(\mathscr P_{18}^{\mathrm{lift}}\) is idempotent on the
18-dimensional displayed subspace and annihilates free-exact cubics. Its
coefficient extraction obeys
\begin{equation}
 \boxed{\|(a(r),b(r))\|_1
 \le\frac{31}{4}\|(\ell(r),\mathscr N r)\|_1.}
 \label{eq:v43-joint-norm}
\end{equation}
For arbitrary input, subtraction preserves its free closure defect and
annihilates both test banks. If \(r\) has a quartic lift, then
\(b(r)=0\), its projected cubic has the explicit lift in
\eqref{eq:v43-old-lifts}, and the unprojected remainder also has a
quartic lift. No converse based solely on zero tests is asserted.
\end{theorem}
\begin{proof}
The combined test matrix is block triangular:
\begin{equation}
 \mathsf A_{43}
 =\begin{pmatrix}\mathsf C&B_0\\0&\mathsf D_{\rm tr}\end{pmatrix},
 \qquad
 \det\mathsf A_{43}=-\frac{2717908992}{15625}.
 \label{eq:v43-block-matrix}
\end{equation}
Its inverse is given by the two equations in \eqref{eq:v43-projector}.
The exact column sum in the inherited coefficient bases is \(31/4\),
verified by rational multiplication and summation. The inverse identities
prove idempotence and test annihilation. Both test banks kill free-exact
inputs. Every element of the projection's range is free closed, so the
remainder has the same \(s_0\)-defect as \(r\).

If \(s_0r_4+s_1r=0\), the annihilator gives \(\mathscr N r=0\), hence
\(b(r)=0\). Subtract
\(\sum_\nu a_\nu(r)s_2n_\nu\) from \(r_4\). Equation
\eqref{eq:v43-old-lifts} shows that the result is a lift of
\(r-\mathscr P_{18}^{\mathrm{lift}}r\). This proves preservation of
the lifting domain without classifying its complement.
\end{proof}

One must not apply the zero-current conclusion to the coefficients from
\emph{an arbitrary previous splitting}. A complementary unknown cubic
class may have nonzero \(\mathscr N\)-evaluation and may cancel that
of the current portion in that splitting. The conclusion concerns the
whole actual cubic and the explicitly defined projection
\eqref{eq:v43-projector}. In particular, this theorem does not set V42's
old nine additional coordinates to zero for a general extendible cubic.

\section{Use in the fixed-gravity comparison and remaining equations}
\label{sec:v43-comparison}

Let \(k=H_g-\widetilde H_g\) be an actual closed critical difference
through arity five. Assume its quadratic coefficient passes V40's exactness
test and let \(Q_2^0=\mathsf S k_2\). Form the complete exact subtraction,
not merely its cubic component:
\begin{align}
 r_3&=k_3-s_1Q_2^0,\\
 r_4&=k_4-s_2Q_2^0,\\
 r_5&=k_5-s_3Q_2^0.
 \label{eq:v43-comparison-residual}
\end{align}
Since the quadratic remainder is zero, closure implies
\[
 s_0r_3=0,\qquad s_0r_4+s_1r_3=0.
\]
Hence \(\mathscr N r_3=0\) is an exact diagnostic of this correctly
formed comparison. A nonzero value would reject the asserted closed
comparison or identify a missing coefficient, not license an extra
normalization parameter.

\begin{corollary}[The known current sector adds no surviving comparison freedom]
\label{cor:v43-comparison}
For \eqref{eq:v43-comparison-residual}, put
\(a=\mathsf C^{-1}\ell(r_3)\). Subtract the full exact germ
\(s_g\sum a_\nu n_\nu\) from the residual. Its cubic component has
\[
 s_0r_3^\perp=0,\qquad
 \ell(r_3^\perp)=0,\qquad
 \mathscr N r_3^\perp=0,
\]
and still possesses a quartic lift. If the cubic class is independently
known to belong to the eighteen-dimensional subspace constructed in V42,
then \(r_3^\perp\) is free-exact. Without that additional reduction,
free exactness of \(r_3^\perp\) remains a separate equation.
\end{corollary}
\begin{proof}
Apply \cref{thm:v43-compatible-projection}. The final conditional statement
uses \cref{thm:v43-eighteen-lifts}, not a completeness claim for the full
cubic quotient.
\end{proof}

This is the specific upstream gain. The eighteen independent \emph{free}
classes do not represent eighteen usable independent normalization
parameters in a closed Einstein comparison. In their computed subspace,
the nine current classes are removed by the next equation; the other nine
are precisely the earlier primitive ambiguity and have an explicit exact
higher completion. An arbitrary cubic not reduced to that subspace is not
being declared solved.

The scalar compensation of any successful comparison remains
\(\operatorname{pr}_{A\le5}tQ\), with the native stress, scalar mass,
current, and antifield descendants retained as in V39--V41. Since this
installment determines no populated \(k_2,k_3\) and inserts no generator,
there is no active scalar or gravitational counterterm update. The
inhomogeneous equations \(Dh_0=\xi_0\) and
\(Dh_1=\xi_1-Mh_0\) are not consequences of the new homogeneous
lifting calculation.

\section{Verification and closing ledgers for V43}
\label{sec:v43-ledgers}

The self-contained checker reconstructs the native free Euler matrix,
the 24-variable Cartan matrix, the two Palatini cofactor variations, the
inherited Weyl/Cotton source factors, and the new constant-ghost source
symbols. It recomputes the three native cubic gradients and all 108
complex translation-source contractions. It checks the 81 selected real
matrix entries, both rational inverse products, the coefficient norm,
and the explicit \(T_{03}\) witness. The eighteen-class block inverse is
checked independently from its two defining blocks.

The verification report counts scalar polynomial assertions as well as
matrix identities; it does not equate that count with independent
mathematical theorems. V42's classification rank and the V41 old action-test
theorem are inherited inputs. The prior native matrices needed by the new
contractions are rechecked, not every previous analytic or cohomological
proof. No new loop integral, full cubic/quartic cohomology matrix, native
critical reference array, or proof-assistant formalization is claimed.

The public utility accepts nine old action evaluations and nine new
quartic evaluations. It returns the rational coordinates
\eqref{eq:v43-projector} and reports whether these tests reject quartic
lifting. It expressly does not report full exactness from zero tests.

\subsection*{Proved}
\addcontentsline{toc}{subsection}{V43: Proved}
A complete constant-translation, three-free-metric evaluation annihilates
all free-exact quartic terms. The native current source and stationary
Einstein cubic give an invertible rank-nine quartic obstruction matrix,
with exact inverse norm \(5/2\). No nonzero combination of V42's nine
current classes has a local quartic lift with zero quadratic part.
The explicit coefficient separation bound is \(64/135\). Within the
specified eighteen-dimensional cubic subspace, the liftable part is
exactly V41's nine images, and an explicit projection with coordinate
norm \(31/4\) preserves quartic liftability. The old-test-invisible
\(T_{03}\) current has an exact nonzero constant-ghost witness.

\subsection*{Inherited}
\addcontentsline{toc}{subsection}{V43: Inherited}
V1--V42, their mathematical bodies, labels, and bibliographies are retained.
The ordinary native master action, its arity identities, unreduced free
Euler matrix, stationary Palatini cubic, coframe chart, V42 current
classification and source factors, V41 quadratic-generator representatives
and action tests, V40 quadratic section, and V39 high-arity tail comparison
keep their original hypotheses. No completed-regulator quantum identity
is inferred from ordinary coefficient nilpotence.

\subsection*{Literature input}
\addcontentsline{toc}{subsection}{V43: Literature input}
The BV deformation interpretation and higher-derivative current discussion
in \cite{V39BDGH}, and the full vacuum-symmetry classification
\cite{V41AndersonTorre}, provide context only. The first obstruction order,
all displayed rank statements, and the projection bound are established
by the native formulas and finite exact certificate. No external
completeness theorem replaces the unresolved cubic image equation.

\subsection*{Conditional}
\addcontentsline{toc}{subsection}{V43: Conditional}
For a correctly formed closed native comparison, its cubic remainder
necessarily passes the new quartic tests. If that remainder is further
reduced to the eighteen-class subspace, its primitive ambiguity is fully
resolved by the old nine-class subtraction. General comparisons still
require the complementary cubic image equation. Matter stress transport
is applied only after an actual complete comparison primitive has been
constructed.

\subsection*{Open}
\addcontentsline{toc}{subsection}{V43: Open}
The populated critical reference boundary \(k_2,k_3,\xi_0,\xi_1\), the
full remaining cubic image, and the inhomogeneous relative scalar--gravity
primitive remain unresolved. The new test does not prove that the full
cubic quotient has dimension eighteen, nor that its entire quartic-liftable
part has dimension nine. Mixed internal-gauge/chiral interfaces,
higher-loop source control, and a realized invariant interacting ESM
EFT class remain distinct tasks. The next useful computation is the
actual normalized boundary, or an exhaustive reduction of its
\emph{quartic-liftable} cubic remainder; another unconstrained free
current census alone would not settle it.

\begin{center}\small
\begin{tabular}{>{\raggedright\arraybackslash}p{.24\textwidth}>{\raggedright\arraybackslash}p{.67\textwidth}}
\toprule
Variable & Scope in V43\\\midrule
Perturbative content & Ordinary classical critical coefficient complex;
cubic-to-quartic lifting. No new quantum loop coefficient.\\
Fields and sources & Ten metric coordinates, four diffeomorphism ghosts,
complete inherited current antifield partners. No TT or gauge-fixed
replacement of the operator complex.\\
Locality & Finite translation-invariant differential polynomials at
weight four; no inverse momentum or mass. Test-data gauge divisions
use fixed nonzero elements of \(\mathbb Q(i)\) only.\\
Cutoff and volume & Finite matrices and coefficient-map bounds have no
mesh, period, lower-reference, or shell-count dependence. No new
many-scale or cutoff estimate is asserted.\\
Order uniformity & Fixed current dimension, arities, derivative weight,
and coefficient norms. No bound uniform in increasing EFT or source
order is asserted.\\
Complex momenta & Algebraic annihilators with exact momentum conservation
and free Euler equations; not integration modes or physical on-shell
scattering claims.\\
Normalization & V40 normalized free convention; physical action
normalization retains the inherited \(2/\chi\) conversion.\\
Active prescription & Fixed gravity, V34, and V38 unchanged. Cocycles and
projections are tests, not newly selected Wilson coefficients.\\\bottomrule
\end{tabular}
\end{center}


\clearpage
\part*{V44: A finite critical gravitational jet has a formal local extension}
\addcontentsline{toc}{part}{V44: Obstruction-degree gap and fixed gravitational normalization}

\section{The missing extension question is at ghost number one}
\label{sec:v44-start}

Sections 1--338 and their labels are preserved. V40--V43 distinguish free
bilinear and cubic classes from the classes that lift through a further
Einstein equation. They do not classify the entire cubic complement. The
present argument goes upstream of that classification. A difference of two
actual critical gravitational primitives is not merely a free cubic
cocycle: it is closed under the \emph{ordinary interacting} differential
through total arity five. We ask whether that finite jet extends to all
arities before trying to classify its individual homogeneous pieces.

The relevant extension obstruction has ghost number one, whereas V39's
high-arity exactness theorem has ghost number zero. Conflating these two
statements would leave a gap. We prove the missing free ghost-one vanishing
at arities at least six by an explicit descent-degree argument. It implies
formal local extension of every critical gravitational five-jet. The
established full pure-Einstein local BRST classification then reduces its
weight-four ghost-zero class to curvature squares, which are Euler-exact
or local total derivatives in four dimensions. It follows that the entire
\emph{closed interacting five-jet}, not every unrestricted free cubic, is
exact.

This resolves the fixed gravitational normalization interface in the
relative matter problem. It does not evaluate a native annular loop array,
and it does not yet construct the absolute critical scalar--gravity
primitive. The distinction between these two tasks is retained. V25's
fixed gravitational normalization, V34's active subcritical scalar
normalization, and V38's logarithmic packet are not changed. All identities
below concern the ordinary local coefficient complex; none asserts
nilpotence of the finite filtered source or a convergent field series.

\subsection{The precise local category and external inputs}

Work on the contractible flat coframe chart with zero cosmological and
independent Holst coefficients, using the normalized ten-component metric
complex of V40. Put
\[
 s_g=s_0+s_1+s_2+\cdots,
 \qquad A(s_jF)=A(F)+j,
 \qquad s_0=\delta_0+\gamma_0.
\]
The ghost and antifields are \(c,\eta,\sigma\), and
\[
 w(q)=0,\quad w(c)=-1,\quad w(\eta)=2,\quad w(\sigma)=3,
 \quad w(\partial)=1.
\]
An integrated density is identified modulo \emph{translation-invariant}
local total derivatives. Explicit coordinate polynomials are not allowed
in counterterms. All component tensors may initially be arbitrary; no
rotation average of an unknown reference coefficient is performed.

The formal completion below means a series in undifferentiated fields,
with polynomial dependence on a bounded collection of derivative and
antifield jets at the fixed weight. It does not mean an analytic norm
limit. At ghost number zero and weight four,
\begin{equation}
 d+n_*=4,
 \qquad n_*=\#\eta+\#\sigma,
 \label{eq:v44-finite-jets}
\end{equation}
so derivative and antifield degree remain bounded even when ordinary field
arity grows. The corresponding ghost-one bound is \(d+n_*=5\).

We use the following established algebraic ingredients, with their roles
kept separate from the new degree argument:
\begin{enumerate}[label=(\roman*)]
\item Free longitudinal cohomology is generated by linearized curvatures,
  antifields and their derivatives, and \(c_\mu\) and
  \(z_{\mu\nu}=\partial_{[\mu}c_{\nu]}\). Ghost derivatives of order
  two and higher are paired jets. This is the reduction used in V20 and
  V39, from \cite[Section 3]{V39BDGH}.
\item In positive antifield number the invariant algebraic Poincare and
  Koszul--Tate reductions remove representatives above antifield number
  two. A nontrivial highest antifield-two coefficient is represented by
  an undifferentiated ghost antifield with constant coefficients. The
  argument in \cite[Sections 4 and Appendix A]{V39BDGH} is an expansion
  in the independent ghost monomials; the elimination does not depend on
  their total ghost number being two. It therefore applies also to the
  ghost-one cocycles used below, whose highest antifield-two term has
  pure ghost number three. These are literature inputs, not new native
  matrix ranks.
\item In the formal, finite-derivative local category the full ordinary
  four-dimensional pure-Einstein ghost-zero cohomology is represented by
  covariant curvature densities; its antifields are cohomologically
  removable. We use \cite[Theorems 1--3]{V25BBHEinstein}, with the
  formal-series convention and normal-theory homological construction
  of \cite[Sections 4 and 10]{V20BBH1994}. The local Lorentz and
  nonminimal quotient is the inherited
  Lemma~\ref{lem:v25-quotient}. The argument below spells out the weight-four
  and translation-invariant restrictions rather than importing a
  finite-truncation conclusion from that classification.
\end{enumerate}
The new proofs in Sections~\ref{sec:v44-lowforms}--\ref{sec:v44-jetexact}
are literature-assisted in precisely these senses. There is no assertion
that the two-derivative interaction classification in the same references
classifies the present four-derivative free vertices.

\section{Two low-form facts that remove the long descents}
\label{sec:v44-lowforms}

Let \(\mathscr R\) denote the linearized Riemann tensor. The form degree
and the weight of a \emph{coefficient} are written separately; in this
section \(dx^\mu\) is not included in the coefficient weight. Thus
horizontal differentiation raises coefficient weight by one.

\begin{lemma}[Translation-invariant invariant forms in degrees zero and one]
\label{lem:v44-lowforms}
Let \(a\) be a polynomial in \(\mathscr R\) and its derivatives,
homogeneous of positive metric degree, with arbitrary spectator tensor
indices. If \(a\) is a horizontally closed zero-form, then \(a=0\).
If \(a=a_\mu dx^\mu\) is a horizontally closed one-form, there is a
translation-invariant curvature polynomial \(f\) such that
\(a=d_xf\). The primitive has the same metric degree and one less
derivative. No inverse momentum is evaluated.
\end{lemma}
\begin{proof}
Polarize the metric arguments and take their formal Fourier symbols.
For positive metric degree \(m\), the total momentum
\(P=p_1+\cdots+p_m\) is a collection of four algebraically independent
linear coordinates after an invertible linear change of momentum
variables. A closed zero-form obeys \(P_\mu a=0\); multiplication by
\(P_\mu\) is injective in the coefficient polynomial module, so
\(a=0\).

For a closed one-form,
\(P_\mu a_\nu-P_\nu a_\mu=0\). The relatively prime coordinates
\(P_0,P_1\) show that \(P_0\) divides \(a_0\), and the other equations
then give \(a_\mu=P_\mu f\) with a polynomial \(f\). Equivalently,
this is the degree-one exactness of the Koszul complex for the four
linear coordinates \(P_\mu\). The quotient polynomial inherits the
polarization symmetries and has one less derivative. This is exact
polynomial division, not a rational symbol evaluated at \(P=0\).

Because \(a\) is strictly invariant,
\(P_\mu\gamma_0 f=\gamma_0 a_\mu=0\). Multiplication by \(P_\mu\)
remains injective in the module with exterior ghost coefficients, so
\(\gamma_0f=0\). At ghost number zero this makes \(f\) a curvature
polynomial by input (i). No explicit coordinate is introduced. Returning
to the jets proves the lemma, with spectator indices carried throughout.
\end{proof}

For the longitudinal reduction, horizontal differentiation induces
\[
 \mathscr D=\mathscr D_0+\mathscr D_1
\]
on the curvature/ghost representatives. Here \(\mathscr D_0\) is ordinary
differentiation of the curvature coefficients; \(\mathscr D_1\) sends an
undifferentiated ghost to its curl times \(dx\), up to the fixed
antisymmetrization convention, and sends a curl to zero. Symmetric first
and all higher ghost derivatives are \(\gamma_0\)-exact. The operator
\(\mathscr D_1\) raises the number of curl factors by one; this number is
bounded in the exterior ghost algebra.

\begin{lemma}[No nontrivial zero- or one-form descent bottom with positive metric degree]
\label{lem:v44-bottom}
Suppose a longitudinal descent has an antifield-free bottom of form
degree zero or one, with positive metric degree and fixed positive ghost
number. The bottom can be removed by translation-invariant
\(\gamma_0\)- and horizontal-exact terms, preserving total arity and
weight. The descent may consequently be shortened.
\end{lemma}
\begin{proof}
Choose a \(\gamma_0\)-closed representative
\(\sum_J\alpha_J(\mathscr R)\omega^J(c,z)\). The equation one level
above the bottom implies
\(\mathscr D\sum_J\alpha_J\omega^J=0\) in longitudinal cohomology.
Order the independent ghost monomials by their number of curls. At the
lowest nonzero curl count, \(\mathscr D_1\) cannot contribute from a
lower coefficient. Hence \(d_x\alpha_J=0\) coefficientwise.

In form degree zero, Lemma~\ref{lem:v44-lowforms} makes these positive-metric
coefficients zero. Repeating at the next curl count removes the whole
representative. In form degree one, write
\(\alpha_J=d_x\beta_J\) by that lemma. Subtracting
\(d_x(\beta_J\omega^J)\), and the corresponding
\(\gamma_0\)-exact correction to the ghost-jet representative, removes
this curl count and changes only higher curl counts. The process
terminates since there are finitely many exterior ghost monomials.

Thus the bottom is \(\gamma_0u+d_xv\). In the preceding descent
equation the \(d_xv\) term has zero derivative, and the
\(\gamma_0u\) term is removed by changing the preceding form by a
signed \(d_xu\). This shortens the descent. Every operation is
homogeneous; the primitive in the one-form step lowers the coefficient
weight by one, exactly as required for its horizontal derivative.
\end{proof}

\section{The free obstruction cohomology vanishes above arity five}
\label{sec:v44-gap}

\begin{theorem}[Native free ghost-one critical gap]
\label{thm:v44-gap}
For the translation-invariant ordinary free spin-two complex in four
dimensions,
\begin{equation}
 \boxed{H^{1,4}_{A=r,w=4}(s_0\mid d_x)=0
 \qquad(r\ge6).}
 \label{eq:v44-gap}
\end{equation}
Here the second superscript is horizontal form degree. More explicitly,
a local integrated ghost-one density of arity \(r\ge6\) and weight four
that is \(s_0\)-closed modulo a divergence has a translation-invariant
local primitive of ghost number zero, the same arity and the same weight.
No assertion is made about the group at arity five or below.
\end{theorem}
\begin{proof}
Use input (ii) to reduce the antifield expansion to degree at most two,
keeping homogeneous arity and weight. A nontrivial highest
antifield-two term is a constant ghost antifield times a ghost monomial
of pure ghost number three. It has arity four and cannot represent an
arity \(r\ge6\) class. Its trivial part is removed, with all lower
antifield terms generated by that removal retained.

The remaining antifield-one term may be made strictly
\(\gamma_0\)-closed. A representative has one metric antifield, two
surviving ghost factors, and \(r-3\) linearized curvatures. Its weight
is at least
\[
 2+2(r-3)-2=2r-6\ge6.
\]
Curl factors or additional derivatives only increase this bound. It
therefore has no representative at weight four and can also be removed.
This uses the invariant reduction, not the supposition that a generic
antifield-dependent density is already invariant.

The residual top form \(a_4\) is antifield-free and obeys
\(\gamma_0a_4+d_xa_3=0\). Descend until a nontrivial bottom \(a_p\),
\(0\le p\le4\), is reached. Positive-arity horizontal cohomology below
form degree four has no constant-form obstruction, so the descent is
available without explicit coordinates. Each coefficient satisfies
\[
 \operatorname{gh}(a_p)=5-p,
 \qquad w(a_p)=p,
 \qquad A(a_p)=r.
\]
At a strictly longitudinally closed bottom, input (i) gives curvature
coefficients and \(5-p\) ghost factors. Hence the number of metric
factors is \(m=r-5+p\).

For \(p\ge2\), even allowing every ghost its most negative weight,
\[
 w(a_p)\ge 2(r-5+p)-(5-p)=2r-15+3p.
\]
For \(r\ge6\) and \(p=2,3,4\), this lower bound exceeds \(p\).
At \(r=6\) the three lower bounds are respectively \(3,6,9\), while
the required weights are \(2,3,4\). Thus no such bottom exists.

At \(p=0\) or \(p=1\), the metric degree is positive:
\(m=r-5\ge1\) or \(m=r-4\ge2\). The bottom is removable by
Lemma~\ref{lem:v44-bottom}. Shortening the descent either reaches a higher
bottom, already excluded, or makes the original form trivial. Reversing
the local reductions gives the required primitive. All the linear
operations, divisions by positive integer homogeneities and polynomial
Koszul sections preserve arity and weight; there is no division by a
physical wave operator or mass.
\end{proof}

The argument is a ghost-one theorem. It is not a repetition of V39's
\(H^0\) tail theorem. It also does not erase V43's explicitly nonzero
arity-four obstruction: the present gap begins at arity six. The bound
six is sufficient for the next construction; no optimal threshold is
claimed.

\section{Every closed critical five-jet extends formally, without changing that jet}
\label{sec:v44-extension}

\begin{theorem}[Formal local extension from the retained gravitational bank]
\label{thm:v44-extension}
Let
\[
 k_{\le N}=\sum_{r=2}^{N}k_r,
 \qquad N\ge5,\qquad
 \operatorname{gh}(k_r)=0,\quad w(k_r)=4,
\]
and suppose
\(\operatorname{pr}_{A\le N}s_gk_{\le N}=0\) modulo local total
derivatives. There is a formal finite-derivative local germ
\(k_\infty=\sum_{r\ge2}k_r\) with
\begin{equation}
 \boxed{s_gk_\infty=0\ \bmod d_x,
 \qquad J_{\le N}k_\infty=k_{\le N}.}
 \label{eq:v44-extension}
\end{equation}
The extension leaves every supplied coefficient through arity \(N\)
unchanged. It is not asserted to converge in an analytic field norm.
\end{theorem}
\begin{proof}
Assume coefficients through \(r-1\) have been chosen with closure through
that arity. The next known obstruction is
\begin{equation}
 O_r=\sum_{j=1}^{r-2}s_jk_{r-j},
 \qquad s_0k_r=-O_r.
 \label{eq:v44-recursion}
\end{equation}
The arity-\(r\) component of
\(s_g^2k_{<r}=0\), using closure at every lower arity, gives
\(s_0O_r=0\) modulo a divergence. Its ghost number is one, its weight
is four, and \(r\ge N+1\ge6\). Apply \cref{thm:v44-gap} to choose
\(k_r\) solving the second equation. Induction constructs all orders.
At any specified arity only finitely many terms of
\eqref{eq:v44-recursion} occur. Equation~\eqref{eq:v44-finite-jets}
bounds all derivative and antifield degrees of the resulting coefficients,
so their formal sum has bounded derivative order and is local in the
formal-series sense. No estimate on the growth of its coefficients as
\(r\to\infty\) is needed or inferred.
\end{proof}

The theorem supplies precisely the missing premise in an unrestricted
truncation argument. In a different filtered complex a degree-zero
arity-five element may have a nonzero arity-six differential; if that
arity-six class has no free primitive, truncation creates a false cycle.
The counterexamples of V25 and V37 remain valid. They do not satisfy the
native ghost-one gap just proved.

\section{The closed interacting critical gravitational bank is exact}
\label{sec:v44-jetexact}

\begin{lemma}[Full ordinary pure-gravity exactness at weight four]
\label{lem:v44-full-exact}
In the formal finite-derivative local category on the contractible
four-dimensional ordinary Einstein chart, every translation-invariant
integrated ghost-zero cocycle of weight four is \(s_g\)-exact modulo a
translation-invariant local divergence.
\end{lemma}
\begin{proof}
Use input (iii) and the inherited local Lorentz/nonminimal quotient. The
classification is applied to the \emph{full formal germ}, not to an
arbitrary truncated array. Its local reductions use the invertible metric
and the regular principal Euler coordinates; they are valid coefficientwise
in the formal neighborhood, as in the normal-theory construction of
\cite{V20BBH1994}. No convergence of that formal series is claimed.

The representative at weight four is a linear combination of curvature
squares and a twice-differentiated curvature. The latter reduces, by
contraction and Bianchi identities, to a divergence such as
\(\sqrt g\nabla^2R\). The parity-even algebraic squares are
\(\sqrt g R^2\), \(\sqrt g R_{ab}R^{ab}\) and
\(\sqrt g R_{abcd}R^{abcd}\). The four-dimensional Euler combination
reduces the last one to the first two modulo the local Euler
Chern--Simons current. The parity-odd Pontryagin density is likewise a
local total derivative. These currents use the local connection and no
explicit coordinate; no global topological coefficient is being removed.

For clarity the remaining Euler-exact step can be written explicitly.
Let the signed, fixed native normalization be
\(\mathcal E_g^{ab}=c_E\sqrt g\,G^{ab}\), \(c_E\ne0\).
Then the local covariant metric variation
\begin{equation}
 F_{ab}=\frac{\beta}{c_E}R_{ab}
        -\frac{\alpha+\beta/2}{c_E}g_{ab}R
 \label{eq:v44-Ricci-lift}
\end{equation}
satisfies
\[
 \mathcal E_g^{ab}F_{ab}
 =\sqrt g\left(\alpha R^2+\beta R_{ab}R^{ab}\right).
\]
In the source convention of V38,
\(s_g\int\eta_g^{ab}F_{ab}\) equals this contraction: its source
remainder \(\mathcal L_cF-DF[\mathcal L_cg]\) vanishes because \(F\)
is a natural tensor. The cotangent coframe conversion is retained.
This proves exactness with weight-four primitives.

Translation invariance is not an extra averaging assumption. In the
full classification the additional constant-horizontal-form classes
come from zero-form BRST classes. At total ghost/form degree four the
only positive-ghost source is the ordinary four-ghost volume class,
whose fourth descendant is the covariant density above; a constant
positive-degree horizontal form cannot multiply a lower positive-ghost
class, since those classes are absent. The constant ghost-zero
four-form itself has coefficient weight zero, not four. The local frame
chart removes the separate global frame-topology classes. Thus the
primitive and current may be chosen with no explicit \(x\), in the
same category as the input.
\end{proof}

\begin{theorem}[Exactness of the complete closed critical jet]
\label{thm:v44-jetexact}
Let \(d_N=\operatorname{pr}_{A\le N}s_g\) on integrated
translation-invariant local coefficients of weight four. For every
\(N\ge5\),
\begin{equation}
 \boxed{\ker(d_N:\mathfrak L^0_{N,4}\to\mathfrak L^1_{N,4})
 =\operatorname{im}(d_N:\mathfrak L^{-1}_{N,4}\to
                         \mathfrak L^0_{N,4}).}
 \label{eq:v44-finite-exact}
\end{equation}
In particular every closed pure-gravity critical normalization difference
through arity five has a primitive through that same arity, including its
bilinear and cubic coefficients.
\end{theorem}
\begin{proof}
Extend a closed finite coefficient by \cref{thm:v44-extension}.
Apply Lemma~\ref{lem:v44-full-exact} to the resulting full formal germ,
obtaining \(k_\infty=s_gQ_\infty\) modulo a local divergence.
The ordinary differential and horizontal differentiation never lower
arity, so projection gives
\(k_{\le N}=d_NJ_{\le N}Q_\infty\). Each retained coefficient is a
finite differential polynomial. The reverse inclusion follows from
\(d_N^2=0\). There is no nonzero integrated arity-one weight-four
term: all derivatives on its single factor give a divergence. It can
therefore be omitted without changing the argument.
\end{proof}

This theorem concerns \emph{interacting closure through five}, not free
closure of a bilinear or cubic in isolation. V40's 43 free bilinear
classes and V42's free current classes remain nonzero. The theorem says
that a nontrivial combination not removable by an ordinary exact germ
must fail one of the retained nonlinear closure equations. It does not
claim that V43's quartic-liftable cubic subspace alone is exhausted by
nine classes: closure through the quintic equation is also used here.

\section{A native finite solver and what its bounds do not claim}
\label{sec:v44-solver}

The exactness theorem eliminates a qualitative image question; it does
not evaluate a reference array or a huge inverse. It has a finite linear
implementation with its actual residual retained.

Choose any rational integrated coefficient bases for the three weight-four
spaces of ghost numbers \(-1,0,1\) and arity at most \(N\). Let
\(A_N,B_N\) be the two actual \(d_N\) matrices. Integration by parts is
imposed exactly, or equivalently current columns and a fixed rational
quotient are retained. Then \(B_NA_N=0\), and
\cref{thm:v44-jetexact} makes
\[
 L_N=A_NA_N^T+B_N^TB_N
\]
positive definite. The two summands have zero product, because
\(B_NA_N=0\), so each commutes with \(L_N^{-1}\). Consequently
\begin{equation}
 \boxed{I=A_NH_N+Q_NB_N,
 \qquad H_N=A_N^TL_N^{-1},\quad Q_N=L_N^{-1}B_N^T.}
 \label{eq:v44-section}
\end{equation}
This is the inherited finite Hodge-section mechanism, now justified on
the \emph{entire} critical interacting bank rather than only its tail.
For a nonclosed array it returns the residual \(Q_NB_Nk\), not a false
primitive.

For an explicit arithmetic bound, stack
\(Z_N=\binom{A_N^T}{B_N}\), with shape \(m\times n\), clear denominators by
an integer \(D_N\), and bound the absolute entries of \(D_NZ_N\) by
\(J_N\ge1\). Put \(L=\max(1,\sqrt{mn}J_N)\). A nonzero full-column
integer minor and Cauchy--Binet give
\begin{equation}
 \|L_N^{-1}\|_2\le D_N^2L^{2(n-1)},\qquad
 \|H_N\|_2\le\|A_N\|_2D_N^2L^{2(n-1)}.
 \label{eq:v44-height}
\end{equation}
The bound has no mesh, period or reference-scale parameter at fixed bank.
It is not a claim of practical conditioning or an evaluated numerical
constant for the unassembled full matrix.

\subsection{Using the already computed low-arity sections}

A more economical prescription consumes V40--V43 rather than discarding
those calculations. For a closed five-jet \(k=k_2+k_3+k_4+k_5\), use
V40's section \(\mathsf S\), V41's representatives \(n_\nu\), its tests
\(\ell\) and invertible matrix \(\mathsf C\). Set
\begin{align}
 r_3&=k_3-s_1\mathsf Sk_2,
 & a&=\mathsf C^{-1}\ell(r_3),\nonumber\\
 Q_2&=\mathsf Sk_2+\sum_\nu a_\nu n_\nu,
 &\widehat r_3&=r_3-\sum_\nu a_\nu s_1n_\nu.
 \label{eq:v44-low-solver}
\end{align}
\begin{corollary}[The remaining cubic range condition is automatic for a closed five-jet]
\label{cor:v44-cubic}
For the input above, \(\Pi k_2=0\) in V40's quotient and
\(\widehat r_3\in\operatorname{im}s_0\) in the full free cubic
complex. Thus one may choose
\begin{equation}
 \begin{aligned}
 Q_3&=(s_0|_{\mathfrak L^{-1}_3})^+\widehat r_3,\\
 Q_4&=H_4\left(k_4-s_2Q_2-s_1Q_3\right),\\
 Q_5&=H_5\left(k_5-s_3Q_2-s_2Q_3-s_1Q_4\right),
 \end{aligned}
 \label{eq:v44-recursive-solver}
\end{equation}
where \(H_4,H_5\) are the tail sections of V39. These choices satisfy
\(d_5(Q_2+Q_3+Q_4+Q_5)=k\).
\end{corollary}
\begin{proof}
By \cref{thm:v44-jetexact}, some complete primitive \(\overline Q\)
exists. Its quadratic equation implies \(\Pi k_2=0\). V40 decomposes
\(\overline Q_2-\mathsf Sk_2\) into a free-exact generator plus the
nine \(n_\nu\). The free-exact term contributes a free-exact cubic.
V41's evaluations annihilate that image and determine the nine coefficients
uniquely, so they equal \(a\) in \eqref{eq:v44-low-solver}. Hence
\(\widehat r_3\) is in the asserted image. After its primitive is
chosen, closure of \(k\) makes each tail right-hand side free-closed.
V39 supplies the last two sections. Directly extracting the arity-two
through arity-five coefficients verifies the stated identity.
\end{proof}

The cubic Moore--Penrose map in \eqref{eq:v44-recursive-solver} is the
rational finite section on its image, not a propagator. Its membership
premise is now proved for closed five-jets. Its native entries and rank
are not reported as evaluated. V40's bound \(145/2\) and V41's bound
\(25/8\) remain available for the first two maps, while the remaining
fixed matrices admit the arithmetic bound used in V25. Arbitrary
free-closed cubics, or arrays that only pass V43's selected evaluations,
are not inputs to this corollary unless their full five-jet closure is
also supplied.

\section{The fixed gravitational primitive no longer requires an independent relative test}
\label{sec:v44-relative}

The application uses the actual ordinary reference equation; it does not
alter a prescription merely to simplify its boundary.

\begin{theorem}[A full formal lift of the fixed finite gravity prescription]
\label{thm:v44-fixed-lift}
Let \(x_g\) be V25's extendable ordinary critical reference, and let
\(H_g\) be its fixed arity-five primitive in the stated normalized
coefficient convention. There is a formal finite-derivative ordinary
primitive \(H_g^{\mathrm{full}}\) of the full reference germ whose
arity-five jet is \emph{exactly} \(H_g\).
\end{theorem}
\begin{proof}
V25 supplies a full ordinary primitive \(\widetilde H_g^{\mathrm{full}}\)
of that reference. The actual difference
\(k=H_g-J_{\le5}\widetilde H_g^{\mathrm{full}}\) is closed under
\(d_5\). By \cref{thm:v44-jetexact}, write \(k=d_5Q_g\) with
\(Q_g\) of arity at most five. Then
\[
 H_g^{\mathrm{full}}
 =\widetilde H_g^{\mathrm{full}}+s_gQ_g
\]
has the same full reference differential and the prescribed finite jet.
The correction is the whole ordinary exact germ, including all source
partners. No finite coefficient of \(H_g\) is changed, and no numerical
value of it is presumed. This is a formal local lift, not an analytic
radius estimate for the fixed prescription.
\end{proof}

Now suppose an \emph{absolute} critical gravity--scalar primitive has
been constructed in the same ordinary local, mass-polynomial and
matter-count convention:
\[
 d_g\widetilde H_g=x_g,
 \qquad d_H\widetilde H_2=x_2-t\widetilde H_g.
\]
As in V37, the two gravity components must solve the \emph{same} actual
reference equation, with the same sector assignment. A scalar vacuum
contribution or a changed reference cannot be hidden in their difference.
The new theorem removes the remaining fixed-gravity qualification:
\begin{corollary}[Automatic fixed-gravity matching for an absolute primitive]
\label{cor:v44-relative}
For any such absolute pair, \(k_g=H_g-\widetilde H_g\) has a local
primitive \(k_g=d_5Q_g\), and
\begin{equation}
 \boxed{H_2=\widetilde H_2+
     \operatorname{pr}_{A\le5}tQ_g,
 \qquad d_HH_2=x_2-tH_g.}
 \label{eq:v44-relative}
\end{equation}
No additional matching assumption on the quadratic or cubic gravity
coefficients is required.
\end{corollary}
\begin{proof}
The same-reference hypothesis makes \(d_5k_g=0\), so apply
\cref{thm:v44-jetexact}. The native triangular relation
\(d_Ht+td_g=0\) then gives
\(d_H(tQ_g)=-tk_g\), proving \eqref{eq:v44-relative}. This is
exactly the stress/current lift of V37, now with its pure-gravity
exactness premise discharged for the complete retained critical bank.
\end{proof}

Every action and source descendant is retained. If
\(Q_g=Q_2+Q_3+Q_4+Q_5\), the correction has
\begin{equation}
 \begin{aligned}
 [tQ_g]_3&=t_{[1]}Q_2,\\
 [tQ_g]_4&=t_{[2]}Q_2+t_{[1]}Q_3,\\
 [tQ_g]_5&=t_{[3]}Q_2+t_{[2]}Q_3+t_{[1]}Q_4.
 \end{aligned}
 \label{eq:v44-matter}
\end{equation}
These are the same source families and the same actual scalar stress and
current already identified in V39--V41. There is no arity-two flat scalar
term. Since \(t(z)=t_0+zt_1\), one generator supplies both the massless
and first-mass correction; the quadratic-in-mass term is unchanged by
this matching. In a full four-matter completion, \(t^2=0\) makes the
corresponding four-matter consistency equation unchanged as well.

The scalar absolute primitive is \emph{not} constructed by this corollary.
V37's complete reference germ, including its four-matter part, is the
appropriate input to that separate problem. Matter fields of weight zero
permit additional high-arity coefficients, so the pure-gravity degree gap
must not be transferred to the enlarged complex by notation. An
application of the tensor-matter classification must still preserve the
stated mass polynomial, source and current category. What has closed is
the independent obstruction attributed to keeping V25's finite gravity
normalization fixed.

\section{Verification, uniformity, and the next gate}
\label{sec:v44-ledgers}

The standalone checker verifies the free gauge/Christoffel and curvature
identities through the derivative orders used in the proof, the
low-form polynomial division, the descent base cases and degree formulas, native
free Euler identities, the curvature-square Euler contraction, finite
extension recursions, the exact section with its nonclosed residual, and
the relative stress signs on finite complexes. It also audits byte-for-byte
preservation of the previous mathematical body. The checker is not a
proof of the imported BRST classifications, a numerical rank computation
of the full arity-five matrix, or an evaluation of a native reference
integral. Its supplied-system utility reports residuals and any harmonic
part instead of assuming the user's matrix has the proved native
exactness property.

\subsection*{Proved}
\addcontentsline{toc}{subsection}{V44: Proved}
The positive-metric zero/one-form descent bottoms are removable without
explicit coordinates. The native free ghost-one weight-four cohomology
vanishes at arities at least six. Every interacting critical gravitational
jet closed through arity \(N\ge5\) extends to a formal local germ with
that jet fixed. Combined with the identified full pure-gravity
classification and explicit curvature-square reduction, every such
closed jet is exact. The fixed V25 primitive has a full formal local
lift, and any absolute same-reference scalar--gravity primitive can be
converted to its gravitational normalization without an extra low-arity
image hypothesis. Fixed-order finite sections and their residuals are
specified; their large native matrices are not evaluated.

\subsection*{Inherited}
\addcontentsline{toc}{subsection}{V44: Inherited}
V1--V43 and all their labels are unchanged. The ordinary master action,
coframe cotangent map, V25 reference germ and primitive, V39 tail sections,
V40 quadratic section and quotient, V41 ambiguity solver, V43 distinction
between quartic lifting and free exactness, and V37 triangular stress lift
retain their stated conventions. Previously nonzero free classes and
finite lifting witnesses remain valid. No old local or analytic result
is silently rewritten as a theorem of the full Einstein--SM theory.

\subsection*{Literature input}
\addcontentsline{toc}{subsection}{V44: Literature input}
The free longitudinal and invariant Koszul--Tate reductions are from
\cite{V39BDGH,V20BBH1994}. The full ordinary four-dimensional pure-Einstein
classification is \cite{V25BBHEinstein}. Its finite-jet applicability is
\emph{not} imported; that is the new obstruction-gap and formal-extension
argument. The curvature-square Euler identity is written explicitly,
and topological densities are discarded only as local divergences on
the declared chart, not as global topological information.

\subsection*{Conditional}
\addcontentsline{toc}{subsection}{V44: Conditional}
The relative matter conversion consumes an absolute scalar--gravity
primitive in the same reference and mass-polynomial category. It does
not produce that primitive or a numerical list of its coefficients.
A computed finite linear section requires a specified integrated basis
and the actual matrix. The arithmetic section bound is uniform only at
the fixed finite bank; it supplies no analytic control of the formal
all-arity extension.

\subsection*{Open}
\addcontentsline{toc}{subsection}{V44: Open}
The remaining critical gate is the absolute ordinary gravity--scalar
reference primitive, with its complete matter-two/four, mass and
antifield data; the fixed finite gravity normalization is no longer an
independent qualitative obstruction. Explicit native coefficients and
practical compressed inverses remain uncomputed. The internal-gauge and
chiral source/measure sectors, higher-loop estimates, and a realized
invariant interacting Einstein--SM EFT class are still separate open
interfaces. No new cutoff rate, quantum measure, or all-order analytic
continuum is asserted.

\begin{center}\small
\begin{tabular}{>{\raggedright\arraybackslash}p{.25\textwidth}>{\raggedright\arraybackslash}p{.66\textwidth}}
\toprule
Variable & Scope in V44\\\midrule
Fields & Ordinary pure gravity: ten metric coordinates, four ghosts,
complete antifields. Scalar fields enter only the stated relative
conversion.\\
Weight and arity & Weight four. Free ghost-one gap for arity at least
six; closed interacting jets through \(N\ge5\). Fixed \(N=5\) is the
programme application.\\
Locality & Translation-invariant finite differential polynomials at
any finite arity. The all-arity extension is formal in undifferentiated
fields, not an asserted convergent analytic function.\\
Topology & Contractible coframe/frame chart and local functionals modulo
local currents. Global characteristic numbers and boundary observables
are not removed.\\
Cutoff and volume & Algebraic section bounds do not depend on mesh,
periods, reference scale or shell count at fixed bank. No new
cutoff-convergence estimate is established.\\
Perturbative order & A local ordinary BV comparison used by the existing
one-loop normalization. No new loop integration or higher-loop result.\\
Normalization & Same-reference comparisons only. Fixed gravity, V34 and
V38 are preserved; no scalar counterterm is silently installed.\\
Computation & Exact finite algebra checks; no computed rank or numerical
primitive for the complete native arity-five matrix.\\\bottomrule
\end{tabular}
\end{center}

\clearpage
\part*{V45: Translation-sensitive scalar currents and the primitive category}
\addcontentsline{toc}{part}{V45: Translation-sensitive scalar currents and the primitive category}

\section{The scalar theorem needs its own counterterm-category audit}
\label{sec:v45-start}

Sections 1--346 and all their labels are preserved. V44 removes the
independent fixed-gravity matching obstruction, conditional on an absolute
scalar--gravity primitive in the \emph{same} local category. The next step
cannot be obtained simply by replacing ``pure gravity'' by ``tensor
matter'' in that argument. The current local bank allows arbitrary constant
component tensors but prohibits explicit spacetime coordinates in both
counterterms and discarded currents. This restriction is relevant when
matter has nontrivial rigid currents.

The tensor-matter classification in \cite[Section 2]{V29BBHMatter}
defines local forms using spacetime forms whose coefficients may depend
on \(x\). Its descent proof uses the spacetime Poincare lemma on such
coefficient forms; see its Section 6, especially equations (6.10)--(6.12).
These are legitimate choices, but they are not a theorem that a primitive
can be selected without explicit \(x\). The source paper's classification
of four-dimensional Lorentz anomaly candidates therefore cannot by itself
settle the more restrictive native primitive problem. This is a difference
of cochain categories, not a disagreement with that classification.

Here that distinction is made concrete on the \emph{same ordinary scalar
action}. We construct 36 independent critical, mass-polynomial,
translation-invariant ghost-one classes before a flavour-singlet
restriction. Each has an elementary primitive linear in \(x\), and none
has a translation-invariant primitive. No internal gauge field or chiral
fermion is involved. These are candidate local breakings, not measured
anomalies. The actual current native reference is flavour invariant, and
we prove that this particular 36-dimensional module cannot occur in its
invariant cohomology component. Thus the example does not obstruct the
native singlet target; it identifies a hypothesis which must be checked,
not silently imported, when constructing its absolute primitive.

\subsection{Fixed carrier, signs, and coefficient category}

Work with the ordinary scalar--gravity master action already used in
V27--V44, in its Euclidean coefficient chart, at zero internal gauge,
Yukawa, and scalar self-interaction couplings. There are four real scalars
with common \(z=m^2\), and their scalar antifields are
\(\upsilon_A\), \(1\le A\le4\). Define
\[
 \rho=\sqrt{\det g},\qquad K^{\mu\nu}=\rho g^{\mu\nu},\qquad
 S_H=\frac12\int
 \left(K^{\mu\nu}\partial_\mu\phi_A\partial_\nu\phi_A
              +\rho z\phi_A^2\right).
\]
Repeated scalar and spacetime indices are summed. In the inherited
ordinary BV convention the relevant transformations are
\begin{equation}
 \begin{split}
 s\phi_A&=c^\mu\partial_\mu\phi_A,\\
 sc^\mu&=c^\nu\partial_\nu c^\mu,\\
 s\upsilon_A&=E_A+\partial_\mu(c^\mu\upsilon_A),\\
 E_A&=-\partial_\mu(K^{\mu\nu}\partial_\nu\phi_A)+\rho z\phi_A.
 \end{split}
 \label{eq:v45-scalar-convention}
\end{equation}
The metric transforms by its Lie derivative. The scalar antifield is a
density. These formulas are not a proposed transformation of the finite
filtered regulator. Algebraic nonminimal doublets are contracted before
the local cohomology test; no nonminimal or measure term is removed from
any measured quantum identity by that contraction.

Let \(\mathfrak L_{\rm TI}\) denote ordinary local integrated coefficient
functionals, formal on the existing coframe chart and polynomial in a
finite collection of jets and in \(z\), modulo
\emph{translation-invariant} total divergences. Counterterms in this
space contain no explicit \(x\). Let \(\mathfrak L_x\) denote the
corresponding comparison category in which polynomial \(x\)-dependence
is allowed. We do not assume the induced map on these integrated
quotients is injective. The examples below distinguish their cohomology
directly. The active category remains \(\mathfrak L_{\rm TI}\).

The relevant weights are
\[
 w(\phi)=0,\quad w(\upsilon)=2,\quad w(c)=-1,\quad
 w(\partial)=1,\quad w(z)=2.
\]
The external tensors introduced below label possible component
coefficients. They are not new dynamical fields or active couplings.

\section{A critical cocycle from each native scalar rotation current}
\label{sec:v45-current}

Fix a real antisymmetric flavour matrix \(R\), and put
\begin{equation}
 a_R=\upsilon_A R_{AB}\phi_B,
 \qquad
 j_R^\mu=K^{\mu\nu}(\partial_\nu\phi_A)R_{AB}\phi_B.
 \label{eq:v45-current}
\end{equation}
The antifield density \(a_R\) is odd and has ghost number minus one;
\(j_R\) is an even vector density. It is the current of an actual
rigid symmetry of this common-mass scalar action, not a symmetry imposed
on a different reference theory.

\begin{lemma}[Exact native current identities]
\label{lem:v45-current}
For every \(z\),
\begin{equation}
 \boxed{
 \partial_\mu j_R^\mu=-E_A R_{AB}\phi_B,\qquad
 sa_R=-\partial_\mu j_R^\mu+\partial_\mu(c^\mu a_R).
 }
 \label{eq:v45-noether}
\end{equation}
Moreover
\[
 sj_R^\mu=c^\nu\partial_\nu j_R^\mu
     -j_R^\nu\partial_\nu c^\mu+j_R^\mu\partial_\nu c^\nu.
\]
\end{lemma}
\begin{proof}
Differentiate \eqref{eq:v45-current}. The term with two scalar first
jets vanishes because \(K\) is symmetric and \(R\) antisymmetric. The
mass term vanishes because \(\phi_A R_{AB}\phi_B=0\). The remaining
term is minus the displayed Euler contraction. Applying
\eqref{eq:v45-scalar-convention} to the ordered odd product
\(\upsilon_A R_{AB}\phi_B\) gives its density Lie derivative and the
Euler term; the minus sign from the odd antifield is essential.
The formula for \(sj_R\) is the tensor-density transformation of the
actual expression \eqref{eq:v45-current}.
\end{proof}

Let \(B_{\mu\nu}=-B_{\nu\mu}\) be any constant spacetime two-form.
Define
\begin{equation}
 f_B=\frac12B_{\mu\nu}c^\mu c^\nu,
 \qquad (\alpha_B)_\mu=B_{\mu\nu}c^\nu,
 \qquad
 A_{B,R}=f_Ba_R+(\alpha_B)_\mu j_R^\mu.
 \label{eq:v45-density}
\end{equation}
All products have their displayed Grassmann order.

\begin{theorem}[An exact critical source-current cocycle]
\label{thm:v45-cocycle}
The full ordinary local functional
\begin{equation}
 \boxed{\mathcal A_{B,R}(z)=z^2\int A_{B,R}}
 \label{eq:v45-anomaly}
\end{equation}
has ghost number one and weight four and satisfies
\(s\mathcal A_{B,R}=0\) modulo a translation-invariant divergence.
More precisely,
\begin{equation}
 \boxed{
 s A_{B,R}
 =\partial_\mu\left(c^\mu A_{B,R}-f_B j_R^\mu\right).
 }
 \label{eq:v45-density-descent}
\end{equation}
It has matter count two. Its antifield term is of type
\(z^2\upsilon\phi cc\); its action-type term is
\(z^2 c\phi\partial\phi\), with its entire metric-density expansion.
\end{theorem}
\begin{proof}
The constant \(B\) obeys
\(sf_B=c^\mu\partial_\mu f_B\) and
\(s(\alpha_B)_\mu=c^\nu\partial_\nu(\alpha_B)_\mu\).
Using \eqref{eq:v45-noether},
\[
 s(f_Ba_R)=\partial_\mu(c^\mu f_Ba_R)
                       -f_B\partial_\mu j_R^\mu.
\]
Since \(\alpha_B\) is odd, its product with the even current instead
satisfies
\[
 s((\alpha_B)_\mu j_R^\mu)
 =\partial_\nu(c^\nu(\alpha_B)_\mu j_R^\mu)
                +(\alpha_B)_\mu j_R^\nu\partial_\nu c^\mu.
\]
Antisymmetry of \(B\) and the odd ghost order give
\((\alpha_B)_\mu\partial_\nu c^\mu=-\partial_\nu f_B\).
Adding the two identities proves \eqref{eq:v45-density-descent}.
The weights of \(f_Ba_R\) and \(\alpha_Bj_R\) are both zero, and
multiplication by \(z^2\) raises the weight to four. No equation of
motion is imposed on the density.
\end{proof}

The lowest action and source arities are three and four. Their higher
metric contacts are supplied by the same \(K^{\mu\nu}\), not by a
new chosen completion. The result is a full-germ identity; the
nontriviality claim below is not a claim that an arbitrary truncated
cochain complex has the same cohomology.

\section{Constant-ghost contraction detects nontriviality}
\label{sec:v45-contraction}

For each coordinate define the odd derivation \(\iota_\mu\) by
\[
 \iota_\mu c^\nu=\delta_\mu^\nu,
 \qquad \iota_\mu\partial_Ic^\nu=0\quad(|I|>0),
\]
and let it vanish on all other jet generators. It is extended by the
left graded Leibniz rule. It contracts a constant diffeomorphism ghost;
it does not set the other ghost derivatives to zero.

\begin{lemma}[Translation homotopy in the actual local category]
\label{lem:v45-translation}
On ordinary minimal local jet functions without explicit coordinates,
\begin{equation}
 \boxed{\{s,\iota_\mu\}=\partial_\mu.}
 \label{eq:v45-cartan}
\end{equation}
Consequently \(s\iota_\mu+\iota_\mu s=0\) on
\(\mathfrak L_{\rm TI}\), and
\(\iota_\nu\iota_\mu\) takes an exact integrated functional to an
exact integrated functional. The map also respects total divergences.
\end{lemma}
\begin{proof}
On a tensor field, the contraction of its Lie derivative selects its
ordinary derivative; derivatives of the ghost are unchanged. On a
diffeomorphism ghost, contraction of \(c\cdot\partial c\) selects
\(\partial_\mu c\). The antifield Euler term contains no ghost;
contraction of its tensor-density Lie derivative again gives
\(\partial_\mu\). The same argument applies to their prolongations.
Both sides are even derivations, so the identity holds on the full jet
algebra. A total derivative integrates to zero precisely in the stated
translation-invariant current convention. Contraction commutes with
horizontal differentiation, completing the proof. This is also the
operator identity used in \cite[Section 6]{V29BBHMatter}; its quotient
consequence is being used here with its explicit-coordinate hypothesis
retained.
\end{proof}

\begin{lemma}[The six scalar rotation charges are not trivial symmetries]
\label{lem:v45-charge}
The charges
\[
 Q_R=\int a_R
\]
are closed ordinary functionals of ghost number minus one. Their six
independent antisymmetric flavour directions are linearly independent
in \(H^{-1}(s\mid d_x)\). The same directions form a torsion-free
module over \(\mathbb R[z]\): a polynomial combination is exact only
when its six coefficient polynomials vanish.
\end{lemma}
\begin{proof}
Closure follows from \eqref{eq:v45-noether}. To check nontriviality,
extract the antifield-linear, ghost-free component of a putative
\(Q_R=sK\), with \(\operatorname{gh}K=-2\). The antifield-two,
ghost-free part of \(K\) can generate only a gauge transformation plus
an Euler-trivial transformation: this follows directly by varying a
ghost antifield, or one factor in a quadratic field-antifield term.
Terms of higher antifield number retain a ghost or an antifield at
this extraction. This is the standard Noether interpretation of
negative-ghost cohomology \cite{V20BBH1994}, with its elementary
implication used explicitly here.

The actual transformation is \(\delta_R\phi=R\phi\). Its nongauge
character has an explicit stationary test requiring no inverse equation.
Complexify a putative real local identity. Choose a nonzero null flavour
vector \(v\in\mathbb C^4\), \(v^Tv=0\), such that \(Rv\) is not
proportional to \(v\), and set
\[
 g_{\mu\nu}=\delta_{\mu\nu},\qquad
 \phi_A(x)=\varepsilon v_Ae^{ip\cdot x},\qquad p^2+z=0.
\]
The scalar equations vanish. Every scalar stress contraction contains
\(v^Tv\), so the metric equations vanish too, together with all their
prolongations. Every local diffeomorphism acts on \(\phi\) along the
single flavour line spanned by \(v\), whereas \(R\phi\) does not.
Lorentz gauge transformations act trivially on the scalar. Such a \(v\)
exists for every nonzero antisymmetric \(R\): if all twelve null lines
\(\mathbb C(e_A\pm i e_B)\) were invariant, every coordinate two-plane
would be invariant, forcing \(R\) to be diagonal and therefore zero.
This establishes the contradiction on exact complex stationary jets.
Complexification tests a real polynomial identity; no complex plane wave
is inserted into a quantum integral. At zero mass a constant scalar and
flat metric provide the simpler test.

For a polynomial combination in \(z\), choose a positive value at which
the combined antisymmetric matrix is nonzero. Evaluation at that mass
would give the forbidden exact fixed-mass charge. Hence every coefficient
polynomial must vanish. No inverse of \(z\) is used.
\end{proof}

\begin{theorem}[Thirty-six independent translation-sensitive critical classes]
\label{thm:v45-independent}
Before imposing flavour invariance, the map
\begin{equation}
 \boxed{
 \Lambda^2(\mathbb R^4)^*\otimes\mathfrak{so}(4)_{\rm fl}
 \longrightarrow H^{1,4}_{w=4}(s\mid d_x;\mathfrak L_{\rm TI}),
 \qquad B\otimes R\longmapsto[\mathcal A_{B,R}]
 }
 \label{eq:v45-injection}
\end{equation}
is injective. In particular, it exhibits a 36-dimensional subspace.
\end{theorem}
\begin{proof}
For \(\mu<\nu\), direct left contraction gives
\begin{equation}
 \boxed{
 \iota_\nu\iota_\mu\mathcal A_{B,R}
   =z^2B_{\mu\nu}Q_R.
 }
 \label{eq:v45-decoder}
\end{equation}
The current term has only one undifferentiated ghost and vanishes after
two contractions. If a linear combination of \(\mathcal A\)'s were
exact in \(\mathfrak L_{\rm TI}\), each contracted charge would be
exact by \cref{lem:v45-translation}. The torsion-free independence in
\cref{lem:v45-charge}, applied to each of the six spacetime pairs,
sets all 36 coefficients to zero.
\end{proof}

The 36 is a lower bound in the stated \emph{unrestricted flavour}
cohomology, not its full dimension and not a count of physical anomalies.
All coefficients are real, mass-polynomial and local. The construction
does not contradict the pure-gravity V44 result: it uses genuine scalar
rigid currents which that theorem does not contain.

\subsection{Coefficient-map bounds}

In the ordered monomial basis each individual
\(\iota_\nu\iota_\mu\) has coefficient norm at most one. A monomial
has at most four distinct undifferentiated ghosts, so the sum over the
six pairs has norm at most six. In the displayed 36-class parameter basis
the decoder \eqref{eq:v45-decoder} is the identity and has norm one.
These statements concern the coefficient maps, not evaluation on a
physical volume.

For a finite coframe Taylor order \(N\), let
\[
 \kappa_N=\max_\mu\sum_\nu
       \|J_{q\le N}K^{\mu\nu}\|_{\mathrm{coeff},1}.
\]
The raw density inclusion satisfies
\begin{equation}
 \|J_{q\le N}(\mathcal A_{B,R}/z^2)\|_{\mathrm{coeff},1}
 \le(2+4\kappa_N)\|B\|_1\|R\|_1,
 \label{eq:v45-inclusion-bound}
\end{equation}
where independent antisymmetric entries are used for both parameter
norms. This follows by counting the two antifield monomials and the two
current components, each containing two flavour terms. The constant is
finite at fixed \(N\) and does not depend on mesh or periods. The
coframe series is not assigned an all-arity norm by this estimate.

\section{The forbidden coordinate primitive can be written exactly}
\label{sec:v45-coordinate}

The preceding classes are trivial in a larger counterterm category.
This can be checked without invoking an anomaly-classification theorem.
Let
\[
 v_\mu(x)=\frac12x^\nu B_{\nu\mu},
 \qquad
 \partial_\nu v_\mu-\partial_\mu v_\nu=B_{\nu\mu}.
\]
The letter \(v_\mu\) here denotes an external coefficient function,
not a dynamical scalar or a new gauge field.

\begin{theorem}[A linear-coordinate primitive and its exact divergence]
\label{thm:v45-coordinate}
Define
\begin{equation}
 \boxed{
 H^x_{B,R}=z^2\int
 v_\mu(x)\bigl(j_R^\mu-c^\mu a_R\bigr).
 }
 \label{eq:v45-xprimitive}
\end{equation}
Then \(sH^x_{B,R}=\mathcal A_{B,R}\) modulo an explicitly
coordinate-dependent local divergence. At density level,
\begin{equation}
 \begin{split}
 s\left[v_\mu(j_R^\mu-c^\mu a_R)\right]
 ={}& A_{B,R}\\
 &+\partial_\nu\left[
 c^\nu v_\mu(j_R^\mu-c^\mu a_R)
       -v_\mu c^\mu j_R^\nu\right].
 \end{split}
 \label{eq:v45-xprimitive-density}
\end{equation}
\end{theorem}
\begin{proof}
Apply the three identities in \cref{lem:v45-current} and retain the
sign in \(s(c^\mu a_R)\). The terms in which a derivative falls on
\(v_\mu\) are
\[
 (\partial_\nu v_\mu)
   \left[c^\nu c^\mu a_R+c^\mu j_R^\nu-c^\nu j_R^\mu\right].
\]
Their antisymmetric part is precisely
\(f_Ba_R+(\alpha_B)_\mu j_R^\mu\); their symmetric part vanishes.
All remaining terms assemble into the displayed divergence.
\end{proof}

This is a sharp category separation: the same candidate is nontrivial
with translation-invariant primitives and trivial when the stated
polynomial \(x\)-dependence is allowed. The coordinate primitive is
\emph{not} inserted into the active theory. It would also not define a
periodic counterterm on the regulator torus.

There is no contradiction with \eqref{eq:v45-cartan}. On a function
with explicit coordinates, \(\{s,\iota_\mu\}\) differentiates the
field jets but does not differentiate its explicit coefficient. After
integration by parts,
\[
 \{s,\iota_\mu\}F=-\partial_\mu^{\mathrm{explicit}}F,
\]
which need not vanish. The extra term is exactly what permits
\eqref{eq:v45-xprimitive}.

The theorem provides a direct audit of the literature interface. The
standard matter-inclusive result in \cite{V29BBHMatter} permits forms
with explicit \(x\), and the associated spacetime primitive above is
admissible there. Its statement that the usual four-dimensional pure
Lorentz chiral anomaly is absent does not imply vanishing of the
translation-invariant group in \eqref{eq:v45-injection}. In particular,
``there is no internal abelian gauge field'' does not remove these
classes when the allowed primitive category is narrowed. They are
not the abelian gauge anomalies classified in that paper.

\section{The actual flavour-invariant reference excludes this entire module}
\label{sec:v45-flavour}

This upstream counterexample is not an obstruction claim for the
populated reference coefficient. The actual branch contains more
information than an arbitrary local scalar--gravity cocycle.

Let \(\mathcal F=(\mathbb Z_2)^4\) act by independent scalar sign
changes,
\[
 \phi_A\longmapsto\epsilon_A\phi_A,
 \qquad\upsilon_A\longmapsto\epsilon_A\upsilon_A,
 \qquad\epsilon_A\in\{\pm1\},
\]
with all gravitational variables and ghosts fixed. This finite subgroup
of the rigid \(O(4)\) already suffices; no spacetime rotation average is
performed. Define
\begin{equation}
 \mathscr P_0=\frac1{16}\sum_{\epsilon\in\mathcal F}\epsilon\cdot,
 \qquad
 \mathscr P_{AB}=\frac1{16}\sum_{\epsilon\in\mathcal F}
       \epsilon_A\epsilon_B\,\epsilon\cdot\quad(A<B).
 \label{eq:v45-flavour-projectors}
\end{equation}
These are projectors on field/source coefficient arrays, not numerical
projections applied after a failed Ward identity.

\begin{proposition}[Exact finite flavour selection]
\label{prop:v45-native-flavour}
The actual ordinary and completed finite action, reference covariances,
minimal source prescription, and scalar measure on the present branch
are \(\mathcal F\)-equivariant. The active gravity prescription, V34
subcritical scalar normalization, and V38 packet preserve this symmetry.
Consequently the complete native reference germ and its actual local
mass jets obey
\begin{equation}
 \boxed{
 \mathscr P_0\mathscr X=\mathscr X,
 \qquad
 \mathscr P_{AB}\mathscr X=0\quad(A<B).
 }
 \label{eq:v45-native-flavour}
\end{equation}
The same statement holds for the relative target
\(\mathscr X_2-tH_g\). All these projectors commute with \(s\),
horizontal differentiation, mass extraction, and constant-ghost
contraction, and have raw coefficient norm at most one.
\end{proposition}
\begin{proof}
Every scalar action coefficient is contracted with the common flavour
metric. The propagator is a scalar operator times the four-dimensional
flavour identity. A sign change conjugates the simultaneous hard Hessian
by a constant orthogonal matrix on the hard scalar factor and transforms
its retained scalar and antifield arguments in the same way. The real
scalar integration measure has absolute Jacobian one. The ghost action
and the coframe contour are unchanged. Thus the measured integral,
normalized logarithm, reference supertrace, and all of their coefficient
extractions obey the symmetry before any local subtraction.

The scalar source \(P((Pc)\cdot D\phi)\) is equivariant. The density
and gravitational normalization have no flavour index; the actual
stress uses the same summed flavour metric. V34 uses summed kinetic and
mass densities, while V38 uses the summed gradient Gram. These facts
establish the claimed symmetry of the active local terms without
assuming that an arbitrary homogeneous addition shares it. Finally the
sign action is diagonal by \(\pm1\) on monomials, proving the norm
bound and commutation assertions.
\end{proof}

For \(R=R_{AB}\), the elementary rotation of the \(AB\)-plane,
\[
 \epsilon\cdot\mathcal A_{B_0,R_{AB}}
 =\epsilon_A\epsilon_B\mathcal A_{B_0,R_{AB}}.
\]
Here \(B_0\) denotes the spacetime two-form, to distinguish it from a
flavour index. Its character is nontrivial.

\begin{theorem}[No native component in the exhibited current module]
\label{thm:v45-singlet}
The 36-dimensional submodule in \eqref{eq:v45-injection} has zero
intersection with the \(\mathcal F\)-invariant cohomology subspace.
In particular the native invariant target has no class in this module.
At coefficient level,
\begin{equation}
 \boxed{
 \mathscr P_{AB}\iota_\nu\iota_\mu\mathscr X=0
 \quad\text{for all }A<B,\ \mu<\nu.
 }
 \label{eq:v45-native-test}
\end{equation}
\end{theorem}
\begin{proof}
Each elementary current pair lies in its nontrivial sign character,
while \(\mathscr P_0\) annihilates that character. Since the finite
group action commutes with the differential, it acts on cohomology with
the same character decomposition. Independence from
\cref{thm:v45-independent} proves that no nonzero combination in the
current module is invariant. Equation \eqref{eq:v45-native-test}
follows directly from \eqref{eq:v45-native-flavour}; it needs no loop
integration or choice of a cohomology representative.
\end{proof}

This conclusion excludes \emph{this module}; it is not a complete
classification of the invariant ghost-one cohomology. A singlet class
outside it is not tested by \eqref{eq:v45-native-test}. Conversely,
adding a nontrivial flavour-valued composite source or internal gauge
coupling can change the coefficient category and its selection rules.
The present result must not be transferred to the chiral/internal ESM
sectors by changing field names.

If a local primitive of the native target is constructed in
\(\mathfrak L_{\rm TI}\), averaging it with \(\mathscr P_0\)
preserves its equation, derivative/mass degrees, and locality, with no
increase of its coefficient norm. This is a valid reduction to the
invariant primitive problem. Averaging a coordinate-dependent primitive,
however, does not remove its explicit \(x\)-dependence.

\section{The mass-lift audit and the remaining absolute equation}
\label{sec:v45-mass}

The three equations of V35 remain unchanged. The new examples test their
logical sufficiency without rederiving the triangular calculus. Let
\(D=s(0)\) and \(M=(\frac12\int\rho\phi_A^2,\cdot)\), so
\(s(z)=D+zM\) and \(DM+MD=0\).
For \(A^0_{B,R}=\int A_{B,R}\),
\begin{equation}
 \boxed{DA^0_{B,R}=0,\qquad MA^0_{B,R}=0.}
 \label{eq:v45-mass-closed}
\end{equation}
The second identity is immediate from
\(M a_R=\rho\phi_A R_{AB}\phi_B=0\); the current has no antifield.

\begin{proposition}[A mass-consistent target with no admissible primitive]
\label{prop:v45-mass-counterexample}
The coefficient array
\[
 \xi_0=0,\qquad \xi_1=0,\qquad \xi_2=A^0_{B,R}
\]
satisfies all three mass-lift consistency equations and the terminal
condition \(M\xi_2=0\). For nonzero \(B\otimes R\), it has no
translation-invariant mass-polynomial critical primitive. The first two
mass equations admit the choices \(h_0=h_1=0\), but this does not
solve the third.
\end{proposition}
\begin{proof}
Consistency is \eqref{eq:v45-mass-closed}. A full solution would be an
admissible primitive of \(z^2A^0_{B,R}\), contrary to
\cref{thm:v45-independent}. The target vanishes when evaluated at
\(z=0\), which also shows why a massless value alone cannot test this
question. The example is excluded by the native flavour selection in
\cref{thm:v45-singlet}; it is not the populated \(\xi_2\).
\end{proof}

The actual remaining absolute problem is therefore
\begin{equation}
 \boxed{
 s(z)H^{[4]}=\mathscr X^{[4]}(z)\pmod{d_x},\qquad
 H^{[4]}\in\mathfrak L_{\rm TI}^{\mathcal F},
 }
 \label{eq:v45-remaining}
\end{equation}
with the full matter-two/four germ, the same reference convention, and
polynomial mass dependence retained. The proven current-module test is
now a zero input to this problem, rather than an unspecified possible
obstruction. It does not prove that the remaining invariant target is
exact. Neither its remaining local cohomology nor a populated finite
primitive has been computed here.

Once \eqref{eq:v45-remaining} is solved in the required category,
V44's fixed-gravity conversion applies with the complete scalar stress
and current. There is no need to redo V40--V44's pure-gravity boundary
analysis. A primitive obtained only in \(\mathfrak L_x\) is not an
input to that conversion. The distinction matters to the programme's
measured, source-local, translation-compatible normalization rather than
merely to terminology.

\section{Verification, uniformity, and status}
\label{sec:v45-ledgers}

The checker uses a finite exterior algebra with independent ghost and
antifield generators. It verifies the two density identities with all
four spacetime directions, reconstructs the double-contraction decoder,
checks the Noether cancellation for all six flavour rotations, and tests
the sixteen sign transformations and their character projectors. It also
checks the mass-consistency example and byte-for-byte preservation of
the previous mathematical body. These exact finite checks do not replace
the proof that rigid rotations are nontrivial local symmetries, classify
the remaining singlet cohomology, or evaluate a native loop integral.

\subsection*{Proved}
\addcontentsline{toc}{subsection}{V45: Proved}
The actual common-mass scalar action supplies six rigid rotation currents.
Their coupling to a constant spacetime two-form gives 36 independent,
critical, mass-polynomial translation-invariant ghost-one classes.
Constant-ghost contraction detects them and has explicit fixed-bank
coefficient bounds. Each has an explicit primitive linear in spacetime
coordinates, demonstrating a genuine change of cohomology when that
larger counterterm category is admitted. The actual finite and ordinary
reference is invariant under independent flavour signs; the entire
exhibited module is excluded from its invariant component. This
selection and its constant-ghost tests are exact, before a loop integral
is evaluated. The mass-lift example proves that consistency and a
coordinate-dependent primitive alone are insufficient for the active
category.

\subsection*{Inherited}
\addcontentsline{toc}{subsection}{V45: Inherited}
V1--V44 and their labels are unchanged. The ordinary scalar action,
minimal source conventions, completed hard carrier, mass grading,
V37 full reference germ, V34 active subcritical scalar prescription,
V38 logarithmic packet, and V44 fixed-gravity lift retain their stated
scope. No prior theorem is replaced or strengthened silently. In
particular the new scalar-current module is not a counterexample to the
pure-gravity hypothesis of V44.

\subsection*{Literature input}
\addcontentsline{toc}{subsection}{V45: Literature input}
The tensor-matter classification of \cite{V29BBHMatter} uses the local
forms and spacetime Poincare lemma specified in its Sections 2 and 6.
The category comparison is made explicitly rather than importing a
translation-invariant vanishing theorem from it. The identification of
negative-ghost classes with rigid variational symmetries is the standard
antifield framework of \cite{V20BBH1994}; the relevant current and
nongauge transformation are displayed directly. No general anomaly
classification is used to assert the new native target is exact.

\subsection*{Conditional}
\addcontentsline{toc}{subsection}{V45: Conditional}
If a translation-invariant primitive of the native target is supplied,
its finite flavour average is an equally local invariant primitive, and
V44 converts an absolute solution to the fixed gravitational
normalization. Existence in the coordinate-dependent category is not
sufficient. The bounds above control fixed coefficient maps and local
jet inclusions, not the inverse of the unresolved full singlet BV
operator or a uniform quantum remainder.

\subsection*{Open}
\addcontentsline{toc}{subsection}{V45: Open}
The remaining critical target is the absolute translation-invariant,
flavour-invariant scalar--gravity primitive in the full stated
mass/source/current category. The 36 current classes computed here are
absent from its invariant component, but other components are not
classified. A proof of their absence or a direct native primitive is
still required. No active counterterm is changed. Internal gauge,
chiral-line, higher-loop and invariant interacting ESM-class interfaces
remain separate. No new cutoff convergence rate, positive quantum
measure, or analytic all-arity result is claimed.

\begin{center}\small
\begin{tabular}{>{\raggedright\arraybackslash}p{.24\textwidth}>{\raggedright\arraybackslash}p{.68\textwidth}}
\toprule
Variable & Scope in V45\\\midrule
Fields & Ordinary ten-coordinate gravity with four equal real scalars,
complete minimal antifields and diffeomorphism ghost; internal gauge,
Yukawa and scalar self-couplings zero.\\
Counterterms & Translation-invariant finite local jets and their formal
coframe germs. Explicit-coordinate primitives are comparison objects
only and are not installed.\\
Flavour & The 36 classes are before flavour restriction. The populated
native reference is invariant under the exact sixteen-element sign
group, which removes this specific module.\\
Mass & Polynomial \(z=m^2\). The cocycles are multiplied by \(z^2\)
and vanish at \(z=0\); their nontriviality is not a massless-value test.
No inverse mass.\\
Weight and arity & Ghost number one, weight four. Full local identities;
fixed finite arity for coefficient bounds. No assertion that 36 is the
cohomology dimension of a truncated bank.\\
Cutoff and volume & Coefficient-map constants are independent of mesh,
periods and shell number at fixed jet order. No new scale-iteration or
continuum estimate.\\
Perturbative role & An algebraic test for a one-loop local primitive.
No measured coefficient of the example module is asserted nonzero.\\
Prescriptions & Fixed gravity, V34 and V38 unchanged. All matrix-valued
source extensions require their own flavour and primitive audit.\\
\bottomrule
\end{tabular}
\end{center}

\clearpage
\part*{V46: Covariant scalar current reduction at critical weight}
\addcontentsline{toc}{part}{V46: Covariant scalar current reduction at critical weight}

\section{A derivative-current question left by V45}
\label{sec:v46-start}

Sections 1--353 and their labels are preserved. The programme remains the
finite-order, source-complete Einstein--Standard Model EFT construction
of the supplied update. This installment stays on its ordinary scalar--gravity
branch with four equal real scalars and vanishing internal gauge, Yukawa,
and scalar self-couplings. It changes no reference, measure, action, or active
counterterm. V44's conditional fixed-gravity conversion and V34's and V38's
normalizations are inherited.

V45 computes the rotation-current module and excludes its flavour-invariant
part. That exclusion is not a classification of derivative-dependent currents.
Here we close a precise additional current sector: all scalar-linear,
metric-fixed, covariantly natural variational symmetries of differential
weight four. They are considered off shell, for \emph{every} metric in the
same local positive coefficient chart. Their scalar characteristics have a
24-dimensional space and an explicit 18-dimensional space of trivial
characteristics. The six remaining classes are exactly the mass-squared
multiples of the rotation charges already treated in V45. In particular,
there is no hidden flavour-invariant class in this larger natural sector.

The qualification ``natural'' is essential. The native finite regulator is
not assumed to be covariant under arbitrary coordinate changes, and its
local breaking is not replaced by a natural ansatz. This is an exact
classification of one possible ordinary conserved-current sector, not a
complete covariantization theorem for the unresolved critical target.
Metric-changing characteristics, nonlinear scalar characteristics, and
other descent sectors are not eliminated by the result below. The new result
must not be used to assert the absolute primitive merely from the word
``singlet.''

\subsection{Operator and source conventions}
Let
\begin{equation}
 P_g=-\Delta_g+z,\qquad z=m^2,\qquad
 S_H=\tfrac12\int\rho\,\phi^TP_g\phi,\qquad
 \rho=\sqrt{\det g},\quad \Delta_g=\nabla^\mu\nabla_\mu.
 \label{eq:v46-reference}
\end{equation}
Thus \(P_g\) is formally self-adjoint for \(\int\rho\); its Euler density
is \(E=\rho P_g\phi\). The transpose below includes flavour transpose and
the formal differential adjoint. Compactly supported tests, or the inherited
periodic integration-by-parts convention, remove boundary terms.

A characteristic \(Q_g\phi\) acts on the scalars and leaves the metric fixed.
Its variational-symmetry equation is exactly
\begin{equation}
 \boxed{P_gQ_g+Q_g^\dagger P_g=0.}
 \label{eq:v46-noether-op}
\end{equation}
No scalar equation of motion is substituted in this identity. A symmetry is
called trivial here when \(Q_g=K_gP_g\), with \(K_g^\dagger=-K_g\).
Its explicit full ordinary BV primitive is given below, so triviality will
not be inferred only from an on-shell slogan.

Assign weight one to a covariant derivative, two to curvature and \(z\),
and zero to the metric and its inverse. Coefficients are real constant
flavour matrices. An operator is \emph{natural in this section} when it is
formed tensorially from these data, the chosen volume form, curvature,
and covariant derivatives, without preferred spacetime tensors, explicit
coordinates, inverse mass, or curvature denominators. It is homogeneous
of weight four and linear in its scalar argument. The allowed parity-odd
term is retained. This definition is narrower than the whole native
component counterterm category.

The source charge \(\mathcal Q_Q=\int\upsilon^TQ_g\phi\) has ghost number
minus one and weight six. This is precisely the weight reached by twice
contracting undifferentiated diffeomorphism ghosts in a critical ghost-one,
weight-four coefficient. Neither that observation nor a projection to one
antifield slot asserts that every such contraction has a representative of
the special metric-fixed form used here.

\section{The complete natural operator inventory at weight four}
\label{sec:v46-inventory}

Write \(R\), \(R_{\mu\nu}\), and \(R_{\mu\nu\rho\sigma}\) for the scalar,
Ricci, and Riemann curvatures. Let
\(\mathfrak p=\tfrac12\varepsilon^{abef}R_{efcd}R_{ab}{}^{cd}\)
be the parity-odd curvature-square scalar. It is not discarded because its
density has a local transgression: as a multiplication \emph{operator}, it
is not a zero function.

\begin{lemma}[Twelve universal scalar operators]
\label{lem:v46-inventory}
Every natural operator in the preceding category has a unique expression
\begin{equation}
 \begin{aligned}
 Q_g={}& A P_g^2+zB P_g+z^2 C+D R P_g
       +E R^{\mu\nu}\nabla_\mu\nabla_\nu
       +F(\nabla^\mu R)\nabla_\mu\\
     &+G\,\Delta_gR+H R^2+I R_{\mu\nu}R^{\mu\nu}
       +J R_{\mu\nu\rho\sigma}R^{\mu\nu\rho\sigma}
       +K zR+L\mathfrak p .
 \end{aligned}
 \label{eq:v46-basis}
\end{equation}
Here \(A,\ldots,L\) are arbitrary real \(4\times4\) flavour matrices;
terms on the second line act by multiplication. Thus the raw operator
space has dimension 192. Without orientation-dependent operators, its
last summand is omitted and the dimension is 176.
\end{lemma}
\begin{proof}
Use normal coordinates at a point. A principal coefficient of order four
has weight zero, so its completely symmetric part is a multiple of the
symmetrized metric square. Reordering its covariant derivatives gives
\(\Delta_g^2\) and lower terms of the types listed. There is no natural
rank-three tensor of weight one, so no independent order-three term.
At order two, the weight-two symmetric tensors are \(zg^{\mu\nu}\),
\(Rg^{\mu\nu}\), and \(R^{\mu\nu}\). At order one the only nonzero
weight-three vector is \(\nabla R\): contracted Bianchi gives
\(\nabla_\mu R^{\mu\nu}=\tfrac12\nabla^\nu R\), and the parity-odd
contractions vanish by the algebraic or differential Bianchi identities.
At order zero the even weight-four scalars are
\(z^2,zR,\Delta R,R^2,|\operatorname{Ric}|^2,|\operatorname{Riem}|^2\),
with the one additional orientation-odd scalar \(\mathfrak p\).
Replacing \(\Delta^2,z\Delta,R\Delta\) by
\(P^2,zP,RP\) is triangular and invertible in this list.

The principal symbols separate the differential orders. Scalar curvature
and trace-free Ricci are independent at a point. The highest metric jets
separate \(\Delta R\), and varying the polynomial parameter \(z\)
separates \(zR\). The remaining four curvature-square polynomials are
independent: scalar curvature, trace-free Ricci, self-dual Weyl, and
anti-self-dual Weyl algebraic tensors give a nonsingular four-row test.
All these jets are realized by smooth metrics near the identity.
These tests also establish uniqueness. No integration-by-parts quotient of
operators is taken in this argument.
\end{proof}

The finite operator inventory is not a basis of all matter--gravity
counterterms. In particular its coefficients depend on the metric, not
on the scalar field, and the metric is not varied by the characteristic.
Those restrictions are part of the theorem, not conclusions about the
uncomputed native primitive.

\section{Universal off-shell symmetries and two explicit curved tests}
\label{sec:v46-classification}

\begin{theorem}[All natural scalar-linear symmetries at this weight]
\label{thm:v46-classification}
The operator \eqref{eq:v46-basis} satisfies \eqref{eq:v46-noether-op}
for every allowed metric and every constant \(z\) if and only if
\begin{equation}
 \boxed{
 Q_g=A P_g^2+zB P_g+z^2 C+D R P_g,
 \qquad A^T=-A,\ B^T=-B,\ C^T=-C,\ D^T=-D.
 }
 \label{eq:v46-answer}
\end{equation}
Consequently the symmetry space has dimension 24, and the reduced
universal determining map has rank 168 on the 192-dimensional inventory.
Its restriction to flavour-sign-invariant operators is injective.
\end{theorem}
\begin{proof}
At a flat metric, \(P=p^2+z\) in Fourier variables. The determining
identity gives
\[
 (A+A^T)P^3+z(B+B^T)P^2+z^2(C+C^T)P=0.
\]
Polynomial independence of \(p^2\) and \(z\) makes \(A,B,C\)
antisymmetric. These three terms solve the equation for every metric,
so subtract them.

The next highest differential symbol separates the scalar and trace-free
Ricci tensors. It forces \(D\) and \(E\) to be antisymmetric. The term
\(DRP\) is already a variational symmetry, because
\[
 P(DRP)+(DRP)^\dagger P=PDRP-PRDP=0.
\]
Subtract it as well. At order three in the remaining determining identity,
the completely trace-free part of
\(\nabla_{(a}R_{bc)}\xi^a\xi^b\xi^c\) has coefficient \(-2E\).
The other order-three terms are multiples of
\((\nabla R\cdot\xi)|\xi|^2\) and cannot cancel it. Hence \(E=0\).
An explicit realizable jet for this conclusion is given below.

For the remaining first-order coefficient, the order-three symbol forces
\(F=F^T\). The order-two symbol then contains
\(-2F\nabla_\mu\nabla_\nu R\,\xi^\mu\xi^\nu\), while multiplication
operators contribute only a multiple of \(|\xi|^2\). A trace-free Hessian
of \(R\) therefore forces \(F=0\).

Only a multiplication matrix \(V_g\) remains. The order-two and order-one
parts of \(PV+V^TP\) respectively give
\(V+V^T=0\) and \(\nabla V=0\). The matrix is zero on a flat open set,
since the already removed \(z^2C\) term is its only curvature-free option.
Take metrics equal to the identity outside a compact set; then \(\nabla
V=0\) forces \(V=0\) everywhere. Arbitrary finite metric jets can be
realized inside such perturbations. Independence in
\cref{lem:v46-inventory} makes \(G,H,I,J,K,L\) vanish.

Conversely every term in \eqref{eq:v46-answer} satisfies the equation by
self-adjointness of \(P\), formal self-adjointness of multiplication by
\(R\), and antisymmetry of its constant flavour matrix. The four
antisymmetric matrices contribute \(4\cdot6=24\) independent coefficients.
There are no additional dependent determining conditions, because this
solution is necessary and sufficient. Finally, a real matrix commuting
with every independent flavour sign is diagonal; an antisymmetric diagonal
matrix is zero.
\end{proof}

\begin{proposition}[Realizable jets separate the two derivative tests]
\label{prop:v46-curved-tests}
The order-three and order-two tests used above can be taken on the same
local metric carrier, without postulating independent curvature variables.
For \(g=e^{2\omega}\delta\), the following jets at the origin suffice:
\begin{enumerate}[label=(\roman*),leftmargin=*]
\item \(\omega=x^0x^1x^2\) gives \(\nabla R=0\), while
\[
 \partial_aR_{ij}\,\xi^a\xi^i\xi^j=-12\xi_0\xi_1\xi_2.
\]
This is a nonzero trace-free cubic.
\item \(\omega=(x^0)^4-(x^1)^4\) gives
\[
 (\nabla^2R)(\xi,\xi)=-144(\xi_0^2-\xi_1^2),\qquad \Delta R=0.
\]
This is a nonzero trace-free quadratic, while the lower curvature jets
needed to isolate it vanish.
\end{enumerate}
A smooth cutoff equal to one near the origin realizes these as bounded
local metric tests. Coordinates in these \emph{tests} are not explicit
coordinates in a counterterm coefficient.
\end{proposition}
\begin{proof}
At these jets \(\omega=\nabla\omega=0\), and the quadratic products of
lower derivatives that enter the displayed curvature orders vanish.
The exact jet formulas are therefore the linear conformal ones,
\(R_{ij}=-2\partial_i\partial_j\omega-\delta_{ij}\Delta\omega\)
and \(R=-6\Delta\omega\). Differentiation gives the displayed tensors.
The first polynomial is harmonic. For the second,
\(R=-72((x^0)^2-(x^1)^2)\) to the relevant degree. The cutoff does not
change any jet under consideration.
\end{proof}

These tests explain why a flat commutant calculation alone is insufficient:
the two curvature-dependent derivative coefficients are detected precisely
by nonconstant background jets. Neither test invokes a metric equation of
motion, a Gaussian inverse, or a cutoff-dependent field amplitude.

\section{An explicit mass-polynomial BV primitive for the derivative currents}
\label{sec:v46-primitive}

Let \(\pi=\rho^{-1}\upsilon\). Its ordinary transformation follows directly
from V45's density convention:
\begin{equation}
 s\pi=P_g\phi+c^\mu\partial_\mu\pi.
 \label{eq:v46-pi}
\end{equation}
For \(Q\) in \eqref{eq:v46-answer}, define
\begin{equation}
 K_g=A P_g+zB+D R,\qquad K_g^\dagger=-K_g,
 \qquad \mathcal T_K=-\tfrac12\int\rho\,\pi^TK_g\pi.
 \label{eq:v46-primitive}
\end{equation}

\begin{theorem}[Exact reduction of every classified current charge]
\label{thm:v46-current-reduction}
In the full ordinary scalar--gravity BV complex, modulo a
translation-invariant divergence,
\begin{equation}
 \boxed{
 \mathcal Q_Q=s\mathcal T_K+z^2\mathcal Q_C,
 \qquad
 \mathcal Q_Q=\int\upsilon^TQ_g\phi,\quad
 \mathcal Q_C=\int\upsilon^TC\phi.
 }
 \label{eq:v46-charge-reduction}
\end{equation}
Both the primitive and its residual are polynomial in \(z\), local on the
same metric chart, and regular at \(z=0\). There are 18 independent
trivial characteristics and six independent residual current classes.
The latter are exactly V45's common-mass rotation charges multiplied by
\(z^2\), not newly found current classes.
\end{theorem}
\begin{proof}
All tensors in \(\rho\pi^TK\pi\) transform naturally; hence its longitudinal
variation is a density divergence. Its Euler variation, with the odd
antifields kept in their displayed order, is
\[
 \delta\left(-\tfrac12\int\rho\pi^TK\pi\right)
 =-\tfrac12\int\rho\left[(P\phi)^TK\pi-\pi^TKP\phi\right]
 =\int\rho\pi^TKP\phi.
\]
The last equality uses \(K^\dagger=-K\) and an even \(P\phi\).
No metric antifield occurs, so no gravitational Euler term is omitted.
Since \(Q=KP+z^2C\), this proves the identity. The three independent
antisymmetric coefficients of \(K\) give 18 parameters. Their multiplication
by \(P\) is injective by the highest differential symbol followed by its
lower orders. V45 proves nontriviality of the six polynomial residual
rotation charges even in the coupled theory; alternatively their flat
symbols are not divisible by \(p^2+z\) in this operator category.
\end{proof}

For example, the natural symmetry \(Q=C\Delta_g^2\) has
\[
 Q=C(P_g^2-2zP_g+z^2),\qquad
 K=C(P_g-2z).
\]
Its charge differs from \(z^2\mathcal Q_C\) by the displayed explicit
quadratic-antifield variation. The curvature characteristic \(DRP_g\)
has zero residual and uses \(K=DR\). Thus curvature-dependent trivial
currents are not accidentally reclassified as new obstructions.
At the massless point, every current in this homogeneous weight-four
operator sector is exact; this statement does not remove the usual
lower-weight massless shift or rotation charges.

\begin{proposition}[A residual-preserving finite compiler]
\label{prop:v46-compiler}
In the basis \eqref{eq:v46-basis}, replace \(A,B,C,D\) by their
antisymmetric parts and set the other eight matrices to zero. Call this
map \(\Pi_{\mathrm{nat}}\). Then
\[
 \Pi_{\mathrm{nat}}^2=\Pi_{\mathrm{nat}},\qquad
 \operatorname{im}\Pi_{\mathrm{nat}}=\ker\mathcal V,
 \qquad \mathcal V(Q)=PQ+Q^\dagger P,
\]
\begin{equation}
 \boxed{
 \mathcal V(Q-\Pi_{\mathrm{nat}}Q)=\mathcal V(Q),\qquad
 \|\Pi_{\mathrm{nat}}\|_{1\to1}=1.
 }
 \label{eq:v46-compiler}
\end{equation}
On the symmetry subspace the maps to \((A,B,D)\) and to \(C\) each have
norm one in the declared coefficient sum. These bounds control finite
coefficient maps, not arbitrary norms of an extensive action.
\end{proposition}
\begin{proof}
Antisymmetrization of a real matrix has induced entry-sum norm one.
The direct product map therefore has the stated norm. The classification
theorem identifies its image. Since that image has zero variational defect,
subtraction preserves the original defect exactly. Coordinate projections
on the four independent antisymmetric matrices have norm one.
\end{proof}

A supplied nonsymmetric input is not declared to solve a Ward identity:
its unremoved eight operator matrices and four symmetric parts are returned.
The reduced determining matrix has 168 independent rows: ten symmetric
conditions for each of four matrices, followed by 128 independent entries
from the remaining eight matrices. Its rational kernel has the 24 explicit
columns in \eqref{eq:v46-answer}. This is an evaluated equivalent determining
matrix, not the full local BV matrix of the Einstein--SM theory.

\section{Flat symbols, flavour selection, and the scope of the exclusion}
\label{sec:v46-flavour}

For comparison, drop covariance in the background but keep a constant flat
metric and arbitrary constant spacetime coefficients. A homogeneous
weight-four scalar-linear variational symmetry then has an antisymmetric
flavour matrix multiplying each monomial
\(p^\alpha z^j\), \(|\alpha|+2j=4\). There are
\(35+10+1=46\) such monomials.

\begin{proposition}[The flat symbol quotient]
\label{prop:v46-flat}
This flat symmetry space has dimension 276. Its Euler-trivial subspace
has dimension 66, and its quotient has dimension 210. The natural curved
current quotient maps into it with dimension six.
A flat polynomial \(f(p,z)\) has the unique division
\begin{equation}
 f=(p^2+z)k+r(p),\qquad \deg_w k=2,\quad \deg r=4,
 \quad r(p)=f(p,-p^2).
 \label{eq:v46-flat-division}
\end{equation}
In the monomial coefficient norm the quotient and remainder maps have
norms five and sixteen, respectively.
\end{proposition}
\begin{proof}
The flat Noether equation makes the flavour matrix of every even-order
monomial antisymmetric. There are six independent such entries, hence
\(6\cdot46=276\). Multiplication by the monic polynomial \(p^2+z\)
is injective on the \(10+1=11\) monomials of weight two. Division in
\(z\) gives all 35 independent quartic remainders, proving the counts.
The input \(z^2\) has quotient \(z-p^2\) of coefficient norm five and
remainder \((p^2)^2\) of norm sixteen; every other basis column has a
smaller norm. The natural residual is \(z^2C\), whose remainder is
\((p^2)^2C\), giving six independent columns. Polynomial quotient
reduction is not a negative-mass evaluation of a propagator.
\end{proof}

This comparison diagnoses two different issues. Flat higher-derivative
currents substantially overcount the currents admitted by a covariant
all-background realization. Conversely, neither classification includes
characteristics that change the metric or are nonlinear in the scalar.
Those cannot be rejected by dropping their gravity or higher-field terms.

Let \(\mathcal F=(\mathbb Z_2)^4\) be the actual independent sign group
of V45. It acts on every coefficient matrix by \(Q\mapsto\epsilon Q\epsilon\).
Both \(\Pi_{\mathrm{nat}}\) and \eqref{eq:v46-charge-reduction} commute
with this action.

\begin{theorem}[No invariant class in the complete classified current sector]
\label{thm:v46-no-singlet}
The flavour-invariant subspace of the 24-dimensional natural symmetry
space is zero. Its six-dimensional current quotient likewise has no
invariant vector. More generally, any invariant cohomology class that
admits a representative among the classified scalar-only natural charges
is zero.
\end{theorem}
\begin{proof}
Invariance under all independent signs makes each flavour matrix diagonal.
All four symmetry matrices are antisymmetric, so they vanish. Alternatively,
each elementary antisymmetric matrix transforms in the nontrivial character
\(\epsilon_A\epsilon_B\). The primitive in
\eqref{eq:v46-charge-reduction} is equivariant. If a cohomology class,
rather than a selected representative, is invariant, average its representative
over the sixteen group elements. This averaging commutes with \(s\), is
still a natural scalar-only charge, represents the same class, and is zero.
\end{proof}

The raw invariant natural inventory has dimension \(12\cdot4=48\);
its determining map has zero kernel and rank 48. Thus there is not even
an invariant off-shell variational characteristic hidden among the
curvature coefficients of this inventory. This extends V45's explicit
rotation calculation to every operator in the stated derivative sector.
It does not assert that the full native invariant critical breaking has
no other descent component.

\section{How the result enters the absolute primitive problem}
\label{sec:v46-interface}

The established identification of negative-ghost local BV cohomology with
rigid variational symmetries provides the interpretation of the current
charges \cite[Sections 7--8]{V20BBH1994}. The operator classification and
the explicit primitive above are proved directly and do not use a
blanket matter-inclusive cohomology-vanishing assertion. The tensor-matter
classification of \cite{V29BBHMatter} remains subject to V45's explicit
counterterm-category audit.

In the translation-invariant integrated category, the actual contraction
identity \(\{s,\iota_\mu\}=\partial_\mu\) is inherited from V45. Hence a
critical cocycle \(\mathscr X\) has closed twice-contracted coefficients
of ghost number minus one and weight six. If a current reduction puts one
of their classes in the natural scalar-linear, metric-fixed sector,
\cref{thm:v46-no-singlet} removes that component using the actual flavour
symmetry. The reduction \emph{to} that sector must itself be proved; it
cannot be inferred from the symmetry of the scalar action alone.

In particular, projecting a full current onto its \(\upsilon\)-linear
part is not in general a chain map. A component with a metric antifield
and two scalar fields can supply the same scalar Euler or stress identity.
Likewise scalar-nonlinear characteristics and the four-matter germ of V37
can contribute. No assertion here removes those components, changes the
allowed physical source tensors, or proves that a possibly non-natural
native array can be replaced by a natural representative while preserving
its locality, mass polynomial, and translation convention.

The exact polynomial source redefinition \(\mathcal T_K\) has ghost
number minus two. It is a primitive for a \emph{current charge}, not by
itself the ghost-number-zero quantum action counterterm requested by the
programme. Multiplication by two ghosts without its complete descent
partners would not turn it into that counterterm. This installment
therefore installs no new action or source normalization.

On a fixed coframe chart with \(\rho\) bounded away from zero,
\(\pi=\rho^{-1}\upsilon\) and the finite-arity expansion of the displayed
primitive have ordinary analytic coefficient bounds. Every fixed source
order and spacetime derivative order has a finite conversion constant,
independent of the number of lattice sites. These are local differential
polynomial bounds. They do not yield a many-scale invariant class, remove
the lower reference scale, or establish a cutoff estimate for the still
uncomputed absolute critical coefficient.

The next remaining issue is the simultaneous invariant current/descent
reduction, including metric-changing and scalar-nonlinear components, or
a direct construction of the native invariant critical primitive. V44's
fixed-gravity conversion remains available only after that absolute
primitive has been obtained in the required category. The present result
removes the derivative-dependent natural scalar-only current sector as
an independent gap, without restating V45's 36 classes as a new theorem.

\section{Verification, uniformity, and status}
\label{sec:v46-ledgers}

The accompanying checker reconstructs the 192-column reduced determining
matrix, its 24-column kernel, the 18 trivial columns and six-class decoder.
It checks the flat polynomial sequence by exact rational arithmetic,
the two conformal metric jets, four independent curvature-square tensors,
the formal adjoint identities, the odd-antifield sign with a finite
noncommuting model, and all sixteen flavour actions. The certificate
separates these checks from the universal naturality proof and from the
unresolved complete critical primitive. No numerical loop quadrature or
proof-assistant verification is asserted.

\subsection*{Proved}
\addcontentsline{toc}{subsection}{V46: Proved}
The twelve-operator natural inventory is complete in its declared category.
Its universal variational symmetries have the explicit 24-dimensional
normal form \eqref{eq:v46-answer}. Two realizable curved metric jets
remove the extra derivative candidates. An explicit quadratic-antifield
primitive reduces 18 directions and leaves exactly the six known
mass-polynomial rotation classes. The full classified current sector has
no invariant flavour-sign component. Exact finite coefficient projectors,
residual preservation, and the flat symbol comparison are evaluated.

\subsection*{Inherited}
\addcontentsline{toc}{subsection}{V46: Inherited}
V1--V45 and their labels, the native scalar action and ordinary BV signs,
the finite flavour symmetry, V45's rotation-class nontriviality and
translation contractions, V34's active scalar normalization, V38's complete
logarithmic packet, and V44's conditional fixed-gravity conversion are
unchanged. No preceding loop or continuum proof is repeated as a new result.

\subsection*{Literature input}
\addcontentsline{toc}{subsection}{V46: Literature input}
The current interpretation of negative-ghost BV cohomology is the
established antifield framework \cite{V20BBH1994}. The broader tensor-matter
classification \cite{V29BBHMatter} is comparison material, not a claim of
vanishing in the smaller translation-invariant source category. Normal
coordinate tensor identities used here are written out in the proofs.

\subsection*{Conditional}
\addcontentsline{toc}{subsection}{V46: Conditional}
A component of the native twice-contracted critical coefficient that
admits a representative in this natural scalar-only class is now known
to vanish as an invariant current class. Existence of such a representative
for every component is not proved. An absolute invariant primitive, once
constructed with all other components, can be converted to the fixed
gravitational prescription by the inherited V44 theorem.

\subsection*{Open}
\addcontentsline{toc}{subsection}{V46: Open}
The full translation-invariant, flavour-invariant, mass-polynomial
scalar--gravity critical primitive remains open. Metric-changing currents,
nonlinear scalar terms, other descent classes, and source-dependent
covariantization must still be controlled. Internal gauge and chiral
sectors, higher loops, and the invariant interacting ESM class remain
separate. No active normalization is changed.

\begin{center}\small
\begin{tabular}{>{\raggedright\arraybackslash}p{.24\textwidth}>{\raggedright\arraybackslash}p{.68\textwidth}}
\toprule
Variable & Scope in V46\\\midrule
Carrier & Same ordinary metric--scalar action; four equal real scalars,
minimal antifields, internal couplings zero.\\
Current sector & Linear in the scalar, metric held fixed by the rigid
characteristic; natural metric coefficients; operator weight four.\\
Mass & Polynomial common \(z=m^2\), no inverse mass. The primitive remains
regular at zero mass; lower-weight massless currents are not reclassified.\\
Coordinates & No explicit coordinates in coefficients or discarded currents.
Polynomial coordinates occur only in realizable jet tests.\\
Cutoff and volume & Fixed coefficient-map norms independent of mesh, periods,
and shell count. No new measured cutoff error or all-scale estimate.\\
Source and field order & Exact local operator identities; fixed arity for
coframe coefficient bounds. No bound uniform in increasing EFT/source order.\\
Perturbative role & Algebraic current input to the existing one-loop
normalization problem; no new loop integral or counterterm installed.\\
\bottomrule
\end{tabular}
\end{center}

\clearpage
\part*{V47: Coupled Ricci-ideal currents and their nonlinear scalar completion}
\addcontentsline{toc}{part}{V47: Coupled Ricci-ideal currents and their nonlinear scalar completion}

\section{The metric-changing interface left by the scalar-only calculation}
\label{sec:v47-start}

Sections 1--360 and their labels are preserved. This installment uses the
same ordinary zero-mesh scalar--gravity action and minimal sources as V46,
with four equal real scalars, a common polynomial mass \(z=m^2\), and zero
internal gauge, Yukawa and scalar self-couplings. The cosmological source
is zero on the inherited flat Einstein branch. The completed finite
regulator, its reference profiles and contour, V34's active subcritical
normalization, V38's logarithmic packet, and the fixed gravitational
prescription are unchanged.

V46 classified scalar-linear variational symmetries with the metric held
fixed. That is not a projection theorem for a simultaneous current: a
metric characteristic of scalar degree two can cancel a scalar-linear
variation through the gravitational Euler tensor. We now compute the
\emph{complete scalar-linear projection} of that larger natural sector,
and give a local mixed-antifield section for every Ricci-dependent
operator in it. The section has a necessary scalar-cubic completion.
Dropping that completion destroys the full variational identity.

The result is a projection and explicit-section theorem, not a
classification of all metric characteristics or of the full invariant
critical BV cohomology. The residual metric-Euler identity and nonlinear
matter equation are kept at the end. In particular, the native,
possibly non-natural reference coefficient is not replaced by the
covariant ansatz in this calculation.

\subsection{The signed native Euler normalization}
Use metric coordinates only as the inherited pointwise coefficient chart,
with the cotangent conversion back to \(q\) retained below. Write
\begin{equation}
 \begin{gathered}
 S=S_g+\tfrac12\int\rho\,\phi^TP\phi,\qquad
 P=-\Delta_g+z,\qquad \rho=\sqrt{\det g},\\
 \rho^{-1}\frac{\delta S_g}{\delta g_{ab}}
   =e^{ab}=c_E G^{ab},\qquad c_E\ne0,\qquad f=P\phi.
 \end{gathered}
 \label{eq:v47-data}
\end{equation}
Here \(c_E\) is the signed normalization already used in
\eqref{eq:v44-Ricci-lift}; it is not fitted here. The normalized scalar
stress with respect to the covariant metric is
\begin{equation}
 \tau^{ab}=\tfrac14g^{ab}(X+zU)-\tfrac12V^{ab},\qquad
 U=\phi^T\phi,\quad
 V^{ab}=\sum_A\nabla^a\phi_A\nabla^b\phi_A,\quad X=g_{ab}V^{ab}.
 \label{eq:v47-stress}
\end{equation}
Thus the complete metric Euler tensor is \(e+\tau\). Tensor pairings
use the full ordered-index contraction, and formal adjoints use
\(\int\rho\). Compact support, or the declared periodic integration-by-parts
convention, is retained.

A natural matter-leading characteristic has
\begin{equation}
 \delta g=F_2(g,\phi)+O(\phi^4),\qquad
 \delta\phi=Q_g\phi+O(\phi^3),
 \label{eq:v47-characteristic}
\end{equation}
where \(F_2\) is homogeneous quadratic in the scalars and their derivatives,
and \(Q_g\) belongs to V46's twelve-operator inventory of weight four.
There is no scalar-independent metric transformation. Naturality means
exactly V46's category: no preferred spacetime tensors, explicit
coordinate coefficients, curvature denominators, or inverse mass.
A completion through scalar degree two satisfies
\begin{equation}
 \int\rho\left(e^{ab}(F_2)_{ab}+f^TQ_g\phi\right)=0
 \quad\hbox{for every local metric and scalar test.}
 \label{eq:v47-leading-identity}
\end{equation}
The construction below also supplies full scalar-polynomial completions,
not merely solutions of this leading equation.

\section{Two Ricci-flat tests of the allowed scalar projection}
\label{sec:v47-vacuum}

\begin{lemma}[Ricci-flat restriction and exact local witnesses]
\label{lem:v47-vacuum}
Every solution of \eqref{eq:v47-leading-identity} obeys
\(PQ+Q^\dagger P=0\) on each Ricci-flat metric. Two positive local
metrics separate the remaining curvature-square operators. In
coordinates \((t,x,y,w)\), \(x>0\), take
\begin{equation}
 g_\epsilon=x(\dd x^2+\dd y^2+\dd w^2)
       +x^{-1}(\dd t+\epsilon y\dd w)^2,\qquad \epsilon=\pm1.
 \label{eq:v47-testmetrics}
\end{equation}
Both have \(\operatorname{Ric}=0\). With the fixed orientation and
\(\mathfrak p=\tfrac12\varepsilon^{abef}R_{efcd}R_{ab}{}^{cd}\),
\begin{equation}
 |\operatorname{Riem}|^2=24x^{-6},\qquad
 \mathfrak p=\epsilon\,24x^{-6}.
 \label{eq:v47-invariants}
\end{equation}
These are local coefficient tests, not new quantum backgrounds.
\end{lemma}
\begin{proof}
On a Ricci-flat metric \(e=0\). Polarizing the real scalar quadratic
form in \eqref{eq:v47-leading-identity} gives the formal operator
identity. For the explicit metrics use the orthonormal coframe
\[
 \omega^0=x^{-1/2}(\dd t+\epsilon y\dd w),\quad
 \omega^1=x^{1/2}\dd x,\quad\omega^2=x^{1/2}\dd y,\quad
 \omega^3=x^{1/2}\dd w.
\]
With \(u=x^{-3/2}\), its structure equations are
\[
 \dd\omega^0=\tfrac u2\omega^0\wedge\omega^1
              +\epsilon u\omega^2\wedge\omega^3,\quad
 \dd\omega^1=0,\quad
 \dd\omega^j=\tfrac u2\omega^1\wedge\omega^j\ (j=2,3),\quad
 \dd u=-\tfrac32u^2\omega^1.
\]
Cartan's equations give the curvature matrix \(x^{-3}M_+\) in the
ordered pair basis \((01,02,03,12,13,23)\), where
\begin{equation}
 M_+=\begin{pmatrix}
 -1&0&0&0&0&-1\\
 0&1/2&0&0&-1/2&0\\
 0&0&1/2&1/2&0&0\\
 0&0&1/2&1/2&0&0\\
 0&-1/2&0&0&1/2&0\\
 -1&0&0&0&0&-1
 \end{pmatrix}.
 \label{eq:v47-curvmat}
\end{equation}
For the opposite sign, \(M_-=SM_+S\), with
\(S=\operatorname{diag}(-1,-1,-1,1,1,1)\).
Ricci contraction vanishes. The Hodge matrices obey
\(\star M_\epsilon=\epsilon M_\epsilon\), and
\(4\sum_{ij}(M_\epsilon)_{ij}^2=24\). This proves
\eqref{eq:v47-invariants}, including its nonconstant dependence on
\(x\). The checker independently reconstructs the Christoffel symbols
and Ricci tensor directly from \eqref{eq:v47-testmetrics}.
At \(x=1,y=0\) the metric equals the identity. A sufficiently small
coordinate neighbourhood is in the inherited positive chart. The
coordinates specify test metrics, not counterterm coefficients.
\end{proof}

\begin{theorem}[Complete admissible scalar-linear projection]
\label{thm:v47-projection}
A scalar operator in V46's inventory admits a natural metric-changing
completion \eqref{eq:v47-leading-identity} if and only if
\begin{equation}
 \boxed{Q=AP^2+zBP+z^2C+Q_{\rm Ric},\qquad
 A^T=-A,\quad B^T=-B,\quad C^T=-C,}
 \label{eq:v47-Qnormal}
\end{equation}
where the seven matrices in
\begin{equation}
 \begin{aligned}
 Q_{\rm Ric}={}&DRP+E R^{ab}\nabla_a\nabla_b
   +F(\nabla^aR)\nabla_a+G\Delta R\\
 &+H R^2+I R_{ab}R^{ab}+KzR
 \end{aligned}
 \label{eq:v47-Qric}
\end{equation}
are arbitrary real \(4\times4\) matrices. In V46's notation the
coefficients \(J\) of \(|\operatorname{Riem}|^2\) and \(L\) of
\(\mathfrak p\) vanish. The admissible scalar projection has dimension
\(3\cdot6+7\cdot16=130\); its determining map has rank 62 on the
192-dimensional inventory.
\end{theorem}
\begin{proof}
At a flat metric, polynomial comparison in momentum and \(z\) forces
\(A,B,C\) to be antisymmetric, exactly as in the inherited scalar
calculation. These three operators solve the scalar variational
identity on every metric. Subtract them. On either Ricci-flat witness,
all seven terms of \eqref{eq:v47-Qric} vanish. The remaining operator
is multiplication by \(V(x)=24x^{-6}(J+\epsilon L)\).
The order-two part of \(PV+V^TP=0\) forces \(V+V^T=0\); its
order-one part then forces \(\nabla V=0\). Since \(x^{-6}\) is
nonconstant, \(J+L=J-L=0\). This is necessary regardless of the
allowed derivative structure of \(F_2\). Sufficiency, including a full
completion, follows from \cref{thm:v47-mixed,thm:v47-section} below.
Independence of the seven Ricci operators is inherited from the original
inventory. No dimension is assigned to the nonunique metric component.
\end{proof}

\section{A local mixed-antifield section for the Ricci-dependent operators}
\label{sec:v47-mixed}

For a symmetric contravariant tensor \(T\), define
\begin{equation}
 \mathfrak r(T)=-c_E^{-1}\operatorname{tr}_gT,\qquad
 \mathfrak r^{ab}(T)=c_E^{-1}
       \left(T^{ab}-\tfrac12g^{ab}\operatorname{tr}_gT\right).
 \label{eq:v47-Ricci-extract}
\end{equation}
In four dimensions, \(\mathfrak r(e)=R\) and
\(\mathfrak r^{ab}(e)=R^{ab}\). Define the operator
\begin{equation}
 \begin{aligned}
 \mathcal K_\phi T={}&D\mathfrak r(T)P\phi
  +E\mathfrak r^{ab}(T)\nabla_a\nabla_b\phi
  +F\nabla^a\mathfrak r(T)\nabla_a\phi
  +G\Delta\mathfrak r(T)\phi\\
 &+HR\mathfrak r(T)\phi
  +IR_{ab}\mathfrak r^{ab}(T)\phi
  +Kz\mathfrak r(T)\phi .
 \end{aligned}
 \label{eq:v47-Kmap}
\end{equation}
It is linear in \(T\) and \(\phi\), differentiates \(T\) at most
twice, and satisfies
\begin{equation}
 \boxed{\mathcal K_\phi e=Q_{\rm Ric}\phi.}
 \label{eq:v47-factor}
\end{equation}
This is trace reversal of the native Euler tensor, not inversion of a
differential field equation.

The adjoint is explicit. For any scalar column \(f\), put
\(a_D=f^TDP\phi\), \(b_{ab}=f^TE\nabla_a\nabla_b\phi\),
\(v^a=f^TF\nabla^a\phi\), and \(a_G=f^TG\phi\). Then
\begin{equation}
 \begin{aligned}
 c_E(\mathcal K_\phi^\dagger f)_{ab}
 ={}&-g_{ab}a_D+b_{ab}-\tfrac12g_{ab}b^c{}_c
      +g_{ab}\nabla_cv^c-g_{ab}\Delta a_G\\
 &-g_{ab}R f^TH\phi
   +\left(R_{ab}-\tfrac12g_{ab}R\right)f^TI\phi
   -g_{ab}z f^TK\phi .
 \end{aligned}
 \label{eq:v47-adjoint}
\end{equation}
Every derivative acts on the whole product. In particular a derivative
on an Euler factor, and the required flavour transpose, are retained.

\begin{theorem}[Exact coupled completion of the Ricci ideal]
\label{thm:v47-mixed}
For arbitrary seven matrices in \eqref{eq:v47-Qric}, the characteristic
\begin{equation}
 \boxed{
 \delta g_{ab}=-(\mathcal K_\phi^\dagger f)_{ab},\qquad
 \delta\phi=\mathcal K_\phi(e+\tau)
       =Q_{\rm Ric}\phi+\mathcal K_\phi\tau
 }
 \label{eq:v47-fullcharacteristic}
\end{equation}
is an exact variational symmetry of the complete ordinary action
\eqref{eq:v47-data}. Its metric component is scalar-quadratic, and its
scalar component is scalar-linear plus scalar-cubic. The full current
charge is BV-exact, with local primitive
\begin{equation}
 \boxed{
 \mathcal T_{\rm mix}
   =-\int\rho\,\pi^T\mathcal K_\phi\vartheta,\qquad
 \pi=\rho^{-1}\upsilon,\quad
 \vartheta^{ab}=\rho^{-1}\eta_g^{ab}.
 }
 \label{eq:v47-mixedprimitive}
\end{equation}
The primitive is polynomial in \(z\), has ghost number minus two and
weight six, and contains no inverse mass, momentum, or cutoff filter.
\end{theorem}
\begin{proof}
The adjoint identity gives
\[
 \delta S=\int\rho\left[-(e+\tau):\mathcal K_\phi^\dagger f
                 +f^T\mathcal K_\phi(e+\tau)\right]=0.
\]
For BV exactness use \(\delta\pi=f\),
\(\delta\vartheta=e+\tau\), with \(\delta\) zero on ordinary
fields. The odd product in \eqref{eq:v47-mixedprimitive} gives
\[
 \delta\mathcal T_{\rm mix}
 =-\int\rho f^T\mathcal K_\phi\vartheta
   +\int\rho\pi^T\mathcal K_\phi(e+\tau)
 =\int\eta_g:\delta g+\int\upsilon^T\delta\phi.
\]
The primitive is a natural scalar density. Its longitudinal variation
is a total divergence, including the transformations of \(\rho\),
the tensor antifield and covariant derivatives. Its full ordinary BV
variation is therefore the displayed charge modulo the allowed
divergence. No completeness theorem for metric gauge generators is
needed for this constructive identity.
\end{proof}

\section{The nonlinear scalar term cannot be omitted}
\label{sec:v47-nonlinear}

The pure scalar-linear projection is not the complete characteristic.
Applying the trace section to the actual stress gives
\begin{equation}
 \boxed{
 \mathfrak r(\tau)=-c_E^{-1}(X/2+zU),\qquad
 \mathfrak r^{ab}(\tau)
 =-c_E^{-1}\left(V^{ab}/2+zUg^{ab}/4\right).
 }
 \label{eq:v47-stress-extract}
\end{equation}
Substitution into \eqref{eq:v47-Kmap} is the entire scalar-cubic
completion. Derivatives on \(X,U,V\), and the actual curvature factors
in the last terms, must remain.

\begin{proposition}[A nonzero four-matter variation if the completion is dropped]
\label{prop:v47-example}
For \(Q_{\rm Ric}=zR\) with identity flavour matrix,
\eqref{eq:v47-fullcharacteristic} becomes
\begin{equation}
 \boxed{
 \delta g_{ab}=\frac z{c_E}g_{ab}\phi^TP\phi,\qquad
 \delta\phi=zR\phi-\frac z{c_E}(X/2+zU)\phi.
 }
 \label{eq:v47-zR}
\end{equation}
Omitting the scalar-cubic term leaves, on a flat constant-scalar test,
the nonzero action variation density \(z^3U^2/c_E\) for \(z>0\),
\(\phi\ne0\). The complete characteristic cancels it exactly.
\end{proposition}
\begin{proof}
Here \(\mathcal K_\phi T=-(z/c_E)\phi\operatorname{tr}T\).
The adjoint and \eqref{eq:v47-stress-extract} give \eqref{eq:v47-zR}.
At flat constant fields, \(R=X=0\), \(P\phi=z\phi\), and
\(\operatorname{tr}\tau=zU\). The metric variation contracts to
\(z^3U^2/c_E\); the scalar-cubic variation contracts to its negative.
No stationary equation is imposed on this off-shell test.
\end{proof}

The construction is also nontrivial at zero mass. For
\(Q_{\rm Ric}=R^{ab}\nabla_a\nabla_b\) with identity flavour matrix,
\begin{equation}
 \begin{aligned}
 \delta g_{ab}&=-c_E^{-1}
 \left[f^T\nabla_a\nabla_b\phi-\tfrac12g_{ab}f^T\Delta\phi\right],\\
 \delta\phi&=R^{ab}\nabla_a\nabla_b\phi
  -c_E^{-1}\left(V^{ab}/2+zUg^{ab}/4\right)
                 \nabla_a\nabla_b\phi .
 \end{aligned}
 \label{eq:v47-Richess-example}
\end{equation}
Both expressions are regular at \(z=0\). They are Euler-trivial
variational characteristics with an explicit primitive, not new
physical rigid symmetries.

The mixed primitive has matter count two; its BV variation contains
matter-count-two and matter-count-four terms because
\(\delta\eta_g\) includes the scalar stress. This evaluates the
triangular bookkeeping on the actual current generators. Dropping the
four-matter source term invalidates \(s\mathcal T_{\rm mix}
=\mathcal Q_{\rm mix}\).

\section{The selected section and its flavour-invariant part}
\label{sec:v47-section}

\begin{theorem}[Full selected completions and their exact directions]
\label{thm:v47-section}
Every projection in \eqref{eq:v47-Qnormal} admits a full natural
scalar-polynomial current completion. A fixed linear choice is
\begin{equation}
 \boxed{
 \mathcal Q_{\rm sel}
 =s\left(\mathcal T_{\rm mix}
     -\tfrac12\int\rho\,\pi^T(AP+zB)\pi\right)
     +z^2\int\upsilon^TC\phi .
 }
 \label{eq:v47-fullsection}
\end{equation}
This section has 124 explicitly exact directions and the six rotation
charges already identified in V45. The count is
\(7\cdot16+2\cdot6=124\). It is not the dimension of the full current
cohomology.

The independent-flavour-sign-invariant admissible scalar projection has
dimension 28. Every one of its selected completions is exact by
\eqref{eq:v47-mixedprimitive}; no residual rotation charge survives.
\end{theorem}
\begin{proof}
The inherited primitive for the skew-adjoint operator \(AP+zB\)
contributes \(AP^2\phi+zBP\phi\) on every metric.
\Cref{thm:v47-mixed} supplies the seven arbitrary Ricci coefficients
and their complete nonlinear partners. The final six terms are the
actual flavour rotation charges multiplied by \(z^2\); their
nontriviality in the translation/mass category is inherited from V45.
These scalar-operator coefficients are independent, proving the section
counts. An invariant flavour matrix is diagonal. Antisymmetric
\(A,B,C\) must vanish, while each of the seven Ricci matrices has four
diagonal entries. The mixed primitive commutes with the flavour action,
so all 28 invariant projections have invariant exact completions.
\end{proof}

V46's metric-fixed invariant symmetry space was zero. The larger sector
has nonzero invariant coupled characteristics such as
\eqref{eq:v47-zR}, but their selected charges are exact. Neither
statement follows by simply deleting the metric characteristic.

\begin{proposition}[Projection with an explicit unremoved defect]
\label{prop:v47-projector}
On the twelve-matrix inventory, let \(\Pi_{\rm mix}\) antisymmetrize
\(A,B,C\), leave \(D,E,F,G,H,I,K\) unchanged, and set \(J,L\) to
zero. In matrix entry-sum coefficient norms,
\begin{equation}
 \Pi_{\rm mix}^2=\Pi_{\rm mix},\qquad
 \|\Pi_{\rm mix}\|_{1\to1}=1.
 \label{eq:v47-pi}
\end{equation}
The discarded symmetric parts and curvature-square matrices are the
unremoved Ricci-flat test defect. A general failed input is not declared
to have a completion. The invariant restriction has dimension 28 inside
the 48-dimensional raw invariant inventory.
\end{proposition}
\begin{proof}
Antisymmetrization has induced entry-sum norm one; the other operations
are coordinate projections or identities. The independent determining
conditions are the ten symmetric entries of each of \(A,B,C\), and
the sixteen entries of each of \(J,L\), giving 62. The preceding
necessity and sufficiency proof identifies their kernel. Restriction to
diagonal matrices gives invariant determining rank 20. Subtracting the
admissible projection leaves these linear determining defects unchanged.
\end{proof}

\section{The residual metric identity and the next matter equation}
\label{sec:v47-interface}

Let a full natural current have the matter-leading form
\eqref{eq:v47-characteristic}. Subtract its selected completion
\eqref{eq:v47-fullsection}, including every scalar-cubic term. The
remaining closed current has zero scalar-linear projection and obeys
\begin{equation}
 \boxed{
 \int\rho\,e^{ab}(\widehat F_2)_{ab}=0
 \quad\text{for arbitrary }g,\phi .
 }
 \label{eq:v47-metric-residual}
\end{equation}
If the rotation projection is retained rather than subtracted, it stays
explicit; it is absent on invariant inputs. At the next matter degree,
\begin{equation}
 \boxed{
 \int\rho\left[
 e^{ab}(\widehat F_4)_{ab}
 +\tau^{ab}(\widehat F_2)_{ab}
 +f^T\widehat Q_3
 \right]=0 .
 }
 \label{eq:v47-next-matter}
\end{equation}
Both statements follow by scaling \(\phi\mapsto\varepsilon\phi\)
after the full exact subtraction. They cannot be replaced by an
independent scalar-only variational identity.

Equation \eqref{eq:v47-metric-residual} is a parameter-dependent
Einstein Noether identity. No local natural section expressing every
such \(\widehat F_2\) as a field-dependent diffeomorphism plus an
Euler-trivial metric operator has been constructed here. Importing an
unrestricted completeness statement for Noether generators would still
require preserving the polynomial mass, translation convention,
naturality, and finite-jet bounds. The nonlinear scalar equation
\eqref{eq:v47-next-matter} is also retained. The present mixed section
removes the Ricci-dependent scalar projection without assuming these
remaining reductions.

Application to the actual critical reference requires a reduction to the
stated category. V45's ghost contractions can retain external tensor
indices, and the finite source need not be natural. Those tensors are
not averaged away. Scalar-independent metric transformations are also
outside \eqref{eq:v47-characteristic}. Neither 130 nor 28 is therefore
a claim of full invariant BV-cohomology vanishing.

Antisymmetric Euler operators and quadratic-antifield primitives are
established BV methodology \cite{V20BBH1994}. The new results are the
seven native Ricci factors, their exact mixed section and adjoint, the
two explicit Ricci-flat exclusions, and the nonlinear scalar completion.
No general cohomology theorem is used to remove the remaining equations.

\section{Coefficient bounds and reproducible checks}
\label{sec:v47-bounds}

In the seven contracted structures of \eqref{eq:v47-Kmap}, the map
\(Q_{\rm Ric}\mapsto c_E\mathcal K_\phi\) has norm one in the sum
of flavour-matrix entry norms. In an orthonormal normal frame, expand
all ordered tensor components and the product derivatives in
\eqref{eq:v47-adjoint}. With the jets of \(\phi,f,T\) and curvature
as labelled coefficient generators, the elementary bounds are
\begin{equation}
 \|\mathcal K_\phi\|_{\mathrm{raw,coeff}}
 \le32|c_E|^{-1}\|Q_{\rm Ric}\|_{\mathrm{coeff}},\qquad
 \|\mathcal K_\phi^\dagger\|_{\mathrm{raw,coeff}}
 \le64|c_E|^{-1}\|Q_{\rm Ric}\|_{\mathrm{coeff}}.
 \label{eq:v47-rawbounds}
\end{equation}
For the \(G\Delta R\) adjoint there are four diagonal outputs, four
Laplacian directions and product-rule coefficients \(1,2,1\), giving
64; the other six structures have no larger sum. These are coefficient
counts, not assumptions that physical field jets have magnitude one.
No bound uniform as \(c_E\to0\) is asserted.

For the coframe coordinate use
\(\eta_g=DG(q)^{-T}\eta_q\) and
\(\delta q=DG(q)^{-1}\delta g\). Derivatives in \(\mathcal K_\phi\)
act on these coordinate matrices and on \(\rho^{-1}\) as well.
On a fixed chart with density and Jacobian bounded away from zero,
each prescribed arity has a finite conversion constant independent of
the number of sites. At fixed external scale \(\mu\), with
\(0\le z\le M\mu^2\), the primitive's two-derivative coefficients
have windowed rooted kernel bounds
\(C_{N,b,M}|c_E|^{-1}\mu^2\), and the characteristics have bounds
\(C_{N,b,M}|c_E|^{-1}\mu^4\). The skew-scalar terms have their own
corresponding bounds without \(c_E^{-1}\). These follow from the
inherited fixed-polynomial Fourier-window estimate, not a new full-zone
or many-shell argument.

The checker reconstructs the curvature of both test metrics, the
projection and flavour matrices, the Ricci/stress trace sections, the
seven formal-adjoint identities, the mixed odd-antifield sign and the
nonlinear omitted-term test. Its finite integration-by-parts model is
an algebraic check, not a finite-lattice diffeomorphism theorem. The
coefficient compiler returns both an admissible scalar projection and
its actual unremoved defect. No loop quadrature or full invariant
cohomology calculation is performed.

The primitive has ghost number minus two, not zero. It is not installed
as an action counterterm and no new field redefinition or measure
Jacobian is performed in the integral. Its use in a ghost-number-one
normalization equation still requires the full descent and source
partners, as in V45--V46.

\section{Closing ledgers and the next interface}
\label{sec:v47-ledgers}

\subsection*{Proved}
\addcontentsline{toc}{subsection}{V47: Proved}
The complete scalar-linear projection of the declared coupled natural
sector has dimension 130. Two exact positive Ricci-flat metrics exclude
the remaining even and odd curvature-square operators. A mixed-antifield
section completes all 112 Ricci-dependent directions, with the adjoint
and nonlinear scalar terms explicit. Together with the inherited
skew-scalar primitive, this selects 124 exact completions and six known
rotation charges. The 28 invariant admissible scalar projections have
invariant exact completions. The residual metric and next-matter
equations and the fixed coefficient bounds are retained.

\subsection*{Inherited}
\addcontentsline{toc}{subsection}{V47: Inherited}
V1--V46 and all labels are unchanged. The ordinary signed Einstein and
scalar actions, V46's natural operator inventory, V45's rotation and
flavour results, the coframe cotangent convention, V34's normalization,
V38's logarithmic packet and V44's conditional fixed-gravity conversion
retain their previous scope. No old loop or cutoff theorem is claimed
as new.

\subsection*{Literature input}
\addcontentsline{toc}{subsection}{V47: Literature input}
The current interpretation and Euler-trivial mechanism are the standard
antifield framework \cite{V20BBH1994}. The mixed primitive and local
metric tensors are evaluated directly here. No classification of all
coupled currents or full invariant anomaly vanishing is imported as a
conclusion.

\subsection*{Conditional}
\addcontentsline{toc}{subsection}{V47: Conditional}
A native component that admits the stated natural matter-leading
representative can have its scalar-linear part removed by the explicit
section, modulo the known rotation projection. In the invariant sector
that projection is absent. This does not solve the residual metric
identity, the nonlinear scalar equation or the naturality reduction.
V44 applies to the fixed gravitational prescription only once the
absolute critical primitive in the full required category exists.

\subsection*{Open}
\addcontentsline{toc}{subsection}{V47: Open}
The next natural-current operands are
\eqref{eq:v47-metric-residual}--\eqref{eq:v47-next-matter}.
The larger problem also retains scalar-independent metric transformations,
non-natural tensor-valued descent components and the populated native
reference coefficient. The full translation-invariant, flavour-invariant,
mass-polynomial scalar--gravity critical primitive remains open.
Internal-gauge/chiral interfaces, higher loops and the invariant
interacting ESM class are separate. No active normalization changes.

\begin{center}\small
\begin{tabular}{>{\raggedright\arraybackslash}p{.24\textwidth}>{\raggedright\arraybackslash}p{.68\textwidth}}
\toprule
Variable & Scope in V47\\\midrule
Carrier & Ordinary Einstein gravity and four equal real scalars on the
inherited positive coefficient chart; the stated internal couplings and
cosmological source are zero.\\
Current category & Natural scalar-linear operator of weight four; metric
characteristic begins at scalar degree two. No full-current cohomology count.\\
Mass & Polynomial \(z=m^2\), regular at zero, with no inverse mass.\\
Coordinates & No explicit coordinate coefficients in charges or primitives;
coordinates occur only in the Ricci-flat test metrics.\\
Cutoff and volume & Fixed coefficient-map bounds independent of mesh,
periods and shell number, for \(|c_E|\ge c_*>0\). No new quantum cutoff rate.\\
Source order & Full local selected completion; fixed arity for coframe
and kernel constants. No uniform all-arity, all-loop or analytic RG bound.\\
Verification & Exact finite algebra and rational metric tests; no native
loop quadrature, full BV matrix or proof-assistant formalization.\\
\bottomrule
\end{tabular}
\end{center}


\clearpage
\part*{V48: Natural metric Noether reductions and the exact mass sectors}
\addcontentsline{toc}{part}{V48: Natural metric Noether reductions and the exact mass sectors}

\section{The residual metric identity, rather than another scalar projection}
\label{sec:v48-start}

Sections 1--368 and their labels are preserved. We use the ordinary
zero-mesh action, signed Euler normalization, and source convention of
\eqref{eq:v47-data}. There are four equal real scalars, \(z=m^2\) is a
polynomial variable, and the internal gauge, Yukawa, scalar self-couplings
and cosmological source remain zero. The finite action, contour, reference
profiles, auxiliary measure, V34 scalar normalization, V38 logarithmic
packet, and fixed gravitational prescription are unchanged.

The selected V47 current has already been subtracted, with its nonlinear
scalar term retained. The leading metric characteristic \(F_2\) of the
remaining current satisfies the actual identity
\begin{equation}
 \int \rho\,e^{ab}(F_2)_{ab}=0
 \quad\hbox{for every local metric and scalar field},\qquad
 e^{ab}=c_EG^{ab},\quad c_E\ne0.
 \label{eq:v48-input}
\end{equation}
This is an off-shell identity with the scalars acting as arbitrary
spectator fields. It is not enough that it hold for one scalar solution,
or only for a flat metric. Compact supports, or the inherited periodic
integration-by-parts convention, are understood.

The category is exactly the natural one used in V46--V47: coefficients are
constant flavour tensors and contractions of the metric, orientation,
curvature and covariant derivatives; they are polynomial in scalar jets
and \(z\), with no preferred spacetime tensor, explicit coordinate,
inverse mass, or curvature denominator. The metric characteristic is
homogeneous of scalar degree two and differential weight four; curvature
and \(z\) have weight two. These restrictions are part of the theorem,
not a naturality claim about every tensor-valued native reference mark.

The new result has three layers. First, the entire positive-mass part of
\eqref{eq:v48-input} is an explicitly generated diffeomorphism. Second,
its massless flat scalar-derivative symbol is also such a transformation,
with a polynomial section that has no inverse total momentum. Subtracting
their \emph{full current completions} leaves a massless curvature-dependent
metric identity. Third, when that remaining characteristic has no scalar
derivatives, its complete natural tensor inventory has one kernel
direction per symmetric flavour pair; this direction has an explicit
skew-Euler completion. The mixed curvature/differentiated-scalar case is
not silently removed by these three statements.

These are local current calculations inside the programme, not a new
quantum cutoff theorem. No theorem classifying all Einstein Noether
identities is invoked to replace the displayed finite constructions.

\section{Every positive-mass residual is an explicit diffeomorphism}
\label{sec:v48-mass}

Write
\[
 F_2=F^{[4]}+zF^{[2]}+z^2F^{[0]},
\]
where brackets here denote the derivative weight after removing the mass
power, not scalar degree. Since \(e\) is independent of \(z\), each
coefficient obeys \eqref{eq:v48-input} separately.

\begin{lemma}[Complete positive-mass tensor inventory]
\label{lem:v48-mass-inventory}
With full ordered metric contractions, the two positive-mass pieces have
unique coefficient form
\begin{align}
 F^{[2]}_{ab}={}&
 \phi^TA\nabla_a\nabla_b\phi
 +g_{ab}\phi^TB\Delta\phi
 +(\nabla_a\phi)^TC\nabla_b\phi
 +g_{ab}(\nabla_c\phi)^TD\nabla^c\phi \nonumber\\
 &+R_{ab}\phi^TE\phi+g_{ab}R\phi^TF\phi,\qquad
 F^{[0]}_{ab}=g_{ab}\phi^TG\phi .
 \label{eq:v48-mass-bank}
\end{align}
Here \(A,B\) are arbitrary real \(4\times4\) matrices and
\(C,D,E,F,G\) are symmetric. The bank has dimensions \(72+10=82\).
Its independent-flavour-sign-invariant restriction has dimension 28.
\end{lemma}
\begin{proof}
At derivative weight two, either both derivatives act on scalars, or a
single curvature multiplies two undifferentiated scalars. Symmetry of the
output indices and commutation of scalar fields give the first four
terms and the stated flavour symmetries. Every contraction of one
curvature to a symmetric two-tensor is a Ricci or scalar-curvature
contraction. A single orientation tensor supplies no additional
symmetric tensor at this weight: antisymmetry annihilates two scalar
momenta and the algebraic Bianchi identity annihilates the curvature
contraction. Weight zero leaves only the last metric multiple. Flavour
sign invariance restricts each matrix to its diagonal. Normal scalar and
curvature jets distinguish the listed coefficients.
\end{proof}

Let \(A_s=(A+A^T)/2\), \(B_s=(B+B^T)/2\), and
\(B_a=(B-B^T)/2\).

\begin{theorem}[Exact positive-mass metric Noether section]
\label{thm:v48-mass}
The positive-mass part of \eqref{eq:v48-input} holds if and only if
\begin{equation}
 \boxed{C=A_s,\qquad B=D=E=F=G=0.}
 \label{eq:v48-mass-conditions}
\end{equation}
For that coefficient,
\begin{equation}
 \boxed{
 zF^{[2]}=\mathcal L_{\xi_z}g,\qquad
 (\xi_z)_b=\frac z2\phi^TA\nabla_b\phi,
 \qquad F^{[0]}=0 .}
 \label{eq:v48-mass-xi}
\end{equation}
Thus the 82-dimensional bank has kernel dimension 16 and determining
rank 66. On the sign-invariant subspace the corresponding dimensions
are \(28,4,24\). This is a characteristic-space count, not a current
cohomology or anomaly count.
\end{theorem}
\begin{proof}
Divide the identity by \(c_E\). The contracted Bianchi identity and
\(\operatorname{tr}G=-R\) reduce its \(z\) coefficient, modulo a
divergence, to
\begin{align}
 &G^{ab}(\nabla_a\phi)^T(C-A_s)\nabla_b\phi
 +R(\nabla_a\phi)^T(B_s-D)\nabla^a\phi \nonumber\\
 &\quad +(\nabla_aR)\phi^TB_a\nabla^a\phi \nonumber\\
 &\quad +\phi^T\left[-\frac12(\Delta R)B_s
       +\left(R_{ab}R^{ab}-\frac12R^2\right)E-R^2F\right]\phi .
 \label{eq:v48-reduced-mass}
\end{align}
The first two terms give the second scalar-derivative coefficients of
the scalar Euler identity. Arbitrary Ricci jets first force
\(C=A_s\) and \(D=B_s\). The first-derivative part left by the third
term is \(2B_a\nabla R\cdot\nabla\phi\); arbitrary curvature-gradient
and scalar jets force \(B_a=0\). At fixed curvature, fourth metric
jets vary \(\Delta R\), forcing \(B_s=0\). Ricci tensors with
zero trace and nonzero square then force \(E=0\), and a nonzero
scalar curvature forces \(F=0\). The \(z^2\) coefficient is
\(-R\phi^TG\phi\), so \(G=0\). These are local off-shell tests;
none divides by curvature or assumes an Einstein solution.

Conversely, differentiating \eqref{eq:v48-mass-xi} gives exactly the
remaining two scalar-derivative terms in \eqref{eq:v48-mass-bank}.
The Bianchi identity makes its integral against \(e\) a divergence.
The free 16 entries of \(A\), or its four diagonal entries, give the
kernel counts. The independent determining map can be written as
\[
 (A,B,C,D,E,F,G)\longmapsto
 (C-A_s,D-B_s,B_a,B_s,E,F,G),
\]
whose output dimension is \(10+10+6+10+10+10+10=66\).
\end{proof}

\begin{proposition}[A mass-sector compiler with a retained error]
\label{prop:v48-mass-compiler}
Let \(\Pi_m\) keep \(A\), replace \(C\) by \(A_s\), and set the
other five matrices to zero. In independent-entry coefficient norms,
\begin{equation}
 \boxed{
 \Pi_m^2=\Pi_m,\qquad
 \|\Pi_m\|_{1\to1}=2,\qquad
 \|I-\Pi_m\|_{1\to1}=1 .}
 \label{eq:v48-mass-norms}
\end{equation}
Its image has identically zero Noether variation. Applying the compiler
to a general array therefore leaves its \emph{actual} determining defect
in the unremoved array. The vector coefficient map \(A\mapsto A/2\)
has norm \(1/2\). None of these constants depends on \(c_E\), the
mesh, mass value, or volume.
\end{proposition}
\begin{proof}
Antisymmetrization and symmetrization are contractions in the expanded
matrix-entry norm. In the declared independent-entry basis a diagonal
entry of \(A\) is copied once into \(C\), giving the column sum two;
o other column is larger. In \(I-\Pi_m\), the \(A\)-row vanishes
and its contribution occurs only in \(C-A_s\), with column sum at
most one. A discarded independent coefficient attains one. Idempotence
and defect preservation follow from \Cref{thm:v48-mass}.
\end{proof}

\section{The massless flat symbol also has a pole-free gauge section}
\label{sec:v48-flat}

Linearize \eqref{eq:v48-input} about the flat metric. Since \(e(\delta)=0\),
self-adjointness of the free native Euler operator implies
\begin{equation}
 E(p+q)\,\mathsf F(p,q)=0,
 \label{eq:v48-flat-test}
\end{equation}
where \(\mathsf F\) is the scalar-bilinear flat symbol of \(F^{[4]}\).
Its two scalar momenta are independent; only the total transfer enters
\(E\). Common powers of \(i\) are suppressed in the following
homogeneous differential-symbol algebra and restored when forming real
operators.

Put \(a=p^2\), \(b=q^2\), and \(c=p\cdot q\). In a symmetric
flavour-pair channel, the complete natural symbol is
\begin{equation}
 \begin{split}
 \mathsf F={}&u(a,b,c)pp^T+u(b,a,c)qq^T
 +v(a,b,c)(pq^T+qp^T)+d(a,b,c)I,\\
 u={}&u_1a+u_2b+u_3c,\qquad
 v=v_1(a+b)+v_2c,\\
 d={}&d_1(a^2+b^2)+d_2ab+d_3c(a+b)+d_4c^2 .
 \end{split}
 \label{eq:v48-even-symbol}
\end{equation}
This has nine coefficients. In an antisymmetric flavour-pair channel,
replace the second term by \(-u(b,a,c)qq^T\) and use
\begin{equation}
 v=v_1(a-b),\qquad d=d_1(a^2-b^2)+d_2c(a-b).
 \label{eq:v48-odd-symbol}
\end{equation}
This has six coefficients. No orientation-odd symmetric two-tensor can
be formed from only the two momenta and the metric; antisymmetry either
leaves an antisymmetric output or repeats a momentum in an orientation
tensor.

\begin{theorem}[Complete flat-spectator symbol reduction]
\label{thm:v48-flat}
Write the three rank-two coefficients in either channel as
\(u,v,w\), with \(w=\pm u(b,a,c)\). Then
\begin{equation}
 \boxed{
 E(p+q)\mathsf F=0
 \quad\Longleftrightarrow\quad
 d=0,\qquad v=\frac{u+w}{2}.}
 \label{eq:v48-flat-kernel}
\end{equation}
The polynomial vector symbol
\begin{equation}
 \boxed{
 \mathsf v(p,q)=\frac12\bigl(u(p,q)p+w(p,q)q\bigr)
 }
 \label{eq:v48-flat-vector}
\end{equation}
satisfies \(\mathsf F=(p+q)\mathsf v^T+\mathsf v(p+q)^T\).
Each flavour channel has three gauge parameters. For four scalars the
raw symbol space has dimension \(9\cdot10+6\cdot6=126\), its kernel
has dimension 48, and its determining rank is 78. The sign-invariant
restriction has dimensions 36 and 12.
\end{theorem}
\begin{proof}
For real \(k\ne0\), choose an orthonormal frame with \(k\) on one
axis. Direct evaluation of the unreduced free Einstein symbol shows
\(\ker E(k)=\{kv^T+vk^T\}\). This uses no gauge-fixed inverse.
For an open set of independent \(p,q\), a vector perpendicular to
both is also perpendicular to \(k=p+q\). Its diagonal contraction
with a pure-gauge tensor is zero, so \(d=0\). The remaining tensor
lies in the \(p,q\) plane. A pure-gauge vector producing it has no
component perpendicular to that plane. Comparing the three independent
coefficients \(pp^T,pq^T+qp^T,qq^T\) gives \(u+w=2v\).
Polynomial identities on this open set hold everywhere, including
\(p+q=0\). Conversely \eqref{eq:v48-flat-vector} is a polynomial
solution on every momentum. It is not obtained by evaluating
\(1/(p+q)^2\).

In the symmetric channel the conditions read
\(v_1=(u_1+u_2)/2\), \(v_2=u_3\), and \(d_j=0\).
In the antisymmetric channel they read
\(v_1=(u_1-u_2)/2\) and \(d_j=0\). The counts follow.
\end{proof}

\begin{proposition}[Natural covariantization and fixed coefficient bounds]
\label{prop:v48-covariantization}
Keeping \(u,w\), replacing \(v\) by \((u+w)/2\), and removing \(d\)
is an idempotent symbol projection. In the displayed symmetry-adapted
bases its norm is at most two and its complementary map has norm one.
Its polynomial vector has a fixed natural covariant realization
\(\xi_0(g,\phi)\), scalar-quadratic and of derivative weight three.
\end{proposition}
\begin{proof}
The column sums of the displayed maps give the bounds. Covariantize each
vector monomial by the fixed ordered replacements
\[
 p_b p^2\rightsquigarrow(\nabla_b\Delta\phi_A)\phi_B,
 \qquad
 p_b q^2\rightsquigarrow(\nabla_b\phi_A)\Delta\phi_B,
 \qquad
 p_b(p\cdot q)\rightsquigarrow
 (\nabla_b\nabla_c\phi_A)\nabla^c\phi_B,
\]
and the exchanged expressions. Restore the common Fourier factors to
obtain real differential operators. Every term is tensorial and its flat
symbol is exactly \eqref{eq:v48-flat-vector}. Covariant derivatives in
this prescription have a specified order; changing that order introduces
curvature terms and is not silently done. Fixed-order coframe conversion
has finite coefficient bounds on the inherited small chart.
\end{proof}

\section{The gauge subtraction includes its nonlinear scalar and ghost-source data}
\label{sec:v48-gauge}

Let \(\xi=\xi_0+\xi_z\) be the preceding natural vector. Its full
ordinary characteristic is
\begin{equation}
 \boxed{
 \delta_\xi g=\mathcal L_\xi g,\qquad
 \delta_\xi\phi=\xi^a\nabla_a\phi .}
 \label{eq:v48-full-gauge}
\end{equation}
The second term is scalar-cubic. It is not optional even though the
leading metric identity \eqref{eq:v48-input} contains no scalar Euler
term.

\begin{theorem}[Exact local gauge completion and mass/flat reduction]
\label{thm:v48-reduction}
In the inherited adjoint-generator convention, the local functional
\begin{equation}
 \boxed{
 \mathcal T_\xi=\int\sigma_a\xi^a(g,\phi),\qquad \sigma=c^*,
 }
 \label{eq:v48-gauge-primitive}
\end{equation}
has ghost number minus two and weight six, and obeys
\begin{equation}
 \boxed{
 s\mathcal T_\xi
 =\int\eta_g:\mathcal L_\xi g
   +\int\upsilon^T\mathcal L_\xi\phi
 \quad\pmod{d_x}.}
 \label{eq:v48-gauge-charge}
\end{equation}
For every \(F_2\) satisfying the stated natural identity,
\begin{equation}
 \boxed{
 F_2=\mathcal L_\xi g+F_{\rm curv},\qquad
 \partial_zF_{\rm curv}=0,\qquad
 F_{\rm curv}|_{g=\delta}=0,
 \qquad
 \int\rho e:F_{\rm curv}=0.}
 \label{eq:v48-final-reduction}
\end{equation}
The residual is therefore a massless natural curvature-dependent
characteristic. Every normal tensor monomial in it contains at least
one curvature tensor or a covariant curvature derivative.
\end{theorem}
\begin{proof}
The action part of \(s\sigma\) is the adjoint of the full ordinary
metric--scalar gauge generator. Pairing with \(\xi\) gives both
terms in \eqref{eq:v48-gauge-charge}. The remaining integrated term is
\(\int\sigma_a(\gamma\xi^a-[c,\xi]^a)\), which vanishes because
\(\xi\) is a natural contravariant vector. This also proves exactness
without treating a field-dependent vector as a field-independent
parameter. No explicit coordinate primitive is introduced.

The native Ward identities are
\(\nabla_ae^{ab}=0\) and
\(\nabla_a\tau^{ab}=\tfrac12f^T\nabla^b\phi\).
They give the direct variational check
\[
 (e+\tau):\mathcal L_\xi g+f^T\mathcal L_\xi\phi
 =2\nabla_a\bigl((e+\tau)^{ab}\xi_b\bigr).
\]
Subtracting the full charge therefore preserves every ordinary current
equation, including the next matter equation of V47.
\Cref{thm:v48-mass} removes all positive mass powers exactly.
\Cref{thm:v48-flat} and its ordered covariantization remove the full
massless flat scalar symbol. The remaining terms in a natural normal
tensor expansion have a curvature factor; scalar-derivative commutators
supply only further such factors. The Bianchi identity shows that the
subtracted metric term itself obeys \eqref{eq:v48-input}, completing
\eqref{eq:v48-final-reduction}.
\end{proof}

The conclusion is about the leading metric characteristic of a full
current. It does not assert that its higher matter coefficients vanish.
In particular, subtracting \(\mathcal L_\xi g\) without the
\(\upsilon\mathcal L_\xi\phi\) partner does not perform the
BV subtraction in \eqref{eq:v48-gauge-charge}. No ghost-number-zero
quantum counterterm is installed by this current reduction.

\section{The scalar-undifferentiated terminal sector has one exact direction}
\label{sec:v48-terminal}

There is a useful complete subcase of the remaining massless identity.
Suppose no derivative acts on either scalar, so
\[
 (F_{\rm curv})_{ab}=\sum_{A,B}\phi_A\phi_B H^{AB}_{ab}(g),
 \qquad H^{AB}=H^{BA},
\]
where \(H\) is a natural symmetric metric tensor of weight four.
Since the scalars are arbitrary, each coefficient satisfies
\(e^{ab}H^{AB}_{ab}=0\) pointwise.

Put
\[
 S_{ab}=R_{ab}-\tfrac14Rg_{ab},\qquad
 C_{abcd}=\text{Weyl tensor},\qquad
 (CS)_{ab}=C_{acbd}S^{cd}.
\]
The star is the Euclidean Hodge star on the first Weyl pair. The complete
natural tensor inventory, including orientation-odd tensors, consists
of the three derivative tensors
\begin{equation}
 \nabla_a\nabla_bR,\qquad g_{ab}\Delta R,\qquad\Delta R_{ab},
 \label{eq:v48-derivative-curv-bank}
\end{equation}
and the eight algebraic tensors
\begin{equation}
 \begin{split}
 RS,&\quad (S^2)_0,\quad gR^2,\quad g|S|^2,\quad
 CS,\quad g|C|^2,\quad ({}^\star C)S,\quad
 g\langle C,{}^\star C\rangle .
 \end{split}
 \label{eq:v48-algebraic-curv-bank}
\end{equation}
Here \((S^2)_0=S^2-\tfrac14g|S|^2\).

\begin{lemma}[Completeness of the terminal inventory]
\label{lem:v48-terminal-bank}
Every such \(H\) is a linear combination of
\eqref{eq:v48-derivative-curv-bank}--\eqref{eq:v48-algebraic-curv-bank}.
The eleven tensors are independent as universal natural tensors.
\end{lemma}
\begin{proof}
At weight four the geometric ingredients are two covariant derivatives
of one curvature, or two curvatures. Contracted differential Bianchi
identities and derivative commutators reduce the first case to the three
displayed derivative tensors plus curvature quadratics. The
orientation-odd double-derivative symbols vanish by the differential
Bianchi identity; commutation leaves only curvature quadratics.
For the second case insert the orthogonal \(C,S,R\) decomposition.
A curvature contraction to two free symmetric indices gives the first
five algebraic tensors and their two orientation-odd counterparts.
The remaining pure Weyl contractions obey the four-dimensional algebraic
identities
\[
 C_{acde}C_b{}^{cde}=\tfrac14g_{ab}|C|^2,
 \qquad
 C_{(a|cde|}{}^\star C_{b)}{}^{cde}
 =\tfrac14g_{ab}\langle C,{}^\star C\rangle.
\]
They follow by antisymmetrization over five indices, and are also
checked on the complete ten-component Weyl basis and its polarizations
in the certificate. This accounts for all metric and orientation
contractions. Independent Ricci and Weyl jets separate the eight
algebraic tensors. Independent fourth metric jets at fixed second jet
separate the three derivative tensors, as made explicit in the next
proof. No integration by parts moves a derivative onto a scalar in this
pointwise inventory.
\end{proof}

\begin{theorem}[Complete terminal Noether kernel]
\label{thm:v48-terminal}
The pointwise equation \(G^{ab}H_{ab}=0\) in this eleven-dimensional
inventory has the one-dimensional kernel
\begin{equation}
 \boxed{
 H_{ab}=\kappa\left(RS_{ab}+g_{ab}|S|^2\right)
 =\kappa\left[R R_{ab}
   +g_{ab}\left(R_{cd}R^{cd}-\tfrac12R^2\right)\right].}
 \label{eq:v48-terminal-H}
\end{equation}
Its determining rank is ten. For four real scalars there are ten
independent symmetric flavour pairs: the corresponding raw dimension
is 110, its kernel is ten, and its determining rank is 100. The
sign-invariant dimensions are \(44,4,40\).
\end{theorem}
\begin{proof}
First vary fourth metric derivatives while fixing the lower metric jet.
The background Einstein tensor is then independent of the variation.
For a quartic perturbation \(h_{ab}(k\cdot x)^4/24\), the highest
symbol of the derivative part is
\[
 a\,k_ak_b\delta R+b\,g_{ab}k^2\delta R
      +c\,k^2\delta R_{ab}.
\]
Take \(k=e_0\). A spatial trace-free \(h\), tested against an
independent spatial trace-free Einstein tensor, first forces \(c=0\).
A spatial trace variation has \(\delta R\ne0\); varying the
background tensor forces \(a=b=0\). The certificate gives the exact
three-by-three test matrix. These are realized fourth metric jets, not
an independent curvature model.

For the algebraic part use coefficients \((a,b,c,d,e,f,g,h)\) in the
order \eqref{eq:v48-algebraic-curv-bank}. Its contraction with
\(G=S-Rg/4\) is
\begin{equation}
 b\operatorname{tr}S^3+(a-d)R|S|^2-cR^3
 +e\langle S,CS\rangle-fR|C|^2
 +g\langle S,{}^\star CS\rangle
 -hR\langle C,{}^\star C\rangle .
 \label{eq:v48-seven-invariants}
\end{equation}
Set \(S=0\) and vary \(R\) and the two independent Weyl
chiralities to get \(c=f=h=0\). Set \(R=C=0\) and use a
trace-free \(S\) with nonzero cubic trace to get \(b=0\).
For \(S=\operatorname{diag}(1,-1,0,0)\), an electric Weyl basis
element gives the pair of contractions \((-2,0)\), while its Hodge
dual gives \((0,-2)\). Thus \(e=g=0\). Finally \(a=d\).
This proves both necessity and the claimed one-dimensional kernel.
Substitution verifies sufficiency. Every algebraic curvature used here
is the curvature of a local metric second jet.
\end{proof}

A residual-preserving coefficient section takes
\(\kappa=(a+d)/2\), discards the other entries, and keeps those
entries in its reported residual. Its coefficient projection and
complement each have induced norm one. This section is valid for an
\emph{independently closed} scalar-undifferentiated input. For a sum
with scalar-derivative terms, those terms can contribute after
integration by parts: the theorem does not authorize projecting them
out and applying a pointwise test to the remaining part alone.

\section{The terminal direction needs a four-matter metric completion}
\label{sec:v48-terminal-primitive}

Lower the Euler tensor with \(g\) when forming the following pointwise
symmetric-tensor operator. For any symmetric tensor \(T\), put
\begin{equation}
 \boxed{
 \Omega_e(T)=g(e:T)-e\operatorname{tr}_gT,
 \qquad \Omega_e^\dagger=-\Omega_e.
 }
 \label{eq:v48-Omega}
\end{equation}
The adjoint uses the same full metric contraction and \(\int\rho\).
For a constant symmetric flavour matrix \(M\), let
\(a_M(\phi)=\phi^TM\phi\).

\begin{theorem}[Exact metric-antifield completion]
\label{thm:v48-terminal-primitive}
The characteristic
\begin{equation}
 \boxed{
 \delta g=\frac{a_M}{c_E^2}\Omega_e(e+\tau),\qquad
 \delta\phi=0
 }
 \label{eq:v48-skew-characteristic}
\end{equation}
has leading scalar degree two equal to
\(a_M(RS+g|S|^2)\). Its degree-four metric term is
\(a_Mc_E^{-2}\Omega_e\tau\). The full charge is BV-exact, with
\begin{equation}
 \boxed{
 \mathcal T_M=-\frac1{2c_E^2}
 \int\rho\,a_M\,\vartheta:\Omega_e\vartheta,
 \qquad \vartheta=\rho^{-1}\eta_g .}
 \label{eq:v48-terminal-BV}
\end{equation}
The primitive has ghost number minus two and weight six. It is local,
translation invariant, natural and polynomial in mass.
\end{theorem}
\begin{proof}
The two elementary identities are
\[
 \langle T_1,\Omega_eT_2\rangle
 =-\langle\Omega_eT_1,T_2\rangle,
 \qquad
 c_E^{-2}\Omega_e e=RS+g|S|^2.
\]
The first proves
\((e+\tau):\delta g=0\) pointwise. For the second, use
\(e=c_E(S-Rg/4)\), \(\operatorname{tr}e=-c_ER\), and
\(|e|^2=c_E^2(|S|^2+R^2/4)\).
The odd product in \eqref{eq:v48-terminal-BV}, with
\(\delta\vartheta=e+\tau\), gives
\(\int\eta_g:\delta g\). Its longitudinal part is a total
divergence because the entire density is natural. No scalar Euler term
has been suppressed; there is no scalar antifield in this primitive.
\end{proof}

The separate matter-degree cancellations are
\begin{equation}
 \begin{gathered}
 e:\Omega_e e=0,\qquad
 e:\Omega_e\tau+\tau:\Omega_e e=0,\qquad
 \tau:\Omega_e\tau=0.
 \label{eq:v48-degree-cancellations}
 \end{gathered}
\end{equation}
They exhibit explicitly why the degree-four metric term cannot be
omitted. At one point take \(c_E=1\), \(a_M=1\),
\(e=\operatorname{diag}(1,0,0,0)\), and a massless scalar gradient
along \(e_1\), so
\(\tau=I/4-e_1e_1^T/2\). These independent local jets give
\begin{equation}
 \boxed{
 \tau:\Omega_e e=\frac14,
 \qquad e:\Omega_e\tau=-\frac14 .}
 \label{eq:v48-nonzero-return}
\end{equation}
The incomplete leading characteristic therefore has a genuine
four-matter variation. The complete primitive cancels it exactly.

\section{Coefficient bounds, source ownership and verification}
\label{sec:v48-bounds}

The mass section has the evaluated norms in
\eqref{eq:v48-mass-norms}. The massless flat symbol projection has norm
at most two in its symmetry-adapted coefficient bases; its complementary
map has norm one. The terminal algebraic section has norm one.
Covariantization and the cotangent transformation back to \(q\) are
fixed finite differential-polynomial maps at each prescribed arity.
They do not change the finite integration measure.

For the terminal operator, in an orthonormal frame,
\begin{equation}
 \|\Omega_e T\|_{\mathrm{HS}}
 \le4\|e\|_{\mathrm{HS}}\|T\|_{\mathrm{HS}},
 \label{eq:v48-Omega-bound}
\end{equation}
because \(\|g\|_{\mathrm{HS}}=2\) and
\(|\operatorname{tr}T|\le2\|T\|_{\mathrm{HS}}\).
Its normalized section is uniformly bounded on every declared jet ball
with \(|c_E|\ge c_*>0\). This is not a bound independent of the
size of arbitrary input curvature jets.

At a fixed Fourier source window, each local polynomial of derivative
degree \(d\) has a weighted rooted kernel bounded by
\(C_{N,b}\mu^d\) times its coefficient norm, with one source in
\(L^1\) and the others in \(L^\infty\). Periodization preserves the
bound. This follows by Fourier integration by parts on the fixed smooth
window, not from an extensive supremum norm of an action. Constants are
independent of mesh and periods under the usual common-window sampling
condition; \(N,b\) and the profile derivative bounds are fixed. No
all-arity or many-shell conclusion is drawn.

The script \texttt{check\_ESM\_metric\_Noether\_V48.py} constructs the
82-column mass map and its invariant restriction, both flat-symbol
sections, the native free-Euler gauge identity, the complete terminal
Weyl contraction tests, and the metric-antifield skew adjoint. It retains
exact rational rank certificates and coefficient-map norms. The mass
compiler accepts all seven input matrices and returns both the gauge
projection and the actual unremoved coefficients. It certifies the
mass-sector identity only if the determining array is zero.

The curvature tests use a complete ten-component algebraic Weyl basis,
independent Ricci jets, and realized highest metric jets. The proofs of
natural-inventory completeness and of the source topology are given in
the text; finite test counts are not substitutes for them. No native
loop integral, full Noether matrix, or uncomputed critical reference
array is evaluated here. General BV/Noether terminology and the role of
the Euler differential are literature background \cite{V20BBH1994};
none of the new kernels is inferred from a blanket completeness claim.

\section{Closing ledgers and the remaining curvature-dependent equation}
\label{sec:v48-ledgers}

\subsection*{Proved}
\addcontentsline{toc}{subsection}{V48: Proved}
The natural positive-mass leading metric identity has exactly the
16-dimensional diffeomorphism kernel in \eqref{eq:v48-mass-conditions};
its invariant restriction has dimension four. The complete massless
flat scalar-bilinear symbol has a polynomial gauge section, with 48
parameters for four scalars. Subtracting the full natural gauge charge
removes every positive-mass term and the flat symbol, leaving a massless
curvature-dependent identity. In the scalar-undifferentiated terminal
sector the complete eleven-tensor inventory has one kernel direction
per symmetric flavour pair. Its exact metric-antifield completion,
including its four-matter metric term, and all stated finite-map bounds
are established.

\subsection*{Inherited}
\addcontentsline{toc}{subsection}{V48: Inherited}
V1--V47 and all their mathematical labels are unchanged. The signed
ordinary action, stress, minimal-source and coframe conventions are
retained from V47. Its selected Ricci-current subtraction is an input,
not repeated as a new result. V34's active scalar prescription, V38's
logarithmic packet and the fixed gravitational normalization are
unchanged. No earlier analytic or cohomological claim is silently
strengthened by this appendix.

\subsection*{Literature input}
\addcontentsline{toc}{subsection}{V48: Literature input}
The distinction between gauge characteristics, Euler-trivial
characteristics and the full BV differential is standard
\cite{V20BBH1994}. The concrete sections here use the Bianchi identity,
tensor contractions and polynomial free symbols. No classification of
all coupled currents or of the invariant anomaly complex is imported.

\subsection*{Conditional}
\addcontentsline{toc}{subsection}{V48: Conditional}
For a full native current that admits V47's natural matter-leading
representative, these exact subtractions preserve its charge class and
remove the specified leading components. A tensor-valued or non-natural
reference component is not assumed to admit that representative.
Likewise the terminal pointwise solver applies only when its complete
input contains no scalar derivatives; it cannot be applied to an
isolated portion of a larger integrated identity.

\subsection*{Open}
\addcontentsline{toc}{subsection}{V48: Open}
The remaining leading metric problem is the massless natural identity
\[
 \int\rho e:F_{\mathrm{curv}}=0,
 \qquad F_{\mathrm{curv}}|_{g=\delta}=0,
\]
with terms involving curvature and differentiated scalars retained together
with the corresponding undifferentiated scalar terms. The previous exact sections
do not prove its general decomposition into field-dependent
diffeomorphisms and skew-Euler characteristics. Higher matter currents,
scalar-independent metric transformations, and reduction of the actual
tensor-valued descent to the declared natural category remain separate.
The critical scalar--gravity primitive must still be constructed with
translation invariance, flavour invariance and polynomial mass dependence.
Internal gauge and chiral interfaces, higher loops, and invariance of the
interacting ESM class are separate unresolved tasks.

\begin{center}\small
\begin{tabular}{>{\raggedright\arraybackslash}p{.24\textwidth}>{\raggedright\arraybackslash}p{.68\textwidth}}
\toprule
Variable & Scope in V48\\\midrule
Carrier & Same ordinary Einstein action and four equal real scalars;
the declared internal couplings and cosmological source remain zero.\\
Mass & Polynomial \(z=m^2\); exact positive-mass leading-metric
reduction; no inverse mass and no infrared-scale limit.\\
Naturality & No preferred spacetime tensors or explicit coordinates in
charges; coordinate polynomials occur only in realized metric-jet tests.\\
Cutoff and volume & Fixed coefficient-map bounds independent of mesh,
periods and number of shells. No new quantum cutoff rate.\\
Input domain & \(|c_E|\ge c_*>0\) and a fixed metric/jet chart for
coframe and pointwise operator bounds. Input coefficients need not be small.\\
Source/perturbative order & Exact local selected completions; fixed arity
for windowed kernel constants. No uniformity in unbounded source order or
loop order.\\
Verification & Exact finite matrices, polynomial identities, tensor
contractions and preservation audit; no native loop quadrature or
proof-assistant formalization.\\
\bottomrule
\end{tabular}
\end{center}

\clearpage
\part*{V49: A complete invariant curvature--spectator Noether section}
\addcontentsline{toc}{part}{V49: A complete invariant curvature--spectator Noether section}

\section{The curvature-dependent leading metric identity is closed in the invariant sector}
\label{sec:v49-start}

Sections 1--376 are preserved. We retain the ordinary positive coefficient
chart, signed Einstein normalization \(e^{ab}=c_EG^{ab}\), and four equal
real scalars of \eqref{eq:v47-data}. The gauge, Yukawa and scalar
self-couplings and the cosmological source remain zero. No action,
reference, contour, measure or active normalization is changed.

The input is precisely the unresolved leading-metric identity of V48,
\begin{equation}
 \int\rho\,e^{ab}F_{ab}(g,\phi)=0\quad\hbox{for every }g,\phi,
 \qquad \partial_zF=0,\qquad F|_{g=\delta}=0 .
 \label{eq:v49-input}
\end{equation}
Here \(F\) is natural, scalar quadratic and of differential weight four.
It is a polynomial in scalar and curvature jets, with no preferred
spacetime tensor, explicit coordinate coefficient, inverse mass or
curvature denominator. This installment solves the complete
\emph{independent-flavour-sign-invariant} input. Such an input is a sum
of four diagonal scalar channels. More generally the same construction
works on each symmetric flavour-pair channel by polarization. No
classification of the antisymmetric flavour channels is claimed.

The principal simplification is to replace a symmetric scalar pair by
its product \(u=\phi^2\), while retaining the separate scalar-gradient
quadratic term. This turns most of the problem into an order-two
operator acting on an arbitrary scalar test. Its first-curvature
identity has a small finite polynomial test. Every surviving principal
coefficient has an exact covariant gauge or skew-Euler completion; thus
no higher-curvature compatibility assumption is needed to lift this
particular linear test. What remains after that subtraction is exactly
the terminal algebraic sector already solved in V48.

The conclusion is a classification of a \emph{leading metric
characteristic}, not of all coupled current cohomology and not of the
critical quantum breaking. It supplies its full local BV subtraction,
including the higher-matter partner, but leaves later current equations
and naturality of the complete native descent to be established.

\section{A thirty-coefficient inventory and the spectator-product reduction}
\label{sec:v49-inventory}

We give one diagonal scalar channel first. All contractions use the
ordinary metric. Put \(S_{ab}=R_{ab}-Rg_{ab}/4\), and let \(C\) be the
Weyl tensor with \(\epsilon_{0123}=1\). For a symmetric tensor \(T\),
define the seven curvature-linear maps
\begin{equation}
\begin{array}{lll}
 M_1(T)=RT,&M_2(T)=gR\operatorname{tr}T,&M_3(T)=\operatorname{Ric}\operatorname{tr}T,\\
 M_4(T)=g\langle\operatorname{Ric},T\rangle,
 &M_5(T)=\operatorname{Ric}\,T+T\operatorname{Ric},
 &M_6(T)_{ab}=C_{acbd}T^{cd},\\
 &&M_7(T)_{ab}=({}^\star C)_{acbd}T^{cd}.
\end{array}
\label{eq:v49-Mbank}
\end{equation}
For a covector \(v\), define
\begin{align}
 N_1(v)_{ab}&=(\nabla_aR)v_b+(\nabla_bR)v_a,
 &N_2(v)_{ab}&=g_{ab}\nabla^cR\,v_c,\nonumber\\
 N_3(v)_{ab}&=(\nabla^cR_{ab})v_c,
 &N_4(v)_{ab}&=(\nabla_aR_b{}^c+\nabla_bR_a{}^c)v_c,\nonumber\\
 N_5(v)_{ab}&=\frac12\nabla^d
 \bigl[({}^\star C)_{acbd}+({}^\star C)_{bcad}\bigr]v^c.
 \label{eq:v49-Nbank}
\end{align}
Let \((D_1,\ldots,D_{11})\) be, in order,
\begin{equation}
\begin{split}
 &\nabla_a\nabla_bR,\quad g_{ab}\Delta R,\quad\Delta R_{ab},\\
 &RS,\quad (S^2)_0,\quad gR^2,\quad g|S|^2,\quad
 CS,\quad g|C|^2,\quad({}^\star C)S,\quad
 g\langle C,{}^\star C\rangle .
\end{split}
\label{eq:v49-Dbank}
\end{equation}
The last list and its universal tensor identities are inherited from
\Cref{lem:v48-terminal-bank}. In particular its scalar-undifferentiated
Noether kernel is inherited, not recomputed as a new result.

\begin{lemma}[Complete curvature inventory for a symmetric scalar pair]
\label{lem:v49-inventory}
Every \(F\) in the one-scalar category of \eqref{eq:v49-input} has
unique coefficients \((a,b,c,d)\in\R^7\oplus\R^7\oplus\R^5\oplus\R^{11}\)
in
\begin{equation}
 F=\sum_{i=1}^7a_i\phi M_i(\nabla^2\phi)
   +\sum_{i=1}^7b_iM_i(\nabla\phi\otimes\nabla\phi)
   +\sum_{j=1}^5c_j\phi N_j(\nabla\phi)
   +\phi^2\sum_{l=1}^{11}d_lD_l .
 \label{eq:v49-raw}
\end{equation}
There are thirty coefficients. The invariant four-scalar sector is the
direct sum of four such inventories.
\end{lemma}
\begin{proof}
Every normal tensor monomial contains curvature and has weight four.
Two scalar derivatives therefore leave exactly one curvature; one
scalar derivative leaves one differentiated curvature; no scalar
derivative leaves weight-four metric tensors. Covariant-derivative
commutators are retained when passing to this normal tensor form.

For two derivatives the symmetric scalar Hessian or gradient product is
the only scalar tensor. The scalar-curvature part has the two maps
\(RT,gR\operatorname{tr}T\). The trace-free Ricci part has three
maps: its product with \(\operatorname{tr}T\), its contraction times
\(g\), and its symmetric product with \(T\). The two Weyl
chiralities supply \(C(T)\) and \({}^\star C(T)\). This gives
\eqref{eq:v49-Mbank}. The absence of further contractions follows also
directly by inserting the Weyl--Ricci--scalar decomposition in one
Riemann tensor. An orientation tensor cannot produce an additional
Ricci/symmetric-scalar map: after symmetrizing the two output indices,
one of the two symmetric input pairs would have to occupy two
antisymmetric slots.

With one scalar derivative, the contracted differential Bianchi
identities reduce every even contraction of \(\nabla\operatorname{Riem}\)
to the four placements in \eqref{eq:v49-Nbank}. In the odd sector,
the dual Bianchi identities leave one symmetrized dual-Cotton
placement, written as \(N_5\); trace and symmetrization remove the
other placements. The full undifferentiated list is
\eqref{eq:v49-Dbank}. Independent scalar jets separate the three
scalar-derivative distributions. Normal metric jets separate their
curvature and curvature-gradient coefficients. Independence of the
fifteen linear-curvature placements is additionally verified by the
explicit free-symbol tests below and the six independent covariant
kernel columns. The eight algebraic curvature terms are the inherited
independent terminal list. Sign invariance prohibits a term containing
one occurrence each of two different scalar flavours, so the four
channels are independent.
\end{proof}

Put \(u=\phi^2\), \(V=\nabla\phi\otimes\nabla\phi\), and
\(r=b-a\). The exact scalar product rules give
\begin{equation}
 F=\frac12\sum_i a_iM_i(\nabla^2u)
   +\frac12\sum_jc_jN_j(\nabla u)
   +u\sum_ld_lD_l+\sum_i r_iM_i(V).
 \label{eq:v49-product}
\end{equation}
No derivative of a curvature coefficient is discarded by this rewriting.

For symmetric tensors \(S,T\), write
\begin{equation}
 \Omega_S(T)=g\langle S,T\rangle-S\operatorname{tr}T,
 \qquad \Omega_S^\dagger=-\Omega_S.
 \label{eq:v49-Omega}
\end{equation}

\begin{lemma}[Pointwise curvature syzygy]
\label{lem:v49-algebraic}
For arbitrary curvature and arbitrary symmetric \(T\),
\begin{equation}
 G:\sum_i r_iM_i(T)=0
 \quad\Longleftrightarrow\quad
 r=\sigma(1,-\tfrac12,0,1,0,0,0).
 \label{eq:v49-alg-syzygy}
\end{equation}
The resulting tensor is \(\sigma\Omega_T(G)\).
\end{lemma}
\begin{proof}
The Weyl and dual-Weyl coefficients vanish independently by varying the
two Weyl chiralities at a fixed nonzero trace-free Ricci tensor, and
then varying \(T\). In the Ricci part, the independent contractions
\(\operatorname{tr}(\operatorname{Ric}^2)\operatorname{tr}T\)
and \(\operatorname{tr}(\operatorname{Ric}^2T)\) force
\(r_3=r_5=0\). The coefficients of
\(R\langle\operatorname{Ric},T\rangle\) and
\(R^2\operatorname{tr}T\) give \(r_4=r_1\) and
\(r_2=-r_1/2\). Conversely the displayed combination is
\(g(G:T)-T\operatorname{tr}G\), whose contraction with \(G\)
is zero. The certificate supplies six independent rational curvature
and tensor tests; their six-column minor has determinant 6. All tested
algebraic curvatures are realized by local metric second jets.
\end{proof}

\begin{proposition}[A principal scalar equation before any metric expansion]
\label{prop:v49-spectator}
If \eqref{eq:v49-input} holds, \(r=b-a\) obeys
\eqref{eq:v49-alg-syzygy}. After subtracting
\(\sigma c_E^{-1}\Omega_Ve\), the remaining characteristic has
the form \(\mathscr D_g u\), where \(\mathscr D_g\) is the
order-two scalar-to-metric operator in the first three terms of
\eqref{eq:v49-product}, and
\begin{equation}
 \mathscr D_g^\dagger e=0
 \quad\hbox{as a local scalar identity}.
 \label{eq:v49-adjoint}
\end{equation}
\end{proposition}
\begin{proof}
Move the derivatives of \(u\) in the integral onto their metric
coefficients. The first three terms become \(\phi^2L(g)\).
The last term is \(K^{cd}(g)\nabla_c\phi\nabla_d\phi\), with
\(K^{cd}=e^{ab}\sum_i r_i(M_i)_{ab}{}^{cd}\).
Its scalar Euler expression has principal coefficient \(-2K^{cd}\).
Arbitrary scalar second jets imply \(K=0\); the remaining zeroth
coefficient gives \(L=0\). Apply \Cref{lem:v49-algebraic}.
This argument justifies the arbitrary-scalar-test identity
\eqref{eq:v49-adjoint} without assuming that every smooth test is the
square of one real scalar.
\end{proof}

\section{A rank-nine principal test with six exact covariant lifts}
\label{sec:v49-principal}

The fifteen one-curvature coefficients of \(\mathscr D_gu\) have
ordered basis
\begin{equation}
 \mathscr L=
 \bigl(M_1(\nabla^2u),\ldots,M_7(\nabla^2u),
 N_1(\nabla u),\ldots,N_5(\nabla u),
 uD_1,uD_2,uD_3\bigr).
 \label{eq:v49-Lbank}
\end{equation}
Denote their vector by
\(t=(a/2,c/2,d_1,d_2,d_3)\in\R^{15}\).
Only this symbol \(t\), not a spacetime coordinate, is used in the
coefficient maps below.

Linearize curvature about the flat metric. For a mode \(h e^{ipx}\),
suppress common Fourier factors and set
\begin{equation}
 R^{(1)}_{abcd}(p,h)=\frac12
 (p_cp_bh_{ad}+p_dp_ah_{bc}-p_cp_ah_{bd}-p_dp_bh_{ac}).
 \label{eq:v49-Rsymbol}
\end{equation}
Let \(\mathscr L_i(p,h;k)\) be the resulting symbol of the
\(i\)-th tensor with scalar-test momentum \(k\). For
\(k=-p-q\), the necessary quadratic metric identity is
\begin{equation}
 \sum_it_i\bigl[
 G^{(1)}(p,h):\mathscr L_i(q,j;k)
 +G^{(1)}(q,j):\mathscr L_i(p,h;k)\bigr]=0.
 \label{eq:v49-principal-test}
\end{equation}
Both metric legs and the differentiated scalar test are retained.

\begin{theorem}[Complete one-curvature kernel]
\label{thm:v49-principal}
The universal test \eqref{eq:v49-principal-test} has rank nine in the
fifteen-dimensional bank. Its six-dimensional kernel consists exactly
of the covariant expressions
\begin{align}
 K_1&=\mathcal L_{R\nabla u}g,&
 K_2&=\mathcal L_{\operatorname{Ric}^{\sharp}\nabla u}g,&
 K_3&=\mathcal L_{u\nabla R}g,\nonumber\\
 K_4&=\Omega_{\nabla^2u}G,&
 K_5&=[u,\Delta]G,&
 K_6&=g[u,\Delta]\operatorname{tr}G.
 \label{eq:v49-six-kernels}
\end{align}
In \(K_5\), \(\Delta\) is the metric rough Laplacian on symmetric
tensors; in \(K_6\) it is the scalar Laplacian. All six expressions
satisfy \(\int\rho G:K_i=0\) on every metric, not only at its
linearized symbol.
\end{theorem}
\begin{proof}
The first three identities are the contracted Bianchi identity. The
fourth is pointwise skewness. The other two use
\begin{equation}
 [u,\Delta]T=-(\Delta u)T-2\nabla^cu\nabla_cT,
 \label{eq:v49-commutator}
\end{equation}
which is formally skew-adjoint. In particular
\begin{align}
 G:[u,\Delta]G
 &=-\nabla_c\bigl[(\nabla^cu)|G|^2\bigr],\nonumber\\
 G:g[u,\Delta]\operatorname{tr}G
 &=-\nabla_c\bigl[(\nabla^cu)(\operatorname{tr}G)^2\bigr].
 \label{eq:v49-divergences}
\end{align}
Their coordinate columns in \eqref{eq:v49-Lbank} are specified by
\begin{align}
 K_1&=2\mathscr L_1+\mathscr L_8,&
 K_2&=\mathscr L_5+\mathscr L_{11},&
 K_3&=\mathscr L_8+2\mathscr L_{13},\nonumber\\
 K_4&=\mathscr L_1-\tfrac12\mathscr L_2+\mathscr L_4,&
 K_5&=\tfrac12\mathscr L_2-\mathscr L_3
       +\mathscr L_9-2\mathscr L_{10},&
 K_6&=\mathscr L_2+2\mathscr L_9.
 \label{eq:v49-columns}
\end{align}
These six columns are independent, so the rank is at most nine.

For the reverse bound, let \(E^{ab}\) have ones in entries \(ab,ba\)
(and one in the diagonal case). The following nine real integer panels
are inserted into \eqref{eq:v49-principal-test}:
\begin{center}\small
\begin{tabular}{ccccc}
\toprule
Row & \(p\) & \(q\) & \(h\) & \(j\)\\\midrule
1&\((0,1,0,0)\)&\((1,0,0,1)\)&\(E^{02}\)&\(E^{03}\)\\
2&\((0,0,1,0)\)&\((1,0,0,1)\)&\(E^{00}\)&\(E^{33}\)\\
3&\((1,1,1,0)\)&\((1,0,1,0)\)&\(E^{12}\)&\(E^{22}\)\\
4&\((1,1,1,0)\)&\((1,1,0,0)\)&\(E^{00}\)&\(E^{11}\)\\
5&\((1,2,0,0)\)&\((1,1,0,0)\)&\(E^{22}\)&\(E^{22}\)\\
6&\((1,0,0,1)\)&\((1,1,1,0)\)&\(E^{12}\)&\(E^{22}\)\\
7&\((1,1,0,0)\)&\((1,2,0,0)\)&\(E^{11}\)&\(E^{01}\)\\
8&\((1,1,1,0)\)&\((1,2,0,0)\)&\(E^{33}\)&\(E^{01}\)\\
9&\((1,2,0,0)\)&\((1,0,1,0)\)&\(E^{02}\)&\(E^{33}\)\\\bottomrule
\end{tabular}
\end{center}
Columns \((1,2,3,5,6,7,12,14,15)\) give the square matrix
{\small
\begin{equation}
\mathsf M_*=
\begin{pmatrix}
0&0&0&0&0&1&-1/2&0&0\\
-1&-6&-3/2&-1&-1/12&0&0&-3&-3/4\\
1&36&9&1&4/3&0&0&10&5/2\\
-3/2&-36&-9/2&-5/2&3/2&0&0&-10&-5/4\\
-7/2&-260&-13/2&-13/2&41/2&0&0&-70&-7/4\\
4&0&0&6&-5/3&13/4&-1&0&0\\
0&104&0&0&-26/3&0&0&28&0\\
18&336&112&30&11&0&0&96&32\\
8&0&36&16&-4/3&0&0&0&14
\end{pmatrix},
\qquad \det\mathsf M_*=2400.
\label{eq:v49-minor}
\end{equation}
}
The displayed symbol, tensor contraction and determinant provide a
rational lower-rank certificate. Common powers of \(i\) and \(c_E\)
do not affect the conclusion. There are no sampled Gaussian histories,
complex stationary solutions, or floating-point rank thresholds in this
test. Together with the six exact covariant lifts this proves the
complete kernel assertion.
\end{proof}

The necessity test uses only the first-curvature part. Its sufficiency
for the subtraction is stronger: \eqref{eq:v49-columns} holds for the
\emph{covariant tensors themselves}. No Ricci commutator is suppressed
when these six representatives are removed. A linearized test alone,
without these exact representatives, would not justify the next step.

\section{The complete invariant leading-metric characteristic has eight parameters}
\label{sec:v49-classification}

Define the formally skew operators
\begin{equation}
 \mathscr C_u(T)=[u,\Delta]T,
 \qquad
 \mathscr C_u^0(T)=g[u,\Delta]\operatorname{tr}T.
 \label{eq:v49-Cops}
\end{equation}
The trace map and multiplication by \(g\) are adjoints and commute
with the metric connection, so the second operator is skew as well.

\begin{theorem}[Complete curvature--spectator Noether decomposition]
\label{thm:v49-main}
For a single scalar, \eqref{eq:v49-input} holds if and only if
\begin{equation}
 \boxed{F=\mathcal L_\xi g+\mathscr A_\phi e,}
 \label{eq:v49-decomp}
\end{equation}
where the eight real coefficients \((\alpha,\beta,\gamma,\delta,
\varepsilon,\zeta,\sigma,\kappa)\) are constant and
\begin{align}
 \xi&=\alpha R\nabla u
       +\beta\operatorname{Ric}^{\sharp}\nabla u
       +\gamma u\nabla R,\nonumber\\
 \mathscr A_\phi
 &=\frac1{c_E}
 \left(\delta\Omega_{\nabla^2u}
 +\varepsilon\mathscr C_u+\zeta\mathscr C_u^0
 +\sigma\Omega_V\right)
 +\frac{\kappa u}{c_E^2}\Omega_e,
 \quad u=\phi^2,\quad V=\nabla\phi\otimes\nabla\phi .
 \label{eq:v49-explicit}
\end{align}
The parameterization is injective as a universal natural
characteristic. The thirty-dimensional determining map has rank 22.
For four independently sign-invariant real scalars its domain, kernel
and rank are respectively \(120,32,88\).
\end{theorem}
\begin{proof}
\Cref{prop:v49-spectator} extracts \(\sigma\Omega_VG\).
The remaining one-curvature coefficient must lie in the six-dimensional
kernel of \Cref{thm:v49-principal}. Subtract the corresponding exact
covariant expressions, obtaining only
\(u\sum_{l=4}^{11}d_lD_l\). The inherited terminal theorem
\Cref{thm:v48-terminal} forces this to be
\[
 \kappa u(RS+g|S|^2)
 =\frac{\kappa u}{c_E^2}\Omega_e e.
\]
All subtractions are literal covariant characteristics satisfying the
integrated identity, so the remaining terminal identity is legitimate.
In particular, it is not an isolated undifferentiated portion of a
larger unprocessed integrated expression.

Conversely each gauge term is killed by the Bianchi identity and each
remaining term is the value of a formally skew operator on \(e\).
This proves sufficiency for arbitrary metric and scalar jets.

The change of variables from \((a,b,c,d)\) to
\((r=b-a,t=(a/2,c/2,d_1,d_2,d_3),d_4,\ldots,d_{11})\)
is invertible. In these three blocks the nullities are respectively
one, six, and one. This proves injectivity and the rank count. Independent
flavour signs split the actual four-scalar invariant sector into four
copies. Symmetric off-diagonal flavour polarizations give the same
formula for each symmetric flavour pair, but antisymmetric flavour
channels are outside this count.
\end{proof}

The new count is not an anomaly or full-current cohomology count.
All eight directions admit the selected exact current completions
below. The mass and flat-symbol transformations that V48 had already
removed are not recounted among these eight curvature-dependent
directions.

\section{A bounded coefficient compiler that preserves failed identities}
\label{sec:v49-compiler}

For \(t\in\R^{15}\), an explicit left inverse of the six columns in
\eqref{eq:v49-columns} is
\begin{equation}
\begin{split}
 \alpha=t_8-\tfrac12t_{13},\qquad
 \beta=t_{11},\qquad \gamma=\tfrac12t_{13},\qquad
 \delta=t_4,\\
 \varepsilon=-\tfrac12t_{10},\qquad
 \zeta=\tfrac12t_9+\tfrac14t_{10}.
\end{split}
\label{eq:v49-six-decoder}
\end{equation}
For a general raw array \(F=(a,b,c,d)\), use these formulas and set
\begin{equation}
 \sigma=b_1-a_1,
 \qquad \kappa=\tfrac12(d_4+d_7).
 \label{eq:v49-other-decoder}
\end{equation}
Let \(H_{49}\) be this eight-coefficient decoder and \(N_{49}\)
the synthesis map in \eqref{eq:v49-decomp}--\eqref{eq:v49-explicit},
expanded back into \eqref{eq:v49-raw}. Put \(\Pi_{49}=N_{49}H_{49}\).

\begin{proposition}[Exact coefficient bounds and residual ownership]
\label{prop:v49-compiler}
In the raw thirty-coefficient \(\ell^1\) norm and the displayed
parameter norm,
\begin{equation}
 \boxed{
 H_{49}N_{49}=I_8,\qquad
 \Pi_{49}^2=\Pi_{49},\qquad
 \|H_{49}\|_{1\to1}=1,
 }
 \label{eq:v49-Hnorm}
\end{equation}
\begin{equation}
 \boxed{
 \|N_{49}\|_{1\to1}=12,\qquad
 \|\Pi_{49}\|_{1\to1}=5,\qquad
 \|I-\Pi_{49}\|_{1\to1}=4.
 }
 \label{eq:v49-Pnorms}
\end{equation}
The same bounds hold on the direct sum of the four invariant channels.
Every synthesized input has zero Noether variation. For any input, the
compiler retains the residual \(F-\Pi_{49}F\), and
\begin{equation}
 \int\rho e:(F-\Pi_{49}F)=\int\rho e:F.
 \label{eq:v49-defect-preserved}
\end{equation}
Within the declared complete inventory, the original identity is
certified exactly when this residual is zero.
\end{proposition}
\begin{proof}
The identities follow by substituting \eqref{eq:v49-six-decoder} and
\eqref{eq:v49-other-decoder} into \eqref{eq:v49-columns}, followed by
\(\nabla^2u=2\phi\nabla^2\phi+2V\) and
\(\nabla u=2\phi\nabla\phi\). These give rational matrices of
sizes \(8\times30\) and \(30\times8\).
The largest absolute column sums are exactly the displayed norms; all
matrices, including the two residual maps, are reconstructed by the
checker. The first-curvature decoder alone has norm one, its synthesis
projection has norm three, and its complement norm two. Defect
preservation follows from \Cref{thm:v49-main}; necessity of zero
residual follows from the injective full parameterization.
\end{proof}

A geometric example records why derivative partners matter. The
characteristic \(F=[u,\Delta]G\) contains both
\(-(\Delta u)G\) and \(-2\nabla u\cdot\nabla G\). Its Noether
contraction is the divergence in \eqref{eq:v49-divergences}.
Keeping only the second term is not a skew operator and does not supply
this identity. The compiler tests the complete raw array; it does not
certify a term after silently adding its missing partners.

These bounds are finite coefficient-map bounds. They imply no smallness
of a cutoff-dependent input array and no uniformity in growing EFT or
source order. The factors \(c_E^{-1},c_E^{-2}\) appear explicitly in
the BV representative; pointwise and differential-operator bounds
require a fixed jet chart and \(|c_E|\ge c_*>0\).

\section{The full coupled current completion and the next matter degree}
\label{sec:v49-BV}

Let \(f=P_g\phi\) and \(\tau\) be the same signed scalar Euler
and metric-stress tensors as in V47--V48. For four scalar channels, sum their vectors and skew operators and use
the total scalar stress; the gauge vector acts on every scalar component.
The current completion is
\begin{equation}
 \boxed{
 \delta g=\mathcal L_\xi g+\mathscr A_\phi(e+\tau),
 \qquad
 \delta\phi=\mathcal L_\xi\phi.
 }
 \label{eq:v49-full-characteristic}
\end{equation}
Although the input leading metric term is massless, the \(\tau\)
partner retains the actual mass dependence of the coupled action.

\begin{theorem}[An exact coupled BV subtraction]
\label{thm:v49-BV}
With \(\vartheta=\rho^{-1}\eta_g\), the local functional
\begin{equation}
 \boxed{
 \mathcal T_F=\int\sigma_a\xi^a
 -\frac12\int\rho\,\vartheta:\mathscr A_\phi\vartheta
 }
 \label{eq:v49-primitive}
\end{equation}
has ghost number minus two and weight six, and satisfies
\begin{equation}
 \boxed{
 s\mathcal T_F
 =\int\eta_g:\delta g+\int\upsilon^T\delta\phi
 \quad\pmod{d_x}.
 }
 \label{eq:v49-exact-charge}
\end{equation}
It is natural, translation invariant and polynomial in the mass. Its
scalar-degree-two metric characteristic is exactly \(F\).
\end{theorem}
\begin{proof}
The first term is the inherited field-dependent gauge generator, whose
complete scalar and ghost-source variation was proved in
\Cref{thm:v48-reduction}. Each term of \(\mathscr A_\phi\) is
formally skew for the volume pairing. For the trace commutator this
uses the adjoint pair \(\operatorname{tr}\) and multiplication by
\(g\); for the rough commutator it uses metric compatibility and
integration by parts. Its coefficients are natural tensors, including
the occurrences of \(e(g)\), so the longitudinal variation of the
complete density is a divergence.

For the Euler part, write \(E_{\mathrm{tot}}=e+\tau\). The odd
product rule and formal adjoint give
\[
 -\frac12\int\rho\bigl[
 E_{\mathrm{tot}}:\mathscr A_\phi\vartheta
 -\vartheta:\mathscr A_\phi E_{\mathrm{tot}}\bigr]
 =\int\rho\,\vartheta:\mathscr A_\phi E_{\mathrm{tot}}.
\]
This is the second metric characteristic in
\eqref{eq:v49-full-characteristic}. There is no scalar antifield in
this term and no scalar Euler contribution has been dropped. The direct
variational check is
\[
 \int\rho\bigl[(e+\tau):\delta g+f^T\delta\phi\bigr]=0:
\]
the gauge terms cancel by the complete matter--gravity Ward identity,
and the remaining pairing with a skew operator is zero. The operator
coefficients have scalar degree two. Thus its \(e\) contribution is
leading metric degree two, its \(\tau\) contribution has scalar
degree four, and the gauge scalar term has scalar degree three.
\end{proof}

One can see the required nonlinear return already in the new
rough-commutator direction. Set \(\xi=0\) and
\(\mathscr A_\phi=c_E^{-1}[u,\Delta]\). Then
\begin{equation}
 \delta g=\frac1{c_E}
 \left[-(\Delta u)(e+\tau)
       -2\nabla^cu\nabla_c(e+\tau)\right],\qquad
 \delta\phi=0.
 \label{eq:v49-explicit-return}
\end{equation}
The mixed matter contribution is not an independently discardable
normalization. For arbitrary symmetric test fields \(T_1,T_2\),
\begin{equation}
 T_1:[u,\Delta]T_2+T_2:[u,\Delta]T_1
 =-2\nabla_c[(\nabla^cu)(T_1:T_2)].
 \label{eq:v49-cross-return}
\end{equation}
Taking \(T_1=e,T_2=\tau\) proves the cancellation of the complete
mixed action row. It need not be a pointwise zero.

Consequently, for a sign-invariant natural current in the V47 category,
first subtract V47's scalar-linear selected completion and V48's mass
and flat gauge charges, and then subtract \(s\mathcal T_F\).
Its leading metric characteristic is now zero. The next unresolved
matter equation has the form
\begin{equation}
 \int\rho\left(e:\widehat F_4+f^T\widehat Q_3\right)=0,
 \label{eq:v49-next}
\end{equation}
with the induced degree-four metric and degree-three scalar changes
already included in \(\widehat F_4,\widehat Q_3\). The previous
\(\tau:\widehat F_2\) term is absent because the complete
\(\widehat F_2\) has been removed, not because the scalar stress
has been put on shell. Equation \eqref{eq:v49-next} is not solved here.

\section{Verification, uniformity, and closing ledgers}
\label{sec:v49-ledgers}

The checker \texttt{check\_ESM\_curvature\_Noether\_V49.py} rebuilds
the curvature and dual-Weyl symbols, every entry of the nine by fifteen
principal matrix, its displayed determinant, the six-row
pointwise syzygy test, the eight-generator thirty-coefficient compiler,
and all stated induced norms. It checks exact skew identities and the
cross-return divergence by polynomial differential tests. It also
checks the decoder on every raw coefficient basis vector and verifies
preservation of the previous mathematical body. The certificate records
the raw data and hashes. No numerical rank threshold, new native loop
integral, full quantum BV matrix or proof-assistant formalization is
used. The natural-inventory reduction and nonlinear covariant lifts are
proved above; a finite assertion count is not substituted for either.

On a fixed metric-jet chart, the skew operator has the explicit
second-order form in \eqref{eq:v49-explicit} and
\eqref{eq:v49-commutator}. For example,
\[
 |[u,\Delta]T|
 \le |\Delta u|\,|T|+2|\nabla u|\,|\nabla T|,
 \qquad |\Omega_ST|\le4|S|\,|T|
\]
in an orthonormal four-frame. With \(|c_E|\ge c_*>0\), these give
fixed-jet operator bounds for the selected primitive. At fixed arity
and a fixed smooth source window, a local polynomial of derivative
degree \(d\) has the inherited weighted rooted-kernel bound
\(C_{N,b}\mu^d\) times its coefficient norm. The constant has no
volume multiplier and is independent of mesh on the common-window
sampling domain. No bound uniform in unbounded source order, arbitrary
curvature jets, shell count or perturbative order is inferred.

\subsection*{Proved}
\addcontentsline{toc}{subsection}{V49: Proved}
The complete invariant massless curvature-dependent leading metric
Noether identity has the eight-parameter decomposition per scalar in
\Cref{thm:v49-main}. The principal map has rank nine, the full
thirty-coefficient map rank 22, and the invariant four-scalar map rank
88 with 32 allowed characteristics. The coefficient decoder and its
residual have the exact norms stated in
\Cref{prop:v49-compiler}. Every allowed characteristic has the full
local ghost-number-minus-two primitive
\eqref{eq:v49-primitive}, including its higher-matter partners.

\subsection*{Inherited}
\addcontentsline{toc}{subsection}{V49: Inherited}
V1--V48 are unchanged. The carrier, signed normalization, scalar action,
source convention and mass/flat reductions are inherited. The
scalar-undifferentiated terminal kernel is
\Cref{thm:v48-terminal}; it is consumed after a justified subtraction,
not repeated as a new classification. Active V34, V38, and gravity
normalizations, and the finite measured reference, remain unchanged.

\subsection*{Literature input}
\addcontentsline{toc}{subsection}{V49: Literature input}
The distinction between gauge transformations, formally skew-Euler
characteristics and the full BV differential is standard
\cite{V20BBH1994}. No general classification of all Einstein gauge
identities, no blanket anomaly-vanishing theorem, and no passage from
coordinate-dependent to translation-invariant primitives is used in the
new kernel argument. Its finite lower-rank certificates and exact
covariant lifts are supplied explicitly.

\subsection*{Conditional}
\addcontentsline{toc}{subsection}{V49: Conditional}
For an invariant full current that admits the declared natural,
scalar-quadratic leading representative, the selected exact
subtractions remove the complete leading metric term without changing
its current class. Naturality of arbitrary tensor-valued reference
components has not been proved. The symmetric-flavour extension follows
by polarization; antisymmetric flavour channels are not classified.

\subsection*{Open}
\addcontentsline{toc}{subsection}{V49: Open}
The leading invariant natural metric identity left by V48 is no longer
an open subproblem. The next equation is \eqref{eq:v49-next}, together
with its higher-matter descendants. Scalar-independent metric
transformations and the reduction of the actual non-natural descent
remain distinct inputs. A ghost-number-minus-two current primitive is
not the ghost-number-zero critical quantum counterterm. The complete
translation-invariant, flavour-invariant, mass-polynomial critical
scalar--gravity primitive remains open, as do the internal-gauge/chiral,
higher-loop and interacting ESM invariant-class interfaces.

\begin{center}\small
\begin{tabular}{>{\raggedright\arraybackslash}p{.24\textwidth}>{\raggedright\arraybackslash}p{.68\textwidth}}
\toprule
Variable & Scope in V49\\\midrule
Theory & Same ordinary Einstein action and four equal real scalars;
internal couplings and the cosmological source remain zero.\\
Current category & Natural, leading metric scalar degree two and
weight four; complete sign-invariant curvature residual, not all
current cohomology.\\
Mass & The input residual is massless. Its full current completion
retains the polynomial mass dependence of \(\tau\) and \(f\).\\
Cutoff and volume & Exact finite coefficient-map bounds independent of
mesh and periods. No new quantum cutoff or many-shell estimate.\\
Jet domain & Fixed chart and \(|c_E|\ge c_*>0\) for differential
and cotangent-coordinate bounds.\\
Source/loop order & Fixed-order kernel estimates; no uniformity in
unbounded source degree or loop number.\\
Verification & Exact rational tensor tests and finite coefficient
maps; neither new loop quadrature nor independent verification of all
inherited results.\\\bottomrule
\end{tabular}
\end{center}

\section*{Append-only continuation marker}
\addcontentsline{toc}{section}{Append-only continuation marker}
\label{sec:v1-continuation-marker}
\textbf{End of V49.} Preserve Sections 1--383 and all active prescriptions.
The invariant natural curvature-dependent leading metric Noether identity
has a complete local gauge/skew-Euler section. Next: the residual
scalar-cubic/metric-quartic current equation with all induced partners,
or a direct construction of the native invariant critical primitive.
Do not identify current exactness with a ghost-number-zero quantum
counterterm or with a new measured cutoff theorem.


\clearpage
\part*{V50: A complete natural scalar-cubic section of the higher-matter current}
\addcontentsline{toc}{part}{V50: A complete natural scalar-cubic section of the higher-matter current}

\section{The scalar-cubic part of the next coupled identity}
\label{sec:v50-start}

Sections 1--383, including their historical closing marker, are preserved.
We retain the ordinary coefficient chart and signed native action
\eqref{eq:v47-data}, the four equal real scalars, and all active
normalizations.  Internal gauge, Yukawa and scalar self-interactions and
the cosmological source are zero.  No finite action, propagator, filter,
contour or measure is changed in this installment.

The input is precisely the next equation retained in
\eqref{eq:v49-next}:
\begin{equation}
 \int\rho\left(e^{ab}(F_4)_{ab}+f^TQ_3\right)=0,
 \qquad e^{ab}=c_EG^{ab},\quad f=P_g\phi,\quad
 P_g=-\Delta_g+z,\quad z=m^2.
 \label{eq:v50-input}
\end{equation}
It holds for arbitrary metric and scalar tests, as a polynomial identity
in \(z\).  The subscripts are scalar degrees: \(F_4\) is quartic and
\(Q_3\) cubic. Both characteristics have differential weight four,
with \(\wt(\nabla)=1\), \(\wt(z)=\wt(\mathrm{curvature})=2\).
They are covariantly natural scalar-jet polynomials in the same category
as V46--V49, with constant flavour tensors, no preferred spacetime
tensor, no explicit coordinate coefficient, and no inverse mass or
curvature. We impose the actual independent flavour-sign invariance.
This is a restriction on the current being analyzed; naturality of the
unreduced native critical reference has not thereby been established.

The result is a section of the \emph{scalar-cubic projection} of
\eqref{eq:v50-input}.  We first construct a scalar-quadratic,
formally skew operator \(\mathcal L_\phi\), then a scalar-cubic
operator \(\mathcal K_\phi\) acting on symmetric tensors, such that
\begin{equation}
 \boxed{Q_3=\mathcal L_\phi f+\mathcal K_\phi e,
 \qquad \mathcal L_\phi^\dagger=-\mathcal L_\phi.}
 \label{eq:v50-target}
\end{equation}
The equality is a curved off-shell identity, not an equality only on
flat solutions.  Both operators are local and use at most two
derivatives beyond their undifferentiated arguments, counting
coefficient derivatives in that total. Their full current completion
will retain \(\mathcal K_\phi\tau\), where \(\tau\) is the inherited
scalar stress \eqref{eq:v47-stress}.

This advances beyond a scalar-only on-shell test without asserting that
all metric characteristics have been classified. After the exact
subtraction the remaining leading identity is an unrestricted natural
quartic-scalar spectator identity for the metric. Its solution is not
assumed here.

\section{A polynomial flat Euler section with no mass-shell division}
\label{sec:v50-Koszul}

At \(g=I\), the gravitational Euler tensor vanishes. Polarize
\(\int f^TQ_3=0\) in four labelled scalar legs with momenta
\(p_1+\cdots+p_4=0\). Let \(q_i\) be the root-normalized cubic
coefficient with the Euler factor on leg \(i\). Scalar permutation
symmetry includes the flavour labels. The exact identity is
\begin{equation}
 W(q):=\sum_{i=1}^4x_iq_i=0,
 \qquad x_i=p_i^2+z.
 \label{eq:v50-flat-ward}
\end{equation}
The overall common polarization factorial has been removed. Each
\(q_i\) has degree two in the scalar products and \(z\).

There is no flat pseudoscalar: four scalar momenta with conserved total
span at most three dimensions. The six Gram entries of three independent
vectors in four dimensions are algebraically independent. This follows,
for example, from the open cone of their positive definite Gram
matrices. Adjoining \(z\) therefore gives a polynomial ring in seven
variables. Introduce the balanced channels
\begin{equation}
 \begin{split}
 S_x&=\sum_i x_i-4z,\\
 a&=(p_1+p_2)^2-S_x/3,\qquad
 b=(p_1+p_3)^2-S_x/3,\\
 (p_1+p_4)^2-S_x/3&=-a-b.
 \end{split}
 \label{eq:v50-balanced}
\end{equation}
Then \((x_1,x_2,x_3,x_4,a,b,z)\) are independent linear coordinates
on that invariant ring. For example,
\[
 p_1\!\cdot p_2=\frac a2-\frac{x_1+x_2}{3}
                 +\frac{x_3+x_4}{6}+\frac z3,
 \qquad p_i^2=x_i-z.
\]
Testing the original identity for real momenta and \(z>0\) determines
this polynomial identity; no complex mass-shell substitution is needed.

\begin{theorem}[Explicit first Euler syzygy and its residual]
\label{thm:v50-koszul}
Decompose a degree-two coefficient vector into its \(x\)-homogeneous
parts \(q=\sum_{m=0}^2q^{(m)}\). Define the skew matrix
\begin{equation}
 \boxed{
 A_{ij}(q)=\sum_{m=0}^2
 \frac{\partial_{x_j}q_i^{(m)}-
       \partial_{x_i}q_j^{(m)}}{m+1}.
 }
 \label{eq:v50-koszul-section}
\end{equation}
For every input, including a nonclosed one,
\begin{equation}
 \boxed{
 q_i-\sum_jA_{ij}(q)x_j
 =\sum_{m=0}^2\frac{\partial_{x_i}W(q^{(m)})}{m+1}.
 }
 \label{eq:v50-koszul-residual}
\end{equation}
In particular \(W(q)=0\) implies \(q=Ax\). The constructed entries
of \(A\) have total invariant degree one, hence differential weight
two. The residual has the same Euler defect as \(q\).
In the coefficient \(\ell^1\) norm, counting only \(i<j\) for the
skew output,
\begin{equation}
 \|A(q)\|_1\le\tfrac23\|q\|_1,\qquad
 \|q-A(q)x\|_1\le\|W(q)\|_1.
 \label{eq:v50-koszul-bounds}
\end{equation}
\end{theorem}
\begin{proof}
Euler's homogeneous identity gives
\[
 \sum_jx_j\partial_{x_j}q_i^{(m)}=m q_i^{(m)},\qquad
 \sum_jx_j\partial_{x_i}q_j^{(m)}
 =\partial_{x_i}W(q^{(m)})-q_i^{(m)}.
\]
Subtract and divide by the positive integer \(m+1\). This proves
\eqref{eq:v50-koszul-residual}. Multiplication by \(x_i\) and summation
annihilates the skew part, proving defect ownership. Different
\(x\)-degrees of \(W\) cannot cancel one another. A monomial of
\(x\)-degree \(m\) contributes at most \(m/(m+1)\) to the upper
skew coefficient sum. On the right of
\eqref{eq:v50-koszul-residual}, the summed derivative norm of an
\(x\)-homogeneous polynomial of degree \(m+1\) is at most
\((m+1)\) times its coefficient norm. The disjoint homogeneous supports
give the two stated bounds.
\end{proof}

The unlabelled algebra here is the elementary degree-two part of the
Koszul resolution for four independent variables. Its vector domain has
dimension \(4\binom82=112\); the image in cubic polynomials consists
of those containing at least one \(x_i\), of dimension
\(\binom93-\binom53=74\). Its kernel consequently has dimension 38.
The skew synthesis domain has dimension \(6\cdot7=42\), with a
four-dimensional relation space. These counts are not the physical
flavour-symmetrized counts below.

If a permutation-equivariant section is desired, average
\eqref{eq:v50-koszul-section} over all 24 leg permutations, including
flavours. A permutation sends \(a,b\) to two members of
\(a,b,-a-b\) and permutes the \(x_i\). Its degree-two norm is at most
four and its degree-one norm at most two. The averaged section therefore
has norm at most \(16/3\). It retains the scalar exchange symmetries
and the skew exchange of the two Euler slots. This averaging is only a
coefficient operation; no finite regulator is rotationally averaged.

The abstract Koszul statement is standard \cite{V50Stacks}; the displayed section and
residual are proved directly and do not use a cohomological existence
claim in place of a coefficient map. See the comparison reference at
the end of this installment.

\section{The complete invariant flat cubic compiler}
\label{sec:v50-flat-compiler}

A natural scalar-cubic characteristic with independent flavour signs
splits into four self channels \(Q_A(\phi_A^3)\) and six unordered
pair channels
\((Q_A(\phi_A\phi_B^2),Q_B(\phi_B\phi_A^2))\).
For three input momenta \(r,s,t\), use the seven variables
\begin{equation}
 (r^2,s^2,t^2,r\cdot s,r\cdot t,s\cdot t,z).
 \label{eq:v50-root-coordinates}
\end{equation}
An explicit coefficient basis is the sum over each distinct permutation
orbit of degree-two monomials. Use all input permutations for a self
channel and only \(s\leftrightarrow t\) for a pair component. These
bases have nine and eighteen elements, respectively. The independent
sign-invariant domain thus has dimension
\begin{equation}
 4\cdot9+6\cdot(18+18)=252.
 \label{eq:v50-domain252}
\end{equation}

\begin{theorem}[Flat scalar-cubic kernel]
\label{thm:v50-flat-classification}
The determining equation \(\int(P_0\phi)^TQ_3=0\), with
\(P_0=-\Delta+z\), has ranks 8 and 27 on the self and unordered-pair
blocks. Their kernel dimensions are 1 and 9. For the four-scalar
sign-invariant domain the rank and kernel dimension are
\begin{equation}
 \boxed{\operatorname{rank}=194,\qquad \dim\operatorname{Ker}=58.}
 \label{eq:v50-flat-counts}
\end{equation}
Every kernel coefficient has the unique representative
\(Q_3=\mathcal L_\phi P_0\phi\) in the following 58-dimensional
skew-operator bank.

For each ordered flavour pair \((A,B)\), include the diagonal
first-order operator
\begin{equation}
 (\mathcal L_\phi)_{AA}
 =\sum_B\lambda_{AB}
 \left(\phi_B\nabla^a\phi_B\nabla_a
       +\tfrac12\nabla_a(\phi_B\nabla^a\phi_B)\right).
 \label{eq:v50-diagonal-L}
\end{equation}
For each \(A<B\), put \(u=\phi_A,v=\phi_B\) and allow
\begin{equation}
 \begin{split}
 (\mathcal L_\phi)_{AB}={}&
 c\,uv\Delta+d\,u\nabla v\cdot\nabla
 +e\,v\nabla u\cdot\nabla\\
 &+f\,u\Delta v+g\,v\Delta u
 +h\,\nabla u\cdot\nabla v+jzuv,\\
 (\mathcal L_\phi)_{BA}&=-(\mathcal L_\phi)_{AB}^{\dagger}.
 \end{split}
 \label{eq:v50-offdiagonal-L}
\end{equation}
The formulas remain formally skew under covariantization with the
volume pairing. The count is \(16+6\cdot7=58\).
\end{theorem}
\begin{proof}
All displayed operators are formally skew. Hence their synthesized
characteristics obey the determining identity on any metric, not only
flat space. The flat scalar symbols of these operators are computed
without a tensor ansatz beyond \eqref{eq:v50-root-coordinates}.

For reproducibility put \(P_r=r^2+z\). The diagonal pair symbol is
\[
 -\frac14P_r\bigl[2r\cdot(s+t)+(s+t)^2\bigr].
\]
The off-diagonal symbol with the target-flavour leg \(r\) is
\begin{align*}
 \frac12\big[&-c(s^2P_s+t^2P_t)-d(s\cdot t)(P_s+P_t)
 -e((r\cdot s)P_s+(r\cdot t)P_t)\\
 &-f(s^2P_t+t^2P_s)-gr^2(P_s+P_t)
 -h((r\cdot s)P_t+(r\cdot t)P_s)+jz(P_s+P_t)\big].
\end{align*}
For the reverse channel replace
\[
 (c,d,e,f,g,h,j)\mapsto
 (-c,e-2c,d-2c,e-c-g,d-c-f,d+e-2c-h,-j).
\]
This is exactly the integration-by-parts adjoint of
\eqref{eq:v50-offdiagonal-L}, with the two flavour roles exchanged.

The common four-leg polynomial maps are
\[
 W_{\rm self}(q)=\sum_{i=1}^4x_iq(p_{j\ne i}),
\]
\begin{align*}
 W_{\rm pair}(q_u,q_v)={}&
 x_1q_u(p_2,p_3,p_4)+x_2q_u(p_1,p_3,p_4)\\
 &+x_3q_v(p_4,p_1,p_2)+x_4q_v(p_3,p_1,p_2).
\end{align*}
Expand them in the 84 degree-three monomials of
\((x_1,x_2,x_3,x_4,a,b,z)\). This defines rational matrices of sizes
\(84\times9\) and \(84\times36\) without sampling a momentum.
The synthesized null columns have full ranks one and nine; their
selected minors are \(-1/2\) and \(1/128\). The determining matrices
have respectively an eight-square minor \(3/16\) and a
27-square minor \(-2^{-21}\). The certificate gives their precise row
and column selections and all entries reconstructed from the preceding
formulas. These lower bounds, and the displayed independent null
columns, prove both ranks. Summing over the disjoint flavour patterns
gives \eqref{eq:v50-flat-counts}. The Gram-variable independence proves
that this is the entire declared natural flat inventory, not merely a
sufficient finite momentum test.
\end{proof}

For one scalar the sole direction has the particularly transparent form
\begin{equation}
 \boxed{
 Q_3=\phi\nabla\phi\cdot\nabla f
  +\tfrac12\bigl(|\nabla\phi|^2+\phi\Delta\phi\bigr)f,
 \qquad f=P_g\phi.
 }
 \label{eq:v50-one-scalar}
\end{equation}
Indeed \(fQ_3=\tfrac12\nabla_a(\phi\nabla^a\phi\,f^2)\).
This is a nonzero cubic characteristic with an exact charge; it is not
a surviving nonlinear internal symmetry.

\begin{proposition}[Evaluated decoder and failed-input ownership]
\label{prop:v50-flat-decoder}
Let \(N\) synthesize the block operators above. Choose the first
independent rows of \(N\) in the recursive monomial ordering recorded
in the certificate, and invert their square minor to define \(H\).
Then \(HN=I\), \(\Pi=NH\) is a projection, and the determining
map obeys \(W(q-\Pi q)=W(q)\). On the full four-flavour direct sum,
\begin{equation}
 \boxed{
 \|H\|_{1\to1}=12,\quad \|N\|_{1\to1}=12,\quad
 \|\Pi\|_{1\to1}=22,\quad \|I-\Pi\|_{1\to1}=21.
 }
 \label{eq:v50-flat-norms}
\end{equation}
Within this inventory the input is variational exactly when
\(q-\Pi q=0\). The public compiler returns the residual and its
original variational defect as well as the selected operator.
\end{proposition}
\begin{proof}
The self block uses row 1 (zero based), whose minor is \(-1/2\).
The pair block uses rows
\((0,1,3,7,8,9,11,13,19)\), whose minor is \(1/128\).
The resulting rational matrices have maximal absolute column sums as
in \eqref{eq:v50-flat-norms}; the self block has respectively
\(2,7/6,7/3,4/3\). Every entry is included in the certificate.
The identities follow from the minor construction and \(WN=0\).
Necessity of zero residual follows from the complete kernel theorem,
not from a claim that \(H\) solves arbitrary failed inputs.
\end{proof}

\section{The curved remainder has no independent Weyl-square projection}
\label{sec:v50-curvature}

Apply the preceding compiler to the flat restriction of \(Q_3\),
and covariantize \eqref{eq:v50-diagonal-L}--\eqref{eq:v50-offdiagonal-L}
with their displayed adjoints. Subtract the complete characteristic
\(\mathcal L_\phi P_g\phi\). Its full current is exact, as verified
below. The residual \(\overline Q_3\) vanishes at every flat metric,
for all \(z\).

\begin{lemma}[Curvature remainder inventory]
\label{lem:v50-curvature-inventory}
Every natural cubic-scalar remainder of weight four with zero flat
restriction is a linear combination of the following nine Ricci
structures (constant flavour coefficients and repeated flavour labels
are understood):
\begin{equation}
\begin{array}{lll}
 R\phi_B\phi_C P_g\phi_D,&
 R^{ab}\phi_B\phi_C\nabla_a\nabla_b\phi_D,&
 R\phi_B\nabla\phi_C\cdot\nabla\phi_D,\\
 R^{ab}\phi_B\nabla_a\phi_C\nabla_b\phi_D,&
 (\nabla^aR)\phi_B\phi_C\nabla_a\phi_D,&
 (\Delta R)\phi_B\phi_C\phi_D,\\
 R^2\phi_B\phi_C\phi_D,&
 R_{ab}R^{ab}\phi_B\phi_C\phi_D,&
 zR\phi_B\phi_C\phi_D,
\end{array}
\label{eq:v50-Ricci-bank}
\end{equation}
plus
\begin{equation}
 |C|^2 Z_A(\phi)+\langle C,{}^\star C\rangle\widetilde Z_A(\phi)
 \label{eq:v50-Weyl-bank}
\end{equation}
in output flavour \(A\), where \(Z,\widetilde Z\) are homogeneous
cubic flavour polynomials. The nine-term list is a spanning normal
form; no unnecessary claim of independent ordered flavour tensors is
made.
\end{lemma}
\begin{proof}
Order scalar covariant derivatives, recording their commutators. The
zero-curvature part is zero by the flat restriction. Each remaining
term has a curvature or a covariant curvature derivative. One
curvature leaves at most two scalar derivatives. A Weyl tensor cannot
be contracted to a scalar with only two derivative indices: any
remaining metric contraction is a trace, and the possible alternating
contraction vanishes by the algebraic Bianchi identity. Thus these
terms reduce to \(R\) and \(R_{ab}\) in the first four positions.
One curvature derivative leaves one scalar derivative, and contracted
Bianchi reduces the resulting vector to \(\nabla R\). Two curvature
derivatives leave no scalar derivatives; their scalar contraction is
\(\Delta R\), modulo curvature-square terms from commuting tensor
derivatives. The curvature-square scalar contractions are
\(R^2,|\operatorname{Ric}|^2,|C|^2,\langle C,{}^\star C\rangle\).
The only mass-bearing curvature term is \(zR\phi^3\). Replacing a
Laplacian by \(-P_g+z\) yields the displayed choice of first term.
All symmetrizations affect only constant flavour coefficients.
\end{proof}

\begin{theorem}[Cubic Weyl potentials are excluded by the off-shell identity]
\label{thm:v50-Weyl-zero}
For an input satisfying \eqref{eq:v50-input}, both cubic maps in
\eqref{eq:v50-Weyl-bank} vanish. Consequently
\(\overline Q_3\) belongs to \eqref{eq:v50-Ricci-bank}.
\end{theorem}
\begin{proof}
The subtraction of \(\mathcal L_\phi f\) preserves the integrated
identity because \(\mathcal L_\phi\) is formally skew. Restrict to
the two exact Ricci-flat metrics of \Cref{lem:v47-vacuum}. All Ricci
terms and \(e:F_4\) vanish. Their two curvature scalars obey
\(|C|^2=24x^{-6}\) and
\(\langle C,{}^\star C\rangle=\pm24x^{-6}\) in the inherited
normalization. For each sign there follows
\[
 \int\rho\,(P_g\phi)_A\,a(x)\,Z_{\pm,A}(\phi)=0,
 \qquad a(x)=24x^{-6}>0,\quad Z_\pm=Z\pm\widetilde Z.
\]
The coefficient of the highest arbitrary scalar Hessian in the Euler
variation of this integral is a nonzero multiple of
\(a(x)(\partial_AZ_{\pm,B}+\partial_BZ_{\pm,A})\).
Metric and \(a\)-derivative terms have lower scalar-jet degree and
cannot cancel it. Therefore each \(Z_\pm\) satisfies the Killing
equation on flat flavour space.

Differentiate that equation and use the three index permutations:
\(
2\partial_A\partial_BZ_C
=\partial_A(\partial_BZ_C+\partial_CZ_B)
+\partial_B(\partial_AZ_C+\partial_CZ_A)
-\partial_C(\partial_AZ_B+\partial_BZ_A)=0.
\)
Thus each map is affine. Homogeneity of scalar degree three forces it
to be zero. Both signs give \(Z=\widetilde Z=0\). This uses arbitrary
scalar tests on a Ricci-flat metric, not the scalar equation of motion.
The existence and curvature of the two metric witnesses are inherited
from V47; they are not reclassified in this version.
\end{proof}

For reference, the finite cubic flavour Killing map has 80 input
coefficients and 100 symmetric-gradient coefficients and is injective.
The checker reconstructs its nonzero 80-square minor. This verifies the
finite instance of the elementary argument, rather than substituting a
rank test for the curved inventory or the Ricci-flat restriction.

\section{An explicit Ricci section with its full adjoint}
\label{sec:v50-Ricci-section}

Use the inherited pointwise trace maps
\begin{equation}
 \mathfrak r(T)=-\frac{\operatorname{tr}_gT}{c_E},\qquad
 \mathfrak r^{ab}(T)=\frac1{c_E}
 \left(T^{ab}-\tfrac12g^{ab}\operatorname{tr}_gT\right).
 \label{eq:v50-trace-maps}
\end{equation}
They obey \(\mathfrak r(e)=R\) and
\(\mathfrak r^{ab}(e)=R^{ab}\). The new use is at scalar degree
three, including two differentiated scalar factors; this is not merely
V47's scalar-linear operator applied independently to each scalar.

Put the Ricci remainder into the explicit normal form
\begin{equation}
 \overline Q_{3,A}
 =R\,U_A+R^{ab}V_{A,ab}
       +(\nabla_aR)W_A^a+(\Delta R)X_A,
 \label{eq:v50-UVWX}
\end{equation}
where
\begin{align}
 U_A={}&a_{A;BCD}\phi_B\phi_C P_g\phi_D
       +c_{A;BCD}\phi_B\nabla\phi_C\cdot\nabla\phi_D
       +(g_{A;BCD}R+i_{A;BCD}z)\phi_B\phi_C\phi_D,
       \nonumber\\
 V_{A,ab}={}&b_{A;BCD}\phi_B\phi_C\nabla_a\nabla_b\phi_D
       +d_{A;BCD}\phi_B\nabla_{(a}\phi_C\nabla_{b)}\phi_D
       +h_{A;BCD}R_{ab}\phi_B\phi_C\phi_D,
       \nonumber\\
 W_A^a={}&e_{A;BCD}\phi_B\phi_C\nabla^a\phi_D,
 \qquad X_A=f_{A;BCD}\phi_B\phi_C\phi_D.
 \label{eq:v50-coefficients}
\end{align}
All these coefficients refer to a fixed ordering of the covariant
normal form, without taking an integration-by-parts quotient of the
characteristic itself.

\begin{proposition}[Cubic scalar Ricci factorization]
\label{prop:v50-K}
Define
\begin{equation}
 \boxed{
 (\mathcal K_\phi T)_A
 =\mathfrak r(T)U_A+\mathfrak r^{ab}(T)V_{A,ab}
 +\nabla_a\mathfrak r(T)W_A^a
 +(\Delta\mathfrak r(T))X_A.
 }
 \label{eq:v50-K}
\end{equation}
Then \(\mathcal K_\phi e=\overline Q_3\). Its full volume adjoint is
\begin{equation}
 \boxed{
 (\mathcal K_\phi^\dagger y)_{ab}
 =\frac1{c_E}\left[
 y_A V_{A,ab}-\tfrac12g_{ab}y_A\operatorname{tr}_gV_A
 -g_{ab}\left(y_AU_A-\nabla_c(y_AW_A^c)+\Delta(y_AX_A)\right)
 \right].
 }
 \label{eq:v50-Kadj}
\end{equation}
In particular the derivatives act on the entire products containing
\(y\); they cannot be held on the scalar coefficients alone.
\end{proposition}
\begin{proof}
The first identity is substitution of
\eqref{eq:v50-trace-maps}. Integrate the first-derivative term once and
the second-derivative term twice in the volume pairing. The adjoint of
\(\mathfrak r\) is \(-g/c_E\), and that of
\(\mathfrak r^{ab}\) is the displayed trace reversal. This gives all
signs and product derivatives in \eqref{eq:v50-Kadj}. No inverse Euler
operator has been used.
\end{proof}

The maps commute with the independent flavour signs. In an ordered
contracted coefficient norm on \((U,V,W,X)\), retaining derivatives of
products as whole slots, the trace map and adjoint have conservative
bounds \(8|c_E|^{-1}\). After expansion into independent tensor and
Leibniz slots a bound \(512|c_E|^{-1}\) suffices: the trace reversal
costs at most eight and each of at most two derivatives on four factors
costs at most sixteen, with the remaining symmetrization covered by the
factor four. This is a bound on this fixed normal-form coefficient map;
conversion from an arbitrary jet basis is an additional fixed rational
linear map. Neither estimate gives smallness of a cutoff-dependent
input.

\section{The complete ordinary BV subtraction and the sixth-matter return}
\label{sec:v50-BV}

Write \(\pi=\rho^{-1}\upsilon\) and
\(\vartheta=\rho^{-1}\eta_g\). The same ordinary differential gives
\(s\pi=f+\mathcal L_c\pi\) and the density/tensor counterpart for
\(\vartheta\). No finite-filter nilpotence is asserted.

\begin{theorem}[Scalar-cubic projection has an exact coupled completion]
\label{thm:v50-main}
Every sign-invariant natural input satisfying
\eqref{eq:v50-input} admits \eqref{eq:v50-target} with the explicitly
constructed \(\mathcal L_\phi\) and \(\mathcal K_\phi\).
The ghost-number-minus-two, weight-six local functional
\begin{equation}
 \boxed{
 \mathcal T_{Q_3}
 =-\tfrac12\int\rho\,\pi^T\mathcal L_\phi\pi
       -\int\rho\,\pi^T\mathcal K_\phi\vartheta
 }
 \label{eq:v50-primitive}
\end{equation}
has full ordinary variation
\begin{equation}
 \boxed{
 s\mathcal T_{Q_3}
 =\int\eta_g:\delta g+\int\upsilon^T\delta\phi
 \pmod{d_x},
 \qquad
 \begin{cases}
 \delta g=-\mathcal K_\phi^\dagger f,\\
 \delta\phi=\mathcal L_\phi f+\mathcal K_\phi(e+\tau)
            =Q_3+\mathcal K_\phi\tau.
 \end{cases}
 }
 \label{eq:v50-full-characteristic}
\end{equation}
The metric characteristic has scalar degree four. The scalar
characteristic has scalar degrees three and five, both of which are
retained. The complete selected current is therefore exact, not merely
its leading scalar projection.
\end{theorem}
\begin{proof}
At flat metric, \eqref{eq:v50-input} supplies the flat variational
identity, so \Cref{thm:v50-flat-classification} constructs
\(\mathcal L_\phi\). Its covariantization is skew for the same volume
pairing. The resulting curved remainder belongs to
\eqref{eq:v50-Ricci-bank} by
\Cref{thm:v50-Weyl-zero}. Proposition~\ref{prop:v50-K} supplies
\(\mathcal K_\phi\), proving \eqref{eq:v50-target}.

The longitudinal variation of each whole natural density in
\eqref{eq:v50-primitive} is a divergence. For the complementary Euler
variation the odd-antifield product rule gives
\[
 -\tfrac12\int\rho
 \bigl(f^T\mathcal L_\phi\pi-\pi^T\mathcal L_\phi f\bigr)
 =\int\upsilon^T\mathcal L_\phi f,
\]
using formal skew-adjointness. In the mixed term it gives
\[
 -\int\rho f^T\mathcal K_\phi\vartheta
 +\int\rho\pi^T\mathcal K_\phi(e+\tau),
\]
which is \eqref{eq:v50-full-characteristic}. The coefficient dependence
on \(\phi,g\) does not create an omitted Euler term: the Euler part
acts on antifields, while the longitudinal coefficient variation was
already retained in the natural density.

As a direct check, the complete action variation is
\[
 \int\rho\left[-(e+\tau):\mathcal K_\phi^\dagger f
 +f^T\mathcal L_\phi f+f^T\mathcal K_\phi(e+\tau)\right]=0.
\]
The first and last terms cancel by the actual adjoint, and the middle
term vanishes by skewness. This holds for arbitrary off-shell fields.
The cotangent conversion from \(\eta_g\) to the inherited coframe
antifield is pointwise and retains the same chart Jacobian convention.
\end{proof}

The stress substitutions needed in the last term are explicit:
\begin{equation}
 \mathfrak r(\tau)=-\frac{X/2+zU}{c_E},\qquad
 \mathfrak r^{ab}(\tau)
 =-\frac{V^{ab}/2+zUg^{ab}/4}{c_E},
 \quad U=\phi^T\phi,\quad V^{ab}=\nabla^a\phi^T\nabla^b\phi.
 \label{eq:v50-stress-substitution}
\end{equation}
Their derivatives in \eqref{eq:v50-K} are also taken. The resulting
\(\mathcal K_\phi\tau\) is a fifth-degree scalar polynomial, with
no singularity at \(z=0\).

For a single scalar choose the Ricci input \(Q_3=zR\phi^3\), so
\(\mathcal L_\phi=0\) and
\(\mathcal K_\phi T=z\phi^3\mathfrak r(T)\). Then
\begin{equation}
 \boxed{
 \delta g_{ab}=\frac z{c_E}g_{ab}\phi^3P_g\phi,
 \qquad
 \delta\phi=zR\phi^3-
       \frac z{c_E}\phi^3\left(\tfrac12|\nabla\phi|^2+z\phi^2\right).
 }
 \label{eq:v50-example}
\end{equation}
On flat metric with constant \(\phi=v\), the leading metric variation
contributes \(z^3v^6/c_E\) to the action density. Its scalar-quintic
partner contributes \(-z^3v^6/c_E\). The sixth-matter cancellation is
therefore nontrivial. Keeping only \(zR\phi^3\) in the scalar
characteristic would fail the full identity, despite the correctness
of its fourth-matter projection.

\section{What the subtraction removes, and what it leaves}
\label{sec:v50-next}

Suppose the residual current before this step has leading
characteristics \((F_4,Q_3)\) and subsequent terms \((F_6,Q_5)\).
Subtract the complete exact charge \(s\mathcal T_{Q_3}\), not merely
its scalar-linear or flat part. The new leading coefficients are
\begin{equation}
 \boxed{
 \widehat Q_3=0,\qquad
 \widehat F_4=F_4+\mathcal K_\phi^\dagger f,\qquad
 \int\rho\,e:\widehat F_4=0.
 }
 \label{eq:v50-remainder}
\end{equation}
The last equality uses \eqref{eq:v50-input} and
\(\int\rho f^T\mathcal L_\phi f=0\). It is an off-shell
spectator identity. At the next matter degree the unchanged logical
requirement is
\begin{equation}
 \boxed{
 \int\rho\left(e:\widehat F_6+
                 \tau:\widehat F_4+f^T\widehat Q_5\right)=0.
 }
 \label{eq:v50-sixth}
\end{equation}
In particular \(\widehat Q_5=Q_5-\mathcal K_\phi\tau\), with the
preceding selections already included. The stress term in
\eqref{eq:v50-sixth} must not be set to zero until the complete quartic
metric characteristic has been removed.

Equation \eqref{eq:v50-remainder} is not a copy of V49's solved
quadratic spectator equation: four scalar arguments allow coefficient
structures not determined by substituting \(u=\phi^2\) into the
old eight directions. No completeness statement for this quartic
metric inventory is made here. Conversely, the scalar-cubic projection
need not be recomputed when addressing that metric identity: it has
just been removed by a full local primitive.

The results are finite coefficient statements in the ordinary local
complex. For fixed metric jets, \(|c_E|\ge c_*>0\), fixed arity, and
fixed external window \(\mu\), the explicit order-two operators yield
polynomial jet bounds from \eqref{eq:v50-flat-norms} and
\eqref{eq:v50-Kadj}. After the standard smooth source window, a local
symbol of derivative degree \(d\) has a rooted kernel bound
\(C_{N,b}\mu^d\) times its coefficient norm, without a total-volume
factor on the inherited sampling domain. These constants are not
uniform in unbounded field degree, arbitrary background jets,
\(c_E\to0\), or loop order. No new bound for a native measured loop,
cutoff mismatch, or iterated RG map follows merely from this section.

Most importantly \(\mathcal T_{Q_3}\) has ghost number minus two.
It is a primitive of a current charge, not a ghost-number-zero quantum
counterterm. Neither the finite source transformations nor the active
normalization are changed. Turning a reduction of a descent current
into the desired critical quantum primitive still requires the full
descent in the translation-invariant, flavour-invariant, mass-polynomial
category. No assumption that an arbitrary non-natural native reference
component belongs to the natural inventory is introduced.

\section{Verification, source audit and closing ledgers}
\label{sec:v50-ledgers}

The checker \texttt{check\_ESM\_scalar\_cubic\_V50.py} rebuilds the
flat polynomial determining matrices, the skew-operator synthesis
matrices, their rational minors and coefficient sections, and all
stated induced norms. The flat source domain has 252 coefficients,
58 kernel directions and rank 194. The two flavour blocks are
reconstructed from the displayed Fourier formulas, not from a stored
numerical rank. The certificate records all coefficient bases, row and
column selections, and rational entries. Public compilation accepts a
nine- or 36-coefficient block and returns both the selected operator and
the original failed-input defect.

The checker also verifies the Koszul residual identity on all 112
ordered vector basis inputs, the balanced-channel permutation bounds,
the cubic flavour Killing map, Einstein trace reversals, mass-stress
substitution, formal-adjoint differential identities, and the explicit
sixth-matter cancellation. These checks do not independently prove the
natural covariant inventory, the inherited Ricci-flat metric formulas,
or every prior theorem. Those inputs and the new analytic arguments are
identified in their proofs. No native critical loop integral or full
quantum coefficient has been evaluated.

\subsection*{Proved}
\addcontentsline{toc}{subsection}{V50: Proved}
The flat scalar-cubic variational equation has an explicit polynomial
Euler section and, in the actual flavour-sign-invariant natural bank,
an evaluated 58-dimensional skew-operator synthesis inside 252 raw
coefficients. Its decoder and residual are bounded. Every natural
scalar-cubic projection of \eqref{eq:v50-input} reduces to a skew scalar
Euler term plus a Ricci factor. The Weyl-square cubic potentials vanish
by the off-shell flavour Killing test. The full mixed adjoint and the
ghost-number-minus-two ordinary BV primitive are explicit and retain
the scalar-quintic partner. The scalar-cubic coefficient is removed
without asserting removal of its metric-quartic residual.

\subsection*{Inherited}
\addcontentsline{toc}{subsection}{V50: Inherited}
V1--V49 are unchanged. The action, Euler signs, independent flavour
symmetry, naturality category, scalar stress, coframe cotangent map,
and periodic/compact-support conventions are inherited. V47 supplies
the two Ricci-flat metric witnesses. V49 supplies the exact prior
leading-metric reduction and the starting higher-matter equation.
V34, V38 and the gravitational normalization, and the measured
reference, remain unchanged.

\subsection*{Literature input}
\addcontentsline{toc}{subsection}{V50: Literature input}
Koszul relations for a regular sequence are established commutative
algebra; here the needed homogeneous relation is proved with its
explicit polynomial section rather than merely invoked. The relation
of Euler-trivial variational symmetries to negative-ghost BV primitives
is standard \cite{V20BBH1994}. No general theorem that all coupled
Einstein--scalar currents are trivial, and no anomaly-vanishing theorem,
is used to settle the remaining current equation.

\subsection*{Conditional}
\addcontentsline{toc}{subsection}{V50: Conditional}
For an invariant full current in the declared natural category whose
lower matter terms have already been removed, this exact subtraction
removes its entire scalar-cubic term. Its full current class is
unchanged. Applying that statement to an unprocessed native descent
requires its representative and the assumed coupled Noether identity
to be established in the same category. There is no corresponding
claim for non-natural external tensors or general interacting SM
couplings.

\subsection*{Open}
\addcontentsline{toc}{subsection}{V50: Open}
The next local equation is the quartic-scalar metric spectator identity
\eqref{eq:v50-remainder}, together with \eqref{eq:v50-sixth}. The
invariant critical quantum primitive still requires the remaining
current/descent reduction and compatible source coefficients.
Scalar-independent metric transformations, naturality of arbitrary
reference components, internal gauge/chiral transport, higher-loop
control and the invariant full interacting ESM shell class remain
separate interfaces. A new finite exact current section is not an
analytic or nonperturbative ultraviolet completion.

\begin{center}\small
\begin{tabular}{>{\raggedright\arraybackslash}p{.23\textwidth}>{\raggedright\arraybackslash}p{.69\textwidth}}
\toprule
Variable & Scope in V50\\\midrule
Theory & Same ordinary Einstein action and four equal free real scalars;
internal couplings and cosmological source zero.\\
Current category & Natural, sign invariant, scalar-cubic/metric-quartic
leading pair of differential weight four; not all current cohomology.\\
Mass & Polynomial in \(z=m^2\); every section is regular at zero mass.
No infrared limit of a quantum functional is taken.\\
Cutoff and volume & Explicit finite coefficient-map bounds independent of
mesh and periods; no new loop or many-shell estimate.\\
Background/source order & Fixed metric-jet chart, nonzero gravitational
normalization margin, finite arity and fixed window.\\
Verification & Rational polynomial maps and differential checks; no
loop quadrature, numerical rank threshold or proof-assistant proof.\\
\bottomrule
\end{tabular}
\end{center}

\begin{thebibliography}{9}
\bibitem{V50Stacks}
The Stacks Project Authors, \emph{Koszul regular sequences},
Section 15.31, especially Lemma 15.31.2, Tag 062D.
\url{https://stacks.math.columbia.edu/tag/062D}.
\end{thebibliography}

\section*{Append-only continuation marker: V50}
\addcontentsline{toc}{section}{Append-only continuation marker: V50}
\label{sec:v50-continuation-marker}
\textbf{End of V50.} Preserve Sections 1--391 and all active
prescriptions. The natural invariant scalar-cubic projection has a
complete local Euler/skew section with its full mixed and
sixth-matter partners. Next: the quartic-scalar metric spectator
identity or a direct construction of the native invariant critical
primitive. Do not identify this current reduction with a
source-complete quantum continuum theorem.


\clearpage
\part*{V51: A polynomial spectator section for the metric Noether identity}
\addcontentsline{toc}{part}{V51: A polynomial spectator section for the metric Noether identity}

\section{The quartic metric identity is part of one polynomial family}
\label{sec:v51-start}

Sections 1--391 and all their prescriptions are retained. The input is the
\emph{metric-only} identity left by \eqref{eq:v50-remainder},
\begin{equation}
 \int\rho\,e^{ab}F_{ab}(g,\phi)=0
 \quad\text{for every }g,\phi,
 \qquad e^{ab}=c_EG^{ab},\quad c_E\ne0.
 \label{eq:v51-input}
\end{equation}
All scalar fields in this test are arbitrary, not stationary fields. We retain
four equal real scalars, zero cosmological source, and zero internal gauge,
Yukawa and scalar self-couplings. The positive ordinary coefficient chart,
volume pairing, signed Einstein normalization and BV convention are those of
V47--V50. Natural tensors are polynomial in scalar and curvature jets, with
metric contractions and the orientation tensor allowed; no preferred
spacetime tensor, explicit coordinate coefficient, inverse mass or curvature
denominator is admitted. Integration means a compactly supported local test
or the corresponding periodic functional, modulo a translation-invariant
local divergence.

The new step is not another independent tensor census at each matter degree.
We prove a section for every fixed positive homogeneous scalar degree
\(r\), at differential weight four, with \(z=m^2\) of weight two. The case
\(r=4\) solves the leading metric problem of V50. The constants may grow with
\(r\); no infinite-field-degree analytic theorem is intended.

The new ingredient is integration of polynomial one-forms in \emph{scalar
value space}. At quartic order such one-forms need not be closed. Their curl
produces a curvature-commutator term which is absent in the diagonal
quadratic-scalar calculation. This term is retained in the full primitive.
The finite curvature identities of V49 and its inherited terminal V48
classification are inputs; their original rank computations are not repeated.

The conclusion is a complete section of \eqref{eq:v51-input}, followed by its
exact coupled-current completion. It is not a proof that an arbitrary native
quantum breaking admits a natural representative, or that the remaining
scalar-quintic current equation has been solved. No active action, source
normalization, measure or propagator is changed.

\section{The mass-dependent identity at arbitrary scalar degree}
\label{sec:v51-mass}

Write \(F=F^{[4]}+zF^{[2]}+z^2F^{[0]}\). Since \(e\) is independent of
\(z\), each coefficient obeys \eqref{eq:v51-input}. Here brackets count
spacetime derivatives and curvature weight, not scalar degree. For scalar
indices \(A,B\), use \(\partial_A=\partial/\partial\phi_A\). The complete
positive-mass inventory is
\begin{align}
 F^{[2]}_{ab}={}&a_A(\phi)\nabla_a\nabla_b\phi_A
 +g_{ab}b_A(\phi)\Delta\phi_A
 +c_{AB}(\phi)\nabla_a\phi_A\nabla_b\phi_B\nonumber\\
 &+g_{ab}d_{AB}(\phi)\nabla_c\phi_A\nabla^c\phi_B
 +E(\phi)R_{ab}+F_0(\phi)g_{ab}R,
 \qquad F^{[0]}_{ab}=g_{ab}H(\phi).
 \label{eq:v51-mass-bank}
\end{align}
The coefficient degrees are respectively \(r-1,r-1,r-2,r-2,r,r,r\);
\(c,d\) are symmetric in scalar indices. Repeated scalar indices are fully
summed. Negative coefficient degrees mean that the corresponding space is
zero. The tensor inventory follows from the same weight-two contraction
argument as V48, now with polynomial scalar coefficients rather than fixed
matrices. There is no additional orientation-odd symmetric tensor at this
weight.

\begin{theorem}[A polynomial mass section]
\label{thm:v51-mass}
For \(r>0\), the positive-mass part of \eqref{eq:v51-input} holds if and only if
\begin{equation}
 \boxed{c_{AB}=\partial_{(A}a_{B)},\qquad
 b=d=E=F_0=H=0.}
 \label{eq:v51-mass-conditions}
\end{equation}
It is exactly \(\mathcal L_{\xi_z}g\), where
\begin{equation}
 \boxed{(\xi_z)_b=\frac z2 a_A(\phi)\nabla_b\phi_A.}
 \label{eq:v51-mass-xi}
\end{equation}
\end{theorem}
\begin{proof}
Subtract \eqref{eq:v51-mass-xi} and divide by \(c_E\). After one integration
by parts and \(\nabla_aG^{ab}=0\), the remaining integrand is
\begin{align*}
 &G^{ab}\bigl(c_{AB}-\partial_{(A}a_{B)}\bigr)
       \nabla_a\phi_A\nabla_b\phi_B
 +R\bigl(\partial_{(A}b_{B)}-d_{AB}\bigr)
       \nabla\phi_A\cdot\nabla\phi_B\\
 &\hspace{12mm}+b_A\nabla R\cdot\nabla\phi_A
 +E\bigl(R_{ab}R^{ab}-R^2/2\bigr)-F_0R^2.
\end{align*}
The principal scalar Euler coefficients, with arbitrary Ricci and scalar
second jets, force the first two symmetric scalar matrices to vanish.
The remaining zeroth scalar-Euler coefficient contains \(-b_A\Delta R\).
Fourth metric jets vary \(\Delta R\) at fixed lower jets, so \(b=0\) and
then \(d=0\). Independent Ricci tensors give \(\partial_AE=
\partial_AF_0=0\). Positive homogeneity in \(\phi\) gives \(E=F_0=0\).
The \(z^2\) coefficient similarly gives \(\partial_AH=0\), hence \(H=0\).
These are off-shell jet tests. Conversely, differentiating
\eqref{eq:v51-mass-xi} gives exactly the surviving inventory, whose
contraction with \(e\) is a divergence.
\end{proof}

For four scalars of degree four with independent sign invariance, the bank
has \(106\) coefficients. There are \(16\) one-form parameters \(a_A\):
\(\phi_A^3\) and \(\phi_A\phi_B^2\), \(B\ne A\), in each component.
Thus the determining rank is \(90\). Keeping \(a\), replacing \(c\) by
\(\partial_{(A}a_{B)}\), and setting the other coefficients to zero defines
an exact projection. In the ordered coefficient norm its synthesis,
projection and complementary norms are respectively \(4,4,3\), and the
vector decoder has norm \(1/2\). The discarded array retains the original
integrated Noether defect. The degree-two restriction recovers V48; the
polynomial coefficient dependence and the degree-four count are new.

\section{A free Einstein section without a scalar-leg restriction}
\label{sec:v51-flat}

After the mass section, linearize the metric in \eqref{eq:v51-input} at
flat space. Polarize the \(r\) scalar arguments and write \(K=\sum_i p_i\)
for their total momentum. The flat characteristic obeys
\begin{equation}
 E(K)\mathsf F(p_1,\ldots,p_r)=0.
 \label{eq:v51-flat-equation}
\end{equation}
The free Einstein matrix is unreduced. Up to the fixed nonzero normalization,
\begin{equation}
 E(K)H=K^2H-K(K^TH)-(HK)K^T+KK^T\operatorname{tr}H
 +I\bigl(K^THK-K^2\operatorname{tr}H\bigr).
 \label{eq:v51-free-symbol}
\end{equation}
It acts only on the total momentum; the independent relative scalar momenta
are polynomial parameters.

\begin{lemma}[Polynomial gauge exactness with arbitrary momentum parameters]
\label{lem:v51-free-module}
Over the real polynomial ring in \(K\in\R^4\) and any finite additional
commuting variables,
\begin{equation}
 \boxed{\ker E(K)=\{K v^T+vK^T:v\text{ polynomial}\}.}
 \label{eq:v51-free-exact}
\end{equation}
The polynomial vector \(v\) is unique.
\end{lemma}
\begin{proof}
For real \(K\ne0\), rotate \(K\) to a coordinate axis in
\eqref{eq:v51-free-symbol}. Its kernel is the gauge image. Consequently the
Saint--Venant double curl of \(H\) vanishes as a polynomial. In particular,
for \(i\ne j\),
\[
 K_i^2H_{jj}-2K_iK_jH_{ij}+K_j^2H_{ii}=0.
\]
Restricting to \(K_i=0\) proves \(K_i\mid H_{ii}\). Define the polynomial
\(v_i=H_{ii}/(2K_i)\) by exact monomial division. Substitution gives
\(2K_iK_j(K_iv_j+K_jv_i-H_{ij})=0\); the polynomial ring is a domain, hence
\(H=Kv^T+vK^T\). The diagonal identities also prove uniqueness. This proof
uses divisibility, not a rational evaluation on \(K_i=0\).
\end{proof}

An equivariant formula avoids even coordinatewise polynomial division.
Decompose a coefficient into homogeneous total-\(K\) pieces \(H^{(n)}\),
with the relative momenta fixed. Set
\begin{equation}
 \boxed{
 (\mathsf H_nH)_j=
 \frac1{n+4}\left(
 \partial_{K_i}H_{ij}
 -\frac{K_j}{2(n+3)}
       \partial_{K_i}\partial_{K_l}H_{il}\right),\qquad n\ge1.
 }
 \label{eq:v51-free-section}
\end{equation}
\begin{proposition}[A bounded soft-transfer-regular section]
\label{prop:v51-free-section}
For \(H=Kv^T+vK^T\) of degree \(n\), \(\mathsf H_nH=v\).
The constant-\(K\) piece in \(\ker E\) is zero. Through total derivative
weight four, the section has coefficient norm at most \(13/14\) in the
full symmetric-entry norm, before changing relative momentum coordinates.
\end{proposition}
\begin{proof}
A direct differentiation gives
\[
 \partial_{K_i}H_{ij}=(n+4)v_j+K_j\operatorname{div}_Kv,
 \qquad
 \partial_{K_j}\partial_{K_i}H_{ij}=2(n+3)\operatorname{div}_Kv.
\]
This proves the formula. At \(n=0\) the polynomial divisibility proof
forces zero. The first derivative has norm at most \(n\), and the second
at most \(n(n-1)\). Multiplication by the vector \(K\) costs at most four
in the summed component norm. Thus
\[
 \|\mathsf H_n\|\le
 \frac{n+2n(n-1)/(n+3)}{n+4}\le\frac{13}{14},\quad 1\le n\le4.
\]
The maximum is a bound, not an assertion that every input attains it.
\end{proof}

Restoring the conventional Fourier factor gives a real differential vector
of weight three. Its uniqueness implies scalar-exchange and orthogonal
covariance whenever the input has them, even though one relative-momentum
coordinate system was used to compute it. Covariantize by one fixed
ordering of derivatives and restore the metric contractions. The resulting
natural \(\xi_0\) has exactly the supplied flat symbol. Subtracting
\(\mathcal L_{\xi_0}g\) preserves \eqref{eq:v51-input}; its remainder is
massless and vanishes at flat metric. Every normal tensor monomial of that
remainder contains curvature. Derivative commutators are retained in it.
For fixed scalar arity, coordinate conversion and polarization cost finite
constants; none depends on a lattice cutoff. This argument includes every
quartic flavour channel, not just a two-scalar symbol.

\section{Polynomial coefficient one-forms and the complete curvature inventory}
\label{sec:v51-curvature-reduction}

Use the seven maps \(M_i\), five maps \(N_j\), and eleven tensors \(D_l\)
from \eqref{eq:v49-Mbank}--\eqref{eq:v49-Dbank}. For scalar degree \(r\), the
symmetric-gradient portion is
\begin{equation}
\begin{aligned}
 F={}&\sum_iM_i\left(a_{iA}(\phi)\nabla^2\phi_A
       +b_{iAB}(\phi)\,\nabla\phi_A\otimes\nabla\phi_B\right)\\
 &+\sum_jN_j\left(c_{jA}(\phi)\,\dd\phi_A\right)
       +\sum_l d_l(\phi)D_l,
\end{aligned}
 \label{eq:v51-curvature-bank}
\end{equation}
where \(b_{iAB}=b_{iBA}\). The degrees are \(r-1,r-2,r-1,r\).
There are also two antisymmetric-gradient maps, which must not be lost:
\begin{equation}
 [\operatorname{Ric},W],\qquad
 [\operatorname{Ric},{}^\star\widetilde W],\qquad
 W=\sum_{A<B}w_{AB}(\phi)\dd\phi_A\wedge\dd\phi_B,
 \label{eq:v51-twoform-bank}
\end{equation}
with \(\deg w_{AB}=r-2\), and the same definition for \(\widetilde W\).
Commutators use the metric to regard these tensors as endomorphisms.
For symmetric \(T\) and antisymmetric \(W\), \([T,W]\) is symmetric and
the map \(T\mapsto[T,W]\) is skew-adjoint on symmetric tensors.

\begin{lemma}[The missing multi-flavour contractions]
\label{lem:v51-inventory}
Equations \eqref{eq:v51-curvature-bank}--\eqref{eq:v51-twoform-bank} form
a complete independent normal tensor inventory. The two additional maps
are identically Einstein-orthogonal.
\end{lemma}
\begin{proof}
At weight four and positive curvature degree there are at most two scalar
derivatives. Their symmetric spacetime part gives precisely the inherited
\(M_i\) contractions. The antisymmetric part is a two-form. Contracting a
Riemann tensor with it and symmetrizing the two output indices kills the
Weyl contribution by pair exchange and the scalar-curvature contribution
by antisymmetry. The remaining Ricci contribution is its commutator with
the two-form. With one orientation tensor the same argument gives the
commutator with its dual. These two maps are independent, for example at
a generic diagonal Ricci tensor, and no further map remains. One and zero
scalar derivatives give the inherited \(N_j,D_l\) inventories. Arbitrary
scalar jets separate the derivative distributions and target coefficients.
Finally \(G\) commutes with \(\operatorname{Ric}\), so
\(G:[\operatorname{Ric},W]=0\), and likewise for the dual map.
\end{proof}

For clarity, the finite metric-current map used in the next proof is
\[
 \mathcal J(a,c)^b
 =\sum_jc_jG^{uv}(N_j)_{uv}{}^b
 -\nabla_a\left(\sum_i a_iG^{uv}(M_i)_{uv}{}^{ab}\right),
 \qquad (a,c)\in\R^{12}.
\]
The coefficients here are constants when defining \(\mathcal J\).
The scalar-linear consequence of V49 is
\begin{equation}
 \ker\mathcal J=\ker(\nabla\cdot\mathcal J)
 =\operatorname{span}\{U_1,U_2,U_4,U_5,U_6\},
 \label{eq:v51-current-kernel}
\end{equation}
where the columns are the first twelve entries of
\eqref{eq:v49-columns}. Indeed a divergence-free \(\mathcal J\) defines
a scalar-linear Noether operator with no undifferentiated-scalar term;
V49 leaves exactly these five columns. Their currents vanish pointwise:
the gauge columns use Bianchi and the other columns are the displayed
pointwise or differential skew identities. Thus no extra covariant
superpotential is assumed to be absent.

\begin{proposition}[Reduction in scalar-value space]
\label{prop:v51-target-homotopy}
Every curvature-dependent solution of \eqref{eq:v51-input} is a sum of
polynomial-coefficient versions of the five columns in
\eqref{eq:v51-current-kernel}, a scalar-function version of
\(\mathcal L_{u\nabla R}g\), the terminal \(u\Omega_e e\) term, and
pointwise symmetric- and antisymmetric-gradient skew terms. The reduction
uses polynomial differentiation and a single homogeneous integration factor
\(1/r\); it has no singularity at \(\phi=0\).
\end{proposition}
\begin{proof}
Remove \eqref{eq:v51-twoform-bank}, which is already pointwise skew.
Integrate the scalar Hessians in \eqref{eq:v51-curvature-bank} once.
The principal scalar Euler coefficient is
\[
 G^{uv}\sum_i(M_i)_{uv}{}^{ab}
 \bigl(b_{iAB}-\partial_{(A}a_{iB)}\bigr).
\]
The pointwise V49 syzygy forces each target pair to be a multiple of
\((1,-1/2,0,1,0,0,0)\). Subtract the corresponding
\(\Omega_{T(\phi,\dd\phi)}G\). The remaining scalar density, modulo a
divergence, is
\[
 K_A^b(g,\phi)\nabla_b\phi_A+\sum_l d_l(\phi)(G:D_l),
 \qquad K_A=\mathcal J(a_{\cdot A},c_{\cdot A}).
\]
Its Euler coefficient linear in scalar gradients is
\(\partial_AK_B-\partial_BK_A\). It must vanish. Consequently the
\(\R^{12}/\ker\mathcal J\)-valued polynomial one-form represented by
\((a_{\cdot A},c_{\cdot A})\) is closed.

Choose a fixed rational complement to \(\ker\mathcal J\). If \(P_A\)
is the projected one-form, put \(p=r^{-1}\phi_AP_A\). For a homogeneous
one-form of coefficient degree \(r-1\),
\begin{equation}
 P_A-\partial_Ap
 =\frac1r\phi_B(\partial_BP_A-\partial_AP_B).
 \label{eq:v51-radial-identity}
\end{equation}
Thus \(P_A=\partial_Ap\). The remaining five coefficient one-forms need
not be closed; they are subtracted using their exact gauge or skew
versions on arbitrary covectors. The latter versions are constructed
explicitly below, including the Ricci commutator required for a covector
with nonzero curl.

After these subtractions, the characteristic is a scalar-linear operator
applied to polynomial scalar functions, plus \(\sum_l d_l(\phi)D_l\).
Integrating its scalar derivatives onto metric coefficients yields
\(\sum_{|I|=r}\phi^I L_I(g)\). Its scalar Euler identity gives every
\(L_I=0\), since \(r>0\). Apply the inherited complete scalar-linear
V49 classification to each monomial. Its five previously used columns,
the \(\mathcal L_{u\nabla R}g\) column and the terminal kernel exhaust
it. No repeated general-matter cohomology theorem is being invoked.
\end{proof}

\section{A complete curvature section and its quartic coefficient matrix}
\label{sec:v51-curvature-section}

Let \(\alpha,\beta,\delta,\varepsilon,\zeta\) be five polynomial one-forms
on scalar-value space, with components of degree \(r-1\). Write
\(v_\alpha=\alpha_A\dd\phi_A\), and similarly for the others.
Let \(\gamma,\kappa\) be homogeneous scalar polynomials of degree \(r\),
let \(\theta_{AB}=\theta_{BA}\) have degree \(r-2\), and define
\[
 T=\delta_A\nabla^2\phi_A+
          \theta_{AB}\nabla\phi_A\otimes\nabla\phi_B.
\]
Use two freely specified two-forms \(W,\widetilde W\) as in
\eqref{eq:v51-twoform-bank}. On symmetric tensors put
\[
 \mathscr C_v(T_0)=-2\nabla_{v^\sharp}T_0-(\nabla\cdot v)T_0,
 \quad
 \mathscr C_v^0(T_0)=g\mathscr C_v(\operatorname{tr}T_0),
 \quad
 \mathscr J_W(T_0)=[T_0,W].
\]
Each is formally skew-adjoint. So is \(\Omega_T\) from V49.

\begin{theorem}[Polynomial curvature--spectator classification]
\label{thm:v51-curvature}
Every curvature-dependent solution of \eqref{eq:v51-input}, and only such
a solution, has the unique parameter form
\begin{equation}
 \boxed{F=\mathcal L_\xi g+\mathscr A_\phi e,}
 \label{eq:v51-curv-main}
\end{equation}
where
\begin{align}
 \xi={}&R\,v_\alpha^\sharp+
 \operatorname{Ric}^{\sharp}v_\beta+\gamma\nabla R,\nonumber\\
 \mathscr A_\phi={}&\frac1{c_E}\left(
 \Omega_T+\mathscr C_{v_\varepsilon}
 +\mathscr C_{v_\zeta}^{0}
 +\mathscr J_{\frac12\dd v_\beta+W}
 +\mathscr J_{{}^\star\widetilde W}\right)
 +\frac\kappa{c_E^2}\Omega_e.
 \label{eq:v51-curv-parameters}
\end{align}
The factor \(\dd v_\beta/2\) is part of the formula, not an optional
homogeneous addition.
\end{theorem}
\begin{proof}
The target-space reduction proves exhaustion once the exact representatives
are checked. For an arbitrary covector \(v\), write
\(\nabla v=S+W_v\), with \(S\) symmetric and \(W_v=\dd v/2\).
Then the literal covariant identity is
\begin{equation}
 \mathcal L_{\operatorname{Ric}^{\sharp}v}g+
 [\operatorname{Ric},W_v]
 =M_5(S)+N_4(v).
 \label{eq:v51-curl-compensation}
\end{equation}
Indeed the derivative terms are
\((\nabla v)\operatorname{Ric}+\operatorname{Ric}(\nabla v)^T\),
whose antisymmetric-derivative contribution is
\(-[\operatorname{Ric},W_v]\). Since \([G,W_v]=
[\operatorname{Ric},W_v]\), the added operator in
\eqref{eq:v51-curv-parameters} gives exactly the needed compensation.
The other four current-kernel columns extend respectively to
\(\mathcal L_{Rv_\alpha^\sharp}g\),
\(\Omega_{\operatorname{sym}\nabla v_\delta}G\),
\(\mathscr C_{v_\varepsilon}G\), and
\(\mathscr C_{v_\zeta}^0G\). Differences between
\(\operatorname{sym}\nabla v_\delta\) and
\(\delta_A\nabla^2\phi_A\) are absorbed into \(\theta\).
All these formulas are gauge or skew-Euler identities on arbitrary metrics.
The terminal representative is inherited from V48. This proves sufficiency
and exhaustion.

For explicit uniqueness and computation, the raw coefficients determine
\begin{align}
 \gamma&=d_1/2,&\kappa&=(d_4+d_7)/2,&
 \alpha_A&=c_{1A}-\partial_A\gamma,\nonumber\\
 \beta_A&=c_{4A},&\delta_A&=a_{4A},&
 \varepsilon_A&=-c_{3A}/2,\nonumber\\
 \zeta_A&=c_{2A}/2+c_{3A}/4,&\theta_{AB}&=b_{4AB}.
 \label{eq:v51-decoder}
\end{align}
The two free two-forms are precisely the two raw antisymmetric-gradient
coefficient arrays. Substitution recovers all parameters, proving
injectivity.
\end{proof}

For reproducibility, put \(S\omega_{AB}=\partial_{(A}\omega_{B)}\).
The nonzero raw synthesis rows are
\begin{align*}
 &(a_1,a_2,a_3,a_4,a_5)
 =(2\alpha+\delta,-\delta/2+\varepsilon/2+\zeta,
       -\varepsilon,\delta,\beta),\\
 &(c_1,c_2,c_3,c_4)
 =(\alpha+\dd_\phi\gamma,\varepsilon+2\zeta,-2\varepsilon,\beta),\\
 &(b_1,b_2,b_3,b_4,b_5)
 =(2S\alpha+\theta,S\varepsilon/2+S\zeta-\theta/2,
       -S\varepsilon,\theta,S\beta),\\
 &d_1=2\gamma,\qquad d_4=d_7=\kappa.
\end{align*}
All unlisted symmetric-gradient rows vanish; the raw two-form rows are
copied. This list already includes the cancellation in
\eqref{eq:v51-curl-compensation}. It does not discard its operator partner.

For four sign-invariant scalars and \(r=4\), the scalar-function,
one-form, symmetric two-tensor and antisymmetric two-tensor coefficient
spaces have dimensions \(10,16,22,6\), respectively. Therefore
\begin{equation}
 \boxed{
 \dim\mathcal V_{\mathrm{curv}}=468,\qquad
 \dim\ker\mathcal N_{\mathrm{curv}}=134,\qquad
 \operatorname{rank}\mathcal N_{\mathrm{curv}}=334.
 }
 \label{eq:v51-quartic-counts}
\end{equation}
The kernel consists of five one-forms, two scalar functions, one symmetric
coefficient tensor and two antisymmetric coefficient tensors:
\(5\cdot16+2\cdot10+22+2\cdot6=134\).

Let \(H_{51}\) be \eqref{eq:v51-decoder} and \(N_{51}\) the synthesis.
Use monomial coefficient norms, summing both ordered entries of symmetric
target tensors (equivalently, weighting each independent off-diagonal
entry by two). The exact sparse matrices obey
\begin{equation}
 \boxed{
 H_{51}N_{51}=I_{134},\quad
 \|H_{51}\|=5/2,\quad \|N_{51}\|=9,\quad
 \|N_{51}H_{51}\|=17,\quad \|I-N_{51}H_{51}\|=16.
 }
 \label{eq:v51-norms}
\end{equation}
The certificate exports all 468 raw labels, 134 parameter labels, and both
rational maps. Rank 334 follows from analytic exhaustion together with the
verified left inverse; it is not a numerical rank fit to sampled metrics.
For every failed input, the complementary coefficient array retains exactly
its integrated Noether defect, because the synthesized characteristic is
always a gauge/skew-Euler identity.

\section{All fixed scalar degrees and an explicit nonclosed one-form test}
\label{sec:v51-all-degrees}

For \(m\) scalars with independent sign symmetry and degree \(r=2k>0\), set
\[
 U_k=\binom{m+k-1}{k},\qquad
 V_k=m\binom{m+k-2}{k-1},\qquad
 A_k=\binom m2\binom{m+k-3}{k-2},
\]
with \(A_1=0\). Symmetric target two-tensors have dimension \(V_k+A_k\).
The proof and decoder above apply coefficientwise and give
\begin{equation}
 \boxed{
 \dim\mathcal V=19V_k+9A_k+11U_k,\qquad
 \dim\ker\mathcal N=6V_k+3A_k+2U_k.
 }
 \label{eq:v51-general-counts}
\end{equation}
For \(k=1,m=4\), these are exactly V49's inherited \(120,32\) counts;
no new quadratic-scalar classification is claimed. For the positive-mass
bank the corresponding dimensions are \(4V_k+2A_k+3U_k\) and \(V_k\).

The coefficient norms have explicit degree dependence. For \(r\ge2\),
\begin{equation}
 \begin{split}
 \|H_r\|&=(r+1)/2,\\
 \|N_r\|&=\max\{2r+1,3r/2+3\},\\
 \|N_rH_r\|&=r^2+1,\\
 \|I-N_rH_r\|&=r^2.
 \end{split}
 \label{eq:v51-general-norms}
\end{equation}
These identities follow by the absolute column sums of the displayed
synthesis and decoder: the \(d_1=\phi_A^r\) column gives \(1+r^2\),
and its complementary column gives \(r^2\). The other columns are bounded
by these values. For the synthesis the largest columns are the \(\alpha\)
and \(\varepsilon\) monomials; for the decoder it is \(d_1\).
The mass projection and complement have norms \(r\) and \(r-1\).
All statements are at a prescribed finite scalar degree.

A concrete new quartic channel is
\[
 v_\beta=\phi_1\phi_2^2\,\dd\phi_1,
 \qquad
 \dd v_\beta=-2\phi_1\phi_2\,\dd\phi_1\wedge\dd\phi_2.
\]
It is invariant under the two independent scalar sign changes but is not
closed. The normal tensor \(M_5(\operatorname{sym}\nabla v_\beta)
+N_4(v_\beta)\) must therefore be completed as
\begin{equation}
 \boxed{
 \mathcal L_{\operatorname{Ric}^{\sharp}v_\beta}g
 +\mathscr J_{\dd v_\beta/2}G.
 }
 \label{eq:v51-example}
\end{equation}
At a point with \(\operatorname{Ric}=\operatorname{diag}(1,2,0,0)\),
\(\phi_1=\phi_2=1\), \(\dd\phi_1=dx^0\), and
\(\dd\phi_2=dx^1\), the added commutator has nonzero \(01\) entry.
Thus tensoring the old exact-gradient formulas with undifferentiated scalar
monomials would miss a literal source term. All metric and scalar jet tests
here are realizable local tests, not selected Gaussian configurations.

\begin{theorem}[Complete polynomial metric-spectator section]
\label{thm:v51-global-section}
At every fixed positive scalar degree, a natural mass-polynomial weight-four
characteristic satisfies \eqref{eq:v51-input} if and only if it can be
written
\begin{equation}
 \boxed{F=\mathcal L_\xi g+\mathscr A_\phi e,
 \qquad \mathscr A_\phi^\dagger=-\mathscr A_\phi,}
 \label{eq:v51-total-section}
\end{equation}
with finite local polynomial scalar coefficients. An explicit section is:
first use \eqref{eq:v51-mass-xi}, next
\eqref{eq:v51-free-section} and its covariantization, and finally
\eqref{eq:v51-decoder} on the remaining curvature inventory. It is regular
at zero scalar field, zero mass and zero external momentum.
\end{theorem}
\begin{proof}
The mass theorem removes all positive-mass pieces. The polynomial free
Einstein section removes the complete flat massless symbol. Naturality and
the weight bound put the residual in
\eqref{eq:v51-curvature-bank}--\eqref{eq:v51-twoform-bank}, to which
\Cref{thm:v51-curvature} applies. Conversely each displayed gauge term is
killed by Bianchi and each skew term is killed in the volume pairing.
Every stage is a finite coefficient operation at fixed scalar degree and
finite derivative weight. No general completeness theorem for all local
Einstein Noether operators is used to replace these steps.
\end{proof}

\section{The full current subtraction and the new remaining equation}
\label{sec:v51-BV}

Let \(\tau\) and \(f=P_g\phi\) be the total scalar metric stress and
scalar Euler expression in V50. For the section
\eqref{eq:v51-total-section}, the full selected characteristic is
\begin{equation}
 \boxed{
 \delta g=\mathcal L_\xi g+\mathscr A_\phi(e+\tau),\qquad
 \delta\phi=\mathcal L_\xi\phi.
 }
 \label{eq:v51-completed-current}
\end{equation}
With \(\vartheta=\rho^{-1}\eta_g\) and \(\sigma=c^*\), put
\begin{equation}
 \boxed{
 \mathcal T_F=\int\sigma_a\xi^a
 -\frac12\int\rho\,\vartheta:\mathscr A_\phi\vartheta.
 }
 \label{eq:v51-BV-primitive}
\end{equation}
\begin{theorem}[Exact coupled completion of the quartic metric section]
\label{thm:v51-BV}
In the inherited ordinary BV convention,
\begin{equation}
 s\mathcal T_F
 =\int\eta_g:\delta g+\int\upsilon^T\delta\phi
 \pmod{d_x}.
 \label{eq:v51-BV-charge}
\end{equation}
For scalar degree four, this removes the complete leading metric-quartic
characteristic and retains a scalar-quintic gauge variation and a
metric-sextic stress return. The primitive has ghost number minus two.
\end{theorem}
\begin{proof}
The gauge part is the full field-dependent diffeomorphism current, with
its ghost-antifield primitive. The longitudinal variation of the second
term is a density divergence because all coefficients, including the
curl and dual-curl terms, are natural. Its Euler variation gives
\(\eta_g:\mathscr A_\phi(e+\tau)\); the two odd antifields give the
sign in \eqref{eq:v51-BV-primitive}. The operator is formally skew even
when its coefficients depend on scalar jets. Directly, the skew part of
the action variation vanishes and the gauge part is the coupled Noether
divergence. This proves the identity without putting either field on shell.
\end{proof}

If \(\widehat F_4\) is V50's remaining characteristic, subtract the whole
charge \eqref{eq:v51-BV-charge}, not just its leading metric term. Then
\begin{equation}
 \boxed{
 \widehat F_4^{\mathrm{new}}=0,\qquad
 \widehat F_6^{\mathrm{new}}=\widehat F_6-\mathscr A_\phi\tau,
 \qquad
 \widehat Q_5^{\mathrm{new}}=\widehat Q_5-\mathcal L_\xi\phi.
 }
 \label{eq:v51-descendants}
\end{equation}
The next actual current equation is
\begin{equation}
 \boxed{
 \int\rho\left(
 e:\widehat F_6^{\mathrm{new}}
 +f^T\widehat Q_5^{\mathrm{new}}
 \right)=0.
 }
 \label{eq:v51-next}
\end{equation}
It is not asserted solved. If a later scalar reduction produces another
metric-only identity at this derivative weight, the all-degree section
above applies without another metric tensor census. A scalar-quintic
variational classification, non-natural components, and scalar-independent
metric symmetries are separate tasks.

For a fixed scalar degree, the maps involve finitely many derivatives,
products and the indicated powers of \(c_E^{-1}\). On a fixed nondegenerate
metric-jet chart and \(|c_E|\ge c_*>0\), their coefficient and physical
jet norms have finite constants. A smooth external momentum window turns
these finite polynomials into rooted kernels with every prescribed
finite moment; repeated integration by parts in all relative momenta
and periodization gives bounds with no total-volume multiplier. The
constants retain their degree, jet-chart, window, and coupling-margin
dependence. No new estimate for loop integrals or a many-shell flow follows
from this finite polynomial statement.

\section{Status, inherited inputs and verification}
\label{sec:v51-ledgers}

\subsection*{Proved}
\addcontentsline{toc}{subsection}{V51: Proved}
The positive-mass metric identity has a polynomial scalar-coefficient
section at every positive scalar degree. The unreduced free Einstein symbol
has an explicit pole-free polynomial gauge section with arbitrary scalar
momenta. The curvature-dependent identity has a complete polynomial
spectator section, including the nonclosed target one-form curl and both
antisymmetric-gradient contractions. In the actual invariant scalar-quartic
curvature bank the dimensions are \(468,134,334\), and the decoder,
synthesis and residual norms are \(5/2,9,16\), with projection norm 17.
The whole metric-quartic residual of V50 admits the exact current subtraction
\eqref{eq:v51-BV-primitive}, with all higher-matter partners retained.

\subsection*{Inherited}
\addcontentsline{toc}{subsection}{V51: Inherited}
V1--V50 are unchanged. The action, ordinary coordinate chart, Euler/BV signs,
flavour symmetry, source interpretation and active prescriptions remain
those already recorded. The finite pointwise curvature syzygy, scalar-linear
current columns and terminal curvature classification are inherited from
V48--V49. Their proofs are used, not silently reverified by the new sparse
matrix checks. V50 supplies the scalar-cubic subtraction and the particular
remaining metric-only identity. V34, V38 and the fixed gravity normalization
remain unchanged.

\subsection*{Literature input}
\addcontentsline{toc}{subsection}{V51: Literature input}
Local gauge and Euler-trivial characteristics and their negative-ghost BV
primitives are established methodology \cite{V20BBH1994}. The polynomial
radial identity used here is proved directly in
\eqref{eq:v51-radial-identity}. No external theorem asserting triviality
of all coupled Einstein--scalar currents or anomaly freedom is used to
close \eqref{eq:v51-next} or the quantum descent.

\subsection*{Conditional}
\addcontentsline{toc}{subsection}{V51: Conditional}
For a full natural invariant current satisfying the preceding coupled
Noether identities, this exact subtraction preserves its class and removes
its metric-quartic part. Applying that conclusion to an unprocessed native
quantum descent still requires its representative and its identity in the
same category. All-degree here means a theorem at each prescribed finite
polynomial scalar degree, not analytic convergence in an infinite field
series or uniformity in unbounded source order.

\subsection*{Open}
\addcontentsline{toc}{subsection}{V51: Open}
The next leading equation is \eqref{eq:v51-next}. The scalar-quintic
projection, scalar-independent metric transformations, naturality of
arbitrary reference components, and completion of the translation-invariant
invariant quantum descent remain separate. Internal gauge/chiral transport,
higher-loop estimates and the invariant measured ESM action class have not
been established by this local current theorem.

\begin{center}\small
\begin{tabular}{p{.23\textwidth}p{.69\textwidth}}
\toprule
Variable & Scope in V51\\\midrule
Theory & Ordinary native Einstein action and four equal free real scalars;
no change to the finite carrier or active normalizations.\\
Field degree & Complete metric-only section at every fixed positive
polynomial scalar degree; explicit invariant quartic matrices.\\
Mass & Polynomial in \(z=m^2\); no inverse mass and no removal of the
positive quantum lower reference scale.\\
Cutoff and volume & Finite coefficient-map bounds independent of mesh and
periods; no new loop or many-scale convergence rate.\\
Naturality & Metric, curvature, scalar jets and orientation only; no claim
for arbitrary external preferred tensors.\\
Verification & Rational sparse sections, polynomial identities and
preservation audit; no loop quadrature or proof-assistant formalization.\\
\bottomrule
\end{tabular}
\end{center}

\section*{Append-only continuation marker: V51}
\addcontentsline{toc}{section}{Append-only continuation marker: V51}
\label{sec:v51-continuation-marker}
\textbf{End of V51.} Preserve Sections 1--399 and all active prescriptions.
The metric-only spectator identity has one finite local section for each
positive polynomial scalar degree, so the metric-quartic residual is no
longer a separate unresolved tensor census. Next: the scalar-quintic/metric-
sextic coupled current equation or a direct construction of the complete
native invariant critical quantum primitive. Do not identify a current
primitive of ghost number minus two with the required quantum action
counterterm.


\clearpage
\part*{V52: An all-degree scalar section and finite-order higher-matter current elimination}
\addcontentsline{toc}{part}{V52: An all-degree scalar section and finite-order higher-matter current elimination}

\section{The scalar-quintic equation belongs to one polarized family}
\label{sec:v52-start}

Sections 1--399 and the V51 closing marker are retained verbatim. The
next input is exactly \eqref{eq:v51-next}. We keep the ordinary native
Einstein--scalar action, its signed normalization, volume pairing,
coframe cotangent convention, and all active prescriptions. No quantum
reference, finite measure, propagator, or interaction is changed. In
particular internal gauge, Yukawa, scalar self-interaction, and
cosmological terms remain zero in this calculation.

For an odd scalar degree \(r\ge3\), consider a natural local identity
\begin{equation}
 \int\rho\bigl(e^{ab}(F_{r+1})_{ab}+f_AQ_{r,A}\bigr)=0,
 \qquad e^{ab}=c_EG^{ab},\quad f=P_g\phi,\quad
 P_g=-\Delta_g+z,\quad z=m^2.
 \label{eq:v52-input}
\end{equation}
It holds for arbitrary off-shell metric and scalar tests and as a
polynomial in \(z\). Both characteristics have differential weight four;
\(\wt(\nabla)=1\) and \(\wt(z)=\wt(\mathrm{curvature})=2\).
The coefficient category is that of V50--V51: polynomial scalar and
curvature jets, metric and orientation contractions, no preferred
spacetime tensors or explicit coordinates, and no negative powers of
mass, fields, or curvature. Four equal real scalars and the actual
independent flavour-sign symmetry are retained. Statements about the
flat labelled algebra below do not need that last symmetry.

V50 solved \(r=3\). V51 supplied the metric-only section at every fixed
positive scalar degree. The new question is whether the scalar section
also persists at every fixed \(r\), beginning with \(r=5\). Two genuine
issues arise: Gram matrices of five or more independent formal momenta
have rank relations in dimension four, and orientation-odd contractions
can occur at sufficiently high scalar degree. Neither issue is removed
by declaring the full Gram ring to be polynomial.

The result will be
\begin{equation}
 \boxed{Q_r=\mathcal L_{r-1}f+\mathcal K_re,
 \qquad \mathcal L_{r-1}^{\dagger}=-\mathcal L_{r-1}.}
 \label{eq:v52-target}
\end{equation}
The subscripts on the operators indicate scalar coefficient degree.
Together with V51, this gives a finite algorithm for the natural
higher-matter current at every prescribed scalar order, rather than a
new separate tensor census at each order. It does not construct a
quantum breaking representative or its ghost-number-zero counterterm.

\section{Euler regularity with arbitrarily many scalar legs}
\label{sec:v52-regularity}

Put \(N=r+1\), take \(p_1,\ldots,p_{N-1}\in\R^4\) as independent
component variables, and set \(p_N=-\sum_{i<N}p_i\). The scalar variation
of \eqref{eq:v52-input} at flat metric gives, after complete labelled
polarization,
\begin{equation}
 W(q):=\sum_{i=1}^N x_iq_i=0,\qquad x_i=p_i^2+z.
 \label{eq:v52-polarized}
\end{equation}
Each root coefficient \(q_i\) has differential weight four. The common
polarization normalization is removed on both sides, as in V50. Flavour
labels travel with their legs. Compactly supported tests determine the
Fourier polynomial on the total-momentum hyperplane; equality for real
momenta and positive \(z\) is a polynomial identity, not a restriction to
stationary scalar waves.

\begin{lemma}[A regular sequence on the actual momentum component ring]
\label{lem:v52-regular}
For every \(N\ge3\), the sequence \(x_1,\ldots,x_N\) is regular in
\[
 \R[z,p_i^\mu:1\le i<N,\ 0\le\mu<4].
\]
In particular, every polynomial relation \eqref{eq:v52-polarized}
admits \(q_i=\sum_j A_{ij}x_j\), with \(A_{ij}=-A_{ji}\).
For homogeneous weight-four \(q\), the entries \(A_{ij}\) can be taken
homogeneous of weight two.
\end{lemma}
\begin{proof}
Start with the domain \(\R[z]\). At each of the first \(N-1\) steps,
adjoin four new variables and impose their squared sum plus \(z\).
Over the fraction field of the preceding domain this quadratic is
irreducible: a product of affine linear polynomials would have a
quadratic homogeneous part of rank at most two, whereas its rank is
four. It is monic in one of the new variables. Polynomial division in
that variable shows that the quotient over the domain embeds into the
corresponding quotient over its fraction field. Thus every intermediate
quotient is a domain and each new polynomial is a nonzerodivisor.

The last polynomial is nonzero in the resulting domain. Indeed, map
all \(p_i\), \(i<N\), to the same vector \(v\) and set \(z=-v^2\).
This is a homomorphism to \(\R[v]\) on the intermediate quotient. The
image of \(x_N\) is \(((N-1)^2-1)v^2\ne0\). This algebraic witness
is not a mass substitution in a quantum coefficient. The common zero
at the origin shows the ideal is proper. Hence the sequence is regular.

For completeness, its first syzygy assertion follows by induction.
Reduce a relation modulo \((x_1,\ldots,x_{N-1})\). Regularity forces
\(q_N\) into that ideal; write \(q_N=\sum_{i<N}b_ix_i\). Subtract the
skew pairs \(A_{Ni}=b_i\), \(A_{iN}=-b_i\). The remaining relation
uses the first \(N-1\) factors and is handled inductively. Homogeneous
components give the stated weight. This is the elementary first stage
of the regular-sequence Koszul argument \cite{V50Stacks}.
\end{proof}

The restriction \(N\ge3\) is real: at \(N=2\), momentum conservation
gives \(x_1=x_2\). The scalar-linear current problem, including its
rotation classes, is not absorbed into this lemma.

\begin{theorem}[No orientation-odd nonlinear flat variational coefficient]
\label{thm:v52-odd}
At differential weight four and \(N\ge3\), the orientation-odd part of
a natural coefficient satisfying \eqref{eq:v52-polarized} is zero as
a momentum polynomial.
\end{theorem}
\begin{proof}
Apply Lemma~\ref{lem:v52-regular} to that part to obtain a weight-two skew
matrix \(A\), initially with arbitrary momentum-component coefficients.
For the sixteen diagonal coordinate reflections \(\epsilon\), apply
\[
 \mathscr P_{\det}A
 =\frac1{16}\sum_{\epsilon\in\{\pm1\}^4}
     \left(\prod_\mu\epsilon_\mu\right)A(\epsilon p,z).
\]
All \(x_i\) are unchanged, and the orientation-odd \(q\) is fixed by
this projector. But a momentum monomial of weight two cannot be odd
in each of four coordinate directions; a factor \(z\) contributes no
coordinate sign. Therefore \(\mathscr P_{\det}A=0\), and hence \(q=0\).
This is a finite coefficient operation, not averaging the regulator.
\end{proof}

For the sign-invariant four-scalar quintic bank the orientation-odd
inventory is already zero: the total six scalar legs have even flavour
multiplicities, hence at most three flavours, while a nonzero scalar
four-gradient alternating product needs four distinct differentiated
flavours. At scalar degree seven the inventory does contain, for example,
\[
 Q_1=\phi_2\phi_3\phi_4\,
 \varepsilon^{abcd}\nabla_a\phi_1\nabla_b\phi_2
                         \nabla_c\phi_3\nabla_d\phi_4.
\]
The theorem excludes this nonzero flat polynomial from a compatible
variational identity; it is not omitted by an inventory assumption.
An explicit labelled test sets the two flavour-1 momenta to \(e_0,2e_1\),
the sums of the flavour-2 and flavour-3 pairs to \(e_2,e_3\), and the
last pair sum to their negative total. Its normalized Euler coefficient
is \((1-4)\det(e_0,2e_1,e_2,e_3)=-6\).

\section{A permutation-equivariant low-degree Gram section}
\label{sec:v52-Gram}

The even coefficient is a polynomial of invariant degree two in scalar
products and \(z\). Its Euler defect has invariant degree three.
There are no Gram relations at those degrees, even with many scalar legs.

\begin{lemma}[The needed stable degree, not an unrestricted Gram ring]
\label{lem:v52-Gram-injective}
A polynomial of degree at most three in the scalar products of any
number of vectors in four dimensions, with polynomial parameter \(z\),
vanishes identically only if its formal Gram polynomial vanishes.
Here degree is the total invariant degree including \(z\).
\end{lemma}
\begin{proof}
Separate the powers of \(z\) and polarize each vector variable. A
homogeneous polynomial of Gram degree \(d\le3\) becomes a linear
combination of pairings of \(2d\) labelled vector slots. Those pairing
tensors are independent in dimension at least \(d\): to test a chosen
pairing, assign one distinct orthonormal basis vector to its two slots,
for each of the \(d\) pairs. All other pairings vanish on that test.
Polarization is injective in characteristic zero, proving the claim.
This also gives independence at Gram degree four; the stated smaller
range is all that is used. The five-by-five Gram determinant is a
nonzero formal relation at degree five, so no assertion about the full
invariant ring is intended.
\end{proof}

Use the abstract symmetric \(N\)-by-\(N\) Gram matrix \(G\) with
\(G\mathbf1=0\). Define \(x_i=G_{ii}+z\),
\(d_i=x_i-z\), \(D=\sum_i d_i\), and
\begin{equation}
 \begin{split}
 G_{ii}&=x_i-z,\\
 G_{ij}&=-\frac{d_i+d_j}{N-2}
       +\frac{D}{(N-1)(N-2)}+h_{ij},\qquad i\ne j,\\
 h_{ii}&=0,\qquad h=h^T,\qquad h\mathbf1=0.
 \end{split}
 \label{eq:v52-balanced}
\end{equation}
The channel space has dimension \(N(N-3)/2\). The variables \(x_i,z\)
and that channel space are independent linear coordinates; the splitting
commutes with every leg permutation. In particular the sextic polarized
identity has six Euler variables, nine channels, and one mass variable.

By Lemma~\ref{lem:v52-Gram-injective}, a physical even variational identity
lifts uniquely to the relevant degree in this polynomial algebra.
Changing the \(x_i\) while holding the channels fixed need not preserve
physical Gram rank. It is a legitimate coefficient operation after the
low-degree injection, not a deformation of the physical momenta.

\begin{theorem}[An all-leg, residual-preserving Euler section]
\label{thm:v52-section}
Write \(q=\sum_{m=0}^2q^{(m)}\), homogeneous in \(x\) with \(h,z\)
fixed. Define
\begin{equation}
 \boxed{A_{ij}(q)=\sum_{m=0}^2
 \frac{\partial_{x_j}q_i^{(m)}-\partial_{x_i}q_j^{(m)}}{m+1}.}
 \label{eq:v52-section}
\end{equation}
Then, for every supplied input,
\begin{equation}
 \boxed{q_i-\sum_jA_{ij}x_j
 =\sum_{m=0}^2\frac{\partial_{x_i}W(q^{(m)})}{m+1}.}
 \label{eq:v52-residual}
\end{equation}
The residual has exactly the original Euler defect. The section is
permutation equivariant and has bounds
\begin{equation}
 \boxed{\|A(q)\|_1\le\frac23\|q\|_1,
 \qquad\|q-A(q)x\|_1\le\|W(q)\|_1.}
 \label{eq:v52-section-bounds}
\end{equation}
The output counts only \(i<j\). The input norm is the coefficient
\(\ell^1\) quotient norm in the independent \(x,z\) variables and the
linearly presented channels. The channel relations do not involve \(x\).
\end{theorem}
\begin{proof}
The homogeneous calculation of V50 works with any number of Euler
variables:
\[
 \sum_jx_j\partial_{x_j}q_i^{(m)}=mq_i^{(m)},\qquad
 \sum_jx_j\partial_{x_i}q_j^{(m)}
   =\partial_{x_i}W(q^{(m)})-q_i^{(m)}.
\]
Subtracting proves \eqref{eq:v52-residual}. The skew part has zero
Euler contraction. A monomial of \(x\)-degree \(m\) contributes at most
\(m/(m+1)\) to the upper skew coefficient norm. The summed derivatives
of its defect have norm at most \(m+1\). Take the quotient norm on the
channels to obtain the same estimates. Permutation equivariance follows
from \eqref{eq:v52-balanced}. The new content is the validity and
orientation control for arbitrary leg count; the four-leg homogeneous
identity is inherited, not reproved as an independent discovery.
\end{proof}

There is also a direct bound in the redundant row-zero Gram coefficient
quotient norm. For \(N\ge4\), a Gram entry in \eqref{eq:v52-balanced}
has coordinate coefficient norm at most \(7/3\). Conversely a channel
entry has a representative of norm at most \(8/3\) in Gram entries,
and \(x_i\) has norm at most two. Thus the complete degree-two input to
degree-one skew map has bound
\begin{equation}
 \frac83\frac23\left(\frac73\right)^2
 =\frac{784}{81}<10.
 \label{eq:v52-Gram-bound}
\end{equation}
Choosing an independent physical jet basis and unpolarizing are separate
finite maps whose constants retain scalar-degree dependence.

Inverse Fourier transformation of \(A\), with its two marked Euler
slots and the other scalar slots retained, defines a local formally
skew operator \(\mathcal L\). Skew exchange of those two slots is
formal adjointness; the spectator permutations are already preserved.
Each entry has weight two, so derivatives on the operator argument
and on its coefficients have total weight at most two. Covariantize
with a fixed ordering and retain the formal skew part
\((\mathcal L^0-(\mathcal L^0)^\dagger)/2\). This has exactly the
constructed flat symbol. It is a polynomial local operator, not an
inverse of any Fourier multiplier.

\section{The complete invariant flat scalar-quintic coefficient bank}
\label{sec:v52-quintic}

The all-degree proof has a small, fully evaluated quintic specialization.
There are six external scalar legs in the polarized action identity.
Independent flavour signs leave the multiplicity patterns
\(6\), \(4+2\), and \(2+2+2\), with respectively four, twelve, and
four choices from the four scalar flavours.

For each root, the five input momenta have fifteen Gram entries and one
mass variable. We use orbit sums of invariant-degree-two monomials under
permutations of identical input flavours. The numbers of root coefficients
are ten for five identical inputs, 31 and 22 for the two roots of a
\(4+2\) channel, and 58 for each of the three roots in a \(2+2+2\)
channel. There is no Gram relation in the degree-three output by
Lemma~\ref{lem:v52-Gram-injective}.

\begin{theorem}[Exact quintic ranks and a finite coefficient compiler]
\label{thm:v52-quintic-ranks}
In these orbit-sum bases, the exact polarized Euler maps have
\begin{center}
\begin{tabular}{lrrrr}
\toprule
Channel & Copies & Input dimension & Kernel & Rank\\\midrule
\(6\) & 4 & 10 & 1 & 9\\
\(4+2\) & 12 & 53 & 11 & 42\\
\(2+2+2\) & 4 & 174 & 41 & 133\\\midrule
Complete bank & & 1372 & 300 & 1072\\\bottomrule
\end{tabular}
\end{center}
For each channel the exported rational matrices \(N,H\) obey
\begin{equation}
 MN=0,\qquad HN=I,\qquad
 \Pi=NH,\qquad M(I-\Pi)=M.
 \label{eq:v52-compiler}
\end{equation}
Here \(M\) is the determining map, \(N\) its kernel synthesis, and
\(H\) extracts the selected free coefficient positions. Their exact
norms are
\begin{center}
\begin{tabular}{lrrrr}
\toprule
Channel & \(\|H\|_1\) & \(\|N\|_1\) & \(\|NH\|_1\) & \(\|I-NH\|_1\)\\\midrule
\(6\)&1&2&2&1\\
\(4+2\)&1&31&31&30\\
\(2+2+2\)&1&80&80&79\\\bottomrule
\end{tabular}
\end{center}
\end{theorem}
\begin{proof}
The coefficient matrix is generated without momentum sampling. Set the
sixth momentum to minus the sum of the first five, expand the root orbit
sums, and multiply each by its actual Euler factor. The output has 816
possible invariant-degree-three monomials. The companion certificate
records rational null vectors with identity free-coordinate blocks and
nonzero maximal minors. With the stated lexicographic ordering their
determinants are respectively
\[
 -3840,\qquad 75497472,\qquad -281474976710656.
\]
The minors prove the lower ranks; the independent null vectors prove the
upper ranks. Every product in \eqref{eq:v52-compiler} and every induced
column norm is checked in rational arithmetic. Weighted sums over the
flavour multiplicities give the total counts. These are variational
coefficient dimensions, not current-cohomology or anomaly dimensions:
the general section supplies skew-Euler operators for the whole kernel.
\end{proof}

The compiler does not insert a counterterm. For an arbitrary array it
returns its selected variational component, its kernel coordinates, and
its unremoved array. Zero residual is equivalent to the original flat
identity in this declared bank. The universal section in
\Cref{thm:v52-section} subsequently constructs the local operator.

For a single scalar, the sole quintic direction can be written
\begin{equation}
 \begin{split}
 \mathcal L_\phi&=\phi^3\nabla^a\phi\,\nabla_a
       +\frac12\nabla_a(\phi^3\nabla^a\phi),\\
 Q_5&=\mathcal L_\phi P_g\phi,\qquad
 fQ_5=\frac12\nabla_a(\phi^3\nabla^a\phi\,f^2).
 \end{split}
 \label{eq:v52-self-example}
\end{equation}
The flat quintic result is not obtained by assuming that every
multiflavour coefficient is a copy of this one-scalar example.

\section{The curved all-degree remainder has an explicit Ricci section}
\label{sec:v52-curved}

Subtract the covariantly skew \(\mathcal L_{r-1}P_g\phi\) from an
input satisfying \eqref{eq:v52-input}. The remaining \(\overline Q_r\)
has zero flat restriction for every mass. Ordering covariant scalar
derivatives, including their commutators, puts it in
\begin{equation}
 \overline Q_{r,A}
 =R U_A+R^{ab}V_{A,ab}+(\nabla_aR)W_A^a+(\Delta R)X_A
 +|C|^2Z_A(\phi)+\langle C,{}^\star C\rangle\widetilde Z_A(\phi).
 \label{eq:v52-curved-inventory}
\end{equation}
All coefficients have homogeneous scalar degree \(r\). More explicitly,
\(U,V\) have at most two scalar derivatives or one further Ricci/mass
factor, \(W\) has one scalar derivative, and \(X,Z,\widetilde Z\)
have none. Their undifferentiated scalar coefficient polynomials have
respectively degrees \(r-1,r-2\), or \(r\), as required by the term.

\begin{lemma}[No new Weyl contraction from undifferentiated spectator fields]
\label{lem:v52-curved-inventory}
Equation \eqref{eq:v52-curved-inventory} is a spanning normal form for
every fixed \(r\) in the stated natural category. For a variational
input, \(Z=\widetilde Z=0\).
\end{lemma}
\begin{proof}
The differential contraction argument in V50 is independent of the
number of undifferentiated scalar factors. A single Weyl tensor leaves
only two available scalar derivative indices at weight four; its
remaining indices force a trace or an algebraic Bianchi contraction.
A differentiated curvature leaves at most one such index and reduces
by contracted Bianchi to \(\nabla R\); two curvature derivatives reduce
to \(\Delta R\) and curvature squares. The two independent Weyl-square
scalars are exactly those displayed. Extra undifferentiated scalars
change their flavour coefficient polynomials, not their tensor types.

On the two Ricci-flat witnesses already established in
Lemma~\ref{lem:v47-vacuum}, the complete identity becomes
\[
 \int\rho\,(P_g\phi)_A\,24x^{-6}
            (Z_A\pm\widetilde Z_A)(\phi)=0.
\]
The principal arbitrary scalar-Hessian coefficient of its scalar Euler
variation forces the flat flavour Killing equation
\(\partial_A Z_{\pm,B}+\partial_B Z_{\pm,A}=0\).
The elementary differentiated Killing identity gives
\(\partial_A\partial_B Z_{\pm,C}=0\). The maps are affine, and
homogeneity \(r>1\) therefore forces them to vanish. No scalar field
is put on shell. The metric witnesses and their curvature calculation
are inherited; the arbitrary-degree application is the present step.
\end{proof}

With the same trace reversal as V50, set
\begin{equation}
 \boxed{
 (\mathcal K_rT)_A
 =\mathfrak r(T)U_A+\mathfrak r^{ab}(T)V_{A,ab}
   +\nabla_a\mathfrak r(T)W_A^a+(\Delta\mathfrak r(T))X_A.
 }
 \label{eq:v52-K}
\end{equation}
Then \(\mathcal K_re=\overline Q_r\), and the full adjoint is
\begin{equation}
 \boxed{
 (\mathcal K_r^\dagger y)_{ab}
 =\frac1{c_E}\left[
 y_AV_{A,ab}-\frac12g_{ab}y_A\operatorname{tr}V_A
 -g_{ab}\{y_AU_A-\nabla_c(y_AW_A^c)+\Delta(y_AX_A)\}
 \right].}
 \label{eq:v52-Kadj}
\end{equation}
Derivatives act on the entire products. These trace and adjoint maps
are the inherited ones, now applied to the arbitrary-degree coefficients
whose completeness has just been proved. In the ordered product-slot
norm they have bound \(8|c_E|^{-1}\); expanding Leibniz factors permits
the conservative bound \(64(r+1)^2|c_E|^{-1}\). Conversion from another
normal tensor basis is a separate fixed finite map.

\begin{theorem}[Scalar section at every prescribed nonlinear degree]
\label{thm:v52-scalar-section}
Every natural mass-polynomial input \eqref{eq:v52-input}, at odd
\(r\ge3\), has the decomposition \eqref{eq:v52-target}. The flat
section, covariant skew completion, and \eqref{eq:v52-K} give finite
local coefficient maps, equivariant under the independent flavour signs.
They remain regular at zero scalar field, zero mass, and zero external
momentum. Their physical-jet constants retain scalar-degree and
nondegenerate metric-chart dependence.
\end{theorem}
\begin{proof}
The flat metric sets \(e=0\). \Cref{thm:v52-odd} removes the alternating
sector, and \Cref{thm:v52-section} constructs its even skew-Euler
operator. Its covariant skew representative has identically zero
integrated scalar variation for every metric. Subtracting it preserves
\eqref{eq:v52-input}. The preceding curved lemma and trace reversal
then produce the remaining term. Every operation is polynomial and
finite at fixed \(r\).
\end{proof}

\section{Exact current completion and finite-order elimination}
\label{sec:v52-BV}

Let \(\tau\) be the total scalar stress, and retain
\(\pi=\rho^{-1}\upsilon\) and \(\vartheta=\rho^{-1}\eta_g\).
For \eqref{eq:v52-target}, define
\begin{equation}
 \boxed{\mathcal T_{Q_r}
 =-\frac12\int\rho\,\pi^T\mathcal L_{r-1}\pi
       -\int\rho\,\pi^T\mathcal K_r\vartheta.}
 \label{eq:v52-current-primitive}
\end{equation}
\begin{theorem}[Complete nonlinear partners at the quintic and higher steps]
\label{thm:v52-complete}
The ordinary variation of \eqref{eq:v52-current-primitive} is the full
current with
\begin{equation}
 \boxed{\delta g=-\mathcal K_r^\dagger f,
 \qquad\delta\phi=\mathcal L_{r-1}f+\mathcal K_r(e+\tau)
                      =Q_r+\mathcal K_r\tau.}
 \label{eq:v52-full-current}
\end{equation}
It retains metric degree \(r+1\) and scalar degrees \(r,r+2\).
The primitive has ghost number minus two and the inherited BV weight
six. The input characteristic still has differential weight four.
\end{theorem}
\begin{proof}
The longitudinal variation of the complete natural density is a
divergence. The complementary Euler variation acts on the two odd
antifields and gives, with the inherited signs,
\[
 -\int\rho f^T\mathcal K_r\vartheta
 +\int\rho\pi^T\mathcal K_r(e+\tau)
 +\int\rho\pi^T\mathcal L_{r-1}f.
\]
The actual adjoint gives \eqref{eq:v52-full-current}. Equivalently the
complete action variation is
\[
 -(e+\tau):\mathcal K_r^\dagger f
 +f^T\mathcal L_{r-1}f+f^T\mathcal K_r(e+\tau),
\]
whose integral vanishes. This is the general-degree use of the V50
identity; no scalar or metric equation has been dropped.
\end{proof}

For example, \(Q_5=zR\phi^5\) has \(\mathcal L=0\) and
\(\mathcal K_5T=z\phi^5\mathfrak r(T)\). Its exact completion is
\begin{equation}
 \boxed{\delta g_{ab}=\frac z{c_E}g_{ab}\phi^5P_g\phi,
 \qquad\delta\phi=zR\phi^5-
 \frac z{c_E}\phi^5\left(\frac12|\nabla\phi|^2+z\phi^2\right).}
 \label{eq:v52-example}
\end{equation}
At flat metric and constant \(\phi=v\), the retained metric and
scalar-septic terms contribute \(z^3v^8/c_E\) and
\(-z^3v^8/c_E\). This example checks the new eighth-matter bookkeeping;
it is not a new loop coefficient.

Apply the theorem to V51's actual remaining \(Q_5\). Subtracting the
whole current first gives
\[
 Q_5'=0,\qquad F_6'=F_6+\mathcal K_5^\dagger f,
 \qquad Q_7'=Q_7-\mathcal K_5\tau.
\]
The identity at that order is \(\int\rho e:F_6'=0\). The inherited
all-degree V51 section therefore writes
\(F_6'=\mathcal L_{\xi_6}g+\mathscr A_6e\), with
\(\mathscr A_6^\dagger=-\mathscr A_6\). Its complete subtraction yields
\begin{equation}
 \boxed{F_6^{\mathrm{new}}=Q_5^{\mathrm{new}}=0,\qquad
 F_8^{\mathrm{new}}=F_8-\mathscr A_6\tau,\qquad
 Q_7^{\mathrm{new}}=Q_7-\mathcal K_5\tau-\mathcal L_{\xi_6}\phi.}
 \label{eq:v52-next-update}
\end{equation}
Both new scalar-septic terms are retained. The remaining equation at
that order is \(\int\rho(e:F_8^{\mathrm{new}}+
 f^TQ_7^{\mathrm{new}})=0\), to which the same scalar and metric
sections apply, rather than a new classification theorem.

\begin{theorem}[Finite-order elimination of the natural higher-matter tail]
\label{thm:v52-finite-tail}
Assume a natural, flavour-invariant, antifield-linear ordinary current
of characteristic weight four satisfies the complete coupled Noether
identity. Suppose its previously treated lower scalar components have
been removed, so its metric component starts at scalar degree six and
its scalar component at degree five. Then for every prescribed even
matter order \(M\), there is a finite polynomial local
\(\mathcal T_{\le M}\), of ghost number minus two, such that
\begin{equation}
 \boxed{\operatorname{pr}_{\mathrm{matter}\le M}
     (\mathcal Q-s\mathcal T_{\le M})=0.}
 \label{eq:v52-finite-tail}
\end{equation}
Matter count assigns one to \(\phi\) and \(\upsilon\), and zero to
the gravitational variables. The choices can be made consistently as
\(M\) increases. Consequently this current is exact in the formal
scalar-degree completion of the stated local category.
\end{theorem}
\begin{proof}
At the first remaining even matter count \(n=r+1\), its identity is
\eqref{eq:v52-input}. Subtract \eqref{eq:v52-current-primitive}. The
remaining leading metric characteristic satisfies the metric-only
identity and is removed by the full V51 primitive. The changes at
matter count \(n+2\) are exactly those in
\eqref{eq:v52-next-update}, with the indices replaced appropriately.
No smaller matter coefficient changes. Iteration through \(M\) is
finite, and each coefficient is a finite local polynomial. Fixing the
sections fixes the lower coefficients of the primitives consistently.
The resulting coefficientwise series has the stated formal exactness.
The proof consumes V51's all-degree metric theorem as an inherited
input; the new scalar theorem is what closes this iteration.
\end{proof}

This theorem does not assert a finite polynomial primitive for the whole
untruncated current, nor convergence of the formal primitive as an
analytic function of the scalar fields. It also does not identify an
arbitrary quantum descent with this natural antifield-linear class.
The scalar-linear and scalar-independent-metric sectors are not swept
into \eqref{eq:v52-finite-tail}.

\section{Coefficient bounds, tests, and the remaining quantum interface}
\label{sec:v52-interface}

The Euler factors used by the flat section are genuine polynomials;
none is inverted. The new all-leg section is bounded by \(2/3\) in its declared
adapted norm, and by \(784/81\) in the redundant Gram quotient norm.
The quintic raw compilers have the separately evaluated norms in
\Cref{thm:v52-quintic-ranks}. Covariantization, unpolarization, and the
Ricci trace maps cost finite constants at every fixed scalar degree.
The combined tail algorithm composes those maps with the degree-dependent
V51 metric section. It has a finite coefficient bound at every fixed
matter order, independent of the lattice cutoff and volume; no bound
uniform in unbounded matter order is asserted.

On a fixed nondegenerate metric-jet chart with \(|c_E|\ge c_*>0\), the
same local polynomials have finite physical jet bounds. A fixed smooth
external momentum window gives rooted source kernels with every fixed
moment and no total-volume multiplier, by integration by parts in the
relative momenta. The constants retain their source degree, window,
metric-jet, and coupling-margin dependence. The present statements are
coefficient-map estimates, not estimates for a new measured loop or
many-shell trajectory.

The companion calculation tests the full quintic raw coefficient maps,
the rational maximal minors and null bases, the residual homotopy on
every sextic labelled degree-two monomial, the finite orientation
projector, the explicit failed orientation-odd example, formal-adjoint
identities, and the nonlinear stress return. It separately audits
byte-for-byte preservation of the previous body. These finite checks
verify the stated identities and matrices; they do not independently
formalize the all-degree proof or the inherited V51 classification.

The remaining programme issue is now different from another scalar
matter-degree census. The natural higher-matter current tail has the
finite-order primitive just constructed. To use that fact in the native
critical quantum equation one still needs the actual invariant descent,
its representative category (including possible non-natural coefficients
and scalar-independent metric transformations), and its ghost-number-zero
local normalization. The distinction between variational current charges
and quantum counterterms is the usual BV distinction
\cite{V20BBH1994}. No primitive of ghost number minus two is installed
as an action counterterm merely by renaming it.

\section{Status, inherited results, and verification scope}
\label{sec:v52-ledgers}

\subsection*{Proved}
\addcontentsline{toc}{subsection}{V52: Proved}
The actual conserved-momentum Euler factors form a regular sequence for
arbitrary leg count at least three. Its weight-four orientation-odd
variational sector vanishes. Stable low Gram degree and the equivariant
Euler/channel splitting give a bounded all-leg scalar section. The
sign-invariant quintic bank has dimensions \(1372,300,1072\), with
explicit rational compilers and rank certificates. The scalar section
extends to every fixed nonlinear degree on curved metrics. Together
with V51 it supplies finite-order elimination of the natural
higher-matter current tail, including every induced stress and scalar
partner.

\subsection*{Inherited}
\addcontentsline{toc}{subsection}{V52: Inherited}
V1--V51 are unchanged. V50 supplies the four-leg homotopy mechanism,
normal curvature contractions, signed adjoint and current completion;
the present proof establishes their all-degree applicability rather
than treating the cubic calculation as new. The Ricci-flat witnesses
are V47's. The complete polynomial metric-only section and its
nonlinear source completion are V51's. The fixed gravity prescription,
V34's active subcritical scalar normalization, and V38's logarithmic
packet remain unchanged.

\subsection*{Literature input}
\addcontentsline{toc}{subsection}{V52: Literature input}
The implication from a regular sequence to its Koszul resolution is
standard \cite{V50Stacks}; the first-syzygy proof and the bounded
low-degree section used here are written out. The interpretation of
Euler-trivial and gauge characteristics in negative-ghost BV degree
is established methodology \cite{V20BBH1994}. No external vanishing
theorem for the full coupled invariant critical cohomology is assumed.

\subsection*{Conditional}
\addcontentsline{toc}{subsection}{V52: Conditional}
The finite-order tail theorem applies to a current satisfying the
complete ordinary identity in the declared natural category, after
the lower components have been treated as stated. Applying it to a
native quantum descent requires those hypotheses to be established
for that descent. Formal all-degree exactness is not analytic
convergence, and fixed-order coefficient bounds do not imply an
invariant quantum RG class.

\subsection*{Open}
\addcontentsline{toc}{subsection}{V52: Open}
The native invariant critical ghost-number-one breaking still needs
its full translation-invariant, mass-polynomial descent and a
compatible ghost-number-zero primitive. Naturality of general reference
components, scalar-independent metric transformations, and any missing
higher-antifield representatives remain separate. The internal gauge,
chiral-line, higher-loop, and invariant measured ESM-class interfaces
are not closed by this ordinary current calculation.

\begin{center}\small
\begin{tabular}{p{.24\textwidth}p{.68\textwidth}}
\toprule
Variable & Scope in V52\\\midrule
Theory & Same ordinary native Einstein action and four equal free real
scalars; no change of finite carrier or active normalization.\\
Derivative weight & Four for characteristics; finite derivative degree
throughout the current primitives.\\
Scalar/source order & Exact at each prescribed finite degree; no
uniform bound or convergence claim as that degree tends to infinity.\\
Mass & Polynomial in \(z=m^2\), regular at \(z=0\). No inverse mass and
no removal of the quantum lower reference scale.\\
Cutoff and volume & New finite coefficient-map bounds are independent
of both; no new quantum cutoff rate is claimed.\\
Verification & Exact polynomial and rational matrix checks, plus
preservation and PDF compilation; no native loop quadrature or
proof-assistant formalization.\\\bottomrule
\end{tabular}
\end{center}

\section*{Append-only continuation marker: V52}
\addcontentsline{toc}{section}{Append-only continuation marker: V52}
\label{sec:v52-continuation-marker}
\textbf{End of V52.} Preserve Sections 1--407. The scalar-quintic step
and the subsequent natural higher-matter tail have finite-order local
current primitives, using the inherited all-degree metric section.
Do not resume scalar degree-by-degree census where this theorem already
applies. Next address the representative/descent interface for the
native invariant critical breaking, with scalar-independent metric and
non-natural sectors retained. No active normalization is changed.


\clearpage
\part*{V53: Translation-compatible ascent of exact current packets}
\addcontentsline{toc}{part}{V53: Translation-compatible ascent of exact current packets}

\section{The next interface is a density descent, not another current census}
\label{sec:v53-start}

Sections 1--407 and the V52 closing marker are retained verbatim. The
ordinary Einstein--scalar action, its four equal real scalars, signed
Euler normalization, coframe cotangent map, and active prescriptions are
unchanged. Internal gauge, Yukawa, scalar self-interaction, and
cosmological terms remain zero in this calculation. All assertions below
concern the ordinary zero-mesh local differential. They are not identities
for the filtered finite-lattice transformation.

V52 supplies a ghost-number-minus-two primitive for the natural
higher-matter current at every prescribed finite scalar degree. That
primitive is not itself an action counterterm. V45, on the other hand,
exhibits a translation-sensitive ghost-one descent from a current by
contracting it with a constant spacetime two-form. The new task is to
connect these two constructions at the level of their actual densities,
including the integration-by-parts current. The connection is explicit:
a matched exact current gives a ghost-number-zero local primitive of
its entire two-form descendant. It requires a boundary-current term,
and a change of the current representative requires a superpotential
term. Neither is obtained by merely changing the name or ghost number
of the V52 primitive.

A second issue is resolved here rather than assumed: every natural,
identically conserved vector current of positive polynomial scalar
degree has a natural, translation-invariant superpotential. An explicit
covariant differential-operator section constructs it. This allows the
ascent to work with the given Noether current, rather than only with a
specially chosen representative. No spacetime-coordinate primitive is
introduced.

The interpretation of negative-ghost cohomology and the distinction
between a charge density and its current are standard BV facts
\cite{V20BBH1994,V20BBH1995}. V45's category counterexample and V52's
current reduction are inherited. The new results are the density-level
ascent, its residual formula, the natural positive-scalar-degree
superpotential, and the applications below. No general vanishing theorem
for the native critical quantum cohomology is invoked.

\subsection{Fixed differential and tensor-density conventions}

Write the ordinary minimal differential as
\[
 s=\delta+\gamma,\qquad
 \delta^2=\gamma^2=0,\qquad
 \delta\gamma+\gamma\delta=0.
\]
Here \(\delta\) is the Euler/adjoint-generator differential, not a
variation of the physical fields. On the existing normalized antifields,
\[
 \pi=\rho^{-1}\upsilon,\quad
 \vartheta^{ab}=\rho^{-1}\eta_g^{ab},\quad
 \delta\pi=f=P_g\phi,\quad
 \delta\vartheta=e+\tau.
\]
The ghost antifield has its full native gauge-adjoint replacement.
The longitudinal differential acts tensorially, with
\(\gamma c^\mu=c^\nu\partial_\nu c^\mu\).
Products use the inherited left graded Leibniz rule.

Use the following ghost-independent natural tensor densities:
\begin{equation}
\begin{array}{c|c|c|c}
\text{symbol}&\text{tensor type}&\text{ghost number}&\text{parity}\\\hline
 t&\text{scalar density}&-2&0\\
 q&\text{scalar density}&-1&1\\
 k^\mu&\text{vector density}&-1&1\\
 j^\mu&\text{vector density}&0&0\\
 \ell^{\mu\nu}=-\ell^{\nu\mu}&\text{bivector density}&0&0
\end{array}
\label{eq:v53-types}
\end{equation}
Ghost independent permits ghost antifields inside \(t\), but no
explicit diffeomorphism ghosts. In the applications, \(t\) is
quadratic in field antifields or linear in the ghost antifield,
\(q,k\) are linear in field antifields, and \(j,\ell\) contain only
physical fields. In particular \(\delta\ell=0\).

For any scalar density \(a\), of either parity,
\(\gamma a=\partial_\mu(c^\mu a)\). In the fixed ordering the vector
and bivector rules are
\begin{align}
 \gamma k^\mu&=c^\nu\partial_\nu k^\mu
   +(\partial_\nu c^\nu)k^\mu-(\partial_\nu c^\mu)k^\nu,
 \label{eq:v53-vector-rule}\\
 \gamma\ell^{\mu\nu}&=c^\alpha\partial_\alpha\ell^{\mu\nu}
 +(\partial_\alpha c^\alpha)\ell^{\mu\nu}
 -(\partial_\alpha c^\mu)\ell^{\alpha\nu}
 -(\partial_\alpha c^\nu)\ell^{\mu\alpha}.
 \label{eq:v53-bivector-rule}
\end{align}
The same vector formula holds for \(j\). All divergences and primitives
in this installment have translation-invariant coefficients.

\section{The current primitive determines an actual Green current}
\label{sec:v53-green}

Suppose the known local charge primitive is \(\int t\), and choose a
rooted charge density \(q\). Retaining the integrations by parts gives
\begin{equation}
 \boxed{\delta t=q+\partial_\mu k^\mu.}
 \label{eq:v53-green-packet}
\end{equation}
The boundary current \(k\) is part of the data. At the integrated level
\(s\int t=\int q\), but that equality alone does not determine a
pointwise representative for the ascent.

\begin{lemma}[A natural finite Green compiler]
\label{lem:v53-green}
For every finite primitive of the V51--V52 type, ordered covariant
integration by parts constructs \(q,k\) in
\eqref{eq:v53-green-packet}. They are natural, translation invariant,
and polynomial in \(z=m^2\), with the complete Einstein--scalar Euler
expressions retained. At any fixed field/derivative order the
construction is a finite coefficient map.
\end{lemma}
\begin{proof}
Apply \(\delta\) to the actual primitive before moving derivatives.
Every resulting term has one remaining field antifield. For an ordered
word with that antifield denoted by \(a\), the elementary identity is
\begin{align}
 U\nabla_{i_1}\cdots\nabla_{i_r}a
 &-(-1)^r(\nabla_{i_r}\cdots\nabla_{i_1}U)a\notag\\
 &=\sum_{h=0}^{r-1}(-1)^h\nabla_{i_{h+1}}
 \left[(\nabla_{i_h}\cdots\nabla_{i_1}U)
       (\nabla_{i_{h+2}}\cdots\nabla_{i_r}a)\right].
 \label{eq:v53-ordered-green}
\end{align}
The absent factors for \(h=0\) or \(h=r-1\) are identities. Tensor
indices and the invariant contractions of the original monomial are
retained. This follows by successive Leibniz rules; no covariant
 derivatives are commuted. Multiplication by \(\rho\) converts the
covariant divergence to the vector-density divergence in
\eqref{eq:v53-green-packet}. Terms from a ghost antifield use the actual
first-order gauge-adjoint replacement in precisely the same way.

Consequently the rooted charge and its boundary are tensorial. Curvature
commutators that arise if one subsequently chooses a normal jet basis
remain in that basis conversion. In the ordered contraction norm one
input word gives one rooted term and at most \(r\) boundary terms.
Expanding the physical Euler tensors and changing to a fixed coframe
Taylor bank costs finite, explicitly computable coefficient maps. No
claim uniform in that bank's unbounded degree is made.
\end{proof}

Let \(j_{\rm can}^\mu=\delta k^\mu\). Nilpotence gives the exact
relations
\begin{equation}
 \boxed{\delta q=-\partial_\mu j_{\rm can}^\mu,
 \qquad\delta j_{\rm can}=0.}
 \label{eq:v53-canonical-current}
\end{equation}
This current includes all Euler and scalar-stress returns. Its
conservation on stationary fields is not used to drop its off-shell
value. For an alternative given physical Noether current \(j\) with
\(\delta q=-\partial j\), the difference
\begin{equation}
 \omega^\mu=j^\mu-\delta k^\mu,
 \qquad \partial_\mu\omega^\mu=0
 \label{eq:v53-current-difference}
\end{equation}
is identically conserved. The next section constructs the required
superpotential when this difference has positive scalar degree.

\section{A covariant superpotential without explicit coordinates}
\label{sec:v53-superpotential}

\begin{lemma}[Divergence-free scalar differential operators are curls]
\label{lem:v53-operator-curl}
Let \(L^\mu[\psi]\) be a finite-order natural linear differential
operator from arbitrary scalar test functions \(\psi_A\) to vector
fields, on the torsion-free metric chart. If
\(\nabla_\mu L^\mu[\psi]=0\) identically for every \(\psi\), there
is a finite natural operator \(K^{\nu\mu}[\psi]=-K^{\mu\nu}[\psi]\)
with
\begin{equation}
 \boxed{L^\mu[\psi]=\nabla_\nu K^{\nu\mu}[\psi].}
 \label{eq:v53-operator-curl}
\end{equation}
An input of test-derivative order \(R\ge1\) requires order at most
\(R-1\) in \(K\). An order-zero divergence-free operator is zero.
\end{lemma}
\begin{proof}
Use symmetrized covariant test derivatives, retaining all curvature
terms at lower test order. For a fixed scalar label, let
\(L_R^\mu(p)\) be the top symbol, homogeneous of degree \(R\).
Identical divergence implies \(p_\mu L_R^\mu=0\). Define
\begin{equation}
 B_R^{\nu\mu}(p)
 =\frac{\partial_{p_\nu}L_R^\mu(p)
        -\partial_{p_\mu}L_R^\nu(p)}{R+1}.
 \label{eq:v53-curl-symbol}
\end{equation}
Euler homogeneity gives, for arbitrary input symbol,
\begin{equation}
 p_\nu B_R^{\nu\mu}
 =L_R^\mu-\frac1{R+1}
       \partial_{p_\mu}(p_\nu L_R^\nu).
 \label{eq:v53-curl-residual}
\end{equation}
Thus covariantizing \(B_R\) supplies the highest symbol of the desired
superpotential. Subtract its full covariant divergence from \(L\).
The test order decreases. Its divergence remains identically zero:
\(\nabla_\mu\nabla_\nu K^{\nu\mu}=0\) for every antisymmetric
bivector, by the Ricci symmetry, or equivalently by the commuting
partial derivatives of the associated density. Iterate. At order zero
one has \(L^\mu=a_A^\mu\psi_A\); the independent first test jets
force \(a_A^\mu=0\). Every step uses only covariant algebra, rational
integers, differentiation of coefficients, and the given metric.
No momentum or coordinate is inverted or integrated.
\end{proof}

In the independent upper-bivector symbol norm,
\begin{equation}
 \|L_R\mapsto B_R\|_{1\to1}=\frac{R}{R+1}
 \qquad(\dim x\ge2).
 \label{eq:v53-curl-norm}
\end{equation}
Indeed an input component monomial contributes the sum of exponents in
the other momentum directions divided by \(R+1\); equality is reached
when its momentum monomial avoids that component's direction. This
bound applies to the principal-symbol step. Full lower-order recursion
also differentiates the coefficient tensors and has fixed finite jet
bounds, not the same norm independently of derivative order.

\begin{theorem}[Positive scalar degree removes the current-improvement gap]
\label{thm:v53-superpotential}
Let \(\omega^\mu(g,\phi,z)\) be a ghost- and antifield-free natural
vector density, polynomial and homogeneous of scalar degree \(n>0\),
with \(\partial_\mu\omega^\mu=0\) identically. There is an explicitly
constructed natural bivector density \(\ell^{\nu\mu}\) such that
\begin{equation}
 \boxed{\omega^\mu=\partial_\nu\ell^{\nu\mu}.}
 \label{eq:v53-improvement}
\end{equation}
It has the same scalar degree and mass polynomial degree, one less
differential weight, and no explicit coordinate dependence or negative
power of a field, momentum, or mass. The construction is equivariant
under the actual flavour-sign group.
\end{theorem}
\begin{proof}
Linearize only the scalar arguments of \(\rho^{-1}\omega\), holding
the metric fixed. This gives a natural differential operator
\[
 L^\mu[\psi]
 =D_\phi(\rho^{-1}\omega^\mu)[\psi]
\]
whose divergence is identically zero for arbitrary \(\psi\). Apply
Lemma~\ref{lem:v53-operator-curl} and evaluate its operator on
\(\psi=\phi\). Scalar homogeneity gives
\(L[\phi]=n\rho^{-1}\omega\). Hence
\[
 \ell^{\nu\mu}=\frac{\rho}{n}K^{\nu\mu}[\phi]
\]
has the stated property. The only new denominator is the positive
integer \(n\). A finite polynomial is treated degree by degree.
Covariant differentiation, the symbol section, scalar linearization,
and evaluation at \(\phi\) commute with the flavour action. We have
not used the scalar equations or a spacetime Poincare primitive.
\end{proof}

The positive scalar-degree hypothesis is essential to this argument.
It makes no assertion about a scalar-independent metric current or an
arbitrary non-natural representative. For the currents of weight five
below, the resulting superpotential has weight four. The coefficient
construction terminates after at most the scalar-jet order of the
current; its constants are independent of mesh and periods at fixed
jet and scalar order.

\section{An exact ghost-number-zero ascent, including improvements}
\label{sec:v53-ascent}

Let \(B_{\mu\nu}=-B_{\nu\mu}\) be any constant numerical two-form.
It is an external component label, not a dynamical field or an additional
geometric structure. Put
\begin{equation}
 f_B=\tfrac12B_{\mu\nu}c^\mu c^\nu,
 \qquad \alpha_\mu=B_{\mu\nu}c^\nu.
 \label{eq:v53-factors}
\end{equation}
As in V45,
\(sf_B=c^\mu\partial_\mu f_B\),
\(s\alpha_\mu=c^\nu\partial_\nu\alpha_\mu\), and
\(\partial_\mu f_B=-\alpha_\nu\partial_\mu c^\nu\).

\begin{theorem}[Translation-compatible ascent of an exact current]
\label{thm:v53-ascent}
Assume the tensor-density data in \eqref{eq:v53-types} obey
\begin{equation}
 \delta t=q+\partial_\mu k^\mu,
 \qquad j^\mu=\delta k^\mu+\partial_\nu\ell^{\nu\mu},
 \qquad\delta\ell=0.
 \label{eq:v53-packet}
\end{equation}
Define the given current descendant and its primitive by
\begin{align}
 \mathcal A_B&=\int(f_Bq+\alpha_\mu j^\mu),
 \label{eq:v53-target}\\
 \mathcal C_B
 &=\boxed{\int\left(f_Bt-\alpha_\mu k^\mu
                   -\tfrac12B_{\mu\nu}\ell^{\mu\nu}\right).}
 \label{eq:v53-counterterm}
\end{align}
Then \(\operatorname{gh}(\mathcal C_B)=0\) and
\begin{equation}
 \boxed{s\mathcal C_B=\mathcal A_B\pmod{d_x}.}
 \label{eq:v53-ascent}
\end{equation}
More precisely, if \(h_B\) denotes the density in
\eqref{eq:v53-counterterm}, then
\begin{equation}
 s h_B=f_Bq+\alpha_\mu j^\mu+
 \partial_\nu\left(
    c^\nu h_B+f_Bk^\nu-\alpha_\mu\ell^{\nu\mu}\right).
 \label{eq:v53-ascent-density}
\end{equation}
Every term in this equality has translation-invariant coefficients.
\end{theorem}
\begin{proof}
Keep the ghost factors to the left as displayed. The scalar-density
rule and \eqref{eq:v53-packet} give
\[
 s(f_Bt)=f_Bq-(\partial_\mu f_B)k^\mu
                  +\partial_\mu(c^\mu f_Bt+f_Bk^\mu).
\]
Since \(\alpha\) and \(k\) are odd, the vector-density rule gives
\[
 s(\alpha_\mu k^\mu)
 =-\alpha_\mu\delta k^\mu-(\partial_\mu f_B)k^\mu
                  +\partial_\nu(c^\nu\alpha_\mu k^\mu).
\]
Subtracting these identities cancels the derivative-of-ghost term.
Finally the bivector-density rule, antisymmetry, and \(\delta\ell=0\)
give
\begin{align*}
 s\left(-\tfrac12B_{\mu\nu}\ell^{\mu\nu}\right)
 ={}&\alpha_\mu\partial_\nu\ell^{\nu\mu}\\
 &+\partial_\nu\left(
 -\tfrac12c^\nu B_{\mu\lambda}\ell^{\mu\lambda}
 -\alpha_\mu\ell^{\nu\mu}\right).
\end{align*}
Adding proves the full density formula. No invariance of the finite
regulator or measure is used.
\end{proof}

The current relations imply \(\delta q=-\partial j\) and
\(\delta j=0\). Directly,
\begin{equation}
 s(f_Bq+\alpha j)
 =\partial_\mu\bigl(c^\mu(f_Bq+\alpha j)-f_Bj^\mu\bigr),
 \label{eq:v53-cocycle}
\end{equation}
in agreement with V45's cocycle construction. That construction is
inherited; the new statement gives its admissible primitive when its
charge is exact, including the current improvement rather than only the
double constant-ghost contraction.

\begin{corollary}[A current representative is not an extra obstruction here]
\label{cor:v53-any-current}
Suppose a natural physical Noether current \(j\) of positive scalar
degree satisfies \(\delta q=-\partial j\), and an actual natural
primitive gives \eqref{eq:v53-green-packet}. Then
Theorem~\ref{thm:v53-superpotential} constructs the \(\ell\) required
by Theorem~\ref{thm:v53-ascent}. Consequently the whole given
current descendant, not merely one canonical representative, has the
translation-invariant primitive \eqref{eq:v53-counterterm}.
\end{corollary}
\begin{proof}
The difference \(j-\delta k\) is identically conserved by
\(\delta^2=0\). It is physical, natural, and of positive scalar
degree, so the superpotential theorem applies. Its superpotential is
physical and therefore \(\delta\)-closed.
\end{proof}

\subsection{An explicit failure residual and coefficient bounds}

For arbitrary supplied data define
\[
 r_q=q-(\delta t-\partial k),\qquad
 r_j=j-(\delta k+\partial\ell).
\]
The same calculation returns the exact unremoved functional
\begin{equation}
 \boxed{\mathcal A_B-s\mathcal C_B
       =\int(f_Br_q+\alpha_\mu r_j^\mu)\pmod{d_x}.}
 \label{eq:v53-residual}
\end{equation}
It does not certify a failed packet. Likewise, for an arbitrary natural
\(q,j\), the closure defect itself is
\[
 s\mathcal A_B
 =\int\left[f_B(\delta q+\partial j)-\alpha_\mu\delta j^\mu\right]
 \pmod{d_x}.
\]

Use independent antisymmetric entries for \(B\) and \(\ell\), and
sum the coefficient norms of vector components. Multiplication in the
free Grassmann algebra gives
\begin{equation}
 \boxed{\|\mathcal C_B\|_{\mathrm{coeff},1}
 \le\|B\|_1\left(\|t\|_1+\|k\|_1+\|\ell\|_1\right),}
 \label{eq:v53-ascent-bound}
\end{equation}
with the same bound using \(r_q,r_j\) for the residual. On the six
unit two-form outputs taken together, the sharper column count is
\begin{equation}
 \sum_{\mu<\nu}\|\mathcal C_{B^{\mu\nu}}\|_1
 \le6\|t\|_1+3\|k\|_1+\|\ell\|_1.
 \label{eq:v53-six-bound}
\end{equation}
These counts are attained in the abstract independent-density basis.
Actual monomial coincidences can lower them, not increase them.

At the critical weight of the current construction,
\[
 \wt(t)=\wt(q)=6,\qquad
 \wt(k)=\wt(j)=5,\qquad \wt(\ell)=4.
\]
Every term of \(\mathcal C_B\) and \(\mathcal A_B\) has weight
four, since \(\wt(c)=-1\). Mass-polynomiality and regularity at
\(m=0\) are preserved. The map changes ghost number by its displayed
factors, not by renaming a charge as an action term.

\section{Two explicit native examples with nonzero boundary contacts}
\label{sec:v53-examples}

\subsection{The scalar skew operator used in V50 and V52}

For one scalar, let \(v^\mu\) be an even natural vector and
\[
 L_v=v^\mu\nabla_\mu+\tfrac12\nabla_\mu v^\mu,
 \qquad f=P_g\phi.
\]
The particular choices \(v=\phi\nabla\phi\) and
\(v=\phi^3\nabla\phi\) are the scalar-cubic and scalar-quintic
representatives already established. With \(\delta\pi=f\), set
\begin{align}
 t&=-\tfrac12\rho\,\pi L_v\pi
   =-\tfrac12\rho\,v^\mu\pi\partial_\mu\pi,\notag\\
 q&=\rho\,\pi L_vf,\qquad
 k^\mu=-\tfrac12\rho\,v^\mu f\pi,\qquad
 j^\mu=-\tfrac12\rho\,v^\mu f^2.
 \label{eq:v53-scalar-packet}
\end{align}
The zero-order part of the first expression vanishes by \(\pi^2=0\),
not by an integration prescription.

\begin{proposition}[Scalar descent with its Euler-square current]
\label{prop:v53-scalar-example}
The packet \eqref{eq:v53-scalar-packet} satisfies
\eqref{eq:v53-packet} with \(\ell=0\). Its counterterm and full
descendant are
\begin{align}
 \mathcal C_B={}&-\frac12\int\rho f_Bv^\mu\pi\partial_\mu\pi
           +\frac12\int\rho\alpha_\mu v^\mu f\pi,
 \label{eq:v53-scalar-counterterm}\\
 \mathcal A_B={}&\int\rho f_B\pi L_vf
           -\frac12\int\rho\alpha_\mu v^\mu f^2,
 \qquad s\mathcal C_B=\mathcal A_B\pmod{d_x}.
 \label{eq:v53-scalar-target}
\end{align}
\end{proposition}
\begin{proof}
The odd product rule gives
\(
 \delta t=-\tfrac12\rho v^\mu
       (f\partial_\mu\pi-\pi\partial_\mu f).
\)
One integration by parts yields precisely \(q+\partial k\).
Moreover \(\delta k=j\), and
\(\delta q=\rho fL_vf=\tfrac12\partial_\mu(\rho v^\mu f^2)
=-\partial j\). Apply Theorem~\ref{thm:v53-ascent}.
\end{proof}

The one-ghost/one-antifield term in
\eqref{eq:v53-scalar-counterterm} is an actual contact. Keeping only
\(f_Bt\) loses it and does not generate
\eqref{eq:v53-scalar-target}. Expanding
\(f=z\phi-\Delta\phi\) retains the mass, gradient, and Euler-square
terms as a polynomial. No scalar equation is imposed. The quintic choice
supplies the same ascent at positive matter count six, including the
metric dependence in the volume and covariant derivatives.

\subsection{A mixed metric--scalar primitive with a curvature derivative}

Use the V47 Ricci section for the scalar-linear operator
\(Q_g\phi=\nabla^\mu R\nabla_\mu\phi\). Put
\[
 \theta=\operatorname{tr}_g\vartheta,\qquad
 E=\operatorname{tr}_g(e+\tau),\qquad\delta\theta=E.
\]
The complete primitive, rooted density, and Green data are
\begin{align}
 t={}&\frac{\rho}{c_E}\pi\nabla^\mu\phi\nabla_\mu\theta,\notag\\
 q={}&-\frac{\rho}{c_E}\left[
  \theta\nabla_\mu(f\nabla^\mu\phi)
    +\pi\nabla^\mu\phi\nabla_\mu E\right],\notag\\
 k^\mu={}&\frac{\rho}{c_E}f\nabla^\mu\phi\,\theta,
 \qquad
 j^\mu=\frac{\rho}{c_E}f\nabla^\mu\phi\,E.
 \label{eq:v53-mixed-packet}
\end{align}
Direct variation gives \(\delta t=q+\partial k\),
\(\delta k=j\), and \(\delta q=-\partial j\). Thus its explicit
ghost-zero primitive is
\begin{equation}
 \boxed{\mathcal C_B^{\mathrm{mix}}
 =\frac1{c_E}\int\rho\left[
 f_B\pi\nabla^\mu\phi\nabla_\mu\theta
 -\alpha_\mu f\nabla^\mu\phi\,\theta\right].}
 \label{eq:v53-mixed-counterterm}
\end{equation}
Its target is \(\int(f_Bq+\alpha j)\), with the \emph{total}
Einstein--scalar Euler trace \(E\), not the vacuum trace alone. The
scalar-quadratic stress inside \(E\) supplies the higher-matter
contacts already required by V47. This example involves a derivative
of a metric antifield, so the Green term is not absent merely because
the original primitive was quadratic in antifields.

\section{Finite-order ascent of V52's higher-matter current tail}
\label{sec:v53-tail}

\begin{theorem}[Action-degree primitive for the natural current-descendant sector]
\label{thm:v53-tail}
Let \(\mathcal Q=\int q\) be a current satisfying the hypotheses of
V52's finite-tail theorem, and let \(j\) be its given natural physical
Noether current. Work at any prescribed finite positive matter order
\(M\), with the same formal coframe and mass-polynomial convention.
For each constant two-form \(B\), the corresponding descendant
\(\mathcal A_B=\int(f_Bq+\alpha j)\) admits a finite, natural-in-its-
current-data, translation-invariant ghost-number-zero primitive
\(\mathcal C_{B,\le M}\) such that
\begin{equation}
 \boxed{\operatorname{pr}_{\mathrm{matter}\le M}
     (\mathcal A_B-s\mathcal C_{B,\le M})=0.}
 \label{eq:v53-tail}
\end{equation}
Its coefficients are polynomial in \(z\), have critical weight four,
and use only the actual V52 current primitive and the two local
constructions proved here. The six sectors can be handled simultaneously.
\end{theorem}
\begin{proof}
Use the inherited finite primitive \(\mathcal T_{\le M}\) for
\(\mathcal Q\). Its chosen density has the type \(t\) above.
Apply the finite Green compiler to all its terms. The complete ordinary
differential cannot lower matter count: Euler replacement of a scalar
antifield preserves it; replacement of a gravitational antifield
preserves or raises it by two. Hence the retained coefficients satisfy
\eqref{eq:v53-green-packet}. The Noether-current difference at each
retained positive scalar degree is identically conserved. Construct its
superpotential using Theorem~\ref{thm:v53-superpotential}, degree by
degree, and use \eqref{eq:v53-counterterm}. The exact density identity
then gives \eqref{eq:v53-tail}. No lower matter coefficient is changed
by increasing \(M\) if the inherited sections and the superpotential
ordering are fixed consistently. All operations are local and finite
at every such order.
\end{proof}

This is a statement about the whole current-derived descendant,
including its current representative. It is stronger than observing
that its double constant-ghost contraction is an exact charge. It does
not assert that every native ghost-one cochain is of the form
\eqref{eq:v53-target}. In particular, components with differentiated
ghosts not belonging to this packet, higher ghost structures, and
non-natural current representatives have not been removed.

For an invariant current the inherited primitive may be averaged over
the sixteen flavour signs. The Green, curl, and ascent maps commute
with that action and the average has coefficient norm at most one.
Thus the resulting action primitive is invariant as well. V45's
nontrivial scalar-rotation currents are not contradictory examples:
they do not possess the charge primitive \(t\), and their nontrivial
flavour character is independently excluded from the current native
reference. No vanishing theorem for all remaining invariant classes
is inferred.

The current-ascent map does not by itself reset an active normalization.
If an actual breaking component is identified with
\(+\mathcal A_B\), its cancelling order-\(\hbar\) action insertion
is \(-\hbar\mathcal C_B\) in the inherited sign convention.
That insertion changes its complete action/source packet. For the
canonical representative \(\ell=0\), it vanishes at zero antifields;
for a given improved current, the required \(-B\ell/2\) term can be
an antifield-free local action contact. It is not discarded to keep
an earlier normalization artificially fixed. The positive-matter
construction has zero pure-gravity restriction. No native breaking
coefficient has been matched to this packet in this installment.

\section{Coefficient control and the remaining native descent test}
\label{sec:v53-interface}

The new principal curl section is bounded by \(R/(R+1)\). Its covariant
lower-order recursion terminates after at most \(R\) steps. At fixed
coframe order, scalar degree, and differential-jet order, all its
coefficient maps have finite constants. If the full covariant divergence
map in the chosen finite operator bank has norm at most \(D_*\), an
inductive bound for the accumulated curl is, for example,
\[
 \|K\|\le\sum_{h=0}^{R-1}(1+D_*)^h\|L\|,
\]
using the principal bound at most one and the same chosen compatible
norms. Scalar homogeneity then multiplies by \(1/n\), not a reciprocal
field. Native Euler substitution and jet conversion retain their own
finite constants and the nonzero \(c_E\) margin. These estimates do
not claim a new sharp raw coframe constant independent of the bank.

Combining them with \eqref{eq:v53-ascent-bound} and the inherited V52
section gives a finite coefficient bound for every prescribed matter
order. It is independent of lattice cutoff, physical periods, and shell
count because the newly constructed maps are ordinary local coefficient
maps. The actual input coefficients can still depend on the regulator.
No bound uniform in unbounded matter order or radius of convergence of
the formal primitive follows.

Fixed smooth external momentum windows give the corresponding rooted
multilinear kernel bounds, with every prescribed moment and no total-
volume multiplier. These follow by integration by parts in all relative
momenta of the finite differential polynomials. Their constants retain
the number of arguments, the window, coefficient norms, metric-jet
chart, and coupling margins. No new quantum loop or shell estimate is
asserted.

The remaining native task is now sharper. For the actual critical
reference/breaking, one must construct its translation-invariant descent
representative and identify its current-derived component with data
\((q,j)\) to which the ordinary current sections apply. In this
positive-matter natural component, an identically conserved current
improvement is no longer an extra unanswered hypothesis: the preceding
superpotential theorem supplies it. The selected current tail then
has the ghost-zero primitive just constructed. What remains includes
the scalar-independent metric sector, non-natural tensor components,
possible higher-antifield or higher-ghost descent pieces, and the
residual after subtracting this identified component.

The public coefficient utility accepts supplied \(t,k,\ell,B\)
arrays and produces the complete local polynomial
\eqref{eq:v53-counterterm}. It does not infer their density identities
or native population from their names. The verification suite checks
those identities on the universal tensor-density algebra and on the
explicit scalar and mixed primitives above. Such checks do not evaluate
a native loop coefficient or classify all quantum anomalies.

\section{Status and verification scope}
\label{sec:v53-ledgers}

\subsection*{Proved}
\addcontentsline{toc}{subsection}{V53: Proved}
An ordered covariant Green-current compiler retains the density data
of the existing current primitives. A natural divergence-free scalar
differential operator has an explicit local curl section, and a
positive-scalar-degree identically conserved physical current has a
natural superpotential without explicit coordinates. The full
translation-sensitive two-form descendant of a matched exact current
has the ghost-number-zero primitive
\(f_Bt-\alpha k-B\ell/2\), with its exact divergence and failed-input
residual. Combining these results with V52 supplies finite-order
action-degree primitives for its natural higher-matter current-derived
descent sector. The scalar and mixed native examples retain nonzero
Green, Euler-square, and stress contacts.

\subsection*{Inherited}
\addcontentsline{toc}{subsection}{V53: Inherited}
V1--V52 remain unchanged. The ordinary master differential, normalized
antifields and density convention are inherited. V45 supplies the
current-descendant construction and the distinction between explicit-
coordinate and translation-invariant primitives. V51--V52 supply the
natural current primitives and their finite-order bounds; no scalar or
metric characteristic census is redone. Fixed gravity, V34's active
subcritical scalar prescription, and V38's logarithmic packet are
unchanged.

\subsection*{Literature input}
\addcontentsline{toc}{subsection}{V53: Literature input}
The BV current/descent framework and interpretation of negative-ghost
cohomology are established methodology \cite{V20BBH1994,V20BBH1995}.
The density ascent and covariant positive-degree superpotential used
here are proved directly. No external theorem is used to remove the
translation-invariance condition or to declare the complete native
invariant critical cohomology zero.

\subsection*{Conditional}
\addcontentsline{toc}{subsection}{V53: Conditional}
The application to V52 consumes that theorem's natural, antifield-linear
current hypotheses and its complete ordinary Noether identity. Applying
the result to the actual quantum breaking requires its representative
and current-derived decomposition. A supplied numerical coefficient
packet without those identities is not certified by the compiler.
Formal finite-order ascent is not convergence of an infinite action or
an exact identity for the finite filtered regulator.

\subsection*{Open}
\addcontentsline{toc}{subsection}{V53: Open}
The complete native invariant critical breaking still needs its
representative reduction and the remaining descent components in the
translation-invariant mass-polynomial category. Scalar-independent
metric symmetries, non-natural preferred-tensor coefficients, and
unaccounted higher-antifield/ghost pieces remain separate. The internal
gauge/chiral, higher-loop, and invariant measured ESM-class interfaces
are not supplied by this ordinary ascent calculation.

\begin{center}\small
\begin{tabular}{p{.24\textwidth}p{.68\textwidth}}
\toprule
Variable & Scope in V53\\\midrule
Theory & Same ordinary Einstein action and four equal free real scalars;
no active finite action, reference, or normalization changes.\\
Local category & Natural current data, constant two-form component
labels, translation-invariant densities and discarded currents; no
explicit-coordinate primitive.\\
Mass and derivatives & Polynomial in \(z=m^2\); weight-four action
and breaking descendants of weight-six charges. No inverse mass or
soft momentum.\\
Scalar/source degree & Exact at every prescribed finite degree using
V52; formal compatibility, not an analytic all-degree theorem.\\
Cutoff and volume & New local coefficient-map constants are independent
of both at fixed bank; no new quantum cutoff rate or shell theorem.\\
Verification & Exact Grassmann density identities, curl-symbol and
variable-operator sections, Green formulas, native example returns,
and preservation audit. No loop quadrature or proof-assistant claim.\\
\bottomrule
\end{tabular}
\end{center}

\section*{Append-only continuation marker: V53}
\addcontentsline{toc}{section}{Append-only continuation marker: V53}
\label{sec:v53-continuation-marker}
\textbf{End of V53.} Preserve Sections 1--415. The natural positive-matter
current-derived two-form descent now has an explicit translation-invariant
ghost-zero primitive, including its Green and superpotential terms.
Next identify that component in the actual invariant critical quantum
descent and treat the remaining representative, scalar-independent, and
higher-antifield/ghost sectors. Do not replace a full descent by its
double constant-ghost contraction or silently alter an active source
normalization.


\clearpage
\part*{V54: Finite-range realization of the complete local source jet}
\addcontentsline{toc}{part}{V54: Finite-range realization of the complete local source jet}

\section{Companion audit and the finite-carrier realization problem}
\label{sec:v54-start}

Sections 1--415, including the V53 continuation marker, are preserved.
The four newly supplied companions have been audited by their version
blocks, terminal proofs, and the earlier definitions used in those
proofs. This is not an independent verification of every theorem in
those cumulative manuscripts. The audit changes the immediate direction:
we address how a prescribed local counterterm/source jet is represented
on the actual finite carrier, rather than performing another census of
natural currents. The current-derived primitive of V53 remains available;
the decomposition of the entire native breaking into that and other
sectors remains open.

\begin{center}\small
\begin{tabular}{p{.24\textwidth}p{.68\textwidth}}
\toprule
Supplied source & Result read and precise transfer decision\\\midrule
YM V79 \cite{V54YM79} &
\srcid{r79:zero}, \srcid{r79:signs}, \srcid{r79:moment},
\srcid{r79:local}: the complete three-orientation contact is not scalar;
midpoint Fourier constants require a real lattice realization. Its
finite-range contact and second-moment method are used as the starting
interface, not proved again. The original compact-law nonlinear
remainder is still unpaid.\\[3pt]
SM source selection V43 \cite{V54SM43} &
\srcid{sel43:thm:renewal}, \srcid{sel43:thm:clock},
\srcid{sel43:thm:physical}: seven accepted events suffice for the
specified count-resolved covariant instrument; the independent original
clock preserves its group. Fine words still fail the fixed-label test.
Our compiler retains every supplied record index; it is not a positive
stopping policy or a proof of physical SM-source occurrence.\\[3pt]
Exact-action bridge V36 \cite{V54Bridge36} &
\srcid{v36:uniform-H}, \srcid{v36:growth-ceiling},
\srcid{v36:slow-real}: the constant-shift Schur return gives an additional
real canonical pair of order inverse amplitude at the fixed three-site
root. No small source sector may be silently removed. This is not a
positive Euclidean covariance or a quantum shell estimate.\\[3pt]
Native stationarity V46 \cite{V54Station46} &
\srcid{v46:sec:coupling}, \srcid{v46:thm:time},
\srcid{v46:prop:accuracy}: two neutral forcing channels are retained;
exact constraints permit original histories through a physical interval
of order amplitude squared. The branch is normally repelling and the
constants are not spatial-refinement uniform. Neither its existence
interval nor its coordinate estimates replace an ESM ultraviolet bound.\\\bottomrule
\end{tabular}
\end{center}

The letters denoting classical preparation amplitude in the last two
sources are unrelated to the lattice spacing \(a\) below. No time
rescaling or preparation result is transferred to the quantum theory.
The files and their SHA-256 digests are recorded in the verification
report. The review is based on the four mounted source files, not on
an identification of the unrelated older snippets with those files.

The YM result supplies a concrete warning. For a contact coefficient
of size \(a^{-4}\), a realization that matches only its value at zero
momentum can introduce a divergent second jet. A smooth Fourier
polynomial on a principal momentum chart is not, merely by its name,
a finite-range operator on a cell complex with differently located
components. We construct a rational, support-explicit section of the
entire prescribed jet, including its source and mass entries. We also
pay for divergent coefficients before claiming a negligible change.

Finite-difference differentiation and accuracy improvement are established
subjects \cite{V54Sadiq,V54Ramos}. No priority for interpolation or
Symanzik improvement is claimed. The contributions here are the
rooted multi-source and offset bookkeeping, a sparse total-degree
section, its use with the actual coefficient growth, and the distinction
between exact jet compatibility and an exact finite-lattice BV identity.
No new interacting loop is evaluated.

\section{Rooted source kernels on sites and shifted cells}
\label{sec:v54-carrier}

Use the same rectangular finite Fourier carrier as in the perturbative
lineage. Its periods satisfy \(L_\mu/a\in\mathbb N\). An argument of type
\(b\) is stored at \(x+a\rho_b\), with fixed \(\rho_b\in[0,1)^4\).
Site fields have \(\rho_b=0\); an edge component may have
\(\rho_b=e_i/2\). Nothing in this construction supplies a new field,
quantum operation, or statistical selection. It acts on already
admitted field/source coefficient types. For charged variables it is
not an exact nonlinear parallel-transport or internal-gauge completion.
The present gravity--scalar application keeps its existing zero internal
coupling convention.

A panel has \(A\ge2\) arguments, one chosen as the root. Put
\(m=4(A-1)\), collect the independent momenta into \(K\in\mathbb R^m\),
and set \(\eta_b=\rho_b-\rho_0\) for each nonroot argument. A finite
shift stencil has the rooted symbol
\begin{equation}
 \widehat T(K)=\sum_{n\in\mathbb Z^m}w_n
                  \exp\bigl(ia(n+\eta)\cdot K\bigr).
 \label{eq:v54-stencil}
\end{equation}
Only finitely many \(w_n\) are nonzero. It is the coefficient of an
actual sum of products of field/source values at the indicated cells,
with the original root sum and measure factor retained. Tensor entries
and ordered Grassmann arguments can be coefficients of \(w_n\).
At a single reciprocal-period shift,
\begin{equation}
 \widehat T(K+2\pi e_j/a)
       =e^{2\pi i\eta_j}\widehat T(K).
 \label{eq:v54-twist}
\end{equation}
This prescribed twist is not a discontinuity of the underlying
site-coordinate operator. Erasing it changes the carrier interpretation.

Let \(z=m_H^2\) denote squared Higgs mass; it is not one of the momentum
coordinates. In the coefficient convention \(\partial\leftrightarrow ik\),
a local polynomial is
\begin{equation}
 C_a(K,z)=\sum_{|\alpha|+2j\le D}
                  c_{\alpha j}(a)(iK)^\alpha z^j.
 \label{eq:v54-local-poly}
\end{equation}
Each matrix entry and each source label is retained separately. The
weighted Taylor map \(\mathbf T_D\) has momentum weight one and mass
weight two, as in V28. The coefficients in \eqref{eq:v54-local-poly}
are the prescribed ones; no value is obtained by fitting the desired
symmetry. A source-only coefficient is not removed by normalizing
\(Z[J]/Z[0]\).

\begin{proposition}[A zeroth-order contact is insufficient for a critical jet]
\label{prop:v54-low-order-failure}
For two distinct midpoint edge orientations, the V79 contact factor
\(\cos(ak_i/2)\cos(ak_j/2)\) has the right zero value and the right
boundary twists. If its coefficient is \(a^{-4}Z_{ij}\), its difference
from the intended constant contact has second jet
\begin{equation}
 -\frac{Z_{ij}}{8a^2}(k_i^2+k_j^2).
 \label{eq:v54-critical-failure}
\end{equation}
Thus its unassigned difference need not tend to zero on a fixed
physical source window. The bounded-coefficient V79 theorem is not
contradicted; the divergence scale has changed.
\end{proposition}
\begin{proof}
Expand the two cosines at zero and multiply by \(a^{-4}Z_{ij}\).
A nonzero displayed second coefficient cannot be part of a vanishing
fixed-window remainder. It must either be assigned as an additional
local coefficient or prevented by a higher-jet realization.
\end{proof}

\section{A sparse rational finite-range section of every prescribed jet}
\label{sec:v54-section}

For formal variables \(u=(u_1,\ldots,u_m)\), let
\([\cdot]_{\le n}\) mean ordinary total-degree truncation. For
\(|\alpha|\le n\) define the finite rational polynomial
\begin{equation}
 P_{n,\eta,\alpha}(u)=
 \left[
 \prod_{j=1}^m(1+u_j)^{-\eta_j}
 \prod_{j=1}^m\log(1+u_j)^{\alpha_j}
 \right]_{\le n}.
 \label{eq:v54-newton-poly}
\end{equation}
Only the finite binomial and logarithm coefficients through degree \(n\)
are used. For rational cell offsets every coefficient is rational. Set
\begin{equation}
 \boxed{
 \mathcal R_{a,n,\eta}[(iK)^\alpha]
 =a^{-|\alpha|}e^{ia\eta\cdot K}
 P_{n,\eta,\alpha}(e^{iaK_1}-1,\ldots,e^{iaK_m}-1).}
 \label{eq:v54-section}
\end{equation}
This is a finite formula, not evaluation of a logarithm across a
Brillouin boundary.

\begin{theorem}[Offset-aware total-degree jet realization]
\label{thm:v54-section}
For every \(|\alpha|\le n\), \eqref{eq:v54-section} is a stencil
\eqref{eq:v54-stencil} supported in
\(\{b\in\mathbb N^m:|b|\le n\}\). It satisfies
\begin{equation}
 \boxed{
 J_{\le n}\mathcal R_{a,n,\eta}[(iK)^\alpha]
        =(iK)^\alpha.}
 \label{eq:v54-jet}
\end{equation}
The number of possible rooted shifts is at most
\(\binom{m+n}{n}\), and their physical radius is at most
\(a(n+\|\eta\|_1)\). No bound depends on the number of lattice cells.
For \(|\eta_j|\le1\), the dimensionless shift weights obey
\begin{equation}
 \boxed{\sum_b|w_{n,\eta,\alpha}(b)|\le4^n2^m.}
 \label{eq:v54-weight-bound}
\end{equation}
The factor \(2^m\) may be omitted when every offset is zero.
\end{theorem}
\begin{proof}
Expand each monomial \(u^\gamma\) in
\eqref{eq:v54-newton-poly} as
\(\prod_j(e^{iaK_j}-1)^{\gamma_j}\). Its translations satisfy
\(0\le b_j\le\gamma_j\), hence \(|b|\le|\gamma|\le n\).
This proves finite support and \eqref{eq:v54-twist}.

In the formal germ at zero, put \(t_j=iaK_j\). Then
\(u_j=e^{t_j}-1\), \(\log(1+u_j)=t_j\), and
\(\prod_j(1+u_j)^{-\eta_j}=e^{-\eta\cdot t}\).
The substitution \(u=t+O(t^2)\) preserves the total-degree filtration.
The omitted terms therefore begin at degree \(n+1\). Multiplication
by \(e^{\eta\cdot t}\) proves \eqref{eq:v54-jet}.

For the coefficient bound, give a power series the absolute coefficient
norm at radius \(1/2\). For \(|\eta_j|\le1\), its binomial factor
has norm at most \(2^{|\eta_j|}\), using the rising-factorial bound
on the absolute binomial coefficients. The logarithm has norm
\(\log2<1\). Consequently the weighted coefficient norm of the
truncated product is at most \(2^m\), or one at zero offset.
Expansion of \(u^\gamma\) into translations costs \(2^{|\gamma|}\).
Relative to its radius-\(1/2\) coefficient weight this costs
\(4^{|\gamma|}\le4^n\). Summing proves
\eqref{eq:v54-weight-bound}. Tensor coefficients obey the same proof
entry by entry in any fixed submultiplicative coefficient norm.
\end{proof}

In one dimension the same weights are the derivatives at zero of the
Lagrange basis on the nodes \(\eta,\eta+1,\ldots,\eta+n\):
\[
 w_j^{(d)}=
 \left.\frac{d^d}{dx^d}
 \prod_{\substack{0\le r\le n\\r\ne j}}
       \frac{x-(\eta+r)}{j-r}\right|_{x=0}.
\]
This is a useful independent check. The multivariable construction
uses a total-degree simplex, not the much larger tensor-product box.
For co-located ESM fields, coordinates not differentiated by a monomial
have \(u_j\)-degree zero and are not averaged at all.

For a nonnegative integer padding \(q\), extend the section to
\eqref{eq:v54-local-poly} by
\begin{equation}
 \boxed{
 \mathcal R^{D,q}_{a,\eta}C_a
 =\sum_{|\alpha|+2j\le D}c_{\alpha j}(a)z^j
       \mathcal R_{a,D+q-2j,\eta}[(iK)^\alpha].}
 \label{eq:v54-weighted-section}
\end{equation}
It obeys \(\mathbf T_{D+q}\mathcal R^{D,q}C_a=C_a\), with zeros
in the extra prescribed jets. Mass differentiation is differentiation
of this same polynomial realization. It is not replaced by independently
realizing each mass derivative with a different stencil order.

\section{Source labels, symmetry, and the limits of finite-jet transport}
\label{sec:v54-sources}

The section acts diagonally on tensor, flavour, record, and ordered
source labels. Coefficient extraction of a labelled source monomial
therefore commutes with the construction. Derivatives of a common
functional retain all its source-only and mixed terms. They are not
independent subtractions of different observables.

\begin{proposition}[Exact graded and record-preserving realization]
\label{prop:v54-symmetry}
Suppose a prescribed jet satisfies a finite set of linear permutation,
conjugation, or reality identities belonging to its labelled source
bank. Averaging the realized physical kernels over those identities,
with the prescribed Grassmann signs, gives a finite-range section
satisfying them exactly and retaining the same weighted jet. A change
of root only enlarges the finite support by a fixed panel-dependent
factor. No scalar replacement of a tensor contact is made.

For a direct sum of retained record sectors the section leaves each
sector separate. If an independently verified clock transform has the
coefficient form
\(\widehat C(s)=h_0(s)\sum_j h(s)^{j-1}C_j\), as in source-selection
V43, the section commutes with that forward transform whenever the
sum converges in the stated coefficient norm. It does not assert
bounded inverse clock reconstruction or covariance of a finer record.
\end{proposition}
\begin{proof}
Each kernel permutation is the corresponding relabelling of actual
cell arguments; conjugation is the adjoint operation on the same
carrier. Such operations take finite stencils to finite stencils.
Their actions on momentum are linear substitutions after the root
momentum is eliminated, so they preserve Taylor degree. The prescribed
jet is fixed by the permitted averaging, hence remains unchanged.
The kernel \(\ell^1\) norm does not increase under signed isometric
permutations; fixed root-coordinate changes have fixed finite bounds.
No averaging over a symmetry that the target fails is allowed.

Record preservation follows because no matrix coefficient of
\(\mathcal R\) acts on that index. Bounded linearity at fixed panel
then permits the stated absolutely convergent forward sum to pass
through the section. This assertion concerns already extracted
coefficients, not the logarithm of a mixture of instruments.
\end{proof}

\begin{proposition}[Nonnegative averaging cannot reproduce a shifted second jet]
\label{prop:v54-positive-no-go}
A real nonnegative averaging kernel supported at half-integer
locations \(x\in\mathbb Z+1/2\), with total weight one, cannot
agree with the constant symbol one through degree two. In fact
\begin{equation}
 -\widehat w''(0)=\sum_xw_xx^2\ge\frac14.
 \label{eq:v54-positive-floor}
\end{equation}
Thus higher-jet midpoint contact realization requires signed weights.
\end{proposition}
\begin{proof}
Every allowed location has \(x^2\ge1/4\). Differentiate the Fourier
sum twice. A constant target has zero second derivative, contradicting
the displayed bound. This concerns nonnegative spatial averaging
weights, not positivity of an operator's Fourier symbol or of the
underlying unmodified physical measure.
\end{proof}

This distinction is exactly why the construction does not supply a
new completely positive operation or stopping law. Local counterterm
coefficients may be signed. Source-selection V43 does not license
identifying them with a positive combination of historical quantum
operations.

\begin{proposition}[Ordinary jet compatibility and an exact finite-defect ledger]
\label{prop:v54-BV}
On a fixed finite source/arity bank, let \(s_\infty\) be the inherited
ordinary local BV differential, with its complete mass grading. Realize
every component through the same required weight plus padding, and
retain all derivative and source partners. Then the realized symbol
has the original ordinary BV image at the prescribed jet order.
In particular, realization neither creates a local primitive nor
removes a nonzero local obstruction.

If \(s_a\) is the actual finite differential and \(C\) an existing
local functional, the full comparison retains
\begin{equation}
 s_a\mathcal RC-s_\infty C
 =s_\infty(\mathcal RC-C)
       +(s_a-s_\infty)\mathcal RC.
 \label{eq:v54-BV-ledger}
\end{equation}
The first term has the declared vanishing ordinary jet. The second
must be evaluated or bounded from the actual finite family; it is
not set to zero by this theorem.
\end{proposition}
\begin{proof}
In the ordinary local germ, the differential acts by finite sums of
polynomial momentum/mass multipliers and momentum substitutions, and
preserves the inherited total local weight. It cannot bring an omitted
higher-weight Taylor term into the retained bank. Applying the exact
jet identity componentwise proves the first statement. Equation
\eqref{eq:v54-BV-ledger} is an identity before taking any projection
or bound. It separates the realization error from the finite regulator's
already recorded symmetry defect.
\end{proof}

The section is not an algebra homomorphism on the full Brillouin zone.
For example, with zero offset, degree two, and \(u=e^{iak}-1\),
\[
 \mathcal R(ik)=a^{-1}(u-u^2/2),\qquad
 \mathcal R((ik)^2)=a^{-2}u^2,
\]
whereas \(\mathcal R(ik)^2=a^{-2}(u^2-u^3+u^4/4)\).
Only the retained jets agree. An exact lattice Leibniz rule, exact
finite BV nilpotence, or an interacting-shell identity cannot be
inferred by replacing all derivatives with this stencil.

\section{A sixth-order midpoint contact with explicit coefficients}
\label{sec:v54-midpoint}

The following realization is a compact symmetric alternative to the
general one-sided section. It directly strengthens the contact
interface supplied by YM V79 without re-evaluating its count law.
Define
\begin{equation}
 \boxed{
 b_4(t)=\frac{75}{64}\cos(t/2)
       -\frac{25}{128}\cos(3t/2)
       +\frac{3}{128}\cos(5t/2).}
 \label{eq:v54-mask}
\end{equation}
Its coefficient mass is \(89/64\), and
\begin{equation}
 b_4(t)=1-\frac5{1024}t^6+O(t^8),\qquad
 b_4(t+2\pi)=-b_4(t).
 \label{eq:v54-mask-series}
\end{equation}
For a real symmetric three-orientation contact matrix \(Z_a\), set
\begin{equation}
 (\widehat{\mathfrak C}^{[4]}_{Z_a})_{ii}= (Z_a)_{ii},\qquad
 (\widehat{\mathfrak C}^{[4]}_{Z_a})_{ij}
      =(Z_a)_{ij}b_4(ak_i)b_4(ak_j)\quad(i\ne j).
 \label{eq:v54-matrix-contact}
\end{equation}
Each off-diagonal entry is supported on the 36 midpoint displacements
\(\pm e_i/2,\pm3e_i/2,\pm5e_i/2\) combined with the analogous
six choices in direction \(j\). These are genuine edges of the
original translated orientation type. No new lattice sites are added.

\begin{theorem}[Critical matrix contact with a vanishing assigned error]
\label{thm:v54-midpoint}
The operator \eqref{eq:v54-matrix-contact} is real, self-adjoint, and
finite range. It has the correct midpoint boundary signs and the
constant matrix jet \(Z_a\) through degree five. For all real principal
momenta,
\begin{equation}
 \boxed{
 \|\widehat{\mathfrak C}^{[4]}_{Z_a}(k)-Z_a\|_{\op}
 \le\frac{725}{32768}\|Z_a\|_{\op}
                       a^6|k|^6.}
 \label{eq:v54-matrix-bound}
\end{equation}
Consequently \(Z_a=a^{-4}Z\) gives an \(O(a^2|k|^6)\) error,
not the divergent second coefficient in
\cref{prop:v54-low-order-failure}.
\end{theorem}
\begin{proof}
The three coefficients in \eqref{eq:v54-mask} solve the moment equations
\(\sum w_j=1\), \(\sum w_j(j+1/2)^2=0\), and
\(\sum w_j(j+1/2)^4=0\), for \(j=0,1,2\).
The signed sixth moment gives \eqref{eq:v54-mask-series}. Taylor's
real remainder bound for cosine and the absolute sixth moment give
\[
 |b_4(t)-1|\le\frac{725}{65536}|t|^6,
 \qquad |b_4(t)|\le\frac{89}{64}.
\]
Let \(A_6=725/65536\), \(B=89/64<2\), and \(t=ak\).
For a fixed row \(i\), the absolute off-diagonal sum of the error
is at most
\[
 \|Z_a\|_{\op}A_6
       \left(2|t_i|^6+B\sum_{j\ne i}|t_j|^6\right)
 \le2A_6\|Z_a\|_{\op}|t|^6.
\]
Symmetry makes this a bound on the spectral norm as well. The cosine
positions prove finite-range realization and the twists. Each
component is even, so the moment equations give degree-five matching.
\end{proof}

The negative coefficients are deliberate and satisfy
\cref{prop:v54-positive-no-go}. The construction retains the entire
matrix, including V79's distinct scalar and traceless orientation
contacts; it cannot convert that matrix into a scalar. It does not
supply an estimate for the compact YM nonlinear remainder, and its
orientation labels are not being identified with physical graviton
or gluon polarizations.

\section{Paying for divergent coefficients and controlling rooted kernels}
\label{sec:v54-bounds}

Let \(\mathcal R\) denote the section in
\eqref{eq:v54-weighted-section}, with permitted source symmetries
restored as in \cref{prop:v54-symmetry}. Fix the panel, offsets,
maximum jet degree, and a smooth compact external window. Set
\(\varepsilon=|K|+\sqrt z\) for \(z\ge0\).

\begin{lemma}[Differentiated realization error]
\label{lem:v54-error}
For every prescribed momentum derivative \(\beta\), a monomial with
\(|\alpha|\le n\) satisfies on a fixed sufficiently small
\(a|K|\) window
\begin{equation}
 |\partial_K^\beta(\mathcal R_{a,n,\eta}[(iK)^\alpha]
                         -(iK)^\alpha)|
 \le C_{n,m,\beta}a^{n+1-|\alpha|}
                   |K|^{n+1-|\beta|}
 \label{eq:v54-monomial-error}
\end{equation}
when \(|\beta|\le n+1\). For higher fixed derivatives the analogous
scaled-window seminorm bound holds, with the window radius in place
of the displayed power of \(|K|\).
\end{lemma}
\begin{proof}
Write the finite stencil as its sum of exponentials. The moment
identities in \cref{thm:v54-section} cancel the complete Taylor
polynomial through degree \(n\). Differentiate the same finite sum
and use Taylor's integral remainder through degree \(n-|\beta|\).
The total node norm is at most \(n+m\) and its absolute weight is
bounded by \eqref{eq:v54-weight-bound}. Each derivative supplies one
factor \(a\). This proves \eqref{eq:v54-monomial-error}, for example
on \(a|K|\le[2(n+m+1)]^{-1}\), with an explicit constant obtained
from those two bounds and the finite factorial. For a derivative above
\(n+1\), use the direct differentiated exponential bound and
\((a\mu)^{|\beta|-n-1}\le1\) in the scaled seminorm.
\end{proof}

\begin{theorem}[Finite-range realization with the full coefficient budget]
\label{thm:v54-budget}
Suppose the actual local coefficients satisfy, for fixed \(D,\delta\)
and fixed integer \(\ell\ge0\),
\begin{equation}
 \sum_{|\alpha|+2j\le D}
     a^{\delta-|\alpha|-2j}|c_{\alpha j}(a)|
 \le B\,L_a,
 \qquad L_a=1+|\log(a\lambda)|^\ell.
 \label{eq:v54-budget}
\end{equation}
Here \(\lambda>0\) is fixed and all coupling normalizations are
included in \(B\). Assume
\(a(|K|+m_H)\le[8A(D+q+m+2)]^{-1}\), so that every
permitted relabelling remains on the small symbol window. With
\(r=|\beta|+2s\le D+q+1\),
\begin{equation}
 \boxed{
 |\partial_K^\beta\partial_z^s(\mathcal RC_a-C_a)|
 \le C B L_a\,a^{D+q+1-\delta}
                   \varepsilon^{D+q+1-r}.}
 \label{eq:v54-budget-error}
\end{equation}
For a window with each external leg bounded by \(\mu\), assume
\(a(A\mu+m_H)\le[8A(D+q+m+2)]^{-1}\). Every fixed rooted
kernel moment of the difference is bounded by
\begin{equation}
 \boxed{
 \|\mathscr K^\Delta_{a,L}\|_{\mathrm{rooted},b}
 \le C_b B L_a\,a^{D+q+1-\delta}
                       (\mu+m_H)^{D+q+1}.}
 \label{eq:v54-rooted}
\end{equation}
There is no total-volume multiplier. The statement is uniform in
periods on \(\mu L_\nu\ge1\), provided the finite stencils can be
read without winding ambiguity (larger periods suffice).
\end{theorem}
\begin{proof}
For the mass coefficient indexed by \(j\), take
\(n_j=D+q-2j\). Equation \eqref{eq:v54-monomial-error} and
\(s\) mass derivatives give its error bounded by a constant times
\[
 |c_{\alpha j}|\,a^{D+q-2j+1-|\alpha|}
 |K|^{D+q-2j+1-|\beta|}z^{j-s}.
\]
If the indicated momentum derivative is above the first residual
order, use the direct bound as in the lemma. Insert the coefficient
budget; the common power of \(a\) is exactly
\(D+q+1-\delta\). Both cases are bounded by the displayed power
of \(\varepsilon\). This uses the finite polynomial mass Taylor
bank, not a new mass expansion of the propagators. Summation proves
\eqref{eq:v54-budget-error}.

For the rooted estimate, multiply the symbol by the fixed smooth
external window in its \(4(A-1)\) independent momenta. All sufficiently
many scaled derivatives obey the same bound; for orders above the
residual order use the higher-derivative part of the lemma. Rescale
these momenta by \(\mu\). Repeated integration by parts gives
arbitrary fixed polynomial decay of the inverse Fourier transform,
with its weighted \(L^1\) norm bounded by the scaled derivative
seminorms. The periodic discrete transform is the periodization of
this Schwartz kernel on the shifted relative lattice. Summing over
one rooted period and then over its periodizing translates is bounded
by the corresponding full-lattice sum. The uniform mesh Riemann-sum
bound holds for \(a\mu\) small and \(\mu L_\nu\ge1\).
Its measure is \(a^{4(A-1)}\), not the number of root cells.
This proves \eqref{eq:v54-rooted} and the one-\(L^1\), other-\(L^\infty\)
multilinear bound. The constants depend on the finite panel and chosen
window/moment order. No unbounded-arity estimate follows.
\end{proof}

For the usual one-loop local source bank \(D=\delta\), choose
\(q=1\). The error is then
\begin{equation}
 \boxed{
 C B a^2[1+|\log(a\lambda)|^\ell]
                (\mu+m_H)^{D+2}.}
 \label{eq:v54-one-loop-rate}
\end{equation}
It tends to zero at fixed lower scale, mass range, and source window.
It does not spoil an already proved \(O(a)\) coefficientwise limit.
An \(O(a^2)\) limit may acquire the displayed logarithm unless the
coefficient bound is sharper or an additional matching order is used.
If only \(q=0\) is used, a marginal logarithm gives an
\(a\log^\ell(a\lambda)\) error. Higher coefficient growth requires
padding according to \(D+q+1-\delta>0\); it is not erased by an
interpolation identity.

\section{Application to the existing normalization and source descent}
\label{sec:v54-application}

For an already computed coefficient write its normalization as
\(\gamma_a^R=\gamma_a-C_a\). The finite-range representative gives
\begin{equation}
 \widetilde\gamma_a^R=\gamma_a-\mathcal RC_a
 =\gamma_a^R+(C_a-\mathcal RC_a).
 \label{eq:v54-matching}
\end{equation}
Thus \cref{thm:v54-budget} is a complete scheme-comparison estimate
for this change of local realization. It uses the original coefficient
and its actual local jet; it does not claim that a missing loop or
cohomology coefficient has been populated. With \(D=\delta,q=1\),
an inherited cutoff error \(\epsilon_a\) becomes at most
\(\epsilon_a+C B a^2L_a(\mu+m_H)^{D+2}\).
For dyadic spacings and fixed bank, these additional errors are
summable. This is not an invariant-class theorem for the full shell.

The metric--Higgs, minimal-antifield, and current-packet components
already supplied in earlier versions can all be realized this way on
their specified coefficient banks. In particular V53's complete
\begin{equation}
 C_B=\int\left(f_Bt-\alpha_\mu k^\mu
                       -\frac12B_{\mu\nu}\ell^{\mu\nu}\right)
 \label{eq:v54-V53-packet}
\end{equation}
has a finite-range representative at each prescribed source degree.
All three terms and all source labels must be realized together.
The common jet still satisfies V53's ordinary ascent equation by
\cref{prop:v54-BV}. The same formula does not assert that a shifted
finite-field expression has an exact ordinary product rule, nor that
its variation under the actual \(s_a\) vanishes. The second term in
\eqref{eq:v54-BV-ledger} remains the regulator-specific test.

This realization rule is an available choice of finite local
representative with an explicit comparison theorem. No coefficient of
the active physical normalization is altered here, and no new primitive
is inserted for the still-unidentified critical breaking. A future
choice to use these stencils in a higher-loop microscopic action must
also retain their hard-leg insertions at that higher loop order.
The present one-loop effective-coefficient comparison is not such a
higher-loop estimate. The Haar factors, contour, retained zero modes,
and auxiliary determinant remain unchanged.

The four companion programmes have therefore supplied one direct
constructive transfer and three concrete restrictions on that transfer.
The complete local source contact is implemented rather than reduced
to a scalar; the record is not coarsened; a signed response or a
shrinking classical interval is not substituted for a quantum bound;
and actual polynomial coefficient growth is paid before taking the
cutoff limit.

\section{Status, audit evidence, and remaining work}
\label{sec:v54-ledgers}

\subsection*{Proved}
The new results are the sparse offset-aware finite-range section for
any prescribed multi-source/mass jet; exact retention of allowed
labelled identities; the nonnegative-half-offset obstruction; the
explicit sixth-order three-orientation contact; and the differentiated
and rooted matching estimate including divergent local coefficients.
The complete current-ascent packet inherits its ordinary local jets
under this realization. An actual finite BV defect has a separate
exact ledger and is not assumed to vanish.

\subsection*{Inherited}
The four reviewed source results are used only in the audit scopes
above. In particular, the YM matrix zero block and second-moment
calculation, the seven-event instrument and clock-injectivity theorem,
the classical inverse-amplitude pair, and the two-neutral-channel
constraint continuation are not presented as newly proved. V53's
ascent and all previous ordinary/quantum coefficient theorems retain
their original hypotheses. The body-preservation audit compares exact
bytes from the first numbered section to the previous terminator.

\subsection*{Conditional}
Applying the new comparison to a supplied coefficient requires its
actual power/logarithm budget. A bound at one loop cannot certify an
analytic expansion of the original compact law. Symmetry restoration
requires the tested finite differential and the whole source packet.
The abstract existence of a finite-range realization is not admission
of a positive historical operation, a new source-selection theorem,
or an all-orders many-shell theorem.

\subsection*{Open}
The native invariant critical breaking still requires the representative
and descent decomposition left open in V53. This installment does not
replace that problem by a proof of locality. The internal gauge/chiral
sector, higher-loop closure, an invariant ESM action class, and the
unpaid original compact YM remainder remain open. The useful next
step is to apply the source-complete local primitive to its actual
populated descent component, carrying the explicit realization and
finite-differential residual together. Another zeroth-order or scalar
contact approximation would not address that step.

\begin{center}\small
\begin{tabular}{p{.24\textwidth}p{.68\textwidth}}
\toprule
Variable & Scope of V54\\\midrule
Carrier & Existing translated sites/cells and full labels; no clock,
source instrument, field content, or measure is synthesized.\\
Locality & Finite support in lattice steps at prescribed jet/arity;
not uniformly finite support as EFT order tends to infinity.\\
Cutoff and volume & Section constants are independent of both at fixed
panel. Matching additionally uses the stated coefficient budget and
fixed physical window/lower scale.\\
Mass and loop order & Polynomial mass jets; source-coefficient
comparison, with a one-loop application. No new loop is integrated.\\
Symmetry & Exact labelled algebraic identities and ordinary retained
jets; finite BV and historical covariance are separate tests.\\
Verification & Exact moment, support, adjoint, mask, mass-jet, budget,
and preservation checks. Analytic rooted estimates are written proofs;
no full native descent array or interval-certified loop is computed.\\\bottomrule
\end{tabular}
\end{center}

\begingroup
\renewcommand{\refname}{References for the V54 source audit and comparison}
\begin{thebibliography}{9}\small
\bibitem{V54YM79} Renewal Geometry research programme,
\emph{Yang--Mills Mass Gap Research Notes V79}, supplied cumulative
LaTeX file; especially \srcid{r79:zero}, \srcid{r79:signs},
\srcid{r79:moment}, \srcid{r79:local}, and \srcid{r79:unpaid}.
\bibitem{V54SM43} Renewal Geometry research programme,
\emph{Historical Standard-Model Source Selection V43}, supplied
\srcid{Renewal_SM_Source_Selection_Research_Notes_V43(1).tex};
\srcid{sel43:thm:renewal}, \srcid{sel43:thm:clock},
\srcid{sel43:thm:physical}, and \srcid{sel43:status}.
\bibitem{V54Bridge36} Renewal Geometry research programme,
\emph{Exact-Action Einstein Bridge Research Notes V36}, supplied
\srcid{Renewal_Exact_Action_Einstein_Bridge_Research_Notes_V36(1).tex};
\srcid{v36:uniform-H}, \srcid{v36:growth-ceiling},
\srcid{v36:slow-real}, and \srcid{v36:direction}.
\bibitem{V54Station46} Renewal Geometry research programme,
\emph{Native Regenerative Stationarity Closure V46}, supplied
\srcid{Renewal_Native_Regenerative_Stationarity_Closure_Notes_V46(1).tex};
\srcid{v46:sec:coupling}, \srcid{v46:thm:time},
\srcid{v46:prop:accuracy}, and \srcid{v46:sec:status}.
\bibitem{V54Sadiq} B.~Sadiq and D.~Viswanath,
\emph{Finite Difference Weights, Spectral Differentiation, and
Superconvergence}, arXiv:1102.3203.
\url{https://arxiv.org/abs/1102.3203}.
Primary-source comparison for differentiation weights; none of the
present record, BV, or cutoff conclusions is imported from that work.
\bibitem{V54Ramos} A.~Ramos and S.~Sint,
\emph{Symanzik improvement of the gradient flow in lattice gauge
theories}, arXiv:1508.05552.
\url{https://arxiv.org/abs/1508.05552}.
Comparison for action and observable improvement; not an estimate
for the present ESM regulator.
\end{thebibliography}
\endgroup

\section*{Append-only continuation marker: V54}
\addcontentsline{toc}{section}{Append-only continuation marker: V54}
\label{sec:v54-continuation-marker}
\textbf{End of V54.} Preserve Sections 1--423. The complete prescribed
local source jet now has a finite-range carrier realization and a
cutoff comparison that pays for its divergent coefficients. This is
not the missing native critical-descent decomposition or a finite
nonlinear symmetry theorem. Keep all tensor, duration, source, and
measure labels, and carry the recorded finite-BV defect when using
the new realization in a measured coefficient.


\clearpage
\part*{V55: Finite BV transport and contraction-aware source realization}
\addcontentsline{toc}{part}{V55: Finite BV transport and contraction-aware source realization}

\section{The two contraction operations left separate in V54}
\label{sec:v55-start}

Sections 1--423, their mathematical text, and their continuation markers
are retained verbatim. V54 constructs finite-range representatives of
prescribed local jets. Its identity \eqref{eq:v54-BV-ledger} deliberately
leaves the finite differential to be tested. This installment separates
two operations which require different estimates: the contraction of a
counterterm with the \emph{tree} master functional, and a coincident
field--antifield trace. The former has no independent loop momentum; the
latter does. High Taylor order by itself controls the former and need
not control the latter.

The first result below bounds the \emph{change} in an actual soft finite
Slavnov row caused by using V54's representative. It does not set the
pre-existing finite breaking to zero. The second result constructs, on
the existing scalar/antifield carrier, a finite-range correction which
preserves the prescribed soft jet and fixes a supplied coincidence trace
exactly. Both are coefficient-level results. No absolute critical
primitive is manufactured and no active normalization is changed.

\subsection*{Inputs and scope}
The ordinary and finite master candidates are those of V18 and V28, with
V16's off-shell nonminimal convention, V24's recorded action/source jets,
and V27's simultaneous scalar completion. The finite source is not
assumed nilpotent. The statements cover fixed panels in the existing
pure-gravity and equal-real-scalar sectors, at zero internal gauge,
Yukawa, and scalar self-interaction order. Extending the field list to
the chiral Standard Model is a further task, not a consequence of a
finite stencil formula.

The V54 section and its coefficient budget are inherited, as are the
prior normalized one-loop rows on their stated domains. The new proofs
use their actual finite formulas; they do not invoke a continuum Ward
identity to estimate the finite differential. The distinction between a
cutoff-modified identity and an unmodified master equation is standard
\cite{V55IIM}; improvement of an action and of its source/contact terms
also has established precedents \cite{V55RS}. Neither reference supplies
the diagonal-restoring section proved here.

Write \(z=m_H^2\), and fix a finite, labelled external panel of maximum
arity \(A_*\). Every source, ghost, antifield, flavour, and record label
is retained. Let \(\mu>0\), \(m_*\ge0\), and
\(\epsilon_* = A_*\mu+m_*\). A scaled symbol seminorm is
\begin{equation}
 \|F\|_{r;\mu,m_*}=
 \max_{|\beta|+2s\le r}
 \mu^{|\beta|}\epsilon_*^{2s}
 \sup_{\substack{\text{each external leg }|p_i|\le\mu\\0\le z\le m_*^2}}
 \|\partial_K^\beta\partial_z^s F(K,z)\|.
 \label{eq:v55-seminorm}
\end{equation}
The rooted kernel norm has the same meaning as
\eqref{eq:v54-rooted}, with one factor of \(a^4\) per relative
four-coordinate and any fixed polynomial spatial moment. All constants
below may depend on the panel, derivative/moment order, fixed profile
seminorms, chart radius, and nonzero coupling margins.

\section{The native soft BV operation has no hidden loop momentum}
\label{sec:v55-soft}

Let \(s_a=(\mathbb S_{a,0},\cdot)_a\) be the ordinary canonical
antibracket with the recorded finite \emph{tree} master candidate.
Here the subscript on the bracket records the lattice volume
normalization, not a propagator. A monomial contraction between a
counterterm panel \(I\) and a tree vertex produces a panel \(J\).
Denote such an incidence by \(\sigma:I\to J\).

\begin{lemma}[Finite soft incidence formula]
\label{lem:v55-incidence}
For every fixed incidence, its momentum coefficient is a finite signed
sum of terms
\begin{equation}
 M_{a,\sigma}(K,z)\,F_I(L_\sigma K,z).
 \label{eq:v55-incidence}
\end{equation}
The linear routing \(L_\sigma\) has integer entries bounded in terms of
\(A_*\). The momentum at the contracted pair is a sum of external
momenta. There is no independent momentum sum and no total-volume
factor. On \(a\epsilon_*\) sufficiently small, the recorded native
multipliers satisfy, for each fixed \(r\),
\begin{equation}
 \|M_{a,\sigma}\|_{r;\mu,m_*}
 \le C_{\sigma,r}\epsilon_*^{\kappa_\sigma},
 \qquad \kappa_\sigma\in\{0,1,2\}.
 \label{eq:v55-multiplier}
\end{equation}
The choice \(\kappa_\sigma\) records the particular replacement:
nonminimal algebraic terms have degree zero, Lie/source terms degree
one, and action Euler terms degree two, including the mass.
\end{lemma}
\begin{proof}
Fourier transformation of the finite functional derivatives produces
one momentum delta from each vertex and one from the antibracket.
After imposing total external conservation, these determine the
contracted momentum. The factor \(a^{-4}\) in the canonical bracket
cancels one of the two action volume factors, leaving one rooted
functional normalization. This is a one-edge tree operation, not the
second-order trace considered below.

For the source multipliers, use \eqref{eq:v18-repaired-source} and
\eqref{eq:v28-scalar-source}. If every external subset sum lies in
\(|ap|<1/64\), each profile is on its plateau. The source is then the
actual ordinary Lie/chart polynomial, with one derivative, and the
nonminimal doublet has the stated algebraic part. Fixed jets of
\(DG(q)^{-1}\) have constants determined by the same pointwise chart;
no inverse wave operator appears.

For the action terms, use \eqref{eq:v16-offshell-action},
\eqref{eq:v17-generated-action}, and
\eqref{eq:v27-completed-Higgs}. The off-shell quadratic Euler row is
\(\widetilde K_a-\chi F^*F\), not the gauge-fixed matrix alone.
Fixed low-chart interaction jets have the two-derivative leading bound.
The native remainders \(a^3O(\epsilon_*^5)\) and
\(a^4O(\epsilon_*^6)\) also obey a constant times
\(\epsilon_*^2\) when \(a\epsilon_*<c\). The sine, hyperbolic-sine,
and finite translated-word formulas give the same bounds after each
prescribed differentiation. The massive scalar vertices include \(z\)
at degree two. Finite products and linear substitutions prove
\eqref{eq:v55-multiplier}. Only already supplied tree jets are consumed;
for a larger panel they must first be supplied at its required degree.
\end{proof}

This proof does not replace a finite difference by a derivation on
arbitrary high lattice modes. In \eqref{eq:v55-incidence}, a translated
counterterm is varied as an actual polynomial in finite field variables.
Its translations give exponential factors evaluated at the routed
\emph{soft} labels. The usual functional Leibniz rule holds in that
finite algebra even when a spatial difference has no continuum product
rule.

\section{Cutoff-faithful transport of an already normalized one-loop row}
\label{sec:v55-transport}

For each input panel let \(C_{a,I}\) satisfy V54's actual coefficient
budget \eqref{eq:v54-budget}, with parameters
\(D_I,\delta_I,B_I,L_a\). Let \(p_I\ge0\) be its padding and put
\(m_I=D_I+p_I+1\). Write
\(E_{a,I}=\mathcal R C_{a,I}-C_{a,I}\).

\begin{theorem}[Finite BV continuity of the realization error]
\label{thm:v55-transport}
On every fixed output panel \(J\),
\begin{equation}
 \boxed{
 \|(s_aE_a)_J\|_{r;\mu,m_*}
 \le C_rL_a\sum_{\sigma:I\to J}
 B_I a^{m_I-\delta_I}
 \epsilon_*^{m_I+\kappa_\sigma}.}
 \label{eq:v55-transport}
\end{equation}
For every fixed rooted kernel moment the same estimate holds, with a
possibly larger constant and sufficiently many symbol derivatives in
the proof. It is uniform in periods on \(\mu L_\nu\ge1\), with the
inherited nonwinding and soft-window restrictions. It contains no
unrestricted source-order or ultraviolet-mode assertion.
\end{theorem}
\begin{proof}
Apply \cref{lem:v55-incidence} to the \emph{difference} \(E_a\).
V54's monomial moment proof bounds it after every required fixed
momentum derivative and mass derivative. Linear pullback to
\(L_\sigma K\) has a panel-dependent finite norm. The product rule in
\eqref{eq:v55-incidence}, together with
\eqref{eq:v55-multiplier}, proves \eqref{eq:v55-transport}; finite
Grassmann signs do not affect this absolute estimate.

Multiply the complete output by the same smooth compact external
window as in V54. After rescaling the independent external momenta by
\(\mu\), fixed symbol derivatives give an integrable inverse transform
with every specified polynomial moment. Periodization on the actual
relative-coordinate torus gives the rooted lattice sum. The factors
from the transform cancel the factors \(a^{4(A_J-1)}\) in that sum.
This proves the source estimate without bounding an extensive action
in supremum norm. The argument requires no covariance or positivity.
\end{proof}

\begin{corollary}[Incremental modified-identity comparison]
\label{cor:v55-one-loop}
Suppose an already established first-order row uses the coefficient
\(\Gamma_{a,1}-C_a\), the fixed tree candidate, the actual lower-reference
insertion, and the same once-counted density convention. Replacing only
\(C_a\) by \(\mathcal RC_a\) changes the row by exactly
\begin{equation}
 \boxed{\widetilde{\mathcal W}_{a,1}-\mathcal W_{a,1}=-s_aE_a.}
 \label{eq:v55-row-change}
\end{equation}
Its additional error is bounded by \eqref{eq:v55-transport}. In
particular \(D_I=\delta_I\), \(p_I=1\) gives an
\(a^2[1+|\log(a\lambda)|^\ell]\) improvement error at fixed scales and
panel, with the displayed physical powers of \(\epsilon_*\).
\end{corollary}
\begin{proof}
The first-order canonical row is linear in its first-order coefficient.
The reference trace and its density operand at this order use only the
unchanged tree action and source. Their change under an order-\(\hbar\)
counterterm would be of the next loop order. Subtraction gives
\eqref{eq:v55-row-change}; apply the theorem. This also proves dyadic
summability of the added error at a fixed bank.
\end{proof}

The conclusion is incremental. The identity
\begin{equation}
 s_a\mathcal RC_a-s_\infty C_a
 =(s_aC_a-s_\infty C_a)+s_aE_a
 \label{eq:v55-incremental}
\end{equation}
retains the first term on the right. If it was never paid, this theorem
does not pay it by renaming it an interpolation error. The existing
current-ascent packet in \eqref{eq:v54-V53-packet} must likewise be
realized and varied as one complete functional.

\section{Arbitrarily accurate soft jets do not control coincidence traces}
\label{sec:v55-trace-defect}

For clarity first use the canonical flat coefficient BV operator on the
actual co-located scalar and scalar-antifield coordinates:
\begin{equation}
 \Delta_a^H=a^{-4}\sum_{x,A}
       \frac{\partial}{\partial\phi_A(x)}
       \frac{\partial^L}{\partial\upsilon_A(x)},
 \qquad \int_a=a^4\sum_x.
 \label{eq:v55-delta}
\end{equation}
Let \(\mathsf T_\mu\) translate one lattice step. For \(L\ge1\), set
\(c_L=\binom{2L}{L}\) and
\begin{equation}
 B_{L,a}=\frac1{4c_L}\sum_{\mu=0}^3
        (2-\mathsf T_\mu-\mathsf T_\mu^{-1})^L,
 \qquad
 b_L(t)=\frac1{4c_L}\sum_{\mu=0}^3(2-2\cos t_\mu)^L.
 \label{eq:v55-bump}
\end{equation}
Its on-site coefficient is one, while its soft Taylor jet vanishes
through degree \(2L-1\).

\begin{proposition}[A jet-null source with a nonzero BV contact]
\label{prop:v55-no-trace-bound}
Let \(J(x)\) be a retained scalar source and use one labelled ghost
component \(c^0(x)\). With one scalar pair define the even, ghost-zero
functional
\begin{equation}
 C_{a,L}=a^{-\delta}\int_a
       \upsilon(x)J(x)c^0(x)(B_{L,a}\phi)(x).
 \label{eq:v55-counterexample}
\end{equation}
It has a zero external Taylor jet through degree \(2L-1\), and its
fixed-window coefficient is \(O(a^{2L-\delta})\). Nevertheless,
\begin{equation}
 \boxed{\Delta_a^HC_{a,L}
       =a^{-\delta-4}\int_a Jc^0.}
 \label{eq:v55-trace-counterexample}
\end{equation}
The statements hold on periodic grids whose lengths in lattice steps
exceed \(2L\). Increasing \(L\) does not improve this trace.
\end{proposition}
\begin{proof}
The Fourier symbol is \(a^{-\delta}b_L(ak_\phi)\). Since
\(2-2\cos t=t^2+O(t^4)\), the stated Taylor and window properties
follow. In \eqref{eq:v55-delta}, the field and antifield derivatives
contract at the same site. Only the on-site coefficient of \(B_{L,a}\)
contributes; it is one. The factor \(a^{-4}\) in \(\Delta_a^H\)
converts the original \(\int_a\) to the stated density. Neither \(J\)
nor \(c^0\) is differentiated by this contraction.
\end{proof}

This is a diagnostic in the allowed component-source algebra, not a
claim that the native critical coefficient contains this particular
monomial or that it is a symmetry. The external \(J\) is retained to
make clear that the contact is not a source-independent vacuum term.
For example \(\delta=4,L=3\) gives a soft \(O(a^2)\) source change
and an \(a^{-8}\) contact. Thus no estimate of \(\Delta_aE_a\) follows
from the soft conclusion of \cref{thm:v55-transport} alone.

\section{A normalized finite-range contact section}
\label{sec:v55-bump}

\begin{lemma}[Exact weights and lower-reference normalization]
\label{lem:v55-bump}
The kernel of \(B_{L,a}\) has support on the \(1+8L\) axial points,
with no extra field sites, and
\begin{equation}
 \boxed{
 b_{L,0}=1,\quad
 J_{\le2L-1}b_L=0,\quad
 \|b_L\|_{\ell^1}=\frac{4^L}{\binom{2L}{L}}=:\gamma_L.}
 \label{eq:v55-bump-properties}
\end{equation}
Its symbol is real, even, hypercubic invariant, and nonnegative. For the
actual lower mask \(\nu_\lambda=1-\Theta_\lambda\), define
\begin{equation}
 \kappa_{L,a,\boldsymbol n}
 =\frac1{\prod_\mu n_\mu}\sum_{p\in\mathrm{BZ}_a}
           \nu_\lambda(p)b_L(ap),
 \qquad n_\mu=L_\mu/a.
 \label{eq:v55-kappa}
\end{equation}
If \(\Theta=0\), then \(\kappa=1\). If
\(0\le\Theta\le1\), \(\operatorname{supp}\Theta\subset\{|p|\le2\lambda\}\),
\(\lambda L_\mu\ge1\), \(a\lambda\le1/4\), and \(n_\mu>2L\), then
\begin{equation}
 \boxed{
 0\le1-\kappa\le
 \frac{4^{L+1}}{c_L}(a\lambda)^{2L+4},
 \qquad \frac12<\kappa\le1.}
 \label{eq:v55-kappa-bound}
\end{equation}
\end{lemma}
\begin{proof}
In one coordinate the coefficients of \((2-T-T^{-1})^L\) are
\((-1)^j\binom{2L}{L+j}\), \(-L\le j\le L\). Averaging four such
axes and dividing by \(c_L\) gives the on-site coefficient one and
\(\ell^1\) norm \(4^L/c_L\). Nonwinding discrete Fourier
orthogonality makes the full-zone average of \(b_L\) equal one.
The Taylor assertion and nonnegativity follow directly from the symbol.

There are at most \(16\lambda^4\prod L_\mu\) momentum labels in the
lower ball: in each direction use
\(1+2\lambda L_\mu/\pi\le2\lambda L_\mu\).
Moreover \(b_L(ap)\le (a|p|)^{2L}/(4c_L)\).
Multiplying this bound on \(|p|\le2\lambda\) by the normalized mode
count proves \eqref{eq:v55-kappa-bound}. At \(a\lambda\le1/4\) its
upper bound is smaller than \(1/2\). This proof uses the finite sum,
not thermodynamic replacement of it.
\end{proof}

The useful fixed choice \(L=3\) is especially small:
\begin{equation}
 \boxed{
 \begin{aligned}
 B_{3,a}={}&I-\frac3{16}\sum_\mu(\mathsf T_\mu+\mathsf T_\mu^{-1})
 +\frac3{40}\sum_\mu(\mathsf T_\mu^2+\mathsf T_\mu^{-2})\\
 &-\frac1{80}\sum_\mu(\mathsf T_\mu^3+\mathsf T_\mu^{-3}),
 \qquad \|B_3\|_{\ell^1}=\frac{16}{5}.
 \end{aligned}}
 \label{eq:v55-L3}
\end{equation}
It uses 25 points and
\(b_3(ak)=a^6\sum_\mu k_\mu^6/80+O(a^8|k|^8)\).
Its weights are signed. Nonnegativity of its Fourier symbol is not a
positive averaging or historical-operation theorem.

\section{Simultaneous source-jet and coincidence matching}
\label{sec:v55-repair}

Consider the scalar-linear transformation-source sector
\begin{equation}
 C[W]=\int_a\sum_{A,B,n}
       \upsilon_A(x)W_{AB,n}(x)\phi_B(x+an).
 \label{eq:v55-source-class}
\end{equation}
The odd coefficients \(W\) have ghost number one and are finite-range
polynomials in retained metric, ghost, and independent composite-source
labels. They contain neither a scalar nor a scalar antifield, and no
other antifields. This includes the linear-scalar transformation
counterterms with prescribed metric marks, not the full multiple-
scalar or multiple-antifield algebra.

For a declared translation-invariant mask define the scalar contraction
\begin{equation}
 \Delta_{\nu,a}^H
 =a^{-4}\sum_{x,y,A}\nu_{xy}
     \frac{\partial}{\partial\phi_A(y)}
     \frac{\partial^L}{\partial\upsilon_A(x)}.
 \label{eq:v55-masked-delta}
\end{equation}
At \(\nu=I\) this is \eqref{eq:v55-delta}. At \(\nu=1-\Theta\)
it is a specified masked coincidence operand. It is not silently
identified with the complete Legendre reference trace, nor combined
with an unmatched canonical bracket to assert a quantum master identity.

The actual contraction of \eqref{eq:v55-source-class} is
\begin{equation}
 \Delta_{\nu,a}^HC[W]=\int_a t_C,
 \qquad
 t_C(x)=a^{-4}\sum_n\nu_{x,x+an}\operatorname{tr}W_n(x).
 \label{eq:v55-diagonal}
\end{equation}
Let \(t_{\mathrm{tar}}\) be a \emph{supplied} finite-range, scalar-free,
antifield-free odd spectator density. It may be zero or a previously
specified finite representative of the required measure/source contact.
There is no rule here setting an unknown native contact to zero.
For \(r\) of this type, define
\begin{equation}
 \boxed{
 \mathcal H_{L,\nu}r
 =\frac{a^4}{4\kappa_{L,a,\boldsymbol n}}\int_a
          \sum_{A=1}^4\upsilon_A(x)r(x)(B_{L,a}\phi_A)(x).}
 \label{eq:v55-contact-section}
\end{equation}
The identity matrix divided by four preserves the flavour action. A
labelled single-flavour version uses no factor four.

\begin{theorem}[Exact contact-restoring section with an unchanged soft jet]
\label{thm:v55-contact}
Set
\begin{equation}
 \boxed{C^{\mathrm{tr}}
 =C[W]+\mathcal H_{L,\nu}(t_{\mathrm{tar}}-t_C).}
 \label{eq:v55-trace-repair}
\end{equation}
Then
\begin{equation}
 \boxed{
 \Delta_{\nu,a}^HC^{\mathrm{tr}}=\int_a t_{\mathrm{tar}},
 \qquad
 J_{\le2L-1}(C^{\mathrm{tr}}-C[W])=0.}
 \label{eq:v55-contact-identity}
\end{equation}
The correction is finite range, retains the full spectator labels, and
has the same Grassmann parity and ghost number as \(C[W]\). The map
at a fixed target is idempotent. At zero target it is a linear projection
onto the specified zero-contact subspace.
\end{theorem}
\begin{proof}
Only the displayed scalar pair is contracted. Its masked diagonal is
\(\kappa\), independently of the spectator location. The four flavour
terms cancel the factor four, and the factor \(a^4\) cancels the
\(a^{-4}\) in the contraction density. Thus
\(\Delta_{\nu,a}^H\mathcal H_{L,\nu}r=\int_ar\).
Substitution proves the first identity and idempotence. In every
coefficient of the correction, the scalar momentum supplies the factor
\(b_L(ak_\phi)\). It has no Taylor term below degree \(2L\); multiplying
by the spectator symbol or a nonnegative mass power cannot lower that
degree. This proves the second identity. The finite support is the
union of the spectator support and the 25-point support at \(L=3\),
or its stated generalization. All source differentiations are of this
one functional.
\end{proof}

With the normalized stencil coefficient \(\ell^1\) norm, deletion of
the contracted scalar pair in \eqref{eq:v55-diagonal} has norm at most
\(a^{-4}\), because \(|\nu_{xy}|\le1\). The section has norm at most
\(a^4\gamma_L/\kappa\). Hence its zero-target projection has norm at
most \(1+\gamma_L/\kappa\). For \(L=3\), this is at most
\(21/5\) for the full canonical trace and \(37/5\) on the lower-mask
branch above. These are convenient upper bounds, not claims of sharp
norms for every spectator basis.

\begin{theorem}[Divergent contact budget and source norm]
\label{thm:v55-contact-budget}
Suppose the finite spectator stencil of
\(a^4(t_{\mathrm{tar}}-t_C)\) has fixed support and finite mass degree,
with coefficient budget
\begin{equation}
 \sum_{j,\,\text{spectator shifts}}
 a^{\delta-2j}|d_{j}(a)|\le B_tL_a.
 \label{eq:v55-trace-budget}
\end{equation}
Then every fixed scaled symbol seminorm and rooted kernel moment of the
correction in \eqref{eq:v55-trace-repair} obeys
\begin{equation}
 \boxed{
 \|C^{\mathrm{tr}}-C[W]\|_{r;\mu,m_*}
 +\|\mathscr K^{\mathrm{tr}}-\mathscr K\|_{\mathrm{rooted},b}
 \le C_{r,b}B_tL_a\,a^{2L-\delta}\epsilon_*^{2L}.}
 \label{eq:v55-contact-budget}
\end{equation}
The same native soft-BV incidence estimate applies to this correction,
with an additional factor \(\epsilon_*^{\kappa_\sigma}\) per output
incidence. The estimates are uniform in the periods on the preceding
nonwinding domain. The constants are not uniform as \(L\), source
arity, or spectator degree tends to infinity.
\end{theorem}
\begin{proof}
The factor \(b_L\) contributes \(a^{2L}\) and its first nonzero soft
weight \(2L\). Its fixed higher derivatives have the same scaled
window bound. The coefficient of \(z^j\) in
\eqref{eq:v55-trace-budget} is at most \(B_tL_a a^{2j-\delta}\);
use \(a^2z\le(a\epsilon_*)^2\le1\). Spectator translations have fixed
scaled derivative bounds, and \(\kappa^{-1}\le2\) is independent of
external momentum and mass. This proves the symbol estimate. Apply the
same external Fourier integration-by-parts and periodization argument
as in \cref{thm:v55-transport}. The BV assertion follows from that
lemma applied to the full correction, not from commutation of
\(B_{L,a}\) with a spatial derivation.
\end{proof}

A V54 stencil with budget \eqref{eq:v54-budget} automatically gives the
corresponding budget for \(a^4t_C\): a derivative monomial contributes
\(a^{-|\alpha|}\) to its stencil weights, exactly cancelling the
\(|\alpha|\) in that budget. The trace takes a bounded linear
combination of these weights. A target must have its own verified
budget. If \(2L>D+p\), the corrected stencil retains V54's whole
prescribed degree-\(D+p\) jet. At \(D=\delta=4,p=1,L=3\), both
improvements have an \(a^2L_a\) fixed-window error and the contact is
matched \emph{exactly}, even when its density grows as \(a^{-8}\).

\section{Measure convention, loop order, and the remaining hard operations}
\label{sec:v55-loop-ledger}

The scalar coefficient measure in \eqref{eq:v27-completed-Higgs} is
Lebesgue. Its geometric Hodge factors are in the action; the auxiliary
amplitude and coordinate Jacobian depend on the metric variables.
Accordingly, a scalar-linear correction of
\eqref{eq:v55-source-class} has no additional scalar-density derivative.
For any declared base density independent of the chosen scalar pair,
the same contraction calculation is unchanged. No new Haar factor or
scalar volume factor is introduced.

Use the canonical operator, or a consistently paired BV operator, when
expanding a master defect. Algebraically,
\begin{equation}
 \begin{aligned}
 \frac12(S_{a,0}+\hbar S_{a,1},S_{a,0}+\hbar S_{a,1})
 -\hbar\Delta_\rho(S_{a,0}+\hbar S_{a,1})
 ={}&\frac12(S_{a,0},S_{a,0})\\
 &+\hbar(s_aS_{a,1}-\Delta_\rho S_{a,0})\\
 &+\hbar^2\left(\frac12(S_{a,1},S_{a,1})
                         -\Delta_\rho S_{a,1}\right).
 \end{aligned}
 \label{eq:v55-loop-counting}
\end{equation}
No term in this equality assumes that the tree defect vanishes.
Thus replacing an order-\(\hbar\) counterterm affects the first loop
through its tree antibracket, and its own BV trace first appears in the
next coefficient. The trace counterexample is not a contradiction of
V54's one-loop effective-coefficient comparison.

For a change \(S_{a,1}\mapsto S_{a,1}+E_a\), the next coefficient
changes by
\begin{equation}
 (S_{a,1},E_a)+\frac12(E_a,E_a)-\Delta_\rho E_a.
 \label{eq:v55-two-loop-change}
\end{equation}
The contact section controls only the specified scalar trace operand.
The remaining brackets, regulated insertions, genuine loop contractions
with a counterterm, source renormalizations, and any required
second-order counterterm are not deleted. Exact trace matching is not
a two-loop master theorem.

For the actual scalar-linear source counterterms with finitely many
metric marks, the construction has the following direct use. First
specify the entire local coefficient and its finite spectator trace
normalization; the latter may be zero only when the source convention
or a calculation establishes it. Apply V54, then
\eqref{eq:v55-trace-repair}. All low jets and their ordinary descent
partners are unchanged, the additional finite first-loop row error is
controlled, and the declared scalar coincidence operand is exact. This
is an available finite realization rule. It is not an instruction to
replace an unknown critical-breaking coefficient or an existing source
normalization by zero.

The contact section does not presently cover arbitrary kernels
containing several scalar antifields, or coefficients that themselves
contain the scalar pair being traced. There the derivative also acts on
those coefficients. One cannot reuse \eqref{eq:v55-contact-identity}
without retaining these further contractions. Nor does it solve the
full nonminimal/chiral measure problem or the representative issue
left open by V53.

\section{Status and reproducible verification}
\label{sec:v55-ledgers}

\subsection*{Proved}
The native finite soft antibracket is a loop-free routed operation with
fixed-panel bounds; V54 realization errors therefore have a complete
incremental first-loop BV estimate. Arbitrarily high matched jets alone
do not bound a coincidence trace. The explicit axial stencil supplies
an exact right inverse for a specified scalar-linear coincidence
operand, preserving arbitrarily many soft jets and paying its
coefficient growth. The same statements include the finite lower-mask
normalization, the scalar measure convention, and the separated
next-loop algebraic ledger.

\subsection*{Inherited}
V54's section, weighted mass convention, and rooted periodization
machinery; the native completed action, source, nonminimal and density
conventions of V14--V33; prior continuum/modified-row estimates only on
their original domains; and V53's complete three-term current-ascent
packet. No companion theorem is reproved or upgraded here.

\subsection*{Literature input}
The cited exact-RG/BV and improvement papers supply comparison context.
No external cohomology classification, anomaly cancellation, or
many-shell theorem is used as a premise of the new finite proofs.

\subsection*{Conditional}
A realization of a populated counterterm requires its actual power/log
budget and its specified trace target. Preserving an already established
finite modified identity uses that identity as input; the estimate does
not repair an unpaid original defect. Applying a masked contraction
inside a full quantum master complex additionally requires its matching
bracket/reference convention. Higher-loop universality needs the actual
hard-leg insertions, not just the soft symbol theorem.

\subsection*{Open}
The native invariant critical descent still needs the representative
and complete coefficient decomposition identified in V53. The present
result supplies an implementable finite symmetry/measure interface,
not that missing local primitive. Multi-antifield coincidence matching,
full internal-gauge and chiral compatibility, an invariant many-shell
ESM class, and higher-loop continuum control remain separate. The useful
next operation is a populated complete descent or hard-contraction
calculation using this source-and-contact realization, rather than
another assertion that Taylor matching alone is a quantum symmetry.

\begin{center}\small
\begin{tabular}{p{.24\textwidth}p{.68\textwidth}}
\toprule
Variable & Scope of V55\\\midrule
Cutoff and volume & Fixed-panel coefficient-map constants are independent
of both; source bounds require the explicit budget, fixed physical
window, and nonwinding periods.\\
Mass and couplings & Nonnegative mass on a fixed retained domain,
polynomial mass jets, and fixed nonzero gravitational coupling margin.
No infrared-scale removal.\\
Source degree & Existing finite action/source panels; trace section
restricted to a scalar-linear one-antifield source with scalar-free
spectators.\\
Symmetry & Incremental native soft-BV control and an exact specified
coincidence operand; no all-frequency nilpotence or full QME claim.\\
Loop order & One-loop change of representative and one explicitly
identified next-loop operand. No new loop integral is evaluated.\\
Verification & Exact rational axial moments, discrete diagonal
identities, source projection, polynomial routing, coefficient budgets,
Grassmann signs, and byte-preservation audit. Analytic source estimates
are written proofs, not certified by finite tests alone.\\\bottomrule
\end{tabular}
\end{center}

\begingroup
\renewcommand{\refname}{References for V55}
\begin{thebibliography}{9}\small
\bibitem{V55IIM} Y.~Igarashi, K.~Itoh, and T.~R.~Morris,
\emph{BRST in the Exact RG}, arXiv:1904.08231 (2019),
\url{https://arxiv.org/abs/1904.08231}.
Comparison for consistently regulated BV and modified reference identities;
not a proof for the present finite completion.
\bibitem{V55RS} A.~Ramos and S.~Sint,
\emph{Symanzik improvement of the gradient flow in lattice gauge theories},
arXiv:1508.05552 (2015),
\url{https://arxiv.org/abs/1508.05552}.
Comparison for separate action, observable, and contact improvements.
\end{thebibliography}
\endgroup

\section*{Append-only continuation marker: V55}
\addcontentsline{toc}{section}{Append-only continuation marker: V55}
\label{sec:v55-continuation-marker}
\textbf{End of V55.} Preserve Sections 1--431. The finite-range local
source section now has an incremental native soft-BV estimate and a
separate exact scalar-coincidence correction. The latter is not inferred
from its high Taylor order. No active critical coefficient is guessed,
no existing contact is reset without its target, and no full quantum
continuum or two-loop symmetry theorem is claimed.


\clearpage
\part*{V56: Multiple-antifield contact matching and a populated current-ascent contraction}
\addcontentsline{toc}{part}{V56: Multiple-antifield contact matching and a populated current-ascent contraction}

\section{The missing contractions in V55}
\label{sec:v56-start}

Sections 1--431 and their continuation markers are preserved verbatim.
V55's coincidence section assumes that the coefficient of its displayed
scalar--antifield pair contains neither member of any scalar pair. That
assumption fails for V53's nonlinear current-ascent counterterms: the
coefficient depends on the scalar, and one term has two scalar
antifields. This installment evaluates a specified finite representative
of that counterterm and constructs a section which includes \emph{all}
scalar contractions of a rooted polynomial density. It does not guess a
coefficient of the unresolved native quantum breaking.

The mechanism is finite separation in the existing carrier. A pair of
linear scalar/antifield combinations is placed at finitely many translated
sites. Its moments vanish, and it is orthogonal, for the actual specified
contraction kernel, to every scalar site occurring in the input density.
Its self-contraction is strictly nonzero. The graded product rule then
gives a contraction homotopy even when the target itself contains scalars
and multiple antifields. Nilpotence imposes a necessary compatibility
condition on a prescribed target; a failed condition is retained as an
explicit residual.

The canonical odd Laplacian and its second-order product rule are standard
finite BV algebra \cite{V56Schwarz}. Compatibility of a cutoff contraction
with its bracket and reference functional is a separate exact-RG issue
\cite{V56IIM}. The new results concern the finite support, moment conditions,
mask-orthogonality, coefficient bounds, and the fully differentiated
current-ascent example. They do not identify an arbitrary masked Laplacian
with the complete measured reference insertion.

\subsection*{Conventions}
The scalar carrier, four real flavours, scalar Lebesgue coefficient
measure, and scalar-independent auxiliary density are unchanged. For an
even real translation-invariant kernel \(\nu\), use exactly the scalar
contraction of V55:
\begin{equation}
 \Delta^H_{\nu,a}
 =a^{-4}\sum_{x,y,A}\nu_{xy}
       \frac{\partial}{\partial\phi_A(y)}
       \frac{\partial^L}{\partial\upsilon_A(x)},
 \qquad \upsilon=\phi^*,\qquad \int_a=a^4\sum_x.
 \label{eq:v56-delta}
\end{equation}
The canonical case is \(\nu=I\); the other case treated below is the
specified hard mask \(\nu=I-\Theta_\lambda\), with
\(0\le\Theta_\lambda(p)\le1\) and Fourier support
\(|p|\le2\lambda\). In either case \((\Delta^H_{\nu,a})^2=0\):
commuting even derivatives and anticommuting left odd derivatives cancel
the double sum, including the repeated-index terms.

A \emph{rooted density} \(c_x\) is a finite super-polynomial whose scalar
and scalar-antifield sites are \(x+aS\), for a fixed finite set
\(S\subset\mathbb Z^4\). Other field, ghost, antifield, mass, and independent
source labels are retained as spectators of \(\Delta^H\). The same
coefficient rule is translated at every root. We work with this supplied
rooted representative, not a quotient by discrete total divergences:
multiplying a divergence by a new pair is not, in general, a divergence.
The integrated identities follow by summing the pointwise-rooted ones.
All polynomial degree bounds, supports, and stencil orders are fixed
before refinement.

\section{A finite pair orthogonal to the complete input support}
\label{sec:v56-mode}

Let \(s=|S|\), fix an integer \(L\ge1\), and choose
\(N=s+L+1\) distinct auxiliary offsets
\(t_j=(R+j)e_0\), \(0\le j<N\), outside \(S\). The integer \(R\)
and the period sizes are chosen so that these offsets and \(S\) are
injective on the finite torus. This is a condition on lattice-step
support, not an additional field species or a new underlying site type.
Let \(T=\{t_j\}\).

Form the real constraint matrix
\begin{equation}
 A_{S,T,L}=
 \begin{pmatrix}
   (R+j)^r\big|_{0\le r<L,\,0\le j<N}\\
   \nu_{S,T}
 \end{pmatrix}.
 \label{eq:v56-constraints}
\end{equation}
It has at most \(L+s=N-1\) independent rows. Let \(P_T\) be its
Euclidean orthogonal kernel projection. Choose the least index maximizing
\((P_T)_{jj}\), and put
\begin{equation}
 w=P_Te_j,\qquad d=w^Tw=(P_T)_{jj},\qquad
 \kappa=w^T\nu_{T,T}w.
 \label{eq:v56-mode-def}
\end{equation}
This prescription needs no square root. For exact rational input matrices,
one may compute \(P_T=Z(Z^TZ)^{-1}Z^T\) from any rational null basis \(Z\).
For a native smooth mask, the defining entries are their actual finite
Fourier sums. Replacing them by rational approximations requires a
separate residual estimate; exact native mask entries are not evaluated
numerically in this installment.

\begin{lemma}[Finite mask margin]
\label{lem:v56-mask}
Assume \(\lambda L_\mu\ge1\) for the four physical periods, and use the
Fourier-normalized matrix of \(\Theta_\lambda\). Then
\begin{equation}
 |(\Theta_\lambda)_{xy}|\le16(a\lambda)^4,
 \qquad
 \|\Theta_{\lambda,T,T}\|_{\op}\le16N(a\lambda)^4.
 \label{eq:v56-mask-bound}
\end{equation}
If \(16N(a\lambda)^4\le1/2\), the mode in
\eqref{eq:v56-mode-def} obeys
\begin{equation}
 \boxed{
 d\ge N^{-1},\qquad \tfrac12d\le\kappa\le d,\qquad
 \frac{\|w\|_1^2}{\kappa}\le2N.}
 \label{eq:v56-mode-bounds}
\end{equation}
For the canonical contraction, \(\kappa=d\) without a mask restriction.
\end{lemma}
\begin{proof}
In one direction, the number of labels with \(|p_\mu|\le2\lambda\)
is at most \(1+2\lambda L_\mu/\pi\le2\lambda L_\mu\). The product
bound, divided by the total number \(\prod_\mu(L_\mu/a)\) of sites,
gives the first estimate. The row-sum estimate gives the second.
The Fourier multiplier \(\Theta\) is positive, so its principal
submatrix is positive as well. It follows that
\((1-16N(a\lambda)^4)I\preceq\nu_{T,T}\preceq I\).
Since \(P_T\) has rank at least one, its maximal diagonal is at least
\(1/N\). The projection identity gives \(\|P_Te_j\|_2^2=(P_T)_{jj}\).
Finally \(\|w\|_1^2\le N\|w\|_2^2=Nd\). These statements prove
all the assertions.
\end{proof}

\begin{proposition}[Orthogonal high-moment mode]
\label{prop:v56-orthogonal}
The scalar symbol
\(b_w(k)=\sum_jw_j\exp(ia t_j\cdot k)\) has zero Taylor jet through
order \(L-1\). Moreover,
\begin{equation}
 \boxed{\nu_{S,T}w=0,\qquad w^T\nu_{T,S}=0.}
 \label{eq:v56-cross-zero}
\end{equation}
For every fixed derivative order, on a fixed soft window,
\(b_w=O(\sqrt d\,a^L|k|^L)\), with the corresponding scaled derivative
bound. Its constants depend on \(N,L,R\), but not on the mesh or periods.
\end{proposition}
\begin{proof}
The first \(L\) rows of \eqref{eq:v56-constraints} cancel the moments.
All offsets lie on the same axis, so this cancels every multivariable
Taylor coefficient of total degree less than \(L\). The remaining rows
give the first contraction equation, and symmetry of \(\nu\) gives the
second. Taylor's integral remainder and
\(\|w\|_1\le\sqrt{Nd}\) give the stated bounds. Fixed derivatives above
order \(L\) are bounded directly using the finite support; multiplication
by the corresponding window scale yields the same estimate for small
\(a\) at fixed window.
\end{proof}

Only a finite matrix kernel has been used. There is no inverse of the
cutoff multiplier \(\nu(p)\), which may vanish on retained modes. The
projection can change discontinuously when a constraint rank changes;
no differentiability in the regulator profile or conditioning estimate
for approximate kernel data is asserted. Its norm and the estimates above
remain uniform for the exact prescribed input.

\section{A contraction homotopy for nonlinear and multiple-antifield densities}
\label{sec:v56-homotopy}

At each root define, using the same \(w\) in every flavour,
\begin{equation}
 \psi_A(x)=\sum_jw_j\phi_A(x+at_j),\qquad
 \zeta_A(x)=\sum_jw_j\upsilon_A(x+at_j),\qquad
 U_x=\frac{a^4}{4\kappa}\sum_{A=1}^4\zeta_A(x)\psi_A(x).
 \label{eq:v56-U}
\end{equation}
The element \(U_x\) is odd and has ghost number minus one. Its two
factors are combinations of existing scalar coordinates, not additional
independent fields. Define \(H_xr_x=U_xr_x\), retaining this order for
all odd coefficients.

\begin{theorem}[Complete scalar-contraction homotopy]
\label{thm:v56-homotopy}
For every super-polynomial density \(r_x\) whose scalar and
scalar-antifield arguments lie in \(x+aS\),
\begin{equation}
 \boxed{
 \Delta^H_{\nu,a}(H_xr_x)+H_x\Delta^H_{\nu,a}r_x=r_x.}
 \label{eq:v56-homotopy}
\end{equation}
The identity allows arbitrary fixed numbers of scalar and scalar-antifield
arguments and arbitrary spectator-source labels. The section is finite
range, translation equivariant, and equivariant under constant orthogonal
flavour transformations, including all independent flavour signs.
Every coefficient of \(H_xr_x\) has a vanishing added pair jet through
total degree \(2L-1\).
\end{theorem}
\begin{proof}
The self-contraction of \(\zeta_A\psi_A\) is
\(a^{-4}w^T\nu_{T,T}w=a^{-4}\kappa\). Summing four flavours and
including the prefactor gives \(\Delta U_x=1\). A contraction between
one factor of \(U_x\) and one scalar or antifield of \(r_x\) vanishes
by \eqref{eq:v56-cross-zero}. This includes derivatives of every
nonlinear scalar coefficient of \(r_x\); none is treated as a frozen
spectator. The second-order product rule for the odd Laplacian therefore
reduces exactly to
\[\Delta(U_xr_x)=(\Delta U_x)r_x-U_x\Delta r_x.\]
The minus sign is the parity of \(U_x\), not an assumed parity of the
input. This proves the identity. In a coefficient Fourier transform the
new factors contribute \(b_w(k_\upsilon)b_w(k_\phi)\), each vanishing to
order \(L\). The support, translation, and flavour statements follow
from the construction.
\end{proof}

In normalized stencil coefficient norms,
\begin{equation}
 \boxed{\|H_x\|_{1\to1}\le2Na^4.}
 \label{eq:v56-H-norm}
\end{equation}
Indeed the four flavours cancel the averaging denominator and
\(\|w\|_1^2/\kappa\le2N\). This bound is independent of how many
antifields the particular input contains; contraction estimates for that
input still depend on its fixed polynomial degree.

\begin{remark}[Other canonical pairs]
\label{rem:v56-other-pairs}
The same identity holds for \(\Delta^H+\Delta^{\rm other}\) if the
second operator differentiates only spectator pairs, has no scalar-pair
coefficient dependence, and is part of the stated nilpotent contraction.
In fact \(\Delta^{\rm other}(U_xr_x)=-U_x\Delta^{\rm other}r_x\).
A scalar-independent bosonic density is compatible with this statement
when its BV density term also acts only on the other pairs. This is an
algebraic extension under explicit hypotheses, not a claim that an
unspecified chiral measure or the complete Legendre reference trace has
that form.
\end{remark}

\section{Compatible contact targets and exact failed-target residuals}
\label{sec:v56-target}

Let \(c_x\) be an actual supplied local source density and let \(t_x\)
be its prescribed contracted density, both rooted on \(S\). A necessary
condition is
\begin{equation}
 \Delta t_x=0.
 \label{eq:v56-target-closed}
\end{equation}
It is a finite polynomial identity; no scalar or antifield appearing in
\(t_x\) is omitted when checking it.

\begin{theorem}[Joint jet and multiple-antifield contact matching]
\label{thm:v56-matching}
Use \(H_x\) from \cref{thm:v56-homotopy} and put
\begin{equation}
 \boxed{
 c_x^{\rm tr}=c_x+H_x(t_x-\Delta c_x).}
 \label{eq:v56-repair}
\end{equation}
If \(\Delta t_x=0\), then
\begin{equation}
 \boxed{
 \Delta c_x^{\rm tr}=t_x,\qquad
 J_{\le2L-1}(c_x^{\rm tr}-c_x)=0.}
 \label{eq:v56-repair-result}
\end{equation}
For fixed compatible target the map is idempotent on its declared input
class. If the input is homogeneous and the target has the degree of
\(\Delta c_x\), the correction has the same total field/source degree
as \(c_x\). It vanishes at zero scalar antifields.

For an arbitrary, possibly incompatible target the exact residual is
\begin{equation}
 \boxed{
 \Delta c_x^{\rm tr}-t_x=-H_x\Delta t_x.}
 \label{eq:v56-invalid-target}
\end{equation}
In particular an input with \(\Delta t_x\ne0\) has no exact solution
of \(\Delta c_x'=t_x\) in any polynomial extension with the same
nilpotent contraction.
\end{theorem}
\begin{proof}
Apply \eqref{eq:v56-homotopy} to
\(r_x=t_x-\Delta c_x\), which is supported on \(S\), and use
\(\Delta^2c_x=0\). This gives the residual formula. If the target is
closed the residual vanishes; a second application has a zero gap and
makes no change. The moment statement is inherited from the two added
mode factors. One contraction removes one scalar and one scalar
antifield, whereas multiplication by \(U_x\) adds precisely that pair.
Nilpotence proves the impossibility assertion. All identities integrate
by the unchanged finite sum \(\int_a\).
\end{proof}

For the zero target no new compatibility assumption is needed:
\(t_x-\Delta c_x=-\Delta c_x\) is automatically closed. This is a
possible contact normalization, not a rule setting every native contact
to zero. For a nonzero physical target its actual polynomial coefficients
and its closure must be supplied. The section does not choose them.

For example, at one canonical scalar pair the target
\(t=\upsilon\phi c^0c^1\) obeys
\(\Delta t=a^{-4}c^0c^1\ne0\), and is therefore inadmissible by itself.
In contrast, for two scalar flavours,
\[c=\upsilon_1\upsilon_2\phi_1\phi_2c^0c^1\]
has the closed, scalar-dependent target
\[\Delta c=a^{-4}(\upsilon_2\phi_2-\upsilon_1\phi_1)c^0c^1.\]
The two terms are a single packet. Prescribing just one would fail
nilpotence. These finite examples illustrate the additional condition
that was automatic in V55's scalar-free target space.

For a degree bound \(d_\phi,d_\upsilon\), deletion of a contracted pair
has coefficient norm at most \(a^{-4}d_\phi d_\upsilon\), using
\(|\nu_{xy}|\le1\). Thus the zero-target projection has bound
\begin{equation}
 \|I-H\Delta\|_{1\to1}\le1+2N d_\phi d_\upsilon.
 \label{eq:v56-projection-norm}
\end{equation}
A nonclosed target retains
\(\|\Delta c^{\rm tr}-t\|_1\le2Na^4\|\Delta t\|_1\).
These are fixed upper bounds, not claimed optimal norms of every source
panel.

\section{The actual nonlinear current-ascent example has a nonzero contact}
\label{sec:v56-populated}

We now differentiate a concrete source functional already supplied in
\cref{sec:v53-start}. This is an evaluated coefficient of that example,
not an invented value for the still-unidentified native critical breaking.
On the flat scalar chart \(\rho=1\), take V53's
\(v^\mu=\phi^r\partial_\mu\phi\), with \(r=1\) or \(3\).
Let
\[f_B=\tfrac12 B_{\mu\nu}c^\mu c^\nu,
  \qquad \alpha_\mu=B_{\mu\nu}c^\nu.\]
They are respectively even and odd, independent of the contracted
scalar pair. The third superpotential term of V53 is zero for this
canonical current representative.

The following \emph{declared source realization} uses fourth-order
finite differences on the existing sites:
\begin{align}
 D_\mu&=
 \frac{-\mathsf T_\mu^2+8\mathsf T_\mu
       -8\mathsf T_\mu^{-1}+\mathsf T_\mu^{-2}}{12a},
 \notag\\
 P_a^{\rm src}&=z+
 \sum_{\mu=0}^{3}
 \frac{30-16(\mathsf T_\mu+\mathsf T_\mu^{-1})
               +(\mathsf T_\mu^2+\mathsf T_\mu^{-2})}{12a^2},
 \qquad f_a=P_a^{\rm src}\phi.
 \label{eq:v56-source-differences}
\end{align}
This does not replace V27's Gaussian scalar operator. It realizes the
\emph{local source polynomial} \(P_0=-\partial^2+z\), as allowed by
V54's coefficient interface. The symbols agree with the ordinary source
through the required weights: \(D=\partial+O(a^4\partial^5)\) and
\(P_a^{\rm src}=P_0+O(a^4\partial^6)\).

Retain both inherited terms, in their displayed Grassmann order:
\begin{equation}
 c_x^{(r)}=
 -\frac12 f_B\sum_\mu\phi^r(D_\mu\phi)\upsilon(D_\mu\upsilon)
 +\frac12\sum_\mu\alpha_\mu\phi^r(D_\mu\phi)f_a\upsilon.
 \label{eq:v56-example-density}
\end{equation}
Its continuum jet is V53's scalar current-ascent density. The first
realization error has weighted degree six. No finite-lattice BRST
identity for this standalone density is asserted.

\begin{theorem}[Complete canonical contact of the inherited source example]
\label{thm:v56-populated-contact}
On a nonwinding grid for \eqref{eq:v56-source-differences}, the canonical
scalar contraction of \eqref{eq:v56-example-density} is
\begin{align}
 \boxed{\begin{aligned}
 \Delta^H_{I,a}c_x^{(r)}
 =\frac{a^{-4}}2\Bigg\{&f_B\left[
       \frac{65}{18a^2}\phi^r\upsilon
       -r\phi^{r-1}\sum_\mu(D_\mu\phi)(D_\mu\upsilon)
                           \right]\\
 &-\sum_\mu\alpha_\mu\left[
       r\phi^{r-1}(D_\mu\phi)f_a
       +\left(z+\frac{10}{a^2}\right)
                       \phi^rD_\mu\phi\right]\Bigg\}.
 \end{aligned}}
 \label{eq:v56-complete-contact}
\end{align}
The density on the right is \(\Delta^H\)-closed. On constant scalar
and scalar-antifield values \(\phi=v\), \(\upsilon=\eta\), its value is
\begin{equation}
 \boxed{\left.\Delta^H_{I,a}c_x^{(r)}\right|_{v,\eta}
       =\frac{65}{36a^6}f_Bv^r\eta.}
 \label{eq:v56-constant-contact}
\end{equation}
The formula holds for every positive integer \(r\); the inherited
examples use odd \(r\) and respect the scalar sign symmetry.
\end{theorem}
\begin{proof}
The stencils have
\[
 (D_\mu)_{xx}=0,\qquad
 \sum_y(D_\mu)_{xy}^2=\frac{65}{72a^2},\qquad
 (P_a^{\rm src})_{xx}=z+\frac{10}{a^2}.
\]
For the first term of \eqref{eq:v56-example-density}, differentiating
\(\upsilon(x)\) first gives
\(-\tfrac r2 f_B\phi^{r-1}D_\mu\phi D_\mu\upsilon\).
Differentiating \(D_\mu\upsilon\) instead gives the opposite left-odd
sign; the scalar derivative then hits \(D_\mu\phi\), producing
\(\tfrac12 f_B\phi^r\upsilon\sum_y(D_\mu)_{xy}^2\).
Both contractions must be included. Their sum over four directions
accounts for the coefficient \(65/18\) inside the braces.

In the second term the left derivative must pass the odd \(\alpha_\mu\),
so it has a minus sign. The subsequent scalar derivative hits
\(\phi^r\) and \(f_a\). Its possible hit on \(D_\mu\phi\) at the
root is zero because \((D_\mu)_{xx}=0\). This gives exactly the second
line of \eqref{eq:v56-complete-contact}. Closure follows either from
nilpotence or by taking another contraction: the two terms containing
one scalar antifield cancel, while the other terms contain none.
Constant fields kill all displayed differences, leaving
\eqref{eq:v56-constant-contact}.
\end{proof}

The source density itself vanishes on constant scalar fields. Its
coincidence trace does not. This is a nonlinear/multiple-antifield
instance not covered by V55's scalar-free coefficient assumption.
It is a fixed finite contact, not an evaluated new loop integral and not
an anomaly assertion. A prescribed zero contact, or any other closed
target on the same support, can now be imposed by
\cref{thm:v56-matching} without changing the source jet through degree
five when \(L=3\). The source remains a descendant of the ordinary
current only at its retained jet; full finite symmetry is not inferred.

\section{Coefficient budgets, source locality, and one-loop transport}
\label{sec:v56-estimates}

Let \(r_x=t_x-\Delta c_x\). Use the same finite mass and source
normalization convention as V55, and assume the \emph{actual} normalized
contact density has budget
\begin{equation}
 a^4r_x=\sum_{j,\omega}d_{j,\omega}(a)z^j\mathsf M_{\omega,x},
 \qquad
 \sum_{j,\omega}a^{\delta-2j}|d_{j,\omega}(a)|
       \le B L_a,
 \quad L_a=1+|\log(a\lambda)|^\ell.
 \label{eq:v56-budget}
\end{equation}
The monomials \(\mathsf M_\omega\) are supplied finite shifted source
words. Their maximum degree and support are fixed. They may contain the
contracted scalars and any fixed number of scalar antifields.

\begin{theorem}[Uniform contact-corrected source comparison]
\label{thm:v56-estimate}
On \(16N(a\lambda)^4\le1/2\), with nonwinding fixed support,
\(\lambda L_\mu\ge1\), a fixed external window \(|p_i|\le\mu\), and
\(0\le z\le m_*^2\), put \(\epsilon_*=A_*\mu+m_*\). If
\(a\epsilon_*\) is sufficiently small for the fixed stencil, then
every prescribed scaled symbol seminorm and rooted kernel moment obeys
\begin{equation}
 \boxed{
 \|c^{\rm tr}-c\|_{r;\mu,m_*}
 +\|\mathscr K^{\rm tr}-\mathscr K\|_{\rm rooted,b}
 \le C_{r,b,S,L}B L_a\,a^{2L-\delta}\epsilon_*^{2L}.}
 \label{eq:v56-bound}
\end{equation}
The constants are independent of mesh and periods under those conditions.
They are not uniform in growing stencil support, polynomial degree,
source order, or in a vanishing gravitational coupling margin included
in \(B\).
\end{theorem}
\begin{proof}
Each inserted Fourier mode has its order-\(L\) moment cancellation.
Their product has bound \(C d a^{2L}\epsilon_*^{2L}\); division by
\(\kappa\ge d/2\) cancels the mode normalization. The factor \(a^4\)
in \(H\) combines with \eqref{eq:v56-budget}. A mass term satisfies
\(a^{2j}z^j\le(a\epsilon_*)^{2j}\), so it cannot worsen the estimate.
Differentiating at fixed order uses the same finite-mode moments and
finite shift bounds. The mode coefficients are independent of external
momenta, mass, and the retained field/source coordinates. Possible
rank changes as the regulator varies do not enter these derivatives.

Multiply by a fixed smooth external window and integrate by parts in
each of the \(4(A-1)\) relative Fourier variables. Taking sufficiently
many of the just established derivatives controls any specified
polynomial position moment. Periodization bounds the finite rooted sum
by the corresponding continuum weighted integral under the same
window/period conditions. Each relative coordinate contributes its
normalizing \(a^4\); there is no free sum over an unrooted volume.
This is the same proved source topology as V54--V55, applied to the
two added high-moment factors rather than a single axial factor.
\end{proof}

For example, \(\delta=4,L=3\) gives unchanged jets through degree five,
an exactly matched compatible contact, and
\begin{equation}
 O\!\left(a^2L_a\epsilon_*^6\right)
 \label{eq:v56-critical-rate}
\end{equation}
in the fixed source norms. This estimate requires the actual budget of
the target, not merely that the target is closed. Unknown native matching
coefficients are not bounded by assigning them a name.

By V55's native loop-free routing theorem, the additional soft tree-BV
variation of this correction satisfies the same bound with the actual
multiplier degree \(\epsilon_*^{\kappa_\sigma}\) on each output panel.
The source arity need not increase: \(\Delta\) deletes the pair that
\(H\) reinserts. All original quantum contacts and all pre-existing
finite-regulator defects retain their own ownership.

\section{Limits of the contraction result and next-loop bookkeeping}
\label{sec:v56-loop}

The construction solves one specified nilpotent contraction equation
on supplied rooted coefficients. It does not solve the classical or
quantum master equation. In particular, at order \(\hbar\) a replacement
of a counterterm changes its tree antibracket; its own odd-Laplacian
contraction enters the following loop coefficient. For a change \(E\)
of the first-loop action coefficient, the remaining next-loop change is
still
\begin{equation}
 (S_{a,1},E)+\tfrac12(E,E)-\Delta_\rho E.
 \label{eq:v56-next-loop}
\end{equation}
Even if the scalar part of the last term is set to its specified target,
the brackets, genuine hard-loop insertions, other measure sectors, and
second-order counterterms are not removed.

The new homotopy is not an algebra homomorphism, a bracket intertwiner,
or a positive stochastic operation. Its scalar-averaged pair preserves
the flavour symmetries. A fixed lattice direction in its support is
regulator data; point-group symmetry may be restored by averaging whole
sections over a finite group preserving the declared support and mask.
Every averaged section still has the same contraction identity and
vanishing pair jet. This does not manufacture a covariance property not
satisfied by the input source.

The use of several translated sites is essential to avoid contracting
into an input scalar coefficient. It is nevertheless finite range on
the original carrier. For the canonical contraction, an especially simple
choice consists of \(L+1\) consecutive auxiliary sites outside \(S\),
with weights \(w_j=(-1)^j\binom Lj\). Its pair normalization is
\(\kappa=\binom{2L}{L}\) and
\begin{equation}
 \frac{\|w\|_1^2}{\kappa}
 =\frac{4^L}{\binom{2L}{L}}.
 \label{eq:v56-binomial}
\end{equation}
At \(L=3\), four existing auxiliary positions and sixteen pair monomials
suffice, with coefficient norm \(16/5\), before multiplying the old
source density. These differ from V55's on-site-antifield/25-point-field
choice: \emph{both} new arguments now have vanishing moments, and neither
can contract into any old argument. For the actual lower mask, the finite
constraint construction replaces spatial disjointness by exact
mask-orthogonality.

At fixed cutoff no section can cover a source supported on every scalar
coordinate while also requiring a new orthogonal pair in that same finite
space. The theorem instead applies to fixed finite local supports on
sufficiently large nonwinding tori, which is the source regime used for
refinement. It supplies no growing-volume top-antifield cohomology claim.
Nor does it transport a target modulo arbitrary changes of rooted
representative without recomputing that representative's contraction.

\section{Status and verification}
\label{sec:v56-ledgers}

\subsection*{Proved}
A finite support/moment/mask construction yields a complete scalar
contraction homotopy for polynomial densities containing multiple scalar
fields and antifields. Compatible targets are characterized by their
nilpotence condition, realized with unchanged prescribed soft jets, and
controlled in the existing coefficient and rooted source norms with their
actual divergence budget. Failed targets retain an exact residual. The
canonical contraction of an explicit finite representative of V53's
nonlinear two-term ascent packet is evaluated, including its nonzero
\(65/(36a^6)\) constant-background coefficient.

\subsection*{Inherited}
V53 supplies the ordinary current-ascent density, not its new coincidence
calculation. V54 supplies the general source-jet realization and its
coefficient conventions. V55 supplies the scalar measure/contraction
convention, the soft native antibracket transport, and the loop-order
ledger. Their original proofs and domains are not reaudited in full or
claimed again as new results. The original quantum action and reference
remain unchanged.

\subsection*{Literature input}
The ordinary finite BV product rule and the distinction between a
cutoff contraction and a consistently regulated master/reference identity
are standard. The cited sources are comparison context; no cohomology
vanishing theorem or all-loop renormalization theorem is imported to
solve the native critical breaking.

\subsection*{Conditional}
Actual source/contact matching requires the supplied rooted local array,
the specified target and its closure, the fixed scalar-independent
measure convention, and the coefficient budget. Applying the optional
other-pair statement requires its explicitly matched nilpotent operator.
The native smooth-mask mode is defined by finite exact Fourier data;
perturbation/conditioning of approximate mode computation is a separate
numerical question. Preservation of an earlier one-loop identity uses
that identity and its original residual estimate as input.

\subsection*{Open}
The native invariant critical descent still requires its representative
and populated coefficient decomposition from V53. The present result
closes V55's multiple-scalar/multiple-antifield coincidence interface,
not that descent. Further local stencil families are not needed before
a concrete remaining master/reference coefficient is supplied. The next
useful operation is that complete coefficient (or a genuine hard-loop
insertion) with its compatible target, brackets, and source labels kept
together. Internal-gauge/chiral transport, the invariant many-shell ESM
class, and higher-loop continuum control remain separate.

\begin{center}\small
\begin{tabular}{p{.23\textwidth}p{.69\textwidth}}
\toprule
Variable & Scope of V56\\\midrule
Carrier and locality & Existing scalar/antifield sites; finite translated
support at fixed input support and moment order. No new field species.\\
Cutoff and volume & Uniform coefficient/source bounds on the stated
mask margin, small external window, and nonwinding period domain.\\
Mass and sources & Polynomial mass dependence; arbitrary fixed scalar
and scalar-antifield degree; every spectator label retained.\\
Symmetry & Exact specified contraction; compatible target required.
No unmatched masked/canonical bracket or full QME claim.\\
Loop order & A populated next-loop contact operand; no new loop
integral and no two-loop cancellation of the entire master defect.\\
Verification & Exact supplied-matrix modes, graded polynomial identities,
multiple-antifield contractions, native-example contact, and preservation
checks. Written analytic estimates are not proved by finite tests.\\
\bottomrule
\end{tabular}
\end{center}

\needspace{18\baselineskip}
\begingroup
\renewcommand{\refname}{References for V56}
\begin{thebibliography}{9}\small
\bibitem{V56Schwarz} A.~Schwarz,
\emph{Geometry of Batalin--Vilkovisky quantization},
arXiv:hep-th/9205088 (1992),
\url{https://arxiv.org/abs/hep-th/9205088}.
Finite odd-Laplacian and measure framework; not the source of the present
mask-orthogonal finite-support section.
\bibitem{V56IIM} Y.~Igarashi, K.~Itoh, and T.~R.~Morris,
\emph{BRST in the Exact RG}, arXiv:1904.08231 (2019),
\url{https://arxiv.org/abs/1904.08231}.
Comparison for matched cutoff BV and modified reference identities;
not a proof of the remaining native ESM descent.
\end{thebibliography}
\endgroup

\section*{Append-only continuation marker: V56}
\addcontentsline{toc}{section}{Append-only continuation marker: V56}
\label{sec:v56-continuation-marker}
\textbf{End of V56.} Preserve Sections 1--439. The complete polynomial
scalar-contact operation now has a finite, high-moment, mask-orthogonal
section, including multiple antifields and derivatives of scalar
coefficients. Targets must pass their actual nilpotence test. The
current-ascent contact above is evaluated but is not substituted for an
unknown quantum breaking. No active normalization, measure, propagator,
or full continuum claim is changed.


\clearpage
\section{The next operation is a measured scalar contraction, not another stencil}
\label{sec:v57-start}

Sections 1--439 and their continuation markers are retained verbatim.
V56 constructs a finite-support section for a specified odd Laplacian.
That operation is not a scalar Gaussian expectation. The latter contracts
\emph{two scalar fields}, with covariance \(\nu_\lambda\widetilde
P_a^{-1}\), whereas V56 contracts a scalar with a scalar antifield,
with kernel \(\nu\) or the canonical kernel. This installment computes
both operations on the same supplied nonlinear source, and then gives
the loop-graded partner needed to transport V56's section through the
actual free scalar Gaussian. No new stencil family is introduced.

The genuinely evaluated insertion is V56's finite representative of
V53's current-ascent density. Its entire Gaussian image is determined by
two finite native covariance sums. One of these has a strictly positive
\(a^{-4}\) leading scale and generates a lower-scalar-degree source term.
It cannot be discarded as a vacuum constant. A scalar integration-by-parts
identity retains the free Euler term, the odd contact, the interacting
scalar force, and the lower-reference term together. This is an exact
conditional identity on the existing carrier, not an assertion of the
full Einstein--SM quantum master equation.

Gaussian moment expansions and finite BV transfer are established
machinery \cite{V57DJP,V57Schwarz}; V1 already proves a finite-order measured
compiler. The new results below are the native covariance operands for the
existing nonlinear density, its explicit lower-degree mixing, and the
conjugated contact section with its actual loop and source budgets.
They do not populate the unidentified invariant critical breaking.

\subsection*{Conventions and the common carrier}
Use the odd periodic four-dimensional grids, physical pairing
\(a^4\sum_x\), four real scalar flavours, and the \emph{same}
completed free scalar operator \(\widetilde P_a\) of
\eqref{eq:v27-P}. Write \(z=m^2\),
\(\nu=I-\Theta_\lambda\), and \(V_L=\prod_\mu L_\mu\).
The coordinate covariance, with its factor of \(\hbar\) removed, is
\begin{equation}
 \boxed{\mathsf G_{xy}
 =\frac1{V_L}\sum_{p\in\mathrm{BZ}_{a,L}}
 e^{ip\cdot(x-y)}\frac{\nu(p)}{\widetilde P_a(p)},
 \qquad \widetilde P_a\mathsf G=a^{-4}\nu.}
 \label{eq:v57-covariance}
\end{equation}
The quotient is zero where \(\nu=0\), including the retained zero mode
at \(z=0\). The Gaussian fluctuation \(\chi\) has covariance
\(\hbar\mathsf G\); scalar antifields \(\upsilon\), ghosts, and
other source labels are retained. The underlying scalar Gaussian is a
factor of the existing reference, not the marginal law of the fully
interacting Einstein--SM theory.

For finite polynomial source coefficients define
\begin{equation}
 \mathscr D_{\mathsf G}
 =\frac12\sum_{A,x,y}\mathsf G_{xy}
       \partial_{\phi_A(x)}\partial_{\phi_A(y)},
 \qquad
 \mathscr W_{\mathsf G}=\exp(\hbar\mathscr D_{\mathsf G}).
 \label{eq:v57-Wick}
\end{equation}
Thus \(\mathscr W_{\mathsf G}F(\varphi,\upsilon)
=\E[F(\varphi+\chi,\upsilon)]\). Its exponential and inverse terminate
at each scalar degree. They do not integrate antifields. All odd
coefficients keep their displayed order.

The domains used for bounds have fixed polynomial/source order, fixed
finite supports and profiles, \(\lambda L_\mu\ge1\),
\(a\lambda\le1/8\), and \(0\le m\le M\lambda\).
They also retain V56's mask margin
\(16N(a\lambda)^4\le1/2\) when its \(N\)-site section is used,
and are intersected with the earlier, more restrictive domains when an
earlier continuum theorem is consumed. Physical external windows and
nonwinding conditions are retained for source-kernel estimates. There is
no removal of the lower Wilsonian scale or uniformity in infinite source
or perturbative order.

\section{Two actual covariance moments control the nonlinear insertion}
\label{sec:v57-moments}

Keep V56's \emph{source} differences \(D_\mu\) and
\(P_a^{\mathrm{src}}\) from \eqref{eq:v56-source-differences}.
They are not substituted for \(\widetilde P_a\) in the Gaussian.
For each root put
\begin{equation}
 \boxed{\begin{aligned}
 u_a&=\mathsf G_{xx}
   =\frac1{V_L}\sum_p\frac{\nu(p)}{\widetilde P_a(p)},\\
 v_a&=(\mathsf G P_a^{\mathrm{src}})_{xx}
   =\frac1{V_L}\sum_p
      \frac{\nu(p)P_a^{\mathrm{src}}(p)}{\widetilde P_a(p)}.
 \end{aligned}}
 \label{eq:v57-uv}
\end{equation}
These are respectively a scalar variance and a scalar--source-Euler
covariance. They are not odd-Laplacian contacts. Translation invariance
makes them independent of the root. Evenness of the covariance and of
\(P_a^{\mathrm{src}}\), and oddness of \(D_\mu\), give exactly
\begin{equation}
 (D_\mu\mathsf G)_{xx}=0,
 \qquad (D_\mu\mathsf G P_a^{\mathrm{src}})_{xx}=0.
 \label{eq:v57-odd-moments}
\end{equation}
No rotational symmetry is needed.

\begin{proposition}[Positive native covariance scales]
\label{prop:v57-moment-bounds}
Let \(C_P\) be the fixed upper constant in
\(z+4|p|^2/\pi^2\le\widetilde P_a\le z+C_P|p|^2\).
Set
\[
 r_- =\min\{1,4/(\pi^2C_P)\},\qquad
 r_+ =\max\{1,\pi^2/3\}.
\]
On the stated domain,
\begin{align}
 \frac{1-16(a\lambda)^4}{4\pi^2 C_P+M^2}\,a^{-2}
 &\le u_a\le C a^{-2},\label{eq:v57-u-bound}\\
 r_-[1-16(a\lambda)^4]a^{-4}
 &\le v_a\le r_+a^{-4}.\label{eq:v57-v-bound}
\end{align}
For fixed profiles, positive physical periods and fixed \(z\),
\(a^2u_a\) and \(a^4v_a\) have strictly positive limits along odd
refinement. Writing
\(p_\star(y)=a^2(\widetilde P_a(y/a)-z)\) and
\(p_s(y)=a^2(P_a^{\mathrm{src}}(y/a)-z)\), which are independent of
\(a,z\), these limits are
\begin{equation}
 \int_{[-\pi,\pi]^4}\frac{\dd^4y}{(2\pi)^4p_\star(y)},
 \qquad
 \int_{[-\pi,\pi]^4}\frac{p_s(y)}{p_\star(y)}
                         \frac{\dd^4y}{(2\pi)^4}.
 \label{eq:v57-leading-integrals}
\end{equation}
Their values depend on the recorded ultraviolet completion. No numerical
values for these integrals are claimed.
\end{proposition}
\begin{proof}
For principal momenta,
\[
 P_a^{\mathrm{src}}(p)-z
 =\frac4{a^2}\sum_\mu\sin^2(ap_\mu/2)
                    [1+\tfrac13\sin^2(ap_\mu/2)]
\]
is between \(4|p|^2/\pi^2\) and \(4|p|^2/3\).
The ratio \(P_a^{\mathrm{src}}/\widetilde P_a\) is therefore
between \(r_-\) and \(r_+\) on nonzero hard modes. V56's finite
count gives
\(\operatorname{av}_{\mathrm{BZ}}\nu\ge1-16(a\lambda)^4\).
Since \(V_L^{-1}\sum_p=a^{-4}\operatorname{av}_{\mathrm{BZ}}\),
this proves \eqref{eq:v57-v-bound}.

The maximum denominator in the principal cube is at most
\(a^{-2}(4\pi^2 C_P+M^2)\) when \(a\lambda\le1\).
This proves the lower bound for \(u_a\). For the upper bound, the
covariance vanishes below \(\lambda\), its symbol is at most
\(C|p|^{-2}\), and the number of periodic momenta below radius
\(R\ge\lambda\), divided by volume, is at most \(CR^4\).
Summation in dyadic annuli up to \(O(a^{-1})\) is at most \(Ca^{-2}\).

For the limits, remove a fixed dimensionless ball of radius \(\epsilon\).
On its complement ordinary Riemann sums converge uniformly as
\(L_\mu/a\to\infty\), the mass becomes \(a^2z\to0\), and the
lower mask becomes one. Inside the ball the first integrand and its sums
have mass at most \(C\epsilon^2\), by the same annular count; the
second ratio is bounded and has mass at most \(C\epsilon^4\), with a
vanishing single-cell correction. First refine and then take
\(\epsilon\downarrow0\). Positivity follows from the ratio bounds
and positivity away from the single zero of \(p_\star\).
\end{proof}

Mass differentiation is at fixed reference mask and uses one mask factor:
\begin{align}
 \partial_z^j u_a
 &=\frac{(-1)^j j!}{V_L}\sum_p
                  \frac{\nu(p)}{\widetilde P_a(p)^{j+1}},\notag\\
 \partial_z^j v_a
 &=\frac{(-1)^{j-1}j!}{V_L}\sum_p
 \frac{\nu(p)[\widetilde P_a(p)-P_a^{\mathrm{src}}(p)]}
      {\widetilde P_a(p)^{j+1}},\qquad j\ge1.
 \label{eq:v57-mass-moments}
\end{align}
The difference in the second numerator is mass independent.
In particular the bounds for \(u_a\) have orders
\(a^{-2}\), \(1+\log(1/(a\lambda))\), and
\(\lambda^{2-2j}\) for \(j=0,1,\ge2\), respectively.
For \(v_a\), the conservative orders are \(a^{-4}\), \(a^{-2}\),
\(1+\log(1/(a\lambda))\), and \(\lambda^{4-2j}\) for
\(j=0,1,2,\ge3\). These follow by the same counting proof.
They are not small-coefficient estimates.

\begin{remark}[Exact mass dependence versus a mass-polynomial EFT bank]
\label{rem:v57-mass-jet}
The fixed-mass Gaussian coefficients \(u_a(z),v_a(z)\) are generally
not polynomials in \(z\). They are regular at zero because every
contracted mode lies above the fixed lower scale. An EFT counterterm
required to be mass polynomial uses one declared finite mass jet of the
\emph{complete} identities below. All those identities remain exact in
\(\mathbb R[z]/(z^{J+1})\) after coherent truncation: multiplication,
scalar differentiation and the odd contraction preserve this ideal.
The coefficients are derivatives of the same massive covariance, not a
replacement by powers of a differently masked propagator. For example,
with \(b=\widetilde P_a(p)|_{z=0}>0\),
\[
 \frac1{b+z}-\sum_{j=0}^J\frac{(-z)^j}{b^{j+1}}
 =\frac{(-z)^{J+1}}{b^{J+1}(b+z)}.
\]
This supplies the ordinary remainder bound on the nonnegative mass
interval. Exact fixed-mass normal ordering and its prescribed finite
mass-polynomial jet are distinguished throughout; neither asserts an
infinite mass-polynomial identity.
\end{remark}

\section{The complete Gaussian image of the inherited ascent density}
\label{sec:v57-populated}

Fix one scalar flavour and retain the exact Grassmann order of
\eqref{eq:v56-example-density}. Abbreviate \(f_a=P_a^{\mathrm{src}}\phi\):
\begin{equation}
 c_x^{(r)}=-\frac12f_B\sum_\mu\phi^r(D_\mu\phi)
                  \upsilon(D_\mu\upsilon)
 +\frac12\sum_\mu\alpha_\mu\phi^r(D_\mu\phi)f_a\upsilon.
 \label{eq:v57-c}
\end{equation}
The inherited examples are \(r=1,3\). Define the moment polynomial
\begin{equation}
 M_r(X,t)=\sum_{j=0}^{\lfloor r/2\rfloor}
 \frac{r!}{2^j j!(r-2j)!}\,t^jX^{r-2j}.
 \label{eq:v57-moment-polynomial}
\end{equation}

\begin{theorem}[All scalar self-contractions of the two-term source]
\label{thm:v57-populated}
For every fixed positive integer \(r\), its complete free scalar
Gaussian image is
\begin{equation}
 \boxed{\begin{aligned}
 \mathscr W_{\mathsf G}c_x^{(r)}
 ={}&-\frac12 f_B\sum_\mu
    M_r(\phi,\hbar u_a)(D_\mu\phi)\upsilon(D_\mu\upsilon)\\
 &+\frac12\sum_\mu\alpha_\mu(D_\mu\phi)
  [f_aM_r(\phi,\hbar u_a)
       +\hbar r v_aM_{r-1}(\phi,\hbar u_a)]\upsilon.
 \end{aligned}}
 \label{eq:v57-full-Wick}
\end{equation}
There is no extra factor of four for a source in one fixed flavour.
For the two inherited densities this becomes
\begin{align}
 \mathscr W c^{(1)}
 &=c^{(1)}+\frac{\hbar v_a}{2}
                      \sum_\mu\alpha_\mu D_\mu\phi\,\upsilon,
 \label{eq:v57-r1}\\
 \mathscr W c^{(3)}
 &=c^{(3)}+3\hbar u_ac^{(1)}
   +\frac{3\hbar v_a}{2}
            \sum_\mu\alpha_\mu\phi^2D_\mu\phi\,\upsilon
   +\frac{3\hbar^2u_av_a}{2}
            \sum_\mu\alpha_\mu D_\mu\phi\,\upsilon.
 \label{eq:v57-r3}
\end{align}
These are exact finite Gaussian coefficients, not only asymptotic
large-momentum formulas.
\end{theorem}
\begin{proof}
At a fixed root use the jointly Gaussian linear coordinates
\(X=\chi_x\), \(Y_\mu=(D_\mu\chi)_x\), and
\(Z=(P_a^{\mathrm{src}}\chi)_x\).
Their covariances include
\(\E X^2=\hbar u_a\), \(\E XZ=\hbar v_a\),
\(\E XY_\mu=\E Y_\mu Z=0\), by
\eqref{eq:v57-odd-moments}. The variance of \(Z\) and the covariances
among the \(Y_\mu\) never occur, because each summand contains only
one \(Y_\mu\) and at most one \(Z\).
Pairing powers of \(X\) gives \(M_r\). In the second source term,
\(Z\) can pair with any one of the \(r\) copies of \(X\), leaving
\(r\hbar v_aM_{r-1}\). All antifields and ghosts are spectators
of these even contractions, so no additional odd reordering is made.
This proves \eqref{eq:v57-full-Wick}; expansion gives the two restrictions.
\end{proof}

The coefficient \(v_a/2\) in \eqref{eq:v57-r1} is strictly positive
and of order \(a^{-4}\). It multiplies a transformation-source term,
not a source-independent vacuum density. The normalization of a scalar
Gaussian partition function cannot remove it. Formula
\eqref{eq:v57-r3} additionally has a double self-contraction of order
\(a^{-6}\). At a single marked insertion all self-contractions are
connected to that same source vertex; no additional normalization graph
has been omitted.

Precisely, if the action is perturbed by
\(\varepsilon\hbar^k\int_a c\), its effective action differentiated
at \(\varepsilon=0\), on the pure scalar reference branch, changes by
\(\hbar^k\int_a\mathscr Wc\). A self-contraction in
\eqref{eq:v57-full-Wick} therefore raises the loop index by one.
This is a coefficient linear in the source insertion, not the complete
higher-loop coefficient of the interacting Einstein--SM action.

\begin{corollary}[The finite normal-ordering partner is fixed]
\label{cor:v57-normal-order}
The unique polynomial input with prescribed Gaussian image
\(c^{(r)}\) is \(\mathscr W^{-1}c^{(r)}\), obtained from
\eqref{eq:v57-full-Wick} by \(\hbar\mapsto-\hbar\).
For \(r=3\), its linear loop term changes sign and its
\(3\hbar^2u_av_a/2\) term does not. It has the same site support,
with a finite tower of lower scalar degrees and no negative powers of
\(\hbar\). This normal ordering is a selectable reference normalization;
it is not installed in the active theory here.
\end{corollary}
\begin{proof}
The constant-coefficient operator \(\mathscr D_{\mathsf G}\) lowers
scalar degree by two. Its exponential has inverse with the opposite
sign on every finite-degree polynomial. All derivatives act on original
arguments, so support cannot enlarge. Every coefficient in the inverse
has the loop index stated above.
\end{proof}

\section{The scalar Euler, contact and reference terms form one exact identity}
\label{sec:v57-balance}

Let
\[
 \delta_P=\sum_{A,x}(\widetilde P_a\phi_A)_x
                    \frac{\partial^L}{\partial\upsilon_A(x)},
 \qquad
 \Delta_\nu=a^{-4}\sum_{A,x,y}\nu_{xy}
 \partial_{\phi_A(y)}\frac{\partial^L}{\partial\upsilon_A(x)}.
\]
The scalar-dependent part of the auxiliary measure is absent, as in V27;
no second scalar volume factor is added. Introduce
\(\mathscr Q_{P,\nu}=\delta_P-\hbar\Delta_\nu\).
This notation denotes the scalar Gaussian integration-by-parts complex.
For \(\nu\ne I\), it is \emph{not} being identified with a full
canonical interacting BV differential or an unmatched masked bracket.

\begin{proposition}[Exact Gaussian--Euler intertwining]
\label{prop:v57-Euler-intertwining}
On finite scalar/source polynomials,
\begin{equation}
 \boxed{
 [\mathscr D_{\mathsf G},\delta_P]=\Delta_\nu,
 \quad [\mathscr W_{\mathsf G},\Delta_\nu]=0,
 \quad
 \mathscr W_{\mathsf G}\mathscr Q_{P,\nu}
       =\delta_P\mathscr W_{\mathsf G}.}
 \label{eq:v57-intertwining}
\end{equation}
In particular \(\mathscr Q_{P,\nu}^2=0\).
These identities hold with multiple antifields and scalar-dependent
coefficients, without taking a continuum limit.
\end{proposition}
\begin{proof}
Only differentiating the linear multiplier
\((\widetilde P_a\phi)_x\) contributes to the commutator.
The resulting matrix is \(\mathsf G\widetilde P_a=a^{-4}\nu\).
All higher commutators vanish because the remaining derivatives have
constant coefficients. The heat operator has even derivatives and
commutes with every derivative in \(\Delta_\nu\). Exponentiation
proves the third identity. The odd derivation \(\delta_P\) squares
to zero, and conjugation by the invertible finite heat operator proves
the nilpotence statement.
\end{proof}

\begin{corollary}[The unit distinguishes the contact and Gaussian complexes]
\label{cor:v57-unit}
The unit belongs to the image of V56's odd contraction on the enlarged
finite support, but not to the image of \(\mathscr Q_{P,\nu}\) on
polynomial insertions. In particular there is no polynomial \(F\) with
\(\mathscr Q_{P,\nu}F=1\).
\end{corollary}
\begin{proof}
The first statement is \(\Delta_\nu U_x=1\). For the second, apply
\(\mathscr W\) and then set the retained scalar field to zero.
Equation \eqref{eq:v57-intertwining} makes the left side zero, because
all multipliers of \(\delta_P\) vanish there. The right side remains
one. This argument is coefficientwise and uses no division by \(\hbar\).
\end{proof}

This calculation has a concrete consequence for V56. Let \(U_x\) be
its pair for this same \(\nu\), with normalization \(\kappa=w^T\nu w\).
Put
\begin{equation}
 \mathcal E_{P,x}=\delta_PU_x
 =\frac{a^4}{4\kappa}\sum_A
       [w\widetilde P_a\phi_A](x)[w\phi_A](x).
 \label{eq:v57-pair-energy}
\end{equation}
Using \(\Delta_\nu U_x=1\), one obtains
\begin{equation}
 \boxed{\mathscr Q_{P,\nu}U_x=\mathcal E_{P,x}-\hbar,
 \qquad
 \mathscr W\mathcal E_{P,x}=\mathcal E_{P,x}+\hbar.}
 \label{eq:v57-pair-balance}
\end{equation}
At zero retained scalar field the two terms cancel exactly in the
measured identity. Thus the unit produced by the \emph{odd} contraction
has an equally sized hard Gaussian Euler return. V56's contraction
homotopy cannot be substituted for a homotopy of the full quantum
operator; dividing it by \(\hbar\) would also leave the admissible
nonnegative-loop packet.

\begin{theorem}[Conditional native scalar Euler/reference balance]
\label{thm:v57-interacting-balance}
Retain the other native fields and sources as parameters, and let
\(V(\phi,\upsilon;\text{spectators})\) be the complete remaining even
interaction in the scalar integral. Denote its physical scalar force by
\(v_i=a^{-4}\partial_{\phi_i}V\), with \(i\) including flavour.
For \(\phi=\varphi+\chi\), define the normalized expectation with
weight \(\exp(-V/\hbar)\) relative to covariance
\(\hbar\mathsf G\). Then, coefficientwise in its finite perturbative
expansion,
\begin{equation}
 \boxed{\begin{aligned}
 \left\langle
 \sum_i[(\widetilde P_a\phi)_i+v_i]\,
                         \partial^L_{\upsilon_i}F
                   -\hbar\Delta_\nu F\right\rangle_V
 ={}&\sum_i(\widetilde P_a\varphi)_i
                        \langle\partial^L_{\upsilon_i}F\rangle_V\\
 &+\left\langle\sum_i(\Theta_\lambda v)_i
                             \partial^L_{\upsilon_i}F\right\rangle_V.
 \end{aligned}}
 \label{eq:v57-conditional-identity}
\end{equation}
The identity is also literal whenever the finite integral and its
integration by parts exist. It remains true inside a subsequent
integration over the spectator species. It is a scalar row, not the
complete diffeomorphism/internal-gauge master equation.
\end{theorem}
\begin{proof}
Apply Gaussian integration by parts to
\(G_i=\partial^L_{\upsilon_i}F\) times the interaction weight:
\[
 \langle\chi_jG_i\rangle_V
 =\hbar\mathsf G_{jk}\langle\partial_{\phi_k}G_i\rangle_V
  -a^4\mathsf G_{jk}\langle v_kG_i\rangle_V.
\]
The differentiating variable is even, so this equation makes no hidden
odd reordering. Multiply by \((\widetilde P_a)_{ij}\), sum, and use
\(\widetilde P_a\mathsf G=a^{-4}\nu\). Add the physical force term.
The difference between it and \(\nu v\) is precisely
\(\Theta_\lambda v\), proving \eqref{eq:v57-conditional-identity}.
The formal proof is an identity of Gaussian moments at each coefficient;
it needs no uniform nonlinear-integral remainder. Source dependence of
\(V\) is retained, since only its scalar derivative was used.
\end{proof}

For \(V=0\), this gives the evaluated balance
\eqref{eq:v57-pair-balance}. For the existing full interaction it records
the scalar force and its actual lower-reference term, rather than
postulating that the scalar odd contact is the entire breaking. The
right-hand side is not the already assembled full supertrace
\(\mathcal X\); it is its independently specified conditional scalar
integration-by-parts counterpart. No unproved equality between these
objects is asserted.

\section{Transporting the complete contact section through the Gaussian}
\label{sec:v57-transported-section}

Fix V56's support \(S\), existing translated mode \(w\), and
contraction kernel \(\nu\). Its odd section is multiplication by
\(U_x\); write it as \(H_x\). The Gaussian covariance is still
\eqref{eq:v57-covariance}. There is generally no reason for
\(\mathsf G_{S,T}w\) to vanish: V56 imposes orthogonality for
\(\nu\), not for \(\nu\widetilde P_a^{-1}\).

Define the odd differential operator
\begin{equation}
 \mathscr V_x r
 =\frac{a^4}{4\kappa}\sum_{A,y\in x+aS}
    \zeta_A(x)(\mathsf Gw_x)_y\,
                \partial_{\phi_A(y)}r,
 \qquad
 (w_x)_y=\sum_jw_j\mathbf1_{y=x+at_j}.
 \label{eq:v57-cross-return}
\end{equation}
The displayed order of \(\zeta_A\) is binding for an odd input.

\begin{theorem}[The unique Gaussian-conjugated contact section]
\label{thm:v57-conjugated}
On densities supported in \(x+aS\),
\begin{equation}
 \boxed{\begin{aligned}
 \mathscr W H_x&=(H_x+\hbar\mathscr V_x)\mathscr W,\\
 H_x^{\mathsf G}:=\mathscr W^{-1}H_x\mathscr W
               &=H_x-\hbar\mathscr V_x,\\
 \Delta_\nu H_x^{\mathsf G}+H_x^{\mathsf G}\Delta_\nu&=I.
 \end{aligned}}
 \label{eq:v57-conjugated}
\end{equation}
All coefficients use the original finite sets \(S,T\).
The covariance is read only at their finitely many differences; it does
not introduce new field sites. The correction is linear in \(\hbar\)
and reduces the scalar/source field degree by two relative to \(H_x\),
so V1's total field-plus-loop weight is preserved. The section is not a
bracket morphism.
\end{theorem}
\begin{proof}
In \([\mathscr D_{\mathsf G},U_x]\), the scalar factor of \(U_x\)
can be differentiated once, giving exactly \(\mathscr V_x\).
There is no scalar second derivative of \(U_x\). The coefficients of
\(\mathscr V_x\) contain only antifields and constants, so
\([\mathscr D_{\mathsf G},\mathscr V_x]=0\). The two exponential
conjugations therefore terminate after this commutator. V56 gives
\(\Delta_\nu H_x+H_x\Delta_\nu=I\) on the old support.
The operator \(\mathscr W\) and its inverse preserve that support
and commute with \(\Delta_\nu\); conjugating proves the last identity.
Uniqueness means uniqueness of the conjugated operator satisfying
\(\mathscr W H_x^{\mathsf G}=H_x\mathscr W\), not uniqueness among
all possible contraction sections.
\end{proof}

\begin{corollary}[Contact matching stable under the actual free scalar step]
\label{cor:v57-joint-matching}
For a supplied density \(c\) and target \(t\) on the old support, put
\(r=t-\Delta_\nu c\) and
\begin{equation}
 c^{\mathsf G}=c+H_x^{\mathsf G}r.
 \label{eq:v57-repair}
\end{equation}
If \(\Delta_\nu t=0\), then
\begin{equation}
 \boxed{\Delta_\nu c^{\mathsf G}=t,
 \qquad
 \mathscr W(c^{\mathsf G}-c)=H_x\mathscr Wr.}
 \label{eq:v57-joint-matching}
\end{equation}
Every coefficient of the \emph{measured added source} has zero added-pair
Taylor jet through degree \(2L-1\). Its odd contact is the transported
target \(\mathscr Wt\), not an unchanged nonlinear \(t\).
If the target fails closure, its exact raw residual is
\(-H_x^{\mathsf G}\Delta_\nu t\). For a separately prescribed measured
target \(t_{\mathrm{out}}\), the raw target must be
\(\mathscr W^{-1}t_{\mathrm{out}}\).
\end{corollary}
\begin{proof}
Apply the last identity of \eqref{eq:v57-conjugated} to \(r\), using
\(\Delta_\nu r=\Delta_\nu t\). The second identity follows by
intertwining. Differentiation cannot create new support in
\(\mathscr Wr\), so the two independent moment-zero factors of
\(H_x\) remain present on every output monomial. Finally
\(\Delta_\nu\mathscr W=\mathscr W\Delta_\nu\) gives the target
statement. The same computation without setting \(\Delta_\nu t=0\)
gives the residual.
\end{proof}

There is a substantive loop-order distinction. V56's multiplication
section has its high-order jet at tree level. Its Gaussian image has
an additional \(\hbar\mathscr V_x\mathscr Wr\) contraction, with
only one guaranteed moment-zero factor. The present section cancels that
specific return. Its raw \(\hbar\)-coefficient can have lower momentum
jets; these are the required lower-field, higher-loop source partners,
not a claim of unchanged raw jets at every loop index. No new mode design
or assumption of covariance orthogonality is used.

\section{Finite-degree source estimates and covariance composition}
\label{sec:v57-bounds}

Let a supplied rooted polynomial have scalar degree at most \(n\),
fixed scalar/antifield support, and coefficient norm \(\|F\|_1\).
Positive semidefiniteness gives
\(|\mathsf G_{xy}|\le u_a\). Counting unordered derivative pairs yields
\begin{equation}
 \boxed{
 \left\|[\hbar^j]\mathscr W^{\pm1}F\right\|_1
 \le
 \frac{n!}{2^j j!(n-2j)!}\,u_a^j\|F\|_1,
 \quad 0\le2j\le n.}
 \label{eq:v57-finite-Wick-bound}
\end{equation}
For an inhomogeneous polynomial, apply this to each homogeneous part
and sum. The same bound includes flavour sums by counting all scalar
occurrences together. Each self-contraction is bounded by an actual
coincidence covariance; an extensive supremum is not introduced.

\begin{proposition}[Bounds for the transported section]
\label{prop:v57-section-bound}
With \(N\) the number of mode sites and \(n\) the maximal scalar degree
of the input,
\begin{equation}
 \|H_xF\|_1\le2Na^4\|F\|_1,
 \qquad
 \|\mathscr V_xF\|_1
 \le\frac{Nn}{2}a^4u_a\|F\|_1
 \le C Nn a^2\|F\|_1.
 \label{eq:v57-H-bound}
\end{equation}
If the raw input obeys V56's budget
\(a^4\|r\|_1\le B a^{-\delta}L_a\), at fixed scalar/source degree,
then the two raw pieces have fixed-window rooted bounds
\begin{equation}
 \|H_xr\|_{\mathrm{src}}
 \le C B L_a a^{2L-\delta}\epsilon_*^{2L},
 \qquad
 \|\mathscr V_xr\|_{\mathrm{src}}
 \le C B L_a a^{L-\delta-2}\epsilon_*^L.
 \label{eq:v57-two-raw-bounds}
\end{equation}
The second bound is not claimed to vanish for \(L=3,\delta=4\).
If \(a^4[\hbar^j]\mathscr Wr\) instead has its own actual budget
\(B_j a^{-\delta_j}L_a\), the measured correction obeys
\begin{equation}
 \boxed{
 \|[\hbar^j]\mathscr W(c^{\mathsf G}-c)\|_{\mathrm{src}}
 \le C_jB_jL_a a^{2L-\delta_j}\epsilon_*^{2L}.}
 \label{eq:v57-measured-bound}
\end{equation}
These estimates hold for every prescribed rooted kernel moment and
scaled momentum derivative, with no total-volume multiplier.
\end{proposition}
\begin{proof}
V56 gives \(\|w\|_1^2/\kappa\le2N\). For \(H\), sum the
four flavour copies. For \(\mathscr V\), the sum of all scalar
derivative coefficient norms is at most \(n\|F\|_1\), and
\(|(\mathsf Gw_x)_y|\le u_a\|w\|_1\). The factor
\(a^4/(4\kappa)\) then proves \eqref{eq:v57-H-bound}.
For soft kernels, \(H\) has two order-\(L\) factors; \(\mathscr V\)
has one. Taylor remainders for their fixed finite supports, followed by
the coefficient budget, give \eqref{eq:v57-two-raw-bounds}.
The exact identity \eqref{eq:v57-joint-matching}, not the triangle
inequality on those two raw pieces, gives
\eqref{eq:v57-measured-bound}. Multiplication by fixed external windows
and integration by parts in all relative momenta give the rooted bounds.
Only fixed finite supports are summed, so periodization introduces no
unrooted volume factor.
\end{proof}

The elementary bound \eqref{eq:v57-finite-Wick-bound} permits
\(\delta_j=\delta+2j\), with explicit degree-dependent constants;
a better populated budget can be used when proved. Keeping the original
power of \(a\) after every self-contraction would be incorrect.
Mass derivatives differentiate the same covariance coefficients, with
\eqref{eq:v57-mass-moments} and the corresponding finite-entry bounds.
No mass-independent approximation to the propagator is made.

\begin{corollary}[A finite loop budget for measured moment matching]
\label{cor:v57-loop-budget}
Suppose the target gap is independent of \(\hbar\), has scalar degree
at most \(n\), and satisfies the raw budget in
\cref{prop:v57-section-bound}. Fix \(J\le\lfloor n/2\rfloor\).
If the inherited mode order is chosen once so that
\(2L\ge\delta+2J+2\), then every coefficient through \(J\) scalar
self-contractions in the measured correction is
\(O(a^2 L_a\epsilon_*^{2L})\), with its own fixed-degree constant.
The target remains exactly matched before transport and is transported
by the same Gaussian afterward. No new sites beyond the finite support
chosen by V56 are required during refinement.
\end{corollary}
\begin{proof}
Use \eqref{eq:v57-finite-Wick-bound} and \(u_a\le Ca^{-2}\) in the
raw budget. This gives \(\delta_j=\delta+2j\). The exponent in
\eqref{eq:v57-measured-bound} is at least two for \(j\le J\).
The source degree, mode order and covariance seminorms remain fixed.
The assertion is about these marked Gaussian coefficients only, not
all interacting diagrams at loop order \(J\).
\end{proof}

For constant covariance matrices
\(\mathsf G=\mathsf G_1+\mathsf G_2\), the two heat operators
commute and
\(\mathscr W_{\mathsf G}=\mathscr W_{\mathsf G_2}
\mathscr W_{\mathsf G_1}\). Consequently the normal-ordering partners
compose in the opposite order, with \emph{all} mixed self-contractions
retained. For example the \(u_av_a\) term in
\eqref{eq:v57-r3} contains
\(u_1v_2+u_2v_1\) as well as \(u_1v_1+u_2v_2\).
Assigning only diagonal shell pairs changes the coefficient. This is an
application of inherited Gaussian composition to the newly evaluated
source, not an interacting many-shell theorem.

\section{Status, nonduplication audit and the remaining master coefficient}
\label{sec:v57-ledgers}

\subsection*{Proved}
The complete free scalar Gaussian image of V56's two-term nonlinear
current-ascent density is evaluated, including its positive order-
\(a^{-4}\) one-contraction mixing and order-\(a^{-6}\) double return.
The two native covariance sums have explicit uniform bounds and positive
leading finite-integral limits. The scalar Euler/contact/reference balance
is proved conditionally inside the existing measured integral without
assuming finite diffeomorphism invariance. V56's contact section has an
explicit one-loop conjugate, which matches its specified raw contact and
preserves the full high-moment pair in the Gaussian-measured correction.
Finite-degree source bounds retain the new covariance growth and every
mixed shell pair.

\subsection*{Inherited}
V53 supplies the source density and its ordinary descent, V54 its
finite-source interpretation, V55 the scalar measure and loop bookkeeping,
and V56 its declared source differences and mask-orthogonal pair. V27's
completed scalar covariance and inverse bounds are unchanged. The general
Wick theorem, Gaussian composition and BV product rule are not new. Earlier
whole-theory one-loop estimates and their local cohomological inputs are
not reaudited or replaced here.

\subsection*{Literature input}
Finite Gaussian/BV transfer and the need for matched cutoff/reference
conventions are comparison context \cite{V57DJP,V57IIM}. No cohomology
vanishing or arbitrary-loop theorem is used to deduce an unknown native
coefficient or a full critical primitive.

\subsection*{Conditional}
The contact target and its nilpotence condition are supplied, not fitted
from the output. The Gaussian-conjugated realization applies to one marked
polynomial insertion in the specified free scalar reference; a Gaussian
step inside an interacting calculation must retain that calculation's
interaction vertices and connected insertions. Equation
\eqref{eq:v57-conditional-identity} keeps them symbolically and exactly,
but supplies no new interacting expectation bound. The native covariance
moments are specified analytically; their profile-dependent numerical
values have not been computed. Higher-loop source locality requires its
actual coefficient budget at that loop order.

\subsection*{Open}
The full native invariant critical breaking still needs the representative
and populated descent decomposition identified in V53. The present
calculation is a genuine hard scalar self-contraction of an inherited
source, not a substitute for that unknown breaking. Inserting a one-loop
source correction changes the next complete master coefficient through
\((S_1,E)+\tfrac12(E,E)-\Delta_\rho E\), and through the genuine
interacting diagrams. Neither a matched odd contact nor a normal-ordered
scalar insertion sets that entire coefficient to zero. Metric/ghost,
internal-gauge/chiral and other measure terms remain required. No active
normalization is changed. Further mode construction is not needed for
this interface; the next step is the populated full interacting
source/reference coefficient with these evaluated scalar returns included.

\begin{center}\small
\begin{tabular}{p{.24\textwidth}p{.68\textwidth}}
\toprule
Variable & Exact scope of V57\\\midrule
Carrier & Same four real scalars, native free \(\widetilde P_a\),
source stencils, contour and spectator labels; no new species.\\
Cutoff and volume & Moment bounds uniform on fixed profile and lower-scale
margins; rooted bounds have no total-volume factor.\\
Mass & \(0\le m\le M\lambda\), including zero; \(\lambda>0\)
retained; differentiated covariances carry one mask.\\
Loops & All self-contractions of a fixed polynomial and one marked
insertion; not all diagrams at the corresponding interacting loop order.\\
Source order & Every fixed scalar/antifield degree; constants retain
support, stencil and degree dependence.\\
Symmetry & Exact scalar conditional integration-by-parts equation;
not a matched full BV fibration or a finite diffeomorphism QME.\\
Verification & Rational polynomial/Grassmann tests, exact model
covariances and native-symbol identities; no numerical native loop value
or proof-assistant formalization.\\
\bottomrule
\end{tabular}
\end{center}

\needspace{16\baselineskip}
\begingroup
\renewcommand{\refname}{References for V57}
\begin{thebibliography}{9}\small
\bibitem{V57DJP} M.~Doubek, B.~Jur\v co and J.~Pulmann,
\emph{Quantum \(L_\infty\) algebras and the homological perturbation lemma},
\emph{Commun. Math. Phys.} \textbf{367} (2019), 215--240,
arXiv:1712.02696. Finite BV/Feynman transfer context, not an estimate of
the present native covariance moments.
\bibitem{V57Schwarz} A.~Schwarz,
\emph{Geometry of Batalin--Vilkovisky quantization},
\emph{Commun. Math. Phys.} \textbf{155} (1993), 249--260,
arXiv:hep-th/9205088. Finite odd-Laplacian framework.
\bibitem{V57IIM} Y.~Igarashi, K.~Itoh and T.~R.~Morris,
\emph{BRST in the Exact RG}, arXiv:1904.08231 (2019).
Comparison for matched quantum and lower-reference identities;
not a complete native Einstein--SM quantum closure theorem.
\end{thebibliography}
\endgroup

\section*{Append-only continuation marker: V57}
\addcontentsline{toc}{section}{Append-only continuation marker: V57}
\label{sec:v57-continuation-marker}
\textbf{End of V57.} Preserve Sections 1--446. The hard scalar
self-contractions of the inherited nonlinear source and the Euler/contact
balance are now explicit. The same V56 pair, with its Gaussian-conjugated
loop partner, transports compatible contact matching without losing the
measured soft moments. The unidentified native critical descent and the
complete interacting master equation remain open. No active action,
normalization, propagator, or measure is changed.


\clearpage
\part*{V58: Geometry-conditioned scalar sources and the masked Euler defect}
\addcontentsline{toc}{part}{V58: Geometry-conditioned scalar sources and the masked Euler defect}

\section{From the free scalar insertion to its actual metric interaction}
\label{sec:v58-start}

Sections 1--446, including their continuation markers, are preserved.
V57 evaluates the nonlinear ascent source on the free scalar reference.
The next interaction already present in the same action is
\(\frac12\langle\phi,[Q_a(q)-\widetilde P_a]\phi\rangle\).
This part includes that interaction before computing the marked source.
It does not add a scalar self-coupling, alter the gravitational contour,
or choose a new reference profile. The resulting conditional integral is
Gaussian only because the retained coupling component has zero internal
gauge, Yukawa and Higgs self-interaction couplings. It still contains all
of the declared metric dependence of the scalar action.

Three distinctions govern the calculation. First, normalizing the scalar
conditional expectation does not permit discarding its metric-dependent
partition factor in the remaining integral. Second, a nonconstant metric
creates scalar--derivative covariances which vanish on the translation-
invariant reference of V57. Third, the physical scalar Euler operator no
longer commutes with the fixed lower mask. The last fact makes the scalar
Euler/contact operator curved: its square is explicitly nonzero. It is
not a new anomaly or an obstruction to the complete modified BV identity.

Gaussian completion and BV transfer are standard machinery
\cite{V58DJP,V58IIM}. The additions here are the complete conditional image
of the existing nonlinear source, an evaluated native metric response,
the first-mark cutoff estimates, and the exact scalar reference and
source-derivative operators. The unidentified invariant critical breaking
is not supplied by this calculation.

\subsection*{Conventions, inherited inputs and scope}
Put \(v=a^4\), \(P=\widetilde P_a\), \(z=m^2\), and
\(\nu=I-\Theta_\lambda\). Matrix products act on coordinate values, whereas
the action pairing is \(v\sum_x\). The scalar covariance as an operator
and as a coordinate kernel is respectively
\[
 C=\nu P^{-1},\qquad \mathsf G=v^{-1}C.
\]
The quotient is zero on retained modes, including the zero mode at
\(z=0\). Let \(Q(q)=Q_a(q)\) be exactly the Hessian of
\eqref{eq:v27-completed-Higgs}; set \(V(q)=Q(q)-P\). Every transpose below
is algebraic, as in the inherited coefficient convention. There are four
identical scalar flavours. Other species, ghosts, antifields and external
source labels remain in their original order.

Identities are literal on the finite real chart where the scalar
quadratic form on its integration support is positive, and analytic
there; their formal metric jets also apply in the declared complex
coefficient chart. This does not establish a globally convergent
nonlinear gravitational integral. Metric jets used in the quantitative
comparison are at \(q=0\), with one arbitrary symmetric metric mark.
For an explicitly geometry-dependent insertion, its own metric derivative
is added in every formula below. The example \(c_x^{(r)}\) is precisely
V56--V57's fixed-coordinate polynomial insertion, not an unrecorded
covariantization of it.

\section{The normalized scalar map and its retained partition factor}
\label{sec:v58-conditional}

Define, without inverting a zero covariance,
\begin{equation}
 \boxed{
 A(q)=(I+CV(q))^{-1},\quad C_q=A(q)C=C A(q)^T,
 \quad m(q)=A(q)\varphi,\quad \mathsf G_q=v^{-1}C_q.}
 \label{eq:v58-resolvents}
\end{equation}
The symbol \(m(q)\) in this section is the conditional mean, not the
scalar mass. The mass appears only as \(z\) or \(m_H=\sqrt z\) below.
For a polynomial insertion \(F\), let
\begin{equation}
 \mathscr T_qF=\frac{\E_{\mathsf G}
 [F(\varphi+\chi,\upsilon)e^{-\frac v{2\hbar}
 \sum_A(\varphi_A+\chi_A)^TV(q)(\varphi_A+\chi_A)}]}
 {\E_{\mathsf G}[e^{-\frac v{2\hbar}
 \sum_A(\varphi_A+\chi_A)^TV(q)(\varphi_A+\chi_A)}]}.
 \label{eq:v58-normalized-map}
\end{equation}
Here \(\chi\) has coordinate covariance \(\hbar\mathsf G\).

\begin{theorem}[Exact conditional source and unnormalized factor]
\label{thm:v58-conditional}
On the analytic germ at \(q=0\),
\begin{equation}
 \boxed{\mathscr T_qF
 =\left[e^{\hbar\mathscr D_{\mathsf G_q}}F(\phi,\upsilon)
 \right]_{\phi=m(q)},\qquad
 \mathscr D_{\mathsf G_q}
 =\frac12\sum_{A,x,y}(\mathsf G_q)_{xy}
 \partial_{\phi_A(x)}\partial_{\phi_A(y)}.}
 \label{eq:v58-Wick}
\end{equation}
\begin{equation}
 Z_H(q,\varphi)
 =\det(I+CV(q))^{-2}
 \exp\left[-\frac v{2\hbar}
 \sum_A\varphi_A^TV(q)A(q)\varphi_A\right].
 \label{eq:v58-partition}
\end{equation}
The determinant branch is fixed by its value one at the origin. The
matrix \(V(q)A(q)\) is symmetric. Both formulas are exact at every
prescribed formal metric and source degree. A sufficient literal domain
is \(\|C^{1/2}VC^{1/2}\|<b<1\) for real symmetric \(V\).
On that domain
\[
 (1+b)^{-1}C\preceq C_q\preceq(1-b)^{-1}C.
\]
No positive lower eigenvalue of \(C\) on the retained space is assumed.
\end{theorem}
\begin{proof}
Write the fluctuation as \(\chi=v^{-1/2}C^{1/2}\xi\) on
\(\operatorname{Ran}C\), with \(\xi\) a standard Gaussian of variance
\(\hbar\). Completing that square gives covariance
\(C^{1/2}(I+C^{1/2}VC^{1/2})^{-1}C^{1/2}=C_q\)
in the operator convention and mean \((I-C_qV)\varphi=A\varphi\).
The determinant on the integration support equals the displayed full
determinant, since the added retained eigenvalues are one. Four flavours
give the power \(-2\). The remaining quadratic form is
\(V-VC_qV=VA\), which is symmetric. Wick's finite polynomial formula
gives \eqref{eq:v58-Wick}. The spectral inequality on the support gives
the covariance inequalities by congruence; analytic continuation and
formal differentiation give the jet statements.
\end{proof}

The partition factor stays in the remaining measured integral:
\begin{equation}
 Z_{\rm full}[F]
 =\int\dd\mu_{\rm rest}\ e^{-S_{\rm rest}/\hbar}
       Z_H(q,\varphi)\,\mathscr T_qF.
 \label{eq:v58-full-integral}
\end{equation}
This is coefficientwise Fubini for the inherited finite expansion, and
literal Fubini whenever the full integral exists. The metric/ghost
measure, contour and once-counted auxiliary density are not removed by
the conditional normalization. The scalar determinant already present
in \eqref{eq:v58-partition} must not be added again.

\begin{proposition}[All three geometric variations]
\label{prop:v58-variation}
For a metric direction \(h\), put \(Q_h=D_qQ(q)[h]\). Then
\begin{equation}
 \boxed{D_hm=-C_qQ_hm,\qquad
 D_hC_q=-C_qQ_hC_q,}
 \label{eq:v58-mean-cov-variation}
\end{equation}
and
\begin{equation}
 D_h(-\hbar\log Z_H)
 =\frac v2\sum_A m_A^TQ_hm_A+2\hbar\Tr(C_qQ_h).
 \label{eq:v58-det-variation}
\end{equation}
For an arbitrary insertion,
\begin{equation}
 D_h\mathscr T_qF
 =\mathscr T_q(D_hF)
 -\frac v{2\hbar}\sum_A
 \Cov_q\bigl(F,\phi_A^TQ_h\phi_A\bigr).
 \label{eq:v58-connected}
\end{equation}
The covariance is the ordered connected scalar expectation with the even
quadratic vertex; no odd coefficients are commuted in this equation.
\end{proposition}
\begin{proof}
Differentiate \((I+CV)A=I\), with \(C\) fixed. This gives
\(D_hA=-C_qQ_hA\) and both identities in
\eqref{eq:v58-mean-cov-variation}. Differentiate the unnormalized
integral before dividing by it. Its logarithmic derivative is the
conditional expectation of the quadratic vertex, which splits into the
mean and covariance terms in \eqref{eq:v58-det-variation}.
The quotient rule yields \eqref{eq:v58-connected}; its second term is
exactly the disconnected determinant contribution subtracted from the
marked insertion.
\end{proof}

For finitely many soft metric marks and scalar inputs whose total
momentum remains strictly below \(\lambda\), every corresponding
nonconstant Taylor coefficient of \(m(q)\) vanishes: its leftmost factor
is \(C\) at a sum of those soft momenta. In particular
\(D_hm|_0=0\) when both \(h\) and \(\varphi\) have momenta at most
\(\mu\) and \(2\mu<\lambda\). This support fact does not remove the
mean shift for general metric configurations, including hard metric
arguments subsequently integrated in \eqref{eq:v58-full-integral}.

\section{The complete nonlinear insertion with a nonconstant metric}
\label{sec:v58-source}

Use the unchanged finite source operators \(D_\mu,P_a^{\rm src}\)
from V56. At the source root \(x\), set
\begin{equation}
 \begin{gathered}
 X=m_x,\qquad Y_\mu=(D_\mu m)_x,\qquad Z=(P_a^{\rm src}m)_x,\\
 u=(\mathsf G_q)_{xx},\quad
 b_\mu=(\mathsf G_qD_\mu^T)_{xx},\quad
 v_s=(\mathsf G_qP_a^{\rm src})_{xx},\quad
 d_\mu=(D_\mu\mathsf G_qP_a^{\rm src})_{xx}.
 \end{gathered}
 \label{eq:v58-four-moments}
\end{equation}
The subscript on \(v_s\) distinguishes this source covariance from the
cell volume \(v=a^4\). The antifields and ghost factors remain at their
original sites. Let \(M_r\) be V57's moment polynomial
\eqref{eq:v57-moment-polynomial}; negative-index terms below are omitted.

\begin{theorem}[All conditional contractions of the actual two-term source]
\label{thm:v58-source}
For \(c_x^{(r)}\) in \eqref{eq:v57-c}, at every fixed integer \(r\ge1\),
\begin{equation}
\boxed{\begin{aligned}
 \mathscr T_qc_x^{(r)}={}&
 -\frac12 f_B\sum_\mu
  [Y_\mu M_r+\hbar r b_\mu M_{r-1}]
                    \upsilon(D_\mu\upsilon)\\
 &+\frac12\sum_\mu\alpha_\mu
 \Big[(Y_\mu Z+\hbar d_\mu)M_r
   +\hbar r(b_\mu Z+v_sY_\mu)M_{r-1}\\
 &\hspace{39mm}
   +\hbar^2r(r-1)b_\mu v_sM_{r-2}\Big]\upsilon,
 \qquad M_j=M_j(X,\hbar u).
\end{aligned}}
\label{eq:v58-full-source}
\end{equation}
This includes every conditional scalar contraction, every number of
metric insertions in the declared germ, and both original source terms.
There is no flavour factor four for one fixed external flavour.
\end{theorem}
\begin{proof}
The three scalar random variables at one summand are \(\phi_x\),
\((D_\mu\phi)_x\), and \((P_a^{\rm src}\phi)_x\). Their means are
\(X,Y_\mu,Z\), and their relevant covariance entries are
\(\hbar u,\hbar b_\mu,\hbar v_s,\hbar d_\mu\).
Each of the last two variables occurs at most once. Differentiating their
joint Gaussian exponential gives
\[
 \E[\phi_x^r(D_\mu\phi)_x]
 =Y_\mu M_r+\hbar b_\mu\partial_XM_r,
\]
\[
 \E[\phi_x^r(D_\mu\phi)_x(P_a^{\rm src}\phi)_x]
 =(Y_\mu Z+\hbar d_\mu)M_r
 +\hbar(b_\mu Z+v_sY_\mu)\partial_XM_r
 +\hbar^2b_\mu v_s\partial_X^2M_r.
\]
Use \(\partial_XM_r=rM_{r-1}\) and preserve the displayed Grassmann
order. No Gaussian contraction acts on an antifield.
\end{proof}

At \(q=0\), \(b_\mu=d_\mu=0\) by the inherited parity, and
\eqref{eq:v58-full-source} reduces exactly to V57. At a nonconstant
metric these two covariances cannot be set to zero. The covariance
\(b_\mu\) creates a two-antifield contact in the first line; \(d_\mu\)
creates a lower-scalar-degree term in the second. They are part of the
same normalized conditional source, not freely adjustable additions.

\subsection*{One metric insertion, evaluated before gravity integration}
Write a dot for \(D_q[\,h\,]|_{q=0}\) and put
\(B_h=D_qQ_a(0)[h]\). All four new coefficients are exact finite sums:
\begin{equation}
\boxed{\begin{aligned}
 \dot u&=-v^{-1}(CB_hC)_{xx},&
 \dot b_\mu&=-v^{-1}(CB_hCD_\mu^T)_{xx},\\
 \dot v_s&=-v^{-1}(CB_hCP_a^{\rm src})_{xx},&
 \dot d_\mu&=-v^{-1}(D_\mu CB_hCP_a^{\rm src})_{xx}.
\end{aligned}}
\label{eq:v58-marked-moments}
\end{equation}
The matrix \(B_h\) is the actual vertex of
\eqref{eq:v27-actual-V}, with the two scalar legs put into
input/output rather than all-incoming convention. There is no scalar
replacement of that vertex. For one soft metric mark and soft scalar
inputs, the mean-shift derivative vanishes as noted above. At \(r=1\),
\begin{equation}
\boxed{\begin{aligned}
 [\hbar]D_h\mathscr T_qc_x^{(1)}\big|_0={}&
 -\frac12 f_B\sum_\mu\dot b_\mu\upsilon D_\mu\upsilon\\
 &+\frac12\sum_\mu\alpha_\mu
 [\dot v_sD_\mu\varphi+\dot b_\mu P_a^{\rm src}\varphi
                     +\dot d_\mu\varphi]\upsilon.
\end{aligned}}
\label{eq:v58-r1-marked}
\end{equation}
This is the complete connected scalar-loop response of this insertion
to one native metric interaction. It is not the entire metric--ghost
master row; no hard metric has yet been integrated here.

\begin{proposition}[A strictly nonzero native conformal source response]
\label{prop:v58-positive}
Take a constant first metric mark \(h=I_4\). Put
\(p_n=P_a^n-z\), \(p_b=P_a^b-z\), and use the same filter \(s\)
as V27. The exact scalar vertex is
\begin{equation}
 B_I(p)=2z+p_b(p)+s(ap)^2[p_n(p)-p_b(p)].
 \label{eq:v58-conformal-vertex}
\end{equation}
Then \(\dot b_\mu=\dot d_\mu=0\), and
\begin{equation}
 \boxed{\dot v_s=-w_a,\qquad
 w_a=\frac1{V_L}\sum_p
 \frac{\nu(p)^2P_a^{\rm src}(p)B_I(p)}{P(p)^2}>0.}
 \label{eq:v58-positive-sum}
\end{equation}
On a sufficiently small fixed \(a\lambda\) domain with
\(\lambda L_\mu\ge1\), fixed profiles and \(0\le\sqrt z\le M\lambda\),
there are positive constants independent of mesh and periods such that
\[
 c a^{-4}\le w_a\le C a^{-4}.
\]
For fixed physical periods, profiles and mass,
\begin{equation}
 a^4w_a\longrightarrow
 \int_{[-\pi,\pi]^4}\frac{\dd^4y}{(2\pi)^4}
 \frac{p_s(y)b_\star(y)}{p_\star(y)^2}>0,
 \quad b_\star=p_b^\star+s^2(p_n^\star-p_b^\star).
 \label{eq:v58-positive-limit}
\end{equation}
Here stars mean dimensionless massless kinetic symbols, exactly as in
V57; this integral is not assigned a numerical value.
\end{proposition}
\begin{proof}
The inherited first coframe variation is \(D\rho[I]=2\) and
\(DW[I]=I\). Differentiate all three interaction terms in
\eqref{eq:v27-completed-Higgs}. The metric filter at zero momentum is
one and the two scalar filters give \(s^2\); the free adjustment has
no metric derivative. This proves \eqref{eq:v58-conformal-vertex}.
It is a positive convex combination of the two kinetic symbols plus
\(2z\). The parity of its diagonal symbol gives the two zeros.
Equation \eqref{eq:v58-marked-moments} gives the finite sum, with
\emph{two} covariance masks. Scalar ellipticity and the source-symbol
bounds of V57 bound its ratio above and below on high momenta. The
fraction of modes in the lower ball is at most \(C(a\lambda)^4\);
outside it \(\nu=1\). The resulting normalized mode count proves the
two-sided estimate. On a fixed dimensionless ball about zero the
integrand ratio is bounded, uniformly as \(a\to0\), and the ball's
normalized mass is \(O(\epsilon^4)\). Riemann convergence on its
complement, then \(\epsilon\downarrow0\), proves the limit.
\end{proof}

In particular the conformal first-mark term in
\eqref{eq:v58-r1-marked} is
\(-w_a\sum_\mu\alpha_\mu D_\mu\varphi\,\upsilon/2\).
Its nonzero value is a calculated interaction-source coefficient. It is
local matching data, not a physical scalar anomalous dimension and not
a nonzero anomaly.

\section{Full-cutoff comparison of the first metric-marked moments}
\label{sec:v58-bounds}

For a Fourier metric mark \(h(x)=H e^{ikx}\), each quantity in
\eqref{eq:v58-marked-moments} is its same Fourier mode times a
coefficient linear in \(H\). Let \(F_{a,L,I}(k,z)[H]\) denote that
coefficient. Assign the following superficial Taylor degrees:
\begin{center}
\begin{tabular}{c c c c}
\toprule
Moment & Symbol order of loop integrand & Taylor degree \(d_I\) & Momentum parity\\\midrule
\(\dot u\)&\(-2\)&2&even\\
\(\dot b_\mu\)&\(-1\)&3&odd\\
\(\dot v_s\)&0&4&even\\
\(\dot d_\mu\)&1&5&odd\\\bottomrule
\end{tabular}
\end{center}
The Taylor operator is the existing weighted operator
\(\mathbf T_{d_I}\) in \((k,z)\), with \(\deg k=1,\deg z=2\).
Every local coefficient is stored. No subtraction is imposed on just
one factor of a later multiple-contraction product.

\begin{theorem}[An interacting first-mark source comparison]
\label{thm:v58-cutoff}
Fix the profiles and source operators above, \(M<\infty\), and a
finite derivative and kernel-moment order. On
\[
 \lambda\ge32768\mu>0,\quad a\lambda\le2^{-24},\quad
 \lambda L_\nu\ge1,\quad |k|\le\mu,\quad0\le\sqrt z\le M\lambda,
\]
put \(\epsilon=|k|+\sqrt z\),
\(r=|\alpha|+2j\), and \(F^R=(I-\mathbf T_{d_I})F\).
Then, for \(r\le d_I+2\),
\begin{equation}
 \boxed{
 |\partial_k^\alpha\partial_z^j(F^R_{a,L,I}-F^R_{0,L,I})[H]|
 \le C_{I,\alpha,j}\|H\|a^2\epsilon^{d_I+2-r}.}
 \label{eq:v58-cutoff}
\end{equation}
All completed Brillouin-zone modes are included. The zero-mesh quantity
on the same torus uses the same lower mask. Its thermodynamic comparison
has the separate bound, for every prescribed \(b>4\),
\begin{equation}
 |\partial_k^\alpha\partial_z^j(F^R_{0,L,I}-F^R_{0,\infty,I})[H]|
 \le C_{I,\alpha,j,b}\|H\|\lambda^{-2}
 \epsilon^{d_I+2-r}(\lambda L_{\min})^{-b}.
 \label{eq:v58-volume}
\end{equation}
Higher fixed derivatives are controlled by the corresponding differentiated
symbol bounds. There is no claim of uniformity in unbounded mark,
perturbative or derivative order, nor removal of \(\lambda\).
\end{theorem}
\begin{proof}
In momentum coordinates each summand contains two actual scalar
covariances and one \(B_h\). The other factors are the displayed fixed
source symbols. With \(R=\lambda+|p|\), an internal or external momentum
derivative lowers the order by one, while a mass derivative lowers it
by two. Proposition~\ref{prop:v27-scalar-reference} and V56's fourth-
order source differences give, on the common interior chart,
\[
 |\partial^{\alpha}_{k,p}\partial_z^j f_{a,I}|+
 |\partial^{\alpha}_{k,p}\partial_z^j f_{0,I}|
 \le C\|H\|R^{d_I-4-|\alpha|-2j},
\]
\[
 |\partial^{\alpha}_{k,p}\partial_z^j(f_{a,I}-f_{0,I})|
 \le C\|H\|a^4R^{d_I-|\alpha|-2j}.
\]
Mass derivatives also act on \(B_h\) and \(P_a^{\rm src}\); these are
included in the order statements. Each individual inverse derivative
retains its own mask once. The two scalar lines are not merged into a
single cutoff factor.

Reflection of every momentum reverses only the odd source differences.
After reflecting the integration variable, this proves the parity in
the table without any rotational averaging. Since \(z\) has even
weight, the local Taylor jet of weight \(d_I+1\) is zero. Thus
\(\mathbf T_{d_I+1}F=\mathbf T_{d_I}F\).

Apply Taylor's formula to \(F(k,m_H^2)\) as a function of \((k,m_H)\),
which is even in \(m_H\). The first possible complementary weight is
\(d_I+2\). At that weight the ordinary differentiated integrand has
order \(R^{-6}\), and the interior difference is \(a^4R^{-2}\).
The lower mask removes the infrared singularity. Normalized periodic
mode counting gives
\[
 a^4\sum_{\lambda\lesssim|p|\lesssim c/a}^{\mathrm{per\ volume}}
 R^{-2}\le Ca^2,\qquad
 \sum_{|p|\gtrsim c/a}^{\mathrm{per\ volume}}R^{-6}\le Ca^2.
\]
The same argument applies after any prescribed internal derivatives.
Smooth periodicity of the completed symbols yields the upper-region
bound: \(d_I+2\) weighted derivatives give \(O(a^6)\) per summand,
against \(O(a^{-4})\) modes per volume. This proves
\eqref{eq:v58-cutoff}. Weighted Taylor integral remainders, or the same
argument after mass differentiation, retain the displayed powers of
\(\epsilon\); no expansion in \(z/p^2\) outside its domain is used.
For the continuum sum--integral comparison, Poisson summation after
arbitrarily many internal derivatives has the same integrable lower-
masked majorant. Its scale is \(\lambda^{-2}\) at the first
complementary weight, proving \eqref{eq:v58-volume}.
\end{proof}

\begin{corollary}[Rooted source kernels and primitive-shell ownership]
\label{cor:v58-rooted}
After multiplication by a fixed external window, the inverse Fourier
kernel in the metric mark's relative coordinate satisfies
\begin{equation}
 a^4\sum_x(1+\mu d_{\rm per}(x,0))^b
 \|\mathcal K^\Delta_{a,L,I}(x)\|
 \le C_{I,b}a^2(\mu+m_H)^{d_I+2}.
 \label{eq:v58-rooted}
\end{equation}
Multiplication by the fixed local scalar, ghost and antifield factors in
\eqref{eq:v58-r1-marked} gives the corresponding full rooted multilinear
bounds, with their source-derivative powers included and no unrooted
volume factor. Expanding both scalar covariances and assigning each word
to its largest primitive shell gives complementary increments of size
\(C\mu^{d_I+2}\Lambda_j^{-2}\) on fixed mass-to-lower-scale domains.
Both mixed shell orientations must be retained.
\end{corollary}
\begin{proof}
The proof just given supplies every fixed external derivative. Rescale
external momenta by \(\mu\), integrate by parts beyond the chosen
moment order plus four, and periodize. Finite source translations have
radius \(O(a)\) and contribute bounded window factors. This proves
\eqref{eq:v58-rooted}. On one highest annulus, the differentiated loop
majorant is \(R^{-6}\); its four-dimensional sum is
\(O(\Lambda_j^{-2})\). Expansion of the product of the two covariance
sums is exact before the largest-label assignment.
\end{proof}

Exact finite Fourier identities refer to admitted grid momenta.
The differentiated coefficient bounds use the same smooth periodic
routing extension as the preceding source comparisons. No exact
finite-volume Slavnov identity is assigned to arbitrary off-grid
interpolated arguments.

The unremoved jets may have the usual bounds
\(Ca^{s-d_I}\) at weight \(s<d_I\) and
\(C[1+\log(1/(a\lambda))]\) at weight \(s=d_I\).
The positive coefficient \(w_a\) is an example of such a retained
power divergence. Products such as \(u\dot v_s\) at the next scalar
contraction order require their common subdivergence assignment; the
four single-loop estimates are not advertised as that higher-loop forest
theorem.

\section{The fixed-mask Euler operator has an evaluated reference curvature}
\label{sec:v58-curvature}

For any matrix \(N\), not necessarily symmetric, abbreviate
\[
 \Delta_N=v^{-1}\sum_{A,i,j}N_{ij}
 \partial_{\phi_A(j)}\partial^L_{\upsilon_A(i)},
 \qquad
 \delta_Q=\sum_{A,i}(Q\phi_A)_i\partial^L_{\upsilon_A(i)}.
\]
The operator under examination is
\(\mathscr Q_{q,\nu}=\delta_Q-\hbar\Delta_\nu\).
It uses the physical scalar Euler operator and the original mask; it is
not called a full interacting BV differential.

\begin{theorem}[Noncommuting cutoff, exact square and closed conditional balance]
\label{thm:v58-reference}
For symmetric \(Q(q)\),
\begin{equation}
 \boxed{
 \mathscr Q_{q,\nu}^{\,2}
 =-\frac{\hbar}{2v}\sum_{A,i,k}[\nu,Q]_{ik}
 \partial^L_{\upsilon_A(i)}\partial^L_{\upsilon_A(k)}.}
 \label{eq:v58-square}
\end{equation}
With the conditional map of \eqref{eq:v58-normalized-map}, the exact
retained-variable identity is
\begin{equation}
 \boxed{
 \mathscr T_q(\delta_Q-\hbar\Delta_\nu)
 =\left[\delta_{QA}+\hbar\Delta_{\Theta_\lambda V C}\right]
             \mathscr T_q,
 \qquad QA=P+\Theta_\lambda VA.}
 \label{eq:v58-effective-reference}
\end{equation}
The coordinate-contact term on the right is present for multiple
antifields and scalar-dependent source coefficients. At \(q=0\) the
curvature and the extra contact vanish, recovering V57.
\end{theorem}
\begin{proof}
The only surviving term in \(\delta_Q\Delta_\nu+
\Delta_\nu\delta_Q\) differentiates \(Q\phi\), giving
\(v^{-1}(\nu Q)_{ik}\partial_{\upsilon_i}^L
\partial_{\upsilon_k}^L\). Anticommutation replaces its coefficient by
\([\nu,Q]/2\), proving \eqref{eq:v58-square}.
For the second identity, conditional Gaussian integration by parts gives
\[
 \mathscr T_q\delta_QF
 =\delta_{QA}\mathscr T_qF
          +\hbar\mathscr T_q\Delta_{QC_q}F.
\]
The exact resolvent identities are
\[
 PC_q=\nu-\nu VC_q,\qquad
 QC_q-\nu=\Theta_\lambda VC_q,\qquad
 C_qA^{-T}=C.
\]
When a source derivative is moved through the mean substitution it
acquires \(A^{-T}\). Therefore
\(\mathscr T_q\Delta_{QC_q-\nu}
=\Delta_{\Theta_\lambda VC}\mathscr T_q\).
Finally \(PA=P-\nu VA\) proves the formula for \(QA\).
No division by \(\nu\) has been made.
\end{proof}

The same equation is the evaluated specialization of V57's conditional
identity to the actual quadratic metric interaction. The force term
\(\Theta_\lambda V\phi\) has conditional mean
\(\Theta_\lambda Vm\) and a contraction term
\(\hbar\Delta_{\Theta_\lambda VC}\). Retaining only its mean would
fail on nonlinear or multiple-antifield insertions.

\begin{proposition}[The first native curvature is not zero]
\label{prop:v58-native-curvature}
For one soft metric Fourier mark,
\begin{equation}
 [\nu,B_h](p,r)=[\nu(p)-\nu(r)]B_h(p,r),\qquad p-r=k.
 \label{eq:v58-mask-commutator}
\end{equation}
If this coefficient is nonzero, both momenta lie in
\(\lambda-|k|\le|p|,|r|\le2\lambda+|k|\), and
\[
 |\partial_p^\alpha\partial_k^\beta[\nu,B_h](p,p-k)|
 \le C_{\alpha,\beta}\|H\|\lambda^{2-|\alpha|-|\beta|}
\]
with the sharper undifferentiated bound
\(C\|H\||k|\lambda\).
For every smooth nonconstant admissible cutoff, the ordinary coefficient
is nonzero for a suitable arbitrarily soft conformal mark. It persists
on sufficiently dense finite momentum grids and small-mesh interior
approximations.
\end{proposition}
\begin{proof}
The first identity is matrix multiplication by the diagonal mask.
If its two values differ, the nearby momenta cannot both be below the
lower plateau or above the upper plateau. The support statement follows
by the triangle inequality. The mean-value theorem bounds their
mask difference by \(C|k|/\lambda\), while the native scalar vertex
on this annulus is bounded by \(C\|H\|\lambda^2\). Differentiation
gives the remaining bounds.
For the ordinary conformal mark \(H=I\), the actual input/output
vertex is \(p\cdot r+2z\), obtained by tracing
\eqref{eq:v27-V0}. Choose two nearby points in a region where the
smooth cutoff changes; their kinetic factor is nonzero. Finite-grid
approximation preserves the nonzero value when that open region is
resolved. No assertion of nonzero sampled curvature on every coarse
finite torus is needed.
\end{proof}

This curvature survives the zero-mesh limit at fixed \(\lambda\).
It is confined to the lower reference region, not an uncontrolled
ultraviolet contribution. It says that the unmatched combination of a
canonical physical Euler replacement and a fixed masked contraction is
not nilpotent at an interacting metric background. It does not contradict
any full cutoff-adapted BV or modified Slavnov identity.

For comparison only, replacing the contact by \(\Delta_{QC_q}\)
gives the exact conjugate of \(\delta_{QA}\) and therefore a nilpotent
scalar operator. This is not a permitted silent change of the active
contact convention. The two contacts differ by the explicit reference
term already retained in \eqref{eq:v58-effective-reference}.

\section{Geometric source derivatives and contact transport must move together}
\label{sec:v58-connection}

\begin{theorem}[Explicit geometry connection for the conditional source map]
\label{thm:v58-connection}
For a fixed metric direction \(h\), define the operator on retained
scalar arguments
\begin{equation}
 \boxed{
 \mathcal A_h^{\rm src}
 =-\sum_A(CQ_hA\varphi_A)\cdot\partial_{\varphi_A}
 -\frac{\hbar}{2v}\sum_A(CQ_hC)_{ij}
            \partial_{\varphi_A(i)}\partial_{\varphi_A(j)}.}
 \label{eq:v58-source-connection}
\end{equation}
Then
\begin{equation}
 \boxed{
 D_h\mathscr T_qF
 =\mathscr T_qD_hF+\mathcal A_h^{\rm src}\mathscr T_qF.}
 \label{eq:v58-connection}
\end{equation}
At every fixed finite scalar degree, \(\mathscr T_q\) is invertible on
its analytic metric germ. Consequently
\(\nabla_h=D_h-\mathcal A_h^{\rm src}\) is a flat connection:
\([\nabla_{h_1},\nabla_{h_2}]=0\) for fixed commuting directions.
This is flatness of a source-coordinate transport, not nilpotence of a
finite-regulator BRST vector field.
\end{theorem}
\begin{proof}
Differentiate the mean and covariance in \eqref{eq:v58-Wick}.
Moving derivatives with respect to the mean to retained coordinates
inserts \(A^{-T}\). Use
\(A^{-1}C_q=C\) and \(C_qA^{-T}=C\).
The drift is \(-CQ_hA\varphi\), and the diffusion is
\(-\hbar CQ_hC/(2v)\), proving the formula. Mean substitution is
invertible since \(A\) is, and the heat exponential has its opposite-
sign finite inverse on polynomials. Thus
\(\nabla_h=\mathscr T_qD_h\mathscr T_q^{-1}\), which proves flatness
without omitting mixed metric derivatives of \(Q\).
\end{proof}

The fixed contraction itself transforms as
\begin{equation}
 \boxed{
 \mathscr T_q\Delta_\nu\mathscr T_q^{-1}
 =\Delta_{\nu A^{-T}}
 =\Delta_{\nu(I+VC)}.}
 \label{eq:v58-contact-transport}
\end{equation}
This follows by moving the scalar derivative through the mean
substitution; the heat operator commutes with constant coefficient
contractions. The transported kernel need not be symmetric. It is
nevertheless an explicitly nilpotent odd contraction because all its
coefficients are scalar independent.

The contact is parallel for the same geometric connection:
\begin{equation}
 \boxed{
 D_h\Delta_{\nu A^{-T}}
 -[\mathcal A_h^{\rm src},\Delta_{\nu A^{-T}}]=0,
 \qquad
 D_h\Delta_{\nu A^{-T}}=\Delta_{\nu Q_hC}.}
 \label{eq:v58-contact-connection}
\end{equation}
Indeed the diffusion term in \(\mathcal A_h^{\rm src}\) commutes with
constant scalar derivatives, and its drift commutator is
\(\Delta_{\nu A^{-T}(CQ_hA)^T}=\Delta_{\nu Q_hC}\).
This uses the actual derivative of the mask after transport, not a
claim that its coefficients remain fixed.

For a supplied raw target \(\Delta_\nu F=t\), the exact measured target
statement is therefore
\begin{equation}
 \Delta_{\nu A^{-T}}\mathscr T_qF=\mathscr T_qt,
 \qquad
 \Delta_\nu\mathscr T_qF-\mathscr T_qt
 =-\Delta_{\nu VC}\mathscr T_qF.
 \label{eq:v58-target}
\end{equation}
V57's simple commutation is recovered only at \(q=0\). Both the
background derivative \eqref{eq:v58-source-connection} and the
contact correction \eqref{eq:v58-contact-transport} must appear in a
combined geometry/antifield identity. Treating its Gaussian covariance
as background dependent while leaving the contact and mean derivative
unchanged is not this pushforward.

No new finite-range contact section is proposed here. Conjugating V56's
section by \(\mathscr T_q\) is an exact formal operation, but its mean
substitution uses \(A(q)\), which is generally spatially nonlocal.
V56's old finite-support theorem cannot simply be transferred to that
conjugate. The first-mark source kernels controlled in
\cref{thm:v58-cutoff} are the quantitative statement proved instead.

\section{Status and the remaining complete measured identity}
\label{sec:v58-ledgers}

\subsection*{Proved}
The native quadratic scalar--metric interaction has been integrated
conditionally with its exact mean, covariance and metric-dependent
partition factor retained. Its action on the full inherited nonlinear
source is evaluated through all scalar self-contractions. The first
metric mark has four explicitly routed covariance coefficients, including
those absent on the free reference, with full-cutoff and rooted
\(O(a^2)\) complementary estimates after their actual local assignments.
The constant conformal response has a strictly nonzero native coefficient
of order \(a^{-4}\). The scalar Euler/contact square, its complete
conditional reference balance, the geometric source connection and the
transported odd contact are exact finite identities.

\subsection*{Inherited}
V27 supplies the completed scalar action, its coframe coefficients and
interior/upper symbol bounds. V56 supplies the source realization and
its finite difference operators. V57 supplies the zero-metric source
image and its normalization conventions. Gaussian completion and general
Wick/BV transport are not new. Earlier gravity, ghost and auxiliary-
measure results remain in the remaining integral and have not been
independently reverified in this part. The new source derivatives do not
rename the previously calculated action determinant as a new theorem.

\subsection*{Literature input}
Finite Gaussian effective actions and cutoff-compatible BV identities
are comparison context \cite{V58DJP,V58IIM}. No external anomaly
classification, cohomology vanishing theorem, or continuum Ward identity
is used to set an unknown native coefficient to zero.

\subsection*{Conditional}
The conditional integral is literal on its stated positive hard-scalar
chart and formal elsewhere in the inherited perturbative convention.
Uniform symbol estimates are established for one soft metric mark, not
for arbitrary finite-amplitude metric fluctuations. The complete source
formula is exact at fixed formal degree, but products of its loop
coefficients require the corresponding common local and subdivergence
normalization. General scalar-dependent contact targets remain supplied,
not fitted. Additional scalar self-interactions or internal gauge/Yukawa
couplings require their own non-Gaussian or enlarged Hessian calculation.

\subsection*{Open}
The next object is the combined metric/ghost/source identity after
\eqref{eq:v58-full-integral}, with
\eqref{eq:v58-source-connection} and
\eqref{eq:v58-effective-reference} inserted on their actual sides.
The lower-mask commutator is a specified source term in that equation,
not a condition to be suppressed by hand. The scalar calculation does
not include the final hard-metric integration, the native invariant
critical descent, the chiral/line contribution, or the complete
higher-loop master coefficient. No active normalization changes.
The programme remains the source-complete measured Einstein--Standard
Model EFT continuum at every prescribed finite order, not a
finite-parameter or nonperturbative ultraviolet completion.

\begin{center}\small
\begin{tabular}{>{\raggedright\arraybackslash}p{.23\textwidth}>{\raggedright\arraybackslash}p{.69\textwidth}}
\toprule
Variable & Exact scope of V58\\\midrule
Carrier & Same four real scalars, \(Q_a(q)\), free \(P\), lower mask,
source differences, and remaining metric/ghost/measure sector.\\
Background & Exact finite scalar conditional germ; quantitative cutoff
comparison at one metric derivative at zero.\\
Cutoff/volume & First-mark bounds uniform on the displayed fixed-window
and profile domains; thermodynamic comparison separate.\\
Mass & \(0\le m_H\le M\lambda\); weighted mass jets of the same
massive inverse, including vertex derivatives.\\
Loops & All scalar contractions of one marked polynomial conditional on
metric; not all full-theory diagrams at those loop orders.\\
Sources & Fixed polynomial and metric-jet order; rooted estimates have
no total-volume multiplier. No unbounded-order uniformity.\\
Symmetry & Exact scalar reference identities and source transport; fixed
Euler/contact operator explicitly curved, not a claimed complete QME.\\
Computation & Exact finite Gaussian/Grassmann tests and native symbol
identities; positive integral limits proved but not numerically fitted.\\\bottomrule
\end{tabular}
\end{center}

\needspace{12\baselineskip}
\begingroup
\renewcommand{\refname}{References for V58}
\begin{thebibliography}{9}\small
\bibitem{V58DJP} M.~Doubek, B.~Jur\v co and J.~Pulmann,
\emph{Quantum \(L_\infty\) Algebras and the Homological Perturbation Lemma},
arXiv:1712.02696, \emph{Commun. Math. Phys.} \textbf{367} (2019),
215--240. General finite BV/Feynman transfer context.
\bibitem{V58IIM} Y.~Igarashi, K.~Itoh and T.~R.~Morris,
\emph{BRST in the Exact RG}, arXiv:1904.08231 (2019).
Comparison for matched Wilsonian and modified-reference identities;
not a native ESM continuum theorem.
\end{thebibliography}
\endgroup

\section*{Append-only continuation marker: V58}
\addcontentsline{toc}{section}{Append-only continuation marker: V58}
\label{sec:v58-continuation-marker}
\textbf{End of V58.} Preserve Sections 1--453. The actual scalar--metric
interaction is now retained in the nonlinear marked source and its
first geometric response. The source mean, determinant, covariance,
contact and lower-reference terms have explicit simultaneous formulas.
The remaining task is their complete metric/ghost/descent assembly, not
another free scalar self-contraction or a new coincidence stencil.
No active action, normalization, propagator or measure is changed.

\clearpage
\part*{V59: Joint scalar pushforward of the complete metric--ghost BV defect}
\addcontentsline{toc}{part}{V59: Joint scalar pushforward of the complete metric--ghost BV defect}

\section{Return to the complete operator before estimating its pieces}
\label{sec:v59-start}

Sections 1--453 are preserved. V58 integrates the actual quadratic scalar
interaction at fixed metric and displays a nonzero square for a physical
scalar Euler replacement paired with only the hard-mask contraction.
The next question is not whether that isolated square can be discarded.
It is how the \emph{complete} field/antifield operator, including metric
contractions and scalar transformation sources, passes through the same
integral. This part answers that finite algebraic question without assuming
that the native finite master defect vanishes.

The new evaluated interface has three parts. The metric derivative of the
conditional scalar covariance contributes to the metric--antifield
contraction. Its apparently higher differential-order terms cancel exactly
against the conditional scalar stress. The remaining stress is the derivative
of the already present scalar determinant and Schur action. The scalar
contact, in turn, must be transported together with its complementary
reference contraction and the associated bracket. Finally, the whole native
master defect and its marked insertions have one ordered pushforward
identity; a nonzero input is retained rather than replaced by zero.

Finite Gaussian BV pushforward and the change of cutoff BV structures are
established machinery \cite{V59DJP,V59IIM,V59Z}. Their general existence is
not claimed as new. Here the formulas are proved with the actual scalar
covariance, its retained null space, the coframe metric variables, and the
source/density conventions of V56--V58. No pure-gravity cohomology theorem,
scalar-current classification, or unknown critical coefficient is used to
infer native master-equation closure.

\subsection*{Carrier and sign conventions}
Put \(v=a^4\), \(P=\widetilde P_a\), \(z=m_H^2\), and
\(\nu=I-\Theta_\lambda\). Retain
\[
 C=\nu P^{-1},\qquad PC=CP=\nu,\qquad
 \mathscr W_C=\exp\left(\frac{\hbar}{2v}
 \sum_{A,i,j}C_{ij}\partial_{\phi_A(i)}\partial_{\phi_A(j)}\right).
\]
There are four real scalar flavours. At \(z=0\), the quotient defining
\(C\) is zero on the retained modes; \(P^{-1}\) is never applied there.
All remaining fields, antifields, and independent composite-source labels
are spectators of \(\mathscr W_C\), not discarded variables. In particular,
the coframe coordinate \(q\), its antifield \(\eta\), the diffeomorphism
ghost and its antifield, and the nonminimal fields remain in the algebra.

For a constant scalar matrix \(N\), set
\begin{equation}
 \Delta_N^H=v^{-1}\sum_{A,i,j}N_{ij}
 \partial_{\phi_A(j)}\partial^L_{\upsilon_A(i)},\qquad
 \delta_P=\sum_{A,i}(P\phi_A)_i\partial^L_{\upsilon_A(i)}.
 \label{eq:v59-scalar-operators}
\end{equation}
Let \(\Delta_{\rm rest}\) be the flat Darboux contraction on the other
coordinate pairs, with their inherited graded signs. It includes
\(v^{-1}\sum_i\partial_{q_i}\partial^L_{\eta_i}\). It is independent of
\(\phi,\upsilon\) and has square zero. Write
\begin{equation}
 \Delta_b=\Delta_{\rm rest}+\Delta_I^H,\qquad
 \Delta_e=\Delta_{\rm rest}+\Delta_{\Theta_\lambda}^H,
 \qquad D_b=\delta_P-\hbar\Delta_b,
 \quad D_e=\delta_P-\hbar\Delta_e.
 \label{eq:v59-two-operators}
\end{equation}
The labels \(b,e\) mean before and after this scalar integration, not a
new ultraviolet completion. Every represented density factor independent
of the scalar is included \emph{once} as \(-\hbar\log\rho_{\rm rest}\)
in the action insertion. It is not also put in a weighted Laplacian.
The same statement in a weighted convention requires conjugating that
entire density convention.

For homogeneous \(F\), the bracket associated to a Laplacian is
\[
 (F,G)_\Delta=(-1)^{|F|}\{
 \Delta(FG)-(\Delta F)G-(-1)^{|F|}F\Delta G\}.
\]
These signs give \((v\phi^TP\phi/2,F)_b=\delta_PF\).
Changing the scalar contraction therefore changes its bracket as well.
The isolated \(\Delta_\nu^H\) of V58 remains an operand:
\(\Delta_I^H=\Delta_\nu^H+\Delta_{\Theta_\lambda}^H\).
It has not silently been replaced by another mask.

\section{The same Gaussian transports the full master defect}
\label{sec:v59-master}

\begin{lemma}[Actual covariance chain relation]
\label[lemma]{lem:v59-chain}
The finite operators in \eqref{eq:v59-two-operators} satisfy
\begin{equation}
 \boxed{D_b^2=D_e^2=0,\qquad
 \mathscr W_CD_b=D_e\mathscr W_C.}
 \label{eq:v59-chain}
\end{equation}
This assertion places no condition on the native interacting action.
\end{lemma}
\begin{proof}
Constant odd contractions square to zero. The anticommutator of
\(\delta_P\) and \(\Delta_N^H\) is a double odd derivative with
coefficient \(NP^T\), hence vanishes for \(N=I,\Theta_\lambda\), since
these matrices commute with the symmetric \(P\). The spectator
contractions anticommute with the scalar Euler replacement.
If \(H_C\) denotes the exponent of \(\mathscr W_C\), then
\[
 [H_C,\delta_P]=\hbar\Delta_\nu^H,
 \qquad [H_C,[H_C,\delta_P]]=0.
\]
All Laplacians commute with \(H_C\). Consequently
\(\mathscr W_CD_b\mathscr W_C^{-1}
=\delta_P-\hbar\Delta_b+\hbar\Delta_\nu^H=D_e\).
The proof uses only \(PC=\nu\), including on its null space.
\end{proof}

Split the actual action as
\[
 S_a=S_{H,0}+I_a,\qquad
 S_{H,0}=\frac v2\sum_A\phi_A^TP\phi_A.
\]
\emph{All} metric, ghost, nonminimal, scalar--metric, minimal-source,
composite-source and recorded density terms are in \(I_a\). In particular,
putting the gravitational quadratic term in \(I_a\) is only bookkeeping
for scalar integration; it does not change the hard gravitational
reference or assume that its finite free BRST block is exactly invariant.
Define the scalar effective interaction and normalized insertions by
\begin{equation}
 e^{-I_e/\hbar}=\mathscr W_Ce^{-I_a/\hbar},\qquad
 \mathcal T_{I_a}F=e^{I_e/\hbar}\mathscr W_C(e^{-I_a/\hbar}F).
 \label{eq:v59-Ieffective}
\end{equation}
These are formal Wick identities at fixed perturbative, field and source
orders. Scalar-independent exponential factors are kept outside the Wick
operation. On the positive quadratic scalar germ, the quadratic part can
instead be resummed literally as in V58; marked sources are then expanded
to their prescribed finite degree. No scalar integration theorem for an
arbitrary nonquadratic positive global measure is asserted.

For an even insertion \(I\), define
\begin{equation}
 \mathfrak A_b(I)=\delta_PI+\tfrac12(I,I)_b-\hbar\Delta_bI,
 \quad
 \mathfrak A_e(I)=\delta_PI+\tfrac12(I,I)_e-\hbar\Delta_eI.
 \label{eq:v59-defects}
\end{equation}
The first is exactly the flat-coordinate full canonical defect of
\(S_{H,0}+I\). The second is its derived cutoff-paired effective
interaction defect. Neither is assumed zero.

\begin{theorem}[Complete scalar pushforward, including a nonzero defect]
\label[theorem]{thm:v59-defect}
For the unchanged finite action and all retained source coefficients,
\begin{equation}
 \boxed{\mathfrak A_e(I_e)
       =\mathcal T_{I_a}\mathfrak A_b(I_a).}
 \label{eq:v59-defect-transfer}
\end{equation}
Let
\(\mathcal L_b=\delta_P+(I_a,\cdot)_b-\hbar\Delta_b\) and
\(\mathcal L_e=\delta_P+(I_e,\cdot)_e-\hbar\Delta_e\).
For every homogeneous marked insertion \(F\), with the defect kept on
its left,
\begin{equation}
 \boxed{\begin{aligned}
 \mathcal L_e\mathcal T_{I_a}F
 ={}&\mathcal T_{I_a}\mathcal L_bF\\
 &-\hbar^{-1}\bigl\{
 \mathcal T_{I_a}(\mathfrak A_bF)
 -(\mathcal T_{I_a}\mathfrak A_b)(\mathcal T_{I_a}F)\bigr\}.
 \end{aligned}}
 \label{eq:v59-marked-transfer}
\end{equation}
No missing connected defect insertion is replaced by a product of
separately transported coefficients.
\end{theorem}
\begin{proof}
The graded second-order product rule gives
\[
 D_b e^{-I_a/\hbar}
 =-\hbar^{-1}e^{-I_a/\hbar}\mathfrak A_b(I_a),
\]
and the identical formula holds on the effective side. Apply
\eqref{eq:v59-chain} to \(e^{-I_a/\hbar}\), then divide by its scalar
pushforward, to prove \eqref{eq:v59-defect-transfer}. For a mark, the
same rule gives
\[
 D_b(e^{-I_a/\hbar}F)
 =e^{-I_a/\hbar}
 [\mathcal L_bF-\hbar^{-1}\mathfrak A_bF].
\]
Apply the chain relation again and use the first assertion. The even
exponential introduces no reordering of \(F\) and the odd defect. This
proves the displayed connected term with its sign.
\end{proof}

The theorem proves invariance of an \emph{established} zero-defect input
under this scalar step. For the native finite filtered transformations,
the nonzero master candidate and its measured source square are still
on the right. The formal Wick map and multiplication by the normalized
exponentials are invertible before restricting to a smaller source
quotient: the formula cannot turn a nonzero complete input into zero
by algebra alone.

\begin{corollary}[Exact conversion back to the canonical bracket]
\label[corollary]{cor:v59-reference}
If the effective expression is tested using the original canonical
bracket and trace instead of its derived effective pair, then
\begin{equation}
 \boxed{\mathfrak A_b(I_e)
 =\mathcal T_{I_a}\mathfrak A_b(I_a)
 +\tfrac12(I_e,I_e)^H_\nu-\hbar\Delta_\nu^H I_e.}
 \label{eq:v59-reference-conversion}
\end{equation}
\end{corollary}
\begin{proof}
Subtract the two definitions in \eqref{eq:v59-defects}, using
\(\Delta_b-\Delta_e=\Delta_\nu^H\) and the corresponding bracket
difference. The free scalar Euler operator is the same in both.
\end{proof}

Equation \eqref{eq:v59-reference-conversion} is an explicit change of
operator bookkeeping, not a claim that the Legendre reference
supertrace calculated in earlier sections equals this expression without
its own field/source transformation. In particular its last two terms
are not asserted to vanish or satisfy an unproved ultraviolet estimate.
No inverse \(\Theta\), hard-mode deletion, or extra phase/density factor
is used in the conversion.

\section{The scalar stress cancels the metric-contact connection exactly}
\label{sec:v59-stress}

To expose the mixed cancellation, return to the zero-mark conditional
map \(\mathscr T_q\) of V58. Put \(V=Q-P\),
\(A=(I+CV)^{-1}\), \(C_q=AC=CA^T\), and \(m=A\varphi\).
For a coordinate metric index \(i\), write \(Q_i=\partial_{q_i}Q\).
Define its scalar force and conditional force by
\begin{equation}
 e_i^H=\tfrac12\sum_A\phi_A^TQ_i\phi_A,
 \qquad
 \bar e_i^H=\tfrac12\sum_A m_A^TQ_im_A
          +\frac{2\hbar}{v}\Tr(C_qQ_i).
 \label{eq:v59-stresses}
\end{equation}
The second trace contains the four-flavour factor. Let
\begin{equation}
 \mathcal A_i=-\sum_A(CQ_im_A)\cdot\partial_{\varphi_A}
 -\frac{\hbar}{2v}\sum_A(CQ_iC)_{jk}
            \partial_{\varphi_A(j)}\partial_{\varphi_A(k)}.
 \label{eq:v59-Ai}
\end{equation}
This is exactly V58's evaluated source connection, now inserted into the
metric--antifield part of the same operator.

\begin{theorem}[Joint stress/contact cancellation]
\label[theorem]{thm:v59-stress}
On every fixed scalar polynomial and metric germ,
\begin{equation}
 \boxed{\begin{aligned}
 \mathscr T_q\left[
  \sum_i e_i^H\partial^L_{\eta_i}
  -\frac{\hbar}{v}\sum_i\partial_{q_i}\partial^L_{\eta_i}
 \right]\mathscr T_q^{-1}
 =\sum_i\bar e_i^H\partial^L_{\eta_i}
  -\frac{\hbar}{v}\sum_i\partial_{q_i}\partial^L_{\eta_i}.
 \end{aligned}}
 \label{eq:v59-stress-cancellation}
\end{equation}
Furthermore
\begin{equation}
 \boxed{\bar e_i^H=v^{-1}\partial_{q_i}I_H^{\rm eff},
 \quad I_H^{\rm eff}=-\hbar\log Z_H
 =\frac v2\sum_A\varphi_A^TVA\varphi_A
       +2\hbar\log\det(I+CV).}
 \label{eq:v59-stress-effective}
\end{equation}
\end{theorem}
\begin{proof}
Gaussian conjugation of a single scalar multiplication is
\[
 \mathscr T_q\phi\mathscr T_q^{-1}
 =m+\frac{\hbar}{v}C\partial_\varphi,
 \qquad
 \mathscr T_q\partial_{q_i}\mathscr T_q^{-1}
 =\partial_{q_i}-\mathcal A_i.
\]
Normal-order the product of two such multiplications. The derivative
of the mean supplies \(\hbar\Tr(C_qQ_i)/(2v)\) per flavour. The
remaining first- and second-scalar-derivative terms are
\[
 \frac{\hbar}{v}(CQ_im)\cdot\partial_\varphi
 +\frac{\hbar^2}{2v^2}(CQ_iC):\partial_\varphi^2
 =-\frac{\hbar}{v}\mathcal A_i.
\]
Thus conjugated stress equals
\(\bar e_i^H-(\hbar/v)\mathcal A_i\), whereas conjugated metric
contact contributes
\(-\hbar\partial_{q_i}/v+(\hbar/v)\mathcal A_i\), each followed by
\(\partial^L_{\eta_i}\). They cancel before a norm or coincident
trace is taken. Differentiating V58's retained partition factor gives
\eqref{eq:v59-stress-effective}.
\end{proof}

The cancellation includes an apparent third-order operator
\((CQ_iC):\partial_\varphi^2\partial^L_{\eta_i}\). Omitting either the
scalar stress or the metric-contact connection leaves it nonzero.
It is not separately absorbed into a scalar source renormalization.
The determinant derivative is not cancelled: it is part of the
remaining effective stress and must stay in the rest-field integral.

\begin{proposition}[A nonzero native operand and its exact cancellation]
\label[proposition]{prop:v59-native}
At the flat metric, the native constant conformal vertex from V58 is
\[
 B_I(p)=2z+p_b(p)+s(ap)^2[p_n(p)-p_b(p)].
\]
On a momentum in the common positive hard interior,
\[
 (CB_IC)(p)
 =\frac{\nu(p)^2B_I(p)}{P(p)^2}>0.
\]
For a polynomial test with one metric antifield and two scalar
arguments in that mode, the second-derivative part of the transformed
stress is nonzero; the transformed metric contraction is its exact
negative. This holds at the finite completed symbol, not only at zero
mesh. At zero metric the surviving one-loop stress is
\(2\hbar v^{-1}\Tr(CB_I)\), with its actual local assignment retained.
\end{proposition}
\begin{proof}
On the interior, \(B_I=p_n+2z>0\) and \(P=p_n+z>0\) away from
zero momentum. Choose \(\nu>0\). The first assertion follows by
matrix multiplication. Two scalar derivatives of the quadratic test
supply the factor two, and \eqref{eq:v59-stress-cancellation} fixes
the opposite sign of the contact. This is an evaluated nonzero
operator coefficient; no full native loop integral is assigned a new
numerical value.
\end{proof}

For clarity, the theorem uses the canonical metric contraction before
any metric-mode pushforward. If a constant matrix \(M_g\) is instead
placed on that contraction, its conjugate has the explicit residual
\[
 \frac{\hbar}{v}\sum_{i,j}[(M_g)_{ij}-\delta_{ij}]
 \mathcal A_j\partial^L_{\eta_i}.
\]
It cannot be discarded. A later metric shell must change its bracket,
reference terms, and source transport together. The present scalar
step neither performs nor assumes that calculation.

\section{The complementary scalar contact and the complete source polynomial}
\label{sec:v59-source}

\begin{proposition}[Where the isolated scalar square goes]
\label[proposition]{prop:v59-complement}
The canonical scalar Euler/contact pair has the exact conditional
image
\begin{equation}
 \boxed{\mathscr T_q(\delta_Q-\hbar\Delta_I^H)
 \mathscr T_q^{-1}
 =\delta_{QA}-\hbar\Delta_{\Theta_\lambda}^H.}
 \label{eq:v59-scalar-complement}
\end{equation}
Both sides square to zero with \(q\) held as a coefficient. This does
not set V58's isolated hard-mask square to zero.
\end{proposition}
\begin{proof}
V58 gives
\(\mathscr T_q(\delta_Q-\hbar\Delta_\nu^H)\mathscr T_q^{-1}
=\delta_{QA}+\hbar\Delta_{\Theta VC}^H\).
Its complementary contact is transported as
\(\mathscr T_q\Delta_\Theta^H\mathscr T_q^{-1}
=\Delta_{\Theta A^{-T}}^H\).
Since \(A^{-T}=I+VC\), subtracting the complement gives
\eqref{eq:v59-scalar-complement}. Equivalently,
\(QC-A^{-T}=PC-I=-\Theta\).
For the square, \(VA\) is symmetric and
\[
 QA=P+\Theta VA,\qquad QA\Theta=P\Theta+\Theta VA\Theta
\]
is symmetric. The general Euler/contact anticommutator therefore
vanishes. On the input side this is simply symmetry of \(Q\) paired
with the identity contraction.
\end{proof}

The separate term involving \([\nu,Q]\) found in V58 is consequently
not an additional anomaly to be cancelled by an arbitrary local term.
It is the square of an incomplete scalar operator. Completing that
operator and transporting its bracket gives the displayed reference
pair. The \emph{full native master defect}, including filtered source
algebra and action breaking, can still be nonzero by
\eqref{eq:v59-defect-transfer}.

At zero internal gauge, Yukawa and scalar self-coupling, the native scalar
transformation source is linear in \(\phi\). Write it in its original
Grassmann order as
\begin{equation}
 J_a=v\sum_A\ell_A^T\phi_A,\qquad
 (\ell_A)_i=\sum_x\upsilon_A(x)
              (\mathscr L_a(c))_{xi},
 \quad
 \mathscr L_a(c)\phi=P_s((P_sc)\cdot D\phi),
 \label{eq:v59-ell}
\end{equation}
where \(P_s=S^2\) is V28's source filter, distinct from the kinetic
operator \(P\). The entries of \(\ell\) are even; they contain an
ordered scalar antifield and ghost. All ghost modes, including those
subsequently integrated, remain present.

\begin{theorem}[Exact conditional scalar action with its antifield source]
\label[theorem]{thm:v59-source}
For the component
\(I_a=I_{\rm rest}+v\phi^TV(q)\phi/2+J_a\), its scalar effective
interaction is
\begin{equation}
 \boxed{\begin{aligned}
 I_e=I_{\rm rest}+I_H^{\rm eff}
 +v\sum_A\ell_A^TA\varphi_A
 -\frac v2\sum_A\ell_A^TC_q\ell_A.
 \end{aligned}}
 \label{eq:v59-source-effective}
\end{equation}
The final term is the complete scalar-exchange two-antifield/two-ghost
contact conditional on all spectator variables. No scalar loop or
extra factor of four is attached to a fixed flavour pair in that term.
For additional marked composite sources \(\mathcal J\), the exact
replacement is
\(I_e=I_{\rm rest}+I_H^{\rm eff}
-\hbar\log\mathscr T_q e^{-(J_a+\mathcal J)/\hbar}\), expanded to
the specified source order.
\end{theorem}
\begin{proof}
The conditional Gaussian has mean \(A\varphi\) and coordinate
covariance \(\hbar C_q/v\). Its generating exponential for the
commuting even source \(\ell\) gives
\[
 \mathscr T_q e^{-v\ell^T\phi/\hbar}
 =\exp[-v\ell^TA\varphi/\hbar
          +v\ell^TC_q\ell/(2\hbar)].
\]
Taking the logarithm proves the formula with the determinant left in
\(I_H^{\rm eff}\). On the finite Grassmann algebra all nilpotent
source coefficients have finite expansions. Marked even composite
sources are handled at their prescribed Taylor order.
\end{proof}

The new formula is a scalar-conditional source polynomial, not a
recalculation of V30's hard metric--ghost--scalar loop. If all source
and ghost momenta entering \(\ell\) lie below the hard plateau, then
\(C\ell=0\); the scalar-exchange term vanishes by support. It must
nevertheless remain before the hard ghost integration, since the
same native \(\ell\) can then carry hard momentum. With two ordered
labels \(\ell_1=\upsilon_1c_1\), \(\ell_2=\upsilon_2c_2\), the
coefficient of \(\upsilon_1\upsilon_2c_1c_2\) in the final term is
\(+v(C_q)_{12}\). The plus sign follows from the exterior reordering,
not a positive-metric replacement for a gravitational propagator.

Every metric or ghost derivative of
\eqref{eq:v59-source-effective} differentiates this same action.
In particular
\(D_hC_q=-C_qQ_hC_q\) and \(D_hA=-C_qQ_hA\) act in the
antifield-quadratic term and the linear source term. Keeping only a
corrected scalar transformation and dropping its quadratic source
partner is not the pushed action whose defect satisfies
\eqref{eq:v59-defect-transfer}.

\section{Products, marked defects, and compatible scalar scale composition}
\label{sec:v59-products}

There is a second reason not to push the terms in the master equation
independently. Conditional Gaussian expectation is not multiplicative.
The defect transfer in \eqref{eq:v59-marked-transfer} keeps its
connected product for this reason.

\begin{proposition}[The exact conditional product kernel]
\label[proposition]{prop:v59-product}
At zero additional scalar source, put
\(\mathcal C(q)=A^{-1}C_qA^{-T}=C+CVC\).
For homogeneous polynomial sources define
\begin{equation}
 f\star_q g
 =m\exp\left[\frac{\hbar}{v}
 \sum_{A,i,j}\mathcal C(q)_{ij}
 \partial_{\varphi_A(i)}\otimes\partial_{\varphi_A(j)}\right](f\otimes g),
 \label{eq:v59-product}
\end{equation}
with the original exterior order and ordinary multiplication \(m\).
Then
\begin{equation}
 \boxed{\mathscr T_q(FG)=(\mathscr T_qF)\star_q(\mathscr T_qG).}
 \label{eq:v59-product-transport}
\end{equation}
The product is associative and graded commutative. It is a description
of the same source expectation, not a change to the finite action.
\end{proposition}
\begin{proof}
Split the scalar heat operator on a product into its two internal
second derivatives and its cross derivatives. All have coefficients
independent of scalar fields and therefore commute. Transforming
mean derivatives to retained coordinates multiplies the cross
covariance by \(A^{-1}\) and \(A^{-T}\). The identity
\(A^{-1}C_qA^{-T}=C(I+VC)\) gives the stated kernel.
Associativity follows either by conjugating ordinary multiplication
or by the three-factor exponential with all pair contractions.
The scalar derivatives are even, so spectator Grassmann order is
unmodified.
\end{proof}

For example,
\(\varphi_i\star_q\varphi_j
=\varphi_i\varphi_j+\hbar\mathcal C_{ij}/v\).
This is why conditional products in a master equation cannot be
replaced by products of conditional means. The kernel has derivative
\(D_h\mathcal C=CQ_hC\), precisely the coefficient entering the
stress/contact cancellation. The third-order ordinary differential
piece in the conjugated metric Laplacian is not absent by itself;
\cref{thm:v59-stress} cancels it only in the complete Euler/contact
operator.

\begin{proposition}[Exact scalar scale composition with the mask updated]
\label[proposition]{prop:v59-composition}
Let \(\Theta_0\succeq\Theta_1\succeq\Theta_2\) be symmetric
Fourier multipliers commuting with \(P\), and use
\(C_{10}=(\Theta_0-\Theta_1)P^{-1}\),
\(C_{21}=(\Theta_1-\Theta_2)P^{-1}\), with the declared zero-mode
convention. With the same spectator contraction at each step,
\[
 \mathscr W_{C_{21}}\mathscr W_{C_{10}}
 =\mathscr W_{C_{20}},\qquad
 \mathscr W_{C_{ji}}D_{\Theta_i}
 =D_{\Theta_j}\mathscr W_{C_{ji}}.
\]
The effective interaction and the ordered defect insertions obtained
by two steps equal the one-step result at every finite perturbative
and source order. The normalized insertion maps obey the corresponding
composition with the \emph{updated interaction} in the second step.
\end{proposition}
\begin{proof}
The constant scalar heat exponents commute and their covariances add.
The proof of \cref{lem:v59-chain} replaces \(I-\Theta\) by
\(\Theta_i-\Theta_j\). Apply the heat semigroup to the complete
unnormalized exponential before forming the effective logarithm or
normalizing marks. The partition factors cancel only in the normalized
map's telescoping quotient, not in the remaining metric integral.
\end{proof}

This algebraic composition retains mixed-shell contractions and all
source/exterior coefficients. It does not prove a stable analytic
class, a uniform number-of-shells bound, or summability of unrenormalized
coincidence coefficients. No new numerical beta coefficient follows
from a heat-semigroup identity.

\section{Fixed-order control and the remaining native estimate}
\label{sec:v59-bounds}

The new mixed cancellation is exact before summation; it has no
cutoff error that needs a separate estimate. Bounds for its surviving
trace stress and marked sources remain those already proved, with
all their local assignments. The complete defect transfer introduces
products of those coefficients; it does not by itself renormalize
such products.

\begin{proposition}[A finite coefficient bound with the actual contraction cost]
\label[proposition]{prop:v59-bound}
Let \(F\) have scalar polynomial degree at most \(n\), in a raw
monomial/exterior coefficient norm. Set
\(g_*=v^{-1}\max_{i,j}|C_{ij}|\). Then
\begin{equation}
 \boxed{\norm{[\hbar^j]\mathscr W_CF}_1
 \le\frac{n!}{2^jj!(n-2j)!}\,g_*^j\norm F_1,
 \qquad 0\le2j\le n.}
 \label{eq:v59-finite-bound}
\end{equation}
The same estimate holds for the opposite-sign heat operator.
It has no explicit number-of-sites factor. The bound for the
normalized conditional map additionally uses the actual finite-degree
mean-substitution norm and \(v^{-1}\max|(C_q)_{ij}|\).
\end{proposition}
\begin{proof}
Choose the \(2j\) scalar occurrences to be contracted, pair them,
and bound each coordinate covariance by \(g_*\). There are at most
\(n!/(2^jj!(n-2j)!)\) such pairings. Summing absolute coefficients
establishes the assertion; coincident variables or odd nilpotence can
only reduce the count. The argument has a fixed number of field
occurrences, not a separate sum over unrelated lattice sites.
\end{proof}

For the native translation-invariant free scalar covariance,
\(g_*\le u_a=O(a^{-2})\) by positivity and V57. Thus every
additional scalar pairing may cost two powers of the cutoff. An
algebraic equality such as \eqref{eq:v59-defect-transfer} supplies no
cutoff-independent norm for its right-hand side without the complete
local counterterms and connected source estimates. In particular it
does not imply that a defect known small only before contraction
remains small afterward.

A useful precise consequence is conditional: if the \emph{complete}
right-hand side of \eqref{eq:v59-marked-transfer}, including its
connected defect term, has a convergent normalized source limit after
the declared counterterms, then the left-hand side has the identical
limit. That is not a newly proved native estimate. V58's single-metric
bounds, or V56's coincidence matching alone, do not supply it for
arbitrary interacting or multiple-insertion coefficients.

\section{Ledgers and next load-bearing coefficient}
\label{sec:v59-ledgers}

\subsection*{Proved}
The same finite scalar covariance gives an exact chain relation between
the before/after contraction--bracket pairs. The complete native master
defect, with all spectator and source terms retained, has the ordered
pushforward identity \eqref{eq:v59-defect-transfer}. Marked insertions
retain their connected defect term. Conditional scalar stress cancels
the entire extra metric-contact connection, including its third-order
ordinary differential piece, leaving the derivative of the actual Schur
action and scalar determinant. The complementary scalar contact resolves
the isolated hard-mask square in its correct canonical/reference
combination. The linear native scalar-antifield source has an exact
quadratic-source completion, and source products and successive scalar
steps have their explicit consistent transport.

\subsection*{Inherited}
V27 supplies the actual scalar action and completed positive reference.
V28 supplies its scalar antifield prescription. V56--V58 supply source
conventions, the conditional mean/covariance and determinant, and the
previous marked-source estimates. V1 already treats general formal BV
pushforward as an input; the present result evaluates the relevant
operator and density interfaces for this scalar step. Earlier claims
about the gravitational/cohomological and chiral sectors are not
independently reverified here.

\subsection*{Literature input}
Gaussian BV effective actions, cutoff-dependent BV structures and
modified reference identities are comparison material
\cite{V59DJP,V59IIM,V59Z}. No external quantum-gravity existence or
anomaly-vanishing theorem is used. The displayed finite identities have
direct algebraic proofs.

\subsection*{Conditional}
The conditional scalar integral is literal on V58's positive quadratic
chart and formal at the specified perturbative and source orders
elsewhere. The output operator pair must be used together; conversion
to the input canonical pair requires
\eqref{eq:v59-reference-conversion}. No identification with a previous
Legendre-source convention is assumed without the corresponding
transformation. Cutoff-uniform norms for complete interacting defect
insertions, products, and the remaining metric/ghost integration are
not furnished by the finite chain relation.

\subsection*{Open}
The smallest unresolved quantity after this step is now an actual
\emph{joint} coefficient of
\(\mathcal T_{I_a}\mathfrak A_b(I_a)\), followed by the remaining
metric/ghost pushforward and the agreed source/reference conversion.
It includes the physical finite action breaking, source-algebra terms,
once-counted density and the connected insertions in
\eqref{eq:v59-marked-transfer}. The current and stencil constructions
of earlier versions do not populate this coefficient. Its invariant
critical local primitive and complete analytic bounds remain open,
as do the internal-gauge/chiral and higher-loop ESM interfaces. No
active counterterm, propagator, contour, density, or physical source
normalization changes in this installment.

\begin{center}\small
\begin{tabular}{>{\raggedright\arraybackslash}p{.24\textwidth}>{\raggedright\arraybackslash}p{.68\textwidth}}
\toprule
Variable & Exact scope of V59\\\midrule
Fields & Same four real scalars and completed kinetic/operator data;
all native metric, ghost, nonminimal and source variables remain spectators
until their later integration.\\
Symmetry & Complete finite defect is transported, not assumed zero;
the scalar output contraction and its bracket are changed by a derived
identity, with the canonical conversion retained.\\
Cutoff and volume & Algebraic identities exact at every finite cutoff;
no new asymptotic error rate. Fixed-degree heat bounds have no explicit
volume factor but retain their actual covariance growth.\\
Mass & Same \(P(z)\) and covariance, including retained zero modes;
identities hold at fixed mass and for common polynomial mass jets.\\
Order & All prescribed finite scalar, metric and source jets;
formal Wick composition. No unbounded-order analytic estimate.\\
Measure & The scalar determinant stays in the effective action and rest
integral. The existing auxiliary density is counted once as specified.\\
Computation & Exact rational and Grassmann fixtures and a native
positive-symbol cancellation test. No populated critical-breaking
integral or full metric/ghost loop is numerically evaluated.\\\bottomrule
\end{tabular}
\end{center}

\needspace{12\baselineskip}
\begingroup
\renewcommand{\refname}{References for V59}
\begin{thebibliography}{9}\small
\bibitem{V59DJP} M.~Doubek, B.~Jur\v co and J.~Pulmann,
\emph{Quantum \(L_\infty\) Algebras and the Homological Perturbation Lemma},
arXiv:1712.02696, \emph{Commun. Math. Phys.} \textbf{367} (2019),
215--240. Finite-dimensional effective action and BV transfer context.
\bibitem{V59IIM} Y.~Igarashi, K.~Itoh and T.~R.~Morris,
\emph{BRST in the Exact RG}, arXiv:1904.08231,
\emph{Prog. Theor. Exp. Phys.} (2019), 103B01.
Cutoff-compatible QME and modified Slavnov comparison.
\bibitem{V59Z} R.~Zucchini,
\emph{Exact renormalization group in Batalin--Vilkovisky theory},
arXiv:1711.01213, \emph{JHEP} \textbf{03} (2018), 132.
Scale-dependent BV structures; not a native ESM continuum theorem.
\end{thebibliography}
\endgroup

\section*{Append-only continuation marker: V59}
\addcontentsline{toc}{section}{Append-only continuation marker: V59}
\label{sec:v59-continuation-marker}
\textbf{End of V59.} Preserve Sections 1--460. The conditional scalar
source now participates in a complete metric/ghost BV defect transfer.
Its metric-contact/stress cancellation, scalar reference complement,
source quadratic return, determinant and connected marked defect are
retained together. The remaining task is the populated interacting
coefficient and its local/analytic control, not another assertion that
an isolated masked Euler operator must be nilpotent. The endpoint remains
the source-complete measured Einstein--Standard Model EFT continuum at
every prescribed finite order, not a finite-parameter or nonperturbative
ultraviolet completion.

\clearpage
\part*{V60: The populated scalar master defect and its vacuum metric row}
\addcontentsline{toc}{part}{V60: The populated scalar master defect and its vacuum metric row}

\section{From a transport identity to its actual scalar operands}
\label{sec:v60-start}

Sections 1--460 are preserved. V59 proves that scalar elimination transports
an arbitrary complete master defect, but leaves its right-hand side as
\(\mathcal T_{I_a}\mathfrak A_b(I_a)\). The present step expands that right-hand
side for the same native scalar-quadratic action and minimal scalar source.
It retains the finite action breaking and the square of the scalar
transformation as different operands. The result is a complete conditional
polynomial in retained scalars and scalar antifields, not an assumption that
either operand vanishes. Its scalar-vacuum restriction gives a new measured
one-loop metric/ghost row on V27's continued and completed scalar reference.

The input action, coframe chart, profiles and scalar source are precisely
\eqref{eq:v27-completed-Higgs}, \eqref{eq:v28-scalar-source} and
\eqref{eq:v59-ell}. Internal gauge, Yukawa and scalar self-couplings are zero
in this component. There are \(n_H=4\) real scalars. All gravitational,
ghost, nonminimal and density data remain in the spectator action. This is
not a new proof of Gaussian BV pushforward, of scalar renormalizability, or
of a current classification. Those frameworks are established
\cite{V60DJP,V60IIM}. No unproved anomaly-vanishing result is used below.

\subsection*{Sign, normalization and source category}
Put \(v=a^4\), \(z=m_H^2\), \(P=\widetilde P_a\),
\(C=\nu P^{-1}\), \(\nu=I-\Theta_\lambda\). As in V59,
\(PC=CP=\nu\), and the coordinate covariance is \(\hbar C/v\).
Inverse symbols are extended by zero on retained modes; no inverse of a
zero cutoff multiplier is used. Set
\begin{equation}
 S=S_R+\frac v2\sum_A\phi_A^TQ(q)\phi_A
       +v\sum_A\upsilon_A^T R(c)\phi_A,
 \qquad R(c)\phi=P_s((P_sc)\cdot D\phi),\quad P_s=S^2.
 \label{eq:v60-action}
\end{equation}
Here \(S_R\) is scalar and scalar-antifield independent, including the
once-counted insertion \(-\hbar\log\rho_R\). It contains the actual native
minimal gravitational sources and nonminimal fields; it need not obey a
master equation. Write
\[
 dF=(S_R,F)_{\mathrm{rest}},\qquad
 \mathfrak A_R=\tfrac12(S_R,S_R)_{\mathrm{rest}}
                         -\hbar\Delta_{\mathrm{rest}}S_R.
\]
The operator \(d\) acts on spectator coefficients, not on \(\phi,\upsilon\).
The Laplacian and its bracket have exactly V59's conventions. In particular,
\(d\) is odd, the entries of \(R\) are odd, and the scalar antifield occurs
\emph{before} the ghost. Algebraic transpose does not complex-conjugate
Fourier coefficients and does not remove Grassmann signs.

The statements hold at every prescribed finite spectator and source jet;
quadratic scalar sources may be absorbed in \(Q\) and linear sources in
the stated linear source when they preserve the assumptions. A general
higher-scalar-degree composite insertion is instead a mark in
\cref{thm:v60-marked}; it is not silently put into the Gaussian action.
The displayed minimal-source degree bound does not apply after arbitrary
nonquadratic interactions are added to \eqref{eq:v60-action}.
This is the native quadratic/minimal-source input component of V59, not
an assertion that all previously selected quantum counterterms have this
form. Counterterms with additional scalar powers, rest-antifield dependence,
or higher scalar-antifield degree must be included as marked perturbations
of the complete action. They are not removed by the displayed conditional
formula. The one-loop vacuum estimate below consumes only the native tree
scalar action and source; its local normalization remains a separate row.


\section{The action-breaking and source-curvature matrices}
\label{sec:v60-operands}

Define the two native matrices
\begin{equation}
 \boxed{\mathcal B=dQ+R^TQ+QR,\qquad
        \mathcal N=dR-R^2.}
 \label{eq:v60-BN}
\end{equation}
The matrix \(\mathcal B\) is odd and symmetric. The even matrix
\(\mathcal N\) is the physical square of the scalar transformation:
\(s^2\phi=\mathcal N\phi\) when its spectator transformation is the one
in \(S_R\). The scalar source has no metric dependence, so on the native
branch \(dR=R(r^c)\); \(dQ=D_qQ[r^q]\). Both are derivatives of the
specified finite action/source, not fitted breaking coefficients.

\begin{theorem}[Complete scalar contribution to the canonical defect]
\label[theorem]{thm:v60-bare}
For \eqref{eq:v60-action}, the full flat-coordinate master defect is
\begin{equation}
 \boxed{\begin{aligned}
 \mathfrak A_b(S)
 ={}&\mathfrak A_R+\frac v2\sum_A\phi_A^T\mathcal B\phi_A
       -v\sum_A\upsilon_A^T\mathcal N\phi_A
       -n_H\hbar\Tr R.
 \end{aligned}}
 \label{eq:v60-bare-defect}
\end{equation}
For the actual native scalar source, \(\Tr R=0\) on every odd finite
grid. No condition \(\mathcal B=0\), \(\mathcal N=0\), or
\(\mathfrak A_R=0\) is assumed.
\end{theorem}
\begin{proof}
The rest bracket acting on the scalar action gives
\(v\phi^T(dQ)\phi/2\). Its action on the ordered scalar source is
\(-v\upsilon^T(dR)\phi\), since the scalar antifield is odd. The scalar
bracket of the quadratic action with the source gives
\(v\phi^T(R^TQ+QR)\phi/2\). Half the source self-bracket is
\(v\upsilon^TR^2\phi\). The only scalar Laplacian of the action is
\(n_H\Tr R\). Cross terms between the scalar quadratic action and the
source under the rest bracket vanish, since neither contains a rest
antifield. These exhaust all terms. The finite trace is V28's paired
Fourier sum: \(\Tr(P_sM_{P_sc}D)\) has diagonal proportional to the
odd sum \(\sum_p ip_\mu s(ap)^2=0\). The rest density remains only
in \(\mathfrak A_R\).
\end{proof}

The sign in front of \(\upsilon\mathcal N\phi\) is fixed by this
bracket calculation. The previously studied insertion
\(\Omega_H=+\langle\upsilon,\mathcal N\phi\rangle\) is still that
insertion; its coefficient in the present signed canonical master defect
is \(-\Omega_H\). Confusing the positive source-square insertion with
the entire signed defect would change \eqref{eq:v60-bare-defect}.

\begin{proposition}[Coherence identities without finite nilpotence]
\label[proposition]{prop:v60-coherence}
The actual matrices obey
\begin{equation}
 \boxed{\begin{aligned}
 d\mathcal B+R^T\mathcal B-\mathcal BR
     &=d^2Q+\mathcal N^TQ+Q\mathcal N,\\
 d\mathcal N-R\mathcal N+\mathcal NR&=d^2R.
 \end{aligned}}
 \label{eq:v60-coherence}
\end{equation}
Thus a source-curvature insertion cannot be treated as an independent
closed action-breaking insertion when the spectator source itself has
curvature.
\end{proposition}
\begin{proof}
Apply the odd Leibniz rule to \eqref{eq:v60-BN}. In the first line the
terms \(R^TdQ\) and \(dQR\) cancel after adding the displayed two
multipliers. For odd matrix entries,
\((R^2)^T=-(R^T)^2\); the remaining terms are precisely
\((dR-R^2)^TQ+Q(dR-R^2)+d^2Q\).
For the second line,
\(d(R^2)=(dR)R-R(dR)\); substituting
\(dR=\mathcal N+R^2\) proves the identity. No use of \(d^2=0\) occurs.
\end{proof}

\section{An evaluated nonzero finite scalar-action defect}
\label{sec:v60-native}

The action-breaking matrix is computable already at zero metric. Let
\(p,r,k\) be outgoing scalar, scalar and ghost momenta with
\(p+r+k=0\), in a nonaliased chart. Write \(\mathcal V_a(p,r)\) for
V27's actual one-metric/two-scalar tensor
\eqref{eq:v27-actual-V}. For a ghost polarization \(c\), the
coefficient in \(\mathcal B(0,c)\) is
\begin{equation}
 \boxed{\begin{aligned}
 \mathcal B_a^{(0)}(p,r;k,c)=iP_s(k)\big\{
 &2k^\mu\mathcal V_{a,\mu\nu}(p,r)c^\nu\\
 &+P(r)P_s(r)(p\cdot c)+P(p)P_s(p)(r\cdot c)\big\}.
 \end{aligned}}
 \label{eq:v60-native-B}
\end{equation}
The two scalar source factors are output filters, not profiles inserted
on an unrelated scalar propagator.

\begin{proposition}[Exact native coefficient and first mesh term]
\label[proposition]{prop:v60-native}
Formula \eqref{eq:v60-native-B} is the complete flat metric
\(c\phi^2\) action-breaking coefficient. On the common interior
plateau,
\begin{equation}
 \mathcal B_a^{(0)}=ia^4\mathcal B_4+O(a^6R^9)\norm c,
 \qquad |p|+|r|+|k|\le R,
 \label{eq:v60-Bfour}
\end{equation}
with all prescribed differentiated estimates. Put
\(P_4(p)=\sum_\mu p_\mu^6/360\). Then
\begin{equation}
 \mathcal B_4=2k^\mu\mathcal V[E_4]_{\mu\nu}c^\nu
              +P_4(r)(p\cdot c)+P_4(p)(r\cdot c),
 \label{eq:v60-Bfour-formula}
\end{equation}
where \(\mathcal V[E_4]=\operatorname{Sym}E_4-	frac12\Tr(E_4)I\),
and the native fourth-order kinetic coefficients are
\begin{align}
 (E_4)_{\mu\nu}
 &=\frac{p_\mu^5r_\nu+p_\mu r_\nu^5}{120}
   +\frac{\epsilon_\mu\epsilon_\nu p_\mu^3r_\nu^3}{36},
             &&\mu\ne\nu,\nonumber\\
 (E_4)_{\mu\mu}
 &=\frac{p_\mu^5r_\mu+p_\mu r_\mu^5}{120}
  +\frac{p_\mu^4r_\mu^2+p_\mu^2r_\mu^4}{48}
  +\frac{p_\mu^3r_\mu^3}{36}.
 \label{eq:v60-Efour}
\end{align}
Here \(\epsilon_0=-1\), \(\epsilon_i=1\) are the inherited
continuation signs. No new correction is added to the action.
\end{proposition}
\begin{proof}
The metric variation is \(R_0(k)c=iP_s(k)(kc^T+ck^T)\), whose
contraction with \(\mathcal V_a\) gives the first term. Varying each
scalar in its actual free quadratic action gives the other two terms,
with the profiles in \eqref{eq:v28-source-symbol}. At zero mesh,
\[
 2k^\mu\mathcal V_{0,\mu\nu}(p,r)
 =-(p^2+z)r_\nu-(r^2+z)p_\nu,
\]
so the ordinary coefficient cancels exactly. Expand the sine/hyperbolic
symbols and the already recorded second-order correction. The off-diagonal
products give the first line of \eqref{eq:v60-Efour}; expanding
\(\kappa_\mu(p)\kappa_\mu(r)c_\mu(p+r)\) gives the second. At
\(r=-p\), \((E_4)_{\mu\mu}=-p_\mu^6/360\), giving \(P_4\).
All source profiles equal one on the plateau. The remaining analytic
Taylor terms have weight at least nine after the ghost derivative, proving
the bound. Mass terms cancel exactly in this extraction.
\end{proof}

\begin{corollary}[A nonzero coefficient with no source-curvature ambiguity]
\label[corollary]{cor:v60-witness}
Take \(p=r=te_j\), \(k=-2te_j\), \(c=e_j\), with all labels on
the plateau. Then
\begin{equation}
 \boxed{\mathcal B_a^{(0)}=-\frac{i}{6}a^4t^7+O(a^6t^9).}
 \label{eq:v60-witness}
\end{equation}
For spatial \(j\), its exact value is
\(i[4t a^{-2}(1-\cos at)^2-a^2t^5]\); for \(j=0\), it is
\(i[-4t a^{-2}(\cosh at-1)^2+a^2t^5]\).
Both are independent of \(z\) and nonzero for sufficiently small
nonzero \(at\). The coefficient is defined for coincident scalar
momenta by labelled polarization, with the action's \(1/2\) retained.
\end{corollary}
\begin{proof}
Only one diagonal kinetic entry is involved. It has fourth mesh
coefficient \(31t^6/360\), while \(P_4(te_j)=t^6/360\).
Substitution in \eqref{eq:v60-Bfour-formula} gives \(-t^7/6\).
The exact expressions follow directly from V27's diagonal symbols and
its actual second-order correction. Their Taylor series prove the claim.
\end{proof}

This is a populated finite action defect, not a quantum anomaly. The same
single-ghost panel contains no two-ghost curvature term. In particular,
\(\mathcal N\) cannot cancel this one-ghost coefficient. Its vanishing
as \(a\to0\) is a soft tree statement, not an estimate after arbitrary
hard contractions.

\section{Complete conditional defect and connected marked insertion}
\label{sec:v60-push}

Use V59's unchanged quadratic scalar conditioning and ordered linear source:
\[
 V=Q-P,\quad A=(I+CV)^{-1},\quad C_q=AC=CA^T,
 \quad (\ell_A)_j=\sum_i\upsilon_A(i)R_{ij}.
\]
The entries of \(\ell\) are even. Write
\begin{equation}
 m_A=A\varphi_A,\qquad b_A=C_q\ell_A,\qquad
 \mu_A=m_A-b_A.
 \label{eq:v60-shifted-mean}
\end{equation}
The full source-conditioned Gaussian has mean \(\mu_A\) and coordinate
covariance \(\hbar C_q/v\); its determinant remains in the effective
action of \cref{thm:v59-source}.

\begin{theorem}[Expanded joint master coefficient after scalar elimination]
\label[theorem]{thm:v60-populated}
For the same minimal scalar action/source component, V59's full effective
master defect is exactly
\begin{equation}
 \boxed{\begin{aligned}
 \mathfrak A_e(I_e)={}&\mathfrak A_R
 +\frac v2\sum_A\mu_A^T\mathcal B\mu_A
 -v\sum_A\upsilon_A^T\mathcal N\mu_A\\
 &+\frac{n_H\hbar}{2}\Tr(C_q\mathcal B)
       -n_H\hbar\Tr R.
 \end{aligned}}
 \label{eq:v60-effective-defect}
\end{equation}
Thus its nonconstant scalar part has the following complete source-degree
expansion:
\begin{align}
  &\frac v2\sum_A m_A^T\mathcal Bm_A,
                      &&(\upsilon\text{-degree }0),\nonumber\\
  &-v\sum_A m_A^T\mathcal B b_A
             -v\sum_A\upsilon_A^T\mathcal Nm_A,
                      &&(\upsilon\text{-degree }1),\nonumber\\
  &\frac v2\sum_A b_A^T\mathcal Bb_A
             +v\sum_A\upsilon_A^T\mathcal N b_A,
                      &&(\upsilon\text{-degree }2).
 \label{eq:v60-source-degrees}
\end{align}
There are no higher scalar-antifield powers in this minimal component.
The quadratic-antifield coefficient has three ghosts. Neither its
\(\mathcal B\) part nor its \(\mathcal N\) part may be omitted.
\end{theorem}
\begin{proof}
V59's defect transport applies to the complete action. Apply its normalized
conditional map to \eqref{eq:v60-bare-defect}. For each scalar flavour,
\[
 \mathcal T\phi_i=\mu_i,\qquad
 \mathcal T(\phi_i\phi_j)=\mu_i\mu_j+\hbar(C_q)_{ij}/v.
\]
The coefficients \(\mathcal B,\mathcal N,\mathfrak A_R,R\) are scalar
independent, so these two moments exhaust the expectation. There is no
Wick contraction of a scalar antifield. Expanding \(\mu=m-b\) proves
\eqref{eq:v60-source-degrees}; \(\mathcal B^T=\mathcal B\) and the
evenness of \(m,b\) justify the single cross term. Only the closed trace
has a factor \(n_H\); a fixed external flavour pair does not.
\end{proof}

This result makes V59's abstract right-hand side explicit in the native
matrices. It leaves \(\mathfrak A_R\) and the subsequent metric/ghost
integral intact. It does not assign an unevaluated local integral a
numerical value or claim that the complete defect is zero.

\begin{theorem}[The connected marked defect is a finite second-order operator]
\label[theorem]{thm:v60-marked}
Let \(F\) be any fixed scalar-polynomial marked insertion with ordered
spectator/exterior coefficients. Put
\(D_{A,i}=(\mathcal B\mu_A)_i-(\upsilon_A^T\mathcal N)_i\).
With the defect on the left, its connected insertion in V59 is
\begin{equation}
 \boxed{\begin{aligned}
 \hbar^{-1}\{\mathcal T(\mathfrak A_bF)
             -(\mathcal T\mathfrak A_b)(\mathcal TF)\}
 ={}&\sum_{A,i,j}D_{A,i}(C_q)_{ij}\mathcal T(\partial_{\phi_A(j)}F)\\
 &+\frac{\hbar}{2v}\sum_{A,i,j,k,l}
 \mathcal B_{ij}(C_q)_{ik}(C_q)_{jl}
                 \mathcal T(\partial_{\phi_A(k)}\partial_{\phi_A(l)}F).
 \end{aligned}}
 \label{eq:v60-connected}
\end{equation}
It is inserted with the minus sign already present in
\eqref{eq:v59-marked-transfer}. The equality holds for odd as well as
even \(F\) without moving a coefficient through \(F\).
\end{theorem}
\begin{proof}
For a Gaussian with mean \(\mu\), the product formula for a quadratic
polynomial \(G\) and a polynomial \(F\) has only the cross-contraction
orders zero, one and two:
\[
 \mathcal T(GF)-(\mathcal TG)(\mathcal TF)
 =\Sigma_{ij}\mathcal T(\partial_iG)\mathcal T(\partial_jF)
 +\tfrac12\Sigma_{ik}\Sigma_{jl}
          \mathcal T(\partial_i\partial_jG)
          \mathcal T(\partial_k\partial_lF).
\]
Here \(\Sigma=\hbar C_q/v\), and the exterior order is that of the
left polynomial followed by the right. For the scalar-dependent part of
\eqref{eq:v60-bare-defect}, the first two derivatives are \(vD\)
and \(v\mathcal B\). Scalar-independent rest and divergence terms
cancel from the connected expression. Substitution gives the result.
\end{proof}

\section{The scalar-vacuum metric row and its lower-reference trace}
\label{sec:v60-vacuum}

There is a particularly useful genuine one-loop restriction of
\eqref{eq:v60-effective-defect}. Set retained scalars and scalar antifields
to zero, but retain metric and ghost arguments. Define
\begin{equation}
 \Gamma_{H,1}(q)=\frac{n_H}{2}\Tr\log(I+CV(q)),\qquad
 X_H(q,c)=n_H\Tr(A(q)\Theta_\lambda R(c)).
 \label{eq:v60-vacuum-functions}
\end{equation}
The common \(\hbar\) is removed from \(\Gamma_{H,1}\) and \(X_H\).
They are not additional determinants on top of V59: they are its existing
scalar determinant and the reference coefficient now extracted together.

\begin{proposition}[Exact finite measured vacuum row]
\label[proposition]{prop:v60-vacuum-row}
At zero scalar backgrounds the extra scalar contribution to the complete
master defect is
\begin{equation}
 \boxed{\frac{n_H}{2}\Tr(C_q\mathcal B)-n_H\Tr R
             =d\Gamma_{H,1}-X_H.}
 \label{eq:v60-vacuum-row}
\end{equation}
This is a finite matrix identity; it does not suppose finite scalar-action
invariance. The first metric/one-ghost coefficient is
\begin{equation}
 \boxed{\begin{aligned}
 D_{H,a}(h,c)={}&\Pi_{H,a}(h,R_0c)+\tau_{H,a}(R_1(h,c)),\\
 X_{H,a}^{(1)}(h,c)={}&-n_H\Tr(CQ_h\Theta_\lambda R(c)),
 \end{aligned}}
 \label{eq:v60-first-row}
\end{equation}
where all vertices and sources are native and
\begin{align}
 \tau_{H,a}(u)&=\tfrac{n_H}{2}\Tr(CQ_u),\nonumber\\
 \Pi_{H,a}(h,u)&=\tfrac{n_H}{2}
   \Tr(CQ_{hu}-CQ_hCQ_u).
 \label{eq:v60-tau-Pi}
\end{align}
In particular, the coframe contact \(Q_{hu}\) and the tadpole/source
contact \(\tau(R_1)\) both remain.
\end{proposition}
\begin{proof}
The resolvent identities give
\(C_qQ=I-A\Theta\), \(QC_q=I-\Theta A^T\). Hence
\[
 \tfrac12\Tr[C_q(R^TQ+QR)]
       =\Tr R-\Tr(A\Theta R).
\]
Every factor except \(R\) is even, so the indicated cyclic moves are
valid without an extra sign. The remaining trace is
\(d\Gamma_{H,1}=n_H\Tr(C_qdQ)/2\). Differentiate once in the
metric, using \(D_hC_q=-CQ_hC\), \(D_hA=-CQ_h\), and the
actual first source derivative \(R_1\), to obtain
\eqref{eq:v60-first-row}--\eqref{eq:v60-tau-Pi}.
\end{proof}

This one-loop row has no metric propagator and hence no additional
\(\chi^{-1}\) factor in the native \(q\) normalization. It is the
scalar-vacuum contribution to the metric/ghost equation, not V27's
mixed scalar--metric self-energy and not V26's pure-gravity loop.

On an external window \(|k|\le\mu\ll\lambda\), every word in
\(X_H^{(1)}\) contains a lower covariance \(\nu\) and a nearby
\(\Theta\) factor. Its loop momentum lies in
\(\lambda/2\le|p|\le3\lambda\). The source profiles there are on
their interior plateau if \(a\lambda\) is sufficiently small. The
factor \(\Theta\) stays between the actual matrices; it is not
commuted through \(Q_h\) or \(R\).

\subsection*{A nonzero native reference coefficient, without rotational averaging}

The retained trace can be evaluated on a useful local projection. This is a
one-loop reference coefficient of the actual scalar determinant, distinct
from the tree example in \cref{cor:v60-witness}.

\begin{theorem}[Positive conformal divergence projection at zero mass]
\label[theorem]{thm:v60-positive-reference}
Let \(h(x)=I_4e^{-ik\cdot x}\) and
\(c(x)=ce^{ik\cdot x}\). At zero mesh and zero mass, the first metric
reference coefficient is the finite integral
\begin{equation}
 \boxed{X_{H,0}^{(1)}(I,c;k)
 =-in_H\int_p\nu(p)\Theta_\lambda(p+k)
       (p\cdot c)\left(1+\frac{p\cdot k}{p^2}\right),}
 \label{eq:v60-conformal-X}
\end{equation}
where \(\int_p=(2\pi)^{-4}\int d^4p\), and the integrand is extended
by zero on the lower ball. Define its averaged divergence coefficient by
\[
 J_{H,0}=\frac1{4i}\sum_{j=0}^{3}
       \left.\partial_{k_j}X_{H,0}^{(1)}(I,e_j;k)\right|_{k=0}.
\]
For every admissible smooth real profile \(0\le\theta\le1\), equal
to one on the unit ball and zero outside the ball of radius two,
\begin{equation}
 \boxed{J_{H,0}
 =\frac{n_H\lambda^4}{4}
   \int\frac{d^4y}{(2\pi)^4}[3\theta(y)-\theta(y)^2]
 \ge\frac{n_H\lambda^4}{64\pi^2}>0.}
 \label{eq:v60-positive-reference}
\end{equation}
Radial symmetry is not required. This is a projection of the actual tensor
coefficient, not an averaging replacement of its other components.
\end{theorem}
\begin{proof}
The ordinary conformal vertex in \eqref{eq:v60-first-row} is
\(Q_I(p,p+k)=p\cdot(p+k)+2z\). The source factor on the trace is
\(i(p\cdot c)\), while the covariance is \(\nu(p)/(p^2+z)\).
At \(z=0\) this proves \eqref{eq:v60-conformal-X}, with the indicated
matrix order and finite support. Differentiating and summing the four
coordinate tests gives
\[
 J_{H,0}=-\frac{n_H}{4}\int_p
   \left[(1-\Theta)p\cdot\nabla\Theta+(1-\Theta)\Theta\right].
\]
Since \(\Theta-\Theta^2/2\) is compactly supported,
\[
 \int_p(1-\Theta)p\cdot\nabla\Theta
       =-4\int_p(\Theta-\Theta^2/2).
\]
Substitution and \(p=\lambda y\) give the equality. Finally,
\(3\theta-\theta^2\ge2\theta\), and the normalized unit
four-ball volume is \(1/(32\pi^2)\). This proves the strict lower
bound without an angular ansatz.
\end{proof}

The native finite counterpart uses exactly V27's vertex in
\eqref{eq:v60-first-row}, not its ordinary conformal replacement. On the
same fixed external-extension convention used for the Taylor jets,
\begin{equation}
 |J_{H,a,L}-J_{H,0}|
 \le C_b n_H\lambda^4\left[(a\lambda)^4
                  +(\lambda L_{\min})^{-b}\right],\qquad b>4.
 \label{eq:v60-positive-reference-error}
\end{equation}
Indeed one external derivative of the reference word has degree zero;
the loop is in the lower transition annulus, and every native inverse and
vertex discrepancy there is fourth relative mesh order. The finite sum
comparison follows by the same differentiated compact-support argument
used in the next section. Positivity persists whenever the displayed
error is below the positive margin; no unevaluated numerical threshold is
claimed. The coefficient is a finite local reference normalization at
fixed \(\lambda\), not an anomaly or a physical beta function.


\begin{lemma}[Ordinary normalized scalar-vacuum identity]
\label[lemma]{lem:v60-ordinary}
At zero mesh and in thermodynamic normalization, assign weights
\(\deg k=1\), \(\deg z=2\). With \(\mathbf T_D\) denoting V28's
common weighted Taylor assignment,
\begin{equation}
 \boxed{(I-\mathbf T_5)(D_{H,0}-X_{H,0}^{(1)})=0.}
 \label{eq:v60-ordinary}
\end{equation}
The lower-reference term is not zero in general.
\end{lemma}
\begin{proof}
The ordinary action is the scalar density action in the same coframe
chart, so its full variation is a total derivative:
\(dQ_0+R_0^{H,T}Q_0+Q_0R_0^H=0\). This is proved by varying that
action, independently of a finite Ward identity. To justify its trace,
first use a smooth even upper guard \(w_\Lambda\) in the covariance.
The reference complement is
\(I-w_\Lambda\nu=\Theta+(1-w_\Lambda)\nu\). The same ordered
matrix calculation as above then gives the lower and upper reference
rows. Finite Fourier boxes can be chosen larger than the upper support
plus all routed external transfers; the only nonzero guarded words are
then identical to their nonaliased ordinary kernels.

The scalar two-point metric coefficient is even under simultaneous
momentum reversal, and its row is odd in the external momentum. After
\(\mathbf T_5\), the first possible weighted degree is seven. Seven
weighted external derivatives make the upper reference contribution
\(O(\Lambda^{-2})\): its undifferentiated integrand has degree one,
the four-dimensional integration adds four, and seven derivatives
remove seven. Finite total-derivative terms and the paired scalar trace
are zero before the guard is removed. The surviving lower trace is
compactly supported in its transition region. Removing the upper guard
therefore proves \eqref{eq:v60-ordinary}. This argument does not use
finite scalar-action invariance to prove its own cutoff estimates.
\end{proof}

At nonzero mass the tadpole is not a single mass-independent constant.
Its complement \((I-\mathbf T_4)\tau_{H,a}(z)\), including its
higher mass jets, is retained. The degree-faithful row is
\[
 D_{H,a}^{R}=\Pi_{H,a}^{R}(h,R_0c)
                    +\tau_{H,a}^{R}(R_1(h,c)).
\]
The weighted assignment intertwines these multipliers because
\(R_0,R_1\) have weight one. Setting the entire massive tadpole to
zero would not give \eqref{eq:v60-ordinary}.

\section{Full-cutoff comparison of the assembled scalar-vacuum row}
\label{sec:v60-estimate}

All estimates in this section concern coefficients per physical volume,
after factoring their momentum-conservation delta, as in V27--V28.
Retain odd grids and fixed completed profiles. Set
\begin{equation}
 \lambda\ge32768\mu>0,\qquad a\lambda\le2^{-24},\qquad
 \lambda L_\nu\ge1,\qquad 0\le m_H\le M\lambda.
 \label{eq:v60-domain}
\end{equation}
Every external leg is in the stated \(\mu\) window. Constants depend
on fixed \(M\), profile seminorms, the coframe chart, derivative and
kernel-moment orders, and the finite source count, not on mesh or volume.
Define
\[
 W_{H,a,L}^{R}=(I-\mathbf T_5)(D_{H,a,L}-X_{H,a,L}^{(1)}),
 \qquad\epsilon=|k|+m_H.
\]

\begin{theorem}[A new scalar contribution with a quantitative vacuum Ward limit]
\label[theorem]{thm:v60-bound}
For \(r=|\alpha|+2j\le7\), the same-volume comparison obeys
\begin{equation}
 \boxed{\left|
 \partial_k^\alpha\partial_z^j
 (W_{H,a,L}^{R}-W_{H,0,L}^{R})(h,c)
 \right|
 \le C_{\alpha,j}n_H a^2\epsilon^{7-r}\norm h\norm c.}
 \label{eq:v60-row-bound}
\end{equation}
The independent thermodynamic comparison obeys, for every fixed \(b>4\),
\begin{equation}
 |\partial_k^\alpha\partial_z^jW_{H,0,L}^{R}(h,c)|
 \le C_{\alpha,j,b}n_H\lambda^{-2}
      (\lambda L_{\min})^{-b}\epsilon^{7-r}\norm h\norm c.
 \label{eq:v60-volume}
\end{equation}
The ordinary thermodynamic coefficient is zero. This is a result for the
assembled scalar-vacuum row, not for the unnormalized full ESM defect.
\end{theorem}
\begin{proof}
We supply the analytic input rather than infer it from the Ward identity.
Each primitive scalar covariance has degree minus two. Each first or
second metric derivative of V27's scalar action has degree two, with its
mass term retained. The two terms of \(\Pi_H\) and the single
\(\tau_H\) term therefore have degree zero. The completed symbols
are smooth periodic in loop momentum. On the common interior their
ordinary difference starts at fourth relative mesh order, since this
loop consumes no gravitational cubic vertex. For any fixed weighted
external derivative of weight \(s\), with \(R_p=\lambda+|p|\),
their complete integrands have bounds
\[
 |\partial^{[s]}f_0|+|\partial^{[s]}f_a|\le C R_p^{-s},
 \qquad
 |\partial^{[s]}(f_a-f_0)|\le C a^4R_p^{4-s}.
\]
These follow by differentiating the actual inverse and vertex symbols;
each mass derivative lowers degree by two and retains its single mask
on each differentiated covariance. Loop-momentum derivatives satisfy
analogous bounds. The estimates include the contact \(Q_{hu}\),
not only the two-vertex loop.

Parity makes every weighted degree in \(\Pi_H\) and \(\tau_H\)
even. After \(\mathbf T_4\), six weighted derivatives suffice. The
ordinary tail above \(c/a\) is bounded by
\(C\int_{c/a}^\infty R^{-3}\,dR=O(a^2)\), and the interior
difference by \(Ca^4\int_\lambda^{c/a}R\,dR=O(a^2)\).
Rescaling the smooth periodic upper region gives the same rate.
The standard cube comparison of positive radial majorants gives these
bounds for normalized finite sums when \(\lambda L_\nu\ge1\).
Weighted Taylor's formula, or the substitution \(z=m_H^2\) followed
by ordinary Taylor's formula in \((k,m_H)\), gives
\[
 |\partial^{[s]}(\Pi_{H,a}^R-\Pi_{H,0}^R)|
 +|\partial^{[s]}(\tau_{H,a}^R-\tau_{H,0}^R)|
 \le Cn_Ha^2\epsilon^{6-s},\qquad s\le6.
\]
For the tadpole, momentum derivatives that do not occur are simply zero;
its mass Taylor remainder is not dropped. Multiplication by \(R_0\)
and \(R_1\) adds one external derivative and yields the row rate.

For \(X_H^{(1)}\), the loop is confined to the lower transition
annulus. Its integrand has degree one. After seven weighted derivatives,
its finite-to-ordinary difference is bounded by
\(Ca^4R_p^{-2}\). Integration over that fixed annulus yields
\(Ca^4\lambda^2\), so
\[
 |\partial^{[r]}(X_{H,a}^{(1),R}-X_{H,0}^{(1),R})|
 \le Cn_Ha^4\lambda^2\epsilon^{7-r}.
\]
This is no larger than the bound in \eqref{eq:v60-row-bound} on
\eqref{eq:v60-domain}. In particular an inverse cutoff multiplier is
never estimated.

For the volume comparison, the differentiated ordinary integrands and
all prescribed loop derivatives are integrable, with scale
\(\lambda^{-2}\) after the seventh row derivative. Repeated Fourier
integration by parts followed by Poisson summation gives
\((\lambda L_{\min})^{-b}\) for every fixed \(b>4\).
\Cref{lem:v60-ordinary} supplies the zero thermodynamic value only after
these analytic estimates have been established. Higher prescribed
external derivatives follow from the same differentiated symbol bounds;
their constants are not uniform in derivative order.
\end{proof}

\begin{corollary}[Rooted sources and scalar-shell ownership]
\label[corollary]{cor:v60-rooted}
After a smooth fixed external window at scale \(\mu\), the same-volume
kernel difference in the two arguments \(h,c\) satisfies
\begin{equation}
 \boxed{a^4\sum_x(1+\mu d_{\mathrm{per}}(x,0))^b
       \norm{\mathcal K_{H,a,L}^{\Delta}(x)}
 \le C_b n_H a^2(\mu+m_H)^7.}
 \label{eq:v60-rooted}
\end{equation}
Its scalar-vacuum determinant parent has complementary shell increments
\(Cn_H(\mu+m_H)^6\Lambda_j^{-2}\); multiplying by the tree
source adds the required derivative. Every scalar covariance in the
contact and double-vertex words participates in the largest-shell
assignment. The complements are summable at fixed order.
\end{corollary}
\begin{proof}
The differentiated symbol bounds in the preceding proof hold to every
fixed order. Integration by parts beyond four relative dimensions gives
the integrable weighted inverse-Fourier bound; periodization gives the
same lattice estimate with the physical factor \(a^4\). This is a
bilinear local-source bound with one argument in \(L^1\) and one in
\(L^\infty\), not a supremum bound on an extensive action. A shell
has normalized four-dimensional mode count \(O(\Lambda_j^4)\), and
the sixth differentiated degree-zero parent has degree minus six.
Expanding each primitive covariance and assigning a word to its largest
label proves the shell rate without a diagonal-shell approximation.
\end{proof}

The local jets may diverge and remain part of the common action/source
normalization. This theorem supplies a nonlocal normalized row; it does
not provide its complete invariant local primitive, nor justify inserting
it in further metric/ghost loops without their own local and connected
estimates. The determinant, reference trace and paired measure convention
must accompany any subsequent integration.

\section{Ledgers, verification and the next interacting coefficient}
\label{sec:v60-ledgers}

\subsection*{Proved}
The minimal scalar contribution to the full native canonical master defect
has the explicit action-breaking/source-curvature decomposition
\eqref{eq:v60-bare-defect}, with its signed antifield term and scalar
divergence. The two coherence equations retain the actual spectator
curvature. The flat finite scalar-action coefficient is evaluated directly
from V27's continued and completed action; it is nonzero at fourth mesh
order on \cref{cor:v60-witness}. Scalar elimination gives the complete
conditional defect through scalar-antifield degree two and a finite
second-order formula for every connected scalar-polynomial mark. Its
scalar-vacuum metric/ghost restriction has the exact finite measured row,
a retained lower-reference trace with the explicitly positive conformal
divergence projection in \cref{thm:v60-positive-reference}, and the full-cutoff and rooted
\(O(a^2)\) complementary comparison of \cref{thm:v60-bound}.

\subsection*{Inherited}
V27 supplies the scalar action, tensor vertices, continuation and smooth
completion. V28 supplies its ordered scalar source. V30's scalar-square
insertion is used with its original meaning; it is not the entire signed
master defect. V58--V59 supply the conditional mean, covariance, determinant
and the complete paired master transport. Their mathematical bodies are
preserved, not independently reverified wholesale. The rest of the
nonminimal, density and gravitational data remain at their previous scope.

\subsection*{Literature input}
Finite Gaussian BV transfer and cutoff-modified Slavnov identities are
comparison material \cite{V60DJP,V60IIM}. Their abstract constructions are
not advertised as new, and neither citation is used to set a native
finite breaking, critical cohomology class, or loop coefficient to zero.

\subsection*{Conditional}
The conditional formulas are literal on the existing positive scalar
quadratic chart and formal in the prescribed spectator/source jets
elsewhere. The reference-derived effective bracket is V59's bracket;
conversion to the original canonical bracket still needs its displayed
additional terms. Further metric/ghost integration requires the full
\(\mathfrak A_R\), the conditional quadratic-antifield terms, and the
connected insertion \eqref{eq:v60-connected}. The new analytic estimate
covers the stated scalar-vacuum row, not all these later contractions.

\subsection*{Open}
The next joint coefficient is the remaining metric/ghost expectation of
\eqref{eq:v60-effective-defect}, with the actual rest breaking and
all scalar-source terms retained. Its interacting local counterterms,
critical-descent decomposition and uniform higher-loop bounds are not
supplied by the finite conditional formula or by the single vacuum row.
Internal gauge/chiral, determinant-line and complete filtered ESM
invariant-class estimates remain separate. No active normalization,
propagator, contour or measure has been changed.

\begin{center}\small
\begin{tabular}{>{\raggedright\arraybackslash}p{.24\textwidth}>{\raggedright\arraybackslash}p{.68\textwidth}}
\toprule
Variable & Scope of V60\\\midrule
Fields and couplings & Four real scalars at zero internal gauge, Yukawa
and scalar self-couplings; native gravitational, ghost and nonminimal
variables retained.\\
Cutoff and volume & Algebraic formulas at every finite odd grid;
\eqref{eq:v60-row-bound} on \eqref{eq:v60-domain}. Separate
thermodynamic comparison; no total-volume loss in rooted bounds.\\
Mass & Same massive propagators; common polynomial jets in \(z=m_H^2\).
No negative-mass loop or inverse mass. Lower scale \(\lambda>0\)
not removed.\\
Source order & Complete conditional minimal scalar polynomial and any
fixed polynomial scalar mark; analytic theorem only for one metric and
one ghost in the scalar-vacuum contribution.\\
Iteration & Summable complementary scalar-loop increments, not a
self-mapping theorem for the complete interacting ESM class.\\
Computation & Exact exterior/rational tests, native trigonometric series
and Ward coefficients. Fixtures distinguished from native numerical
loop sums; no full critical primitive evaluated.\\\bottomrule
\end{tabular}
\end{center}

The checker independently compares the complete bare bracket with
\eqref{eq:v60-bare-defect}, verifies the odd-matrix coherence equations,
tests the shifted Gaussian and marked coefficients, reconstructs the
native flat defect and its fourth-order coefficient, and checks the
vacuum trace relation with noncommuting matrices and retained zero modes.
These finite checks do not certify all inherited analytic or cohomological
arguments, numerical values of native profile integrals, or a proof-assistant
formalization.

\begingroup
\renewcommand{\refname}{References for V60}
\begin{thebibliography}{9}\small
\bibitem{V60DJP} M.~Doubek, B.~Jur\v co and J.~Pulmann,
\emph{Quantum \(L_\infty\) Algebras and the Homological Perturbation Lemma},
\emph{Commun. Math. Phys.} \textbf{367} (2019), 215--240,
arXiv:1712.02696. Gaussian BV transfer comparison.
\bibitem{V60IIM} Y.~Igarashi, K.~Itoh and T.~R.~Morris,
\emph{BRST in the Exact RG}, \emph{Prog. Theor. Exp. Phys.}
(2019), 103B01, arXiv:1904.08231. Paired cutoff/reference identities;
not a native ESM continuum theorem.
\end{thebibliography}
\endgroup

\section*{Append-only continuation marker: V60}
\addcontentsline{toc}{section}{Append-only continuation marker: V60}
\label{sec:v60-continuation-marker}
\textbf{End of V60.} Preserve Sections 1--467. The scalar part of the
joint master defect and its connected insertions are now explicit in
native action and source matrices. A nonzero finite action coefficient
and a quantitative scalar-vacuum metric row have been evaluated. The
remaining task is the interacting metric/ghost expectation with its
local normalization and source/reference conversion, not another
standalone assertion of Gaussian pushforward. The endpoint remains
source-complete measured Einstein--Standard Model EFT continuum control
at every prescribed finite order.


\clearpage
\part*{V61: The first scalar--metric two-loop contraction and its overlapping subgraphs}
\addcontentsline{toc}{part}{V61: The first scalar--metric two-loop contraction and its overlapping subgraphs}

\section{From a conditional scalar identity to a second hard loop}
\label{sec:v61-start}

Sections 1--467, including all previous continuation markers, are preserved.
V60 evaluates the conditional scalar master defect and a scalar-vacuum
metric/ghost row before the remaining metric integration. The next operation
is not another application of scalar Gaussian conditioning: it contracts
metric legs of the scalar determinant. That operation produces two-loop
graphs with proper ultraviolet subgraphs. A fixed-soft-window bound on the
first loop does not control its contraction at arbitrary hard momentum.

This installment performs that additional contraction in one concrete
component of the unchanged theory. It retains four real scalars, the V27
completed scalar action, the V14/V24 gravitational reference and tree jets,
and the same odd periodic grids, contour and density convention. Internal
gauge, Yukawa and scalar self-couplings are zero. The new component has no
retained scalar, ghost or antifield, but its ordinary metric sources are
retained through degree two. It is the part of the connected negative
logarithm that contains a scalar loop, at loop order two. A common factor
\(\hbar^2\) is removed from its coefficient. We first compute the bare
coefficient. Insertions of the existing order-\(\hbar\) counterterms are
separate terms of the renormalized two-loop expression, not silently
included in the bare graphs.

The new results are: the complete two-loop core and its source derivatives;
an evaluated six-power vacuum coefficient in an explicitly declared
ordinary comparison regulator; an exact hard-momentum growth formula for
an already soft-renormalized one-loop subgraph; and a full-cutoff comparison
for the all-comparable-scale part of the two-metric coefficient. The
scale-separated, overlapping-subgraph normalization is identified but not
claimed completed. In particular, the all-comparable projection is not an
approximation that discards the remaining modes: it names one part of a
fixed decomposition whose complement remains unpaid.

The native field action is never replaced by a continuum Feynman rule.
Continuum tensors below are the ordinary limits of the same native
vertices, and every comparison regulator is labelled as such. The
coefficient of a vacuum functional is distinguished from a normalized
observable. Differentiation in geometry is performed before evaluating the
flat background; an unmarked vacuum constant alone would cancel from
\(\log Z[J]-\log Z[0]\).

\section{The complete scalar-containing two-loop core}
\label{sec:v61-core}

Use finite orthonormal coordinates for the physical field inner products.
In these coordinates a Gaussian contraction is \(\hbar G\); conversion to
site coordinates reintroduces the physical cell factors as in V58--V60.
Let \(G_g(q)\) be the actual background hard metric Wick covariance and
\(G_H(q)\) the actual hard scalar covariance, each with \(\hbar\) removed.
The former retains the signed trace sector of the declared contour. It is
not replaced by a positive metric tensor. Write
\[
 Q_I(q)=D_qQ_a(q)[e_I],\qquad
 Q_{IJ}(q)=D_q^2Q_a(q)[e_I,e_J],
\]
where \(Q_a\) is the completed scalar operator and \(e_I\) ranges over
all hard metric directions. These derivatives include the coframe
coordinate contact. On the background germ,
\[
 D_hG_H=-G_HQ_hG_H,\qquad
 D_hG_g=-G_gB_h^gG_g,
\]
with \(B_h^g\) the derivative of the actual completed gravitational
Hessian. The hard-support inverses mean covariance resolvents; a retained
zero mode is never inverted.

\begin{lemma}[Bridge exclusion after, not before, connectedness]
\label[lemma]{lem:v61-bridge}
Consider a connected vacuum graph with at most two external metric marks,
each of momentum at most \(\mu\), and hard covariances supported above
\(\lambda>2\mu\). If an internal line is a bridge, the graph is zero.
This applies also when either component cut off by the bridge contains a
closed scalar, metric or ghost loop or an order-\(\hbar\) density vertex.
\end{lemma}
\begin{proof}
Removing a bridge separates the graph into two components. Summing momentum
conservation over one component cancels every internal momentum in it.
The bridge momentum is the negative sum of the external marked momenta
on that component, hence has norm at most \(2\mu\). Its covariance is
zero. The argument is unaffected by the species of the loops or their
Wick signs. On the small nonaliased external window the same statement
holds on the periodic carrier. This does not exclude graphs with a cycle
connecting the components, nor graphs with additional hard external
sources. It is a support proof for the stated component only.
\end{proof}

\begin{theorem}[Native two-loop scalar core and its two-metric restriction]
\label[theorem]{thm:v61-core}
The scalar-containing two-loop coefficient through two ordinary retained
metric marks is the corresponding background jet of
\begin{equation}
 \boxed{\mathcal G_{a,2}^H(q)
 =\frac{n_H}{4}\sum_{I,J}(G_g(q))_{IJ}
 \Tr\!\left[
 G_H(q)Q_{IJ}(q)-G_H(q)Q_I(q)G_H(q)Q_J(q)
 \right],\qquad n_H=4.}
 \label{eq:v61-core}
\end{equation}
The first term is a metric--scalar contact graph with two loops; the second is
the three-propagator graph with one metric line and two scalar lines.
The formula includes all primitive internal momenta and all metric
polarizations. No closed ghost loop or additional auxiliary determinant
is missing from this scalar-containing coefficient at this loop order.
The latter assertion is restricted to the stated zero-ghost/zero-antifield
external sector and uses \cref{lem:v61-bridge}.
\end{theorem}
\begin{proof}
The mixed cubic and quartic hard vertices are, in these coordinates,
\[
 S_{H,3}=\frac12\sum_{A,I}x_I\phi_A^TQ_I\phi_A,
 \qquad
 S_{H,4}=\frac14\sum_{A,I,J}x_Ix_J\phi_A^TQ_{IJ}\phi_A.
\]
The connected negative logarithm contributes
\(\langle S_{H,4}\rangle_0-\langle S_{H,3}^2\rangle_{0,c}/(2\hbar)\)
at the order in question. The quartic expectation gives the first term
of \eqref{eq:v61-core}. Of the two scalar pairings between the cubic
vertices, each gives the same ordered trace by cyclicity, producing the
second term with coefficient \(-n_H/4\). Pairing the scalar legs
within each cubic instead gives
\[
 -\frac{n_H^2}{8}\sum_{I,J}(G_g)_{IJ}
     \Tr(G_HQ_I)\Tr(G_HQ_J).
\]
This is a bridge graph and vanishes on the declared source window.
It is not discarded as a disconnected diagram: it is connected and needs
the support argument.

For an independent completeness check, let every external metric mark be
absorbed into a background derivative of a propagator or vertex. A
connected two-loop vacuum core has
\(\sum_v(\deg_h(v)-2)=2\), after quadratic hard vertices are resummed.
The possibilities are one hard-degree-four vertex or two hard-degree-three
vertices. A scalar-containing quartic is precisely \(x^2\phi^2\).
Two scalar-containing cubics are \(x\phi^2\) and give the two pairings
just analyzed. A scalar cubic joined to a purely gravitational or
metric--ghost cubic connects its scalar tadpole to the other loop by a
bridge. Two scalar loops joined by one metric edge do likewise. An
order-\(\hbar\) scalar-independent density or ghost-determinant insertion
couples to the scalar determinant at two-loop order only by its linear
bridge term. Every such contribution vanishes by the lemma; higher
vertices or non-bridge connections would raise the loop order. Hard-linear
tree branches themselves vanish by the same soft-transfer argument.

Equivalently, the scalar determinant is
\(\Gamma_{H,1}=n_H\Tr\log(I+C_HV_H)/2\). Its normalized hard metric
expectation at this order is \(G_g:D_q^2\Gamma_{H,1}/2\), with the
bridge terms treated above. Differentiating the determinant gives
\eqref{eq:v61-core}. Expanding the background covariances and vertices
through two metric marks supplies all such marked graphs. This uses only
scalar jets through four metric derivatives and gravity jets through the
quartic vertex, which are already part of the completed input.
\end{proof}

There is an important distinction from V59: \eqref{eq:v61-core} is an
actual second-loop contraction, not only an identity transporting a defect
through a scalar Gaussian. It is nevertheless not the entire two-loop
master coefficient. Quantum counterterm insertions and ghost/antifield
marks require their own terms.

\begin{proposition}[The seven terms of the first geometric derivative]
\label[proposition]{prop:v61-source}
Put \(C=G_H\) and
\(U_{IJ}=\Tr(CQ_{IJ}-CQ_ICQ_J)\). The first derivative of
\eqref{eq:v61-core} is \(n_H/4\) times
\begin{align}
 &\sum_{I,J}(-G_gB_h^gG_g)_{IJ}U_{IJ}\notag\\
 &\quad+\sum_{I,J}(G_g)_{IJ}\Tr\bigl(
 CQ_{IJh}-CQ_hCQ_{IJ}-CQ_{Ih}CQ_J-CQ_ICQ_{Jh}\notag\\
 &\hspace{53mm}+CQ_hCQ_ICQ_J+CQ_ICQ_hCQ_J\bigr).
 \label{eq:v61-seven}
\end{align}
Its second derivative is obtained by differentiating every factor of this
single expression. Both scalar covariances, the metric covariance, and
all three contact derivatives are included.
\end{proposition}
\begin{proof}
Apply the two inverse-derivative formulas and the ordinary product rule.
The positive last two terms differentiate the two scalar inverse factors
of \(-CQ_ICQ_J\). The separate first line differentiates the metric
inverse. None of these matrices is commuted through another factor.
Differentiating again is a finite labelled product rule and accounts for
every distribution of the two external marks.
\end{proof}

The formula concerns arbitrary off-shell geometric source polarizations.
It is not a transverse-traceless extraction. Although the Gaussian metric
inverse is gauge fixed, a subsequent Ward equation must still use its
proper off-shell classical Euler multiplier; no replacement of that
multiplier by a gauge-fixed one is authorized here.

\section{An evaluated ordinary two-loop coefficient and its cutting factors}
\label{sec:v61-value}

This section computes an ordinary comparison coefficient to fix the Wick
weights and show nonvanishing. It does not install a new finite regulator.
At zero retained metric and zero mass put the same nonnegative even
upper/lower weight \(w\) on each ordinary propagator:
\[
 G_g(p)=\frac{2w(p)}{\chi p^2}J,\qquad
 G_H(p)=\frac{w(p)}{p^2}.
\]
The inherited native tensor contraction is
\(V_0(p,q):JV_0(p,q)=p^2q^2/2\), and the inherited coframe contact
has \(K_{\rm con}=\operatorname{diag}(-1/2,2,2,2)\). These are
\cref{prop:v27-numerator,eq:v27-contact}, not new vertex calculations.
Let \(\int_p=(2\pi)^{-4}\int\dd^4p\).

\begin{proposition}[Guarded massless two-core coefficient]
\label[proposition]{prop:v61-vacuum}
For radial \(w\), the coefficient per physical volume is
\begin{equation}
 \boxed{
 \mathcal E_2^H[w]
 =\frac{11n_H}{16\chi}
     \left(\int_p\frac{w(p)}{p^2}\right)
     \left(\int_qw(q)\right)
 -\frac{n_H}{4\chi}
     \int_{p,q}\frac{w(p)w(q)w(p+q)}{(p+q)^2}.}
 \label{eq:v61-vacuum}
\end{equation}
All terms are formed under one convergent guard. No separate routing is
assigned to a divergent summand.
\end{proposition}
\begin{proof}
The contact trace is the closure of V27's contact self-energy with weight
\(n_H/2\). Since
\(\operatorname{av}_{S^3}p^TK_{\rm con}p/p^2
 =\Tr K_{\rm con}/4=11/8\), it is the first term displayed.
The two-cubic term of \eqref{eq:v61-core} has weight \(-n_H/4\).
The metric covariance supplies \(2/\chi\), and the signed numerator
supplies \(p^2q^2/2\); cancel the two scalar denominators before
integrating. This gives the second term.
\end{proof}

The different cutting factors can be stated independently of radiality.
If \(\Pi_g^{\rm con}\) and \(\Pi_g^{\rm bub}\) denote V27's fixed-flavour
one-loop scalar two-point coefficients with the same primitive weights,
then
\begin{equation}
 \mathcal E_2^H
 =\frac{n_H}{2}\Tr(G_H\Pi_g^{\rm con})
  +\frac{n_H}{4}\Tr(G_H\Pi_g^{\rm bub}).
 \label{eq:v61-cutting}
\end{equation}
The contact graph has one scalar edge; the two-cubic graph has two.
Using one common factor when closing the full scalar self-energy counts
the two-cubic graph incorrectly. Equation~\eqref{eq:v61-cutting} is
not a prescription to renormalize its two terms independently: all
subgraph assignments and source derivatives must subsequently be
retained.

\begin{theorem}[A nonzero six-power ordinary comparison]
\label[theorem]{thm:v61-gaussian}
Take \(w_\Lambda(p)=\exp(-p^2/\Lambda^2)\). Then
\begin{equation}
 \boxed{\mathcal E_2^H[w_\Lambda]
   =\frac{31n_H}{12288\pi^4\chi}\Lambda^6.}
 \label{eq:v61-gaussian-value}
\end{equation}
With the fixed smooth lower mask included,
\(w_{\lambda,\Lambda}=\nu_\lambda w_\Lambda\), and
\(\Lambda\ge8\lambda\), the same leading term has error
\(O(n_H\chi^{-1}\lambda^2\Lambda^4)\).
\end{theorem}
\begin{proof}
The Gaussian integrals are
\[
 \int_p\frac{e^{-p^2/\Lambda^2}}{p^2}
   =\frac{\Lambda^2}{16\pi^2},\qquad
 \int_pe^{-p^2/\Lambda^2}=\frac{\Lambda^4}{16\pi^2}.
\]
For the second integral use \(r=p+q\) and translate \(p\) by \(r/2\).
Then
\[
 p^2+q^2+(p+q)^2=2(p-r/2)^2+\tfrac32r^2,
\]
and the integral equals
\[
 \frac{\Lambda^4}{64\pi^2}
 \frac{\Lambda^2}{24\pi^2}
 =\frac{\Lambda^6}{1536\pi^4}.
\]
Inserting these values into \eqref{eq:v61-vacuum} gives
\(11/4096-1/6144=31/12288\).

For the lower-mask error, telescope the product of the three weights.
A removed metric momentum \(r\) is confined below \(2\lambda\);
its Coulomb integral is \(O(\lambda^2)\), while the remaining Gaussian
integral is \(O(\Lambda^4)\). A removed scalar momentum has volume
\(O(\lambda^4)\), and the other Coulomb Gaussian is bounded by
\(C\Lambda^2\), uniformly for that small shift. The contact product
obeys the same or a smaller error bound. This proves the claim without
requiring radiality of the lower profile.
\end{proof}

With a lower mask the bridge argument applies and this is the full
stated bare scalar-containing component. Without that mask,
\(\mathcal E_2^H[w_\Lambda]\) denotes only the infrared-integrable
sum of the two cores. No massless zero-momentum bridge or unnormalized
Gaussian partition is defined by that comparison integral.

The coefficient in \eqref{eq:v61-gaussian-value} is not a universal
vacuum-energy counterterm: the power divergence depends on the regulator
and the stated coframe/gauge convention. The native finite coefficient is
\eqref{eq:v61-core}, not this Gaussian value. This exact comparison fixes
the signs and symmetry factors and shows that the two primitive graphs
need not cancel. The unmarked constant disappears from the normalized
source functional, whereas its metric derivatives must be formed from
\eqref{eq:v61-core} before setting the background to zero.

\section{A soft-renormalized subgraph still has a divergent hard closure}
\label{sec:v61-hardclosure}

V27's ordinary massless bubble has the already established nonlocal
representative
\begin{equation}
 \Pi^{\rm bub,R}_0(k)=\frac{k^2}{\chi}
       \bigl[J_\lambda(k)-\mathsf T_2J_\lambda(k)\bigr],\qquad
 J_\lambda(k)=\int_p\frac{\nu_\lambda(p)\Theta_\lambda(p+k)}{p^2}.
 \label{eq:v61-inherited-J}
\end{equation}
The integral is finite for every \(k\). Here its extension outside the
retained window is the explicitly specified primitive two-point graph;
we do not identify it there with every term of a Wilsonian effective
functional, whose support-excluded branches may change. This is the
amplitude that becomes a proper subgraph when its external scalar line is
contracted in the next loop. Its normalization is the inherited Taylor
assignment, not a new one. A soft estimate does not assert smallness at
these new hard subgraph momenta.

Restrict for this diagnostic to an admissible smooth radial lower profile
\(\theta(r)\), equal to one for \(r\le1\) and zero for \(r\ge2\).
Define
\[
 a_\theta=\int_y\frac{\theta(|y|)}{|y|^2},\quad
 b_\theta=\int_y\frac{\theta(|y|)^2}{|y|^2},\quad
 M_\theta=\int_y\theta(|y|),\quad
 E_\theta=\int_1^2r\,\theta'(r)^2\dd r.
\]
These profile moments are not new coupling choices.

\begin{theorem}[Exact hard-momentum polynomial of the inherited subgraph]
\label[theorem]{thm:v61-hardbubble}
Put
\(\kappa_\theta=(E_\theta-1)/(64\pi^2)\). Then
\begin{equation}
 \boxed{\kappa_\theta\ge
       \frac{1/\log2-1}{64\pi^2}>0.}
 \label{eq:v61-kappa}
\end{equation}
For \(|k|>4\lambda\), the inherited renormalized coefficient is exactly
\begin{equation}
 \boxed{
 \Pi^{\rm bub,R}_0(k)
 =\frac1\chi\left[
       \lambda^4M_\theta
       -\lambda^2(a_\theta-b_\theta)k^2
       -\kappa_\theta(k^2)^2
 \right].}
 \label{eq:v61-hard-polynomial}
\end{equation}
\end{theorem}
\begin{proof}
At zero momentum,
\(J_\lambda(0)=\lambda^2(a_\theta-b_\theta)\).
Radiality makes the Hessian of \(J\) scalar. The four-dimensional
radial integral gives
\[
 \Delta_kJ_\lambda(0)=\frac1{8\pi^2}
 \int_1^2r(1-\theta)(\theta''+3\theta'/r)\dd r
 =\frac{E_\theta-1}{8\pi^2}.
\]
Indeed the integrand is the derivative of
\(r(1-\theta)\theta'\), plus
\(2(1-\theta)\theta'+r(\theta')^2\); the integrated first term
vanishes and the second integrates to \(-1\). Hence
\(\mathsf T_2J=\lambda^2(a_\theta-b_\theta)+\kappa_\theta k^2\).
Cauchy--Schwarz applied to
\(1=-\int_1^2\theta'(r)\dd r\) gives
\(1\le E_\theta\log2\), proving \eqref{eq:v61-kappa}.

For \(|k|>4\lambda\), make the change \(y=p+k\) in \(J\).
On the support of \(\Theta_\lambda(y)\), \(\nu_\lambda(y-k)=1\).
The function \(|k-y|^{-2}\) is harmonic in \(y\) on the ball
\(|y|<|k|\). Its mean on every sphere centered at zero is \(|k|^{-2}\).
Radial integration therefore gives exactly
\(J_\lambda(k)=\lambda^4M_\theta/k^2\). Substitution in
\eqref{eq:v61-inherited-J} proves the formula. No asymptotic expansion
of a divergent integral has been used.
\end{proof}

\begin{corollary}[A calculated failure of soft-only iteration]
\label[corollary]{cor:v61-hardclosure}
Close the two scalar legs of this already renormalized subgraph with the
ordinary hard scalar weight \(\nu_\lambda(k)e^{-k^2/\Lambda^2}\), retaining
the correct factor \(n_H/4\). For \(\Lambda\ge8\lambda\),
\begin{align}
 \frac{n_H}{4}\int_k
   \frac{\nu_\lambda(k)e^{-k^2/\Lambda^2}}{k^2}
                    \Pi^{\rm bub,R}_0(k)
 ={}&-\frac{n_H\kappa_\theta}{32\pi^2\chi}\Lambda^6\notag\\
 &-\frac{n_H\lambda^2(a_\theta-b_\theta)}{64\pi^2\chi}\Lambda^4
 +\frac{n_H\lambda^4M_\theta}{64\pi^2\chi}\Lambda^2
 +O(n_H\chi^{-1}\lambda^6).
 \label{eq:v61-closed-divergence}
\end{align}
In particular this contraction is not bounded as the upper scale grows,
despite the subgraph's established soft degree-six remainder.
\end{corollary}
\begin{proof}
Insert \eqref{eq:v61-hard-polynomial} outside the radius-\(4\lambda\)
ball. Extending its three polynomial terms over all momentum uses
\(\int_k e^{-k^2/\Lambda^2}/k^2=\Lambda^2/(16\pi^2)\),
\(\int_k e^{-k^2/\Lambda^2}=\Lambda^4/(16\pi^2)\), and
\(\int_k k^2e^{-k^2/\Lambda^2}=\Lambda^6/(8\pi^2)\).
The alteration inside the fixed ball is \(O(\lambda^6/\chi)\),
as is the integral of the actual smooth subgraph there. Multiplying by
\(n_H/4\) proves the displayed formula.
\end{proof}

This is not an obstruction to renormalization. Its local divergences
belong to the next EFT assignment, and other subgraph counterterms must
be included in the complete graph. It is a direct counterexample to
inferring an iterative hard-contraction bound from the already proved
fixed-soft-window norm alone. In particular, the ordinary massless
simplification in V27 cannot be used to declare the subsequent scalar
closure finite. The diagnostic radial convention is not imposed on the
active finite regulator.

\section{A full-cutoff estimate on the all-comparable two-loop part}
\label{sec:v61-comparable}

We now return to the actual finite regulator. Work at \(z=0\), with two
external metric marks of momenta \(k,-k\), and retain the full matrix of
polarizations. Expand every covariance into the inherited smooth dyadic
pieces at scales \(\Lambda_j\asymp2^j\lambda\), including the periodic
upper pieces. Background derivatives split a covariance into its actual
primitive factors before assigning scales. The contact and two-cubic
cores have at most five such factors after two metric derivatives.

Fix an integer \(b_*\ge2\). The all-comparable part is the sum of the
assigned marked graphs for which the largest and smallest primitive
covariance labels differ by at most \(b_*\). This selection is by labels,
not by deleting a singular momentum region after integration. Denote its
per-volume coefficient by \(\gamma_{a,L,\mathrm{cmp}}(k)\), and set
\[
 \gamma^R_{a,L,\mathrm{cmp}}
   =(I-\mathsf T_6)\gamma_{a,L,\mathrm{cmp}}.
\]
The degree-six local polynomial is stored. It is a component assignment,
not yet a claimed symmetry-restoring common two-loop counterterm.

\begin{lemma}[Native common-scale degree]
\label[lemma]{lem:v61-common-degree}
At a largest common scale \(\Lambda_j=R\), the differentiated ordinary
and finite integrands obey, for every fixed multi-index \(\alpha,\beta\),
\begin{align}
 |\partial_k^\alpha\partial_{p,r}^\beta f_0|
 +|\partial_k^\alpha\partial_{p,r}^\beta f_a|
 &\le C_{\alpha,\beta}\frac{n_H}{\chi}
             R^{-2-|\alpha|-|\beta|},\notag\\
 |\partial_k^\alpha\partial_{p,r}^\beta(f_a-f_0)|
 &\le C_{\alpha,\beta}\frac{n_H}{\chi}
             a^3R^{1-|\alpha|-|\beta|}
 \label{eq:v61-common-symbol}
\end{align}
on the common interior plateau. The normalizations include every
primitive covariance and both loop integrations are still to be taken.
The finite upper pieces satisfy the first bound in their dimensionless
periodic coordinates. Constants depend on \(b_*\), the fixed finite
source degree, profiles and inverse margins, not on \(a,L,j\).
\end{lemma}
\begin{proof}
For the two-cubic graph, three propagators contribute degree minus six
and the two scalar vertices degree four, giving degree minus two.
For the contact, its two propagators contribute minus four and its vertex
degree two, with the same result. A metric derivative can split a scalar
or metric covariance, adding one degree-minus-two inverse and its actual
degree-two Hessian vertex. Differentiating a vertex preserves its degree.
Thus the result is unchanged by either external metric mark.

All routed momenta in a nonzero primitive factor are at scales comparable
to \(R\). Their inverse and derivative bounds are consequently those of
V14 and V27, with no small unprotected denominator. Each external or loop
momentum derivative lowers this degree by one. Native scalar vertices and
free inverses agree to fourth relative mesh order on the common plateau.
A retained metric mark can also consume the corrected gravitational cubic
or quartic Hessian, whose conservative comparison is third order. A finite
product difference therefore has the second bound. The finite completed
vertices and covariance inverses have smooth periodic scaled symbols;
this proves the upper-region estimate. At no point is the signed metric
covariance replaced by a positive tensor: only its finite matrix norm is
bounded. The native local contact is included in this argument.
\end{proof}

\begin{theorem}[All-comparable two-metric source comparison]
\label[theorem]{thm:v61-comparable}
On
\begin{equation}
 \lambda\ge32768\mu>0,\qquad
 a\lambda\le2^{-24},\qquad
 \lambda L_\nu\ge1,\qquad |k|\le\mu,\qquad z=0,
 \label{eq:v61-domain}
\end{equation}
the all-comparable coefficients have an ordinary thermodynamic limit
after their actual degree-six assignment. For \(r=|\alpha|\le7\),
\begin{equation}
 \boxed{
 |\partial_k^\alpha(
     \gamma^R_{a,L,\mathrm{cmp}}
     -\gamma^R_{0,L,\mathrm{cmp}})(h_1,h_2;k)|
 \le C_\alpha\frac{n_H}{\chi}
      a|k|^{7-r}\norm{h_1}\norm{h_2}.}
 \label{eq:v61-comparable-limit}
\end{equation}
The ordinary finite-volume discrepancy is bounded by
\begin{equation}
 C_{\alpha,b}\frac{n_H}{\chi}\lambda^{-1}
       (\lambda L_{\min})^{-b}|k|^{7-r}
       \norm{h_1}\norm{h_2},\qquad b>8.
 \label{eq:v61-comparable-volume}
\end{equation}
These statements are for the declared comparable-scale component, not
the entire two-loop coefficient.
\end{theorem}
\begin{proof}
For a fixed largest label there are only a bounded number of primitive
label assignments, since their number and the allowed spread \(b_*\)
are fixed. Both loop momenta are confined to balls of radius \(CR\).
Their normalized lattice mode count is at most \(CR^8\) on the stated
volume domain. At the seventh external derivative, the first bound of
\eqref{eq:v61-common-symbol} is \(CR^{-9}\), so its integral or sum
is \(CR^{-1}\). The interior difference is \(Ca^3R^{-6}\),
whose normalized two-loop sum is \(Ca^3R^2\).

Sum dyadically over common scales below a fixed multiple of \(a^{-1}\).
The two geometric sums obey
\[
 Ca^3\sum_{R\le c/a}R^2\le Ca,
 \qquad C\sum_{R\ge c/a}R^{-1}\le Ca.
\]
Only a bounded number
of scaled upper labels remain between the common plateau and the edge
of the finite Brillouin zone; their seventh-derivative estimates each
give \(Ca\). Smooth periodicity includes routed wraps at that edge.
Thus no finite upper mode of the comparable component is omitted.
Taylor's integral remainder of order seven gives
\eqref{eq:v61-comparable-limit}. Higher prescribed derivatives follow
from the same estimates; they are needed for the source kernel below.

For the ordinary volume comparison, all seven-times-differentiated
integrands, with any fixed additional loop derivatives, are integrable.
On each scale their Fourier transforms have rapid decay in the two
four-dimensional loop-position variables. Poisson summation in eight
variables bounds the nonzero images by
\(CR^{-1}(RL_{\min})^{-b}\); their sum is bounded by
\(C\lambda^{-1}(\lambda L_{\min})^{-b}\). Taylor's formula gives
\eqref{eq:v61-comparable-volume}. This estimate compares the same
canonical off-grid interpolation and does not assign an exact finite-torus
Ward identity.
\end{proof}

\begin{corollary}[Rooted source norm and summable comparable increments]
\label[corollary]{cor:v61-rooted}
After a fixed smooth external window,
\begin{equation}
 \boxed{
 a^4\sum_x(1+\mu d_{\rm per}(x,0))^b
       \norm{\mathcal K^\Delta_{\mathrm{cmp}}(x)}
 \le C_b\frac{n_H}{\chi}a\mu^7.}
 \label{eq:v61-rooted}
\end{equation}
The complementary common-scale increments have bounds
\(C(n_H/\chi)\mu^7\Lambda_j^{-1}\), hence are summable. One metric
source may be in \(L^1\) and the other in \(L^\infty\), with no
total-volume factor.
\end{corollary}
\begin{proof}
Differentiate the windowed symbol to any fixed order exceeding the four
relative dimensions and the desired moment. The preceding estimates
control each term after multiplication by the window derivatives.
Fourier integration by parts and periodization give
\eqref{eq:v61-rooted}. The scale bound was proved before summing in
\cref{thm:v61-comparable}. Every covariance in every inverse-derivative
return receives a label; using only equal scalar and metric labels would
not give this component.
\end{proof}

\section{The three overlapping channels cannot be replaced by one scalar subtraction}
\label{sec:v61-forest}

Label the internal edges of the two-cubic core by
\(e_g,e_1,e_2\), with the last two scalar. Its three proper one-loop
subgraphs are
\begin{equation}
 \gamma_g=\{e_1,e_2\},\qquad
 \gamma_{s,1}=\{e_g,e_1\},\qquad
 \gamma_{s,2}=\{e_g,e_2\}.
 \label{eq:v61-three-subgraphs}
\end{equation}
They are, respectively, a scalar contribution to the metric two-point
kernel and two metric--scalar contributions to the scalar two-point
kernel. V60 and V27 provide the corresponding primitive operands and
soft one-loop comparisons. Every pair shares an internal edge and neither
member contains the other. They are overlapping, not a nested forest.

\begin{proposition}[Proper-subgraph bookkeeping for the native two-cubic core]
\label[proposition]{prop:v61-forest}
The proper forests of the core are the empty forest and the three
singletons in \eqref{eq:v61-three-subgraphs}. If \(t_\gamma\) denotes the
actual local subgraph assignment through degree four, a graphwise
forest subtraction has the form
\begin{equation}
 \boxed{
 R\mathcal I_\Gamma
 =(1-t_{\Gamma,6})
   (1-t_{\gamma_g,4}-t_{\gamma_{s,1},4}
                         -t_{\gamma_{s,2},4})\mathcal I_\Gamma.}
 \label{eq:v61-forest}
\end{equation}
The formula is an exact finite subtraction prescription. It is not, by
itself, a convergence estimate. Geometric marks must be distributed on
each subgraph and its complement before their source derivatives are
assigned. The distinct two-contact-loop core has its own tadpole and
composite-contact assignments and is not replaced by
\eqref{eq:v61-forest}.
\end{proposition}
\begin{proof}
A proper loop needs two of the three parallel core edges. These give
exactly the three displayed subsets. Each pair shares one edge, so no
forest may contain both. The forest sum therefore has the stated four
terms, followed by the overall degree-six assignment. A background mark
splits an edge into its actual differentiated chain or differentiates a
vertex; it does not justify dropping any of the three core channels.
Each assigned local term must be inserted on the complementary edge with
its original source labels and symmetry factor. The proof is
combinatorial and makes no estimate of the subtracted integral.
\end{proof}

A separated assignment with the metric line at a low scale and both
scalar lines high uses \(\gamma_g\). A separated assignment with one
scalar line low and the other two high uses one of
\(\gamma_{s,1},\gamma_{s,2}\). Renormalizing the scalar determinant
before performing the second loop treats the first channel but does not
automatically implement the other two. This is the concrete analytic
interface missing from a repetition of V60's fixed-soft-window theorem.

The local subgraph momenta include the momentum of the complementary
hard line. They are not restricted by \(\mu\). To finish the full
comparison, one must estimate the signed forest sum in these scale
hierarchies, together with the tadpole/composite-contact subtraction and
the prescribed one-loop finite normalizations. The all-comparable theorem
does not pay that estimate. Conversely, \cref{cor:v61-hardclosure} is
not an unrenormalizability conclusion: it evaluates exactly the term that
makes the further assignment necessary.

The ordinary six-derivative bank and its associated geometry-dependent
source contacts are permitted by the programme's EFT filtration. It is
not legitimate to force them into the one-loop degree-four bank. The
current result does not determine an independent invariant curvature
basis or the symmetry-restoring coefficients at weight six. Those
questions require the complete two-loop action and source complex.

\section{Which two-loop identity would have to be assembled}
\label{sec:v61-master}

Let \(S=S_0+\hbar S_1+\hbar^2S_2+\cdots\) be a finite source-complete
master candidate in one fixed contraction/bracket convention, with the
density assigned once. For a nilpotent canonical coefficient contraction
\(\Delta\), write
\[
 \mathcal A_0=\tfrac12(S_0,S_0),\quad
 \mathcal A_1=(S_0,S_1)-\Delta S_0,\quad
 \mathcal A_2=(S_0,S_2)+\tfrac12(S_1,S_1)-\Delta S_1.
\]
The standard BV Jacobi/product identities give, without assuming
\(\mathcal A_0=0\),
\begin{equation}
 (S_0,\mathcal A_2)+(S_1,\mathcal A_1)
 +(S_2,\mathcal A_0)-\Delta\mathcal A_1=0.
 \label{eq:v61-order-two-consistency}
\end{equation}
This equation is recalled to specify the task, not advertised as a new
abstract BV theorem. In the present finite carrier, V60 already supplies
nonzero pieces of \(\mathcal A_0\). In a lower-reference convention,
the paired extra terms of V59 must be included as well.

Equation~\eqref{eq:v61-core} supplies one previously uncontracted part
of \(S_2\), including its metric-source derivatives. It supplies neither
the entire \(S_2\) nor the two other contributions to \(\mathcal A_2\).
The ghosts and scalar antifields that vanish on the present action panel
still participate in the full identity. In particular a contraction of
the scalar determinant in isolation is not a test of
\eqref{eq:v61-order-two-consistency}.

A complete forest computation must use the actual previously chosen
\(S_1\), not independently refit its metric and scalar local terms.
The forest formula with zero finite additions is only a comparison
renormalization prescription. Adopting it would be a change of scheme,
which has not been made. The fixed gravitational normalization, V34's
subcritical scalar prescription, V38's logarithmic packet, and the
once-counted measure remain unchanged in this installment.

Lattice power counting and local forest renormalization are established
frameworks \cite{V61ReiszPower,V61ReiszRen}. The scale organization of high
subgraphs is likewise standard \cite{V61KRT}. These references are used
for comparison and terminology only. Their hypotheses are not declared
verified for the complete native two-loop forest on the strength of the
all-comparable estimate. In particular, the new result is not advertised
as a first construction of perturbative scalar gravity or of its EFT
counterterm tower.

\section{Ledgers, checks and the next overlapping-scale estimate}
\label{sec:v61-ledgers}

\subsection*{Proved}
The complete bare scalar-containing two-loop core through two metric marks
is \eqref{eq:v61-core}; bridge exclusions, all seven first-derivative
terms and scalar cutting factors are explicit. The ordinary guarded
vacuum coefficient is analytically evaluated, and the hard closure of
the inherited soft-renormalized massless bubble has the nonzero
six-power term \eqref{eq:v61-closed-divergence}. The actual finite
all-comparable two-metric component has a degree-six local assignment,
an \(O(a)\) full-cutoff comparison, a separate volume bound and a rooted
source estimate. The three proper overlapping channels of the two-cubic
core and their finite forest prescription are specified.

\subsection*{Inherited}
V14/V24 supply the completed gravitational covariance, native tree jets,
contour and density convention. V27 supplies the completed scalar action,
the signed mixed numerator and contact tensor; its finite source bounds
are consumed at their stated scope. V58--V60 supply scalar elimination,
the determinant, source/reference bookkeeping and populated one-loop
vacuum row. The old results are not independently certified wholesale by
the new checker. No formerly open local cohomology coefficient is assigned
a value here.

\subsection*{Literature input}
Reisz's lattice power-counting and renormalization theorems and the
multiscale high-subgraph framework provide comparison, not an unproved
shortcut for this regulator. The core Wick calculation, hard external
formula and all-comparable estimate are derived above. No external
anomaly cancellation or pure-gravity cohomology theorem is invoked.

\subsection*{Conditional}
The core is a literal finite Gaussian coefficient on the declared scalar
chart and a coefficientwise Wick expression on the inherited metric
contour. Its physical meaning beyond this coefficient order uses the same
perturbative setup as the earlier notes. A complete two-loop renormalized
functional would additionally require the scale-separated forest bounds,
the actual one-loop finite additions and their marked insertions, and the
complete master/reference identity. An overall local polynomial by itself
does not renormalize those subgraphs.

\subsection*{Open}
The next load-bearing analytic calculation is the common signed forest
comparison in the three separated-scale channels of
\eqref{eq:v61-three-subgraphs}, together with the two-contact-loop
composite assignment, using the same one-loop source normalization. The
remaining metric/ghost master-defect integration and the critical local
primitive remain open. The new massless two-loop action component does
not prove an invariant interacting ESM class, an all-orders continuum,
the chiral line theorem, or removal of the positive lower scale.

\begin{center}\small
\begin{tabular}{>{\raggedright\arraybackslash}p{.24\textwidth}>{\raggedright\arraybackslash}p{.68\textwidth}}
\toprule
Variable & Scope in V61\\\midrule
Loop and source order & Bare scalar-containing two-loop action; zero
retained scalar/ghost/antifield, through two metric marks. Counterterm
insertions and the full master row are not claimed included.\\
Cutoff and volume & Exact finite core on odd grids. The analytic
comparison is on \eqref{eq:v61-domain} for the all-comparable component,
with fixed label spread. No total-volume factor in rooted bounds.\\
Mass and couplings & Analytic two-loop theorem and hard-momentum
diagnostic at zero scalar mass; fixed \(\chi>0\), four equal scalars;
internal gauge, Yukawa and scalar self-couplings zero.\\
Scheme & Existing local choices are retained. Gaussian upper guard and
radial lower profile are labelled ordinary diagnostics, not new native
profiles. The compared common-scale component is not a separate physical
observable.\\
Uniformity & Fixed metric/source order, shell spread, derivative order,
profile seminorms and inverse margins. No bound uniform in growing EFT
order or arbitrary numbers of interacting RG steps.\\
Numerical status & Exact rational/symbolic combinatorics and analytic
ordinary integrals. Native finite profile integrals and the complete
renormalized forest have not been numerically evaluated.\\\bottomrule
\end{tabular}
\end{center}

The self-contained checker verifies the full finite Gaussian cumulant
against the core plus its bridge term, also with signed metric covariance;
tests the mixed gravitational-cubic bridge; checks both scalar cutting
weights; reconstructs the seven-term background derivative with
noncommuting matrices; verifies the inherited signed vertex contraction,
the ordinary Gaussian integral constants and profile integration identity;
and checks the overlap graph and common-scale geometric sums. The append
and label audits are separate. These checks are not a proof-assistant
formalization or a verification of all earlier analytic claims.

\begingroup
\renewcommand{\refname}{References for V61}
\begin{thebibliography}{9}\small
\bibitem{V61ReiszPower} T.~Reisz,
\emph{A power counting theorem for Feynman integrals on the lattice},
Commun. Math. Phys. \textbf{116} (1988), 81--126,
doi:10.1007/BF01239027.
\bibitem{V61ReiszRen} T.~Reisz,
\emph{Renormalization of Feynman integrals on the lattice},
Commun. Math. Phys. \textbf{117} (1988), 79--108,
doi:10.1007/BF01228412.
\bibitem{V61KRT} T.~Krajewski, V.~Rivasseau and A.~Tanasa,
\emph{Combinatorial Hopf algebraic description of the multiscale
renormalization in quantum field theory}, arXiv:1211.4429.
\end{thebibliography}
\endgroup

\section*{Append-only continuation marker: V61}
\addcontentsline{toc}{section}{Append-only continuation marker: V61}
\label{sec:v61-continuation-marker}
\textbf{End of V61.} Preserve Sections 1--475. The second hard loop has
been taken in the scalar-containing action sector, with its actual
contact and three-propagator terms and metric derivatives. The
all-comparable source limit is controlled, while a calculated hard
closure shows why the scale-separated forest must still be paid. The
next step is that common forest estimate and its prescribed source/master
insertions, not another conditional scalar Gaussian identity. The endpoint
remains source-complete measured Einstein--Standard Model EFT continuum
control at every prescribed finite order.


\clearpage
\section{The separated forest: locality must be tested after contraction}
\label{sec:v62-start}

Sections 1--475 and all earlier continuation markers are preserved. This
installment addresses the three overlapping proper subgraphs of
\cref{eq:v61-three-subgraphs}. It does not repeat the Wick calculation,
the hard-closure counterexample, or the bare all-comparable estimate of
V61. The new object is the \emph{signed forest of the three-propagator
core}, including its zero, first and second ordinary metric derivatives.
The nontrivial complementary estimate is stated for two metric marks of
momenta $k,-k$. Four equal real scalars, zero scalar mass, the actual
V14/V24 metric covariance and signed contour, the V27 scalar action,
and the once-counted measure are retained. Internal gauge, Yukawa and
scalar self-couplings remain zero. A common factor $\hbar^2$ is removed.

The important extra region is not a region of the bare graph. A Taylor
subtraction can survive on a scale assignment on which the original
three-line product is zero. Dropping those assignments would change the
fixed local subtraction into a scale-restricted prescription. The proof
below keeps them and does not assume that a subgraph's remaining external
momentum is in the original source window.

The result is a graphwise renormalization theorem, not a complete two-loop
Einstein--SM theorem. The distinct two-loop contact core in
\cref{eq:v61-core}, its composite-contact normalization, and the remaining
master/reference rows are not included. Local Taylor coefficients and
finite additions are retained as data. Zero finite additions will be used
only to define a comparison forest; the active common one-loop
normalization is not reset. A separate theorem controls any specified
bounded finite local additions in this core.

\subsection*{Inherited symbol input and terminology}
The completed propagator and tree-vertex estimates are those of
\cref{lem:v27-symbols,lem:v61-common-degree}. In particular, a propagator
has order $-2$, a scalar or metric Hessian vertex has order at most two,
and an inserted metric vertex paired with its extra covariance preserves
that degree. On the common interior the conservative relative mismatch
is $O(a^3R^3)$; the stronger scalar-only fourth-order comparison is not
needed. The estimates hold with any fixed number of momentum derivatives.
No continuum Ward identity enters them. A factor $n_H/\chi$ remains after
all primitive contractions and source insertions of this core.

Lattice local forest subtraction and high-subgraph organization are
established frameworks \cite{V61ReiszRen,V61KRT}. The proof below is a
two-loop verification using the existing smooth native symbol bounds,
not an invocation of those general theorems for an unchecked regulator.

\section{Assigned forests and the counterterm-only region}
\label{sec:v62-assignments}

Expand every background derivative in \cref{eq:v61-core} before assigning
scales. In the two-cubic core, an edge differentiated in a metric mark
becomes a chain of actual covariances and Hessian vertices. There are
three core chains, denoted $g,1,2$, with at most five primitive covariance
factors in total. A derivative on an end vertex does not create a new
cycle. The three proper cycles are still the two-chain subsets
$\gamma_g,\gamma_{s,1},\gamma_{s,2}$.

Use a smooth dyadic decomposition with scales $R_j\asymp2^j\lambda$,
$j\ge0$. It can be chosen identical for the finite and ordinary kernels
on the common interior; the bounded number of periodic upper pieces
have scale comparable to $a^{-1}$. This is a decomposition of the existing
covariances, not a new Gaussian or ultraviolet cutoff. A fixed support
constant controls the annular thickness and the top periodic pieces.

On the fixed soft panel the primitive labels in one chain differ by at
most a fixed integer $b_0$. Indeed, their momenta differ only by sums of
at most two soft external momenta. Terms violating this condition vanish,
also after the external Taylor derivatives used below. Group each chain
by a representative scale and absorb the bounded multiplicity into the
constants. Choose a fixed separation integer $b_*$ larger than the
support and routing constants needed below. Constants can depend on this
choice, but never on the number of ultraviolet shells.

For one assigned differentiated graph let $\mathcal I_\mu(k;p,r)$
include its original sign, Wick factor, tensor contractions, and both loop
measures. When displaying integrand bounds, the loop measures will be
written separately. For a proper cycle $\gamma$, factor its internal
word from the complementary chain before Taylor expansion. Its independent
external variables $U_\gamma$ are one boundary momentum and all metric
momenta attached to the cycle. One external boundary momentum is eliminated
by conservation. The complementary hard momentum is therefore a variable
of $U_\gamma$, not a soft parameter.

Let $t_{\gamma,4}$ be the total Taylor polynomial through degree four in
$U_\gamma$, acting on this complete internal word only. It leaves the
complementary chain unchanged. Its coefficients are differentiated finite
loop sums, not fitted constants. Metric marks and coframe contacts are
included before this operation. Define the fixed comparison forest by
\begin{equation}
 \boxed{\mathcal F_{a,L}(k)
 =(1-\mathsf T_{6,k})\sum_\mu
 \left(1-\sum_{\gamma\in\{\gamma_g,\gamma_{s,1},\gamma_{s,2}\}}
                  t_{\gamma,4}\right)\mathcal I_{a,\mu}(k).}
 \label{eq:v62-forest-definition}
\end{equation}
There are no products of overlapping proper Taylor operators. At finite
cutoff this is the prescription of \cref{prop:v61-forest}, with source
routing now specified. The polynomial insertion first uses the common
unwrapped local momentum chart; \cref{sec:v62-realization} separately
constructs a periodic finite-range representative and bounds its effect.
For the ordinary theory start with a finite shell cutoff and take its
limit only after the signed forest is formed.

\begin{lemma}[Support after localization]
\label[lemma]{lem:v62-support}
For every assigned marked core, the following statements hold on the
fixed soft window, including all external derivatives on that window.
\begin{enumerate}[label=\textup{(\roman*)},leftmargin=2.8em]
\item A nonzero bare graph has its two largest core-chain scales
comparable within a fixed support constant.
\item A nonzero $t_{\gamma,4}\mathcal I_\mu$ has all the internal
chain scales of $\gamma$ comparable. No comparison between those scales
and the complementary-chain scale is forced.
\item If two chains have scale $R$ and the third scale $r\le2^{-b_*}R$,
the signed proper forest is exactly $(1-t_{\gamma,4})\mathcal I_\mu$,
where $\gamma$ is the high two-chain cycle. The other proper Taylor terms
vanish.
\item If just one chain has scale $R$ and the other two have comparable
scale $r\le2^{-b_*}R$, the bare graph is zero, but the forest can equal
$-t_{\gamma,4}\mathcal I_\mu$, with $\gamma$ the low two-chain cycle.
If no two scales are comparable, all terms vanish.
\end{enumerate}
\end{lemma}
\begin{proof}
At an original vertex the three core momenta sum to external momenta of
size at most $2\mu$. The triangle inequality, applied in the periodic
momentum group with its nearest-representative distance, excludes one
large momentum balanced by two much smaller ones. Annular support converts
this into (i). Along a marked chain the same argument uses differences
of size at most $2\mu$.

For (ii), evaluate the proper Taylor coefficient before reinserting its
polynomial. At $U_\gamma=0$, every internal momentum of the one-loop
cycle is $p$ or $-p$. Momentum derivatives do not enlarge the support of
a smooth shell profile. Its internal shell labels must therefore overlap
at this common magnitude. The complementary momentum is not set to zero
in the complementary chain, which is outside the Taylor operation.
Statements (iii) and (iv) follow by applying (i) and (ii) to all three
cycles. A finite buffer $b_*$ handles the annular boundaries, the bounded
chain-label spreads and the periodic upper pieces. Assignments in that
buffer are included in the comparable class. No off-grid change of loop
variables by a nonlattice external shift is needed.
\end{proof}

\begin{proposition}[An actual scalar-vertex test of subtraction-created support]
\label[proposition]{prop:v62-created-support}
The conformal scalar bubble inside $\gamma_g$ can have zero bare assigned
integrand and a nonzero assigned Taylor counterterm on the same pair of
loop labels.
\end{proposition}
\begin{proof}
The ordinary conformal vertex from V60 is $p\cdot(p+q)$ at zero mass.
Select $p=e_0$ and $q=t e_1$. Before shell factors its two-vertex scalar
word is
\[
 \frac{[p\cdot(p+q)]^2}{p^2(p+q)^2}=\frac1{1+t^2}.
\]
Choose the two scalar shells with a plateau at $|p|=1$ and support below
$|p|=2$, and a complementary metric shell at $t=64$. The bare word is
zero because $|p+q|>2$. At $q=0$ both shell profiles are identically one
near the differentiated point. The fourth Taylor polynomial along this
line is $1-t^2+t^4$, nonzero at $t=64$. Multiplication by the nonzero
complementary metric component leaves a nonzero subtraction term. This
is a component/support witness built from the native limiting scalar
vertex, not a numerical value of a full two-loop amplitude.
\end{proof}

In particular, the scale split in \cref{lem:v62-support} is applied to
\emph{all four signed forest terms}, not only to assignments supported
by the original graph. This distinction is what retains the fixed
one-loop subtraction rather than silently changing to running high-part
subtractions only.

\section{Two-scale Taylor bounds with hard boundary momenta}
\label{sec:v62-two-scale}

Set $g_*=n_H/\chi$ and suppress the product of the two metric-polarization
norms until the final statements. Norms of tensor factors are taken only
after the indicated proper-forest cancellation. The constants below
include the finite number of tensor components and source placements.
All loop sums are normalized per physical volume.

\begin{lemma}[Native two-scale operands]
\label[lemma]{lem:v62-operands}
Let a cycle have comparable internal scale $R$ and external variables
$|U|\le\varepsilon_0R$, where $\varepsilon_0$ is fixed by $b_*$. Write
its unintegrated word as $H_{a,R}(p;U)$. After stripping the appropriate
fixed coupling/flavour factor, for every fixed derivative order,
\begin{align}
 |\partial_U^\alpha\partial_p^\beta H_{a,R}|+
 |\partial_U^\alpha\partial_p^\beta H_{0,R}|
 &\le C R^{-|\alpha|-|\beta|},\\
 |\partial_U^\alpha\partial_p^\beta(H_{a,R}-H_{0,R})|
 &\le C a^3R^{3-|\alpha|-|\beta|}
 \label{eq:v62-high-word}
\end{align}
on the common interior. The derivative support has normalized volume
at most $CR^4$. A complementary chain at scale $r$ has order $-2$,
with one inverse power of $r$ for every momentum derivative and a
conservative relative mismatch $Ca^3r^3$.

The integrated degree-$d$ Taylor coefficient of a cycle at scale $r$
obeys
\begin{equation}
 |a_{a,\gamma,r}^{[d]}|+|a_{0,\gamma,r}^{[d]}|
 \le Cr^{4-d},\qquad
 |a_{a,\gamma,r}^{[d]}-a_{0,\gamma,r}^{[d]}|
 \le Ca^3r^{7-d},\quad 0\le d\le4.
 \label{eq:v62-local-shell-coefficients}
\end{equation}
The finite upper pieces obey the corresponding scale bounds without
requiring an interior comparison there.
\end{lemma}
\begin{proof}
A two-chain cycle has two order-minus-two propagators and two order-two
end vertices. Its unintegrated degree is zero. An inserted metric mark
on a chain adds an order-minus-two propagator and its native order-two
Hessian vertex; a mark on an end vertex preserves its order. These are
exactly the differentiated factors of V61. This proves the degree with
at most two marks. In the declared high-to-external ratio, every internal
argument remains at scale $R$ along $U\mapsto tU$, $0\le t\le1$.
The differentiated native bounds used in V27/V61 consequently apply
with $U$ as a variable throughout this larger window, not merely with
$|U|\le\mu$. This proves the first line. Telescoping the product and
using the conservative third relative mesh order gives the second.

The complementary chain has one fewer order-two vertex than
propagators, hence degree $-2$. All its arguments are within $2\mu$
of its loop momentum. Mode counting gives $CR^4$ and $Cr^4$ on
$\lambda L_\nu\ge1$. Differentiation of the high word at $U=0$ and
integration proves \eqref{eq:v62-local-shell-coefficients}. The upper
bounds follow instead from the inherited smooth periodic dimensionless
symbols. This is a direct use of the native vertex estimates; no
hard-momentum estimate is inferred from a previously soft-renormalized
one-loop functional.
\end{proof}

\begin{lemma}[Cancellation estimate in the two-high region]
\label[lemma]{lem:v62-high-remainder}
For a high cycle of scale $R$ and complementary scale
$\lambda\le r\le2^{-b_*}R$, the integrated signed forest contribution
$A^H_{a,R,r}(k)$ satisfies, for $|\alpha|=7$,
\begin{align}
 |\partial_k^\alpha A^H_{a,R,r}|+
 |\partial_k^\alpha A^H_{0,R,r}|&\le Cg_*R^{-1},\\
 |\partial_k^\alpha(A^H_{a,R,r}-A^H_{0,R,r})|
 &\le Cg_*a^3R^2
 \label{eq:v62-high-remainder}
\end{align}
on the common interior. Each additional fixed external or slow-loop
momentum derivative can be paid by $r^{-1}$, and each high-loop
derivative by $R^{-1}$. The generic first bound also holds on finite
upper pieces.
\end{lemma}
\begin{proof}
Form $H(U)-\mathsf T_{4,U}H(U)$ before estimating. Taylor's integral
formula, followed by differentiation, gives for $t$ derivatives in the
slow/external variables
\[
 \left|\partial_U^t(1-\mathsf T_4)H_{0,R}\right|
 \le CR^{-5}r^{5-t}.
\]
For $t>5$, the direct bound $CR^{-t}$ implies this weaker bound because
$r\le R$. Derivatives of the Taylor polynomial above its degree are
zero. The same argument for the difference gives
$Ca^3R^{-2}r^{5-t}$. Here $U$ includes the complementary loop momentum
and the marks on the cycle, so $|U|\le Cr$ along the Taylor segment.

Distribute seven external derivatives between the high remainder and
the complementary chain. Their product is bounded by
\[
 (R^{-5}r^{5-t})(r^{-2-(7-t)})=R^{-5}r^{-4}.
\]
The normalized two-loop count $R^4r^4$ gives $R^{-1}$. If the high
factor is a lattice difference it gives $a^3R^2$ after both counts.
If the complementary factor is the difference, the result is
$a^3r^3/R\le a^3R^2$. The estimates are unchanged by the finite
Leibniz sum. Additional derivatives are distributed identically.
The cancellation in this proof is the actual high-cycle subtraction;
bounding the bare and Taylor terms separately would leave $R^4$.
\end{proof}

\begin{lemma}[The subtraction-created region is also summable]
\label[lemma]{lem:v62-counterterm-only}
Let the sole high chain have scale $R$ and the Taylor-collapsed cycle
scale $r\le2^{-b_*}R$. For its degree-$d$ local coefficient,
$0\le d\le4$, the signed term $A^L_{a,R,r,d}$ satisfies
\begin{align}
 |\partial_k^\alpha A^L_{a,R,r,d}|+
 |\partial_k^\alpha A^L_{0,R,r,d}|
 &\le Cg_*R^{-1}(r/R)^{4-d},\\
 |\partial_k^\alpha(A^L_{a,R,r,d}-A^L_{0,R,r,d})|
 &\le Cg_*a^3R^2(r/R)^{4-d},\\[-3pt]
 &\hspace{22mm}|\alpha|=7.
 \label{eq:v62-counterterm-only}
\end{align}
The second bound is an interior bound; the first controls upper pieces.
\end{lemma}
\begin{proof}
The integrated cycle coefficient is bounded by $Cr^{4-d}$.
Its reinserted polynomial, times the complementary order-minus-two
chain, has degree $d-2$ in its scale-$R$ arguments. After seven
external derivatives and the remaining loop integration, its bound is
$CR^{d-5}$. The product is $R^{-1}(r/R)^{4-d}$.
The inner mismatch replaces $r^{4-d}$ by $a^3r^{7-d}$; the outer
mismatch contributes $a^3R^3$. Both are bounded by the second displayed
quantity. This proof estimates a \emph{counterterm-only} term and does
not need the bare graph to have support on the same assignment.
\end{proof}

All-comparable forest terms satisfy the same bounds with $r\asymp R$.
For the bare term this is \cref{lem:v61-common-degree}. For each proper
Taylor term it follows from $R^{4-d}R^{d-5}=R^{-1}$ and the same
product-difference calculation. This additionally pays the comparable
\emph{counterterm} pieces, which were not part of the bare comparable
coefficient in V61.

\begin{remark}[Why the fourth local degree is retained]
The $d=4$ term in \eqref{eq:v62-counterterm-only} has no geometric
suppression in $r/R$. There is no general license to delete it. For
example, away from the ends of an ordinary scalar shell the conformal
bubble has the angular expansion
\[
 \operatorname{av}_{\omega\in S^3}
 \frac{[\omega\cdot(\omega+tk)]^2}{|\omega+tk|^2}
 =1-\frac34t^2|k|^2+\frac14t^4|k|^4+O(t^6|k|^6).
\]
Its fourth homogeneous density is nonzero. This is an elementary vertex
check, not a recomputation of a full native renormalization constant.
The logarithm in the bound below comes from counting potentially
marginal scale separations; no logarithmic improvement is assumed.
\end{remark}

\section{Full-cutoff convergence of the three-propagator forest}
\label{sec:v62-continuum}

Write $\mathfrak L_a=1+\log(1/(a\lambda))$ on the small-cutoff
domain. The same constants apply to the finite list of placements of
two labelled metric marks. The root convention for all off-grid
interpolations is held fixed.

\begin{theorem}[Forest-complete two-metric estimate for the three-line core]
\label[theorem]{thm:v62-forest-limit}
On
\begin{equation}
 \boxed{\lambda\ge32768\mu>0,\qquad
 a\lambda\le2^{-24},\qquad \lambda L_\nu\ge1,
 \qquad |k|\le\mu,\qquad m_H=0,}
 \label{eq:v62-domain}
\end{equation}
let $\mathcal F_{a,L}$ be \eqref{eq:v62-forest-definition}.
The ordinary finite-shell sequence converges in all retained source
seminorms to $\mathcal F_{0,L}$. For $j=|\alpha|\le7$,
\begin{equation}
 \boxed{
 |\partial_k^\alpha(\mathcal F_{a,L}-\mathcal F_{0,L})
             (h_1,h_2;k)|
 \le C_\alpha\frac{n_H}{\chi}
 a\mathfrak L_a|k|^{7-j}\norm{h_1}\norm{h_2}.}
 \label{eq:v62-forest-rate}
\end{equation}
The ordinary thermodynamic comparison is
\begin{equation}
 |\partial_k^\alpha(\mathcal F_{0,L}-\mathcal F_{0,\infty})
             (h_1,h_2;k)|
 \le C_{\alpha,b}\frac{n_H}{\chi}\lambda^{-1}
 (\lambda L_{\min})^{-b}|k|^{7-j}\norm{h_1}\norm{h_2},
 \quad b>8.
 \label{eq:v62-forest-volume}
\end{equation}
Every prescribed higher derivative is controlled. This theorem includes
all three proper cycles, every relative scale assignment, and the full
periodic ultraviolet region of the three-propagator core. It does not
include the separate two-contact-loop core or a complete master equation.
\end{theorem}
\begin{proof}
Use \cref{lem:v62-support} to divide the \emph{signed} forest into
comparable, two-high, and counterterm-only terms. Outside these classes
it vanishes. Bounded intra-chain and high-pair label multiplicities are
absorbed in $C$. At largest scale $R_j\asymp2^j\lambda$, there are
at most $C(j+1)$ lower-scale choices. The preceding lemmas give the
seventh-derivative majorant
\begin{equation}
 Cg_*(j+1)R_j^{-1}
 \label{eq:v62-shell-bound}
\end{equation}
for the ordinary forest and the generic finite bound. The sharper
$(r/R)^{4-d}$ factor is available in the counterterm-only regions but
is not needed for $d=4$.

The dyadic sum converges: with $R_j=2^j\lambda$,
\[
 \sum_{j=0}^\infty(j+1)R_j^{-1}=4\lambda^{-1},\qquad
 \sum_{j>J}(j+1)R_j^{-1}
 =\lambda^{-1}(J+3)2^{-J}.
\]
This proves convergence of the differentiated forest, without ever
assigning a separate value to a divergent infinite local subgraph jet.
The finite counterterms are formed at finite ultraviolet cutoff first.

For scales wholly in the common interior, the difference is bounded by
$Cg_*a^3(j+1)R_j^2$. If $R_J\asymp a^{-1}$,
\[
 a^3\sum_{j\le J}(j+1)R_j^2
 \le Ca^3(J+1)R_J^2\le Ca(1+J).
\]
A fixed number of top finite labels contributes at most
$Ca(1+J)$ by \eqref{eq:v62-shell-bound}; the ordinary tail contributes
the same order by the exact tail sum. Interior and upper pieces may
be separated with an identical common smooth partition. It is not
necessary to compare their symbols pointwise outside that common region.
Thus no routed Brillouin-zone corner or subtraction-created assignment
is omitted. Taylor's integral remainder in $k$ of order seven proves
\eqref{eq:v62-forest-rate}.

For higher derivatives, the two-scale lemmas pay additional inverse
powers of the lower scale. Summation gives the required fixed-window
bounds; for example the order-$j>7$ mesh difference is bounded by
$Ca\mathfrak L_a\lambda^{7-j}$. These estimates also justify every
fixed source-kernel differentiation.

For volume comparison work first with each ordinary smooth assigned
forest term after the high cancellation. Integer unimodular changes of
the two loop coordinates can place the high and low loop directions
in separate variables. Soft external shifts remain arguments, not
changes of the summation lattice. Arbitrarily many loop derivatives
have bounds $R^{-1}r^{-b}$ in the separated terms. Poisson summation
in the two four-dimensional loop variables gives a summable estimate
$CR^{-1}(rL_{\min})^{-b}$; in the comparable case use $r\asymp R$.
Now sum first over $R\ge r$ and then over $r\ge\lambda$.
This is bounded by $C\lambda^{-1}(\lambda L_{\min})^{-b}$ for $b>8$.
Applying the external Taylor remainder proves
\eqref{eq:v62-forest-volume}. This is a comparison of the same
interpolation convention, not an exact finite-torus Ward identity.
\end{proof}

\begin{corollary}[Rooted sources and the largest-scale remainder]
\label[corollary]{cor:v62-rooted}
After a fixed smooth external window,
\begin{equation}
 \boxed{a^4\sum_x(1+\mu d_{\rm per}(x,0))^b
 \norm{\mathcal K^\Delta_{a,L}(x)}
 \le C_b\frac{n_H}{\chi}a\mathfrak L_a\mu^7.}
 \label{eq:v62-rooted}
\end{equation}
The renormalized largest-scale coefficient is bounded by
$C(n_H/\chi)\mu^7(j+1)\Lambda_j^{-1}$. It is summable and its
ultraviolet tail above $\Lambda_J$ is bounded by
$C(n_H/\chi)\mu^7\Lambda_J^{-1}
[1+\log(\Lambda_J/\lambda)]$.
\end{corollary}
\begin{proof}
Use the high external derivatives in \cref{thm:v62-forest-limit},
rescale by the fixed source window, Fourier integrate by parts and
periodize. This is the same rooted construction as V61, now applied to
the entire signed three-line forest. The physical relative measure is
$a^4$, not the number of root cells. The asserted shell and tail bounds
are \eqref{eq:v62-shell-bound} and its geometric sum. One metric
argument may be in $L^1$ and the other in $L^\infty$.
\end{proof}

\section{Finite local normalizations and a periodic contact realization}
\label{sec:v62-realization}

A graphwise Taylor forest with zero finite additions is not the active
common Einstein--scalar normalization. The following statements keep
the finite additions explicit and show what may be supplied without
changing the preceding limit.

For each proper channel let $f_\gamma(U)$ be a prescribed tensor-valued
local polynomial of degree at most four, including the specified metric
marks. Let $\mathcal J_{a,\gamma}[f_\gamma]$ be its insertion in the
complementary chain, followed by $1-\mathsf T_{6,k}$. The multiplicity
and cutting weights are exactly those of its forest term; the two scalar
channels are not identified before those weights are applied. Normalize
the coefficient norm by stripping the native channel coupling factor,
and set
\[
 M_f=\sum_\gamma\sum_{d=0}^4
              \lambda^{d-4}\norm{f_\gamma^{[d]}}.
\]

\begin{proposition}[Prescribed finite additions have controlled insertions]
\label[proposition]{prop:v62-finite-additions}
For $M_f<\infty$, each $\mathcal J_{0,\gamma}[f_\gamma]$ exists in
the same source topology. Its seventh derivative is bounded by
$C(n_H/\chi)M_f\lambda^{-1}$. Using the same finite polynomial in
both regulators gives an $O((n_H/\chi)M_f a|k|^{7-j})$ comparison
for $j\le7$. If $f_a$ and $f_0$ differ, add
\[
 C\frac{n_H}{\chi}\lambda^{-1}M_{f_a-f_0}|k|^{7-j}.
\]
The rooted and finite-volume estimates have the corresponding form.
\end{proposition}
\begin{proof}
At complementary scale $R$, a degree-$d$ polynomial times the marked
chain has integrated seventh derivative bounded by $CR^{d-5}$.
Since $d\le4$, its sum is $C\lambda^{d-5}$. The interior difference
is bounded by $Ca^3R^{d-2}$, whose sum through $a^{-1}$, plus the
upper and ordinary tail, is at most $Ca\lambda^{d-4}$ on $a\lambda$
small. Multiplication by the stated coefficient norm gives the result.
Linearity gives the additional mismatch term. No coefficient of a
previously uncomputed common counterterm is assigned a numerical value.
\end{proof}

This proposition applies when the channel restrictions of the common
order-$\hbar$ action/source functional have actually been supplied.
It does not declare that all its other insertions belong to this graph,
and it does not construct a missing critical primitive. V34's scalar
normalization, the fixed gravitational normalization and V38's packet
remain unchanged.

\subsection*{Counterterms at hard momenta require their own realization bound}
A polynomial symbol in a proper Taylor coefficient is not by itself a
smooth periodic multiplier on the finite momentum torus. For a genuinely
finite-range representative, use V54's simultaneous source-jet section
with matching degree six. For a monomial $U^\beta$, $d=|\beta|\le4$,
write its realization as $\mathfrak r_{a,6}[U^\beta]$. Its weights and
support are fixed once the finite panel and offsets are fixed; all
permutation and reality identities are imposed on the complete array.

\begin{lemma}[Global estimate for the fixed finite stencil]
\label[lemma]{lem:v62-global-stencil}
On $|aU|\le C_0$, for $0\le j\le7$,
\begin{equation}
 |\partial_U^j(\mathfrak r_{a,6}[U^\beta]-U^\beta)|
 \le C_{C_0,\beta}a^{7-d}|U|^{7-j}.
 \label{eq:v62-global-stencil}
\end{equation}
The inequality is in the unwrapped comparison chart. The realized
multiplier itself has the prescribed periodic or offset-periodic
carrier convention.
\end{lemma}
\begin{proof}
V54's exact moment identities make the dimensionless difference vanish
to order seven at the origin. On a fixed dimensionless ball it is a
finite sum of exponentials minus a polynomial. Taylor's integral
remainder bounds it and all seven derivatives by the fixed support and
coefficient norm. Rescaling gives \eqref{eq:v62-global-stencil}.
No smallness of the hard argument relative to $a^{-1}$ is asserted.
\end{proof}

\begin{theorem}[Periodic local insertions have the same forest limit]
\label[theorem]{thm:v62-periodic-insertions}
Realize every proper degree-four Taylor coefficient in
\eqref{eq:v62-forest-definition} by the same degree-six section,
and retain the overall local degree-six assignment. The resulting
finite-range proper-counterterm forest differs from $\mathcal F_{a,L}$
by at most
\begin{equation}
 C_\alpha\frac{n_H}{\chi}
 a\mathfrak L_a|k|^{7-|\alpha|}\norm{h_1}\norm{h_2},
 \qquad |\alpha|\le7.
 \label{eq:v62-realization-rate}
\end{equation}
It therefore has the same limiting three-line source coefficient.
The same construction realizes any fixed finite addition of the
preceding proposition with an $O(a)$ insertion error. If an overall
local degree-six counterterm is retained as an action coordinate rather
than removed by $1-\mathsf T_6$, it may separately be realized by V54's
coefficient-budget theorem with sufficient padding.
\end{theorem}
\begin{proof}
Do not infer this statement from a soft stencil error alone. The actual
integrated degree-$d$ proper coefficient satisfies, by summing
\eqref{eq:v62-local-shell-coefficients} through the cutoff,
\[
 \norm{a_{a,\gamma}^{[d]}}
 \le C\begin{cases}
 a^{d-4},&d<4,\\
 1+\log(1/(a\lambda)),&d=4.
 \end{cases}
\]
After reinsertion at complementary scale $R$, distribute seven
external derivatives between the realization difference and the
complementary chain. Equation~\eqref{eq:v62-global-stencil} gives
an integrated bound $Ca^{7-d}R^2$ times the proper coefficient.
Summing $R\lesssim a^{-1}$ gives $Ca^{5-d}$ times that coefficient.
It is $Ca$ for $d<4$ and $Ca\mathfrak L_a$ for $d=4$.
No relation between the proper and complementary scales is used in
this estimate: every hard contraction of the realization error is paid.
Taylor's integral remainder, higher fixed derivatives, and the rooted
Fourier argument finish the proof. A bounded finite addition gives
$Ca^{5-d}\norm{f_d}\le Ca\lambda^{d-4}\norm{f_d}$.
This theorem is an implementation comparison, not an exact finite BV
identity or a new physical choice of finite normalizations.
\end{proof}


\begin{corollary}[Matched-completion independence for this graph]
\label[corollary]{cor:v62-matching}
Two admitted smooth ultraviolet completions with the same ordinary
vertices, lower-reference profile and supplied finite local additions
have the same limiting three-propagator forest coefficient. Their
finite-cutoff source difference is bounded by the sum of the two
estimates \eqref{eq:v62-forest-rate} and, when used, the two
realization estimates \eqref{eq:v62-realization-rate}.
\end{corollary}
\begin{proof}
Both converge to the same ordinary forest defined at finite shell
cutoff before taking its limit. The common finite additions have the
same limiting one-loop insertions by
\cref{prop:v62-finite-additions}. Apply the triangle inequality to
these already renormalized coefficients, not to divergent bare graphs.
\end{proof}

\section{What is and is not now iterative}
\label{sec:v62-iteration}

The complete three-propagator forest can now be integrated through all
relative scales at the stated order. The analytic result consumes actual
one-loop subgraph jets at hard complementary momenta. It does not apply
a soft one-loop theorem at an argument outside its domain. The
subtraction-created regions are indispensable for keeping the fixed
local scheme: retaining only genuine high subgraphs without their
corresponding local running data would be a different prescription.

The logarithmic factor in \eqref{eq:v62-forest-rate} has a transparent
bound-level origin. At fixed largest scale, the two-high Taylor remainder
has size $R^{-1}$ independently of the lower scale after seven external
derivatives. The marginal local term in a counterterm-only region has
the same property. There can be $O(1+\log(R/\lambda))$ admissible
lower scales. Their sum is still ultraviolet summable. No assertion is
made that this bound is sharp for the full signed tensor amplitude.

The graphwise local data through degree six are not declared an
independent invariant curvature basis. Nor does the forest theorem imply
that the already selected common one-loop action/source coefficients
satisfy every two-loop symmetry constraint. The complete order-two
master coefficient remains the expression in
\cref{eq:v61-order-two-consistency}, with its action, bracket, contact,
reference and measure terms on one common convention.

The remaining contact core has two loops attached to the same native
coframe vertex. Its tadpole/composite assignments must be made on that
full vertex, including translated high--high contacts and differentiated
profiles. It is not obtained by calling the three parallel-edge forest
on a different edge list. After this component is paid, the combined
finite additions from the actual common $S_1$ and the remaining action
and source sectors still have to be assembled. The present theorem is
therefore a genuine two-loop analytic interface, but not an invariant
interacting RG-class theorem.

\section{Ledgers and verification scope}
\label{sec:v62-ledgers}

\subsection*{Proved}
The fixed three-channel proper forest splits into comparable, two-high,
and subtraction-created scale regions. The last region can be nonzero
where the original native scalar-vertex word is zero. Explicit native
two-scale Taylor estimates control every region, giving the
$O(a[1+\log(1/(a\lambda))])$ full-cutoff and rooted comparison of the
three-propagator two-metric coefficient. The ordinary forest has a
separate thermodynamic estimate. Bounded prescribed finite local
additions have controlled complementary insertions. Degree-six
finite-range realization of the proper degree-four jets has the same
limit, with the divergent inner coefficients included in its hard-loop
estimate.

\subsection*{Inherited}
The bare two-loop graph, its original signs and scalar cutting factors,
the full two-metric differentiation, and the source window come from
V61. The scalar and metric symbol estimates come from V14/V24/V27 and
are used at the specific comparable internal arguments established
above. V54 supplies the finite stencil section; its hard insertion is
newly bounded here rather than inferred from its earlier soft estimate.
The once-counted density and active local normalization choices are
unchanged. No inherited analytic theorem is claimed independently
reverified in its entirety by the new checker.

\subsection*{Literature input}
Reisz's lattice forest framework and Krajewski--Rivasseau--Tanasa's
scale-assigned high-subgraph formalism supply context and terminology
\cite{V61ReiszRen,V61KRT}. The elementary support, Taylor and summation
proofs above verify the particular native three-line graph. Neither
an all-orders lattice theorem nor a quantum-gravity anomaly theorem is
invoked as an unchecked premise.

\subsection*{Conditional}
Insertion of a particular physical one-loop normalization uses its
actual local array and the explicit finite-addition bound. The theorem
does not compute a still-unknown critical coefficient. The finite graph
is a coefficientwise Wick expression on the inherited signed metric
contour, not a nonperturbative positive gravitational measure.
Constructing the complete renormalized two-loop action additionally
requires the contact core and the complete common one-loop insertion
ledger. A complete source/master theorem requires all other fields and
source components at that order.

\subsection*{Open}
The smallest remaining analytic graph is the native two-contact-loop
core with its shared-vertex composite and tadpole assignments, followed
by the common $S_1$ insertion matching across the full scalar-containing
two-loop functional. The populated critical local descent, combined
internal/chiral identities, higher-loop closure and filtered interacting
ESM invariant class remain open. None of these follows merely from the
convergence of this three-line forest.

\begin{center}\small
\begin{tabular}{>{\raggedright\arraybackslash}p{.24\textwidth}>{\raggedright\arraybackslash}p{.68\textwidth}}
\toprule
Variable & Scope in V62\\\midrule
Loop/source order & Scalar-containing three-propagator two-loop core,
through two ordinary metric marks; no retained scalar, ghost or antifield.
All source placements in this component are included.\\
Mass/couplings & Zero scalar mass; four equal scalars, fixed $\chi>0$;
internal gauge, Yukawa and scalar self-couplings zero.\\
Cutoff/volume & Whole completed Brillouin zone, all relative scales, and
the explicit domain \eqref{eq:v62-domain}. Rooted bounds have no total
volume multiplier. The volume comparison is independent of the mesh
comparison.\\
Scheme & Fixed Taylor forest is a comparison convention. Any supplied
bounded finite additions are covered without refitting them. Finite-range
proper insertions preserve their prescribed jets and the limiting graph,
not an exact finite-cutoff BV algebra.\\
Uniformity & Fixed source degree, moment/derivative order, shell support
constants, profile seminorms, interpolation convention and inverse margins.
No bound uniform in arbitrary EFT order or interacting RG steps.\\
Computation & Exact Taylor, support/combinatorial, power-counting and
dyadic-sum checks; no numerical native two-loop integral, full contact
forest, or master coefficient is reported.\\\bottomrule
\end{tabular}
\end{center}

The checker verifies exact mixed-variable Taylor remainders, three-channel
forest combinatorics, the subtraction-created support example, the
conformal angular fourth jet, the scale exponents and dyadic sums, a
finite ordered marked-word factorization, and the global finite-stencil
moment identities. These finite tests support the explicit algebra;
they are not proof-assistant formalization of the analytic bounds.
The append-only and old-label audits are independent and record hashes.

\section*{Append-only continuation marker: V62}
\addcontentsline{toc}{section}{Append-only continuation marker: V62}
\label{sec:v62-continuation-marker}
\textbf{End of V62.} Preserve Sections 1--482. The separated overlapping
forest of the three-propagator scalar--metric core is now controlled,
including the counterterm-only scale regions and hard realization of its
local jets. The remaining contact core and common action/source matching
are explicit next tasks. This is a finite-order source-complete EFT
step, not a full interacting Einstein--Standard Model continuum claim.


\clearpage
\part*{V63: Shared-root contact subtraction and the two-loop scalar--metric coefficient}
\addcontentsline{toc}{part}{V63: Shared-root contact subtraction and the two-loop scalar--metric coefficient}

\section{The shared-vertex graph requires a composite normalization}
\label{sec:v63-start}

Sections 1--482, including their continuation markers, are preserved. The
remaining graph in \cref{eq:v61-core} is
\begin{equation}
 \mathcal C_a(q)=\frac{n_H}{4}\sum_{I,J}(G_g(q))_{IJ}
                   \Tr\bigl(G_H(q)Q_{IJ}(q)\bigr),\qquad n_H=4.
 \label{eq:v63-bare-contact}
\end{equation}
Its metric and scalar loops meet at one scalar--metric vertex. The
three-overlapping-cycle formula of V62 does not define its subtraction.
The appropriate upstream object is the pair of rooted composite kernels
carried by that vertex. This part constructs their subtraction and proves
convergence of the resulting complete contact component, with zero, one
and two ordinary metric marks. The nontrivial source estimate concerns
two marks of momenta $k,-k$. A common factor $\hbar^2$ is removed.

The input is unchanged: four equal real massless scalars, the completed
V27 scalar action, the V14/V24 metric covariance and its signed contour,
the same coframe chart and physical pairing, and the once-counted
measure. Internal gauge, Yukawa and scalar self-couplings are zero in
this component. No independent ghost or antifield mark is added here.
The exclusion of bridge graphs and the two-loop Wick factor are inherited
from \cref{lem:v61-bridge,thm:v61-core}; they are not proved again.

The new subtraction is an explicit \emph{rooted-composite comparison
scheme}. It includes local prescriptions for both tadpole tensors.
A supplied common one-loop normalization can be compared with it, but
an unspecified common normalization is not declared equal to it. In
particular, the theorem does not decide the still-open critical BV
primitive. Locality of a subtraction, existence of a graphwise limit,
and consistency of one common action/source counterterm are separate
requirements of the supplied EFT programme.

Normal products and forest subtraction are established methods
\cite{V63ZimmermannNormal,V63Blaschke}. The new argument below uses a
factorization of this particular completed native vertex and proves the
required bounds directly. It does not invoke a general forest theorem
for an unchecked finite regulator. No additional finite-range
realization or change to an active counterterm is made.

\section{Exact factorization on the native root, including its filters}
\label{sec:v63-factorization}

Use the pointwise tensor $W(q)=\det(e(q))G(q)^{-1}$ of
\cref{sec:v27-scalar}. In the massless action its three interaction
channels are
\[
 \sigma\in\{b,nS,bS\},\qquad
 \varepsilon_\sigma=(1,1,-1),\qquad
 F_b=I,\quad F_{nS}=F_{bS}=S.
\]
The scalar bilinear tensors are $D^b$, $\widehat D^n$, and $D^b$,
respectively, with the signs in \cref{eq:v27-native-scalar-action}.
Let $q_\sigma=F_\sigma q$. The covariance in \emph{every} channel is
the same actual background covariance $G_H(q)$, not a covariance
reconstructed from that individual channel.

\Needspace{7\baselineskip}
In the following kernels $I,J$ label the ten metric components; their
site indices have been summed into the coordinate kernels. For a root $x$, define
\begin{align}
 M_{a,\sigma}^{IJ}(q;x)
   &=\bigl[(F_\sigma\otimes F_\sigma)G_g^{IJ}(q)\bigr](x,x),\notag\\*
 H_{a,\sigma}^{\mu\nu}(q;x)
   &=\left[\mathscr D_{a,\sigma}^{\mu\nu}
          (F_\sigma\otimes F_\sigma)G_H(q)\right](x,x),
 \label{eq:v63-rooted-kernels}
\end{align}
where the bilinear operator $\mathscr D$ has Fourier symbol
$-D^b_{\mu\nu}(p,r)$ or $-\widehat D^n_{\mu\nu}(p,r)$ as applicable.
Covariance kernels in this formula are converted to physical coordinate
kernels: one loop sum is $V_L^{-1}\sum_p$, not an unnormalized count.
All tensor contractions are algebraic. Complex conjugation is used only
in estimates.

Set
\[
 B_{\sigma,IJ;\mu\nu}(q;x)
       =\partial_I\partial_JW^{\mu\nu}(q_\sigma(x)).
\]
The hard-leg filters are already included in $M$ and $H$; there is no
second application to $B$. Background marks differentiating $B$ do
include their own $F_\sigma$.

\begin{theorem}[Native contact factorization]
\label[theorem]{thm:v63-factorization}
The complete finite coefficient \eqref{eq:v63-bare-contact} equals
\begin{equation}
 \boxed{\mathcal C_a(q)
   =\frac{n_H}{4}\int_a\sum_{\sigma,I,J,\mu,\nu}
      \varepsilon_\sigma B_{\sigma,IJ;\mu\nu}(q;x)
         M_{a,\sigma}^{IJ}(q;x)H_{a,\sigma}^{\mu\nu}(q;x).}
 \label{eq:v63-factorization}
\end{equation}
The two loop momenta are independent. This is a finite identity before
subtraction or integration of further fields. It includes the complete
coframe Hessian and every native upper-profile factor.
\end{theorem}
\begin{proof}
Differentiate the single action \eqref{eq:v27-completed-Higgs} twice in
hard metric directions. Its free term is metric independent. In each
of the three interaction channels, both metric derivatives act on the
pointwise $W(q_\sigma(x))$ and give $B_\sigma$, with one
$F_\sigma$ on each differentiated field. Wick pairing those two fields
gives $M_\sigma$. Pairing the two scalar fields in the same vertex gives
$H_\sigma$. There is no contraction joining a scalar and a metric, so
the sums factor at the root. The coefficient $n_H/4$ is exactly the
inherited quartic Wick factor.

The scalar bilinear is a finite sum of separated one-leg operators.
Off-diagonal components of \eqref{eq:v27-Dnative} already factor. The
diagonal cosine and hyperbolic cosine split by their addition formulas;
every monomial in $a^2E_2$ also factors. Thus no scalar-loop momentum
is hidden in $M$, or metric-loop momentum in $H$. In the filtered
channels the profiles act individually on the legs, exactly as in the
native action. No profile at the sum of a metric and scalar hard
momentum is inserted. The root imposes only the existing total momentum
conservation. This proves the finite identity, including temporal
continuation and translated scalar factors.
\end{proof}

\begin{proposition}[Labelled geometric differentiation]
\label[proposition]{prop:v63-marks}
For a labelled mark set $E$ with $|E|\le2$, the derivative of each
summand in \eqref{eq:v63-factorization} is
\begin{equation}
 \sum_{E_0\sqcup E_g\sqcup E_H=E}
 B_\sigma^{[E_0]}M_{a,\sigma}^{[E_g]}H_{a,\sigma}^{[E_H]}.
 \label{eq:v63-partitions}
\end{equation}
The first derivative has three terms and the second has nine labelled
terms. In differentiating $M$ and $H$, respectively, use
$D_hG_g=-G_gB_h^gG_g$ and $D_hG_H=-G_HQ_hG_H$ with the actual completed
vertices; their second derivatives include both inverse returns and the
second Hessian derivative.
\end{proposition}
\begin{proof}
Apply the labelled product rule to the exact finite identity. Assigning
each mark to one of three factors gives $3^{|E|}$ terms. Inverse
differentiation is an identity on the hard-support resolvent, not an
inverse on retained zero modes. The endpoint filters are fixed; their
momentum arguments still participate in every momentum differentiation.
\end{proof}

On the soft window all external mark profiles equal one. Consequently
$B_\sigma^{[E_0]}$ is the same local, momentum-independent coframe
tensor in all three channels. It is not replaced by a scalar trace.
This last fact is used only for the local-polynomial comparison below;
the internal endpoint filters remain present.

\section{Two independent rooted tadpole sections}
\label{sec:v63-tadpoles}

For each marked coefficient through metric degree two, use all its
independent external momenta, including the momentum of an independent
root source. Taylor subtraction is performed before identifying these
momenta with those of the other factor. Write
\begin{equation}
 \begin{aligned}
 P_{g,a,\sigma}&=\mathsf T_2M_{a,\sigma},&
 U_{g,a,\sigma}&=(1-\mathsf T_2)M_{a,\sigma},\\
 P_{H,a,\sigma}&=\mathsf T_4H_{a,\sigma},&
 U_{H,a,\sigma}&=(1-\mathsf T_4)H_{a,\sigma}.
 \end{aligned}
 \label{eq:v63-one-loop-sections}
\end{equation}
The root-source notation merely keeps the local tensor components
independent until the final contraction; it adds no new integration
variable. Unmarked coefficients are constants and are removed completely
by these zero-finite-addition sections.

The respective superficial degrees are $d_g=2$ and $d_H=4$: a metric
covariance contributes minus two derivatives, whereas the scalar
bilinear differentiates its covariance by total order two. The finite
$a^2E_2$ correction is included in the latter tensor; it is not assigned
an independent divergent graph with its compensating terms missing.
Let $\eta_g=\chi^{-1}$ and $\eta_H=1$.

\begin{lemma}[Assigned rooted composite bounds]
\label[lemma]{lem:v63-symbols}
For $s\in\{g,H\}$, a rooted coefficient with at most two metric marks,
and any prescribed number of momentum derivatives, its shell at scale
$R\ge\lambda$ obeys
\begin{equation}
 |\partial_K^\alpha L_{s,a,R}(K)|
       \le C_\alpha\eta_sR^{d_s-|\alpha|}.
 \label{eq:v63-shell-base}
\end{equation}
On the common interior its finite-to-ordinary difference obeys
\begin{equation}
 |\partial_K^\alpha(L_{s,a,R}-L_{s,0,R})(K)|
       \le C_\alpha\eta_s a^2R^{d_s+2-|\alpha|}.
 \label{eq:v63-shell-difference}
\end{equation}
Here the normalized one-loop sum is already included. The corresponding
integrands have four fewer powers of $R$. The finite upper pieces satisfy
\eqref{eq:v63-shell-base} in their periodic scaled coordinates.
\end{lemma}
\begin{proof}
At fixed source degree every background derivative of an inverse adds
one order-minus-two covariance with one actual order-two vertex, so it
preserves the degree. Endpoint filtering also preserves degree. The
scalar bilinear has effective order two on the finite zone, including
$a^2E_2$; on its support $a^2R^4\le CR^2$. All internal momenta in a
nonzero rooted chain differ by sums of at most two soft marks, so their
scales are comparable within a fixed support constant. The normalized
one-loop count is bounded by $CR^4$. These statements, together with the
inherited native inverse and vertex estimates, prove
\eqref{eq:v63-shell-base}.

The bare periodic scalar tensor has second-order consistency, the
continued corrected scalar tensor has fourth-order consistency, and
the metric vertices have the conservative third-order consistency of
V61. On $aR<c$, all are bounded by the common relative error
$Ca^2R^2$. Endpoint profiles are one there. The three channels may
therefore be compared separately using this conservative estimate;
no cancellation between their norms is asserted. A product difference
inserts exactly one such discrepancy, giving
\eqref{eq:v63-shell-difference}. The grouped tensor includes the
compensating scalar mesh terms. Smooth periodic completion gives the
upper estimate. Signed metric contractions are assembled before taking
matrix norms.
\end{proof}

\begin{theorem}[Renormalized rooted composites]
\label[theorem]{thm:v63-one-loop}
On the domain
\begin{equation}
 \lambda\ge32768\mu>0,\quad a\lambda\le2^{-24},\quad
 \lambda L_\nu\ge1,\quad |k_i|\le\mu,\quad m_H=0,
 \label{eq:v63-domain}
\end{equation}
every $U_{s,a,\sigma}$ of \eqref{eq:v63-one-loop-sections} has an
ordinary limit. For $j=|\alpha|\le d_s+1$,
\begin{align}
 |\partial_K^\alpha U_{s,0,L}(K)|
 &\le C\eta_s\lambda^{-1}|K|^{d_s+1-j},\notag\\
 |\partial_K^\alpha(U_{s,a,L}-U_{s,0,L})(K)|
 &\le C\eta_s a|K|^{d_s+1-j}.
 \label{eq:v63-one-loop-low}
\end{align}
For every fixed $j>d_s+1$, the corresponding bounds are
$C\eta_s\lambda^{d_s-j}$ and $C\eta_sa\lambda^{d_s+1-j}$.
The ordinary volume error has the same first bound multiplied by
$(\lambda L_{\min})^{-b}$, for any fixed $b>4$; the higher-derivative
version has that factor as well.
\end{theorem}
\begin{proof}
After $d_s+1$ external derivatives,
\eqref{eq:v63-shell-base} is $C\eta_sR^{-1}$ and
\eqref{eq:v63-shell-difference} is $C\eta_sa^2R$.
Dyadic summation gives
\[
 \sum_{R\ge\lambda}R^{-1}\le C\lambda^{-1},\qquad
 a^2\sum_{R\le c/a}R\le Ca,\qquad
 \sum_{R\ge c/a}R^{-1}\le Ca.
\]
There are a bounded number of upper periodic pieces at scale $a^{-1}$;
each contributes $C\eta_sa$. Taylor's integral remainder proves
\eqref{eq:v63-one-loop-low}. Extra derivatives lower each scale power.
The borderline sum $a^2\log(1/(a\lambda))$ is at most $Ca/\lambda$;
all other sums are geometric. This proves the stated higher-derivative
bounds. Poisson summation in the single four-dimensional loop, after
these derivatives have made the integrand integrable, gives the volume
bound on each scale and then its summation. At no stage is an isolated
divergent infinite local coefficient assigned a value.
\end{proof}

For later multiplication define a finite scaled seminorm
\[
 \|F\|_{d,\lambda;m}
   =\sum_{|\alpha|\le m}\lambda^{|\alpha|-d}
      \sup_{|K|\le c_0\lambda}|\partial_K^\alpha F(K)|,
\]
where $c_0>0$ is fixed below the support margin. The theorem implies
$\|U_{s,0}\|_{d_s,\lambda;m}\le C\eta_s$ and
$\|U_{s,a}-U_{s,0}\|_{d_s,\lambda;m}\le C\eta_sa\lambda$.
These norms are on fixed multilinear source panels, not an extensive
supremum norm on an action.

The stored local coefficient of degree $d\le d_s$ satisfies
\begin{equation}
 |P_{s,a}^{[d]}|\le C\eta_s
 \begin{cases}
 a^{d-d_s},&d<d_s,\\
 1+\log(1/(a\lambda)),&d=d_s.
 \end{cases}
 \label{eq:v63-local-budget}
\end{equation}
Thus the shared-root local product can contain powers through $a^{-6}$
and products of logarithms. It is not uniformly small.

\section{The shared-root subtraction and its local double term}
\label{sec:v63-subtraction}

Suppress tensor and channel indices temporarily. Let $f_g,f_H$ be
\emph{supplied} local tensor polynomials of degrees at most two and four,
respectively. They specify a comparison normalization for the two
composites. Set
\[
 A_g=P_g-f_g,\qquad A_H=P_H-f_H,\qquad
 M^\sharp=M-A_g=U_g+f_g,\quad
 H^\sharp=H-A_H=U_H+f_H.
\]
All these are jets of one rooted functional through the retained metric
degree, not unrelated values for individual observables.

\begin{theorem}[Product preparation and shared local assignment]
\label[theorem]{thm:v63-preparation}
For the complete tensor contraction in \eqref{eq:v63-factorization},
put
\begin{equation}
 \mathcal C^{\sharp,R}_{a,L}
  =(1-\mathsf T_6)\frac{n_H}{4}
       \int_a\sum_\sigma\varepsilon_\sigma B_\sigma
                              M_\sigma^\sharp H_\sigma^\sharp.
 \label{eq:v63-ren-contact}
\end{equation}
On the declared source window this is exactly
\begin{equation}
 \boxed{(1-\mathsf T_6)\frac{n_H}{4}
       \int_a\sum_\sigma\varepsilon_\sigma B_\sigma
       (M_\sigma H_\sigma-A_{g,\sigma}H_\sigma
                                      -M_\sigma A_{H,\sigma}).}
 \label{eq:v63-three-term}
\end{equation}
Before the overall assignment, the product preparation has the additional
local term
\begin{equation}
 \mathcal D_a=\frac{n_H}{4}
       \int_a\sum_\sigma\varepsilon_\sigma B_\sigma
                                  A_{g,\sigma}A_{H,\sigma}.
 \label{eq:v63-double-local}
\end{equation}
It has total derivative degree at most six and must be retained in the
local bookkeeping, although $(1-\mathsf T_6)\mathcal D_a=0$.
\end{theorem}
\begin{proof}
Expand $(M-A_g)(H-A_H)$ in the given tensor order. This gives the three
terms of \eqref{eq:v63-three-term} plus $A_gA_H$. Each local composite
coefficient is a differential polynomial in the retained metric, of
degree at most two or four. The root tensor is algebraic in that metric.
On soft marks its external filters are exactly one. Their product
therefore has degree at most six at each of the zero, one and two mark
panels, and is fixed by $\mathsf T_6$. Applying $1-\mathsf T_6$ proves
the equality. This is a coefficient identity at finite cutoff; it needs
neither a convention about forests that share vertices nor separate
convergence of the divergent local factors.
\end{proof}

The graph in \eqref{eq:v63-bare-contact} is connected. Its algebraic
factorization is not permission to remove it as a disconnected vacuum
piece. Conversely, \eqref{eq:v63-double-local} is not an additional
bare Feynman graph. It is the local cross term that makes two successive
composite replacements equal to their simultaneous product. If the
prepared three-term expression is used, its overall counterterm is
$-\mathsf T_6$ of that three-term expression. If the product preparation
is used, its different overall local polynomial includes
\eqref{eq:v63-double-local}. Both yield the same complementary source
coefficient. This fixes the shared-vertex ambiguity without claiming
that the numerical local polynomial is zero.

For $f_g=f_H=0$, the unmarked renormalized composites vanish. In the
coefficient with two labelled metric marks only the two split-mark
terms remain:
\begin{equation}
 \boxed{\frac{n_H}{4}\sum_\sigma\varepsilon_\sigma B_\sigma(0)
 \left[U_{g,\sigma}^{[h_1]}U_{H,\sigma}^{[h_2]}
       +U_{g,\sigma}^{[h_2]}U_{H,\sigma}^{[h_1]}\right].}
 \label{eq:v63-split-marks}
\end{equation}
Each product vanishes through degree seven since its factors start at
orders three and five. The overall sixth jet therefore vanishes in
this particular zero-finite-addition comparison. This does not set the
raw local graph, its one-loop subtractions, or a physical finite
normalization to zero. Terms with both marks on one loop return when
its partner has a nonzero prescribed finite constant.

\section{All relative scales of the contact component converge}
\label{sec:v63-limit}

Assume that the supplied finite composite tensors have a common ordinary
limit and a fixed scaled coefficient budget
\begin{equation}
 \sum_{b=0}^{2}\sum_{d=0}^{2}
       \chi\lambda^{d-2}\|f_g^{[b,d]}\|
 +\sum_{b=0}^{2}\sum_{d=0}^{4}
       \lambda^{d-4}\|f_H^{[b,d]}\|\le M_f.
 \label{eq:v63-finite-budget}
\end{equation}
For the first assertion below the same polynomial tensors are used in
both regulators. Denote the similarly normalized discrepancy between
different supplied arrays by $\varepsilon_f(a,L)$.

\begin{theorem}[Complete contact comparison]
\label[theorem]{thm:v63-limit}
On \eqref{eq:v63-domain}, the two-metric coefficient of
\eqref{eq:v63-ren-contact} has an ordinary thermodynamic limit. For
$j=|\alpha|\le7$,
\begin{equation}
 \boxed{\left|\partial_k^\alpha
 (\mathcal C^{\sharp,R}_{a,L}-\mathcal C^{\sharp,R}_{0,L})
 (h_1,h_2;k)\right|
 \le C_\alpha\frac{n_H}{\chi}(1+M_f)^2
           a|k|^{7-j}\|h_1\|\|h_2\|.}
 \label{eq:v63-main-bound}
\end{equation}
If the supplied finite tensors differ, add
$C(n_H/\chi)(1+M_f)\lambda^{-1}\varepsilon_f|k|^{7-j}
\|h_1\|\|h_2\|$. The ordinary volume discrepancy is bounded by
\begin{equation}
 C_{\alpha,b}\frac{n_H}{\chi}(1+M_f)^2\lambda^{-1}
 (\lambda L_{\min})^{-b}|k|^{7-j}\|h_1\|\|h_2\|,
 \qquad b>8.
 \label{eq:v63-volume}
\end{equation}
All relative metric/scalar loop scales and all completed upper modes are
included. No unproved restriction to comparable scales is made.
\end{theorem}
\begin{proof}
At any fixed marked panel, \cref{thm:v63-factorization} is a finite sum
of contractions of one degree-two composite and one degree-four
composite with a dimensionless algebraic root tensor. Product rules in
the scaled seminorms give
\[
 \|B M^\sharp H^\sharp\|_{6,\lambda;m}
           \le C\chi^{-1}(1+M_f)^2,
\]
\[
 \|B(M_a^\sharp H_a^\sharp-M_0^\sharp H_0^\sharp)
                         \|_{6,\lambda;m}
           \le C\chi^{-1}(1+M_f)^2a\lambda.
\]
Use $M_aH_a-M_0H_0=(M_a-M_0)H_a+M_0(H_a-H_0)$ after both factors
have been subtracted. The local divergent $P_s$ are absent from these
bounds because the signed preparation was formed first. The finitely
many marked tensors $B$ are bounded on the common coframe chart and
are identical between regulators on the soft external window. Each
of the three scalar-action channels has the same ordinary rooted
kernels and root tensor; its signs sum to $1+1-1=1$. This identifies
the comparison with the ordinary native contact, rather than three
independent continuum vertices. If the supplied finite additions differ
between channels, their ordinary values combine as the signed sums
$\sum_\sigma\varepsilon_\sigma f_{s,\sigma}$. The difference between
summing the products of finite additions and multiplying these two
sums is local of degree at most six and is removed by the same overall
assignment; no new ordinary vertex is needed.

The seventh derivative of the difference is thus bounded by
$C\chi^{-1}a$: the seminorm has supplied $\lambda^{6-7}$ times
$a\lambda$. Taylor's integral remainder proves
\eqref{eq:v63-main-bound}, including its differentiated versions.
Higher fixed derivatives follow from the same product estimate.
A differing finite tensor gives the displayed $\varepsilon_f$ term.
Apply the volume statement of \cref{thm:v63-one-loop} to one factor
at a time to obtain \eqref{eq:v63-volume}.

For explicit scale accounting, write the two remainders as sums over
\emph{independent} primitive shell chains. After their respective
third and fifth Taylor remainders their bounds are proportional to
$R^{-1}$ and $r^{-1}$. Both series are absolutely summable. Their
product contains every $(R,r)$, including $R\gg r$, $r\gg R$, and
comparable scales. The finite difference is estimated by one error
series and one convergent series; no marginal count of the number of
lower scales is required. This proves why the contact estimate has
no extra ultraviolet logarithm. Prescribed polynomial finite additions
multiply only one convergent series at a time, and the seventh overall
derivative gives the same result. This argument includes all labelled
inverse-return factors and is not a diagonal-shell approximation.
\end{proof}

\begin{corollary}[Rooted topology and a sharper comparison convention]
\label[corollary]{cor:v63-rooted}
After a fixed smooth source window,
\begin{equation}
 \boxed{a^4\sum_x(1+\mu d_{\rm per}(x,0))^b
        \|\mathcal K^\Delta_{a,L}(x)\|
 \le C_b\frac{n_H}{\chi}(1+M_f)^2a\mu^7.}
 \label{eq:v63-rooted}
\end{equation}
For zero finite additions, \eqref{eq:v63-split-marks} instead gives
$C_b(n_H/\chi)a\lambda^{-1}\mu^8$. One metric source may be in
$L^1$ and the other in $L^\infty$, without a total-volume factor.
\end{corollary}
\begin{proof}
Window derivatives and the higher fixed derivative bounds in
\cref{thm:v63-limit} give an integrable Fourier majorant with any fixed
moment in the four relative coordinates. Periodization gives
\eqref{eq:v63-rooted}. For zero finite additions use
$|\Delta U_g|\le C\chi^{-1}a|k|^3$ and
$|U_H|\le C\lambda^{-1}|k|^5$, and the opposite product, directly in
\eqref{eq:v63-split-marks}. The additional $a^2$ product is absorbed
since $a\lambda$ is bounded. The same differentiated argument gives
the sharper windowed estimate.
\end{proof}

\section{What must match before the two core graphs are combined}
\label{sec:v63-matching}

The theorem proves the complete contact limit in a declared composite
comparison scheme. It does not identify that scheme with an unknown
solution of a local BV equation. Its compatibility with a common
one-loop action can be stated exactly at finite cutoff.

Introduce independent rooted quadratic sources for the metric pair and
for the scalar kinetic pair at the shared vertex. Their local one-loop
jets are $P_g$ and $P_H$ in \eqref{eq:v63-one-loop-sections}. The
corresponding insertions in the contact coefficient are the actual
$-BA_gH$ and $-BMA_H$ terms of \eqref{eq:v63-three-term}. Each retains
the parent sign, the $n_H/4$ weight and all metric derivatives. These
are subgraph components of a common one-loop normalization; they are
not generally the entire one-loop self-energy, which also has the
second graph of \cref{eq:v61-core}.

\begin{proposition}[Finite normalized comparison and nonlocal dependence]
\label[proposition]{prop:v63-scheme}
Changing the supplied finite tensors by $\delta f_g,\delta f_H$
changes the contact remainder by exactly
\begin{equation}
 (1-\mathsf T_6)\frac{n_H}{4}\int_a\sum_\sigma
 \varepsilon_\sigma B_\sigma
 \left(\delta f_g H^\sharp+M^\sharp\delta f_H
                              +\delta f_g\delta f_H\right).
 \label{eq:v63-scheme-change}
\end{equation}
The last term is local and is removed by the overall assignment. The
first two need not be local and cannot in general be absorbed into one
new overall gravitational polynomial. They have the bounds of
\cref{thm:v63-limit} with the corresponding finite coefficient budgets.
\end{proposition}
\begin{proof}
Expand the product in \eqref{eq:v63-ren-contact}; tensor order and
source allocation are unchanged. Degrees two and four make the last
product a degree-six local tensor. The other two contain renormalized
one-loop functions, so locality is not automatic. The earlier seminorm
product proof supplies their estimates.
\end{proof}

This formula is the explicit condition needed before combining this
contact with V62's three-propagator forest. A proposed common $S_1$
must supply the same local rooted tensor restrictions, or supply their
finite differences in \eqref{eq:v63-scheme-change} and the corresponding
insertions in the other graph. No source tensor is averaged to make
this condition hold. The profiles on internal legs are those of the
native filtered composite realization. No claim of strict finite range
for these smooth profile kernels is made. Replacing them by another
realization inside a hard contraction would require its own matched
estimate; V54's soft-jet theorem alone would not suffice.

\begin{corollary}[Both primitive action graphs, with compatible supplied data]
\label[corollary]{cor:v63-combine}
Suppose the finite one-loop data supplied to V62 and to
\eqref{eq:v63-ren-contact} are the restrictions of one common
normalization with the stated bounded budgets and compatible
realizations. Their sum, through two ordinary metric marks, has a
cutoff comparison bounded by
\begin{equation}
 C\frac{n_H}{\chi}
 a\left[1+\log\frac1{a\lambda}\right]|k|^{7-j}
 \|h_1\|\|h_2\|,\qquad j\le7,
 \label{eq:v63-combined}
\end{equation}
with the separate volume and rooted estimates obtained by addition.
This combines the two primitive scalar-containing action cores only.
It does not add a missing common normalization or a master identity.
\end{corollary}
\begin{proof}
Apply \cref{thm:v63-limit} and the inherited V62 forest estimate. The
scalar vacuum constant is discarded only after geometric
differentiation and the stipulated local assignments, as in V61.
Additional one-loop counterterm insertions outside the two specified
subgraph ledgers remain separate terms if the common action contains
them. Existence of a source-complete $S_1$ with all required BV and
physical normalization conditions is not a premise established here.
\end{proof}

\section{Ledgers, scope, and the next common-action calculation}
\label{sec:v63-ledgers}

\subsection*{Proved}
The native shared-root contact factorizes exactly into its full signed
coframe tensor, rooted metric covariance and rooted scalar kinetic
tensor. Both tadpole sections have complete source and cutoff bounds.
Their simultaneous product gives an explicit shared-vertex subtraction,
including the stored local double term. The resulting two-metric
contact has an $O(a)$ full-cutoff comparison, a separate volume error
and rooted source bounds over all relative internal scales. Supplied
finite composite changes have an exact, generally nonlocal insertion
formula.

\subsection*{Inherited}
V27 supplies the single completed scalar action and its derivative
symbols. V14/V24 supply the signed metric reference and needed native
tree jets. V61 supplies the contact Wick factor, connected-graph
completeness and bridge exclusion. V62 supplies the other core graph's
complete overlapping forest estimate. These inputs are consumed under
the same massless, two-metric restrictions; they are not promoted to
all-order ESM statements.

\subsection*{Literature input}
Composite normal products and forest subtraction are comparison
methodology \cite{V63ZimmermannNormal,V63Blaschke}. No external theorem
is used to supply a missing native symbol bound, choose a finite
counterterm, or enforce a Ward identity after subtraction.

\subsection*{Conditional}
Combining both core limits in a specified physical scheme requires the
actual common one-loop action/source arrays and their graphwise
restrictions, as stated in \cref{cor:v63-combine}. The particular
zero-finite-addition composite scheme is a comparison normalization,
not a silent replacement of the fixed gravitational or scalar
prescription. No numerical native integral is asserted evaluated.

\subsection*{Open}
The next load-bearing task is the common $S_1$ insertion ledger and its
coupled two-loop source/reference consistency equation, not another
estimate of the now controlled contact topology. The complete
coefficient still contains $(S_0,S_2)+(S_1,S_1)/2-\Delta S_1$ and the
actual finite tree defect. Internal gauge/chiral sectors, the populated
critical primitive, higher-loop induction and the filtered interacting
ESM class remain open. A graphwise continuum limit does not solve those
interfaces.

\begin{center}\small
\begin{tabular}{>{\raggedright\arraybackslash}p{.23\textwidth}>{\raggedright\arraybackslash}p{.69\textwidth}}
\toprule
Variable & Scope in V63\\\midrule
Loop/source order & Two-loop shared-root contact through two ordinary metric
marks; no retained scalar, ghost or antifield.\\
Mass/couplings & Four massless real scalars; fixed $\chi,\lambda>0$;
internal gauge, Yukawa and scalar self-couplings zero.\\
Cutoff/volume & All completed upper modes and all relative scales on
\eqref{eq:v63-domain}; no total-volume factor in rooted bounds.\\
Normalization & Divergent jets retained; finite tensors supplied under
\eqref{eq:v63-finite-budget}; common-action matching remains separate.\\
Uniformity & Fixed source/derivative order, profiles, inverse margins, coframe
chart and shell constants; no arbitrary-order or interacting-step bound.\\
Verification & Exact finite algebra and preservation audits; no numerical
native integral or proof-assistant certification.\\
\bottomrule
\end{tabular}
\end{center}

\begin{thebibliography}{99}
\bibitem{V63ZimmermannNormal} W.~Zimmermann,
Normal products and the short distance expansion in the perturbation
 theory of renormalizable interactions,
\emph{Annals of Physics} \textbf{77} (1973), 570--601.
\bibitem{V63Blaschke} D.~N.~Blaschke, F.~Gieres, F.~Heindl, M.~Schweda
and M.~Wohlgenannt, BPHZ renormalization and its application to
non-commutative field theory, \emph{European Physical Journal C}
\textbf{73} (2013), 2566; arXiv:1307.4650.
\end{thebibliography}

\section*{Append-only continuation marker: V63}
\addcontentsline{toc}{section}{Append-only continuation marker: V63}
\label{sec:v63-continuation-marker}
\textbf{End of V63.} Preserve Sections 1--489. The shared-root contact
has a complete graphwise composite subtraction and full-cutoff source
limit. Common one-loop normalization and the full two-loop
master/reference coefficient remain the next interfaces. No claim of
finite-parameter or nonperturbative ultraviolet gravity is made.


\clearpage
% ============================================================================
% V64 -- common one-loop action, its two-loop insertion, and reference response
% ============================================================================
\section{V64: from independent subgraph data to one counterterm action}
\label{sec:v64-start}

V62 and V63 supply complementary estimates for the two primitive
scalar-containing two-loop action graphs. Their common-normalization
condition is not an additional graph estimate: the local tensors inserted
on different cut lines must be derivatives of the same order-$\hbar$
functional. This installment computes that insertion, including its
cutting multiplicities and geometric derivatives, and separates its
ordinary action-integrability test from the still required BV equation.
A finite scalar kinetic normalization provides an explicitly evaluated
native example: at fixed reference it produces a nonzero lower-reference
metric coefficient at the next loop order. It is not merely a vacuum
constant.

No action, covariance, contour, measure or chosen finite normalization is
changed. In particular the fixed gravity prescription, V34's active
subcritical scalar normalization and V38's complete logarithmic packet
are retained. The parameter introduced in the kinetic example describes
a proposed comparison, not a newly fitted native coefficient. The
unknown critical counterterm is not populated by the constructions below.

\paragraph{Scope and conventions.}
All finite expressions use the same physical cell pairing and normalized
traces as V61--V63. Covariances and Hessian operators below have their
factors of $\hbar$ removed. Antifields are retained sources, not Gaussian
integration variables. The integrated hard fields are denoted by $\Phi$;
their supertrace includes the actual metric, scalar, ghost and nonminimal
blocks. On a scalar flavour-diagonal block the factor is $n_H=4$.
The signed metric covariance is never replaced by a positive tensor.
The local coordinate density is inserted once as $-\hbar\log\rho$
when the flat Darboux Laplacian is used.

The quantitative example uses four massless real scalars. Internal gauge,
Yukawa and scalar self-interaction couplings are zero. The inherited window is
\begin{equation}
 \lambda\ge32768\mu>0,\qquad a\lambda\le2^{-24},\qquad
 \lambda L_\nu\ge1,\qquad |k_i|\le\mu.
 \label{eq:v64-domain}
\end{equation}
The purely algebraic coefficient identities also allow the finite mass
polynomials and spectator sources already present in a supplied action.
A trace written before division by physical volume is an extensive
functional; a bound on a rooted coefficient below is not a bound on that
extensive trace.

The standard relation between counterterm vertices and forest preparation
is used as comparison methodology \cite{V64BPHZ}. Neither that relation
nor the finite BV product rule supplies the actual missing finite
normalization. The distinction between a fixed reference and a simultaneous
change of reference is essential \cite{V64IIM}.

\section{The common one-loop functional has one two-loop insertion}
\label{sec:v64-insertion}

Write the bare perturbative input as
\[
 S_a=S_{a,0}+\hbar C_{a,1}+\hbar^2C_{a,2}+O(\hbar^3).
\]
Here $C_{a,1}$ is the entire supplied local action-and-source functional,
including any once-counted density contribution. It is not just one
self-energy block. On a fixed soft background let $\mathbb G_a$ be the
actual hard tree covariance resolvent, and put
\[
 H_{a,1}=D_\Phi^2C_{a,1}.
\]
All derivatives in this formula are coefficient derivatives on the
finite carrier. For odd hard fields the ordered Berezin Hessian is used.

\begin{theorem}[Counterterm insertion on the native hard carrier]
\label[theorem]{thm:v64-insertion}
Through two ordinary metric marks and zero retained ordinary scalar,
ghost and antifield, the additional order-$\hbar^2$ effective action
caused by $C_{a,1}$ is
\begin{equation}
 \boxed{\mathcal I_a[C_{a,1}]
       =\frac12\operatorname{STr}(\mathbb G_aH_{a,1}).}
 \label{eq:v64-insertion}
\end{equation}
The same formula holds coefficientwise on a prescribed finite
source panel whenever the hard-linear bridge exclusion holds for every
external argument of that panel. Its source factors are not integrated
out of the Hessian. The complete order-two action also contains
$C_{a,2}$ and the bare two-loop graphs; \eqref{eq:v64-insertion} is the
counterterm contribution, not their replacement.
\end{theorem}
\begin{proof}
Expand the same finite connected negative logarithm in explicit powers
of $\hbar$. A connected graph of topological loop number $\ell$ with
$r_1$ order-one counterterm vertices and $r_2$ order-two vertices has
order $\hbar^{\ell+r_1+2r_2}$. At order two the possibilities involving
a counterterm are: one $C_1$ with one loop; two $C_1$ vertices and no
loop; or one $C_2$ with no loop. The last gives the retained local
$C_2$. Every nontrivial tree in the other cases contains a bridge.
Its momentum is the sum of external momenta on one side, hence is
below the hard support in the declared panel. These terms vanish by
\cref{lem:v61-bridge}, not because they are disconnected.

In the one-loop case, a hard-linear vertex likewise forces a bridge.
For a connected unicyclic graph with every hard valence at least two,
$\sum_v(\deg_hv-2)=0$ forces every vertex to have hard valence two.
Exactly one is the Hessian of $C_1$. The other quadratic vertices are
precisely the tree-background returns resummed in $\mathbb G_a$.
Cyclic Wick counting gives $1/2$ for a real even pair, and ordered
Berezin contraction gives the supertrace signs. A ghost--antighost
block therefore gives $-\Tr(G_cM_1)$, not an additional positive half
trace. The argument retains every source derivative. It never inverts
a retained zero mode and does not replace the native action by an
independent continuum vertex.
\end{proof}

For example, on the zero-ghost/zero-scalar background write
\[
 C_1=g_1(q)+\frac12\sum_A\langle\phi_A,W_1(q)\phi_A\rangle
       +\langle\bar c,M_1(q)c\rangle+C_{1,\mathrm{other}}.
\]
Then \eqref{eq:v64-insertion} contains
\begin{equation}
 \frac12\Tr(G_gD_q^2g_1)
 +\frac{n_H}{2}\Tr(G_HW_1)-\Tr(G_cM_1)
 +\frac12\operatorname{STr}(\mathbb G_aH_{1,\mathrm{other}}).
 \label{eq:v64-blocks}
\end{equation}
The last term retains the nonminimal or other actually supplied hard
blocks. A minimal-antifield counterterm can vanish in this action
extraction and still be essential to the master coefficient below.
No physical zero is assigned to an unspecified block.

Let $K_h$ and $K_{hk}$ be derivatives of the native tree hard Hessian,
so $G_h=-GK_hG$. For even metric marks, the whole first and second
source derivatives are
\begin{align}
 D_h\mathcal I[C_1]
 &=\tfrac12\operatorname{STr}
       (G H_{1,h}-GK_hG H_1),\label{eq:v64-first}\\
 D_kD_h\mathcal I[C_1]
 &=\tfrac12\operatorname{STr}\bigl[
 G H_{1,hk}-GK_hG H_{1,k}-GK_kG H_{1,h}\notag\\
 &\hspace{10mm}+GK_hGK_kG H_1+GK_kGK_hG H_1
                         -GK_{hk}G H_1\bigr].
 \label{eq:v64-second}
\end{align}
\begin{proposition}[One insertion determines all six second-mark terms]
\label[proposition]{prop:v64-marks}
Equations \eqref{eq:v64-first}--\eqref{eq:v64-second} are the full
geometric derivatives of \eqref{eq:v64-insertion}, with no commuting
of matrix factors. At this action/source order, $g_1$ is consumed
through four metric derivatives, $W_1$ through two, and the ghost and
other supplied blocks through their corresponding Hessian marks.
\end{proposition}
\begin{proof}
Differentiate $G H_1$ twice and use
$G_{hk}=GK_hGK_kG+GK_kGK_hG-GK_{hk}G$.
These are six labelled product-rule terms. The two hard derivatives
already present in $D_q^2g_1$ explain its fourth metric jet. The
result is a finite coefficient identity also when the matrices have
nilpotent source coefficients; it is not a norm estimate obtained by
dropping their signs.
\end{proof}

\section{Opening a primitive graph fixes the counterterm multiplicities}
\label{sec:v64-cuts}

This section evaluates the compatibility factors connecting V62 and V63.
Use their $Q_I,Q_{IJ}$, $G_g$ and $C=G_H$, and put
\begin{align}
 \mathcal C&=\frac{n_H}{4}(G_g)_{IJ}\Tr(CQ_{IJ}),\notag\\
 \mathcal S&=-\frac{n_H}{4}(G_g)_{IJ}\Tr(CQ_ICQ_J).
 \label{eq:v64-cores}
\end{align}
Repeated hard metric indices are summed. These are the two already
computed bare cores; they are not evaluated anew here.
Their open one-loop two-point tensors, with the external scalar flavour
fixed in the scalar tensors, are
\begin{align}
 (\Sigma_H^{\mathrm c})_{IJ}&=\frac{n_H}{2}\Tr(CQ_{IJ}),&
 (\Sigma_H^{\mathrm s})_{IJ}&=-\frac{n_H}{2}\Tr(CQ_ICQ_J),\notag\\
 \Sigma_g^{\mathrm c}&=\frac12(G_g)_{IJ}Q_{IJ},&
 \Sigma_g^{\mathrm s}&=-(G_g)_{IJ}Q_ICQ_J.
 \label{eq:v64-open-tensors}
\end{align}
These formulas include the actual coframe contact and all endpoint
profiles through the vertices and covariances.

\begin{theorem}[Common-action cut multiplicities]
\label[theorem]{thm:v64-cuts}
The ordinary one-counterterm contraction of
\eqref{eq:v64-open-tensors} satisfies the four exact identities
\begin{equation}
 \boxed{\begin{array}{ll}
 \frac12(G_g)_{IJ}(\Sigma_H^{\mathrm c})_{IJ}=\mathcal C,
 &\frac{n_H}{2}\Tr(C\Sigma_g^{\mathrm c})=\mathcal C,\\[1mm]
 \frac12(G_g)_{IJ}(\Sigma_H^{\mathrm s})_{IJ}=\mathcal S,
 &\frac{n_H}{2}\Tr(C\Sigma_g^{\mathrm s})=2\mathcal S.
 \end{array}}
 \label{eq:v64-cut-identities}
\end{equation}
Consequently the two proper contact subtractions and the three proper
parallel-edge subtractions are generated by the same unit coefficient
of an action Hessian insertion. The factor two in the last entry
counts the two scalar edges. It is not a second independent scalar
normalization.
\end{theorem}
\begin{proof}
Insert \eqref{eq:v64-open-tensors} and contract, preserving trace order.
For the last identity there are two indistinguishable scalar cut edges
in the original core; equivalently its displayed expression is twice
$\mathcal S$. The two-scalar-vertex expansion with scalar legs paired
within their own vertex gives an additional connected bridge term;
it is absent here by the support lemma, not absorbed into a symmetry
factor. The scalar contact has one scalar cut edge, so its coefficient
is different. This verifies the graph multiplicities directly from
the finite native quadratic and cubic vertices.
\end{proof}

A local proper-subgraph assignment can be applied to each open tensor
before closing its complementary edge. Since it is linear, the same
cut multiplicities persist, with the boundary hard momentum and all
geometric marks still arguments of that assignment. The two scalar
cuts may carry different labelled external transfers; they are summed
with their own routing before being identified by a symmetry. One must
not Taylor-subtract only on the soft physical source window.

For the contact, the two insertions obtained this way are exactly the
$-BA_gH$ and $-BMA_H$ terms of V63 after the specified rooted tensors
are contracted with the native scalar vertex. The degree-six local
product $BA_gA_H$ remains the overall-assignment conversion of
\cref{eq:v63-three-term}; it is not generated as an additional bare
order-two graph. For the parallel-edge core, the same prescription
gives the three single-cycle terms of V62. This identifies their
common action-level multiplicities, but does not choose their finite
local tensor values.

An implementation that inserts $\tfrac12\Sigma_g^{\mathrm s}$ merely
to reconstruct one bare vacuum graph would miss one scalar subtraction
channel. Conversely counting two independent copies of the complete
scalar counterterm would double its contact contribution as well.
The shared action insertion \eqref{eq:v64-insertion} makes neither
mistake. The identities remain valid after the complete labelled
geometric differentiation of \cref{prop:v64-marks}.

\section{An action-integrability test for the supplied local Hessians}
\label{sec:v64-hessian}

Graphwise local arrays are often supplied with two distinguished hard
legs and several soft legs. A common action cannot distinguish which
two of its identical-type legs will later be contracted. The finite
condition is therefore stronger than symmetry in the two hard legs.
This section concerns the even metric/scalar action slots. It does not
replace the graded BV or physical normalization conditions.

Let $X^a$ be the actual finite even coordinates, or a fixed translated
local coefficient bank including field type and momentum labels. Write
\[
 H_{ab}(X)=\sum_{n=0}^{m}\frac1{n!}
                  h_{ab\mid i_1\ldots i_n}X^{i_1}\cdots X^{i_n},
 \qquad H_{ab}=H_{ba}.
\]
The last $n$ indices are symmetric. A coefficient label includes its
field type and derivative/momentum assignment; symmetrization moves
these labels together and never converts one species into another.
Constant and linear action coefficients are supplied independently.

\begin{theorem}[Local Hessian section with an explicit failed-input residual]
\label[theorem]{thm:v64-hessian}
There is a polynomial action $F$, through degree $m+2$, with $D^2F=H$
if and only if every $h_{ab\mid I}$ is symmetric in all $n+2$ labels.
Equivalently,
\begin{equation}
 \partial_cH_{ab}-\partial_aH_{cb}=0.
 \label{eq:v64-helmholtz}
\end{equation}
Let $\operatorname{Sym}_{n+2}$ denote full labelled averaging and set
\begin{equation}
 \boxed{F_H(X)=\int_0^1(1-t)X^aH_{ab}(tX)X^b\,dt.}
 \label{eq:v64-hessian-section}
\end{equation}
For arbitrary symmetric input, not necessarily integrable,
$D^2F_H=\Pi_HH$ is the full-symmetry projection. In particular
\begin{equation}
 \boxed{\Pi_H^2=\Pi_H,\qquad
 H=D^2F_H+H^{\rm rem},\qquad
 \mathfrak c(H^{\rm rem})=\mathfrak c(H),}
 \label{eq:v64-hessian-residual}
\end{equation}
where $\mathfrak c(H)_{cab}=\partial_cH_{ab}-\partial_aH_{cb}$.
For a supplied compatible input the unique action with given constant
and linear anchors is $F_H$ plus those anchors.
\end{theorem}
\begin{proof}
Third and higher derivatives of an action are fully symmetric, proving
necessity. Conversely, differentiating a fully symmetric $(n+2)$-tensor
in the action coefficient $X^{a_1}\cdots X^{a_{n+2}}/(n+2)!$ gives the
stated Hessian coefficient. For a general tensor, contraction with the
commuting $X$ variables averages over all labels. The radial integral
has coefficient $1/((n+1)(n+2))$, since
$\int_0^1(1-t)t^n\,dt=1/((n+1)(n+2))$. Thus its Hessian is exactly
full symmetrization. This proves the projection identity and its range.
The curl of a Hessian is zero, proving the residual formula. Two
polynomials with the same Hessian differ by an affine polynomial.
All arguments are coefficientwise, so fixed momentum polynomials,
mass powers and spectator labels are retained. On finite translated
kernels the permutations can reroot a leg but preserve finite range up
to a fixed arity-dependent factor. No spatial antiderivative or inverse
momentum occurs.
\end{proof}

In the ordered divided-tensor norm
$\|H^{[n]}\|=\sum_{a,b,I}|h_{ab\mid I}|/n!$, full averaging has norm
at most one. Hence
\begin{equation}
 \boxed{\|F_H^{[n+2]}\|\le
 \frac{\|H^{[n]}\|}{(n+1)(n+2)},\quad
 \|\Pi_H\|\le1,\quad\|I-\Pi_H\|\le2.}
 \label{eq:v64-hessian-norms}
\end{equation}
These bounds refer to this redundant labelled-tensor convention.
Changing to a compressed monomial or integration-by-parts basis has
its own finite conversion constants. The theorem does not assert these
same numerical norms in every such basis.

For example, specifying $H_{\phi\phi}=2q$ and declaring
$H_{q\phi}=H_{qq}=0$ fails \eqref{eq:v64-helmholtz}:
$\partial_qH_{\phi\phi}-\partial_\phi H_{q\phi}=2$.
If only the scalar slot was supplied and the mixed slot was unknown,
the action $q\phi^2$ would complete it with $H_{q\phi}=2\phi$.
An unspecified coefficient is not a zero constraint. On the fully
specified inconsistent array, \eqref{eq:v64-hessian-section} instead
returns $q\phi^2/3$ and a nonzero residual. It does not certify the
original data.

In particular, for a local scalar-bilinear action
$\tfrac12\phi^TW(q)\phi$, the $q\phi$ and $qq\phi^2$ vertices are
$W_q\phi$ and $\tfrac12\phi^TW_{qq}\phi$. They cannot be normalized
independently while keeping $W$ fixed. A common action Taylor assignment
made before choosing the hard legs has this symmetry automatically;
separate rooted assignments must be checked or compared with it.
Total divergences are zero as finite periodic functionals. Working at
the functional-kernel level before selecting density representatives
avoids silently changing an integration-by-parts convention.

This is ordinary action integrability, not BV exactness. Ghost source
slots require their actual graded identities, and an integrable action
can still violate the master equation. Applying $\Pi_H$ is a diagnostic
with a recorded residual, not permission to alter the fixed counterterm.

\section{A common kinetic normalization has a nonzero reference response}
\label{sec:v64-kinetic}

Consider the explicit comparison direction
\begin{equation}
 F_\zeta=\zeta S_{H,a}
 =\frac\zeta2\sum_A\langle\phi_A,Q_a(q)\phi_A\rangle,
 \label{eq:v64-kinetic-direction}
\end{equation}
with a supplied constant $\zeta$. It is inserted with coefficient
$\hbar$. On the ordinary massless branch it is the full covariant
scalar kinetic functional, so it is a permitted homogeneous action
direction if the scalar normalization is allowed to change. It is not
an allowed change while insisting on V34's fixed numerical kinetic
condition. Here it is only a comparison test; no value of $\zeta$ is
chosen. At finite cutoff its first-order BV defect is the actual
$\zeta s_aS_{H,a}$, including V60's nonzero finite action breaking.

Keep the reference fixed:
\[
 P=\widetilde P_a,\quad \nu=I-\Theta_\lambda,\quad C=\nu P^{-1},
 \quad V=Q_a(q)-P,\quad A=(I+CV)^{-1},\quad C_q=AC.
\]
In particular, $P$ is not replaced by $(1+\hbar\zeta)P$.
By \cref{thm:v64-insertion}, the induced two-loop action coefficient
on a purely metric background is
\[
 \mathcal J_{a,\zeta}(q)=\frac{\zeta n_H}{2}\Tr(C_qQ_a(q)).
\]

\begin{theorem}[Exact lower-reference normalization response]
\label[theorem]{thm:v64-kinetic}
After removing only the field-independent value at $q=0$,
\begin{equation}
 \boxed{\mathcal J_{a,\zeta}(q)-\mathcal J_{a,\zeta}(0)
 =\frac{\zeta n_H}{2}\Tr\bigl[(I-A(q))\Theta_\lambda\bigr].}
 \label{eq:v64-kinetic-reference}
\end{equation}
This is an exact finite identity, including the retained zero modes.
Its first and second metric jets are
\begin{align}
 D_h\mathcal J(0)&=\frac{\zeta n_H}{2}\Tr(CQ_h\Theta_\lambda),
 \label{eq:v64-kinetic-first}\\
 D_kD_h\mathcal J(0)&=\frac{\zeta n_H}{2}\Tr\bigl[
 CQ_{hk}\Theta_\lambda-CQ_hCQ_k\Theta_\lambda
                         -CQ_kCQ_h\Theta_\lambda\bigr].
 \label{eq:v64-kinetic-second}
\end{align}
Each source term retains the actual reference in its displayed matrix
position. The $Q_{hk}$ coframe contact cannot be deleted.
\end{theorem}
\begin{proof}
Since $CP=\nu$,
\[
 C_qQ=A(\nu+CV)=I-A\Theta_\lambda.
\]
At zero background this is $\nu$. Subtracting its trace gives
\eqref{eq:v64-kinetic-reference}; no separately divergent infinite
trace is assigned a value. Differentiate $A$ at the origin:
$A_h=-CQ_h$ and
$A_{hk}=CQ_hCQ_k+CQ_kCQ_h-CQ_{hk}$. This proves the displayed jets.
They are exactly the same result as differentiating both the scalar
covariance and the counterterm Hessian in \eqref{eq:v64-first}--
\eqref{eq:v64-second}. Keeping only the covariance derivatives would
miss the $CQ_h$ and $CQ_{hk}$ counterterm vertices.
\end{proof}

A constant conformal mark is $h=I_4$ in the inherited tangent chart.
At $z=0$, V58's actual vertex is
\[
 Q_I(p)=p_b(p)+s(ap)^2[p_n(p)-p_b(p)].
\]
On the lower-reference support, all completion and source profiles lie
on their unit plateaux; hence $Q_I(p)=p_n(p)=P(p)$ there. This uses
the native continued action, not an independently selected conformal
coupling.

\begin{proposition}[A populated nonzero finite-normalization coefficient]
\label[proposition]{prop:v64-positive}
On \eqref{eq:v64-domain}, the conformal derivative per physical volume
is exactly
\begin{equation}
 \boxed{\frac1{V_L}D_I\mathcal J_{a,\zeta}(0)
 =\frac{\zeta n_H}{2V_L}\sum_p
       \Theta_\lambda(p)[1-\Theta_\lambda(p)].}
 \label{eq:v64-positive-finite}
\end{equation}
Its ordinary thermodynamic value is
\begin{equation}
 \boxed{\frac{\zeta n_H\lambda^4}{2}\,m_\theta,
 \qquad
 m_\theta=\int\frac{d^4y}{(2\pi)^4}\theta(y)[1-\theta(y)]>0.}
 \label{eq:v64-positive-limit}
\end{equation}
Here $\theta$ is any admitted smooth profile with $0\le\theta\le1$,
unit plateau on the unit ball and support in the radius-two ball.
No radial symmetry is required. The uniform upper bound is
$ m_\theta\le15/(128\pi^2)$; there is no profile-independent positive
lower bound without an additional transition-width condition.
\end{proposition}
\begin{proof}
Use $C=\nu/P$ and $Q_I=P$ in \eqref{eq:v64-kinetic-first}.
The zero mode has $\Theta\nu=0$, so no omitted-zero-mode contact
appears. Rescaling $p=\lambda y$ gives the integral. Continuity from a
unit plateau to the vanishing exterior supplies an open set on which
$0<\theta<1$, proving strict positivity. The integrand is at most
$1/4$ and is supported in the annulus of volume $15\pi^2/2$, giving
the upper bound after division by $(2\pi)^4$. Smooth transitions can
be arbitrarily narrow, so these assumptions alone give no common
positive lower bound. At a finite torus the sum is nonnegative and is
strictly positive only if a momentum samples the transition; the
thermodynamic limit is strictly positive. Poisson summation gives
an error at most $C_b|\zeta|n_H\lambda^4
(\lambda L_{\min})^{-b}$ for every fixed $b>4$.
\end{proof}

For fixed physical periods the expression
\eqref{eq:v64-positive-finite} is independent of further sufficiently
fine mesh refinement, since all its supported momenta are already
inside the finite zone. It is a finite lower-reference coefficient,
not an ultraviolet divergence or anomaly. Subtracting
$\mathcal J(0)$ cannot remove it: metric differentiation must precede
any vacuum normalization. Changing the reference together with the
kinetic coefficient is a different comparison and must transport the
measure, determinant and source convention as well.

\begin{proposition}[Two-mark lower-reference comparison]
\label[proposition]{prop:v64-kinetic-bounds}
For the normalized coefficients in
\eqref{eq:v64-kinetic-first}--\eqref{eq:v64-kinetic-second}, the loop
momentum lies in $\lambda/2\le |p|\le3\lambda$ after a fixed routing.
For any prescribed external derivative order $j$ and $b>4$,
\begin{equation}
 |\partial_K^j(\mathcal J_{a,L}-\mathcal J_{0,\infty})|
 \le C_{j,b}|\zeta|n_H\lambda^{4-j}
 \left[(a\lambda)^2+(\lambda L_{\min})^{-b}\right]
 \prod_i\|h_i\|.
 \label{eq:v64-kinetic-bound}
\end{equation}
After fixed source windows the same estimate gives every fixed rooted
kernel moment, with no total-volume multiplier. The conservative
second-order mesh power is used; no new optimal rate is asserted.
\end{proposition}
\begin{proof}
Every displayed word contains a $\Theta_\lambda$ and at least one
covariance with the complementary factor $\nu$. The routed momenta
differ by at most two soft marks. Both factors are therefore supported
in the stated lower annulus. On that fixed annulus the scalar inverse
and the grouped native vertices through two marks have the inherited
second-order comparison used in \cref{lem:v63-symbols}; differentiating
preserves it with the corresponding powers of $\lambda$. A product
difference inserts one relative factor $O(a^2\lambda^2)$ into a
normalized four-dimensional integral of degree $4-j$. This proves the
mesh term. Smooth compact support and Poisson summation give the volume
term. Arbitrarily many fixed loop and source derivatives are bounded
on the annulus, so integration by parts in the windowed external
Fourier transform gives every prescribed rooted moment. This is a
single finite-reference loop, not another estimate of the two-loop
ultraviolet cores. The exact conformal identity above is stronger for
its particular coefficient.
\end{proof}

\section{The induced two-loop BV condition retains the first-order defect}
\label{sec:v64-BV}

It is useful to state exactly what a successful common-action matching
would still have to prove. Work in one fixed Darboux density/bracket
convention, as in V61, and write
\[
 \mathcal A_0=\tfrac12(S_0,S_0),\quad
 \mathcal A_1=(S_0,S_1)-\Delta S_0,\quad
 s_0=(S_0,\cdot).
\]
Here $S_i$ are coefficients of the complete master candidate at the
stage being tested. They need not be the bare counterterms $C_i$ alone.
The finite $\mathcal A_0$ and $\mathcal A_1$ are not set to zero.

\begin{theorem}[Relative order-two obstruction with the full finite defects]
\label[theorem]{thm:v64-relative-BV}
For an even supplied change $S_1\mapsto S_1+F$ and
$S_2\mapsto S_2+G$, define
\begin{equation}
 b_F=(S_1,F)+\tfrac12(F,F)-\Delta F.
 \label{eq:v64-bF}
\end{equation}
Then the exact changes of master coefficients are
\begin{equation}
 \boxed{\delta\mathcal A_1=s_0F,\qquad
 \delta\mathcal A_2=s_0G+b_F,}
 \label{eq:v64-relative-defect}
\end{equation}
and
\begin{equation}
 \boxed{s_0b_F=-(F,\mathcal A_1)
 -(S_1+F,s_0F)+\Delta(s_0F).}
 \label{eq:v64-relative-consistency}
\end{equation}
In particular, a first-order closed change gives a second-order cocycle
only after the remaining term $-(F,\mathcal A_1)$ is controlled. If
$\mathcal A_0\ne0$, $s_0$ is not itself a differential and one cannot
invoke its cohomology without an additional ordinary/quotient passage.
\end{theorem}
\begin{proof}
Expand the bilinear master bracket to order two. For the consistency
identity use the graded Jacobi rule
\[
 s_0(S_1,F)=(s_0S_1,F)-(S_1,s_0F),\quad
 s_0\tfrac12(F,F)=-(F,s_0F),
\]
and the BV product identity
$s_0\Delta F+\Delta s_0F=(\Delta S_0,F)$.
Their sum is \eqref{eq:v64-relative-consistency}.
Moreover $s_0^2G=(\mathcal A_0,G)$; together with the graded
antisymmetry of the bracket this reproduces the exact changed
order-two consistency equation of V61, including $(G,\mathcal A_0)$.
No finite classical master equation was used.
\end{proof}

At the bare-action level a change $C_1\mapsto C_1+F$ also induces the
loop coefficient \eqref{eq:v64-insertion} and any supplied local
order-two change. At the effective level the corresponding $G$ must
include those induced terms and all other sectors of that same
pushforward. Choosing $G$ to be an arbitrary local polynomial while
omitting $\mathcal I_a[F]$ tests a different effective action.
The finite-reference coefficient in
\cref{prop:v64-positive} is a concrete example of the omitted term.

For a smooth external parameter, derivatives of
$\tfrac12(S_1,S_1)$ include $(S_{1,h},S_{1,k})$ at second order, in
addition to $(S_1,S_{1,hk})$. The contact $\Delta S_1$ is differentiated
in the same convention. The action-only Hessian test in
\cref{thm:v64-hessian} cannot determine these antifield/source
contributions. V55--V60 already showed why a soft Taylor estimate is
not a bound on an unknown coincidence contact.

After the scalar pushforward, this theorem is used with the paired
effective bracket and Laplacian of V59. Re-expressing it in the original
canonical convention requires precisely the reference conversion in
that version. A reference force, scalar determinant or geometric source
term is not deleted on the strength of the algebraic identities above.
The new result specifies the relative equation for the common insertion;
it does not solve its populated local range problem.

\section{A finite common-action test and the remaining continuum interface}
\label{sec:v64-matching}

At a fixed EFT/source order define $\mathsf R C_1$ to be the labelled
collection of open local vertices used in V62 and V63, including their
geometric derivatives and the other requested action/source slots.
This is a finite linear differentiation-and-restriction map. Its entries
are fixed by the native field types, derivative labels, routing and
Wick weights; it is not a fitted collection of independent graph
constants. The bosonic full-Hessian part of its range is tested by
\cref{thm:v64-hessian}. Missing slots can be completed; prescribed
contradictory slots produce a residual and cannot be silently changed.
BV and physical normalization conditions are additional constraints.

\begin{proposition}[Common-action matching is stronger than two graph limits]
\label[proposition]{prop:v64-common-matching}
Suppose a supplied $C_{a,1}$ has the stated local realization and its
restriction $\mathsf RC_{a,1}$ agrees with the actual proper-subgraph
counterterms, including the supplied finite parts, in both primitive
cores. Then the sum of the counterterm terms generated by
\eqref{eq:v64-insertion}, restricted to those cores, is exactly the
proper preparation in V62 and V63. The contact double-local difference
is absorbed only in the same degree-six overall assignment. Subject to
their inherited coefficient budgets, the two primitive remainders have
the combined $O(a[1+\log(1/(a\lambda))])$ limit already proved there.
If the restrictions differ, the differences must be inserted linearly
on every corresponding cut, with the multiplicities
\eqref{eq:v64-cut-identities}. In general these are nonlocal one-loop
functions and are not removable by an overall local $C_2$ alone.
\end{proposition}
\begin{proof}
Apply the common Hessian contraction to each labelled proper
counterterm. \cref{thm:v64-cuts} identifies its opening/closing weights
with the forest terms. The two scalar cuts of the parallel graph are
both retained. For the shared vertex use the exact local-product
conversion of V63. These are finite equalities before passing to a
limit. The analytic conclusion is the sum of the inherited V62 and
V63 estimates, not a new graph bound. For mismatched finite data,
linearity of \eqref{eq:v64-insertion} gives their remaining insertions.
V63's exact scheme-change formula already contains a finite tensor
times a nonpolynomial one-loop function, and
\eqref{eq:v64-kinetic-reference} supplies an additional explicit
common-action example. No assertion that an arbitrary supplied
$C_1$ solves its local BV or normalization conditions is required for
these identities, or follows from them.
\end{proof}

For finite additions with the bounded budgets of V62/V63, their
insertion comparisons are inherited directly. A divergent local
$C_{a,1}$ outside those matched assignments has no automatic uniform
bound from \eqref{eq:v64-insertion}; its local stock and proper forest
must be accounted for. The source integrability map has fixed
coefficient bounds but does not supply such an analytic budget.
No unknown native $S_1$ array is reconstructed from a vacuum amplitude.

The next concrete input is therefore one complete finite coefficient
array for the common $S_1$, with its proper graph restrictions,
source/antifield coefficients and prescribed finite anchors. The
insertion functional, its six second-mark terms, its cut
multiplicities and its order-two consistency residual are now fixed.
No further freedom to choose them independently remains. A successful
range test for that populated array, followed by the paired two-loop
local master/reference equation, is still required.

\section{Ledgers and verification boundary}
\label{sec:v64-ledgers}

\subsection*{Proved}
The common one-loop functional has the single Hessian insertion
\eqref{eq:v64-insertion} at the next order on the native soft panels.
Its complete geometric derivatives and the contact/parallel-edge cut
multiplicities are evaluated. The finite bosonic Hessian compatibility
map has an explicit local section and failed-input residual. A common
scalar kinetic comparison gives the exact lower-reference response
\eqref{eq:v64-kinetic-reference}, its strictly positive conformal
coefficient and fixed-order kernel estimates. The relative order-two
BV defect and its consistency equation retain the actual first-order
and tree defects.

\subsection*{Inherited}
V27 supplies the common completed scalar action and native vertices.
V58 supplies its conformal derivative and hard covariance convention.
V61 supplies the two primitive cores and bridge exclusion. V62 and V63
supply their graphwise forest, volume and rooted bounds. V59 supplies
the paired scalar pushforward and canonical-reference conversion.
These results are used under their original fixed-order assumptions;
their analytic proofs have not been independently re-certified here.

\subsection*{Literature input}
Forest counterterms, finite Gaussian/BV effective actions and modified
reference identities are established methodology
\cite{V64BPHZ,V64IIM,V64DJP}. The finite Hessian test is elementary
coefficient integrability, not a new inverse-variational classification.
No external result supplies the missing populated critical primitive.

\subsection*{Conditional}
Combining the primitive limits in the chosen physical scheme requires
the actual common $S_1$ restrictions and the inherited budgets.
A finite change of scalar kinetic normalization is evaluated but not
installed. The complete $S_2$ contains other graphs and source sectors
besides the two primitive scalar-containing action cores.

\subsection*{Open}
The common source-complete critical $S_1$ coefficient array and the
paired populated two-loop local master/reference equation remain open.
The new action-integrability section is not a BV solution, and the
nonzero finite normalization response is not an anomaly. Internal
gauge/chiral sectors, arbitrary-loop closure and an invariant
interacting EFT class remain separate programme gates.

\begin{center}\small
\begin{tabular}{>{\raggedright\arraybackslash}p{.23\textwidth}>{\raggedright\arraybackslash}p{.69\textwidth}}
\toprule
Variable & Scope in V64\\\midrule
Loop/source order & One order-$\hbar$ counterterm inserted at total order
$\hbar^2$, through two ordinary metric marks; finite source algebra
retained in the operator identities.\\
Native integration & Same signed metric, scalar and ghost references;
no new loop integral of the critical unknown array is evaluated.\\
Concrete comparison & Four massless real scalars, fixed lower scale and
supplied $\zeta$; no active kinetic normalization changes.\\
Cutoff/volume & Exact conformal finite sum; two-mark lower-annulus
comparison on \eqref{eq:v64-domain}; rooted bounds have no total-volume
factor.\\
Uniformity & Fixed jet/source order, profile seminorms, chart and coupling
margins. No estimate uniform in unbounded order or RG steps.\\
Verification & Exact matrix, polynomial, Wick and finite BV tests;
rational fixtures are not numerical native loop coefficients.\\\bottomrule
\end{tabular}
\end{center}

\begin{thebibliography}{99}
\bibitem{V64BPHZ} D.~N.~Blaschke, F.~Gieres, F.~Heindl, M.~Schweda
and M.~Wohlgenannt, BPHZ renormalization and its application to
non-commutative field theory, \emph{Eur. Phys. J. C} \textbf{73}
(2013), 2566; arXiv:1307.4650.
\bibitem{V64IIM} Y.~Igarashi, K.~Itoh and T.~R.~Morris,
BRST in the exact renormalization group,
\emph{Prog. Theor. Exp. Phys.} (2019), 103B01; arXiv:1904.08231.
\bibitem{V64DJP} M.~Doubek, B.~Jur\v co and J.~Pulmann,
Quantum $L_\infty$ algebras and the homological perturbation lemma,
\emph{Commun. Math. Phys.} \textbf{367} (2019), 215--240;
arXiv:1712.02696.
\end{thebibliography}

\section*{Append-only continuation marker: V64}
\addcontentsline{toc}{section}{Append-only continuation marker: V64}
\label{sec:v64-continuation-marker}
\textbf{End of V64.} Preserve Sections 1--497. The common-action
insertion, cut weights, Hessian compatibility test and relative
order-two defect are explicit. The populated common critical
normalization and full two-loop source/reference closure remain open.
No finite-parameter or nonperturbative ultraviolet completion is claimed.

\clearpage
% ============================================================================
% V65 -- quantum completion of the populated V38 logarithmic normalization
% ============================================================================
\section{V65: continue the populated logarithmic packet, not an unknown array}
\label{sec:v65-start}

V64 fixes the insertion of a common $C_1$ but does not populate its complete
critical coefficient array. There is nevertheless a nonzero part already
known: V38's logarithmic scalar--gravity packet
\eqref{eq:v38-log-counterterm}. Its next-order completion was expressly left
unevaluated in Section~\ref{sec:v38-sources}. We now compute that completion
for the specified field-redefinition path. In particular, a six-scalar
six-derivative coefficient and a nonzero auxiliary-density contact are
obtained analytically. This uses known input rather than choosing the
uncomputed remainder of $C_1$.

The general quantum canonical transformation is established methodology
\cite{V65BLT,V65CPV}. The new content is its evaluated application to the
native tensor \eqref{eq:v38-native-F}, the actual coframe Jacobian and
auxiliary measure, and the explicit interface with V64's insertion.
The construction is a comparison between source-complete schemes. It is
not a new choice of the unspecified finite Wilson coefficients, and no
active normalization is changed in this installment.

\paragraph{Two distinct settings.}
Finite identities below hold in the full canonical odd-cotangent algebra
of the existing carrier, before a scalar or metric hard-mode pushforward.
They use the actual finite tree master \emph{candidate}, with no claim
that its master defect vanishes. A finite realization $F_a(\phi)$ is
independent of the metric and uses a fixed linear difference/filter on
each scalar occurrence. In the ordinary local coefficient calculations
$F=F^{\mathcal K}(\partial\phi)$ is exactly V38's metric-coordinate
polynomial. These calculations do not evaluate a lattice propagator on a
continuum wave. A comparison between the two settings retains its actual
finite-jet discrepancy explicitly.

The base density, denoted $\rho_a$, is the inherited leading auxiliary
and coordinate density, independent of antifields. It is not the scalar
volume density $\sqrt{\det g}$. In the convention of
\eqref{eq:v15-measure-once},
\begin{equation}
 \Delta_{\rho_a}=\Delta_0+(\log\rho_a,\cdot),\qquad
 \mathfrak A_{\rho_a}(S)=\tfrac12(S,S)-\hbar\Delta_{\rho_a}S.
 \label{eq:v65-measure-convention}
\end{equation}
Additional perturbative density coefficients, when present, remain in
the corresponding action coefficients. Equivalently one may put
$-\hbar\log\rho_a$ in the action and use $\Delta_0$; the two conventions
are never added together. The canonical cell pairing is $v=a^4$ and
$\Delta_0$ contains $v^{-1}$ in coefficient coordinates. Antifields
remain external sources.

\section{An exact metric shear and its quantum master transport}
\label{sec:v65-flow}

Use the same metric-coordinate generator as V38,
\begin{equation}
 Q_F=\int_a\eta_g:F_a(\phi),\qquad
 \mathcal D_F=(\,\cdot\,,Q_F),\qquad
 \varepsilon\in\R\text{ formal}.
 \label{eq:v65-generator}
\end{equation}
Here $Q_F$ is odd of ghost number $-1$, while $\mathcal D_F$ is even.
The ordinary generator is obtained by replacing the prescribed finite
scalar derivative by $\partial$. All transposes below are the actual
cell-pairing adjoints.

\begin{proposition}[Finite shear, cotangent source, and base Jacobian]
\label[proposition]{prop:v65-shear}
In metric coordinates the flow generated by $\mathcal D_F$ is exactly
\begin{equation}
 \boxed{g_\varepsilon=g+\varepsilon F_a,\quad
 \phi_\varepsilon=\phi,\quad
 (\eta_g)_\varepsilon=\eta_g,\quad
 \upsilon_\varepsilon=\upsilon-\varepsilon( D_\phi F_a)^\dagger\eta_g.}
 \label{eq:v65-shear}
\end{equation}
Other canonical pairs are unchanged. The base-field Jacobian is one in
metric coordinates, but the weighted divergence is
\begin{equation}
 \boxed{\Delta_{\rho_g}Q_F=\mathcal D_F\log\rho_g.}
 \label{eq:v65-metric-divergence}
\end{equation}
In the existing coframe coordinate $g=G(q)$, with $A(q)=DG(q)$ and
$J_G(q)=\prod_x\det A(q_x)$, the corresponding base map is
\begin{equation}
 q_\varepsilon=G^{-1}(G(q)+\varepsilon F_a),\qquad
 \operatorname{Jac}(q_\varepsilon,\phi;q,\phi)
       =\frac{J_G(q)}{J_G(q_\varepsilon)}.
 \label{eq:v65-coframe-flow}
\end{equation}
For the actual density $\rho_q=(\rho_g\circ G)J_G$,
\begin{equation}
 \boxed{
 \frac{\rho_q(q_\varepsilon)\operatorname{Jac}(q_\varepsilon;q)}
      {\rho_q(q)}
 =\frac{\rho_g(G(q)+\varepsilon F_a)}{\rho_g(G(q))}.}
 \label{eq:v65-density-ratio}
\end{equation}
Thus the coframe Jacobian cancels its coordinate-density partner, not the
physical auxiliary-density derivative.
\end{proposition}
\begin{proof}
The four Hamilton equations in \eqref{eq:v65-shear} follow directly
from the canonical bracket. Their right-hand sides are constant along
the flow because $\phi$ and $\eta_g$ do not move. Equivalently, the
canonical one-form is unchanged:
$\eta_g\,d(g+\varepsilon F_a)+
 (\upsilon-\varepsilon DF_a^\dagger\eta_g)d\phi
 =\eta_g\,dg+\upsilon\,d\phi$.
The base Jacobian has triangular blocks $(I,\varepsilon DF_a;0,I)$.
Also $\Delta_0Q_F=0$: differentiating its metric antifield leaves $F_a$,
whose metric derivative is zero, and there is no scalar-antifield
summand in $Q_F$. Scalar dependence in $F_a$ does not create an extra
diagonal contraction. This proves \eqref{eq:v65-metric-divergence}.
Differentiating $G(q_\varepsilon)=G(q)+\varepsilon F_a$ in $q$ gives
$D_qq_\varepsilon=A(q_\varepsilon)^{-1}A(q)$ and hence
\eqref{eq:v65-coframe-flow}. The density factors cancel as in
\eqref{eq:v65-density-ratio}. The coframe antifield is transported by
$\eta_{q,\varepsilon}=A(q_\varepsilon)^TA(q)^{-T}\eta_q$; the scalar
antifield has the same shift as in \eqref{eq:v65-shear} after this
cotangent conversion. No spacetime inverse is involved.
\end{proof}

\begin{theorem}[Quantum shear transports the actual defect]
\label[theorem]{thm:v65-quantum-flow}
Let $\rho$ be an antifield-independent nonzero base density in the
specified finite chart. For an arbitrary even master candidate $S$, set
\begin{equation}
 \boxed{S_\varepsilon=e^{\varepsilon\mathcal D_F}S
       -\hbar\int_0^\varepsilon
              e^{t\mathcal D_F}\Delta_\rho Q_F\,dt.}
 \label{eq:v65-quantum-flow}
\end{equation}
Then
\begin{equation}
 \boxed{\mathfrak A_\rho(S_\varepsilon)
             =e^{\varepsilon\mathcal D_F}\mathfrak A_\rho(S).}
 \label{eq:v65-defect-flow}
\end{equation}
In metric coordinates the integral in \eqref{eq:v65-quantum-flow} is
$\log\rho_g(g+\varepsilon F_a)-\log\rho_g(g)$. No zero-defect hypothesis
is needed.
\end{theorem}
\begin{proof}
An antifield-independent base density gives $\Delta_\rho^2=0$ in the
canonical algebra. The BV product identity implies
$[\Delta_\rho,\mathcal D_F]H=-(\Delta_\rho Q_F,H)$.
Equation \eqref{eq:v65-quantum-flow} solves
$\dot S_\varepsilon=\mathcal D_FS_\varepsilon-\hbar\Delta_\rho Q_F$.
Differentiating its master defect, using the commutator identity and
$\Delta_\rho^2Q_F=0$, gives
$\dot{\mathfrak A}_\rho=\mathcal D_F\mathfrak A_\rho$.
The initial value proves \eqref{eq:v65-defect-flow}. The last assertion
uses \eqref{eq:v65-metric-divergence} and the fundamental theorem of
calculus. This proof also applies to formal local coefficient series.
It does not assert an unbounded real-chart or contour equivalence for a
nonlinear finite integral. Such an equivalence additionally transports
the integration domain, observables and the source arguments.
\end{proof}

For a polynomial realization every retained coefficient of this identity
is a finite algebraic statement. For an analytic coframe germ it holds
inside its declared chart, or formally to every fixed coefficient order.
The finite parameter cannot be kept uniformly small as $a\to0$ merely
because it is proportional to a loop-counting variable.

\section{The specified first-order packet has a determined second-order continuation}
\label{sec:v65-order-two}

At a fixed cutoff write $S=S_0+\hbar S_1^-+\hbar^2S_2^-+O(\hbar^3)$.
For the already evaluated V38 coefficient put
\begin{equation}
 b_a=\frac{\log(1/(a\lambda))}{8\pi^2\chi^2},\qquad
 \varepsilon=\hbar b_a.
 \label{eq:v65-b}
\end{equation}
A minus superscript denotes the comparison scheme before the specified
packet, not a new prescription for the unknown part of $S_1$.

\begin{theorem}[Complete order-two relative action]
\label[theorem]{thm:v65-order-two}
For the exact canonical realization of the packet,
\begin{align}
 S_1^+&=S_1^-+b_a\mathcal D_FS_0,\label{eq:v65-S1}\\
 S_2^+-S_2^-
 &=\boxed{b_a\mathcal D_FS_1^-
   +\frac{b_a^2}{2}\mathcal D_F^2S_0
   -b_a\Delta_\rho Q_F.}\label{eq:v65-S2}
\end{align}
If $\mathfrak A_j^\pm$ denote coefficients of the actual master defects,
\begin{align}
 \mathfrak A_0^+&=\mathfrak A_0^-,\notag\\
 \mathfrak A_1^+&=\mathfrak A_1^-+b_a\mathcal D_F\mathfrak A_0^-,\notag\\
 \mathfrak A_2^+
 &=\boxed{\mathfrak A_2^-+b_a\mathcal D_F\mathfrak A_1^-
              +\frac{b_a^2}{2}\mathcal D_F^2\mathfrak A_0^-.}
 \label{eq:v65-order-defect}
\end{align}
The cross term in \eqref{eq:v65-S2} consumes the complete previous
one-loop functional, including its sources and any separately assigned
density coefficients. It is not determined by the known packet alone.
\end{theorem}
\begin{proof}
Expand both sides of \eqref{eq:v65-quantum-flow} and
\eqref{eq:v65-defect-flow} in $\hbar$. The integral contributes
$-\hbar^2b_a\Delta_\rho Q_F$ at this order; its next term starts at
$\hbar^3$. The remaining coefficients give the formulas. In particular,
no term involving $\mathfrak A_0^-$ or $\mathfrak A_1^-$ is discarded.
\end{proof}

In a flat-measure representation let
$\widehat S_1^-=S_1^--\log\rho$. Equation \eqref{eq:v65-S2} is exactly
\begin{equation}
 S_2^+-S_2^-
 =b_a\mathcal D_F\widehat S_1^-
 +\frac{b_a^2}{2}\mathcal D_F^2S_0-b_a\Delta_0Q_F.
 \label{eq:v65-once}
\end{equation}
This identity is the once-counted measure rule, not a second density
counterterm.

\begin{proposition}[Literal comparison with the previously retained jet]
\label[proposition]{prop:v65-realization}
Let $L_a^{38}$ be a specified finite realization of the retained V38
local packet. Set
\begin{equation}
 E_a=L_a^{38}-b_a\mathcal D_{F_a}S_{a,0}.
 \label{eq:v65-realization-error}
\end{equation}
There is no need to assume $E_a=0$. To leave the literal first-order
coefficient $C_1^{\rm rest}+L_a^{38}$ unchanged in the comparison, use
$S_1^-=C_1^{\rm rest}+E_a$ in \eqref{eq:v65-S1}. Its next-order cross
term contains $b_a\mathcal D_{F_a}E_a$. This term and its insertions must
be retained until their own finite realization and hard-contraction
bounds are proved.
\end{proposition}
\begin{proof}
Substitution in \eqref{eq:v65-S1} gives the stated coefficient exactly.
Substitution in \eqref{eq:v65-S2} displays the additional term. Agreement
of retained ordinary jets does not imply an estimate on every hard
insertion of $E_a$, as V61 already demonstrates. This prescription is
algebraic comparison data; it neither fixes $C_1^{\rm rest}$ nor declares
the comparison scheme to satisfy its master equation.
\end{proof}

For the ordinary minimal metric/scalar sources, $\mathcal D_F^2S_{0,\rm
min.src}=0$. Indeed their first variation is V38's
$\int\eta_g:\mathcal R_F(c,\phi)$, independent of $g$ and $\upsilon$;
its bracket with $Q_F$ is zero. The same statement holds for a finite
minimal source affine in $g$ and a metric-independent scalar source.
Gauge-fixing action, nonminimal action coefficients and quantum sources
are still included in the full derivatives in \eqref{eq:v65-S2}; no
claim about their vanishing is made.

\section{An evaluated six-scalar coefficient and the retained gravitational partner}
\label{sec:v65-sextic}

The ordinary scalar action is the already fixed
\begin{equation}
 S_H(g,\phi)=\frac12\int\sqrt{\det g}
      \left[\operatorname{tr}(g^{-1}\mathsf V)+zU\right],\qquad
 \mathsf V_{ab}=\sum_A\partial_a\phi_A\partial_b\phi_A,
 \quad U=\sum_A\phi_A^2.
 \label{eq:v65-SH}
\end{equation}
$\mathsf V$ and $U$ do not vary along the metric shear. For any symmetric
metric-independent $F$, put $f_1=\operatorname{tr}F$ and
$f_2=\operatorname{tr}(F^2)$.

\begin{theorem}[Second scalar action coefficient of the native shear]
\label[theorem]{thm:v65-sextic}
At $g=I$, the coefficient $\tfrac12\mathcal D_F^2S_H$ is
\begin{equation}
 \boxed{\int\left[\mathcal P_6(\mathsf V)+zU\mathcal Z_4(\mathsf V)\right],}
 \label{eq:v65-PZ}
\end{equation}
where
\begin{align}
 \mathcal P_6
 &=\tfrac12\operatorname{tr}(F^2\mathsf V)
  -\tfrac14 f_1\operatorname{tr}(F\mathsf V)
  +\left(\tfrac1{16}f_1^2-\tfrac18f_2\right)
       \operatorname{tr}\mathsf V,\notag\\
 \mathcal Z_4&=\tfrac1{16}f_1^2-\tfrac18f_2.
 \label{eq:v65-general-PZ}
\end{align}
For precisely $F=F^{\mathcal K}$, write
$\mathsf V=\left(\begin{smallmatrix}t&b^T\\b&S\end{smallmatrix}\right)$,
$s=\operatorname{tr}S$, $u=\operatorname{tr}(S^2)$,
$w=\operatorname{tr}(S^3)$, $\beta=b^Tb$, and $\gamma=b^TSb$. Then
\begin{equation}
 \boxed{\begin{aligned}
 8192\mathcal P_6={}&1536t^3-1920t^2s+1947ts^2+503s^3\\
 &-3465(t+s)u+4900w
       +5152t\beta-5824s\beta+10976\gamma,\\
 8192\mathcal Z_4={}&-768t^2+704ts+283s^2-1225u-1568\beta.
 \end{aligned}}
 \label{eq:v65-native-PZ}
\end{equation}
There is no additional flavour multiplicity: $\mathsf V$ and $U$
already contain all four scalar components.
\end{theorem}
\begin{proof}
The pointwise expansions are
\[
 (I+\varepsilon F)^{-1}=I-\varepsilon F+\varepsilon^2F^2+O(\varepsilon^3),
 \quad
 \sqrt{\det(I+\varepsilon F)}
 =1+\tfrac\varepsilon2f_1
   +\varepsilon^2(\tfrac18f_1^2-\tfrac14f_2)+O(\varepsilon^3).
\]
Multiplying these in \eqref{eq:v65-SH} proves
\eqref{eq:v65-general-PZ}. Insert the three blocks of
\eqref{eq:v38-native-F}. Expanding the traces and collecting the ten
independent cubic Gram contractions gives \eqref{eq:v65-native-PZ};
off-diagonal entries are counted twice in the trace. The accompanying
checker verifies this polynomial identity in all ten independent
symmetric Gram entries, without a rotational averaging assumption.
\end{proof}

For one temporal scalar gradient, $\mathsf V=\operatorname{diag}(t,0,0,0)$,
\begin{equation}
 \boxed{\mathcal P_6=\frac3{16}t^3,\qquad
        \mathcal Z_4=-\frac3{32}t^2.}
 \label{eq:v65-temporal}
\end{equation}
Thus the specified canonical continuation has a nonzero
$\hbar^2 b_a^2\mathcal P_6$ six-derivative scalar-six coefficient.
It is fixed by the known first-order tensor along this path, not a
calculated new two-loop Feynman integral or a universal Wilson coefficient.
Additional second-order generators and physical finite normalizations
can change the full $S_2$.

The gravitational action supplies a different partner. On its flat
stationary branch,
\begin{equation}
 \boxed{\left[\tfrac12\mathcal D_F^2S_g\right]_{g=I}
        =\tfrac12\langle F,H_g^{\rm cl}F\rangle.}
 \label{eq:v65-gravity-four}
\end{equation}
This has four scalar arguments and six derivatives. $H_g^{\rm cl}$ is
the inherited signed \emph{Euler} Hessian, not a positive or gauge-fixed
inverse. It can be nonzero for varying scalar jets. Constant-gradient
tests make $F$ constant and isolate \eqref{eq:v65-PZ}; they do not justify
omitting \eqref{eq:v65-gravity-four} from the general functional.

The previous arity-five bank is therefore not large enough to record
this full second-order continuation: the scalar-six term must be added
when that external order is claimed. This is a finite EFT enlargement,
not an infinite-parameter claim at a single prescribed order.

\begin{proposition}[The populated packet has a nonzero self-bracket]
\label[proposition]{prop:v65-self-bracket}
Let $C_F=(S_{0,\infty},Q_F)$ be V38's ordinary complete packet, and
isolate the scalar-six component of $\tfrac12(C_F,C_F)$. For one scalar
with constant first jet $v=e_0+e_1$, flat metric, and an affine rotation
ghost with $\partial_1c^0=1$, $\partial_0c^1=-1$, its integrated-by-parts
local density has coefficient
\begin{equation}
 \boxed{\left[\tfrac12(C_F,C_F)\right]_{\phi^6c}
                  =-\frac{201}{512}.}
 \label{eq:v65-self-bracket}
\end{equation}
The coefficient of the second scalar action term transforms by
$+201/512$ on the same jet, so the two cancel in the ordinary
order-two identity. This is a coefficient of the packet's self-interaction,
not the value of the full two-loop master defect.
\end{proposition}
\begin{proof}
Write $\Omega^0{}_1=1$, $\Omega^1{}_0=-1$ and
$\dot{\mathsf V}=\Omega^T\mathsf V+\mathsf V\Omega$. If
$F=\mathsf B\mathsf V$, V38's source tensor on these jets is
$R_F=\Omega^TF+F\Omega-\mathsf B\dot{\mathsf V}$.
The scalar-six self-bracket comes solely from contracting the metric
antifield in $\int\eta_g:R_F$ against the metric derivative of
$\int E_H:F$. It is $D_g^2S_H[F,R_F]$. At $g=I$, $z=0$, its density is
\begin{align*}
 &\tfrac12\operatorname{tr}[(FR_F+R_FF)\mathsf V]
 -\tfrac14\{\operatorname{tr}F\operatorname{tr}(R_F\mathsf V)
            +\operatorname{tr}R_F\operatorname{tr}(F\mathsf V)\}\\
 &\hspace{10mm}
 +\{\tfrac18\operatorname{tr}F\operatorname{tr}R_F
              -\tfrac14\operatorname{tr}(FR_F)\}
                    \operatorname{tr}\mathsf V.
\end{align*}
Substituting $\mathsf V=vv^T$ gives $-201/512$ by rational matrix
multiplication. Independently, differentiate the explicit polynomial
$\mathcal P_6(\mathsf V)$ of \eqref{eq:v65-native-PZ} in direction
$\dot{\mathsf V}$ to obtain $+201/512$.
More generally these two polynomials sum to zero for every symmetric
$\mathsf V$, by direct polarization of
\eqref{eq:v65-general-PZ}. This is also the scalar-six restriction of
$ (S_0,\mathcal D_F^2S_0)+(\mathcal D_FS_0,\mathcal D_FS_0)=0$
on the ordinary classical master branch. The first-jet density is a
nonzero homogeneous polynomial of degree six for one scalar; it cannot
be a scalar total divergence by itself, since a first-order one-scalar
null Lagrangian has zero Hessian in its gradient variables. The affine
jets are local polynomial tests, not periodic fields. Additional
counterterm sectors and the measure contact have their own scalar,
parameter and loop degrees and are not assigned values by this test.
\end{proof}

\section{The previously deferred auxiliary-density contact is nonzero}
\label{sec:v65-density}

The physical metric-coordinate measure is not flat. V15 evaluated its
restriction to the lapse line:
\begin{equation}
 g=\operatorname{diag}(n^2,1,1,1),\qquad
 \rho_g=n^{-7},\quad J_G=n,\quad \rho_q=n^{-6},
 \quad q_{00}=2(n-1).
 \label{eq:v65-lapse-density}
\end{equation}
This input includes the actual stationary auxiliary/Haar terms on that
line. We do not infer the density away from it from this formula.

\begin{theorem}[A populated order-two measured contact]
\label[theorem]{thm:v65-density}
For a metric shear tangent to V15's spatially constant lapse line,
$n=n(x^0)$ and $F=f(x^0)\,e_0e_0^T$ with $f$ independent of $g$,
the complete weighted divergence is
\begin{equation}
 \boxed{\Delta_{\rho_a}Q_F
             =-\frac72\sum_x\frac{f(x)}{n(x)^2}}
 \label{eq:v65-lapse-div}
\end{equation}
on the inherited lapse restriction whenever it is defined, and in
particular for its constant-background coefficient. At the flat
constant-background coefficient of the \emph{native} tensor, the Gram
test
\begin{equation}
 \mathsf V=\operatorname{diag}(t,4t,4t,4t),\qquad
 F^{\mathcal K}=\tfrac92t\,e_0e_0^T
 \label{eq:v65-lapse-Gram}
\end{equation}
gives the order-two measure contribution per physical volume
\begin{equation}
 \boxed{
 \left[-b_a\Delta_{\rho_a}Q_{F^{\mathcal K}}\right]_{
       g=I,\,\mathsf V=\operatorname{diag}(t,4t,4t,4t)}
       /V_L=\frac{63b_a}{4a^4}t.}
 \label{eq:v65-density-populated}
\end{equation}
This is nonzero for $t>0$ and $a\lambda<1$. The evaluation is a local
constant-jet coefficient test, not a claim that affine scalar fields are
periodic configurations or that this contact is an anomaly.
\end{theorem}
\begin{proof}
In metric coordinates use
$\partial_{g_{00}}\log(n^{-7})=-7/(2n^2)$ and
\eqref{eq:v65-metric-divergence}. One contraction at each site produces
\eqref{eq:v65-lapse-div}. Equivalently, in the active coordinate,
$R_q=f/n$, so
\[
 \partial_{q_{00}}R_q=-\frac{f}{2n^2},\qquad
 R_q\partial_{q_{00}}\log\rho_q=-\frac{3f}{n^2}.
\]
Their sum is $-7f/(2n^2)$, not $-3f/n^2$. This explicitly retains the
coordinate Jacobian. Substitution in \eqref{eq:v38-native-F} verifies
\eqref{eq:v65-lapse-Gram}. At flat background the integrated density's
one-metric coefficient is translation invariant; its constant lapse
coefficient is fixed by \eqref{eq:v65-lapse-density}. The Gram test
selects precisely that component. The site density is $a^{-4}$ per
physical volume, proving \eqref{eq:v65-density-populated}. No value of an
unknown density derivative in a spatial-metric direction is used.
\end{proof}

For the same constant-jet test, the scalar action coefficients are
$\mathcal P_6=-729t^3/64$ and $\mathcal Z_4=-81t^2/64$.
They occur with $b_a^2$ and have different scalar and cutoff degree from
\eqref{eq:v65-density-populated}; none cancels the other by a vacuum
normalization. The cross term $b_a\mathcal D_FS_1^-$ still has to be
included before any claim about the complete next-order coefficient.

The density formula is not a new density choice. If the inherited
$-\log\rho_q$ is kept in the flat-measure $S_1^-$, its directional
derivative together with $-\Delta_0Q_F$ produces exactly the same
contact by \eqref{eq:v65-once}. Counting both representations would
double it. Nor does a vanishing flat-metric-coordinate Jacobian imply
that the physical weighted contact vanishes.

\section{Observable transport, insertion bookkeeping, and bounded local operations}
\label{sec:v65-interface}

For a supplied geometric observable $\mathcal O(g,\phi)$, the same path
requires
\begin{equation}
 \boxed{\mathcal O^+
       =\mathcal O+\hbar b_a\mathcal D_F\mathcal O
         +\frac{\hbar^2b_a^2}{2}\mathcal D_F^2\mathcal O
         +O(\hbar^3).}
 \label{eq:v65-observables}
\end{equation}
This is an identity for transformed source coefficients, not an assertion
of equality of generating functions at unchanged external sources.
For volume at flat metric the second contact is explicitly
\begin{equation}
 \boxed{\tfrac12\mathcal D_F^2\sqrt{\det g}\big|_{g=I}
       =2\mathcal Z_4(\mathsf V).}
 \label{eq:v65-volume}
\end{equation}
For the temporal test \eqref{eq:v65-temporal}, the first volume variation
is zero but the second coefficient equals $-3t^2/16$. Therefore even a
vanishing first contact does not justify omitting the next source term.
Curvature and stress sources follow by differentiating their full native
coefficient functionals twice with fixed $F$, including every derivative
on $F(\partial\phi)$. They are specified by \eqref{eq:v65-observables};
no uncomputed component formula is claimed evaluated.

\begin{proposition}[Insertion ledger for the selected completion]
\label[proposition]{prop:v65-insertion-ledger}
On V64's bridge-excluded panels, a first-order change
$\delta C_1=b_a\mathcal D_FS_0$ changes the order-two effective action by
\begin{equation}
 \boxed{\delta\Gamma_2
     =\delta C_2+\mathcal I_a[b_a\mathcal D_FS_0],\qquad
 \mathcal I_a[F]=\tfrac12\operatorname{STr}(\mathbb G_aD_\Phi^2F).}
 \label{eq:v65-insertion-ledger}
\end{equation}
Here $\delta C_2$ is \eqref{eq:v65-S2}, with its complete $C_1^-$ cross
term. The source functions are transported by \eqref{eq:v65-shear} and
\eqref{eq:v65-observables}. The insertion is not included a second time
inside $\delta C_2$. If the literal implementation has the discrepancy
\eqref{eq:v65-realization-error}, its corresponding cross terms and
insertion are included as well.
\end{proposition}
\begin{proof}
This is V64's insertion theorem applied to the now specified first-order
change, plus the tree insertion of its second-order change. The graph
counting belongs to that inherited result. Equation \eqref{eq:v65-S2}
was derived at the bare action level and contains no hard Gaussian
integration, so it cannot replace $\mathcal I_a[\delta C_1]$. Conversely
adding the insertion inside both coefficients would count it twice.
\end{proof}

The exact master transport in \eqref{eq:v65-defect-flow} uses the
canonical bracket/Laplacian pair. After a hard-mode pushforward the
reference pair and source map must be transported together as in V59.
A scalar-dependent metric shear mixes field blocks; one cannot merely
retain a fixed scalar mask and substitute the transformed scalar Euler
operator. V58's explicit noncommuting-mask calculation remains relevant.
No equality for a partially integrated functional with untransported
retained fields, antifields or contour is asserted here.

The local computations have fixed-order coefficient bounds. From V38,
$\|F^{\mathcal K}\|_{\rm HS}\le\tfrac32\|\mathsf V\|_{\rm HS}$.
The explicit trace formula therefore gives
\begin{equation}
 \boxed{|\mathcal P_6|\le\tfrac{63}{16}\|\mathsf V\|_{\rm HS}^3,
 \qquad |\mathcal Z_4|\le\tfrac{27}{32}\|\mathsf V\|_{\rm HS}^2.}
 \label{eq:v65-pointwise-bounds}
\end{equation}
Indeed $|\operatorname{tr}F|\le2\|F\|_{\rm HS}$,
$|\operatorname{tr}\mathsf V|\le2\|\mathsf V\|_{\rm HS}$ and
Cauchy--Schwarz bound the three terms in $\mathcal P_6$ by
$(1/2+1/2+3/4)\|F\|_{\rm HS}^2\|\mathsf V\|_{\rm HS}$.
The bound on $\mathcal Z_4$ is
$(1/4+1/8)\|F\|_{\rm HS}^2$. This proves both constants.

At each fixed local arity the ordinary $\mathcal D_F$ operation raises
derivative weight by two and raises minimum arity by one. It is a bounded
finite tensor map after the coframe chart and coefficient basis are
fixed. With source windows of scale $\mu$, the scalar-six term has
rooted multilinear bound $C_b|b_a|^2\mu^6$, the gravitational
four-scalar term has bound $C_b\chi|b_a|^2\mu^6$, and the displayed
measure contact has coefficient size $C|b_a|a^{-4}\mu^2$.
These follow by polarizing the local polynomials and differentiating
the fixed smooth windows; their relative-coordinate kernels have no
total-volume multiplier. They are \emph{counterterm} bounds, not small
remainders: the logarithms and the $a^{-4}$ remain. General weighted
density derivatives and hard contractions require their own bounds.

\section{Status and the next common-action coefficient}
\label{sec:v65-ledgers}

\subsection*{Proved}
The specified scalar-dependent metric shift has an exact finite cotangent
flow, with its actual coframe Jacobian. The density-corrected quantum flow
transports a nonzero master defect rather than assuming it away. Its
complete second-order coefficient is given in \eqref{eq:v65-S2}.
The native tensor produces the evaluated scalar-six/mass polynomial
\eqref{eq:v65-native-PZ}, the retained signed gravitational Hessian
partner, and the nonzero auxiliary-density coefficient
\eqref{eq:v65-density-populated}. Observable and insertion bookkeeping
are recorded in the same comparison.

\subsection*{Inherited}
V38 supplies the tensor, its logarithmic coefficient, the ordinary
first-order packet and its coframe realization. V15 supplies the actual
lapse-line density and once-counted measure convention. V64 supplies the
one-counterterm hard insertion and its bridge hypothesis. V59 supplies
the paired hard-pushforward convention. None of their analytic proofs is
claimed independently verified by the present algebraic checker.

\subsection*{Literature input}
Quantum anticanonical transformations and higher-order field-redefinition
matching are established methods \cite{V65BLT,V65CPV,V65Schwarz}. Their
general existence is not the new theorem. The finite tensor coefficients,
measured lapse test, and retained source/counterterm ledger are derived
here from the inherited carrier and packet.

\subsection*{Conditional}
The formulas select one canonical second-order continuation, with no new
order-two generator. It is not the only admissible physical normalization.
The full $b_a\mathcal D_FS_1^-$ term requires the actual common $S_1^-$
array. A literal finite implementation requires its discrepancy
$E_a$ and hard insertions. Counterterm addition preserves already
restored identities only under the actual defect hypotheses, as shown
explicitly by \eqref{eq:v65-order-defect}.

\subsection*{Open}
The nonredundant critical part of the common source-complete $S_1$, the
remaining populated local descent, and the complete two-loop
master/reference equation remain open. This computation removes the
previously unpaid next-order completion of one \emph{populated}
first-order direction; it does not fill the unknown coefficient array.
No two-loop native integral, chiral-line theorem, invariant many-shell
class or convergence of a full formal series is claimed.

\begin{center}\small
\begin{tabular}{>{\raggedright\arraybackslash}p{.24\textwidth}>{\raggedright\arraybackslash}p{.68\textwidth}}
\toprule
Variable & Scope in V65\\\midrule
Action/source order & Order-$\hbar$ populated packet and its order-$\hbar^2$
canonical continuation; includes scalar arity six and EFT weight six.\\
Native versus ordinary & Finite master/density identities are exact;
explicit scalar polynomial is the ordinary local coefficient of the
same action. Their realization difference is retained.\\
Mass/flavour & Four equal real scalars; polynomial in $z=m_H^2$;
internal gauge, Yukawa and scalar self-couplings zero in evaluated tests.\\
Volume & Coefficient and rooted bounds have no extra total-volume
factor; the complete functional itself is extensive.\\
Cutoff & $b_a$ is logarithmic, the evaluated density contact is
$a^{-4}b_a$. No new vanishing-cutoff error is asserted.\\
Uniformity & Fixed arity, jet order, chart, source window and nonzero
coupling margin. No uniform analytic shear at fixed $\hbar$ as $a\to0$.\\
Verification & Exact tensor polynomials, finite BV identities, coordinate
Jacobian and preservation tests; no numerical native loop evaluation.\\\bottomrule
\end{tabular}
\end{center}

\begin{thebibliography}{99}
\bibitem{V65BLT} I.~A.~Batalin, P.~M.~Lavrov and I.~V.~Tyutin,
Finite anticanonical transformations in field--antifield formalism,
\emph{Eur. Phys. J. C} \textbf{75} (2015), 270; arXiv:1501.07334.
\bibitem{V65CPV} J.~C.~Criado and M.~P\'erez-Victoria,
Field redefinitions in effective theories at higher orders,
\emph{JHEP} \textbf{03} (2019), 038; arXiv:1811.09413.
\bibitem{V65Schwarz} A.~Schwarz,
Geometry of Batalin--Vilkovisky quantization,
\emph{Commun. Math. Phys.} \textbf{155} (1993), 249--260;
arXiv:hep-th/9205088.
\end{thebibliography}

\section*{Append-only continuation marker: V65}
\addcontentsline{toc}{section}{Append-only continuation marker: V65}
\label{sec:v65-continuation-marker}
\textbf{End of V65.} Preserve Sections 1--504. The known logarithmic
packet has an explicit canonical order-two action, measure and source
continuation. The unknown common critical normalization and full
interacting two-loop master/reference closure remain separate. No
finite-parameter or nonperturbative ultraviolet completion is claimed.


% ============================================================================
% V66: the logarithmic scalar-determinant block of the common one-loop action
% ============================================================================
\clearpage
\section{V66: extract a populated common-action block before completing the descent}
\label{sec:v66-start}

V65 completes one already known redundant direction. It does not determine
all of the common one-loop action. We return instead to an actual term in
that action: the closed scalar loop in the metric background. Its full
finite determinant is present in V60, but its universal critical local
logarithm has not been evaluated in this lineage. Computing this block
supplies actual input to V64's common-Hessian insertion, rather than another
section acting on an unspecified coefficient array.

The calculation below concerns the same four equal real scalar components,
with internal gauge, Yukawa and scalar self-interactions zero in this
component. There is no hard graviton or ghost line inside this determinant.
The metric is a retained source. The real positive ordinary metric germ
is used for the elliptic comparison; its finite jets extend holomorphically
to the inherited continued coframe chart. The full Einstein--SM problem still
contains those other lines and the scalar-antifield sectors. In particular,
a scalar determinant does not stand for a simultaneous graviton--matter
determinant at nonzero retained scalar background.

\paragraph{New result and external input.}
The native matching argument is proved from the finite symbols in
\eqref{eq:v27-completed-Higgs}. The scalar heat coefficient and the
logarithmic residue identities are established external inputs
\cite{V66Vassilevich,V66Scott}. They evaluate an ordinary elliptic comparison
operator, not a substitute finite measure. We show why that comparison has
the same logarithmic coefficient as the actual completed finite determinant.
A direct scalar-loop calculation independently verifies its entire
massless bilinear metric coefficient. No claim of a first derivation of
scalar-induced curvature counterterms is made.

Set $z=m_H^2$, $n_H=4$, $\rho=\sqrt{\det g}$ and $g=G(q)=e(q)^Te(q)$
on the inherited coefficient chart. The geometric density $\rho$ is not
the auxiliary integration density $\rho_a$ of V65. Write
\begin{equation}
 \Gamma^H_{a,1}(q,z)=\frac{n_H}{2}
 \Tr\log\bigl(I+C_aV_a(q,z)\bigr),\quad
 C_a=\nu_\lambda\widetilde P_a^{-1},\quad
 V_a=Q_a(q,z)-\widetilde P_a(z).
 \label{eq:v66-native-determinant}
\end{equation}
This is the existing normalized scalar determinant; its common factor
$\hbar$ is removed. The trace, momentum routing, coframe contacts and odd-grid
normalization are unchanged. A scalar coefficient measure is not multiplied
by a second factor of $\rho$.

For a fixed $N\le5$, let $\mathsf J_N$ denote metric arity at most $N$ and
let $\mathsf P_4$ select total external weight exactly four, with
$\deg K=1$ and $\deg z=2$. All statements concern these local coefficient
jets, not a uniform expansion at unbounded metric arity. The domain is
\begin{equation}
 \lambda\ge32768\mu>0,\qquad a\lambda\le2^{-24},\qquad
 \lambda L_\nu\ge1,\qquad 0\le m_H\le M\lambda.
 \label{eq:v66-domain}
\end{equation}
Mass coefficients are derivatives at $z=0$ on the hard reference, not
values of a massless uncut determinant. The positive lower scale is retained.

\section{The native critical coefficient has an ordinary residue}
\label{sec:v66-native-residue}

At fixed metric arity, expand the actual trace as the finite sum
\begin{equation}
 \mathsf J_N\Gamma^H_{a,1}
 =\frac{n_H}{2}\mathsf J_N
   \sum_{r=1}^{N}\frac{(-1)^{r+1}}r\Tr(C_aV_a)^r.
 \label{eq:v66-finite-words}
\end{equation}
Each vertex is at least linear in $q$. This expression includes single-vertex
coframe contacts as well as multi-vertex loops. A derivative of a covariance
with respect to $z$ retains its single $\nu_\lambda$ factor.

\begin{lemma}[One-loop critical logarithm on the completed carrier]
\label{lem:v66-native-log}
Fix a tensor component $I$ of \eqref{eq:v66-finite-words}, including a metric
arity and a Taylor coefficient with $|\alpha|+2j=4$. There is an ordinary
homogeneous loop coefficient $h_I(p)$ of degree $-4$ such that its native
critical coefficient has the form
\begin{equation}
 g_{a,L,I}^{[4]}=\ell_a\,\mathfrak r_I+B_{a,L,I},\qquad
 \ell_a=\log\frac1{a\lambda},\qquad
 \mathfrak r_I=\frac1{(2\pi)^4}\int_{S^3}h_I(\omega)\,d\omega,
 \label{eq:v66-critical-asymptotic}
\end{equation}
where $|B_{a,L,I}|\le C_I$ uniformly in $a,L$ on
\eqref{eq:v66-domain}, for fixed profiles, source order and inverse margins.
The factor $n_H/2$ can either be included in $h_I$ or restored at the end;
it is included in \eqref{eq:v66-critical-asymptotic}.
The coefficient $\mathfrak r_I$ does not depend on the admitted smooth upper
completion or the lower transition profile. It is the corresponding
coefficient of
\begin{equation}
 \frac{n_H}{2}\operatorname{res}
 \log\bigl(P_0^{-1}Q_0(g,z)\bigr),\qquad
 Q_0=-\partial_a(\rho g^{ab}\partial_b)+z\rho,\quad
 P_0=-\partial^2+z.
 \label{eq:v66-residue-target}
\end{equation}
The residue and equality are for integrated local functionals. An
arbitrary extra multiplicative composite source is not included.
\end{lemma}
\begin{proof}
Differentiate the complete routed words before estimating. An ordinary
vertex has total momentum/mass weight two, and an ordinary covariance has
weight minus two. At critical Taylor weight four the complete loop
integrand is therefore homogeneous of degree minus four wherever every
lower mask is one. The terms in which a derivative hits a lower mask are
supported in a fixed multiple of the lower annulus. Scaling $p=\lambda y$
bounds their integrals and normalized sums independently of $a,L$.

Use a fixed $c>0$ small enough that all routed native scalar labels are on
the common interior plateau for $|p|<c/a$ at the prescribed arity. The
inherited scalar symbol estimates imply, for the critical coefficient,
\[
 |\partial_p^\beta(f_{a,I}^{[4]}-h_I)(p)|
       \le C_{I,\beta}a^2|p|^{-2-|\beta|},
 \qquad C_N\lambda<|p|<c/a.
\]
The fourth-order scalar improvement permits a better interior estimate,
but the displayed conservative second-order estimate suffices. Its integral
is bounded by $Ca^2\int^{c/a}r\,dr\le C$. This also controls the finite
normalized sums by the inherited shell counting bound. On the upper region,
rescale $p=\theta/a$. Four external weighted derivatives supply $a^4$;
the loop count costs $a^{-4}$. All completed symbols and their fixed
rescaled derivatives are bounded there, so that region contributes $O(1)$.
This includes derivatives of the upper profiles and all surviving upper
interactions.

The remaining homogeneous annular integral is
$\mathfrak r_I\log(1/(a\lambda))+O(1)$. A nonround upper boundary changes
only the angular finite constant, not this logarithmic coefficient.
To compare finite sums with the integral, insert a smooth dyadic partition.
At radius $R$ the Poisson-summation error for a degree-minus-four shell is
at most $C_b(RL_{\min})^{-b}$ after $b>4$ loop derivatives. Summing the
shells from $\lambda$ upwards is bounded by
$C_b(\lambda L_{\min})^{-b}$. The last periodic regions were already
bounded without assigning them a continuum value. Hence the residue is
independent of periods and the remainder has the asserted uniform bound.

On an ordinary central region the trace words are precisely the symbol
expansion of the formal logarithm in \eqref{eq:v66-residue-target}.
The integrated degree-minus-four coefficient is its residue. Equivalently,
differentiate with respect to the metric amplitudes, use cyclicity of the
residue to obtain $\operatorname{res}(Q_0^{-1}\delta Q_0)$, and integrate
back from $q=0$. All these operations are finite at each prescribed metric
arity. This identifies the same logarithm without assigning a value to the
unsubtracted infinite determinant.
\end{proof}

\begin{remark}[What the bounded remainder does not assert]
\label{rem:v66-finite-not-fixed}
$B_{a,L,I}$ is actual finite matching data. Lemma~\ref{lem:v66-native-log}
does not set it to zero, identify it with a covariant subtraction, or prove
that the complete critical Slavnov target has a primitive. Power-divergent
weights below four remain in their previous local assignment. The theorem
concerns a populated ghost-number-zero logarithm, not a ghost-number-one
anomaly coefficient.
\end{remark}

\section{Density matching and the evaluated curvature logarithm}
\label{sec:v66-density-heat}

The ordinary scalar Hessian $Q_0$ acts in the flat coefficient pairing.
The geometric scalar operator $\mathscr P_g=-\Delta_g+z$ instead acts in
$L^2(\rho\,dx)$. Define its flat-pairing representative
\begin{equation}
 \widehat{\mathscr P}_g=\rho^{1/2}\mathscr P_g\rho^{-1/2},
 \qquad Q_0=\rho^{1/2}\widehat{\mathscr P}_g\rho^{1/2}.
 \label{eq:v66-density-congruence}
\end{equation}
These operators are not declared equal.

\begin{proposition}[The density congruence changes no integrated logarithmic residue]
\label{prop:v66-density-residue}
For the same positive metric germ,
\begin{equation}
 \operatorname{res}\log Q_0
 =\operatorname{res}\log\widehat{\mathscr P}_g
 =\operatorname{res}\log\mathscr P_g.
 \label{eq:v66-density-residue}
\end{equation}
This equality concerns the integrated logarithmic coefficient only. It is
not equality of finite determinants, measures or pointwise residue densities.
\end{proposition}
\begin{proof}
Let $f=\log\rho$ and
$Q_t=e^{tf/2}\widehat{\mathscr P}_ge^{tf/2}$. The residue trace and its
logarithmic differentiation rule \cite[Theorems 1.7--1.8]{V66Scott} give
\[
 \frac d{dt}\operatorname{res}\log Q_t
 =\operatorname{res}(Q_t^{-1}\dot Q_t)
 =\tfrac12\operatorname{res}(Q_t^{-1}fQ_t+f)
 =\operatorname{res}f=0.
\]
Multiplication by $f$ has no homogeneous symbol of degree minus four.
Similarity invariance gives the second equality. Work first at $z>0$;
finite-rank removal of any zero eigenspace changes no residue. The local
coefficient is polynomial in $z$, so the same result holds at its mass
Taylor origin. No actual retained zero mode is inverted. Cyclicity used
here is integrated cyclicity. With an additional independent source it
would produce commutator/source-derivative terms, which are not discarded.
\end{proof}

\paragraph{Established heat input.}
With curvature convention positive on the round sphere, the scalar
Laplace-type heat coefficient on a closed four-manifold is
\begin{equation}
 a_4(\mathscr P_g)
 =\frac1{(4\pi)^2}\int\rho\left[
 \frac{z^2}{2}-\frac{zR}{6}+\frac{R^2}{72}
 -\frac{R_{ab}R^{ab}}{180}
 +\frac{R_{abcd}R^{abcd}}{180}+\frac{\Delta_gR}{30}\right].
 \label{eq:v66-a4}
\end{equation}
This is the specialization $E=-z$, $\Omega=0$ of the established formula
\cite[Eq.~(4.28)]{V66Vassilevich}, together with
$\operatorname{res}\log\mathscr P_g=-2a_4(\mathscr P_g)$
\cite[Theorem 1.8 and Eq.~(1.22)]{V66Scott}. The zero-eigenspace correction
in that theorem is included; it is why the local identity remains valid
at $z=0$. These are literature inputs, not conclusions of the finite
checker.

Let
\begin{equation}
 \mathscr L_H[g,z]=\int\rho\left[
 \frac{z^2}{2}-\frac{zR}{6}+\frac{R^2}{72}
 -\frac{|\operatorname{Ric}|^2}{180}
 +\frac{|\operatorname{Riem}|^2}{180}\right],\quad
 \Delta\mathscr L_H=\mathscr L_H[g,z]-\mathscr L_H[I,z].
 \label{eq:v66-LH}
\end{equation}
The integrated Laplacian in \eqref{eq:v66-a4} is zero on the periodic
manifold or for compactly supported local variations. Its pointwise
coefficient is recorded, not assigned zero for independently sourced
composites.

\begin{theorem}[Native scalar-vacuum critical logarithm]
\label{thm:v66-native-logarithm}
For every fixed metric arity $N\le5$,
\begin{equation}
 \boxed{\mathsf J_N\mathsf P_4\Gamma^H_{a,1}
 =-\frac{n_H}{16\pi^2}\ell_a\,
       \mathsf J_N\Delta\mathscr L_H+\mathscr B_{a,L,N}.}
 \label{eq:v66-main}
\end{equation}
In the normalized local tensor coefficient basis,
$\|\mathscr B_{a,L,N}\|_{4,\lambda}\le C_Nn_H$ on
\eqref{eq:v66-domain}. Here this norm means division of the local polynomial
by $\lambda^4$ after $K=\lambda\kappa$, $z=\lambda^2\zeta$; at exact weight
four it is simply the sum of the normalized coefficient magnitudes.
The sign is the sign of $+n_H\Tr\log Q/2$ in the effective action.
The counterterm that cancels this logarithmic block has the opposite sign.
\end{theorem}
\begin{proof}
Apply Lemma~\ref{lem:v66-native-log} to every coefficient of the finite
metric jet. Proposition~\ref{prop:v66-density-residue} converts its residue
to \eqref{eq:v66-a4}. Multiplication by $n_H/2$ cancels the factor two in
$\operatorname{res}\log\mathscr P_g=-2a_4$. Subtraction of the flat
reference gives \eqref{eq:v66-main}. There are finitely many coefficients
at each $N$, so their proven uniform bounds sum. All coframe derivatives
are taken by substituting the same $g=G(q)$ in \eqref{eq:v66-LH}.
\end{proof}

The flat mass check fixes both sign and normalization: the $z^2$ term of
$\frac{n_H}{2}\int_p\log(p^2+z)$ is
$-\frac{n_Hz^2}{4}\int_p p^{-4}$, whose logarithm is
$-n_Hz^2\log\Lambda/(32\pi^2)$, as in \eqref{eq:v66-main}.
No factor $\chi^{-1}$ is present, because this loop contains no metric
propagator.

With $E_4=|\operatorname{Riem}|^2-4|\operatorname{Ric}|^2+R^2$ and Weyl
tensor $C_g$, the same density has the useful equivalent forms
\begin{align}
 a_H(g,z)
 &\equiv \frac{z^2}{2}-\frac{zR}{6}
       +\frac{R^2}{120}+\frac{|\operatorname{Ric}|^2}{60}
       +\frac{E_4}{180},\label{eq:v66-Ricci-form}\\
 &\equiv \frac{z^2}{2}-\frac{zR}{6}
       +\frac{R^2}{72}+\frac{|C_g|^2}{120}-\frac{E_4}{360}.
 \label{eq:v66-Weyl-form}
\end{align}
Here $a_H$ is defined by $\mathscr L_H=\int\rho a_H$. The Euler term is retained as a topological density.
Its integrated local metric derivatives vanish on the stated periodic
chart; no statement about boundary observables or global topology is made.

\section{Independent bilinear and conformal-source checks}
\label{sec:v66-bilinear}

For a metric-difference direction $h$ and Fourier momentum $k$, put
\begin{equation}
 \begin{split}
 \mathcal R(h)&=k^2\operatorname{tr}h-k^Thk,\\
 \mathcal R_{ab}(h)&=\tfrac12\left[
 k^2h_{ab}+k_ak_b\operatorname{tr}h-k_a(hk)_b-k_b(hk)_a\right].
 \end{split}
 \label{eq:v66-linear-Ricci}
\end{equation}
These are the actual linearized curvatures, not a transverse--traceless
projection.

\begin{proposition}[Complete massless metric Hessian]
\label{prop:v66-bilinear}
The degree-four bilinear Hessian of \eqref{eq:v66-LH} is
\begin{equation}
 \boxed{\mathscr H_H(h,u;k)
 =\frac1{60}\mathcal R(h)\mathcal R(u)
  +\frac1{30}\mathcal R_{ab}(h)\mathcal R^{ab}(u).}
 \label{eq:v66-HH}
\end{equation}
It annihilates a free gauge leg $h=kv^T+vk^T$. For $h=u=I$ it is
$(k^2)^2/4$. For $k=te_0$ and $h=u=e_1e_2^T+e_2e_1^T$ it is $t^4/60$.
The logarithmic scalar contribution to the metric two-point coefficient is
$-n_H\ell_a\mathscr H_H/(16\pi^2)$.
\end{proposition}
\begin{proof}
Use \eqref{eq:v66-Ricci-form}. At $z=0$ the background curvatures vanish,
so only the products of linearized curvatures contribute at quadratic
order. The Euler integral has zero second variation. Polarization gives
\eqref{eq:v66-HH}. Substitution gives the stated gauge and tensor tests.
\end{proof}

This Hessian can also be obtained without using the heat coefficient.
Take $k=xe_0$, write a unit loop direction as $n\in S^3$, and set
\[
 v_0=\tfrac12\operatorname{tr}h-n^Thn,\qquad
 v_1=\tfrac12\operatorname{tr}h\,n_0-(hn)_0.
\]
The native limiting scalar vertex is $v_0+xv_1$ after its radial factor
is removed. The fourth Taylor coefficient of the two-vertex word is
\begin{equation}
 A_h(n)=v_0^2(16n_0^4-12n_0^2+1)
       +2v_0v_1(-8n_0^3+4n_0)
       +v_1^2(4n_0^2-1).
 \label{eq:v66-angular-polynomial}
\end{equation}
Using
\begin{equation}
 \langle n_0^{2a_0}\cdots n_3^{2a_3}\rangle_{S^3}
 =\frac{\prod_i(2a_i-1)!!}
        {4\cdot6\cdots(4+2\sum_i a_i-2)}
 \label{eq:v66-angular-moments}
\end{equation}
(with odd moments zero), direct polynomial contraction gives
$\langle A_h\rangle=\mathscr H_H(h,h;e_0)$.
The companion checker verifies every one of the 55 quadratic tensor
coefficients, not only the conformal and shear examples. The one-vertex
ordinary massless contact has insufficient external derivative degree to
supply a fourth-order logarithm, although it remains essential at lower
local weights and finite matching orders. Thus this is an independent
check of the entire massless bilinear block and its sign.


\begin{proposition}[Mass terms and the required tadpole partner]
\label{prop:v66-massive-bilinear}
Let $\mathcal E_1(u)=\mathcal R_{ab}(u)-\frac12\delta_{ab}\mathcal R(u)$.
The complete critical Hessian in metric-difference coordinates is
\begin{equation}
 \boxed{\begin{aligned}
 \mathscr H_H(h,u;k,z)={}&\mathscr H_H(h,u;k,0)
   +\frac z6 h:\mathcal E_1(u)\\
 &+z^2\left[\frac{\operatorname{tr}h\operatorname{tr}u}{8}
                          -\frac{\operatorname{tr}(hu)}4\right].
 \end{aligned}}
 \label{eq:v66-massive-Hessian}
\end{equation}
The corresponding linear metric coefficient is
$\frac{z^2}{4}\int\operatorname{tr}h$. The $z^2$ Hessian is not by itself
annihilated by a free gauge leg. Its contribution to the one-ghost/one-metric
identity cancels against the nonlinear variation of this linear coefficient.
\end{proposition}
\begin{proof}
At flat metric,
$D^2\int\rho R[h,u]=-\int h:\mathcal E_1(u)$, while
$D^2\rho[h,u]=\operatorname{tr}h\operatorname{tr}u/4
-\operatorname{tr}(hu)/2$. Insert these into \eqref{eq:v66-LH}.
For the last assertion take $h$ at momentum $-k$, $c$ at $k$.
The free-gauge leg is $R_0c=i(kc^T+ck^T)$, and the constant Fourier
coefficient of the nonlinear Lie term is
\[
 R_1(h,c)=i\left[-(k\cdot c)h+k(hc)^T+(hc)k^T\right].
\]
The $z^2$ part of
$\mathscr H_H(h,R_0c;k,z)+(z^2/4)\operatorname{tr}R_1(h,c)$ is
exactly zero. Thus the tadpole is required even though the integrated
linear curvature term is a divergence.

There is also a direct native loop check. At $k=xe_0$ introduce
$y=z/R_p^2$ for a loop of radius $R_p$, and replace the two-vertex word by
\[
 \frac{(v_0+xv_1+y\operatorname{tr}h/2)^2}
      {(1+y)(1+2n_0x+x^2+y)}.
\]
Its $x^2y$ coefficient has average $h:\mathcal E_1(h)/6$.
For the $y^2$ coefficient subtract the actual one-vertex contribution
$n^TW''[h,h]n-\rho''[h,h]$, where
$W''=\rho''I-(\operatorname{tr}h)h+2h^2$.
The average is $(\operatorname{tr}h)^2/8-\operatorname{tr}(h^2)/4$.
The checker verifies all 55 tensor coefficients at each of the two mass
weights, in addition to the 55 massless coefficients. This retains the
coframe-independent density origin of the tadpole; conversion to the
active coframe uses the full chain rule, including $D^2G(q)$.
\end{proof}

\begin{proposition}[Full conformal source curve]
\label{prop:v66-conformal}
On $g=e^{2\sigma}\delta$, modulo a periodic divergence and the unchanged
Euler integral,
\begin{equation}
 \boxed{\Delta\mathscr L_H
 =\int\left[
 \frac12\bigl(\partial^2\sigma+|\partial\sigma|^2\bigr)^2
 -z e^{2\sigma}|\partial\sigma|^2
 +\frac{z^2}{2}(e^{4\sigma}-1)\right].}
 \label{eq:v66-conformal}
\end{equation}
In particular the massless cubic and quartic source contacts are respectively
$\int(\partial^2\sigma)|\partial\sigma|^2$ and
$\frac12\int|\partial\sigma|^4$. They are derivatives of the same action,
not independently chosen finite tensors.
\end{proposition}
\begin{proof}
For this metric $C_g=0$,
$R=-6e^{-2\sigma}(\partial^2\sigma+|\partial\sigma|^2)$ and
$\rho=e^{4\sigma}$. Insert these in \eqref{eq:v66-Weyl-form}.
Integration by parts gives
$\int ze^{2\sigma}\partial^2\sigma
=-2\int ze^{2\sigma}|\partial\sigma|^2$.
This proves \eqref{eq:v66-conformal}. It is an ordinary local source curve,
not a change of the finite Gaussian contour.
\end{proof}

\section{The finite part is matching data, not an invariant replacement}
\label{sec:v66-finite-matching}

A more explicit comparison is useful when the remaining finite local array
is eventually supplied. Extend an ordinary smooth upper-cutoff loop to the
same Brillouin cube, with its cutoff zero near the cube boundary and one
near its centre. Put its own cutoff on every primitive covariance, and
retain exactly the same ordinary vertices, coframe chart and lower mask.
This is a comparison calculation, not an active change of the regulator.
For a critical component $I$ let $F_I^{A}(\theta)$ and $F_I^{B}(\theta)$
be the dimensionless critical integrands for the native and comparison
schemes, with $a=1,z=0$ and the lower mask removed after differentiation.
All geometric and momentum derivatives are taken before extracting these
functions.

\begin{proposition}[An absolutely convergent critical matching quadrature]
\label{prop:v66-finite-matching}
The difference
\begin{equation}
 \mathfrak m_I(A,B)
 =\int_{[-\pi,\pi]^4}\frac{d^4\theta}{(2\pi)^4}
       \bigl(F_I^A(\theta)-F_I^B(\theta)\bigr)
 \label{eq:v66-matching-integral}
\end{equation}
is absolutely convergent. It is the refinement limit of the difference
of the two actual critical local coefficients, with the same lower
reference, at fixed physical periods. It is independent of that lower
reference and those periods. It need not vanish and is not in general
separately diffeomorphism invariant.
\end{proposition}
\begin{proof}
At the origin the leading degree-minus-four functions are equal by
Lemma~\ref{lem:v66-native-log}. The differentiated interior symbol estimates
bound their difference by $C|\theta|^{-2}$; away from the origin it is
bounded on a compact set. Both bounds are integrable. Terms in which a
derivative hits the lower mask lie in $|\theta|=O(a\lambda)$; their
native-minus-ordinary discrepancy is bounded after integration by
$C(a\lambda)^2$. On the rest of a small dimensionless ball the normalized
sums and integrals of the difference are $O(\epsilon^2)$, uniformly in
refinement. Outside that ball the rescaled grid spacing tends to zero,
and ordinary Riemann convergence applies. First remove the cutoff at fixed
$\epsilon$, then let $\epsilon\downarrow0$. This yields
\eqref{eq:v66-matching-integral} and its independence of the lower reference.
The proof determines a finite quadrature, not its numerical value.
\end{proof}

\begin{remark}[Power contacts, topology and independent composite sources]
\label{rem:v66-source-boundary}
The density congruence in \eqref{eq:v66-density-congruence} has a pointwise
finite determinant Jacobian if performed as a finite change of all scalar
coordinates. That is a separate local volume contact; in the hard split it
also transports the reference. Neither operation is performed here.
Equality of logarithmic residues is not permission to remove the
auxiliary density, to add another scalar $\rho$ measure, or to equate the
two finite determinants.

Likewise $\int\rho\Delta R=0$ and the local constancy of $\int\rho E_4$
refer to the integrated action on the stated boundary-free chart. Multiplying
by a new independent source destroys these simplifications and introduces
source-derivative or topological-current contacts. This installment covers
the geometric derivatives of one integrated action. It does not purport to
renormalize all independently sourced curvature composites.
\end{remark}

\section{A Ricci counterterm carries a definite matter partner}
\label{sec:v66-matter-partner}

The scalar-loop logarithm cannot be discarded merely because curvature
squares are redundant in vacuum Einstein gravity. The common action has the
full metric Euler tensor
\begin{equation}
 \mathcal E^{ab}=e^{ab}+\tau^{ab},\qquad e^{ab}=c_EG^{ab},\qquad
 \tau^{ab}=\frac14g^{ab}(X+zU)-\frac12V^{ab},
 \quad X=\operatorname{tr}_g V,\quad U=\phi^T\phi,
 \label{eq:v66-stress}
\end{equation}
with $V^{ab}=\sum_A\nabla^a\phi_A\nabla^b\phi_A$ and
$c_E\ne0$ the signed native normalization. No value for $c_E$ is changed.

\begin{theorem}[Exact full-Euler conversion of the evaluated residue]
\label{thm:v66-matter-partner}
Set
\begin{equation}
 \begin{split}
 F_{H,ab}={}&\frac{z}{6c_E}g_{ab}
  +\frac{1}{120c_E^2}\operatorname{tr}(e-\tau)g_{ab}
  +\frac{1}{60c_E^2}(e-\tau)_{ab},\\
 \mathcal M_H={}&\frac{z^2}{2}
  -\frac{zX}{12c_E}-\frac{z^2U}{6c_E}
  +\frac1{c_E^2}\left[
    \frac{X^2}{480}+\frac{\operatorname{tr}(V^2)}{240}
    +\frac{zUX}{80}+\frac{z^2U^2}{80}\right].
 \end{split}
 \label{eq:v66-FM}
\end{equation}
Then the ordinary local functional satisfies the exact off-shell identity
\begin{equation}
 \boxed{\mathscr L_H
 =\int\rho\left(\mathcal E^{ab}F_{H,ab}+\mathcal M_H
                       +\frac{E_4}{180}\right).}
 \label{eq:v66-matter-conversion}
\end{equation}
The equality uses no field equation. The previously recorded integrated
$\Delta R$ is a separate divergence. The metric-field generator
$Q_H=\int\eta_g^{ab}F_{H,ab}$ is local and natural; for the ordinary ungauge-fixed minimal
Einstein--scalar master representative,
\begin{equation}
 (S_0,Q_H)=\int\rho\,\mathcal E^{ab}F_{H,ab}\pmod{d_x}.
 \label{eq:v66-BV-conversion}
\end{equation}
Its coframe realization uses the same pointwise inverse $DG(q)^{-1}$.
The identity is stated before gauge fixing; adjoining the nonminimal
doublets and applying the inherited gauge-fixing canonical transformation
transports both sides together, rather than dropping a gauge-fixing
contribution from a mixed representation.
\end{theorem}
\begin{proof}
In four dimensions $\operatorname{tr}e=-c_ER$ and
$|\operatorname{Ric}|^2=|e|^2/c_E^2$. Thus, after keeping the Euler term,
the integrand in \eqref{eq:v66-Ricci-form} is
\[
 \frac{z^2}{2}+\frac{z}{6c_E}\operatorname{tr}e
  +\frac1{c_E^2}\left[
       \frac{(\operatorname{tr}e)^2}{120}+\frac{|e|^2}{60}\right].
\]
Subtract this quadratic polynomial evaluated at $e=-\tau$. Difference of
squares factors it exactly as $(e+\tau):F_H$. The stress has
\[
 \operatorname{tr}\tau=X/2+zU,\qquad
 |\tau|^2=\tfrac14\operatorname{tr}(V^2)+\tfrac14zUX+\tfrac14z^2U^2.
\]
Substitution yields \eqref{eq:v66-FM}. Since $F_H$ is a natural metric
tensor, the longitudinal variation of $\eta_g:F_H$ is a density divergence;
its Euler variation uses $e+\tau$. This proves
\eqref{eq:v66-BV-conversion}. In particular a vacuum-Euler replacement
alone would miss the entire difference involving $\tau$.
\end{proof}

At zero mass the evaluated matter density is
\begin{equation}
 \boxed{\mathcal M_H\big|_{z=0}
 =\frac{X^2+2\operatorname{tr}(V^2)}{480c_E^2}.}
 \label{eq:v66-massless-matter}
\end{equation}
For real Euclidean gradients $V\succeq0$ this is strictly positive unless
all gradients vanish; for one scalar it equals $X^2/(160c_E^2)$.
Consequently the scalar-loop logarithm, when represented with the Ricci
terms removed by this complete Euler conversion, includes
\[
 -\frac{n_H\ell_a}{16\pi^2}
 \int\rho\,\frac{X^2+2\operatorname{tr}(V^2)}{480c_E^2}.
\]
This is a populated matter partner of this loop contribution. We do not
claim that it exhausts an independent operator basis of the full coupled
BV quotient. Other loops and finite physical normalizations remain.
It is distinct from V38's anisotropic four-scalar metric-loop packet.

\section{Insertion into the common action and normalization ledger}
\label{sec:v66-insertion}

If the scalar-loop logarithm has not already been included in a supplied
common counterterm, its required addition has the sign
\begin{equation}
 C_{H,1}^{\log}
 =\frac{n_H}{16\pi^2}\ell_a\,
       \mathsf J_N\Delta\mathscr L_H.
 \label{eq:v66-counterterm}
\end{equation}
This is an identified component of $C_1$, not an additional copy to add
when $-\mathsf P\Gamma_1$ already contains it. It contains no scalar
antifield, and $s_\infty\mathscr L_H=0$ modulo a divergence for the
ordinary diffeomorphism differential. This homogeneous identity does
not solve the inhomogeneous finite critical reference equation.

\begin{proposition}[Populated metric-channel logarithm and its common insertion]
\label{prop:v66-insertion}
The metric, two-metric and higher metric vertices of
\eqref{eq:v66-counterterm} are derivatives of this single local functional.
They pass V64's Hessian integrability test without independent adjustment.
Their order-two hard insertion, on V64's bridge-excluded action panels, is
\begin{equation}
 \boxed{\mathcal I_a[C_{H,1}^{\log}]
 =\frac{n_H\ell_a}{32\pi^2}
   \Tr\left(G_{g,a}D_q^2\mathsf J_N\Delta\mathscr L_H\right).}
 \label{eq:v66-hard-insertion}
\end{equation}
It is the scalar-loop contribution to the proper metric two-point
counterterm channel, with its scalar-vacuum coframe contacts included.
No second factor $n_H$ or $1/2$ is inserted when closing that channel.
\end{proposition}
\begin{proof}
Mixed derivatives of a single polynomial/analytic metric germ commute.
V64's formula $\frac12\operatorname{STr}(\mathbb G D_\Phi^2C_1)$ has only
the metric block for the displayed representation of this counterterm;
its coefficient is \eqref{eq:v66-hard-insertion}. Using the Euler-reduced
representation \eqref{eq:v66-matter-conversion} instead requires the
complete field/source redefinition and its induced matter and measure
partners. The two presentations cannot be mixed graph by graph without
those partners. V65 fixes the general next-order canonical bookkeeping.
\end{proof}

The new local action also has fixed-window multilinear bounds. A term with
$d$ derivatives and $z^j$, $d+2j=4$, has a rooted symbol/kernel bound
$C_{N,b}\,n_H|\ell_a|(\mu+m_H)^4$ at fixed arity and kernel moment,
with one source in $L^1$ and the others in $L^\infty$. This follows by
polarizing the finite metric jet and integrating the fixed smooth windows;
there is no total-volume multiplier. The logarithm is retained. These are
local counterterm bounds, not an additional vanishing-cutoff estimate.

In particular V34's finite weight-two scalar normalization is untouched.
V38 and V65's packet is also untouched. No value is assigned to the bounded
finite tensor $\mathscr B_{a,L,N}$, to the mixed graviton--scalar logarithm,
to a ghost or fermion loop, or to the finite critical source-breaking
primitive. The common array is now populated in this scalar-vacuum
logarithmic block, not in all of its blocks.

\section{Status: the scalar logarithm is known, the finite critical descent is not}
\label{sec:v66-ledgers}

\subsection*{Proved}
The actual completed scalar determinant has the logarithm
\eqref{eq:v66-main} at every retained metric arity through five, with a
cutoff/volume-uniform bounded critical remainder. Its full curvature and
mass polynomial is evaluated by the stated residue/heat inputs after the
native matching and density arguments. The full 165-coefficient bilinear check, including the two mass
weights and its tadpole, is independent. The finite completion comparison has the convergent quadrature
\eqref{eq:v66-matching-integral}. The complete off-shell Euler conversion
has the evaluated scalar-quartic/mass partner \eqref{eq:v66-FM}; the common
metric-channel insertion is fixed by \eqref{eq:v66-hard-insertion}.

\subsection*{Inherited}
V27 supplies the completed scalar action, floor, coframe vertices and
interior/upper symbol estimates. V60 supplies the scalar determinant and
its master/reference role. V64 supplies the common-counterterm insertion
and bridge hypotheses. V65 supplies the full next-order bookkeeping for
an optional canonical conversion. Their analytic claims are treated as
inherited inputs, not reverified by the current finite algebraic tests.

\subsection*{Literature input}
Equation~\eqref{eq:v66-a4} and the integrated logarithmic residue trace
identities are established theorems \cite{V66Vassilevich,V66Scott}. The
present work does not prove the general heat-kernel theorem. It verifies
its application to the actual density-normalized scalar operator and
matches that residue to the native finite cutoff. Boundary and global
topological extensions of those theorems are not claimed.

\subsection*{Conditional}
Counterterm insertion in a renormalized two-loop action requires the same
common $C_1$ coefficients, source convention and finite anchors in all
channels. A full Euler-reduced scheme must transport the scalar and metric
observables and the measure; \eqref{eq:v66-matter-conversion} is not a
license to drop them. The finite matching quadratures have not been
numerically evaluated. The logarithm alone does not fix their values.

\subsection*{Open}
The finite inhomogeneous critical source primitive and the remaining
mixed, ghost, internal-gauge and chiral contributions to the common
one-loop array remain open. The complete two-loop master/reference
coefficient and its interacting uniform estimates remain open. The present
result is not an anomaly-cancellation theorem, a covariant finite regulator,
or an invariant many-shell Einstein--SM class.

\begin{center}\small
\begin{tabular}{>{\raggedright\arraybackslash}p{.23\textwidth}>{\raggedright\arraybackslash}p{.69\textwidth}}
\toprule
Variable & Scope in V66\\\midrule
Loop/source order & One closed real-scalar loop; integrated metric action and
its jets through arity five; no independent curvature-composite source.\\
Derivative/mass order & Critical weight four, polynomial through $z^2$;
positive lower reference retained, including at $z=0$.\\
Native/ordinary & Finite symbols determine the same integrated logarithmic
residue; finite and power-divergent pieces are not equated.\\
Volume/cutoff & Bounded critical remainder at fixed profiles and
$\lambda L_\nu\ge1$; logarithmic counterterm retained; no new full forest
or higher-loop estimate.\\
Symmetry & Ordinary action logarithm is diffeomorphism invariant; finite
reference breaking and density contacts remain separate.\\
Verification & Exact heat-polynomial specialization, 165 bilinear loop
coefficients, Euler conversion, source-curve and density tests; no numerical
native integral, full cohomology matrix or proof-assistant claim.\\\bottomrule
\end{tabular}
\end{center}

\begin{thebibliography}{99}
\bibitem{V66Vassilevich} D.~V.~Vassilevich,
Heat kernel expansion: user's manual,
\emph{Phys. Rept.} \textbf{388} (2003), 279--360;
arXiv:hep-th/0306138, especially Eqs.~(2.2), (4.26)--(4.28).
\bibitem{V66Scott} S.~Scott,
The residue determinant,
\emph{Comm. Partial Differential Equations} \textbf{30} (2005), 483--507;
arXiv:math/0406268, especially Theorems~1.7--1.8 and Eq.~(1.22).
\end{thebibliography}

\section*{Append-only continuation marker: V66}
\addcontentsline{toc}{section}{Append-only continuation marker: V66}
\label{sec:v66-continuation-marker}
\textbf{End of V66.} Preserve Sections 1--512. The native scalar-vacuum
logarithmic block is now evaluated as one common local action, with its
density distinction and coupled-matter Euler conversion. The bounded
finite matching arrays and the inhomogeneous critical source primitive
remain distinct unresolved inputs. No active normalization changes.


% ============================================================================
% V67: a locally sourced scalar determinant, rather than an unsourced residue
% ============================================================================
\clearpage
\section{V67: companion audit and a source-dependent logarithmic block}
\label{sec:v67-start}

The three newly supplied companions do not provide the missing critical
quantum primitive. They do sharpen which source operation must be estimated.
We therefore continue V66 at its explicitly excluded interface: local
composite sources, not merely metric differentiation of an integrated
unsourced action. The new calculation evaluates a matrix-valued kinetic
source together with the quadratic scalar source. It includes the local
connection generated by a noncommuting kinetic matrix. This connection is
an output of the source normalization, not an additional physical gauge
field. At zero added source the action, measure, covariance and all active
normalizations are exactly those of V66.

\paragraph{Targeted companion audit.}
The following entries refer to the terminal developments and the definitions
and preceding results those developments consume. They are not independent
revalidations of the entire cumulative companions.
\begin{center}\small
\begin{tabular}{>{\raggedright\arraybackslash}p{.24\textwidth}>{\raggedright\arraybackslash}p{.68\textwidth}}
\toprule
Companion & Exact input and limit of the transfer\\\midrule
SM source selection V52~\cite{V67SM52} &
\srcid{sel52:thm:global} proves the one-call maxima
$(3-2\sqrt2)/3$ and $(2-\sqrt3)/3$ in its two specified preparation--steering
classes. \srcid{sel52:prop:bank} exhibits a spanning invariant preparation
bank with zero exact-ray capacity: signed reconstruction is not positive
preparation. No quantum field source is certified executable by that
operator-algebraic result.\\
Native stationarity V54~\cite{V67Station54} &
\srcid{v54:thm:descriptor} retains the rank-two derivative-source map
$ZG_0^{-1}Z^T$. \srcid{v54:thm:lift} controls the restored coefficient
problem with a $C^3$ covector norm and two-scale preparation; the nonlinear
original remainder is still unestimated. This motivates retaining source
derivatives here, but supplies no spatial quantum bound.\\
Exact-action bridge V45~\cite{V67Bridge45} &
\srcid{v45:weak-rank} resolves the weak rank $24$ and the five zero
modes; \srcid{v45:canonical-theorem} identifies four intermediate
oscillatory roots. \srcid{v45:frequency-formula} still contains the
unevaluated next Poisson coefficient. These are fixed-$N=3$ amplitude
statements, not positivity or ultraviolet estimates for our Gaussian.\\\bottomrule
\end{tabular}
\end{center}
The amplitude denoted $a$ in the last two companions is not the mesh $a$
in this lineage. In particular their positive intermediate pencil is not
used to replace the signed gravitational covariance.

\paragraph{Nonduplication and provenance.}
V9 already defines a matrix-valued scalar-potential source and estimates
its Hodge module. V54--V56 already distinguish signed counterterm synthesis
from a positive physical instrument. V66 already evaluates the unsourced
scalar logarithm. None is claimed again. The new source below extends the
\emph{active V27 completion}, rather than silently switching back to V9's
unimproved Hodge regulator. Its source derivatives are a specified formal
observable family on the existing field carrier. Their independent physical
preparation or readout is not inferred from the companion source theorem.
The general heat coefficient and residue trace remain explicit literature
inputs \cite{V67Vassilevich,V67Scott}; the source normal form, its native
matching, and the displayed coefficient evaluations are proved below.

\section{One declared finite action for kinetic and potential sources}
\label{sec:v67-finite-family}

Work on the same four odd periodic grids and coefficient measure as V27.
Let $\Sigma(x)$ and $Y(x)$ be real symmetric $n_H\times n_H$ sources,
$n_H=4$, with weights zero and two respectively. Set
\begin{equation}
 B=e^\Sigma,\qquad Z=B^2=e^{2\Sigma},\qquad
 \mathcal M=B(zI+Y)B,\qquad z=m_H^2.
 \label{eq:v67-source-matrices}
\end{equation}
$B$ is a pointwise positive matrix on the real source chart. Its inverse
is a pointwise analytic matrix, not an inverse scalar field or propagator.
Source Taylor coefficients extend to the same complexified metric chart.

For $\alpha=n,b$, define $S_{H,a}^{\alpha}(q,\Sigma,Y,\phi)$ by the actual
scalar bilinear of V27 with the replacements
\begin{equation}
 W^{\mu\nu}(q)I\ \longmapsto\ W^{\mu\nu}(q)Z,
 \qquad z\rho(q)I\ \longmapsto\rho(q)\mathcal M.
 \label{eq:v67-native-replacements}
\end{equation}
Explicitly, the $(A,B)$ scalar Fourier coefficient is
\[
 \frac12\phi_A(p)\phi_B(r)
 \left\{\widehat{\rho\mathcal M}_{AB}(-p-r)
  -\widehat{W^{\mu\nu}Z}_{AB}(-p-r)
                  D^{\alpha}_{\mu\nu}(p,r)\right\}.
\]
Here $D^n$ means V27's corrected $\widehat D^n$, and $D^b$ its periodic
bare tensor. All products at the root and all factors on the two scalar
legs retain that convention. Write
$I^\alpha=S_{H,a}^\alpha(q,\Sigma,Y,\phi)-S_{H,a}^\alpha(0,0,0,\phi)$.
With the same filter $S=s(aD)$ as V27, define
\begin{equation}
 \boxed{\begin{aligned}
 S_{H,a}^{\rm src}
 ={}&\frac12\langle\phi,\widetilde P_a\phi\rangle
       +I^b(q,\Sigma,Y,\phi)\\
 &+I^n(Sq,S\Sigma,SY,S\phi)-I^b(Sq,S\Sigma,SY,S\phi).
 \end{aligned}}
 \label{eq:v67-finite-action}
\end{equation}
This is an explicit source completion. It is exactly
\eqref{eq:v27-completed-Higgs} when $\Sigma=Y=0$. All upper interactions
and the temporal continuation are retained. The reference remains
$C_a=\nu_\lambda\widetilde P_a^{-1}I_{n_H}$ and is independent of the
added sources. Let $Q_a^{\rm src}$ be the Hessian of
\eqref{eq:v67-finite-action} and define
\begin{equation}
 \Delta\Gamma_{a,1}
 =\frac12\Tr\log(I+C_a[Q_a^{\rm src}-\widetilde P_a])
  -\frac12\Tr\log(I+C_a[Q_a(q)-\widetilde P_a]).
 \label{eq:v67-delta-det}
\end{equation}
The trace already includes the real flavour space. There is no additional
factor $n_H$ outside it. The common $\hbar$ is removed. In particular the
normalization at zero sources does not remove source-only terms.

\begin{lemma}[Native source jets and the unchanged reference]
\label{lem:v67-source-symbols}
At any fixed total metric/source arity $N\le5$, the differentiated vertices
of \eqref{eq:v67-finite-action} have the same order-two estimates as V27.
A potential-source argument removes two powers of hard momentum. On the
common interior plateau their discrepancy from the ordinary vertices is
$O(a^4R^6)$ before assigning these potential weights, with one factor
$R^{-1}$ per momentum derivative. On the periodic upper regions they
have the corresponding rescaled bounds. The ordinary operator is
\begin{equation}
 Q_0^{\rm src}
 =-\partial_\mu(\rho g^{\mu\nu}B^2\partial_\nu)
             +\rho B(zI+Y)B.
 \label{eq:v67-ordinary-Q}
\end{equation}
At fixed positive lower scale a common sufficiently small hard-source
chart exists; no claim about an unintegrated zero mode at $z=0$ is needed.
\end{lemma}
\begin{proof}
Each derivative of the matrix exponentials is a finite sum of ordered
products of its polarization matrices. At a fixed arity these products
have fixed norm bounds. They multiply exactly the scalar difference
tensors in \eqref{eq:v27-Dcorrected}, not a new momentum symbol.
The interior Taylor estimate therefore follows from the same corrected
sine/hyperbolic-sine expansion. Each potential vertex is linear in $Y$
and has no hard derivative. The completed upper symbols are finite
products of the same smooth profiles and bounded root coefficients;
rescaling gives their derivative bounds. Root products may shift the
momentum on a scalar line but do not change its differential order.
The lower hard support has $\widetilde P_a\ge c\lambda^2$ there.
Cauchy--Schwarz in the quadratic forms bounds its relative perturbation
by $C_N(\|q\|_{\rm W}+\|\Sigma\|_{\rm W}+\lambda^{-2}\|Y\|_{\rm W})$
on a fixed small chart. Neumann inversion then gives the hard analytic
germ, with all retained covariance zeros unchanged. Finally, setting the
sources to zero in the defining action gives V27 term by term.
\end{proof}

Under a constant orthogonal change of flavour frame,
$\Sigma,Y\mapsto O\Sigma O^T, OYO^T$, and $\phi\mapsto O\phi$.
The finite family is equivariant, not invariant with an arbitrary
nonscalar source held fixed. No local internal-gauge assertion is made
for this zero-gauge-field source family.

\section{The source normal form has a determined matrix connection}
\label{sec:v67-normal-form}

The ordinary computation is performed in the geometric scalar pairing
$L^2(\rho\,dx)$, solely to evaluate a logarithmic residue. It does not
change the finite scalar coefficient measure. Define
\begin{equation}
 L_\mu=B^{-1}\partial_\mu B,\quad
 R_\mu=(\partial_\mu B)B^{-1},\quad
 H_\mu=\frac{L_\mu+R_\mu}{2},\quad
 \omega_\mu=\frac{L_\mu-R_\mu}{2}.
 \label{eq:v67-MC}
\end{equation}
Thus $H_\mu^T=H_\mu$ and $\omega_\mu^T=-\omega_\mu$.
Let $\mathcal D_\mu=\nabla_\mu+\omega_\mu$ on flavour vectors and use its
commutator action on matrix tensors. Put
\begin{equation}
 \mathcal V_\Sigma=\mathcal D^\mu H_\mu+H_\mu H^\mu,
 \qquad U=Y+\mathcal V_\Sigma.
 \label{eq:v67-potential}
\end{equation}

\begin{theorem}[Exact source-dependent Laplace normal form]
\label{thm:v67-normal-form}
The operator \eqref{eq:v67-ordinary-Q} satisfies
\begin{equation}
 \boxed{B^{-1}\rho^{-1}Q_0^{\rm src}B^{-1}
           =-\mathcal D^2+zI+U.}
 \label{eq:v67-Laplace}
\end{equation}
Its connection curvature is exactly
\begin{equation}
 \boxed{\Omega_{\mu\nu}
   =\partial_\mu\omega_\nu-\partial_\nu\omega_\mu
      +[\omega_\mu,\omega_\nu]
   =-[H_\mu,H_\nu].}
 \label{eq:v67-curvature}
\end{equation}
The potential $U$ is symmetric. A noncommuting source is not, in general,
equivalent to several independently normalized scalar sources.
\end{theorem}
\begin{proof}
Write $\psi=B\phi$. Then $B\nabla_\mu\phi=
\nabla_\mu\psi-R_\mu\psi$. The geometric adjoint of this first-order
operator gives
\[
 (\nabla-R)^\dagger(\nabla-R)
 =-\Delta_g+(R_\mu-L_\mu)\nabla^\mu
                 +\nabla^\mu R_\mu+L_\mu R^\mu.
\]
Substitute $L=H+\omega$, $R=H-\omega$, and complete
$-(\nabla+\omega)^2$. The remaining multiplier is
$\nabla^\mu H_\mu+[\omega^\mu,H_\mu]+H_\mu H^\mu$, giving
\eqref{eq:v67-Laplace}. Symmetry follows either from the displayed
adjoint product or from $H^T=H$, $\omega^T=-\omega$.
The left and right Maurer forms obey
$dL=-L\wedge L$ and $dR=R\wedge R$. Hence
\[
 d\omega+\omega\wedge\omega
 =-\tfrac14(L+R)\wedge(L+R)=-H\wedge H.
\]
This proves \eqref{eq:v67-curvature} without diagonalizing $B$ or
assuming that its values commute at different points.
\end{proof}

\begin{proposition}[A fixed source-jet bound]
\label{prop:v67-source-bound}
In an orthonormal metric frame assume $\|\Sigma\|_{\op}\le r$ and put
$p_1=\sum_\mu\|\nabla_\mu\Sigma\|_{\op}$,
$p_2=\sum_{\mu,\nu}\|\nabla_\mu\nabla_\nu\Sigma\|_{\op}$. Then
\begin{equation}
 \|\mathcal V_\Sigma\|_{\op}
 \le e^{2r}p_2+(e^{2r}+4e^{4r})p_1^2,
 \qquad
 \|\Omega_{\mu\nu}\|_{\op}
 \le2e^{4r}\|\nabla_\mu\Sigma\|_{\op}
                 \|\nabla_\nu\Sigma\|_{\op}.
 \label{eq:v67-pointwise-bounds}
\end{equation}
These are bounds in a fixed metric/source-jet chart, not in the source
amplitude norm alone.
\end{proposition}
\begin{proof}
The differentiated exponential is
$\partial B=\int_0^1e^{(1-t)\Sigma}(\partial\Sigma)e^{t\Sigma}dt$.
Together with $\|B^{-1}\|\le e^r$ this bounds $L,R,H,\omega$ by
$e^{2r}\|\partial\Sigma\|$. The second differentiated exponential
has two ordered integrals whose total simplex volume is one; hence its
norm is at most $e^r(\|\nabla^2\Sigma\|+\|\nabla\Sigma\|^2)$.
Differentiating $B^{-1}$ adds an $e^{4r}$ product. In
$\mathcal D\cdot H+H^2$ the commutator and square cost at most three
further such products. Summing gives the first bound. The second is
\eqref{eq:v67-curvature} and $\|[A,B]\|\le2\|A\|\|B\|$.
\end{proof}

\section{The native logarithm and complementary source estimates}
\label{sec:v67-logarithm}

Define the source-only integrated local functional
\begin{equation}
 \boxed{\mathscr L_{\Sigma,Y}
 =\int\rho\,\tr\left[
     \frac12U^2+\left(z-\frac R6\right)U
                   +\frac1{12}\Omega_{\mu\nu}\Omega^{\mu\nu}
                         \right].}
 \label{eq:v67-source-L}
\end{equation}
Here $\tr$ is the flavour trace, not the spatial functional trace.
It already counts four real scalars when the sources are scalar matrices.
The contraction $\tr(\Omega_{\mu\nu}\Omega^{\mu\nu})$ is nonpositive
on a real Euclidean chart; it is not replaced by its negative.

\begin{theorem}[Populated logarithm of the native composite-source family]
\label{thm:v67-native-log}
Fix total metric/source arity $N\le5$. With
$\ell_a=\log(1/(a\lambda))$ and total weight
$\deg K=1$, $\deg z=\deg Y=2$, the determinant
\eqref{eq:v67-delta-det} obeys
\begin{equation}
 \boxed{\mathsf J_N\mathsf P_4\Delta\Gamma_{a,1}
  =-\frac{\ell_a}{16\pi^2}\mathsf J_N\mathscr L_{\Sigma,Y}
                                  +\mathscr B_{a,L,N}.}
 \label{eq:v67-main-log}
\end{equation}
The critical coefficient norm of $\mathscr B_{a,L,N}$ is bounded uniformly
in mesh and periods on \eqref{eq:v66-domain}, at fixed profiles, source
order, polarization norms and inverse margins. It is not assigned zero.
The logarithmic counterterm is the opposite sign of the displayed term.
No extra power of $\chi^{-1}$ occurs in this scalar loop.
\end{theorem}
\begin{proof}
Expand the finite trace logarithm into its ordered words before
extracting its source jets. Lemma~\ref{lem:v67-source-symbols} gives the
same intermediate-region argument as V66: at critical total weight each
ordinary integrand has degree minus four, while a conservative native
mismatch is $O(a^2|p|^{-2})$. Its integrated contribution is bounded.
The lower-mask derivatives occupy a fixed annulus; the completed upper
region is bounded after rescaling $p=\theta/a$. Dyadic Poisson summation
at degree minus four gives a bounded finite-volume remainder. This proves
native equality to the ordinary integrated logarithmic residue for every
specified coefficient, including multiple source contacts.

To compute that residue, put $\mathcal P=-\mathcal D^2+zI+U$ and
$\widehat{\mathcal P}=\rho^{1/2}\mathcal P\rho^{-1/2}$.
Equation~\eqref{eq:v67-Laplace} gives
$Q_0^{\rm src}=M\widehat{\mathcal P}M$ with
$M=\rho^{1/2}B$. Along the positive multiplicative congruence
$M_t=e^{t\log M}$, cyclicity and logarithmic differentiation of the
residue give
$\frac d{dt}\operatorname{res}\log(M_t\widehat{\mathcal P}M_t)
=2\operatorname{res}\log M=0$.
This is the residue theorem of \cite{V67Scott}; it is not a finite
Gaussian change of measure. Work first at positive $z$ on a spectral
cut and then take mass jets; finite-rank zero-mode corrections do not
change the residue.

For the Laplace-type operator the established integrated heat coefficient
\cite[Eq.~(4.28)]{V67Vassilevich} is
\[
 \int\rho\,\tr\left[
 I\left(\frac{R^2}{72}-\frac{|\operatorname{Ric}|^2}{180}
                    +\frac{|\operatorname{Riem}|^2}{180}\right)
 -\frac R6(zI+U)+\frac12(zI+U)^2+\frac1{12}\Omega^2
 \right].
\]
The full density also contains $n_H\Delta_gR/30-
\tr(\mathcal D^2(zI+U))/6$. Their integrals vanish on the declared
periodic or compactly supported comparison. Subtracting the zero-added-source
coefficient leaves \eqref{eq:v67-source-L}. The real-scalar determinant
sign and factor are those of V66, giving \eqref{eq:v67-main-log}.
\end{proof}

\begin{remark}[No multiplicative-source shortcut]
\label{rem:v67-residue-density}
Equation~\eqref{eq:v67-main-log} concerns the determinant of the declared
source-dependent operator. It does not follow by multiplying V66's
unsourced heat density by a source. In particular a further independent
multiplicative curvature-composite insertion would retain its own
integration-by-parts and residue-density contacts. Equality of integrated
residues is not equality of finite determinants, power divergences, measures,
or pointwise heat densities. The new source functional is integrated
\emph{after} its source-dependent normal form has been computed.
\end{remark}

\begin{theorem}[Full-cutoff complement on the new fixed source panels]
\label{thm:v67-complement}
Let $I$ contain $s\le2$ potential-source arguments and any fixed number
of kinetic/metric arguments, with total arity $A\le5$ and at least one
added source. Put $d_I=4-2s$ and
$\gamma^R_{a,L,I}=(1-\mathbf T_{d_I})\gamma_{a,L,I}$, where the Taylor
assignment acts jointly on external momenta and $z$. Then, with
$r=|\alpha|+2j\le d_I+2$,
\begin{equation}
 \boxed{
 |\partial_K^\alpha\partial_z^j
 (\gamma^R_{a,L,I}-\gamma^R_{0,L,I})|
 \le C_{I,\alpha,j}\mathfrak n_I
       a^2(|K|+m_H)^{d_I+2-r}.}
 \label{eq:v67-complement}
\end{equation}
Every fixed higher derivative is controlled in the corresponding window
seminorm. The rooted kernel bound is
\begin{equation}
 \|\mathcal K_{a,L,I}^{\Delta}\|_{\rm rooted,b}
 \le C_{I,b}\,a^2(\mu+m_H)^{d_I+2},
 \label{eq:v67-rooted}
\end{equation}
with all $4(A-1)$ relative coordinates retained. Constants are independent
of periods and mesh on the fixed domain, not of unbounded source order.
The ordinary thermodynamic error has the separate bound
$C\lambda^{-2}(\lambda L_{\min})^{-b}(|K|+m_H)^{d_I+2-r}$ for $b>4$.
\end{theorem}
\begin{proof}
Before derivatives the primitive scalar loop has degree $-2s$.
Simultaneous reversal of every external and loop momentum leaves each
scalar difference tensor, root profile and covariance unchanged. Matrix
ordering does not affect that parity. Thus the full coefficient is even
in external momentum; its next odd weighted Taylor degree vanishes.
After $d_I+2$ weighted derivatives the ordinary integrand is $O(R^{-6})$.
The improved native interior discrepancy is $O(a^4R^{-2})$ by
Lemma~\ref{lem:v67-source-symbols}. The normalized four-dimensional sums
are bounded by
$a^4\int^{c/a}R\,dR=O(a^2)$ and
$\int_{c/a}^{\infty}R^{-3}dR=O(a^2)$.
Smooth completed upper pieces have the same scaling. The weighted Taylor
integral, or its equivalent even-mass dilation, restores the stated
external powers; mass differentiation retains one lower mask in each
propagator derivative. No comparison is made with an uncut massless
integral at zero momentum.

Additional loop-momentum derivatives and Poisson summation yield the
volume estimate shell by shell. Smooth external windows and sufficiently
many external derivatives, followed by Fourier integration by parts,
give every fixed rooted moment as in V26. The maximal relative
coordinate dimension is sixteen; the number of derivatives is chosen
accordingly. Multiplication by polarization tensors and one $L^1$ with
remaining $L^\infty$ source arguments introduces no total volume.
All signed loop words and contacts are formed before norms are taken.
\end{proof}

\section{Commuting local sources have linked derivative contacts}
\label{sec:v67-commuting}

Take $\Sigma=\sigma I$, $Y=jI$. Then $\omega=\Omega=0$ and
$\mathcal V_\Sigma=b_\sigma I$ with
$b_\sigma=\Delta_g\sigma+|\nabla\sigma|^2$.
Equation~\eqref{eq:v67-source-L} becomes exactly
\begin{equation}
 \boxed{\mathscr L_{\sigma,j}
 =n_H\int\rho\left[
    \frac{j^2}{2}+\left(z-\frac R6\right)j
    +\left(z+j-\frac R6\right)b_\sigma
    +\frac{b_\sigma^2}{2}\right].}
 \label{eq:v67-commuting-L}
\end{equation}
For example $\int\rho j\Delta_g\sigma=-\int\rho\nabla j\cdot\nabla\sigma$.
A local mass source and a local kinetic source therefore have a fixed
mixed derivative contact. Holding both sources constant misses it.

\begin{proposition}[Complete flat two-source logarithmic Hessian]
\label{prop:v67-quadratic}
At $g=I$, $\sigma=j=0$, the Hessian of
\eqref{eq:v67-commuting-L}, in the source order $(\sigma,j)$, is
\begin{equation}
 \boxed{n_H
 \begin{pmatrix}
 (k^2)^2+2zk^2&-k^2\\-k^2&1
 \end{pmatrix}.}
 \label{eq:v67-source-Hessian}
\end{equation}
The native logarithmic Hessian is
$-\ell_a/(16\pi^2)$ times this matrix, up to its bounded critical
remainder. At $z=0$, it has rank one for every $k$, and the source
combination is $j(k)-k^2\sigma(k)$.
\end{proposition}
\begin{proof}
Expand \eqref{eq:v67-commuting-L} to degree two and integrate the
Laplacian of $\sigma$. The result is
$n_H\int[\frac12(j+\Delta\sigma)^2+z|\nabla\sigma|^2]$,
with the separate linear coefficient $n_Hz\int j$ retained.
For an independent loop check, use the actual ordinary source vertices
$Q_\sigma(p+k,p)=2(p\cdot(p+k)+z)$,
$Q_j=1$, $Q_{\sigma\sigma}(p,p)=4(p^2+z)$ and
$Q_{\sigma j}=2$. Insert them in
$\frac12\Tr(CQ_{ab}-CQ_aCQ_b)$.
The angular coefficients of $|p|^{-4}$, divided by $n_H$, are respectively
\[
 -\tfrac12 k^2(k^2+2z),\qquad
 +\tfrac12 k^2,\qquad -\tfrac12
\]
for $\sigma\sigma$, $\sigma j$, and $jj$. In dimension four
$\langle t^2\rangle=1/4$, $\langle t^4\rangle=1/8$ and
$(2\pi)^{-4}|S^3|=1/(8\pi^2)$, giving the stated Hessian.
The $jj$ restriction checks the inherited potential-source sector;
the two coefficients involving $\sigma$ and their prescribed mixed
contacts are the new block.
\end{proof}

\begin{proposition}[An amplitude-only bound cannot control the source logarithm]
\label{prop:v67-amplitude}
On the flat $(2\pi)^4$ torus at $z=j=0$, take
$\sigma_N=N^{-1}\cos(Nx^1)$. Then
\begin{equation}
 \boxed{\|\sigma_N\|_\infty=N^{-1}\longrightarrow0,
 \qquad
 \frac{\mathscr L_{\sigma_N,0}}{\operatorname{Vol}}
     =n_H\left(\frac{N^2}{4}+\frac3{16}\right).}
 \label{eq:v67-amplitude}
\end{equation}
\end{proposition}
\begin{proof}
Here $b_{\sigma_N}=-N\cos(Nx^1)+\sin^2(Nx^1)$.
The cross term has zero average, while the other two averages are
$N^2/2$ and $3/8$. Substitute into
$n_H\int b_{\sigma_N}^2/2$.
\end{proof}
This is a genuine source-derivative test, not a counterexample inside a
fixed Fourier window: the sequence leaves every such window.
It makes the derivative dependence of the source norms necessary, much
as the companion descriptor retains derivatives of its covector.
The two statements concern different problems and no classical estimate
has been substituted for \eqref{eq:v67-complement}.

\section{A noncommuting four-source coefficient is evaluated}
\label{sec:v67-noncommuting}

Use the first two real scalar components and extend the following matrices
by zero on the other two:
\[
 A=\begin{pmatrix}1&0\\0&-1\end{pmatrix},\qquad
 B_0=\begin{pmatrix}0&1\\1&0\end{pmatrix},\qquad
 \Sigma=\varepsilon\bigl(s(x)A+t(x)B_0\bigr).
\]
The symbol $B_0$ here is a fixed test polarization, distinct from
$B=e^\Sigma$. On the flat $(2\pi)^4$ torus choose
$s=\sin x^1$, $t=\sin x^2$, $Y=z=0$.

\begin{theorem}[Complete periodic source coefficient with its curvature term]
\label{thm:v67-noncommuting}
In the declared integrated representative,
\begin{equation}
 \boxed{
 \frac{\mathscr L_{\Sigma,0}}{\operatorname{Vol}}
   =\varepsilon^2+\frac{43}{12}\varepsilon^4
                         +O(\varepsilon^6).}
 \label{eq:v67-periodic-coefficient}
\end{equation}
The connection-curvature summand contributes
$-\varepsilon^4/3$ to the volume-normalized fourth-order coefficient.
It is nonzero and cannot be omitted from the trace formula.
These are source Taylor coefficients of the specified operator; they are
not new physical gauge couplings or an anomaly statement.
\end{theorem}
\begin{proof}
Write $T=sA+tB_0$, $P_\mu=\partial_\mu T$. Direct differentiated
exponential series give
\[
 H_\mu=\varepsilon P_\mu+
       \frac{\varepsilon^3}{6}[T,[T,P_\mu]]+O(\varepsilon^5),\qquad
 \omega_\mu=-\frac{\varepsilon^2}{2}[T,P_\mu]+O(\varepsilon^4).
\]
Consequently
\begin{align*}
 \mathcal V_1&=\Delta T,&
 \mathcal V_2&=\sum_\mu P_\mu^2,\\
 \mathcal V_3&=\frac23\sum_\mu[P_\mu,[T,P_\mu]]
                      +\frac16[T,[T,\Delta T]],&
 \Omega_{\mu\nu}^{[2]}&=-[P_\mu,P_\nu].
\end{align*}
Since $\Delta T=-T$, the last double commutator in $\mathcal V_3$
vanishes. With $s'=\cos x^1$, $t'=\cos x^2$,
\[
 \mathcal V_1=-sA-tB_0,\quad
 \mathcal V_2=(s'^2+t'^2)I_2,\quad
 \mathcal V_3=-\frac83(ts'^2B_0+st'^2A).
\]
Now $\tr[A,B_0]^2=-8$. The second-order integrand is $s^2+t^2$;
the third is zero. The fourth-order integrand is
\[
 s'^4+t'^4+\frac23s'^2t'^2
       +\frac{16}{3}(s^2t'^2+t^2s'^2).
\]
The first part comes from $\tr(\mathcal V_2^2)/2$, the last from
$\tr(\mathcal V_1\mathcal V_3)$, and the change of the mixed coefficient
from $2$ to $2/3$ is precisely $\tr(\Omega^2)/12$.
Using the averages $\langle\sin^2\rangle=\langle\cos^2\rangle=1/2$
and $\langle\cos^4\rangle=3/8$ gives $43/12$ and the stated $-1/3$
curvature contribution. The functional is even in $\varepsilon$ on
this test: translation by $\pi$ in both active coordinates changes
$T$ to $-T$. Thus the next possible integrated term is order six.
The checker independently expands $e^{\varepsilon T}$ before forming
$L,R,H,\omega$, rather than assuming the displayed potential coefficients.
\end{proof}

The determinant contribution to these coefficients is
$-\ell_a\mathscr L_{\Sigma,0}/(16\pi^2)$. There is no additional factor
four: only the two displayed flavour directions contribute and their
trace is already evaluated. An independent local density source would
also retain the total-derivative heat terms; the periodic integral here
is not a pointwise heat-density assertion.

\section{One source counterterm, its Ward identity, and matching data}
\label{sec:v67-common-source}

The new logarithmic functional is an explicit common action/source block:
\begin{equation}
 C_{\Sigma,Y,1}^{\log}
  =\frac{\ell_a}{16\pi^2}\mathsf J_N\mathscr L_{\Sigma,Y}.
 \label{eq:v67-counterterm}
\end{equation}
It is to be included only if not already contained in the common local
assignment $-\mathsf P\Gamma_1$. It vanishes at $\Sigma=Y=0$ and
therefore changes no active zero-source normalization. The pure metric
logarithm of V66 is not counted a second time.

\begin{proposition}[External-source Ward and integrability identities]
\label{prop:v67-ward}
For the ordinary simultaneous variations
$\delta g=\mathcal L_cg$, $\delta\Sigma=c\cdot\partial\Sigma$ and
$\delta Y=c\cdot\partial Y$, one has
\begin{equation}
 \boxed{
 \int\left[
 \frac{\delta\mathscr L}{\delta g}:\mathcal L_cg
 +\tr\left(\frac{\delta\mathscr L}{\delta\Sigma}
                            c\cdot\partial\Sigma\right)
 +\tr\left(\frac{\delta\mathscr L}{\delta Y}
                            c\cdot\partial Y\right)
 \right]=0.}
 \label{eq:v67-Ward}
\end{equation}
All mixed source/metric Hessians are derivatives of this same functional.
The metric term alone need not vanish when the added sources vary.
\end{proposition}
\begin{proof}
$H$ is a matrix covector, $\omega$ transforms as a matrix covector under
spacetime diffeomorphisms, $\Omega$ as a two-form, and $U$ as a matrix
scalar. Every term in \eqref{eq:v67-source-L} is consequently a scalar
density. Its Lie variation is a divergence. Differentiate the identity
in the actual source amplitudes; ordinary mixed derivatives commute.
The source terms are not deleted when computing a metric response.
\end{proof}

This is an external-source identity of the ordinary logarithmic block,
not a proof of finite filtered BV restoration, internal-gauge invariance
at fixed sources, or the unknown inhomogeneous source primitive. A formal
source coordinate cannot be interpreted as an actually executed preparation
without the additional occurrence data distinguished by SM V52.

At every fixed source arity, derivatives of the matrix exponential,
pointwise inverse and normal form are bounded on the declared chart.
Polarization and smooth windows give a logarithmic local bound
$C_{N,b}|\ell_a|(\mu+m_H)^4$, with potential-source weights included
and one source in $L^1$, the others in $L^\infty$. There is no total-volume
factor. These local coefficients are not vanishing remainders.

The finite critical matching data have the same convergent dimensionless
quadrature as V66 after replacing its vertices by
\eqref{eq:v67-native-replacements}. For two specified completions with the
same ordinary source germ, their critical integrands differ by
$O(|\theta|^{-2})$ near $\theta=0$ and are integrable on the completed
zone. This computes a matching \emph{prescription}, not its unevaluated
numerical value. Different nonlinear source completions with the same
first source derivative can have different higher contacts. The present
normal form and coefficients belong specifically to
\eqref{eq:v67-finite-action}; that choice is not claimed unique.

\section{Status, verification, and the remaining continuum interface}
\label{sec:v67-ledgers}

\subsection*{Proved}
The explicitly declared source extension of the V27 scalar action has
an exact ordinary matrix Laplace normal form and determined curvature.
Its native critical logarithm is \eqref{eq:v67-main-log}, with bounded
finite critical remainder and the stated source-complement estimates.
The scalar kinetic/potential Hessian, its direct loop check, the
amplitude-only counterexample and the noncommuting periodic fourth-source
coefficient are evaluated. All mixed source/metric logarithms come from
one local functional with the external-source Ward identity.

\subsection*{Inherited}
V9 supplies the earlier potential-source interpretation, not a replacement
regulator. V27 supplies the completed scalar difference tensors, floor,
profiles and improved symbols. V66 supplies the zero-source determinant
and the residue-matching strategy. V54--V56 supply the separate
finite-range/contact distinctions. The companion audit uses only the
identified terminal claims and does not import classical rate estimates
into quantum ultraviolet bounds.

\subsection*{Literature input}
The general matrix Laplace heat coefficient and residue-trace rules are
\cite{V67Vassilevich,V67Scott}. They are not proved by the accompanying
finite checker. The application is to the derived operator
\eqref{eq:v67-Laplace}, after matching its integrated residue to the
native finite source determinant. No first general theorem about
variable couplings or scalar composite renormalization is claimed.

\subsection*{Conditional}
A physical interpretation of the added source labels requires its own
readout/preparation specification. Using the resulting block in later
hard-metric insertions requires the same common counterterm and all
source contacts. Finite critical matching quadratures remain unevaluated.
There is no source-independent suppression of the connection term by
assuming simultaneous diagonalization. No stronger convergence follows
from a source-amplitude bound without derivative control.

\subsection*{Open}
The bounded inhomogeneous critical source primitive, other matter and
chiral contributions, full two-loop master/reference assembly and
invariant many-shell EFT class remain open. The next missing local data
are not the source logarithms computed here but the finite matched
source array and its combined symmetry equation. The companion exact
preparation and classical descriptor theorems do not supply those data.

\begin{center}\small
\begin{tabular}{>{\raggedright\arraybackslash}p{.22\textwidth}>{\raggedright\arraybackslash}p{.70\textwidth}}
\toprule
Variable & Scope in V67\\\midrule
Loop/fields & One closed four-real-scalar determinant; metric and added
sources retained; no integrated gauge, ghost or graviton line.\\
Source/EFT degree & Total prescribed arity at most five; critical weight
four. Complement theorem displayed for zero, one or two potential marks.\\
Cutoff/volume & Fixed lower scale and inverse/profile margins; bounded
critical finite part and $O(a^2)$ complementary comparison, with separate
volume error and rooted moments.\\
Mass/chart & Polynomial mass jets at the positive hard reference, including
$z=0$; pointwise positive source chart for the ordinary normal form.\\
Symmetry/measure & Complete external-source transformation; no change of
finite scalar measure, active action or normalization.\\
Verification & Exact matrix, source-series, angular and periodic-coefficient
checks; no numerical native matching integral or proof-assistant claim.\\\bottomrule
\end{tabular}
\end{center}

\begin{thebibliography}{99}
\bibitem{V67SM52} A.~Pelissier, \emph{Renewal SM Source Selection Research
Notes}, V52, supplied companion, terminal labels
\srcid{sel52:thm:global}, \srcid{sel52:thm:prepface},
\srcid{sel52:prop:bank}.
\bibitem{V67Station54} A.~Pelissier, \emph{Renewal Native Regenerative
Stationarity Closure Notes}, V54, supplied companion, terminal labels
\srcid{v54:thm:descriptor}, \srcid{v54:lem:green},
\srcid{v54:thm:lift}, \srcid{v54:eq:native-defect}.
\bibitem{V67Bridge45} A.~Pelissier, \emph{Renewal Exact Action Einstein Bridge
Research Notes}, V45, supplied companion, terminal labels
\srcid{v45:weak-rank}, \srcid{v45:canonical-theorem},
\srcid{v45:frequency-formula}.
\bibitem{V67Vassilevich} D.~V.~Vassilevich,
Heat kernel expansion: user's manual,
\emph{Phys. Rept.} \textbf{388} (2003), 279--360;
arXiv:hep-th/0306138, Eqs.~(2.1)--(2.4), (4.25), (4.28).
\bibitem{V67Scott} S.~Scott, The residue determinant,
\emph{Comm. Partial Differential Equations} \textbf{30} (2005), 483--507;
arXiv:math/0406268, Proposition 1.2 and Theorems 1.7--1.8.
\end{thebibliography}

\section*{Append-only continuation marker: V67}
\addcontentsline{toc}{section}{Append-only continuation marker: V67}
\label{sec:v67-continuation-marker}
\textbf{End of V67.} Preserve Sections 1--520. Local kinetic and potential
source logarithms are now specified on an explicit extension of the active
scalar carrier. Their derivative and noncommuting contacts are retained.
The finite critical source normalization and full quantum descent remain
separate. No active zero-source normalization changes.

% ============================================================================
% V68 -- finite source normalization beyond the logarithmic residue
% ============================================================================
\clearpage
\section{V68: a populated finite source-matching problem}
\label{sec:v68-start}

V67 computes the logarithmic local functional, but equality of logarithmic
residues does not determine finite critical source contacts. This installment
computes such contacts for the \emph{same} source family and proves a local
matching theorem for one explicitly specified comparison family. The comparison
has the same zero-source action, covariance and measure. It is not installed
as the active source convention. The new identity is the central kinetic-source
rescaling identity, not the diffeomorphism or complete ESM master equation.

The input is exactly \eqref{eq:v67-finite-action}, with the free scalar symbol
\eqref{eq:v27-P} and the two different upper profiles
\eqref{eq:v14-profiles}. No heat coefficient is recalculated. No profile
is replaced by a sharp cutoff, no new gauge field is introduced, and no
companion classical continuation estimate is used as a quantum bound.
We keep the four-real-scalar trace, including its noncommuting matrix sources.
The familiar need to transport reference and source data together under field
redefinitions is comparison material \cite{V68CPV,V68IIM}; all finite identities
used here are proved directly.

Throughout, $z=m_H^2$, $P=\widetilde P_a$, $C=\nu_\lambda P^{-1}$ and
$\nu_\lambda=I-\Theta_\lambda$. The quotient $\nu_\lambda/P$ is defined to
be zero on retained covariance-zero modes, including at $z=0$. Define
\begin{equation}
 K_a=P_a^b+s(aD)^2(P_a^n-P_a^b),\qquad
 D_a=P-K_a=(\zeta-s^2)(aD)(P_a^n-P_a^b).
 \label{eq:v68-defect}
\end{equation}
$K_a$ is the scalar operator obtained by a \emph{constant kinetic-source}
variation of the completed interactions. It is not the chosen free operator
$P$. The notation $D_a$ in V68 denotes this even operator, not a first-order
scalar difference or the ordinary BV differential.

The new local problem is unusually concrete: $D_a$ has support strictly away
from zero dimensionless momentum, but its contribution to a critical
one-loop coefficient need not tend to zero. Source-logarithm universality
alone cannot choose that finite coefficient.

\section{The actual finite source-rescaling identity has two reference terms}
\label{sec:v68-rescaling}

In dimensionless momentum $x=ap$, put
\begin{align}
 p_b(x)&=4\sum_{\mu=0}^3\sin^2(x_\mu/2),\nonumber\\
 p_n(x)&=4\sinh^2(x_0/2)+4\sum_{i=1}^3\sin^2(x_i/2)
                  +\frac{-x_0^4+\sum_i x_i^4}{12},\nonumber\\
 p_*(x)&=p_b(x)+\zeta(x)[p_n(x)-p_b(x)],\qquad
 k_*(x)=p_b(x)+s(x)^2[p_n(x)-p_b(x)],\nonumber\\
 d_*(x)&=p_*(x)-k_*(x).
 \label{eq:v68-dimensionless}
\end{align}
Thus $P=a^{-2}p_*+z$, $K_a=a^{-2}k_*+z$, and $D_a=a^{-2}d_*$.
The last operator is independent of the mass parameter.

\begin{lemma}[Location and sign of the native defect]
\label{lem:v68-positive-defect}
One has $0\preceq D_a\prec P$ on nonzero modes. Its symbol vanishes for
$|ap|\le1/64$ and $|ap|\ge1/8$. On $1/32<|ap|<1/16$ it is strictly
positive. In particular the two profile placements cannot be identified.
\end{lemma}
\begin{proof}
The support conditions imply $\zeta\ge s^2$: where $s$ is nonzero,
$\zeta=1$. Direct expansion gives
\[
 p_n-p_b=\frac1{12}\sum_{\mu=0}^3x_\mu^4
                  +4\sum_{j\ge2}\frac{x_0^{4j}}{(4j)!}.
\]
It is positive at every nonzero momentum in the central chart. Moreover
$k_*\ge p_b>0$ at every nonzero principal momentum, so $d_*<p_*$.
The declared plateaux and supports prove the remaining assertions.
\end{proof}

Let $\mathcal Q_a(q,\Sigma,Y)$ denote V67's source Hessian and write
$B=e^\Sigma$. Let $\mathscr S$ be differentiation of the \emph{constant
central source shift} $\Sigma(x)\mapsto\Sigma(x)+tI$ at $t=0$.
It does not act on $q,Y,z,C$ or the coefficient measure. Define
\[
 A=(I+C(\mathcal Q_a-P))^{-1},\qquad C_{\mathcal Q}=AC.
\]
All identities below are finite coefficient identities; the inverses can
also be read on the common sufficiently small source chart.

\begin{theorem}[Exact central source identity on the native completion]
\label{thm:v68-native-Ward}
For arbitrary retained local metric and matrix sources,
\begin{equation}
 \boxed{\mathscr S\mathcal Q_a=2(\mathcal Q_a-D_a).}
 \label{eq:v68-Q-shift}
\end{equation}
For the one-loop determinant with its factor $\hbar$ removed,
\begin{equation}
 \boxed{\mathscr S\Gamma_a
       =\Tr I-\Tr(A\Theta_\lambda)-\Tr(C_{\mathcal Q}D_a).}
 \label{eq:v68-native-Ward}
\end{equation}
The trace includes site and real-flavour indices. Normalizing $\Gamma_a$
at zero added source, at the same retained metric, does not change this
identity: that subtracted functional is independent of the added source.
\end{theorem}
\begin{proof}
Expand the three terms of \eqref{eq:v67-finite-action} before varying.
Their source-independent subtraction is precisely $K_a$, whereas the
retained Gaussian is $P$. Hence $\mathcal Q_a=D_a+\overline Q_a$, where
all three full source bilinears in $\overline Q_a$ scale by $e^{2t}$.
The filtered source also shifts by $tI$ because $S1=1$. This proves
\eqref{eq:v68-Q-shift}, including nonconstant $q,\Sigma,Y$ and their
noncommuting flavour values. Finite logarithmic differentiation gives
$\mathscr S\Gamma_a=\Tr(C_{\mathcal Q}(\mathcal Q_a-D_a))$.
Since $CP=\nu_\lambda$,
\[
 C_{\mathcal Q}\mathcal Q_a
 =A[\nu_\lambda+C(\mathcal Q_a-P)]=I-A\Theta_\lambda.
\]
This proves \eqref{eq:v68-native-Ward} without inverting $C$.
\end{proof}

The first term of \eqref{eq:v68-native-Ward} is the finite coefficient
Jacobian of a scalar rescaling. Its derivatives in other sources vanish.
The next terms have different support: the retained lower reference and
the ultraviolet mismatch of the free and source quadratic forms.
Neither term is an anomaly conclusion. Keeping $D_a$ is necessary even
though it vanishes on every fixed low-momentum scalar mode for sufficiently
small $a$.

\section{An explicitly recorded comparison, not a silent source replacement}
\label{sec:v68-comparison}

Consider, solely as a matching reference, the finite source Hessian
\begin{equation}
 \boxed{\mathcal Q_a^\sharp
       =\mathcal Q_a+B D_a B-D_a.}
 \label{eq:v68-sharp}
\end{equation}
Here $B$ is pointwise multiplication, and $D_a$ acts on the scalar carrier.
The new term is symmetric and its coefficients are spatially quasilocal,
not asserted strictly finite range. It is zero at $\Sigma=0$, for every
$q,Y$. Thus the zero-added-source action and all its old metric derivatives
are exactly unchanged. This is a \emph{declared alternative source
completion}, not a change to the active V67 family.

\begin{proposition}[Comparison identity and agreement of the lower row]
\label{prop:v68-sharp-Ward}
Set $A^\sharp=(I+C(\mathcal Q_a^\sharp-P))^{-1}$. Then
\begin{equation}
 \mathscr S\mathcal Q_a^\sharp=2\mathcal Q_a^\sharp,
 \qquad
 \boxed{\mathscr S\Gamma_a^\sharp=\Tr I-\Tr(A^\sharp\Theta_\lambda).}
 \label{eq:v68-sharp-Ward}
\end{equation}
For a coefficient with at most $N$ soft source/metric marks, each of
momentum at most $\mu$, assume
\begin{equation}
 2\lambda+N\mu<\frac1{64a}.
 \label{eq:v68-support-domain}
\end{equation}
Then its lower-reference coefficients agree exactly:
\begin{equation}
 \boxed{[\Tr((A-A^\sharp)\Theta_\lambda)]_{\le N}=0.}
 \label{eq:v68-lower-equality}
\end{equation}
\end{proposition}
\begin{proof}
$\mathscr S(BD_aB)=2BD_aB$, so \eqref{eq:v68-Q-shift} gives the first
identity. The previous determinant proof gives the second. For the last
claim use the finite resolvent identity for $A-A^\sharp$ and expand only
to the prescribed source degree. Each word contains $BD_aB-D_a$ and a
factor $\Theta_\lambda$. At zero source the former vanishes; every
nonzero differentiated word contains one $D_a$ on a routed scalar
momentum. A loop based at the $\Theta_\lambda$ factor has momentum at
most $2\lambda$. Every other line differs by a partial sum of at most
$N$ external momenta. It cannot reach the support of $D_a$ under
\eqref{eq:v68-support-domain}. All words therefore vanish individually.
The argument uses finite source order; it is not an unrestricted statement
about arbitrarily high Fourier modes of $e^\Sigma$.
\end{proof}

In particular, if $\mathcal D_a=\Gamma_a-\Gamma_a^\sharp$, then on these
coefficients
\begin{equation}
 \boxed{\mathscr S\mathcal D_a=-\Tr(C_{\mathcal Q}D_a).}
 \label{eq:v68-difference-Ward}
\end{equation}
This equality determines the extra finite normalization equation instead
of assuming that the ultraviolet term is absent. A source shift lowers
the number of kinetic marks by one; checking a row through arity $N$
requires the determinant and counterterm through arity $N+1$.

\section{Exact constant-source critical coefficients, including mass}
\label{sec:v68-constant}

Set $q=0$ and take $\Sigma,Y$ constant symmetric matrices. They need not
commute. At a momentum $p$ the native Hessian is exactly
\begin{equation}
 \boxed{\mathcal Q_a=D_a+B(K_a+Y)B.}
 \label{eq:v68-constant-Q}
\end{equation}
The comparison is $\mathcal Q_a^\sharp=B(P+Y)B$. These equations are
consequences of the source placements, not continuum replacements.

Write $p_0=P|_{z=0}$, $k_0=K_a|_{z=0}$, $\nu=1-\Theta$, and put
\[
 h=1-\nu k_0/p_0,\quad E=e^{-2\Sigma},\quad
 F_h=(I+h(E-I))^{-1},\quad Z_\Sigma=\nu I+\Theta E.
\]
At retained zero modes the following expressions are read by cancellation
or as zero contributions to the source determinant. No mass expansion of
a physical zero-mode propagator is made.

\begin{theorem}[Full derivative-free critical kinetic dependence]
\label{thm:v68-constant-critical}
Let $\mathfrak F_{a,4}(\Sigma,Y,z)$ be the term of total degree two in
$(z,Y)$ in the per-volume difference
$\Gamma_a(\Sigma,Y,z)-\Gamma_a(0,Y,z)$. Then
\begin{equation}
 \boxed{\mathfrak F_{a,4}
 =-\frac1{4V_L}\sum_p p_0(p)^{-2}
 \tr\left\{
 [F_h(zZ_\Sigma+\nu Y)]^2-[zI+\nu Y]^2
 \right\}.}
 \label{eq:v68-full-critical}
\end{equation}
Only two disjoint momentum bands contribute. On the lower band $D_a=0$,
$h=\Theta$ and $F_hZ_\Sigma=I$, so its summand is
\begin{equation}
 -\frac1{4p_0^2}\left\{
 2z\nu\tr((F_\Theta-I)Y)
 +\nu^2\tr(F_\Theta YF_\Theta Y-Y^2)\right\}.
 \label{eq:v68-lower-generating}
\end{equation}
On the upper band $\Theta=0$, $h=D_a/p_0$, and the summand is
\begin{equation}
 \boxed{-\frac1{4p_0^2}
 \tr\{F_h T F_h T-T^2\},\qquad T=zI+Y.}
 \label{eq:v68-upper-generating}
\end{equation}
The last band is exactly the native-minus-comparison critical coefficient.
All intermediate momenta contribute zero to this kinetic-dependent
critical difference.
\end{theorem}
\begin{proof}
Multiplication of $I+C(\mathcal Q_a-P)$ by the scalar $P=p_0+z$,
and congruence by $B^{-1}$, give
\[
 B^{-1}\{P I+\nu(\mathcal Q_a-P)\}B^{-1}
 =p_0\{I+h(E-I)\}+zZ_\Sigma+\nu Y.
\]
Finite-dimensional determinant multiplication is legitimate here; its
$2\tr\Sigma$ and $\log P$ terms have no uncancelled contribution to the
kinetic-dependent degree-two difference. Expanding the \emph{finite}
trace logarithm to degree two yields \eqref{eq:v68-full-critical}; the
factor is $\tfrac12(-\tfrac12)=-\tfrac14$. Since lower and upper
supports are disjoint, $D_a=0$ on the lower band and $\Theta=0$ on the
upper. The two displayed specializations follow. Where both vanish,
$F_h=Z_\Sigma=I$ and $\nu=1$, so the difference is zero. On covariance-zero
modes $F_hZ_\Sigma=I$ and $\nu=0$, with the same conclusion.
\end{proof}

This theorem computes the entire zero-spacetime-derivative critical
source array by finite expansion in $\Sigma$. For source arity at most
five retain precisely the words with that arity; $z$ carries weight two
but is not a source occurrence. It includes the $zY$ and $z^2$ contacts,
not only two potential insertions. It excludes the pure-$Y$ logarithm
already assigned in V67.

\begin{proposition}[Exact matrix-potential Hessian]
\label{prop:v68-potential-Hessian}
At fixed $\Sigma$ and $z=0$, define $c=\nu/p_0$. Then
\begin{equation}
 \boxed{D_Y^2\Gamma_a(\Sigma,0)[Y_1,Y_2]
    =-\frac1{2V_L}\sum_p c(p)^2
                     \tr(F_hY_1F_hY_2).}
 \label{eq:v68-Hessian}
\end{equation}
The derivative of its kinetic-dependent part at $\Sigma=0$ is
\begin{equation}
 -\frac1{V_L}\sum_pc^2h\,
                   \tr\bigl(S(Y_1Y_2+Y_2Y_1)\bigr).
 \label{eq:v68-mixed-third}
\end{equation}
\end{proposition}
\begin{proof}
The finite covariance attached to a potential vertex is $cF_h$ after
cyclically moving the two factors $B$. Differentiate the finite trace
logarithm twice. Next use $D_\Sigma F_h|_0[S]=2hS$. Trace cyclicity,
not commutation of $S,Y_1,Y_2$, gives \eqref{eq:v68-mixed-third}.
\end{proof}

\section{Nonzero finite contacts and an evaluated noncommuting source tower}
\label{sec:v68-moments}

Let $\int_x=(2\pi)^{-4}\int d^4x$. Define the finite upper moments
\begin{equation}
 U_j=\int_{[-\pi,\pi]^4}\frac{\dd^4x}{(2\pi)^4}
       \frac1{p_*(x)^2}\left(\frac{d_*(x)}{p_*(x)}\right)^j,
 \qquad j\ge1.
 \label{eq:v68-upper-moments}
\end{equation}
The integrands are defined to be zero near $x=0$. They are smooth and
supported in $1/64\le|x|\le1/8$. For the lower profile
$\Theta_\lambda(p)=\theta(p/\lambda)$, put
\begin{equation}
 L_j=\int_y\frac{(1-\theta(y))^2\theta(y)^j}{|y|^4},\qquad
 N_j=\int_y\frac{(1-\theta(y))\theta(y)^j}{|y|^4},\qquad
 M_j=U_j+L_j.
 \label{eq:v68-lower-moments}
\end{equation}
Only $1\le|y|\le2$ contributes. These are specified finite coefficient
integrals, not numerical values inferred without choosing the profiles.

\begin{theorem}[Strict native finite coefficient and its comparison rate]
\label{thm:v68-positive-contact}
All $U_j,L_j,N_j$ are strictly positive, and decrease with $j$. The nested
native upper profiles imply the explicit profile-independent lower bound
\begin{equation}
 \boxed{U_1\ge\frac1{2^{21}\pi^2}>0.}
 \label{eq:v68-positive-bound}
\end{equation}
The native kinetic-dependent coefficient of $\tr(\Sigma Y^2)$ is
$-M_1$; the native-minus-comparison coefficient is $-U_1$. Therefore
neither is determined by its vanishing logarithmic coefficient.

Let $U_{j,a,L},L_{j,a,L}$ be the actual finite sums corresponding to
\eqref{eq:v68-upper-moments}--\eqref{eq:v68-lower-moments}, using $p_0$
in the lower sum. For each fixed $j$ and $b>4$,
\begin{align}
 |U_{j,a,L}-U_j|&\le C_{j,b}(a/L_{\min})^b,\nonumber\\
 |L_{j,a,L}-L_j|+|N_{j,a,L}-N_j|
 &\le C_{j,b}\left[(a\lambda)^4+
                           (\lambda L_{\min})^{-b}\right].
 \label{eq:v68-moment-errors}
\end{align}
\end{theorem}
\begin{proof}
Positivity of the upper moments follows from Lemma~\ref{lem:v68-positive-defect}.
A smooth transition from one to zero has a nonempty open set with
$0<\theta<1$, proving positivity of the lower moments. On the annulus
$1/32\le r=|x|\le1/16$, one has $\zeta=1,s=0$,
$p_n-p_b\ge r^4/48$ and $p_n\le2r^2$. The last elementary inequality
follows from the sine/hyperbolic-sine series for $r\le1/8$. Hence
$d_*/p_*^3\ge1/(384r^2)$ there. Since the normalized four-dimensional
radial measure is $r^3dr/(8\pi^2)$,
\[
 U_1\ge\frac1{3072\pi^2}
           \int_{1/32}^{1/16}r\,dr
       =\frac1{2^{21}\pi^2}.
\]
The coefficient claim follows from $F_h=I+2h\Sigma+O(\Sigma^2)$ in
\eqref{eq:v68-lower-generating} and \eqref{eq:v68-upper-generating}.
For the upper finite sum, scaling $p=x/a$ gives exactly a normalized
Riemann sum of a fixed smooth periodic function. Its Fourier coefficients
decay to every fixed order, giving the first estimate by Poisson summation.
On the lower band, $p_0=p^2+O(a^4|p|^6)$ by the inherited corrected
symbol. Thus $p_0^{-2}-|p|^{-4}=O(a^4)$ there. Integrating over a band
of volume $O(\lambda^4)$ and separately applying Poisson summation at
scale $\lambda$ proves the second estimate. No divergent integral is
split into separately assigned infinite constants.
\end{proof}

The ultraviolet lower bound does not require a lower bound on the width
of the lower-reference transition. A fixed finite grid can have a smaller moment;
\eqref{eq:v68-moment-errors} states exactly how its positive continuum
coefficient is reached. In the Hessian normalization,
\[
 D_\Sigma D_Y^2(\Gamma_a-\Gamma_a^\sharp)|_0[I,Y_1,Y_2]
 \longrightarrow-2U_1\tr(Y_1Y_2).
\]
For $Y_1=Y_2=I_4$ this is $-8U_1$, not $-2U_1$ times an additional
unrecorded scalar-loop multiplicity.

\begin{proposition}[The two-potential tower through total source arity five]
\label{prop:v68-tower}
At zero mass, the kinetic-dependent critical quadratic-in-$Y$ density,
through cubic degree in $\Sigma$, is
\begin{equation}
 \boxed{\begin{aligned}
 \mathfrak F_{Y^2}={}&-M_1\tr(\Sigma Y^2)
 +(M_1-2M_2)\tr(\Sigma^2Y^2)-M_2\tr(\Sigma Y\Sigma Y)\\
 &+(-4M_3+4M_2-\tfrac23M_1)\tr(\Sigma^3Y^2)\\
 &+(2M_2-4M_3)\tr(\Sigma^2Y\Sigma Y).
 \end{aligned}}
 \label{eq:v68-tower}
\end{equation}
The corresponding finite matching to the comparison is obtained by replacing
$M_j$ by $U_j$ and changing the sign when used as an additive counterterm.
\end{proposition}
\begin{proof}
The ordered inverse has the expansion
\[
 F_h=I+2h\Sigma+(4h^2-2h)\Sigma^2
             +(8h^3-8h^2+\tfrac43h)\Sigma^3+O(\Sigma^4).
\]
Insert it in $-\tfrac14\tr(F_hYF_hY-Y^2)$. Trace cyclicity identifies
only cyclic words; in particular $\tr(\Sigma Y\Sigma Y)$ is not replaced
by $\tr(\Sigma^2Y^2)$. Collecting the powers of $h$ proves the formula.
For arbitrary fixed degree the scalar generating coefficients are
\[
 a_r(h)=\frac{(-2)^r}{r!}
          \sum_{k=1}^r(-h)^k k!\left\{\begin{matrix}r\\ k\end{matrix}\right\},
 \quad r\ge1,\qquad a_0=1,
\]
where braces denote Stirling numbers of the second kind. The matrix series
$F_h=\sum_ra_r(h)\Sigma^r$ therefore gives all the retained mass/source
coefficients in Theorem~\ref{thm:v68-constant-critical} by finite arithmetic.
\end{proof}

The quartic part of \eqref{eq:v68-tower} can also be written
\begin{equation}
 (M_1-3M_2)\tr(\Sigma^2Y^2)
                 -\frac{M_2}{2}\tr[\Sigma,Y]^2.
 \label{eq:v68-commutator}
\end{equation}
It is a derivative-free \emph{finite} contact. It is distinct from V67's
curvature $-[H_\mu,H_\nu]$, which vanishes for constant sources.
For two flavour directions take
$\Sigma=u\operatorname{diag}(1,-1)$ and
$Y=v\left(\begin{smallmatrix}0&1\\1&0\end{smallmatrix}\right)$,
with zeros on the other two directions. The $u^2v^2$ coefficient is
$2(M_1-M_2)$. Replacing $Y$ by $v\operatorname{diag}(1,-1)$ gives
$2(M_1-3M_2)$. Their difference is exactly $4M_2>0$.
Both configurations have constant sources and identical V67 logarithmic
kinetic dependence (zero). Their finite contacts are nevertheless different.

\section{One local matching functional and its full-cutoff error}
\label{sec:v68-matching}

At fixed total arity $N\le5$, expand the \emph{difference of the two
finite determinants} in the same ordered field/source labels. For a panel
$I$ with $s_I\le2$ potential arguments, let
$d_I=4-2s_I$ and define
\begin{equation}
 M_{a,I}=\mathbf T_{d_I}(\Gamma_a-\Gamma_a^\sharp)_I,\qquad
 \Gamma_a^{\rm match}=\Gamma_a-M_a.
 \label{eq:v68-matching-functional}
\end{equation}
$M_a$ is one local scalar action/source functional, reconstructed from the
finite Taylor coefficients of a single determinant difference. It includes
its power-divergent parts as well as its finite critical part. These are
not independently fitted Hessians. It has at least one kinetic-source
argument and vanishes at $\Sigma=0$, with all metric marks retained.

\begin{theorem}[Native finite matching on the local source panels]
\label{thm:v68-matching}
On the inherited domain
\begin{equation}
 \lambda\ge32768\mu>0,\quad a\lambda\le2^{-24},\quad
 \lambda L_\nu\ge1,\quad 0\le m_H\le M\lambda,
 \label{eq:v68-domain}
\end{equation}
and at the stated fixed source order,
\begin{equation}
 \boxed{\begin{aligned}
 &\left|\partial_K^\alpha\partial_z^j
 (\Gamma_a^{\rm match}-\Gamma_a^\sharp)_I\right|\\
 &\hspace{12mm}\le C_{I,\alpha,j}\mathfrak n_I
       a^2(|K|+m_H)^{d_I+2-r},\qquad
       r=|\alpha|+2j\le d_I+2.
 \end{aligned}}
 \label{eq:v68-matching-error}
\end{equation}
Every prescribed higher derivative is controlled in its fixed-window
seminorm. The rooted kernel satisfies
\begin{equation}
 \boxed{\|\mathcal K^\Delta_{a,L,I}\|_{\rm rooted,b}
    \le C_{I,b}a^2(\mu+m_H)^{d_I+2}.}
 \label{eq:v68-rooted}
\end{equation}
The critical coefficients of $M_a$ have finite ultraviolet limits given
by smooth dimensionless integrals. The derivative-free part of those
limits is exactly \eqref{eq:v68-upper-generating}; the limits of the
remaining metric/derivative coefficients are specified by the same
finite determinant, but not numerically evaluated here. No new logarithm
is generated by this comparison.
\end{theorem}
\begin{proof}
Use the exact interpolation
$\mathcal Q_t=\mathcal Q_a+t(BD_aB-D_a)$ at finite cutoff. Then
\begin{equation}
 \mathcal D_a=-\frac12\int_0^1
   \Tr\{(I+C(\mathcal Q_t-P))^{-1}C(BD_aB-D_a)\}\,dt.
 \label{eq:v68-interpolation}
\end{equation}
Every fixed-order differentiated word has at least one $D_a$. Each loop
momentum is within $N\mu$ of its momentum. Therefore all routed loop
momenta lie in a fixed dimensionless annulus bounded away from zero;
all lower masks on the loop equal one. This holds for the native upper
vertices too, including their distinct scalar-leg and free profiles.

After $p=x/a$, a coefficient with $s_I$ potential marks has the form
$a^{2s_I-4}$ times a smooth dimensionless finite sum depending on
$(aK,a^2z)$. Its derivatives are bounded uniformly at every fixed order.
The same remains true on the interpolating analytic germ, since at zero
added sources $\mathcal Q_t$ is the unchanged $Q_a(q)$ and the expansion
contains only finitely many ordered words. Simultaneous reversal of all
momenta leaves the scalar kernels and the $D_a$ term unchanged. The next
odd Taylor weight is absent. The weighted remainder after degree $d_I$
therefore begins at $d_I+2$ and is bounded by
$a^{2s_I-4}a^{d_I+2}(|K|+m_H)^{d_I+2}=a^2(|K|+m_H)^{d_I+2}$.
Differentiating gives \eqref{eq:v68-matching-error}.

For a critical local coefficient, the total weight is
$|\alpha|+2j+2s_I=4$. Scaling leaves no power of $a$. At zero external
momenta its dimensionless integrand is fixed, smooth, periodic and zero
near the origin. Fourier decay yields convergence of its finite sum to
its dimensionless integral, with error $C_b(a/L_{\min})^b$. Subcritical
coefficients carry their explicit powers $a^{w-4}$ and are stored rather
than bounded as small errors. The annular support precludes a logarithm.

Finally multiply by the fixed smooth windows in every independent
external momentum. Integrate by parts in those variables to order exceeding
$4(A-1)+b$, using the just-proved higher-derivative bounds. The resulting
lattice Fourier kernel is summable with the stated root weight. This is
the same physical Fourier normalization as \eqref{eq:v67-rooted}; it gives
one $L^1$ argument and the other arguments in $L^\infty$, without a volume
factor. No assertion about arbitrarily large source or derivative order
is used.
\end{proof}

This theorem fixes a concrete source-normalization comparison. It does
not assert that $M_a$ is separately diffeomorphism invariant, or that the
new source completion satisfies a full finite BV identity. Its signs are
fixed: $M_a$ represents native minus comparison, so an additive
counterterm transporting the native answer to this comparison is $-M_a$.
In particular its upper $\tr(\Sigma Y^2)$ term is $+U_1\tr(\Sigma Y^2)$.

\section{The finite source identity after matching and the remaining anchors}
\label{sec:v68-common-identity}

\begin{corollary}[Restored central-source identity to the prescribed jets]
\label{cor:v68-restoration}
Use one additional kinetic mark for the constant shift, and retain all
finite local weights in \eqref{eq:v68-matching-functional}. Through the
remaining arity $N-1$, the extra ultraviolet term in
\eqref{eq:v68-native-Ward} is removed at the assigned Taylor degrees:
\begin{equation}
 \boxed{\mathbf T\left\{
 \mathscr S\Gamma_a^{\rm match}
 -\Tr I+\Tr(A^\sharp\Theta_\lambda)\right\}=0.}
 \label{eq:v68-restored-jet}
\end{equation}
The unprojected source-window difference satisfies the differentiated
$O(a^2)$ bound of Theorem~\ref{thm:v68-matching}. The lower-reference
term is retained and agrees with the native one on these soft coefficients.
\end{corollary}
\begin{proof}
Subtract \eqref{eq:v68-sharp-Ward} from the derivative of
$\Gamma_a^{\rm match}=\Gamma_a-M_a$. The result is
$\mathscr S(\mathcal D_a-M_a)$. The Taylor projector commutes with the
constant source derivative when the extra input mark is supplied. Its
jet is zero by definition, and its complementary estimate is the
corresponding marked case of Theorem~\ref{thm:v68-matching}. Equation
\eqref{eq:v68-lower-equality} gives the last assertion.
\end{proof}

Even after this comparison, the lower part in
\eqref{eq:v68-lower-generating} remains. At $z=0$ its cubic source
coefficient is $-L_1\tr(\Sigma Y^2)$, which is strictly nonzero. Setting
that coefficient to zero would require a \emph{further} local source
normalization and the corresponding change of the reference-rescaling
identity. It is not achieved by the ultraviolet matching just proved.

The logarithmic source functional of V67 is unchanged: the determinant
difference has no logarithm, and its ordinary source germ is zero. The
present finite matrix contacts are therefore additional data, not extra
copies of the logarithmic counterterm. At zero sources both descriptions
have the same original action. All higher contacts of the chosen
comparison are derivatives of \eqref{eq:v68-sharp}, not arbitrary
eigenvalue-by-eigenvalue source renormalizations.

The scalar rescaling studied here acts simultaneously on a kinetic source
and the integration variable. It is a normalization identity for the
source-dependent scalar integral, not a physical local gauge symmetry.
Neither \eqref{eq:v68-restored-jet} nor positivity of $U_1$ proves the
inhomogeneous diffeomorphism/internal-gauge/chiral master equation. The
latter still requires its complete source and antifield array. A physical
preparation interpretation of these formal labels also remains subject
to the occurrence conditions audited in V67.

\section{Status and verification of the finite normalization step}
\label{sec:v68-ledgers}

\subsection*{Proved}
The active finite source completion has the exact additional ultraviolet
term \eqref{eq:v68-native-Ward}. An explicitly declared comparison removes
that term without changing the zero-source action or covariance. Its lower
reference agrees coefficientwise on the stated soft panels. The derivative-
free finite critical matrix array is evaluated by
\eqref{eq:v68-full-critical}--\eqref{eq:v68-upper-generating}; its
$\Sigma Y^2$ coefficient has a proved nonzero native ultraviolet part.
The ordered two-potential tower through arity five and its noncommuting
contact are explicit. One local finite matching functional gives the
$O(a^2)$ complementary and rooted estimates and the central-source jet
identity, with all lower-reference terms retained.

\subsection*{Inherited}
V14 supplies the distinct upper profiles. V27 supplies the actual corrected
scalar action, positivity bounds and momentum conventions. V67 supplies the
nonlinear matrix source definition and its logarithmic counterterm; these
are not rederived. The source-window Fourier topology and finite-volume
normalization are unchanged. No classical companion estimate is imported
into the quantum matching bound.

\subsection*{Literature input}
The broad distinction between a field-redefinition comparison and a
reference-preserving quantum symmetry is discussed in
\cite{V68CPV,V68IIM}. The new proof needs only finite determinant identities,
the supplied scalar symbols, Taylor estimates and Fourier summation.
No external anomaly/cohomology-vanishing theorem is used to infer a missing
primitive, and no heat-coefficient evaluation is claimed again.

\subsection*{Conditional}
Choosing the comparison as a normalization would mean adding exactly
$-M_a$, together with its source and metric derivatives, where those terms
are not already present in the common assignment. It is not installed here.
A different finite reference anchor changes the corresponding scalar
normalization identity. The profile moments are determined by the actual
profiles but have not been given numerical values. Other derivative/metric
critical coefficients are specified by convergent matching quadratures;
no full numerical array is claimed.

\subsection*{Open}
The complete bounded inhomogeneous critical BV/source primitive, its
internal-gauge and chiral sectors, and the common higher-loop/many-shell
assembly remain open. This installment resolves one finite, previously
unevaluated normalization block and its scalar calibration equation. It
does not convert that equation into the complete quantum descent.

\begin{center}\small
\begin{tabular}{>{\raggedright\arraybackslash}p{.22\textwidth}>{\raggedright\arraybackslash}p{.70\textwidth}}
\toprule
Variable & Scope of V68\\\midrule
Loop and fields & One closed real-scalar determinant; no integrated
metric, ghost, internal gauge or chiral line.\\
Source and derivative order & Total input arity at most five; at most two
potential marks for the displayed complementary theorem; one extra kinetic
mark is required for each central-shift identity.\\
Cutoff and volume & Uniform at fixed order, profiles, soft windows and
inverse margins. Upper matching has no new logarithm; power local terms
are retained. No total-volume multiplier in rooted source bounds.\\
Mass & Polynomial mass jets at the positive hard reference, including
$z=0$; physical massive propagators are not replaced before differentiation.\\
Normalization & Active V67 source family unchanged. The comparison is
explicit, agrees at $\Sigma=0$, and has its own finite matching contacts.\\
Verification & Finite rational operator tests and exact ordered source
series; moment positivity is an analytic proof. No numerical native moment,
complete BV array or proof-assistant certification.\\\bottomrule
\end{tabular}
\end{center}

\begin{thebibliography}{99}
\bibitem{V68CPV} J.~C.~Criado and M.~P\'erez-Victoria,
Field redefinitions in effective theories at higher orders,
\emph{JHEP} \textbf{03} (2019), 038; arXiv:1811.09413.
\bibitem{V68IIM} Y.~Igarashi, K.~Itoh and T.~R.~Morris,
BRST in the Exact RG, \emph{Prog. Theor. Exp. Phys.} (2019), 103B01;
arXiv:1904.08231.
\end{thebibliography}

\section*{Append-only continuation marker: V68}
\addcontentsline{toc}{section}{Append-only continuation marker: V68}
\label{sec:v68-continuation-marker}
\textbf{End of V68.} Preserve Sections 1--528. The finite kinetic-source
matching contains an explicitly nonzero upper-reference matrix contact.
Its complete prescribed scalar-rescaling comparison is now fixed. The
lower-reference contacts, the remaining finite source anchors, and the
complete ESM master problem are not deleted. No active normalization changes.


% ============================================================================
% V69 -- local frame closure, nonabelian reference identity and a populated row
% ============================================================================
\clearpage
\section{V69: from central rescaling to a closed local source operation}
\label{sec:v69-start}

V68 supplies finite matching for the constant central kinetic-source shift.
It does not show that the two-source family is closed under local matrix
changes of scalar variables. We address that upstream question before
attempting to identify a more general source Ward equation. The result is
an explicit missing first-order current coefficient, an exactly closed
finite source operation on the same scalar carrier, and a populated
nonzero lower-reference two-source coefficient. No new heat coefficient,
critical cohomology classification or interacting shell theorem is used.

The added operation is a \emph{declared external frame source}. It is not
identified with a physical preparation, with the internal Standard-Model
gauge group, or with a new integrated field. Its value is the identity at
zero frame source, so it changes none of the earlier action, measure,
propagator or finite normalization. All four real scalar components are
retained. The counterterm comparison of V68 remains a comparison only.
The relevance of complete source and reference transport under field
redefinitions is established in \cite{V69CPV,V69IIM}; the identities below
are derived directly for the finite operators used here.

Put $P=\widetilde P_a$, $\Theta=\Theta_\lambda$, $\nu=I-\Theta$ and
$C=\nu P^{-1}$ with the inherited hard-support convention. Thus
$PC=CP=\nu$, including at $z=m_H^2=0$, where no inverse of a retained
zero mode is taken. $Q$ denotes the actual symmetric scalar Hessian,
including the source placements of \eqref{eq:v67-finite-action}. For a
pointwise invertible real matrix $T(x)$ near $I_4$ set
\begin{equation}
 \boxed{\begin{aligned}
 Q_T&=T^TQT,\qquad A_T=(I+C(Q_T-P))^{-1},\qquad G_T=A_TC,\\
 \Gamma_T&=\frac12\Tr\log(I+C(Q_T-P)).
 \end{aligned}}
 \label{eq:v69-orbit}
\end{equation}
The transpose includes the coefficient pairing and site labels.
Multiplication by $T$ is at the two \emph{actual scalar endpoints}; it is
not moved through a difference or a smooth profile. Consequently this is
an exactly specified quasilocal finite source family, not a replacement
of the finite vertices by continuum rules. All following finite identities
hold on its invertible analytic germ, or coefficientwise to any prescribed
finite source order. No claim about a large-source Gaussian chart is made.

The loop factor $\hbar$ is removed from $\Gamma$. Normalization is at zero
sources, not separately at each value of $T$: a frame-dependent constant
is a source functional and cannot be discarded as a vacuum normalization.

\section{Why the kinetic and potential sources alone do not close}
\label{sec:v69-closure}

We first test closure on the ordinary flat limit, for which the source
family is known explicitly. At $\Sigma=Y=0$ it is $P_0=-\partial^2+z$.
A symmetric matrix test $X(x)$ has the congruence variation
\begin{equation}
 X P_0+P_0X
 =-2X\partial^2-2(\partial_\mu X)\partial_\mu
   -(\partial^2X)+2zX.
 \label{eq:v69-sym-tangent}
\end{equation}
This is represented at this point by
$\delta\Sigma=X$, $\delta Y=-\partial^2X$ in V67's family.

\begin{proposition}[A necessary local current direction]
\label{prop:v69-nonclosure}
The ordinary $(\Sigma,Y)$ family is not invariant under all local scalar
congruences, even though every symmetric congruence has the tangent
\eqref{eq:v69-sym-tangent} at its origin. The Lie closure of these
operations contains a skew first-order coefficient that neither source
can represent at that origin.
\end{proposition}
\begin{proof}
For fixed even matrix functions $X,Y$, the vector field
$\mathscr W_XQ=X^TQ+QX$ satisfies
$[\mathscr W_X,\mathscr W_Y]=\mathscr W_{[X,Y]}$ by multiplication.
If $X,Y$ are symmetric, $\Omega=[X,Y]$ is skew and
\begin{equation}
 \mathscr W_\Omega P_0=[P_0,\Omega]
   =-2(\partial_\mu\Omega)\partial_\mu-\partial^2\Omega.
 \label{eq:v69-skew-tangent}
\end{equation}
Its principal second-order variation is zero. A tangent within the old
family must therefore have $\delta\Sigma=0$, leaving only symmetric
zeroth-order $\delta Y$, which cannot reproduce its first-order term.
For a periodic two-flavour witness embedded into the four scalars take
\[
 X=\begin{pmatrix}1&0\\0&-1\end{pmatrix},\qquad
 Y=\sin x^1\begin{pmatrix}0&1\\1&0\end{pmatrix},\qquad
 J=\begin{pmatrix}0&1\\-1&0\end{pmatrix}.
\]
Then $\Omega=2\sin x^1J$ and the missing operator is
$-4\cos x^1J\partial_1+2\sin x^1J$. It is nonzero already at its
first-order coefficient at $x^1=0$. The claim is a failure of invariance,
not a claim that arbitrary chosen coordinate lifts on the old source
space have the same commutator without that invariance hypothesis.
\end{proof}

The necessary coefficient has the usual current interpretation:
$\int \phi^T J\partial_\mu\phi$ is not zero for skew $J$. At one point,
$\partial_\mu\Omega$ can prescribe any of the six skew-flavour matrices
independently for each of the four values of $\mu$. Thus a general
second-order quadratic-source operator must retain 24 such first-order
\emph{jet coefficients}. This is not a count of new physical particles,
independent preparation records, or an exhaustive quantum source algebra.

There is an explicit ordinary normal form for this necessary enlargement.
Let $B\succ0$ be symmetric. A self-adjoint operator with principal symbol
$\rho B^2g^{\mu\nu}\xi_\mu\xi_\nu$, in the flat coefficient pairing, has a
unique representation
\begin{equation}
 Q=\rho B\bigl[-(\nabla+\omega)^2+zI+U\bigr]B,
 \qquad \omega_\mu^T=-\omega_\mu,\quad U^T=U.
 \label{eq:v69-normal-form}
\end{equation}
Here $\rho=\sqrt{\det g}$. Uniqueness refers to fixed $B,g$.
Indeed the normalized operator has principal part $-\Delta_gI$;
self-adjointness makes its additional first-order coefficient skew, fixing
$\omega$, and the remaining multiplication operator fixes $U$.
V67 is the subfamily with its already derived $\omega_B$ and
$U=Y+\mathcal D^\mu H_\mu+H_\mu H^\mu$; these are not recalculated.

\begin{theorem}[Polar source transport without diagonalizing the matrices]
\label{thm:v69-polar}
For any pointwise real $T$ of positive determinant, define
\begin{equation}
 B_T=(T^TB^2T)^{1/2},\qquad O=BTB_T^{-1}.
 \label{eq:v69-polar}
\end{equation}
Then $O^TO=I$, and the normal form of $T^TQT$ has
\begin{equation}
 \boxed{\omega_{T,\mu}=O^T\omega_\mu O+O^T\partial_\mu O,
 \qquad U_T=O^TUO.}
 \label{eq:v69-polar-law}
\end{equation}
Its curvature is $\Omega_{T,\mu\nu}=O^T\Omega_{\mu\nu}O$.
\end{theorem}
\begin{proof}
The definition gives $T^TB= B_TO^T$ and $BT=OB_T$. Substitute these
into \eqref{eq:v69-normal-form} before differentiating. The identity
$O^T(\nabla+\omega)O=\nabla+\omega_T$ gives the result; squaring it
retains every derivative of $O$. Symmetry of $U_T$ and skewness of
$\omega_T$ follow from orthogonality. No simultaneous eigenbasis of
$B,U$ is used. All inverses are pointwise source-coordinate inverses.
\end{proof}

This ordinary normal form is explanatory; the finite source remains
\eqref{eq:v69-orbit}. In particular an independent continuum connection
is not discretized and then substituted for the existing finite scalar
operator. The integrated logarithmic density of V67, expressed in $U$
and $\Omega$, is invariant under \eqref{eq:v69-polar-law} by cyclicity of
the flavour trace. This consumes the previously evaluated heat density,
not a new logarithm calculation. Finite reference terms need not vanish.

\section{The finite local-frame Ward equation and its consistency}
\label{sec:v69-Ward}

Let $\mathscr W_X$ now mean differentiation of $T\mapsto T e^{tX}$ at
$t=0$, with $X(x)$ any fixed real matrix test. It acts on $Q_T$ but not
on $P,C,\Theta$ or the scalar coefficient measure.

\begin{theorem}[Exact finite nonabelian reference identity]
\label{thm:v69-Ward}
On the actual finite source family,
\begin{equation}
 \boxed{\mathscr W_X\Gamma_T=\Tr X-\mathfrak R_T(X),
 \qquad \mathfrak R_T(X)=\Tr(XA_T\Theta).}
 \label{eq:v69-Ward}
\end{equation}
Moreover,
\begin{equation}
 \boxed{\mathscr W_X\mathfrak R_T(Y)
       -\mathscr W_Y\mathfrak R_T(X)
       =\mathfrak R_T([X,Y]).}
 \label{eq:v69-consistency}
\end{equation}
These equations do not require $X$ to commute with $C$, with the other
source matrices, or with spatial translations.
\end{theorem}
\begin{proof}
The resolvent $G_T=A_TC=CA_T^T$ is symmetric. Therefore finite determinant
differentiation gives
\[
 \mathscr W_X\Gamma_T
 =\tfrac12\Tr\bigl[G_T(X^TQ_T+Q_TX)\bigr]
 =\Tr(XG_TQ_T).
\]
Using $CP=\nu$ yields $G_TQ_T=I-A_T\Theta$, which proves
\eqref{eq:v69-Ward}. The source vector fields obey
$[\mathscr W_X,\mathscr W_Y]=\mathscr W_{[X,Y]}$; apply this identity
to the same finite determinant and use $\Tr[X,Y]=0$. Alternatively,
insert $\mathscr W_XA_T=-G_T(X^TQ_T+Q_TX)A_T$ in both sides of
\eqref{eq:v69-consistency}. All factors keep their displayed order.
\end{proof}

$\Tr X$ is the finite scalar change-of-variables Jacobian. It includes
site and flavour indices and is not an additional auxiliary density. The
reference term is not declared zero. In a full change of variables the
scalar field and antifield transform oppositely,
$\phi\mapsto T^{-1}\phi$, $\upsilon\mapsto T^T\upsilon$; a linear
source must also transform on its actual side. The scalar antibracket is
preserved. A fixed hard covariance, however, is not invariant under an
arbitrary local $T$, which is why \eqref{eq:v69-Ward} retains its reference
term. The operation is a source-reparametrization identity, not the
unproved diffeomorphism/internal-gauge master equation.

\section{A complete two-frame loop collapses to the lower reference}
\label{sec:v69-two-frame}

We populate one new source row rather than treating the reference term
as an unspecified functional. Set the old metric and added sources to
zero, so $Q=P$. Take a real skew matrix $J$ and a real scalar test $\omega$,
and use $T=e^{t\omega J}$. Set $\kappa_J=-\operatorname{tr}J^2>0$ for
$J\ne0$. There is no separate $n_H$ multiplier: $\kappa_J$ is the actual
flavour contraction.

\begin{theorem}[Native two-frame coefficient with its contact vertex]
\label{thm:v69-two-frame}
Let $\Omega$ be the multiplication operator $\omega J$. Then
\begin{equation}
 \boxed{\left.\frac{d^2\Gamma_{e^{t\Omega}}}{dt^2}\right|_0
 =\frac12\Tr\left(C[[P,\Omega],\Omega]
             -C[P,\Omega]C[P,\Omega]\right).}
 \label{eq:v69-Hessian}
\end{equation}
For lattice Fourier arguments, its per-volume momentum kernel, with the
flavour factor removed, is
\begin{equation}
 \boxed{\mathcal K_{a,L}(k,z)
 =\frac1{V_L}\sum_p
       \nu(p)\Theta(p+k)\,
       \frac{P(p+k)-P(p)}{P(p)}.}
 \label{eq:v69-kernel}
\end{equation}
The quotient is zero where the inherited covariance is zero. With Fourier
coefficients $\omega_k$ normalized by $\int_a1=V_L$, the second variation
per physical volume equals
$\kappa_J\sum_k\omega_{-k}\mathcal K_{a,L}(k,z)\omega_k$.
\end{theorem}
\begin{proof}
$Q'(0)=[P,\Omega]$ and $Q''(0)=[[P,\Omega],\Omega]$ follow from the
actual endpoint conjugation. Differentiation of one half log determinant
gives \eqref{eq:v69-Hessian}, including its one-vertex contact.
Diagonalize the translation-invariant $P,C,\Theta$, writing their scalar
symbols as $P_i,C_i,\Theta_i$. Skew adjointness of $\Omega$ gives
\[
 \Gamma''(0)=-\frac12\sum_{i,j}
 \left[(C_i-C_j)(P_i-P_j)+C_iC_j(P_i-P_j)^2\right]
 |\Omega_{ij}|^2.
\]
The coefficient in brackets is exactly
\begin{equation}
 (P_i-P_j)(C_i\Theta_j-C_j\Theta_i),
 \label{eq:v69-cancellation}
\end{equation}
because $P_iC_i=\nu_i$ and $\Theta_i=1-\nu_i$. Symmetrize the two
summands before taking bounds, route $j=p+k$, and contract the flavour
indices. This gives \eqref{eq:v69-kernel}. The algebra is valid when a
$P_i$ vanishes because it was performed with $C_i$, not its inverse.
\end{proof}

This is the complete one-loop coefficient on this two-frame orbit, not
only its two-vertex part. For $\nu=I$ and an invertible full Gaussian the
answer vanishes, as required by $\det(e^{t\Omega})=1$. At the fixed lower
reference it need not vanish. For $|k|\le\mu<\lambda/2$, every nonzero
term of \eqref{eq:v69-kernel} has
\begin{equation}
 \lambda\le|p|\le2\lambda+\mu.
 \label{eq:v69-annulus}
\end{equation}
Hence it has no ultraviolet tail after the signed cancellation. The
periodic continuation used for momentum derivatives is defined by the
right-hand side of \eqref{eq:v69-kernel}; equality to the finite Fourier
functional is asserted at its actual lattice source momenta.

\section{A strictly nonzero local frame-response coefficient}
\label{sec:v69-negative}

At zero mesh put $P_0(p)=p^2+z$ and use the same lower profile, not a newly
chosen one. In thermodynamic normalization,
\begin{equation}
 \mathcal K_0(k,z)=\int\frac{d^4p}{(2\pi)^4}
 \frac{\nu_\lambda(p)\Theta_\lambda(p+k)
       (2p\cdot k+k^2)}{p^2+z}.
 \label{eq:v69-continuum-kernel}
\end{equation}
It is even in $k$ and zero at $k=0$ if the inherited profile is even.
Let $\mathcal T_{ij}(z)=\tfrac12\partial_{k_i}\partial_{k_j}
\mathcal K_0(0,z)$.

\begin{theorem}[A nonzero reference tensor without rotational averaging]
\label{thm:v69-negative}
For $z\ge0$,
\begin{equation}
 \mathcal T_{ij}(z)=\delta_{ij}\int_p\frac{\nu\Theta}{p^2+z}
 +\int_p\frac{\nu(p_i\partial_j\Theta+p_j\partial_i\Theta)}{p^2+z},
 \label{eq:v69-Tij}
\end{equation}
and its trace is
\begin{equation}
 \boxed{\operatorname{tr}\mathcal T(z)
 =-2\int_p\frac{p^2\Theta(p)^2+2z\Theta(p)}{(p^2+z)^2}<0.}
 \label{eq:v69-negative-mass}
\end{equation}
In particular, at $z=0$ and $\Theta_\lambda(p)=\theta(p/\lambda)$,
\begin{equation}
 \boxed{\frac14\operatorname{tr}\mathcal T(0)
 =-\frac{\lambda^2}{2}\,b_\theta,
 \qquad b_\theta=\int\frac{d^4y}{(2\pi)^4}\frac{\theta(y)^2}{|y|^2}
 \ge\frac1{16\pi^2}.}
 \label{eq:v69-negative-coefficient}
\end{equation}
No radial assumption is needed for this trace identity. If the profile is
radial, $\mathcal T_{ij}(0)=-\lambda^2b_\theta\delta_{ij}/2$.
\end{theorem}
\begin{proof}
Differentiate \eqref{eq:v69-continuum-kernel} under its compact annular
integral to obtain \eqref{eq:v69-Tij}. Set
$F=\Theta-\Theta^2/2$, so that $\nu\partial_i\Theta=\partial_iF$.
The trace is $4\int\nu\Theta/(p^2+z)+2\int p\cdot\partial F/(p^2+z)$.
Integration by parts uses
\[
 \partial_{p_i}\frac{p_i}{p^2+z}
 =\frac4{p^2+z}-\frac{2p^2}{(p^2+z)^2}.
\]
Collecting terms gives \eqref{eq:v69-negative-mass}. At $z=0$, the inner
sphere boundary term is $O(\varepsilon^2)$ and vanishes; no delta contact
is omitted. The unit plateau gives
$\int_{|y|\le1}d^4y/((2\pi)^4|y|^2)=1/(16\pi^2)$.
For nonradial profiles only the trace, not every eigenvalue, has the proved
sign. Rotation invariance in the radial subcase gives the last formula.
\end{proof}

For the $2\times2$ generator $J$ of Proposition~\ref{prop:v69-nonclosure},
$\kappa_J=2$. Thus the new row is strictly nonzero in the source family
forced by local closure. It is a lower-reference quadratic frame contact,
not a physical scalar instability, anomaly or anomalous dimension. Global
constant frames still give $\mathcal K(0,z)=0$ exactly. Numerical values
of $b_\theta$ for the native profile have not been evaluated.

The negative coefficient is compatible with positivity of the finite
scalar quadratic form on its small chart: the frame is an external
coordinate, and the fixed hard split is not invariant under a local
orthogonal change of that coordinate. Renormalizing this coefficient to
zero would require a new finite source prescription and the corresponding
change of \eqref{eq:v69-Ward}; it is not achieved by declaring the full
Gaussian Jacobian to be one.

\section{Cutoff and source estimates after the finite cancellation}
\label{sec:v69-estimates}

We prove an analytic comparison only for the two-frame row just populated.
The all-source Ward equation is exact algebra but is not thereby an
all-source quasilocal estimate. Fix the inherited domain
\begin{equation}
 \lambda\ge32768\mu>0,\quad a\lambda\le2^{-24},\quad
 \lambda L_\nu\ge1,\quad |k|\le\mu,\quad z=0.
 \label{eq:v69-domain}
\end{equation}
Massive identities above are exact; the estimates stated here are the
massless ones. Set
$\mathcal K_{a,L}^R=(1-\mathsf T_{4,k})\mathcal K_{a,L}$, retaining the
actual quadratic and quartic coefficients as local source data.

\begin{theorem}[A fourth-order mesh comparison for the reference row]
\label{thm:v69-estimates}
For $j=|\alpha|\le6$,
\begin{equation}
 \boxed{\left|\partial_k^\alpha
 (\mathcal K_{a,L}^R-\mathcal K_{0,L}^R)(k)\right|
 \le C_\alpha a^4\lambda^2|k|^{6-j}.}
 \label{eq:v69-cutoff}
\end{equation}
For every $b>4$ its separate ordinary volume comparison is
\begin{equation}
 \boxed{\left|\partial_k^\alpha
 (\mathcal K_{0,L}^R-\mathcal K_{0,\infty}^R)(k)\right|
 \le C_{\alpha,b}\lambda^{-2}(\lambda L_{\min})^{-b}|k|^{6-j}.}
 \label{eq:v69-volume}
\end{equation}
Every fixed higher external derivative is bounded in the associated
scaled symbol norms. The difference of quadratic Taylor tensors obeys
\begin{equation}
 \norm{\mathcal T_{a,L}-\mathcal T_{0,\infty}}
 \le C_b\left[a^4\lambda^6+
                 \lambda^2(\lambda L_{\min})^{-b}\right].
 \label{eq:v69-local-error}
\end{equation}
\end{theorem}
\begin{proof}
Use the assembled formula \eqref{eq:v69-kernel}, not separate absolute
bounds on the two terms of \eqref{eq:v69-Hessian}. Its support is
\eqref{eq:v69-annulus}. In this annulus the native free scalar operator is
exactly the interior corrected symbol, with
$P_a-P_0=O(a^4|p|^6)$ and the corresponding differentiated estimates.
No oriented gravitational vertex is used. The symbol ratio in
\eqref{eq:v69-kernel} therefore differs by
$O(a^4\lambda^{4-j})$ after $j$ external derivatives. Its normalized
four-dimensional sum has size $O(\lambda^4)$, uniformly when
$\lambda L_\nu\ge1$. Thus the unsubtracted difference has derivative
bound $Ca^4\lambda^{8-j}$. It is even and zero at the origin. After its
fourth Taylor jet is removed, the fifth jet also vanishes. The sixth
order Taylor integral proves \eqref{eq:v69-cutoff}. The same estimates
for $j>6$ give the asserted scaled norms.

The ordinary integrand is smooth with compact support away from the
origin. Differentiated Fourier decay followed by Poisson summation gives
$C_b\lambda^{4-j}(\lambda L_{\min})^{-b}$; the image-lattice sum
converges for $b>4$. Apply the same Taylor remainder to obtain
\eqref{eq:v69-volume}. For \eqref{eq:v69-local-error}, take exactly two
external derivatives before evaluating at zero. The possible retained
zero scalar mode never occurs in the annulus.
\end{proof}

In particular the strictly negative trace is stable once the displayed
error is smaller than its continuum margin. No numerical refinement
threshold is asserted without the profile constants. For a fixed smooth
window on each external frame mark, the usual Fourier integration by
parts gives the rooted estimate
\begin{equation}
 \boxed{a^4\sum_x(1+\mu d_{\rm per}(x,0))^b
 |\mathcal K^{\Delta,R}(x)|
 \le C_b a^4\lambda^2\mu^6.}
 \label{eq:v69-rooted}
\end{equation}
Polarization inserts the factor
$|\operatorname{tr}(J_1J_2)|\le\norm{J_1}_{\rm HS}\norm{J_2}_{\rm HS}$.
One frame test can be in $L^1$ and the other in $L^\infty$, without a
volume multiplier. The analogous rooted ordinary-volume bound has factor
$\lambda^{-2}(\lambda L_{\min})^{-b'}\mu^6$, with sufficiently large
$b'$. Constants depend on fixed derivative and window orders and profiles,
not on mesh or periods. This is a reference-supported one-loop row, not
a many-shell interacting estimate.

\section{V68 matching extends to the full finite frame orbit}
\label{sec:v69-matching}

Let $Q$ and $Q^\sharp=Q+BD_aB-D_a$ be precisely the native and comparison
source Hessians in \eqref{eq:v68-sharp}; do not choose new finite
coefficients. Apply the same frame $T$ to both. Write
$\mathcal D_T=\Gamma_T-\Gamma_T^\sharp$ and use total source/metric/frame
arity $N$, with every external momentum at most $\mu$.

\begin{theorem}[No additional ultraviolet matching on the frame orbit]
\label{thm:v69-matching}
If $2\lambda+N\mu<1/(64a)$, every Taylor coefficient of
$\mathcal D_T-\mathcal D_I$ through total arity $N$ is zero on the
specified Fourier panels:
\begin{equation}
 \boxed{[\mathcal D_T]_{\le N}=[\mathcal D_I]_{\le N}.}
 \label{eq:v69-frame-independence}
\end{equation}
Equivalently, every such coefficient containing at least one frame mark
vanishes in the native-minus-comparison difference. This is not only a
zero-momentum Taylor-jet statement. It holds for all allowed soft momenta
in that finite source bank. In particular V68's single matching
functional $M_a$ is sufficient on this enlarged orbit; no new independent
frame matching coefficients are chosen.
\end{theorem}
\begin{proof}
The exact difference of Ward equations is
\begin{equation}
 \mathscr W_X\mathcal D_T
 =-\Tr\bigl[X(A_T-A_T^\sharp)\Theta\bigr].
 \label{eq:v69-difference-Ward}
\end{equation}
A resolvent identity expands the difference through any fixed source
arity into finitely many words, each containing a differentiated
$T^T(BD_aB-D_a)T$ and a factor $\Theta$. Every nonzero differentiated
word has one $D_a$ on a routed scalar momentum. Start the trace at its
$\Theta$ factor, whose momentum has size at most $2\lambda$. All other
momenta differ by partial sums of the source momenta, including the
extra $X$ mark. Their total count is at most $N$. They cannot reach the
support of $D_a$ under the displayed hypothesis. Each word vanishes.

Parametrize the local frame germ by $T=e^H$ and differentiate along
$t\mapsto e^{tH}$. At each finite source order the right-hand side is
zero by the preceding argument. Integrate this coefficient identity in
$t$ from zero to one. This proves \eqref{eq:v69-frame-independence}; it
requires no convergence in unbounded source order. The normalization at
zero old sources cancels between $Q$ and $Q^\sharp$, since they coincide
there at every $T$. The local matching $M_a=\mathbf T\mathcal D_I$ is
therefore independent of $T$ on the specified bank, as claimed.
\end{proof}

This extends the comparison to noncommuting local frame operations while
retaining the original kinetic/potential labels. It does not state that
$M_a$ is separately diffeomorphism invariant. It is invariant under the
\emph{added frame coordinate}, on this prescribed orbit and finite soft
bank. The intrinsic $(\Sigma,Y)$ transformations need not close without
that coordinate, by Proposition~\ref{prop:v69-nonclosure}.

The new reference response of Section~\ref{sec:v69-negative} is common to
both completions at zero old sources and is \emph{not} removed by this
ultraviolet matching. Thus the two conclusions are compatible: the frame
matching difference is zero, whereas a frame-dependent lower-reference
coefficient is strictly nonzero. Dropping the latter would change the
same exact Ward equation that enabled the matching proof.

\section{Status, verification and the remaining combined symmetry problem}
\label{sec:v69-ledgers}

\subsection{Proved}
The kinetic/potential family has an explicit local congruence-closure
failure. Its necessary skew current direction and ordinary polar normal
form are given without a commuting-matrix assumption. The declared finite
frame orbit obeys one exact nonabelian reference Ward identity and its
consistency equation. The complete two-frame loop, including its
second-order source contact, reduces to a lower annulus. Its quadratic
Taylor tensor has a strictly negative nonzero trace, and its subtracted
coefficient has fourth-order mesh, independent volume, and rooted source
bounds. Finally the original V68 ultraviolet matching is independent of
the added local frame coordinate through each declared soft arity, with
no separately fitted current matching tensor.

\subsection{Inherited}
V27 supplies the same completed scalar kinetic operator and its interior
fourth-order comparison. V67 supplies the old matrix source action,
its ordinary normal form and logarithmic functional. V68 supplies $D_a$,
its support, the specified comparison Hessian and its matching functional.
The ordinary field-redefinition interpretation does not replace any of
these finite inputs. No classical companion estimate is used as a
quantum source bound.

\subsection{Literature input}
Source and reference transport under field redefinitions and cutoff
modified identities are comparison material \cite{V69CPV,V69IIM}.
All determinant, polar, contact, sign and support identities used in this
installment are proved here. The cutoff proof consumes the previously
stated native symbol estimates, not an external continuum-renormalizability
theorem. No new anomaly/cohomology theorem is claimed.

\subsection{Conditional}
The frame operation is a formal composite-source extension on the existing
carrier. Physical occurrence of those source operations is not supplied.
Adopting $-M_a$ is still V68's explicit normalization choice and has not
been performed here. Its agreement on frame-marked panels does not
establish agreement of arbitrary separately specified external current
sources outside the orbit. Positivity of a native Gaussian is used only
on its declared small chart; the finite algebra also makes sense
coefficientwise. Numerical values of the profile moments remain unevaluated.

\subsection{Open}
The full common inhomogeneous BV/source primitive, diffeomorphism and
internal-gauge/chiral compatibility, and the source-complete interacting
many-shell construction remain open. The present local-frame Ward equation
must not be relabelled as that combined master equation. The next source
comparison must retain the first-order current directions just identified;
central rescaling alone cannot supply their algebra. A general physical
gauge-connection source sector outside the congruence orbit requires its
own native vertices, counterterms, and uniform estimates.

\begin{center}\footnotesize
\begin{tabular}{>{\raggedright\arraybackslash}p{.22\textwidth}>{\raggedright\arraybackslash}p{.70\textwidth}}
\toprule
Variable & Scope of V69\\\midrule
Loop/fields & One real-scalar determinant; no integrated graviton, ghost,
internal gauge, or chiral line. The displayed flavour norm counts its
actual active scalar directions.\\
Source order & Finite frame identity exact on its matrix germ. Comparison
through every stated finite total arity satisfying the support inequality;
new analytic estimates concern two frame marks at zero old sources.\\
Cutoff/volume & Fourth-order mesh comparison and independent Poisson
volume estimate for that lower-reference row; no total-volume loss in
rooted bounds. No unbounded-order or many-step constant is asserted.\\
Mass & Exact two-frame and trace formulas for $z\ge0$; new quantitative
cutoff theorem stated at $z=0$. No inverse retained zero mode.\\
Normalization & Zero-frame source family, active counterterms, covariance,
contour and once-counted density unchanged. No local coefficient is
silently deleted.\\
Verification & Exact finite matrix, polynomial, Lie-bracket and contraction
tests. Analytic support and estimates are proofs, not conclusions of
numerical sampling. No numerical native profile integral.\\\bottomrule
\end{tabular}
\end{center}

\begin{thebibliography}{99}
\bibitem{V69CPV} J.~C.~Criado and M.~P\'erez-Victoria,
Field redefinitions in effective theories at higher orders,
\emph{JHEP} \textbf{03} (2019), 038; arXiv:1811.09413.
\bibitem{V69IIM} Y.~Igarashi, K.~Itoh and T.~R.~Morris,
BRST in the Exact RG, \emph{Prog. Theor. Exp. Phys.} (2019), 103B01;
arXiv:1904.08231.
\end{thebibliography}

\section*{Append-only continuation marker: V69}
\addcontentsline{toc}{section}{Append-only continuation marker: V69}
\label{sec:v69-continuation-marker}
\textbf{End of V69.} Preserve Sections 1--536. The local-frame source and
reference identities, populated two-frame response, and V68 matching
extension hold at their stated scope. The complete interacting ESM master
problem remains open. Active normalizations are unchanged.



\clearpage
\part*{V70: the first Higgs-quartic insertion on the finite source orbit}
\addcontentsline{toc}{part}{V70: the first Higgs-quartic insertion on the finite source orbit}

\section{From a quadratic source identity to an interacting Higgs coefficient}
\label{sec:v70-start}

Sections 1--536 are retained. V69 closes the matrix-frame identity for a
quadratic scalar integral. Repeating that determinant calculation would not
control the first interaction insertion. We instead turn on the coefficient
of the local Higgs quartic already specified in
\eqref{eq:v9-Higgs-Hessian}, at flat retained metric and zero internal
connection. The result consumes the existing Gaussian and frame formulas
but computes a different perturbative component: one quartic interaction,
its one-loop mass insertion, and its two-loop vacuum/source coefficient.
The potentially divergent products are cancelled at the finite cutoff
before any estimates are taken.

The quartic convention in this installment is
\begin{equation}
 V_4(\phi)=g_H\int_a(\phi^T\phi)^2,
 \qquad n=4,
 \qquad \int_a=a^4\sum_x.
 \label{eq:v70-quartic-convention}
\end{equation}
Here $g_H$ is the coefficient in this real-field normalization. It is not
identified with a differently normalized textbook quartic coupling. The
quadratic part of $g_H(\phi^T\phi-v_H^2)^2$ belongs to the mass parameter;
the calculation isolates the quartic coefficient and does not reset a
prescribed physical vacuum scale. There is no integrated graviton, gauge
boson, ghost, or fermion in the new diagrams. The other couplings are zero
in this component. No gravitational propagator, and hence no extra
$\chi^{-1}$, appears.

We use the exact scalar reference of V27 and the exact endpoint frame of
V69. Write
\begin{equation}
 \begin{gathered}
 v=a^4,\quad P=\widetilde P_a(z),\quad z=m_H^2,\quad
 \Theta=\Theta_\lambda,\quad \nu=I-\Theta,\quad C=\nu P^{-1},\\
 PC=CP=\nu,\qquad
 Q_T=T^TPT,\quad A_T=(I+C(Q_T-P))^{-1},\quad G_T=A_TC.
 \end{gathered}
 \label{eq:v70-reference}
\end{equation}
The covariance is zero on the retained modes. There is no inverse of a
retained zero mode at $z=0$. Each $T(x)$ is an invertible real matrix near
$I_n$; its transpose in \eqref{eq:v70-reference} includes site labels.
The ordinary operator $P_0=p^2+z$ is used only for the comparison limit.
The finite Gaussian still uses the continued, corrected, periodically
completed operator, not a trigonometric replacement.

The source operation acts on the \emph{entire} local potential:
$\phi(x)\mapsto y(x)=T(x)\phi(x)$, not just on the kinetic Hessian.
This is V69's declared source reparametrization extended to one already
present interaction. It is not an assertion of physical occurrence of a
new preparation or an independent gauge connection. The frame-free action
at tree level is unchanged.

\subsection*{The normalization selected for this component}
Let
\begin{equation}
 u_a(z)=\frac1{V_L}\sum_{p\in\mathrm{BZ}_{a,L}}
                     \frac{\nu(p)}{P(p,z)}.
 \label{eq:v70-u}
\end{equation}
The coordinate covariance of one scalar is $\hbar C/v$, so $u_a$ is its
coincident variance with $\hbar$ removed. V57's bounds give
$u_a\asymp a^{-2}$ on the stated fixed-lower-scale domain. We select the
following explicit \emph{normal-ordering comparison} in the coefficient
linear in $g_H$:
\begin{equation}
 \boxed{\mathcal V^{\mathrm{NO}}_{a,T}
 =\int_a\left[
 (y^Ty)^2-2\hbar(n+2)u_a\,y^Ty
       +\hbar^2 n(n+2)u_a^2\right].}
 \label{eq:v70-NO-action}
\end{equation}
The coefficient $g_H$ is outside this display. Its mass and vacuum pieces
are thus fixed by the \emph{same} local functional. For $n=4$ the mass
coefficient in the convention $\frac12\delta m^2\phi^T\phi$ is
$\delta m^2=-24g_H\hbar u_a$, and the vacuum insertion is
$24g_H\hbar^2u_a^2$.

This is a newly specified comparison in the $g_H$ component, not an
additional term to count again if the full common $S_1,S_2$ already
contain those assignments. V34's $g_H=0$ scalar normalization and V38's
metric-loop packet are not altered. The quantum subtractions are not
resummed into $P$. At fixed mass their coefficients are functions of
$z$; they are generally not polynomials in $z$. A prescribed
mass-polynomial EFT jet is obtained by Taylor-extracting the \emph{whole}
identity below, including both powers of $u_a(z)$, before truncation.
Replacing just one occurrence of $u_a$ by an unrelated mass jet is not
part of the scheme.

Normal products, finite renormalizations of Wick powers, and their
background dependence are established subjects \cite{V70KMM}. What is
proved here is a finite native realization of this subtraction, its
interacting frame coefficient, and its cutoff estimate. No Wick-power
classification or anomaly theorem is substituted for the calculation.

\section{One common subtraction evaluates the complete marked quartic insertion}
\label{sec:v70-Wick}

For a retained scalar $\varphi$, V58's normalized conditional Gaussian has
mean $m=A_T\varphi$ and covariance $\hbar G_T/v$. Define at each root
\begin{equation}
 \begin{split}
 M(x)&=(T A_T\varphi)(x),\\
 K_T(x)&=v^{-1}[T G_T T^T]_{xx},\qquad
 D_T(x)=K_T(x)-u_a I_n.
 \end{split}
 \label{eq:v70-dressed}
\end{equation}
The letter $D_T$ here denotes a covariance difference, not a lattice
derivative or the ultraviolet source mismatch $D_a$ in V68. Those
objects are not used interchangeably. A fixed even linear scalar source
can be included by replacing $m$ with $A_T\varphi-G_T\ell$; all formulas
remain polynomial identities in its ordered even coefficients. No
complete antifield master identity is inferred from that substitution.

\begin{theorem}[Exact first-quartic coefficient with its source partners]
\label{thm:v70-Wick}
The coefficient linear in $g_H$ of the normalized effective action
$-\hbar\log Z$, with insertion \eqref{eq:v70-NO-action}, is
\begin{equation}
 \boxed{\begin{aligned}
 \Gamma_{a}^{[g_H]}[T,\varphi]
 =g_H\int_a\Big\{&(M^TM)^2
  +2\hbar\operatorname{tr}D_T\,M^TM
  +4\hbar M^TD_TM\\
  &+\hbar^2\big[(\operatorname{tr}D_T)^2
                         +2\operatorname{tr}(D_T^2)\big]\Big\}.
 \end{aligned}}
 \label{eq:v70-Wick-result}
\end{equation}
This is the entire coefficient linear in the quartic coupling, through
its maximum of two scalar self-contractions, at every fixed number of
frame marks. The trace is over the actual $n$ real scalar components.
There is no further flavour multiplicity outside it.
\end{theorem}
\begin{proof}
For an $n$-component Gaussian with mean $M$ and covariance $\hbar K$,
finite differentiation of its generating exponential gives
\begin{align*}
 \E(y^Ty)&=M^TM+\hbar\operatorname{tr}K,\\
 \E((y^Ty)^2)&=(M^TM)^2
       +2\hbar\operatorname{tr}K\,M^TM+4\hbar M^TKM\\
 &\hspace{12mm}+\hbar^2\big[(\operatorname{tr}K)^2
                                      +2\operatorname{tr}(K^2)\big].
\end{align*}
Substitute $K=u_a I_n+D_T$ and the three terms of
\eqref{eq:v70-NO-action}. The $u_aM^TM$ terms cancel, the
$u_a\operatorname{tr}D_T$ terms cancel, and the constant $u_a^2$ terms
cancel. What remains is \eqref{eq:v70-Wick-result}. All operations are
finite polynomial contractions on the existing carrier. Differentiating
$-\hbar\log Z$ with respect to $g_H$ at zero coupling gives the normalized
Gaussian expectation of the insertion, with positive sign. At
$T=I,\varphi=0$ that expectation is zero. Hence normalizing the partition
function at that same frame-free point introduces no additional
order-$g_H$ constant. A single quartic vertex has at most two scalar
contractions, so the displayed expansion is complete at this coupling
order. The frame vertices are already resummed in $G_T$ and $M$.
\end{proof}

In particular the vacuum/source component is
\begin{equation}
 \boxed{\Gamma_a^{[g_H\hbar^2]}[T,0]
      =g_H\hbar^2\mathcal R_a[T],\qquad
 \mathcal R_a[T]=\int_a\left[
       (\operatorname{tr}D_T)^2+2\operatorname{tr}(D_T^2)\right].}
 \label{eq:v70-R}
\end{equation}
It contains the quartic two-loop graph, the one-loop insertion of the
order-$\hbar$ mass counterterm, and the order-$\hbar^2$ vacuum counterterm.
It is not merely the bare double contraction. On a real positive scalar
chart $D_T$ is symmetric, so $\mathcal R_a[T]\ge0$. This is a statement
about this coefficient after its specified subtractions, not positivity
of the full signed gravitational perturbation theory.

\begin{proposition}[Vacuum subtraction alone fails at the next frame mark]
\label{prop:v70-tadpole-required}
At a constant central frame let $D_T=dI_n$. Removing only the frame-free
bare vacuum constant leaves
\begin{equation}
 n(n+2)\bigl(2u_a d+d^2\bigr)
 \label{eq:v70-unpaid}
\end{equation}
instead of $n(n+2)d^2$. If $d$ tends to a nonzero finite
lower-reference coefficient, the first term diverges like $a^{-2}$.
The mass insertion in \eqref{eq:v70-NO-action} cancels exactly that term.
\end{proposition}
\begin{proof}
The bare expectation is $n(n+2)(u_a+d)^2$. Subtract its value at $d=0$
and use $u_a\asymp a^{-2}$. The mass insertion is
$-2n(n+2)u_a(u_a+d)$ and the chosen vacuum insertion is
$n(n+2)u_a^2$. Their sum with the bare term is $n(n+2)d^2$ exactly.
No bound on an isolated tadpole is used in place of this identity.
\end{proof}

\section{The dressed covariance difference is a lower-reference object}
\label{sec:v70-annulus}

The Gaussian covariance $G_T$ itself still has ultraviolet modes. The
object appearing after the complete subtraction is instead $D_T$.
The following identity locates its differentiated coefficients without
replacing $P$ or treating the frame as commuting with it.

\begin{theorem}[Exact ordered frame derivative]
\label{thm:v70-covariance-Ward}
Let $\mathscr W_X$ differentiate $T\mapsto Te^{tX}$. Then
\begin{equation}
 \boxed{\mathscr W_X(TG_TT^T)
 =T\left[A_T\Theta XG_T+G_TX^T\Theta A_T^T\right]T^T.}
 \label{eq:v70-covariance-Ward}
\end{equation}
The identity is valid at zero mass with the inherited hard-support
convention. Every positive-order frame coefficient of $D_T$ has a
finite representation by ordered words containing at least one actual
$\Theta$ and at least one hard covariance $C$.
\end{theorem}
\begin{proof}
Put $B_X=X^TQ_T+Q_TX$. Resolvent differentiation gives
$\mathscr W_XG_T=-G_TB_XG_T$. Therefore the left side of
\eqref{eq:v70-covariance-Ward}, with the outer $T,T^T$ removed, is
\[
 XG_T+G_TX^T-G_TX^TQ_TG_T-G_TQ_TXG_T.
\]
Use the finite identities
\[
 G_TQ_T=I-A_T\Theta,\qquad
 Q_TG_T=I-\Theta A_T^T.
\]
They reduce this expression to the right side in the stated matrix
order. The covariance $G_T=A_TC$ retains a primitive $C$. For a frame
germ $T=e^H$, integrate \eqref{eq:v70-covariance-Ward} along
$e^{tH}$, $0\le t\le1$, and subtract the value $C$ at the identity.
Expansion to any prescribed source order terminates in a finite set of
resolvent words. Every word retains the displayed $\Theta$ and at least
one $C$. This is a coefficientwise argument, not an assertion of
convergence in unbounded source order.
\end{proof}

\begin{lemma}[All positive-order coincidence coefficients have annular support]
\label{lem:v70-support}
Give every frame mark Fourier support $|k_i|\le\mu$, and retain at most
$N$ external marks. If $N\mu\le\lambda/2$, every loop momentum in the
representation of a positive-order derivative of $D_T(x)$ can be routed
within
\begin{equation}
 \lambda/2\le |p|\le 3\lambda.
 \label{eq:v70-annulus}
\end{equation}
The statement concerns the assembled words of
\eqref{eq:v70-covariance-Ward}, not the separate terms obtained by
expanding $TG_TT^T-C$ before the cancellation.
\end{lemma}
\begin{proof}
Starting at the $\Theta$ factor gives one momentum of norm at most
$2\lambda$. Nonzero $C$ requires one momentum outside the lower plateau,
of norm at least $\lambda$. All routed momenta differ by partial sums
of the at most $N$ external momenta. Their total possible displacement
is at most $N\mu$. This gives the lower and upper bounds simultaneously.
Taking the coincidence kernel closes the scalar momentum sum but does
not remove either support factor. Derivatives of a cutoff profile are
supported in its transition and satisfy the same bounds.
\end{proof}

At first order this word formula is particularly simple. Write
$T=e^{\epsilon X(x)}$. Translation invariance of $P,C,\Theta$ at the
source origin gives
\begin{equation}
 \boxed{D_T^{[X]}(x)
   =\sum_k e^{ikx}\,\beta_{a,L}(k,z)
                          \bigl[X(k)+X(k)^T\bigr],}
 \label{eq:v70-first-D}
\end{equation}
where the Fourier coefficients are normalized by $\int_a1=V_L$, and
\begin{equation}
 \boxed{\beta_{a,L}(k,z)
 =\frac1{V_L}\sum_p
       \frac{\nu(p)\Theta(p+k)}{P(p,z)}.}
 \label{eq:v70-beta}
\end{equation}
It is even in $k$ on the actual symmetric Fourier lattice. In particular
all pointwise skew frame directions have $D_T^{[X]}=0$. For such a path,
$\mathcal R_a[e^{tX}]=O(t^4)$; the order-$g_H$ vacuum coefficient has no
quadratic or cubic frame term. This does not remove V69's distinct
order-$g_H^0$ determinant response.

The massless ordinary function in \eqref{eq:v70-beta} is exactly the
$J_\lambda(k)$ already used in V27 and evaluated in V61. Its integral
is not advertised as new. The new quantity below is the complete
interacting coefficient in which \emph{two} copies occur after the
common mass/vacuum subtraction.

\section{A populated nonzero two-loop frame kernel}
\label{sec:v70-twoframe}

\begin{theorem}[Full two-frame response at first quartic order]
\label{thm:v70-twoframe}
Let $X,Y$ be real symmetric flavour polarizations, with frame Fourier
momenta $k,-k$. Remove the common $g_H\hbar^2$ from the vacuum coefficient
\eqref{eq:v70-R}, and divide by physical volume. Its polarized Hessian is
\begin{equation}
 \boxed{\mathcal H_{a,L}(X,Y;k,z)
   =8\beta_{a,L}(k,z)^2
       \left[\operatorname{tr}X\operatorname{tr}Y
                            +2\operatorname{tr}(XY)\right].}
 \label{eq:v70-twoframe}
\end{equation}
There is no unrenormalized ultraviolet tadpole in this formula. Both
independent scalar loop momenta are retained by the product of the two
finite sums. For general real $X,Y$, replace them by their symmetric
parts.
\end{theorem}
\begin{proof}
At the frame origin $D_I=0$, so the second variation of
\eqref{eq:v70-R} has no term involving a second derivative of $D_T$.
It is
\[
 2\operatorname{tr}D^{[X]}\operatorname{tr}D^{[Y]}
                  +4\operatorname{tr}(D^{[X]}D^{[Y]}).
\]
Insert \eqref{eq:v70-first-D} and use the evenness of $\beta$. For
symmetric polarizations $D^{[X]}=2\beta X$ and
$D^{[Y]}=2\beta Y$. This yields \eqref{eq:v70-twoframe}, including
its factor eight. The cancellation theorem already includes the
one-loop mass insertion. Using only the quartic vacuum graph would
instead leave the $u_a$ term in \eqref{eq:v70-unpaid}.
\end{proof}

For a constant central frame $T=e^{tI_n}$ there is also an exact
finite-amplitude formula, valid on its Gaussian germ:
\begin{equation}
 \begin{split}
 D_T(x)&=d_a(t,z)I_n,\\
 d_a(t,z)&=\frac1{V_L}\sum_p\frac{\nu(p)}{P(p,z)}
       \frac{\Theta(p)(e^{2t}-1)}{\Theta(p)+e^{2t}\nu(p)},\\
 \frac{\mathcal R_a[e^{tI_n}]}{V_L}&=n(n+2)d_a(t,z)^2.
 \end{split}
 \label{eq:v70-central}
\end{equation}
The denominator is bounded away from zero for bounded small real $t$.
All nonzero summands of $d_a$ lie in the lower transition. The same
statement does not hold for the unsubtracted variance $u_a$.

\begin{proposition}[Strict positivity and a nonuniversal profile coefficient]
\label{prop:v70-positive}
In ordinary thermodynamic normalization define
\begin{equation}
 \beta_\lambda(z)=\int\frac{d^4p}{(2\pi)^4}
                  \frac{\nu_\lambda(p)\Theta_\lambda(p)}{p^2+z}.
 \label{eq:v70-beta0}
\end{equation}
For every admissible smooth transition and $z\ge0$,
$\beta_\lambda(z)>0$. At $z=0$,
\begin{equation}
 0<\beta_\lambda(0)
   =\lambda^2\int\frac{d^4y}{(2\pi)^4}
            \frac{\theta(y)(1-\theta(y))}{|y|^2}
   \le\frac{3\lambda^2}{64\pi^2}.
 \label{eq:v70-beta-bound}
\end{equation}
The four-real-scalar central Hessian is
\begin{equation}
 \boxed{\left.\frac1{V_L}\frac{d^2}{dt^2}
 \Gamma^{[g_H\hbar^2]}[e^{tI_4},0]\right|_{t=0}
       =192g_H\hbar^2\beta_\lambda(z)^2}
 \label{eq:v70-192}
\end{equation}
in the continuum comparison. It is nonzero for nonzero $g_H$, with sign
fixed by $g_H$. The finite coefficient is the same formula with
$\beta_{a,L}(0,z)$; strict positivity at a particular finite torus
requires at least one momentum in the smooth transition.
\end{proposition}
\begin{proof}
The integrand in \eqref{eq:v70-beta0} is nonnegative and is positive on
an open part of the transition. Such a part exists because a smooth
profile connects the value one on the unit ball to zero outside the
radius-two ball. The upper bound follows from
$\theta(1-\theta)\le1/4$ and
\[
 \frac{1}{4}\int_{1\le|y|\le2}\frac{d^4y}{(2\pi)^4|y|^2}
 =\frac3{64\pi^2}.
\]
No profile-independent positive lower bound follows: the transition
may be arbitrarily narrow. Finally $d_a'(0,z)=2\beta_{a,L}(0,z)$.
The second derivative of $n(n+2)d_a^2$ is
$8n(n+2)\beta_{a,L}^2$, which is 192 times that square at $n=4$.
\end{proof}

\begin{corollary}[The populated symmetric-frame block has rank ten]
\label{cor:v70-frame-rank}
At zero external momentum, identify real matrix frame tangents with
$M_4(\mathbb R)$ in the Hilbert--Schmidt pairing. After removal of
$g_H\hbar^2\beta_\lambda(z)^2$, the Hessian of
\eqref{eq:v70-twoframe} has eigenvalues
\[
 0\quad(6\text{ skew directions}),\qquad
 16\quad(9\text{ symmetric trace-free directions}),\qquad
 48\quad(1\text{ scalar direction}).
\]
Its rank is exactly ten. This is a rank of the displayed interacting
source coefficient, not the rank of a gauge or anomaly operator.
\end{corollary}
\begin{proof}
The bilinear form acts only on the symmetric part of each argument.
Its corresponding operator is $X\mapsto16X+8\operatorname{tr}(X)I_4$
on symmetric matrices. The trace-free and scalar subspaces give the
displayed eigenvalues. Strict positivity of $\beta_\lambda$ proves the
rank assertion for a nonzero coupling coefficient. At a finite cutoff
the same statement holds whenever $\beta_{a,L}(0,z)\ne0$.
\end{proof}

For the radial subclass, V61's already evaluated second Taylor
coefficient sharpens the momentum dependence. With its notation,
\begin{equation}
 \beta_\lambda(k,0)
 =\beta_\lambda(0)+\kappa_\theta k^2+O(|k|^4/\lambda^2),\qquad
 \kappa_\theta=\frac{E_\theta-1}{64\pi^2}>0.
 \label{eq:v70-radial-consumed}
\end{equation}
Thus the scalar factor $8\beta_\lambda(k,0)^2$ has the new interacting
quadratic coefficient $16\beta_\lambda(0)\kappa_\theta$. This consumes
\cref{thm:v61-hardbubble}; it is not a new proof of that one-loop profile
identity. No radiality assumption enters \eqref{eq:v70-twoframe} or
\eqref{eq:v70-beta-bound}.

The contact in \eqref{eq:v70-192} is lower-reference normalization data.
It is not a Higgs beta function, a gauge anomaly, or a prediction for a
physical scattering amplitude. Setting it to zero would specify another
finite composite-source normalization and change its source identity.

\section{All fixed source orders have a complete cutoff comparison}
\label{sec:v70-bounds}

We now estimate the already subtracted coefficient, not its individual
bare diagrams. Fix a maximum total external arity $N$, a derivative
order, and a rooted kernel moment. Every external frame and retained
scalar test has Fourier support $|k_i|\le\mu$. It suffices to work on
\begin{equation}
 \boxed{\lambda\ge32768N\mu>0,\qquad
 a\lambda\le2^{-24},\qquad
 \lambda L_\nu\ge1,\qquad 0\le z\le M^2\lambda^2.}
 \label{eq:v70-domain}
\end{equation}
Here $M,N$ and the scalar profiles are fixed. The formulas apply
coefficientwise at the source origin. The external momenta are not
allowed to grow with the cutoff while being estimated by the same
window constant.

Let $D^{[I]}_{a,L}(K,z)$ denote a positive-order frame coefficient of the
rooted difference \eqref{eq:v70-dressed}, with the frame polarization
norms factored into $\mathfrak n_I$. For $r=|\alpha|+2j$ the following
bounds are in the inherited physical Fourier convention.

\begin{lemma}[Uniform marked covariance bounds after the reference cancellation]
\label{lem:v70-cov-bounds}
For each fixed coefficient and differentiation order,
\begin{align}
 |\partial_K^\alpha\partial_z^jD_{0,L}^{[I]}|
 &\le C_{I,\alpha,j}\mathfrak n_I\lambda^{2-r},\nonumber\\
 |\partial_K^\alpha\partial_z^j(D_{a,L}^{[I]}-D_{0,L}^{[I]})|
 &\le C_{I,\alpha,j}\mathfrak n_I a^4\lambda^{6-r},
 \label{eq:v70-cov-bounds}\\
 |\partial_K^\alpha\partial_z^j(D_{0,L}^{[I]}-D_{0,\infty}^{[I]})|
 &\le C_{I,\alpha,j,b}\mathfrak n_I
       \lambda^{2-r}(\lambda L_{\min})^{-b},\qquad b>4.\nonumber
\end{align}
All original upper modes were included before the cancellation; none is
removed by a new regulator.
\end{lemma}
\begin{proof}
Use the finite word representation in
\cref{thm:v70-covariance-Ward}. At the source origin $G=C$ and $A=I$.
Every additional resolvent derivative inserts one order-two frame
vertex and one covariance. The net scalar-loop symbol has order $-2$,
regardless of the number of frame marks. The outer endpoint frames
have order zero. On the annulus in \cref{lem:v70-support}, each
momentum derivative costs at most $\lambda^{-1}$ and each mass
 derivative costs at most $\lambda^{-2}$. The normalized loop count is
$O(\lambda^4)$, uniformly when $\lambda L_\nu\ge1$. This proves the
first estimate.

On the same annulus all momentum arguments satisfy $|ap|<1/64$.
\Cref{prop:v27-scalar-reference} gives the differentiated native
comparison $P_a-P_0=O(a^4\lambda^6)$. The covariances differ by
$O(a^4\lambda^2)$, and an order-two source vertex differs by
$O(a^4\lambda^6)$, with the corresponding differentiated powers.
Telescope each finite ordered word. Replacing either one covariance or
one source vertex gives a symbol error of order $a^4\lambda^2$ before
integration; the loop count gives $a^4\lambda^6$. The source order is
fixed, so there are finitely many such replacements and rooted tree
paths. This proves the second estimate with every prescribed derivative.
For mass derivatives use the actual formula
$\partial_z^j C=(-1)^jj!\nu P^{-j-1}$; no extra mask factor is inserted.

Finally the ordinary assigned integrands are smooth and compactly
supported away from zero, by \cref{lem:v70-support}. After rescaling by
$\lambda$, Fourier integration by parts of arbitrarily prescribed order
bounds the nonzero Poisson images by
$C_b(\lambda L_{\min})^{-b}$. This proves the independent volume
comparison. No singular zero-mode term is hidden in that step.
\end{proof}

\begin{theorem}[First Higgs-quartic source continuum with common subtractions]
\label{thm:v70-continuum}
Let $k_s\in\{0,2,4\}$ be the number of retained scalar arguments in
\eqref{eq:v70-Wick-result}; it is not a momentum component. Remove the
common factor $g_H\hbar^{(4-k_s)/2}$ and write the resulting labelled
coefficient as $\mathcal F^{[k_s]}_{a,L,I}$. For $k_s=0,2$ and every
fixed $r=|\alpha|+2j$,
\begin{equation}
 \boxed{\begin{aligned}
 |\partial_K^\alpha\partial_z^j
 (\mathcal F^{[k_s]}_{a,L,I}-\mathcal F^{[k_s]}_{0,L,I})|
 &\le C_{I,\alpha,j}\mathfrak n_I
                 a^4\lambda^{8-k_s-r},\\
 |\partial_K^\alpha\partial_z^j
 (\mathcal F^{[k_s]}_{0,L,I}-\mathcal F^{[k_s]}_{0,\infty,I})|
 &\le C_{I,\alpha,j,b}\mathfrak n_I
 \lambda^{4-k_s-r}(\lambda L_{\min})^{-b}.
 \end{aligned}}
 \label{eq:v70-main-bound}
\end{equation}
A fixed $b>8$ suffices for the second line. The four-scalar tree
coefficient has no mesh discrepancy on these nonaliased soft panels.
All local frame coefficients of the displayed renormalized expression,
not only a further Taylor complement, are covered by this theorem.
\end{theorem}
\begin{proof}
First, $A_T\varphi=\varphi$ through the declared external order on this
soft panel. Every nonidentity term in the resolvent expansion has an
outer covariance $C$ carrying a sum of the external scalar and frame
momenta. That sum is smaller than $\lambda$, so the term is zero.
The same statement holds in the ordinary comparison. Thus $M=T\varphi$
for these coefficient extractions; this is not a statement for hard
retained scalar arguments or for arbitrary finite frame configurations.

For $k_s=2$, the one-loop terms of \eqref{eq:v70-Wick-result} contain
one $D_T$ and two exact local factors $T\varphi$. Apply
\eqref{eq:v70-cov-bounds} and the finite labelled product rule. For
$k_s=0$, differentiate the two products in \eqref{eq:v70-R}. Each
factor has positive frame degree because $D_I=0$. The cutoff difference
is a sum of one covariance difference of size $a^4\lambda^6$ times
one uniformly bounded covariance coefficient of size $\lambda^2$.
Leibniz and the differentiated versions of these estimates give
$a^4\lambda^{8-r}$. The product of two difference terms is no larger
on \eqref{eq:v70-domain}. The volume proof uses one Poisson error at a
time and the same product bounds. The two scalar loops are independent
rooted contractions, so the argument includes all their relative scale
assignments, rather than only diagonal scales. At $k_s=4$ only
$(T\varphi)^4$ remains, which is the same local multiplication in both
regulators before any aliasing can occur.
\end{proof}

\begin{corollary}[Rooted bounds without a total-volume multiplier]
\label{cor:v70-rooted}
After the same fixed smooth external window has been applied in both
regulators, the source kernels for $k_s=0,2$ obey
\begin{equation}
 \boxed{\|\mathscr K_{a,L,I}^{\Delta,[k_s]}\|_{\mathrm{rooted},b}
       \le C_{I,b}a^4\lambda^{8-k_s}.}
 \label{eq:v70-rooted}
\end{equation}
All $4(A-1)$ relative coordinates of an $A$-argument functional are
included. This gives a multilinear bound with one source in $L^1$ and
the remaining sources in $L^\infty$. The ordinary volume error has the
analogous factor
$\lambda^{4-k_s}(\lambda L_{\min})^{-b'}$, with sufficiently large
fixed $b'$.
\end{corollary}
\begin{proof}
\Cref{thm:v70-continuum} supplies every fixed external derivative.
Multiply the coefficient by the common smooth compact external window
and integrate by parts separately in each independent external momentum.
The powers $\mu^{|\alpha|}$ introduced by scaling the kernel are bounded
by the corresponding $\lambda^{|\alpha|}$ since $\mu\le\lambda$.
Choose enough derivatives to dominate the prescribed moment and the
$4(A-1)$-dimensional relative-coordinate summation. The Fourier
coefficients are normalized per physical volume, so this is a sum over
relative coordinates, not an estimate of an extensive action in sup
norm. Poisson comparison is applied before this same operation. Local
retained scalar factors and their frame multipliers do not introduce an
additional volume sum.
\end{proof}

For the two symmetric frame marks of \eqref{eq:v70-twoframe} the new
bound is, in particular, $O(a^4\lambda^8)$ for the entire two-loop
kernel at fixed $\lambda$. This rate is not inferred from two separately
renormalized hard integrals: the cancellation to $D_TD_T$ is essential.
Numerical evaluation of $u_a$ is not needed to prove the cancellation,
but an executable finite counterterm uses its actual finite sum.

\section{The interacting reference row and its normalization remain coupled}
\label{sec:v70-reference-row}

The matrix polynomial in \eqref{eq:v70-R} also determines the reference
response of this interaction. Define
\[
 Z_T(x)=\operatorname{tr}D_T(x)\,I_n+2D_T(x).
\]
\begin{proposition}[An explicit interacting frame response]
\label{prop:v70-interacting-Ward}
The source derivative of the complete two-loop coefficient is
\begin{equation}
 \boxed{\mathscr W_X\mathcal R_a[T]
 =2\int_a\operatorname{tr}
                \bigl(Z_T\mathscr W_XD_T\bigr),}
 \label{eq:v70-interacting-Ward}
\end{equation}
where $\mathscr W_XD_T$ is the diagonal block of
\eqref{eq:v70-covariance-Ward}, divided by $v$. It contains the actual
lower-reference factors in their original positions. If
$\mathfrak R_T^{[1]}(X)=-\mathscr W_X\mathcal R_a[T]$, then
\begin{equation}
 \mathscr W_X\mathfrak R_T^{[1]}(Y)
 -\mathscr W_Y\mathfrak R_T^{[1]}(X)
 =\mathfrak R_T^{[1]}([X,Y]).
 \label{eq:v70-consistency}
\end{equation}
\end{proposition}
\begin{proof}
Differentiate the two traces of squares. Symmetry is not used to
commute different source matrices; cyclicity alone gives
\eqref{eq:v70-interacting-Ward}. Then use the exact right-frame Lie
identity $[\mathscr W_X,\mathscr W_Y]=\mathscr W_{[X,Y]}$ on the one
finite functional $\mathcal R_a$. This is the source-row consistency of
the declared interaction and subtractions. It is not a proof that an
independently specified row has those coefficients.
\end{proof}

Restoring loop factors, the normalized vacuum source functional through
first quartic order is
\[
 \Gamma_a(g_H,T;0)
 =\hbar\Gamma_{a,0}(T)
       +g_H\hbar^2\mathcal R_a[T]+O(g_H^2).
\]
The first term is the V69 determinant. Its source identity has reference
$\mathfrak R_T^{[0]}(X)=\Tr(XA_T\Theta)$; the new reference coefficient
is exactly the $g_H\hbar$ correction
$\mathfrak R_T^{[1]}(X)$. Thus the finite Ward response and the mass
insertion are generated by the same action. The higher metric, gauge,
ghost and chiral identities are not obtained by relabelling this source
Lie algebra as the ESM gauge algebra. Cutoff-modified identities retain
reference terms even when a physical theory is anomaly free
\cite{V70IIM}; here those terms have been evaluated in this component.

\subsection*{Finite normalization changes are not pure vacuum changes}
Adding a supplied order-$g_H\hbar$ mass coefficient
$\frac12 g_H\hbar\delta m^2\int_a y^Ty$ changes the order-$g_H\hbar^2$
vacuum source functional, after its zero-frame value is removed, by
\begin{equation}
 \frac12g_H\hbar^2\delta m^2\int_a\operatorname{tr}D_T.
 \label{eq:v70-finite-anchor}
\end{equation}
It can therefore change a nonlocal finite frame coefficient. A matching
prescription must transport this term rather than changing only an
overall vacuum polynomial. No full common $S_1$ is presumed to have the
normal-ordering finite part; \eqref{eq:v70-finite-anchor} records the
conversion when the appropriate mass restriction of that action is
supplied. More general source-dependent finite quadratic tensors require
their own complete contraction with $K_T$, just as in V64.

\subsection*{All primitive scalar shell labels remain}
If the reference covariance is decomposed as $C=C_1+C_2$ with coincident
variances $u_1,u_2$, the normal-ordering operator factors exactly:
\[
 e^{-\hbar(u_1+u_2)\Delta_\phi/2}
 =e^{-\hbar u_1\Delta_\phi/2}
  e^{-\hbar u_2\Delta_\phi/2}.
\]
The vacuum coefficient contains $(u_1+u_2)^2$, not $u_1^2+u_2^2$.
Keeping the common mass subtraction but deleting the mixed vacuum term
would leave
$-2n(n+2)\hbar^2u_1u_2$ at the frame origin. In
\eqref{eq:v70-twoframe}, both independent sums likewise retain every
cross-shell product. This is a finite composition identity. It is not a
uniform many-step estimate for the full interacting action.

\subsection*{Which problem is and is not removed}
The previously uncomputed $g_H$ tadpole/mass insertion and its first
source-dependent double contraction are now one explicit comparison
packet. The $a^{-2}$ interaction-source error exposed by
\cref{prop:v70-tadpole-required} is cancelled. At order $g_H^2$ there
are multiple interaction vertices, connected contractions between them,
and overlapping subgraphs not represented by one local covariance
square. Turning on retained metric variations also changes the
coincident reference variance beyond a pure frame orbit. Neither problem
is eliminated by the present lower-annulus identity.

\section{Status, verification, and the next joint coefficient}
\label{sec:v70-ledgers}

\subsection*{Proved}
The first quartic-coupling contribution is evaluated as one finite
quartic/mass/vacuum insertion on V69's source orbit. Its transported
covariance difference has an exact lower-reference word representation.
The complete two-frame, two-loop coefficient is
\eqref{eq:v70-twoframe}, with a strictly nonzero central continuum
coefficient and the stated finite-torus qualification. Every prescribed
finite source order of this first-coupling packet has cutoff, volume,
and rooted estimates after the actual finite tadpole cancellation.
The same functional supplies its interacting source response and Lie
consistency. The single-vertex shell products are retained explicitly.

\subsection*{Inherited}
V9 specifies the real Higgs quartic normalization being consumed.
V27 supplies the corrected completed scalar reference, its positive
hard-domain bounds and fourth-order interior symbol comparison. V58
supplies conditional Gaussian source transport; V69 supplies the exact
finite endpoint frame and determinant identity. V57's variance growth
and V61's optional radial $J_\lambda$ coefficient are used at their
stated scope. None is claimed as a new result. No classical renewal
propagation estimate is used as a quantum bound.

\subsection*{Literature input}
Wick renormalization and cutoff-modified BV identities provide context
\cite{V70KMM,V70IIM}. The finite polynomial contractions, common
subtractions, covariance derivative, source Hessian and support argument
are proved here. The new computation is not a claim of first
renormalization of an $O(4)$ scalar theory. No external beta function,
cohomology classification, or anomaly-cancellation assertion supplies an
unknown native coefficient.

\subsection*{Conditional}
The normal-ordering mass anchor is an explicitly specified comparison in
the first $g_H$ component. Incorporating it into the common ESM action
requires agreement with that action's finite mass/source conditions or
the indicated finite conversion. The source frame has the same formal
status as in V69; no physical preparation protocol is inferred. The
finite formulas make sense on the positive scalar chart and
coefficientwise on its analytic germ. No infinite source expansion or
coupling series is asserted to converge.

\subsection*{Open}
The source-complete critical ESM primitive and the combined master
identity remain open. In particular the other scalar, gauge, fermion,
ghost and gravitational couplings and the first $g_H^2$ connected
coefficient require their own common subtractions. The useful next
calculation is a multiple-vertex or mixed-metric coefficient with this
\emph{populated} mass restriction of $S_1$ inserted, not another
independent choice of a tadpole constant. The existing finite graph
estimates do not by themselves establish the filtered interacting shell
self-map or the complete coefficientwise ESM continuum.

\begin{center}\footnotesize
\begin{tabular}{>{\raggedright\arraybackslash}p{.22\textwidth}>{\raggedright\arraybackslash}p{.70\textwidth}}
\toprule
Variable & Scope of V70\\\midrule
Coupling/loops & Exactly linear in $g_H$; zero, one and two scalar
contractions of that vertex. No integrated graviton, internal gauge,
ghost or chiral line.\\
Sources & Every prescribed finite frame/retained-scalar arity $N$ in
\eqref{eq:v70-domain}; estimates at the source origin. General finite
frames are used only in the finite algebraic formulas.\\
Cutoff/volume & Fourth-order comparison after common tadpole subtraction;
independent Poisson volume error; no total-volume multiplier in rooted
bounds. Constants depend on fixed $N$, derivatives, moments and profiles.\\
Mass & Exact fixed-mass algebra with the actual $u_a(z)$, including $z=0$.
A finite mass-polynomial EFT assignment is extracted jointly; an
untruncated $u_a(z)$ is not called a polynomial.\\
Normalization & Existing $g_H=0$ prescriptions, covariance and measure
unchanged. A new normal-ordering comparison in the $g_H$ component is
recorded; its finite conversion is not silently set to zero.\\
Verification & Exact rational finite covariance and polynomial tests;
no numerical native profile moment, complete master coefficient, or
proof-assistant verification of inherited analytic results.\\\bottomrule
\end{tabular}
\end{center}

\begin{thebibliography}{99}
\bibitem{V70KMM} I.~Khavkine, A.~Melati and V.~Moretti,
\emph{On Wick polynomials of boson fields in locally covariant algebraic
QFT}, arXiv:1710.01937, version 2 (2018).
\bibitem{V70IIM} Y.~Igarashi, K.~Itoh and T.~R.~Morris,
\emph{BRST in the Exact RG}, Prog. Theor. Exp. Phys. (2019), 103B01;
arXiv:1904.08231.
\end{thebibliography}

\section*{Append-only continuation marker: V70}
\addcontentsline{toc}{section}{Append-only continuation marker: V70}
\label{sec:v70-continuation-marker}
\textbf{End of V70.} Preserve Sections 1--543. The first Higgs-quartic
source coefficient now includes its common mass and vacuum subtractions,
its nonzero lower-reference response, and its fixed-order cutoff limit.
The complete interacting ESM normalization and many-shell closure remain
open. No zero-coupling action or active earlier normalization was reset.


\clearpage
\part*{V71: two quartic vertices, a generated marginal, and the retained six-field return}
\addcontentsline{toc}{part}{V71: two quartic vertices and the retained six-field return}

\section{A multiple-vertex test of the actual scalar blocking class}
\label{sec:v71-start}

Sections 1--543 and the V70 continuation marker are preserved. V70
normal-orders each quartic vertex with the actual finite scalar variance.
That removes self-contractions at one vertex, not contractions between
vertices. The present installment computes the full second cumulant of
that same insertion and renormalizes its one-loop four-scalar component.
It then tests its composition under a split of the \emph{actual}
reference covariance. The resulting marginal update consumes a generated
six-field kernel: discarding that kernel because it vanishes on the final
soft six-field panel changes the next quartic coefficient.

The endpoint is a coefficientwise statement in the scalar part of the
ESM action, not a new construction of scalar quantum field theory. The
renormalization of $O(n)$ quartic models and their perturbative flow are
established, including at much higher orders \cite{V71KP}. The question
here is whether this coefficient is produced, normalized and composed by
the inherited finite measured map. No continuum beta coefficient is
inserted as an assumption. The other ESM couplings and integrated metric,
ghost and fermion fluctuations are zero in this component.

Use exactly V70's convention, with $n=4$ in the Higgs application:
\begin{equation}
 v=a^4,\qquad P=\widetilde P_a(z),\quad z=m_H^2,\qquad
 c_{a,\lambda}(p,z)=\frac{\nu_\lambda(p)}{P(p,z)},\qquad
 \nu_\lambda=1-\Theta_\lambda.
 \label{eq:v71-reference}
\end{equation}
The quotient is zero on retained modes, including the zero mode at
$z=0$. Set
\begin{equation}
 c(x-y)=\frac1{V_L}\sum_p e^{ip\cdot(x-y)}c_{a,\lambda}(p,z),
 \qquad u=c(0),\qquad U(y)=(y^Ty)^2.
 \label{eq:v71-coordinate}
\end{equation}
Thus the scalar coordinate covariance is $\hbar c(x-y)I_n$.
Both the local quartic and its already specified comparison subtractions
are retained:
\begin{equation}
 :U:_u(y)=U(y)-2\hbar(n+2)u\,y^Ty
                    +\hbar^2n(n+2)u^2.
 \label{eq:v71-NO}
\end{equation}
The input interaction is $g_H\int_a:U:_u(T\phi)$. No subtraction is
resummed into $P$. The exact fixed-mass expressions below can be
Taylor-extracted jointly in $z$ when a mass-polynomial EFT bank is
required. The four-field local anchor will be specified at $z=0$.

For a frame $T$, retain V70's $M=T A_T\varphi$ and define the complete
coordinate covariance and its diagonal difference by
\begin{equation}
 K_T(x,y)=v^{-1}[T G_TT^T]_{xy},\qquad
 D_x=K_T(x,x)-uI_n.
 \label{eq:v71-frame-data}
\end{equation}
These are not replaced by their coincident values when two different
vertices are contracted. The scalar determinant remains in the
$g_H$-independent part of the action. Normalizing the total partition
function at $T=I,\varphi=0$ removes only the value of the resulting
vacuum coefficient there, not its source derivatives.

\section{The entire second cumulant with nonlinear frame sources}
\label{sec:v71-cumulant}

\begin{theorem}[Finite connected two-vertex formula]
\label{thm:v71-cumulant}
Let $s,t$ be two dummy real flavour vectors. Define
\begin{equation}
 \begin{split}
 \mathfrak C_T(x,y)
 =\Big[&\exp\!\left\{\frac\hbar2D_x:\partial_s^2
                         +\frac\hbar2D_y:\partial_t^2\right\}\\
 &\quad\times
 \left(\exp\{\hbar K_T(x,y)_{AB}\partial_{s_A}\partial_{t_B}\}-1\right)
 U(s)U(t)\Big]_{s=M(x),\,t=M(y)}.
 \end{split}
 \label{eq:v71-connected}
\end{equation}
All exponential series in this formula terminate on $U(s)U(t)$.
The complete effective-action coefficient with two copies of the
specified quartic/mass/vacuum packet is
\begin{equation}
 \boxed{\Gamma_{a,2}^{[g_H^2]}[T,\varphi]
       =-\frac{g_H^2}{2\hbar}\int_a\!dx\int_a\!dy\,
                         \mathfrak C_T(x,y).}
 \label{eq:v71-second}
\end{equation}
The formula includes all self- and cross-contractions and all ordered
matrix-frame derivatives. It does not include counterterms first
introduced at order $g_H^2$.
\end{theorem}
\begin{proof}
The coefficient of $g_H^2$ in $-\hbar\log\E\exp(-g_HV/\hbar)$ is
$-(2\hbar)^{-1}\Cov(V,V)$. The Gaussian expectation on the two dummy
fields is generated by their two self-covariances and the cross
covariance. Normal ordering at each vertex subtracts $uI_n$ from its
self-covariance; the remaining matrices are exactly $D_x,D_y$.
The product of the two one-vertex expectations is the term with no cross
contraction. Subtracting it produces the minus one in
\eqref{eq:v71-connected}. All differential operators have coefficients
independent of $s,t$ and commute. This argument also applies at $x=y$:
the dummy fields distinguish the vertices, not the underlying Gaussian
variable. Every operation is on the original finite covariance. It uses
neither independent diagonalizations of $K_T$ nor an additional flavour
multiplicity.
\end{proof}

At $T=I$, $D_x=0$, $M=\varphi$ and $K_T(x,y)=c(x-y)I_n$.
Write $U_x=\varphi_x^T\varphi_x$ and
$W_{xy}=\varphi_x^T\varphi_y$.
\begin{proposition}[All four connected edge multiplicities]
\label{prop:v71-four-channels}
In this restriction, \eqref{eq:v71-second} is exactly
\begin{align}
 \Gamma_{a,2}^{[g_H^2]}={}&
 -8g_H^2\int_{a,x,y}c_{xy}\,U_xU_yW_{xy}\nonumber\\
 &-4g_H^2\hbar\int_{a,x,y}c_{xy}^2
                   \big[(n+4)U_xU_y+4W_{xy}^2\big]\nonumber\\
 &-16g_H^2\hbar^2(n+2)\int_{a,x,y}c_{xy}^3W_{xy}\nonumber\\
 &-4g_H^2\hbar^3n(n+2)\int_{a,x,y}c_{xy}^4.
 \label{eq:v71-components}
\end{align}
The rows contain respectively six, four, two and zero retained scalar
fields. They have zero, one, two and three loops. Only the final row is
removed at the frame-free normalization point.
\end{proposition}
\begin{proof}
Put $A=s^Ts$, $B=t^Tt$, $W=s^Tt$ and
$\mathcal D=\sum_A\partial_{s_A}\partial_{t_A}$. Direct differentiation
gives
\begin{align*}
 \mathcal D(U(s)U(t))&=16ABW,\\
 \mathcal D^2(U(s)U(t))&=16[(n+4)AB+4W^2],\\
 \mathcal D^3(U(s)U(t))&=192(n+2)W,\\
 \mathcal D^4(U(s)U(t))&=192n(n+2).
\end{align*}
Multiply the $r$th row by $(\hbar c)^r/r!$ and by
$-g_H^2/(2\hbar)$. Higher derivatives vanish. This determines the
symmetry factors without diagrammatic multiplicity conventions.
\end{proof}

\begin{lemma}[The bridge is absent only after the final soft restriction]
\label{lem:v71-bridge}
If all six scalar momenta have norm at most $\mu$ and $3\mu<\lambda$,
the first row of \eqref{eq:v71-components} vanishes. With finitely many
frame marks the analogous one-cross-edge terms vanish whenever every
momentum sum entering that bridge remains below the hard plateau.
This statement does not hold for a field that is to be integrated in a
subsequent Gaussian step.
\end{lemma}
\begin{proof}
The momentum in the single covariance is the sum of the external momenta
on one side of the bridge. There is no free loop momentum on that edge.
Its multiplier $\nu_\lambda$ is zero. For frame-marked words the same
argument applies after expanding every endpoint and resolvent: every
partial momentum differs by a sum of the declared soft labels. A later
Gaussian field is not restricted to that final source window, so that
last assertion cannot be inferred for it.
\end{proof}

\section{One common quartic counterterm and its source dependence}
\label{sec:v71-fish}

Set
\begin{equation}
 B_{a,\lambda,L}(k,z)=\frac1{V_L}\sum_p
     c_{a,\lambda}(p,z)c_{a,\lambda}(p+k,z),\qquad
 b_{a,\lambda,L}=B_{a,\lambda,L}(0,0).
 \label{eq:v71-bubble}
\end{equation}
All momentum additions in a finite sum use the existing periodic
completion. At nonzero external momenta this is the two-line scalar
bubble, not the one-propagator $J_\lambda$ of V27/V61/V70.

\begin{theorem}[The populated local quartic coefficient]
\label{thm:v71-local}
The momentum-and-mass weight-zero part of the four-scalar row is
\begin{equation}
 \boxed{-4(n+8)\hbar g_H^2 b_{a,\lambda,L}
                        \int_a(\varphi^T\varphi)^2.}
 \label{eq:v71-local}
\end{equation}
It is cancelled by the opposite local coefficient. A source-consistent
extension is obtained by replacing $\varphi$ by $T\phi$ in that same
quartic counterterm. Choosing its complete V70 normal-ordered version
also prescribes lower-field, higher-loop partners; those do not by
themselves renormalize the two- and zero-field rows of
\eqref{eq:v71-components}.
\end{theorem}
\begin{proof}
For a local projection the fields at the two vertices have the same
value. The flavour polynomial in the four-field row becomes
$(n+8)(\varphi^T\varphi)^2$. The integral of $c_{xy}^2$ over $y$ is
$B(0,z)$ by the finite Fourier normalization. Its weight-zero
assignment evaluates $z=0$, proving \eqref{eq:v71-local}. The same
coefficient multiplies one action monomial, so its source derivatives
are not independent adjustments. The higher-loop assertion follows
from the explicit powers of $\hbar$ in \eqref{eq:v71-NO}.
\end{proof}

For nonconstant frames write $H_U(m)=4(m^Tm)I_n+8mm^T$.
On final soft panels excluding one-cross-edge bridges, the entire
one-loop, four-scalar coefficient is
\begin{equation}
 \boxed{-\frac{g_H^2\hbar}{4}\int_{a,x,y}
 \operatorname{tr}\big[H_U(M_x)K_T(x,y)
                       H_U(M_y)K_T(y,x)\big].}
 \label{eq:v71-frame-fish}
\end{equation}
It is the two-cross-edge term of \eqref{eq:v71-connected}. At this
field and loop degree no $D_x$ insertion can be added to it. A one-cross-
edge term with one $D_x$ is excluded by \cref{lem:v71-bridge}, not
silently omitted from the general cumulant.
For any fixed total soft arity $N$ with $N\mu<\lambda$, the mean obeys
$A_T\varphi=\varphi$ through that arity. Indeed every correction has a
leftmost hard covariance at a sum of soft labels. Thus $M=T\varphi$ in
this assertion.

V70's ordered derivative for $TG_TT^T$ applies before taking a diagonal
block. Hence every positive-order frame coefficient of
$K_T-cI_n$ has a lower-reference word representation. In
\eqref{eq:v71-frame-fish}, any term containing this difference has its
one loop near $\lambda$. The ultraviolet logarithm therefore comes
only from the $cI_n,cI_n$ product and its endpoint fields $T\varphi$.
This proves that the one common counterterm in
\cref{thm:v71-local} cancels the logarithm at every fixed soft frame
arity. It does not prove a full ESM Slavnov equation.

\begin{proposition}[A finite interacting frame anchor is not zero]
\label{prop:v71-frame-anchor}
For a constant central frame $T=e^{tI_n}$ at zero mass, put
\begin{equation}
 c_t(p)=\frac{e^{2t}\nu(p)}{P(p)[\Theta(p)+e^{2t}\nu(p)]},
 \quad b_t=\frac1{V_L}\sum_pc_t(p)^2,
 \quad j_{a,L}=\frac1{V_L}\sum_p\frac{\nu(p)^2\Theta(p)}{P(p)^2}.
 \label{eq:v71-central}
\end{equation}
After the quartic normalization in \eqref{eq:v71-local}, the local
four-field coefficient at this frame, with $g_H^2\hbar$ removed, is
\begin{equation}
 -4(n+8)e^{4t}(b_t-b_0).
 \label{eq:v71-central-ren}
\end{equation}
Its first derivative is $-16(n+8)j_{a,L}$. In thermodynamic ordinary
normalization
\begin{equation}
 j_{a,L}\longrightarrow
 j_\theta=\int\frac{d^4y}{(2\pi)^4}
               \frac{(1-\theta(y))^2\theta(y)}{|y|^4}>0.
 \label{eq:v71-jtheta}
\end{equation}
The moment in \eqref{eq:v71-jtheta} is V68's existing $L_1$; the new
quantity is its coefficient in the interacting four-field normalization.
\end{proposition}
\begin{proof}
At constant $t$, $M=e^t\varphi$ on the soft scalar support, and the
covariance is diagonal in momentum with symbol $c_t$. Apply
\eqref{eq:v71-frame-fish} and subtract its common local counterterm.
Since $\partial_tc_t|_0=2\Theta c$, $\partial_tb_t|_0=4j_{a,L}$.
The difference $b_t-b_0$ has support only in the lower transition, so
its ordinary limit is a finite integral. Smoothness, a unit plateau
and compact support imply an open set with $0<\theta<1$, proving
strict positivity. No profile-independent positive lower bound follows.
The finite sum is strictly positive only when a lattice momentum meets
that transition. The scalar determinant response of V69 and the first-
$g_H$ response of V70 are different coefficients and are not counted here.
\end{proof}

\section{Full-cutoff comparison of the renormalized four-field kernel}
\label{sec:v71-bounds}

Define the joint weight-zero subtraction
\begin{equation}
 B^R_{a,\lambda,L}(k,z)
       =B_{a,\lambda,L}(k,z)-B_{a,\lambda,L}(0,0).
 \label{eq:v71-BR}
\end{equation}
For the ordinary coefficient the same subtraction is made in the
integrand before summing or integrating. Because that integrand is not
pointwise odd at first external order, it may equivalently be averaged
at $k,-k$ before estimating; this changes no finite symmetric sum.
No value is assigned to either divergent ordinary unrenormalized term.

\begin{theorem}[Scalar-bubble continuum in the source topology]
\label{thm:v71-bound}
Use V27's completed scalar reference and fix $\lambda>0$, $M<\infty$,
a finite source order and finitely many derivatives and kernel moments.
It suffices that
\begin{equation}
 \lambda\ge32768N\mu>0,\qquad a\lambda\le2^{-24},\qquad
 \lambda L_\nu\ge1,\qquad |k|\le N\mu,\qquad
 0\le z\le M^2\lambda^2.
 \label{eq:v71-domain}
\end{equation}
For $r=|\alpha|+2j\le2$ and $\epsilon=|k|+\sqrt z$,
\begin{align}
 |\partial_k^\alpha\partial_z^j
 (B^R_{a,\lambda,L}-B^R_{0,\lambda,L})|
 &\le C_{\alpha,j}a^2\epsilon^{2-r},
 \label{eq:v71-rate}\\
 |\partial_k^\alpha\partial_z^j
 (B^R_{0,\lambda,L}-B^R_{0,\lambda,\infty})|
 &\le C_{\alpha,j,b}\lambda^{-2}(\lambda L_{\min})^{-b}
                              \epsilon^{2-r},\quad b>4.
 \label{eq:v71-volume}
\end{align}
The right sides at $(k,z)=(0,0)$ are interpreted by continuity and the
stored jets. All completed ultraviolet modes are included. Higher fixed
external derivatives obey the differentiated symbol estimates in the
proof. For the massless four-field action, its localized rooted kernel
has cutoff error at most
\begin{equation}
 C_{N,b}\,|g_H|^2\hbar\,a^2\mu^2
 \label{eq:v71-rooted}
\end{equation}
after the flavour and source polarizations are factored out. There is
no total-volume multiplier.
\end{theorem}
\begin{proof}
Write $R_p=\lambda+|p|$. On every fixed set of derivatives the product
of two actual scalar covariances has ordinary degree $-4$:
\[
 |\partial_p^\beta\partial_k^\alpha\partial_z^j
             (c(p,z)c(p+k,z))|
 \le C R_p^{-4-|\beta|-|\alpha|-2j}.
\]
Cutoff derivatives are in the lower transition and satisfy the same
scaled bound. By \cref{prop:v27-scalar-reference}, on a common
interior region the finite/ordinary difference has the improved bound
$Ca^4R_p^{-|\beta|-|\alpha|-2j}$. Thus two weighted external
derivatives give ordinary order $-6$ and discrepancy $Ca^4R_p^{-2}$.
On a dyadic shell the normalized four-dimensional mode count is at
most $CR^4$, uniformly for $\lambda L_\nu\ge1$. Hence
\[
 a^4\sum_{R\lesssim a^{-1}}R^2\le Ca^2,
 \qquad
 \sum_{R\gtrsim a^{-1}}R^{-2}\le Ca^2.
\]
Smooth periodic upper pieces, including a covariance whose shifted
argument crosses a face, obey the same $a^2$ bound after rescaling.
These bounds also apply to one $z$ derivative. Momentum inversion
makes the summed coefficient even in $k$. Integrating the bounded
second momentum derivatives and first mass derivative from $(0,0)$
proves \eqref{eq:v71-rate}, including its first derivatives. This is a
joint subtraction, not a subtraction at a different mass in one factor.

For the ordinary finite-volume comparison, apply Poisson summation on
each smooth dyadic piece of those differentiated integrands. Integration
by parts in $p$ supplies arbitrary fixed powers of $(RL_{\min})^{-1}$;
the leading sum is $C\lambda^{-2}(\lambda L_{\min})^{-b}$. Integrate the
same Taylor remainder to obtain \eqref{eq:v71-volume}. At finite volume
external interpolation is fixed by the same smooth routed integrand;
no off-grid Ward identity is inferred from it.

Additional momentum derivatives improve the ordinary ultraviolet
degree. The corresponding discrepancy sums are still $O(a^2)$ at
fixed $\lambda$; for derivative order four a harmless
$a^4\log(1/(a\lambda))$ is bounded by $Ca^2\lambda^{-2}$. These bounds
control all scaled derivatives after localization to the fixed source
window. Fourier integration by parts gives the weighted $\ell^1$
rooted norm, including the three independent relative coordinates of
the four-field panel. The local vertex delta kernels can be convolved
with these window kernels without a volume factor. Equivalently the
rooted norm gives one scalar argument in $L^1$ and the other three in
$L^\infty$. Mass and derivative-window constants are retained; no
uniformity in an unbounded source order is asserted.
\end{proof}

\begin{proposition}[Fixed frame derivatives and the retained finite contact]
\label{prop:v71-frame-bound}
For any prescribed number of soft frame marks, the renormalized
four-scalar coefficient of \eqref{eq:v71-frame-fish} is well defined
with the same local quartic anchor. Its massless windowed cutoff error
is bounded by
\begin{equation}
 C_{N,b}|g_H|^2\hbar
       \bigl[a^2\mu^2+(a\lambda)^4\bigr]
 \label{eq:v71-frame-bound}
\end{equation}
in the corresponding rooted norm, with a separate volume error
$C_{N,b}|g_H|^2\hbar[(\mu/\lambda)^2+1]
(\lambda L_{\min})^{-b}$. The second term in
\eqref{eq:v71-frame-bound} is not set to zero at zero external momentum.
\end{proposition}
\begin{proof}
Expand $K_T=cI_n+(K_T-cI_n)$ using the ordered derivative identity of
V70 before taking norms. The $c,c$ term is the bubble just estimated,
with endpoint fields $T\varphi$ and their complete local counterterm.
Every other differentiated term contains a $\Theta$ and a $C$ on its
single loop. At fixed source arity its loop is in
$\lambda/2\le|p|\le3\lambda$. Its total degree is zero; V27's
fourth-order free-symbol comparison therefore gives
$C(a\lambda)^4$, and its differentiated bounds give the corresponding
windowed rooted estimate. The same annular words have Poisson error
$C(\lambda L_{\min})^{-b}$. Both signs and matrix orders are retained.
\cref{prop:v71-frame-anchor} shows why an unqualified
$O(a^2|K|^2)$ assertion for the entire marked row would be unjustified:
a local lower-reference coefficient survives at $K=0$ and its finite
comparison need not vanish exactly there.
\end{proof}

\section{The marginal logarithm is generated by this finite coefficient}
\label{sec:v71-marginal}

\begin{theorem}[Native critical logarithm and finite-scale update]
\label{thm:v71-marginal}
In the ordinary thermodynamic normalization,
\begin{equation}
 b_{a,\lambda,\infty}
 =\frac1{8\pi^2}\log\frac1{a\lambda}
      +\kappa_{P,\theta}+O((a\lambda)^4).
 \label{eq:v71-blog}
\end{equation}
Here $\kappa_{P,\theta}$ is a finite completion/profile-dependent
constant; it is not assigned a numerical value. For the four-field
coefficient, at bidegree $g_H^2\hbar$, a fixed zero-mass local anchor at
$\lambda_*$ gives
\begin{equation}
 \boxed{g_H(\lambda)
   =g_H(\lambda_*)+
       \frac{n+8}{2\pi^2}\hbar g_H(\lambda_*)^2
                         \log\frac{\lambda}{\lambda_*}.}
 \label{eq:v71-running}
\end{equation}
Equality means equality of the indicated formal bidegree. For $n=4$
the coefficient is $6/\pi^2$. This is not an analytic trajectory at
arbitrary ultraviolet scale.
\end{theorem}
\begin{proof}
At $z=0$ the native symbol has the form $P(p,0)=a^{-2}p_*(ap)$,
where $p_*(x)=|x|^2+O(|x|^6)$ near zero and has the inherited positive
floor elsewhere in the principal cube. Therefore
$p_*^{-2}-|x|^{-4}=O(1)$ near zero. Separate a small radial ball contained
in the common interior plateau from its complement. On the complement
the lower mask equals one for small $a\lambda$ and the integral is
finite. In the ball the radial integral of $|x|^{-4}$ supplies
$\operatorname{area}(S^3)/(2\pi)^4=1/(8\pi^2)$ times the logarithm.
The transition-dependent constant is finite ray by ray, without a radial
assumption on $\theta$. The omitted bounded difference on a ball of
radius $O(a\lambda)$ is $O((a\lambda)^4)$. This proves
\eqref{eq:v71-blog}.

More directly, for any fixed pair of lower scales,
\begin{equation}
 \int\frac{d^4p}{(2\pi)^4}
 \frac{\nu_{\lambda_1}(p)^2-\nu_{\lambda_2}(p)^2}{p^4}
       =\frac1{8\pi^2}\log\frac{\lambda_2}{\lambda_1}.
 \label{eq:v71-Frullani}
\end{equation}
For each unit direction set $F(r)=(1-\theta(r\widehat p))^2$.
It is zero below one and one above two. Differentiating the scaled
integral in $\log\lambda$ gives $-\int_0^\infty F'(r)dr=-1$;
compact transition support justifies differentiation even without
monotonicity. Integrate this identity and then the sphere. The finite
native difference has error $C(a\lambda_{\max})^4$ for a fixed scale
ratio and a separate $C_b(\lambda_{\min}L_{\min})^{-b}$ error.
Finally the local coefficient of the measured effective action is
$g_b-4(n+8)\hbar g_b^2b_{a,\lambda,L}$. Eliminate the same bare
coefficient using its value at $\lambda_*$, and use
\eqref{eq:v71-Frullani}. This gives \eqref{eq:v71-running} with its
sign. The mass counterterms and other field degrees have not been
projected away from the full action. A change of the first-order mass
anchor can contribute to this four-field bidegree only through a
single hard bridge joining a quadratic vertex to a quartic vertex.
That bridge vanishes on the final soft panel by
\cref{lem:v71-bridge}. Thus comparing these projected couplings does
not require silently changing the underlying bare mass convention. All further orders are outside this
identity.
\end{proof}

The familiar scalar coefficient is thus reproduced by the native
finite bubble and its actual Wick multiplicity, not imported from
\cite{V71KP}. With the conventional interaction
$\lambda_4(\phi^T\phi)^2/4!$, $\lambda_4=24g_H$, the same equation reads
$\partial_{\log\lambda}\lambda_4=(n+8)\hbar\lambda_4^2/(48\pi^2)$
at this order. It is positive rather than asymptotically free. The
positive lower scale, finite coupling-order interpretation, and ESM
Wilson-coefficient tower remain part of the claim.

\section{Why the generated six-field kernel is necessary for that update}
\label{sec:v71-mixed-shell}

\begin{theorem}[Exact two-step composition in the quartic row]
\label{thm:v71-composition}
Split an actual finite positive scalar covariance as $c=c_1+c_2$,
with coincident variances $u_1,u_2$. Integrate two independent Gaussian
innovations in that order, starting with normal ordering at $u_1+u_2$.
The one-loop four-field coefficient equals the one-step coefficient
\eqref{eq:v71-components} with $c_{xy}^2$ replaced by
\begin{equation}
 \boxed{c_{1,xy}^2+2c_{1,xy}c_{2,xy}+c_{2,xy}^2.}
 \label{eq:v71-cross-shell}
\end{equation}
The middle term is generated by the six-field coefficient of the first
step. It is not obtained by adding only the two single-step quartic
local corrections.
\end{theorem}
\begin{proof}
Gaussian convolution and conditional cumulance give the exact identity
\[
 \Cov_{1+2}(V,V)=\E_2\Cov_1(V,V)
                              +\Cov_2(\E_1V,\E_1V).
\]
Normal ordering changes from $u_1+u_2$ to $u_2$ after the first
expectation. In its first covariance, the one-cross-edge six-field term
is
$-8g_H^2c_{1,xy}U_xU_yW_{xy}$, with its local polynomial normally
ordered in the remaining self-covariance $u_2$. One cross-contraction
from the second innovation is $\hbar c_{2,xy}\mathcal D$.
Directly,
\[
 \mathcal D(U_xU_yW_{xy})=(n+4)U_xU_y+4W_{xy}^2.
\]
Its resulting coefficient is $-8g_H^2\hbar c_1c_2$ times the same
four-field polynomial. The two-cross-edge first-step term gives
$-4g_H^2\hbar c_1^2$ and the second conditional covariance gives
$-4g_H^2\hbar c_2^2$. Self-contractions at both stages cancel in the
common normal-product convention; no independently chosen tadpole
constant is introduced. This proves the formula. The identity is at
fixed finite carrier and does not assert a uniform many-step norm.
\end{proof}

\begin{proposition}[A finite marginal error if the six-field return is dropped]
\label{prop:v71-missing}
Take $\lambda_1\ge2\lambda_2$ and
\[
 c_1(p)=\frac{1-\Theta_{\lambda_1}(p)}{P(p,0)},\qquad
 c_2(p)=\frac{\Theta_{\lambda_1}(p)-\Theta_{\lambda_2}(p)}{P(p,0)}.
\]
Both are nonnegative, even without radial monotonicity of the admissible
profile. Their sum is $c_{a,\lambda_2}(p,0)$. The mixed local integral is
\begin{equation}
 \frac1{V_L}\sum_pc_1(p)c_2(p)
 =\frac1{V_L}\sum_p
 \frac{\nu_{\lambda_1}(p)\Theta_{\lambda_1}(p)}{P(p,0)^2}
 \longrightarrow s_\theta,
 \label{eq:v71-mixed-moment}
\end{equation}
where
\begin{equation}
 \boxed{0<s_\theta=
 \int\frac{d^4y}{(2\pi)^4}
       \frac{\theta(y)(1-\theta(y))}{|y|^4}
       \le\frac{\log2}{32\pi^2}.}
 \label{eq:v71-stheta}
\end{equation}
Deleting the generated six-field coefficient misses the local
four-field contribution
\begin{equation}
 \boxed{-8(n+8)\hbar g_H^2s_\theta.}
 \label{eq:v71-missing}
\end{equation}
It does not tend to zero with the mesh. No profile-independent strictly
positive lower bound is asserted.
\end{proposition}
\begin{proof}
Below $\lambda_1$, the first covariance is zero. Above $\lambda_1$,
$\Theta_{\lambda_2}=0$, so the product is exactly as displayed.
The same support separation proves positivity of the second
covariance. The mixed loop lies in the transition
$\lambda_1<|p|<2\lambda_1$. Ordinary rescaling proves the limit with
error $C(a\lambda_1)^4+C_b(\lambda_1L_{\min})^{-b}$. An admissible
smooth transition has a nonempty open set with $0<\theta<1$.
The upper bound follows from $\theta(1-\theta)\le1/4$ and the
four-dimensional radial measure $dr/(8\pi^2r)$. Finally the middle
term of \eqref{eq:v71-cross-shell} has the local flavour multiplier
$n+8$, giving \eqref{eq:v71-missing}.
\end{proof}

For a fixed ratio $r=\lambda_1/\lambda_2\ge2$, the ordinary integral of
$c_2^2$ alone is
\begin{equation}
 \int_pc_2(p)^2=\frac{\log r}{8\pi^2}-2s_\theta,
 \label{eq:v71-shell-alone}
\end{equation}
whereas the complete incremental quartic coefficient consumes
$\log r/(8\pi^2)$. Thus a coupling-only update based on the second
shell's bubble would acquire a spurious profile dependence. Equation
\eqref{eq:v71-running} is a generated marginal only when the
complementary six-field return is carried to the next integration.
This explicitly realizes, in one populated component, the distinction
between a final soft source norm and the hard-domain norm needed for
an invariant blocking state. General combinatorial formulations of
Wilsonian composition are comparison material \cite{V71KM}; they do
not replace this native coefficient and its mixed contribution.

\section{Status and the next common-action problem}
\label{sec:v71-ledgers}

\subsection*{Proved}
The complete finite second cumulant of the V70 interaction is given with
all frame, mass and normal-product partners. Its one-loop four-scalar
row has one populated quartic counterterm, an evaluated logarithm, a
finite lower-reference frame contact, and a cutoff/volume/rooted
comparison. The indicated scalar marginal runs with the coefficient
obtained from that finite bubble. Two-step composition is evaluated,
and a nonzero cutoff-stable marginal error is exhibited when its
six-field return is dropped. These are new calculations in this
lineage, not first renormalizability results for a quartic scalar model.

\subsection*{Inherited}
V27 supplies the completed scalar carrier, symbol floor, fourth-order
interior comparison and smooth periodic upper bounds. V70 supplies the
local quartic normalization, its actual tadpole/vacuum packet and the
ordered frame covariance. V69 supplies the finite source frame, and
V68 supplies the previously identified profile moment appearing in the
constant-frame test. Those results are used at their stated scope and
are not independently reproved in their entirety.

\subsection*{Literature input}
Wick powers, scalar perturbative renormalization, Wilsonian composition
and cutoff-modified BRST identities are established
\cite{V71KMM,V71KP,V71KM,V71IIM}. They supply context and convention
checks, not an unknown coefficient or a uniform native estimate.
The cumulant contractions, logarithmic shell difference and
cross-shell mechanism are derived here. No outside anomaly theorem
or all-orders scalar construction is invoked as the ESM endpoint.

\subsection*{Conditional}
The normal-product finite mass anchor is V70's comparison, not an
assertion about an unpopulated common ESM $S_1$. The newly specified
quartic finite anchor and any normal-ordered higher-loop partners must
be converted if the common normalization chooses different finite
coefficients. Analytic estimates apply at the source origin to each
fixed number of frame and scalar marks. The general matrix-frame
cumulant is an exact finite algebraic formula; higher-loop estimates
for all of its components have not been supplied. The running equation
is coefficientwise at $g_H^2\hbar$, not an infinite-scale small-coupling
trajectory. No active action, covariance, contour or density is reset.

\subsection*{Open}
The two-field row in \eqref{eq:v71-components} contains the genuine
three-line scalar sunset; the zero-field row contains the four-line
vacuum graph. Their proper subtractions, including the new coupling
restriction and the already fixed mass restriction, must be assembled
before claiming a full $g_H^2$ continuum. Independent metric marks and
the mixed internal-gauge, Yukawa and chiral sectors require their
native vertices and common source equations. The explicit six-field
return identifies a component the iterative class must retain; it does
not establish its full interacting self-map or all-step source bounds.
The next useful calculation is the complete renormalized scalar
sunset/marked two-point coefficient with this supplied quartic
subtraction, not another change of tadpole constant.

\begin{center}\footnotesize
\begin{tabular}{p{.23\textwidth}p{.69\textwidth}}
\toprule
Variable & Scope of this installment\\\midrule
Coupling and loops & Exact second cumulant at $g_H^2$; new continuum
estimate for its one-loop four-field component. Two- and three-loop
coefficients are algebraically recorded but not renormalized here.\\
Cutoff and volume & $O(a^2)$ four-field complement; fixed-frame terms
also retain $O((a\lambda)^4)$; separate Poisson volume bounds and no
total-volume multiplier in rooted source estimates.\\
Mass and sources & $0\le m_H\le M\lambda$ for the bubble comparison;
zero-mass coupling anchor and flow; every fixed soft frame arity has
its own constants. General finite frames appear in the algebraic
formula, not in an unbounded-source analytic assertion.\\
RG steps & Exact two-step coefficient composition; no uniform estimate
in an unlimited number of steps and no analytic ultraviolet scalar
trajectory.\\
Verification & Exact polynomial and rational finite-covariance tests,
including a mixed-shell return; analytic symbol arguments remain
proofs in the text, not numerically certified estimates.\\\bottomrule
\end{tabular}
\end{center}

\begin{thebibliography}{99}
\raggedright
\bibitem{V71KP} M.~V.~Kompaniets and E.~Panzer,
\emph{Minimally subtracted six loop renormalization of $O(n)$-symmetric
$\phi^4$ theory and critical exponents}, Phys. Rev. D \textbf{96}
(2017), 036016; arXiv:1705.06483v2.
\bibitem{V71KM} T.~Krajewski and P.~Martinetti,
\emph{Wilsonian renormalization, differential equations and Hopf
algebras}, Contemp. Math. \textbf{539} (2011), 187;
arXiv:0806.4309v2.
\bibitem{V71KMM} I.~Khavkine, A.~Melati and V.~Moretti,
\emph{On Wick polynomials of boson fields in locally covariant
algebraic QFT}, arXiv:1710.01937v2 (2018).
\bibitem{V71IIM} Y.~Igarashi, K.~Itoh and T.~R.~Morris,
\emph{BRST in the Exact RG}, Prog. Theor. Exp. Phys. (2019),
103B01; arXiv:1904.08231v3.
\end{thebibliography}

\section*{Append-only continuation marker: V71}
\addcontentsline{toc}{section}{Append-only continuation marker: V71}
\label{sec:v71-continuation-marker}
\textbf{End of V71.} Preserve Sections 1--550. The two-quartic
coefficient supplies its common one-loop quartic subtraction and a
coefficientwise marginal update. A generated six-field return is
necessary for exact composition and has an explicitly nonzero local
feedback. The remaining sunset, vacuum, combined BV and interacting
ESM continuum problems are not declared solved.



\clearpage
\part*{V72: the renormalized scalar sunset and the generated kinetic coefficient}
\addcontentsline{toc}{part}{V72: the renormalized scalar sunset and the generated kinetic coefficient}

\section{The two-field row after the quartic subtraction}
\label{sec:v72-start}

Sections 1--550 and the V71 continuation marker are preserved.  V71
identified the complete two-field term at order $g_H^2\hbar^2$ but left its
overlapping one-loop subtractions and overall two-point normalization open.
This installment closes that scalar two-point interface at zero Higgs mass.
It uses the quartic counterterm already fixed in V71 and does not choose a
second independent four-field normalization.

Write, with the same completed scalar covariance,
\begin{equation}
 c_a(p)=\frac{\nu_\lambda(p)}{P_a(p)},\qquad
 \mathcal S_{a,L}(k)=\frac1{V_L^2}\sum_{p,q}
 c_a(p)c_a(q)c_a(k-p-q),
 \label{eq:v72-sunset}
\end{equation}
where the quotient is zero on retained modes.  The scalar-two-point row of
\eqref{eq:v71-components} is therefore
\begin{equation}
 \boxed{
 \Gamma^{(2),\mathrm{sun}}_{a,L}(k)
 =-16(n_H+2)g_H^2\hbar^2\,\mathcal S_{a,L}(k),
 \qquad n_H=4.}
 \label{eq:v72-two-point-row}
\end{equation}
The displayed coefficient multiplies
$\sum_A\phi_A(-k)\phi_A(k)$; no extra factor $1/2$ is hidden in
\eqref{eq:v72-two-point-row}.

There are three proper logarithmic four-point subgraphs, obtained by deleting
one of the three scalar lines.  Each is normalized by the \emph{same} V71
quartic action coefficient.  Because the proper subtraction has degree zero
in all boundary momenta, closing its remaining scalar line produces a local
mass insertion.  Denote their sum, including the chosen V70/V71 normal-product
mass convention, by $\mathfrak m_{a,L}^{\rm sub}$.  It is independent of the
external two-point momentum.

\begin{proposition}[Proper quartic subtractions are local in the final two-point row]
\label{prop:v72-proper-local}
Let $t_{\gamma_i,0}$ be the degree-zero Taylor operation on the complete
four-point subgraph $\gamma_i$, $i=1,2,3$, with the V71 common quartic
normalization inserted before the complementary line is closed.  Then
\begin{equation}
 \boxed{
 \sum_{i=1}^{3}t_{\gamma_i,0}\mathcal S_{a,L}(k)
 =\mathfrak m_{a,L}^{\rm sub},\qquad
 \partial_k\mathfrak m_{a,L}^{\rm sub}=0.}
 \label{eq:v72-proper-local}
\end{equation}
Consequently, if $\mathbf T_{2,k}$ stores the complete even two-point Taylor
polynomial through momentum degree two,
\begin{equation}
 \boxed{
 (1-\mathbf T_{2,k})
 \left(1-\sum_i t_{\gamma_i,0}\right)\mathcal S_{a,L}
 =(1-\mathbf T_{2,k})\mathcal S_{a,L}.}
 \label{eq:v72-complement-equality}
\end{equation}
The proper subtractions are nevertheless required before scale-wise absolute
bounds are taken.
\end{proposition}
\begin{proof}
A proper subgraph contains two vertices and two scalar lines.  Its independent
external variables are the physical two-point momentum together with the
momentum of the omitted scalar line.  The degree-zero Taylor operation sets all
of these boundary variables to zero inside the subgraph.  The resulting local
four-scalar vertex is precisely the V71 common quartic counterterm tensor.
Closing the omitted line gives a tadpole at one local vertex and therefore no
remaining dependence on the physical two-point momentum.  The three flavour
channels are already combined in the common action coefficient; replacing
them by three copies of a scalar bubble would give the wrong $O(n_H)$
multiplicity.  Equation~\eqref{eq:v72-complement-equality} follows because
$\mathbf T_{2,k}$ contains the constant local mass slot.  The last statement is
important: on a separated scale assignment the bare high bubble and its local
Taylor term need not have the same support, so their signed difference must be
formed before estimating it, exactly as in the V62 overlap calculation.
\end{proof}

Thus this installment uses the forest representative
\begin{equation}
 \mathcal S^{R}_{a,L}
 =(1-\mathbf T_{2,k})
 \left(1-\sum_{i=1}^{3}t_{\gamma_i,0}\right)\mathcal S_{a,L},
 \label{eq:v72-forest}
\end{equation}
while \eqref{eq:v72-complement-equality} may be used only after the forest has
been assembled.  A finite change of the V71 quartic anchor changes
$\mathfrak m_{a,L}^{\rm sub}$ and hence the local two-point mass coefficient,
but it does not change the nonlocal complement or the kinetic logarithm below.

\section{Universal kinetic logarithm of the massless sunset}
\label{sec:v72-log}

The wave-function logarithm can be extracted without evaluating a cutoff-
dependent quadratic divergence.  Consider first the ordinary massless
integral in $d=4-2\varepsilon$ dimensions,
\begin{equation}
 I_d(k)=\mu^{4\varepsilon}
 \int\frac{d^dp}{(2\pi)^d}\frac{d^dq}{(2\pi)^d}
 \frac1{p^2q^2(p+q+k)^2}.
 \label{eq:v72-Id}
\end{equation}

\begin{lemma}[Exact massless sunset integral]
\label{lem:v72-sunset-gamma}
For $3<\Re d<4$ and by meromorphic continuation elsewhere,
\begin{equation}
 \boxed{
 I_d(k)=
 \frac{\mu^{4\varepsilon}(k^2)^{d-3}}{(4\pi)^d}
 \frac{\Gamma(d/2-1)^3\Gamma(3-d)}{\Gamma(3d/2-3)}.}
 \label{eq:v72-gamma}
\end{equation}
At $d=4-2\varepsilon$,
\begin{equation}
 \boxed{
 I_d(k)=-\frac{k^2}{1024\pi^4}\frac1\varepsilon+O(1).}
 \label{eq:v72-pole}
\end{equation}
\end{lemma}
\begin{proof}
Integrate first in $q$.  The standard one-loop massless convolution gives
\[
 \int_q\frac1{q^2(q+p+k)^2}
 =(4\pi)^{-d/2}\frac{\Gamma(2-d/2)\Gamma(d/2-1)^2}
 {\Gamma(d-2)}[(p+k)^2]^{d/2-2}.
\]
Apply the same convolution formula to the remaining $p$ integral with powers
$1$ and $2-d/2$.  The factors $\Gamma(2-d/2)$ and $\Gamma(d-2)$ cancel,
leaving \eqref{eq:v72-gamma}.  Since
$\Gamma(3-d)=\Gamma(-1+2\varepsilon)=-1/(2\varepsilon)+O(1)$ and
$\Gamma(3d/2-3)=\Gamma(3-3\varepsilon)=2+O(\varepsilon)$,
\eqref{eq:v72-pole} follows.  No infrared scale is introduced in the finite
programme by this auxiliary calculation; it is used only to identify the
homogeneous ultraviolet residue.
\end{proof}

The conversion from the pole to the native cutoff logarithm is calibrated by
the one-loop bubble already evaluated in V71.  In the same dimensional
convention the massless bubble has pole $1/(16\pi^2\varepsilon)$, whereas the
native shell coefficient is $(8\pi^2)^{-1}\log(1/(a\lambda))$.  Hence
$1/\varepsilon$ corresponds to $2\log(1/(a\lambda))$ for the common
homogeneous residue.

\begin{theorem}[Native wave-function logarithm]
\label{thm:v72-wave-log}
For the completed scalar regulator and the forest prescription
\eqref{eq:v72-forest}, the local kinetic part of the massless sunset is
\begin{equation}
 \boxed{
 \mathbf T_{2,k}\mathcal S_{a,\infty}^{\rm forest}(k)
 =\mathfrak m_{a}^{(2)}
 -\frac{k^2}{512\pi^4}\log\frac1{a\lambda}
 +k_i k_j\,\mathfrak z^{\rm fin}_{a,ij}.}
 \label{eq:v72-sunset-local}
\end{equation}
Here $\mathfrak m_a^{(2)}$ is the complete local mass coefficient, including
the proper-subgraph and normal-product convention, and
$\mathfrak z^{\rm fin}_{a,ij}$ is a bounded finite matching tensor after the
logarithm is removed.  Neither is assigned a numerical value.

Consequently the two-point effective action contains the universal local term
\begin{equation}
 \boxed{
 +\frac{n_H+2}{32\pi^4}g_H^2\hbar^2
 \log\frac1{a\lambda}
 \sum_A\int (\partial_\mu\phi_A)(\partial_\mu\phi_A),}
 \label{eq:v72-loop-kinetic}
\end{equation}
and the corresponding counterterm has the opposite sign.
Equivalently, in the convention
$\frac12Z_H\sum_A\int(\partial\phi_A)^2$,
\begin{equation}
 \boxed{
 \delta Z_H^{\log}
 =-\frac{n_H+2}{16\pi^4}g_H^2\hbar^2
 \log\frac1{a\lambda}.}
 \label{eq:v72-Z}
\end{equation}
For $\lambda_4=24g_H$, its magnitude is
$(n_H+2)\lambda_4^2/[36(16\pi^2)^2]$, in the standard $O(n_H)$ convention.
\end{theorem}
\begin{proof}
The logarithmic coefficient is determined by the homogeneous ultraviolet
residue.  The completed native symbol agrees with the ordinary scalar symbol
through the order needed to leave the logarithmic residue unchanged; its
upper completion and the fixed lower reference change only the finite local
matching tensor.  Lemma~\ref{lem:v72-sunset-gamma}, together with the V71
one-loop calibration, therefore gives the coefficient
$-1/(512\pi^4)$ in \eqref{eq:v72-sunset-local}.  Proper quartic subtractions
are momentum independent after the complementary line is closed, by
Proposition~\ref{prop:v72-proper-local}, so they cannot change this coefficient.
Multiplying by the exact row factor in \eqref{eq:v72-two-point-row} gives
$16/512=1/32$ and proves \eqref{eq:v72-loop-kinetic}.  The factor two between
that action coefficient and \eqref{eq:v72-Z} comes only from the standard
$\frac12Z_H$ normalization.  No continuum anomalous dimension has been inserted
as an assumption.
\end{proof}

The logarithm is isotropic.  The finite tensor $\mathfrak z^{\rm fin}_{a,ij}$
need not be.  A physical normalization must therefore specify its finite kinetic
anchor rather than silently replacing it by its rotational average.

\section{All-relative-scale cutoff estimate for the sunset complement}
\label{sec:v72-estimate}

The local logarithm does not by itself prove that the remaining two-point
coefficient has a continuum limit.  The required estimate uses the signed
proper subtractions on scale-separated assignments.

Let $R$ be the larger internal shell scale and $r$ the smaller.  For a proper
high bubble, parity removes its linear boundary term.  After its degree-zero
subtraction, its differentiated shell amplitude obeys
\begin{equation}
 \left|\partial_U^t(1-t_{\gamma,0})H_R(U)\right|
 \le C R^{-2-t}r^{2-t},\qquad 0\le t\le2,
 \label{eq:v72-highbubble}
\end{equation}
when $|U|\lesssim r$; for higher $t$ use the corresponding homogeneous
$R^{-t}$ bound.  The finite-to-ordinary difference gains the inherited scalar
fourth-order factor and satisfies, at the two-derivative level,
\begin{equation}
 \left|\partial_U^2\left[(1-t_{\gamma,0})H_{a,R}
 -(1-t_{\gamma,0})H_{0,R}\right]\right|
 \le C a^4R^2.
 \label{eq:v72-highbubble-diff}
\end{equation}
These estimates are statements about the signed high bubble after subtraction;
estimating its bare and Taylor pieces separately loses the scale gain.

\begin{theorem}[Full sunset complement has an $O(a^2\log)$ cutoff comparison]
\label{thm:v72-sunset-bound}
Fix the V71 scalar reference, the common quartic normalization, a positive lower
scale $\lambda$, and finitely many source derivatives and rooted moments.  It
suffices that
\begin{equation}
 \lambda\ge32768\mu>0,\qquad a\lambda\le2^{-24},\qquad
 \lambda L_\nu\ge1,\qquad m_H=0.
 \label{eq:v72-domain}
\end{equation}
Put
$\mathfrak L_a=1+\log(1/(a\lambda))$.  For $j=|\alpha|\le4$,
\begin{equation}
 \boxed{
 \left|\partial_k^\alpha
 \left(\mathcal S^R_{a,L}-\mathcal S^R_{0,L}\right)(k)\right|
 \le C_\alpha a^2\mathfrak L_a|k|^{4-j}.}
 \label{eq:v72-cutoff-bound}
\end{equation}
The separate ordinary finite-volume comparison is
\begin{equation}
 \boxed{
 \left|\partial_k^\alpha
 \left(\mathcal S^R_{0,L}-\mathcal S^R_{0,\infty}\right)(k)\right|
 \le C_{\alpha,b}\lambda^{-2}(\lambda L_{\min})^{-b}|k|^{4-j},
 \qquad b>8.}
 \label{eq:v72-volume-bound}
\end{equation}
All completed Brillouin-zone modes and all relative internal scales are included.
\end{theorem}
\begin{proof}
Assign a shell to every primitive scalar covariance.  When all three scales are
comparable to $R$, four external derivatives make the ordinary integrand
$O(R^{-10})$; the two-loop mode count is $O(R^8)$, hence the shell is
$O(R^{-2})$.  The inherited scalar-symbol mismatch is fourth relative mesh
order.  At the same derivative order it contributes $O(a^4R^{-6})$ before the
mode count and therefore $O(a^4R^2)$ after it.  Summing comparable scales up to
$R\asymp a^{-1}$ gives $O(a^2)$, and the ordinary ultraviolet tail gives the
same order.

Now suppose two lines form a high bubble at scale $R\gg r$, while the
complementary line has scale $r$.  Form the signed bubble subtraction before
absolute values.  Distribute the four external derivatives, with $t$ landing on
the high remainder.  The high integrated remainder contributes
$R^{-2}r^{2-t}$, while the differentiated low covariance and its normalized
loop count contribute $r^{t-2}$.  Their product is $CR^{-2}$, independently of
$r$.  There are at most $C(j+1)$ lower-shell choices at largest shell
$R_j\asymp2^j\lambda$, so the ordinary sum is
$C\sum_j(j+1)R_j^{-2}<\infty$.

For the finite-to-ordinary difference, \eqref{eq:v72-highbubble-diff} and the
same low-chain derivative count give $Ca^4R^2$ per admissible lower scale.
Thus
\[
 a^4\sum_{j\le J}(j+1)R_j^2
 \le Ca^2(1+J),\qquad R_J\asymp a^{-1}.
\]
The finite upper region and the ordinary tail obey the same
$a^2(1+J)$ bound.  This proves \eqref{eq:v72-cutoff-bound}.  Scale regions in
which a proper Taylor term survives while the bare assigned graph vanishes are
included by the same signed estimate; they are the scalar analogue of the V62
counterterm-only regions.  No product of two proper subtractions appears because
the three bubble subgraphs overlap pairwise.

For finite volume, apply Poisson summation after the four derivatives to each
smooth shell word.  Eight integrations by parts are sufficient because there
are two loop momenta; arbitrary fixed $b>8$ gives the displayed bound.  Integrate
the stored Taylor remainder back to order zero.  The subtraction-created local
mass terms are part of $\mathbf T_{2,k}$ and do not enter the complement.
\end{proof}

\begin{corollary}[Rooted two-point source estimate]
\label{cor:v72-rooted}
For the effective two-point row \eqref{eq:v72-two-point-row}, after factoring its
flavour polarization,
\begin{equation}
 \boxed{
 a^4\sum_x(1+\mu d_{\rm per}(x,0))^b
 \|\mathcal K^\Delta_{a,L}(x)\|
 \le C_b(n_H+2)|g_H|^2\hbar^2
 a^2\mathfrak L_a\mu^4.}
 \label{eq:v72-rooted}
\end{equation}
There is no total-volume multiplier.  The estimate is uniform in mesh and
periods on \eqref{eq:v72-domain}, at fixed derivative/source order and profile
bounds; it is not uniform in unbounded perturbative order or in the number of
RG steps.
\end{corollary}
\begin{proof}
Fourier integration by parts converts the differentiated bound of
Theorem~\ref{thm:v72-sunset-bound} into the stated weighted rooted norm.  The
panel has one independent relative coordinate after translation invariance.  The
local degree-two polynomial is stored separately and therefore contributes no
nonlocal kernel tail.
\end{proof}

\section{The common kinetic counterterm and its marked source completion}
\label{sec:v72-source-completion}

The logarithm in \eqref{eq:v72-Z} belongs to one action coefficient.  Its source
contacts cannot be normalized independently.

For the local flavour frame used in V67--V71, and for the retained metric source,
define the ordinary kinetic functional
\begin{equation}
 \mathscr K_H[g,T,\phi]
 =\frac12\sum_A\int\rho\,g^{\mu\nu}
 \partial_\mu(T\phi)_A\partial_\nu(T\phi)_A.
 \label{eq:v72-kinetic-functional}
\end{equation}
The universal logarithmic kinetic counterterm is
\begin{equation}
 \boxed{
 C_{H,2}^{\rm kin,log}
 =-\frac{n_H+2}{16\pi^4}g_H^2\hbar^2
 \log\frac1{a\lambda}\,\mathscr K_H.}
 \label{eq:v72-kinetic-counterterm}
\end{equation}
At $g=I$, $T=I$, this is exactly the opposite of
\eqref{eq:v72-loop-kinetic}.

\begin{proposition}[One common kinetic source packet]
\label{prop:v72-source-packet}
Every retained metric or frame mark of the logarithmic two-point normalization is
obtained by differentiating \eqref{eq:v72-kinetic-functional}.  In particular,
for $T=e^\Sigma$ the first local frame derivative contains both
\begin{equation}
 \partial_\mu\phi^T\Sigma\,\partial_\mu\phi
 \quad\text{and}\quad
 \partial_\mu\phi^T(\partial_\mu\Sigma)\phi,
 \label{eq:v72-frame-contacts}
\end{equation}
and higher marks contain the corresponding noncommuting matrix contacts.  A
finite change of the quartic V71 anchor changes only the local mass part of the
sunset insertion at this order; it does not change the coefficient of
\eqref{eq:v72-kinetic-counterterm}.
\end{proposition}
\begin{proof}
Differentiate $\partial(T\phi)=T\partial\phi+(\partial T)\phi$ before setting
the source to zero.  The two terms in \eqref{eq:v72-frame-contacts} are therefore
part of one action Hessian and cannot be separated by source normalization.  The
second statement follows from Proposition~\ref{prop:v72-proper-local}: the
quartic finite addition closes into a momentum-independent tadpole in the final
two-point row, annihilated by the kinetic projection.
\end{proof}

The bounded finite kinetic tensor remains common matching data.  V72 does not set
it to zero.  If the programme later adopts a different finite wave-function
normalization, all of its metric/frame/source derivatives must be converted with
it.

\section{What the sunset closes, and what it does not}
\label{sec:v72-closure}

The sunset is the first $g_H^2$ coefficient that requires both an overlapping
one-loop subtraction and a new overall kinetic normalization.  The preceding
results give one complete scalar-sector statement:

\begin{theorem}[Renormalized scalar sunset through the marked two-point panel]
\label{thm:v72-combined}
At zero Higgs mass and at bidegree $g_H^2\hbar^2$, the V70/V71 scalar action has a
renormalized two-point coefficient
\begin{equation}
 \boxed{
 \Gamma^{(2),R}_{a,L}
 =-16(n_H+2)g_H^2\hbar^2\mathcal S^R_{a,L}
 +\Gamma^{(2),\rm loc}_{a,L},}
 \label{eq:v72-combined}
\end{equation}
where $\Gamma^{(2),\rm loc}$ contains the supplied local mass normalization and
the common kinetic counterterm \eqref{eq:v72-kinetic-counterterm}.  Its nonlocal
part has the cutoff, volume and rooted estimates of
Theorem~\ref{thm:v72-sunset-bound} and Corollary~\ref{cor:v72-rooted}.
The kinetic logarithm is independent of the finite V71 quartic anchor; the local
mass coefficient is not.

Every fixed metric/frame derivative of the logarithmic kinetic block is generated
by the same action functional.  This theorem includes all relative scalar loop
scales and the common quartic subdivergence.  It does not include internal gauge,
Yukawa, ghost or gravitational loops.
\end{theorem}
\begin{proof}
Combine Proposition~\ref{prop:v72-proper-local},
Theorems~\ref{thm:v72-wave-log} and \ref{thm:v72-sunset-bound}, and
Proposition~\ref{prop:v72-source-packet}.  The proper forest fixes the separated
scales, the overall degree-two assignment stores the remaining local mass and
kinetic data, and the same action functional supplies every marked derivative.
\end{proof}

The zero-field row of V71 still contains the four-line vacuum graph.  Normalizing
$Z[0]$ removes its unmarked value but not its metric/frame derivatives.  Those
marked derivatives require their own common forest with the populated quartic,
mass and kinetic restrictions.  The scalar two-point theorem above does not
supply that forest automatically.

\section{Status and next load-bearing step}
\label{sec:v72-ledgers}

\subsection*{Proved}
The complete $g_H^2\hbar^2$ scalar sunset is assembled with the V71 common
quartic subdivergence and a degree-two overall assignment.  The universal
wave-function logarithm is evaluated from the massless two-loop integral and
matched to the native cutoff convention.  The all-relative-scale nonlocal
remainder has an $O(a^2[1+\log(1/(a\lambda))])$ cutoff bound, separate volume
control, and a rooted source estimate.  The logarithmic kinetic counterterm is
written as one marked action functional, fixing its metric/frame contacts.

\subsection*{Inherited}
V27 supplies the completed scalar covariance and its fourth-order interior symbol
comparison.  V70 supplies the common one-vertex quartic/mass/vacuum packet.  V71
supplies the complete two-vertex cumulant, the common quartic local coefficient,
and the exact mixed-shell return needed for iterative composition.  V54--V57
supply finite-range and measured-source bookkeeping used only at their stated
scope; none is reinterpreted as a two-loop scalar theorem.

\subsection*{Literature input}
The massless convolution formula, BPHZ/forest organization and the standard
$O(n)$ scalar comparison are established.  The numerical kinetic residue used
here is nevertheless derived explicitly in Lemma~\ref{lem:v72-sunset-gamma} and
calibrated against V71's native one-loop logarithm.  No continuum anomalous
dimension is inserted as an assumption.

\subsection*{Conditional}
The finite local mass coefficient and finite anisotropic kinetic tensor are common
normalization data and have not been numerically populated.  The theorem is
massless for the nonlocal sunset estimate; a full mass-polynomial two-point bank
requires the joint $z$-Taylor assignment.  Marked source bounds beyond fixed
source order require their own constants.  No estimate is uniform in unbounded
perturbative order or an unlimited number of RG steps.

\subsection*{Open}
The V71 four-line zero-field graph still needs its marked vacuum forest.  More
importantly for the full ESM programme, the scalar sector now contains populated
quartic and kinetic running data but still lacks a proved invariant iterative
state space carrying them, their six-field return, and the source transports
through arbitrarily many shells.  Internal gauge, Yukawa, fermionic/chiral and
gravity sectors must enter the same common action before a full ESM theorem can be
stated.

The next useful move is therefore upstream: use the now-populated scalar examples
(quartic marginal, generated six-field irrelevant kernel, and two-loop kinetic
coordinate) to formulate and prove the first filtered perturbative state-space
self-map, rather than continuing indefinitely with isolated scalar graphs.

\begin{center}\footnotesize
\begin{tabular}{p{.23\textwidth}p{.69\textwidth}}
\toprule
Variable & Scope of V72\\\midrule
Coupling order & $g_H^2$: quartic one-loop and scalar two-point two-loop rows are now both normalized; the three-loop vacuum row is not.\\
Cutoff & Sunset nonlocal complement: $O(a^2[1+\log(1/(a\lambda))])$ at fixed lower scale.\\
Volume & Separate $\lambda^{-2}(\lambda L_{\min})^{-b}$ two-loop comparison, $b>8$.\\
Sources & One common kinetic action fixes every fixed metric/frame mark of the logarithm; finite source order only.\\
RG steps & No uniform many-step theorem; V71 proves exact two-step quartic composition only.\\
Verification & Exact convolution residue, normalization factors and append-only preservation are checked; analytic shell estimates remain proofs in the text.\\
\bottomrule
\end{tabular}
\end{center}

\begin{thebibliography}{99}
\raggedright
\bibitem{V72KP} M.~V.~Kompaniets and E.~Panzer,
\emph{Minimally subtracted six loop renormalization of $O(n)$-symmetric
$\phi^4$ theory and critical exponents}, Phys. Rev. D \textbf{96}
(2017), 036016; arXiv:1705.06483v2.
\bibitem{V72Z} W.~Zimmermann,
\emph{Convergence of Bogoliubov's method of renormalization in momentum
space}, Commun. Math. Phys. \textbf{15} (1969), 208--234.
\end{thebibliography}

\section*{Append-only continuation marker: V72}
\addcontentsline{toc}{section}{Append-only continuation marker: V72}
\label{sec:v72-continuation-marker}
\textbf{End of V72.} Preserve Sections 1--557.  The scalar sunset is now
renormalized through its common quartic subdivergence and overall kinetic/mass
assignment, with the universal kinetic logarithm and an all-scale cutoff estimate.
The next load-bearing task is upstream: construct the filtered iterative ESM action
class that retains the generated quartic, kinetic and six-field coordinates
rather than adding another isolated scalar graph.

\clearpage
\part*{V73: A filtered scalar ESM spine and an exact two-loop shell self-map}
\addcontentsline{toc}{part}{V73: Filtered scalar ESM spine and exact two-loop shell self-map}

\section{Upstream pivot: from isolated graphs to an iterative finite-order state}
\label{sec:v73-pivot}

V70--V72 populate three coefficients that cannot be treated independently in an
iterated shell calculation: the quartic marginal, the generated six-field return,
and the two-loop kinetic row.  The purpose of this installment is to package them
into the first explicit filtered state on which a shell step closes.  This is the
scalar $O(n_H)$ spine of the Einstein--Standard-Model (ESM) programme, not a
replacement for the full ESM state.

For the physical specialization we set
\begin{equation}
 n:=n_H=4
 \label{eq:v73-n}
\end{equation}
in the physical specialization, but retain symbolic $n$ in the combinatorial
identities.  We truncate at coupling degree $g_H^2$, loop degree $\hbar^2$, and
local EFT dimension six.  Connected coefficients with at least two retained
scalar legs are included.  The $g_H^2\hbar^3$ zero-field graph of V71 is outside
this loop truncation; its marked vacuum forest therefore remains an independent
open problem rather than being silently set to zero.

The one-vertex quartic interaction is always understood in V70's common
normal-product packet.  Thus its tadpole mass restriction and vacuum restriction
move with the quartic coefficient.  They are not regenerated as independent
parameters at every shell.  The retained scalar field is not canonically rescaled
between shells in this installment.  Wave-function renormalization is stored as a
local kinetic coefficient, so the scalar source coordinates themselves have the
identity adjacent transport on this spine.

\subsection{The local and complementary coordinates}

Let $F_m$ denote a connected scalar vertex with $m$ retained scalar legs.  At the
present order only $m=6,4,2$ occurs in the genuinely new two-vertex sector.  Let
$\mathsf T_d$ denote the joint even external-momentum Taylor polynomial of total
degree at most $d$ at the origin.  We use the exact decompositions
\begin{align}
 F_6&=\underbrace{\mathsf T_0F_6}_{\text{local dimension six}}
      +\underbrace{(1-\mathsf T_0)F_6}_{R_6},
 \label{eq:v73-split6}\\
 F_4&=\underbrace{\mathsf T_0F_4}_{\text{quartic physical local}}
      +\underbrace{(\mathsf T_2-\mathsf T_0)F_4}_{\text{local dimension six}}
      +\underbrace{(1-\mathsf T_2)F_4}_{R_4},
 \label{eq:v73-split4}\\
 F_2&=\underbrace{\mathsf T_2F_2}_{\text{mass/kinetic physical local}}
      +\underbrace{(\mathsf T_4-\mathsf T_2)F_2}_{\text{local dimension six}}
      +\underbrace{(1-\mathsf T_4)F_2}_{R_2}.
 \label{eq:v73-split2}
\end{align}
The first terms in \eqref{eq:v73-split4}--\eqref{eq:v73-split2} are the usual
renormalizable scalar coordinates.  The three middle/local entries
\begin{equation}
 w^{\rm loc}_6
 =\bigl(\mathsf T_0F_6,
        (\mathsf T_2-\mathsf T_0)F_4,
        (\mathsf T_4-\mathsf T_2)F_2\bigr)
 \label{eq:v73-zloc}
\end{equation}
are the first scalar dimension-six EFT bank.  They are local Wilson
coordinates and therefore belong to the two-boundary local bank, not to the
ultraviolet remainder.  The $R_m$ are the complete quasilocal complements.  Nothing is deleted by the Taylor maps: recombining all
entries in \eqref{eq:v73-split6}--\eqref{eq:v73-split2} recovers the original
vertices exactly.

The physical local coordinate is
\begin{equation}
 u=(m_H^2,Z_H,g_H),
 \label{eq:v73-u}
\end{equation}
with the V70 vacuum restriction carried as a triangular normalization coordinate
when zero-field quantities are requested.  The finite filtered scalar state is
therefore
\begin{equation}
 \boxed{
 \mathfrak S^H_\lambda
 =\bigl(u_\lambda,w^{\rm loc}_{6,\lambda},
        R_{6,\lambda},R_{4,\lambda},R_{2,\lambda};\mathcal J_\lambda\bigr),}
 \label{eq:v73-state}
\end{equation}
where $\mathcal J_\lambda$ is any fixed finite scalar source panel obtained by
differentiating the same action coefficients.  In the unrescaled-field convention
used here its adjacent coordinate map is the identity.  Pure background-source
vacuum coefficients beyond loop degree two are not included in
\eqref{eq:v73-state}.

For later scale estimates define the dimensionless rescaled $m$-point vertex
\begin{equation}
 \widehat F_{m,\lambda}(q_1,\ldots,q_{m-1})
 :=\lambda^{m-4}F_m(\lambda q_1,\ldots,\lambda q_{m-1}),
 \qquad |q_i|\le 2^{-7}.
 \label{eq:v73-rescaled-vertex}
\end{equation}
A local operator of canonical dimension six has coefficient dimension $-2$;
therefore its dimensionless coefficient scales as $\lambda^2$.  This gives the
same excess degree two for the three entries in \eqref{eq:v73-zloc}.

\begin{proposition}[Finite scalar filter at $(g_H^2,\hbar^2,\Delta=6)$]
\label{prop:v73-finite-filter}
At coupling degree at most two and loop degree at most two, the connected
$O(n)$ scalar sector generated by the quartic interaction and containing at least
two retained scalar legs is contained in the finite state
\eqref{eq:v73-state}.  At local dimension at most six, no additional scalar tensor
type is produced by Gaussian contraction: the new local types are exactly
$\phi^6$, a two-derivative four-field type, and a four-derivative two-field type,
together with the renormalizable mass, kinetic and quartic coordinates.
\end{proposition}
\begin{proof}
Two quartic vertices have eight scalar legs.  A connected two-vertex contraction
uses at least one line, hence leaves at most six external scalar legs.  Through
loop degree two it uses at most three contraction lines and therefore leaves
respectively six, four or two retained legs.  Each contraction removes two fields
of total canonical dimension two and contributes one scalar covariance of
canonical dimension minus two, so the total local canonical dimension remains six
for the first nonrenormalizable descendant.  At arity six this is the zero-momentum
jet; at arity four it is the degree-two jet; at arity two it is the degree-four
jet.  Higher Taylor degrees have dimension at least eight and remain in $R_m$.
The one-vertex normal-product packet contributes only the already stored
renormalizable mass/quartic/vacuum restrictions.  This proves the finite list.
\end{proof}

\section{The exact contraction cascade}
\label{sec:v73-cascade}

The key closure is combinatorial and finite.  It is useful to display it before
any scale estimate is taken.

For two spacetime points $x,y$, write
\begin{equation}
 U_x=\phi(x)^T\phi(x),\qquad
 U_y=\phi(y)^T\phi(y),\qquad
 W_{xy}=\phi(x)^T\phi(y),
 \label{eq:v73-UW}
\end{equation}
and define
\begin{align}
 P_6&=U_xU_yW_{xy},
 \label{eq:v73-P6}\\
 P_4&=(n+4)U_xU_y+4W_{xy}^2,
 \label{eq:v73-P4}\\
 P_2&=W_{xy}.
 \label{eq:v73-P2}
\end{align}
Let
\begin{equation}
 \mathscr D_{xy}
 =\sum_{A=1}^{n}
 \frac{\partial^2}{\partial\phi_A(x)\,\partial\phi_A(y)}
 \label{eq:v73-cross-D}
\end{equation}
be the cross-contraction differential.  Same-point contractions are already
accounted for by the common V70 normal-product restrictions.

\begin{lemma}[The three invariant contraction identities]
\label{lem:v73-contractions}
The polynomials \eqref{eq:v73-P6}--\eqref{eq:v73-P2} satisfy
\begin{equation}
 \boxed{
 \mathscr D_{xy}P_6=P_4,\qquad
 \mathscr D_{xy}P_4=12(n+2)P_2,\qquad
 \mathscr D_{xy}^2P_6=12(n+2)P_2.}
 \label{eq:v73-contractions}
\end{equation}
\end{lemma}
\begin{proof}
For the first identity differentiate $U_xU_yW$ once at each endpoint.  The
terms in which both derivatives hit $W$ give $nU_xU_y$; the terms in which one
hits $W$ and one hits a quadratic factor give $4U_xU_y$; the terms in which the
derivatives hit the two quadratic factors give $4W^2$.  This yields
$(n+4)U_xU_y+4W^2$.

For the second identity,
$\mathscr D_{xy}(U_xU_y)=4W$, while
$\mathscr D_{xy}(W^2)=2(n+1)W$.  Hence
\[
 \mathscr D_{xy}P_4
 =4(n+4)W+8(n+1)W
 =12(n+2)W.
\]
The third identity is the composition of the first two.
\end{proof}

For a translation-invariant covariance kernel $c_{xy}$ define the connected
second-cumulant functional, modulo field-independent terms of loop degree at
least three, by
\begin{equation}
 \boxed{
 \begin{aligned}
 \mathfrak G[c]
 ={}&-8g_H^2\int_{x,y}c_{xy}P_6\\
 &-4g_H^2\hbar\int_{x,y}c_{xy}^2P_4\\
 &-16(n+2)g_H^2\hbar^2\int_{x,y}c_{xy}^3P_2.
 \end{aligned}}
 \label{eq:v73-Gc}
\end{equation}
This is exactly the six-, four- and two-field part of V71.
For an independent shell covariance $d$, define its normal-product cross transport
on a two-point polynomial by
\begin{equation}
 \mathsf H_d^\times
 :=\exp\!\left(\hbar d_{xy}\mathscr D_{xy}\right),
 \label{eq:v73-Hd}
\end{equation}
truncated consistently at loop degree two.

\begin{theorem}[Exact two-loop Gaussian semigroup on the filtered scalar spine]
\label{thm:v73-semigroup}
For any two independent translation-invariant scalar covariances $c,d$,
\begin{equation}
 \boxed{
 \mathfrak G[c+d]
 =\mathsf H_d^\times\mathfrak G[c]+\mathfrak G[d]
 \quad\bmod\bigl(g_H^3,\hbar^3,\text{field-independent terms}\bigr).}
 \label{eq:v73-semigroup}
\end{equation}
Equivalently, through loop degree two the two-vertex sector is closed under a
second Gaussian blocking step if and only if all three arity rows in
\eqref{eq:v73-Gc} are retained.
\end{theorem}
\begin{proof}
Use Lemma~\ref{lem:v73-contractions}.  The six-field row contributes under
\eqref{eq:v73-Hd}
\[
 -8g_H^2cP_6
 -8g_H^2\hbar cdP_4
 -48(n+2)g_H^2\hbar^2cd^2P_2.
\]
The four-field row contributes
\[
 -4g_H^2\hbar c^2P_4
 -48(n+2)g_H^2\hbar^2c^2dP_2,
\]
and the two-field row contributes itself through loop degree two.  Adding the
fresh-shell functional $\mathfrak G[d]$ gives the coefficients
\[
 c+d,
 \qquad c^2+2cd+d^2,
 \qquad c^3+3c^2d+3cd^2+d^3,
\]
which are exactly the three powers in $\mathfrak G[c+d]$.  No estimate or
continuum identity is used.
\end{proof}

\begin{corollary}[Both higher-field rows are load-bearing]
\label{cor:v73-load-bearing}
If the six-field row is discarded after the first shell, the next four-field
coefficient misses
\begin{equation}
 \boxed{-8g_H^2\hbar\,c_{xy}d_{xy}P_4.}
 \label{eq:v73-missing4}
\end{equation}
At two-point level it also misses the $cd^2$ contribution
\begin{equation}
 \boxed{-48(n+2)g_H^2\hbar^2c_{xy}d_{xy}^2P_2.}
 \label{eq:v73-missing2a}
\end{equation}
If the four-field nonlocal row is discarded, the next two-point coefficient
misses
\begin{equation}
 \boxed{-48(n+2)g_H^2\hbar^2c_{xy}^2d_{xy}P_2.}
 \label{eq:v73-missing2b}
\end{equation}
Thus a coupling-only state cannot compose even at this fixed perturbative order.
\end{corollary}
\begin{proof}
These are respectively the first contraction of the six-field row, its second
contraction, and the first contraction of the four-field row in the proof of
Theorem~\ref{thm:v73-semigroup}.  Their sum is exactly the mixed part of the
binomial square and cube.  V71 established nonvanishing of the first term for
actual separated shells; the identities here show that the same obstruction
propagates into the two-point row.
\end{proof}

Theorem~\ref{thm:v73-semigroup} is the finite-order form of conditional
cumulance for this populated sector.  Its importance here is not the abstract
probability identity but the explicit list of coordinates that must survive the
first shell for the second shell to be correct.

\section{The exact filtered one-step map}
\label{sec:v73-one-step}

Let $\lambda_{j+1}=b\lambda_j$, $b>1$, and let $d_j$ be one positive scalar
innovation between the two lower reference scales.  Denote by
$\mathsf{Rec}_{\lambda}$ the exact recombination of
\eqref{eq:v73-split6}--\eqref{eq:v73-split2}, and by
$\mathsf{Split}_{\lambda}$ the corresponding exact decomposition.
At coupling degree two define
\begin{equation}
 \boxed{
 \mathcal R_j^H
 :=\mathsf{Split}_{\lambda_j}\circ
 \left(\mathsf H_{d_j}^\times\circ
       \mathsf{Rec}_{\lambda_{j+1}}+\mathfrak G[d_j]\right),}
 \label{eq:v73-Rj}
\end{equation}
with the V70 one-vertex normal-product packet and the pre-existing tree/local
quadratic and quartic terms carried in the obvious triangular lower-degree rows.
The notation in \eqref{eq:v73-Rj} means addition of functionals before the final
split, not addition of coordinate vectors with incompatible arities.

\begin{theorem}[Algebraic one-step closure of the scalar filtered state]
\label{thm:v73-algebraic-selfmap}
At $(g_H^2,\hbar^2,\Delta=6)$, the map \eqref{eq:v73-Rj} sends every state of
the form \eqref{eq:v73-state} to another state of the same form.  In particular:
\begin{enumerate}[label=\textup{(\roman*)},leftmargin=2.8em]
\item no generated scalar interaction of arity greater than six occurs;
\item every generated local operator of dimension at most six is stored in
      $u$ or $w^{\rm loc}_6$;
\item every remaining generated term is in one of the complete complements
      $R_6,R_4,R_2$;
\item the local quartic and two-point updates are outputs of the same shell map,
      not externally supplied running data;
\item every fixed scalar source derivative of these nonzero-field coefficients is
      obtained by differentiating the same recombined functional before
      \eqref{eq:v73-split6}--\eqref{eq:v73-split2} are taken.
\end{enumerate}
The local projectors are therefore renormalization coordinates, not deletion
operators.
\end{theorem}
\begin{proof}
Theorem~\ref{thm:v73-semigroup} proves closure of the complete two-vertex
functional under another Gaussian innovation.  Gaussian contraction can only
lower arity by an even number.  Proposition~\ref{prop:v73-finite-filter}
classifies every resulting local type through dimension six.  The decompositions
\eqref{eq:v73-split6}--\eqref{eq:v73-split2} are exact linear identities, so their
complements retain all terms not assigned to the finite local bank.  The one-vertex
normal-product packet is already closed by V70 and is triangular in coupling
order.  Source differentiation commutes with finite Gaussian expectation and with
the finite Taylor maps when it is performed before the source is set to zero.
Thus no source contact generated in the stated nonzero-field panel is lost.
\end{proof}

\subsection{The internally generated local logarithms}

On the canonical generated trajectory, take a massless ordinary shell whose scale
ratio is $b$.  V71 and V72 give the logarithmic pieces of the local output of
\eqref{eq:v73-Rj}.  With $j+1$ the more ultraviolet scale and $j$ the result after
eliminating that shell,
\begin{align}
 \left[\Delta g_H\right]_{\log}
 &=-\frac{n+8}{2\pi^2}\,\hbar g_H^2\log b,
 \label{eq:v73-dg}\\
 \left[\Delta Z_H^{\rm loop}\right]_{\log}
 &=+\frac{n+2}{16\pi^4}\,\hbar^2g_H^2\log b.
 \label{eq:v73-dZ}
\end{align}
Here $\Delta Z_H^{\rm loop}$ denotes the loop coefficient in the effective
kinetic action; a counterterm defined to hold a fixed physical endpoint has the
opposite increment.  The power-sensitive mass coefficient remains a physical
local two-boundary coordinate.  Equations~\eqref{eq:v73-dg}--\eqref{eq:v73-dZ}
are not assumptions in \eqref{eq:v73-Rj}; they are evaluations of two of its local
projections.

For $n=4$,
\begin{equation}
 \left[\Delta g_H\right]_{\log}
 =-\frac6{\pi^2}\hbar g_H^2\log b,
 \qquad
 \left[\Delta Z_H^{\rm loop}\right]_{\log}
 =\frac3{8\pi^4}\hbar^2g_H^2\log b.
 \label{eq:v73-n4logs}
\end{equation}
The sign of the first line is the infrared step corresponding to the positive
ultraviolet scalar beta coefficient found in V71.  No analytic ultraviolet
completion of the scalar theory is inferred.

\section{A scale norm and the first quantitative self-map}
\label{sec:v73-norm}

We now add the estimate needed to turn the algebraic closure into a filtered
one-step class.  The only shell input used here is the native scalar bound already
available for the completed covariance: after rescaling to its physical shell,
for every prescribed derivative order $r$,
\begin{equation}
 |\partial_p^\beta d_j(p)|
 \le C_{\beta}\lambda_j^{-2-|\beta|},
 \qquad
 \operatorname{supp}d_j\subset
 \{c_0\lambda_j\le |p|\le c_1\lambda_{j+1}\},
 \label{eq:v73-shell-bound}
\end{equation}
with fixed $0<c_0<c_1<\infty$.  Smooth endpoint filters are included in the
constant.

For a rescaled vertex \eqref{eq:v73-rescaled-vertex}, let
$\|\widehat F_{m,\lambda}\|_{r}$ be the maximum of its momentum derivatives
through order $r$ on the fixed dimensionless window, together with the
corresponding rooted weighted kernel norm.  The exact choice of finite $r$ is
fixed once the source panel is fixed.  Define
\begin{equation}
 \|w_6\|_{\lambda}
 :=\sum_{\text{three local dim-6 jets}}|w_I|
 \label{eq:v73-z-norm}
\end{equation}
and
\begin{equation}
 \|R\|_{6,\lambda}
 :=\|R_6\|_{r,\lambda}
  +\|R_4\|_{r,\lambda}
  +\|R_2\|_{r,\lambda},
 \label{eq:v73-R-norm}
\end{equation}
where the Taylor zeros in \eqref{eq:v73-split6}--\eqref{eq:v73-split2} are part
of the definition of the three remainder spaces.

\begin{lemma}[Local growth versus complementary decay]
\label{lem:v73-excess}
Transport from scale $\lambda_{j+1}=b\lambda_j$ to the smaller external
window at $\lambda_j$ has two different weights:
\begin{equation}
 \boxed{
 \|w_{6,j}^{\rm carry}\|_{\lambda_j}
 \le b^{-2}\|w_{6,j+1}\|_{\lambda_{j+1}},
 \qquad
 \|R_j^{\rm carry}\|_{6,\lambda_j}
 \le b^{-4}\|R_{j+1}\|_{6,\lambda_{j+1}}.}
 \label{eq:v73-bminus2}
\end{equation}
Equivalently, reconstruction of a fixed physical dimension-six Wilson
coefficient toward the ultraviolet has inverse growth $b^2$, whereas the
complement after the dimension-six Taylor jets are stored has two additional
powers of decay.
\end{lemma}
\begin{proof}
For an $m$-field homogeneous local monomial with external momentum degree $r$,
the coefficient dimension is $4-m-r$.  The three stored dimension-six cases have
$m+r=6$, hence coefficient dimension $-2$.  Under
\eqref{eq:v73-rescaled-vertex} the corresponding dimensionless coefficient is
$\lambda^2$ times its physical coefficient.  Replacing $\lambda_{j+1}$ by
$\lambda_j=b^{-1}\lambda_{j+1}$ therefore multiplies it by $b^{-2}$.

After the dimension-six Taylor jet is removed, $R_6,R_4,R_2$ begin respectively
at momentum degrees $2,4,6$.  In all three cases the first omitted local
monomial has canonical dimension eight.  The integral Taylor formula therefore
supplies two additional powers of the ratio of external scale to shell scale.
Together with the dimension-six ancestry weight this gives the conservative
$b^{-4}$ carry bound in the fixed dimensionless remainder norm.  This is exactly
the growth--decay gap required by the filtered V6 theorem; it is not a claim that
the remainder has lost its dimension-six EFT ancestry.
\end{proof}

\begin{lemma}[Bounded contraction maps on the scalar filtered norm]
\label{lem:v73-contraction-bound}
Under \eqref{eq:v73-shell-bound}, at every fixed source/derivative order there is
$C_r<\infty$, independent of $j$, the periods and the ultraviolet endpoint, such
that the normal-product contractions obey
\begin{align}
 \|\mathscr C_{d_j}F_6\|_{4,j}
 &\le C_r\hbar\|F_6\|_{6,j+1},
 \label{eq:v73-C64}\\
 \|\mathscr C_{d_j}F_4\|_{2,j}
 &\le C_r\hbar\|F_4\|_{4,j+1},
 \label{eq:v73-C42}\\
 \|\mathscr C_{d_j}^2F_6\|_{2,j}
 &\le C_r\hbar^2\|F_6\|_{6,j+1}.
 \label{eq:v73-C62}
\end{align}
The fresh two-vertex shell rows in \eqref{eq:v73-Gc}, after the split
\eqref{eq:v73-split6}--\eqref{eq:v73-split2}, satisfy
\begin{equation}
 \|w^{\rm fresh}_{6,j}\|_{\lambda_j}
 +\|R^{\rm fresh}_{j}\|_{6,\lambda_j}
 \le C_r|g_H|^2.
 \label{eq:v73-fresh-bound}
\end{equation}
\end{lemma}
\begin{proof}
After $p=\lambda_j q$, \eqref{eq:v73-shell-bound} turns every shell covariance
and each prescribed derivative into a uniformly bounded dimensionless multiplier
on a fixed annulus.  A contraction closes two scalar legs and one loop momentum;
canonical dimensions cancel exactly, leaving the dimension-six normalization of
Lemma~\ref{lem:v73-excess}.  Young's inequality in the rooted weighted kernel norm
and the finite product rule for momentum derivatives give
\eqref{eq:v73-C64}--\eqref{eq:v73-C62}.  There are only finitely many flavour
contractions at fixed $n$ and source order.

For the fresh rows, the six-field kernel contains one shell covariance, the
four-field row two, and the two-field row three.  Their explicit coefficients in
\eqref{eq:v73-Gc} have precisely the powers needed to make the rescaled amplitudes
dimensionless.  Taylor extraction is bounded on the fixed dimensionless window,
and the complementary integral Taylor remainders have the same uniform symbol
bounds.  This proves \eqref{eq:v73-fresh-bound}.
\end{proof}

\begin{theorem}[First quantitative filtered scalar self-map]
\label{thm:v73-selfmap}
Fix $b>1$ and a finite source/derivative norm.  Under
\eqref{eq:v73-shell-bound}, the exact finite-order map
\eqref{eq:v73-Rj} obeys, for sufficiently small perturbative scalar coupling,
\begin{align}
 \|w_{6,j}\|_{\lambda_j}
 &\le b^{-2}\|w_{6,j+1}\|_{\lambda_{j+1}}
      +C_*b^{-4}\|R_{j+1}\|_{6,\lambda_{j+1}}
      +C_*|g_{H,j+1}|^2,
 \label{eq:v73-w6-recursion}\\
 \|R_j\|_{6,\lambda_j}
 &\le b^{-4}\|R_{j+1}\|_{6,\lambda_{j+1}}
      +C_*|g_{H,j+1}|^2.
 \label{eq:v73-stable-recursion}
\end{align}
The physical local rows $(m_H^2,Z_H,g_H)$ receive only finite Taylor
projections of the same recombined functional; at the present perturbative order
their forcing is a polynomial in lower-order local data plus the displayed
complement-to-local contractions.  In particular, the shell transformation maps
the finite coefficient space \eqref{eq:v73-state} to itself with constants
uniform in the periods and shell index on the declared profile class.

The local dimension-six block has inverse ultraviolet growth exponent $s=2$,
whereas the complete remainder has decay exponent $t=4$.  These exponents are
uniform at this fixed perturbative/EFT/source order.  No claim is made of a
cutoff-independent small ball for arbitrary physical dimension-six Wilson
coefficients; those coefficients are two-boundary local data.
\end{theorem}
\begin{proof}
Recombine the incoming state before applying any estimate.  Lemma
\ref{lem:v73-excess} gives $b^{-2}$ on the carried local dimension-six block and
$b^{-4}$ on the carried complement.  Lemma~\ref{lem:v73-contraction-bound}
bounds every additional contraction.  A contraction of a carried complement may
feed a stored local jet; this is precisely the allowed complement-to-local block
and gives the second term of \eqref{eq:v73-w6-recursion}.  A stored local
polynomial cannot generate a same-order nonlocal remainder under a local Gaussian
contraction, so there is no same-degree local-to-remainder block.  The fresh shell
functional has bounded local and complementary pieces by
\eqref{eq:v73-fresh-bound}.  This proves
\eqref{eq:v73-w6-recursion}--\eqref{eq:v73-stable-recursion}.

The physical mass, kinetic and quartic rows are bounded finite Taylor projections
of the same recombined functional.  Their dependence on the incoming local
six-dimensional block is the expected EFT power mixing and is retained in the
local coefficient map rather than estimated as a remainder.  Thus the map closes
on the stated finite coefficient space without assuming a running coupling.
\end{proof}

The inequalities \eqref{eq:v73-w6-recursion}--\eqref{eq:v73-stable-recursion}
are the first native interacting realization in this lineage of V6's abstract
``local growth versus complementary decay'' split.  The local dimension-six
coefficients remain owned by the EFT boundary data; the six-field row is not an
error term.

\section{The scalar two-boundary consequence}
\label{sec:v73-two-boundary}

The preceding theorem can now be inserted into the filtered two-boundary theorem
of V6 without supplying the scalar stable trajectory externally.
At fixed perturbative degree two, the renormalizable local block has inverse
exponent zero, while the largest inverse growth comes from the dimension-six
local Wilson coordinates in \eqref{eq:v73-w6-recursion}.  Thus
$s=2$.  Equation~\eqref{eq:v73-stable-recursion} gives complementary exponent
$t=4$.  One may therefore choose, for example, the V6 weight $\sigma=3$:
\begin{equation}
 \boxed{s=2\le\sigma=3<t=4.}
 \label{eq:v73-window}
\end{equation}
The finite source panel has identity coordinate transport in the present
unrescaled-field convention, and terminal Gaussian evaluation is the inherited
finite Wick evaluation.

\begin{corollary}[Filtered two-boundary trajectory for the scalar spine]
\label{cor:v73-two-boundary}
Fix the truncation $(g_H^2,\hbar^2,\Delta=6)$, a finite connected scalar source
panel with at least two retained scalar legs, and the shell class of
Theorem~\ref{thm:v73-selfmap}.  Prescribe the physical local scalar
renormalization data at the lower endpoint and polynomially bounded ultraviolet
stable boundary data at a finite endpoint $J$.  Then the scalar spine has a unique
finite two-boundary coefficient trajectory.  As $J\to\infty$, the stable boundary
influence at fixed $j$ is bounded by
\begin{equation}
 \boxed{
 C(1+J)^q\theta^{J-j},
 \qquad b^{-1}<\theta<1,}
 \label{eq:v73-boundary-loss}
\end{equation}
and the connected scalar source coefficients in the declared finite panel converge
coefficientwise at this perturbative/EFT order.

The local quartic logarithm and the two-loop kinetic logarithm along this
coefficient trajectory are the internally generated values
\eqref{eq:v73-dg}--\eqref{eq:v73-dZ}.  This corollary is a coefficientwise
finite-order continuum statement.  It does not assert an analytic ultraviolet
trajectory for positive scalar $g_H$, whose one-loop beta coefficient has the
Landau sign.
\end{corollary}
\begin{proof}
Theorem~\ref{thm:v73-algebraic-selfmap} gives the finite triangular state,
Theorem~\ref{thm:v73-selfmap} gives its local-growth/complementary-decay
bounds and the strict gap \eqref{eq:v73-window}, and the terminal Gaussian/source
conditions are inherited from the measured scalar shell construction.  These are
the scalar-spine specialization of V6's weighted R1--R4 hypotheses.  Apply
Theorem~\ref{thm:v6-filtered-trajectory}.  Its ultraviolet boundary-loss estimate
gives \eqref{eq:v73-boundary-loss}.  The local logarithms are the already evaluated
projections of the same shell map, so no supplied marginal trajectory enters.
\end{proof}

\begin{remark}[What has and has not been instantiated from V6]
\label{rem:v73-v6-scope}
The corollary verifies the filtered theorem only on the scalar connected spine just
defined.  It does not verify the combined diffeomorphism--Lorentz--internal-gauge
BV complex, chiral determinant line, mixed gravity--matter shell bounds, or a full
ESM source bank.  Those sectors cannot be imported merely because the scalar
state has the same abstract triangular shape.
\end{remark}

\section{Anti-circularity audit and the next ESM interface}
\label{sec:v73-audit}

The construction closes a specific loop in the earlier programme.
The quartic marginal update is extracted only after the six-field return has been
retained.  The two-point row is extracted only after both the six- and four-field
returns have been retained.  The stable estimate is proved on the resulting
state, and only then is the V6 two-boundary theorem invoked.  Thus none of the
following is assumed in order to prove itself:
\begin{enumerate}[label=\textup{(A\arabic*)},leftmargin=3em]
\item no running quartic coupling is supplied to derive
      \eqref{eq:v73-dg};
\item no coupling-only truncation is used to prove shell composition;
\item no ultraviolet stable boundary is set to zero before its loss is proved;
\item no continuum Ward identity is used in the scalar contraction algebra;
\item no all-ESM filtered self-map is inferred from the scalar one.
\end{enumerate}

The next load-bearing task is no longer another pure-scalar coefficient.  The
natural merger is the first \emph{mixed gravity--Higgs filtered shell}: combine
V63--V66's populated metric/scalar action blocks and common one-loop insertion
with the state architecture above, and test whether its local dimension-six bank
and quasilocal complement satisfy an analogue of
\eqref{eq:v73-stable-recursion}.  That is the point at which the free gravitational
reference of V6 would first consume the interacting scalar self-map proved here.

\section{Status ledger for V73}
\label{sec:v73-ledgers}

\subsection*{Proved}
A finite scalar perturbative/EFT state has been defined at
$(g_H^2,\hbar^2,\Delta=6)$.  The six-, four- and two-field rows obey an exact
Gaussian semigroup, and both higher-field rows are proved necessary for shell
composition.  The local/complement decomposition is exact and retains the first
three scalar dimension-six EFT coordinates as local Wilson data.  Under the native
scalar shell symbol bounds, those coefficients have inverse ultraviolet growth
$b^2$ while the complete complement has decay $b^{-4}$, giving the strict
filtered window $2<4$.  The V6 filtered two-boundary theorem is thereby
instantiated on this scalar spine,
with coefficientwise ultraviolet boundary loss and internally generated quartic
and kinetic logarithms.

\subsection*{Inherited}
V70 supplies the common quartic/mass/vacuum normal-product packet and the exact
one-vertex source transport.  V71 supplies the complete two-vertex cumulant and
the nonzero six-field return, including its first mixed-shell contraction.  V72
supplies the overlapping sunset forest, kinetic logarithm and all-relative-scale
two-point estimate.  V6 supplies the already proved abstract weighted
R1--R4/two-boundary theorem; V73 verifies its scalar-spine hypotheses rather than
reproving that theorem.

\subsection*{Literature input}
No new literature theorem is used in the shell-composition proof.  Gaussian
conditional cumulance and Taylor/BPHZ organization are standard background, but
the contraction coefficients, finite state and scale recursion are derived from
the populated native scalar rows.  The external scalar anomalous-dimension and
beta-function literature is not used to set \eqref{eq:v73-dg} or
\eqref{eq:v73-dZ}.

\subsection*{Conditional}
The quantitative self-map uses a fixed finite scalar source/derivative panel and
the native annular shell symbol estimate \eqref{eq:v73-shell-bound}.  The physical
mass and finite kinetic/quartic anchors remain lower-endpoint normalization data.
Pure background-source vacuum coefficients at loop degree three are outside the
truncation.  No estimate is uniform in increasing perturbative order, EFT degree,
source order or an unbounded number of interacting ESM sectors.

\subsection*{Open}
The smallest load-bearing ESM theorem is now a mixed gravity--Higgs one-shell
self-map on a common filtered state, including the actual gravitational ghost and
measure restrictions and the scalar local/complement coordinates above.  Beyond
that remain internal-gauge/Yukawa/chiral merger, combined BV closure, coherent
line transport, arbitrary-order filtered induction and the full source-complete
ESM continuum theorem.

\begin{center}\footnotesize
\begin{tabular}{p{.23\textwidth}p{.69\textwidth}}
\toprule
Variable & Scope of V73\\\midrule
Perturbative order & Coupling degree $g_H^2$, loop degree $\hbar^2$.\\
EFT degree & Local scalar bank through canonical dimension six; higher Taylor jets remain in the quasilocal complement.\\
Sources & Fixed finite connected scalar source panel with at least two retained scalar legs; nonzero-field frame/metric derivatives generated from the same coefficients.\\
Cutoff & One-shell constants uniform in shell index and periods on the declared native profile class; scalar two-boundary ultraviolet boundary loss follows at this fixed order.\\
RG steps & Coefficientwise two-boundary trajectory at the stated finite order; no analytic scalar UV completion.\\
Full ESM & Not proved. Mixed gravity--Higgs, gauge/Yukawa/chiral and combined BV/line sectors remain open.\\
\bottomrule
\end{tabular}
\end{center}

\section*{Append-only continuation marker: V73}
\addcontentsline{toc}{section}{Append-only continuation marker: V73}
\label{sec:v73-continuation-marker}
\textbf{End of V73.} Preserve Sections 1--563.  The first interacting filtered
scalar state is now closed through $g_H^2\hbar^2$: its six-, four- and two-field
rows compose exactly, its dimension-six local Wilson bank and quasilocal
complement are retained, and the native shell realizes the strict filtered gap
between local inverse growth $b^2$ and complementary decay $b^{-4}$.  The next load-bearing step is the mixed gravity--Higgs
filtered shell, not another isolated scalar graph.



% =============================================================================
% V74 APPEND-ONLY CONTINUATION
% =============================================================================
\part*{V74: The first mixed gravity--Higgs filtered row from one joint hard shell}
\addcontentsline{toc}{part}{V74: The first mixed gravity--Higgs filtered row from one joint hard shell}
\label{part:v74}

\section{Why the first mixed row must contain a hard scalar line}
\label{sec:v74-scope}

V73 closed the interacting scalar spine through
$(g_H^2,\hbar^2,\Delta=6)$.  The next interface is not obtained by inserting
that soft scalar action into the soft current exchange of V6.  The propagating
gravitational covariance used there is a genuine hard shell.  If all scalar
legs entering one linear stress vertex are retained in the soft window, the
momentum carried by the hard graviton is a sum of soft external momenta and the
term vanishes by support.  A nonzero mixed scalar two-point coefficient first
appears when the same shell contains a hard scalar line.

We therefore keep the exact V6 reduced gravitational reference and the V27
completed scalar reference simultaneously.  Let $d^g_j$ be one positive hard
TT graviton innovation between the scales $\lambda_{j+1}=b\lambda_j$ and
$\lambda_j$, and let $d^H_j$ be the corresponding scalar innovation.  Their
supports are contained in fixed annuli
\begin{equation}
 c_0\lambda_j\le |p|\le c_1\lambda_{j+1},
 \label{eq:v74-shell-support}
\end{equation}
with the inherited smooth endpoint completions.  The scalar covariance is
flavour diagonal.  The graviton covariance is the signed/physical TT covariance
obtained after V6's exact linear constraint--ghost reduction; no positive
replacement of that tensor is made.

Expand the native scalar Hessian in the normalized hard gravitational variable
$q$ as
\begin{equation}
 Q_a(q)=P_a+q^I Q_{a,I}+\frac12q^Iq^JQ_{a,IJ}+O(q^3).
 \label{eq:v74-Q-expansion}
\end{equation}
The multi-index $I$ denotes the actual propagating gravitational component after
the linear reduction.  The factors of the gravitational expansion parameter are
retained in the vertices.  When useful we write
\begin{equation}
 Q_{a,I}=\kappa_g\,\overline Q_{a,I},\qquad
 Q_{a,IJ}=\kappa_g^2\,\overline Q_{a,IJ}.
 \label{eq:v74-kappa-tag}
\end{equation}
This tag is part of the EFT ancestry.  In particular, a mass or kinetic
counterterm generated by the graph below is not reclassified as a primitive
renormalizable scalar interaction merely because Gaussian contraction has
lowered its canonical field degree.

The present V74 component has zero internal-gauge and Yukawa couplings.  The
quartic scalar interaction remains in the common V73 state, but it does not
enter the mixed one-loop row computed below.  At the displayed order the hard
mixed graph is therefore a statement about the common Gaussian gravity--scalar
reference with arbitrary fixed scalar source/frame marks inherited from V73.

\begin{proposition}[Hard-support bridge exclusion]
\label{prop:v74-hard-bridge}
Let both scalar legs of a linear $q\phi\phi$ vertex lie in the retained source
window $|k|\le\mu$, with $N\mu<c_0\lambda_j$ for the prescribed finite source
arity $N$.  Then a single hard graviton line cannot close two such vertices
without at least one hard scalar or other hard field line.  Equivalently, the
purely soft tree exchange of the hard TT shell vanishes on this panel.
\end{proposition}
\begin{proof}
At a linear scalar--graviton vertex, momentum conservation identifies the
momentum on the graviton line with a signed sum of the scalar and source
momenta entering that vertex.  With at most $N$ retained marks its magnitude is
strictly below $c_0\lambda_j$.  The covariance $d_j^g$ is zero there by
\eqref{eq:v74-shell-support}.  A hard scalar line changes the routed momentum by
an annular momentum and removes this support obstruction.  The statement uses
the actual support, not an assertion that a connected bridge is ``small.''
\end{proof}

\section{Exact finite mixed scalar self-energy}
\label{sec:v74-selfenergy}

Split one scalar flavour as $\phi=\varphi+\eta$, with $\varphi$ retained and
$\eta$ distributed with covariance $\hbar d_j^H$.  The two interaction pieces
that contribute to the retained scalar quadratic action at order
$\kappa_g^2\hbar$ are
\begin{equation}
 V_{11}=\langle\varphi,q^IQ_{I}\eta\rangle,
 \qquad
 V_{20}=\frac14\langle\varphi,q^Iq^JQ_{IJ}\varphi\rangle.
 \label{eq:v74-V11-V20}
\end{equation}
All pairings are the finite complex-bilinear continuations of the real field
pairing.  The factors in \eqref{eq:v74-V11-V20} follow directly from
$\frac12\langle\phi,Q(q)\phi\rangle$ and
\eqref{eq:v74-Q-expansion}.

For a finite operator word $A(p,k)$ let $\operatorname{Ctr}_p A$ denote its normalized finite
sum over the routed shell momentum, including the actual endpoint filters and
matrix contractions.  Define the retained-scalar kernel
\begin{equation}
 \boxed{\begin{aligned}
 \Sigma_{a,j}^{gH}(k)
 &= \frac{\hbar}{2}\sum_{I,J}
 \operatorname{Ctr}_p\!\bigl[d^g_{j,IJ}(p)\,Q_{a,IJ}(p,-p;k,-k)\bigr]\\
 &\quad-\hbar\sum_{I,J}\operatorname{Ctr}_p\!\bigl[
 d^g_{j,IJ}(p)Q_{a,I}(k,p)d^H_j(p+k)
 Q_{a,J}(p+k,-k)\bigr].
 \end{aligned}}
 \label{eq:v74-selfenergy}
\end{equation}
The tensor order in the second term is the displayed order; no commutation of
native vertex factors is assumed.

\begin{theorem}[Exact one-shell mixed two-point formula]
\label{thm:v74-exact-selfenergy}
Through order $\kappa_g^2\hbar$ and quadratic order in the retained scalar,
the simultaneous finite Gaussian elimination of $(q,\eta)$ contributes
\begin{equation}
 \boxed{
 \Delta\mathscr S^{gH,(2)}_{a,j}[\varphi]
 =\frac12\langle\varphi,\Sigma_{a,j}^{gH}\varphi\rangle.}
 \label{eq:v74-effective-quadratic}
\end{equation}
The first term of \eqref{eq:v74-selfenergy} is the hard-graviton contact and
the second is the graviton--scalar loop.  Both are restrictions of the same
native scalar Hessian \eqref{eq:v74-Q-expansion}.

The exact linear gravitational ghost determinant of V6 contributes no additional
scalar two-point vertex to \eqref{eq:v74-effective-quadratic}: at this reference
order it cancels the normal constraint Jacobian before the TT covariance is
formed.  Nonlinear ghost, connection, lapse/shift and measure vertices are not
included in the theorem and remain separate rows of the full ESM problem.
\end{theorem}
\begin{proof}
The first cumulant of $V_{20}$ gives
\[
 \mathbb E V_{20}
 =\frac{\hbar}{4}\left\langle\varphi,
    \sum_{I,J}\operatorname{Ctr}_p[d^g_{j,IJ}Q_{IJ}]\varphi\right\rangle.
\]
Writing the resulting quadratic action as
$\frac12\langle\varphi,\Sigma\varphi\rangle$ yields the coefficient
$\hbar/2$ in the first term of \eqref{eq:v74-selfenergy}.

The only second-cumulant contribution with two retained scalar legs is
$V_{11}^2$.  Independence of $q$ and $\eta$ gives
\[
 \mathbb E(V_{11}^2)
 =\hbar^2\left\langle\varphi,
   \sum_{I,J}\operatorname{Ctr}_p[d^g_{j,IJ}Q_I d^H_jQ_J]\varphi\right\rangle.
\]
The logarithm of the normalized Gaussian integral contributes
$-(2\hbar)^{-1}\mathbb E(V_{11}^2)$, which becomes the second term of
\eqref{eq:v74-selfenergy} after again extracting the outer factor $1/2$.
All other terms have either the wrong retained-scalar arity or higher powers of
$\kappa_g$ or $\hbar$.  V6's constraint--ghost theorem has already been used in
defining $d_j^g$, so reinserting that determinant would double count the linear
reduction.
\end{proof}

\begin{remark}[Why this is not obtained from V63 by opening a line]
The contact and three-propagator vacuum graphs of V61--V64 have their own vacuum
symmetry factors.  Closing or opening an external scalar line changes those
factors.  Equation~\eqref{eq:v74-selfenergy} is therefore derived directly from
the finite cumulant rather than inferred by cutting a vacuum diagram.  The
common native vertices are the same, but the combinatorics is not replaced by a
line-opening mnemonic.
\end{remark}

\section{The ancestry-preserving local split}
\label{sec:v74-split}

The self-energy has canonical mass and kinetic descendants as well as a genuine
dimension-six scalar operator.  The descendants must remain tagged by the
$\kappa_g^2$ ancestry in \eqref{eq:v74-kappa-tag}.  This is the mixed analogue
of V6's $\phi^6\mapsto\phi^4\mapsto\phi^2$ warning.

The inherited finite symbols imply, after rescaling the joint shell to unit
annuli, that for every prescribed $r\le6$,
\begin{equation}
 \boxed{
 \left\|\partial_k^\alpha
   \bigl(\hbar^{-1}\kappa_g^{-2}\Sigma_{a,j}^{gH}\bigr)(k)
 \right\|
 \le C_{\alpha}\lambda_j^{4-|\alpha|},
 \qquad |\alpha|\le6,
 \quad |k|\le N\mu.}
 \label{eq:v74-symbol-bound}
\end{equation}
The bound is deliberately conservative.  Individual massless continuum terms
can vanish or cancel, but no such cancellation is assumed in the filtered norm.
The native reality and momentum-reversal symmetries make the quadratic action
even in the external momentum after symmetrization.

Let $\mathbf T_{4,k}^{\rm even}$ denote the joint even Taylor polynomial through
momentum degree four.  Put
\begin{equation}
 L^{gH}_{a,j}:=\mathbf T_{4,k}^{\rm even}\Sigma_{a,j}^{gH},
 \qquad
 R^{gH}_{a,j}:=(1-\mathbf T_{4,k}^{\rm even})\Sigma_{a,j}^{gH}.
 \label{eq:v74-local-remainder}
\end{equation}
The three local Taylor levels have the schematic interpretation
\begin{equation}
 \boxed{
 k^0:\ \kappa_g^2\phi^2,
 \qquad
 k^2:\ \kappa_g^2(\partial\phi)^2,
 \qquad
 k^4:\ \kappa_g^2\phi\,\partial^4\phi.}
 \label{eq:v74-descendants}
\end{equation}
The first two are lower-canonical descendants but remain in the same gravity
EFT ancestry row.  The last is a genuine dimension-six local Wilson operator.
At the physical lower endpoint the descendants may mix with the ordinary scalar
mass and kinetic normalization; the state keeps the ancestry label until that
physical matching is imposed.

\begin{proposition}[Six-derivative remainder and rooted locality]
\label{prop:v74-remainder}
On the shell domain \eqref{eq:v74-shell-support}, for $|\alpha|\le6$,
\begin{equation}
 \boxed{
 \|\partial_k^\alpha R^{gH}_{a,j}(k)\|
 \le C_{\alpha}\hbar\kappa_g^2
      \lambda_j^{-2}|k|^{6-|\alpha|}.}
 \label{eq:v74-remainder-bound}
\end{equation}
For every fixed rooted source order and polynomial tree weight there is a
corresponding position-space bound with no total-volume multiplier.  The same
statement holds after finitely many scalar frame/source derivatives, with the
constant enlarged by their fixed product-rule budget.
\end{proposition}
\begin{proof}
Equation~\eqref{eq:v74-symbol-bound} with $|\alpha|=6$ and the integral Taylor
formula give the momentum estimate.  The odd fifth Taylor tensor vanishes in
the symmetrized scalar quadratic action.  For rooted kernels, differentiate the
finite shell symbols before Fourier inversion.  Each momentum derivative gains
a factor $\lambda_j^{-1}$, while the smooth annular support removes boundary
terms.  The V6 gravitational shell estimate and the V27 scalar shell estimate
therefore give arbitrary fixed polynomial moments.  Products of the two kernels
are handled by the same rooted-tree inequality used in V6 and V73.  The argument
never estimates an extensive action in supremum norm.
\end{proof}

\section{A first mixed filtered state and its scale gap}
\label{sec:v74-filtered-state}

Let $\mathfrak S_j^H$ be the V73 scalar state.  Enlarge it by one ancestry-tagged
mixed two-point row
\begin{equation}
 \boxed{
 \mathfrak S_j^{H\oplus gH,\mathrm{ref}}
 =\bigl(\mathfrak S_j^H;\ u^{gH}_{j,<6},\ w^{gH}_{6,j},
                 R^{gH}_{j};\ \mathcal J_{g,j}\bigr).}
 \label{eq:v74-state}
\end{equation}
Here $u^{gH}_{j,<6}$ stores the mass and kinetic descendants in
\eqref{eq:v74-descendants} with their $\kappa_g^2$ ancestry tag,
$w^{gH}_{6,j}$ stores the degree-four Taylor tensor, and $R^{gH}_{j}$ is the
complete complement.  $\mathcal J_{g,j}$ denotes the already declared finite
gravitational source normalization on this reference panel.  The zero-retained-
scalar mixed vacuum/source rows of V61--V66 may be adjoined as independent
already-populated local coordinates; the theorem below does not re-estimate their
forests.

Use the same scale convention $\lambda_{j+1}=b\lambda_j$.  In the dimensionless
mixed remainder norm, transport of a fixed physical dimension-six Wilson tensor
from $j+1$ to $j$ gains $b^{-2}$, while the complement after the degree-four
Taylor jet has first omitted canonical dimension eight and gains $b^{-4}$.
Thus
\begin{equation}
 \boxed{
 \|w^{gH,\mathrm{carry}}_{6,j}\|
 \le b^{-2}\|w^{gH}_{6,j+1}\|,
 \qquad
 \|R^{gH,\mathrm{carry}}_j\|
 \le b^{-4}\|R^{gH}_{j+1}\|.}
 \label{eq:v74-carry}
\end{equation}
The lower descendants in $u^{gH}_{<6}$ are local two-boundary coordinates and
are not placed in the stable complement.

The fresh shell has the dimensionless bound
\begin{equation}
 \boxed{
 \|u^{gH,\mathrm{fresh}}_{j,<6}\|
 +\|w^{gH,\mathrm{fresh}}_{6,j}\|
 +\|R^{gH,\mathrm{fresh}}_j\|
 \le C_*\hbar(\kappa_g\lambda_j)^2\,\mathfrak z_{H,j}^2,}
 \label{eq:v74-fresh}
\end{equation}
where $\mathfrak z_{H,j}$ is a fixed finite stock of the scalar quadratic
normalization and source vertices.  On the V73 coefficient trajectory this stock
has at most polynomial growth at the present perturbative order.

\begin{theorem}[First quantitative mixed gravity--Higgs reference self-map]
\label{thm:v74-selfmap}
Fix the combined truncation consisting of the V73 scalar spine through
$(g_H^2,\hbar^2,\Delta=6)$ and the mixed reference row
$(\kappa_g^2\hbar,\Delta=6)$ of
\eqref{eq:v74-selfenergy}.  For sufficiently small finite perturbative
coefficients at the physical endpoint, the one-shell map sends
\eqref{eq:v74-state} to a state of the same form and obeys
\begin{align}
 \|w^{gH}_{6,j}\|
 &\le b^{-2}\|w^{gH}_{6,j+1}\|
   +C_*b^{-4}\|R^{gH}_{j+1}\|
   +C_*\hbar(\kappa_g\lambda_j)^2\mathfrak z_{H,j+1}^2,
 \label{eq:v74-w-rec}\\
 \|R^{gH}_{j}\|
 &\le b^{-4}\|R^{gH}_{j+1}\|
   +C_*\hbar(\kappa_g\lambda_j)^2\mathfrak z_{H,j+1}^2.
 \label{eq:v74-R-rec}
\end{align}
The mass and kinetic Taylor descendants are updated by finite local projections
of the same $\Sigma^{gH}_{a,j}$ and are retained in $u^{gH}_{<6}$.
No same-degree local-to-remainder block occurs.

The mixed local Wilson block has inverse ultraviolet growth exponent $s=2$ and
the complement has carry exponent $t=4$.  The fresh forcing has exponential
weight $\gamma=2$ coming from $(\kappa_g\lambda_j)^2$; hence the V6 normalized
forcing condition is satisfied by the same choice
\begin{equation}
 \boxed{s=2\le\sigma=3<t=4,\qquad \gamma=2<\sigma.}
 \label{eq:v74-window}
\end{equation}
All constants are uniform in the periods and shell index on the declared native
profile class at fixed source and derivative order.
\end{theorem}
\begin{proof}
The exact finite formula is Theorem~\ref{thm:v74-exact-selfenergy}.  Split it by
\eqref{eq:v74-local-remainder}.  Carrying an incoming state gives
\eqref{eq:v74-carry}.  A carried complement may contribute to a stored local
Taylor tensor after restriction to the smaller source window, producing the
second term of \eqref{eq:v74-w-rec}; this is the allowed complement-to-local
block.  A local polynomial remains local under restriction and cannot create a
same-order nonlocal complement.

For a fresh shell, rescale the two internal momenta and native symbols to the
fixed annuli.  The two gravitational vertices contribute $\kappa_g^2$ and the
four-dimensional shell counting leaves the dimensionless factor
$(\kappa_g\lambda_j)^2$.  Proposition~\ref{prop:v74-remainder} controls the
complement, while finite Taylor projection controls the local rows.  This gives
\eqref{eq:v74-fresh} and hence
\eqref{eq:v74-w-rec}--\eqref{eq:v74-R-rec}.

The V73 scalar rows already satisfy their own $s=2<t=4$ estimates.  In the mixed
row the local degree-six coefficient again has inverse growth $b^2$ and the first
omitted dimension-eight term has $b^{-4}$ carry.  Multiplying the fresh forcing
by the V6 weight $b^{-3j}$ leaves $b^{-j}$ times the polynomial scalar stock.
This is exactly the normalized-forcing condition of V6.  No ultraviolet
trajectory for the dimensionless Newton coupling is assumed or inferred.
\end{proof}

\begin{remark}[Why the ancestry tag matters]
If the $k^0$ or $k^2$ Taylor descendants in
\eqref{eq:v74-descendants} were merged immediately with the primitive scalar
mass and kinetic rows and their $\kappa_g^2$ ancestry discarded, the shell
forcing $\hbar\kappa_g^2\lambda_j^4$ or
$\hbar\kappa_g^2\lambda_j^2$ would appear as an unexplained lower-dimensional
UV growth.  The filtered bank instead records that these are contraction
descendants of one dimension-six gravitational EFT row.  Physical matching at
the lower endpoint may combine their numerical values, but the iterative state
must retain the origin needed for power counting.
\end{remark}

\section{Two-boundary consequence on the mixed reference spine}
\label{sec:v74-two-boundary}

The direct sum of the V73 scalar system and the new mixed row has the same maximal
local inverse exponent $s=2$ and the same minimal complementary carry exponent
$t=4$.  The mixed forcing satisfies the grade inequality
\eqref{eq:v74-window}.  The terminal finite Wick evaluation and scalar source
transport are inherited, while the linear gravitational source map is the V6
reference map.

\begin{corollary}[Coefficientwise mixed gravity--Higgs reference trajectory]
\label{cor:v74-two-boundary}
Fix the finite truncation
\begin{equation}
 (g_H^2,\hbar^2,\Delta=6)
 \ \oplus\
 (\kappa_g^2\hbar,\Delta=6)
 \label{eq:v74-truncation}
\end{equation}
and a finite scalar/gravitational source panel satisfying the hard-support
condition above.  Prescribe the physical scalar and mixed local renormalization
data at the lower endpoint and polynomially bounded ultraviolet complementary
data at a finite endpoint $J$.  Then the reference mixed spine has a unique
finite two-boundary coefficient trajectory.  Its ultraviolet complementary
boundary influence at every fixed $j$ obeys
\begin{equation}
 \boxed{
 \|x_j^{(J)}-x_j^{(\infty)}\|
 \le C b^{3j}(1+J)^q\theta^{J-j},
 \qquad b^{-1}<\theta<1,}
 \label{eq:v74-boundary-loss}
\end{equation}
with the corresponding fixed-endpoint connected source estimate.

The statement is coefficientwise in the gravitational EFT expansion.  It does
not assert that $\kappa_g\lambda_j$ remains small on an infinite analytic UV
trajectory; its exponential growth is precisely what the weighted two-boundary
normalization in \eqref{eq:v74-window} absorbs at this finite EFT order.
\end{corollary}
\begin{proof}
Combine Theorem~\ref{thm:v73-selfmap} with
Theorem~\ref{thm:v74-selfmap}.  Their direct-sum local/complement map is
triangular at the declared perturbative degrees and satisfies V6's weighted
R1--R4 bounds with $\sigma=3$.  The normalized mixed forcing is exponentially
decaying and the remaining lower-order coefficients grow only polynomially.
Theorem~\ref{thm:v6-filtered-trajectory} therefore gives the unique trajectory
and \eqref{eq:v74-boundary-loss}.  The argument is coefficientwise and makes no
asymptotic-safety assumption.
\end{proof}

\section{Interface with the earlier scalar--metric loops}
\label{sec:v74-interface}

V74 fills the first retained-scalar row that was missing from the earlier
zero-field scalar--metric calculations.  The relations are as follows.

\begin{enumerate}[label=\textup{(I\arabic*)},leftmargin=3em]
\item V61--V64 evaluate and renormalize the zero-retained-scalar contact and
      three-propagator scalar--metric vacuum/source cores.  Their common-action
      counterterm insertion remains the correct insertion when the scalar
      two-point row is also retained.
\item V66 evaluates the universal logarithmic scalar determinant on a geometric
      background and therefore populates curvature/mass local coefficients in
      the same finite ESM bank.  V74 does not replace those geometric rows by a
      scalar self-energy.
\item The new two-point row uses the same $Q_I,Q_{IJ}$ vertices but has different
      graph symmetry factors.  Those factors are fixed by
      Theorem~\ref{thm:v74-exact-selfenergy}, not by opening a V61 vacuum graph.
\item At this reference order the V6 linear constraint--ghost cancellation has
      already been consumed in $d_j^g$.  The nonlinear gravitational ghost,
      coframe/connection measure and combined BV restoration are still absent.
\end{enumerate}

Thus the combined populated bank now contains a scalar interacting spine, a
scalar determinant geometric block, zero-field scalar--metric two-loop cores,
and a retained-scalar mixed one-loop row.  What is still missing is not another
copy of these Gaussian contractions but the nonlinear gravitational shell that
makes all of them components of one symmetry-complete map.

\begin{proposition}[The TT reference does not close the full gravity--Higgs EFT]
\label{prop:v74-tt-not-full}
The self-map of Theorem~\ref{thm:v74-selfmap} cannot be promoted to the full
mixed ESM shell by notation alone.  In particular, lapse/shift or conformal
constraint exchange couples to scalar potential stress that is annihilated by
the pure TT projector.  At order $\kappa_g^2g_H^2$ such nonpropagating exchange
can generate local or quasilocal eight-scalar descendants before scalar
contraction.  Those directions are not contained in the V74 TT reference row
and require the actual nonlinear constraint/measure completion before their
coefficients can be assigned.
\end{proposition}
\begin{proof}
The quartic scalar potential varies with the volume density and therefore has a
nonzero pure-trace geometric stress.  The TT projector annihilates a pure trace,
so this direction is absent from the reference exchange used above.  A
nonpropagating scalar gravitational response need not annihilate it.  Pairing two
quartic potential stresses gives a term of scalar arity eight and overall
$\kappa_g^2g_H^2$ ancestry.  Whether and how it is cancelled by the native
constraint, measure and ghost sectors is a question about the full finite shell,
not the TT reference.  Hence it cannot be deleted or assigned a coefficient from
V74 alone.
\end{proof}

This proposition identifies the next genuinely upstream interface.  Before a
full mixed self-map is claimed, the nonlinear gravitational constraint and ghost
completion must determine whether the arity-eight ancestor survives, is a local
coboundary, or cancels in the common measured action.

\section{Anti-circularity audit and status ledger for V74}
\label{sec:v74-ledgers}

\subsection*{Proved}
The first retained-scalar mixed gravity--Higgs shell coefficient has been derived
directly from the native finite Gaussian cumulant.  Its hard-graviton contact and
graviton--scalar loop have the exact relative factors in
\eqref{eq:v74-selfenergy}.  Hard support proves that the corresponding purely
soft TT tree bridge vanishes.  After the even Taylor jet through momentum degree
four is stored, the mixed two-point complement starts at degree six and obeys the
rooted quasilocal estimate \eqref{eq:v74-remainder-bound}.  The resulting
ancestry-preserving mixed state has local inverse growth $b^2$, complementary
carry $b^{-4}$ and fresh gravitational forcing of weight $b^{2j}$, so the V6
window $2\le3<4$ closes on the direct sum with the V73 scalar spine.  A
coefficientwise mixed two-boundary trajectory follows at the finite truncation
\eqref{eq:v74-truncation}.

\subsection*{Inherited}
V6 supplies the exact linear constraint--ghost reduction, the TT hard covariance,
its quasilocal shell bound and the weighted two-boundary theorem.  V27 supplies
the completed scalar covariance and the native metric derivatives
$Q_I,Q_{IJ}$.  V61--V64 supply the zero-retained-scalar mixed contact/sunset
bookkeeping and common counterterm insertion.  V66 supplies the scalar-determinant
geometric logarithmic block.  V73 supplies the interacting scalar filtered state
and its internally generated scalar local trajectory.

\subsection*{Literature input}
No new external coefficient is used.  Gaussian cumulance, Taylor subtraction and
EFT ancestry are standard background, but the factors, support statement and
filtered recursion are derived from the native finite vertices and covariances.
No continuum graviton--scalar anomalous dimension is inserted as an assumption.

\subsection*{Conditional}
The quantitative mixed self-map is restricted to the V6 flat reduced TT reference,
one hard scalar shell, fixed finite source/derivative order and the declared
coefficient truncation.  The constants use the inherited native symbol and
annular-support bounds.  The already-populated V61--V66 zero-field rows are
compatible local coordinates but their complete filtered shell recursion is not
reproved here.  The physical finite mass and kinetic matching conditions remain
lower-endpoint data.

\subsection*{Open}
The smallest load-bearing theorem is now the nonlinear gravitational
constraint/ghost/measure completion of this mixed shell.  It must decide the
constraint-mediated scalar-potential directions identified in
Proposition~\ref{prop:v74-tt-not-full}, construct the required local ancestor
envelope (potentially including arity eight at
$\kappa_g^2g_H^2$), and prove a common filtered estimate together with the
combined diffeomorphism--Lorentz BV restoration.  Only after that step can the
V74 reference spine be promoted to a full mixed gravity--Higgs one-shell
self-map.

\begin{center}\footnotesize
\begin{tabular}{p{.23\textwidth}p{.69\textwidth}}
\toprule
Variable & Scope of V74\\\midrule
Perturbative order & V73 scalar spine through $g_H^2\hbar^2$ plus the mixed reference row $\kappa_g^2\hbar$.\\
EFT degree & Mixed scalar two-point local bank through effective dimension six; lower mass/kinetic descendants retain the $\kappa_g^2$ ancestry tag.\\
Sources & Fixed finite scalar/frame and reduced gravitational source panel; complete nonlinear lapse/shift, ghost and antifield panels are not yet included.\\
Cutoff & One-shell constants uniform in shell index and periods on the declared smooth native profile class.\\
RG steps & Coefficientwise two-boundary trajectory at the displayed finite orders; no analytic gravitational UV trajectory.\\
Full ESM & Not proved. Nonlinear gravity constraints/measure, combined BV, gauge/Yukawa/chiral merger and arbitrary-order closure remain open.\\
\bottomrule
\end{tabular}
\end{center}

\section*{Append-only continuation marker: V74}
\addcontentsline{toc}{section}{Append-only continuation marker: V74}
\label{sec:v74-continuation-marker}
\textbf{End of V74.} Preserve Sections 1--570.  The first retained-scalar
mixed gravity--Higgs reference row is now inside the filtered two-boundary
architecture: its contact and graviton--scalar loop are derived from one finite
joint shell, its lower canonical descendants retain their gravitational EFT
ancestry, and its dimension-six local/complement split realizes the same strict
$2<4$ growth--decay gap as the scalar spine.  The next load-bearing step is the
nonlinear constraint/ghost/measure completion, because the TT reference alone
cannot decide the scalar-potential arity-eight ancestor.



\clearpage
\part*{V75: Soft-constraint retention and the locality obstruction to all-scale TT reduction}
\addcontentsline{toc}{part}{V75: Soft-constraint retention and the locality obstruction to all-scale TT reduction}

\section{Why the next step must go upstream of the TT reference}
\label{sec:v75-upstream}

Sections 1--570 and all of their labels are unchanged.  V74 identified a
possible scalar-potential direction that is invisible to the pure TT reference.
The first question is not yet its numerical coefficient.  It is whether the
operation that would generate it belongs inside a local Wilsonian shell map in
the first place.

The answer is already visible in the exact linear constraint geometry of V6.
Solving the Hamiltonian constraint at \emph{all} momenta produces an
inverse-Laplacian response.  For a source-complete local RG state this is the
wrong operation: the soft constraint fields are massless geometric variables
and should remain on the retained side of the shell.  Only their hard Fourier
components may be eliminated in a local Wilsonian step.

This version makes that distinction precise.  It proves three statements.
First, global TT reduction sends a local Higgs-potential source to a nonlocal
inverse-Laplacian field response, and one explicit volume-coupling contribution
then has scalar arity eight.  Second, after the constraint/ghost complex is
split at the same smooth shell as the propagating graviton, a purely soft Higgs
potential has no hard constraint component, so the V6 normal/ghost cancellation
remains exact and no such arity-eight term is generated by the hard shell.
Third, the linear BRST differential commutes with this shell split.  Thus the
correct next reference is a \emph{soft-constraint-retained} mixed state, not a
further reduction to TT variables at all scales.

No nonlinear coframe/connection measure is inferred from this argument.  The
nonlinear shell remains the next open theorem.

\section{Global constraint elimination is nonlocal on a Higgs-potential source}
\label{sec:v75-nonlocal-reduction}

Work in the V6 flat reduced canonical chart at a nonzero spatial momentum.
Write
\[
 \lambda=\lambda_a(k)>0,\qquad \tau=\operatorname{tr}u,
\]
and retain the V6 gauge condition
\[
 F_i=\delta_j u_{ij}-\tfrac12\delta_i\tau=0.
\]
In Fourier variables it implies
\begin{equation}
 \kappa_j u_{ij}=\frac12\kappa_i\tau,
 \qquad
 \kappa^*u\kappa=\frac12\lambda\tau.
 \label{eq:v75-gauge-trace}
\end{equation}
Hence the V6 scalar constraint reduces exactly to
\begin{equation}
 \boxed{C_0^{\rm grav}(k)=\frac12\lambda_a(k)\tau(k).}
 \label{eq:v75-C0-trace}
\end{equation}

Let
\[
 \mathcal V_H(x)=g_H\bigl(\phi(x)^T\phi(x)\bigr)^2
\]
be the quartic Higgs-potential density.  The common minimally coupled action
has a nonzero matter-energy contribution to the Hamiltonian constraint and a
nonzero first variation of the spatial volume density.  Since V6 did not fix
their common numerical normalization in the reduced quantum reference, keep
those two inherited nonzero normalizations explicit:
\begin{align}
 C_0^{\rm tot}(k)
 &=\frac12\lambda_a(k)\tau(k)
   +\kappa_g\alpha_V\widehat{\mathcal V_H}(k)+\cdots,
 \label{eq:v75-matter-constraint}\\
 S_{V}^{[1]}
 &=\frac{\kappa_g\beta_V}{2}
   \sum_{k\ne0}\tau(-k)\widehat{\mathcal V_H}(k)+\cdots.
 \label{eq:v75-volume-coupling}
\end{align}
Here the dots contain the other geometric and matter pieces of the common
action.  The theorem below uses only the displayed linear source map.  In the
usual metric normalization the factor $1/2$ in
\eqref{eq:v75-volume-coupling} is the derivative of the spatial volume density;
we do not need to identify $\alpha_V$ with $\beta_V$.

\begin{theorem}[Locality obstruction to an all-scale TT reduction]
\label{thm:v75-global-reduction-obstruction}
On the nonzero-mode branch, solving the displayed Hamiltonian constraint and
V6 gauge condition gives
\begin{equation}
 \boxed{
 \tau(k)=-2\kappa_g\alpha_V\,
              \lambda_a(k)^{-1}\widehat{\mathcal V_H}(k)+\cdots .}
 \label{eq:v75-trace-solution}
\end{equation}
Consequently the displayed volume-coupling piece alone contributes
\begin{equation}
 \boxed{
 S_{V,\rm red}^{[8]}
 =-\kappa_g^2\alpha_V\beta_V
   \sum_{k\ne0}
   \frac{\widehat{\mathcal V_H}(-k)\widehat{\mathcal V_H}(k)}
        {\lambda_a(k)}
   +\cdots .}
 \label{eq:v75-nonlocal-eight}
\end{equation}
It has scalar arity eight and an inverse-Laplacian kernel.  In particular the
field-elimination map \eqref{eq:v75-trace-solution} is not a finite local Taylor
operator at $k=0$.

Equation \eqref{eq:v75-nonlocal-eight} is not asserted to be the complete
reduced action coefficient: shifted gravitational, ghost or measure terms may
modify or cancel this particular contribution.  However, such a cancellation
would not make the elimination map \eqref{eq:v75-trace-solution} local.
Therefore a source-complete Wilsonian state based on local/quasilocal kernels
cannot globally eliminate the soft Hamiltonian-constraint sector unless a
separate source-complete cancellation theorem removes every occurrence of this
inverse-Laplacian response.
\end{theorem}
\begin{proof}
Equation \eqref{eq:v75-C0-trace} follows directly from
\eqref{eq:v75-gauge-trace}.  Setting the displayed part of
\eqref{eq:v75-matter-constraint} to zero gives
\eqref{eq:v75-trace-solution}.  Substitution into
\eqref{eq:v75-volume-coupling} gives \eqref{eq:v75-nonlocal-eight}.

For the locality statement, on the ordinary branch
$\lambda_a(k)=|k|^2+O(a^4|k|^6)$.  The multiplier
$\lambda_a(k)^{-1}$ therefore has a second-order pole at $k=0$ and has no
finite Taylor expansion there.  A finite local differential operator has a
polynomial Fourier symbol; hence the displayed elimination is not local.  The
statement concerns the source-complete elimination map and is unaffected by a
possible cancellation in one particular action coefficient.
\end{proof}

\begin{remark}[The arity-eight question is relocated, not answered by fiat]
The theorem sharpens Proposition~\ref{prop:v74-tt-not-full}.  A pure potential
exchange obtained by solving a massless soft constraint is naturally a
retained soft-gravity effect, not a new local hard-shell counterterm.  If one
first integrates out the soft constraint field, the inverse-Laplacian kernel
appears.  A Wilsonian shell should instead keep that field and integrate only
its hard component.  The possible nonlinear hard-shell descendants are a
different calculation and remain open.
\end{remark}

\section{A shell-local split of the constraint--ghost complex}
\label{sec:v75-shell-split}

Use the same spatial shell profile as V6.  Let
\[
 H_j(k)=\chi\!\left(\lambda_a(k)/\Lambda_j^2\right),
 \qquad 0\le H_j\le1,
\]
with $\operatorname{supp}\chi\subset(1,4)$.  Thus every hard gravitational
mode obeys
\begin{equation}
 \Lambda_j<\sqrt{\lambda_a(k)}<2\Lambda_j.
 \label{eq:v75-hard-annulus}
\end{equation}
The soft gravitational and ghost variables are not removed.  The hard normal
variables, hard TT variables and hard ghost variables are integrated in the
same Fourier shell.  As in V6, the zero spatial mode is retained.

Let all retained scalar source momenta obey $|p|\le\mu_j$, and assume the
nonaliasing/support separation
\begin{equation}
 \boxed{
 4\mu_j<\Lambda_j,
 \qquad
 2\Lambda_j+4\mu_j<c_{\rm BZ}a^{-1},}
 \label{eq:v75-support-domain}
\end{equation}
where the second inequality keeps all routed momenta inside the native
nonaliased Brillouin region.  The numerical value of $c_{\rm BZ}$ is fixed by
the inherited profile support; no estimate below depends on the periods.

\begin{lemma}[A purely retained quartic potential has no hard constraint mode]
\label{lem:v75-soft-potential-support}
If each scalar Fourier argument lies in $|p|\le\mu_j$, then
\begin{equation}
 \operatorname{supp}\widehat{\mathcal V_H}
 \subseteq\{|k|\le4\mu_j\}.
 \label{eq:v75-potential-support}
\end{equation}
Under \eqref{eq:v75-support-domain},
\begin{equation}
 \boxed{H_j\widehat{\mathcal V_H}=0.}
 \label{eq:v75-no-hard-V}
\end{equation}
The same statement holds for every polarization of the four retained scalar
legs.
\end{lemma}
\begin{proof}
The Fourier transform of $(\phi^T\phi)^2$ is a four-fold convolution.  The
sum of four momenta of norm at most $\mu_j$ has norm at most $4\mu_j$.
The hard shell starts above $\Lambda_j$, proving
\eqref{eq:v75-no-hard-V}.  The second condition in
\eqref{eq:v75-support-domain} prevents wraparound from invalidating this
support argument.
\end{proof}

The V6 linear BRST differential is translation invariant and momentum
diagonal.  Denote it by $s_0$.  It sends the metric and momentum variables to
linear gauge variations at the same Fourier momentum, sends antighosts to the
auxiliary field at that momentum, and sends the linear ghosts to zero.

\begin{proposition}[The linear BRST complex splits exactly into hard and soft subcomplexes]
\label{prop:v75-brst-split}
On the V6 flat reference,
\begin{equation}
 \boxed{[s_0,H_j]=0.}
 \label{eq:v75-brst-filter-commute}
\end{equation}
Consequently the hard and soft Fourier sectors are separately $s_0$-invariant.
On every nonzero hard mode the V6 identity
\begin{equation}
 \delta(C)\delta(F)|\det\{F,C\}|\,\dd u\,\dd\pi
 =\dd u_{\rm TT}\,\dd\pi_{\rm TT}
 \label{eq:v75-hard-ghost-cancel}
\end{equation}
continues to hold with the same oriented hard ghost determinant.  The soft
constraint, gauge, ghost and auxiliary variables remain explicit coordinates
of the retained state.
\end{proposition}
\begin{proof}
Every entry of the linear BRST operator is a polynomial in the translation
symbols $\kappa_i(k)$ and acts at fixed $k$.  The shell multiplier $H_j$ is a
scalar function of $\lambda_a(k)$, so it commutes with all such entries.  This
proves \eqref{eq:v75-brst-filter-commute}.  Restricting V6's finite-dimensional
constraint/gauge calculation to the hard Fourier blocks then gives
\eqref{eq:v75-hard-ghost-cancel} mode by mode.  No soft determinant is divided
out because the soft modes have not been integrated.
\end{proof}

\section{The candidate arity-eight term is absent from the hard reference shell}
\label{sec:v75-eight-absent}

We can now answer the specific ambiguity left by V74 at the level of the
linear gravitational reference.

\begin{theorem}[No pure-soft arity-eight forcing from the hard constraint shell]
\label{thm:v75-no-hard-eight}
Consider the shell-local reference of Section~\ref{sec:v75-shell-split} with
the Higgs quartic potential retained in the common action.  At perturbative
ancestry $\kappa_g^2g_H^2$ and with all eight external scalar legs in the
retained window $|p_i|\le\mu_j$, the hard normal/ghost elimination produces
no coefficient whose only internal line is a hard Hamiltonian-constraint,
shift, or gauge-normal mode.  In particular the candidate pure-soft
constraint-mediated arity-eight row is exactly zero in the hard shell.

The nonlocal expression \eqref{eq:v75-nonlocal-eight} appears only if the soft
constraint field is also eliminated.  It is therefore a retained soft-gravity
exchange (or a globally reduced representation of one), not a fresh local
hard-shell Wilson coefficient.
\end{theorem}
\begin{proof}
A hard constraint or gauge-normal line carries one Fourier momentum in the
annulus \eqref{eq:v75-hard-annulus}.  A vertex made only from four retained
scalar legs contributes the potential source
$\widehat{\mathcal V_H}(k)$ at that same momentum.  By
Lemma~\ref{lem:v75-soft-potential-support} this source is zero on the hard
annulus.  Hence each endpoint of a proposed one-line hard exchange vanishes.

Equivalently, in the delta-function form of the hard normal integral, the
matter shift of the hard Hamiltonian constraint is
$H_j\widehat{\mathcal V_H}=0$.  The hard constraint remains homogeneous, so
Proposition~\ref{prop:v75-brst-split} reduces the normal/ghost measure exactly
to the hard TT measure with no source-dependent residue.  Eliminating the
soft constraint as well would instead use
\eqref{eq:v75-trace-solution}; that is a different, nonlocal operation.
\end{proof}

\begin{corollary}[Where a nonzero nonlinear constraint contribution can first come from]
\label{cor:v75-hard-descendant-origin}
At the same finite shell, a nonzero constraint/ghost contribution involving
the Higgs potential must contain at least one of the following:
\begin{enumerate}[label=\textup{(\roman*)},leftmargin=2.7em]
\item a hard matter line or another hard field whose momentum can route the
      potential source into the hard annulus;
\item a nonlinear gravitational/ghost vertex carrying more than one shell
      momentum;
\item a source profile outside the retained window used in
      \eqref{eq:v75-support-domain}.
\end{enumerate}
The first two are genuine interacting shell diagrams and must be assigned by
the full nonlinear measured action.  They cannot be inferred from the global
soft constraint solve.
\end{corollary}

\section{The first soft-constraint-retained filtered state}
\label{sec:v75-state}

The previous theorem changes the correct iterative state.  The V73 scalar
spine and V74 TT mixed row remain valid hard-shell coordinates, but the soft
normal/ghost sector must be retained as part of the same local functional.
At the present finite order define schematically
\begin{equation}
 \boxed{
 \mathfrak S_j^{\rm scr}
 =\bigl(
 u_j,\ w_{6,j},\ R_j;
 \mathcal N_{<,j},\mathcal G_{<,j};
 \mathcal J_j,\mathcal T_j
 \bigr).}
 \label{eq:v75-state}
\end{equation}
Here $u_j,w_{6,j},R_j$ contain the V73 scalar and V74 TT local/complement
coordinates, $\mathcal N_{<,j}$ denotes the retained soft lapse/shift,
longitudinal/trace and auxiliary source coordinates, and $\mathcal G_{<,j}$
denotes their soft ghost/antighost/antifield packet.  These are functional
coordinates, not new independent physical couplings.

The local bank at fixed $(N,\Delta)$ contains the finite set of monomials in
these soft fields allowed by the inherited combined symmetry and EFT weight.
For the present Higgs-potential panel it includes the unreduced local
constraint/source vertex itself, rather than the nonlocal kernel obtained by
solving it.  The exact field inventory is intentionally larger than the TT
bank and smaller than the final nonlinear coframe--connection BV bank.

\begin{theorem}[Reference self-map with soft constraints retained]
\label{thm:v75-reference-selfmap}
Fix the V73 scalar truncation through $(g_H^2,\hbar^2,\Delta=6)$ and the V74
mixed TT row through $(\kappa_g^2\hbar,\Delta=6)$.  Enlarge the state by the
soft linear constraint/ghost packet in \eqref{eq:v75-state}.  Under
\eqref{eq:v75-support-domain}, one reference shell has the exact factorization
\begin{equation}
 \boxed{
 \mathcal R_j^{\rm ref}
 =\mathcal R_{j,\rm TT+H}^{\rm V73/V74}
   \otimes \operatorname{id}_{\mathcal N_<,\mathcal G_<}
 }
 \label{eq:v75-reference-factorization}
\end{equation}
through the displayed perturbative orders, modulo field-independent hard
normal/ghost normalization already canceled in V6.

In particular:
\begin{enumerate}[label=\textup{(\alph*)},leftmargin=2.7em]
\item the candidate pure-soft $\kappa_g^2g_H^2\phi^8$ hard-shell forcing is
      absent;
\item the V73/V74 local dimension-six coordinates still have inverse UV
      exponent $s=2$;
\item their stored complement still carries with exponent $t=4$;
\item the same admissible filtered window
      \begin{equation}
       \boxed{2\le\sigma=3<4}
       \label{eq:v75-window}
      \end{equation}
      remains valid for this enlarged linear reference.
\end{enumerate}
\end{theorem}
\begin{proof}
Proposition~\ref{prop:v75-brst-split} gives an exact hard/soft tensor product
for the linear constraint--ghost complex.  The hard normal/ghost factor reduces
to the TT hard measure by \eqref{eq:v75-hard-ghost-cancel}.  The soft factor is
not integrated.  Theorem~\ref{thm:v75-no-hard-eight} shows that adding the
purely retained quartic potential does not spoil this factorization on the
specified panel.  The remaining hard scalar and TT coefficients are precisely
the already computed V73/V74 maps.  Their growth and complementary carry
exponents are unchanged, giving \eqref{eq:v75-window}.
\end{proof}

\begin{corollary}[The V74 two-boundary trajectory survives the locality repair]
\label{cor:v75-two-boundary}
On the finite source panel and perturbative truncation of
Theorem~\ref{thm:v75-reference-selfmap}, the coefficientwise two-boundary
trajectory of Corollary~\ref{cor:v74-two-boundary} remains valid after the
soft linear constraint/ghost packet is adjoined.  Its ultraviolet boundary
loss has the same form
\begin{equation}
 \|x_j^{(J)}-x_j^{(\infty)}\|
 \le C b^{3j}(1+J)^q\theta^{J-j},
 \qquad b^{-1}<\theta<1.
 \label{eq:v75-boundary-loss}
\end{equation}
The statement now has a local source-complete interpretation at the linear
gravity reference: soft constraint variables are retained rather than hidden
inside inverse-Laplacian reduced kernels.
\end{corollary}

\section{What remains for the nonlinear gravitational shell}
\label{sec:v75-next}

The locality repair is exact only for the linear flat gravitational complex.
The nonlinear BRST differential does not commute with a scalar Fourier shell:
products of fields and ghosts mix hard and soft momenta.  Likewise the native
coframe/connection density, contorsion determinant and nonlinear
Faddeev--Popov matrix need not factor into V6's hard determinant times a soft
spectator.

The correct next theorem is therefore not an arity-eight coefficient extracted
from a globally reduced TT action.  It is a \emph{shell-local nonlinear BV
reduction theorem}.  In schematic form, with
\[
 s=s_0+s_1+s_2+\cdots,
\]
one must construct a local hard elimination for which
\begin{equation}
 \boxed{
 \mathcal R_j\circ s
 =s_{j+1}\circ\mathcal R_j+\mathcal E_j,
 \qquad
 \|\mathcal E_j\|_{N,\Delta}
 \le C_{N,\Delta}\,\text{(summable/filtered shell weight)},}
 \label{eq:v75-next-target}
\end{equation}
while retaining the soft lapse/shift, Lorentz, constraint and ghost source
coordinates.  Its local bank must contain every hard-matter descendant of the
constraint sector, including any genuine arity-eight local jet generated by
hard internal lines.

The distinction is now sharp:
\begin{equation}
 \boxed{\text{soft massless constraint exchange is retained as a field}}
 \label{eq:v75-soft-principle}
\end{equation}
rather than being integrated into a nonlocal Wilson term, whereas
\begin{equation}
 \boxed{\text{hard constraint/ghost/matter subgraphs are integrated locally}}
 \label{eq:v75-hard-principle}
\end{equation}
with their local Taylor data stored in the filtered EFT bank.

This is the shell organization required before the nonlinear measure can be
meaningfully compared with the filtered ESM theorem.  It also aligns the
quantum shell with the programme's instruction to use the same finite common
Einstein--SM carrier rather than replace it by a globally reduced nonlocal
matter theory.

\section{Anti-circularity audit and status ledger for V75}
\label{sec:v75-ledgers}

\subsection*{Proved}
The V6 gauge equations imply the exact trace identity
\eqref{eq:v75-C0-trace}.  A local quartic Higgs-potential contribution to the
Hamiltonian constraint therefore produces the inverse-Laplacian global response
\eqref{eq:v75-trace-solution}; one displayed volume-coupling contribution has
the arity-eight form \eqref{eq:v75-nonlocal-eight}.  This proves a locality
obstruction to eliminating the soft constraint sector inside a source-complete
local RG state.  Splitting the linear constraint/ghost complex with the same
smooth shell as the graviton gives exact hard and soft BRST subcomplexes.  On a
retained scalar window satisfying \eqref{eq:v75-support-domain}, the quartic
potential has no hard Fourier component, so the V6 hard normal/ghost cancellation
is unchanged and the candidate pure-soft arity-eight hard-shell coefficient is
exactly absent.  The V73/V74 filtered reference self-map and its strict $2<4$
growth--decay window therefore survive after soft constraints and ghosts are
retained.

\subsection*{Inherited}
V6 supplies the linear constraints, gauge functions, exact normal/ghost measure
cancellation, TT covariance and filtered two-boundary theorem.  The common
finite ESM carrier supplies the local Higgs potential and its geometric volume
coupling.  V7--V8 establish that nonlinear mixed BV completion cannot be inferred
from the linear ghost algebra.  V73 supplies the interacting scalar filtered
spine, and V74 supplies the first retained-scalar TT mixed row and its local
ancestry bookkeeping.

\subsection*{Literature input}
No external numerical coefficient is used.  The only general background is the
standard Wilsonian distinction between retained massless fields and integrated
hard modes.  Every support and constraint identity used here is derived from the
native finite V6 symbols.  No continuum gauge-fixed graviton propagator is
substituted for the measured finite shell.

\subsection*{Conditional}
The nonzero constants $\alpha_V,\beta_V$ in
\eqref{eq:v75-matter-constraint}--\eqref{eq:v75-volume-coupling} encode the
normalization of the inherited common action in the reduced V6 chart; their
numerical values are not evaluated here.  The arity-eight expression
\eqref{eq:v75-nonlocal-eight} is one displayed contribution to a globally
reduced action, not the complete coefficient.  The hard-shell factorization is
proved only for the linear flat gravitational BRST complex and the specified
soft source window.

\subsection*{Open}
The smallest load-bearing theorem is now the shell-local \emph{nonlinear}
constraint/Lorentz/ghost/measure elimination on the unreduced finite
Einstein--SM carrier.  It must control the hard--soft mixing of the nonlinear
BRST differential, the coframe/connection density and Faddeev--Popov/contorsion
factors, retain all soft constraint sources, and place every genuine hard-matter
descendant into a finite filtered local bank.  Only after that theorem can the
V75 reference self-map be promoted to the full mixed ESM one-shell map and the
combined BV identity iterated.

\begin{center}\footnotesize
\begin{tabular}{p{.23\textwidth}p{.69\textwidth}}
\toprule
Variable & Scope of V75\\\midrule
Perturbative order & V73 scalar spine and V74 mixed TT row, plus the linear constraint/ghost reference needed to interpret the next mixed order.\\
EFT degree & Local scalar/mixed bank through dimension six; no local arity-eight coefficient is assigned from a soft constraint solve.\\
Sources & Fixed finite scalar and linear gravitational constraint/ghost panels with $4\mu_j<\Lambda_j$; zero spatial modes remain retained.\\
Cutoff & Exact support and BRST-split statements at every shell satisfying the nonaliasing domain; constants are period independent.\\
RG steps & The V74 coefficientwise two-boundary trajectory survives on the enlarged linear reference.\\
Full ESM & Not proved. Nonlinear coframe/connection gauge fixing, hard--soft BV mixing, measure/Jacobian factors, internal gauge/Yukawa/chiral merger and arbitrary-order closure remain open.\\
\bottomrule
\end{tabular}
\end{center}

\section*{Append-only continuation marker: V75}
\addcontentsline{toc}{section}{Append-only continuation marker: V75}
\label{sec:v75-continuation-marker}
\textbf{End of V75.} Preserve Sections 1--577.  The potential arity-eight
ambiguity has been moved to its correct logical location: global elimination of
a soft massless constraint is nonlocal and is not a Wilsonian hard-shell
operation.  Retaining the soft constraint/ghost sector while eliminating only
its hard linear component preserves the exact V6 ghost cancellation, removes
the spurious pure-soft arity-eight hard forcing, and leaves the V73/V74 strict
$2<4$ filtered reference window intact.  The next theorem must construct this
hard/soft organization for the nonlinear finite coframe--connection BV shell.



\clearpage
\part*{V76: A shell-local nonlinear BV connection and bounded composed source transport}
\addcontentsline{toc}{part}{V76: Shell-local nonlinear BV connection and composed source transport}

\section{The exact shell object is the transported BV operator, not a projected vector field}
\label{sec:v76-start}

Sections 1--577 are preserved verbatim.  V75 isolates the correct locality
principle for the gravitational constraints: the soft lapse/shift and their
soft ghost/source partners must remain among the retained variables, whereas
only a hard shell is eliminated.  The remaining obstruction was phrased as a
failure of the nonlinear BRST differential to commute with a Fourier shell.
That phrasing is useful diagnostically but is not the correct object to
iterate.  The exact measured shell does not project a nonlinear symmetry
vector field and then ask it to remain unchanged.  It transports the complete
field--antifield operator through the same normalized hard pushforward that
transports the action and every marked source.

This installment makes that statement finite and quantitative on the first
nontrivial common bank for which all required ingredients have already been
proved.  The bank is the V26 minimal gravitational arity-five source complex,
its V27 metric--Higgs two-field extension and the V28 scalar-transformation
row, together with V75's retained soft linear constraint/ghost coordinates.
The propagating metric and ghost upper completion, coframe-coordinate density,
local Lorentz quotient and auxiliary connection amplitude are the ones already
fixed in V10--V14 and V24--V28.  No measure, contour, covariance or physical
normalization is changed.

The result has three logically separate levels.
\begin{enumerate}[label=\textup{(\roman*)},leftmargin=2.7em]
\item At every finite cutoff and fixed perturbative/source degree, the
      normalized hard pushforward conjugates the complete BV operator exactly.
      Its square, or its nonzero defect, is transported with it.
\item The difference between this induced operator and the declared retained
      local operator is a \emph{shell BV connection}.  At one loop on the
      stated bank its local Taylor jet is precisely of the type already solved
      by the V24--V25 restoring complex, while its complementary kernel has a
      summable annular bound.
\item After the local jet is absorbed into the common counterterm/source
      coordinates, the adjacent source map is identity plus a summable
      operator.  Thus the composed source transport, which earlier programmes
      treated as an assumption, is a theorem on this finite one-loop
      Einstein--Higgs panel.
\end{enumerate}

The result is deliberately not called a complete nonlinear Einstein--Standard
Model BV theorem.  Internal gauge, Yukawa and chiral-line sectors are not added
by analogy, and the two-loop master coefficient is not inferred from a
one-loop transport estimate.

\subsection*{Finite algebra and conventions}
Let \(\mathscr A_j\) be the finite polynomial field--antifield algebra at one
Wilsonian stage, truncated only after a fixed perturbative, field/source and
EFT degree has been declared.  Write the retained/hard split collectively as
\[
 \Phi=\Phi_<+\xi_j,
\]
where \(\Phi_<\) contains the soft gravitational constraint and ghost packet
of V75 and all other retained coordinates, while \(\xi_j\) contains the
propagating metric/ghost and matter variables assigned to the hard shell.  The
nonpropagating auxiliary connection is treated by its already established
local stationary amplitude and is not counted once per propagating shell.

Let \(S_j\) denote the complete once-counted action-density representative at
this stage and let \(\Delta_j\) be the canonical finite BV contraction in the
same coordinates.  Put
\begin{equation}
 D_j=(S_j,\,\cdot\,)_j-\hbar\Delta_j .
 \label{eq:v76-Dj}
\end{equation}
No assertion \(D_j^2=0\) is made.  The exact defect is retained as
\(\mathfrak Q_j=D_j^2\).

Let \(\mathscr T_j\) be the normalized hard pushforward acting on marked
insertions,
\begin{equation}
 \mathscr T_jF
 =e^{S_{j+1}/\hbar}\,
   \mathbb E_{h,j}\!\left[e^{-S_j/\hbar}F\right],
 \qquad
 e^{-S_{j+1}/\hbar}
 =\mathbb E_{h,j}\!\left[e^{-S_j/\hbar}\right],
 \label{eq:v76-Tj}
\end{equation}
with the Gaussian/Berezin expectation and density convention of the actual
shell.  At every fixed formal degree, \(\mathscr T_j=I+O(\hbar)\) on the
retained polynomial algebra and is therefore algebraically invertible order by
order.  V59 is the evaluated scalar instance of this same normalization.

\begin{theorem}[Exact induced BV operator and defect transport]
\label{thm:v76-induced}
On every fixed finite polynomial truncation on which \(\mathscr T_j\) is
invertible, define
\begin{equation}
 \boxed{\widehat D_{j+1}:=\mathscr T_jD_j\mathscr T_j^{-1}.}
 \label{eq:v76-induced-D}
\end{equation}
Then
\begin{equation}
 \boxed{
 \widehat D_{j+1}^{\,2}
   =\mathscr T_j\mathfrak Q_j\mathscr T_j^{-1}.}
 \label{eq:v76-induced-square}
\end{equation}
In particular, a vanishing complete input defect is transported exactly, while
a nonzero defect cannot be removed by the shell algebra alone.

For every marked insertion \(F\), the induced action of
\(\widehat D_{j+1}\) is the normalized connected pushforward of the complete
input action/source variation.  No commutation of the nonlinear BRST vector
field with the hard projector is required.
\end{theorem}
\begin{proof}
Equation~\eqref{eq:v76-induced-square} is the algebraic identity
\((TDT^{-1})^2=TD^2T^{-1}\).  The nontrivial point is that the normalized
hard map used here is the same map that transports the action and the marked
insertions.  At fixed formal degree, Gaussian/Berezin Wick expansion is a
finite triangular map with identity zeroth term, so its inverse exists.  The
second-order BV product rule gives the marked connected term exactly as in
Theorem~\ref{thm:v59-defect}; replacing a scalar-only expectation there by the
present finite product hard expectation changes neither the algebraic
argument nor the ordering requirement.  All field-independent density factors
are counted once in \(S_j\), so no additional weighted Laplacian is inserted.
The formula therefore transports the actual input defect rather than a
projected continuum Ward identity.
\end{proof}

This theorem is the upstream correction to the tempting equation
\([s,H_j]=0\) beyond the linear reference.  The latter is true for V75's
\(s_0\), but it is neither needed nor expected for the nonlinear differential.

\section{The shell BV connection and its exact curvature identity}
\label{sec:v76-connection}

Let \(D_{<,j+1}\) be the declared local retained operator in the same source
coordinates after the shell.  It contains the retained soft constraint and
ghost fields, the ordinary nonlinear coframe-coordinate transformation jets,
the off-shell nonminimal doublets and the same lower-reference convention.
Define the shell BV connection
\begin{equation}
 \boxed{
 \mathcal A_j^{\rm BV}
 :=\widehat D_{j+1}-D_{<,j+1}.}
 \label{eq:v76-A}
\end{equation}
It is an odd operator on the retained polynomial/source bundle.  It records
both the familiar nonlinear hard--soft transformation leakage and the quantum
contractions generated by eliminating the hard shell.

\begin{proposition}[Exact curvature identity for the shell connection]
\label{prop:v76-curvature}
Without assuming nilpotence of either operator,
\begin{equation}
 \boxed{
 D_{<,j+1}\mathcal A_j^{\rm BV}
 +\mathcal A_j^{\rm BV}D_{<,j+1}
 +(\mathcal A_j^{\rm BV})^2
 =\mathscr T_j\mathfrak Q_j\mathscr T_j^{-1}
  -D_{<,j+1}^{\,2}.}
 \label{eq:v76-curvature}
\end{equation}
Consequently the shell connection is flat exactly when the transported input
and retained output defects agree.  At first order in \(\hbar\), the quadratic
term in \(\mathcal A_j^{\rm BV}\) does not contribute.
\end{proposition}
\begin{proof}
Expand
\((D_<+\mathcal A)^2-D_<^2\) and apply
Theorem~\ref{thm:v76-induced}.  The grading signs are the ordinary signs for
the anticommutator of two odd operators.  Since
\(\mathcal A_j^{\rm BV}=O(\hbar)\) for the one-loop quantum transport after
the classical retained transformation has been put in \(D_<\), its square
starts at \(O(\hbar^2)\).
\end{proof}

Equation~\eqref{eq:v76-curvature} is useful because it separates two problems
which were previously easy to conflate.  Local restoration chooses the
coordinate representative of \(D_<\) and its counterterms.  Hard-shell
transport determines \(\mathcal A_j^{\rm BV}\).  One cannot set the latter to
zero by choosing a convenient local representative unless the measured
connection itself is exact in that representative.

\subsection{No hidden inverse soft constraint in the connection}
The shell expectation in \eqref{eq:v76-Tj} integrates only the hard variables.
The soft constraint fields of V75 remain explicit arguments of
\(\mathscr T_jF\).  Thus no factor \(\lambda_a(k)^{-1}\) associated with a
soft Hamiltonian-constraint solve is introduced into
\(\mathcal A_j^{\rm BV}\).  Hard propagator inverses are permitted: their
momenta lie on the shell and are the objects whose local Taylor/complement
splitting is estimated below.

\section{One-loop locality of the induced connection on the common minimal bank}
\label{sec:v76-locality}

We now restrict to the first bank for which every required operand has already
been populated.  Let \(\mathfrak B_{5,4}^{g}\) be V26's minimal pure-gravity
arity-five, weight-four source bank, including its ghost and multiple-antifield
rows.  Adjoin the V27 two-Higgs metric panels and V28 source panels
\begin{equation}
 \mathfrak B_{5,4}^{H}
 =\{\upsilon q^b\phi c:0\le b\le2\}
 \quad\hbox{and their measured descendants.}
 \label{eq:v76-bank}
\end{equation}
The complete bank is
\(\mathfrak B_{5,4}^{EH}=\mathfrak B_{5,4}^{g}\oplus
\mathfrak B_{5,4}^{H}\), with the V10--V12 local density and nonminimal
operands included once.

For a panel \(I\in\mathfrak B_{5,4}^{g}\), let \(n(I)\) be its number of
minimal antifields, as in V26, and define the dimensionless source norm
\begin{equation}
 \|F_I\|_{*,r}
 :=\mu^{-(4-n(I))}
   \sum_{|\alpha|\le r}\frac{\mu^{|\alpha|}}{\alpha!}
   \sup_{|K|\le\mu}|\partial_K^\alpha F_I(K)| .
 \label{eq:v76-star-norm}
\end{equation}
For a V28 scalar-source panel use the analogous mass-graded normalization,
with \(\epsilon=|K|+m_H\) and its canonical weight three.  Fixed flavour,
source and polarization norms are suppressed.

\begin{lemma}[Hard-degree reduction of a one-loop induced BV coefficient]
\label{lem:v76-hard-degree}
At fixed external arity at most five, every connected one-loop coefficient of
\(\mathcal A_j^{\rm BV}\) can be represented as a finite sum of the following
already measured hard words:
\begin{enumerate}[label=\textup{(\alph*)},leftmargin=2.7em]
\item background-complete metric/ghost resolvent words of V26;
\item the simultaneous metric--Higgs Schur words of V27;
\item the scalar-antifield three-propagator words of V28;
\item local auxiliary-connection, Haar and coordinate-density insertions of
      V10--V12, assigned once rather than once per propagating shell.
\end{enumerate}
No additional one-loop topology is created solely by retaining the soft
constraint fields of V75.
\end{lemma}
\begin{proof}
Expand the normalized hard expectation and the operator
\(D_j\) to first loop order.  A connected hard word with a vertex carrying
exactly one hard leg has that hard momentum equal to a sum of at most five
external retained momenta.  On the source panels used for the continuum
coefficient this lies strictly below the hard support by the inherited
nonaliasing window, so the corresponding covariance vanishes.  Thus every
surviving propagating vertex has at least two hard legs.

For a connected graph with one independent hard loop,
\(\sum_v(\deg_h v-2)=2(\ell-1)=0\).  Hence every surviving vertex has exactly
two hard legs.  Its background dependence is therefore obtained by
successively differentiating a hard Hessian, a ghost-action Hessian or a
linear transformation/source Hessian.  Those are exactly the dressed words
listed in V26--V28.  The retained soft constraint coordinates can mark these
vertices, but they are not integrated and therefore create no new hard line.
The auxiliary connection has a pointwise inverse rather than a propagating
shell covariance; V10 already assigns its loop density locally.  The Haar and
coframe-coordinate factors are ultralocal action-density terms.  This exhausts
the one-loop possibilities on the stated bank.
\end{proof}

The lemma is an ownership statement, not a claim that the nonlinear symmetry
has closed at arbitrary field degree.  At the next source arity new hard
vertices must be supplied before the same argument can be repeated.

\begin{theorem}[Local-plus-summable decomposition of the one-loop shell BV connection]
\label{thm:v76-local-connection}
Let \(\Lambda_j=b^j\lambda\), \(b>1\), and keep the V26/V28 source windows
with \(\lambda\ge8192\mu\), \(a\lambda\ll1\) and
\(\lambda L_\nu\ge1\).  After applying the common V24 local projection to
every component of \(\mathcal A_j^{\rm BV}\), write
\begin{equation}
 \mathcal A_j^{\rm BV}
 =\mathsf P\mathcal A_j^{\rm BV}
  +\mathcal E_j^{\rm BV}.
 \label{eq:v76-A-split}
\end{equation}
Then, at one loop on \(\mathfrak B_{5,4}^{EH}\), the complementary adjacent
operator obeys the dimensionless bound
\begin{equation}
 \boxed{
 \|\mathcal E_{0,j}^{\rm BV}\|_{*\to *}
 \le C\frac{\mu}{\Lambda_j}.}
 \label{eq:v76-shell-connection-bound}
\end{equation}
For the finite regulator and the ordinary completed reference,
\begin{equation}
 \boxed{
 \|\mathcal E_{a,j}^{\rm BV}-\mathcal E_{0,j}^{\rm BV}\|_{*\to *}
 \le C a^3\mu\Lambda_j^2}
 \label{eq:v76-shell-connection-mismatch}
\end{equation}
for interior shells; the final smooth upper cap has the corresponding V26/V28
full-zone bound.  Rooted source-kernel versions have the same scale factors and
no total-volume multiplier.

The V28 unmarked scalar-source row has the sharper ordinary shell decay
\(C(\mu/\Lambda_j)^2\), but this improvement is not used in the common bound.
\end{theorem}
\begin{proof}
For the pure-gravity slots, Proposition~\ref{prop:v26-shells} gives
\[
 \|(I-\mathsf T_{4-n})\Delta_j\Gamma_{1,I}\|
 \le C\mu^{5-n}\Lambda_j^{-1}.
\]
The finite BV/source operator is a fixed sum of local multipliers and Euler
insertions.  In the scaled norm \eqref{eq:v76-star-norm}, one additional soft
derivative or the corresponding canonical weight converts the right-hand side
to \(C\mu/\Lambda_j\), uniformly over the finite bank.  This is the same
count used in the proof of \eqref{eq:v26-W-rate}, now before summing the shell
labels.

For the mixed scalar-source slots, the proof of
Theorem~\ref{thm:v28-cutoff} gives, after the required weighted derivatives,
a continuum hard integrand of order \(R^{-5}\).  Restricting it to one annulus
therefore gives \(C\Lambda_j^{-1}\) before restoring the canonical external
weight, hence again \(C\mu/\Lambda_j\) in the dimensionless bundle norm.  The
unmarked odd row has order \(R^{-6}\) and gives the stated stronger decay.
V27's mixed action coefficient has the same or better fifth-derivative
majorant.

The interior finite--ordinary differences in V26 have the bound
\(a^3\mu^{5-n}\Lambda_j^2\).  Applying the same fixed local operator and
normalizing by \(\mu^{4-n}\) proves
\eqref{eq:v76-shell-connection-mismatch}.  V28's corresponding difference is
no worse on the common bank.  The auxiliary and density factors are local and
are not repeated in this annular estimate.  Their fixed source derivatives
were already proved volume-uniform in V10--V12.  Fourier integration by parts
in the independent external variables gives the rooted versions.
\end{proof}

\section{The local connection is absorbed by the existing common restoring bank}
\label{sec:v76-restoration}

The preceding theorem deliberately leaves the local Taylor jet
\(\mathsf P\mathcal A_j^{\rm BV}\) untouched.  That jet is not an error: it is
part of the running local EFT/source coordinates.  The question is whether it
requires a new independent symmetry structure after the V75 split.

\begin{theorem}[Shell-local use of the V24--V25 restoring primitive]
\label{thm:v76-local-primitive}
On the pure-gravity component of \(\mathfrak B_{5,4}^{EH}\), the local jet of
the one-loop shell BV connection lies in the same finite degree-closed complex
used in V24--V25.  Consequently there is a shell-local ghost-number-zero
coefficient \(C^{\rm BV}_{j,\rm loc}\) in that bank such that
\begin{equation}
 \boxed{
 \mathsf P\mathcal A_{j}^{\rm BV}
 =D_{<,0}\,C^{\rm BV}_{j,\rm loc}
   +\mathcal X^{\rm ref}_{j,\rm loc}}
 \label{eq:v76-local-primitive}
\end{equation}
with the same lower-reference term and normalization convention as V25--V26.
Its finite-dimensional section norm is bounded independently of the shell
number after canonical scaling.

On the mixed \(c\phi^2\) component, no pure-gravity cohomology theorem is
imported.  Instead the six-term measured identity of V28 fixes the local
combination of metric source, scalar source, matter stress and lower-reference
trace.  Thus its local jet is stored as one common mixed action/source
coordinate rather than independently normalized pieces.
\end{theorem}
\begin{proof}
V24 proves that the ordinary arity-five, weight-four gravitational source
space is degree closed under the ordinary BV differential and that its common
Taylor projector is a chain map.  V25 proves exactness of the extendable
critical ghost-number-one germ in the stated local coframe chart and supplies
a bounded finite-dimensional primitive.  The local jet in
Theorem~\ref{thm:v76-local-connection} has exactly this field, ghost, antifield
and derivative content by Lemma~\ref{lem:v76-hard-degree}; retaining the soft
constraint variables changes its coefficient arguments but does not create a
new inverse momentum or new local tensor type.  Applying the same finite
section shell by shell proves \eqref{eq:v76-local-primitive}.

For the Higgs row, Theorem~\ref{thm:v28-finite-row} is the relevant measured
identity.  Its right side contains the actual classical breaking insertion and
reference trace, and its source-loop coefficient includes all six terms of the
metric/scalar transformation.  The local Taylor operation is linear and acts
on the complete identity.  Hence the shell-local mixed jet is one constrained
coordinate.  Declaring its separate pieces to be independently adjustable
would violate that identity.
\end{proof}

\subsection{Measure ownership}
Three measure contributions remain distinct throughout the construction.
\begin{enumerate}[label=\textup{(\alph*)},leftmargin=2.7em]
\item V10's auxiliary-connection determinant and Lorentz-Haar density are
      assigned once as a local action-density amplitude; they are not added
      once per propagating metric shell.
\item V12's metric-to-coframe-coordinate Jacobian stays paired with the
      nonlinear transformation divergence.  Its local jet belongs to the same
      restoring functional and is not a new ghost determinant.
\item Propagating metric/ghost determinants and their source derivatives are
      part of the V26 shell words and therefore are assigned by maximal
      internal shell label exactly once.
\end{enumerate}
This prevents the shell-local primitive from canceling a local coefficient
whose measure partner has accidentally been counted a second time.

\section{Bounded composed source transport on the nonlinear one-loop panel}
\label{sec:v76-transport}

After adding the local primitive of
Theorem~\ref{thm:v76-local-primitive}, use the corresponding renormalized
source coordinates at each scale.  The remaining adjacent map on the fixed
source bundle is
\begin{equation}
 \boxed{
 \mathsf U_{j+1,j}=I+E_j,
 \qquad
 \|E_j\|_{*\to *}\le C\frac{\mu}{\Lambda_j}.}
 \label{eq:v76-adjacent}
\end{equation}
The identity term includes the canonical rescaling already built into the
star norm; all finite wave-function/source normalizations are in the local
coordinate choice.

\begin{theorem}[Uniform composed source transport]
\label{thm:v76-composed}
For \(\Lambda_j=b^j\lambda\), define
\begin{equation}
 \mathsf U_{n,j}
 =\mathsf U_{n,n-1}\cdots\mathsf U_{j+1,j}.
 \label{eq:v76-composed}
\end{equation}
Then, on the fixed one-loop bank \(\mathfrak B_{5,4}^{EH}\),
\begin{equation}
 \boxed{
 \sup_{n\ge j}\|\mathsf U_{n,j}\|_{*\to *}
 \le
 \exp\!\left[
  \frac{Cb}{b-1}\frac{\mu}{\Lambda_j}
 \right].}
 \label{eq:v76-composed-bound}
\end{equation}
In particular the composed transports are uniformly bounded in the number of
ultraviolet shell steps.

Let \(\mathsf U^{(a)}\) and \(\mathsf U^{(0)}\) denote the finite and ordinary
completed transports with the same local normalization.  Up to the common
upper cap and the separate finite-volume term,
\begin{equation}
 \boxed{
 \sup_{n\le J}
 \|\mathsf U^{(a)}_{n,0}-\mathsf U^{(0)}_{n,0}\|_{*\to *}
 \le C'e^{C\mu/\lambda}\,a\mu,}
 \label{eq:v76-composed-mismatch}
\end{equation}
whenever \(\Lambda_J\le c/a\).  The upper-cap contribution obeys the inherited
full-zone vanishing bound, and the thermodynamic comparison is independent.
\end{theorem}
\begin{proof}
From \eqref{eq:v76-adjacent},
\[
 \|\mathsf U_{n,j}\|
 \le\prod_{k=j}^{n-1}(1+\|E_k\|)
 \le\exp\!\left(\sum_{k=j}^{n-1}\|E_k\|\right).
\]
The geometric series is
\[
 \sum_{k=j}^{\infty}\frac{\mu}{\Lambda_k}
 =\frac{\mu}{\Lambda_j}\frac{b}{b-1},
\]
which proves \eqref{eq:v76-composed-bound}.

For the finite--ordinary difference, telescope the product once at every
adjacent factor.  The factors on both sides of the replaced term are bounded
by \eqref{eq:v76-composed-bound}.  Theorem~\ref{thm:v76-local-connection}
gives
\[
 \sum_{k=0}^{J}a^3\mu\Lambda_k^2
 \le C_ba^3\mu\Lambda_J^2
 \le C_ba\mu.
\]
This proves \eqref{eq:v76-composed-mismatch} for the interior shells.  The
smooth upper completion is one grouped factor controlled by the existing
full-zone comparison, not by extending the interior estimate beyond its
support.  The independent volume errors compose by the same telescoping
argument because their fixed-shell sum is absolutely convergent.
\end{proof}

This is the first place in the ESM lineage where a composed source-transport
bound of the form required by the perturbative continuum programme is obtained
from the actual shell coefficients rather than postulated.  Its scope is still
the fixed one-loop Einstein--Higgs panel, not the chiral Standard Model source
bundle.

\section{One-loop retained-constraint Einstein--Higgs self-map}
\label{sec:v76-selfmap}

Combine V75's soft-constraint-retained state with the local and complementary
coordinates above.  Denote the resulting one-loop source sector by
\begin{equation}
 \mathfrak C_{1,\le5}^{EH}
 =\bigl(
  u_{\rm loc},\ C_{\rm BV,loc},\ R_{\rm ql};
  \mathcal N_<,\mathcal G_<;\mathcal J
 \bigr),
 \label{eq:v76-state}
\end{equation}
where \(R_{\rm ql}\) is the quasilocal complementary source kernel and
\(\mathcal J\) uses the renormalized coordinates of
Section~\ref{sec:v76-transport}.  The V73/V74 dimension-six action spine is a
compatible direct summand at its already proved orders; it is not needed to
prove the one-loop BV statement below.

\begin{theorem}[First nonlinear one-loop self-map with soft constraints retained]
\label{thm:v76-selfmap}
Fix the common completed regulator, the V26/V28 source windows and a finite
period torus satisfying the inherited nonaliasing and lower-scale margins.
At one loop and on the bank \(\mathfrak B_{5,4}^{EH}\), one hard shell maps
\begin{equation}
 \boxed{
 \mathcal R_j^{(1)}:
 \mathfrak C_{1,\le5}^{EH}
 \longrightarrow
 \mathfrak C_{1,\le5}^{EH}.}
 \label{eq:v76-selfmap}
\end{equation}
More precisely:
\begin{enumerate}[label=\textup{(\roman*)},leftmargin=2.7em]
\item soft constraint, Lorentz and nonminimal source coordinates remain
      retained; no inverse soft wave operator is introduced;
\item every hard local Taylor coefficient lands in the finite V24--V25
      gravitational bank or in the common V28 mixed Higgs source coordinate;
\item every nonlocal one-loop source coefficient lands in \(R_{\rm ql}\) with
      adjacent norm \(O(\mu/\Lambda_j)\);
\item the complete finite BV defect is transported by
      Theorem~\ref{thm:v76-induced}, while the difference between the induced
      and declared local representatives is the controlled shell connection;
\item the composed renormalized source transports obey
      \eqref{eq:v76-composed-bound}, uniformly in the number of shell steps;
\item finite-to-ordinary transport differences vanish with
      \eqref{eq:v76-composed-mismatch}, together with the inherited independent
      volume and upper-cap errors.
\end{enumerate}
Thus the nonlinear hard--soft BRST mixing that was open in V75 is closed at
this one-loop, finite-source level without requiring the nonlinear BRST vector
field itself to commute with the shell projector.
\end{theorem}
\begin{proof}
The exact operator transport is Theorem~\ref{thm:v76-induced}.  Lemma
\ref{lem:v76-hard-degree} identifies the complete set of one-loop hard words
on the bank.  Theorem~\ref{thm:v76-local-connection} splits them into the
common local bank and quasilocal complement.  The local part is absorbed by
Theorem~\ref{thm:v76-local-primitive}, with the density ownership fixed there.
The soft fields were never integrated, so no inverse soft constraint appears.
The complementary coordinate and source transport estimates are precisely
Theorem~\ref{thm:v76-composed}.  These operations retain all terms rather than
projecting a noninvariant answer after the shell.  Hence every output component
is again one of the coordinates in \eqref{eq:v76-state}.
\end{proof}

\begin{corollary}[Compatibility with the V73--V75 filtered action spine]
\label{cor:v76-filtered}
Take the direct sum of \(\mathfrak C_{1,\le5}^{EH}\) with the V73 scalar
\((g_H^2,\hbar^2,\Delta=6)\) spine and the V74 mixed
\((\kappa_g^2\hbar,\Delta=6)\) TT row, interpreted with V75 soft constraints
retained.  At the perturbative orders already proved in those versions, the
action coordinates retain the strict filtered window
\begin{equation}
 \boxed{2\le\sigma=3<4,}
 \label{eq:v76-window}
\end{equation}
while the one-loop source bundle has the independently summable transport
\eqref{eq:v76-composed-bound}.  No inequality between the action remainder
exponent and the source-transport exponent is required: they control different
coordinates of the perturbative state.
\end{corollary}

\section{Anti-circularity audit and status ledger for V76}
\label{sec:v76-ledgers}

\subsection*{Proved}
The normalized hard pushforward defines an exact induced BV operator by
conjugation, and transports the complete finite defect rather than assuming it
vanishes.  Its difference from the declared retained local differential is a
shell BV connection satisfying the exact curvature identity
\eqref{eq:v76-curvature}.  On the populated V26--V28 Einstein--Higgs one-loop
bank every induced hard coefficient is one of the already measured
metric/ghost, metric--Higgs, scalar-source or local auxiliary/density words.
After the common local projection, the complementary adjacent connection has
the dimensionless bound \(C\mu/\Lambda_j\), with finite--ordinary interior
mismatch \(Ca^3\mu\Lambda_j^2\).  The local gravitational jet is absorbed by
the existing V24--V25 primitive, while the mixed Higgs jet remains one common
measured V28 coordinate.  Consequently the composed renormalized source
transports are uniformly bounded in the number of ultraviolet shells and
their finite-to-ordinary difference is \(O(a\mu)\) up to the separately
controlled cap and volume terms.  This gives the first nonlinear one-loop
retained-constraint Einstein--Higgs self-map on the stated source bank.

\subsection*{Inherited}
V10--V12 supply the auxiliary connection amplitude, Lorentz-Haar/coframe
coordinate density and nonlinear transformation divergence.  V24--V25 supply
the degree-closed pure-gravity local complex and its bounded critical
restoring primitive.  V26 supplies the full arity-five pure-gravity one-loop
functional, modified identity, shell ownership and annular complement bound.
V27 supplies the simultaneous metric--Higgs hard block, and V28 supplies the
scalar transformation source and its finite measured identity.  V59 supplies
the exact normalized BV pushforward logic with a nonzero defect retained.
V73--V75 supply the filtered scalar/mixed action spine and the requirement that
soft massless constraints remain retained rather than globally eliminated.

\subsection*{Literature input}
No new cohomology or quantum-gravity existence theorem is invoked.  The only
external cohomological input remains the one already isolated in V25 for the
ordinary local pure-gravity critical germ.  The present shell transport and
composed bounds are derived from the finite measured coefficients and shell
estimates already established in this lineage.

\subsection*{Conditional}
The exact conjugation theorem holds at every fixed formal degree, but the
quantitative self-map is proved only at one loop on the arity-five
Einstein--Higgs source bank.  The mixed local Higgs row uses V28's measured
identity rather than a classification of the enlarged gravity--matter BV
cohomology.  Constants are uniform in volume, cutoff and shell number only in
the displayed fixed source, derivative, profile and coupling domains.  No
uniformity in unbounded source order or perturbative order is asserted.

\subsection*{Open}
The smallest load-bearing theorem is now the \emph{two-loop} retained-constraint
BV forest.  It must combine the populated one-loop counterterm insertion of
V64, the logarithmic canonical packet of V65, the scalar sunset and quartic
subdivergences of V71--V72, and the mixed gravity--Higgs shell of V74 inside one
common equation for
\[
 (S_0,S_2)+\tfrac12(S_1,S_1)-\Delta S_1,
\]
with overlapping forests and the shell BV connection transported together.
Only after that can the one-loop self-map above be promoted to a two-loop
filtered Einstein--Higgs BV self-map.  Internal gauge, Yukawa, chiral-line and
arbitrary-order ESM closure remain further gates; no asymptotic-safety,
finite-parameter quantum-gravity UV completion or nonperturbative continuum is
claimed.

\begin{center}\footnotesize
\begin{tabular}{p{.23\textwidth}p{.69\textwidth}}
\toprule
Variable & Scope of V76\\\midrule
Loop order & One-loop nonlinear BV/source transport; action spines from V73--V75 retained at their separately proved orders.\\
Local bank & Pure-gravity arity at most five and weight four, plus the populated two-Higgs metric/source rows; auxiliary density local jets counted once.\\
Soft sector & Lapse/shift, constraint, Lorentz/nonminimal and their soft source/ghost coordinates are retained; no soft inverse wave operator.\\
Cutoff & Completed finite regulator with the inherited lower-scale and nonaliasing margins; interior shell mismatch sums to $O(a\mu)$ and upper cap is treated separately.\\
Volume & Rooted source norms and shell connection bounds have no total-volume multiplier; thermodynamic errors are separate.\\
Shell count & Composed source transports are uniformly bounded for arbitrarily many ultraviolet shell steps on this fixed panel.\\
Full ESM & Not proved. Two-loop BV forests, internal gauge/Yukawa/chiral-line merger and arbitrary-order filtered closure remain open.\\
\bottomrule
\end{tabular}
\end{center}

\section*{Append-only continuation marker: V76}
\addcontentsline{toc}{section}{Append-only continuation marker: V76}
\label{sec:v76-continuation-marker}
\textbf{End of V76.} Preserve Sections 1--584.  The nonlinear shell problem
has been reformulated in the correct transported form: the hard pushforward
conjugates the complete BV operator, while its difference from the chosen
retained local representative is a shell connection.  On the populated
one-loop Einstein--Higgs bank, that connection has a common local restoring
jet and a summable quasilocal complement, yielding a bounded composed source
transport with no assumption that the nonlinear BRST vector field commutes
with the shell projector.  The next theorem is the common two-loop BV forest
on this retained-constraint state.



% =============================================================================
% V77 -- weight-six two-loop BV forest and finite local obstruction
% =============================================================================
\part*{V77: The weight-six two-loop BV forest and the finite local obstruction}
\addcontentsline{toc}{part}{V77: Weight-six two-loop BV forest and finite local obstruction}
\label{sec:v77-start}

V76 closes the nonlinear retained-constraint source transport at one loop, but
its local BV bank is still the weight-four bank of V24--V29.  That bank cannot
be used unchanged in the order-two master equation.  The obstruction is already
visible without evaluating another graph: the antibracket of two weight-four
one-loop functionals has weight six.  V65 supplies an explicit nonzero
self-bracket, so this term cannot be removed by a bookkeeping convention.  The
purpose of V77 is therefore upstream.  We construct the weight-six coupled local
complex, assemble the common proper-subgraph forest from the already chosen
one-loop action, and reduce the remaining two-loop restoration problem to one
finite local range test.  The soft constraint and ghost sectors remain retained
as in V75--V76.

No new value is assigned to an unevaluated pure-gravity two-loop coefficient.
The result is a filtered finite-order reduction theorem, not a claim of a
complete two-loop Einstein--Standard-Model calculation.

\section{Why the two-loop source-complete bank must reach weight six}
\label{sec:v77-weight-six}

Use the coupled ordinary bookkeeping already fixed in V24 and V29.  Give
\(q,\phi\) weight zero, \(c\) weight \(-1\), the metric and scalar antifields
\(\eta=q^*\), \(\upsilon=\phi^*\) weight two, and \(\sigma=c^*\) weight
three.  Ordinary derivatives have weight one and \(z=m_H^2\) has weight two.
The nonminimal canonical pairs retain the V24 weight-one convention.  Thus every
canonical field--antifield pair has total weight two.

For a local monomial let \(n_*\) be the total number of minimal antifields,
let \(g\) be its ghost number, and let \(d_w\) be its derivative--mass weight,
with every power of \(z\) counted twice.  Then
\begin{equation}
 w=d_w+n_*-g.
 \label{eq:v77-weight}
\end{equation}
This is V24's grading with the V29 scalar pair adjoined.

\begin{lemma}[The antibracket raises the one-loop cap from four to six]
\label{lem:v77-bracket-weight}
For homogeneous local functionals \(F,G\), every nonzero canonical contraction
obeys
\begin{equation}
 \boxed{w((F,G))=w(F)+w(G)-2.}
 \label{eq:v77-bracket-weight}
\end{equation}
Consequently
\begin{equation}
 w(F),w(G)\le4\quad\Longrightarrow\quad w((F,G))\le6.
 \label{eq:v77-44-to-6}
\end{equation}
Moreover the classical differential \(s_0=(S_0,\cdot)\) preserves \(w\),
because every term of the ordinary tree master action has weight two.
\end{lemma}
\begin{proof}
A functional derivative with respect to a field of weight \(a\) removes weight
\(a\); the paired antifield derivative removes weight \(2-a\).  Their product
therefore removes exactly two from the sum of the two input weights.  This gives
\eqref{eq:v77-bracket-weight}.  Taking \(G=S_0\) gives preservation of weight.
The argument applies to the metric, scalar, ghost and nonminimal canonical
pairs separately and is unchanged by the tensor and Grassmann exchange
symmetries.
\end{proof}

The finite measured BV Laplacian contains one canonical contraction at
coincidence.  Removing the canonical pair lowers the local weight by two; the
four-dimensional coincidence trace has superficial weight at most four.
Hence, for a one-loop local functional of weight at most four,
\begin{equation}
 \boxed{w(\Delta S_1)\le6.}
 \label{eq:v77-laplacian-cap}
\end{equation}
This is only a superficial local-bank statement.  It does not replace the
actual finite coincidence coefficient or its once-counted density convention.

\begin{proposition}[Weight four is genuinely insufficient at order two]
\label{prop:v77-weight4-insufficient}
Suppose the one-loop source-complete action has a nonzero weight-four component.
Then the order-two master coefficient
\begin{equation}
 \mathcal A_2=(S_0,S_2)+\tfrac12(S_1,S_1)-\Delta S_1
 \label{eq:v77-A2}
\end{equation}
requires a local bank through weight six unless every weight-six contribution is
proved to cancel.  Such a cancellation cannot be assumed in the present lineage:
V65 evaluates a nonzero component of \(\tfrac12(S_1,S_1)\), and cancels it only
after the corresponding second-order canonical completion is retained.
\end{proposition}
\begin{proof}
The first statement follows from Lemma~\ref{lem:v77-bracket-weight} and
\eqref{eq:v77-laplacian-cap}.  V65's explicit scalar-six self-bracket is a
nonzero coefficient of the quadratic term.  Its cancellation by the V65
second-order action is an example of the required weight-six completion, not a
reason to omit the weight-six bank.
\end{proof}

This is the precise upstream repair anticipated in V61: the one-loop
weight-four complex and the two-loop weight-six complex are distinct filtered
objects.

\section{The coupled weight-six local BV complex and common Taylor projection}
\label{sec:v77-complex}

Fix an arity cap \(N\).  Let \(\mathfrak L^g_{N,6}\) be the finite space of
translation-invariant local polynomial functionals of ghost number \(g\), total
arity at most \(N\), and weight \(w\le6\), modulo total derivatives and the
actual tensor/Grassmann symmetries.  The scalar mass variable is counted with
weight two.  This is an EFT/source bank at one prescribed order; no all-order
uniform dimension is asserted.

For a component \(I\) with ghost number \(g(I)\) and \(n_*(I)\) minimal
antifields, define
\begin{equation}
 d_I^{(6)}=6+g(I)-n_*(I),
 \qquad
 (\mathsf P_{N,6}F)_I=\mathsf T_{d_I^{(6)}}F_I,
 \label{eq:v77-projector}
\end{equation}
with \(\mathsf T_d=0\) when \(d<0\).  The Taylor operator acts jointly on
external momenta and \(z\), counting \(z\) twice.

\begin{theorem}[Degree-closed coupled BV complex at weight six]
\label{thm:v77-complex}
Let \(s_\infty=(\mathbb S_\infty,\cdot)\) be the ordinary coupled
Einstein--Higgs classical BV differential in the same coframe chart as V24 and
V29, including the nonminimal doublet after the fixed gauge-fermion
transformation.  Then
\begin{equation}
 \boxed{
 d_{N,6}=\operatorname{pr}_{A\le N}s_\infty,
 \qquad d_{N,6}^2=0,
 \qquad
 d_{N,6}:\mathfrak L^g_{N,6}\to\mathfrak L^{g+1}_{N,6}.}
 \label{eq:v77-complex}
\end{equation}
The common Taylor projection is a chain retraction,
\begin{equation}
 \boxed{\mathsf P_{N,6}d_{N,6}=d_{N,6}\mathsf P_{N,6}.}
 \label{eq:v77-chain-projection}
\end{equation}
In the scaled coefficient-jet norm its retraction norm is at most one, while
inclusion of the local polynomial in the corresponding smooth norm is bounded
by \(2^{N+6}\).  The constants are independent of the mesh, periods and number
of shell steps at fixed \(N\).
\end{theorem}
\begin{proof}
The proof is the V24/V29 weight argument with the cap four replaced by six.
Every term of the ordinary coupled tree master action has weight two and arity
at least two.  By Lemma~\ref{lem:v77-bracket-weight}, bracketing with that action
preserves \(w\) and cannot lower arity.  The ordinary classical master identity
and Jacobi therefore give square zero before quotienting, while the arity ideal
and total derivatives are invariant.

For the projection, every contribution to the differential is a homogeneous
polynomial in the external momenta times an input coefficient whose arguments
are linear subset sums.  Homogeneous Taylor extraction commutes with those
linear substitutions, and the weight-preserving multiplication shifts the
allowed Taylor degree by exactly the degree of the polynomial.  Hence
\eqref{eq:v77-chain-projection}.  The coefficient-norm estimate is the same
binomial bound as V24, with maximal Taylor degree at most \(N+6\).
\end{proof}

\begin{remark}[Action and source filtrations remain distinct]
\label{rem:v77-two-filtrations}
V73's scalar EFT operators of canonical dimension six all lie inside this
weight-six envelope, but the two gradings are not identified.  The BV weight
tracks derivatives, ghosts and antifields so that the local complex closes;
the EFT ancestry tags retain powers of \(\kappa_g\), \(g_H\), and the physical
Wilson normalization.  Both labels are kept in the iterative state.
\end{remark}

\section{One common one-loop action owns every populated proper subtraction}
\label{sec:v77-forest}

Let \(S_{1,j}^{\rm loc}\) denote the actual one-loop local functional already
chosen at shell \(j\), including its gravity, scalar, mixed, source and
once-counted density components.  Its sign is the physical counterterm sign;
no graphwise refit is made below.  Let
\begin{equation}
 \mathcal I_j[S_{1,j}^{\rm loc}]
 =\frac12\operatorname{STr}\!\left(
       \mathbb G_jD_\Phi^2S_{1,j}^{\rm loc}\right)
 \label{eq:v77-I}
\end{equation}
be V64's one-counterterm hard insertion, retaining every labelled external
source argument.

Denote by \(\mathfrak G_{\rm pop}^{(2)}\) the two-loop graph family already
populated in this lineage: the V62 three-propagator scalar--metric core, the V63
shared-root contact core, and the V71--V72 two-quartic scalar family including
the scalar sunset.  The family does not include the still unevaluated complete
pure-gravity two-loop skeletons.

\begin{theorem}[Common forest ownership on the populated two-loop family]
\label{thm:v77-forest-ownership}
For every graph in \(\mathfrak G_{\rm pop}^{(2)}\), the sum of its proper
one-loop local subtraction insertions, with the finite parts fixed by the
previous normalization, is exactly the restriction of
\(\mathcal I_j[S_{1,j}^{\rm loc}]\) to its labelled cuts.  In particular:
\begin{enumerate}[label=\textup{(\roman*)},leftmargin=2.7em]
\item the two contact subtractions and the three overlapping scalar--metric
      subtractions have the V64 cut multiplicities, including the factor two
      from the two scalar cuts of the parallel-edge graph;
\item the three proper one-loop four-point subgraphs of the scalar sunset all
      use the single V71 quartic counterterm and its fixed finite part;
\item V63's shared degree-six double-local product remains an overall-assignment
      conversion and is not a new bare graph;
\item V65's second-order canonical packet is not counted as a proper forest
      subtraction; it is a separate order-two action term.
\end{enumerate}
Thus no independent graphwise finite counterterm remains on this populated
family once the common \(S_{1,j}^{\rm loc}\) is fixed.
\end{theorem}
\begin{proof}
Items (i) and (iii) are the exact opening/closing identities of V64 together
with V63's shared-root conversion.  For the scalar family, V71 supplies one
local quartic action coefficient and V72 identifies the three proper one-loop
four-point channels of the sunset.  Opening any one of those labelled channels
therefore restricts the same quartic Hessian; closing the complementary scalar
line is precisely the corresponding term of \eqref{eq:v77-I}.  Linearity
preserves the already chosen finite part.  The V65 term arises from a finite
canonical continuation and is already quadratic in the order-one generator;
it is not a subgraph counterterm.  These are finite contraction identities
before any cutoff estimate is taken.
\end{proof}

Define the prepared known two-loop action on this family by
\begin{equation}
 S_{2,j}^{\rm known}
 =\Gamma_{2,j}^{\rm bare,pop}
  +\mathcal I_j[S_{1,j}^{\rm loc}]_{\rm pop}
  +S_{2,j}^{\rm can},
 \label{eq:v77-known-S2}
\end{equation}
where \(S_{2,j}^{\rm can}\) contains the already specified V65 canonical
completion and any previously fixed order-two local normalization.  Equation
\eqref{eq:v77-known-S2} is a bookkeeping identity for the populated family; it
does not assign the missing pure-gravity primitive graphs a value.

\section{The order-two master row becomes one finite weight-six range problem}
\label{sec:v77-range}

Work first in the ordinary local reference where \(d_{N,6}\) is a genuine
differential.  Write
\begin{equation}
 S_2=S_2^{\rm known}+C_{2}^{\rm loc},
 \label{eq:v77-S2split}
\end{equation}
where \(C_2^{\rm loc}\in\mathfrak L^0_{N,6}\) is the still adjustable common
order-two local functional.  Put
\begin{equation}
 B_2
 =\mathsf P_{N,6}\left
 \{
 (S_0,S_2^{\rm known})
 +\tfrac12(S_1,S_1)-\Delta S_1
 \right\}.
 \label{eq:v77-B2}
\end{equation}
The actual lower-reference and finite defects are restored below; they are not
set to zero in this definition.

\begin{theorem}[Finite local range criterion for the two-loop master row]
\label{thm:v77-range}
On the ordinary local bank the projected order-two master equation is exactly
\begin{equation}
 \boxed{
 \mathsf P_{N,6}\mathcal A_2
 =d_{N,6}C_2^{\rm loc}+B_2.}
 \label{eq:v77-range-equation}
\end{equation}
Let \(P_{\rm coker}\) denote orthogonal projection onto
\((\operatorname{Ran}d_{N,6})^\perp\) in the chosen finite coefficient norm.
Then a local order-two restoring functional exists if and only if
\begin{equation}
 \boxed{P_{\rm coker}B_2=0.}
 \label{eq:v77-coker-test}
\end{equation}
When this holds, the minimum-norm choice is
\begin{equation}
 \boxed{C_{2,\min}^{\rm loc}=-d_{N,6}^{\dagger}B_2,}
 \label{eq:v77-min-primitive}
\end{equation}
and, if \(\sigma_{N,6}>0\) is the least nonzero singular value of
\(d_{N,6}\),
\begin{equation}
 \|C_{2,\min}^{\rm loc}\|
 \le\sigma_{N,6}^{-1}\|B_2\|.
 \label{eq:v77-primitive-bound}
\end{equation}
Thus the formerly qualitative ``critical two-loop primitive'' problem is a
finite, explicitly populated range test once the remaining primitive graph
coefficients are supplied.
\end{theorem}
\begin{proof}
Substitute \eqref{eq:v77-S2split} into \eqref{eq:v77-A2}, use
\((S_0,C_2^{\rm loc})=d_{N,6}C_2^{\rm loc}\), and commute the common Taylor
projection through the ordinary differential by
Theorem~\ref{thm:v77-complex}.  This proves
\eqref{eq:v77-range-equation}.  The remaining assertions are the finite
Moore--Penrose range criterion and its singular-value bound.
\end{proof}

The weight decomposition sharpens this test.  If
\(S_1=\sum_{w\le4}S_1^{[w]}\), then the critical weight-six quadratic term is
\begin{equation}
 \boxed{
 \left[\tfrac12(S_1,S_1)\right]_{[6]}
 =\tfrac12\bigl(S_1^{[4]},S_1^{[4]}\bigr)_{[6]}.}
 \label{eq:v77-critical-bracket}
\end{equation}
Lower-weight pairs cannot reach six.  This is the only source of weight six
coming from the quadratic master term.

\begin{proposition}[The populated canonical packet removes its own quadratic obstruction]
\label{prop:v77-canonical-cancel}
For V65's canonical logarithmic packet, the contribution of its specified
\(S_1^F=b\mathcal D_FS_0\) and its specified second-order continuation to the
relative order-two defect is
\begin{equation}
 \boxed{
 \delta\mathfrak A_2
 =b\mathcal D_F\mathfrak A_1
  +\frac{b^2}{2}\mathcal D_F^2\mathfrak A_0.}
 \label{eq:v77-canonical-defect}
\end{equation}
Consequently, on the ordinary branch where the lower-order defects vanish, the
nonzero V65 self-bracket is cancelled inside \(B_2\) by the same packet's
second-order action and measure terms.  It must not be presented again as a new
cohomological obstruction.
\end{proposition}
\begin{proof}
This is V65's exact order-two defect-transport formula, restricted to the
packet.  Its explicit nonzero self-bracket and the corresponding second scalar
action coefficient are two terms in the same identity.
\end{proof}

The pure scalar action counterterms of V71--V72 are antifield independent; their
mutual antibrackets therefore vanish.  Their role in \(B_2\) is through the
common action, its diffeomorphism/source completion, and the proper forests of
Theorem~\ref{thm:v77-forest-ownership}, not through an independently fitted
scalar self-bracket.

\section{Two-loop shell-connection curvature: the quadratic one-loop term is summable}
\label{sec:v77-connection}

V76's exact shell connection provides a second, operator-level view of the same
quadratic structure.  Expand
\begin{equation}
 \mathcal A_j^{\rm BV}
 =\hbar\mathcal A_{1,j}+\hbar^2\mathcal A_{2,j}+O(\hbar^3),
 \qquad
 \widehat D_{j+1}^{\,2}-D_{<,j+1}^{\,2}
 =\hbar\mathcal K_{1,j}+\hbar^2\mathcal K_{2,j}+O(\hbar^3).
 \label{eq:v77-connection-expansion}
\end{equation}
Here the square is the graded operator square in the same density/bracket
convention as V76.

\begin{theorem}[Exact coefficient equations for the two-loop shell connection]
\label{thm:v77-connection-equations}
The V76 curvature identity gives
\begin{align}
 \boxed{D_<\mathcal A_{1,j}+\mathcal A_{1,j}D_<
       =\mathcal K_{1,j},}
 \label{eq:v77-connection-one}\\
 \boxed{D_<\mathcal A_{2,j}+\mathcal A_{2,j}D_<
       +\mathcal A_{1,j}^2
       =\mathcal K_{2,j}.}
 \label{eq:v77-connection-two}
\end{align}
After the V76 local one-loop primitive is removed, write
\(\mathcal A_{1,j}=\mathcal E_{1,j}\) for the complementary part.  Then
\begin{equation}
 \boxed{
 \|\mathcal E_{1,j}^2\|_{*\to *}
 \le C\left(\frac{\mu}{\Lambda_j}\right)^2.}
 \label{eq:v77-square-bound}
\end{equation}
Thus the quadratic one-loop connection is absolutely summable over ultraviolet
shells.  It cannot be the obstruction to two-loop composed source transport.
The remaining analytic input is the local/range decomposition and quasilocal
bound for the genuinely new coefficient \(\mathcal K_{2,j}\).
\end{theorem}
\begin{proof}
Expand V76's exact curvature equation
\(D_<\mathcal A+\mathcal AD_<+\mathcal A^2
 =\widehat D^2-D_<^2\) in powers of \(\hbar\).  This gives
\eqref{eq:v77-connection-one}--\eqref{eq:v77-connection-two}.  V76 proves
\(\|\mathcal E_{1,j}\|_{*\to *}\le C\mu/\Lambda_j\); submultiplicativity
gives \eqref{eq:v77-square-bound}.  The geometric series
\(\sum_j(\mu/\Lambda_j)^2\) converges.
\end{proof}

\begin{corollary}[Conditional two-loop composed source transport]
\label{cor:v77-two-loop-transport}
Suppose the new two-loop connection coefficient admits the common split
\begin{equation}
 \mathcal K_{2,j}
 =\mathsf P_{N,6}\mathcal K_{2,j}
  +(1-\mathsf P_{N,6})\mathcal K_{2,j},
 \end{equation}
its local part passes the finite range test \eqref{eq:v77-coker-test}, and after
that primitive is installed
\begin{equation}
 \|(1-\mathsf P_{N,6})\mathcal K_{2,j}\|_{*\to *}
 \le C_2\frac{\mu}{\Lambda_j}.
 \label{eq:v77-K2-bound}
\end{equation}
Then the renormalized adjacent source map through order \(\hbar^2\) has
\begin{equation}
 \mathsf U_{j+1,j}
 =I+\hbar E_{1,j}+\hbar^2E_{2,j}+O(\hbar^3),
 \qquad
 \|E_{2,j}\|_{*\to *}\le C'_2\frac{\mu}{\Lambda_j},
 \label{eq:v77-adjacent2}
\end{equation}
and its composed products are uniformly bounded in the number of ultraviolet
shells on every fixed perturbative \(\hbar\)-domain.
\end{corollary}
\begin{proof}
Equation~\eqref{eq:v77-connection-two} shows that the only nonlinear forcing
from the already controlled one-loop connection is
\(-\mathcal E_{1,j}^2\), which is bounded by
\eqref{eq:v77-square-bound}.  The new complementary forcing is
\eqref{eq:v77-K2-bound}.  Their sum is bounded by
\(C\mu/\Lambda_j\), since \(\mu/\Lambda_j<1\).  The product estimate is the
same logarithmic/product argument as V76, applied coefficientwise through
order two.
\end{proof}

This corollary isolates the open theorem sharply: one does not need a new
mechanism for composing source transports.  One needs the genuinely new
\(\mathcal K_{2,j}\) local range and quasilocal estimate.

\section{Filtered action closure on the populated two-loop spine}
\label{sec:v77-filtered}

On the antifield-free action restriction the populated weight-six local
operators are exactly of the kind already encountered in V61--V74: six-derivative
pure/mixed local terms and the scalar dimension-six bank of V73.  Their local
Wilson coordinates have inverse ultraviolet growth exponent two.  After the
complete weight-six Taylor assignment, the parity-even action complement starts
two derivative weights higher and has decay exponent four.  Thus the populated
action spine retains
\begin{equation}
 \boxed{s_2=2\le\sigma_2=3<t_2=4.}
 \label{eq:v77-action-window}
\end{equation}

\begin{theorem}[Two-boundary closure for the populated two-loop action family]
\label{thm:v77-populated-trajectory}
Take the direct sum of the V73 scalar spine, the V74 mixed retained-scalar row,
the V62--V63 scalar-containing two-loop metric cores, V72's sunset, and V65's
specified canonical order-two completion.  Use the common V71/V64 one-loop
counterterms in every proper subgraph as in
Theorem~\ref{thm:v77-forest-ownership}.  At their already proved coupling,
source and arity orders, the resulting antifield-free action coefficients form
a finite triangular filtered system satisfying the V6 window
\eqref{eq:v77-action-window}.  Hence fixed physical lower-endpoint local data
and polynomially bounded ultraviolet complementary data determine a unique
coefficientwise two-boundary trajectory on this populated family, with
\begin{equation}
 \boxed{
 \|x_{j}^{(J)}-x_j^{(\infty)}\|
 \le C b^{3j}(1+J)^q\theta^{J-j},
 \qquad b^{-1}<\theta<1.}
 \label{eq:v77-populated-rate}
\end{equation}
The theorem does not include unevaluated pure-gravity two-loop skeletons or the
unresolved local BV range condition \eqref{eq:v77-coker-test}.
\end{theorem}
\begin{proof}
V73 and V74 prove the local/complement scaling \(b^{-2}\) versus \(b^{-4}\)
for the scalar and mixed dimension-six rows.  V61--V63 use the same degree-six
overall assignment and V62 supplies the separated-scale forest estimate.  V72
supplies the scalar-sunset forest and its rooted cutoff bound.  V65's specified
canonical completion is local weight six and does not introduce a new
complementary ultraviolet boundary.  Theorem~\ref{thm:v77-forest-ownership}
identifies the proper counterterms with the same one-loop action, so no extra
same-degree graph coordinate is introduced.  The finite direct sum therefore
has the V6 triangular form with \(s_2=2\), \(t_2=4\), and the already used
choice \(\sigma_2=3\).  Apply the V6 filtered two-boundary theorem.
\end{proof}

The distinction is important.  The antifield-free populated action family now
has a two-loop coefficient trajectory.  A source-complete two-loop
Einstein--Higgs theorem additionally requires the finite weight-six master range
test and the genuinely new shell-connection coefficient of the preceding
section.

\section{Anti-circularity audit and status ledger for V77}
\label{sec:v77-ledgers}

\subsection*{Proved}
The one-loop weight-four BV bank is insufficient for the order-two master
equation; the native grading forces a weight-six bank because an antibracket
removes weight two, and V65 provides an explicit nonzero quadratic witness.  The
ordinary coupled Einstein--Higgs BV complex and the common Taylor retraction
extend from weight four to weight six with the same mesh-, volume-, and
shell-independent finite coefficient bounds at fixed arity.  On every populated
two-loop graph family, the proper one-loop subtractions are restrictions of one
common previously chosen \(S_1\); there are no independent graphwise finite
counterterms.  The projected order-two master equation is reduced to the finite
range test \eqref{eq:v77-coker-test}, with an explicit Moore--Penrose primitive
and singular-value bound when the test passes.  V65's canonical packet cancels
its own populated self-bracket relative to the lower-order defects.  The exact
V76 shell-connection curvature splits at orders one and two as
\eqref{eq:v77-connection-one}--\eqref{eq:v77-connection-two}; the quadratic
one-loop complementary term is \(O((\mu/\Lambda_j)^2)\) and is summable.  The
antifield-free populated two-loop action spine retains the strict filtered
window \(2\le3<4\) and therefore has a two-boundary coefficient trajectory.

\subsection*{Inherited}
V24 and V29 supply the weight/arity bookkeeping and the coupled ordinary
classical BV differential.  V59 supplies exact normalized BV pushforward with
finite defects retained.  V62--V64 supply the overlapping scalar--metric forests,
shared-root contact subtraction, common one-loop counterterm insertion and cut
multiplicities.  V65 supplies the exact quantum-canonical second-order packet and
its nonzero self-bracket.  V71--V72 supply the common quartic normalization,
scalar sunset and its full separated-scale forest estimate.  V73--V75 supply the
filtered scalar/mixed action spine and retained soft constraints.  V76 supplies
the exact shell BV connection and the one-loop composed source-transport bound.

\subsection*{Literature input}
No new external two-loop gravity coefficient or cohomology classification is
used in the proved statements above.  The finite range test is deliberately
left as a finite native calculation.  Established two-loop gravity EFT results
remain comparison information only; they are not substituted for the missing
native pure-gravity shell coefficients.

\subsection*{Conditional}
The two-loop composed source-transport conclusion requires the genuinely new
coefficient \(\mathcal K_{2,j}\) to pass the weight-six local range test and to
obey the quasilocal bound \eqref{eq:v77-K2-bound}.  The populated action
trajectory excludes pure-gravity two-loop primitive skeletons not yet evaluated
in this regulator.  All estimates remain at fixed perturbative order, source
order, arity, profile class and coupling domain.  No uniformity in an unbounded
EFT tower is asserted.

\subsection*{Open}
The smallest load-bearing problem is now finite and explicit: populate the
remaining native pure-gravity and retained-constraint contributions to \(B_2\)
and \(\mathcal K_{2,j}\), then evaluate the cokernel coordinate
\(P_{\rm coker}B_2\) in the weight-six coupled complex.  If it vanishes, install
\(-d_{N,6}^{\dagger}B_2\) and prove \eqref{eq:v77-K2-bound}; this would promote
V76 to a genuine two-loop retained-constraint Einstein--Higgs BV self-map.  If
it does not vanish, the nonzero cokernel coordinate is the precise local
obstruction and the programme must enlarge the symmetry/source content rather
than fit graphwise counterterms.  Internal gauge, Yukawa, chiral-line and
arbitrary-order ESM closure remain downstream gates.  No asymptotic-safety or
nonperturbative quantum-gravity conclusion is made.

\begin{center}\footnotesize
\begin{tabular}{p{.23\textwidth}p{.69\textwidth}}
\toprule
Variable & Scope of V77\\\midrule
Loop order & Complete algebraic/order-two BV assembly on the populated family; missing native pure-gravity primitive coefficients remain open.\\
Local BV bank & Coupled Einstein--Higgs ghost/source complex at fixed arity and weight at most six; common Taylor projection is a chain retraction.\\
Action forest & V62--V65 and V71--V74 populated two-loop/action pieces with one common one-loop normalization.\\
Soft sector & V75 retained lapse/shift, Lorentz/nonminimal and ghost/source coordinates; no inverse soft constraint operator.\\
Action UV window & Populated weight-six local/complement coordinates satisfy $s=2$, $\sigma=3$, $t=4$.\\
Source transport & One-loop composition proved; two-loop composition reduced to the new $\mathcal K_2$ local/quasilocal estimate, with the quadratic $\mathcal A_1^2$ term already summable.\\
Cutoff/volume & Existing graphwise cutoff and rooted estimates are inherited at their stated domains; no new bound is claimed for the missing pure-gravity two-loop skeletons.\\
Full ESM & Not proved. Internal gauge/Yukawa/chiral-line merger and arbitrary-order filtered closure remain open.\\
\bottomrule
\end{tabular}
\end{center}

\section*{Append-only continuation marker: V77}
\addcontentsline{toc}{section}{Append-only continuation marker: V77}
\label{sec:v77-continuation-marker}
\textbf{End of V77.} Preserve Sections 1--591.  The two-loop problem is no
longer posed on the one-loop weight-four bank.  The coupled local complex has
been enlarged minimally to weight six, every populated proper subtraction is
owned by one common $S_1$, and the full local order-two master problem is the
finite range test \eqref{eq:v77-coker-test}.  The quadratic one-loop shell
connection is already summable.  The next theorem must populate the remaining
native pure-gravity/retained-constraint coefficient, evaluate its cokernel
coordinate, and prove the new two-loop quasilocal connection bound.



% =============================================================================
% V78 -- weight-six pure-gravity range theorem and Wilson/anomaly split
% =============================================================================
\part*{V78: The weight-six pure-gravity range theorem and the Wilson/anomaly split}
\addcontentsline{toc}{part}{V78: Weight-six pure-gravity range theorem and Wilson/anomaly split}
\label{sec:v78-start}

V77 reduces the order-two source-complete master equation to a finite
weight-six range test, but leaves the remaining pure-gravity and
retained-constraint coefficients unevaluated.  Before computing those
coefficients graph by graph, there is a sharper upstream question: can a
closed pure-gravity ghost-one coefficient at this weight carry a genuine local
anomaly, or do the missing graphs determine only a restoring primitive and the
independent ghost-zero EFT Wilson data?  The two possibilities must not be
conflated.

This installment answers that question on the bosonic pure-gravity restriction
of the retained-constraint state.  First we extend V44's free ghost-one
arity gap from weight four to weight six.  This permits a closed seven-jet to
be extended to a full formal local cocycle.  The same four-dimensional
pure-Einstein anomaly classification already used in V25 then makes the full
ghost-one cocycle exact modulo a translation-invariant divergence.  Hence the
pure-gravity component of V77's cokernel coordinate vanishes without knowing
the numerical native two-loop graph array.  This does \emph{not} make the
weight-six action bank exact: its ghost-zero cohomology contains genuine EFT
Wilson directions.  In particular the two-loop curvature-cubic direction must
be retained rather than absorbed into the BV restoring primitive.

We finally prove the shell-local analytic statement needed on this sector.  A
complete two-loop hard forest, after all proper one-loop subtractions and the
weight-six local Taylor assignment, has a quasilocal complement of size
$O(\mu/\Lambda_j)$ in the V76 source norm.  Soft constraints remain retained,
so no inverse soft wave operator enters this estimate.  Combined with V77's
summable quadratic one-loop connection, this closes the pure-gravity plus
retained-nonminimal part of the two-loop source transport.  Matter-dependent
weight-six anomaly/range questions are kept separate.

No numerical Goroff--Sagnotti coefficient is imported into the active native
normalization.  The established two-loop result is used only to exhibit a
nontrivial physical Wilson direction in the weight-six ghost-zero quotient.

\section{The free ghost-one gap moves from arity six to arity eight}
\label{sec:v78-gap}

Use V44's ordinary free spin-two complex and the same weights
\[
 w(q)=0,\qquad w(c)=-1,\qquad w(\eta)=2,\qquad
 w(\sigma)=3,\qquad w(\partial)=1.
\]
The coefficient weight is independent of horizontal form degree.  For a
four-form coefficient of weight six, a descent bottom of form degree $p$ has
coefficient weight $p+2$ and ghost number $5-p$.

\begin{theorem}[Free ghost-one weight-six arity gap]
\label{thm:v78-gap}
For the translation-invariant ordinary free spin-two complex in four
spacetime dimensions,
\begin{equation}
 \boxed{H^{1,4}_{A=r,w=6}(s_0\mid d_x)=0\qquad(r\ge8).}
 \label{eq:v78-gap}
\end{equation}
Equivalently, every local integrated ghost-one density of arity $r\ge8$ and
coefficient weight six that is $s_0$-closed modulo a translation-invariant
local divergence has a translation-invariant ghost-zero primitive of the same
arity and weight.  No inverse momentum, mass, or wave operator is used.
\end{theorem}
\begin{proof}
Repeat V44's invariant antifield reduction, which is homogeneous in arity and
weight.  A possible highest antifield-two representative contains one
undifferentiated ghost antifield, three surviving ghost factors, and, at arity
$r$, at least $r-4$ linearized curvatures.  Its minimum coefficient weight is
\[
 3-3+2(r-4)=2r-8.
\]
At weight six this requires $r\le7$.  Thus no antifield-two representative
survives for $r\ge8$.

An antifield-one representative contains one metric antifield, two surviving
ghost factors and at least $r-3$ curvatures.  Its minimum weight is
\[
 2-2+2(r-3)=2r-6>6\qquad(r\ge7),
\]
so it is absent in particular for $r\ge8$.

The remaining top form is antifield-free.  Descend to a nontrivial bottom
$a_p$, $0\le p\le4$.  Its ghost number is $5-p$, its required coefficient
weight is $p+2$, and the number of positive-metric-degree curvature factors is
at least $m=r-5+p$.  Even assigning every surviving ghost its most negative
weight, an invariant representative has
\begin{equation}
 w(a_p)\ge 2(r-5+p)-(5-p)=2r-15+3p.
 \label{eq:v78-bottom-bound}
\end{equation}
For $r\ge8$ and $p\ge1$ this is strictly larger than $p+2$.  At $p=0$ the
metric degree is $r-5\ge3$, so V44's translation-invariant zero-form lemma
removes the bottom.  Shortening the descent either reaches a forbidden higher
bottom or trivializes the original top form.  Reversing the local reductions
gives the stated primitive.  All divisions are the polynomial homogeneity and
Koszul sections already used in V44; none is a physical propagator.
\end{proof}

The threshold eight is sufficient for the order-two local bank.  It is not
claimed to be optimal.  The two additional allowed arities relative to V44 are
exactly the room created by the two additional units of local weight.

\section{Closed weight-six seven-jets extend, and the pure-gravity anomaly class vanishes}
\label{sec:v78-extension}

Write the ordinary interacting pure-gravity differential as
$s_g=s_0+s_1+s_2+\cdots$, with $s_j$ raising ordinary field arity by $j$.

\begin{theorem}[Formal extension of a closed weight-six seven-jet]
\label{thm:v78-extension}
Let
\[
 a_{\le N}=\sum_{r=2}^{N}a_r,\qquad N\ge7,
 \qquad \operatorname{gh}(a_r)=1,\quad w(a_r)=6,
\]
and suppose
\[
 \operatorname{pr}_{A\le N}s_ga_{\le N}=0\quad\bmod d_x.
\]
Then there is a formal finite-derivative local germ
$a_\infty=\sum_{r\ge2}a_r$ such that
\begin{equation}
 \boxed{s_ga_\infty=0\quad\bmod d_x,
 \qquad J_{\le N}a_\infty=a_{\le N}.}
 \label{eq:v78-extension}
\end{equation}
The extension changes none of the supplied coefficients through arity $N$.
\end{theorem}
\begin{proof}
Assume the coefficients have been chosen through arity $r-1$.  The next
obstruction is
\[
 O_r=\sum_{j=1}^{r-2}s_ja_{r-j},\qquad s_0a_r=-O_r.
\]
Nilpotence of $s_g$ and the lower equations give $s_0O_r=0$ modulo a
divergence.  For $r\ge N+1\ge8$, Theorem~\ref{thm:v78-gap} supplies a local
$a_r$.  Induction gives the formal extension.  Fixed weight bounds derivative
and antifield degree at every arity, exactly as in V44.
\end{proof}

The remaining step is all-arity ghost-one cohomology rather than finite-jet
linear algebra.  We use the same established four-dimensional pure-Einstein
local BRST classification already invoked in V25 and V44
\cite{V20BBH1994,V25BBHEinstein}.  In particular there is no local pure
gravitational anomaly in four dimensions in this ordinary bosonic complex.
The translation-invariant restriction requires one additional weight check.

\begin{lemma}[No translation-sensitive descent exception at weight six]
\label{lem:v78-no-exception}
In the pure-gravity descent of a ghost-one four-form of coefficient weight six,
the constant one-form exception that could obstruct a translation-invariant
primitive is absent.
\end{lemma}
\begin{proof}
Give $dx^\mu$ weight $-1$, as in V25.  The top four-form has total weight
$6-4=2$.  The only constant one-form exception in the imported descent has the
form $dx\,\widehat\Theta L(\mathcal T)$, with four undifferentiated ghosts.
The one-form and four ghosts contribute total weight $-5$, so preservation of
total weight two would require a pure-gravity Lorentz scalar $L$ of coefficient
weight seven.  A scalar built from $r$ curvatures and $m$ covariant derivatives
has weight $2r+m$ and $4r+m$ tensor indices.  Weight seven forces
$m=7-2r$, hence $4r+m=7+2r$ is odd and cannot be completely contracted with
the metric and the four-index alternating tensor.  Thus no such $L$ exists.
\end{proof}

\begin{theorem}[Pure-gravity weight-six ghost-one range theorem]
\label{thm:v78-pg-range}
Let $d_{N,6}^{\rm pg}$ be the ordinary interacting pure-gravity BV differential
on translation-invariant integrated local coefficients of arity at most
$N\ge7$ and weight at most six, with the inherited local Lorentz and
nonminimal doublets contracted as in V25.  On the homogeneous ghost-one,
weight-six subspace,
\begin{equation}
 \boxed{
 \ker\!\left(d_{N,6}^{\rm pg}:\mathfrak L^1_{N,6}\to
                                 \mathfrak L^2_{N,6}\right)
 =
 \operatorname{im}\!\left(d_{N,6}^{\rm pg}:\mathfrak L^0_{N,6}\to
                                     \mathfrak L^1_{N,6}\right).}
 \label{eq:v78-pg-range}
\end{equation}
Consequently the pure-gravity restriction of every closed V77 two-loop local
breaking has zero cokernel coordinate.
\end{theorem}
\begin{proof}
Extend the closed finite jet by Theorem~\ref{thm:v78-extension}.  The established
all-arity pure-Einstein anomaly classification makes the resulting formal
ghost-one cocycle $s_g$-exact modulo a local divergence.  Lemma~\ref{lem:v78-no-exception}
keeps the primitive translation invariant at the present weight.  Projecting
back to arity at most $N$ gives the displayed finite range identity because
$s_g$ and $d_x$ do not lower arity.  Algebraic nonminimal pairs and the retained
constraint gauge-fixing doublets do not change local BRST cohomology, by the
same contraction used in V25.  Hence the result applies to the pure-gravity
restriction of the V75--V77 retained-constraint bank.
\end{proof}

This is a statement about the ghost-one master range.  It is deliberately not
an assertion that every ghost-zero weight-six action coefficient is redundant.

\section{The order-two consistency identity makes the native pure-gravity row closed}
\label{sec:v78-consistency}

Let
\[
 \mathfrak A(S)=\frac12(S,S)-\hbar\Delta S
 =\mathfrak A_0+\hbar\mathfrak A_1+\hbar^2\mathfrak A_2+\cdots .
\]
The BV identity gives
\begin{equation}
 \boxed{(S,\mathfrak A)-\hbar\Delta\mathfrak A=0.}
 \label{eq:v78-quantum-consistency}
\end{equation}
Its order-two coefficient is
\begin{equation}
 (S_0,\mathfrak A_2)+(S_1,\mathfrak A_1)
 +(S_2,\mathfrak A_0)-\Delta\mathfrak A_1=0.
 \label{eq:v78-order2-consistency}
\end{equation}

\begin{proposition}[Closure of the restored ordinary two-loop row]
\label{prop:v78-B2-closed}
On the ordinary branch on which the tree and one-loop local master defects have
been restored,
\[
 \mathfrak A_0=0,\qquad \mathfrak A_1=0,
\]
the weight-six projected order-two breaking satisfies
\begin{equation}
 \boxed{d_{N,6}B_2=0.}
 \label{eq:v78-B2-closed}
\end{equation}
In the finite filtered regulator the corresponding formula retains the
transported lower-order defects on the right-hand side; they are not set to
zero by this proposition.
\end{proposition}
\begin{proof}
Equation~\eqref{eq:v78-order2-consistency} reduces to
$(S_0,\mathfrak A_2)=0$.  The common Taylor map of V77 is a chain map, so its
weight-six local projection is $d_{N,6}$-closed.  The finite-regulator caveat
is exactly V59's defect-retention convention.
\end{proof}

Combining Proposition~\ref{prop:v78-B2-closed} with
Theorem~\ref{thm:v78-pg-range} proves
\begin{equation}
 \boxed{P_{\rm coker}^{\rm pg}B_2=0.}
 \label{eq:v78-pg-coker-zero}
\end{equation}
Thus the unevaluated native pure-gravity two-loop graphs can change the
coefficients of the restoring primitive, but they cannot create a new local
pure-gravity anomaly in this sector.

\section{Finite pure-gravity restoring section and the separate Wilson quotient}
\label{sec:v78-hodge}

Choose any fixed rational coefficient bases for the pure-gravity weight-six
spaces at arity at most $N\ge7$.  Let
\[
 D_0:\mathfrak L^0_{N,6}\to\mathfrak L^1_{N,6},\qquad
 D_1:\mathfrak L^1_{N,6}\to\mathfrak L^2_{N,6}
\]
be the actual matrices of $d_{N,6}^{\rm pg}$ after quotienting total
divergences.  Then $D_1D_0=0$.

\begin{theorem}[Finite minimum-norm pure-gravity two-loop restoration]
\label{thm:v78-finite-section}
On the closed ghost-one subspace the matrix
\[
 L_1=D_0D_0^T+D_1^TD_1
\]
is positive definite.  The operator
\begin{equation}
 \boxed{H_1=D_0^TL_1^{-1}}
 \label{eq:v78-H1}
\end{equation}
is a right inverse of $D_0$ on $\ker D_1$.  Hence the pure-gravity local
restoring coefficient can be chosen canonically as
\begin{equation}
 \boxed{C_{2,\rm rest}^{\rm pg}=-H_1B_2^{\rm pg},\qquad
 D_0C_{2,\rm rest}^{\rm pg}=-B_2^{\rm pg}.}
 \label{eq:v78-restorer}
\end{equation}
At fixed bank its coefficient norm obeys
\[
 \|C_{2,\rm rest}^{\rm pg}\|
 \le \sigma_{\rm pg,6}^{-1}\|B_2^{\rm pg}\|,
\]
where $\sigma_{\rm pg,6}$ is the least nonzero singular value of $D_0$ on the
orthogonal complement of its kernel.  These constants are independent of mesh,
periods and shell number at fixed arity and local weight.
\end{theorem}
\begin{proof}
Theorem~\ref{thm:v78-pg-range} gives
$\ker D_1=\operatorname{im}D_0$.  On $\mathfrak L^1$ the two positive summands
of $L_1$ have orthogonal ranges because $D_1D_0=0$.  Thus $L_1$ has no kernel
and the standard finite Hodge identity gives
$D_0D_0^TL_1^{-1}=I$ on $\ker D_1$.  The singular-value estimate is immediate.
\end{proof}

The same local bank has a different ghost-zero question.  Put
\begin{equation}
 \mathcal H^0_{N,6}
 :=\ker D_0/\operatorname{im}D_{-1},
 \label{eq:v78-H0}
\end{equation}
where $D_{-1}$ is the ghost-minus-one field-redefinition differential.  Choose
the orthogonal harmonic realization
\[
 \mathbf H^0_{N,6}
 =\ker(D_{-1}D_{-1}^T+D_0^TD_0).
\]

\begin{proposition}[Restoration and physical Wilson data are different coordinates]
\label{prop:v78-Wilson-split}
Every closed local ghost-zero weight-six coefficient has the unique orthogonal
split
\begin{equation}
 \boxed{X=D_{-1}Q+W,\qquad W\in\mathbf H^0_{N,6}.}
 \label{eq:v78-Wilson-split}
\end{equation}
The first term is a redundant field/source redefinition.  The harmonic
coefficient $W$ is an independent EFT Wilson coordinate and is not fixed by the
ghost-one restoration equation \eqref{eq:v78-restorer}.
\end{proposition}
\begin{proof}
This is the finite Hodge decomposition of the cochain complex at ghost number
zero, restricted to $\ker D_0$.  The coexact component vanishes on a closed
vector, leaving the exact and harmonic summands.
\end{proof}

The distinction is physically nonempty.  In four-dimensional pure gravity the
parity-even curvature-cubic density
\begin{equation}
 \mathcal O_{R^3}
 =\int\!\sqrt g\,
 R_{\alpha\beta}{}^{\gamma\delta}
 R_{\gamma\delta}{}^{\mu\nu}
 R_{\mu\nu}{}^{\alpha\beta}
 \label{eq:v78-R3}
\end{equation}
represents a nonredundant weight-six direction after the Ricci/Euler-redundant
terms are removed.  The classic two-loop calculations of Goroff--Sagnotti and
van~de~Ven exhibit a nonzero coefficient in this on-shell class
\cite{V6GS,V6vdV}.  We use that literature only to certify that
$\mathbf H^0_{N,6}$ is not the zero space and that an $R^3$-type Wilson
coordinate must be retained.  Its native Renewal-regulator coefficient is not
assigned here.

\section{A shell-local two-loop forest bound for the pure-gravity BV connection}
\label{sec:v78-forest-bound}

The local range theorem does not by itself control composed source transport.
We now address the genuinely new analytic part left by V77.  Work on one shell
$j$ with hard scale $\Lambda_j=b^j\lambda$ and fixed source window
$|K|\le\mu\ll\Lambda_j$.  Soft lapse/shift, Lorentz and nonminimal variables
remain retained as in V75; every integrated internal propagator therefore
contains the hard shell profile.

At fixed arity $N$ and loop order two there are finitely many connected
pure-gravity/ghost graph words built from the native local vertices.  For a graph
$G$, let $\omega_G$ be its superficial local Taylor degree in the V77 weight-six
bank.  Denote by $R_G$ the complete forest expression in which every proper
one-loop divergent subgraph is subtracted by the restriction of the already
chosen common $S_1$, and the overall local Taylor polynomial through degree
$\omega_G$ is removed.

\begin{theorem}[Uniform hard-annulus forest remainder]
\label{thm:v78-annulus}
For every fixed graph word $G$ in the preceding family and every prescribed
source derivative multi-index $\alpha$ with
$|\alpha|\le\omega_G+1$,
\begin{equation}
 \boxed{
 |\partial_K^\alpha R_{G,j}(K)|
 \le C_{G,\alpha}\,
 \Lambda_j^{-1}|K|^{\omega_G+1-|\alpha|}.}
 \label{eq:v78-annulus}
\end{equation}
After the fixed source normalization used in V76, summing the finite graph bank
gives
\begin{equation}
 \boxed{
 \|(1-\mathsf P_{N,6})\mathcal K^{\rm pg}_{2,j}\|_{*\to *}
 \le C_{N,2}\frac{\mu}{\Lambda_j}.}
 \label{eq:v78-K2-bound}
\end{equation}
The constants are independent of volume, mesh and shell number on the declared
profile/inverse-margin class.  The same estimate holds for the counterterm-only
forest assignments on which a bare shell graph vanishes but a proper Taylor
counterterm survives.
\end{theorem}
\begin{proof}
Rescale every hard loop momentum by $p=\Lambda_jr$.  On interior shells the
hard profiles restrict $r$ to a fixed compact annulus and the inherited native
propagator symbols and all required derivatives are uniformly bounded after
factoring their engineering powers.  Near the lattice ultraviolet cap use
$x=ap$ instead; the completed Brillouin-zone symbols live on a fixed compact
set with the same inverse margins.  Thus all dimensionless differentiated
integrands are uniformly bounded before the external Taylor remainder is
formed.

Apply the Zimmermann forest formula with the actual V64/V77 ownership rule for
proper subgraphs.  For each forest term, take the overall Taylor remainder in
the external source variables.  The integral remainder formula supplies
$\omega_G+1-|\alpha|$ powers of $K$ and the same number of external derivatives.
Relative to the superficial degree, one additional derivative leaves exactly
one inverse hard scale, giving \eqref{eq:v78-annulus}.  Proper-subgraph Taylor
operators do not spoil the estimate: their coefficients are hard-annulus
integrals with the already subtracted superficial degrees.  If the bare assigned
graph vanishes by shell support while a Taylor counterterm remains, the same
argument is applied to that local coefficient before closing the complementary
hard line; this is the general version of the counterterm-only region treated
explicitly in V62.

There are finitely many graph words at fixed loop order, arity and weight.
The V76 rooted/source norm multiplies each coefficient by the corresponding
power of the fixed source scale $\mu$; \eqref{eq:v78-annulus} therefore gives
one factor $\mu/\Lambda_j$ uniformly after summation.  No soft inverse propagator
appears because V75 keeps the soft constraint sector explicit.
\end{proof}

The theorem is a shell-local BPHZ estimate, not a numerical evaluation of the
pure-gravity two-loop amplitudes.  It consumes the common one-loop forest and
the native hard-symbol bounds; it does not assume continuum locality of the
unsubtracted graph.

\section{Pure-gravity retained-constraint two-loop BV self-map}
\label{sec:v78-selfmap}

Let $B_{2,j}^{\rm pg}$ denote the pure-gravity restriction of V77's local
order-two breaking, and let $W_{6,j}^{\rm pg}\in\mathbf H^0_{N,6}$ denote the
independent physical weight-six Wilson coordinates.  Install the restoring
primitive \eqref{eq:v78-restorer}, but leave $W_{6,j}^{\rm pg}$ as an explicit
coordinate determined by physical lower-endpoint data and the native shell
projection.

\begin{theorem}[Two-loop self-map on the pure-gravity retained-constraint sector]
\label{thm:v78-selfmap}
Fix $N\ge7$ and the weight-six local bank.  On the bosonic pure-gravity
restriction of the V75--V77 retained-constraint state, the order-two measured
shell map closes after:
\begin{enumerate}[label=\textup{(\roman*)},leftmargin=2.8em]
\item the common V64/V77 one-loop forest subtractions;
\item the local restoring primitive $-H_1B_{2,j}^{\rm pg}$;
\item retention of the independent harmonic Wilson coordinate
      $W_{6,j}^{\rm pg}$ rather than deleting it;
\item retention of the soft constraint/nonminimal variables themselves.
\end{enumerate}
The new complementary two-loop shell connection satisfies
\begin{equation}
 \boxed{
 \|E^{\rm pg}_{2,j}\|_{*\to *}
 \le C\frac{\mu}{\Lambda_j},}
 \label{eq:v78-E2}
\end{equation}
where V77's quadratic one-loop contribution is separately
$O((\mu/\Lambda_j)^2)$.  Consequently
\begin{equation}
 \boxed{
 \sup_{n\ge j}\|\mathsf U^{\rm pg,[\le2]}_{n,j}\|_{*\to *}<\infty}
 \label{eq:v78-composed}
\end{equation}
uniformly in the number of ultraviolet shell steps on every fixed perturbative
$\hbar$-domain.
\end{theorem}
\begin{proof}
Proposition~\ref{prop:v78-B2-closed} and
Theorem~\ref{thm:v78-pg-range} give the local restoring primitive.
Proposition~\ref{prop:v78-Wilson-split} keeps the independent ghost-zero
cohomology in the state rather than treating it as a breaking.  Theorem
\ref{thm:v78-annulus} gives the genuinely new complementary order-two bound.
V77 gives the $O((\mu/\Lambda_j)^2)$ square of the renormalized one-loop shell
connection.  Equation~\eqref{eq:v77-connection-two} therefore yields
\eqref{eq:v78-E2}.  Since
$\sum_j\mu/\Lambda_j<\infty$, the same product estimate as V76 proves
\eqref{eq:v78-composed} coefficientwise through order two.
\end{proof}

This closes the local-anomaly and source-transport interfaces on the
pure-gravity retained-constraint sector without evaluating its physical
weight-six Wilson coefficients.  Those coefficients are EFT data, not anomaly
coordinates.

\section{Anti-circularity audit and status ledger for V78}
\label{sec:v78-ledgers}

\subsection*{Proved}
The free pure-gravity ghost-one cohomology at coefficient weight six vanishes
for arity at least eight.  Hence every closed weight-six seven-jet extends to a
formal all-arity local cocycle.  Using the same established four-dimensional
pure-Einstein anomaly classification already consumed in V25, the entire
translation-invariant pure-gravity ghost-one weight-six finite jet is in the
BV range for every arity cap $N\ge7$.  The order-two BV consistency identity
makes the restored ordinary native coefficient closed, so its pure-gravity
cokernel coordinate vanishes without numerical evaluation of the missing
native graphs.  A finite minimum-norm restoring section with cutoff-independent
coefficient-map bounds is constructed.  Ghost-zero weight-six local data split
instead into redundant exact directions and independent harmonic EFT Wilson
coordinates; the curvature-cubic class is retained as a nontrivial example.
Finally, the complete two-loop hard forest on the pure-gravity retained-constraint
sector has a shell-local $O(\mu/\Lambda_j)$ quasilocal remainder after the common
weight-six subtraction.  Together with V77's summable one-loop-square term this
gives bounded composed source transport through order two on this sector.

\subsection*{Inherited}
V25 supplies the translation-invariant pure-gravity anomaly bridge and finite
local solver at weight four.  V44 supplies the free high-arity descent mechanism
and formal-extension strategy.  V64 and V77 supply common one-loop forest
ownership and the order-two local master assembly.  V75 keeps soft constraints
and nonminimal variables explicit; V76 supplies the one-loop shell BV connection
and its source norm; V77 supplies the weight-six complex and the quadratic
connection equation.  V62 supplies the counterterm-only forest warning used in
the generic hard-annulus proof.

\subsection*{Literature input}
The absence of a local four-dimensional pure gravitational anomaly in the
ordinary bosonic Einstein complex and the all-arity local BRST classification
are the same inputs already cited in V25 and V44
\cite{V20BBH1994,V25BBHEinstein}.  The Goroff--Sagnotti/van~de~Ven result
\cite{V6GS,V6vdV} is used only to witness that the weight-six ghost-zero Wilson
quotient is nontrivial; its numerical coefficient is not imported into the
Renewal regulator.  The forest organization is standard BPHZ methodology,
with the native ownership and shell estimates supplied internally by prior
versions.

\subsection*{Conditional}
The pure-gravity self-map does not determine the native numerical values of its
weight-six Wilson coordinates.  It also does not prove that the full
Einstein--Higgs ghost-one cokernel vanishes: positive-matter current/descent
classes and the genuinely mixed two-loop shell connection remain separate.
The hard-annulus theorem is at fixed loop order, arity, weight, source order,
profile class and coupling domain; no bound uniform in an unbounded EFT tower
is asserted.

\subsection*{Open}
The smallest load-bearing local problem is now the positive-matter part of
V77's weight-six cokernel together with its mixed retained-constraint shell
coefficient.  The native two-loop pure-gravity graphs are still needed to
populate physical $R^3$-type Wilson coordinates, but not to decide anomaly
restorability.  The next efficient attack is therefore to project the mixed
Einstein--Higgs $B_2$ onto the current/descent sectors already constructed in
V46--V53 and determine whether any flavour-invariant weight-six class survives.
If that projection vanishes, the same hard-annulus estimate promotes the full
bosonic Einstein--Higgs source map to two loops.  Internal gauge, Yukawa and
chiral-line sectors remain downstream and are not inferred from the bosonic
result.  No asymptotic-safety or nonperturbative quantum-gravity conclusion is
made.

\begin{center}\footnotesize
\begin{tabular}{p{.24\textwidth}p{.68\textwidth}}
\toprule
Variable & Scope of V78\\\midrule
Loop order & Two-loop local master/source closure on the pure-gravity retained-constraint restriction; mixed positive-matter cokernel remains open.\\
Local anomaly bank & Translation-invariant pure-gravity ghost-one weight-six complex, arity cap $N\ge7$; exact range theorem and finite restoring section.\\
Physical EFT bank & Independent ghost-zero harmonic weight-six coordinates retained explicitly; $R^3$ is a nontrivial witness, not assigned a native coefficient.\\
Soft sector & Lapse/shift, Lorentz/nonminimal and ghost/source variables remain retained; no inverse soft constraint operator is introduced.\\
Quasilocality & Complete pure-gravity two-loop hard forests after common proper subtractions obey $O(\mu/\Lambda_j)$ in the V76 source norm.\\
Source transport & Bounded through order two on the pure-gravity retained-constraint sector, uniformly in the number of UV shell steps at fixed perturbative domain.\\
Cutoff/volume & Local coefficient-map bounds are mesh/volume independent at fixed bank; shell remainder constants are uniform on the declared native profile and inverse-margin class.\\
Full ESM & Not proved. Mixed positive-matter range, internal gauge/Yukawa, chiral line and arbitrary-order filtered closure remain open.\\
\bottomrule
\end{tabular}
\end{center}

\section*{Append-only continuation marker: V78}
\addcontentsline{toc}{section}{Append-only continuation marker: V78}
\label{sec:v78-continuation-marker}
\textbf{End of V78.} Preserve Sections 1--598.  The pure-gravity part of the
weight-six order-two master breaking is now known to have zero local cokernel,
and its shell BV complement is summable after the common forest subtraction.
The independent ghost-zero weight-six Wilson quotient has been separated from
the anomaly-restoring primitive and must be retained as EFT data.  The next
load-bearing theorem is the positive-matter/mixed Einstein--Higgs weight-six
range problem, using the current/descent machinery of V46--V53 rather than
recomputing the already closed pure-gravity sector.



% =============================================================================
% V79 -- bosonic Einstein--Higgs weight-six range closure
% =============================================================================
\part*{V79: Bosonic Einstein--Higgs weight-six range closure and the two-loop source self-map}
\addcontentsline{toc}{part}{V79: Bosonic Einstein--Higgs weight-six range closure and two-loop source self-map}
\label{sec:v79-start}

V78 removes the pure-gravity part of the weight-six ghost-one obstruction but
leaves the positive-matter component of V77's local range test open.  The next
step should not be another graphwise coefficient calculation.  There is a
stronger upstream question: in the ordinary bosonic Einstein--Higgs BV complex,
can a flavour-invariant positive-scalar-degree ghost-one local class exist at
all?

The answer is no in the present branch.  The local BRST classification of
Einstein--Yang--Mills theory with general matter fields applies to matter fields
that transform as spacetime scalars after converting world indices to Lorentz
indices.  With no internal gauge factor in the present Einstein--Higgs
restriction, there is no free abelian gauge field capable of coupling a
conserved matter current into a new antifield-dependent anomaly.  The remaining
candidate anomalies are the ordinary chiral/topological classes.  The current
branch contains only real bosonic scalars, hence no chiral matter anomaly, and
four spacetime dimensions are not among the exceptional dimensions supporting
the topological candidates in that classification.  Consequently the complete
positive-matter ghost-one local class vanishes.

The literature theorem is used only for this ordinary local cohomology
statement.  V46--V53 remain essential internally: they supply explicit
translation-compatible current sections, control the positive-matter
antifield-dependent representatives at every prescribed scalar degree, and
show how current improvements are lifted without introducing inverse fields or
soft momenta.  Their role below is constructive and finite-order; the external
classification is not used as a substitute for the native measured shell map.

Combining this range theorem with V78's pure-gravity result gives zero local
bosonic Einstein--Higgs ghost-one cokernel at weight six.  The ghost-zero
cohomology remains nontrivial and is retained as the mixed gravitational EFT
Wilson bank.  Finally the hard-annulus forest argument of V78 extends to the
mixed scalar--gravity graph family because every surviving integrated line is
hard after V75's soft-constraint retention.  This closes the bosonic
Einstein--Higgs retained-constraint source map through two loops at fixed
weight, arity and perturbative order, without assigning numerical values to
unevaluated Wilson coefficients.

\section{The ordinary bosonic Einstein--Higgs anomaly complex}
\label{sec:v79-complex}

Work in the ordinary local reference used in V77--V78.  The fields are the
coframe/metric and its diffeomorphism/local-Lorentz ghosts, the nonminimal
doublets, and four equal real scalar fields \(\phi_A\), together with their BV
antifields.  The scalars are neutral under any internal gauge group in this
restriction.  Their only gauge transformation is the diffeomorphism action;
in a coframe formulation they are Lorentz scalars.  The independent flavour
sign group
\begin{equation}
 \mathfrak F=(\mathbb Z_2)^4,
 \qquad \phi_A\longmapsto\epsilon_A\phi_A,
 \qquad \epsilon_A=\pm1,
 \label{eq:v79-flavour-signs}
\end{equation}
is retained as in V34 and V45--V53.

Let \(\mathfrak L^{g,+}_{6}\) denote the translation-invariant local
ghost-number-\(g\), weight-at-most-six subspace of positive scalar degree in
the formal scalar-degree completion of the ordinary Einstein--Higgs local BV
algebra, modulo total derivatives.  The mass variable \(z=m_H^2\) is a
spectator parameter of weight two.  Let
\begin{equation}
 s_{\rm EH}=(S_{0,\rm EH},\cdot)
 \label{eq:v79-sEH}
\end{equation}
be the ordinary classical differential.  It preserves the weight, scalar-sign
invariance and mass-polynomial grading.

\begin{theorem}[No bosonic positive-matter anomaly class in four dimensions]
\label{thm:v79-no-matter-anomaly}
On the preceding ordinary Einstein--Higgs branch,
\begin{equation}
 \boxed{
 H^{1,4}_{\rm loc,+}(s_{\rm EH}\mid d_x)=0.}
 \label{eq:v79-H1zero}
\end{equation}
More explicitly, if a translation-invariant, flavour-sign-invariant local
four-form \(B^+\) has ghost number one, positive scalar degree, finite weight
and satisfies
\begin{equation}
 s_{\rm EH}B^+ + d_x M=0,
 \label{eq:v79-cocycle}
\end{equation}
then there are translation-invariant local forms \(C^+\), \(N\) with the same
mass-polynomial and flavour-sign symmetries such that
\begin{equation}
 \boxed{B^+=s_{\rm EH}C^+ + d_xN.}
 \label{eq:v79-exact}
\end{equation}
At fixed homogeneous weight, \(C^+\) may be chosen at that same weight.
\end{theorem}
\begin{proof}
The external input is the local BRST classification for
Einstein--Yang--Mills theory with general matter fields
\cite{V20BBH1995}, together with the general antifield framework
\cite{V20BBH1994}.  That analysis permits matter fields which transform as
spacetime scalars after converting world indices to Lorentz indices and treats
all ghost numbers.  In the absence of free abelian gauge factors, candidate
anomalies can be represented without antifields; the extra current-dependent
anomaly classes arise only when such abelian factors are present.  The present
restriction has none.

The remaining candidates are the ordinary chiral and exceptional topological
classes.  Four real scalar fields carry no chiral fermion measure, and the
exceptional topological anomaly dimensions listed by the classification do
not include four.  Hence there is no nontrivial ghost-one four-form class on
this bosonic branch.  Decomposing by scalar degree, mass degree and the V77
weight is legitimate because the classical differential preserves each of
those filtrations.  Averaging a primitive over the finite group
\(\mathfrak F\) preserves the equation and does not increase the chosen
coefficient norm.  The classification is formulated in covariant local
variables and introduces no explicit spacetime coordinate, so the primitive
may be kept in the translation-invariant local category used here.
\end{proof}

\begin{remark}[What this theorem does not cover]
\label{rem:v79-scope}
The theorem is deliberately bosonic.  It does not include the Standard-Model
hypercharge factor, chiral fermions, determinant/Pfaffian orientation, or
possible global anomalies.  Those are precisely the sectors for which the
full ESM programme keeps the anomaly and line-transport questions separate.
No conclusion of this section is transferred to them.
\end{remark}

\section{V46--V53 give a constructive finite-order section of the matter-current part}
\label{sec:v79-current-section}

The cohomology theorem proves existence, but the finite measured programme also
needs control of the representative and of current improvements.  For the
positive-matter antifield-dependent component this is already available
internally.

Let \(\mathfrak C^{1,+}_{M,6}\subset\mathfrak L^{1,+}_{6}\) be the natural,
flavour-invariant current-derived subspace through scalar matter count at most
\(M\), in the category of V46--V53.  It includes the complete metric/scalar
Euler returns and the translation-sensitive two-form descendants of the
corresponding Noether currents.

\begin{theorem}[Finite-order current/descent section]
\label{thm:v79-current-section}
For every prescribed finite positive matter order \(M\), there is a linear
local map
\begin{equation}
 \mathsf H^{\rm cur}_{M,6}:\mathfrak C^{1,+}_{M,6}
       \longrightarrow \mathfrak L^{0,+}_{6}
 \label{eq:v79-Hcur}
\end{equation}
such that, modulo total derivatives,
\begin{equation}
 \boxed{
 s_{\rm EH}\mathsf H^{\rm cur}_{M,6}B=B
 \qquad(B\in\mathfrak C^{1,+}_{M,6},\ s_{\rm EH}B=0).}
 \label{eq:v79-Hcur-identity}
\end{equation}
The map is translation compatible, polynomial in \(z\), equivariant under
\(\mathfrak F\), and has a finite coefficient bound independent of mesh,
periods and shell number at fixed \(M\), arity and metric chart.
\end{theorem}
\begin{proof}
V46 removes the complete natural scalar-linear invariant current sector after
the flavour-singlet test.  V47--V49 supply the metric-changing Ricci and
curvature-current sections with their full matter-stress returns.  V50--V52
iterate the scalar and metric sections to every prescribed positive matter
order without dividing by a scalar field, mass or soft momentum.  V53 then
lifts the exact current packet through the translation-sensitive descent,
including the Green current and the superpotential needed for an improved
physical Noether current.  Composing those fixed sections gives
\eqref{eq:v79-Hcur}; their stated finite-order coefficient estimates multiply
to a finite constant.  Every step is flavour equivariant, so finite sign
averaging preserves the bound.
\end{proof}

This theorem is not needed to infer the vanishing of
\eqref{eq:v79-H1zero}; rather, it gives the explicit internal section on the
part of the native mixed breaking which V46--V53 already identify.  The
residual antifield-independent anomaly candidate is what
Theorem~\ref{thm:v79-no-matter-anomaly} removes.

\section{The positive-matter component of V77's two-loop cokernel vanishes}
\label{sec:v79-cokernel}

Return to V77's ordinary weight-six bank.  Let
\begin{equation}
 B_2=B_2^{\rm pg}+B_2^+
 \label{eq:v79-Bsplit}
\end{equation}
be the decomposition by scalar degree, where \(B_2^{\rm pg}\) is the pure
scalar-degree-zero part treated in V78 and \(B_2^+\) contains every positive
scalar-degree coefficient.  Proposition~\ref{prop:v78-B2-closed} gives
\begin{equation}
 d_{N,6}B_2=0
 \label{eq:v79-Bclosed}
\end{equation}
on the restored ordinary branch.  Since the scalar-degree filtration is
preserved,
\begin{equation}
 d_{N,6}B_2^+=0.
 \label{eq:v79-Bplusclosed}
\end{equation}

The finite arity cap requires one bookkeeping point.  The actual order-two
master functional exists before the arity projection.  Apply the common local
weight-six Taylor map first in the formal scalar-degree completion and denote
its positive-matter coefficient by \(\widehat B_2^+\).  The exact order-two
consistency identity makes \(\widehat B_2^+\) a full local cocycle.  Its
finite arity projection is the \(B_2^+\) above.  Because
\(s_{\rm EH}\) does not lower arity, projecting any full primitive back to
arity at most \(N\) preserves the finite range equation.

\begin{theorem}[Mixed Einstein--Higgs weight-six range theorem]
\label{thm:v79-range}
For every fixed arity cap \(N\) and every fixed finite scalar/mass degree
appearing in the order-two bosonic Einstein--Higgs coefficient,
\begin{equation}
 \boxed{P_{\rm coker}^{+}B_2=0.}
 \label{eq:v79-cokerzero}
\end{equation}
Equivalently there is a finite local restoring functional
\(C_{2,\rm rest}^{+}\in\mathfrak L^0_{N,6}\) such that
\begin{equation}
 \boxed{d_{N,6}C_{2,\rm rest}^{+}=-B_2^+.}
 \label{eq:v79-restorer}
\end{equation}
In a fixed rational coefficient basis the minimum-norm choice is
\begin{equation}
 C_{2,\rm rest}^{+}=-(D_0^+)^\dagger B_2^+,
 \qquad
 \|C_{2,\rm rest}^{+}\|
 \le (\sigma_{N,6}^{+})^{-1}\|B_2^+\|,
 \label{eq:v79-minrestorer}
\end{equation}
where \(\sigma_{N,6}^{+}\) is the least nonzero singular value of the
positive-matter block of \(d_{N,6}\).  The bound is independent of mesh,
periods and shell number at fixed bank.
\end{theorem}
\begin{proof}
The full positive-matter local cocycle \(\widehat B_2^+\) is exact by
Theorem~\ref{thm:v79-no-matter-anomaly}.  Project its primitive to weight six,
the required mass/flavour sector and arity at most \(N\).  Weight and flavour
projection commute with the differential.  Since the differential does not
lower arity, no term of primitive arity greater than \(N\) can contribute back
to \(B_2^+\).  This gives \eqref{eq:v79-restorer}.  The Moore--Penrose formula
and singular-value estimate are finite-dimensional linear algebra.  On the
current-derived block one may choose instead the explicit section of
Theorem~\ref{thm:v79-current-section}; both sections differ only by a
closed ghost-zero local term.
\end{proof}

Combining with V78 gives the entire bosonic ghost-one result
\begin{equation}
 \boxed{P_{\rm coker}^{\rm EH,bos}B_2=0.}
 \label{eq:v79-full-cokerzero}
\end{equation}
No numerical two-loop gravity or mixed-matter amplitude is needed to decide
this range question.

\section{Ghost-zero mixed EFT data remain independent Wilson coordinates}
\label{sec:v79-wilson}

As in V78, vanishing ghost-one cohomology does not make the order-two action
bank redundant.  Let
\begin{equation}
 D_{-1}^{\rm EH}:\mathfrak L^{-1}_{N,6}	o\mathfrak L^0_{N,6},
 \qquad
 D_0^{\rm EH}:\mathfrak L^0_{N,6}	o\mathfrak L^1_{N,6}
 \label{eq:v79-D0}
\end{equation}
be the bosonic Einstein--Higgs matrices, and define
\begin{equation}
 \mathbf H^{0,\rm EH}_{N,6}
 =\ker D_0^{\rm EH}\cap(\operatorname{Ran}D_{-1}^{\rm EH})^\perp.
 \label{eq:v79-H0}
\end{equation}
Then every closed local action coefficient has the orthogonal decomposition
\begin{equation}
 \boxed{X=D_{-1}^{\rm EH}Q+W,
 \qquad W\in\mathbf H^{0,\rm EH}_{N,6}.}
 \label{eq:v79-Wsplit}
\end{equation}

The positive-matter part of \(W\) includes the ordinary mixed
curvature--Higgs EFT directions permitted at dimension/weight six.  Depending
on the chosen integration-by-parts and equation-of-motion basis, representatives
may be taken from combinations of curvature with Higgs kinetic or potential
insertions and from higher-derivative scalar operators.  V79 does not declare a
particular basis canonical and does not set any of their native coefficients to
zero.

\begin{proposition}[Restoration and Wilson matching are independent problems]
\label{prop:v79-wilson-separate}
Installing the ghost-one restoring functional
\(C_{2,\rm rest}^{\rm EH}\) fixes the order-two master row but leaves the
harmonic coordinate \(W_{6,j}^{\rm EH}\in\mathbf H^{0,\rm EH}_{N,6}\)
undetermined.  Changing the latter within a BRST-closed physical
renormalization changes EFT Wilson data without reopening the local anomaly
condition.  Conversely a nonzero Wilson coefficient is not evidence of a
nonzero ghost-one cokernel.
\end{proposition}
\begin{proof}
The restoring equation concerns the range of \(D_0^{\rm EH}\), whereas
\eqref{eq:v79-H0} is orthogonal to that range after quotienting redundant
\(D_{-1}^{\rm EH}\) directions.  The two coordinates therefore occupy
different cohomological degrees.  The statement is the finite Hodge splitting
of the local complex.
\end{proof}

This distinction is required by the EFT programme: the bosonic two-loop shell
may generate new mixed gravity--matter Wilson coefficients even though its
master breaking is locally restorable.

\section{The hard-annulus forest estimate extends to mixed scalar--gravity graphs}
\label{sec:v79-forest}

It remains to control the genuinely mixed two-loop shell connection rather than
only its local range.  Fix loop order two, arity \(N\), weight at most six and
an external source window \(|K|\le\mu\ll\Lambda_j\).  The V75 state retains
the soft constraint/nonminimal variables.  Therefore a surviving integrated
constraint, graviton, ghost or scalar line in the shell calculation always
carries the hard shell profile.  Purely soft constraint exchange is not
integrated and never produces an inverse soft Laplacian inside the Wilsonian
remainder.

For each connected mixed graph word \(G\), let \(\omega_G\) be its overall
local Taylor degree and let \(R^{\rm EH}_G\) be the complete forest expression
with every proper one-loop divergent subgraph subtracted by the restriction of
the already chosen common \(S_1\), followed by the overall weight-six Taylor
subtraction.

\begin{theorem}[Mixed hard-annulus forest remainder]
\label{thm:v79-annulus}
For every fixed mixed Einstein--Higgs two-loop graph word and every
\(|\alpha|\le\omega_G+1\),
\begin{equation}
 \boxed{
 |\partial_K^\alpha R^{\rm EH}_{G,j}(K)|
 \le C_{G,\alpha}\Lambda_j^{-1}
       |K|^{\omega_G+1-|\alpha|}.}
 \label{eq:v79-annulus}
\end{equation}
After the V76 source normalization and summation of the finite graph bank,
\begin{equation}
 \boxed{
 \|(1-\mathsf P_{N,6})\mathcal K^{\rm EH,+}_{2,j}\|_{*\to *}
 \le C_{N,2}^{+}\frac{\mu}{\Lambda_j}.}
 \label{eq:v79-K2}
\end{equation}
The constants are independent of volume, mesh and shell number on the fixed
profile, inverse-margin, coupling and source-order domain.  The same estimate
holds on counterterm-only shell assignments for which the bare graph vanishes
but a proper Taylor counterterm survives.
\end{theorem}
\begin{proof}
The proof is the V78 annulus argument with the mixed native symbol bounds added.
V27, V62--V64, V72 and V74 provide the required uniformly differentiated hard
scalar, mixed scalar--metric and contact symbols; V76 supplies the nonlinear
constraint/ghost/density vertices at one loop.  After rescaling every hard loop
momentum by \(p=\Lambda_jr\), all differentiated propagators and local vertices
live on a fixed compact annulus with uniform inverse margins.  Near the finite
ultraviolet cap the change \(x=ap\) gives the same compact-symbol control.

Apply the common Zimmermann forest with V77's ownership rule.  The overall
Taylor remainder contributes \(\omega_G+1-|\alpha|\) powers of the soft
external variables.  One derivative beyond the superficial degree leaves one
inverse hard scale, producing \eqref{eq:v79-annulus}.  Proper-subgraph Taylor
coefficients have already had their own superficial degrees removed and obey
the same annular estimate before the complementary hard line is closed.  V75
excludes any inverse soft constraint propagator from the integrated graph.
There are finitely many words at the stated loop order, arity and weight, so the
V76 rooted/source norm gives \eqref{eq:v79-K2} after summation.
\end{proof}

The theorem is structural and coefficientwise.  It does not numerically evaluate
the mixed two-loop Wilson coefficients.

\section{Two-loop bosonic Einstein--Higgs retained-constraint self-map}
\label{sec:v79-selfmap}

Let
\begin{equation}
 C_{2,\rm rest}^{\rm EH}
 =C_{2,\rm rest}^{\rm pg}+C_{2,\rm rest}^{+}
 \label{eq:v79-fullrestorer}
\end{equation}
combine V78's pure-gravity primitive with the positive-matter primitive above.
Retain separately the complete physical ghost-zero Wilson coordinate
\(W_{6,j}^{\rm EH}\in\mathbf H^{0,\rm EH}_{N,6}\).

\begin{theorem}[Two-loop bosonic Einstein--Higgs source/BV self-map]
\label{thm:v79-selfmap}
Fix the loop order through \(\hbar^2\), the local weight cap six, a finite arity
and source order, and the V75 retained soft-constraint state.  On the bosonic
Einstein--Higgs branch with no internal gauge or Yukawa couplings, the measured
shell map closes after:
\begin{enumerate}[label=\textup{(\roman*)},leftmargin=2.8em]
\item the common one-loop forest subtractions fixed by \(S_1\);
\item the local order-two restoring functional
      \(-C_{2,\rm rest}^{\rm EH}\) in the sign convention of V77;
\item retention of all independent ghost-zero weight-six Einstein--Higgs EFT
      Wilson coordinates;
\item retention of the soft constraint, local-Lorentz, nonminimal and source
      variables of V75 rather than global elimination.
\end{enumerate}
The renormalized order-two adjacent source map has
\begin{equation}
 \boxed{
 \|E^{\rm EH}_{2,j}\|_{*\to *}
 \le C\frac{\mu}{\Lambda_j},}
 \label{eq:v79-E2}
\end{equation}
while V77's quadratic one-loop term remains
\(O((\mu/\Lambda_j)^2)\).  Hence
\begin{equation}
 \boxed{
 \sup_{n\ge j}
 \|\mathsf U^{\rm EH,[\le2]}_{n,j}\|_{*\to *}<\infty}
 \label{eq:v79-composed}
\end{equation}
uniformly in the number of ultraviolet shell steps on every fixed bosonic
perturbative/EFT domain.

On the antifield-free action spine, the local dimension-six Wilson coordinates
and their quasilocal complements retain the previously proved filtered window
\begin{equation}
 \boxed{2\le\sigma=3<4,}
 \label{eq:v79-window}
\end{equation}
so the corresponding coefficientwise two-boundary trajectory survives the
source/BV completion.
\end{theorem}
\begin{proof}
V78 and Theorem~\ref{thm:v79-range} give zero local ghost-one cokernel and the
common restoring primitive.  Proposition~\ref{prop:v79-wilson-separate} keeps
physical ghost-zero EFT data as state coordinates.  Theorem~\ref{thm:v79-annulus}
gives the new mixed complementary order-two shell bound, while V77 supplies the
summable square of the one-loop shell connection.  The exact connection-curvature
identity of V76 then yields \eqref{eq:v79-E2}.  Since
\(\sum_j\mu/\Lambda_j<\infty\), the same product argument as V76 and V78 proves
\eqref{eq:v79-composed}.  The action-side filtered window is inherited from
V73--V74 and is not altered by the source restoring primitive, which lives in
the separate BV/source coordinate block.
\end{proof}

This is the first two-loop source-complete self-map for the 
\emph{bosonic Einstein--Higgs restriction} of the programme.  It establishes
existence, locality and transport control of the finite-order state; it does not
supply numerical values for every two-loop Wilson coefficient and does not yet
include the internal gauge, Yukawa or chiral determinant-line sectors.

\section{Anti-circularity audit and status ledger for V79}
\label{sec:v79-ledgers}

\subsection*{Proved}
The ordinary four-dimensional bosonic Einstein--Higgs ghost-one local cohomology
has no positive-matter anomaly class on the branch with neutral real scalars and
no internal gauge factor.  V46--V53 provide an explicit finite-order section of
the current-derived positive-matter component and its translation-sensitive
descent.  Applying the ordinary order-two consistency identity before the
finite arity restriction shows that the positive-matter part of V77's weight-six
breaking lies in the BV range at every prescribed finite bank.  Together with
V78, the complete bosonic Einstein--Higgs local ghost-one cokernel therefore
vanishes at two loops.  The independent ghost-zero mixed EFT Wilson quotient is
retained rather than confused with the restoring primitive.  The V78
hard-annulus forest estimate extends to the mixed scalar--gravity two-loop graph
family, yielding an \(O(\mu/\Lambda_j)\) complementary shell connection.  Hence
composed bosonic Einstein--Higgs source transports are bounded through order two
uniformly in the number of ultraviolet shell steps at fixed perturbative/EFT
order.

\subsection*{Inherited}
V46--V53 supply the natural matter-current classifications, all-degree scalar
section, finite-order current elimination and translation-compatible ascent.
V59 supplies exact finite defect retention.  V64--V65 and V71--V74 supply the
common one-loop forest ownership and populated scalar/mixed action pieces.
V75 retains the soft gravitational constraints.  V76 supplies the shell BV
connection and one-loop source norm.  V77 supplies the weight-six coupled
complex and two-loop range equation.  V78 supplies the pure-gravity range theorem,
independent Wilson quotient and pure-gravity hard-annulus estimate.

\subsection*{Literature input}
The decisive external input is the local BRST cohomology classification of
Einstein--Yang--Mills theory with matter fields
\cite{V20BBH1995}, together with the general antifield framework
\cite{V20BBH1994}.  It is used only to identify the ordinary local ghost-one
cohomology on the present bosonic branch: no free abelian gauge field is present,
real scalars generate no chiral anomaly, and four dimensions is not an exceptional
topological-anomaly dimension in that classification.  No native loop coefficient,
cutoff estimate or Wilson coefficient is imported from that literature.

\subsection*{Conditional}
The two-loop self-map is at fixed local weight, arity, source order, profile
class and bosonic coupling domain.  Numerical values of the new mixed and
pure-gravity ghost-zero Wilson coordinates remain native matching data.  The
hard-annulus theorem assumes the inherited uniform hard-symbol/inverse-margin
bounds and the V75 hard/soft support separation.  It is not uniform in an
unbounded EFT tower or perturbative order.

\subsection*{Open}
The smallest load-bearing theorem has now moved out of the neutral bosonic
Einstein--Higgs sector.  The next extension must adjoin the internal
\(SU(3)\times SU(2)\times U(1)_Y\) gauge/ghost complex and the Higgs gauge
representation to the same weight-six filtered state.  The semisimple
nonabelian factors can use the existing V1--V4/V20 gauge machinery, but the
abelian hypercharge factor must retain its current-dependent BRST possibilities
rather than inheriting the bosonic theorem above.  After that, Yukawa/chiral
fermions and the determinant/Pfaffian line must be merged with the local anomaly
condition before a full ESM two-loop self-map can be claimed.  No analytic
ultraviolet trajectory for gravity or for the full Standard Model is asserted.

\begin{center}\footnotesize
\begin{tabular}{p{.24\textwidth}p{.68\textwidth}}
\toprule
Variable & Scope of V79\\\midrule
Loop order & Source-complete bosonic Einstein--Higgs closure through two loops at fixed weight six; Wilson coefficients need not be numerically evaluated.\\
Local anomaly bank & Diffeomorphism/local-Lorentz plus four neutral real scalars; positive-matter and pure-gravity ghost-one cokernels both vanish.\\
Current/descent & V46--V53 supply constructive finite-order sections for the current-derived positive-matter component; external cohomology is used for the residual ordinary class.\\
Physical EFT bank & Pure and mixed ghost-zero harmonic weight-six coordinates are retained as Wilson data, not restoring terms.\\
Soft sector & Lapse/shift, Lorentz/nonminimal and source/ghost coordinates remain retained; no inverse soft wave operator.\\
Quasilocality & Complete mixed two-loop hard forests after common proper subtractions obey $O(\mu/\Lambda_j)$ in the V76 source norm.\\
Source transport & Bounded through order two uniformly in the number of UV shell steps on the fixed bosonic perturbative/EFT domain.\\
Full ESM & Not proved. Internal gauge, hypercharge-current, Yukawa/chiral and determinant/Pfaffian-line sectors remain open.\\
\bottomrule
\end{tabular}
\end{center}

\section*{Append-only continuation marker: V79}
\addcontentsline{toc}{section}{Append-only continuation marker: V79}
\label{sec:v79-continuation-marker}
\textbf{End of V79.} Preserve Sections 1--605.  The weight-six local ghost-one
cokernel now vanishes on the complete bosonic Einstein--Higgs restriction, and
the mixed two-loop shell BV complement is summable after the common forest
subtraction.  Independent ghost-zero mixed gravity--Higgs Wilson coordinates
remain explicit EFT data.  The next load-bearing extension is the internal gauge
complex, especially the abelian hypercharge/current sector, before Yukawa/chiral
and determinant-line transport are merged.


\clearpage
\part*{V80: Bosonic internal-gauge extension and the hypercharge anomaly projection}
\addcontentsline{toc}{part}{V80: Bosonic internal-gauge extension and the hypercharge anomaly projection}

\section{The next obstruction is the standard gauge-anomaly projection, not a free-abelian current class}
\label{sec:v80-start}

Sections 1--605 and all their labels are retained unchanged.  V79 closes the
weight-six ghost-one range problem for diffeomorphism/local-Lorentz gravity
coupled to four neutral real Higgs components.  The next extension must turn
those four real components into the actual charged Higgs representation and
adjoin the internal
\(SU(3)\times SU(2)\times U(1)_Y\) gauge/ghost complex.  The first question is
therefore not yet a new Feynman integral.  It is whether the abelian factor
creates a new antifield-dependent obstruction which would invalidate V79's
range argument before any numerical coefficient is evaluated.

We work on the bosonic Standard-Model branch only.  The retained fields are the
V75 gravitational/nonminimal packet, the colour and weak gauge fields
\(A^3,A^2\), the hypercharge gauge field \(B\), their ghosts and antifields, and
the Higgs doublet \(H\in\C^2\), equivalently four real scalar components.
There are no quarks, leptons, Yukawa matrices, Majorana terms or fermionic
Pfaffian line in this installment.  The Higgs hypercharge is a fixed nonzero
number \(y_H\); no result below depends on its normalization.

At the ordinary local level write the internal part of the BRST differential
as
\begin{align}
 s_{\rm int}A^3_\mu&=D_\mu c_3,
 &s_{\rm int}c_3&=-\tfrac12[c_3,c_3],\\
 s_{\rm int}A^2_\mu&=D_\mu c_2,
 &s_{\rm int}c_2&=-\tfrac12[c_2,c_2],\\
 s_{\rm int}B_\mu&=\partial_\mu c_Y,
 &s_{\rm int}c_Y&=0,\\
 s_{\rm int}H&=-c_2\cdot T_HH-i y_Hc_YH,
 &s_{\rm int}H^\dagger&=H^\dagger c_2\cdot T_H+i y_Hc_YH^\dagger.
 \label{eq:v80-brst}
\end{align}
The colour action on \(H\) is trivial.  Diffeomorphism and local-Lorentz
pieces are the same retained transformations used in V79 and commute with the
internal action in the declared product representation.

\begin{proposition}[Hypercharge is not a free abelian factor on the Higgs branch]
\label{prop:v80-not-free}
In the terminology of the Einstein--Yang--Mills local-cohomology
classification \cite{V20BBH1995}, the hypercharge generator is not a
``free'' abelian generator on the present bosonic branch.  Equivalently, it is
not an abelian gauge generator acting trivially on all matter fields.
\end{proposition}
\begin{proof}
The abelian generator acts on the Higgs by
\(\delta_YH=-iy_H\epsilon_YH\).  Since \(y_H\ne0\), this action is nonzero.
In real four-component notation it is an antisymmetric generator
\(Y_H=y_HJ_H\ne0\).  Thus the condition
\(\delta_Iy^m=0\) for every matter field, used to define the free abelian
indices in \cite{V20BBH1995}, fails for \(I=Y\).  The same fact is visible in
the nonzero gauge current
\begin{equation}
 J_Y^\mu
 =iy_H\bigl(H^\dagger D^\mu H-(D^\mu H)^\dagger H\bigr),
 \label{eq:v80-hypercurrent}
\end{equation}
which is not identically zero as a local polynomial.  No equation of motion or
vacuum expectation value is used.
\end{proof}

This removes one possible circularity but does not prove anomaly freedom.  The
extra antifield-dependent classes which occur only for free abelian fields are
absent, whereas the ordinary four-dimensional gauge-anomaly classes associated
with an abelian factor remain possible cohomology representatives.  The native
breaking must therefore be projected onto that finite standard anomaly space.

\section{The weight-six bosonic Einstein--gauge--Higgs complex and its anomaly quotient}
\label{sec:v80-complex}

Let \(\mathfrak L^{g,{\rm EGH}}_{N,6}\) be V77's weight-six local complex
enlarged by the internal gauge fields, ghosts, antifields and nonminimal pairs,
with the Higgs in the representation just specified.  Canonical weight is
assigned in the same derivative/mass convention as V77; the internal ghosts
have weight zero and one unit of ghost number, and the Yang--Mills curvatures
have weight two.  The local projector \(\mathsf P^{\rm EGH}_{N,6}\) is the
componentwise Taylor projector determined by
\begin{equation}
 d_I^{(6)}=6+g(I)-n_*(I),
 \label{eq:v80-degree}
\end{equation}
followed by the finite gauge-covariant tensor projection on the declared colour,
weak and hypercharge representations.

\begin{theorem}[Finite weight-six EGH complex]
\label{thm:v80-complex}
At every fixed arity \(N\), the ordinary coupled differential
\begin{equation}
 d^{\rm EGH}_{N,6}
 =\mathsf P^{\rm EGH}_{N,6}
 \bigl((S_0^{\rm EGH},\cdot)\bigr)
 \mathsf P^{\rm EGH}_{N,6}
 \label{eq:v80-d}
\end{equation}
is nilpotent on the projected local bank and
\begin{equation}
 \boxed{
 (d^{\rm EGH}_{N,6})^2=0,
 \qquad
 \mathsf P^{\rm EGH}_{N,6}d^{\rm EGH}_{N,6}
 =d^{\rm EGH}_{N,6}\mathsf P^{\rm EGH}_{N,6}.}
 \label{eq:v80-chain}
\end{equation}
The nonminimal pairs remain contractible.  The coefficient norms of the
projection and inclusion maps are finite constants depending on \(N\), the
fixed representation matrices and weight six, but not on mesh, volume or shell
number.
\end{theorem}
\begin{proof}
Before projection the ordinary Einstein--Yang--Mills--Higgs BV action obeys the
classical master equation, so its Hamiltonian vector field squares to zero.
The internal gauge transformations in \eqref{eq:v80-brst} preserve canonical
weight and the tensor type of every local coefficient.  The same paired-jet
argument as V44 and V77 shows that Taylor projection through the declared
weight commutes with the differential after all descendants at that weight are
retained.  The internal representation projection is an intertwiner for the
colour, weak and hypercharge actions and hence also commutes with the
differential.  Tensoring the V77 nonminimal contraction with the standard
internal antighost--auxiliary doublets proves the last assertion.  All spaces
at fixed \(N\) and weight are finite dimensional.
\end{proof}

The local ghost-one cohomology is now not zero.  What is needed for the
measured theory is a finite projection onto the ordinary anomaly representatives.
Denote by \(\mathfrak A_{\rm std}\) the span, in the present four-dimensional
product gauge group, of the independent consistent descendants associated with
\begin{equation}
 \operatorname{tr}(F_3^3),\qquad
 F_Y\operatorname{tr}(F_3^2),\qquad
 F_Y\operatorname{tr}(F_2^2),\qquad
 F_Y^3,\qquad
 F_Y\operatorname{tr}(\mathcal R^2).
 \label{eq:v80-anomaly-polys}
\end{equation}
Here the notation records the six-form invariant polynomial whose descent gives
the ghost-one four-form; it does not assert a nonzero native coefficient.  The
pure \(SU(2)^3\) cubic invariant is absent, and mixed products containing one
traceless generator from each nonabelian simple factor vanish by factorization.

Choose once and for all a coefficient projection
\begin{equation}
 \mathsf A_{\rm std}:\ker d^{\rm EGH}_{N,6}|_{g=1}
 \longrightarrow \mathfrak A_{\rm std}
 \label{eq:v80-anomaly-proj}
\end{equation}
by evaluating the corresponding independent flat-background curvature/ghost
jets.  Any two such choices differ only by an invertible finite change of basis
on \(\mathfrak A_{\rm std}\).

\begin{theorem}[Cohomological reduction of the bosonic EGH range problem]
\label{thm:v80-cohom-reduction}
On the present branch, the free-abelian antifield-dependent anomaly families of
\cite{V20BBH1995} are absent by Proposition~\ref{prop:v80-not-free}.  Modulo
local divergences and the exceptional topological classes which do not occur in
four dimensions, every remaining ghost-number-one obstruction is represented
by the standard anomaly space \(\mathfrak A_{\rm std}\).  Consequently a
closed populated coefficient \(B\) lies in the local BV range whenever
\begin{equation}
 \boxed{\mathsf A_{\rm std}B=0.}
 \label{eq:v80-range-test}
\end{equation}
At finite arity the corresponding restoring primitive may be chosen by the
Moore--Penrose inverse on the orthogonal complement of
\(\mathfrak A_{\rm std}\).
\end{theorem}
\begin{proof}
The matter-inclusive Einstein--Yang--Mills classification
\cite{V20BBH1995} separates the special antifield-dependent classes associated
with free abelian generators from the ordinary trace/chiral anomaly classes.
The former are absent because hypercharge acts nontrivially on the Higgs.  In
four dimensions the exceptional vielbein-topological anomaly dimensions of that
classification are absent.  What remains at ghost number one is therefore the
finite ordinary gauge-anomaly quotient represented by
\eqref{eq:v80-anomaly-polys}, together with exact terms.  Restriction to the
finite weight-six/arity bank commutes with the differential by
Theorem~\ref{thm:v80-complex}; the same finite-jet extension argument used in
V78 then identifies the finite cokernel with the restriction of this ordinary
quotient.  Orthogonal finite-dimensional Hodge decomposition gives the stated
primitive when its projection vanishes.
\end{proof}

The theorem is deliberately a range criterion, not yet an assertion that the
native bosonic shell has zero anomaly projection.

\section{The native bosonic hard shell has zero standard gauge-anomaly projection}
\label{sec:v80-native-zero}

We now use the finite measured regulator rather than a continuum anomaly
coefficient.  The internal gauge shell inherited from the gauge--matter line is
chosen gauge covariantly.  At each fixed cutoff the hard bosonic Hessian and
ghost operator transform by conjugation on their finite hard carriers.  Write
schematically
\begin{equation}
 \delta_\epsilon\mathbb H=[\mathbb R_\epsilon,\mathbb H],
 \qquad
 \delta_\epsilon\mathbb M_{\rm gh}
 =[\mathbb R^{\rm gh}_\epsilon,\mathbb M_{\rm gh}].
 \label{eq:v80-conjugation}
\end{equation}
The hard projectors and reference covariances are the inherited gauge-covariant
ones, so the same identities hold for the supported operators entering the
normalized shell determinant.

\begin{lemma}[Finite bosonic determinant has no internal gauge Jacobian]
\label{lem:v80-det-zero}
For every finite supported hard block on the present bosonic branch,
\begin{equation}
 \delta_\epsilon\left[
 \frac12\operatorname{Tr}\log\mathbb H
 -\operatorname{Tr}\log\mathbb M_{\rm gh}
 \right]=0.
 \label{eq:v80-det-inv}
\end{equation}
The same is true for the Higgs Gaussian determinant and for every coefficient
obtained by differentiating it with respect to retained bosonic sources before
setting the sources to zero.
\end{lemma}
\begin{proof}
All carriers are finite at fixed regulator.  Hence
\begin{align*}
 \delta_\epsilon\operatorname{Tr}\log\mathbb H
 &=\operatorname{Tr}\bigl(
 \mathbb H^{-1}[\mathbb R_\epsilon,\mathbb H]\bigr)=0,\\
 \delta_\epsilon\operatorname{Tr}\log\mathbb M_{\rm gh}
 &=\operatorname{Tr}\bigl(
 \mathbb M_{\rm gh}^{-1}[\mathbb R^{\rm gh}_\epsilon,
 \mathbb M_{\rm gh}]\bigr)=0
\end{align*}
by cyclicity.  The Higgs operator transforms by the same finite similarity on
its real four-component carrier.  Differentiating a finite conjugation identity
with respect to commuting retained bosonic sources preserves the identity.
No infinite trace or regulator-removal interchange is used.
\end{proof}

The coefficient measure is equally important.  Compact-link Haar measure is
exactly invariant under internal gauge multiplication.  The Higgs representation
is unitary; on the real four-component carrier its connected transformation is
orthogonal with determinant one.  The adjoint ghost representation is
unimodular, and the hypercharge ghost has trivial adjoint action.  Thus the
finite bosonic coefficient measure has no internal-gauge Berezinian/Jacobian
term.

\begin{theorem}[Zero native bosonic anomaly projection]
\label{thm:v80-native-zero}
Let \(B^{\rm EGH,bos}_{\ell,j}\) be the complete internally gauge-covariant
bosonic Slavnov/BV breaking generated by one normalized hard shell, including
its gauge, ghost, Higgs, gravitational, Haar, coordinate-density and reference
terms in the same convention.  For every fixed perturbative order for which the
finite hard expansion is formed,
\begin{equation}
 \boxed{
 \mathsf A_{\rm std}B^{\rm EGH,bos}_{\ell,j}=0.}
 \label{eq:v80-native-anomaly-zero}
\end{equation}
In particular this holds for the two-loop weight-six coefficient consumed by
V77--V79.
\end{theorem}
\begin{proof}
Form the normalized finite hard integral before local Taylor subtraction.  The
classical bosonic action and the hard/soft blocking data transform covariantly
under the inherited internal gauge action.  Lemma~\ref{lem:v80-det-zero}
shows that all Gaussian/Berezin determinant contributions are invariant.  Haar
and coefficient measures are invariant by the preceding paragraph.  Therefore
the complete normalized hard pushforward has zero internal gauge variation as a
finite identity, with the reference contribution included once.  Extracting a
fixed loop coefficient and then a fixed local source jet preserves this identity.
The anomaly projection \(\mathsf A_{\rm std}\) tests precisely the non-exact
ordinary internal-gauge part of that local breaking, so every one of its
coordinates vanishes.

This argument is specific to the bosonic shell.  It does not apply to a chiral
fermion determinant/Pfaffian line, where a phase can transform nontrivially even
when the formal classical action is gauge invariant.
\end{proof}

A useful corollary is that the abelian factor creates a nontrivial \emph{space of
possible} anomaly representatives but no populated bosonic anomaly coefficient.
This is stronger and more precise than incorrectly declaring the abelian
cohomology itself zero.

\section{Bosonic Einstein--gauge--Higgs two-loop range theorem}
\label{sec:v80-range}

Let \(B^{\rm EGH}_{2,j}\) denote the complete weight-six local two-loop row after
all proper one-loop subgraphs are subtracted by the same common \(S_{1,j}\), and
split its internal-gauge anomaly projection from the orthogonal exact part.
The diffeomorphism/local-Lorentz positive-matter and pure-gravity components are
the ones already treated in V78--V79.

\begin{theorem}[Bosonic EGH weight-six ghost-one range]
\label{thm:v80-range}
On the bosonic
\(\mathrm{Einstein}\!\times SU(3)\times SU(2)\times U(1)_Y\!+\!H\)
branch described above,
\begin{equation}
 \boxed{
 P_{\rm coker}^{\rm EGH,bos}B^{\rm EGH}_{2,j}=0.}
 \label{eq:v80-coker-zero}
\end{equation}
Hence there is a local two-loop restoring functional
\begin{equation}
 \boxed{
 C_{2,j}^{\rm rest,bos}
 =-\bigl(d^{\rm EGH}_{N,6}\bigr)^\dagger B^{\rm EGH}_{2,j}}
 \label{eq:v80-restorer}
\end{equation}
modulo the independent ghost-zero cohomology.  The minimum-norm representative
obeys
\begin{equation}
 \|C_{2,j}^{\rm rest,bos}\|
 \le (\sigma^{\rm EGH}_{N,6})^{-1}\|B^{\rm EGH}_{2,j}\|
 \label{eq:v80-rest-bound}
\end{equation}
with a finite constant at fixed \(N\), representation packet and weight.
\end{theorem}
\begin{proof}
V78--V79 remove the pure gravitational and neutral positive-matter
cokernels for the diffeomorphism/local-Lorentz complex.  The only new ordinary
ghost-one quotient introduced by the internal gauge extension is the standard
anomaly quotient of Theorem~\ref{thm:v80-cohom-reduction}; the special
free-abelian family is absent by Proposition~\ref{prop:v80-not-free}.  The
native coefficient has zero projection onto the standard quotient by
Theorem~\ref{thm:v80-native-zero}.  Therefore the full local row is in the
range.  Finite-dimensional Hodge decomposition gives the stated representative
and bound.
\end{proof}

The independent ghost-zero weight-six quotient is retained as physical EFT data.
In addition to V78--V79's pure/mixed gravitational coordinates it contains, in a
convenient redundant generating bank, representatives of the forms
\begin{equation}
 \begin{gathered}
 R\,\operatorname{tr}F_3^2,
 \quad R\,\operatorname{tr}F_2^2,
 \quad RF_Y^2,
 \quad R|D H|^2,
 \quad R(H^\dagger H)^2,\\
 (H^\dagger H)\operatorname{tr}F_3^2,
 \quad (H^\dagger H)\operatorname{tr}F_2^2,
 \quad (H^\dagger H)F_Y^2,
 \quad |D H|^4,
 \quad (H^\dagger H)^3,
 \end{gathered}
 \label{eq:v80-wilson-bank}
\end{equation}
together with curvature contractions, integration-by-parts partners and
equation-of-motion and other redundant directions required by the chosen basis.  No claim
that the displayed list is already an independent minimal Warsaw-type basis is
made.  The quotient construction, not this schematic list, owns the physical
Wilson coordinates.

\section{Hard-annulus forests and the two-loop source transport survive the gauge extension}
\label{sec:v80-annulus}

The internal gauge extension adds hard gauge, Higgs and ghost propagators but does
not alter the support architecture established in V75.  Soft gravitational
constraints remain retained.  The gauge shell is the inherited measured
hard/soft gauge block; every integrated propagator is therefore supported on a
fixed annulus at scale \(\Lambda_j\), after its proper lower reference factors
are included.

\begin{theorem}[Bosonic EGH two-loop hard-annulus remainder]
\label{thm:v80-annulus}
Fix perturbative order two, weight six, source arity \(N\), the representation
packet and the profile/inverse-margin bounds.  For every connected two-loop
bosonic Einstein--gauge--Higgs graph after all proper one-loop subgraphs are
subtracted by the common \(S_1\) and the complete weight-six local Taylor
assignment is removed,
\begin{equation}
 \boxed{
 |\partial_K^\alpha R^{\rm EGH}_{G,j}(K)|
 \le C_{G,\alpha}\Lambda_j^{-1}
 |K|^{\omega_G+1-|\alpha|}.}
 \label{eq:v80-graph-bound}
\end{equation}
Consequently the new two-loop shell BV connection satisfies
\begin{equation}
 \boxed{
 \|(1-\mathsf P^{\rm EGH}_{N,6})
 \mathcal K^{\rm EGH}_{2,j}\|_{*\to *}
 \le C_{N,2}\frac{\mu}{\Lambda_j}.}
 \label{eq:v80-K2}
\end{equation}
The corresponding rooted kernels carry one \(L^1\) source and the remaining
sources in \(L^\infty\), with no total-volume multiplier.
\end{theorem}
\begin{proof}
The proof is the V78--V79 forest argument with the additional gauge/Higgs/ghost
line types.  Their hard propagators have the same large-momentum order
\(|p|^{-2}\) in the gauge-fixed bosonic reference.  Vertices are finite
polynomials in shell momenta at the prescribed EFT weight.  All proper divergent
subgraphs are removed by restrictions of the same one-loop action \(S_1\);
independent graphwise finite constants are not introduced.  After the full
weight-six Taylor assignment, one additional external momentum remains in the
complementary integrand.  Scale ordering and the finite number of two-loop
forests then give \(\Lambda_j^{-1}\) after shell rescaling.  V75 prevents a
soft lapse/shift inverse from entering the integrated graph.  Gauge-covariant
source rooting and integration by parts give the stated local kernel norm.
\end{proof}

Let
\(\mathsf U^{\rm EGH,[\le2]}_{j+1,j}=I+E^{\rm EGH}_{1,j}+E^{\rm EGH}_{2,j}\)
be the normalized adjacent source map, with loop order displayed.  V76 and V79
supply the gravitational/Higgs first- and second-order pieces, while
Theorem~\ref{thm:v80-annulus} supplies the new internal-gauge part.  Hence
\begin{equation}
 \|E^{\rm EGH}_{1,j}\|_{*\to *}
 \le C_1\frac{\mu}{\Lambda_j},
 \qquad
 \|E^{\rm EGH}_{2,j}\|_{*\to *}
 \le C_2\frac{\mu}{\Lambda_j},
 \label{eq:v80-source-step}
\end{equation}
with the V77 square term separately bounded by
\(C(\mu/\Lambda_j)^2\).

\begin{corollary}[Uniform composed bosonic EGH source transport through two loops]
\label{cor:v80-source-product}
At fixed perturbative/EFT order and source arity,
\begin{equation}
 \boxed{
 \sup_{n\ge j}
 \|\mathsf U^{\rm EGH,[\le2]}_{n,j}\|_{*\to *}
 \le
 \exp\!\left[
 C\sum_{m\ge j}\frac{\mu}{\Lambda_m}
 \right]
 \le
 \exp\!\left[
 \frac{Cb}{b-1}\frac{\mu}{\Lambda_j}
 \right].}
 \label{eq:v80-composed}
\end{equation}
This is uniform in the number of ultraviolet shell steps on the declared fixed
bosonic perturbative/EFT domain.
\end{corollary}

\section{First bosonic Einstein--gauge--Higgs two-loop filtered self-map}
\label{sec:v80-selfmap}

The complete finite-order state now contains
\begin{equation}
 \mathfrak S_j^{\rm EGH,bos}
 =\bigl(
 u_j^{\rm grav},g_{3,j},g_{2,j},g_{Y,j},g_{H,j};
 w_{6,j},R_j;
 \mathcal N_{<,j},\mathcal G_{<,j};
 \mathcal J_j,\mathcal T_j
 \bigr),
 \label{eq:v80-state}
\end{equation}
where \(w_{6,j}\) includes every local ghost-zero weight-six Wilson coordinate
selected by the finite quotient, not merely the examples in
\eqref{eq:v80-wilson-bank}.  The gauge couplings are physical lower-weight
coordinates.  Their detailed beta functions are not recomputed here; the
present theorem concerns closure of the state which carries them.

\begin{theorem}[Bosonic Einstein--gauge--Higgs two-loop self-map]
\label{thm:v80-selfmap}
On the fixed bosonic perturbative domain, after imposing one common set of
lower-weight physical normalization conditions and retaining the complete
weight-six Wilson quotient, the normalized measured shell defines
\begin{equation}
 \boxed{
 \mathcal R_j^{[\le2]}:
 \mathfrak C^{\rm EGH,bos}_{N,6}
 \longrightarrow
 \mathfrak C^{\rm EGH,bos}_{N,6}.}
 \label{eq:v80-selfmap}
\end{equation}
The local ghost-one row is restored by
\eqref{eq:v80-restorer}; all independent ghost-zero EFT coordinates remain in
\(w_{6,j}\); the quasilocal source/BV complement obeys
\eqref{eq:v80-K2}; and the composed source transport obeys
\eqref{eq:v80-composed}.  On the action side the dimension-six filtered window
remains
\begin{equation}
 \boxed{2\le\sigma=3<4.}
 \label{eq:v80-window}
\end{equation}
Thus the coefficientwise two-boundary trajectory already established in
V73--V79 extends to the bosonic internal-gauge/Higgs state at this fixed order.
\end{theorem}
\begin{proof}
The local range statement is Theorem~\ref{thm:v80-range}.  The independent
physical Wilson coordinates are kept rather than projected out.  Theorem
\ref{thm:v80-annulus} and Corollary~\ref{cor:v80-source-product} close the
quasilocal/source block.  The gauge/Higgs marginal coordinates belong to the
lower-weight local bank and are transported by the inherited gauge-shell
normalization.  Weight-six local coefficients have the same canonical
\(b^{-2}\) ultraviolet carry factor as in V73--V74, while the complementary
Taylor remainder begins two dimensions higher and carries \(b^{-4}\).  Hence
one may keep the same intermediate weight \(\sigma=3\).  The filtered
V6 two-boundary theorem then applies coefficientwise at the present finite
order.
\end{proof}

This is not yet a bosonic Standard-Model continuum theorem to arbitrary order.
It is the first two-loop source-complete self-map in which the actual three
internal gauge factors and the charged Higgs occupy the same filtered state as
the retained gravitational sector.

\section{Anti-circularity audit and status ledger for V80}
\label{sec:v80-ledgers}

\subsection*{Proved}
Hypercharge is not a free abelian gauge factor on the charged-Higgs branch in
the precise sense relevant to the Einstein--Yang--Mills BRST classification.
The weight-six finite EGH bank is a chain complex.  Its ghost-one cokernel
reduces to the ordinary standard gauge-anomaly quotient rather than the
free-abelian current family.  The complete native bosonic hard shell has zero
projection onto that quotient because its finite bosonic/ghost Hessians transform
by similarity and its Haar/Lebesgue/Berezin coefficient measures are internally
gauge invariant.  Combining this with V78--V79 gives zero populated local
ghost-one cokernel for the complete bosonic Einstein--gauge--Higgs two-loop row.
The independent ghost-zero gauge/gravity/Higgs Wilson quotient remains in the
state.  The mixed gauge/Higgs/gravity two-loop hard forests obey the same
\(O(\mu/\Lambda_j)\) complementary source bound, yielding a composed two-loop
source transport uniform in the number of ultraviolet shell steps on the fixed
bosonic domain and a filtered two-loop self-map.

\subsection*{Inherited}
V1--V4 and the later gauge-shell interfaces supply the gauge-covariant measured
blocking, compact Haar ownership, gauge/ghost reference and lower-weight Slavnov
normalizations.  V20 and V29 record the relevant local-cohomology distinction for
abelian factors.  V73 supplies the scalar filtered action spine; V74 the first
mixed gravity--Higgs hard row; V75 the retained soft constraints; V76 the shell
BV connection; V77 the weight-six two-loop range equation; V78--V79 the
pure-gravity and neutral-matter range/forest estimates.

\subsection*{Literature input}
The local cohomological classification used to identify the finite ordinary
anomaly quotient is the matter-inclusive Einstein--Yang--Mills analysis
\cite{V20BBH1995}, with the general antifield framework
\cite{V20BBH1994}.  In that classification a free abelian generator is an
abelian generator acting trivially on every matter field, and additional
antifield-dependent families occur in that special case.  The standard
four-dimensional gauge-anomaly representatives themselves are literature
structure.  No anomaly coefficient is imported: their vanishing in the present
bosonic shell is proved by the finite native conjugation/Jacobian identities.

\subsection*{Conditional}
The theorem is at fixed loop order two, weight six, source/field arity, profile
class and perturbative coupling domain.  Numerical beta functions and finite
Wilson matching coefficients are not supplied here.  The hard-annulus estimate
uses the inherited gauge-fixed inverse bounds and V75 support separation.  The
statement is coefficientwise/EFT and not uniform in an unbounded perturbative or
canonical-dimension tower.

\subsection*{Open}
The smallest load-bearing interface is now the chiral fermion/Yukawa sector.
Once quarks and leptons are inserted, the standard anomaly projection need not
vanish term by term; it must cancel only after the actual Standard-Model
representation and hypercharge assignments are assembled.  The local anomaly
condition must then be tied to the determinant/Pfaffian-line transport rather
than treated as sufficient by itself.  The next version should therefore adjoin
the chiral matter representations and Yukawa block, compute the finite anomaly
trace vector in the common convention, and test whether the same shell map
transports one coherent determinant/Pfaffian orientation.  No nonperturbative
fermion topology change or Standard-Model ultraviolet completion is asserted.

\begin{center}\footnotesize
\begin{tabular}{p{.24\textwidth}p{.68\textwidth}}
\toprule
Variable & Scope of V80\\\midrule
Loop/EFT order & Bosonic EGH closure through two loops at fixed local weight six and source arity.\\
Gauge group & $SU(3)\times SU(2)\times U(1)_Y$ with charged Higgs; no fermions or Yukawa couplings.\\
Abelian issue & Hypercharge is not free abelian because $y_H\ne0$; free-abelian current anomaly family is absent, standard abelian anomaly quotient retained and tested.\\
Native anomaly & Standard anomaly projection vanishes for the complete bosonic hard shell by finite conjugation and measure invariance.\\
Physical EFT bank & Complete ghost-zero weight-six quotient retained; schematic $RF^2$, $(H^\dagger H)F^2$, $R|DH|^2$, $|DH|^4$, $(H^\dagger H)^3$ representatives included.\\
Soft sector & Gravitational constraints/nonminimal variables remain retained as in V75; no inverse soft wave operator.\\
Source transport & Bounded through order two uniformly in the number of UV shell steps on the fixed bosonic domain.\\
Chiral line & Not included.  Standard-Model fermion anomaly cancellation and determinant/Pfaffian coherence remain open.\\
\bottomrule
\end{tabular}
\end{center}

\section*{Append-only continuation marker: V80}
\addcontentsline{toc}{section}{Append-only continuation marker: V80}
\label{sec:v80-continuation-marker}
\textbf{End of V80.} Preserve Sections 1--612.  The charged Higgs removes the
free-abelian hypercharge exception, while the standard internal-gauge anomaly
quotient is retained explicitly and shown to have zero coefficient in the native
bosonic shell.  Together with V78--V79 this closes the bosonic
Einstein--$SU(3)\times SU(2)\times U(1)_Y$--Higgs source/BV self-map through two
loops at weight six.  The next load-bearing extension is the chiral
fermion/Yukawa and determinant/Pfaffian-line sector, where anomaly cancellation
must be proved in the populated Standard-Model representation rather than
inferred from bosonic gauge invariance.


\end{document}
